% Options for packages loaded elsewhere
\PassOptionsToPackage{unicode}{hyperref}
\PassOptionsToPackage{hyphens}{url}
\PassOptionsToPackage{dvipsnames,svgnames,x11names}{xcolor}
%
\documentclass[
  10pt,
]{article}
\usepackage{amsmath,amssymb}
\usepackage{iftex}
\ifPDFTeX
  \usepackage[T1]{fontenc}
  \usepackage[utf8]{inputenc}
  \usepackage{textcomp} % provide euro and other symbols
\else % if luatex or xetex
  \usepackage{unicode-math} % this also loads fontspec
  \defaultfontfeatures{Scale=MatchLowercase}
  \defaultfontfeatures[\rmfamily]{Ligatures=TeX,Scale=1}
\fi
\usepackage{lmodern}
\ifPDFTeX\else
  % xetex/luatex font selection
  \setmainfont[]{DejaVu Serif}
  \setsansfont[]{DejaVu Sans}
  \setmonofont[]{DejaVu Sans Mono}
  \setmathfont[]{Latin Modern Math}
\fi
% Use upquote if available, for straight quotes in verbatim environments
\IfFileExists{upquote.sty}{\usepackage{upquote}}{}
\IfFileExists{microtype.sty}{% use microtype if available
  \usepackage[]{microtype}
  \UseMicrotypeSet[protrusion]{basicmath} % disable protrusion for tt fonts
}{}
\makeatletter
\@ifundefined{KOMAClassName}{% if non-KOMA class
  \IfFileExists{parskip.sty}{%
    \usepackage{parskip}
  }{% else
    \setlength{\parindent}{0pt}
    \setlength{\parskip}{6pt plus 2pt minus 1pt}}
}{% if KOMA class
  \KOMAoptions{parskip=half}}
\makeatother
\usepackage{xcolor}
\usepackage[margin=22mm]{geometry}
\usepackage{longtable,booktabs,array}
\usepackage{calc} % for calculating minipage widths
% Correct order of tables after \paragraph or \subparagraph
\usepackage{etoolbox}
\makeatletter
\patchcmd\longtable{\par}{\if@noskipsec\mbox{}\fi\par}{}{}
\makeatother
% Allow footnotes in longtable head/foot
\IfFileExists{footnotehyper.sty}{\usepackage{footnotehyper}}{\usepackage{footnote}}
\makesavenoteenv{longtable}
\setlength{\emergencystretch}{3em} % prevent overfull lines
\providecommand{\tightlist}{%
  \setlength{\itemsep}{0pt}\setlength{\parskip}{0pt}}
\setcounter{secnumdepth}{-\maxdimen} % remove section numbering
\usepackage{amsmath,amssymb,amsthm}
\usepackage{fvextra}
\usepackage{newunicodechar}
\providecommand{\bigboxtimes}{\mathop{\boxtimes}\displaylimits}
\providecommand{\PsiAt}{\operatorname{PsiAt}}
\newunicodechar{𝓞}{\ensuremath{\mathcal O}}
\newunicodechar{ℓ}{\ensuremath{\ell}}
\newunicodechar{∎}{\ensuremath{\blacksquare}}
\newunicodechar{⋯}{\ensuremath{\cdots}}
\newunicodechar{⋊}{\ensuremath{\rtimes}}
\newunicodechar{𝔣}{\ensuremath{\mathfrak f}}
\newunicodechar{𝜎}{\ensuremath{\sigma}}
\newunicodechar{ℒ}{\ensuremath{\mathcal L}}
\newunicodechar{〈}{\ensuremath{\langle}}
\newunicodechar{〉}{\ensuremath{\rangle}}
\DefineVerbatimEnvironment{verbatim}{Verbatim}{breaklines,breakanywhere,fontsize=\small}
\setlength{\emergencystretch}{3em}
\setcounter{tocdepth}{2}
\ifLuaTeX
  \usepackage{selnolig}  % disable illegal ligatures
\fi
\usepackage{bookmark}
\IfFileExists{xurl.sty}{\usepackage{xurl}}{} % add URL line breaks if available
\urlstyle{same}
\hypersetup{
  pdftitle={ALP\_MACHINE: Mixed CM theta sections, lift costs, and arithmetic comparison gates},
  pdfauthor={GPT-6 Astra Pro research team},
  colorlinks=true,
  linkcolor={blue},
  filecolor={Maroon},
  citecolor={Blue},
  urlcolor={blue},
  pdfcreator={LaTeX via pandoc}}

\title{ALP\_MACHINE: Mixed CM theta sections, lift costs, and arithmetic comparison gates}
\author{GPT-6 Astra Pro research team}
\date{21 November 2026 --- active research version}

\begin{document}
\maketitle

\section{Status and conventions}\label{status-and-conventions}

The requested shorter actual mixed-type lift chain, universal negative lift-cost theorem, arithmetic MainX pairing value, and Kummer Hecke RH theorem have not been established. This is the integrated active paper, containing the full proofs written so far. Checked findings and work awaiting whole-proof review are separated below. The dated register gives the exact status of each original question and each completed subproblem.

Alpöge's theorem and archive are accepted as stipulated. No Lean build or independent validity audit of that theorem is claimed. Settled prior findings retain the supplied intake statuses. A literature theorem is an input, not a newly proved finding. All new proofs have stable identifiers; the event log records who read a complete proof and distinguishes those readings from numerical diagnostics and exact computations.

The notation \nolinkurl{M} in the source arithmetic Level is distinct from the number of iterated lifts when both occur in a discussion. The actual Level requires an odd integer. Three theta germs must also be kept distinct: the value-one even-characteristic germ F1 on the D135 line, its normalized Mumford translate Sigma\_tr on the trace-one line, and the trace-one germ sigma\_q with quadratic term u1*u3-u2\^{}2. AM-R-006 computes the scalar between the latter two. Later findings refine the limitations recorded in earlier contributions; none silently replaces a selected automorphic image by an unrestricted formal function.

ALP-LIFT-01 and AM-R-001 are two independently written versions of the same elementary changed-family accounting identity. They are retained with their check histories but count as one finding in novelty summaries.

\tableofcontents
\clearpage
\part{Checked findings and their full proofs}

\clearpage

\section{ALP\_MACHINE root contribution: actual comparison and endpoint normalization}\label{alp_machine-root-contribution-actual-comparison-and-endpoint-normalization}

Author: GPT-6 Astra Pro, \nolinkurl{/root}; reasoning effort is not exposed by the runtime. Run: ALP\_MACHINE, 21 November 2026. This is an application of accepted DP\_MULT results and directly read source definitions. It is not a Lean execution or a reproof of Alpöge's theorem.

\subsection{AM-R-001. Comparison of two rebuilt families}\label{am-r-001.-comparison-of-two-rebuilt-families}

Let \(\mathscr D\) and \(\mathscr D'\) be two fully specified families of lengths \(M\) and \(M'<M\). Fix corresponding family members, but do not suppose their endpoints or coefficient fields are equal. Use DP44's scalar definitions on each family, with primes denoting the second family. Thus

\[
 Q_M=\sum_{j=1}^{M}(g_j-e_j)-T_\tau-MH+K_u,
\qquad
 Q'_{M'}=\sum_{j=1}^{M'}(g'_j-e'_j)-T'_\tau-M'H'+K'_{u'}.
\]

Then the exact difference is

\[
\begin{split}
Q'_{M'}-Q_M={}&-\sum_{j=M'+1}^{M}(g_j-e_j-H)\\
&+\sum_{j=1}^{M'}\bigl((g'_j-g_j)-(e'_j-e_j)\bigr)
-M'(H'-H)\\
&-(T'_\tau-T_\tau)+(K'_{u'}-K_u).
\end{split}\tag{R1}
\]

If the compared quantity is \(Q-B\), add \(B_M-B'_{M'}\) to the right side. No sign of \(H\), or equality between the two period normalizations, is assumed.

\textbf{Proof.} Subtract the two displayed defining expressions and split the unprimed sum into retained and deleted indices. The unprimed main-term difference is \((M-M')H-M'(H'-H)\). Combining its first term with the deleted sum gives (R1). This short proof gives the changed-family form that DP44 did not write as an equation; the frozen-family formula DP44.3 is carried, not rederived as a new finding.

\textbf{Application and limitation.} A proposed calculation of a lift saving must specify every term of (R1). The supplied curve, prime, formal-group slopes and generic rank records do not supply those numbers. \nolinkurl{SizeData₂.gain}, \nolinkurl{SizeData₂.E}, \nolinkurl{SizeData₂.tau}, the weights and level are additional source data tied to a family through the global links. \nolinkurl{LinkG_at} includes a coarse upper bound as well as the ramp inequality; that ramp inequality is only a lower certificate and cannot be reversed into a sharp upper bound on true gains. Thus the local slope upper obstructions in DP42/DP51 do not supply the missing upper bound in (R1). This is a scope statement about the available implications, not a theorem that no shorter family exists. In particular, even a negative difference alone is not equivalent to failure of the full global theorem's inequalities.

\textbf{Locators.} DP\_MULT DP44(d), equations DP44.1--DP44.4; original source \nolinkurl{Formalism/LocalRingCardinality/DataParams₂.lean}, \nolinkurl{SizeData₂} lines 76--181, \nolinkurl{sz} 203--206, \nolinkurl{gainD} 212--213, \nolinkurl{Ga} 220, \nolinkurl{Ee} 222, \nolinkurl{logOm} 224, \nolinkurl{MAIN} 232--234; \nolinkurl{Formalism/SharpRestriction/GlobalTheoremInputs.lean}, \nolinkurl{LinkDIV_at} 77--89 and \nolinkurl{LinkG_at} 122--136. Source bytes were recovered from the intake-pinned mirror and matched to the intake's original-author archive SHA256 values.

\subsection{AM-R-002. The arithmetic and archimedean cost of an endpoint detector}\label{am-r-002.-the-arithmetic-and-archimedean-cost-of-an-endpoint-detector}

This statement makes quantitative the distinction retained in DP55 between a nonzero complex pairing and an integral, controlled detector. It is an elementary application of the exact source pairing, not a claim that ideal content or dual norms are new general concepts.

Let \(K\) be a number field, \(\mathcal R\) a full projective \(\mathcal O_K\)-lattice in the row coefficient space \(W\), and \(B:V\times W\to K\) the algebraic endpoint pairing, including its stabilizer and coefficient normalization. Let \(0\ne f\in V(K)\), and suppose the functional \(\ell_f=B(f,-)\) is nonzero. Define its nonzero fractional content ideal

\[
 I_f=\ell_f(\mathcal R)\subset K.\tag{R2}
\]

At a finite place \(\mathfrak p\), write \(I_f\mathcal O_{K,\mathfrak p}=\pi_{\mathfrak p}^{s_{\mathfrak p}}\mathcal O_{K,\mathfrak p}\), with the valuation normalized by \(v_{\mathfrak p}(\pi_{\mathfrak p})=1\). Then:

\begin{enumerate}
\def\labelenumi{\arabic{enumi}.}
\tightlist
\item
  There is a detector \(z\in\mathcal R_{\mathfrak p}\) with \(B(f,z)=1\) if and only if \(s_{\mathfrak p}\le0\).
\item
  The least integer \(d\ge0\) for which such a detector exists in \(\pi_{\mathfrak p}^{-d}\mathcal R_{\mathfrak p}\) is \[d_{\min}(f,\mathfrak p)=\max\{0,s_{\mathfrak p}\}.\tag{R3}\]
\item
  If \(r_1,\ldots,r_N\) is a local lattice basis and \(b_i=B(f,r_i)\), then \[s_{\mathfrak p}=\min_i v_{\mathfrak p}(b_i).\tag{R4}\]
\end{enumerate}

For a fixed embedding \(\sigma:K\hookrightarrow\mathbb C\), put a positive Hermitian norm on \(W_\sigma\), and let \(\beta_\sigma(f)=\|\ell_{f,\sigma}\|_*\). Then

\[
 \inf_{B_\sigma(f,z)=1}\|z\|_\sigma=\beta_\sigma(f)^{-1},\tag{R5}
\]

and the infimum is attained. If the norm on rows is transported through DP55's representing map \(J\), then \(\beta_\sigma(f)=\|f\|_{\rm cl,\sigma}\), provided that transported form is a norm; alternatively apply (R5) to the quotient by \(\ker J\).

For \(a\in K^\times\),

\[
 I_{af}=aI_f,\qquad \beta_\sigma(af)=|\sigma(a)|\beta_\sigma(f).\tag{R6}
\]

Consequently the signed normalization balance

\[
 \mathcal B(f)=\frac1{[K:\mathbb Q]}
 \left(\sum_{\mathfrak p}s_{\mathfrak p}\log N\mathfrak p
       -\sum_{\sigma:K\hookrightarrow\mathbb C}\log\beta_\sigma(f)\right)\tag{R7}
\]

is invariant under \(f\mapsto af\). All complex embeddings are counted; this convention automatically counts a conjugate pair twice. It is the negative Arakelov degree of the quotient image line \(I_f\), equipped with the quotient norm \(\|1\|_\sigma=\beta_\sigma(f)^{-1}\).

\textbf{Full proof.} Since \(\mathcal R\) is a finite lattice, its image under a nonzero \(K\)-linear map is a nonzero finitely generated \(\mathcal O_K\)-submodule of \(K\), hence a fractional ideal. Localization commutes with its image. The image of \(\pi^{-d}\mathcal R_{\mathfrak p}\) is exactly \(\pi^{s_{\mathfrak p}-d}\mathcal O_{K,\mathfrak p}\); it contains 1 exactly when \(d\ge s_{\mathfrak p}\). This proves (R3), including negative \(s_{\mathfrak p}\), and assertion 1. A basis generates the image ideal, and a finitely generated ideal over a DVR has valuation equal to the least valuation of its generators, proving (R4).

By the finite-dimensional Riesz representation theorem, there is a unique nonzero vector \(y\in W_\sigma\) such that \(\ell_{f,\sigma}(z)=\langle y,z\rangle\), in the convention conjugate-linear in the first argument. Its norm is \(\beta_\sigma(f)\). Cauchy--Schwarz gives \(1\le\beta_\sigma(f)\|z\|\) for every detector. The vector \(z=y/\|y\|^2\) is a detector and attains equality. This proves (R5). DP55's identity \(B(f,z)=\langle Jz,f\rangle_{\rm cl}\), with \(J\) conjugate-linear and onto, gives the stated transported-norm specialization by the same argument; if necessary take its injective quotient first.

Bilinearity gives the ideal equality in (R6), and the definition of operator norm gives the norm equality. The change in the numerator of (R7) is \[\sum_{\mathfrak p}v_{\mathfrak p}(a)\log N\mathfrak p-\sum_\sigma\log|\sigma(a)|=0\] by the product formula. Finally, relative to the lattice \(I_f\), the rational section 1 has divisor coefficient \(-s_{\mathfrak p}\); its quotient norm is \(\beta_\sigma(f)^{-1}\). Substituting these two facts into the definition of arithmetic degree gives \(-\mathcal B(f)\). This proves every assertion.

\subsubsection{The exact source coefficient matrix}\label{the-exact-source-coefficient-matrix}

DP55's source form is the bilinear normalized Fock form

\[
 B_\lambda(x,y)=D_\lambda^{-1}\sum_\alpha x_\alpha y_\alpha\alpha!,
\quad
 D_\lambda=\prod_\tau\sum_\alpha (\operatorname{ext}_{\lambda_\tau})_\alpha^2\alpha! .\tag{R8}
\]

At class \(c\) it is multiplied by \(1/a_c\), where \(a_c\) is the actual finite stabilizer order. Choose bases of the actual invariant coefficient lattices, not a larger unconstrained polynomial module. Write \(G_c\) for their raw Fock matrix and \(f_c\) for the point coordinates. When the row lattice is expressed as the direct sum of these invariant class lattices, (R4) becomes the explicit formula

\[
 s_{\mathfrak p}
 =\min_{c,b}v_{\mathfrak p}\left(
       \frac{(f_c^{\mathsf T}G_c)_b}{a_cD_\lambda}\right).\tag{R9}
\]

If integral coordinates are tied together by an additional lattice inclusion \(C\), replace the row of all displayed coefficients by its product with \(C\). Formula (R9) then applies to generators of that actual image lattice. Thus no arbitrary basis is silently assumed unimodular.

\textbf{Proof of the application.} The source period is the finite sum of the classwise forms (R8), weighted by \(a_c^{-1}\). Evaluating it on a basis row supported on one class gives exactly the coefficient inside (R9). Formula (R4) then applies. If rows are restricted to an integral submodule, evaluating on its generators multiplies by its inclusion matrix. This proves the corrected formula for that case too.

\textbf{Locators.} DP\_MULT DP55, (ED3), (ED7), (ED11), and its explicit statement that the arbitrary detecting row has no proved denominator or norm bound. Directly read original-source declarations: \nolinkurl{Formalism/LocalToGlobal/Columns/CoefModB.lean}, \nolinkurl{fockB} line 18, \nolinkurl{extNormB} 55--56, \nolinkurl{pairB} 58--59, \nolinkurl{pairBK} 61--62, \nolinkurl{PairBDual} 113--119; \nolinkurl{Columns/Record.lean}, \nolinkurl{perCl} 315--329 and \nolinkurl{periodC} 1568--1571; \nolinkurl{Columns/TiedPer.lean}, \nolinkurl{per_jMCL_eq} 243--244 and \nolinkurl{per_jMC_onto} 250--252. Perfectness over the field does not assert that the two chosen integral lattices are mutually dual.

\textbf{Consequence for Q3e.} A unit-valued complex period cannot be inserted at zero arithmetic cost merely by rescaling. For a specified algebraic terminal lift vector, the endpoint denominator problem is exactly the content computation (R9), and the archimedean problem is the dual norm (R5). These are computable once the actual kernel/image and lattice matrices are given. This result does not construct that terminal lift vector or identify its lattice, and therefore does not decide the sign in (R1).

\subsection{AM-R-003. Powers of the genus-three Jacobian do not repair an elliptic Serre-tensor requirement}\label{am-r-003.-powers-of-the-genus-three-jacobian-do-not-repair-an-elliptic-serre-tensor-requirement}

Let \(J/\overline{\mathbb Q}\) be the Jacobian of \(y^2=x^7-1\), with its accepted CM type \(\Phi=\{1,2,3\}\) on \(E=\mathbb Q(\zeta_7)\). For every \(N>0\), no nonzero elliptic-power abelian variety is isogenous to an abelian subvariety or quotient of \(J^N\). In particular there is no nonzero isogeny, nor a nonzero homomorphism, in either direction between \(J^N\) and \(C^d\), for any elliptic curve \(C\) and \(d>0\).

This remains true after a finite extension of the base field, after twisting the displayed CM action by automorphisms, and for the geometric product \(J^P\times\bar J^Q\) used by DP65, where \(\bar J\) is the same abelian variety with conjugate \(E\)-action. Consequently an actual Serre tensor of an elliptic curve by a finite projective module cannot become such a geometric power of \(J\) through a field enlargement or an isogeny. This excludes only attempts to retain the source's elliptic \nolinkurl{IsTen} requirement. It is not a no-go theorem for a newly defined mixed-type theta correspondence.

\textbf{Proof.} The only proper CM subfield of the cyclic degree-six field \(E\) is its imaginary quadratic field \(K=\mathbb Q(\sqrt{-7})\). The two induced types from \(K\) are \(\{1,2,4\}\) and \(\{3,5,6\}\). The actual type \(\{1,2,3\}\) is neither, hence primitive. Milne, \emph{Complex Multiplication} (2020), Proposition 3.13, identifies primitive CM types with simple complex CM abelian varieties up to isogeny. Thus \(J\) is geometrically simple of dimension three.

A nonzero homomorphism \(J\to C\) would have an abelian image of positive dimension, hence image \(C\). Its connected kernel would have dimension two, contradicting simplicity. A nonzero map \(C\to J\) would give a positive-dimensional abelian subvariety of dimension at most one, again a contradiction. Every homomorphism between products is a matrix of homomorphisms between factors, so all maps \(J^N\leftrightarrow C^d\) vanish. By Poincare complete reducibility, a nonzero elliptic isogeny factor of a subvariety or quotient of \(J^N\) would give such a nonzero map, which is impossible.

The argument is over \(\overline{\mathbb Q}\), so a finite base extension changes nothing. Twisting a named endomorphism action does not change the underlying variety. For a finite projective module \(M\) over an order acting on an elliptic curve, \(M\otimes C\) is, after clearing denominators in a module basis, isogenous to an elliptic power. Hence it cannot be isogenous to a nonzero variety in the \(J\)-isotypic category. This proves the statement.

\textbf{Locators and scope.} The CM type and the earlier type-package obstruction are carried DP22, DP25, DP49, DP52(d), DP57. The exact old \nolinkurl{IsTen} input is \nolinkurl{Development/AbelianVarietyTheory/Definitions/CMLattices/IsDiffOf.lean}, lines 73--77, as identified in DP57 §7. The new assertion applies the simple-isogeny-factor invariant to arbitrary powers and field enlargements; it does not reprove the accepted local derivative results. In motives it says that the weight-one simple constituent \(h^1(J)\) cannot be replaced by an elliptic \(h^1\) constituent. Exterior and tensor powers of motives are a different operation; no vanishing of all higher-weight correspondences is claimed.

\clearpage

\section{ALP\_MACHINE: lift accounting and explicit CM metric calculations}\label{alp_machine-lift-accounting-and-explicit-cm-metric-calculations}

Author: /root/lift\_accounting, GPT-6 Astra Pro. Reasoning effort: not exposed. Started 2026-11-21 19:35 UTC. This working contribution is not a resolution of DP\_MULT Q3e. Accepted DP findings are cited and are not re-proved.

The later, complete companion findings are:

\begin{itemize}
\item
  \nolinkurl{hodge_lattice_and_cm_basis.md}: ALP-LIFT-05, the actual stable Hodge lattices at 2 and 7, full determinant metric/height, and a primitive induced CM eigenbasis with Gram matrix \(G_B I\).
\item
  \nolinkurl{directional_budget_cancellation.md}: ALP-LIFT-06, the exact directionwise prime constants and the arbitrary-order-vector cancellation against the actual diagonal Cauchy cost.
\item
  \nolinkurl{finite_theta_null_normalization.md}: ALP-LIFT-07, the exact finite theta-null normalization, explicit zero characteristic, trace-one principal line and the surviving two-torsion difference at source odd levels.
\item
  \nolinkurl{trace_one_leading_jet.md}: ALP-LIFT-08, the unique actual trace-one section with leading jet \(u_1u_3-u_2^2\), primitive at every \(p\nmid14\), and its complete finite scalar contents at 2 and 7.
\item
  \nolinkurl{actual_torus_first_lift.md}: ALP-LIFT-09, the norm-one torus class calculation, an explicit honest global splitting, actual nonzero first U(1) to U(2) theta lift, and the exact finite-index cost of a common fixed-weight first-image projector.
\item
  \nolinkurl{permitted_orbit_first_image.md}: ALP-LIFT-10, the exact DP59 half-torsion finite-input family, common actual first-lift input, exact index/norm and exceptional-prime projection costs. AM-K018 closes the ordinary-power source gauge; explicit DP59 harmonic sharp values now occur in an actual first image. Ordinary mixed-line values use a separate multiweight projector.
\item
  \nolinkurl{fixed_conductor_type_repair.md}: ALP-LIFT-11, an explicit global torus character with arbitrary integral infinity type and fixed conductor29, its exact determinant twist of the actual first lift, and removal of the specified next-edge real-type obstruction. Its fixed bilinear repair works for all degrees; the common source-character pins remain separate.
\item
  second\_definite\_lift\_nonvanishing.md: ALP-LIFT-12, an actual nonzero U(2) to U(4) lift of every nonzero repaired first-image vector, including automorphic multiplicities. The proof checks all finite local occurrence, the explicit compact-real Rallis factor, and absolute nonvanishing of the standard L-value. AM-R009 gives the concrete same-host corollary over E-prime. The second Schwartz function is unrestricted; its prescribed arithmetic row and cost data are not constructed.
\item
  fixed\_section\_sharp\_first\_images.md: ALP-LIFT-13, actual sharp first-image coefficients at3 and5 from the single fixed section in AM-K020. One fixed finite vacuum and one unit-norm source character realize every nonzero witness; the exact Fock norm and common period are retained.
\end{itemize}

\nolinkurl{cmdata_theta_constructor_audit.md} records the directly checked input and output of the generic six-identity constructor. Its principal distinction is already accepted DP52 background, not a new finding.

Each companion proof has been read in full and its actual checks are recorded in \nolinkurl{lift_accounting_CHECK_EVENTS.tsv}. None identifies the complete selected \nolinkurl{SizeData₂} or constructs the requested lift chain.

\subsection{1. Source boundary and the actual cost variables}\label{source-boundary-and-the-actual-cost-variables}

The attached DP\_MULT return contains the paper, proof snapshots and manifests, but no raw Lean modules. The present run recovered the pinned mirror modules listed in lift\_sources/FETCH\_MANIFEST.json. Their hashes are compared separately with the intake's original-author manifest; a source read is not a Lean build.

The exact primary declarations are:

\begin{itemize}
\tightlist
\item
  \nolinkurl{Formalism/LocalRingCardinality/DataParams₂.lean,} SizeData₂, lines 79--183: the fields gain, F, E, tau, supp, aT, logDerivativeAtOne and the size windows.
\item
  The same file, sz, gainD, Ga, Ee, logOm and MAIN, lines 203--234.
\item
  \nolinkurl{Formalism/LocalRingCardinality/DataHost.lean,} rdE and v, lines 140--145: \(v=\operatorname{rd}(E)/2\), not \(\sqrt d/2\).
\item
  \nolinkurl{Formalism/SharpRestriction/GlobalTheoremInputs.lean,} LinkDIV\_at, lines 77--89; flG, loG, LinkG\_at, lines 100--136.
\item
  \nolinkurl{Formalism/LocalRingCardinality/SpanPowerMembership/NZRowAsm.lean,} ExtraNZIn, lines 61--86.
\end{itemize}

In particular, a Jacobian, a prime, a CM type and an admissible Ranks record do not determine numerical gain or error fields. The two-sided LinkG\_at conditions include the coarse bound \[
 0\le {\rm gain}_{j,w}\le
 4\,|{\rm TT}|\,a/p,
\] as well as the lower ramplike certificate. The latter is not an equality. LinkDIV\_at controls each \(e_j\) within \(C B_0\) of \[
 \left(\sum_t A_t\right)\log(g(0)g(1)).
\] Its actual value still depends on the selected geometric and integral data. The exact period scalar of the old field package is \[
 \log\Omega=
 \gamma+\log(2\pi)-\log\operatorname{rd}(E)
 -\frac{L'(1,\chi_K)}{L(1,\chi_K)}.
 \tag{L0}
\] Consequently the old \(H\), before any mixed-type replacement, is \[
 H=\sum_t A_t\left[
 \log A_t-1-\gamma-\log(2\pi)+\log\operatorname{rd}(E)
 +\frac{L'(1,\chi_K)}{L(1,\chi_K)}
 \right].
 \tag{L1}
\] This is an exact substitution in source definitions, not a value of a newly constructed mixed theta family.

DP44 and DP52 already settle rank reconstruction, the generic Hermitian spaces, the conditional six-identity induction and the endpoint-period interface. DP65 supplies actual CM-labelled local rows, but explicitly does not supply one SizeData₂ package or membership in a prescribed theta image. Its final statement must not be substituted for either of those missing identifications.

\subsection{2. ALP-LIFT-01: the comparison for two rebuilt families}\label{alp-lift-01-the-comparison-for-two-rebuilt-families}

Let the old family have length \(M\), and the new family length \(M'<M\). Primes, fields and norms may differ; define their scalar totals separately: \[
 Q_M=\sum_{j=1}^{M}(g_j-e_j)-T_\tau-MH+K,\qquad
 Q'_{M'}=\sum_{j=1}^{M'}(g'_j-e'_j)-T'_\tau-M'H'+K'.
\] No identification of local place sets is needed to compare these real numbers. Then \[
\begin{split}
 Q'_{M'}-Q_M={}&-\sum_{j=M'+1}^{M}(g_j-e_j-H)\\
 &+\sum_{j=1}^{M'}[(g'_j-g_j)-(e'_j-e_j)]\\
 &+(T_\tau-T'_\tau)+(K'-K)+M'(H-H').
\end{split}
\tag{L2}
\] If an additional allowance is subtracted, the right side acquires \(B_M-B'_{M'}\).

\textbf{Proof.} Split the old gain and error sums at \(M'\), subtract the two definitions, and use \(MH-M'H'=(M-M')H+M'(H-H')\). This accounts for every displayed term and proves (L2). Adding the allowance is direct subtraction. No positivity assertion is used. \(\square\)

When the retained data and all global terms are equal, (L2) is precisely DP44.3. For a genuinely changed CM type it is the second through fourth lines, rather than the deleted tail alone, that must be computed as well. Moreover, \(Q'_{M'}<Q_M\) alone would not prove failure of a global threshold condition: both quantities could exceed that threshold. This last observation is elementary logic about inequalities; it is not an example of an admissible geometric family.

\subsection{3. ALP-LIFT-02: exact stable CM-height penalty for the proposed Jacobian}\label{alp-lift-02-exact-stable-cm-height-penalty-for-the-proposed-jacobian}

Put \(E=\mathbb Q(\zeta_7)\), \(K=\mathbb Q(\sqrt{-7})\). Let \(J\) have its actual integral CM action of type \(\Phi=\{1,2,3\}\), as in DP22 and DP49. Choose a CM elliptic curve \(B\) with full ring \(\mathcal O_K\), and define \[
 A_{\mathrm{ind}}=\mathcal O_E\otimes_{\mathcal O_K} B
\] using the embedding of \(K\) for which the induced type is \(\Psi=\{1,2,4\}\). In any common standard normalization of stable Faltings height for three-dimensional abelian varieties, \[
\boxed{\;
 h_F(J)-h_F(A_{\mathrm{ind}})
 =\log\frac{\Gamma(2/7)\Gamma(4/7)^3}
                  {\Gamma(3/7)^3\Gamma(5/7)}
 =:\delta_7>0.\;}
\tag{L3}
\] The exact rational computation accompanying this proof certifies \[
       0.05568702<\delta_7<0.05568705.
\tag{L4}
\] The higher precision decimal \(0.0556870335243028288211564027305\ldots\) is a numerical cross-check, not the rational certificate.

\subsubsection{Proof: integral comparison variety}\label{proof-integral-comparison-variety}

Let \(S=\zeta_7+\zeta_7^2+\zeta_7^4\), \(T=-1-S\). The roots \(\zeta_7,\zeta_7^2,\zeta_7^4\) give \[
 \zeta_7^3-S\zeta_7^2+T\zeta_7-1=0.
\] Since \(\mathcal O_K=\mathbb Z[S]\), \(\mathcal O_E=\mathbb Z[\zeta_7]\), and \([E:K]=3\), the elements \(1,\zeta_7,\zeta_7^2\) are an \(\mathcal O_K\)-basis of \(\mathcal O_E\). The Serre tensor construction therefore identifies \(A_{\mathrm{ind}}\) with \(B^3\) as an abelian variety (over a common field of definition), not merely up to isogeny. Thus \(h_F(A_{\mathrm{ind}})=3h_F(B)\). The \(\mathcal O_E\)-action has exactly the embeddings extending the chosen \(K\)-embedding, namely \(\Psi\). No isogeny-invariance assertion for general Faltings heights is being used.

\subsubsection{Proof: the normalized Colmez class function}\label{proof-the-normalized-colmez-class-function}

Let \(\sigma=\sigma_3\) generate \({\rm Gal}(E/\mathbb Q)\), and let \(\chi(\sigma)=\exp(2\pi i/6)\). Then \(\chi^3=\chi_{-7}\), and the exponent sets of \(\Phi\) and \(\Psi\) are respectively \(\{0,1,2\}\) and \(\{0,2,4\}\).

Use Haar probability measure in convolution: \[
 A_\Phi=\mathbf1_\Phi*\mathbf1_{\Phi^{-1}}.
\] For \(\sigma^k\), \(k=0,\ldots,5\), its values are \[
 (1/2,\,1/3,\,1/6,\,0,\,1/6,\,1/3).
\] The corresponding induced-type values are \((1/2,0,1/2,0,1/2,0)\). Discrete Fourier expansion gives \[
 A_\Phi=\frac14\mathbf1+
             \frac19(\chi+\bar\chi)+\frac1{36}\chi^3,\qquad
 A_\Psi=\frac14\mathbf1+\frac14\chi^3.
\tag{L5}
\] These expansions can also be checked at the six group elements, using \(\chi+\bar\chi=(2,1,-1,-2,-1,1)\) and \(\chi^3=(1,-1,1,-1,1,-1)\).

The abelian Colmez theorem applies to both integral CM types. Its normalization is explicit in Yang (2010), Section 3, pp.~464--466, Conjecture 3.1 and the elliptic example (3.3); the abelian case is supplied by Colmez (1993), Théorème 0.5 and Théorème III.2.9, with the possible dyadic ambiguity removed by Obus (2013), Theorem 3.9 in arXiv:1107.0684v2. Write \(Z_\eta=L'(0,\eta)/L(0,\eta)\). In the difference, the trivial character and metric constants cancel. All three nontrivial characters here have conductor \(7\), and their coefficient difference sums to zero. The conductor term therefore cancels as well. Multiplication by \([E:\mathbb
Q]=6\), which converts the type-normalized height to the actual three-dimensional height, gives \[
 h_F(J)-h_F(A_{\mathrm{ind}})
   =-\frac23(Z_\chi+Z_{\bar\chi}-2Z_{\chi^3})
   =\frac43(Z_{\chi^3}-\operatorname{Re}Z_\chi).
\tag{L6}
\]

\subsubsection{Proof: reduction to gamma values}\label{proof-reduction-to-gamma-values}

For any nontrivial Dirichlet character \(\eta\) modulo \(7\), the Hurwitz formula, together with DLMF 25.11.13 and 25.11.18, gives \[
 L(0,\eta)=-\frac17\sum_{a=1}^{6}a\eta(a),\qquad
 L'(0,\eta)=-\log7\,L(0,\eta)+
                 \sum_{a=1}^{6}\eta(a)\log\Gamma(a/7).
\tag{L7}
\] Here the constant from \(\zeta'(0,x)=\log\Gamma(x)-\tfrac12\log(2\pi)\) vanishes because the character sum is zero.

Put \(\ell_a=\log\Gamma(a/7)\). The character values give \(L(0,\chi^3)=1\) and \[
 L(0,\chi)=\frac{4+2\sqrt3\,i}{7},\qquad
 L(0,\chi)^{-1}=1-\frac{\sqrt3}{2}i.
\] Consequently \[
 Z_{\chi^3}=-\log7+\ell_1+\ell_2-\ell_3+\ell_4-\ell_5-\ell_6
\] and \[
 \operatorname{Re}Z_\chi=
 -\log7+\ell_1+\frac14\ell_2+\frac54\ell_3
                    -\frac54\ell_4-\frac14\ell_5-\ell_6.
\] Substitution in (L6) proves the equality in (L3).

\subsubsection{Proof: strict sign and exact rational certificate}\label{proof-strict-sign-and-exact-rational-certificate}

Let \(f(x)=\log\Gamma(x)\), \(h=1/7\). The expression in (L3) is \(-\Delta_h^3 f(2/7)\). By differentiating DLMF 5.15.1, \[
 f'''(x)=-2\sum_{n\ge0}(n+x)^{-3}<0 \quad(x>0).
\] Repeated integration of a finite difference therefore gives \[
 \delta_7=-\int_{[0,h]^3} f'''(2/7+s+t+u)\,ds\,dt\,du>0.
\] This proves the strict sign without numerical approximation.

For a reproducible rational bound, the Weierstrass product (or the same integral identity after integration term by term) yields \[
 \delta_7=\sum_{n=0}^{\infty}
 \log\frac{(7n+3)^3(7n+5)}{(7n+2)(7n+4)^3}
 =\sum_{n=0}^{\infty}\log(1+x_n),
\quad
 x_n=\frac{14n+7}{(7n+2)(7n+4)^3}.
\] For \(t=x/(2+x)\), the six-term sum \[
 2\sum_{k=0}^{5}\frac{t^{2k+1}}{2k+1}
\] is a lower bound on \(\log(1+x)\); the omitted positive tail is at most \(2t^{13}/(13(1-t^2))\). Bound the first \(100\) terms with exact fractions. Each remaining logarithm is a third finite difference of \(\log x\), and lies between \(2/(7n+5)^3\) and \(2/(7n+2)^3\). Integral comparison bounds the tail starting at \(N=100\) between \[
 \frac1{7(7N+5)^2},\qquad \frac1{7(7N-5)^2}.
\] The script colmez\_height\_certificate.py implements exactly these rational inequalities and verifies (L4). This completes the proof. \(\square\)

\subsubsection{Independent normalization cross-check against Wilms}\label{independent-normalization-cross-check-against-wilms}

In Yang's height convention, (L5) gives \[
 h_F^Y(J)=-\frac43\operatorname{Re}Z_\chi
               -\frac16Z_{\chi^3}-\frac34\log7
          =-2.0840659166599491401\ldots,
\] \[
 h_F^Y(A_{\mathrm{ind}})
          =-\frac32Z_{\chi^3}-\frac34\log7
          =-2.1397529501842519690\ldots.
\] Wilms (2026), Theorem 1.1, specialized to \(n=7,g=3\), gives \[
 h_F^W(J)=\log7-\frac32\log\pi+\frac3{14}\log2
     -\sum_{j=1}^{3}
       \log\frac{\Gamma((2j-1)/14)}{\Gamma((j+3)/7)}.
\tag{L8}
\] His Section 7, pp.~12--13, explicitly records the normalization \(h'_F=h_F^W+(g/2)\log\pi\). The duplication formula and \(\prod_{a=1}^6\Gamma(a/7)=(2\pi)^3/\sqrt7\) show algebraically that (L8) equals \(h_F^Y(J)-(3/2)\log\pi\). The independent numerical evaluation is \(-3.8011607454340494014\ldots\). This check catches the degree-six, dimension-three and metric factors; it is not used to prove (L3).

\subsubsection{Exact implication and limit}\label{exact-implication-and-limit}

The actual Jacobian has a strictly larger base CM stable height than the specified integral induced-type product. Thus replacing that product by J does not provide a free saving in this complete invariant of the base Hodge determinant. It does not follow that \(H'\ge H\): the source's \(H\) also incorporates a particular derivative normalization, root discriminant, unequal embedding weights, gauges and other fields. The identification of that source metric with this determinant line is a separate equation that must be proved in a mixed-type construction. There is no asserted sign for \(Q'_{M'}-Q_M\).

In Grothendieck's language, (L3) compares the arithmetic degrees of the metrized determinants of the Hodge bundles \(e^*\Omega^1_{\mathcal J}\) and \(e^*\Omega^1_{\mathcal A_{\mathrm{ind}}}\) on a common semistable integral base. The CM idempotents are those acting on the \(H^1\) motive; (L5) is the explicit Artin-character expression entering its period height. FGA's Picard-scheme theorem supplies the Jacobian, as already used in DP22; the exact FGA locators are carried from that finding. No new manuscript cote or scan is asserted to contain the Colmez formula.

\subsection{4. ALP-LIFT-03: full complex Hodge metric in the actual CM basis}\label{alp-lift-03-full-complex-hodge-metric-in-the-actual-cm-basis}

For \(C:y^2=x^7-1\), put \(\omega_i=x^{i-1}dx/y\) for \(i=1,2,3\). Equip \(H^0(C,\Omega^1)\) with \[
 (\omega,\eta)=\frac{i}{2}\int_{C(\mathbb C)}\omega\wedge\bar\eta.
\] Then \[
 (\omega_i,\omega_j)=0\quad(i\ne j),\qquad
\boxed{\;
 \|\omega_i\|^2=\frac{2\pi}{7}
 \frac{\Gamma(i/7)\Gamma(1/2-i/7)}
      {\Gamma(1-i/7)\Gamma(1/2+i/7)}.\;}
\tag{L9}
\] The norm squares are approximately \(9.47604253508638,\ 8.00728744152725,\ 15.34580357802465\).

\textbf{Proof.} The automorphism \(x\mapsto\zeta_7 x\) preserves the integral and multiplies \(\omega_i\wedge\bar\omega_j\) by \(\zeta_7^{i-j}\). For \(i\ne j\) this scalar is not one, proving orthogonality. Branch points have measure zero. Over the affine line, the map from C has degree two and the two sheets contribute equally, so \[
 \|\omega_i\|^2=
 2\int_{\mathbb C}
       \frac{|x|^{2i-2}}{|x^7-1|}\,dA(x).
\] Substitute \(z=x^7\), accounting for its degree seven and Jacobian \(49|x|^{12}\). The expression becomes \[
 \frac27\int_{\mathbb C}|z|^{2a-2}|1-z|^{-1}\,dA(z),
 \qquad a=i/7.
\tag{L10}
\] For \(a,b>0\) and \(a+b<1\), the complex beta integral is \[
 \int_{\mathbb C}|z|^{2a-2}|1-z|^{2b-2}\,dA(z)
 =\pi\frac{\Gamma(a)\Gamma(b)\Gamma(1-a-b)}
            {\Gamma(1-a)\Gamma(1-b)\Gamma(a+b)}.
\tag{L11}
\] For completeness, write each negative power as its positive gamma integral, apply Tonelli, and integrate the Gaussian in z: \[
 \frac{\pi}{\Gamma(1-a)\Gamma(1-b)}
 \int_{s,t>0}\frac{s^{-a}t^{-b}}{s+t}
                  e^{-st/(s+t)}\,ds\,dt.
\] With \(s=uv,t=u(1-v)\), this becomes \[
 \frac{\pi\Gamma(1-a-b)}{\Gamma(1-a)\Gamma(1-b)}
 \int_0^1 v^{b-1}(1-v)^{a-1}\,dv.
\] The ordinary beta identity proves (L11). All changes of variables are valid on positive integrands; the stated inequalities give convergence. Set \(b=1/2\) in (L11) and substitute into (L10). This proves (L9). \(\square\)

This supplies the full complex metric for arbitrary unequal embedding weights once the geometric differential basis is tied to the chosen theta connection and source gauge. It does not remove that tie or the integral corrections at 2 and 7. In particular the Hodge norms are not silently declared to be the source's \(\Omega\).

Geometrically, (L9) gives the Hermitian metric on the Hodge subbundle of the de Rham realization of \(h^1(J)\), split by the \(\mathcal O_E\) idempotents. Its dual metric controls invariant tangent derivations. The connection to the Poincaré biextension and its theta sections still requires the arithmetic comparison asked for in THETA3 G1.

\subsection{5. ALP-LIFT-04: the exact constant in the saturated bare-budget difference}\label{alp-lift-04-the-exact-constant-in-the-saturated-bare-budget-difference}

This result uses the accepted DP64 definition (CB1); it does not repeat the cyclic-display or global-integrality proofs behind that definition. It determines the linear term left undetermined by DP64's general \(O(N)\) comparison for this particular pair of types and equal orders.

Continue with \(\Phi=\{1,2,3\}\), \(\Psi=\{1,2,4\}\). For a rational prime \(p\ne7\), define the \textbf{formal slope totals} \[
 S_\Phi(p)=\sum_{i\in\Phi}\sum_{r\ge1}
       \frac{\mathbf1_{p^r i\notin\Phi}}{p^r},\qquad
 \Delta S(p)=S_\Phi(p)-S_\Psi(p),
\] where labels are multiplied modulo 7. Let \(S\) be a fixed finite set of primes containing 7. Set \[
 \delta_{7,S}=\delta_7-
       \sum_{p\in S\setminus\{7\}}\Delta S(p)\log p.
 \tag{L12}
\] Then the following are exact statements about these explicitly specified arithmetic certificates:

\begin{enumerate}
\def\labelenumi{\arabic{enumi}.}
\tightlist
\item
  The local difference is \[
  \boxed{\quad
  \Delta S(p)=
  \begin{cases}
  \displaystyle\frac{2p(p+1)}{p^3-1},&p\equiv2,4\pmod7,\\[4pt]
  \displaystyle-\frac{2p(p-1)}{p^3+1},&p\equiv3,5\pmod7,\\[4pt]
  0,&p\equiv1,6\pmod7.
  \end{cases}\quad}
  \tag{L13}
  \]
\item
  In ordinary increasing rational-prime order, \[
  \boxed{\quad
  \lim_{X\longrightarrow\infty}
    \sum_{\substack{p\le X\\p\notin S}}
           \Delta S(p)\log p=\delta_{7,S}.
  \quad}
  \tag{L14}
  \] This convergence is conditional. No rearrangement of the primes is authorized by (L14).
\item
  Put equal derivative order \(n\) on each of the three CM directions in both types, and use the DP64 integer certificates \[
  D_S(n)=\sum_{p\notin S}\log p
      \left(\sum_{i\in\Phi}E^\Phi_{i,p}(n)
                  -\sum_{i\in\Psi}E^\Psi_{i,p}(n)\right).
  \] The sum is finite, since a term can be nonzero only when \(p\le n\). For some constants \(c>0,C_S>0\), \[
  \boxed{\quad
      D_S(n)=n\delta_{7,S}
                 +O_S\!\left(n e^{-c\sqrt{\log n}}\right)
      \qquad(n\longrightarrow\infty).
  \quad}
  \tag{L15}
  \] In particular, with total order \(N=3n\), the exact coefficient of \(N\) is \(\delta_{7,S}/3\).
\item
  The same linear term holds before saturation. If \(D_S(n;X)\) denotes the definition of \(D_S(n)\) with \(p\le X\), then, uniformly for \(n\ge4\), \(X\ge2\), with some fixed \(c>0\), \[
  D_S(n;X)=n\delta_{7,S}
    +O_S\!\left(n e^{-c\sqrt{\log\min\{X,\sqrt n\}}}\right).
  \tag{L15a}
  \] Consequently \(D_S(n;X(n))/n\to\delta_{7,S}\) for \emph{every} \(n\to\infty\) and \(X(n)\to\infty\), including cutoffs much smaller than the derivative order. The three directions still have the same order \(n\), fixed as \(p\) varies.
\end{enumerate}

The source-supported common good-model range in DP49/DP64/DP65 excludes 2 and 7. It is therefore the specialization \[
 \boxed{\quad
 \delta_{7,\{2,7\}}=
       \delta_7-\frac{12}{7}\log2<0,
 \quad}
 \tag{L16}
\] that applies to that range. The identity with \(\delta_7\) itself uses the \emph{formal} cyclic slope at 2; it does not assert an available integral good-model certificate at 2 for the specified curve. Every additional discarded prime has the explicit correction in (L12). The strict sign in (L16) follows from (L4) and the elementary inequality \(\log2>1/2\).

\subsubsection{Proof: exact local class function}\label{proof-exact-local-class-function}

For \(a\in(\mathbb Z/7)^\times\), put \[
 b(a)=\#\{i\in\Phi:ai\notin\Phi\}
                  -\#\{i\in\Psi:ai\notin\Psi\}.
\] The Haar normalization already checked in (L5) gives \[
 b(a)=-6(A_\Phi-A_\Psi)(\sigma_a).
\] Thus \(b(1)=b(6)=0\), \(b(2)=b(4)=2\), and \(b(3)=b(5)=-2\). Write \(b(0)=0\). The function has mean zero over the six nonzero residues, and \[
 b(a)=-6\sum_{\eta\in\{\chi,\bar\chi,\chi^3\}}d_\eta\eta(a),
 \qquad(d_\chi,d_{\bar\chi},d_{\chi^3})=(1/9,1/9,-2/9).
 \tag{L17}
\] Directly from the DP64 definitions, \[
 \Delta S(p)=\sum_{r\ge1}\frac{b(p^r)}{p^r}.
 \tag{L18}
\] For residues 2 or 4 the period-three numerator is \((2,2,0)\), giving \(2(p^2+p)/(p^3-1)\). For residues 3 or 5 the period-six numerator is \((-2,2,0,2,-2,0)\). Its rational sum is \[
 \frac{-2p^5+2p^4+2p^2-2p}{p^6-1}
   =-\frac{2p(p-1)}{p^3+1}.
\] Residues 1 and 6 give zero at every power. This proves (L13), including \(\Delta S(2)=12/7\), \(\Delta S(3)=-3/7\), and \(\Delta S(11)=132/665\).

\subsubsection{Proof: ordinary prime-order convergence and the constant}\label{proof-ordinary-prime-order-convergence-and-the-constant}

The known fixed-modulus prime number theorem, as stated in Rudnick (2015), Section 2, printed p.~2 immediately before equation (1), implies for some fixed \(c_0>0\) \[
 B_1(x):=\sum_{p\le x}b(p)\log p
           =O\!\left(xe^{-c_0\sqrt{\log x}}\right).
 \tag{L19}
\] Indeed the six residue-class main terms cancel because the mean of \(b\) is zero. Replacing the displayed von Mangoldt sum in that theorem by its prime-only part changes it by \(O(\sqrt{x}\log^2x)\), which is absorbed by the stated error after possibly decreasing \(c_0\). Partial summation now shows convergence of \(\sum_p b(p)\log p/p\): its tail is at most a constant times \[
 (1+\sqrt{\log X})e^{-c_0\sqrt{\log X}}.
\] The contribution of \(r\ge2\) in (L18) is absolutely convergent, since \[
 \sum_p\log p\sum_{r\ge2}\frac{|b(p^r)|}{p^r}
 \le2\sum_p\frac{\log p}{p(p-1)}<\infty.
\] This proves the convergence assertion in ordinary prime order before any use of analytic continuation. Conversely, (L13) and the same fixed-modulus prime theorem give \[
 \sum_{p\le X}|\Delta S(p)|\log p
           =\frac43\log X+O(1),
 \tag{L20}
\] so the series is not absolutely convergent.

For real \(s>1\), the absolutely convergent Euler product, DLMF 25.15.2, and (L17) give \[
 \sum_{p\ne7}\log p\sum_{r\ge1}\frac{b(p^r)}{p^{rs}}
       =6\sum_\eta d_\eta\frac{L'(s,\eta)}{L(s,\eta)}.
 \tag{L21}
\] The left side tends as \(s\downarrow1\) to the ordinary prime sum just proved convergent. To justify the passage without a Tauberian premise, apply partial summation to (L19): the prime, first-power part is an integral with integrand bounded, uniformly for \(1\le s\le2\), by a constant times \(e^{-c_0\sqrt{\log x}}/x\), which is integrable at infinity. The higher-power part admits the absolute majorant above. Dominated convergence therefore applies. Nonvanishing of the three \(L(1,\eta)\) is the classical Dirichlet nonvanishing theorem (DLMF 25.15.9).

All three characters are odd, primitive and of conductor 7. The functional equation (DLMF 25.15.5), logarithmically differentiated at \(s=1\), yields \[
 \frac{L'(0,\eta)}{L(0,\eta)}
   +\frac{L'(1,\bar\eta)}{L(1,\bar\eta)}
       =\gamma+\log(2\pi/7).
 \tag{L22}
\] For clarity, in the DLMF normalization the sine factor has logarithmic derivative \((\pi/2)\cot(\pi s/2)=0\) at 1, while \(\Gamma'(1)/\Gamma(1)=-\gamma\); the Gauss sum is constant. The conjugate in (L22) is essential. Since the \(d_\eta\) are conjugation-invariant and sum to zero, (L21) at 1 equals \[
 6\sum_\eta d_\eta\frac{L'(1,\eta)}{L(1,\eta)}
 =-6\sum_\eta d_\eta\frac{L'(0,\eta)}{L(0,\eta)}
 =\delta_7,
\] where the last equality is the already-proved (L6). Removing the fixed finite set of prime summands proves (L12) and (L14).

\subsubsection{Proof: the actual integer floors have the same linear constant}\label{proof-the-actual-integer-floors-have-the-same-linear-constant}

First take \(S=\{7\}\). Define an arithmetic sequence \(a(d)=b(d)\Lambda(d)\), where \(\Lambda\) is the von Mangoldt function and \(b(d)\) is read modulo 7. Thus \(a(d)=0\) except at prime powers, and \[
 D_{\{7\}}(n)=\sum_{d\le n}a(d)\left\lfloor\frac nd\right\rfloor.
 \tag{L23}
\] The fixed-modulus prime theorem directly gives \[
 A(x):=\sum_{d\le x}a(d)
       =2\psi(x;7,2)+2\psi(x;7,4)
          -2\psi(x;7,3)-2\psi(x;7,5)
       =O\!\left(xe^{-c_0\sqrt{\log x}}\right).
 \tag{L24}
\] Also \(\sum_{d\le y}|a(d)|\le2\psi(y)=O(y)\), by the same known prime theorem. Partial summation proves that \(\sum_da(d)/d\) converges, with a tail \(O((1+\sqrt{\log y})e^{-c_0\sqrt{\log y}})\). Its value is \(\delta_7\): the prime-first-power sum is the same one as before, and every prime-power term of exponent at least two belongs to an absolutely convergent series. This last absolute convergence justifies the change from increasing primes to increasing prime powers.

Put \(y=\lfloor\sqrt n\rfloor\). Expanding the floor as a count of positive integers gives the exact finite identity \[
 D_{\{7\}}(n)=
 \sum_{d\le y}a(d)\left\lfloor\frac nd\right\rfloor
 +\sum_{k\le n/y}\bigl(A(n/k)-A(y)\bigr).
 \tag{L25}
\] Here the second sum enumerates exactly the pairs \((d,k)\) with \(d>y\), \(dk\le n\); if \(n/k=y\), its summand is zero. The first sum is \[
 n\sum_{d\le y}\frac{a(d)}d+O(y)
 =n\delta_7+
 O\!\left(n(1+\sqrt{\log y})e^{-c_0\sqrt{\log y}}+y\right).
\] For the second sum, (L24), \(n/k\ge y\), and the harmonic-sum bound give \[
 \sum_{k\le n/y}|A(n/k)-A(y)|
 \ll n(1+\log n)e^{-c_0\sqrt{\log y}}.
\] As \(y\asymp\sqrt n\), these errors together are \(O(n e^{-c\sqrt{\log n}})\) for some smaller \(c>0\). This proves (L15) for \(S=\{7\}\).

For one fixed omitted prime \(p\ne7\), the floor error and the geometric tail give \[
 \sum_{r\ge1}b(p^r)\left\lfloor\frac n{p^r}\right\rfloor
       =n\Delta S(p)+O_p(1+\log n).
\] Indeed there are \(O_p(1+\log n)\) nonzero floors, each rounding error has absolute value at most 2, and the omitted geometric tail after \(p^r>n\) contributes \(O_p(1)\). Removing finitely many such contributions proves (L15) and (L16). This completes the proof. \(\square\)

\subsubsection{Proof of the unsaturated-cutoff refinement (L15a)}\label{proof-of-the-unsaturated-cutoff-refinement-l15a}

It suffices first to treat \(S=\{7\}\). Split the sum into prime first powers and powers with exponent at least two. Write \[
 C_1(X)=\sum_{p\le X}b(p)\log p/p,\qquad
 C_2(X)=\sum_{p\le X}\log p\sum_{r\ge2}b(p^r)/p^r.
\] The argument above gives \[
 C_1(X)+C_2(X)=\delta_7+
             O\!\left(e^{-c_1\sqrt{\log X}}\right)
 \tag{L26}
\] for some smaller \(c_1>0\); also \(C_2(\infty)-C_2(X)=O((1+\log X)/X)\).

The second-power-and-higher rounding error is uniform in \(X\). Indeed, let \[
 H_2(t)=\sum_{\substack{p^r\le t\\r\ge2}}\log p
            =O(\sqrt t\log t).
\] This bound follows from \(\theta(u)=O(u)\), by summing \(\theta(t^{1/r})\) over at most \(\log_2t\) exponents. The floor errors for \(p^r\le n\) have total absolute value at most \(2H_2(n)\). The linear geometric terms with \(p^r>n\) have total absolute value at most \(2n\int_{(n,\infty)}t^{-1}\,dH_2(t)\), which is \(O(\sqrt n\log n)\) by integration by parts. These bounds remain valid after imposing \(p\le X\). Thus that whole contribution equals \[
                    nC_2(X)+O(\sqrt n\log n)
 \tag{L27}
\] uniformly in \(X\).

Put again \(y=\lfloor\sqrt n\rfloor\). If \(X\le y\), the first-power contribution is \(nC_1(X)+O(\theta(X))\), hence (L26) and (L27) give (L15a). If \(X\ge y\), its exact hyperbola decomposition is \[
 \sum_{p\le X}b(p)\log p\left\lfloor\frac np\right\rfloor
 =\sum_{p\le y}b(p)\log p\left\lfloor\frac np\right\rfloor
  +\sum_{k\le n/y}
       \bigl(B_1(\min\{X,n/k\})-B_1(y)\bigr).
 \tag{L28}
\] In each summand, \(y\le\min\{X,n/k\}\le n/k\). Applying (L19) and the harmonic-sum bound makes the second sum \(O(n(1+\log n)e^{-c_0\sqrt{\log y}})\). The first sum is \(nC_1(\infty)+O(n e^{-c_1\sqrt{\log y}}+y)\). Together with (L27) and \(|C_2(\infty)-C_2(X)|\ll(1+\log y)/y\), this proves (L15a) after decreasing \(c\). In both cases \(\sqrt n\log n\) is absorbed by \(n e^{-c\sqrt{\log\min\{X,\sqrt n\}}}\).

For a fixed finite exceptional set, once \(X\) contains it, the omitted-prime estimate following (L25) applies unchanged. If it does not yet contain it, the missing finite correction is bounded by a constant depending on \(S\) times \(n\), while \(X\) stays below the fixed largest exceptional prime. Enlarging the implied constant therefore proves the stated uniform assertion for all \(X\ge2\). \(\square\)

\subsubsection{Consequence, and precise limitation}\label{consequence-and-precise-limitation}

For the formal all-unramified-prime slope comparison, the finite part of the mixed-type gain equals the complete stable CM-height difference in (L3). For the actual common good-model range, the finite part is (L16), with the missing formal 2-adic amount explicitly visible. The new statement reaches the order-\(n\) term of the integer certificates both at saturation and along any growing cutoff, so it is stronger than a comparison of leading \(n\log n\) terms or an Abel-regularized identity.

It is still a theorem about the bare lower certificates, with equal orders held fixed over all primes. It does not identify those certificates with \(g_j\), does not give an upper bound on a selected theta value, and does not replace the source bad-prime, gauge, orbit, connection, row-lattice or division-error terms. In particular the height/slope equality cannot be inserted into \(Q_M\) until the coefficient with which the actual source metric occurs is proved. No conclusion of minimal chain length or impossibility of a shorter mixed-type chain follows from (L12)--(L16).

In motivic terms the same Artin class function \(A_\Phi-A_\Psi\) controls this finite part of the Frobenius-weighted bare certificate and the CM Hodge determinant height, through the functional equation. The exact identification here is a calculation for these two specified invariants; it supplies no new comparison isomorphism between the abelian theta kernel and an automorphic kernel.

\subsection{6. Approaches pursued and present Q3e boundary}\label{approaches-pursued-and-present-q3e-boundary}

\begin{enumerate}
\def\labelenumi{\arabic{enumi}.}
\tightlist
\item
  Re-read DP44's exact comparison and the actual SizeData₂ declarations. The numerical arrays are inputs and the lower certificate is not their exact value. No actual mixed family was found in the supplied return.
\item
  Re-read DP52's six-identity induction and DP65's sharp full-row witnesses. These occupy different interfaces; the latter does not instantiate the former's arithmetic theta family.
\item
  Tested whether field/type or rank restrictions alone give minimality. DP44 already rules out rank minimality; DP25 and DP52 exclude only reuse of the unchanged induced-type wrappers. They do not prohibit a new theta correspondence.
\item
  Examined the old logOm and root-discriminant normalization. This produces exact expressions (L0)--(L1), but not a mixed-family \(H\).
\item
  Calculated the actual complete base CM height, with its strict positive comparison (L3), and the complete complex Hodge Gram (L9). These remove vague dependence on the base CM period as far as these specified metrics are concerned.
\item
  No admissible rescaling counterexample is asserted. Changing gain/E fields in a bare record would not preserve the geometric ties.
\end{enumerate}

The remaining concrete problem is to identify the classical mixed-type theta kernel, its invariant differential operators and its integral lattices with one source-compatible kernel package, then evaluate its gain/error arrays and the gauge change from (L9) in that same package. The full difference is (L2). Neither its strict negativity nor a shorter valid chain has been proved in this contribution.

\subsection{7. Sources with exact locators}\label{sources-with-exact-locators}

\begin{itemize}
\tightlist
\item
  Colmez, P. (1993). Périodes des variétés abéliennes à multiplication complexe. Annals of Mathematics, 138(3), 625--683. Théorèmes 0.3, 0.5 and III.2.9. \url{https://webusers.imj-prg.fr/~pierre.colmez/artin.pdf}
\item
  Obus, A. (2013). On Colmez's product formula for periods of CM-abelian varieties. Mathematische Annalen, 356, 401--418. Theorem 3.9 in the inspected arXiv:1107.0684v2, Section 3.3. \url{https://arxiv.org/abs/1107.0684}
\item
  Yang, T. (2010). The Chowla--Selberg formula and the Colmez conjecture. Canadian Journal of Mathematics, 62(2), 456--472. Section 3, pp.~464--466, definition of \(A_\Phi\), Conjecture 3.1, and equation (3.3); p.~457 states the abelian case. \url{https://doi.org/10.4153/CJM-2010-028-x}
\item
  Wilms, R. (2026). On the Faltings height of the curve \(y^2=x^n-1\). arXiv:2601.15271v1. Theorem 1.1; Section 7, pp.~12--13; Remark 10.2 (the cyclotomic CM type). \url{https://arxiv.org/abs/2601.15271}
\item
  NIST Digital Library of Mathematical Functions. Equations 5.15.1, 25.11.13, 25.11.18. The multiplication and duplication identities used for the cross-check are 5.5.6 and 5.5.5. For ALP-LIFT-04: 25.15.2 (Euler product), 25.15.5 (primitive functional equation, including the conjugate character), and 25.15.9 (nonvanishing at 1). \url{https://dlmf.nist.gov/25.15}
\item
  Rudnick, Z. (2015, May 26). Primes in arithmetic progressions: Fixed modulus. Course notes, Section 2, printed p.~2, quantitative prime number theorem immediately before equation (1). The theorem is used as known; no proof of it is claimed here. \url{https://www.math.tau.ac.il/~rudnick/courses/sieves2015/fixed\%20progressions.pdf}
\item
  DP\_MULT checkpoint 06: DP22, DP25, DP44, DP49, DP52, DP65, each cited at statement level with the intake status retained. ALP-LIFT-04 additionally uses DP64 (CB1), with its explicit fixed exceptional-set convention and distinction from theta-selected values.
\end{itemize}

Web retrieval references for parent reuse: Yang turn507172355view2 and turn815074106view2; Wilms turn507172355view1 and turn815074106view1; Obus turn477052200view3; DLMF turn477052200view0 and turn477052200view2. The parent must itself open a source before citing its web reference. ALP-LIFT-04 primary retrievals: Rudnick turn953517703view0; DLMF functional equation/Euler product/nonvanishing turn953517703view1.

\subsection{ALP-LIFT-14: actual selected even-rank chain}\label{alp-lift-14-actual-selected-even-rank-chain}

Full proof: \nolinkurl{selected_even_chain_nonvanishing.md}, E1-E28. Starts from AM-R-012 on the common Eprime host, uses literal bilinear source convention and equal stage characters. Exact standard factors reduce to one absolutely convergent initial factor and nonzero angular Hecke factors at 1 through d; spherical local occurrence and compact real Schur tests give an actual nonzero chain with unrestricted later finite Schwartz vectors. Whole-proof author check logged 2026-11-21T22:39:18.688282+00:00. Raw Xi membership, later integral restricted-image matrices, six identities and DP44 costs remain open.

\subsection{ALP-LIFT-15: one finite unit-norm translate per edge}\label{alp-lift-15-one-finite-unit-norm-translate-per-edge}

Full proof: \nolinkurl{single_finite_heisenberg_chain.md}, H1-H10. Every later actual nonzero selected lift in LIFT14 can use one finite Heisenberg translate of the actual vacuum, norm exactly1 and coefficient1. A depth-N full window has an orthonormal basis of N\^{}(2g) such vectors and selection bound N\^{}g. Required depth and support are uncontrolled; the family is larger than DP59 and does not identify raw Xi. Author whole-proof PASS 2026-11-21T22:49:14.726742+00:00.

\subsection{ALP-LIFT-16: 29-only single translates and exact level masses}\label{alp-lift-16-29-only-single-translates-and-exact-level-masses}

Full proof: \nolinkurl{spherical_tests_and_29_only_chain.md}, S1-S27. Exact spherical-vacuum transforms, including residue2, refine LIFT15 to one norm-one translate supported only at29 on every later edge. A sufficient rank-n principal level is29\^{}(2e), with exact index I\_n=29\^{}(12 e n\^{}2) product\_(i=1)\^{}n(1-29\^{}(-i))\^{}6 and all stabilizers trivial. Its v11 is6(r+v11(r!)), r=floor(n/10), independent of unknown depth e. Every-vector and multiplicity details are retained. Whole-proof author PASS 2026-11-21T23:04:38.479472+00:00. Depth bounds, raw Xi and the restricted arithmetic image remain open. Finite arithmetic certificate: \nolinkurl{spherical_test_level_certificate.py}.

\subsection{ALP-LIFT-17: prescribed local column types exclude the spherical chain}\label{alp-lift-17-prescribed-local-column-types-exclude-the-spherical-chain}

Full proof: \nolinkurl{prescribed_column_local_type_obstruction.md}, F1-F19. For1\textless=j\textless m the actual colLev\_j has nontrivial determinant-block characters and occurs in no unramified local representation. Forj\textgreater=2, distinct block characters also exclude every tame determinant twist of an unramified representation. Thus the29-only chain cannot itself supply these explicit column modules at the prescribed G places. The variable ChIdx family and raw Xi remain separate. Any genuine source IsMu0 realization has conductor norm(g0*g1)\^{}12 on Eprime; global existence is not asserted. Author whole-proof PASS 2026-11-21T23:20:39.359309+00:00. Finite certificate: \nolinkurl{column_type_flag_certificate.py}.

\subsection{ALP-LIFT-18: exact literal-G even characters}\label{alp-lift-18-exact-literal-g-even-characters}

Full proof: \nolinkurl{literal_g_character_compatibility.md}, G1-G21. The source G half has Norm\_Eprime/Kq(zeta28)=1, so literal residue characters impose exactly sum\_(a in Phi) a e\_i,a=0 mod28; class-one norm-one torus descent proves necessity and sufficiency for IsMu0. Uniform e\_i divisible7 works. Explicit primes113/281, coefficient fieldQ(zeta560), conductor norm31753\^{}12, and numerical Hyp3 family k6,m2,t1795,C46539458987 are supplied. Odd etaT and subsequent prescribed local kernel/type nonvanishing remain open. Author whole-proof PASS 2026-11-21T23:30:15.183105+00:00. Exact certificate: \nolinkurl{literal_g_character_certificate.py}.

\subsection{ALP-LIFT-19: actual ramified theta constituents with prescribed nonterminal finite types}\label{alp-lift-19-actual-ramified-theta-constituents-with-prescribed-nonterminal-finite-types}

Full proof: \nolinkurl{ramified_column_type_chain.md}, T1-T24. At any degree-one split G place, an initial finite Steinberg type plus distinct literal tame stage characters gives irreducible actual local theta descendants and the exact inverse adelic column type at every j\textless m. Explicit finite Mackey vectors are primitive and have norm squared q\footnote{k(k-1)/2+d\emph{j}(j-1)/2+epsilon\emph{k}j} divided by the Gaussian multinomial index for(k,d\^{}j), with an exact11-valuation formula. DO033 supplies an old-host actual first input; LIFT14 extends analytically with unrestricted later finite Schwartz. Terminal HXm, raw Xi and global integral image remain open. Author whole-proof PASS 2026-11-21T23:49:36.760890+00:00. Exact certificate: \nolinkurl{ramified_column_type_certificate.py}. Independent DO033 review: \nolinkurl{review_steinberg_first_input.md}.

\subsection{ALP-LIFT-20: exact raw first G kernel and Steinberg extraction}\label{alp-lift-20-exact-raw-first-g-kernel-and-steinberg-extraction}

Full proof: \nolinkurl{raw_first_g_steinberg_bridge.md}, RG1-RG24. For m=2 and the actual allowed class y=q\^{}(-1)e\_k, both source PhiG/PhiGS give the primitive generic-unipotent residue function f\_U. Its full Steinberg component is (\textbar U\textbar\^{}2/\textbar GL\_k\textbar) f\_St, with exact compact-filter and first theta-extraction formulas. The raw finite Heisenberg coefficient absolute sum is1; raw norm squared is q\^{}(-bk)\textbar U\textbar{} and all11 averaging denominators are retained. Other raw T/Q factors, terminal PhiTop, later images and full DP44 costs remain separate. Author whole-proof PASS 2026-11-21T23:59:38.030210+00:00. Exact certificate: \nolinkurl{raw_first_g_steinberg_certificate.py}.

\subsection{ALP-LIFT-21: actual raw terminal G quotient preserves the preceding column}\label{alp-lift-21-actual-raw-terminal-g-quotient-preserves-the-preceding-column}

Full proof: \nolinkurl{raw_terminal_g_theta_quotient.md}, TG1-TG28. For m=2, the literal terminal PhiTop/ftY factor maps every nonzero vector in the actual LIFT19 preceding column into the actual small theta quotient. A compact full-rank integral and the primary Minguez boundary filtration prove the quotient identification. Exact identity value is vol(J)v (sharp vol(J)/q\^{}k); graph-plane norms are computed, and removing that volume gives a saturated injection of the natural induced11-lattices. The global automorphic scalar and exact row-lattice comparison remain open. Author whole-proof PASS 2026-11-22T00:26:12.923878+00:00. Exact finite diagnostic: \nolinkurl{raw_terminal_g_certificate.py}.

\subsection{ALP-LIFT-22: the two fixed sharp raw G edges compose on the actual middle level}\label{alp-lift-22-the-two-fixed-sharp-raw-g-edges-compose-on-the-actual-middle-level}

Full proof: \nolinkurl{composed_raw_sharp_g_chain.md}, CG1-CG21. Unwinding actual ChIdx.Ok forces the m2 sharp middle level Hsharp and its inverse character; midGSD uses this chF level, not generic ColV1. The first PhiGS forces the unique G class qy=e\_k. Its actual quotient image is the identity-parabolic section already in that middle type, and the fixed terminal sharp kernel preserves it. Their composition is injective on the initial generic type, with exact Haar scalar, composed norms, and saturated natural11-lattice image after scalar removal. Global quotient pins, other places/Y sum, and row-lattice comparison remain open. Author whole-proof PASS 2026-11-22T00:34:08.173546+00:00. Certificate: \nolinkurl{composed_raw_sharp_g_certificate.py}.

\subsection{ALP-LIFT-23: exact full coupled-Q Fourier contraction}\label{alp-lift-23-exact-full-coupled-q-fourier-contraction}

Full proof: raw\_coupled\_q\_fourier.md, Q1-Q20. At an unramified degree-one Q place the entire literal m2 Y sum collapses to the cofactor-supported Fourier function q\^{}(-d)G(chi,psi)chi(-c)\^{}(-1). The auxiliary depth cancels exactly. Joint Schwartz norm, primitive11 content, and scalar height (d-1/2)logq are explicit. This is a raw coupled-array identity; actual second theta-image survival and the global row scalar remain separate. Author whole-proof PASS 2026-11-22T00:47:32Z. Certificate: \nolinkurl{raw_coupled_q_fourier_certificate.py,10496} exact transforms.

\subsection{ALP-LIFT-24: actual fixed second T period on the first image}\label{alp-lift-24-actual-fixed-second-t-period-on-the-first-image}

Full proof: raw\_t\_projective\_period.md, MT1-MT33. The nontrivial projective-line source type is sigma\_spherical times one tame character. Whole-Schwartz GJ continuation gives the actual small theta quotient, and the raw T incidence vector has exact induced-model scalar L((d+1)/2,sigma1 eta1\^{}(-1))/(q\^{}n-1). Literal compact covariance transports the first T character into that same one-dimensional middle type. Every nonzero actual first T vector therefore survives the unchanged second T kernel. Haar denominator, L-factor valuation and global scalar comparison remain explicit. Author whole-proof PASS 2026-11-22T01:00:19Z. Certificate: \nolinkurl{raw_t_projective_period_certificate.py,58} compact vector checks and63958 Schur-complement checks.

\subsection{ALP-LIFT-25: actual selected compact highest kernels and their preceding-image contraction}\label{alp-lift-25-actual-selected-compact-highest-kernels-and-their-preceding-image-contraction}

Full proof: selected\_archimedean\_composition.md, AR1-AR26. The actual first highest polynomial outputs h1; the literal selected determinant insertion identifies the second highest source functional with h1 dual. The two compact contractions therefore preserve the actual matched line, including every multiplicity copy. Literal Haar coefficient factors are1/(d0 d1); the separate unitary quotient model gives squared factorA0²A1²/(d0 d1), with all Gaussian/CM/hook norms explicit. Exact11 Haar denominator valuations are1806 on K031 and14448 on K032. Global pins, nonhighest sharp witnesses and terminal row comparison remain separate. Author whole-proof PASS2026-11-22T01:14:48Z. \nolinkurl{Certificate:selected_archimedean_composition_certificate.py,810} exact compact projections plus polynomial and actual rank valuation checks.

\subsection{ALP-LIFT-26: full fixed sharp coupled two-edge automorphic nonvanishing}\label{alp-lift-26-full-fixed-sharp-coupled-two-edge-automorphic-nonvanishing}

Proof: \nolinkurl{fixed_global_coupled_second_lift.md}, GS1--GS31. Author whole-proof PASS 2026-11-22T01:30:35Z. On actual K031/K032 hosts and the unchanged R017 source, the complete sharp G/T/Q Y-coupled operator with LIFT25 selected highest real kernels is nonzero, retaining maximal vacua elsewhere. GQT is used in its anisotropic convergent case n=1802\textless m=3562/3436\textless2n. The exact partial standard factor is the initial degree-a Euler product at (d+1)/2, absolutely convergent with margin860/734, times Hecke L(l,chi0/chi1),1\textless=l\textless=d; the l1 boundary argument is the carried angular-type lemma. No global automorphic base change is assumed.

New assembly: local Howe multiplicity one and automorphic semisimplicity factor each actual global theta map as an injective multiplicity map tensor the product of local small quotients. Nonvanishing for every actual embedding makes that multiplicity map injective. The literal CRT Y quotient yields a tensor product of coupled local arrays. LIFT22, LIFT24, DO037 and LIFT25 give actual local left inverses on the initial types; LIFT16 supplies basic lines elsewhere. Distinct initial constituents cannot cancel after two edges by local Howe injectivity. The exact norm is a sum of positive constants c0*c1 times the squared local-tensor output norm; the constants remain the actual Rallis and quotient-mass costs.

Scope: full Y-sum nonzero, hence some admitted matched Y chain survives; the specific earlier Y0 is not asserted to survive. Nonzero G2 output does not prove its literal brC/tilQ/PtilFF restricted endpoint period, and the detecting Q graph lies outside that subgroup. Global pins, integral prior-image content, terminal branching detection, MainX mass and DP44 gain comparison remain open. The failed shortcut was inferring fixed-vector positivity directly from GQT11.6; the proof uses global factorization instead.

\subsection{ALP-LIFT-27: unchanged raw coupled operator is injective on the entire initial Mram}\label{alp-lift-27-unchanged-raw-coupled-operator-is-injective-on-the-entire-initial-mram}

Proof: full\_ramified\_raw\_coupled\_injectivity.md, MR1--MR33. Author whole-proof PASS 2026-11-22T02:01:09Z. This removes the finite Steinberg restriction at G by proving whole-Schwartz Godement--Jacquet holomorphy for every unitary initial representation and its actual first theta image. The same raw sharp G composition is injective on the entire prescribed inverse-generic space. The literal unweighted T norm kernel is the direct sum of all nontrivial projective dilation characters, each handled by the existing actual raw T theorem. The fixed full Y tensor consequently gives an injective full-G2 map on every actual initial Mram vector, including all multiplicities and rational stabilizer relations.

The exact complete-row coevaluation is d0\emph{d1 times the source-normalized highest-kernel operator, or d0}d1\emph{Lambda0}Lambda1 times its Chern version. Actual strict-covector/unipotent levels, stabilizer masses and weighted projective kernel averages are retained; coarsening does not make raw summands identical. K038's completed comparison applies after the literal terminal tilQ restriction and retains its own v0*v1 and period factors. Thus this does not prove injectivity/nonvanishing of that restricted arithmetic row or its brC/PtilFF period. It removes the need to identify a future actual bottom/Pu witness with the single R017 input. Uniform effective image costs and the actual bottom detector remain open.

\subsection{ALP-LIFT-28: fixed-level holomorphic bottom jets detect the full initial row space}\label{alp-lift-28-fixed-level-holomorphic-bottom-jets-detect-the-full-initial-row-space}

Proof: holomorphic\_bottom\_jet\_detector.md, HJ1--23. Author whole-proof PASS 2026-11-22T02:29:00+00:00. For either actual mixed CM host, the literal initial strict G/T/basic level extends to an explicit level on U(a,1), with inverse finite character for the bilinear period. Its compact U(a) packet is embedded with exactly the initial rational stabilizers, because the extra norm-one scalar is a root of unity forced to1 by the G unipotent residue level. The actual bottom HC list has selected initial K-type at the unique Pieri degree N\_tau=sum u\_tau-sum v\_tau+b0\_tau-a/2, independent of b1. Increasing all b1 by (a+1)ell changes the dual-minimal automorphic bundle by an ample canonical power and a flat determinant correction trivial on every arithmetic component. Serre vanishing on the Baily--Borel compactification therefore gives all compatible selected jets at the same finite level for sufficiently large ell in a fixed progression. PBW principal symbols transfer them to actual higher-K-type restrictions; their normalization factors are retained.

The actual cusp restriction is surjective onto the whole finite initial coefficient space, uniformly over any fixed finite word set. Every nonzero initial Mram row P thus has a bottom cusp test with exact period (v0/d\_lambda) sum\_x \textbar\textbar P\_x\textbar\textbar²/\textbar Gamma\_x\textbar. The test norm and sufficient weight are not bounded. The constructed detector has not been identified with source rJ/rChi/RealPack/Pu, and its cusp test has not been placed in the prescribed downward theta image. Consequently no literal MainX endpoint or DP44 cost comparison is claimed. The argument permits narrow-cloud interlacing profiles, but does not fit its ineffective weight increase into full SizeData2 windows.

\subsection{ALP-LIFT-29: exact source-shaped bottom level and T translation}\label{alp-lift-29-exact-source-shaped-bottom-level-and-t-translation}

Proof: bottom\_dat0\_translation\_detector.md, BT1--14. Author whole-proof PASS 2026-11-22T02:34:29+00:00. The actual rS=0 source PsharpSym(r=a+1,0,a) is exactly diag(U\_a,1), with generic U\_a character; dat0 has maximal compacts everywhere else. At T the explicit lower unipotent tau={[}{[}I,0{]},{[}varpi\^{}-1 u,1{]}{]} conjugates the maximal bottom compact so its intersection with H is exactly the original strict covector stabilizer. The translated packet x*tau has precisely the original classes/stabilizers. LIFT28 therefore gives large-weight cusp detection at this standard-frame source-shaped level, with no added T refinement and the unchanged initial mass. Translation is an ambient Petersson isometry, not an integral-lattice/height assertion.

The prescribed tS and betaY0 pins, actual two vertical kernels, and same-family downward image have not been identified with this concrete candidate. The new proof matches the finite level formula and initial-level intersection, not a completed RealPack or MainX endpoint.

\subsection{ALP-LIFT-30: primitive negative line and the finite/real bottom transport}\label{alp-lift-30-primitive-negative-line-and-the-finitereal-bottom-transport}

Proof: primitive\_negative\_line\_and\_bottom\_pin.md, NL1--14. Author whole-proof PASS 2026-11-22T02:48:20Z. The actual DO040 norm generator pi\emph{pi\_bar=T gives an explicit rational SU(a,1) matrix on a positive/negative plane, with columns (A,B)/pi and (B,A)/pi\_bar, A=(T+1)/2 and B=(T−1)/2. It sends pi\_bar times the negative coordinate to the primitive integral vector (B e1,A) of norm minus T. At T its inverse factors exactly as the LIFT29 lower unipotent for covector pi\_bar e1}, times an integral determinant-one compact element; away T it is integral and at G it is identity modulo the level. An arbitrary actual primitive T covector is obtained by an explicit integral finite H-coordinate transport, with unchanged strict class masses and Mram.

The actual identity of automorphic restrictions also transports the real vector by sigma(g\_infinity); omitting that factor would be false. Its Hilbert norm is unchanged, while the ambient matrix has singular values sqrt(T),1/sqrt(T), so algebraic height or coefficient costs are not set to zero. At 11 this matrix and inverse are integral. Exact Q(sqrt(-7)) replay for both actual DO040 primes passed24 matrix/congruence checks. Source audit distinguishes the full row compact pin tS from this Bessel embedding: rowK2 is maximal away G, and RowLevelMatchB/C further fixes the full residue action and spectral idempotent. None of these source ties, the idempotent-selected detector, or the fixed downward image is asserted complete.

\clearpage

\section{ALP-LIFT-05: stable Hodge lattices and a normalized induced CM eigenbasis}\label{alp-lift-05-stable-hodge-lattices-and-a-normalized-induced-cm-eigenbasis}

Run ALP\_MACHINE; authors /root/lift\_accounting and /root, GPT-6 Astra Pro. Reasoning effort is not exposed. The construction at 7 and the positive trace polarization were suggested independently by /root and checked here. This is an actual geometric normalization theorem. It does not identify an arithmetic theta lift with an automorphic lift, or determine a SizeData₂ gain or division-error field.

\subsection{Statement}\label{statement}

Write \(\zeta=\zeta_7\), \(E=\mathbb Q(\zeta)\), \(K=\mathbb Q(\sqrt{-7})\), and \[
 C:y^2=x^7-1,\qquad J=\operatorname{Jac}(C),\qquad
 \omega_i=x^{i-1}dx/y\quad(1\le i\le3).
 \tag{H1}
\] For a finite place above a rational prime \(p\), use the valuation normalized by \(v_p(p)=1\). Fix algebraic numbers \[
 \iota^2=-1,\quad \alpha^7=-4,\quad
 \lambda=\zeta-1,\quad \beta^2=\lambda^7,\quad \rho^4=7.
 \tag{H2}
\] After adjoining these numbers, the following assertions hold.

\textbf{(a) Actual stable Hodge lattices.} At a place above 2, the good model and its integral differential basis are \[
 Y^2+Y=X^7,\qquad
 \eta_i^{(2)}=X^{i-1}dX/(2Y+1),\qquad
 \omega_i=\alpha^i\iota^{-1}\eta_i^{(2)}.
 \tag{H3}
\] At a place above 7, they are \[
 Y^2=f_\lambda(X):=\frac{(1+\lambda X)^7-1}{\lambda^7},\qquad
 \eta_j^{(7)}=X^{j-1}dX/Y,
 \tag{H4}
\] with \[
 \omega_i=\frac{\lambda}{\beta}
       \sum_{k=0}^{i-1}\binom{i-1}{k}\lambda^k\eta_{k+1}^{(7)}.
 \tag{H5}
\] Here and below an integral differential basis means a basis of \(e^*\Omega^1_{\mathcal J/R}\), identified through the Abel map with \(H^0(\mathcal C,\Omega^1_{\mathcal C/R})\), for the good model over the local valuation ring \(R\). Consequently the signed determinant valuations of \(\omega_1\wedge\omega_2\wedge\omega_3\) are exactly \[
 \begin{array}{c|c}
 p&v_p(\omega_1\wedge\omega_2\wedge\omega_3)
       \text{ relative to an integral generator}\\ \hline
 2&12/7\\
 7&-3/4\\
 p\nmid14&0.
 \end{array}
 \tag{H6}
\] The negative entry is a valuation of a rational section, not a positive index of an inclusion of integral lattices.

\textbf{(b) Complete determinant metric and height.} Give the determinant Hodge line the normalization \[
 \|s\|_{L^2}^{\,2}
   =\frac{\mathrm i^{g^2}}{2^g}\int_{A(\mathbb C)}s\wedge\overline{s}
 \quad(g=\dim A).
 \tag{H7}
\] Let \(G=\operatorname{diag}(G_1,G_2,G_3)\) be the Gram matrix of (H1), already computed in ALP-LIFT-03: \[
 G_i=\frac{2\pi}{7}
 \frac{\Gamma(i/7)\Gamma(1/2-i/7)}
      {\Gamma(1-i/7)\Gamma(1/2+i/7)}.
 \tag{H8}
\] Then \[
 \boxed{h_{F,L^2}(J)=\frac{12}{7}\log2-\frac34\log7
                       -\frac12\log\det G.}
 \tag{H9}
\] This equals the Wilms normalization used in ALP-LIFT-02. Numerically it is \(-3.80116074543404940138079328403117913069\ldots\); the decimal is a cross-check of the exact expression.

\textbf{(c) An actual comparison variety and compatible principal polarization.} Take the elliptic curve of full CM by \(\mathcal O_K\) \[
 B:y^2+xy=x^3-x^2-2x-1,\qquad \psi=dx/(2y+x),
 \tag{H10}
\] and put \(G_B=(\mathrm i/2)\int_{B(\mathbb C)}\psi\wedge\overline\psi\), where \(\mathrm i\) is the usual positive imaginary unit in each complex metric, not the image of the algebraic symbol \(\iota\) under an arbitrary embedding. This is LMFDB 49.a4, Cremona 49a1; the full CM discriminant \(-7\) is the only catalogued fact used about it. Its Weierstrass coefficients, discriminant, local changes of variables and all following metrics are computed below.

The basis \(1,\zeta,\zeta^2\) identifies the Serre tensor variety \(A_{\mathrm{ind}}=\mathcal O_E\otimes_{\mathcal O_K}B\) with \(B^3\) as in ALP-LIFT-02. Put \[
 S=\zeta+\zeta^2+\zeta^4=(-1+\sqrt{-7})/2,\qquad
 u=-(\zeta-\zeta^{-1})^{-2}.
\] The matrix of the positive Hermitian trace form \(\operatorname{Tr}_{E/K}(u\overline x y)\) in that basis is \[
 H=\begin{pmatrix}
 2&-1&1+S\\
 -1&2&-1\\
 -S&-1&2
 \end{pmatrix},\qquad \det H=1.
 \tag{H11}
\] It defines an \(\mathcal O_E\)-compatible principal polarization on \(A_{\mathrm{ind}}\). The product polarization is not silently used as that compatible polarization.

Let \(\psi_a\), \(a=0,1,2\), denote the product invariant differentials on \(B^3\). For \(i\in\Psi=\{1,2,4\}\), set \[
 \vartheta_i=(\zeta^i-\zeta^{-i})^{-1}
                     \sum_{a=0}^{2}\zeta^{ia}\psi_a.
 \tag{H12}
\] These are CM eigenforms of type \(\Psi\). Their change-of-basis matrix \(W\), with rows in the order \(1,2,4\), has determinant \(-1\), and both \(W\) and \(W^{-1}\) have entries integral away from 7. They are therefore a primitive integral CM eigenbasis at every prime \(p\nmid7\). For the principal polarization (H11), their Hodge Gram matrix is exactly \[
                         G_B I_3.
 \tag{H13}
\] Their determinant section has normalized valuation \(-3/4\) at 7 and zero at every other prime, even though their individual lattice at 7 is not asserted to be the primitive Hodge lattice. In particular \[
 3h_{F,L^2}(B)=-\frac34\log7-\frac32\log G_B.
 \tag{H14}
\]

\textbf{(d) Exact comparison in these CM coordinates.} With \(\delta_7\) from ALP-LIFT-02 and \(\delta_{7,\{2,7\}}=\delta_7-(12/7)\log2\) from ALP-LIFT-04, \[
 \boxed{\delta_{7,\{2,7\}}
       =-\frac12\log\frac{G_1G_2G_3}{G_B^3}.}
 \tag{H15}
\] Thus the order-linear difference of the equal-direction Cauchy bounds in these actual diagonal CM coordinates is exactly the order-linear difference of the bare prime certificates over \(p\nmid14\). The precise analytic statement and its scope are given after the proof.

\subsection{Proof of (a): the good model at 2}\label{proof-of-a-the-good-model-at-2}

Make the change \(x=\alpha X\), \(y=\iota(2Y+1)\). Since \(\alpha^7=-4\), the original equation becomes exactly \(Y^2+Y=X^7\). Complete this affine chart by the chart \[
 U=X^{-1},\qquad V=YU^4,\qquad V^2+U^4V=U.
 \tag{H16}
\] On their common locus the changes are inverse to one another. Both charts are finite free of rank two over their respective affine charts of \(\mathbb P^1\), with bases \(1,Y\) and \(1,V\). They therefore form a finite flat double cover of \(\mathbb P^1\), which is proper. The algebra generators transform by \(V=U^4Y\), so the underlying rank-two pushforward is \(\mathcal O\oplus\mathcal O(-4)\); its Euler characteristic is \(1-3=-2\), giving arithmetic genus three in every fiber.

In residue characteristic 2 the finite equation has derivative 1 with respect to \(Y\). In (H16) the derivative with respect to \(U\) of \(V^2+U^4V-U\) is \(-1\). Thus the whole special fiber is smooth, including the point at infinity \((U,V)=(0,0)\). The special fiber is geometrically integral: if \(Y^2+Y-X^7\) were reducible over \(\overline k(X)\), then \(X^7=h^2+h\) for a rational function \(h\). Any finite pole of \(h\) would produce a finite pole of \(h^2+h\), and the pole at infinity of a nonconstant polynomial \(h^2+h\) has even order. Both alternatives contradict \(X^7\). The generic fiber is the original smooth curve. This proves that (H16) is a smooth proper genus-three model with an infinity section over the local valuation ring, and remains so after any finite base extension.

The residue differential on a plane curve \(F(X,Y)=0\) is \(dX/F_Y=-dY/F_X\) on the respective smooth open subsets. Thus \(\eta_i^{(2)}\) is regular on the finite chart. On the other chart it is \[
                  -U^{3-i}\frac{dU}{2V+U^4}.
 \tag{H17}
\] The same residue description shows regularity at infinity: when \(2V+U^4\) is not a unit, the alternate denominator \(4U^3V-1\) is a unit. The power \(U^{3-i}\) is nonnegative for all three indices. Reduction on the finite chart gives \(X^{i-1}dX\). These three differentials are independent over the residue field, since the Artin--Schreier function-field extension is separable and \(dX\ne0\).

For completeness, their being an integral basis follows without any unproved saturation assumption. They are a basis on the generic fiber because the genus is three. Express any global regular differential as \(\sum a_i\eta_i^{(2)}\) with coefficients in the fraction field. If some coefficient has negative valuation, multiply by the smallest positive power of a uniformizer that makes all coefficients integral and at least one a unit. This product, as a multiple of the original integral differential, reduces to zero, contradicting the independence just proved. Hence all coefficients are integral. This also proves the claim after further ramified base change.

The infinity section supplies the Abel map to the relative Jacobian. The smooth proper curve has an abelian relative Picard scheme, and the Abel map identifies its invariant differentials with those of the curve. The precise references are Milne (2021), Proposition 2.2, pp.~6--7, and Theorem 8.1 with the following paragraph, pp.~27--28; the relative assertion follows also by applying the field assertion on every fiber and Nakayama's lemma to the rank-three modules. This is the same FGA input already accepted in DP22, not a new Picard representability claim.

Finally, direct substitution gives (H3). Since \(v_2(\alpha)=2/7\), the determinant valuation is \((1+2+3)(2/7)=12/7\).

\subsection{Proof of (a): the good model at 7 and the remaining primes}\label{proof-of-a-the-good-model-at-7-and-the-remaining-primes}

The local cyclotomic uniformizer has \(v_7(\lambda)=1/6\). Write \(f_\lambda(X)=\sum_{k=1}^{7}f_kX^k\). Its top coefficient is 1, and for \(1\le k\le6\), \[
 f_k=\binom7k\lambda^{k-7},\qquad v_\lambda(f_k)=k-1.
\] Here \(v_\lambda(7)=6\) and the binomial coefficients have exactly one factor of 7. The cyclotomic identity \(\Phi_7(1+\lambda)=0\) can be rearranged as \[
 \frac7{\lambda^6}
 (1+3\lambda+5\lambda^2+5\lambda^3+3\lambda^4+\lambda^5)=-1.
\] It follows that \(f_1\equiv-1\), while \(f_2,\ldots,f_6\equiv0\). The reduction of (H4) is therefore \(Y^2=X^7-X\). Its affine chart is smooth because \(d(X^7-X)/dX=-1\) in characteristic 7. The polynomial has simple roots and is not a square, so this special fiber is geometrically integral.

The infinity chart is \[
 U=X^{-1},\quad V=YU^4,\qquad
 V^2=U+f_6U^2+f_5U^3+\cdots+f_1U^7.
 \tag{H18}
\] Away from \(U=0\) it is already covered by the finite chart. At \(U=0\) its only point is \(V=0\), and its equation has derivative \(-1\) with respect to \(U\). Thus all of this fiber is smooth. As before, the two monic charts give a proper finite flat double cover with pushforward \(\mathcal O\oplus\mathcal O(-4)\), genus three and an infinity section. Its generic fiber is the original curve by \(x=1+\lambda X\), \(y=\beta Y\).

The three differentials in (H4) are regular, since 2 is a unit and \(dX/Y\) is twice the plane residue differential. On the infinity chart they become \(-U^{3-j}dU/V\); again this is a unit multiple of the residue differential times a nonnegative power of \(U\). Their reductions are independent on the separable hyperelliptic function field. The saturation argument above proves that they form the full integral Hodge basis. Substitution of the change of variables gives (H5). Its triangular diagonal entries are \(\lambda^i/\beta\), so its determinant is \[
                      \lambda^6/\beta^3.
\] Its normalized valuation is \(6/6-3(7/12)=-3/4\), proving the second row of (H6). An explicit inverse formula, useful for the lattice rather than only its determinant, is \[
 \eta_j^{(7)}=\beta\lambda^{-j}
    \sum_{a=0}^{j-1}\binom{j-1}{a}(-1)^{j-1-a}\omega_{a+1}.
 \tag{H19}
\]

At a prime not dividing 14, the original two-chart hyperelliptic model is smooth: 2 is a unit and its branch polynomial has squarefree reduction. The same residue and independence argument gives precisely the basis (H1), proving the last row of (H6). These local models can also be made global over a common number field containing (H2). Model (H3) is smooth over \(\mathcal O_M[1/7]\); model (H4) is smooth over \(\mathcal O_M[1/2]\). On their overlap the specified changes of variables give an isomorphism, preserving infinity. Gluing yields a smooth proper curve over \(\mathcal O_M\) and hence an abelian scheme with generic fiber \(J_M\).

\subsection{Proof of (b): the determinant degree}\label{proof-of-b-the-determinant-degree}

The Hodge line is \(\det e^*\Omega^1_{\mathcal J/\mathcal O_M}\). Use its nonzero rational section \(s=\omega_1\wedge\omega_2\wedge\omega_3\). By definition of the degree of a metrized line, \[
 \frac{\widehat{\deg}(\overline\omega_{\mathcal J})}{[M:\mathbb Q]}
 =\frac1{[M:\mathbb Q]}
   \left(\sum_{v\text{ finite}}\operatorname{ord}_v(s)\log Nv
              -\sum_{\sigma:M\hookrightarrow\mathbb C}
                           \log\|s\|_\sigma\right).
 \tag{H20}
\] The local orders from (H6) give exactly \((12/7)\log2-(3/4)\log7\) after division by the degree: \(\sum_{v\mid p}e_vf_v=[M:\mathbb Q]\). The original curve and the three differentials are rational, so their complex metric is the same under every \(\sigma\). The Riemann bilinear relation identifies (H7) with \(\det G\) on this section; this is precisely the determinant convention fixed in ALP-LIFT-03. Thus the archimedean term is \(-\tfrac12\log\det G\), proving (H9). It is the stable height because the chosen model is an abelian scheme, and further base extension leaves the normalized degree unchanged. This proof does not use isogeny invariance of height.

\subsection{Proof of (c): the comparison lattice}\label{proof-of-c-the-comparison-lattice}

The Weierstrass coefficients of (H10) give \(c_4=105\), \(c_6=1323\), and discriminant \(-343\); these follow by the usual polynomial discriminant identities and can also be checked by completing the square as follows. Put \(y'=y+x/2\), \(z=x-1/4\). The equation becomes \[
                   y'^2=z^3-\frac{35}{16}z-\frac{49}{32}.
\] At 7 make the substitution \(z=\rho^2 Z\), \(y'=\rho^3 T\), where \(\rho^4=7\). The resulting equation is \[
                  T^2=Z^3-\frac5{16}Z-\frac{\rho^2}{32}.
 \tag{H21}
\] Its discriminant is exactly \(-1\), either by dividing the old discriminant by \(\rho^{12}=7^3\), or by direct substitution in \(-16(4a^3+27b^2)\). It is a smooth proper elliptic model; its special fiber is \(T^2=Z^3+Z\). The differential transforms as \[
                        \psi=\rho^{-1}\frac{dZ}{2T}.
 \tag{H22}
\] The invariant differential on the right after removing \(\rho^{-1}\) is primitive, so \(v_7(\psi)=-1/4\). At every prime different from 7 the original integral Weierstrass model has unit discriminant, and \(\psi\) is primitive. Applying (H20) in dimension one proves \(h_{F,L^2}(B)=-(1/4)\log7-(1/2)\log G_B\), and thus (H14). The value of \(G_B\) is approximately \(4.9445046002825467364981969682\); this agrees with the independent LMFDB period-area entry, but the proof of (H14) does not use that numerical entry.

\subsection{Proof of (c): the compatible polarization and eigenforms}\label{proof-of-c-the-compatible-polarization-and-eigenforms}

The integral basis \(1,\zeta,\zeta^2\) was proved in ALP-LIFT-02; it gives the specified product coordinates. Each real embedding of \(E^+\) sends \(u=-(\zeta-\zeta^{-1})^{-2}\) to \(1/(4\sin^2(2\pi i/7))>0\). Direct cyclotomic reduction gives \[
 u=\frac{6+3\zeta^2+\zeta^3+\zeta^4+3\zeta^5}{7}.
\] Taking the trace over the embeddings \(\zeta\mapsto\zeta^i\), \(i=1,2,4\), gives (H11). Its entries belong to \(\mathcal O_K\); its leading principal minors are \(2,3,1\), which proves both positive definiteness and determinant one. Multiplication by an element of \(E\) is adjoint to multiplication by its complex conjugate for this trace form. In particular, multiplication by \(\zeta\) preserves it.

Let \(\lambda_0\) be the product principal polarization on \(B^3\). The integral Hermitian matrix \(H\) is an endomorphism of that product; the homomorphism \(\lambda_H=\lambda_0\circ H\) has integral positive Riemann form, and is therefore a polarization. This standard polarization criterion is the Riemann-form criterion in Milne (2022), Section 16, p.~31. Since \(H^{-1}\) has integral entries by the adjugate formula and \(\det H=1\), \(H\) is an automorphism, so \(\lambda_H\) is principal. Its Rosati involution restricts to conjugation on \(E\) by the trace identity just proved. Thus it is the required compatible principal polarization. The polarization homomorphism is defined over a number field after the CM endomorphisms are; an ample line realizing it can be chosen after a further finite extension if necessary. No special theta trivialization is asserted at this step.

For the metric calculation, let \[
 V_{i,a}=\zeta^{ia}\quad(i=1,2,4;\ a=0,1,2),\qquad
 U=\operatorname{diag}\bigl(\sigma_i(u)\bigr).
\] Then \[
                         H=V^*UV.
 \tag{H23}
\] In product differential coordinates, the Hodge Gram for \(\lambda_H\) is \(G_B H^{-1}\). One may check this directly: the normalized product Kähler form has coefficient matrix \(G_B^{-1}H^t\); integrating \((\mathrm i/2)\psi_a\wedge\overline\psi_b\) against its second power divided by \(2!\) gives the corresponding cofactor, hence \(G_B(\det H)H^{-1}\). Here \(\det H=1\). This computation fixes the transpose convention and does not identify covectors with tangent vectors without taking the inverse metric.

The rows of \(V\) are eigen-covectors for multiplication by \(\zeta\), with eigenvalues \(\zeta,\zeta^2,\zeta^4\). Indeed evaluation at \(\zeta^i\) sends multiplication by \(\zeta\) on \(E\) to multiplication by \(\zeta^i\). It follows from (H23) that \[
 V(G_BH^{-1})V^*=G_B U^{-1}
     =G_B\operatorname{diag}(|\zeta^i-\zeta^{-i}|^2).
 \tag{H24}
\] Dividing the corresponding eigenforms by \(\zeta^i-\zeta^{-i}\) therefore proves (H13).

The Vandermonde determinant and the cyclotomic product are \[
 \det V=-\sqrt{-7},\qquad
 \prod_{i=1,2,4}(\zeta^i-\zeta^{-i})=\sqrt{-7},
\] with the same square-root convention \(2S+1=\sqrt{-7}\). Their quotient gives \(\det W=-1\). Each denominator is a unit away from 7, so all entries of \(W\) are integral there; its adjugate gives the same conclusion for its inverse. This proves primitivity at those primes. At 7 the determinant section equals \(-\psi_0\wedge\psi_1\wedge\psi_2\), exactly, so its valuation is still \(-3/4\) by (H22), regardless of the nonintegral entries of \(W\). There has been no assertion that a general change of coordinates transforms a product of derivatives by a determinant. The CM eigenbasis (H12) is chosen first, and the derivatives are then its dual coordinate directions.

Finally subtract (H14) from (H9). ALP-LIFT-02 identifies the left side with \(\delta_7\), and the finite 7-adic contributions cancel exactly. This gives \[
 \delta_7=\frac{12}{7}\log2-
                    \frac12\log\frac{\det G}{G_B^3},
\] which is (H15) and completes the proof. \(\square\)

\subsection{Exact differential rescaling and the 2-adic limitation}\label{exact-differential-rescaling-and-the-2-adic-limitation}

Let \(D_i\) be the invariant derivations dual to \(\omega_i\), and let \(E_i\) be dual to \(\eta_i^{(2)}\). Formula (H3) gives the exact operator identities \[
 D_i=\iota\alpha^{-i}E_i,\qquad
 D_1^{n_1}D_2^{n_2}D_3^{n_3}
  =\iota^{n_1+n_2+n_3}\alpha^{-(n_1+2n_2+3n_3)}
                         E_1^{n_1}E_2^{n_2}E_3^{n_3}.
 \tag{H25}
\] The associated valuation shift is \(-\sum_i(2i/7)n_i\). The same identity holds for covariant derivatives when the connection and local frame are transported along the same isomorphism: connections are linear in a vector field, the scaling factors are constant, and the individual operators are composed in the displayed order. This is a statement about the same section and connection, not a comparison of two independently normalized theta frames.

For the abstract DP64 slope certificate at 2, the three directions of \(\Phi=\{1,2,3\}\) have slopes respectively \(2/7,4/7,6/7\); the induced type \(\Psi\) has slope zero in each direction. These numerical slopes equal the valuations of the coefficients in (H3). The equality gives a concrete geometric interpretation of the omitted \((12/7)\log2\) in ALP-LIFT-04. It does not extend the accepted good-model theta theorem to 2. In particular the trace-one endomorphism \(w_0=-(\zeta+\zeta^2+\zeta^4)\) has \(w_0+w_0^\dagger=1\) but \(w_0w_0^\dagger=2\), so one of its 2-adic conjugates is a nonunit. The odd-prime unit-path hypotheses used in DP49 do not follow there. The primitive theta frame, gauge and selected gain still need their own comparison.

\subsection{Precise equal-direction Cauchy consequence}\label{precise-equal-direction-cauchy-consequence}

This paragraph states exactly what part of a theta estimate is fixed by the calculation. Let \(F_J\) and \(F_A\) be holomorphic functions in the universal-cover coordinates dual to \(\omega_i\) and \(\vartheta_i\), respectively, satisfying \[
 |F_J(z)|\le C_J\exp\!\left(\frac\pi2
                       \sum_{i=1}^3|z_i|^2/G_i\right),\qquad
 |F_A(t)|\le C_A\exp\!\left(\frac\pi{2G_B}
                       \sum_{i=1}^3|t_i|^2\right)
 \tag{H26}
\] for positive constants independent of the derivative order. These are the Chern-frame growth inequalities for bounded-norm sections of the respective principal theta lines; /root's AM-R004 constructs the specific analytic germ on \(J\) and supplies its constant separately. Define the logarithmic Cauchy upper bounds, with equal order \(n\ge1\) in the three chosen CM directions, by \[
 \mathcal C_J(n)=\log C_J+3\log(n!)+
               \frac{3n}{2}\log(\pi e/n)-\frac n2\log\det G,
\] and the same expression with \(C_A\) and \(G_B^3\) for \(\mathcal C_A(n)\). Cauchy's formula on a polydisk of radii \(R_i^2=nG_i/\pi\) proves the bound for \(J\); replacing the three \(G_i\) by \(G_B\) proves the other. Subtraction and (H15) give the exact difference of these \textbf{upper-bound expressions}: \[
 \mathcal C_J(n)-\mathcal C_A(n)
     =\log(C_J/C_A)+n\delta_{7,\{2,7\}}.
 \tag{H27}
\] Therefore ALP-LIFT-04 implies, for its bare integer certificate \(D_{\{2,7\}}(n;X)\) and every joint limit \(n\to\infty\), \(X\to\infty\), \[
 \frac{D_{\{2,7\}}(n;X)
              -\bigl(\mathcal C_J(n)-\mathcal C_A(n)\bigr)}n
                         \longrightarrow0.
 \tag{H28}
\] The same identity holds after multiplying both displayed budgets by the source's common factor 2. A tensor level \(\kappa\) in the exponent replaces \(\pi\) by \(\kappa\pi\); the added \((3n/2)\log\kappa\) cancels between the two bounds.

Neither (H27) nor (H28) is an equality between actual derivative values. The constants can change under arithmetic normalization of a theta section, and no primitive integral frame or selected automorphic row is identified by choosing an analytic sup norm. Thus this closes the order-linear \textbf{bare certificate versus specified diagonal Cauchy metric} comparison in actual CM coordinates, while leaving the source gain, division error, bad-prime theta normalization, row content, and kernel matching as separate inputs to DP44. It gives no shorter valid lift chain and no universal sign assertion for its full cost.

\subsection{Grothendieck and source scope}\label{grothendieck-and-source-scope}

The smooth models above give an actual relative Picard scheme and its Hodge bundle, not merely a characteristic-zero CM type. The Picard representability and smoothness locators are those already accepted in DP22: FGA V, Bourbaki 232, Theorem 3.1, pp.~148--149; FGA VI, Bourbaki 236, Corollary 1.7, p.~229, Corollary 2.3, p.~232, and Proposition 2.10(ii), p.~236. Those source results are used with their accepted status, not re-proved or claimed newly scan-checked. The Hodge line is the determinant of the degree-one de Rham realization of this Picard scheme. Its metric and finite lattices are evaluated explicitly in (H6)--(H15).

The principal polarization is the homomorphism to the dual Picard scheme determined by its Poincaré biextension. The trace identity in (H11) makes the \(E\)-action and its conjugate adjoint, so it is compatible with that biextension. This is the concrete content relevant to SGA 7 I, Exposés VII--VIII and the previously accepted DP49 trace-one construction; no further comparison theorem is ascribed to those references. No new claim about a lambda-ring, an unrelated K-theory class, or weights of an automorphic sheaf is needed for this Hodge determinant calculation.

\subsection{References and exact retrievals}\label{references-and-exact-retrievals}

\begin{itemize}
\tightlist
\item
  Milne, J. S. (2021, June 12). \emph{Jacobian varieties}. Proposition 2.2, printed pp.~6--7; Theorem 8.1 and following paragraph, printed pp.~27--28. \url{https://www.jmilne.org/math/xnotes/JVs.pdf} Primary text opened this run; current retrievals turn963921859view0, turn794617041view1, turn794617041view2.
\item
  Milne, J. S. (2022, January 2). \emph{Abelian varieties}. Section 16, printed p.~31, nondegenerate positive Riemann-form criterion; Section 13, printed p.~25, principal-polarization convention. \url{https://www.jmilne.org/math/xnotes/AVs.pdf} Current retrievals turn794617041view0 and turn892289938view0.
\item
  The LMFDB Collaboration. (2026). \emph{Elliptic curve 49.a4: Data}. The ec\_curvedata row identifies coefficients {[}1,−1,0,−2,−1{]}, CM discriminant −7 and Cremona label 49a1; the ec\_mwbsd row contains the independent period-area decimal. Retrieved November 21, 2026, from \url{https://www.lmfdb.org/EllipticCurve/Q/data/49.a4} Current retrieval turn963921859view1.
\item
  ALP-LIFT-02, ALP-LIFT-03 and ALP-LIFT-04 in \nolinkurl{work/agents/lift_accounting.md}, with their exact cited Colmez--Obus, Wilms, DLMF and fixed-modulus prime-theorem inputs.
\item
  DP\_MULT DP22, DP49 and DP64, with intake status unchanged. The actual theta gain is not identified with the abstract DP64 floor certificate, at 2 or at any other prime.
\end{itemize}

\clearpage

\section{ALP-LIFT-06: directional prime constants and unequal-order Cauchy cancellation}\label{alp-lift-06-directional-prime-constants-and-unequal-order-cauchy-cancellation}

ALP\_MACHINE; /root/lift\_accounting, GPT-6 Astra Pro, with the directional question proposed by /root. Reasoning effort not exposed. This refines ALP-LIFT-04 using the actual CM eigenbases and polarization constructed in ALP-LIFT-05. It concerns the bare DP64 certificate and a specified Chern metric. It does not identify actual theta-selected gain or source division error.

\subsection{Statement}\label{statement-1}

Use \(\Phi=\{1,2,3\}\), \(\Psi=\{1,2,4\}\) in \((\mathbb Z/7\mathbb Z)^\times\). For \(i=1,2,3\), define \[
 b_i(a)=\mathbf1_{ai\notin\Phi}-\mathbf1_{a\notin\Psi},
 \qquad
 d_i(p)=\sum_{r\ge1}\frac{b_i(p^r)}{p^r}\quad(p\ne7).
 \tag{D1}
\] The second indicator is the same for every selected direction of \(\Psi\), since \(\Psi\) is a subgroup. Thus it is the appropriate comparison for any pairing of the three chosen \(\Phi\)-directions with the three \(\Psi\)-directions. Let \(\ell_a=\log\Gamma(a/7)\) and \[
 Q=\ell_1+\ell_2-\ell_3+\ell_4-\ell_5-\ell_6.
\]

\textbf{(a) Directional constants in ordinary prime order.} The limits \[
 C_i=\lim_{X\longrightarrow\infty}
                  \sum_{\substack{p\le X\\p\ne7}}d_i(p)\log p
 \tag{D2}
\] exist, with the primes in their ordinary increasing order, and satisfy \[
 \boxed{C_i=\frac12\bigl(Q-2\ell_i+2\ell_{7-i}
                              -\ell_{7-2i}+\ell_{2i}\bigr).}
 \tag{D3}
\] In particular \[
\begin{split}
 C_1&=-\tfrac12\ell_1+\ell_2-\tfrac12\ell_3
                         +\tfrac12\ell_4-\ell_5+\tfrac12\ell_6,\\
 C_2&=\tfrac12\ell_1-\tfrac12\ell_2-\ell_3
                         +\ell_4+\tfrac12\ell_5-\tfrac12\ell_6,\\
 C_3&=\tfrac12\ell_2-\tfrac32\ell_3
                         +\tfrac32\ell_4-\tfrac12\ell_5
      =C_1+C_2=\tfrac12\delta_7.
\end{split}
 \tag{D4}
\] Let \(G_i\), \(G_B\) be the exact metrics in ALP-LIFT-05. Then \[
 \boxed{C_i-\frac{2i}{7}\log2
               =-\frac12\log(G_i/G_B)\qquad(i=1,2,3).}
 \tag{D5}
\] The subtraction removes exactly the formal prime-2 contribution; it leaves the actual common good-prime range \(p\nmid14\).

\textbf{(b) Integer certificates with an arbitrary growing cutoff.} For an integer \(n\ge0\), real \(X\ge2\), and fixed finite set of rational primes \(S\) containing 7, put \[
 D_{i,S}(n;X)=\sum_{\substack{p\le X\\p\notin S}}\log p
                     \sum_{r\ge1}b_i(p^r)\left\lfloor n/p^r\right\rfloor,
 \qquad
 C_{i,S}=C_i-\sum_{p\in S\setminus\{7\}}d_i(p)\log p.
 \tag{D6}
\] There is \(c>0\) such that, uniformly in \(n\ge4\), \(X\ge2\), \[
 D_{i,S}(n;X)=nC_{i,S}
    +O_S\!\left(n\exp\{-c\sqrt{\log\min(X,\sqrt n)}\}\right).
 \tag{D7}
\] The constant can be chosen for all three directions at once. Consequently, for any vector of nonnegative integer orders \(\boldsymbol n=(n_1,n_2,n_3)\), any common cutoff \(X\), and \(N=n_1+n_2+n_3\), \[
 \sum_{i=1}^3D_{i,\{2,7\}}(n_i;X)
       =-\frac12\sum_{i=1}^3n_i\log(G_i/G_B)+o(N)
 \tag{D8}
\] in every joint limit \(N\to\infty\), \(X\to\infty\). No fixed positive limiting ratio between the three orders is required.

\textbf{(c) Unequal-order Cauchy comparison.} In the actual diagonal CM coordinates of ALP-LIFT-05, let the two functions satisfy its growth inequalities (H26), with fixed constants \(C_J^\theta,C_A^\theta>0\). Pair \((\omega_1,\omega_2,\omega_3)\) with \((\vartheta_1,\vartheta_2,\vartheta_4)\). Define the logarithmic polydisk Cauchy bounds \[
 \mathcal C_J(\boldsymbol n)=\log C_J^\theta+
   \sum_{i=1}^3\left[
     \log(n_i!)+\frac{n_i}{2}\log(\pi e/n_i)
                          -\frac{n_i}{2}\log G_i\right],
 \tag{D9}
\] where a term with \(n_i=0\) is interpreted by continuity as zero, and define \(\mathcal C_A\) by replacing \(G_i\) by \(G_B\) and the leading constant by \(C_A^\theta\). Then \[
 \mathcal C_J(\boldsymbol n)-\mathcal C_A(\boldsymbol n)
 =\log(C_J^\theta/C_A^\theta)
               +\sum_{i=1}^3n_i C_{i,\{2,7\}}.
 \tag{D10}
\] It follows from (D8) that \[
 \sum_iD_{i,\{2,7\}}(n_i;X)
       -\bigl(\mathcal C_J(\boldsymbol n)-
                                 \mathcal C_A(\boldsymbol n)\bigr)
                                                        =o(N)
 \tag{D11}
\] for the same arbitrary joint limit. Multiplying both budgets by the source's factor 2 leaves this cancellation unchanged. The assertion is about the displayed bounds, not equality of actual theta values.

\subsection{Proof of the local functions and convergence}\label{proof-of-the-local-functions-and-convergence}

Direct evaluation of the six residue classes gives \[
 b_1=\mathbf1_{\{4\}}-\mathbf1_{\{3\}},\qquad
 b_2=\mathbf1_{\{2\}}-\mathbf1_{\{5\}},\qquad
 b_3=b_1+b_2.
 \tag{D12}
\] They have mean zero and are odd: \(b_i(-a)=-b_i(a)\). Geometric summation of the powers of a residue class gives the following exact local table:

\begin{longtable}[]{@{}llll@{}}
\toprule\noalign{}
\(p\bmod7\) & \(d_1(p)\) & \(d_2(p)\) & \(d_3(p)\) \\
\midrule\noalign{}
\endhead
\bottomrule\noalign{}
\endlastfoot
\(1,6\) & \(0\) & \(0\) & \(0\) \\
\(2\) & \(p/(p^3-1)\) & \(p^2/(p^3-1)\) & \(p(p+1)/(p^3-1)\) \\
\(4\) & \(p^2/(p^3-1)\) & \(p/(p^3-1)\) & \(p(p+1)/(p^3-1)\) \\
\(3\) & \(-p^2/(p^3+1)\) & \(p/(p^3+1)\) & \(-p(p-1)/(p^3+1)\) \\
\(5\) & \(p/(p^3+1)\) & \(-p^2/(p^3+1)\) & \(-p(p-1)/(p^3+1)\) \\
\end{longtable}

For example, in class 3 the nonzero terms of \(b_1(p^r)\) are \(-1\) at \(r\equiv1\pmod6\) and \(+1\) at \(r\equiv4\pmod6\); hence \((-p^{-1}+p^{-4})/(1-p^{-6})=-p^2/(p^3+1)\). In class 2 the period is three, and the nonzero \(b_1\)-term is at exponent 2, giving \(p^{-2}/(1-p^{-3})=p/(p^3-1)\). Interchanging classes 2 and 4, or 3 and 5, interchanges \(b_1\) and \(b_2\); (D12) gives the last column. These observations prove every entry. In particular \[
                  (d_1(2),d_2(2),d_3(2))=(2/7,4/7,6/7).
 \tag{D13}
\]

For clarity about convergence, the quantitative prime theorem for fixed modulus 7, already cited in ALP-LIFT-04, gives \[
 B_i(x):=\sum_{p\le x}b_i(p)\log p
                   =O(xe^{-c_0\sqrt{\log x}}).
 \tag{D14}
\] The terms at 7 are defined to be zero. The main terms cancel because \(\sum_{a=1}^6b_i(a)=0\). The prime-power version of the same theorem implies this prime version as well, since powers of exponent at least two contribute \(O(\sqrt x\log x)\). Partial summation now proves convergence of \(\sum_p b_i(p)\log p/p\), with tail \(O((1+\sqrt{\log X})e^{-c_0\sqrt{\log X}})\). The series of the remaining terms is absolutely convergent, since \[
 \sum_p\log p\sum_{r\ge2}|b_i(p^r)|p^{-r}
       \le\sum_{m\ge2}\frac{\log m}{m(m-1)}<\infty.
\] This proves (D2) in the asserted ordinary-prime order. It also justifies absorbing the harmless polynomial factor in the tail by decreasing \(c_0\).

\subsection{Proof of the exact constants}\label{proof-of-the-exact-constants}

Let \(\chi(3)=\exp(\pi\mathrm i/3)\), and write \(Z_\eta=L'(0,\eta)/L(0,\eta)\). In the ordered odd character basis \((\chi,\chi^3,\overline\chi)\), the coefficients of the three functions in (D12) are \[
 \begin{array}{c|ccc}
 &\chi&\chi^3&\overline\chi\\\hline
 b_1&(-1+\mathrm i\sqrt3)/6&1/3&(-1-\mathrm i\sqrt3)/6\\
 b_2&(-1-\mathrm i\sqrt3)/6&1/3&(-1+\mathrm i\sqrt3)/6\\
 b_3&-1/3&2/3&-1/3.
 \end{array}
 \tag{D15}
\] This follows either by substituting the six group elements into the table, or by the normalized Fourier coefficient \(6^{-1}\sum_ab_i(a)\overline\eta(a)\). The sum of each row of coefficients is zero, since \(b_i(1)=0\).

For real \(s>1\), the Euler product gives \[
 \sum_{p\ne7}\log p\sum_{r\ge1}b_i(p^r)p^{-rs}
        =-\sum_\eta\widehat b_i(\eta)
                                      \frac{L'(s,\eta)}{L(s,\eta)}.
 \tag{D16}
\] Partial summation using (D14) permits the limit \(s\downarrow1\) for the first powers by an integrable bound independent of \(s\in[1,2]\); the higher powers converge absolutely there. The functional equation for an odd primitive character of conductor 7 gives \[
 \frac{L'(1,\eta)}{L(1,\eta)}+Z_{\overline\eta}
                         =\gamma+\log(2\pi/7).
 \tag{D17}
\] The constant vanishes after summation because the coefficient row sums in (D15) are zero. Hence \[
 C_i=\sum_\eta\widehat b_i(\eta)Z_{\overline\eta}.
 \tag{D18}
\] The conjugation in this formula is essential. Write \(Z=Z_\chi\), \(Z_3=Z_{\chi^3}\). Then (D15)--(D18) give \[
 C_1=(Z_3-\operatorname{Re}Z+\sqrt3\operatorname{Im}Z)/3,
 \quad
 C_2=(Z_3-\operatorname{Re}Z-\sqrt3\operatorname{Im}Z)/3,
 \quad C_3=2(Z_3-\operatorname{Re}Z)/3.
 \tag{D19}
\]

The Hurwitz special values used in ALP-LIFT-02 yield \[
 Z_3=-\log7+Q,
\] \[
 \operatorname{Re}Z=-\log7+\ell_1+\tfrac14\ell_2
          +\tfrac54\ell_3-\tfrac54\ell_4-\tfrac14\ell_5-\ell_6,
\] and the same finite character calculation yields \[
 \sqrt3\operatorname{Im}Z=-\tfrac32\ell_1+\tfrac94\ell_2
             +\tfrac34\ell_3-\tfrac34\ell_4-\tfrac94\ell_5
             +\tfrac32\ell_6.
 \tag{D20}
\] For example, the last two formulas follow by multiplying \(\sum_a\chi(a)\ell_a\) by \(L(0,\chi)^{-1}=1-\mathrm i\sqrt3/2\), whose value was computed exactly in ALP-LIFT-02. Substituting (D20) into (D19) gives all three expansions in (D4), and these are exactly (D3).

To identify the actual metric, the elliptic case of the same normalized Colmez formula, together with (H22), gives \[
                         \log G_B=\log(\pi/7)+Q.
 \tag{D21}
\] Here is the normalization explicitly. ALP-LIFT-02's Yang convention gives \(h_F(B)=-\tfrac12Z_3-\tfrac14\log7\); passing to its L² convention subtracts \(\tfrac12\log\pi\). On the other hand (H22) gives \(h_{F,L^2}(B)=-\tfrac14\log7-
\tfrac12\log G_B\). Equating these expressions proves (D21). No isogeny invariance is used.

The gamma duplication formula, applied at \(i/7\) and \(1/2-i/7\), rewrites (H8) as \[
 \log G_i=\log(\pi/7)+\frac{4i}{7}\log2
          +2\ell_i-2\ell_{7-i}+\ell_{7-2i}-\ell_{2i}.
 \tag{D22}
\] All gamma arguments are positive for \(i=1,2,3\). Subtracting (D21) from (D22), then multiplying by \(-1/2\) and using (D3), proves (D5). The explicit prime-2 values (D13) show that its left side is exactly \(C_{i,\{2,7\}}\). This proves part (a).

\subsection{Proof of the uniform integer estimate}\label{proof-of-the-uniform-integer-estimate}

We give the estimates rather than identifying a formal slope with a theta value. First let \(S=\{7\}\), suppress the index \(i\), and put \(b=b_i\). Decompose (D6) into first powers and higher powers. Define \[
 P_1(X)=\sum_{p\le X}b(p)\log p/p,\qquad
 P_2(X)=\sum_{p\le X}\log p\sum_{r\ge2}b(p^r)/p^r.
\] The convergence proof gives \[
 P_1(X)+P_2(X)=C_i+O(e^{-c_1\sqrt{\log X}}),\quad
 P_1(X)-P_1(\infty)=O(e^{-c_1\sqrt{\log X}}),
 \tag{D23}
\] and \(P_2(X)-P_2(\infty)=O((1+\log X)/X)\).

The higher-power floor errors are uniformly \(O(\sqrt n\log n)\). Indeed \(H_2(t)=\sum_{p^r\le t,r\ge2}\log p=O(\sqrt t\log t)\) by \(\theta(t)=O(t)\). Floors for \(p^r\le n\) produce an error at most \(H_2(n)\); the tail of the linear approximation is at most \(n\int_{(n,\infty)}t^{-1}dH_2(t)
=O(\sqrt n\log n)\). Restricting to \(p\le X\) only reduces these absolute bounds. Thus the higher-power contribution is \(nP_2(X)+O(\sqrt n\log n)\).

Put \(y=\lfloor\sqrt n\rfloor\). If \(X\le y\), the first-power floor error is at most \(\theta(X)=O(X)\). Together with (D23) this proves (D7) in that range. If \(X\ge y\), expanding the floors gives the exact identity \[
 \sum_{p\le X}b(p)\log p\left\lfloor n/p\right\rfloor
 =\sum_{p\le y}b(p)\log p\left\lfloor n/p\right\rfloor
 +\sum_{k\le n/y}
       \{B_i(\min(X,n/k))-B_i(y)\}.
 \tag{D24}
\] In the second sum, \(y\le\min(X,n/k)\le n/k\). Formula (D14) and the harmonic-sum bound show that its absolute value is at most \(O(n(1+\log n)e^{-c_0\sqrt{\log y}})\). The first sum is \(nP_1(\infty)+O(ne^{-c_1\sqrt{\log y}}+y)\). Combining with the higher powers and \(P_2(X)-P_2(\infty)=O((1+\log y)/y)\) proves (D7) for \(X\ge y\), after reducing \(c>0\). In both ranges the \(O(\sqrt n\log n)\) error is absorbed by the displayed exponential bound. The replacement of \(y\) by \(\sqrt n\) only changes a fixed multiplicative constant for \(n\ge4\).

For a fixed omitted prime \(p\ne7\), the corresponding floor sum differs from \(nd_i(p)\) by \(O_p(1+\log n)\): only \(O_p(1+\log n)\) powers have a nonzero floor, and the remaining linear geometric tail is \(O_p(1)\). Its \(\log p\) factor is fixed. Once \(X\) contains the exceptional set, subtracting these finitely many terms proves (D7) for \(S\). When it does not, \(X\) stays bounded by the largest exceptional prime, and an extra \(O_S(n)\) is absorbed by the right side after enlarging its constant. This establishes the full uniform assertion.

To deduce (D8) without a restriction on relative orders, fix \(\varepsilon>0\). Choose \(K\) large enough that the error factor in (D7) is below \(\varepsilon\) whenever \(n_i\ge K\) and \(X\ge\sqrt K\). The sum of errors in these directions is at most a constant times \(\varepsilon N\). In directions with \(n_i<K\), both their floor sums and their linear approximations are bounded by a constant depending only on \(K\), since every nonzero floor has \(p\le n_i<K\), independent of \(X\). Divide by \(N\), let \(N,X\to\infty\), then let \(\varepsilon\to0\). Using (D5) now proves (D8).

\subsection{Proof of the Cauchy statement and its scope}\label{proof-of-the-cauchy-statement-and-its-scope}

Apply polydisk Cauchy to the bounds in (H26). For a positive directional order choose \(R_i^2=n_iG_i/\pi\); for order zero take \(R_i\downarrow0\). The resulting bound is (D9). Perform the same calculation in the normalized induced eigenbasis (H12), whose three metric values are all \(G_B\). The factorial terms, \(\pi e\) terms and the individual \(n_i\log n_i\) terms cancel exactly between the two expressions. The remaining metric terms give (D10) by (D5); (D8) proves (D11). A common tensor level \(\kappa\) changes \(\pi\) to \(\kappa\pi\) and hence adds \((N/2)\log\kappa\) to both expressions, so it cancels as well. This is the level dependence whose source differential normalization is treated in AM-R004. This completes the proof. \(\square\)

The input in this theorem is an actual diagonal CM eigenbasis for both principally polarized abelian varieties, primitive at every common good prime. Thus the metric term is not a hypothetical period parameter. The lower prime certificate remains precisely the accepted DP64 object. The theorem proves that assigning different orders to the three fixed CM directions does not create an order-linear advantage within this paired bare-certificate and Cauchy-bound calculation.

At 7 the normalized eigenbasis need not be integral one direction at a time, although its determinant has the exact valuation proved in ALP-LIFT-05. This theorem does not bound an unequal-order 7-adic theta error. Nor does it show that the Cauchy bounds are attained, that selected values equal their universal lower certificates, or that the same arithmetic normalization gives the constants \(C_J^\theta,C_A^\theta\). The arithmetic kernel, row content, bad-prime theta normalization and division losses are still required before applying the full DP44 comparison.

\subsection{Numerical cross-checks and citations}\label{numerical-cross-checks-and-citations}

The accompanying exact-algebra/numerical script is \nolinkurl{work/agents/hodge_lattice_certificate.py}; it checks the gamma, Fourier and metric identities to 90 decimal digits as a cross-check of the symbolic proof. Its current approximate constants are:

\begin{longtable}[]{@{}
  >{\raggedright\arraybackslash}p{(\columnwidth - 4\tabcolsep) * \real{0.2727}}
  >{\raggedleft\arraybackslash}p{(\columnwidth - 4\tabcolsep) * \real{0.3636}}
  >{\raggedleft\arraybackslash}p{(\columnwidth - 4\tabcolsep) * \real{0.3636}}@{}}
\toprule\noalign{}
\begin{minipage}[b]{\linewidth}\raggedright
direction
\end{minipage} & \begin{minipage}[b]{\linewidth}\raggedleft
all \(p\ne7\): \(C_i\)
\end{minipage} & \begin{minipage}[b]{\linewidth}\raggedleft
good \(p\nmid14\): \(C_{i,\{2,7\}}\)
\end{minipage} \\
\midrule\noalign{}
\endhead
\bottomrule\noalign{}
\endlastfoot
1 & −0.127202946857769870150102196626 & −0.325244998446325672840739945614 \\
2 & 0.155046463619921284560680397991 & −0.241037639557190320820595099985 \\
3 & 0.0278435167621514144105782013653 & −0.566282638003515993661335045599 \\
\end{longtable}

External known inputs and exact locators are unchanged from ALP-LIFT-02/04/05: the abelian Colmez theorem in Yang (2010), Section 3, pp.~464--466, with Colmez (1993) and Obus (2013), arXiv v2 Theorem 3.9; Rudnick (2015), fixed-modulus prime theorem, Section 2, printed p.~2 before equation (1); DLMF 25.15.2, 25.15.5 and 25.15.9 for the Euler product, primitive functional equation and nonvanishing at 1, 25.11.13 and 25.11.18 for Hurwitz special values, and 5.5.5 for gamma duplication. The smooth Picard/Hodge and compatible-biextension interpretation is precisely that in ALP-LIFT-05, with the accepted DP22 FGA locators; it requires no additional weight or comparison-isomorphism claim.

\clearpage

\section{AM-R-004 --- an actual CM period lattice, theta section and uniform archimedean jet bound}\label{am-r-004-an-actual-cm-period-lattice-theta-section-and-uniform-archimedean-jet-bound}

Author: GPT-6 Astra Pro, ALP\_MACHINE \nolinkurl{/root}, 21 November 2026 UTC; reasoning effort is not separately exposed. This is a geometric normalization theorem on the actual Jacobian. It is not a construction of the requested global automorphic lift chain. The period-matrix theorem itself is known: the exact source is Tadokoro (2014), Proposition 3.2, Proposition 4.1 and Theorem 4.4, arXiv:1211.6910v5, pp.~9--10 and 13--15. We cite that theorem and calculate the particular basis, polarization, sections and bounds needed here.

\subsection{Statement}\label{statement-2}

Let C be the smooth projective curve y²=x⁷−1, J its principally polarized Jacobian, ζ=exp(2πi/7), E=Q(ζ), and Φ=\{1,2,3\}. Put

\[
 \omega_i=x^{i-1}dx/y,\qquad
 c_i=-\frac{2i_{\mathbb C}}7 B(i/7,1/2)(\zeta^i-1),\quad
 D=\operatorname{diag}(c_1,c_2,c_3).
 \tag{R4.1}
\]

Here i\_C is the imaginary unit, distinct from the index i. Define the 3 by 3 matrices, with i=1,2,3 and j=0,1,2,

\[
 A_{ij}=\zeta^{2ij},\qquad
 B_{ij}=\sum_{k=j}^{2}\zeta^{i(2k+1)},\qquad P=DA,
 \qquad \tau=A^{-1}B.
 \tag{R4.2}
\]

The following normalizations are explicit.

\begin{enumerate}
\def\labelenumi{\arabic{enumi}.}
\tightlist
\item
  The actual period lattice in the integral coordinates u\_i=∫ω\_i is
\end{enumerate}

\[
 \Lambda=D\Phi(\mathcal O_E)=P(\mathbb Z^3+\tau\mathbb Z^3).
 \tag{R4.3}
\]

Its principal Riemann form in O\_E coordinates is

\[
 \mathcal E(x,y)=\operatorname{Tr}_{E/\mathbb Q}
       \left(\frac{\zeta-\zeta^{-1}}7x\bar y\right).
 \tag{R4.4}
\]

In particular this is a concrete mixed-type CM torus with its actual principal polarization, not a torus induced from the quadratic type \{1,2,4\}.

\begin{enumerate}
\def\labelenumi{\arabic{enumi}.}
\setcounter{enumi}{1}
\tightlist
\item
  The normalized matrix is the following symmetric element of the Siegel upper half-space:
\end{enumerate}

\[
\tau=\begin{pmatrix}
-\zeta^5-1&-\zeta^5-\zeta^4-\zeta^2-2&\zeta^5+\zeta^3+\zeta\\
-\zeta^5-\zeta^4-\zeta^2-2&\zeta^5-\zeta^4+2\zeta^3+\zeta&\zeta^5+\zeta^3+\zeta^2+1\\
\zeta^5+\zeta^3+\zeta&\zeta^5+\zeta^3+\zeta^2+1&\zeta^2
\end{pmatrix}.
\tag{R4.5}
\]

Write Y=Im τ, and let G=diag(G\_1,G\_2,G\_3) be the Hodge Gram matrix computed in ALP-LIFT-03. Then

\[
 G=P Y\bar P^{\mathsf T},\qquad
 \operatorname{covol}_{\mathbb R^6}(\Lambda)=\det G
 =\frac{7^{5/2}}8\prod_i|c_i|^2.
 \tag{R4.6}
\]

The Euclidean volume here refers to the displayed u coordinates. It is not a level-group volume or a \nolinkurl{MainX} covolume.

\begin{enumerate}
\def\labelenumi{\arabic{enumi}.}
\setcounter{enumi}{2}
\tightlist
\item
  For every integer κ≥1, define
\end{enumerate}

\[
 \theta_\kappa(z)=\sum_{n\in\mathbb Z^3}
     \exp(\pi i\kappa n^{\mathsf T}\tau n+2\pi i\kappa n^{\mathsf T}z),
 \quad
 F_\kappa(u)=
 \frac{\exp\!\left(\frac{\pi\kappa}{2}z^{\mathsf T}Y^{-1}z\right)
       \theta_\kappa(z)}{\theta_\kappa(0)},\quad z=P^{-1}u.
 \tag{R4.7}
\]

These are actual sections of L\^{}κ, for the symmetric principal theta line L attached to τ, written in its local Chern-normalized holomorphic frame. Their rigidified germs satisfy F\_κ(0)=1. Explicit rational bounds prove

\[
 \lambda_{\min}(Y)>6/25,\qquad
 |\theta_\kappa(0)-1|<3/8,\qquad
 |\theta_\kappa(0)|>5/8.
 \tag{R4.8}
\]

Uniformly for every u∈C³ and κ≥1,

\[
 |F_\kappa(u)|<49\exp\!\left(
       \frac{\pi\kappa}{2}\sum_{i=1}^3\frac{|u_i|^2}{G_i}\right).
 \tag{R4.9}
\]

Consequently, for every multi-index α=(α\_1,α\_2,α\_3) of nonnegative integers,

\[
 |\partial^\alpha F_\kappa(0)|
 \le49\,\alpha!\prod_{\alpha_i>0}
       \left(\frac{\pi e\kappa}{\alpha_iG_i}\right)^{\alpha_i/2}.
 \tag{R4.10}
\]

For the source-shaped level normalization κ=M² and N=\textbar α\textbar{} this gives the exact cancellation

\[
 \left|M^{-N}\partial^\alpha F_{M^2}(0)\right|
 \le49\,\alpha!\prod_{\alpha_i>0}
       \left(\frac{\pi e}{\alpha_iG_i}\right)^{\alpha_i/2}.
 \tag{R4.11}
\]

This is a uniform archimedean inequality for an explicit family of sections. It supplies neither a prescribed nonzero positive-degree jet nor its required integral finite-place lattice.

\begin{enumerate}
\def\labelenumi{\arabic{enumi}.}
\setcounter{enumi}{3}
\tightlist
\item
  With probability Haar measure dx dt on J(C), z=x+τt, the metric in the usual theta frame is exp(−2πκ t\^{}T Yt). In this metric the exact squared L² norm of the rigidified section is
\end{enumerate}

\[
 \|\theta_\kappa/\theta_\kappa(0)\|_{L^2}^2
   =\frac{\det(2\kappa Y)^{-1/2}}{|\theta_\kappa(0)|^2}.
 \tag{R4.12}
\]

Thus even in this explicit example, normalization to value 1 and normalization to L² norm 1 are different operations, with the difference recorded rather than discarded.

\subsection{Proof of the lattice and polarization}\label{proof-of-the-lattice-and-polarization}

Use Tadokoro's paths λ\_j from (0,i\_C) to (0,−i\_C) through the branch point ζ\^{}j, and write γ\_j=λ\_\{j+1\}−λ\_j, j=0,\ldots,5. Tadokoro's Proposition 3.2 states that

\[
 a_i^{\rm T}=-\gamma_{2i-1},\qquad
 b_i^{\rm T}=\sum_{k=0}^{i-1}\gamma_{2k},\quad i=1,2,3
\]

are a symplectic basis. This fixes the topological orientation; it is not inferred from a numerical period calculation. Set

\[
 a_j=\gamma_{2j},\qquad b_j=\sum_{k=j}^{2}\gamma_{2k+1},
       \quad j=0,1,2.
\]

These are integral invertible triangular combinations of the previous basis: a\_j=b\^{}T\_\{j+1\}−b\^{}T\_j with b\^{}T\_0=0, and b\_j=−Σ\_\{k=j\}²a\^{}T\_\{k+1\}. Therefore they form a Z basis. Directly their intersections are

\[
 a_i\cdot b_j=\mathbf1_{i\ge j}-\mathbf1_{i-1\ge j}=\delta_{ij},
 \qquad a_i\cdot a_j=b_i\cdot b_j=0.
\]

On the outward segment x=ζ\^{}j t, y=i\_C√(1−t⁷), the beta substitution s=t⁷ gives

\[
 \int_{\lambda_j}\omega_i
   =-\frac{2i_{\mathbb C}}7 B(i/7,1/2)\zeta^{ij}.
\]

The factor 2 includes the return segment on the opposite sheet. Subtracting adjacent paths gives ∫γ\_j ω\_i=c\_i ζ\^{}\{ij\}. Hence the periods on a\_j and b\_j are DA and DB. The six γ\_j themselves are a Z basis, so their period columns give DΦ(Z{[}ζ{]})=DΦ(O\_E), proving (R4.3). The Vandermonde matrix A is invertible because ζ²,ζ⁴,ζ⁶ are distinct.

For ξ=(ζ−ζ\^{}−1)/7, the cyclotomic trace formula is

\[
 \operatorname{Tr}(\zeta^h)=
 \begin{cases}6,&7\mid h,\\-1,&7\nmid h.\end{cases}
\]

It follows that \(\operatorname{Tr}(\xi\zeta^j\overline{\zeta^k})\) is 1 for k=j+1, −1 for j=k+1, and zero otherwise for 0≤j,k≤5. These are exactly the intersections of the γ basis determined above. Thus (R4.4) is the actual principal Riemann form. Its positivity also follows directly from Im σ\_i(ξ)\textgreater0 for i=1,2,3. The ideal ξ O\_E is the codifferent: ζ−ζ\^{}−1 generates the prime above 7, whereas 7 has valuation 6 there, so its exponent is −5, the inverse of the different of Q(ζ\_7). This is consistent with unimodularity and does not replace the explicit intersection check.

Multiplying A by the matrix in (R4.5), reducing by 1+ζ+⋯+ζ⁶=0, gives B. This proves the displayed exact matrix; symmetry is visible. Positive definiteness follows either from the Riemann bilinear relations for the symplectic basis, or from the independent positive lower bound below.

The bilinear relations give \(G=P Y\bar P^T\). Taking determinants gives det G=\textbar det P\textbar²det Y, which is the real covolume of the six columns P and Pτ. The standard complex-embedding lattice of O\_E has covolume 2\^{}−3√\textbar disc E\textbar=7\^{}\{5/2\}/8. Multiplication by D scales real volume by ∏\textbar c\_i\textbar², proving (R4.6). The discriminant formula here is also the determinant of the full cyclotomic embedding matrix; it can be checked from the polynomial discriminant of 1+X+⋯+X⁶.

\subsection{Exact nonvanishing certificate}\label{exact-nonvanishing-certificate}

Put a=sin(π/7), b=sin(2π/7), c=sin(3π/7). Formula (R4.5) gives

\[
 Y=\begin{pmatrix}
 c&a&a+b-c\\
 a&3a+b-c&a\\
 a+b-c&a&c
 \end{pmatrix}.
 \tag{R4.13}
\]

The supplied reproducible script \nolinkurl{scripts/jacobian_theta_certificate.py} uses only exact \nolinkurl{Fraction} arithmetic for its proof assertions. It brackets π by Machin's identity 16 arctan(1/5)−4 arctan(1/239), with 40-term alternating-series bounds, and brackets the three sines with 24-term alternating-series bounds. Outward rounding then gives rational intervals with denominator 10⁶. The smallest Gershgorin lower bound of this rational enclosure is strictly greater than 6/25; all signs of the off-diagonal entries are positive. This proves the first part of (R4.8).

Here is the full finite reduction behind the remaining assertions. For t≥0 define

\[
 U(t)=\left(\sum_{k=0}^{40}\frac{t^k}{k!}\right)^{-1}\ge e^{-t}.
\]

The code rounds U upward to a multiple of 10\^{}−12. For each nonzero n in the cube \textbar n\_i\textbar≤3 it uses the rational interval lower endpoint q\_n for n\^{}T Yn, choosing the lower or upper matrix endpoint according to the sign of n\_i n\_j. Thus q\_n\textgreater0 and

\[
 \sum_{\substack{n\ne0\\|n_i|\le3}}e^{-\pi n^TYn}
 \le\sum_{\substack{n\ne0\\|n_i|\le3}}
       U((314159/100000)q_n).
\]

For the complement of the cube, π λ\_min(Y)\textgreater3/4. Set a\_0=3/4, q=e\^{}−a\_0. Then

\[
 \sum_{k\in\mathbb Z}e^{-a_0 k^2}\le1+\frac{2q}{1-q},\qquad
 \sum_{|k|>3}e^{-a_0 k^2}
 \le\frac{2e^{-16a_0}}{1-e^{-9a_0}}.
\]

The second inequality uses (4+j)²≥16+9j. A union bound on which of the three coordinates lies outside the cube bounds the complement by

\[
 3\frac{2e^{-16a_0}}{1-e^{-9a_0}}
       \left(1+\frac{2e^{-a_0}}{1-e^{-a_0}}\right)^2.
\]

Replacing the three exponentials by U with upward rounding gives a rational upper bound. The final rational comparison is strictly less than 3/8 (its separately labelled decimal approximation is 0.366129006419). All numerators and denominators and all inequalities are retained in \nolinkurl{evidence/theta_certificate.json}. This proves Σ\_\{n≠0\}exp(−πn\^{}TYn)\textless3/8. Increasing κ decreases every absolute summand. Consequently \textbar θ\_κ(0)−1\textbar\textless3/8 for all κ≥1, proving (R4.8). No floating-point approximation is a premise in this certificate.

\subsection{Sections, Chern jets and the uniform bound}\label{sections-chern-jets-and-the-uniform-bound}

The series θ\_κ converges normally on compact sets. Direct reindexing gives, for a,b∈Z³,

\[
 \theta_\kappa(z+a+\tau b)
  =e^{-\pi i\kappa b^T\tau b-2\pi i\kappa b^Tz}\theta_\kappa(z).
\]

This is the factor of automorphy of L\^{}κ, so these are genuine sections on the actual abelian variety. The line L is a symmetric representative of its principal polarization. We have not asserted that this representative or its particular section is fixed by the order-seven action. That distinction matters for a prescribed CM-labelled automorphic row.

In the usual holomorphic theta frame the metric is h(z)=exp(−2πκ(Im z)\^{}T Y\^{}−1 Im z). Its Chern connection has (1,0) coefficient \(\pi\kappa Y^{-1}(z-\bar z)\). Put Q(z)=πκ z\^{}TY\^{}−1z/2. Replacing the frame e by e′=e\^{}−Q e makes the connection coefficient \(-\pi\kappa Y^{-1}\bar z\) and the metric \(\exp(-\pi\kappa\bar z^TY^{-1}z)\). Thus on the formal holomorphic directions at 0, covariant derivatives are exactly the ordinary derivatives of e\^{}Qθ\_κ. Normalization by the nonzero θ\_κ(0) proves the claimed interpretation of F\_κ.

For y=Im z, completing the square yields

\[
 |\theta_\kappa(z)|
 \le e^{\pi\kappa y^TY^{-1}y}
      \sum_{n\in\mathbb Z^3}
        e^{-\pi\kappa(n+Y^{-1}y)^TY(n+Y^{-1}y)}.
\]

For every λ\textgreater0 and real t, the elementary unimodal Gaussian bound is

\[
 \sum_{n\in\mathbb Z}e^{-\pi\lambda(n+t)^2}
      \le1+\lambda^{-1/2}.
\]

To see it, keep a nearest-to-zero sample as a term at most 1. Every sample to its left or right is at most the Gaussian integral on the adjoining unit interval toward that nearest sample; these intervals partition the real line. The remaining sum is at most the Gaussian integral λ\^{}−1/2. Applying this in three coordinates, using κ≥1 and λ\_min(Y)\textgreater6/25, bounds the displayed three-dimensional sum by (1+5/√6)³\textless(25/8)³.

Multiplication by \textbar e\^{}Q\textbar{} changes the exponential to

\[
 \exp\left(\frac{\pi\kappa}{2}\bar z^T Y^{-1}z\right).
\]

Since \(G=P Y\bar P^T\) is real diagonal, \(\bar z^TY^{-1}z=\sum_i|u_i|^2/G_i\). Dividing by \textbar θ\_κ(0)\textbar\textgreater5/8 leaves a constant smaller than (25/8)³(8/5)=3125/64\textless49. This proves (R4.9).

Cauchy's integral inequality on a polydisc of radii R\_i gives

\[
 |\partial^\alpha F_\kappa(0)|
 \le49\alpha!\prod_i
      \left(e^{\pi\kappa R_i^2/(2G_i)}R_i^{-\alpha_i}\right).
\]

For α\_i\textgreater0 minimize at R\_i²=α\_i G\_i/(πκ); for α\_i=0 take the limit R\_i→0. This proves (R4.10). With κ=M², the factor κ\^{}\{\textbar α\textbar/2\} is precisely M\^{}N, proving (R4.11). This calculation is an explicit instance of the source normalization issue already emphasized in DP50; it does not alter DP50's carried proof status.

Finally integrate the squared section norm with z=x+τt. Integration in x over {[}0,1{]}³ kills all cross terms n≠m. The remaining integrand is Σ\_n exp(−2πκ(n+t)\^{}TY(n+t)). Integration in t unfolds the unit cubes into R³, whose Gaussian integral is det(2κY)\^{}−1/2. Division by \textbar θ\_κ(0)\textbar² proves (R4.12). All integrations are justified by normal convergence or positive-term Tonelli integration. This completes the proof.

\subsection{What remains before this is a lift-count input}\label{what-remains-before-this-is-a-lift-count-input}

The statement now fixes an actual mixed-type CM lattice, polarization, complex covolume, nonzero section family, metric, Chern gauge and level-uniform archimedean jet bound. The new exact finite Hodge lattices in ALP-LIFT-05 complement it. These are concrete geometric data that can be used in a replacement arithmetic constructor.

They do not identify F\_κ with the primitive trace-one Poincaré kernel of the selected source row. Nor do they establish integral normalized coefficients, the allowable-translate expansion with controlled algebraic coefficient height, membership in the chosen automorphic theta image, all-degree \nolinkurl{ClsK.Core}, \nolinkurl{SizeData₂}, or the six identities for that tied actual family. In particular, θ\_κ(0)≠0 proves nonvanishing of the section at the origin; it is not controlled nonvanishing of the prescribed high-degree harmonic jet or terminal lift. Passing from a nonzero complex section to such a family remains a substantive arithmetic step.

In Grothendieck's language, (R4.3--6) fix the Betti lattice and polarization of the Picard variety; (R4.7--12) fix its complex metrized theta line and a Hodge-splitting jet calculation. The Poincaré biextension framework is SGA 7 I, Exposés VII--VIII; the integral crystalline realization and its filtration remain the objects of cote 54 scans 3 and 5--7 and cote 4 scan 3, (8.1.1--2). No assertion in those manuscripts is replaced by the complex analytic bounds here.

\subsection{References}\label{references}

Tadokoro, Y. (2014). The period matrix of the hyperelliptic curve w²=z\^{}(2g+1)−1. \emph{Tsukuba Journal of Mathematics, 38}(2), 137--158. The version actually consulted is arXiv:1211.6910v5, Proposition 3.2, Proposition 4.1, Theorem 4.4, pp.~9--10, 13--15. \url{https://arxiv.org/abs/1211.6910}.

Supplied programme sources: DP\_MULT DP49, DP50, DP57, DP64, DP65 and DP67; THETA3 Part II G1. Their statuses are carried without upgrade. The exact new rational certificate and the symbolic period check are documented separately.

\clearpage

\section{AM-R-005: a uniform bound for every nonsingular even theta characteristic}\label{am-r-005-a-uniform-bound-for-every-nonsingular-even-theta-characteristic}

Author: GPT-6 Astra Pro, \nolinkurl{/root}; ALP\_MACHINE, 21 November 2026 UTC. Reasoning effort is not separately exposed. This extends the one-characteristic archimedean calculation AM-R-004 to all algebraic conjugates and to normalized Mumford translates. It does not assert finite-place integrality of the reduced formal frame, membership in the selected arithmetic rows, or a lift-cost sign.

\subsection{Statement}\label{statement-3}

Retain the actual curve, differentials, period matrix and metric of AM-R-004: \[
C:y^2=x^7-1,\quad E=\mathbf Q(\zeta_7),\quad
\omega_i=x^{i-1}dx/y,\quad u=Pz,\quad Y=\operatorname{Im}\tau,
\quad P Y\overline P^{\,t}=\operatorname{diag}(G_1,G_2,G_3).
\tag{A1}
\] For a three-element subset \(S\subset\{0,\ldots,6\}\), let \(D_S=\sum_{j\in S}W_j-\infty\), with \(W_j=(\zeta_7^j,0)\). Let \(\delta_S=[a_S;b_S]\) be its nonsingular even half-characteristic in this period basis, with \(a_S,b_S\in\{0,1/2\}^3\). Put \[
F_S(u)=\exp\!\left(\frac\pi2 z^tY^{-1}z\right)
       \frac{\theta[\delta_S](z;\tau)}{\theta[\delta_S](0;\tau)}.
\tag{A2}
\] The classification and nonsingularity of these 35 characteristics are AM-DO-012. Their exact identification includes \(S_0=\{1,3,5\}\), whose characteristic is zero, in ALP-LIFT-07.

For all such \(S\), \[
\left|\frac{\theta[\delta_S](0;\tau)}{\theta[0;0](0;\tau)}\right|
   >\frac18,
\qquad
\boxed{|F_S(u)|<392\exp\!\left(\frac\pi2
                           \sum_{i=1}^3 |u_i|^2/G_i\right).}
\tag{A3}
\] In particular \(|\theta[\delta_S](0;\tau)|>5/64\). The constant is deliberately a simple certified uniform bound, not an optimum. For every multi-index \(\alpha\in\mathbf N^3\), with the usual convention for a zero component, \[
|\partial^\alpha F_S(0)|
 \le392\,\alpha!\prod_{\alpha_i>0}
             \left(\frac{\pi e}{\alpha_iG_i}\right)^{\alpha_i/2}.
\tag{A4}
\] The same bound holds for every normalized Mumford translate of \(F_S\), even when the translated section vanishes at the origin. For its pullback under \([M]\), the level-normalized derivative \(M^{-|\alpha|}\partial^\alpha\) again satisfies (A4).

Each \(F_S\) has coefficients in \(E\). If \(\sigma_a(\zeta_7)=\zeta_7^a\), then coefficientwise conjugation in the rational differential coordinates satisfies \[
                    \sigma_a(F_S)=F_{aS}.
\tag{A5}
\] Consequently (A4) bounds every conjugate of each algebraic derivative. For the Mumford translates, every embedding of their finite torsion field is bounded in the same way, after the corresponding torsion translation; the possible change of a period representative is a factor of absolute value one.

\subsection{Proof of the theta-null comparison}\label{proof-of-the-theta-null-comparison}

The only external theta-constant input is the classical first Thomae formula in the precise form stated by Eilers (2018), Theorem 2.1, equation (2.17), with the partition convention of equations (2.12)--(2.16). It applies to an odd-degree hyperelliptic equation with the branch point at infinity omitted from the Vandermonde products. Taking absolute values and a ratio removes both its common period factor and its eighth-root-of-unity ambiguity. Thus \[
\left|\frac{\theta[\delta_S](0)}{\theta[\delta_{S_0}](0)}\right|^4
 =\frac{\displaystyle
     \prod_{i<j\in S}|\zeta_7^i-\zeta_7^j|
     \prod_{i<j\notin S}|\zeta_7^i-\zeta_7^j|}
    {\displaystyle
     \prod_{i<j\in S_0}|\zeta_7^i-\zeta_7^j|
     \prod_{i<j\notin S_0}|\zeta_7^i-\zeta_7^j|}.
\tag{A6}
\] Every numerator and denominator has exactly \(\binom32+\binom42=9\) factors.

For distinct indices every factor is between \(2\sin(\pi/7)\) and 2. The lower bound is strictly greater than \(4/5\). One elementary verification uses \(\pi>3\), the monotonicity of sine on \([0,\pi/2]\), and its alternating Taylor lower bound: \[
 \sin(\pi/7)>\sin(3/7)\ge 3/7-(3/7)^3/6
                 =285/686>2/5.
\] It follows from (A6) that its right side is greater than \((2/5)^9\). The exact integer inequality \(2^{21}=2097152>1953125=5^9\) says \((2/5)^9>1/8^4\). Taking the positive fourth root gives the first inequality of (A3). AM-R-004 gives \(|\theta[0;0](0)|>5/8\), proving the asserted absolute lower bound.

\subsection{Proof of the global growth and derivative bounds}\label{proof-of-the-global-growth-and-derivative-bounds}

For real half-characteristic vectors \(a,b\) and \(z=x+iy\), completing the real quadratic form in the absolutely convergent theta series gives \[
 e^{-\pi y^tY^{-1}y}\,|\theta[a;b](z;\tau)|
 \le\sum_{n\in\mathbf Z^3}
 e^{-\pi(n+a+Y^{-1}y)^tY(n+a+Y^{-1}y)}.
\tag{A7}
\] The phase containing \(b\) has absolute value one. AM-R-004 proves \(\lambda_{\min}(Y)>6/25\) and the shifted one-dimensional Gaussian bound used there. Hence the right side of (A7) is at most \[
                         C_0=(1+5/\sqrt6)^3.
\] The same exact inequalities in AM-R-004 give \(C_0/|\theta[0;0](0)|<49\). Combining this with (A6) bounds \(C_0/|\theta[\delta_S](0)|\) by 392. Finally \[
 \operatorname{Re}(z^tY^{-1}z)/2+y^tY^{-1}y
                         =z^tY^{-1}\overline z/2,
\] and (A1) turns this into \(\tfrac12\sum_i|u_i|^2/G_i\). This proves the second inequality of (A3).

Apply polydisk Cauchy with \(R_i^2=\alpha_iG_i/\pi\) when \(\alpha_i>0\). For a zero component let its radius decrease to zero. The resulting product is exactly (A4). No equality of actual derivative size with this upper bound is claimed.

Let \(H(u,v)=\sum_i u_i\overline v_i/G_i\). A normalized translation by a period-cover representative \(t\) has the analytic form \[
 (U_tF)(u)=\varepsilon_t
       e^{-\pi H(u,t)-\pi H(t,t)/2}F(u+t),\qquad |\varepsilon_t|=1.
\tag{A8}
\] This is Bannai--Kobayashi's normalized translation convention; its Mumford version has the explicitly specified finite theta-group factor. Since \[
 H(u+t,u+t)-2\operatorname{Re}H(u,t)-H(t,t)=H(u,u),
\] the exponential cross terms cancel exactly in the absolute-value bound. Thus (A3) and (A4) hold for \(U_tF_S\) with the same constant. Substitution of \(Mu\) and the chain rule prove the level-normalized assertion. Replacing \(t\) by \(t+\lambda\), \(\lambda\in\Lambda\), changes the normalized translate only by its semicharacter and the exponential of a purely imaginary pairing, both of absolute value one.

\subsection{Algebraicity and the conjugation assertion}\label{algebraicity-and-the-conjugation-assertion}

AM-DO-011 applies Bannai--Kobayashi (2010), arXiv v4 Proposition 2.1, to these symmetric algebraic divisors, their admissible polarization, the algebraic differential coordinates, and value 1 at the origin. The same hypotheses hold for every \(D_S\), so \(F_S\in E[[u_1,u_2,u_3]]\).

We justify (A5), since algebraicity alone is not a statement about the analytic growth of conjugates. The coefficient construction in the proof of that proposition is Galois equivariant in this situation. More explicitly choose \(b\in2\mathcal O_E\) there so that every positive-degree multi-index has unequal eigenvalue products for \(b\) and \(\overline b\). The normalized periodic function \[
 R_{S,b}(u)=F_S(\iota(b)u)/F_S(\iota(\overline b)u)
\tag{A9}
\] is the unique rational function of value 1 at the origin with divisor \(\iota(b)^*\Theta_S-\iota(\overline b)^*\Theta_S\). Its Taylor coefficients are algebraic, and (A9) recursively determines every coefficient of \(F_S\) from the earlier ones. The nonzero coefficient multiplying the next unknown is exactly the difference of the two displayed eigenvalue products. These facts are the proof of the cited proposition, not a new replacement for its algebraicity theorem.

The curve and \(\omega_i\) are rational, while \(\sigma_a\iota(b)=\iota(\sigma_ab)\) because the generating endomorphism is \((x,y)\mapsto(\zeta_7x,y)\). Conjugation commutes with \(\sigma_a\), and \(\sigma_aD_S=D_{aS}\). Therefore \(\sigma_aR_{S,b}\) is exactly \(R_{aS,\sigma_ab}\), by its normalized rational divisor description. The transformed eigenvalues are the corresponding eigenvalues of \(\iota(\sigma_ab)\) on the same rational differentials. The recursive equations and initial value 1 now agree; induction proves (A5).

For torsion translations use the canonical algebraic isomorphisms in Bannai--Kobayashi Proposition 2.3 and Theorem 2.9. Formation of those isomorphisms commutes with coefficient-field embeddings. Analytically they are the translations (A8) by the embedded torsion point, up to the specified finite-order factor. The period-representative observation following (A8) removes that harmless ambiguity from absolute values. This proves the final assertion and completes the proof. \(\square\)

\subsection{Scope for the main question and Grothendieck's frameworks}\label{scope-for-the-main-question-and-grothendiecks-frameworks}

This closes the uniform archimedean normalization subgate for the 35 normalized even theta germs and their canonical algebraic torsion translations on this fixed Jacobian. It retains the exact metric \(G_i\), instead of substituting an arbitrary complex-coordinate norm. Together with ALP-LIFT-06 it shows that passing to a different such theta characteristic changes this certified bound by a fixed additive constant, not by an order-linear metric advantage. This does not bound the height of a general varying linear combination, the finite denominators of its theta-group projector, or its actual vanishing and divisibility. Those data remain in Q1 and THETA3 G1.

Geometrically the characteristic change is translation of the relative Picard theta line by a section of its finite etale two-torsion scheme over \(\mathbf Z[\zeta_7,1/14]\). The compatible translation isomorphisms belong to the symmetric-line/theta-group structure and the Poincare biextension, not to an arbitrary identification of functions. The Hodge metrics and determinant lines are the ones in ALP-LIFT-05; the precise FGA, cote54 and SGA7 dictionary is recorded separately. No new claim about number-field L-function zeros or finite-field weights follows from the uniform analytic norm bound.

\subsection{References and check provenance}\label{references-and-check-provenance}

\begin{itemize}
\tightlist
\item
  Eilers, K. (2018). Rosenhain--Thomae formulae for higher genera hyperelliptic curves. \emph{Journal of Nonlinear Mathematical Physics, 25}(1), 86--105. \url{https://doi.org/10.1080/14029251.2018.1440744}. Exact input: Theorem 2.1, equation (2.17), and the odd-degree partition convention (2.12)--(2.16). The primary article text was read at \url{https://www.atlantis-press.com/journals/jnmp/125950502/view}.
\item
  Bannai, K., \& Kobayashi, S. (2010). Algebraic theta functions and the p-adic interpolation of Eisenstein--Kronecker numbers. \emph{Duke Mathematical Journal, 153}(2), 229--295. \url{https://doi.org/10.1215/00127094-2010-024}. Exact text used: arXiv:math/0610163v4, Definition 1.8, Proposition 2.1 and its proof, Proposition 2.3, and Theorem 2.9. No ordinary-integrality result is extended to a nonordinary prime here.
\item
  AM-R-004; AM-DO-011--013; ALP-LIFT-06--07. These dependencies retain their separate whole-proof-check events. The elementary bound above is a new corollary of the cited known Thomae formula, whose proof is not repeated.
\end{itemize}

\clearpage

\section{ALP-LIFT-07: the actual theta characteristic, its finite normalization, and the source line}\label{alp-lift-07-the-actual-theta-characteristic-its-finite-normalization-and-the-source-line}

Run ALP\_MACHINE. Author: \nolinkurl{/root/lift_accounting}, GPT-6 Astra Pro; reasoning effort is not separately exposed. The good models and Hodge bases are the proved ALP-LIFT-05 inputs. The homology basis is AM-R-004. The branch half-period table also occurs independently in AM-DO-012; credit it once in a final synthesis. The new finite calculation uses the determinant of cohomology, so no unspecified Thomae scalar enters it.

\subsection{Statement and normalization}\label{statement-and-normalization}

Put \[
 C:y^2=x^7-1,\quad J=\operatorname{Jac}(C),\quad
 P_a=(\zeta_7^a,0)\ (0\leq a\leq6),\quad P_\infty=\infty.
 \tag{T1}
\] For a three-element subset \(S\subset\{0,\ldots,6\}\), set \[
 E_S=\sum_{a\in S}P_a,\qquad D_S=E_S-P_\infty.
 \tag{T2}
\] At each finite place work over a complete DVR \(R\), after a finite extension large enough for ALP-LIFT-05's good curve model and these sections. Write \(v_p(p)=1\). On the relative degree-two Picard scheme take its canonical determinant theta divisor, pull it back to \(\mathcal J\) by \(x\mapsto\mathcal O(D_S)\otimes x\), and denote the resulting line and canonical section by \((\mathcal L_S,s_S)\). The section has no vertical divisor: on each geometric fibre it is the ordinary nonzero theta section. This is the primitive integral theta model meant here. The valuation of its nonzero origin value means its valuation relative to any primitive generator of the free rank-one module \(e^*\mathcal L_S\); this is independent of the generator's integral-unit ambiguity.

\textbf{(a) Exact finite value.} All 35 choices have \(s_S(0)\ne0\) on the generic fibre, and \[
 \boxed{v_p(s_S(0))=
 \begin{cases}6/7,&p=2,\\0,&p\text{ odd}.\end{cases}}
 \tag{T3}
\] The entry at 7 is zero on the good model after extension. Normalization of the value to 1 therefore contributes the fixed scalar valuation \(-6/7\) at 2 and zero at every odd prime, in these local integral origin frames.

\textbf{(b) Characteristic of the explicit analytic germ.} In AM-R-004's symplectic basis \[
 a_j=\gamma_{2j},\qquad b_j=\sum_{k=j}^2\gamma_{2k+1}
 \quad(0\leq j\leq2),
 \tag{T4}
\] write a half-period as \((m+\tau n)/2\), with \(m,n\in\mathbb F_2^3\). The Riemann constant for infinity is \[
 K_\infty=\frac{(1,0,1)^t+\tau(1,1,1)^t}{2}
 =\operatorname{AJ}(P_1+P_3+P_5-3P_\infty).
 \tag{T5}
\] Thus \(D_{135}=P_1+P_3+P_5-P_\infty\) corresponds to the \textbf{zero analytic characteristic}. Its theta divisor is the zero locus of \(\theta(z,\tau)\), and AM-R-004's value-one germ \[
 F_1(u)=\exp\!\left(\frac\pi2z^t(\operatorname{Im}\tau)^{-1}z\right)
 \frac{\theta(z,\tau)}{\theta(0,\tau)},\qquad u=Pz,
 \tag{T6}
\] belongs to this precise principal line. AM-DO-011 identifies it as the reduced theta in the specified CM differential coordinates and applies Bannai--Kobayashi's algebraicity theorem. Those are separate carried findings; (T3) is an additional integral origin calculation.

\textbf{(c) The trace-one line is a different principal line.} Let \(a\) be induced by \(x\mapsto\zeta_7x\), let \(w_0=-(a+a^2+a^4)\), and put \(L_{w_0}=(1,w_0)^*\mathcal P_\lambda\), as in accepted DP49. Let \(L_\infty\) be the theta line obtained from \(2P_\infty\). Superscript \(\mathrm{rig}\) means tensoring a line with the inverse of its origin fibre, giving a rigidified line. On the good models as well as generically, \[
 L_{w_0}\simeq L_\infty^{\mathrm{rig}}.
 \tag{T7}
\] Its unique principal theta section vanishes to order \textbf{two} at zero. In contrast, \[
 L_{135}^{\mathrm{rig}}\simeq L_{w_0}\otimes\mathcal P_T,
 \qquad T=K_\infty\in J[2]\setminus\{0\},
 \tag{T8}
\] using the canonical polarization to identify \(J\) with its dual. For every positive even integer \(M\) there is an integral rigidified isomorphism \[
 [M]^*L_{135}^{\mathrm{rig}}\simeq L_{w_0}^{\otimes M^2}.
 \tag{T9}
\] The pulled-back primitive section retains (T3). Its normalized reduced germ is \(F_1(Mu)\), and \[
 M^{-|\alpha|}\partial_u^\alpha F_1(Mu)|_{u=0}
       =\partial_u^\alpha F_1(u)|_{u=0}.
 \tag{T10}
\] This construction is \textbf{outside the present source \nolinkurl{Level} record}, which requires \nolinkurl{Odd M}. For odd \(M\) the nonzero two-torsion twist survives pullback. An odd-order theta-group translation within \(H^0(J,L_{w_0}^{M^2})\) cannot change its underlying line into that of \([M]^*s_{135}\). This is a line distinction, not a claim that allowed odd-order translates have zero origin value.

\subsection{Proof of the determinant normalization}\label{proof-of-the-determinant-normalization}

\subsubsection{1. The determinant theta model and its dual complex}\label{the-determinant-theta-model-and-its-dual-complex}

The degree \(g-1\) theta divisor has its standard determinant-of- cohomology description: a perfect two-term complex of equal-rank vector bundles computes the universal degree \(g-1\) line's cohomology, and the determinant of its differential is the theta section. See Casalaina-Martin and Kass (2010, Definition 3.1 and the ensuing paragraphs, printed pp.~9--10). Invariance of this determinant element under a change of perfect representative is the determinant functor formalism (Stacks Project, Tag 0FJI). We use this construction over the DVR; no equal-characteristic arc theorem is asserted to cover the present mixed-characteristic family.

Write \(\omega=\Omega^1_{\mathcal C/R}\). The exact sequence \[
 0\to\omega(P_\infty-E_S)\to\omega(P_\infty)
   \to\omega(P_\infty)|_{E_S}\to0
 \tag{T11}
\] gives a simple complex. On every geometric fibre, \(H^1(\omega(P_\infty))=0\) by Serre duality and \(h^0(\omega(P_\infty))=3\) by Riemann--Roch. The inclusion \(H^0(\omega)\to H^0(\omega(P_\infty))\) is an isomorphism on each fibre: both have dimension three, or a differential cannot have a single nonzero simple-pole residue. Cohomology and base change and Nakayama therefore identify these integral rank-three modules. Since \(E_S\) is finite flat of degree three, its restricted invertible sheaf has free rank three over \(R\). Hence \[
 r_S:H^0(\omega(P_\infty))\longrightarrow
 H^0(E_S,\omega(P_\infty)|_{E_S})
 \tag{T12}
\] is a map between free rank-three modules. Its complex in degrees zero and one computes \(R\Gamma(\omega\otimes\mathcal O(-D_S))\).

Relative Serre duality identifies this with the dual shifted complex of \(R\Gamma(\mathcal O(D_S))\); see Stacks Project, Tag 0E6N. For a free complex \([A\xrightarrow d B]\) in degrees zero and one, the dual shifted complex is \([B^*\xrightarrow{\pm d^t}A^*]\), which has the same determinant valuation. Thus the theta section at \(D_S\) has valuation \(v_p(\det r_S)\). Equivalently, Smith normal form says that the determinant order of this generically acyclic complex is the sum of its elementary-divisor orders, the length of its torsion cokernel. This elementary argument works over every DVR. We have not used an unverified integral identification of \(\omega(-D_S)\) with \(\mathcal O(D_S)\).

The canonical relative theta section is primitive because its restriction to every geometric Picard fibre is a proper effective theta divisor, not that entire fibre. The translated pullback is primitive for the same reason. The determinant in (T12) therefore computes the stated model, with no arbitrary vertical scalar.

\subsubsection{2. Every odd-prime value is a unit}\label{every-odd-prime-value-is-a-unit}

For \(p\nmid14\), the original good model has seven distinct finite branch points, separate from infinity. At 7 use ALP-LIFT-05's good model \[
 Y^2=\frac{(1+\lambda X)^7-1}{\lambda^7},\qquad
 \lambda=\zeta_7-1,
 \tag{T13}
\] whose special fibre is \(Y^2=X^7-X\). The branch coordinates are \[
 X_a=(\zeta_7^a-1)/\lambda\equiv a\pmod\lambda,
 \tag{T14}
\] so they are again distinct. On each smooth odd-characteristic fibre, \(K_C\sim4P_\infty\) and \(H^0(\mathcal O(5P_\infty))=\langle1,X,X^2\rangle\). Riemann--Roch and evaluation at three distinct coordinates give \[
 h^0(D_S)=h^0(5P_\infty-E_S)=0:
 \tag{T15}
\] the quadratic-polynomial evaluation matrix is an invertible Vandermonde. This also proves the generic nonvanishing. Since the degree is \(g-1\), first cohomology vanishes as well. The matrix (T12) is therefore invertible after reduction, proving all the odd entries of (T3). No parity assertion in characteristic 2 was used.

\subsubsection{3. The dyadic collision algebra}\label{the-dyadic-collision-algebra}

Use the good model and integral Hodge basis from ALP-LIFT-05: \[
 Y^2+Y=X^7,\quad x=\alpha X,\quad y=\iota(2Y+1),\quad
 \alpha^7=-4,\quad\iota^2=-1,\qquad
 \eta_i=X^{i-1}\frac{dX}{2Y+1}\ (1\leq i\leq3).
 \tag{T16}
\] On the infinity chart \(U=1/X\), \(V=YU^4\), the equation is \(V^2+U^4V=U\). The rational function \[
 t=X^3/Y=U/V=V+U^4
 \tag{T17}
\] is a uniformizer at the infinity section. Indeed \(U\) is a formal power series in \(V\) beginning with \(V^2\), so \(t=V+O(V^8)\). The completed ring is \(R[[t]]\), and \(dt/t\) is an integral generator of \(\omega(P_\infty)\) on this disk. At \(P_a\), \[
 X_a=\zeta_7^a/\alpha,\quad Y_a=-1/2,\quad
 t_a=-2\alpha^{-3}\zeta_7^{3a}.
 \tag{T18}
\] All seven sections specialize to infinity. Since \(v_2(\alpha)=2/7\), \[
 v_2(t_a)=1/7,\qquad v_2(t_a-t_b)=1/7\quad(a\ne b).
 \tag{T19}
\] The differences of distinct seventh roots of unity are dyadic units: their reductions are distinct since \(X^7-1\) is separable in characteristic 2. The actual finite flat collision algebra is thus \[
 H^0(E_S,\mathcal O_{E_S})
   =R[[t]]\big/\!\prod_{a\in S}(t-t_a).
 \tag{T20}
\] The product is a distinguished monic cubic, so Weierstrass division gives the integral basis \(1,t,t^2\). Inverting the uniformizer splits the algebra into the three point evaluations. That split evaluation basis is not integral: its determinant relative to \(1,t,t^2\) is a Vandermonde of valuation \(3/7\).

Write \(\eta_i=f_i(t)dt/t\). Differentiating (T17) and using \(X^7=Y^2+Y\) yields the exact identity \[
 dt=-\frac{X^2(Y+4)}{Y(2Y+1)}dX,\qquad
 \frac{\eta_i}{dt}=-\frac{YX^{i-3}}{Y+4}.
 \tag{T21}
\] The ratio extends at every \(P_a\). Substitution gives \[
 f_i(t_a)=t_a\frac{\eta_i}{dt}(P_a)=-\frac27X_a^i.
 \tag{T22}
\] Up to the ordering sign, this split evaluation determinant is \[
 \det(f_i(t_a))=
 \left(-\frac27\right)^3
 \left(\prod_{a\in S}X_a\right)
 \prod_{a<b\text{ in }S}(X_b-X_a).
 \tag{T23}
\] Each \(X_a\) and each pairwise difference has valuation \(-2/7\). Thus (T23) has valuation \[
 3-6/7-6/7=9/7.
 \tag{T24}
\] The conversion matrix from the integral target basis \((1,t,t^2)dt/t\) to split evaluation has determinant \(\prod_{a<b}(t_b-t_a)\), of valuation \(3/7\). Consequently the actual integral restriction map in (T12) has \[
 v_2(\det r_S)=9/7-3/7=6/7.
 \tag{T25}
\] This proves the dyadic entry of (T3) for every triple. It also explains why using only split point values would give the incorrect answer \(9/7\). The special-fibre class \(D_S\) is \(2P_\infty\), whose space of sections has dimension two, consistent with the nonunit theta value.

\subsection{Proof of the characteristic and source-line assertions}\label{proof-of-the-characteristic-and-source-line-assertions}

\subsubsection{4. Half-periods and the Riemann constant}\label{half-periods-and-the-riemann-constant}

In this subsection only, write \(t_a=\operatorname{AJ}(P_a-P_\infty)\); this is distinct from the parameter values in (T18). Each \(t_a\) is two-torsion because \(\operatorname{div}(x-\zeta_7^a)=2P_a-2P_\infty\). The adjacent-arc relation for the paths of AM-R-004 is \[
 t_a-t_0=\frac12\sum_{j=0}^{a-1}\gamma_j
                    \pmod{H_1(C,\mathbb Z)}.
 \tag{T26}
\] This is the hyperelliptic doubled-arc calculation: the two sheets give opposite integrals between the ramification endpoints, so the closed cycle is twice the Abel integral. It also follows from AM-DO-012's independently evaluated branch half-periods. Since \(\operatorname{div}(y)=\sum_aP_a-7P_\infty\), one has \(\sum_at_a=0\). Summing (T26), using two-torsion, and reducing cycle coefficients modulo two gives \[
 t_0=(\gamma_1+\gamma_3+\gamma_5)/2=b_0/2.
 \tag{T27}
\] Substitution of \(\gamma_{2j}=a_j\), \(\gamma_1=b_0-b_1\), \(\gamma_3=b_1-b_2\), and \(\gamma_5=b_2\) now gives

\begin{longtable}[]{@{}rcc@{}}
\toprule\noalign{}
\(a\) & \(m_a\) & \(n_a\) in \(t_a=(m_a+\tau n_a)/2\) \\
\midrule\noalign{}
\endhead
\bottomrule\noalign{}
\endlastfoot
0 & 000 & 100 \\
1 & 100 & 100 \\
2 & 100 & 010 \\
3 & 110 & 010 \\
4 & 110 & 001 \\
5 & 111 & 001 \\
6 & 111 & 000 \\
\end{longtable}

The Riemann constant is determined from this table and parity, rather than assumed from the labels. The theta characteristic \(2P_\infty\) has \(h^0=2\). Its twist by any \(t_a\) is \(P_\infty+P_a\), with \(h^0=1\): Riemann--Roch identifies this dimension with the space of sections of \(3P_\infty-P_a\), a codimension-one subspace of \(\langle1,x\rangle\). The standard Riemann theta correspondence identifies algebraic and analytic characteristic parity; equivalently use Riemann singularity and analytic theta parity. These classical inputs are recalled by Mumford (1971, pp.~181--182) and Casalaina-Martin and Kass (2010, p.~1).

In the convention of DLMF 21.2.5--21.2.7, the analytic characteristic of \((m+\tau n)/2\) is \([n/2;m/2]\) and has parity \(n\cdot m\). This follows directly by replacing the theta summation index by \(-k-n\). If \((m_K,n_K)\) is the characteristic of \(2P_\infty\), its parity is even and all seven twists above are odd. Therefore \[
 n_a\cdot m_a+n_K\cdot m_a+m_K\cdot n_a=1
                         \quad\text{in }\mathbb F_2.
 \tag{T28}
\] The rows \(a=0,1,2,3,4,5\), in succession, force \[
 m_K=(1,0,1),\qquad n_K=(1,1,1).
 \tag{T29}
\] The last row is consistent and this characteristic is even. The table also gives \(t_1+t_3+t_5=K_\infty\). The divisor for \(2P_\infty\) is analytically \(\theta(z+K_\infty,\tau)=0\) up to a nowhere-vanishing factor. Translating its degree-two class by \(t_1+t_3+t_5=K_\infty\) changes it to \(D_{135}\) and adds another \(K_\infty\) to the theta argument. Since \(2K_\infty\) is a period, the divisor becomes \(\theta(z,\tau)=0\). This proves (T5)--(T6). All parity arguments here take place over \(\mathbb C\), not in residue characteristic 2.

\subsubsection{5. Identification of the invariant principal line}\label{identification-of-the-invariant-principal-line}

Use the accepted DP49 identities \(a^\dagger=a^{-1}\), \(w_0+w_0^\dagger=1\), \(w_0w_0^\dagger=2\), and its principal polarization conclusion for \(L_{w_0}\). Symmetry and multiplication isomorphisms are also in accepted DP59, proof of (TS7)--(TS8). The same Poincaré functoriality applies to each good abelian scheme in ALP-LIFT-05. In particular \[
 a^*L_{w_0}\simeq(1,a^\dagger w_0a)^*\mathcal P_\lambda
                 \simeq L_{w_0}.
 \tag{T30}
\] The theta line \(L_\infty\) is symmetric as well: relative duality sends \(2P_\infty\otimes x\) to \(2P_\infty\otimes x^{-1}\), since \(\omega\simeq\mathcal O(4P_\infty)\). This canonical- divisor identity holds on the good models, where the first integral differential has divisor exactly \(4P_\infty\). The automorphism fixes infinity and preserves effectivity, so it preserves this theta divisor.

Two rigidified principal theta lines with the same polarization differ by a rigidified line in \(\operatorname{Pic}^0\). Because both are symmetric, that difference is two-torsion; because both are \(a\)-invariant, it is fixed by the dual \(a\)-action. But \(a-1\) is invertible on two-torsion. Indeed \(\Phi_7(a)=0\), and \[
 (T-1)(T^5+2T^4+3T^3+4T^2+5T+6)=\Phi_7(T)-7
 \tag{T31}
\] shows that \(a^5+a^3+a\) is an inverse to \(a-1\) on \(J[2]\). The dual action has the same property. This is a group-scheme endomorphism identity, so it remains valid on two-torsion in residue characteristic 2. The invariant difference is therefore trivial, proving (T7) integrally with rigidifications.

The space of sections of \(2P_\infty\) is exactly \(\langle1,x\rangle\), so Riemann singularity gives theta multiplicity two at zero. Translation by \(T=K_\infty\) and the definition of the polarization homomorphism give (T8). The explicit table shows \(T\ne0\).

\subsubsection{6. Pullbacks and the exact source limitation}\label{pullbacks-and-the-exact-source-limitation}

For a rigidified symmetric line, \([M]^*L\simeq L^{M^2}\); for a rigidified degree-zero line, \([M]^*\mathcal P_T\simeq\mathcal P_{MT}\). These are the theorem-of-the-cube identities; for \(L_{w_0}\) the first is exactly the integral Poincaré identity in DP59 (TS7). Applied to (T8), they prove (T9) when \(M\) is even.

Choose a primitive local generator of \(e^*\mathcal L_{135}\) to transfer its canonical section to the rigidified line. Pullback by \([M]\) retains the origin value. It retains primitivity as well: multiplication on an abelian scheme is finite faithfully flat, so a nonzero section on the special fibre cannot pull back to zero. The isomorphism (T9) is integral and rigidified and introduces no nonunit origin scalar. This proves its claimed origin valuation. In reduced analytic frames the pullback is composition with \(u\mapsto Mu\); normalization gives \(F_1(Mu)\). The chain rule proves (T10), including the exact factor \(M^{-|\alpha|}\).

For odd \(M\), instead \(MT=T\ne0\), and hence \[
 [M]^*L_{135}^{\mathrm{rig}}
 \simeq L_{w_0}^{M^2}\otimes\mathcal P_T
       \not\simeq L_{w_0}^{M^2}.
 \tag{T32}
\] The directly recovered source file \nolinkurl{ReciprocityTheory/ReductiveGroups/ArithmeticGeometry/BirationalGeometry/ToroidalCompactification/TwoWeightEqOneAddRoot.lean}, lines 53--60, defines \nolinkurl{Level} with \nolinkurl{odd : Odd M} at line 56. Accepted DP59/DP62 likewise use odd levels with theta translations in the specified subgroup of \(J[M^2]\). A theta-group translation acts within the section space of its given line, so it cannot erase the nontrivial difference in (T32). This is the stated limitation. It does not imply that allowed odd torsion translates of the singular principal section have zero origin value; they can have nonzero value, and DP59 already supplies jet-spanning results for their prescribed span.

\subsection{What this closes, and what remains}\label{what-this-closes-and-what-remains}

This identifies the exact algebraic principal line for AM-R-004's nonvanishing analytic germ, computes its entire finite origin normalization on the stable good models, and identifies the different principal line specified by the trace-one construction. Thus this explicit analytic model has neither an unspecified characteristic nor an unspecified finite theta-null scalar.

It does not make the value-one Chern/reduced formal frame integral at nonordinary primes. A unit value at 11 does not imply integral Taylor coefficients or identify the selected gain. Conversion from an integral algebraic frame to the reduced CM frame remains a separate local problem, as AM-DO-013--014 explain. At 2 the norm \(w_0w_0^\dagger=2\) is a nonunit: neither (T3) nor (T7) verifies DP49's unit-parameter primitive-kernel hypotheses. The even-level construction (T9) is outside \nolinkurl{Level.odd}. Finally \(F_1(Mu)\) is a pullback theta, distinct from AM-R-004's other family \(\theta(M^2z,M^2\tau)\). No all-polynomial arithmetic/automorphic kernel identity, full \nolinkurl{MainX}, selected \(g_j,e_j\), source period \(H\), or shorter lift chain is asserted.

\subsection{References and exact source locators}\label{references-and-exact-source-locators}

Casalaina-Martin, S., \& Kass, J. L. (2010). \emph{A Riemann singularity theorem for integral curves} (arXiv:0907.0212v2, 1 June 2010). \url{https://arxiv.org/pdf/0907.0212} . The smooth Riemann singularity statement is recalled on printed p.~1. The determinant/Fitting construction is Definition 3.1 and the ensuing paragraphs, pp.~9--10. Proposition 3.3, p.~11, concerns an equal-characteristic test arc; its elementary determinant idea is extended here by the written DVR argument. No mixed-characteristic scope is attributed to that result.

Mumford, D. (1971). Theta characteristics of an algebraic curve. \emph{Annales scientifiques de l'École Normale Supérieure}, 4(2), 181--192. \url{https://www.numdam.org/article/ASENS_1971_4_4_2_181_0.pdf} . Parity and its quadratic law are on pp.~181--182, formula (*); the proof is in §3, pp.~188--190. Only complex/characteristic-not-two parity is used.

National Institute of Standards and Technology. (n.d.). \emph{Digital Library of Mathematical Functions}, §21.2(ii), equations 21.2.5--21.2.7. \url{https://dlmf.nist.gov/21.2} . These fix characteristic and translation conventions. The parity identity follows directly from the displayed defining series.

The Stacks Project Authors. (n.d.). \emph{Determinants of complexes}, Tag 0FJI. \url{https://stacks.math.columbia.edu/tag/0FJI} . The rank-zero two-term determinant element and its quasi-isomorphism invariance are used. \emph{The relative dualizing sheaf}, Tag 0E6N. \url{https://stacks.math.columbia.edu/tag/0E6N} . Relative duality and its base-change compatibility are used.

Accepted intake: \nolinkurl{RETURN_DP_MULT_checkpoint_06_20261121T0420Z/DP_MULT.tex}, DP49 statement and trace-one polarization proof; DP59 proof (TS7)--(TS8), symmetry and rigidified multiplication; DP62, odd-level canonical translates. These are carried at their accepted status, not newly proved. \nolinkurl{Level.odd} was read directly in the recovered pinned source module named above; no Lean execution is claimed.

\clearpage

\section{ALP-LIFT-08: the normalized actual trace-one section and its complete finite scalar content}\label{alp-lift-08-the-normalized-actual-trace-one-section-and-its-complete-finite-scalar-content}

Run ALP\_MACHINE; \nolinkurl{/root/lift_accounting}, GPT-6 Astra Pro; reasoning effort not separately exposed. This uses ALP-LIFT-05's good models and integral Hodge bases and ALP-LIFT-07's integral identification of the trace-one line with the theta characteristic \(2\infty\). It is a theorem about the actual principal section and its first nonzero jet. It does not assert integral reduced-theta coefficients from an integral first jet.

\subsection{Statement}\label{statement-4}

Let \(C:y^2=x^7-1\), \(J=\operatorname{Jac}(C)\), and \(\omega_i=x^{i-1}dx/y\), \(1\leq i\leq3\). Let \(L\) be the symmetric rigidified trace-one line \(L_{w_0}\) of DP49, identified integrally by ALP-LIFT-07 with the rigidified principal theta line for \(2\infty\). In logarithmic differential coordinates \(u_i\), with \(du_i=\omega_i\) at the origin, put \[
 q_\omega(u)=u_1u_3-u_2^2.
 \tag{Q1}
\] There is a unique algebraic section \[
 s_q\in H^0(J,L)
 \tag{Q2}
\] whose initial homogeneous jet in any origin-normalized formal frame is precisely \(q_\omega\). A formal frame is origin-normalized if its value at the identity agrees with the given rigidification. Changing such a frame cannot alter this first nonzero degree-two jet. The section is defined over \(\mathbb Q\) after using the canonical identification with the \(2\infty\) theta line, and in particular over \(E=\mathbb Q(\zeta_7)\) in the trace-one description.

At a finite place after the good-model extensions of ALP-LIFT-05, let \(s_{\mathrm{prim}}\) be any primitive generator of the rank-one module \(H^0(\mathcal J,\mathcal L)\). Define the scalar content by \[
 c_p(s_q)=v_p(a),\qquad s_q=a\,s_{\mathrm{prim}},\qquad v_p(p)=1.
 \tag{Q3}
\] This is independent of the generator. Then \[
 \boxed{c_2(s_q)=8/7,\qquad c_7(s_q)=-1/2,
             \qquad c_p(s_q)=0\quad(p\nmid14).}
 \tag{Q4}
\] The negative entry means a pole relative to the primitive integral section module. These are finite scalar contents of the entire section, not just valuations of two selected Taylor coefficients.

More explicitly, write \(v_i\) for tangent/logarithmic coordinates dual to the primitive local Hodge basis \(\eta_i\) of ALP-LIFT-05. Only their linear terms are needed. At 2 and 7 respectively, \[
 q_\omega=\frac{\alpha^4}{\iota^2}(v_1v_3-v_2^2),
 \qquad
 q_\omega=\frac{\lambda^4}{\beta^2}(v_1v_3-v_2^2)
          =\lambda^{-3}(v_1v_3-v_2^2),
 \tag{Q5}
\] where \(\alpha^7=-4\), \(\iota^2=-1\), \(\lambda=\zeta_7-1\), and \(\beta^2=\lambda^7\).

If \(\sigma_q\) denotes the reduced CM germ of \(s_q\), its initial term is again \(q_\omega\), since a normalized reduced frame has value one at the origin. The comparison of this germ with the classical hyperelliptic sigma, and integrality of its Hurwitz coefficients, are separate assertions proved by the Kummer agent. Likewise the moved germ called \(\Sigma_{\mathrm{tr}}\) in AM-DO-015 is a scalar multiple \(c\sigma_q\); it is not normalized by (Q1). Root's derivative-Thomae calculation determines that separate scalar. These notations must not be interchanged.

\subsection{Full proof}\label{full-proof}

\subsubsection{\texorpdfstring{1. The integral cohomology at the characteristic \(2\infty\)}{1. The integral cohomology at the characteristic 2\textbackslash infty}}\label{the-integral-cohomology-at-the-characteristic-2infty}

Let \(R\) be one of the complete good-model DVRs. Write \(\mathcal B=\mathcal O_{\mathcal C}(2\infty)\), to distinguish the curve line from the theta line on its Jacobian. On each good model, let \(X\) be its degree-two hyperelliptic coordinate. The modules \[
 U=H^0(\mathcal C,\mathcal B),\qquad
 W=H^0(\mathcal C,\omega_{\mathcal C/R}\otimes\mathcal B^{-1})
 \tag{Q6}
\] have integral bases \[
 U=(1,X),\qquad W=(\eta_1,\eta_2),
 \tag{Q7}
\] where \(\eta_2=X\eta_1\) and \(\eta_3=X^2\eta_1\). Indeed \(X\) has its only pole, of order two, at infinity, and the other curve generator has pole order seven. Thus the functions with pole at most two are precisely \(\langle1,X\rangle\) on every fibre. The integral Hodge forms have respective orders four, two, and zero at infinity, as verified in ALP-LIFT-05; vanishing by at least two selects precisely \(\eta_1,\eta_2\). Cohomology and base change, or reduction followed by Nakayama, proves these are the full integral modules. Relative Serre duality gives \[
 H^1(\mathcal C,\mathcal B)=W^*.
 \tag{Q8}
\] Both cohomology modules are consequently free of rank two, including at 2 and 7.

\subsubsection{2. A two-by-two determinant with integral first-order matrix}\label{a-two-by-two-determinant-with-integral-first-order-matrix}

The determinant construction of the relative theta divisor is the same one used in ALP-LIFT-07: on the completed Picard space at the class of \(\mathcal B\), a perfect two-term free complex computes the universal line's cohomology, and its determinant is a local theta equation. This standard construction is recalled in Casalaina-Martin and Kass (2010, Definition 3.1, pp.~9--10), with its determinant invariance in Stacks Project Tag 0FJI.

At the origin, the cohomology in (Q6)--(Q8) is free. The specialized complex therefore splits into these two cohomology modules, with zero differential, plus an acyclic free summand. To see that this splitting is integral, the images, kernels and cohomology are free over the DVR, so every short exact sequence with free quotient splits. Lifting the chosen decompositions to the completed base and eliminating the invertible acyclic block by row/column operations leaves a two-by-two matrix \(d\), all of whose entries belong to the origin ideal. The eliminated block has unit determinant. Thus the initial quadratic theta equation is, up to an integral unit, the determinant of the linear part \(d_1\).

We give the first-order identification explicitly. A tangent vector to the Picard scheme is a class \(\xi\in H^1(\mathcal O_C)\). If transition functions \(g_{ij}\) for \(\mathcal B\) are deformed to \(g_{ij}(1+\epsilon\xi_{ij})\), then a section \(s\) lifts to first order precisely when the Čech cocycle \(\xi_{ij}s_j\) is a coboundary (a simultaneous overall sign has no effect on the two-by-two determinant). The first differential on cohomology is therefore cup product \[
 U\longrightarrow H^1(\mathcal B),\qquad s\longmapsto\xi\smile s.
 \tag{Q9}
\] Under the Serre duality (Q8), its matrix entries are the pairings of \(\xi\) with the products \(s_a t_b\in H^0(\omega)\), for the bases in (Q7). Those products are exactly \[
 \begin{pmatrix}\eta_1&\eta_2\\\eta_2&\eta_3\end{pmatrix}.
 \tag{Q10}
\] Since \(H^1(\mathcal O_C)\) is integrally dual to \(H^0(\omega)\), in the tangent coordinates \(v_1,v_2,v_3\) dual to these forms the first-order matrix and its determinant are \[
 d_1=\begin{pmatrix}v_1&v_2\\v_2&v_3\end{pmatrix},\qquad
 \det d_1=q_\eta=v_1v_3-v_2^2.
 \tag{Q11}
\] This polynomial has an integral-unit coefficient at every place, including residue characteristic two. In particular it is nonzero there; no division by 2 or characteristic-not-two theorem was used. The argument is the usual determinant/cup-product mechanism of the Riemann--Kempf tangent cone, with the present integral case established directly rather than inferred from a complex tangent-cone statement.

The origin fibre of the determinant theta line identifies with \(\det U^{-1}\otimes\det W^{-1}\). The bases in (Q7) give it a primitive generator. Passing to its rigidification and using the integral isomorphism with \(\mathcal L_{w_0}\) from ALP-LIFT-07, one sees that the degree-two jet map on the actual principal section module has image exactly \[
 R\,q_\eta\subset\operatorname{Sym}^2 e^*\Omega^1_{\mathcal J/R}.
 \tag{Q12}
\] Here a remaining integral unit in the determinant presentation can be absorbed in the primitive section generator. The principal theta section module is free of rank one and formation commutes with base change, by the standard principal-polarization cohomology theorem used in DP59. Therefore (Q12) is a saturated rank-one image; it determines the scalar content of the \emph{entire section} from its quadratic jet.

\subsubsection{3. Algebraic normalization and finite contents}\label{algebraic-normalization-and-finite-contents}

Over the original characteristic-zero field use \(U=(1,x)\), \(W=(\omega_1,\omega_2)\). The same calculation gives the jet line spanned by (Q1). Its generator is defined over \(\mathbb Q\), as are the curve, infinity, the rigidified \(2\infty\) theta line, and this quadratic tensor. The nonzero degree-two jet map from a one-dimensional section space has a unique inverse image of \(q_\omega\). This proves the existence, uniqueness and field assertion in (Q2). A change of origin-normalized formal frame multiplies a local representative by a power series with constant term one, and cannot change its first nonzero degree-two term.

At 2, ALP-LIFT-05 gives \(\omega_i=\alpha^i\iota^{-1}\eta_i\). Thus at the tangent level \(u_i=\alpha^i\iota^{-1}v_i\), and \[
 q_\omega=\alpha^4\iota^{-2}q_\eta,
 \qquad v_2(\alpha^4\iota^{-2})=8/7.
 \tag{Q13}
\] Equation (Q12) proves \(c_2(s_q)=8/7\).

At 7, ALP-LIFT-05 gives \[
 \begin{aligned}
 u_1&=(\lambda/\beta)v_1,\\
 u_2&=(\lambda/\beta)(v_1+\lambda v_2),\\
 u_3&=(\lambda/\beta)(v_1+2\lambda v_2+\lambda^2v_3).
 \end{aligned}
 \tag{Q14}
\] Substitution cancels the \(v_1^2\) and \(v_1v_2\) terms and gives \[
 q_\omega=(\lambda^4/\beta^2)q_\eta
           =\lambda^{-3}q_\eta,
 \qquad v_7(\lambda^{-3})=-1/2.
 \tag{Q15}
\] Again (Q12) proves the stated content. At every other prime the original basis is already primitive, so (Q12) gives content zero. This proves (Q4)--(Q5).

There is also an independent determinant-frame check of the two scalings. At 2, relative to (Q7), the determinant valuations of \((1,x)\) and \((\omega_1,\omega_2)\) are \(2/7\) and \(6/7\). At 7 they are respectively \(1/6\) and \(v_7(\lambda^3/\beta^2)=-2/3\). Their sums are exactly \(8/7\) and \(-1/2\). These sums agree with (Q13) and (Q15), as required by the determinant line \(\det U^{-1}\otimes\det W^{-1}\).

\subsection{Consequence at the actual prime 11, and scope}\label{consequence-at-the-actual-prime-11-and-scope}

The section \(s_q\) is now an explicitly normalized generator of the actual integral principal section module at 11. It can be used as the fixed primitive base section in accepted DP59 and DP62. This removes an unspecified scalar from their application to this Jacobian. Their permitted subgroup, square level, theta-group lift and selected-jet conditions are unchanged. Their integral-span statements are cited as accepted results, not re-proved here.

The conclusion does not infer full reduced-theta integrality from the quadratic term. The scalar content at 7 is measured in the primitive stable Hodge/line model. It is compatible with separate Hurwitz integrality in the original \(\omega\)-coordinates, because those coordinates are not primitive at 7. The dyadic content also says nothing by itself about the norm-two, nonunit trace-one kernel parameter. Finally a principal section and its allowed translates still require an actual theta-lift row/column identification and all global costs before any integer chain improvement follows.

\subsection{Literature and status}\label{literature-and-status}

The determinant and relative-duality inputs are the same directly checked primary sources as in ALP-LIFT-07: Casalaina-Martin and Kass (2010), arXiv:0907.0212v2, Definition 3.1, pp.~9--10; Stacks Project Tags 0FJI and 0E6N. The Čech first-order argument and the integral splitting argument needed here are written in full above.

The classical tangent-cone theorem is due to Kempf, G. (1973), \emph{On the geometry of a theorem of Riemann}, Annals of Mathematics, 98(1), 178--185, \url{https://doi.org/10.2307/1970910} . The original publication metadata was checked on the Annals site; its paywalled full text was not inspected. The related cup-product matrix description was read in Hitching and Hoff, arXiv:1610.03456v2, §3.1, p.~11. Their stated Petri-injective setup is not asserted for this hyperelliptic curve; our integral proof does not use that hypothesis or their reconstruction theorem. No whole-proof check of those literature results is claimed.

\clearpage

\section{AM-R-006: the exact scalar joining the translated theta and the normalized trace-one section}\label{am-r-006-the-exact-scalar-joining-the-translated-theta-and-the-normalized-trace-one-section}

Author: \nolinkurl{/root}, GPT-6 Astra Pro; reasoning effort not externally exposed; run ALP\_MACHINE, 21 November 2026. This finding uses AM-R-004, AM-DO-015--016 and ALP-LIFT-07--08. It removes the remaining scalar between their two actual sections of the same principal theta line. It is not an identification of an iterated automorphic lift image.

\subsection{Statement}\label{statement-5}

Use the period coordinates of AM-R-004, with \(u=Pz\), \(Y=\operatorname{Im}\tau\), and \(Q(z)=\pi z^tY^{-1}z/2\). Put \[
 \theta_0=\theta[0;0](0;\tau),\qquad
 \theta_K(z)=\theta[\tfrac12(1,1,1);\tfrac12(1,0,1)](z;\tau).
 \tag{R6.1}
\] By AM-DO-015 the exact normalized Mumford translate is \[
 \Sigma_{\rm tr}(u)=-i\,e^{Q(P^{-1}u)}
                   \frac{\theta_K(P^{-1}u)}{\theta_0}.
 \tag{R6.2}
\] Let \(\sigma_q\) be the reduced germ of the actual trace-one section \(s_q\) of ALP-LIFT-08, whose first nonzero homogeneous term is \(q(u)=u_1u_3-u_2^2\). Define the nonzero scalar by the exact convergent analytic expression \[
 c=-\frac{i}{\theta_0}
     \left.\partial_{u_1}\partial_{u_3}
         \theta_K(P^{-1}u)\right|_{u=0}.
 \tag{R6.3}
\] This fixes its phase, not merely an unspecified root. Then \[
 \Sigma_{\rm tr}=c\,\sigma_q.
 \tag{R6.4}
\] For \(S=\{1,3,5\}\), \(S^c=\{0,2,4,6\}\), set \[
 R_S=\prod_{a\in S}\prod_{b\in S^c}(\zeta_7^a-\zeta_7^b)
 =49(-\zeta_7^5-\zeta_7^4-2\zeta_7^2+\zeta_7-2).
 \tag{R6.5}
\] The scalar is algebraic and satisfies the exact phase-independent identity \[
 c^8=\frac{R_S^2}{4^8},\qquad
 |c|=\frac{|R_S|^{1/4}}4.
 \tag{R6.6}
\] In normalized valuations \(v_p(p)=1\), at every place above the indicated rational prime, \[
 v_2(c)=-2,\qquad v_7(c)=\frac12,\qquad
 v_p(c)=0\quad(p\ne2,7).
 \tag{R6.7}
\] The primitive integral scalar contents of the translated section are therefore \[
 c_2(\Sigma_{\rm tr})=-\frac67,\qquad
 c_p(\Sigma_{\rm tr})=0\quad(p\text{ odd}).
 \tag{R6.8}
\] Here content is relative to the integral rigidified trace-one line of ALP-LIFT-08. Equation (R6.8) concerns the entire algebraic section, not the coefficient integrality of a Chern formal frame.

The following explicit global estimate now holds for the actual quadratic-normalized trace-one germ: \[
 |\sigma_q(u)|<\frac{6125}{16}
       \exp\!\left(\frac\pi2\sum_{j=1}^3\frac{|u_j|^2}{G_j}\right).
 \tag{R6.9}
\] Consequently, for every multiindex \(\alpha\), \[
 |\partial^\alpha\sigma_q(0)|
 <\frac{6125}{16}\,\alpha!
   \prod_{j:\alpha_j>0}
      \left(\frac{\pi e}{\alpha_jG_j}\right)^{\alpha_j/2}.
 \tag{R6.10}
\] For any positive integer \(M\), substitution of \(Mu\) multiplies the right side before level normalization by \(M^{|\alpha|}\). Thus the source factor \(M^{-|\alpha|}\) cancels this increase exactly; for the actual arithmetic \nolinkurl{Level} one retains its oddness condition. The section \(\sigma_q(Mu)\) is in the correct trace-one line for every such odd \(M\).

\subsection{Proof}\label{proof}

\subsubsection{The derivative Thomae comparison with all differential factors retained}\label{the-derivative-thomae-comparison-with-all-differential-factors-retained}

We use the derivative Thomae formula as a literature theorem, rather than reprove it. Bernatska (2020), arXiv v3, Corollary 2, equations (34)--(37), printed p.~18, and Example 4, equations (38)--(39), printed p.~19, identify the Hessian of the unique singular even theta characteristic in genus three. Her differentials are \[
 d v_1=-\frac{x^2dx}{2y},\qquad
 d v_3=-\frac{x dx}{2y},\qquad
 d v_5=-\frac{dx}{2y};
\] this follows from her equation (2), printed p.~2, with \(f=-y^2+x^7-1\). Hence \[
 (v_1,v_3,v_5)=-\tfrac12(u_3,u_2,u_1).
 \tag{R6.11}
\] The cited Hessian formula says that the degree-two term of \(\theta_K\), in these variables, is \(C_3(v_1v_5-v_3^2)\). The derivative normalization of \(C_3\) is stated immediately after (39). In the original \(u\) variables the term is therefore exactly \[
 \frac{C_3}{4}(u_1u_3-u_2^2).
 \tag{R6.12}
\] An eighth-root phase in the homology/characteristic convention is included in the actual derivative value \(C_3\); equation (R6.3) selects it in our fixed period convention. The quadratic exponential in (R6.2) has constant term one, so it cannot alter (R6.12).

Both \(\Sigma_{\rm tr}\) and \(\sigma_q\) are nonzero sections of the same rigidified principal line, represented in its same reduced frame. Its section space has dimension one. Comparing their first nonzero jets proves (R6.4) with \(c=-iC_3/(4\theta_0)\). This argument does not assume a global algebraic trivialization of that ample line.

Bernatska's Remark 3 and equation (30), printed p.~16, give \(C_3\) as the product of the square root of the period determinant factor and the fourth root of the Vandermonde in all seven finite branch points. The first Thomae formula for the nonsingular characteristic belonging to \(S\) gives the same period determinant factor times the fourth root of the two within-part Vandermondes. In our period convention that nonsingular characteristic is precisely the zero characteristic, by ALP-LIFT-07. Dividing the two formulas cancels the period determinant. The result is \[
 \frac{C_3}{\theta_0}=\epsilon R_S^{1/4},\qquad \epsilon^8=1,
 \tag{R6.13}
\] with the ordering sign absorbed in the root convention. This is also the displayed consequence of first Thomae immediately after (30). Taking absolute values and eighth powers proves (R6.6). In particular the scalar defined analytically in (R6.3) is algebraic. There is no need to select an arbitrary fourth root in place of that exact analytic definition. The cyclotomic reduction in (R6.5) is ordinary polynomial multiplication modulo \(1+T+\cdots+T^6\), and is independently reproduced by the accompanying script.

\subsubsection{Finite valuations and integral section content}\label{finite-valuations-and-integral-section-content}

For distinct seventh roots of unity, their difference is a unit away from 7 and has normalized valuation \(1/6\) at 7. For example \(\zeta^a-\zeta^b=\zeta^b(\zeta^{a-b}-1)\); the quotient of \(\zeta^{a-b}-1\) by \(\zeta-1\) is a unit and \(v_7(\zeta-1)=1/6\). There are twelve factors in (R6.5), so \(v_7(R_S)=2\), with zero valuation elsewhere. Applying valuations to (R6.6) yields \[
 8v_p(c)=2v_p(R_S)-8v_p(4),
\] which proves (R6.7) at every place after any further finite extension. Roots of unity have valuation zero, including at 2, so the phase convention causes no exception.

Scalar contents add under multiplication of a section by a scalar. ALP-LIFT-08 gives content \(8/7\) at 2, \(-1/2\) at 7, and zero elsewhere for \(s_q\). Adding (R6.7) proves (R6.8). This also checks compatibility with the finite normalization of the original value-one section in ALP-LIFT-07: the translated section has exactly the same dyadic content \(-6/7\), and is primitive at 11 and at every other odd prime on the stated good models. The equality of these contents is a computed conclusion; it was not assumed from integrality of a two-torsion translation in characteristic two.

\subsubsection{The archimedean estimate and derivatives}\label{the-archimedean-estimate-and-derivatives}

Every difference of two distinct seventh roots of unity has absolute value at least \(2\sin(\pi/7)>4/5\). The rational trigonometric certificate in AM-R-004 proves this elementary inequality. Thus \[
 |R_S|>(4/5)^{12},\qquad |c|>\frac14(4/5)^3=\frac{16}{125}.
 \tag{R6.14}
\] AM-DO-015's normalized Mumford translation preserves the AM-R-004 Gaussian estimate, including its constant 49. Divide that estimate by (R6.14), and use \(49\cdot125/16=6125/16\), to obtain (R6.9). This bound is valid on all of \(\mathbb C^3\), not only a compact set or a formal neighborhood.

Apply the multivariable Cauchy estimate on a polydisc of radii \(r_j>0\). For \(\alpha_j>0\), minimize \(e^{\pi r_j^2/(2G_j)}r_j^{-\alpha_j}\) at \(r_j^2=\alpha_jG_j/\pi\). For zero entries let \(r_j\) tend to zero after applying Cauchy. This gives (R6.10). Composition with \([M]\) acts on the universal cover by \(u\mapsto Mu\), and the chain rule gives the final level-normalization assertion.

\subsection{Consequence and limits for the main question}\label{consequence-and-limits-for-the-main-question}

The source-line section can now be used with a fixed algebraic quadratic normalization, exact finite scalar content, and an explicit archimedean Cauchy bound. Its relation to the literal finite Schwartz input of AM-DO-016 is \[
 e^{Q(z)}\Theta(\varphi_t;z,\tau)
       =\theta_0 c\,\sigma_q(Pz).
 \tag{R6.15}
\] The exact scalar is retained on the right. This closes the scalar comparison for that actual CM fibre. It neither makes the reduced formal frame integral nor computes the automorphic quotient volume, the terminal mass, a selected lift's actual gain/error, or a common all-polynomial \nolinkurl{MainX} identity. Those are separate questions.

In the Picard and biextension formulation this is an equality of sections of the same rigidified line after the two-torsion change of characteristic; the scalar valuations measure its integral section lattice. In determinant-of-cohomology terms, the chosen quadratic jet is the cup-product determinant of ALP-LIFT-08. The comparison concerns the metric and integral lattice of a line, not an identity of higher cohomological weights. See the exact FGA/SGA7/cote54 locators in \nolinkurl{GROTHENDIECK_DICTIONARY.md}.

\subsection{Primary reference and computation}\label{primary-reference-and-computation}

Bernatska, J. (2020). General derivative Thomae formula for singular half-periods. \emph{Letters in Mathematical Physics, 110}, 2983--3014. \url{https://doi.org/10.1007/s11005-020-01315-8} . The directly inspected version is \url{https://arxiv.org/pdf/1904.09333v3} , equation (2), Remark 3/equation (30), Corollary 2/equations (34)--(37), and Example 4/equations (38)--(39), printed pp.~2, 16, 18--19.

\nolinkurl{scripts/trace_one_theta_scalar.py} performs exact cyclotomic multiplication and an independent 60-digit theta-series diagnostic of the Hessian shape, scalar modulus and eighth-power identity. The latter is explicitly numerical, not an interval or formal proof. The theorem above uses the exact literature formulas and elementary algebra, independently of the numerical diagnostic.

\clearpage

\section{ALP\_MACHINE --- primitive CM reduced theta and the integral-frame boundary}\label{alp_machine-primitive-cm-reduced-theta-and-the-integral-frame-boundary}

Author: GPT-6 Astra Pro, \nolinkurl{/root/definite_oscillator}, inherited effort not separately exposed, 21 November 2026 UTC. New consequences below are labelled AM-DO-011--016. Root's AM-R-004 and the supplied THETA3/DP\_MULT results retain their own proof statuses. This note does not instantiate \nolinkurl{PinX}, \nolinkurl{ClsK.Core}, a primitive trace-one kernel, or the requested lift family.

\subsection{Literature actually read}\label{literature-actually-read}

Bannai, K., \& Kobayashi, S. (2010). Algebraic theta functions and the p-adic interpolation of Eisenstein--Kronecker numbers. \emph{Duke Mathematical Journal, 153}(2), 229--295. \url{https://doi.org/10.1215/00127094-2010-024}. The mathematical text read in this run is arXiv:math/0610163v4 (11 December 2007), \url{https://arxiv.org/pdf/math/0610163}, specifically Definition1.8 pp10--11; §2.1 and Proposition2.1 pp19--21; Theorem2.9 pp25--26; §2.3, condition(28), Lemma2.13 and Propositions2.14--2.15 pp27--29. Published bibliographic metadata was checked on Bannai's own publication page.

The used conclusions are: a reduced theta germ on a CM abelian variety has algebraic coefficients in algebraic eigen-differential coordinates, if its divisor and normalized leading value are algebraic and its line is admissible; prime-to-p torsion translations have an algebraic comparison; formal-parameter integrality requires the stated CM ordinarity condition and an arithmetic unit normalization. Those are cited theorems, not new theorems rederived in this run.

The supplied THETA3 C11 already provides an actual reduced-theta counterexample to automatic integral formal coefficients, and C19 treats integral de Rham CM splittings and derivative estimates. They are carried, not re-proved. In particular C11's y²=x³−x example at good supersingular p=7 has normalized reduced theta coefficient1411/2016 at formal degree9, of valuation−1. This is a concrete warning against promoting a generic algebraic/Chern comparison to an integral one.

\subsection{AM-DO-011. The exact R004 completion is a reduced theta, with algebraic coefficients}\label{am-do-011.-the-exact-r004-completion-is-a-reduced-theta-with-algebraic-coefficients}

Carry AM-R-004's actual Jacobian J of y²=x⁷−1, E=Q(ζ7), Φ=\{1,2,3\}, algebraic differential coordinates u\_i=∫ω\_i, u=Pz, normalized period matrix τ, and Y=Imτ. For κ≥1 write

\begin{verbatim}
θκ(z)=Σ_(n∈Z³) exp(πiκ nᵗτn+2πiκ nᵗz),
Fκ(u)=exp(πκ zᵗY⁻¹z/2) θκ(z)/θκ(0).
\end{verbatim}

The carried quantitative result gives θκ(0)≠0. Define

\begin{verbatim}
Hκ(z,w)=κ zᵗY⁻¹ conjugate(w),
ακ(a+τb)=(−1)^(κ aᵗb),  a,b∈Z³.
\end{verbatim}

Then Fκ is exactly the reduced theta function for the Appell--Humbert data(Hκ,ακ), normalized to value1 at0. It is not merely proportional to an unspecified reduced theta after an omitted quadratic exponential.

\subsubsection{Full calculation of the multiplier}\label{full-calculation-of-the-multiplier}

Direct reindexing gives

\begin{verbatim}
θκ(z+a+τb)/θκ(z)
  =exp(−πiκ bᵗτb−2πiκ bᵗz).
\end{verbatim}

For λ=a+τb, multiplication by the completion factor changes the linear coefficient in z to

\begin{verbatim}
πκ zᵗ(Y⁻¹λ−2ib)=πκ zᵗY⁻¹ conjugate(λ)=πHκ(z,λ).
\end{verbatim}

The difference between the resulting constant term and πHκ(λ,λ)/2 is

\begin{verbatim}
(πκ/2)λᵗY⁻¹(λ−conjugate(λ))−πiκ bᵗτb
   =πiκ λᵗb−πiκ bᵗτb=πiκ aᵗb.
\end{verbatim}

Therefore

\begin{verbatim}
Fκ(u+Pλ)=ακ(λ) exp(πHκ(z,λ)+πHκ(λ,λ)/2) Fκ(u),
\end{verbatim}

which is precisely Bannai--Kobayashi Definition1.8. Since ακ²=1, its semicharacter meets their finite-order hypothesis. The completion's constant is1 at the origin, so its value normalization is the stated one. This proves the assertion. ∎

\subsubsection{Admissibility and algebraicity}\label{admissibility-and-algebraicity}

In the first-argument-linear Appell--Humbert convention used here, Im Hκ is minus κ times the R004 alternating form

\begin{verbatim}
E_R4(x,y)=Tr_(E/Q)((ζ−ζ⁻¹)x conjugate(y)/7).
\end{verbatim}

The sign can be checked on the symplectic periods: Im Hκ(a\_i,b\_j)=−κδ\_ij, while E\_R4(a\_i,b\_j)=δ\_ij. The sign convention has no effect on admissibility. For every a∈E, direct substitution gives E\_R4(ax,y)=E\_R4(x,conjugate(a)y), and therefore the same equality for Im Hκ. This verifies Φ-admissibility in the exact sense of Bannai--Kobayashi §2.1. Its algebraic invariant differentials are the specified ω\_i=x\^{}(i−1)dx/y, and the CM action x↦ζx is defined over E.

For κ=1 the associated symmetric theta line and its divisor are defined over E: AM-DO-012 below identifies every possible symmetric representative with one of the E-rational hyperelliptic theta characteristics. With F1(0)=1, Proposition2.1 applies directly and gives

\begin{verbatim}
F1(u1,u2,u3) ∈ E[[u1,u2,u3]].               (DO11.1)
\end{verbatim}

Here the conclusion concerns the coefficients in the algebraic differential coordinates u, including the ordinary factorial denominators inherent in a Taylor expansion. It does not identify the theta-null period θ1(0) itself as algebraic.

For κ\textgreater1 one must also descend the selected section line, since an arbitrary complex section of L\^{}κ need not be algebraic up to scalar. In this case the descent is explicit. Let Hκ⊂J{[}κ{]} be the subgroup represented by (κ⁻¹Z³)/Z³. Translation z↦z+a/κ fixes θκ. These translations have commuting finite-order lifts to the theta group of L\^{}κ, and the corresponding joint character space in H⁰(J,L\^{}κ) has dimension1. For example the usual κ³ theta-series basis has joint eigencharacters indexed by j∈(Z/κZ)³, and θκ is its j=0 vector. Choose the matching algebraic lifts after a finite torsion-field extension. Such lifts are algebraic: starting from any algebraic isomorphism for a torsion translation, requiring its finite power to equal1 only adjoins a root of an algebraic number; the remaining exact finite commutator relations are the same finite theta-group relations. The chosen simultaneous eigenspace is the kernel of algebraic linear operators and hence is an algebraic line. Its divisor is therefore defined over a finite number field Kκ. Proposition2.1 gives Fκ∈Kκ{[}{[}u{]}{]}. This argument gives a finite descent for each κ, not a level-uniform field-degree or coefficient-height bound.

\subsection{AM-DO-012. A nonzero even theta value is a unit at every good prime for this Jacobian}\label{am-do-012.-a-nonzero-even-theta-value-is-a-unit-at-every-good-prime-for-this-jacobian}

Work over R=Z{[}ζ7,1/14{]}, with the smooth projective hyperelliptic curve C:y²=x⁷−1. Let W\_j=(ζ\^{}j,0), j=0,\ldots,6, and let ∞ be its unique point at infinity. The branch sections are pairwise disjoint over R because only7 ramifies in their differences, and2 is inverted. The standard relative hyperelliptic model is smooth and proper over R. For a subset S of the seven finite branch points with s=\textbar S\textbar≤3, put

\begin{verbatim}
D_S=Σ_(j∈S) W_j+(2−s)∞.
\end{verbatim}

Every geometric fiber has the following exact theta-characteristic classification and cohomology:

\begin{verbatim}
S=∅:             h⁰(O(D_S))=2,
|S|=1 or2:       h⁰(O(D_S))=1,
|S|=3:           h⁰(O(D_S))=0.              (DO12.1)
\end{verbatim}

There are1+7+21+35=64 such characteristics, and they exhaust all of them. In particular the35 non-effective characteristics remain non-effective in every residue characteristic other than2 and7.

\subsubsection{Full proof of the classification and cohomology calculation}\label{full-proof-of-the-classification-and-cohomology-calculation}

The differential dx/y has divisor4∞, so K\_C∼4∞. Also div(x−ζ\^{}j)=2W\_j−2∞. Hence2D\_S∼4∞ and every D\_S is a theta characteristic.

For distinct S,T with cardinalities at most3, the difference of their associated two-torsion points is represented, using2(W\_j−∞)∼0, by

\begin{verbatim}
Σ_(j∈S△T) W_j−|S△T|∞.
\end{verbatim}

Suppose this is principal for a nonempty U=S△T of size at most6. Its divisor is invariant under the hyperelliptic involution ι. If f has that divisor, then ι(f)/f is constant and its square is1. Since the characteristic differs from2, either f∈k(x) or f=y r(x) with r∈k(x). A rational function in x has even order at every branch point, contradicting the required order1 at j∈U. A function y r(x) has odd order at every one of the seven finite branch points, contradicting the required order0 at a point outside U. Thus all64 points are distinct. A genus-three Jacobian has64 geometric two-torsion points in characteristic different from2, and the theta characteristics form their torsor by addition to2∞; the list is complete.

Now compute H⁰(O(D\_S)) directly. The space is stable under ι, so every permitted function decomposes into its even and odd parts f=u(x)+y v(x), each still permitted. The even part has no finite poles, since a pole at W\_j would have even order at least2 while D\_S permits only1. Thus u is a polynomial. Its pole order2deg u at∞ must be at most2−s. This gives dimension2 when s=0, dimension1 when s=1 or2, and dimension0 when s=3.

For the odd part all finite poles must occur among S, and each possible denominator x−ζ\^{}j occurs to exponent at most1. Therefore

\begin{verbatim}
y v(x)=y Q(x)/∏_(j∈S)(x−ζ^j)
\end{verbatim}

for a polynomial Q, unless the odd part is zero. Its order at∞ is2s−7−2deg Q. To be a section of O(D\_S), this order must be at least s−2. It would follow that2deg Q≤s−5, impossible for s≤3. The odd part is therefore zero. This proves(DO12.1) over every indicated geometric fiber. ∎

\subsubsection{The actual algebraic unit rigidification}\label{the-actual-algebraic-unit-rigidification}

For \textbar S\textbar=3, Riemann--Roch gives h¹(O(D\_S))=0 too, because deg D\_S=g−1=2 and h⁰=0. There is an actual relative theta section on the Jacobian attached to D\_S, and its value at the identity is a unit over R. One can see the unit without choosing a transcendental theta null.

Let p:C×\emph{R J→J, let P be the rigidified Poincaré line on C×J, and let M=P⊗pr\_C\^{}\emph{O(D\_S). Take the effective divisor B=6∞. The line M(B) has fiber degree8, hence H¹=0 on every fiber. Thus p\_}M(B) and p}*(M(B)\textbar\_B) are locally free of the same rank6, and the determinant of their restriction map is the usual determinant-of-cohomology theta section. At the identity, its kernel and cokernel are H⁰(O(D\_S)) and H¹(O(D\_S)), both zero on every fiber. The restriction map there is an isomorphism. Its determinant is therefore a unit and supplies the canonical rigidification of the theta line at0 in which the section has value1. This constructs the integral section and its unit value; no assertion that a Chern formal trivialization is integral has been used.

The symmetric principal theta line in R004 must correspond to one of these theta characteristics. Its theta value at0 is nonzero, by R004's exact estimate, so its characteristic has h⁰=0 and therefore belongs to the35 triples. It can consequently be modeled by the corresponding D\_S over R with this unit rigidification. An explicit S is unnecessary for the conclusion, although it can be recovered from the finite level-two theta characteristic. The isomorphism to the complex R004 representative is fixed by the principal polarization and the rigidification; the resulting normalized reduced germ remains exactly F1.

\subsubsection{Explicit branch half-period data in the R004 basis}\label{explicit-branch-half-period-data-in-the-r004-basis}

This calculation makes the finite level-two comparison concrete. Use root's basis a0=γ0,a1=γ2,a2=γ4 and b0=γ1+γ3+γ5,b1=γ3+γ5,b2=γ5. Then the branch points have the following half-periods, with a point denoted(ε;δ) when its z coordinate is(ε+τδ)/2:

\begin{longtable}[]{@{}lll@{}}
\toprule\noalign{}
Point & ε & δ \\
\midrule\noalign{}
\endhead
\bottomrule\noalign{}
\endlastfoot
W0 & 000 & 100 \\
W1 & 100 & 100 \\
W2 & 100 & 010 \\
W3 & 110 & 010 \\
W4 & 110 & 001 \\
W5 & 111 & 001 \\
W6 & 111 & 000 \\
∞ & 000 & 000 \\
\end{longtable}

To prove the first row, integrate on the real path∞→W0 with y=√(x⁷−1):

\begin{verbatim}
∫∞^W0 ω_i=−B(1/2−i/7,1/2)/7.
\end{verbatim}

The gamma reflection identities give

\begin{verbatim}
B(1/2−t,1/2)/B(t,1/2)=tan(πt).
\end{verbatim}

Divide by root's c\_i=−2i\_C B(i/7,1/2)(ζ\^{}i−1)/7. The quotient is−1/{[}2(1+ζ\^{}i){]}=(ζ\^{}i+ζ\^{}\{3i\}+ζ\^{}\{5i\})/2, using(1+ζ)\^{}(−1)=−ζ−ζ³−ζ⁵. Thus AJ(W0−∞)=b0/2. The identity AJ(Wj−W0)=½Σ\_(k\textless j)γ\_k follows by taking half the difference of the two sheet-symmetric branch paths λ\_j−λ\_0. Expanding γ0=a0, γ1=b0−b1, γ2=a1, γ3=b1−b2, γ4=a2, γ5=b2 modulo2 gives every row. This proves the table. No unverified Riemann-constant convention is used to name the triple S of θ{[}0;0{]}.

\subsection{AM-DO-013. An actual pullback family with fixed normalization and ordinary integrality}\label{am-do-013.-an-actual-pullback-family-with-fixed-normalization-and-ordinary-integrality}

Let L,s be the integral symmetric theta line and unit-rigidified section from AM-DO-012 matching F1. For every M≥1, the section {[}M{]}\^{}\emph{s lives in {[}M{]}\^{}}L≃L\^{}(M²), using the rigidified symmetric theorem-of-the-cube isomorphism. Its exact normalized reduced germ is

\begin{verbatim}
G_M(u)=F1(Mu).                              (DO13.1)
\end{verbatim}

It has value1 at0, lies in E{[}{[}u{]}{]}, and satisfies

\begin{verbatim}
M^(−|α|) ∂^α G_M(0)=∂^α F1(0).             (DO13.2)
\end{verbatim}

From AM-R-004's bound at κ=1 one therefore obtains

\begin{verbatim}
|G_M(u)|<49 exp((πM²/2)Σ_i |u_i|²/G_i),
\end{verbatim}

and exactly the same level-normalized derivative bound as R4.11. This is a different actual section family from θ\_(M²)(z); the equality(DO13.2) states its chosen normalization and makes no claim that it equals a primitive trace-one kernel.

Proof. Pullback along{[}M{]} acts on the universal cover by u↦Mu. Its Appell--Humbert Hermitian form is M²H1, and its semicharacter is α1(Mλ)=α1(λ)\^{}(M²), so the reduced multiplier is that of the specified symmetric tensor power. The germ is therefore the displayed pullback, with no additional nonconstant factor; the value normalization fixes the constant. The chain rule proves(DO13.2), and substitution in the carried κ=1 estimate proves the bound. ∎

At a good prime p≠2,7 satisfying Bannai--Kobayashi condition(28), let λ(t) be the formal logarithm in an integral parameter system. The section's arithmetic divisor avoids0, and its normalized value is1. Proposition2.14, applied with local defining function f=1, gives

\begin{verbatim}
F1(λ(t)) ∈ O_v[[t1,t2,t3]]^×.              (DO13.3)
\end{verbatim}

For all M, the formal multiplication map\href{t}{M} has integral coefficients, including when p divides M. Since λ(\href{t}{M})=Mλ(t), composition gives

\begin{verbatim}
G_M(λ(t))=F1(λ([M](t)))∈O_v[[t]]^×.
\end{verbatim}

Thus this family has exact uniform formal integrality at the stated ordinary primes. Prime-to-p torsion translations have the further explicit algebraic/integral comparison of Bannai--Kobayashi Theorem2.9 and Proposition2.15; their theta-group lifts, local divisor functions and torsion field are still part of the datum.

For the primitive type Φ=\{1,2,3\} of Q(ζ7), condition(28) at an unramified rational prime is equivalent to p≡1(mod7). Indeed let D=

⊂(Z/7Z)\^{}× be its decomposition group. The ideals occurring in Φ(p) and conjugate(Φ)(p) are disjoint exactly when ΦD and conjugate(Φ)D are disjoint. Since their union is the whole six-element group, this is equivalent to Φ being a union of D-cosets, or DΦ=Φ. Its stabilizer is trivial: multiplication by2 sends3 to6; by3 sends2 to6; by4 sends1 to4; by5 sends1 to5; by6 sends1 to6. Hence D=\{1\}, equivalently p≡1(mod7). This calculation concerns the exact hypothesis of the cited theorem; it does not assert non-integrality at every other prime.

\subsection{AM-DO-014. What the CM Hodge splitting does and does not settle}\label{am-do-014.-what-the-cm-hodge-splitting-does-and-does-not-settle}

Over a number field containing E and the endomorphisms, H¹\_dR(J) is the sum of six one-dimensional E-eigenlines. Its Hodge subspace is the sum of the three Φ-eigenlines. The sum of the other three eigenlines is the unique E-stable complement. Indeed any other complement is the graph of an E-linear map from that sum to the Hodge subspace, and such a map is zero because their eigencharacters are disjoint. After base change to C the anti-holomorphic Hodge summand has precisely those complementary characters, so it is this same complement.

At a good prime not dividing7, after a finite base change splitting O\_E⊗O\_v, the same idempotents are integral, because the CM algebra is étale there. The Hodge subbundle is saturated in the de Rham lattice. Each of its eigenparts has either full rank1 or rank0, so saturation makes it respectively the entire eigenline lattice or zero. Thus the complementary sum gives an integral E-stable module splitting as well. These assertions prove the proposed comparison of the linear Hodge splittings.

They do not prove that the generic formal horizontal section of a connection, the formal exponential of a Lie splitting, or a chosen Chern/reduced theta frame is an integral formal object. The supplied C11 and C19 already establish this distinction for actual CM elliptic theta functions. It is therefore not enough to invoke uniqueness of the E-linear Hodge complement to discharge all integral-frame hypotheses of the mixed-type arithmetic constructor. The generic algebraicity needed here is independently justified by the exact reduced multiplier and Bannai--Kobayashi Proposition2.1; the ordinary integrality is the narrower result(DO13.3).

For the requested bridge to \nolinkurl{I.J.V.thMoved} and \nolinkurl{PinX}, the remaining tasks are exact and separate: identify the actual primitive/trace-one section and its fixed theta-group translates with this CM model; match the selected local coordinates, determinant roots, additive character and periods; establish the prescribed finite-place derivative or divided-power bounds for the tied family at the nonordinary primes actually used; and prove the all-polynomial kernel identity with the specified volumes and terminal rows. None is implied by a unit theta value at0 or by an algebraic normalized germ alone.

\subsection{AM-DO-015. The exact line-class and parity gate, and its torsion-translation repair}\label{am-do-015.-the-exact-line-class-and-parity-gate-and-its-torsion-translation-repair}

For this paragraph carry ALP-LIFT-07's identification of the R004 zero characteristic with D135=W1+W3+W5−∞ and of the trace-one line Lw0 with the unique ζ-invariant symmetric principal line L0 attached to2∞. Let t be their nonzero two-torsion displacement. In root's period coordinates it has

\begin{verbatim}
t=(ε+τδ)/2,  ε=(1,0,1),  δ=(1,1,1).        (DO15.1)
\end{verbatim}

Thus L1=L\_D135=t\_t\^{}\emph{L0 and t\_t\^{}}L1=L0. The source's actual arithmetic \nolinkurl{Level} requires M odd: \nolinkurl{ReciprocityTheory/ReductiveGroups/ArithmeticGeometry/BirationalGeometry/ToroidalCompactification/TwoWeightEqOneAddRoot.lean}, lines53--62, field \nolinkurl{odd}, line56. Its \nolinkurl{SAB.prim}, line82, also fixes the trace-one line by the cube identity. The actual file was fetched and read in this run; this is not inferred only from the convenient example M=29\^{}e.

Let N=L1⊗L0\^{}(−1), a nontrivial element of Pic⁰(J){[}2{]}. For every M,

\begin{verbatim}
[M]^*N=N^M,
[M]^*L1 ≃ L0^(M²) ⊗ N^M.                 (DO15.2)
\end{verbatim}

Consequently the value-one pullback family F1(Mu) lies in L0\^{}(M²) for even M, but remains in a different line for odd M. Theta-group actions on L0\^{}(M²), including those in \nolinkurl{IsVtOver}, preserve that line and cannot by themselves change its Pic⁰ class. Even M is therefore a valid geometric way to remove this particular twist, but it does not instantiate the current arithmetic \nolinkurl{Level}.

Proof. Pullback by{[}M{]} on Pic⁰ is multiplication by M, so{[}M{]}\^{}\emph{N=N\^{}M. The symmetric rigidified multiplication isomorphism gives{[}M{]}\^{}}L0=L0\^{}(M²). Their tensor product proves(DO15.2). As N has exact order2, N\^{}M is trivial exactly for even M. A theta-group action is, by definition, an automorphism of sections of its specified line via an isomorphism of a translate of that same line. It does not identify two different Pic⁰ classes. This proves the stated obstruction to that direct substitution. ∎

\subsubsection{Repair by changing the base section to the actual trace-one line}\label{repair-by-changing-the-base-section-to-the-actual-trace-one-line}

There is a precise algebraic torsion translation that fixes the line issue without changing the odd arithmetic level. Let H(z,w)=zᵗY⁻¹ conjugate(w), Q(z)=πzᵗY⁻¹z/2, and write F1 also for its expression in z. Define the normalized reduced translation

\begin{verbatim}
Σ(z)=exp(−πH(z,t)−πH(t,t)/2) F1(z+t).       (DO15.3)
\end{verbatim}

This is Bannai--Kobayashi's U\^{}M\_t convention, equations(14) and following pp14--15. Since t\_t\^{}*L1=L0, it is a nonzero section of the actual trace-one line L0. It vanishes at0, with the expected effective even theta characteristic; it is not asserted to have value1 there.

Its exact theta-series formula is

\begin{verbatim}
Σ(z)=−i · exp(Q(z))
      θ[δ/2; ε/2](z;τ) / θ[0;0](0;τ).    (DO15.4)
\end{verbatim}

To check the constant, substitute t=(ε+τδ)/2. The linear part of Q(z+t)−πH(z,t)−Q(z) is πiδᵗz. Its constant part after subtracting πH(t,t)/2 is(πi/4)(δᵗτδ+εᵗδ). By the definition of a theta characteristic,

\begin{verbatim}
θ[δ/2;ε/2](z;τ)
  =exp(πiδᵗτδ/4+πiδᵗz+πiδᵗε/2) θ(z+t;τ).
\end{verbatim}

The quotient of these two exponentials is exp(−πiεᵗδ/4)=−i because εᵗδ=2. This proves(DO15.4), including its scalar.

Bannai--Kobayashi Theorem2.9 applies to this two-torsion translation and shows that the coefficients of Σ in the algebraic u=Pz coordinates are algebraic over E(J{[}16{]}); an extra root-of-unity convention for U\_t rather than U\^{}M\_t changes only the stated finite torsion scalar. This is a generic algebraic statement. We do not infer that every coefficient, or its chosen integral-frame scalar, is a unit at a nonordinary prime.

The SAME Gaussian constant49 is preserved. Indeed R004 gives \textbar F1(z+t)\textbar\textless49 exp(πH(z+t,z+t)/2). Taking the absolute value of the exponential in(DO15.3) subtracts πReH(z,t)+πH(t,t)/2; the cross terms cancel exactly, leaving

\begin{verbatim}
|Σ(u)|<49 exp((π/2)Σ_i |u_i|²/G_i).         (DO15.5)
\end{verbatim}

For every allowed odd M=29\^{}e from carried DP59, set Σ\_M(u)=Σ(Mu). This is literally{[}M{]}\^{}*Σ in L0\^{}(M²). The standard rigidified multiplication isomorphism, with t′=t=0 and the identity theta action, verifies the generic \nolinkurl{IsVtOver} clauses, exactly as in DP59(c). Its normalized derivatives satisfy

\begin{verbatim}
M^(−|α|) ∂^αΣ_M(0)=∂^αΣ(0),
\end{verbatim}

and the R004 level-independent Cauchy bound holds with the same49. This gives an actual algebraic base section, in the correct trace-one line, and an actual generic permitted pullback family. It still does not prove membership of the selected row in an iterated automorphic theta image, any prescribed nonzero positive-degree jet, or the global \nolinkurl{PinX.mv} identity for every moved lattice.

\subsection{AM-DO-016. A literal finite Schwartz input and the retained CM period}\label{am-do-016.-a-literal-finite-schwartz-input-and-the-retained-cm-period}

The trace-one section in(DO15.4) has a concrete theta input. Fix the classical additive-character convention with ψ\_∞(x)=exp(2πix), so ψ\_f(r)=exp(−2πir) on rational r. Put α=δ/2 and β=ε/2, and define the finite Schwartz function on A\_f³ by

\begin{verbatim}
φ_t(q_f)=−i·1_(α+Zhat³)(q_f) ψ_f(−q_fᵗβ). (DO16.1)
\end{verbatim}

This is locally constant and compactly supported. Its values on rational points q=n+α are−i exp(2πi(n+α)ᵗβ), hence algebraic fourth roots of unity. Its rational theta sum is exactly

\begin{verbatim}
Σ_(q∈Q³) φ_t(q) exp(πi qᵗτq+2πi qᵗz)
   =−i θ[α;β](z;τ).                        (DO16.2)
\end{verbatim}

Multiplying by the actual Chern completion gives

\begin{verbatim}
exp(Q(z)) Θ(φ_t;z,τ)=θ[0;0](0;τ)·Σ(z).     (DO16.3)
\end{verbatim}

All derivatives of(DO16.3) follow from normally convergent theta series. Thus in this explicit one-fiber Schrödinger/Chern convention the analytic-to-normalized-algebraic comparison constant is precisely

\begin{verbatim}
c0=θ[0;0](0;τ),  |c0|>5/8.                (DO16.4)
\end{verbatim}

It has not been set to1. The phase−i is part of the finite input, so no phase is omitted from c0. The finite input is a phase and a half-lattice translation of the standard input1\_Zhat³, hence belongs to its finite Heisenberg orbit. It is supported in (1/2)Zhat³ and invariant under2Zhat³, so its finite Schrödinger support/period quotient is (1/2)Z³/2Z³. This is a modulus4 statement for the input function, not a claimed level4 stabilizer for the full metaplectic group action.

Proof. Compact support and local constancy follow from the coset of the compact open subgroup Zhat³ and the finite-conductor additive character. At a rational point the stated product formula for ψ gives its displayed value. Its support on Q³ is precisely α+Z³, so writing q=n+α identifies the sum term by term with the defining series for the theta characteristic. This proves(DO16.2), and substitution of(DO15.4) proves(DO16.3). The inequality in(DO16.4) is carried R004. ∎

This is a genuine fiberwise theta equality with an explicit finite input and an explicit retained period. It is more than an assertion of abstract nonvanishing. It is still less than the required common arithmetic \nolinkurl{ClsK.Core} for the whole family. In particular the recovered Alpöge convention is \nolinkurl{IsStdChar ψ (-1)}; a comparison with it must track the conjugate additive character and the associated polarization or coordinate convention. Nor have the other moved CM lattices, their covolumes, the fixed all-degree branching maps or the terminal mass matrix been identified by(DO16.3). A common scalar for one fixed Heisenberg module does not supply these further identities.

\clearpage

\section{AM-K012--016 --- actual reduced theta germs and derivative bounds on y²=x⁷−1}\label{am-k012016-actual-reduced-theta-germs-and-derivative-bounds-on-yuxb2xux20771}

Author: GPT-6 Astra Pro, ALP\_MACHINE agent /root/kummer\_extension, 21 November 2026 UTC. This note supplies actual germs on the specified Jacobian. The accepted Alpöge theorem is not re-examined. The distinction between the translated non-effective characteristic and the source trace-one line is maintained.

\subsection{0. Statements and conventions}\label{statements-and-conventions}

Let \(E=\mathbb Q(\zeta)\), \(\zeta=\zeta _7\), and \[
C:y^2=x^7-1,\quad J=\operatorname{Jac}(C),\quad
\omega_i=x^{i-1}dx/y,\quad u_i=\int\omega_i\quad(1\leq i\leq3).
\] Write \(P_a=(\zeta^a,0)\). The translated characteristic is \(D=P_1+P_3+P_5-\infty\). Its value-one Chern-reduced germ is AM-R-004's \[
F(u)=e^{\pi z^tY^{-1}z/2}\frac{\theta(z,\tau)}{\theta(0,\tau)},
\qquad z=P^{-1}u,\quad Y=\operatorname{Im}\tau.
\tag{K12.1}
\] ALP-LIFT-07 identifies this characteristic exactly and proves that the value-one normalization is primitive at every odd good prime. The source trace-one line is instead \(L_\infty=L_{2\infty}\). Let \(s_q\) denote its theta section normalized by the determinant-of-cohomology tangent cone \[
q(u)=u_1u_3-u_2^2.
\] Its Chern-reduced germ is denoted \(\sigma_q(u)\). ALP-LIFT-08 identifies this exact scalar normalization with an actual primitive section at every \(p\nmid14\).

A ``Hurwitz coefficient'' means \(\partial^\alpha f(0)\), so \(f=\sum_\alpha\partial^\alpha f(0)u^\alpha/\alpha!\). It does \textbf{not} mean the divided Taylor coefficient. Valuations satisfy \(v_{\mathfrak p}(p)=1\).

The results are:

\begin{enumerate}
\def\labelenumi{\arabic{enumi}.}
\tightlist
\item
  \textbf{AM-K012.} The canonical Schiffer bidifferential equals the explicit algebraic Baker bidifferential. Consequently \(F\) is precisely the Baker even sigma function for \(D\), with no unrecorded quadratic exponential. An algebraic recursion determines all its Hurwitz coefficients, which belong to \(\mathcal O_E[1/14]\). In particular \[
  \partial_1\partial_3F(0)=-\zeta^2/4.
  \tag{K12.2}
  \]
\item
  \textbf{AM-K013.} At every prime above 11, \[
  v_{\mathfrak p}(\partial_i^nF(0))
  \geq\max\left\{0,c_in-\frac{s_{11}(n)}{10}\right\},
  \quad c_1=c_3=\frac1{20},\quad c_2=\frac1{220}.
  \tag{K13.1}
  \] Here \(s_p(n)\) is the base-\(p\) digit sum. The proof uses a stable integral ring for the actual logarithmic Hessian and three finite exact coefficient checks.
\item
  \textbf{AM-K014.} The actual source-line germ is \[
  \sigma_q(u)=4\sigma(u_1/2,u_2/2,u_3/2),
  \tag{K14.1}
  \] where \(\sigma\) is the algebraically normalized \((2,7)\) sigma function in half-differential coordinates. An explicit integer determinant formula determines every coefficient. Thus \(\partial^\alpha\sigma_q(0)\in2^{2-|\alpha|}\mathbb Z\). It is even, has CM weight 4 modulo 7, and \[
  \partial_1\partial_3\sigma_q(0)=1,\quad
  \partial_2^2\sigma_q(0)=-2,\quad
  \partial_1^4\sigma_q(0)=2,\quad
  \partial_3^6\sigma_q(0)=1.
  \tag{K14.2}
  \]
\item
  \textbf{AM-K015.} The bounds (K13.1) hold for \(\sigma_q\) too, by the actual integral algebraic ratio \[
  \sigma_q(2u)=r(u)F(2u).
  \tag{K15.1}
  \] Pullback by \([2]\) trivializes the two-torsion difference of the lines. This is a comparison of germs; the final section belongs to \(L_\infty\), and no even value is inserted into the source's odd Level.M.
\item
  \textbf{AM-K016.} Each actual pure restriction \(f_i(t)=\sigma_q(te_i)\) is an entire function of order 2 and exact type \(\pi/(2G_i)\). Its nonzero coefficient orders form an infinite computable sequence \(n_{i,k}\) satisfying \(n_{i,k+1}/n_{i,k}\to1\), within the necessary congruence class in (K14.11). The coefficient limsup is \[
  \limsup_{n\to\infty}
    \left(\frac{|\partial_i^n\sigma_q(0)|}{\sqrt{n!}}\right)^{1/n}
    =\sqrt{\pi/G_i}.
  \tag{K16.1}
  \] The support and growth assertions are archimedean; they do not bound finite-place valuations above.
\end{enumerate}

These statements do not establish nonvanishing in every admissible degree, identify a prescribed automorphic harmonic row, or prove the terminal every-\(\varepsilon\) window.

\subsection{1. AM-K012: identify the actual reduced frame}\label{am-k012-identify-the-actual-reduced-frame}

Set \(f(x)=x^7-1\) and \(\mathcal F(x,z)=x^3z^3(x+z)-2\). The Baker fundamental bidifferential is \[
\Omega_B(P,Q)=
\frac{\mathcal F(x,z)+2y(P)y(Q)}
 {4(x-z)^2y(P)y(Q)}\,dx\,dz.
\tag{K12.3}
\] The \(2y(P)y(Q)\) term removes a spurious pole on the conjugate diagonal. This symmetric differential has diagonal leading term \(dt_Pdt_Q/(t_P-t_Q)^2\) and no other poles. Exact general formulas: Buchstaber, Enolskii, and Leykin (1997), §1.2, equations (1.7)--(1.9); Edelstein, Gómez-Reino, and Mariño (2001), equations (3.3)--(3.6).

Let \(B_0\) be the normalized fundamental bidifferential with zero \(a\)-periods and let \(\nu=P^{-1}\omega\) be the column of normalized holomorphic differentials. The Schiffer bidifferential is \[
\Omega_S=B_0-\pi\nu(P)^tY^{-1}\nu(Q).
\tag{K12.4}
\] It is independent of the symplectic marking. Its period relation pairs its periods with those of the intrinsic Bergman reproducing kernel and characterizes it without a marking. Shramchenko (2004), §4, printed p.1468, gives (K12.4) and these period relations. Thus every curve automorphism preserves \(\Omega_S\).

The automorphism \(a:(x,y)\mapsto(\zeta x,y)\) preserves (K12.3): \(\mathcal F\) has only terms of degree 7 and a constant, and \(dx\,dz/(x-z)^2\) is invariant. Moreover \(a^*\omega_i=\zeta^i\omega_i\). The difference \(\Omega_B-\Omega_S\) is a global holomorphic bidifferential, hence a linear combination of \(\omega_i(P)\omega_j(Q)\). Every term has weight \(i+j\in\{2,3,4,5,6\}\), none divisible by 7. Invariance forces every coefficient to vanish: \[
\boxed{\Omega_B=\Omega_S.}
\tag{K12.5}
\] This uniqueness proof makes no inference about finite-place integrality from a Hodge splitting.

The quadratic normalization includes a factor of two. If Baker's sigma theta expression has exponential \(e^{u^tKu}\), its normalized bidifferential is \[
B_0=\Omega_B+2\omega(P)^tK\omega(Q).
\] This follows by taking \(a\)-periods of the second-kind differential identity, equivalently Edelstein et al., equations (3.14)--(3.18). Comparison gives \[
2K=\pi P^{-t}Y^{-1}P^{-1}.
\tag{K12.6}
\] Hence the exponential is exactly \(e^{\pi z^tY^{-1}z/2}\). The characteristic is zero by ALP-LIFT-07, and normalization at zero divides by \(\theta(0,\tau)\). This proves (K12.1). The same calculation applies to the characteristic \(2\infty\).

\subsection{2. Actual algebraic logarithmic Hessian and recursion}\label{actual-algebraic-logarithmic-hessian-and-recursion}

Let \(a_1=\zeta,a_2=\zeta^3,a_3=\zeta^5\). Near the ordered divisor \(((a_1,0),(a_2,0),(a_3,0))\), put \[
U(X)=\prod_j(X-x_j),\qquad
\ell_j(X)=\frac{\prod_{m\ne j}(X-x_m)}{U'(x_j)}
         =\sum_{i=1}^3q_{ji}X^{i-1}.
\] The matrix \(q\) is the inverse of \(V_{ij}=x_j^{i-1}\). The invariant derivatives are \[
D_ix_j=y_jq_{ji},\qquad D_iy_j=\tfrac72x_j^6q_{ji}.
\tag{K12.7}
\] The first formula inverts \(\sum_jx_j^{k-1}dx_j/y_j=du_k\); the second differentiates \(y_j^2=x_j^7-1\). They are regular at the branch points. The \(D_i\) commute because they lift the commuting invariant derivatives on \(J\).

Let \(s_1=\sum x_j,s_2=\sum_{j<m}x_jx_m,s_3=x_1x_2x_3\). The symmetric Kleinian matrix \(W\) is determined by \[
W_{33}=s_1/4,\quad W_{23}=-s_2/4,\quad W_{13}=s_3/4
\tag{K12.8}
\] and \[
\sum_{i,k}W_{ik}x_j^{i-1}x_m^{k-1}
=\frac{\mathcal F(x_j,x_m)-2y_jy_m}{4(x_j-x_m)^2},
\qquad j<m.
\tag{K12.9}
\] These are Buchstaber et al.~(1997), Corollary 2.1.2, equation (2.9), and Theorem 2.2, equations (2.10)--(2.14). Their canonical model has leading coefficient 4. Passing from \(Y^2=4(x^7-1)\), \(du^{\rm can}=dx/Y\), to our \(u=2u^{\rm can}\) divides the logarithmic Hessian by 4 and gives exactly these factors.

After the last row is inserted, the unknowns \(W_{11},W_{12},W_{22}\) have coefficient rows \((1,x_j+x_m,x_jx_m)\). The determinant is a signed product of the three differences. Here is explicit elimination: \[
\begin{split}
R_{jm}={}&\frac{\mathcal F(x_j,x_m)-2y_jy_m}{4(x_j-x_m)^2}
-W_{13}(x_j^2+x_m^2)\\
&-W_{23}(x_j^2x_m+x_jx_m^2)-W_{33}x_j^2x_m^2,\\
t_1={R_{13}-R_{12}\over x_3-x_2},\qquad
t_2={R_{23}-R_{12}\over x_3-x_1}.
\end{split}
\] Then \[
W_{22}={t_1-t_2\over x_1-x_2},\quad W_{12}=t_1-x_1W_{22},\quad
W_{11}=R_{12}-(x_1+x_2)W_{12}-x_1x_2W_{22}.
\tag{K12.10}
\]

For \(h=\log F\), \(h(0)=0\), the characteristic translation gives \[
\partial_i\partial_jh(u)=-W_{ij}(a+u),
\quad a=\operatorname{AJ}(P_1+P_3+P_5-3\infty).
\tag{K12.11}
\] This is the Hessian of the specified actual theta germ. Translating changes an affine exponential only, whose Hessian is zero; evenness and \(F(0)=1\) fix the first and zeroth derivatives. Thus (K12.11) determines \(F\) uniquely.

\subsubsection{Stable integral ring}\label{stable-integral-ring}

Let \(R_0=\mathcal O_E[1/14]\), and \[
\mathcal R=R_0[x_1,x_2,x_3,y_1,y_2,y_3,
(x_1x_2x_3\prod_{j<m}(x_j-x_m))^{-1}]
/(y_j^2-x_j^7+1)_{j=1}^3.
\tag{K12.12}
\] Every \(D_i\) preserves this ring, by (K12.7), the quotient rule, and differentiation of the defining relations. Every \(W_{ij}\) belongs to it. Evaluation \(e:\mathcal R\to R_0\) at \((x_j,y_j)=(a_j,0)\) is defined: these roots are units, and root differences have prime divisors only above 7.

The Abel map from this ordered-divisor open set to \(J\) is étale. In a branch chart the local parameters \(y_j\) give the differential matrix \[
\left(2x_j^{i-1}/(7x_j^6)\right)_{ij},
\] with invertible determinant. At nonbranch points use \(x_j\); the same Vandermonde criterion applies. This concerns algebraic formal coordinates on \(J\); it is \textbf{not} a claim that the Abelian logarithms \(u_i\) are integral formal parameters.

For \(|\alpha|\geq2\), choose \(i,j\) with \(\alpha-e_i-e_j\geq0\). Then \[
h_\alpha:=\partial^\alpha h(0)
=-e(D^{\alpha-e_i-e_j}W_{ij})\in R_0.
\tag{K12.13}
\] The result is independent of the choice by the Kleinian Hessian identity. Also \(h_0=h_{e_i}=0\). If \(S_\alpha=\partial^\alpha F(0)\), then \[
S_0=1,\qquad
S_\alpha=\sum_{0\leq\beta\leq\alpha-e_i}
{\alpha-e_i\choose\beta}h_{\beta+e_i}S_{\alpha-e_i-\beta}
\quad(\alpha_i>0).
\tag{K12.14}
\] This is \(\partial_iF=F\partial_ih\). It involves integer binomial coefficients and no division by an order, proving \(S_\alpha\in R_0\).

\subsubsection{Initial matrix and checks}\label{initial-matrix-and-checks}

Let \(Q(X)=\prod_{a\in\{1,3,5\}}(X-\zeta^a)\), \(R(X)=f(X)/Q(X)\). Bolza's formula is \[
\sum_{i,j}W_{ij}(a)X^{i-1}Z^{j-1}
=\frac{\mathcal F(X,Z)-Q(X)R(Z)-Q(Z)R(X)}{4(X-Z)^2}.
\tag{K12.15}
\] See Edelstein et al.~(2001), equation (3.24), reversing their descending basis. Explicitly, \(4W(a)\) is \[
\begin{pmatrix}
-\zeta^5-\zeta^4-2\zeta^2+\zeta-2&\zeta^3+1&\zeta^2\\
\zeta^3+1&-2\zeta^4+\zeta^3-\zeta^2+\zeta-2&
\zeta^5+\zeta^3+\zeta^2+1\\
\zeta^2&\zeta^5+\zeta^3+\zeta^2+1&\zeta^5+\zeta^3+\zeta
\end{pmatrix}.
\tag{K12.16}
\] This proves (K12.2). Exact polynomial division modulo \(\Phi_7\) is checked in \nolinkurl{jacobian_theta_hessian_check.py.}

A separate numerical diagnostic, not a proof premise, used 65-digit independent theta summation with AM-R-004's periods. Its Hessian agrees with \(-W(a)\) to \(6.85\cdot10^{-66}\). Pure jets of orders 2,4,6,8 agree with the independent flow/Bell recurrence, largest observed discrepancy \(4.55\cdot10^{-64}\). Results and diagnostic scope are retained in jacobian\_theta\_hessian\_check.json.

\subsection{3. AM-K013: invariant-derivative nilpotence and actual logarithmic bounds}\label{am-k013-invariant-derivative-nilpotence-and-actual-logarithmic-bounds}

Fix an odd prime \(p\ne7\) and complete \(R_0\) at a prime above \(p\). The reduction is a smooth abelian variety with the order-seven automorphism. The \(D_i\) form an integral invariant basis, with distinct CM labels \(i\in\{1,2,3\}\).

In characteristic \(p\), the \(p\)-fold composition of a derivation is a derivation, because the intermediate binomial coefficients in the Leibniz formula vanish. Translation invariance is retained. The CM weight is multiplied by \(p\), so a label \(pi\bmod7\) outside \(\{1,2,3\}\) forces \(D_i^p=0\). If the label is inside this set, \(D_i^p\) is a scalar multiple of its uniquely labelled invariant derivative. Iteration gives \[
D_i^{q_i}=0\text{ on the special fibre},\qquad q_i=p^{h_i},
\tag{K13.2}
\] with these sufficient exponents:

\begin{longtable}[]{@{}lrrr@{}}
\toprule\noalign{}
\(p\bmod7\) & \(h_1\) & \(h_2\) & \(h_3\) \\
\midrule\noalign{}
\endhead
\bottomrule\noalign{}
\endlastfoot
2 & 2 & 1 & 1 \\
3 & 3 & 1 & 2 \\
4 & 1 & 2 & 1 \\
5 & 1 & 3 & 2 \\
6 & 1 & 1 & 1 \\
\end{longtable}

For \(p\equiv1\bmod7\), this argument supplies no exponent. The table is obtained by multiplying a label successively by \(p\bmod7\) until it leaves \(\{1,2,3\}\); no classification of the Newton polygon is needed.

Because the Abel map in §2 is étale, the zero derivation on the target has only the zero lift. Thus (K13.2) holds on \(\mathcal R/p\mathcal R\), and \[
D_i^{q_i}\mathcal R\subset p\mathcal R.
\tag{K13.3}
\] This is an operator assertion on a stable integral ring. Since \(D_i\mathcal R\subset\mathcal R\) and \(D_i(p)=0\), it may genuinely be iterated: \[
D_i^m\mathcal R\subset p^{\lfloor m/q_i\rfloor}\mathcal R.
\tag{K13.4}
\] For \(\ell_n=\partial_i^n\log F(0)\), applying this to \(W_{ii}\) gives \[
v_{\mathfrak p}(\ell_n)\geq
\max\{0,\lfloor(n-2)/q_i\rfloor\}\quad(n\geq2).
\tag{K13.5}
\] Also \(\ell_n=0\) for odd \(n\).

\subsubsection{Bell--Legendre lemma, with all losses}\label{belllegendre-lemma-with-all-losses}

Let \(p\) be odd, \(q=p^h\), and \[
H(t)=\sum_{j\geq2}\ell_jt^j/j!,\qquad
e^{H(t)}=\sum_{n\geq0}S_nt^n/n!.
\] Suppose \(H\) is even, all \(\ell_j\) are integral, and (K13.5) holds. Then \[
v(S_n)\geq
\max\left\{0,\frac{n}{(p-1)(q+1)}-\frac{s_p(n)}{p-1}\right\}.
\tag{K13.6}
\] If additionally \(v(\ell_{q+1})\geq1\), then \[
v(S_n)\geq
\max\left\{0,\frac{n}{2p^{h-1}(p-1)}-\frac{s_p(n)}{p-1}\right\}.
\tag{K13.7}
\]

To prove this, a complete Bell-polynomial term, indexed by \(m_j\geq0\), \(\sum jm_j=n\), is \[
\frac{n!}{\prod_jm_j!(j!)^{m_j}}\prod_j\ell_j^{m_j}.
\] Only even \(j\) occur. Legendre's formula bounds its valuation below by \[
-\frac{s_p(n)}{p-1}
+\sum_jm_j\left(v(\ell_j)+\frac{s_p(j)-1}{p-1}\right)
+\frac{\sum_js_p(m_j)}{p-1}.
\tag{K13.8}
\] The final summand is nonnegative.

For (K13.6), even \(j>0\) has \(s_p(j)\geq2\). If \(j\leq q+1\), then \((p-1)v(\ell_j)+s_p(j)-1\geq1\geq j/(q+1)\). If \(j\geq q+3\), put \(k=\lfloor(j-2)/q\rfloor\geq1\). Since \(j\leq(k+1)q+1\), \[
(p-1)v(\ell_j)+s_p(j)-1\geq(p-1)k+1\geq j/(q+1).
\] Insert this into (K13.8). Integrality supplies the maximum with zero.

For (K13.7), if \(2\leq j<q\), then \(j\leq p^{h-1}s_p(j)\), while \(s_p(j)-1\geq s_p(j)/2\). Thus the lower bound is \(j/(2p^{h-1})\). The sole even gap is \(j=q+1\): the additional assumption gives a lower bound at least \(p\), larger than \((q+1)/(2p^{h-1})\). For \(j\geq q+3\), retain \(k=\lfloor(j-2)/q\rfloor\). If \(k=1\), evenness gives \(j\leq2q\), so the lower bound \(p\) suffices. If \(k\geq2\), \[
(p-1)k+1-\frac{(k+1)q+1}{2p^{h-1}}
=\left(\frac p2-1\right)k+1-\frac p2-\frac1{2p^{h-1}}\geq0.
\] This holds for \(p\geq3,k\geq2,h\geq1\); its smallest case \(p=3,k=2,h=1\) is zero. Substitution into (K13.8) proves (K13.7).

Therefore (K13.6) is already an unconditional directional lower bound for the actual translated theta at every residue class in the table, including \(p\equiv3,5\bmod7\).

\subsubsection{Exact finite certificate for the improved 11-adic constants}\label{exact-finite-certificate-for-the-improved-11-adic-constants}

At \(p=11\), \(q_1=q_3=11\) and \(q_2=121\). The only additional checks needed are \[
\partial_1^{12}\log F(0)\equiv0,\qquad
\partial_3^{12}\log F(0)\equiv0,\qquad
\partial_2^{122}\log F(0)\equiv0\pmod{11}.
\tag{K13.9}
\] All three were verified by exact arithmetic in \[
B=\mathbb F_{11}[Z]/(1+Z+\cdots+Z^6),
\] the product of the two residue fields above 11. Zero in \(B\) proves zero at both primes.

The reproducible recurrence is \nolinkurl{jacobian_theta_direction_jets.py.} Its inputs are (K12.7)--(K12.10). To avoid power-series division, add \(z_{jm}=(x_j-x_m)^{-1}\) with \[
D_iz_{jm}=-z_{jm}^2(D_ix_j-D_ix_m).
\] This is a polynomial differential system over \(B\). The \(n\)-th Hurwitz coefficient of a product is \(\sum_{k=0}^n\binom nk a_kb_{n-k}\). The coefficient of order \(n\) in a variable's prescribed derivative gives its coefficient of order \(n+1\), with no division by \(n+1\). Initial branch-difference inverses are obtained by Euclidean polynomial inversion modulo \(\Phi_7\). Finally \(\ell_{n+2}=-D_i^nW_{ii}|_a\).

The executed command was:

\begin{verbatim}
python work/agents/jacobian_theta_direction_jets.py --prime 11 --order 160 --output jacobian_theta_direction_jets_mod11.json
\end{verbatim}

It retains all nonzero coefficient vectors. In each of (K13.9), all six coordinates in \(1,Z,\ldots,Z^5\) are zero. The companion \nolinkurl{jacobian_theta_11_certificate.py} checks the saved tables, checks the unit status of the retained coefficients, and regenerates the three exceptional coefficients using the stated recurrence.

The operator proof (K13.3)--(K13.5), the three exact checks, and (K13.7) now prove (K13.1) for every \(n\). The finite computation is used only for these three coefficients; no pattern is extrapolated.

There are exact all-order mod-11 consequences as well. In the reduced one-variable Hurwitz algebra, \(b^{11}=b(0)^{11}\): a positive divided-power monomial has zero 11-th power, as follows from \(v_{11}((11k)!/(k!)^{11})=s_{11}(k)>0\). The universal Bell/Jacobson identity is \[
\partial^{11}e^H/e^H=(H')^{11}+\partial^{11}H.
\] Iteration gives \(\partial^{11^h}e^H/e^H=\partial^{11^h}H\), since each constant \(H^{(11^j)}(0)\) is zero by evenness. Formula (K13.3) kills the logarithmic derivatives of orders at least \(q_i+2\) modulo 11; (K13.9) and oddness supply the two preceding orders. Hence \[
\partial_1^nF(0)\equiv\partial_3^nF(0)\equiv0\ (n\geq11),
\qquad\partial_2^nF(0)\equiv0\ (n\geq121).
\tag{K13.10}
\] This establishes a first congruence only. Repeated valuation gain in (K13.1) comes instead from the stable-ring estimate (K13.4) and the Bell lemma.

\subsection{4. AM-K014: an explicit determinant formula for the actual source sigma}\label{am-k014-an-explicit-determinant-formula-for-the-actual-source-sigma}

The standard theta divisor based at infinity is \(W_2\), hence its line is \(L_{2\infty}\). Buchstaber et al.~(1997), Definition 1 and equation (2.2), printed pp.7--8, normalize its Baker sigma by initial homogeneous term \(v_1v_3-v_2^2\) in the canonical half-differential coordinates. The construction below fixes the whole series and its sign, including the non-leading terms.

Use the integral arithmetic parameter \(t=-x^3/y\) at infinity. There is a unique \(A\in1+t^{14}\mathbb Z[[t^{14}]]\) satisfying \[
A=1-t^{14}A^7,\qquad x=t^{-2}A^{-1},\qquad y=-t^{-7}A^{-3}.
\tag{K14.3}
\] Substitution verifies both the curve equation and \(t=-x^3/y\). Coefficients of \(A\) are recursively determined without division, and all inverse powers \(A^{-m}\) are integral. With \(T=t^{14}\), \[
[T^0]A^{-m}=1,\qquad
[T^k]A^{-m}=\frac m k(-1)^{k-1}{7k-m-1\choose k-1}\quad(k\geq1).
\tag{K14.4}
\] For \(m=0,k>0\) this is zero. Formula (K14.4) is Lagrange inversion applied to \(1-A=T A^7\); its expression is an integer because \(A^{-m}\in\mathbb Z[[T]]\). The division in this closed expression may be avoided entirely by the recursion (K14.3).

Let the increasing nongaps be \[
n_1,n_2,\ldots=0,2,4,6,7,8,9,\ldots .
\] For even \(n=2k\), take \(\phi_n=x^k\); for odd \(n=2k+7\), take \(\phi_n=-x^ky\). Then \[
t^2\phi_n=t^{2-n}A^{-\lfloor n/2\rfloor}.
\tag{K14.5}
\] Thus the basis of \(t^2H^0(C-\{\infty\},\mathcal O_C)\) has leading exponents \[
2,0,-2,-4,-5,-6,\ldots ,
\] corresponding to the partition \(\lambda=(3,2,1)\).

Here is a finite algorithm for every coefficient. For any partition \(\mu\), choose \(L\geq\max(3,\operatorname{length}\mu)\), append zeros to \(\mu\), and define \[
\xi_\mu=
\det\left([t^{\,\mu_i-i-(2-n_j)}]\,
A(t)^{-\lfloor n_j/2\rfloor}\right)_{1\leq i,j\leq L}.
\tag{K14.6}
\] A negative power or a power not divisible by 14 has coefficient zero. Enlarging \(L\) adds an upper triangular identity tail and leaves the determinant unchanged. Each \(\xi_\mu\) is therefore an explicitly computable integer. It vanishes unless \(\mu\geq(3,2,1)\) and \(|\mu|=6+14k\) for some \(k\geq0\); \(\xi_{321}=1\).

Define polynomials \(p_m\) by \[
\sum_{m\geq0}p_m(v)z^m
=\exp(v_3z+v_2z^3+v_1z^5),\quad p_m=0\ (m<0),
\] and put \[
S_\mu(v)=\det(p_{\mu_i-i+j}(v))_{1\leq i,j\leq L}.
\] Then the actual Baker sigma is \[
\boxed{\sigma(v_1,v_2,v_3)=\sum_\mu\xi_\mu S_\mu(v).}
\tag{K14.7}
\] This is Nakayashiki's tau/sigma theorem, specialized explicitly: Nakayashiki (2010), §§4--5, equation (7), and Theorem 8; for the arithmetic parameter version, Ônishi (2015), Theorem 6.5. The affine-ring embedding is shifted by \(g-1=2\), exactly as in (K14.5), not by an unspecified convention. The monomial frame is upper unitriangularly equivalent to the normalized frame and has identical Plücker coordinates. Formula (K14.6) is its finite determinant.

We verify why there is no omitted exponential or coordinate correction in this specialization. The half-differentials \(x^{i-1}dx/(2y)\) have leading terms \(t^{6-2i}dt\); all later terms have exponents shifted by multiples of 14. Hence the gap-coordinate submatrix \(B\), at gaps \(1,3,5\), is the identity after the displayed reversal. The linear correction coefficients \(c_w\) vanish at these gaps because the square-root logarithm of the normalized highest differential contains powers only \(t^{14k}\). Likewise the regular coefficients \(q_{ij}\) of \(\Omega_B-dt_1dt_2/(t_1-t_2)^2\) vanish unless \(14\mid i+j\). Indeed the combination of \(a\) and the hyperelliptic involution acts on \(t\) by the primitive fourteenth root \(-\zeta^3\) and preserves \(\Omega_B\). For \(i,j\in\{1,3,5\}\), \(i+j\leq10\), so every retained \(q_{ij}\) vanishes. These are exactly the exponential corrections in the tau/sigma theorem.

Ônishi and Nakayashiki sometimes use \(-x^{i-1}dx/f_y\) rather than \(+x^{i-1}dx/(2y)\). Applying their convention to the identical model with \(Y=-y\) makes these differentials exactly ours. Equivalently a simultaneous sign of all three coordinates has no effect because the genus-three sigma here is even. The leading Schur polynomial fixes the remaining overall sign.

Every desired monomial has one finite weight \(5\alpha_1+3\alpha_2+\alpha_3\), so only finitely many partitions contribute to its coefficient in (K14.7). This is a checkable all-order construction of this particular function.

Each Schur polynomial in exponential times is Hurwitz integral: equivalently its coefficient of \(\prod_jT_j^{m_j}/m_j!\) is the integral symmetric-group character at cycle type \(1^{m_1}2^{m_2}\cdots\). One may instead use Ayano (2024), Lemma 7.2, or Ônishi (2015), Theorem 1.8. Since the \(\xi_\mu\) are integers, (K14.7) is Hurwitz integral over \(\mathbb Z\). Changing to \(u=2v\) and multiplying by 4 proves \[
\partial^\alpha\sigma_q(0)\in2^{2-|\alpha|}\mathbb Z.
\tag{K14.8}
\] This includes 7-adic integrality in the original, nonprimitive at 7, \(\omega\)-coordinates. It does not contradict ALP-LIFT-08's finite scalar content in primitive coordinates.

By (K12.5), the Baker gauge is the actual Chern-reduced gauge for this characteristic. The leading quadratic is \(q\); the principal theta space is one-dimensional. Thus (K14.1) agrees with the actual \(q\)-normalized section of ALP-LIFT-08, not merely with an abstractly isomorphic line.

\subsubsection{Leading coefficients and selection}\label{leading-coefficients-and-selection}

The weight-six term of (K14.7) is \[
S_{321}(v)=v_1v_3-v_2^2-\tfrac13v_2v_3^3+\tfrac1{45}v_3^6.
\tag{K14.9}
\] All subsequent weights are \(6+14k\). At weight 20, the coefficient of \(v_1^4\) is \(1/3\). This follows directly from (K14.6), or from the first-order frame shifts \(r\mapsto r+14\). In the latter calculation the only nonzero contributions after setting \(v_2=v_3=0\) are \[
(-1)(-1/6)+3(-1/2)+(-6)(-1/2)+8(-1/6)=1/3.
\] These terms come respectively from the partitions \((15,4,1)\), \((10,4,3,2,1)\), \((5,4,3,2,1^6)\), and \((3,2,2,2,1^{11})\). The script \nolinkurl{jacobian_sigma_determinant_certificate.py} checks these values without using the shift shortcut. It evaluates (K14.6) for all 627 partitions of 20, checks two extra identity-tail sizes for every nonzero coordinate, and computes the Schur determinants exactly. The retained JSON contains all 13 nonzero weight-20 Plücker coordinates and the four nonzero contributions above. It also checks (K14.9) symbolically and all four jets in (K14.2).

Consequently the beginning contains \[
\sigma_q(u)
=u_1u_3-u_2^2-\frac1{12}u_2u_3^3
+\frac1{720}u_3^6+\frac1{12}u_1^4
+\text{other terms of weights at least 20}.
\tag{K14.10}
\] In this notation \(u_1,u_2,u_3\) have weights \(5,3,1\); the displayed \(u_1^4\) is the indicated weight-20 term, and the last phrase excludes it. This proves (K14.2).

Every monomial is even in total degree, and its CM weight \(\alpha_1+2\alpha_2+3\alpha_3\) is \(4\bmod7\), since \[
\alpha_1+2\alpha_2+3\alpha_3
\equiv3(5\alpha_1+3\alpha_2+\alpha_3)\equiv18\equiv4\pmod7.
\] Thus the necessary pure-direction congruences are \[
n\equiv4\pmod{14}\ (i=1),\qquad
n\equiv2\pmod{14}\ (i=2),\qquad
n\equiv6\pmod{14}\ (i=3).
\tag{K14.11}
\] These are selection rules, not assertions that every permitted coefficient is nonzero.

\subsection{5. AM-K015: transfer to the actual source section}\label{am-k015-transfer-to-the-actual-source-section}

Fix \(p\nmid14\). On the good abelian scheme, the rigidified line difference is \[
L_D^{\rm rig}=L_\infty^{\rm rig}\otimes\mathcal P_T,\qquad 2T=0.
\] The canonical rigidified Poincaré isomorphism gives \([2]^*\mathcal P_T\simeq\mathcal P_{2T}\simeq\mathcal O\), integrally. Hence the two pulled-back theta sections can be divided to give the actual rational function \(r\) of (K15.1). The denominator has unit origin value by ALP-LIFT-07. The numerator is integral in the \(q\)-normalization by ALP-LIFT-08. Therefore \(r\) is regular and integral on an affine neighborhood \(\mathcal U\) of the origin.

The identity (K15.1) uses exactly the Chern-reduced gauges. The metrics of the two lines have the same translation-invariant curvature; their quotient is the flat metric on \(\mathcal P_T\). The normalized rigidified trivialization of \([2]^*\mathcal P_T\) is parallel and has value one at the origin. Thus it introduces neither a linear nor a quadratic exponential. Equivalently, the ratio of the two analytic reduced multipliers at \(2u\) is periodic: the difference semicharacter has order two. This confirms that the analytic ratio is the same algebraic \(r\), with no unknown scalar.

The invariant derivatives preserve \(\mathcal A=\Gamma(\mathcal U,\mathcal O)\): regular derivations preserve a localization of a regular function algebra. The CM-label proof (K13.2) applies on this actual abelian scheme, so \[
D_i^{q_i}\mathcal A\subset p\mathcal A,\qquad
v_{\mathfrak p}(D_i^kr(0))\geq\lfloor k/q_i\rfloor.
\tag{K15.2}
\] This is a separate stable-ring application. It does not assert that a Chern frame itself is an integral formal frame.

Suppose a bound already proved for \(F\) is \[
v_{\mathfrak p}(\partial_i^nF(0))
\geq c n-s_p(n)/(p-1).
\] Both constants used in §3 satisfy \[
cq_i\leq1,\qquad c\leq 1/(p^{h_i-1}(p-1)).
\tag{K15.3}
\] Differentiate (K15.1) \(n\) times along \(u_i\). Since \(2\) is a \(p\)-adic unit, the \(k\)-th Leibniz term has valuation at least \[
\begin{split}
&v_p{n\choose k}+\lfloor k/q_i\rfloor
+c(n-k)-s_p(n-k)/(p-1)\\
&=cn-s_p(n)/(p-1)
+\lfloor k/q_i\rfloor+s_p(k)/(p-1)-ck.
\end{split}
\tag{K15.4}
\] Write \(k=aq_i+b\), \(0\leq b<q_i=p^{h_i}\). Then \[
s_p(k)=s_p(a)+s_p(b),\qquad
b\leq p^{h_i-1}s_p(b).
\] Using (K15.3), \[
\lfloor k/q_i\rfloor+s_p(k)/(p-1)
\geq a+b/(p^{h_i-1}(p-1))\geq c(aq_i+b)=ck.
\] Thus every term satisfies the claimed lower bound. The sum does too. All coefficients of \(\sigma_q\) are integral at \(p\) by (K14.8), giving the maximum with zero.

This proves (K13.1) for the actual source-line germ, and also transfers the generic bound (K13.6) at every prime in the table. In particular, it gives explicit positive linear lower slopes with a logarithmic digit-sum loss at every \(p\equiv3,5\bmod7\).

The argument respects the source's odd level condition. The auxiliary map \([2]\) is used only to compare two already constructed sections and their derivatives. The section whose bounds are concluded is \(s_q\in H^0(J,L_\infty)\). A later odd-\(M\) pullback of this section has the exact normalization \[
M^{-n}\partial_i^n\sigma_q(Mu)|_{u=0}=\partial_i^n\sigma_q(0).
\] No translated \(D\)-line is substituted for \(L_\infty\).

The first congruence (K13.10) transfers too. Modulo 11, \(D_i^{q_i}\) is a derivation and kills the actual integral function \(r\). Applying it to (K15.1) and using the corresponding identity for \(F\) gives \[
\partial_1^n\sigma_q(0)\equiv\partial_3^n\sigma_q(0)\equiv0\ (n\geq11),
\qquad \partial_2^n\sigma_q(0)\equiv0\ (n\geq121).
\tag{K15.5}
\] The even order sets \(0,2,\ldots,10\) and \(0,2,\ldots,120\) found to consist of units for \(F\) are not asserted to be unit sets for \(\sigma_q\). Its different characteristic has its own selection rules.

\subsection{6. AM-K016: exact archimedean growth and controlled nonvanishing}\label{am-k016-exact-archimedean-growth-and-controlled-nonvanishing}

Fix \(i\), and put \[
f(t)=\sigma_q(te_i)=\sum_{n\geq0}a_nt^n,
\qquad \alpha=\frac{\pi}{2G_i}>0,\qquad
M(R)=\max_{|t|\leq R}|f(t)|.
\] The exact metric matrix is AM-R-004's diagonal \(G\). AM-R-006 gives the explicit bound \[
|f(t)|\leq C e^{\alpha|t|^2},\qquad C=6125/16.
\tag{K16.2}
\] The initial jets (K14.2) prove that \(f\) is nonzero for each \(i\). We now prove statements about this actual function, rather than all entire functions satisfying (K16.2).

\subsubsection{A lower growth bound from the compact theta metric}\label{a-lower-growth-bound-from-the-compact-theta-metric}

The Chern norm of the section is a continuous function on the compact torus \(J(\mathbb C)\). After a fixed constant normalization it pulls back along \(t\mapsto te_i\) as \[
g(t)=|f(t)|e^{-\alpha|t|^2}.
\tag{K16.3}
\] Restrict temporarily to real \(t\), and let \(K\) be the closure in \(J(\mathbb C)\) of this real one-parameter subgroup. It is a compact subgroup. Some real \(t_0\) has \(f(t_0)\ne0\), since otherwise the identity theorem would make \(f\) zero. Choose \(\varepsilon>0\) such that the relatively open set \[
V=\{x\in K:\|s_q(x)\|>\varepsilon\}
\] is nonempty, with the same fixed norm normalization as (K16.3).

The nonnegative real orbit is dense in \(K\). Here is an elementary justification. Its closure \(S\) is a compact subsemigroup of the torus. For any \(x\in S\), a pigeonhole argument applied to integer multiples of \(x\) in a successively finer finite cover produces positive multiples approaching zero. These return multiples may be chosen unbounded: if a bounded subsequence already returns arbitrarily close, one of its positive indices is an exact period and its multiples give unbounded exact returns. If the return indices are \(m_j\to\infty\), then \((m_j-1)x\to-x\), so \(-x\in S\). Thus \(S\) is a subgroup and must equal the closure \(K\) of the full real orbit. Positive translates of every \(x\in K\) therefore meet \(V\).

For each \(x\in K\), choose \(t_x\geq0\) with \(x+t_xe_i\in V\). The open sets \(V-t_xe_i\) cover \(K\); compactness supplies a finite subcover. Let \(L\) be the largest time used. For every \(r\geq0\), at least one \(t\in[r,r+L]\) has \(g(t)>\varepsilon\). Therefore \[
\varepsilon e^{\alpha(R-L)^2}\leq M(R)\leq Ce^{\alpha R^2}
\qquad(R\geq L).
\tag{K16.4}
\] In particular, \[
\log M(R)=\alpha R^2+O(R),
\tag{K16.5}
\] and \(f\) has order 2 and exact type \(\alpha\). It cannot be a polynomial, so it has infinitely many nonzero coefficients. The constants \(L,\varepsilon\) and the resulting thresholds are not made effective by this compactness argument.

\subsubsection{Exact coefficient limsup}\label{exact-coefficient-limsup}

Optimizing Cauchy's bound at radius \(\sqrt{n/(2\alpha)}\) gives \[
|a_n|\leq C\left(\frac{2e\alpha}{n}\right)^{n/2}\quad(n>0).
\tag{K16.6}
\] Consequently \(\limsup n|a_n|^{2/n}\leq2e\alpha\). If this inequality were strict, choose \(0<\beta<\alpha\) with \(|a_n|\leq(2e\beta/n)^{n/2}\) for all sufficiently large \(n\). With \(N=2\beta R^2\), a term at radius \(R\) is then at most \[
\exp\left(\frac N2\varphi(n/N)\right),
\qquad \varphi(x)=x(1-\log x),\qquad \varphi(1)=1.
\] The first \(\lceil e^2N\rceil\) terms sum to \(O(N+1)e^{N/2}\), and the remaining terms have a bounded geometric tail because the derivative of their logarithmic bound is at most \(-1\). The finitely many excluded terms form a fixed polynomial. Hence \(\log M(R)\leq\beta R^2+O(\log(R+2))\), contradicting (K16.5). We have proved \[
\limsup_{n\to\infty}n|a_n|^{2/n}=2e\alpha=\pi e/G_i.
\tag{K16.7}
\] Since \(\partial_i^n\sigma_q(0)=n!a_n\), Stirling's formula converts (K16.7) exactly to (K16.1). This is a limsup over all orders, with zero terms allowed. It does not assert that every nonzero derivative attains the upper bound.

\subsubsection{Every sufficiently large multiplicative window contains support}\label{every-sufficiently-large-multiplicative-window-contains-support}

Fix \(0<\delta<1\). We claim that every sufficiently large real \(N\) has a nonzero coefficient with index in \[
[(1-\delta)N,(1+\delta)N].
\tag{K16.8}
\] Suppose instead that the whole window is zero and put \(R=\sqrt{N/(2\alpha)}\). By (K16.6), each other term is bounded by \(C\exp(N\varphi(n/N)/2)\). Define \[
\kappa_\delta=1-\max\{\varphi(1-\delta),\varphi(1+\delta)\}>0.
\] There are at most \(N+1\) terms below the window, each bounded by \(C\exp((1-\kappa_\delta)N/2)\). Above the window the derivative of the logarithmic bound with respect to \(n\) is at most \(-\tfrac12\log(1+\delta)\), so those terms sum to a geometric tail with ratio at most \((1+\delta)^{-1/2}\). The constant term is zero for this source section; even if it were present it would satisfy the same estimate for large \(N\). Thus \[
M(R)\leq C\left(N+1+\frac1{1-(1+\delta)^{-1/2}}\right)
e^{(1-\kappa_\delta)N/2}.
\tag{K16.9}
\] But (K16.4) gives \(\log M(R)\geq N/2-O(\sqrt N)\), a contradiction for large \(N\). This proves (K16.8).

List the nonzero orders increasingly as \(n_1<n_2<\cdots\). If \(n_{k+1}/n_k\) stayed at least \(1+\eta\) along a subsequence, the interval between these two indices would contain a multiplicative window (K16.8) of a fixed positive relative width: take its midpoint and any \(\delta<\eta/(2+\eta)\). This contradicts (K16.8). Hence \[
\frac{n_{k+1}}{n_k}\longrightarrow1.
\tag{K16.10}
\] These orders can be enumerated exactly by (K14.6)--(K14.7): for a proposed pure order only one finite weight and finitely many partitions occur, and all coefficients are rational. The enumeration is therefore an actual algorithm with provable infinitely many outputs. Neither bounded additive gaps nor a nonzero coefficient at every allowed order is proved.

\subsection{7. What the new bounds do and do not supply to the arithmetic gate}\label{what-the-new-bounds-do-and-do-not-supply-to-the-arithmetic-gate}

The actual all-order finite lower bound at a prime in the table is \[
v_{\mathfrak p}(\partial_i^n\sigma_q(0))
\geq\max\{0,c_i(p)n-s_p(n)/(p-1)\},
\quad c_i(p)=\frac1{(p-1)(p^{h_i}+1)}.
\tag{K17.1}
\] At 11 the three finite checks replace these constants by the larger constants in (K13.1). For the two inert orbit classes specifically requested, the generic constants are

\begin{longtable}[]{@{}
  >{\raggedright\arraybackslash}p{(\columnwidth - 6\tabcolsep) * \real{0.2500}}
  >{\raggedright\arraybackslash}p{(\columnwidth - 6\tabcolsep) * \real{0.2500}}
  >{\raggedright\arraybackslash}p{(\columnwidth - 6\tabcolsep) * \real{0.2500}}
  >{\raggedright\arraybackslash}p{(\columnwidth - 6\tabcolsep) * \real{0.2500}}@{}}
\toprule\noalign{}
\begin{minipage}[b]{\linewidth}\raggedright
Prime class
\end{minipage} & \begin{minipage}[b]{\linewidth}\raggedright
\(c_1(p)\)
\end{minipage} & \begin{minipage}[b]{\linewidth}\raggedright
\(c_2(p)\)
\end{minipage} & \begin{minipage}[b]{\linewidth}\raggedright
\(c_3(p)\)
\end{minipage} \\
\midrule\noalign{}
\endhead
\bottomrule\noalign{}
\endlastfoot
\(p\equiv3\pmod7\) & \(1/((p-1)(p^3+1))\) & \(1/((p-1)(p+1))\) & \(1/((p-1)(p^2+1))\) \\
\(p\equiv5\pmod7\) & \(1/((p-1)(p+1))\) & \(1/((p-1)(p^3+1))\) & \(1/((p-1)(p^2+1))\) \\
\end{longtable}

These are uniform actual-section bounds. They are not claimed sharp. In particular, the accepted DP45/DP53 bare integral-function bounds cannot simply be assigned to this Chern-reduced germ. In the fixed global direction order, their comparison at 11 is

\begin{longtable}[]{@{}
  >{\raggedright\arraybackslash}p{(\columnwidth - 6\tabcolsep) * \real{0.2000}}
  >{\raggedleft\arraybackslash}p{(\columnwidth - 6\tabcolsep) * \real{0.2667}}
  >{\raggedleft\arraybackslash}p{(\columnwidth - 6\tabcolsep) * \real{0.2667}}
  >{\raggedleft\arraybackslash}p{(\columnwidth - 6\tabcolsep) * \real{0.2667}}@{}}
\toprule\noalign{}
\begin{minipage}[b]{\linewidth}\raggedright
Direction
\end{minipage} & \begin{minipage}[b]{\linewidth}\raggedleft
Accepted bare sharp slope from DP53
\end{minipage} & \begin{minipage}[b]{\linewidth}\raggedleft
Actual theta lower slope proved here
\end{minipage} & \begin{minipage}[b]{\linewidth}\raggedleft
Difference of the two stated constants
\end{minipage} \\
\midrule\noalign{}
\endhead
\bottomrule\noalign{}
\endlastfoot
\(\omega_1\) & \(121/1330\) & \(1/20\) & \(109/2660\) \\
\(\omega_2\) & \(11/1330\) & \(1/220\) & \(109/29260\) \\
\(\omega_3\) & \(132/1330\) & \(1/20\) & \(131/2660\) \\
\end{longtable}

The supplied exact locators are \nolinkurl{RETURN_DP_MULT_checkpoint_06_20261121T0420Z/research/DP45_integrated_fragment.tex}, the DP45 theorem, and \nolinkurl{DP53_integrated_fragment.tex}, displayed global label table and its proof paragraph beginning ``For a=3''. Their accepted \(p\equiv3,5\) triples are respectively \((p,p^3,p^2)/((p-1)(p^3+1))\) and \((p^3,p,p^2)/((p-1)(p^3+1))\). The numerical differences in the table are not proved losses of the actual theta section: its optimal finite slopes are still unknown.

There is also a uniform, explicit but much weaker upper bound on each \textbf{nonzero} actual pure jet. For \(n\geq2\), (K14.8) makes \(I_n=2^{n-2}\partial_i^n\sigma_q(0)\) a nonzero rational integer. For every odd prime \(p\), \(p^{v_p(I_n)}\leq|I_n|\). Using (K16.6) gives \[
v_p(\partial_i^n\sigma_q(0))
\leq\frac{\log(6125/16)+(n-2)\log2+\log(n!)
+\frac n2\log\!\left(\frac{\pi e}{nG_i}\right)}{\log p}.
\tag{K17.2}
\] Thus the finitely computable nonzero pure derivatives have both the actual linear lower bounds and an actual \(O(n\log n)\) upper bound. Formula (K17.2) applies at every prime above \(p\) because these derivatives are rational. It supplies no matching linear upper bound, optimal finite slope, or finite-prime analogue of (K16.1).

All mixed actual derivatives are determined by the Sato determinant and are integral away from 2. The positive linear finite bounds proved here concern the three pure directions. AM-K016 supplies nonzero pure derivatives with support ratios tending to one, and an exact archimedean limsup, but it does not ensure nonvanishing of a prescribed mixed harmonic row or charge. No transfer from one pure sequence to a source-selected harmonic image is claimed. Matching linear finite upper bounds, effective support thresholds, and the prescribed-row image comparison remain additional gates.

The full finite data for the translated \(F\) have useful but limited scope: modulo 11 its even derivatives through order 10 in directions 1 and 3, and through order 120 in direction 2, are units at both primes above 11. The exact certificate checks this by polynomial gcd with \(\Phi_7\); it does not convert those finite lists into an asymptotic finite upper bound for \(F\), much less for \(\sigma_q\).

\subsection{8. Primary sources, exact locators, and reproducibility}\label{primary-sources-exact-locators-and-reproducibility}

The literature supplies the general bidifferential, Kleinian, and tau/sigma identities cited here. The CM uniqueness argument, stable-ring estimates, Bell--Legendre constants, algebraic comparison, and growth/support arguments are proved in this note. No claim that integral Hodge eigenvectors give an integral Chern frame is used; the supplied THETA3 C11 and C19 counterexamples and actual-theta method remain consistent with every assertion above.

\begin{enumerate}
\def\labelenumi{\arabic{enumi}.}
\tightlist
\item
  Buchstaber, V. M., Enolskii, V. Z., \& Leykin, D. V. (1997). Hyperelliptic Kleinian functions and applications. \emph{American Mathematical Society Translations, Series 2, 179}, 1--33. \href{https://homepage.mi-ras.ru/~buchstab/download/advmathsci.2.179.1.pdf}{Author-hosted primary PDF}. Used: §1.2, equations (1.7)--(1.9); §2.1, Definition 1 and equations (2.1)--(2.2), pp.7--8; Corollary 2.1.2, equation (2.9), p.10; Theorem 2.2, equations (2.10)--(2.14), p.11. Local retrieved source: \nolinkurl{theta_sources/BEL1997.pdf} and \nolinkurl{.txt}.
\item
  Edelstein, J. D., Gómez-Reino, M., \& Mariño, M. (2001). \emph{Blowup formulae in Donaldson--Witten theory and integrable hierarchies} (arXiv:hep-th/0006113v3). \href{https://arxiv.org/pdf/hep-th/0006113}{Primary PDF}. Used: equations (3.3)--(3.6) for the explicit polar and bidifferential; (3.14)--(3.18) for the quadratic period normalization; (3.24), p.12, for the branch-partition Hessian. Local source: \nolinkurl{theta_sources/Edelstein2001.pdf} and \nolinkurl{.txt}.
\item
  Shramchenko, V. (2004). New Frobenius structures on Hurwitz spaces in terms of Schiffer and Bergmann kernels. \emph{Proceedings of the Institute of Mathematics of the NAS of Ukraine, 50}(3), 1466--1473. \href{https://www.slac.stanford.edu/econf/C0306234/papers/shramchenko.pdf}{Primary proceedings PDF}. Used: §4, p.1468, Schiffer formula and its intrinsic period characterization.
\item
  Nakayashiki, A. (2010). Sigma function as a tau function. \emph{International Mathematics Research Notices, 2010}(3), 373--394. \href{https://arxiv.org/pdf/0904.0846v2}{Author preprint, arXiv:0904.0846v2}. The cited v2 was revised in 2012; the journal article is from 2010. Used: §4, Plücker and Schur definitions and equation (7); §5, the shift by \(g-1\) and leading partition; §7, Theorem 4; §8, Theorem 8, normalized tau/sigma equality. All page references in this note concern the arXiv version.
\item
  Ônishi, Y. (2015). \emph{Hurwitz integrality of power series expansion of the sigma function for a plane curve} (arXiv:1510.03002v1). \href{https://arxiv.org/pdf/1510.03002}{Primary PDF}. Used: equation (2.6) for the differential sign convention; Theorem 1.8 for the general Hurwitz result; Theorem 6.5, printed p.21, for the arithmetic-parameter determinant and exponential corrections; Proposition 6.7 for weights. Local source: \nolinkurl{theta_sources/Onishi2015.pdf} and \nolinkurl{.txt}. The stronger integral statement in this specialization is proved directly from its integer frame and vanishing corrections, not imported from the general theorem without checking powers of 2.
\item
  Ayano, T. (2025). Hurwitz integrality of the power series expansion of the sigma function for a telescopic curve. \emph{Journal of the Mathematical Society of Japan, 77}(3), 833--867. \url{https://doi.org/10.2969/jmsj/90129012}. \href{https://arxiv.org/pdf/2208.10079v2}{Primary preprint, arXiv:2208.10079v2}, revised 14 April 2024. Used only as corroboration: Lemma 7.2 for Schur Hurwitz integrality and Theorem 7.5 for the tau/sigma relation. The reference in §4 to ``Ayano (2024)'' denotes this exact preprint version.
\end{enumerate}

Reproducible artifact paths, all relative to this report's directory:

\begin{longtable}[]{@{}
  >{\raggedright\arraybackslash}p{(\columnwidth - 2\tabcolsep) * \real{0.5000}}
  >{\raggedright\arraybackslash}p{(\columnwidth - 2\tabcolsep) * \real{0.5000}}@{}}
\toprule\noalign{}
\begin{minipage}[b]{\linewidth}\raggedright
Artifact
\end{minipage} & \begin{minipage}[b]{\linewidth}\raggedright
Actual check and scope
\end{minipage} \\
\midrule\noalign{}
\endhead
\bottomrule\noalign{}
\endlastfoot
\nolinkurl{jacobian_theta_direction_jets.py} / \nolinkurl{.json} & Exact algebraic ODE and Bell recurrence; rational pure \(F\) jets through order 8. \\
\nolinkurl{jacobian_theta_direction_jets_mod11.json} & Exact coefficient vectors modulo 11 through order 160 in all three directions. \\
\nolinkurl{jacobian_theta_11_certificate.py} / \nolinkurl{.json} & Regenerates the three exceptional log derivatives, checks saved coefficients and finite unit lists at both residue primes; PASS. \\
\nolinkurl{jacobian_theta_hessian_check.py} / \nolinkurl{.json} & Exact Bolza polynomial identity plus an independent finite theta-sum diagnostic for Hessians and initial \(F\) and \(\sigma_q\) jets; numerical comparisons are not proof premises. \\
\nolinkurl{jacobian_sigma_determinant_certificate.py} / \nolinkurl{.json} & Exhaustive exact weight-20 Plücker determinants, Schur polynomial, and four source-section jets; PASS. \\
\end{longtable}

The imported normalization statements are ALP-LIFT-07 in \nolinkurl{finite_theta_null_normalization.md}, ALP-LIFT-08 in \nolinkurl{trace_one_leading_jet.md}, AM-R-004 in \nolinkurl{../research/jacobian_period_theta.md}, and AM-R-006 in \nolinkurl{../research/trace_one_theta_scalar.md}. This report's own checks are recorded separately in \nolinkurl{kummer_extension_CHECK_EVENTS.tsv}; an author's whole-proof recheck is distinguished there from the parent's independent review.

\clearpage

\section{ALP\_MACHINE: definite comparison and oscillator contribution}\label{alp_machine-definite-comparison-and-oscillator-contribution}

Author: GPT-6 Astra Pro, agent \nolinkurl{/root/definite_oscillator}; effort inherited, not separately exposed. Run start 2026-11-21 19:35 UTC. These are work products for root integration, not a claim that all assigned tasks are finished.

\subsection{Source conventions and carried results}\label{source-conventions-and-carried-results}

Use the supplied run names and statement numbers. The current run carries THETA3 A7, B9, B13, B15, B17, B19, B24 and GCT\_PLETH T-052; it does not rerun their proofs. Exact local source paths follow below. The supplied source archive has reports and source manifests. In this run the needed Lean bodies were recovered centrally from the intake-pinned mirror, with hashes compared against the intake original-author archive metadata wherever that metadata contains the path. Later additional Rallis and TLiftK files have no author-hash entry in that comparison manifest and are identified as pinned-mirror reads in the events. The relevant definitions and proof bodies below were then read directly; no Lean execution or complete dependency audit is claimed.

\begin{enumerate}
\def\labelenumi{\arabic{enumi}.}
\tightlist
\item
  \nolinkurl{RETURN_THETA3_checkpoint_07_20261121T0604Z/documentation/paper_A7.tex}, equation (A7-lift), final two paragraphs: the multiplicity space P is totally definite exactly when the first-list complete CM type functions agree. In the balanced 2n-factor construction this forces all 2n type functions to agree, and the selected balanced tensor is divisor-generated.
\item
  \nolinkurl{documentation/paper_B.tex}, B9, equations (B9a) and following coefficient calculation: the canonical right coefficient on U(2,1) has compact type det(A)d\^{}\{-2\}, equivalently lowest compact weight (1,1;−2). For probability Haar and the compatible compact det test line, the LOCAL rank-zero Bessel integral is one. This is not a global period evaluation.
\item
  Same file, B13: actual MainX relation transfer comes from LawsX. Common geometry, tuple validity, finite support, common real action, double and last integral laws, kernel reduction, terminal factorization and row covariance are distinct hypotheses.
\item
  Same file, B15: literal one-active-place rank3 geometry has global Witt index zero, so abstract MainX length m=2; the remaining definite-side pattern is A2 → S → X2=T⊥Lz, with Lz of opposite sign to L. The small choice S2,T1 gives A2→S2→X2, equal ranks on the definite side. The relevant kernel identity is (B15m).
\item
  \nolinkurl{documentation/paper_B17.tex}, (B17a)--(B17g): the adjoint double lift and its positive inverse are exact analytic constructions. Its upper space is A⊥H of signature (3,1) at every real place, even if P is totally definite. It is therefore not an instantiation of the arithmetic definite upper tower.
\item
  \nolinkurl{documentation/paper_B19.tex}, (B19a)--(B19k) and final paragraph: actual CM-jet model and all period factors; the arithmetic functional is the fixed \nolinkurl{abelian₂}, not a freely chosen linear functional. The proved StdPair kernel constructor requires both ranks even and strictly increasing.
\item
  \nolinkurl{documentation/paper_B24.tex}, (B24a)--(B24e): in the repeated-type definite sector the actual integral Gram inverse is one specified positive scalar times an algebraic matrix. This is stronger than mere analytic existence but does not identify the MainX kernel or terminal masses.
\item
  \nolinkurl{RETURN_GCT_PLETH_checkpoint_09_20261121T0546Z/work/theta/oscillator_charge_report.md}, T-052, (OC.1)--(OC.12), (OC.20)--(OC.22): r=m−2, pair U(r,1)×U(m), polynomial Fock space C{[}X\_ai,Y\_i{]}, and all vacuum shifts.
\end{enumerate}

Primary literature for the last convention: Basile, T., Joung, E., Mkrtchyan, K., \& Mojaza, M. (2020), \emph{Dual pair correspondence in physics: Oscillator realizations and representations}, arXiv:2006.07102v2, §5.2, equations (5.21)--(5.30), printed pp.29--31. This source was opened anew at \url{https://arxiv.org/html/2006.07102v2} and \url{https://arxiv.org/pdf/2006.07102}. The new obstruction below is a calculation from the carried T-052 operators, not a claim of a theorem appearing in Basile et al.

\subsection{AM-DO-001. Explicit repeated-type definite geometry over Q(ζ7)}\label{am-do-001.-explicit-repeated-type-definite-geometry-over-qux3b67}

Let ζ=exp(2πi/7), E=Q(ζ), F=Q(ζ+ζ\^{}\{-1\}), and a=ζ+ζ\^{}\{-1\}. Index the three real embeddings by a↦2cos(2jπ/7), j=1,2,3. The signs of a are (+,−,−). In fact its minimal polynomial is t³+t²−2t−1, so it is an algebraic integer and a unit. Put δ=ζ−ζ\^{}\{-1\}. Its imaginary parts at the complex extensions σ1,σ2,σ3 are all positive.

Carry the CM type Φ=\{1,2,3\} of J=Jac(y²=x⁷−1) from DP\_MULT DP45/DP49. Let v be its odd type function on (Z/7Z)×. The inverse type d(σ)=v(σ\^{}\{-1\}) has signs (+,−,−) on σ1,σ2,σ3, since 1\^{}\{-1\}=1, 2\^{}\{-1\}=4, 3\^{}\{-1\}=5 modulo7. Therefore the rank-one Hermitian form h\_P=−a has exactly the required signs −d(σ)sign Imσ(δ) from THETA3 A6. Take both first-list types equal to v and

\begin{verbatim}
V=(E³,diag(1,1,−a)),  L=Ee3,  A=L⊥=(E²,I2),
P=(E²,δ(−a)I2) as a skew-Hermitian space.
\end{verbatim}

Then V has signature (2,1) at σ1 and (3,0) at the other real places; A is totally positive definite and U(P) is compact at all real places. Its Hermitian sign model is (0,2) at σ1 and (2,0) elsewhere. V and P are globally anisotropic because each has a definite real completion. These are actual global forms over E/F, not separately prescribed incoherent local signatures.

For any nonzero ξ0∈E with ξ0\^{}c=−ξ0 put q=N\_\{E/F\}(ξ0)\textgreater0 at every real place and Lz=−q\^{}\{-1\}L. Its Gram entry is a/q. Taking T=Lz gives X2=T⊥Lz with Gram (a/q)I2: positive at σ1 and negative at the other real places, hence totally definite. For the explicit ξ0=√−7/7, q=1/7, its Gram is 7a I2. Here √−7 can be taken as ζ+ζ²+ζ⁴−ζ³−ζ⁵−ζ⁶; its square is−7 by reduction modulo Φ7. The exact field identities and inverse-type signs are also recorded by \nolinkurl{definite_oscillator_geometry.py} and its JSON receipt. The identity line map satisfies G\_L=−q G\_Lz exactly, as required by the abstract last seesaw relation.

Proof: all signatures follow from the displayed cosines and the inverse type calculation; the line relation is −a=−q(a/q). The compactness criterion and the interpretation as the repeated-type A7 multiplicity space are carried A7 inputs. The polynomial identity for a follows by ζ⁷=1 and ζ≠1. No global theta nonvanishing or arithmetic MainX law is claimed merely from this geometry.

Scope: A7 supplies nonzero canonical wedges in this repeated-type family after suitable inputs and level. If a balanced last pair is requested, it must have the same type v, so this family addresses divisor-generated lines. The alternatives (adjoint T=V versus arithmetic T=Lz) remain different. The former has indefinite upper X4 and an actual analytic inverse; the latter has definite X2 but still requires a global reverse lift, actual finite classes, kernel equations, and terminal masses. This construction does not give their numerical values.

\subsection{AM-DO-002. A canonical compact-weight obstruction for the charged oscillator}\label{am-do-002.-a-canonical-compact-weight-obstruction-for-the-charged-oscillator}

\subsubsection{Precise statement}\label{precise-statement}

Let m≥3, r=m−2, and let ω\_m be the polynomial oscillator of T-052, with the fixed positive-energy convention and dual pair U(r,1)×U(m). On G=U(r,1) let K=U(r)×U(1). Let π\_can be the irreducible canonical holomorphic discrete-series Harish--Chandra module with lowest K type

\begin{verbatim}
κ_can(A,d)=det(A)d^(−r).
\end{verbatim}

For r=2 this is exactly THETA3 B9/B15's canonical ball-surface coefficient, with lowest weight (1,1;−2). There is no nonzero (g,K)-map from ω\_m or any of its algebraic subquotients to π\_can, nor from π\_can to any such subquotient. The same holds after arbitrary determinant character twists of either representation and after a simultaneous contragredient convention. In particular no charge-respecting intertwiner can identify a T-052 component with this canonical 032 component.

For odd m the bare oscillator may live on the metaplectic preimage of K. In that case every displayed map means a (g, K̃)-map on a common cover, with π\_can pulled back; equivalently the obstruction can first be read at the complexified Lie algebra and the traceless compact circle. Its generator H has the integral spectrum computed below, independently of the individual half-vacuum weights. This convention is already part of T-052 and does not assume an unprovided descent of the bare representation.

The statement concerns the canonical automorphic component, not arbitrary maps between entire oscillator representations after changing dual pairs, groups, coefficient systems, or taking additional tensor factors. When comparing precisely the ball surface U(2,1), the T-052 matching rank is m=4,r=2.

\subsubsection{Full proof}\label{full-proof-1}

T-052 (OC.7) gives

\begin{verbatim}
A_ab=Σ_i X_ai ∂_X_bi+(m/2)δ_ab,
A_00=−Σ_iY_i∂_Y_i−m/2.
\end{verbatim}

Let h=diag(1,\ldots,1,−r) in the complexified compact Lie algebra. Its imaginary multiple generates the compact circle c(z)=diag(zI\_r,z\^{}(−r))⊂SU(r,1). Its infinitesimal oscillator action is

\begin{verbatim}
H=Σ_a A_aa−r A_00=N_X+rN_Y+rm.
\end{verbatim}

All monomials in the polynomial Fock model have N\_X,N\_Y≥0. Thus the H spectrum lies in rm+Z\_\{≥0\}. More precisely it is rm+N\_X+rN\_Y. The action is diagonal on a monomial basis. Its restriction to every invariant subspace remains semisimple with this same lower bound; every algebraic quotient has only these eigenvalues as well, since it is the direct sum of its circle weight quotients.

The canonical character on c(z) is

\begin{verbatim}
det(zI_r)(z^(−r))^(−r)=z^{r(r+1)}.
\end{verbatim}

But rm=r(r+2)\textgreater r(r+1). Therefore the lowest K type of π\_can does not occur in ω\_m, or in any subquotient. A nonzero homomorphism from the irreducible π\_can is injective, contradicting absence of its lowest K type. A nonzero homomorphism to π\_can is surjective; exactness of taking an isotypic component for the compact circle (or simply its semisimple weight decomposition on K-finite vectors) would then force its lowest-circle-weight vector to have a preimage of that weight. Again impossible.

Every determinant character is trivial on c(z), since det c(z)=1. Changing the unitary splitting only adds such determinant characters to the fixed pair, so cannot change H. A unitary equivalent realization obtained by changing the polarization preserves the action and its spectrum. Simultaneously passing to contragredients negates both spectra; the distinguished weight −r(r+1) still lies outside the new spectrum, bounded above by −rm. The same absence argument applies. This completes the proof.

\subsubsection{Exact charge calculation}\label{exact-charge-calculation}

On the simultaneous charge-k subspace of T-052, with Q\_i=Σ\_aX\_ai∂\_X\_ai−Y\_i∂\_Y\_i=k, one has N\_X−N\_Y=mk. If b=N\_Y, then

\begin{verbatim}
H=m(k+r)+(r+1)b.
\end{verbatim}

The canonical equality would require m k+(r+1)b=−r, impossible for the T-052 parameters k,b≥0. For m=4,r=2 this reads 4k+3b+8=6. The contradiction is already present before charge selection.

One can also restrict the standard U(r,1) factor to a standard U(2,1) subblock for any m≥4. Its traceless compact circle generator has spectrum N\_X,selected+2N\_Y+2m≥2m≥8, whereas the canonical ball-surface coefficient has weight6. Thus this standard restriction does not evade the obstruction.

\subsubsection{Check scope and consequence}\label{check-scope-and-consequence}

This is a local archimedean obstruction. It survives global choices of finite Schwartz vectors, Haar measures, finite levels, and CM periods. It does not show that every conceivable comparison with Alpöge is impossible. It specifically excludes a route which first identifies T-052's actual charged module with the canonical 032 component while preserving the indicated group action. Q-ALPOGE and Q-ALPOGE-METHOD remain open for changed weights, other coefficient systems, auxiliary oscillator factors, or comparisons not passing through that component.

\subsection{AM-DO-003. Partial Fourier comparison with a definite dual pair, and its exact category obstruction}\label{am-do-003.-partial-fourier-comparison-with-a-definite-dual-pair-and-its-exact-category-obstruction}

\subsubsection{Algebraic statement}\label{algebraic-statement}

Over Z let A be the Weyl algebra on coordinates X\_ai,Y\_i and derivations ∂\_X\_ai,∂\_Y\_i, with {[}∂\_u,u{]}=1. Let A' be the Weyl algebra on X\_ai,Z\_i. The assignments

\begin{verbatim}
X_ai↦X_ai,  ∂_X_ai↦∂_X_ai,
Y_i↦∂_Z_i,  ∂_Y_i↦−Z_i
\end{verbatim}

define a Z-algebra isomorphism F:A→A'. (For the half-integer oscillator constants extend coefficients to Z{[}1/2{]}.) It carries the T-052 complex Lie algebra operators into the definite compact-pair oscillator operators for U(r+1)×U(m), conjugating the first matrix action by diag(I\_r,−1):

\begin{verbatim}
F(A_ab)=Σ X_ai∂_X_bi+(m/2)δ_ab,
F(A_00)=Σ Z_i∂_Z_i+m/2,
F(A_a0)=−Σ X_ai∂_Z_i,
F(A_0a)=−Σ Z_i∂_X_ai,
F(B_ij)=Σ X_ai∂_X_aj+Z_i∂_Z_j+((r+1)/2)δ_ij.
\end{verbatim}

The exact transformed charge is

\begin{verbatim}
F(Q_i)=Σ_a X_ai∂_X_ai+Z_i∂_Z_i+1.
\end{verbatim}

Thus the scalar shift1 is mandatory; charge does not become the unshifted compact column degree.

Integrally, ``Weyl algebra'' here means precisely the algebra generated by coordinate multiplication and the ordinary derivations. It is not the full Grothendieck ring of differential operators over Z. Accordingly the target below is an integral Weyl-algebra module; after tensoring with Q it is the usual algebraic delta D-module. No divided-power differential-operator action, stratification, or crystalline Frobenius is supplied. EGA IV4, Theorem16.11.2 p54 and Remarks16.11.3(ii) p55, gives the exact distinction: the integral differential-operator basis consists of divided derivatives, and ordinary derivations do not generate them in general. This primary text was read anew.

\subsubsection{Full proof of the algebra identities}\label{full-proof-of-the-algebra-identities}

The only nontrivial defining relation is {[}−Z\_i,∂\_Z\_j{]}=δ\_ij; all mixed generators commute as required. The inverse sends Z\_i↦−∂\_Y\_i and ∂\_Z\_i↦Y\_i. For A\_00 use ∂Z\_i·Z\_i=Z\_i∂Z\_i+1. For B\_ij use ∂Z\_j·Z\_i=Z\_i∂Z\_j+δ\_ij. The resulting constants are −m/2+m=m/2 and (r−1)/2+1=(r+1)/2. The off-diagonal signs are both negative and are exactly conjugation by diag(I\_r,−1). The charge formula follows from the same normal ordering. These calculations prove all displayed identities.

\subsubsection{The actual module and why it is not the definite Fock space}\label{the-actual-module-and-why-it-is-not-the-definite-fock-space}

Let Mδ be the A'-module with cyclic vector δ satisfying Z\_iδ=0 and ∂\_X\_aiδ=0. As a Z-module it has basis X\^{}u∂Z\^{}vδ. The map

\begin{verbatim}
U:Z[X,Y]→Mδ,  X^uY^v↦X^u∂Z^vδ
\end{verbatim}

is a bijection and satisfies U(a f)=F(a)U(f). Indeed, multiplication by Y becomes ∂Z, and −Z\_i∂Z\^{}vδ=v\_i∂Z\^{}\{v−e\_i\}δ reproduces ∂Y\_i. X and ∂X are unchanged. This proves exact equivariance and shows that U preserves the ORIGINAL charges, using their transformed shifted operators. In particular Z\_i∂Z\_i acts on ∂Z\^{}vδ by −(v\_i+1), so the transformed charge is deg\_X,column−v\_i as required.

Over C, Mδ is the algebraic D-module generated by the delta section on the subspace Z=0. It is not the ordinary holomorphic polynomial Fock module C{[}X,Z{]}. There is NO nonzero linear map L:C{[}X,Y{]}→C{[}X,Z{]} satisfying L(a f)=F(a)L(f) for every Weyl generator. For f=1, ∂Y\_i1=0 forces −Z\_iL(1)=0, and multiplication by Z\_i is injective on C{[}X,Z{]}, so L(1)=0. The source is cyclic under coordinate multiplications, so L=0.

It is also not a unitary oscillator intertwiner. In positive Fischer/Fock structures Y\_i†=∂Y\_i and (∂Z\_i)†=Z\_i, whereas F(∂Y\_i)=−Z\_i. Thus F fails the adjoint relation. It is a complex symplectic transformation of the algebraic Weyl algebra, not a real metaplectic equivalence of these positive-energy representations. No splitting character changes this Heisenberg adjoint obstruction.

\subsubsection{Consequence for the actual problem}\label{consequence-for-the-actual-problem}

There is an explicit algebraic comparison between the formulas underlying T-052 and a totally definite compact dual pair. Its landing object is a delta D-module and its charge has the displayed shift. It does not map the prescribed Schwartz/Fock theta inputs to the corresponding definite theta inputs, and therefore does not prove any of Alpöge's lift identities, quantitative nonvanishing, or the plethysm complexity claims. Conversely it prevents the overly strong statement that there is no comparison of any kind: there is a precise comparison after changing category.

\subsection{Current state of THETA3 G1}\label{current-state-of-theta3-g1}

The exact arithmetic equation, inherited B15m, is

\begin{verbatim}
volLev(X_j,U_{j,α}) θ_{K,j}(h_{j,α,a},h_{j+1,γ,a'};Φ_{j,P})
  = c_{K,j} abelian₂_{j,α,γ,a,a'}(mvZ(1,s_{j,τ}I)(castP P)).
\end{verbatim}

The functional \nolinkurl{abelian₂} is fixed by the specified arithmetic theta section, its character factors, and its period frame. It cannot be defined post hoc from the analytic left side. The terminal condition \nolinkurl{ClsT.Is.tail} includes top-level covolumes and reciprocal stabilizer masses. Its exact recovered statement has now been read directly in \nolinkurl{OnFramesRdApply.lean}, \nolinkurl{ClsT.Is.tail}:

\begin{verbatim}
volLev(T,U_T) Σ_{β,a} stabOrd(U_T,h_a)^(−1)
   ell_{β,a,G}(low(rd_m(ψ) evaluated on emb(h_a,1)))
  = cP·tail(ψ,G).
\end{verbatim}

The separate \nolinkurl{top} clause says that the translated terminal automorphic function at h\_a k is the functional ell\_\{β,a,G\} applied to the fixed archimedean coefficient J(v,k). Thus the terminal equation must match the actual top representation and its branching map; it is not merely the total mass of the finite class set.

Paths: \nolinkurl{Development/AbelianVarietyTheory/Definitions/LatticeFrames/OnFramesRdApply.lean}, \nolinkurl{ClsK.Core} lines108--113, class masses115--121, \nolinkurl{ClsT.Is.tail}147--150, \nolinkurl{ClsPinD}152--153; \nolinkurl{LevChiEqOne.lean}, \nolinkurl{abelian₂}73--77; \nolinkurl{SplitPart.lean}, \nolinkurl{StdPair}38--49; \nolinkurl{KernelPartLawsOfColumnPins.lean}, \nolinkurl{clsKL_is}22--52; \nolinkurl{ThetaKernelSplitPlacesEqJetClassicalTheta.lean}, \nolinkurl{thetaKernel_eq_jet_thC}46--71.

Approaches inspected:

\begin{itemize}
\tightlist
\item
  Repeated CM type supplies an actual totally definite P, and AM-DO-001 makes the forms explicit over the same cyclotomic field as the main Jacobian. It leaves the finite class matrices and global reverse lift unconstructed.
\item
  The exact adjoint double lift B17 has an indefinite upper space. Repeated types do not fix that signature; this is a carried source conclusion.
\item
  The B15 rank-one terminal alternative has the required upper signature and locally compatible archimedean representations, but neither finite local occurrence for the chosen global data nor the global reverse realization is supplied.
\item
  B24 makes the actual kernel inverse algebraic up to one retained scalar in this sector. Algebraicity/proportionality of an inverse is not equality with the fixed MainX functional, and cannot set terminal masses.
\item
  The StdPair constructor only covers strictly increasing even ranks. It cannot directly be applied to the natural rank2/3 analytic pair, nor to the small equal-rank definite side. This is a limitation of that constructor, not a proof that abstract MainX is impossible.
\item
  Arbitrarily rescaling a test vector to make one Bessel period equal1 does not solve the common kernel equations; that would lose the fixed input and uniformity.
\end{itemize}

Thus question2 is not settled by this contribution. Its current gate is simultaneous construction of the chosen global two-step theta representation and finite arithmetic class/jet matrices, with exact covolumes and terminal masses, for the concrete repeated-type data. No pairing value satisfying MainX has been computed here.

\subsection{AM-DO-004. A simultaneous sign obstruction to the common positive arithmetic pin}\label{am-do-004.-a-simultaneous-sign-obstruction-to-the-common-positive-arithmetic-pin}

\subsubsection{Statement with all hypotheses}\label{statement-with-all-hypotheses}

Let E/F be a CM extension and let X0,X1,X2 be nonzero Hermitian spaces, each definite at every real place, with signs ε\_j(τ)∈\{+1,−1\}. On the oriented matrix space W\_j=Hom\_E(X\_j,X\_\{j+1\}) use the exact source form

\begin{verbatim}
h_j(x,y)=tr(G_j^(−1) y* G_{j+1} x),
B_j(x,y)=Tr_{E/F}(ξ0 h_j(x,y)),  ξ0^c=−ξ0.
\end{verbatim}

A \textbf{label-preserving positive transport} means: identify each abstract matrix group with the corresponding positive definite one, replace all node Gram matrices by positive definite ones, use one common ξ0 and one common additive character ψ\_b(x)=ψ(bx) (b∈F×), and identify each oriented matrix oscillator space by an E-linear map, intertwining the fixed two group actions and carrying its Heisenberg form to the new form with that common additive character. No conjugation of the E action, exchange with a contragredient, or separate additive character for different edges is allowed in this definition.

A necessary condition for such a transport is

\begin{verbatim}
ε_j(τ) ε_{j+1}(τ)=sign τ(b)   for j=0,1,
\end{verbatim}

and hence ε0(τ)=ε2(τ) at every real place. More generally in a chain all even-indexed node signs must agree and all odd-indexed node signs must agree.

For the literal rank-zero MainX geometry of a one-active-place compact 032 ball surface, this necessary condition fails at every compact real place. Consequently its signed definite side cannot be transported to the common all-positive \nolinkurl{PinX}/\nolinkurl{StdPair} geometry by label-preserving positive transports, even after one common change of additive character. This does not refute the abstract signed \nolinkurl{MainX} construction.

\subsubsection{Exact source locator}\label{exact-source-locator}

\nolinkurl{Formalism/LocalToGlobal/IndefiniteReal/LocalNewformFunction.lean}: \nolinkurl{hO} lines25--26, \nolinkurl{symplO}28--29, \nolinkurl{heisO}34--35, \nolinkurl{actO}37--38, and \nolinkurl{IsWeilDatumO.om_heis}62--63. The actual formula uses the SOURCE Gram inverse; this is why a scalar edge ratio is q\_\{j+1\}/q\_j rather than its product.

\nolinkurl{Development/AbelianVarietyTheory/Definitions/LatticeFrames/SplitPart.lean}: \nolinkurl{StdPair.gS}, \nolinkurl{gT}, \nolinkurl{cc}, \nolinkurl{xi}38--49. \nolinkurl{LevChiEqOne.lean}: \nolinkurl{PinX.xs} and \nolinkurl{vc} specify Cartan involution1, orthonormal frame1 and standard lattice for every node, all with the same ψ. The abstract \nolinkurl{DefSp.Is.pm}, in \nolinkurl{JointSmooth.lean}, instead allows signs+1 or−1 separately at every real place. \nolinkurl{PlaceLaws.lean}, \nolinkurl{WeilOK.tau}165--166, imposes the opposite-line Gram equation.

\subsubsection{Full proof}\label{full-proof-2}

Fix a real place τ. Write G\_j=ε\_j P\_j with P\_j positive definite. For x≠0,

\begin{verbatim}
h_j(x,x)=ε_j ε_{j+1} tr(P_j^(−1) x* P_{j+1} x).
\end{verbatim}

The trace on the right is positive: it is the sum of squared absolute values of the matrix P\_\{j+1\}\^{}\{1/2\} x P\_j\^{}\{−1/2\}. Thus h\_j is definite of sign ε\_j ε\_\{j+1\} on the complex vector space W\_j.

Let h\_j\^{}+ be the Hermitian form for the transported positive pair and T\_j the required complex-linear map. Equality of the real symplectic forms with additive-character scaling says

\begin{verbatim}
Tr_{C/R}(ξ0 h_j(x,y))
  = b Tr_{C/R}(ξ0 h_j^+(T_jx,T_jy)).
\end{verbatim}

Because T\_j is complex linear, replacing x by ix also gives equality of the complementary real part. Since ξ0 is nonzero purely imaginary, the two real equalities imply

\begin{verbatim}
h_j(x,y)=b h_j^+(T_jx,T_jy).
\end{verbatim}

For x=y≠0 the right side has sign sign b. Therefore sign b=ε\_j ε\_\{j+1\}. Multiplying the two edge equalities gives ε0ε2=(sign b)\^{}2=1, proving the necessary condition. Repeating this argument along the chain proves the parity assertion. The same proof covers arbitrary positive Gram matrices; it is not limited to scalar forms.

For the ball surface, write V=A⊥L and X2=T⊥Lz in the literal m=2 last-step pattern. At a compact real place τ, V is definite, so A and L have the same sign. The exact line relation

\begin{verbatim}
τ_line* G_L τ_line=−N_{E/F}(ξ0) G_Lz
\end{verbatim}

has a totally positive norm factor, and therefore Lz has the opposite sign to L. Since X2 is definite and contains Lz, it has Lz's sign. Hence X0=A and X2 have opposite signs at τ. This contradicts the necessary condition. Such a compact real place exists because the 032 construction under discussion has {[}F:Q{]}\textgreater1 and only one active place.

This argument does not say that U(qI) and U(I) are different algebraic groups; they are literally the same matrix group. It says that the E-linear symplectic embeddings and their common Heisenberg character cannot simultaneously be the positive ones. Splitting-character twists multiply group actions by determinant characters but do not change the Heisenberg commutator form, so they cannot remove this obstruction. A common change ξ0↦cξ0 is already absorbed into b. Changing the middle sign multiplies both edge signs by−1 and still leaves them opposite. Rescaling the original definite completion of V changes both A and L together and again leaves the endpoint sign opposition unchanged. This completes the proof.

\subsubsection{The scalar equations after fixing a Hermitian middle-space conversion}\label{the-scalar-equations-after-fixing-a-hermitian-middle-space-conversion}

The endpoints of AM-DO-001 have q0=1 and q2=a/q, where q=Nξ0\textgreater0. Its middle object P is initially skew-Hermitian, with Gram δ(−a)I2. A comparison with the oriented Hermitian \nolinkurl{hO} model must first specify the skew-to-Hermitian conversion, the source-dual convention and the group labels. In particular one must not silently declare the middle Hermitian Gram to be −aI2 while independently choosing ξ0. Suppose that actual conversion gives a scalar Hermitian Gram q1=s∈F×. The edge ratios are then

\begin{verbatim}
r0=s,   r1=a/(qs).
\end{verbatim}

An E-linear scalar matrix map x↦c\_jx into the standard positive model would have to satisfy

\begin{verbatim}
b N_{E/F}(c_0)=s,
b N_{E/F}(c_1)=a/(qs).
\end{verbatim}

Their product gives b²N(c\_0c\_1)=a/q, independently of s. The left side is totally positive, while a/q has two negative real embeddings. Thus there is no solution. If a compatible convention separately proves s=−a, the displayed equations specialize to bN(c0)=−a and bN(c1)=−1/q; that specialization is not an additional claim here. The obstruction itself uses only the endpoints and the necessary real sign of a norm, so neither the unresolved middle conversion nor possible finite norm obstructions affect this exclusion.

\subsubsection{Meaning for the open questions}\label{meaning-for-the-open-questions}

This refutes one concrete attempt to solve THETA3 G1: transport both edges of its literal two-step definite side into the already-built positive-identity arithmetic kernel constructor while keeping the E labels, charges and a common additive character. It leaves the signed arithmetic kernel/covolume equations themselves open. They would require a new signed CM-type comparison (or a different global realization), with explicit conjugations, characters and periods. The theorem does not obstruct Q1's already-positive abstract definite tower; it is specific to importing the 032 one-active-place geometry into that same positive pin.

\subsubsection{Extension allowing whole-edge E/F conjugation}\label{extension-allowing-whole-edge-ef-conjugation}

The same route still fails if each edge identification is allowed to be either E-linear or E-conjugate-linear globally over F, with the corresponding conjugation of its group labels. Write η\_j=+1 in the former case and η\_j=−1 in the latter. Each η\_j is independent of the real place: the only nontrivial E/F automorphism is complex conjugation at all of them. The necessary condition becomes

\begin{verbatim}
sign τ(b)=η_j ε_j(τ) ε_{j+1}(τ),
ε0(τ) ε2(τ)=η0η1.
\end{verbatim}

To prove the first equality, the E-linear case is already proved. In the conjugate-linear case, put T(ix)=−iT(x) in the two symplectic equalities. They give h\_j(x,y)=−b overline(h\_j\^{}+(T x,T y)), so the sign is −sign b. Multiplying the two edge conditions proves the second equality.

For the one-active-place ball geometry, ε0ε2 is +1 at the active place because A and Lz are both positive, and −1 at every compact place because Lz has the opposite sign to A there. It is therefore not constant, contrary to η0η1. Thus whole-edge global conjugation does not repair the obstruction either. Independent conjugations at selected real places are not E/F-semilinear maps of the rational oscillator space; no claim about constructions that replace that global space and its rational theta lattice is made here.

\subsection{Additional source and proof checks}\label{additional-source-and-proof-checks}

After recovering the actual Lean files, the agent read all of \nolinkurl{OnFramesRdApply.lean} through the ClsK/ClsT definitions and conversion lemmas; all of \nolinkurl{ValLastX.lean}; \nolinkurl{StdPair} and \nolinkurl{xiN} in \nolinkurl{SplitPart.lean}; \nolinkurl{TwDat}, \nolinkurl{abelian₂}, \nolinkurl{PinX}, \nolinkurl{PinColT} in \nolinkurl{LevChiEqOne.lean}; \nolinkurl{DefSp}, \nolinkurl{KerDat}, \nolinkurl{BesAt}, \nolinkurl{JacAt} in \nolinkurl{JointSmooth.lean}; \nolinkurl{WeilOK} in \nolinkurl{PlaceLaws.lean}; and the complete proof bodies \nolinkurl{clsKL_is} and \nolinkurl{thetaKernel_eq_jet_thC}. The dependency chain has not been compiled or audited in full.

Root's AM-R-002 was independently whole-read and checked: content ideal and local detector denominator, dual norm and quotient by kernel, scaling/product-formula invariant, arithmetic-degree sign, actual normalized coefficient matrix and stabilizer masses all pass. Its formula must be applied to the actual terminal lift vector and actual invariant lattice; it does not supply them for the present signed seesaw.

\subsection{AM-DO-005. An explicit charge-preserving local bicanonical comparison}\label{am-do-005.-an-explicit-charge-preserving-local-bicanonical-comparison}

This gives a nonzero comparison after changing the target coefficient weight. It is included to make the exact scope of AM-DO-002 clear; it does not turn the canonical 032 class into an Alpöge arithmetic column.

\subsubsection{Data, splitting and polarization}\label{data-splitting-and-polarization}

Set m=4, r=2 in T-052. The real dual pair is U(2,1)×U(4) acting on the real symplectic space underlying C\^{}\{2,1\}⊗C⁴, with the T-052 convention ω=Im(h\_\{2,1\}⊗h\_4). Use its positive compatible complex structure i on the positive summand and −i on the negative summand. The polynomial Fock space is C{[}X\_ai,Y\_i{]}, a=1,2 and i=1,\ldots,4. On K\_G×U(4) the bare vacuum character has square

\begin{verbatim}
χ0²=(det A)^4 d^(−4) det g,
\end{verbatim}

so it is (det A)²d\^{}(−2)χ\_W(g), with χ\_W²=det g on the indicated metaplectic pullback. These are the exact T-052 conventions, not newly inferred global splitting data.

Twist the G=U(2,1) action by det\_G\^{}(−2), and untwist the compact U(4) factor by χ\_W\^{}(−1), just as T-052 untwists its ordinary column action. Both combined actions then descend in this selected model. The former is an honest determinant character. In the archive's \nolinkurl{splitFacO} convention, changing the character attached to the compact member by \(\kappa(z)=(z/\bar z)^2\) produces this det\_G\^{}(−2) factor: if \(x/\bar x=\det g\), then κ(x)\^{}(−1)=det(g)\^{}(−2). This states the splitting CHANGE relative to the fixed model; it does not invent a global absolute character or its finite components. The polarization is unchanged.

\subsubsection{Explicit lowest vector and intertwiner}\label{explicit-lowest-vector-and-intertwiner}

Put Δij=X\_\{1i\}X\_\{2j\}−X\_\{1j\}X\_\{2i\} and f=Δ12Δ34. Then

\begin{verbatim}
C_1 f=C_2 f=0,    Q_i f=f for every i,
(N_X,N_Y)=(4,0),  α=(2,2), k=1, b=0.
\end{verbatim}

Each column occurs once and f contains no Y, which proves the first two assertions directly. Under U(2), f transforms as det²; with the vacuum shifts, its bare lowest K\_G weight is (4,4;−2). After the determinant twist it becomes (2,2;−4), the bicanonical compact type from THETA3 B14. Its restriction to the traceless compact circle has weight12, consistently with H=8+N\_X+2N\_Y. This does not contradict AM-DO-002, which excludes weight6.

The polynomial f expands into four distinct squarefree monomials, each with coefficient±1. Under the Fischer norm its squared norm is therefore4. Let e=f/2. The G-cyclic module generated by e lies in the T-052 irreducible lowest-weight Howe summand indexed by α=(2,2),b=0. This irreducibility is a carried T-052 (OC.6),(OC.20)--(OC.22) input, not an assertion made for an arbitrary lowest vector. After the determinant twist it is the irreducible bicanonical holomorphic discrete series π\_2.

Choose a unit lowest vector v\_2 in a fixed π\_2 model. The explicit intertwiner on algebraic vectors is

\begin{verbatim}
I(u·e)=u·v_2,   u∈U(gl_3(C)),
\end{verbatim}

where the action on the source includes the det\^{}(−2) twist. This is well-defined: the carried irreducible lowest-weight identification gives the same annihilator ideal of the fixed lowest type in the two modules; equivalently both are the same irreducible lowest-weight quotient. It is a bijection. Its unit-vector normalization fixes its scalar uniquely, and invariance of the positive forms makes it unitary on completion. Because the G action commutes with every Q\_i, its whole domain has charge(1,1,1,1). Thus the comparison preserves the prescribed charge sector.

An actual nonzero coherent target may be obtained from any nonzero canonical wedge ξ supplied by A7: ξ² is a nonzero section of K\_X², and its automorphic coefficient generates a copy of π\_2. This observation specifies the archimedean recipient. No compatibility with the finite Weil action or membership in the global U(4) theta image follows from it.

\subsubsection{The multiplicity factor cannot be discarded}\label{the-multiplicity-factor-cannot-be-discarded}

The charge-one lowest space is the multilinear component of S\_(2,2)(C⁴). It has basis

\begin{verbatim}
f1=Δ12Δ34,  f2=Δ13Δ24,
\end{verbatim}

with the third pairing f3=Δ14Δ23 satisfying f1−f2+f3=0. Their Fischer Gram matrix is {[}{[}4,2{]},{[}2,4{]}{]}. The ordinary S4 action is its irreducible Specht module {[}2,2{]}, of dimension2. This is also T-052 (OC.4) at n=1,λ=(2,2,0,0), with its vacuum untwist. Consequently the full charge-selected local comparison retains

\begin{verbatim}
π_2 ⊗ [2,2],
\end{verbatim}

rather than mapping the S4 multiplicity to a scalar. For example choose the two orthonormal lowest vectors e1=f1/2 and e2=(f2−f1/2)/√3 and send u·e\_j to (u·v\_2)⊗w\_j, where w1,w2 are the corresponding orthonormal basis of {[}2,2{]}. This fixes the whole G×S4 intertwiner; S4 acts through its exact Gram-compatible matrices on the second factor.

For all n≥1 this is the T-052 parameter family λ=(n+1,n+1,n−1,n−1), since α=(2,2),k=1,b=0. The ordinary S4 comparison may include the source's sign\^{}(n+1), but {[}2,2{]} is self-conjugate and hence unchanged up to its specified Specht isomorphism. This parameter observation is carried from (OC.1) and (OC.4), not a new plethysm computation.

\subsubsection{A normalization with an actual local value}\label{a-normalization-with-an-actual-local-value}

For the compact subgroup U(2)⊂U(2,1), e1 has character det² in the twisted representation. With probability Haar measure and the compact test representation det²,

\begin{verbatim}
∫_{U(2)} <e1,π_2(diag(A,1))e1> det(A)^(−2) dA=1.
\end{verbatim}

The integrand is identically1. This is precisely the q=2 specialization of the carried THETA3 B9/B14 local calculation. It is not the global kernel/covolume or terminal-mass value requested in question2.

\subsubsection{Consequence}\label{consequence}

A blanket assertion that no charge-preserving comparison with a ball-quotient representation exists would be false: this is one, with the dual pair, covering/splitting adjustment, polarization, vector, multiplicity factor and normalization given explicitly. It uses a bicanonical coherent section. THETA3 B14 shows why that differs from the ordinary canonical cohomology class and why a diagonal cup-product argument does not return it to that class. Finite-place compatibility, actual global theta occurrence, arithmetic columns and controlled period costs remain separate. Q-ALPOGE and Q-ALPOGE-METHOD therefore remain open in their original global form.

\subsection{AM-DO-006. Full charge torus obstruction to a polarized compact tensor comparison}\label{am-do-006.-full-charge-torus-obstruction-to-a-polarized-compact-tensor-comparison}

\subsubsection{Exact scope and statement}\label{exact-scope-and-statement}

Keep the T-052 positive polarization and m≥3, r=m−2. Let T=(U(1))\^{}m be the individual column-charge torus. Its action on the creation space E\_m has character multiset

\begin{verbatim}
E_m = ⊕_{i=1}^m (C_{e_i}^{⊕r} ⊕ C_{−e_i}),
\end{verbatim}

where e\_i are the independent coordinate characters. In particular dim\_C E\_m=m(m−1). Vacuum and splitting characters do not occur in this formula: these are the commutator weights on creation operators, so a scalar twist of the whole oscillator cancels.

A \textbf{polarized full-charge tensor comparison} to a totally definite dual pair U(a)×U(b) means a real symplectic isomorphism of the full Heisenberg spaces, preserving the specified positive compatible polarizations, together with a torus homomorphism T→U(a)×U(b), such that the isomorphism intertwines the individual charge torus. The two compact factors may be dualized, and either common orientation of the symplectic form may be used. These changes only dualize one or both defining representations and do not change their tensor-product form. A comparison of this kind requires ab=m(m−1).

There is no such comparison with a,b≥2. In particular, even when dimensions happen to agree, it cannot identify the T-052 full charged oscillator with a single one of Alpöge's nonzero even-rank, totally definite tensor-pair oscillators while keeping all individual charges and the indicated positive polarizations.

This statement does not concern a map only on a chosen simultaneous-charge eigenspace, a map of Harish--Chandra modules which does not intertwine the Heisenberg action, changes of the real form or polarization not identified in the hypothesis, extra oscillator factors, or a rank-one compact factor. It is a precise obstruction to one proposed direct comparison, not a proof that Q-ALPOGE has no solution by any method.

\subsubsection{Elementary support lemma}\label{elementary-support-lemma}

Let m≥3 and S=\{±e\_1,\ldots,±e\_m\}⊂R\^{}m. If A and B are nonempty finite sets of vectors with A+B=S, then A or B is a singleton.

Proof. Suppose each set has at least two elements. Choose distinct α,α′∈A and β,β′∈B and form

\begin{verbatim}
x=α+β,  y=α′+β,  z=α+β′,  w=α′+β′.
\end{verbatim}

All lie in S, and x+w=y+z. Adjacent vertices in this displayed rectangle differ, because the selected α's and β's differ. If x=w, then x=(y+z)/2; y≠z, since otherwise x=y would follow. This contradicts the fact that every element of S is an extreme point of its convex hull. The same argument excludes y=z. Thus x,y,z,w are four distinct elements of S.

The sum of two distinct elements of S is either zero, when they are opposite, or a vector with exactly two nonzero coordinates equal to ±1. In the latter case the sum uniquely determines the unordered pair of vertices. Since the two pairs \{x,w\} and \{y,z\} are disjoint, their common sum must therefore be zero. Hence

\begin{verbatim}
α+α′+β+β′=0.
\end{verbatim}

Now fix α,α′ and β and vary β′ over B\{β\}. This equality determines β′ uniquely. It follows that B has at most two elements. By symmetry A has at most two elements. But then A+B has at most four elements, whereas S has 2m≥6. This contradiction proves the lemma. The proof uses no lattice or integrality hypothesis. ∎

\subsubsection{Full proof of the representation obstruction}\label{full-proof-of-the-representation-obstruction}

Preservation of the positive polarization identifies the complex creation spaces equivariantly for T. Restrict the defining a- and b-dimensional compact-factor representations along the claimed torus homomorphism. Compactness makes each a direct sum of characters. Write their character multisets as \{α\_1,\ldots,α\_a\} and \{β\_1,\ldots,β\_b\}, absorbing any contragredient into the signs of these characters. The creation space of the definite tensor pair then has character multiset

\begin{verbatim}
{α_p+β_q : 1≤p≤a, 1≤q≤b}.
\end{verbatim}

All multiplicities are nonnegative, so its support is exactly A+B, where A and B are the two character supports. The required isomorphism forces A+B=S. By the support lemma one factor, say A, consists of one character α. That factor representation therefore consists of a copies of α. Every character multiplicity in its tensor product is divisible by a. But every negative character −e\_i in E\_m has multiplicity1. Consequently a=1. If B is the singleton, the same argument gives b=1. Either conclusion contradicts a,b≥2.

This proof also shows why scalar splitting shifts do not help. They only multiply an entire oscillator representation by a character; the adjoint weights of the creation operators remain unchanged. If one allows an affine common shift of all creation weights in an intermediate notation, it can be absorbed into one factor's α's before applying the same argument. Passing to a common cover to define half-integral vacuum characters similarly leaves the proof intact. ∎

\subsubsection{Algebraic geometry and K-theory form}\label{algebraic-geometry-and-k-theory-form}

For the split torus T\_Z=G\_m\^{}m the creation module is the free T\_Z-equivariant vector bundle

\begin{verbatim}
E_m = ⊕_i (L_{e_i}^{⊕r} ⊕ L_{−e_i})
\end{verbatim}

on Spec Z, with character class rΣ\_i t\_i+Σ\_i t\_i\^{}(−1) in K\_0\^{}\{T\_Z\}(Spec Z)=Z{[}t\_1\^{}±1,\ldots,t\_m\^{}±1{]}. The theorem says that this positive character cannot be a product of two positive character classes of ranks at least2. The adjective positive means actual representations or vector bundles; virtual K-theory factorizations with negative coefficients are not excluded. Thus an ordinary Grothendieck-group equality is weaker than the tensor-bundle identification needed for the positive oscillator construction.

The exact Grothendieck-framework locator for the weight decomposition is SGA3, Exposé I, Proposition4.7.3; Corollary4.7.3.1 gives its graded-algebra/equivariant-affine-scheme form. Both were read anew in the French recomposed text. The relevant archive framework in the supplied materials is cote133 scan006 (pre-λ Grothendieck groups and universal identities in classifying topoi), while the exact mathematical input and new statement are T-052 (OC.3), (OC.7)--(OC.12) and AM-DO-006 here. No assertion that SGA3 or cote133 proves this factorization obstruction is made.

\subsection{AM-DO-007. Canonical rank cap for a compact middle source}\label{am-do-007.-canonical-rank-cap-for-a-compact-middle-source}

\subsubsection{Statement}\label{statement-6}

For the real dual pair U(n,1)×U(d), let the second factor be totally definite and n≥1, d≥1, q≥1. In the positive-energy orientation, if the irreducible holomorphic coefficient representation of the qth canonical power occurs as a subquotient of the oscillator, with arbitrary determinant splitting twists, then

\begin{verbatim}
d≤q(n+1).
\end{verbatim}

The same necessary bound holds in the simultaneous contragredient convention. In the opposite oscillator orientation with the original holomorphic target, no occurrence is possible. As in AM-DO-002, the comparison is made on a common metaplectic cover when necessary.

For the unchanged canonical coefficient of the 032 ball surface, n=2 and q=1, the compact source must have d≤3. Therefore the first definite edge A\_2→S\_d in a literal MainX realization cannot at the same time satisfy the existing \nolinkurl{StdPair.rk} conditions d\textgreater2 and Even d.~Enlarging the middle totally definite rank cannot put that canonical realization into this particular arithmetic constructor.

\subsubsection{Full proof}\label{full-proof-3}

Here n and d are independent; no T-052 relation between them is assumed. The general U(n,1)×U(d) oscillator formulas are the same creation-annihilation formulas carried in T-052 (OC.7), with a=1,\ldots,n and i=1,\ldots,d.~Equivalently they are Basile et al.~(2020), §5.2, equations (5.21)--(5.30), specialized to q=1 negative direction and d compact columns. The positive compact diagonal operators have vacuum shift d/2 each, and the negative diagonal operator is −N\_Y−d/2. Consequently the generator of diag(zI\_n,z\^{}(−n))⊂SU(n,1) is

\begin{verbatim}
H=N_X+nN_Y+nd,
\end{verbatim}

with every eigenvalue at least nd. The qth canonical coefficient has compact character det(A)\^{}q d\_0\^{}(−nq), where d\_0 denotes the last U(1) entry, and therefore has the weight qn(n+1) for this circle. Its lowest vector must appear in any nonzero subquotient realization. Exactness of compact weight decomposition, as proved in AM-DO-002, gives nd≤qn(n+1). Division by n\textgreater0 proves the bound.

Every determinant splitting twist is trivial on this traceless circle. Reversing the entire oscillator orientation replaces its spectrum by eigenvalues at most−nd, which excludes the unchanged positive canonical lowest weight. Simultaneously dualizing the target negates its lowest weight too and gives the same inequality. These cover the two definite signs of the auxiliary Hermitian space with fixed additive-character convention.

For the ball, d≤3 follows. The actual source \nolinkurl{SplitPart.lean}, \nolinkurl{StdPair.rk} lines43--44, requires 0\textless rank A\textless rank S and both ranks even; since rank A=2, it requires d≥4. The inequalities are incompatible. This proves the last assertion. ∎

\subsubsection{Consequence and boundary}\label{consequence-and-boundary}

This obstruction is additional to the simultaneous sign obstruction AM-DO-004. Even at the active real place alone, the unchanged canonical coefficient prevents using a larger even middle group to meet the strict kernel-constructor range. A signed equal-rank kernel theorem, a different arithmetic constructor, or a change of coherent coefficient is therefore needed for that route to G1. The bound is necessary and not sufficient: it does not prove finite local occurrence, global nonvanishing, the kernel identity, or the terminal mass equation when d≤3. At q=2, d=4 is allowed by this bound and the explicit AM-DO-005 lowest vector realizes the corresponding local bicanonical component.

\subsection{T-053, the Fourier comparison, and the precise geometric category}\label{t-053-the-fourier-comparison-and-the-precise-geometric-category}

Carry T-053, \nolinkurl{work/theta/charged_incidence_rees_report.md}, (IR.1)--(IR.5), (IR.12)--(IR.14). Its incidence ring is B=Q{[}X,Y{]}/(R\_a), R\_a=Σ\_iX\_aiY\_i, and its actual frame Rees algebra is ⊕\emph{\{k≥0\}B}(k,\ldots,k). The source gives the generator maps, multiplication and quotient scheme; these are settled results and are not rederived here. In particular its equal-charge construction is an algebraic torus-weight operation in SGA3 I4.7.3 and4.7.3.1.

The partial Fourier map AM-DO-003 has the exact additional formulas

\begin{verbatim}
F(R_a)=Σ_i X_ai∂_Zi,
F(C_a)=−Σ_i Z_i∂_Xai,
\end{verbatim}

where C\_a=Σ\_i∂\_Xai∂\_Yi. These follow by substitution into the already-proved Weyl-algebra isomorphism. They exchange the displayed multiplication constraints and differential constraints with off-diagonal compact Lie operators. They do not carry them to multiplication by polynomials in the ordinary positive definite coordinates X,Z: indeed

\begin{verbatim}
[F(R_a), Z_j]=X_aj≠0,
\end{verbatim}

whereas every such coordinate multiplication operator commutes with Z\_j. This proves that this specific Fourier map does not give a morphism to the ordinary definite polynomial incidence scheme preserving the positive coordinate polarization. It remains an exact Weyl-module comparison into the delta realization, and the original Rees algebra still has its own carried scheme interpretation. No assertion that an arbitrary algebraic comparison between selected Rees pieces is impossible follows.

Likewise AM-DO-006 concerns the full creation-space tensor factorization, not a factorization of a single equal-charge quotient B\_(k,\ldots,k). It therefore does not refute T-053, its positive row-weight counts, or its exact remaining highest-type and Specht projectors. It excludes the proposed direct identification of the whole charged oscillator with one nontrivial compact tensor-pair oscillator under the named polarization and charge requirements.

Primary framework locators read in this run: Grothendieck, EGA IV4 (1967), Theorem16.11.2 and Remarks16.11.3(ii), pp54--55, \url{https://www.iut-theory.org/divers/ega44.pdf}; Demazure--Grothendieck, SGA3, Exposé I, Proposition4.7.3 and Corollary4.7.3.1, French recomposed edition, \url{https://www.iut-theory.org/divers/SGA3I.pdf}. The former distinguishes full differential operators from the integral Weyl algebra; the latter identifies actual torus-equivariant modules and schemes with gradings. Neither identifies these realizations with an integral theta lift, a Dieudonné crystal, an ℓ-adic sheaf with a specified Frobenius, or a motive. Those require additional structures and actual comparison maps.

\subsection{AM-DO-008. The actual kernel equation as period-rescaling and lattice-membership tests}\label{am-do-008.-the-actual-kernel-equation-as-period-rescaling-and-lattice-membership-tests}

This is a reduction of the already displayed MainX equation, not a computation of its missing theta entries. It makes precise which finite calculations would certify a selected arithmetic column, and which additional identity is needed for the source's full \nolinkurl{ClsK.Core}.

\subsubsection{Polynomial and period equations}\label{polynomial-and-period-equations}

Fix one edge j, all finite charts and class representatives, the actual arithmetic \nolinkurl{abelian₂} functionals, and the corresponding analytic kernel. Index a finite list of tests by i=(α,γ,a,a′,µ), where P\_µ is a monomial polynomial. Put

\begin{verbatim}
Θ_i = volLev(X_j,U_j,α) θ_K,j(h_j,α,a,h_j+1,γ,a′; Φ_j,Pµ),
A_i = abelian₂_j,α,γ,a,a′(castP P_µ),
d_i,τ = degree of P_µ in the variables at τ.
\end{verbatim}

The exact source equation restricted to these tests is

\begin{verbatim}
Θ_i = c_j (∏_τ s_j,τ^(d_i,τ)) A_i,              (DO8.1)
\end{verbatim}

where c\_j=cK\_j≠0 and every s\_j,τ=sc\_j,τ≠0. No covolume has been omitted from Θ\_i. If the c\_j or s\_j,τ are already fixed by the arithmetic pin, (DO8.1) must use those values.

For freely allowed complex rescalings c\_j,s\_j,τ, a finite system (DO8.1) is solvable if and only if:

\begin{enumerate}
\def\labelenumi{\arabic{enumi}.}
\item
  A\_i=0 exactly when Θ\_i=0;
\item
  on the nonzero indices, with r\_i=Θ\_i/A\_i and u\_i=(1,(d\_i,τ)\_τ), one has

\begin{verbatim}
∏_i r_i^(v_i)=1
\end{verbatim}

  for every integer relation Σ\_i v\_i u\_i=0.
\end{enumerate}

In particular all ratios belonging to the same multidegree must be identical. For a single homogeneous coefficient space, every two-by-two minor Θ\_i A\_l−Θ\_l A\_i must vanish, together with the stated support condition; a nonzero analytic period alone does not impose these identities.

Proof. The scalar matrix substitution in \nolinkurl{mvZ(1,sI)} multiplies a monomial by ∏\_τs\_τ\^{}(d\_i,τ), proving (DO8.1). Nonzero scalars prove condition1; multiplying the equations according to an integer relation proves condition2. Conversely let M be the integer matrix with rows u\_i for the nonzero indices. Choose integer unimodular matrices U,V with UMV=D in Smith form, whose nonzero diagonal entries are δ1,\ldots,δh\textgreater0. Replacing target ratios by r′\_k=∏\_i r\_i\^{}(U\_ki) and the torus variables by x\_l=∏\_k y\_k\^{}(V\_lk) changes the system into

\begin{verbatim}
y_k^(δk)=r′_k  (k≤h),   1=r′_k  (k>h).
\end{verbatim}

The zero rows are integer relations and therefore their equations hold by condition2. Each nonzero complex number has a δkth root, so choose those roots and arbitrary nonzero unused variables. Undoing the unimodular substitutions solves the system. The equal-multidegree criterion is exactly the specialization where all u\_i agree. This proves every assertion. Over a number field K the same proof additionally requires the indicated r′\_k to be δkth powers in K×; it does not silently adjoin those roots or identify periods with algebraic numbers. ∎

\subsubsection{Integral finite-input equations}\label{integral-finite-input-equations}

If a finite-dimensional ansatz for the finite Schwartz input has been specified, the analytic left side of (DO8.1) is linear in its input coefficients. After choosing actual integral generators, fixing the period scalars, and making a common algebraic denominator clearing that has itself been justified, write a localized system as

\begin{verbatim}
Cx=b,   C∈Mat_(q×s)(O_v),  b∈O_v^q.          (DO8.2)
\end{verbatim}

Here b is the specified jet vector including c\_j and the period scales. Let π be a uniformizer. Suppose UCV has nonzero Smith diagonal π\^{}(a\_1),\ldots,π\^{}(a\_h), with U,V invertible over O\_v, and put b′=Ub. Then an integral solution exists exactly when

\begin{verbatim}
b′_i=0 for i>h, and v(b′_i)≥a_i for i≤h.     (DO8.3)
\end{verbatim}

Assuming the zero-row conditions hold, the least d≥0 admitting a solution x∈π\^{}(−d)O\_v\^{}s is

\begin{verbatim}
max(0, max_(i≤h)(a_i−v(b′_i))),              (DO8.4)
\end{verbatim}

with v(0)=+∞ and an empty maximum contributing no positive bound.

Proof. Smith form over a DVR follows by choosing a matrix entry of least valuation, moving it to the top left, and clearing its row and column using divisibility by that entry. Repeat on the smaller submatrix; multiplication by units turns every nonzero diagonal entry into π to its valuation. All row and column operations are invertible over O\_v. With y=V\^{}(−1)x, (DO8.2) becomes π\^{}(a\_i)y\_i=b′\_i for i≤h and 0=b′\_i afterwards. This proves (DO8.3). Permitting y\_i∈π\^{}(−d)O\_v changes the first inequalities to v(b′\_i)≥a\_i−d, giving (DO8.4). These basis changes preserve the lattice, so the criterion does not depend on the selected Smith bases. ∎

This is a conditional integral criterion for a \emph{specified} finite-input span. Allowing a larger Schwartz space may change C; choosing arbitrary analytic coordinates and treating them as integral does not instantiate (DO8.2). For the terminal detector, the analogous rank-one image-ideal and archimedean-norm costs are precisely AM-R-002, applied to the actual terminal vector and the fixed mass-weighted row lattice. Its stabilizer and Fock denominators must be retained.

\subsubsection{The remaining all-degree and family requirements}\label{the-remaining-all-degree-and-family-requirements}

The actual \nolinkurl{OnFramesRdApply.lean} definition \nolinkurl{ClsK.Core}, lines108--113, quantifies over every polynomial P. Thus a finite matrix check certifies its restriction to the chosen monomial span, not the entire predicate. A full witness needs the equality for every degree, or a proved equality of the relevant analytic theta germs from which every jet follows. The finite class sets do not make the polynomial quantifier finite.

The source's \nolinkurl{ClsK.Kval}, lines115--121, subsequently uses fixed finite weight spaces and the reciprocal stabilizer masses. The terminal \nolinkurl{ClsT.Is.top/tail}, lines143--150, must hold for all the stated terminal vectors and rows, with one fixed branching map, level and top covolume. AM-R-002 supplies the cost of normalizing one already constructed detector; it cannot substitute for those simultaneous identities. Finally \nolinkurl{PinColT.cst} in \nolinkurl{LevChiEqOne.lean} fixes cK\_j from the actual covolume, finite datum multiplier, vacuum scalar and determinant-root characters, so the free rescalings in the first criterion represent the most permissive test. Failure of that test is an obstruction; success does not override the specified normalization.

For the repeated-type geometry AM-DO-001 the actual matrices Θ,A,C,b and terminal mass matrix have not been calculated. Consequently no global value for the requested G1 pairing follows from this reduction. The concrete remaining calculation is the signed equal-rank first kernel, with its all-degree CM-jet identity, followed by the compatible second kernel and terminal row equations. AM-DO-004 and AM-DO-007 show why the existing common positive, strictly increasing constructor does not supply these entries.

\subsection{AM-DO-009. Exact tensor-to-oriented-Hom comparison, including the source dual}\label{am-do-009.-exact-tensor-to-oriented-hom-comparison-including-the-source-dual}

\subsubsection{The global symplectic and group map}\label{the-global-symplectic-and-group-map}

Use the first-argument-linear Hermitian convention of 032, §4.2, equation(4.4), printed pp14--15. Let V and P have Gram matrices G\_V and G\_P, and let δ and ξ0 be nonzero purely imaginary elements of E. The 032 space is

\begin{verbatim}
W032 = V⊗_E P,
B032(x,y)=Tr_{E/F}(δ (h_V⊗h_P)(x,y)).
\end{verbatim}

Define S to be the algebraic E-dual P\^{}∨, with Hermitian Gram

\begin{verbatim}
G_S=(ξ0/δ) G_P^(−t).                         (DO9.1)
\end{verbatim}

The ratio ξ0/δ belongs to F×. In the dual basis of S and the given basis of V, identify a tensor with its coefficient matrix:

\begin{verbatim}
I: V⊗P → Hom_E(S,V),
I(Σ_(i,a) X_ia v_i⊗p_a)=X.                 (DO9.2)
\end{verbatim}

Then I is E-linear and an exact isometry of the F-valued alternating spaces. It intertwines the tensor action of U(V)×U(P) with the oriented action of U(S)×U(V) under

\begin{verbatim}
(g,k) ↦ (k^(−t),g),
I((g⊗k)x)=g I(x) k^t.                      (DO9.3)
\end{verbatim}

This is an explicit underlying rational oscillator comparison. It is not a comparison with the all-positive \nolinkurl{StdPair} pin, and it does not alone establish any MainX kernel or covolume identity.

\subsubsection{Full proof}\label{full-proof-4}

The dual Hermitian Gram G\_P\^{}(−t) is Hermitian: this follows by transposing and conjugating the identity G\_P\emph{=G\_P. Scaling it by ξ0/δ∈F× preserves Hermitian symmetry and nondegeneracy. If k}G\_Pk=G\_P, transpose and invert this equality to obtain

\begin{verbatim}
(k^(−t))* G_P^(−t) k^(−t)=G_P^(−t).
\end{verbatim}

Thus k\^{}(−t) lies in U(S), proving that the group map is well-defined. The ordinary action on the algebraic dual is k\^{}(−t), so its inverse acts on the right of Hom matrices as k\^{}t. This proves(DO9.3).

Write x=Σ X\_ia v\_i⊗p\_a and y=Σ Y\_jb v\_j⊗p\_b. The first-argument-linear convention gives

\begin{verbatim}
(h_V⊗h_P)(x,y)
  =Σ_(i,a,j,b) conjugate(Y_jb)(G_V)_ji(G_P)_ba X_ia
  =tr(G_P^t Y* G_V X).
\end{verbatim}

By(DO9.1), G\_S\^{}(−1)=(δ/ξ0)G\_P\^{}t. The exact source definition \nolinkurl{hO} therefore gives

\begin{verbatim}
hO(S,V)(I(x),I(y))
  =tr(G_S^(−1)Y*G_VX)
  =(δ/ξ0)(h_V⊗h_P)(x,y).
\end{verbatim}

Multiplying by ξ0 and taking Tr\_\{E/F\} proves

\begin{verbatim}
symplO(S,V)(I(x),I(y))=B032(x,y).
\end{verbatim}

The same equality remains true after Tr\_\{F/Q\} and at every completion. The tensor-to-Hom map is bijective on the displayed bases, completing the proof. ∎

\subsubsection{Polarization, additive character, splitting and measure scope}\label{polarization-additive-character-splitting-and-measure-scope}

Transport 032's actual positive Lagrangian through I\_C. At the active place its two blocks are exactly the images of L⊗P\_μ and Q\emph{⊗P}\_μ from032(4.5); at the compact places they are the images of the specified embedding blocks. The positivity equation is preserved because the alternating forms are equal. This describes the chosen polarization rather than replacing it by an all-positive coordinate pin.

Choose the SAME global additive character and transport one rational Schrödinger position space and its self-dual measures through I. The associated Heisenberg groups then identify by (x,t)↦(I(x),t), with central coordinate and character unchanged. On the transported Schrödinger function spaces the intertwiner is explicitly f↦f∘I\^{}(−1) on the position spaces. The Weil actions coincide when both are pulled back from the same metaplectic representation and its same lifts. This assertion needs no change of additive character and no partial complex Fourier transform. The predicate \nolinkurl{IsStdChar D ψ (-1)} of the recovered source must separately be compared with 032's fixed theta multiplier; equality of alternating forms does not discharge that predicate. The map also does not identify separately chosen lattice measures or finite theta structures; those must be transported and compared.

The group relabelling in(DO9.3) changes determinant-root labels. Choose idele roots t,y with

\begin{verbatim}
t/c(t)=det k,  y/c(y)=det g.
\end{verbatim}

Then the source root for k\^{}(−t) is t\^{}(−1), and Alpöge's oriented character factor is

\begin{verbatim}
splitFacO(χ_S,χ_V)(t^(−1),y)
  =χ_V(t) χ_S(y).                          (DO9.4)
\end{verbatim}

This follows from \nolinkurl{LocalNewformFunction.lean}, line51. The root condition is \nolinkurl{TiedWeil.lean}, \nolinkurl{IsDetRoot}, lines64--66. At a split place \nolinkurl{IsDetAt}, lines77--78, fixes the particular root t=e·det k+(1−e); the same inverse substitution applies. The arguments of these characters are the determinant roots, not the determinants. Moreover \nolinkurl{IsWeilDatumO.om_split} retains a separate positive real factor r; (DO9.4) describes only its character factor.

Any proposed identification with already fixed 032 circle splittings must match these roots and this character factor. Transporting the existing metaplectic lifts gives one consistent comparison; matching arbitrary independently prescribed Hecke-character splittings additionally requires their determinant-character ratio. On SU(V) the ratio of two group splittings is trivial, as carried032 Lemma4.4, printed pp20--21, states. This does not set an arbitrary χ\_S(y) to1 for every possible root y. No unprovided finite components of a global Hecke character, positive normalization r or universal Weil-index scalar are assigned here.

\subsubsection{The explicit Q(ζ7) middle form}\label{the-explicit-qux3b67-middle-form}

For AM-DO-001, G\_P=−a I2. Formula(DO9.1) gives

\begin{verbatim}
G_S=s I2,  s=−ξ0/(δa).                      (DO9.5)
\end{verbatim}

With ξ0=√−7/7 and δ=ζ−ζ\^{}(−1), the signs of ξ0/δ at σ1,σ2,σ3 are(+,+,−): the quadratic Gauss sum √−7 is fixed by σ2 and negated by σ3, whereas all three imaginary parts of δ are positive. As sign(a)=(+,−,−), the signs of s are(−,+,−). In particular S is still totally definite. The scalar Gram −aI2 would give the wrong exact symplectic normalization for this independently fixed ξ0; it is not used here.

The literal definite chain has endpoint forms G\_A=I2 and G\_X2=(a/q)I2, q=Nξ0. Its actual oriented edge ratios are s and a/(qs). Their product is a/q, so AM-DO-004's simultaneous positive-pin obstruction applies unchanged. The conversion solves the rational symplectic and source-label issue; it does not solve the finite CM-jet, character, covolume, all-degree kernel or terminal-mass gates.

\subsection{AM-DO-010. The evaluation-point CM tori are explicit primitive-J powers}\label{am-do-010.-the-evaluation-point-cm-tori-are-explicit-primitive-j-powers}

Carry B19a's exact rational E-Hodge decomposition for one multiplicity line at the rational negative line x={[}Ee3{]}:

\begin{verbatim}
H^1(A_x,Q) ≃ U(d)^⊕2 ⊕ U(dflip),
A_x ∼_E B(d)^2 × B(dflip).                 (DO10.1)
\end{verbatim}

The action is the specified geometric E-action of B19a. In AM-DO-001, the positive set of d is\{1,4,5\}; changing the active pair gives the positive set of dflip as\{4,5,6\}.

For s∈(Z/7Z)×, write J\^{}(s) for the same abelian variety J with its E-action precomposed by σ\_s. Its H\^{}\{1,0\} type is sΦ, where Φ=\{1,2,3\} is the carried type of J. Therefore

\begin{verbatim}
B(d) ∼_E J^(4),   B(dflip) ∼_E J^(6),
A_x ∼_E (J^(4))^2 × J^(6).                 (DO10.2)
\end{verbatim}

Forgetting the E-action, A\_x is isogenous to J\^{}3. The two repeated multiplicity inputs of AM-DO-001 consequently use CM abelian varieties whose product is isogenous to J\^{}6.

Proof. On a type-τ eigenspace, replacing ι(a) by ι(σ\_s(a)) changes its eigencharacter to τ∘σ\_s. The Galois group is abelian and identified with multiplication modulo7, so the new type is sΦ. Direct multiplication gives

\begin{verbatim}
4Φ={4,1,5}={1,4,5},
6Φ={6,5,4}={4,5,6}.
\end{verbatim}

Each rank-one rational E-Hodge structure is determined by this type. Thus its E-Hodge structure is isomorphic to that of the indicated twist of J. The carried equivalence between rational Hodge maps in weight one and homomorphisms of complex abelian varieties up to isogeny, already used to prove B19a, supplies the two E-isogenies. Inserting them into(DO10.1) proves(DO10.2). The underlying abelian variety of every J\^{}(s) is J, giving the final two assertions. ∎

This is an isogeny and a type identification, not an identification of the actual polarization, symmetric theta line, finite theta structure, integral lattice, or algebraic tangent frame. The exact B19 comparison still retains c\_i, p\_i, N\_i, M\_i and κ\_x from(B19c)--(B19f). In particular a model by J powers makes the primitive CM source explicit but does not turn its theta-null period into1 or its arithmetic functionals into freely chosen linear maps. Root's AM-R-003 explains why these simple primitive-J powers cannot simultaneously be the old elliptic Serre-tensor model.

\clearpage

\section{ALP\_MACHINE --- question 3: the Kummer character and the precise transfer gate}\label{alp_machine-question-3-the-kummer-character-and-the-precise-transfer-gate}

Author/checker: GPT-6 Astra Pro, ALP\_MACHINE child \nolinkurl{/root/kummer_extension}; inherited reasoning effort is not separately exposed. Date: 21 November 2026. This is a mathematical contribution for integration into the main paper, not a claim that the requested RH extension has been completed.

\subsection{0. Outcome and scope}\label{outcome-and-scope}

The source has a field-independent analytic theorem that can be applied to the Kummer character. \textbf{AM-K005} below proves all of its pole-data hypotheses and identifies a single, explicit remaining arithmetic window estimate. For \(r=3\), \textbf{AM-K004} also realizes the character as a rank-one summand in a tensor construction from two explicit CM elliptic curves. These are constructive reductions. They do not prove the missing window estimate.

The literal existing character host cannot receive the Kummer splitting field or its induced representation. This is stronger than a missing notation adapter: \textbf{AM-K002} proves three exact obstructions---non-CM splitting field, failure of norm pullback from a rational Dirichlet character, and failure of a finite Artin-formalism factorization into rational Dirichlet characters. Each obstruction is explicitly scoped; none asserts that a different CM-geometric construction is impossible.

Alpöge's theorem and archive are used as accepted theorems, as directed. CX-ALG-024, CX-R02, CX-G06, DP6, DP32, DP34, DP42 and the supplied source-interface results are carried with their existing status. No check of the validity of the accepted Dirichlet GRH theorem is undertaken. Masullo's consequences remain conditional. No Lean compilation was performed in this contribution; the archive checks below are checks of actual declarations and their types.

\subsubsection{0.1 Source provenance and what was actually read}\label{source-provenance-and-what-was-actually-read}

The primary note is Alpöge, \emph{Remark}, pp.~3--4, numbered timeline steps 2--20, \url{https://alpo.ge/remark.pdf}. The note calls this a rough timeline. Its historical step numbering is not itself a one-to-one module index.

Actual Lean text was recovered at pinned mirror commit \nolinkurl{7cbf7a0c32fca1badf05931879a37c3f8a2b03b4} of \nolinkurl{CocoSgt/riemann-hypothesis-lean}, under \nolinkurl{src/}. The local copy is \nolinkurl{work/agents/kummer_sources/}; \nolinkurl{manifest.json} records each downloaded file, byte count and SHA-256. Two files checked against the supplied original-author source manifest already match: \nolinkurl{Formalism/SizeWindowBrackets/GrowthBoundGivesRiemannHypothesis.lean} and \nolinkurl{Development/AutomorphicArithmetic/ThetaCorrespondence/ThetaPairLifts/ThirdPlaceSeries/SiegelEisensteinSecondTermHypotheses.lean}. The remaining copies are pinned mirror text, corroborated where applicable by the supplied original-archive map; this contribution does not claim every mirror file has independently been compared with the original archive.

The supplied map is \nolinkurl{RETURN_THETA3_checkpoint_07_20261121T0604Z/supplied_sources/ALPOGE_LEAN_MAP_20261120.md}, especially §§3--4. Its low-confidence keyword correspondences for steps 11, 14, 17 and 19 are not promoted to a complete dependency audit here.

\subsection{\texorpdfstring{1. Setup and the \(p=r\) edge case}{1. Setup and the p=r edge case}}\label{setup-and-the-pr-edge-case}

Let \(r\) be an odd prime, \(p\) a rational prime, \(\alpha=p^{1/r}>0\), and

\[
K=\mathbf Q(\zeta_r),\qquad F=\mathbf Q(\alpha),\qquad E=K(\alpha).
\]

The user assumes \(r\nmid(p-1)\). The algebraic results below do not need this restriction. It matters for the surrounding finite-field construction. Let \(\chi=\chi_{p,r}:G_K\to\mu_r\) be \(\sigma\mapsto\sigma(\alpha)/\alpha\), viewed by class field theory as the primitive finite-order Hecke character with its actual conductor. Write \(L(s)=L_K(s,\chi)\), retaining all primitive local factors.

\subsubsection{\texorpdfstring{AM-K000. The edge \(p=r\) is permitted and is algebraically nondegenerate}{AM-K000. The edge p=r is permitted and is algebraically nondegenerate}}\label{am-k000.-the-edge-pr-is-permitted-and-is-algebraically-nondegenerate}

\textbf{Statement.} The displayed condition permits \(p=r\). For this case \([E:K]=r\), the same affine Galois group and character analysis used below apply, and the finite-field task of constructing \(\mathbf F_{p^p}\) has the elementary irreducible polynomial \(T^p-T-1\in\mathbf F_p[T]\). Its use requires no RH premise.

\textbf{Proof.} The polynomial \(T^r-r\) is Eisenstein over \(\mathbf Q\), so \([F:\mathbf Q]=r\). Because \([K:\mathbf Q]=r-1\), the intersection \(F\cap K\) is \(\mathbf Q\), and \([E:K]=r\). Equivalently, at the unique prime \(\mathfrak r\) of \(K\) above \(r\), the normalized valuation satisfies \(v_{\mathfrak r}(r)=r-1\not\equiv0\pmod r\), so \(r\) is not an \(r\)-th power in \(K\). The statement is not that \(T^r-r\) is Eisenstein over \(K\); its constant-term valuation there is \(r-1\). The automorphisms \(\alpha\mapsto\zeta_r^b\alpha\), \(\zeta_r\mapsto\zeta_r^u\) give the full affine group of order \(r(r-1)\).

For the finite-field assertion, let \(a\) be a root of \(T^p-T-1\) in an algebraic closure of \(\mathbf F_p\). Then \(a^p=a+1\), and induction gives \(a^{p^m}=a+m\) for every nonnegative integer \(m\). The least positive \(m\) with \(a^{p^m}=a\) is \(p\). Its Frobenius orbit therefore has length \(p\), so its minimal polynomial has degree \(p\) and equals the given degree-\(p\) polynomial. ∎

The previously accepted CX-ALG-024 was stated for \(p\ne r\). The finite-field observation separates the omitted wild case from that statement without changing it.

\subsection{2. The exact target representation and L-function}\label{the-exact-target-representation-and-l-function}

\subsubsection{AM-K001. A single real L-function for all nontrivial Kummer powers}\label{am-k001.-a-single-real-l-function-for-all-nontrivial-kummer-powers}

\textbf{Statement.} Set \(N=\operatorname{Gal}(E/K)\simeq C_r\), \(G=\operatorname{Gal}(E/\mathbf Q)\simeq\operatorname{AGL}_1(\mathbf F_r)\), and let \(P\) be the permutation representation of \(G\) on the \(r\) roots of \(T^r-p\). For every nontrivial character \(\vartheta\) of \(N\),

\[
\rho:=\operatorname{Ind}_N^G\vartheta\simeq P-\mathbf 1.
\tag{K1}
\]

Consequently, for \(a=1,\ldots,r-1\),

\[
L_K(s,\chi^a)=L_{\mathbf Q}(s,\rho)
=\frac{\zeta_F(s)}{\zeta_{\mathbf Q}(s)}=:L(s).
\tag{K2}
\]

The identities hold with ramified local factors as well. The resulting Dirichlet series has real coefficients, \(L(\bar s)=\overline{L(s)}\), and its representation is self-dual. Under the accepted Dirichlet GRH theorem, RH for this \(L\) is equivalent to RH for \(\zeta_F\) and for \(\zeta_E\).

\textbf{Proof of the additional representation statement.} Write an affine map as \(x\mapsto ux+b\). The permutation character has values \(r\), \(0\), \(1\), respectively, at the identity, a nonidentity translation, and an element with \(u\ne1\). Its augmentation character therefore has values \(r-1,-1,0\).

Since \(N\) is normal, the induced character vanishes off \(N\). On \(N\), its restriction is the sum of the conjugates of \(\vartheta\) under \(\mathbf F_r^\times\). This group acts transitively and freely on all nontrivial characters of \(N\). Their sum has value \(r-1\) at the identity and \(-1\) at every other translation. The two characters agree, proving (K1). This also proves that inducing any nontrivial power gives the same representation.

Artin induction compatibility for L-functions now gives (K2); the standard permutation representation identity gives \(L(P)=\zeta_F\), while \(L(\mathbf1)=\zeta_{\mathbf Q}\). These identities are local identities on inertia invariants, so they include ramified primes. Since \(P-\mathbf1\) is defined over \(\mathbf R\) and self-dual, the asserted reality and duality follow. Equivalently, conjugation of the Euler product sends \(\chi\) to \(\chi^{-1}\), which has the same L-function by (K1).

For the final equivalence, use the already checked CX-ALG-024 identity \(\zeta_E=\zeta_K L^{r-1}\); it is not rederived here. Inside \(0<\Re s<1\), both \(\zeta_{\mathbf Q}L\) and \(\zeta_K L^{r-1}\) are products of holomorphic functions, so their zero divisors add with nonnegative multiplicities. Dirichlet GRH puts every zero of \(\zeta_{\mathbf Q}\) and \(\zeta_K\) on the critical line. The equivalence follows. ∎

\textbf{Sources.} CX\_GRH/GRH\_NIGHT, CX-ALG-024, in \nolinkurl{GRH_NIGHT.tex}, subsection containing lines 14723--14826 of the supplied checkpoint. For induction, Milne (2020a), VIII §10, pp.~261--262. Milne uses ``Dirichlet'' also for ray-class characters of general number fields in that discussion; this does not make those characters rational Dirichlet characters of \((\mathbf Z/N\mathbf Z)^\times\).

\subsection{3. Three exact obstructions to literal reuse}\label{three-exact-obstructions-to-literal-reuse}

\subsubsection{AM-K002. No cyclotomic host, norm pullback, or rank-one Dirichlet factorization}\label{am-k002.-no-cyclotomic-host-norm-pullback-or-rank-one-dirichlet-factorization}

\textbf{Statement.} For odd \(r\):

\begin{enumerate}
\def\labelenumi{\arabic{enumi}.}
\tightlist
\item
  \(E\) is not a CM field and cannot embed into a CM field. In particular it cannot be the field supplied to the literal \nolinkurl{IsChiField} record.
\item
  \(\chi\) is not the pullback of a rational Dirichlet character along the idelic norm \(N_{K/\mathbf Q}\).
\item
  \(\rho\) is not a virtual finite sum of one-dimensional rational Galois characters. Thus a finite factorization of its L-function into rational Dirichlet L-functions cannot be obtained by equality of Artin Euler factors, even after enlarging the common finite Galois extension.
\end{enumerate}

These statements do not exclude CM abelian varieties defined over \(K\), tensor constructions involving their twists, or a new Hecke-character version of Alpöge's arithmetic host.

\textbf{Proof.} Let \(\sigma(\alpha)=\zeta_r\alpha\) and let \(c\) be ordinary complex conjugation. Since \(c(\alpha)=\alpha\) and \(c(\zeta_r)=\zeta_r^{-1}\),

\[
c\sigma c^{-1}=\sigma^{-1}\ne\sigma.
\tag{K3}
\]

Complex conjugation is central in the Galois group of a normal CM field: its intrinsic involution commutes with every automorphism. Thus \(E\), which is normal, is not CM. There is also a direct containment obstruction. In a CM field embedded in \(\mathbf C\), the subfield of elements fixed by ordinary complex conjugation is totally real. The positive real number \(\alpha\) cannot be in that subfield, because its minimal polynomial \(T^r-p\) has nonreal conjugates. If an embedding of the normal field \(E\) into a CM field existed, its image in \(\mathbf C\) would again be \(E\) and would contain \(\alpha\), a contradiction. This proves (1). The central-involution characterization is Milne (2020b), Proposition 1.4 and Corollary 1.5, p.~10.

For (2), conjugation acts on the Kummer character by \(\chi^c=\chi^{-1}\), from (K3). A character pulled back through \(N_{K/\mathbf Q}\) is \(c\)-invariant, since \(N(c x)=N(x)\). It would follow that \(\chi=\chi^{-1}\). An order-\(r\) character for odd \(r>1\) cannot satisfy this equality.

For (3), the commutator of a multiplier \(h_u:x\mapsto ux\) with a translation \(\sigma_b:x\mapsto x+b\) is translation by \((u-1)b\). Choosing \(u\ne1\) shows that every translation is a commutator. Since \(G/N\) is abelian, \([G,G]=N\). Every one-dimensional character is trivial on \(N\); every virtual sum of such characters therefore has the same trace at the identity and at any translation. But (K1) gives trace \(r-1\) at the identity and \(-1\) at a nonidentity translation. This proves the representation obstruction, extending the accepted cubic example CX-R02 to the displayed family.

For completeness, enlargement of the finite quotient cannot remove it. Suppose a finite collection of rational Dirichlet characters and integer exponents gave the same unramified Euler factors as \(L(s,\rho)\). Form a finite Galois extension containing \(E\) and all their cyclotomic fields. At every prime outside a finite set, equality of the coefficient of \(T\) in the Euler factors gives equality of the corresponding virtual traces. Chebotarev gives this equality on every conjugacy class of the enlarged Galois group. The induced map on commutator subgroups surjects onto \([G,G]=N\), whereas each one-dimensional character is trivial on the enlarged commutator. Lifting a nonidentity translation contradicts its trace \(-1\ne r-1\). The same argument applies to rational exponents after clearing a common denominator. ∎

\textbf{Exact source boundary.} \nolinkurl{CharacterBoxes/ExistsPageN.lean}, definition \nolinkurl{IsChiField}, line 19, explicitly requires a \(\mathbf Q\)-algebra embedding into \nolinkurl{CyclotomicField N ℚ}. It also uses a fixed field determined by rational congruences modulo \(2^{e+2}\), a primitive rational character \(\psi\), and two Jacobi symbols. \nolinkurl{Formalism/SharpRestriction/RamificationFrobenius/TotientGaloisClasses/Analytic.lean}, \nolinkurl{HostRaw.dir}, line 102, explicitly constructs a \nolinkurl{DirichletCharacter ℂ H.condN} from the rational residue map \nolinkurl{H.art}. This is the literal specialization excluded here.

\textbf{Scope.} The conclusion is not that RH for this L-function is logically independent of Dirichlet GRH. It rules out the stated finite rank-one Artin factorization and these literal host substitutions. In particular, Brauer induction permits characters over other base fields; it does not turn those bases into \(\mathbf Q\).

\subsection{4. The cubic case with its actual conductor and completed function}\label{the-cubic-case-with-its-actual-conductor-and-completed-function}

\subsubsection{AM-K003. Cubic arithmetic data and the target Euler coefficients}\label{am-k003.-cubic-arithmetic-data-and-the-target-euler-coefficients}

Let \(r=3\), \(K=\mathbf Q(\sqrt{-3})\), \(\lambda=1-\zeta_3\), and assume \(3\nmid p-1\). Thus \(p=3\) or \(p\equiv2\pmod3\). Let \(\mathfrak f\) be the primitive conductor of \(\chi_{p,3}\).

\[
\begin{array}{c|c|c|c}
\text{case}&D_F&\mathfrak f&N\mathfrak f\\ \hline
p\ne3,\ p\equiv8\pmod9&-3p^2&(p)&p^2\\
p\ne3,\ p\equiv2,5\pmod9&-27p^2&(p)\lambda^2&9p^2\\
p=3&-243&\lambda^4&81
\end{array}
\tag{K4}
\]

The analytic conductor factor is \(|D_K|N\mathfrak f=|D_F|\). With \(\Gamma_{\mathbf C}(s)=2(2\pi)^{-s}\Gamma(s)\),

\[
\Lambda(s)=|D_F|^{s/2}\Gamma_{\mathbf C}(s)L(s),
\qquad \Lambda(s)=\Lambda(1-s).
\tag{K5}
\]

In particular the root number is \(+1\); the functional equation does not force a central zero.

\textbf{Justification and specialization.} Harron (2016), §2.1, Lemma 2.1, p.~5, gives the integral bases and discriminants of pure cubic fields. Its Type II condition specializes to \(p\equiv\pm1\pmod9\), and under \(p\equiv2\pmod3\) this is \(p\equiv8\pmod9\). Otherwise the discriminant is \(-27p^2\), including \(p=3\).

The Artin conductor of an induced character is the discriminant of the base times its relative conductor norm; in the present quadratic induction case see, for example, the formula in the proof of Hida (2025), Corollary 6.18, p.~45. One can verify the normalization directly here without any local convention ambiguity. The quotient of the standard completed Dedekind functions \(\Lambda_F/\Lambda_{\mathbf Q}\) is \(|D_F|^{s/2}\Gamma_{\mathbf C}(s)L(s)\), because \(F\) has signature \((1,1)\). It satisfies the functional equation with sign \(+1\). The standard primitive Hecke completion is \((3N\mathfrak f)^{s/2}\Gamma_{\mathbf C}(s)L(s)\). Its dual has the same L-function by AM-K001. Dividing its functional equation by that of the Dedekind quotient gives \(C^{s/2}=W C^{(1-s)/2}\), where \(C=3N\mathfrak f/|D_F|>0\). On an open set where the functions are nonzero this identity forces \(C=1\) and \(W=1\), and then holds by continuation. This proves the asserted conductor normalization.

Only primes above \(3p\) can ramify, from the polynomial discriminant. If \(p\ne3\), \(p\equiv2\pmod3\) is inert in \(K\), and \(T^3-p\) is Eisenstein at \((p)\). The resulting local degree-three extension is tame, so every nontrivial local character has conductor exponent one. The remaining conductor exponent at \(\lambda\), whose norm is \(3\), is determined by the displayed norm identity. This gives (K4), including \(\lambda^4\) when \(p=3\). ∎

At a rational prime \(\ell\nmid3p\), use the factorization of \(T^3-p\) modulo \(\ell\). With \(T=\ell^{-s}\), the reciprocal Euler polynomials of \(L(s)\) are

\[
\begin{array}{c|c|c}
\text{factorization}&\det(1-\rho(\mathrm{Frob}_\ell)T)&
\operatorname{tr}(\rho(\mathrm{Frob}_\ell)^m)\\ \hline
1+1+1&(1-T)^2&2\\
1+2&1-T^2&0\ (m\text{ odd}),\ 2\ (m\text{ even})\\
3&1+T+T^2&-1\ (3\nmid m),\ 2\ (3\mid m)
\end{array}
\tag{K6}
\]

These formulas follow directly from the eigenvalues \(1,1\), \(1,-1\), and \(\zeta_3,\zeta_3^{-1}\) of the augmentation representation. Thus the rational prime-power coefficient in \(-L'/L\) is the last column times \(\log\ell\). This is the concrete coefficient system that a new geometric family must produce. It is not a periodic rational Dirichlet-character coefficient system.

The same completion argument in arbitrary odd prime degree gives

\[
\Lambda_{p,r}(s)=|D_F|^{s/2}\Gamma_{\mathbf C}(s)^{(r-1)/2}L(s),
\qquad\Lambda_{p,r}(s)=\Lambda_{p,r}(1-s),
\tag{K7}
\]

using the signature \((1,(r-1)/2)\) of \(F\) and (K2). This statement needs no explicit general formula for \(D_F\).

\subsection{5. An explicit geometric realization in the cubic case}\label{an-explicit-geometric-realization-in-the-cubic-case}

\subsubsection{AM-K004. The Kummer character is a Hom line between two CM elliptic motives}\label{am-k004.-the-kummer-character-is-a-hom-line-between-two-cm-elliptic-motives}

\textbf{Statement.} Over \(K=\mathbf Q(\zeta_3)\), take

\[
A_0:y^2=x^3+1,\qquad A_p:y^2=x^3+p^2.
\]

For an auxiliary prime \(\ell\), use a coefficient field \(C_\ell\subset\overline{\mathbf Q}_\ell\) containing \(\zeta_3\). Let \(V_0,V_p\) be the \(\zeta_3\)-eigenlines of the endomorphism \([\zeta_3]:(x,y)\mapsto(\zeta_3x,y)\) acting on \(H^1_{\mathrm{\acute et}}\). Then, with this eigenline convention,

\[
V_p\simeq V_0\otimes\chi_{p,3},\qquad
\operatorname{Hom}_{C_\ell}(V_0,V_p)\simeq\chi_{p,3}.
\tag{K8}
\]

This is realized by algebraic correspondences: the Artin line is a direct factor of \(h^1(A_p)\otimes h^1(A_0)^\vee\), with coefficients containing \(K\). It has weight zero. The realization does not provide Alpöge's required quantitative theta-lift family or the window estimate of §6.

\textbf{Proof.} Put \(\alpha=p^{1/3}\). Over \(E=K(\alpha)\) the map

\[
g:A_p\longrightarrow A_0,\qquad (x,y)\longmapsto(\alpha^{-2}x,p^{-1}y)
\tag{K9}
\]

is an isomorphism, since \(\alpha^6=p^2\). For \(\sigma\in G_K\), the ratio \(\sigma(\alpha)/\alpha=\chi(\sigma)\) lies in \(\mu_3\), and \(\chi^{-2}=\chi\). Therefore

\[
\sigma(g)=[\chi(\sigma)]\circ g.
\tag{K10}
\]

On cohomology, functoriality gives

\[
\rho_p(\sigma)g^*\rho_0(\sigma)^{-1}
=(\sigma g)^*=g^*[\chi(\sigma)]^*.
\]

The two \([\zeta_3]^*\)-eigenvalues on \(H^1\) are \(\zeta_3,\zeta_3^{-1}\). This can be seen on the differential \(dx/y\), where the eigenvalue is \(\zeta_3\), and its complex-conjugate Hodge summand; the characteristic polynomial is unchanged in étale realization. The projector

\[
e_{\zeta_3}=\frac13\sum_{j=0}^2\zeta_3^{-j}[\zeta_3]^{j*}
\tag{K11}
\]

is an algebraic-correspondence projector over the coefficient field. Restricting the displayed functoriality identity to its image yields

\[
\rho_p(\sigma)g^*=\chi(\sigma)g^*\rho_0(\sigma),
\]

which is precisely (K8). The line in the Hom motive becomes the trivial line over \(E\) with descent cocycle \(\chi\); this identifies it with the Artin motive defined by that character, not merely with an unrelated equal-dimensional representation. As a tensor of weight-one cohomology with its dual, it has weight zero. ∎

\textbf{What this does and does not change.} Both elliptic curves have CM by \(K\), are isomorphic over \(\overline{\mathbf Q}\), and have the same geometric CM type and stable Faltings height. At a good prime their formal groups become isomorphic after extension. At good rational primes inert in \(K\), these formal groups have height two, so the existing height-two derivative mechanism applies after suitable base extension; DP6 supplies its known exponent. Split good primes have height one and give no extra height-two contribution. A finite-order descent twist is actual arithmetic information, but it is not a new geometric height budget.

The invariant differentials in these equations satisfy

\[
g^*(dx/y)_{A_0}=\alpha\,(dx/y)_{A_p}.
\tag{K12}
\]

Indeed \(\alpha^{-2}/p^{-1}=p\alpha^{-2}=\alpha\). For transported complex cycles, \(\Omega_0=\alpha\Omega_p\). These algebraic factors must be retained in the integral lattices and period normalization. Removing them by consistent normalization cannot itself supply a favorable extra lift cost. No equality of \nolinkurl{MainX} kernel, covolume, mass, or terminal-period records has been inferred from this isomorphism.

This construction is also why the non-CM statement in AM-K002 must not be overstated as ``the character cannot occur in CM motives.'' It does occur in this explicit tensor category over \(K\).

\subsection{6. A completed analytic transfer and the surviving arithmetic gate}\label{a-completed-analytic-transfer-and-the-surviving-arithmetic-gate}

\subsubsection{AM-K005. Field-independent Landau reduction for the Kummer L-function}\label{am-k005.-field-independent-landau-reduction-for-the-kummer-l-function}

Define \(\Lambda_K(\mathfrak a)=\log N\mathfrak q\) when \(\mathfrak a=\mathfrak q^m\) for a prime ideal \(\mathfrak q\) and \(m\ge1\), and zero otherwise. Extend the primitive character by zero at ideals not coprime to its conductor. Put \(c_1=L'(1)/L(1)\) and, for \(x\ge1\),

\[
G_\chi(x)=\sum_{N\mathfrak a\le x}\Lambda_K(\mathfrak a)\chi(\mathfrak a)
\left(\frac{x}{N\mathfrak a}-1\right)+xc_1.
\tag{K13}
\]

\textbf{Theorem.} This is real-valued. Fix \(0<\theta<1\) and \(\delta>0\). Suppose the following single window condition holds:

\begin{quote}
\textbf{\(\mathcal W_\chi(\theta,\delta)\).} For every \(\eta>0\) and every \(Y_0\ge1\), there exist \(Y\ge Y_0\) and a measurable \(R:[Y,Y^{1+\delta}]\to\mathbf R\) such that \(0\le R(x)\), \(G_\chi(x)\le R(x)\), \(R(x)x^{-\theta-1}\) is integrable, and \[
\int_Y^{Y^{1+\delta}}R(x)x^{-\theta-1}\,dx\le Y^\eta.
\tag{K14}
\]
\end{quote}

Then \(L(s)\) has no zero with \(\Re s>\theta\). Consequently it suffices to prove \(\mathcal W_\chi(1/2+\epsilon,1/2)\) for every \(0<\epsilon<1/2\) to prove RH for \(L\), and hence the CX-ALG-024 consequences under accepted Dirichlet GRH.

\textbf{Proof.} We verify the actual \nolinkurl{Deduction.PoleData} interface and invoke the accepted source theorem \nolinkurl{Deduction.landau_window}. The analytic continuation, functional equation, Euler product for \(\Re s>1\), and nonvanishing \(L(1)\ne0\) of a nontrivial finite-order Hecke character are classical. In the Artin description (K2), the relevant holomorphy is the classical monomial case, obtained by its identification with this nontrivial Hecke L-function; it is not inferred from an unproved cancellation in a quotient of zeta functions. The standard nonvanishing at \(1\) is Milne (2020a), V Theorem 2.4(b), p.~156; its Euler product and continuation conventions are V §2 and VIII §10.

By AM-K001, the coefficients of \(L\), and therefore those of \(-L'/L\), are real. In particular \(c_1\in\mathbf R\). Grouping the ideal terms by rational norm writes

\[
-\frac{L'}L(s)=\sum_{n\ge1}\frac{b(n)}{n^s},\qquad
G_\chi(x)=\sum_{n\le x}b(n)(x/n-1)+xc_1,
\tag{K15}
\]

with \(b(n)\in\mathbf R\). At a rational prime \(\ell\), the local factor is the determinant of \(1-\rho(\mathrm{Frob}_\ell)\ell^{-s}\) on inertia invariants. All its nonzero eigenvalues are roots of unity, and there are at most \(d=r-1\) of them. Thus

\[
|b(\ell^m)|\le d\log\ell,
\qquad b(n)=0\quad\text{if }n\text{ is not a prime power}.
\tag{K16}
\]

This bound includes ramified primes and proves the absolute convergence needed below. It also gives the elementary growth estimate, for \(x\ge1\),

\[
|G_\chi(x)|\le dx\sum_{2\le n\le x}\frac{\log n}{n}+|c_1|x
\ll x(1+\log x)^2\ll x^2.
\tag{K17}
\]

At a threshold \(x=n\) the newly appearing term is zero. Hence \(G_\chi\) is continuous and piecewise affine. Set \(g(u)=G_\chi(e^u)\) for \(u\ge0\), and \(g(u)=0\) for \(u<0\). It is measurable and \(|g(u)|\le C e^{2u}\) for \(u\ge0\).

For \(\Re s>1\), direct integration gives

\[
\int_n^\infty (x/n-1)x^{-s-1}\,dx=\frac{n^{-s}}{s(s-1)}.
\]

The absolute integral of the summands is bounded by the corresponding expression with \(|b(n)|\) and \(\Re s>1\), which is finite by (K16). Termwise integration is therefore legitimate and yields

\[
\Phi(s):=\int_1^\infty G_\chi(x)x^{-s-1}\,dx
=-\frac{L'(s)}{s(s-1)L(s)}+\frac{c_1}{s-1}
=\frac{-L'/L(s)+s c_1}{s(s-1)}.
\tag{K18}
\]

At \(s=1\) the numerator vanishes. Define \(\Phi(1)\) by the removable continuation; it equals \(c_1-(L'/L)'(1)\). At zeros of \(L\), regard \(\Phi\) as meromorphic and assign arbitrary values if a total function is required by the formal interface. At points with positive real part outside those zeros it is holomorphic, including \(1\). Under \(x=e^u\), (K18) is the Laplace identity

\[
\Phi(s)=\int_0^\infty g(u)e^{-su}\,du,
\]

in particular on the source's required half-plane \(\Re s>2\).

Let \(Z=\{\rho:\Re\rho>0,\ L(\rho)=0\}\). The function \(L\) is entire, not identically zero, and \(L(1)\ne0\). Its zeros are isolated. If \(\rho\in Z\) has multiplicity \(m\), write \(L(s)=(s-\rho)^m a(s)\), with \(a\) nonvanishing near \(\rho\). Then

\[
\Phi(s)=\frac{m}{\rho(1-\rho)}\frac1{s-\rho}+h_\rho(s),
\tag{K19}
\]

where \(h_\rho\) is holomorphic near \(\rho\). The residue is nonzero since \(\rho\ne0,1\). If \(\rho\) is real, the Euler product excludes \(\rho>1\), and \(L(1)\ne0\), so \(0<\rho<1\). The residue \(m/[\rho(1-\rho)]\) is then strictly positive and real. These are exactly the positivity and simple-pole conditions of \nolinkurl{PoleData}; the pole is simple even when the zero of \(L\) has higher multiplicity.

All fields of the source record have now been supplied: measurability, exponential growth, Laplace representation, positive real parts of \(Z\), holomorphy away from \(Z\), isolation, and the required nonzero residues with positive sign at real zeros. No Hecke-family extension of its Dirichlet-specialized \nolinkurl{CharFamily} is needed for this application.

It remains to translate (K14). Given the \(U_0\in\mathbf R\) required by \nolinkurl{WindowHypU}, choose \(Y_0=\max(1,e^{U_0})\). Put \(U=\log Y\) and \(\widetilde R(u)=R(e^u)\). Then \(U\ge U_0\), \(\widetilde R\ge0\), \(g\le\widetilde R\), and the change of variables \(x=e^u\) gives

\[
\int_U^{(1+\delta)U}\widetilde R(u)e^{-\theta u}\,du
=\int_Y^{Y^{1+\delta}}R(x)x^{-\theta-1}\,dx
\le e^{\eta U}.
\]

This is precisely \nolinkurl{WindowHypU g θ δ}. The accepted theorem \nolinkurl{landau_window} gives \(\Re\rho\le\theta\) for every \(\rho\in Z\).

If the condition is known for every \(\theta=1/2+\epsilon\), then no zero has real part \(>1/2\). The self-dual functional equation (K7) maps a nontrivial zero of real part \(<1/2\) to one of real part \(>1/2\), while the gamma factor is nonzero and holomorphic in the critical strip. This proves RH under the asserted window conditions. ∎

\textbf{Actual Lean locators.} With common prefix \nolinkurl{Development/AutomorphicArithmetic/DedekindLFunctions/}:

\begin{longtable}[]{@{}
  >{\raggedright\arraybackslash}p{(\columnwidth - 4\tabcolsep) * \real{0.3333}}
  >{\raggedright\arraybackslash}p{(\columnwidth - 4\tabcolsep) * \real{0.3333}}
  >{\raggedright\arraybackslash}p{(\columnwidth - 4\tabcolsep) * \real{0.3333}}@{}}
\toprule\noalign{}
\begin{minipage}[b]{\linewidth}\raggedright
Record/theorem
\end{minipage} & \begin{minipage}[b]{\linewidth}\raggedright
File and locator
\end{minipage} & \begin{minipage}[b]{\linewidth}\raggedright
Exact transfer use
\end{minipage} \\
\midrule\noalign{}
\endhead
\bottomrule\noalign{}
\endlastfoot
\nolinkurl{Deduction.PoleData} & \nolinkurl{LandauZeroFreeCriterion.lean:15} & Arbitrary \(g:\mathbf R\to\mathbf R\), \(\Phi:\mathbf C\to\mathbf C\), \(Z\subset\mathbf C\); no character or base-field parameter. \\
\nolinkurl{Deduction.WindowHypU} & \nolinkurl{LandauZeroFreeCriterion.lean:25} & Quantifiers and weighted window integral translated in (K14). \\
\nolinkurl{Deduction.landau_window} & \nolinkurl{LandauTheorems/landauWindow.lean:83} & The accepted analytic implication used above. \\
\nolinkurl{poleData_abstract} & \nolinkurl{LandauTheorems/poleData.lean:22} & Confirms a generic analytic layer precedes the Dirichlet-family specialization at line 248. The proof above directly verifies the record instead of claiming a new Lean instantiation. \\
ramp Mellin kernel & \nolinkurl{LandauTheorems/MellinRamp.lean:10}; \nolinkurl{mellinGchiOfNeOne.lean:65,116} & Generic scalar kernel and arithmetic-sequence integration before Dirichlet specialization. \\
\end{longtable}

\textbf{A sufficient stronger estimate.} If, for every \(\epsilon>0\), \(G_\chi(x)\le x^{1/2+\epsilon}\) for all sufficiently large \(x\), take \(R(x)=x^{1/2+\epsilon}\). The integral in (K14) is \(\delta\log Y\), which is at most \(Y^\eta\) for all sufficiently large \(Y\). Thus a one-sided pointwise bound suffices; absolute-value control is unnecessary for this deduction. No such bound is proved here.

\textbf{Current gate, Q3-KUMMER-WINDOW.} Construct a relative character host and an actual arithmetic theta-lift family whose period, divisibility, covolume, integral-structure, and cost equations prove (K14), for each \(\epsilon>0\), with \(\theta=1/2+\epsilon\), already for the explicit cubic data (K4), (K8), (K13). This is an explicit analytic obligation left by the arithmetic construction. The source's generic Landau theorem eliminates a separate analytic-generalization obligation.

\subsection{7. Steps 2--20: actual interface audit and replacement obligations}\label{steps-220-actual-interface-audit-and-replacement-obligations}

The short titles in the following table identify the historical timeline. The right-hand column distinguishes a completed replacement from a required new instance. A source anchor is a declaration whose scope was actually checked, or a clearly labelled carried source interface. It is not a claim to have exhaustively inspected every proof body contributing to that step.

\begin{longtable}[]{@{}
  >{\raggedright\arraybackslash}p{(\columnwidth - 4\tabcolsep) * \real{0.3333}}
  >{\raggedright\arraybackslash}p{(\columnwidth - 4\tabcolsep) * \real{0.3333}}
  >{\raggedright\arraybackslash}p{(\columnwidth - 4\tabcolsep) * \real{0.3333}}@{}}
\toprule\noalign{}
\begin{minipage}[b]{\linewidth}\raggedright
Step
\end{minipage} & \begin{minipage}[b]{\linewidth}\raggedright
Q-sensitive interface and replacement over \(K\)
\end{minipage} & \begin{minipage}[b]{\linewidth}\raggedright
State and exact anchor
\end{minipage} \\
\midrule\noalign{}
\endhead
\bottomrule\noalign{}
\endlastfoot
2: local derivations & The local divided-power/height statement is already local arithmetic geometry. For \(r=3\), the CM elliptic curves of AM-K004 give explicit good-reduction inputs after extension. All conductor primes remain in the global bookkeeping. & DP6, DP32, DP34 and DP42 are carried; no new exponent is asserted. AM-K004 supplies the twist, not a better divisibility budget. The supplied map §4 labels its correspondence to \nolinkurl{BlockTheory/HeightDefect} and CM-good-curve files as low confidence. \\
3: quadric and biextension & Poincaré bundles and the relevant quadratic constructions can be formed over a number field. What must be newly proved is that the resulting character-weighted arithmetic functional equals the norm-indexed \(G_\chi\) in (K13), with its actual integral lattices. & The generic geometry does not manufacture this equality. Carried source families: \nolinkurl{Formalism/LocalToGlobal/GlobalTheory}; directly checked assembly record \nolinkurl{ThetaIntegralsAgree} and \nolinkurl{LocalMatching} in \nolinkurl{SiegelEisensteinSecondTermHypotheses.lean:44,69}. \\
4: determinant positivity & Positivity of a fixed Hermitian determinant is field-independent linear algebra. Identifying it with the target arithmetic sum is a separate global construction. & No new positivity theorem is needed at this level; no target determinant identity has been proved. The supplied map found no exact one-to-one source module for this historical step. \\
5: smoothed explicit formula & Replace rational integer coefficients \(\Lambda(n)\psi(n)\) by \(\Lambda_K(\mathfrak a)\chi(\mathfrak a)\), and rational conductor normalization by (K4). & Completed analytic replacement in AM-K005. Literal definitions are \nolinkurl{Analytic.lean:9,11,18,23} (\nolinkurl{abelian₂}, \nolinkurl{logDerivativeAtOne}, \nolinkurl{sumLamChi}, \nolinkurl{gChi}). \\
6: zero-free deduction & Apply a base-field-free pole theorem to the real induced L-function; do not substitute a Hecke character into \nolinkurl{PChar}. & Completed conditional deduction AM-K005. \nolinkurl{LandauZeroFreeCriterion.lean:15,25}; \nolinkurl{LandauTheorems/landauWindow.lean:83}. \\
7: periods and weights & The underlying CM/unitary setting is not restricted to \(\mathbf Q\). The new task is to choose the target ray-class character, test vectors and interlacing data so its period is exactly the coefficient functional needed in step 5. & Generic CM Galois row is visible in \nolinkurl{WeightsConjBis.lean:69} (\nolinkurl{Row0b}); it does not state a Hecke RH conclusion. Carried period anchor \nolinkurl{Formalism/PeriodIntegrals/BesselPeriodGGP.lean}. Target period identity open. \\
8: repeated definite lifts & Build an actual chain for the new ray-class family, with the six carried lift identities, controlled nonvanishing and correct rank costs. Keeping \(K\) CM is possible; using non-CM \(E\) as a CM base is excluded by AM-K002. & AM-K004 gives initial geometric descent data only. Direct assembly anchor \nolinkurl{nonempty_hyps_definite_side_of_closed_bound}, \nolinkurl{SiegelEisensteinSecondTermHypotheses.lean:44}, requires linked norm bounds and the other records. \\
9: norms and masses & Replace level/discriminant and local indices by those of the relative conductor, and include the unit and archimedean normalizations over \(K\). The cubic conductor is explicit in (K4). & Generic mass identities may be reused with hypotheses, but their evaluation for the new chain is open. Carried \nolinkurl{Formalism/PeriodIntegrals/UnitaryGroupMassFormula}; \nolinkurl{CovolumeBoundOfX} in the checked assembly theorem. \\
10: global divisibility & Local DP bounds remain available; the global sum must be normalized by \(\log N\mathfrak q\), residue degrees, bad-prime lattices and all embeddings. & No equality of the new global gain to the original gain is proved. Carried divisibility interfaces \nolinkurl{Formalism/SharpRestriction/Divisibility} and \nolinkurl{Formalism/LocalToGlobal/Divisibility}. \\
11: auxiliary tower & A relative ray-class tower over \(K\) must replace the rational cyclotomic fixed-field construction. The Kummer layer need not be CM as an absolute number field. & Literal substitution fails by AM-K002. \nolinkurl{ExistsPageN.lean:19,27,39} (\nolinkurl{IsChiField}, \nolinkurl{WinPair}, \nolinkurl{PageN}) explicitly require rational congruence and cyclotomic data. New tower with the necessary character separation and size bounds is open. \\
12: seesaw and terminal period & Seesaw/Rallis formalism over CM \(K\) is available at the generic level. The target family must satisfy the specified character independence and terminal-period normalization. & The explicit twist identity is not that terminal period identity. Actual checked bridge: \nolinkurl{ThetaIntegralsAgree}, \nolinkurl{LocalMatching}, \nolinkurl{SiegelEisensteinSecondTermHypotheses.lean:44,69}. Carried \nolinkurl{ThetaCorrespondence/RallisInnerProduct}. \\
13: slope inequality & The Arakelov slope inequality is already a number-field statement. Supply the actual metrized lattices, normalize degrees by the field degree, and include residue degrees in finite contributions. & New lattice/cost instance open. Carried source \nolinkurl{Formalism/LatticeSlopes/LatticeSlopesProofs/SlopeInequality.lean}; source assembly still needs \nolinkurl{CovolumeBoundOfX}. \\
14: cost evaluation & Compute every local and global term for the new lifts, using the actual conductor and period scaling (K12). The algebraic twist alone does not determine those terms or their cancellation. & No cost improvement or RH follows from the same geometric CM type. The supplied step-to-source map identifies \nolinkurl{UnitaryGroupMassFormula/UnipotentCost} only with low confidence; it is not treated as a complete formula for the target. \\
15: common integral structure & Replace rational primitive-character sets by finite-order ray-class characters over \(K\), and prove common coefficient-ring, saturation and congruence properties for their sections. & New host adapter required. \nolinkurl{Analytic.lean:102,106,145,149,154} (\nolinkurl{HostRaw.dir}, \nolinkurl{IsPrimOf} and analytic bounds) all quantify rational Dirichlet characters. AM-K002 excludes the naive norm-pullback route. \\
16: trace/orthogonality bound & The relative trace identity needs an actual ideal/ray-class indexing, stabilizers, unit factors and quantitative orbital count for the new host. & A formal orthogonality equality alone does not supply the needed scale-uniform bound. Carried \nolinkurl{Formalism/PeriodIntegrals/OrbitalBoundsConclusion}; no target version is proved. \\
17: cohomology control & The weighted Galois-cohomology size predicate \nolinkurl{VVle E t B} is generic in a number field \(E\). Its available supply theorem is indexed by rational \nolinkurl{PChar} and \nolinkurl{IsChiField}. & \nolinkurl{ExistsPageN.lean:12,41,44,49}; \nolinkurl{SupplyChi.lean:8}. A relative version over the new tower, with the same quantitative allowance, remains open. It is not supplied just by Dirichlet GRH. \\
18: covolumes and indices & Compute the actual determinant and finite-index lattice for the new global family; the conductor table fixes only part of its arithmetic input. & \nolinkurl{CovolumeBoundOfX} is a required hypothesis at \nolinkurl{SiegelEisensteinSecondTermHypotheses.lean:44,69}. No \nolinkurl{MainX} instance or covolume equality for \(\chi\) has been constructed. \\
19: combinatorial tower choice & Any abstract finite-set combinatorics must be realized with the new ray-class/tower parameters and the necessary cost and independence properties. & The supplied map §4 explicitly marks the use of \nolinkurl{Formalism/BehrendSets} in the final chain as uncertain. No extra Q-specific theorem is invented from that historical label. The precise checked tower restrictions remain \nolinkurl{IsChiField} and \nolinkurl{WinPair}. \\
20: final theorem & The literal endpoint is \nolinkurl{EveryCharacter}, a statement about \nolinkurl{DirichletCharacter ℂ N}. The replacement endpoint is the conclusion of AM-K005 once the target window is proved. & \nolinkurl{CharacterBoxes/EveryCharacter.lean:13}; \nolinkurl{ThirdPlaceSeries/TopPlug.lean:12,16}. The new arithmetic window remains unproved, so the requested Hecke RH conclusion is not established. \\
\end{longtable}

Two additional interface observations make the gate concrete. First, \nolinkurl{CharacterBoxes/Supply.lean:10} defines a generic request type \nolinkurl{FReq}, but its actual target request at line 42 is

\[
\texttt{GChi ψ}:
\quad\texttt{EAtH (Level.cyc ψ) H.toHostRaw TT s.x s.eps 4}
\ \land\ \texttt{VVle H.E t (sqrt s.b)}.
\]

The ambient request formalism is not the obstruction; the specific character/host supply is. Second, \nolinkurl{CharacterBoxes/ZeroFreeAllOfTopG.lean:12,26} separates the generic design/top-level family machinery from its application to \nolinkurl{PChar}. A relative version must produce a new target request and prove its supply, then connect it to (K13). Merely adding a new character type without those proofs would relocate the gap.

\subsection{8. Grothendieck-language formulation}\label{grothendieck-language-formulation}

\subsubsection{AM-K006. The finite étale and motivic dictionaries preserve the gate}\label{am-k006.-the-finite-uxe9tale-and-motivic-dictionaries-preserve-the-gate}

Let \(U=\operatorname{Spec}\mathbf Z[1/(rp)]\), \(U_K=\operatorname{Spec}\mathcal O_K[1/(rp)]\). The normalization \(U_E\to U_K\) is a finite étale \(\mu_r\)-torsor. Indeed \(p\) and \(r\) are units there; the derivative \(rT^{r-1}\) is invertible on \(T^r=p\), and \(K\) already contains \(\mu_r\). The character \(\chi\) is its Kummer local system. For coefficients containing \(\mu_r\), the idempotent

\[
e_\chi=\frac1r\sum_{g\in C_r}\chi(g)^{-1}g
\]

in the finite-cover correspondence algebra projects its rank-one Artin summand (with the chosen action convention). Every geometric Frobenius eigenvalue on that summand is a root of unity. Thus it is pointwise pure of weight zero. This is a statement about each closed fibre, not about zeros of the global L-function.

For the finite étale cover \(U_F\to U\), the permutation local system satisfies

\[
[\rho]=[h^0(U_F/U)]-[\mathbf1]
\]

in the Grothendieck group of finite-monodromy local systems, equivalently in the representation ring of \(G\). Its characteristic-polynomial identity is

\[
\det(1-gT\mid\rho)=\frac{\det(1-gT\mid P)}{1-T}.
\]

Exterior-power operations \(\lambda^j\) encode these determinants. AM-K001 and the carried CX-ALG-024 are therefore identities in the Artin-motivic \(K_0\) / representation-ring setting. AM-K002 says that this particular class is outside the subring generated by rational abelian Artin lines: every virtual object in that subring has trivial commutator action, whereas \(\rho\) does not. Adding tensor products, duals, exterior powers, symmetric powers or characteristic-zero direct summands of those abelian lines does not change that action. This is the elementary finite-group content of the carried GRH\_NIGHT AG-LAMBDA discussion, not a new global purity theorem.

For \(r=3\), AM-K004 gives the more useful relative realization over \(K\): the finite-monodromy line is the Hom of two CM elliptic eigensummands. The elliptic curves and their Picard duals extend as abelian schemes after deleting their bad primes. Their Poincaré biextensions give the usual tensor/duality pairings. The finite descent cocycle changes the Artin factor while leaving the geometric CM type and local height-two formal-group isomorphism unchanged after extension. The arithmetic integral lattice and descent factors remain essential for the cost.

\textbf{Exact framework locators.}

\begin{longtable}[]{@{}
  >{\raggedright\arraybackslash}p{(\columnwidth - 4\tabcolsep) * \real{0.3333}}
  >{\raggedright\arraybackslash}p{(\columnwidth - 4\tabcolsep) * \real{0.3333}}
  >{\raggedright\arraybackslash}p{(\columnwidth - 4\tabcolsep) * \real{0.3333}}@{}}
\toprule\noalign{}
\begin{minipage}[b]{\linewidth}\raggedright
Language
\end{minipage} & \begin{minipage}[b]{\linewidth}\raggedright
Locator and status
\end{minipage} & \begin{minipage}[b]{\linewidth}\raggedright
Meaning here
\end{minipage} \\
\midrule\noalign{}
\endhead
\bottomrule\noalign{}
\endlastfoot
Finite étale covers and finite G-sets & SGA 1, Exposé V, Théorème 4.1; primary reedition \url{https://arxiv.org/pdf/math/0206203} & The covering/local-system/representation dictionary used above. \\
Kummer torsors & SGA 1, Exposé XI §6, Proposition 6.1, Définition 6.2, Corollaires 6.3 and 6.5; reedition PDF pp.~248--250, original marginal pp.~303--305 & \(T^r=p\) on the base with \(rp\) inverted and its \(\mu_r\)-character. \\
\(\lambda\)-rings and \(K_0\) & SGA 6, Exposé V §§2 and 7, especially (7.1.1); carried GRH\_NIGHT AG-LAMBDA & Exterior-power determinant operations on Artin classes; no analytic RH follows from the class identity. \\
CM motives & Milne (2020b), §5, p.~45, for the framework; explicit new construction (K8)--(K11) here & The Artin Hom line belongs to a relative CM tensor category despite the non-CM splitting field. \\
Formal groups and crystals & Carried DP\_MULT DP10, DP6, DP32/DP34; Grothendieck cote 54 scans 5--7; cote 4 scans 5--6, 8.4.2 and Lemma 8.4.3, and scans 7--10 for torsion corrections & The local crystalline/Cartier-duality interpretation of the formal groups. No new monodromy result for cote 4 scan 12, §8.6, is asserted. The boxed zero restricted power in cote 54 concerns the vector kernel, as DP10 already established. \\
Picard schemes & FGA, Bourbaki exposé 232, Théorème 3.1; exposé 236, Corollaires 1.7 and 2.3, Proposition 2.10(ii), carried through DP\_MULT & Picard/dual abelian schemes for the elliptic curves; this gives geometric objects, not a \nolinkurl{MainX} arithmetic identity. \\
Biextensions & SGA 7 I, Exposés VII--VIII; carried DP\_MULT/THETA3 framework & Poincaré biextension and Cartier duality organize the pairings; no finer unverified proposition number is fabricated. \\
Weights versus global zeros & Carried CX\_GRH CX-G06 & Weight zero of a finite-monodromy sheaf does not imply RH over a number field. The missing input remains (K14). \\
\end{longtable}

The archive cote locators above are carried from the supplied DP\_MULT record, not newly reread scans in this child contribution. Their exact statements are not enlarged.

\subsection{9. Tried approaches and the remaining work}\label{tried-approaches-and-the-remaining-work}

\begin{longtable}[]{@{}
  >{\raggedright\arraybackslash}p{(\columnwidth - 4\tabcolsep) * \real{0.3333}}
  >{\raggedright\arraybackslash}p{(\columnwidth - 4\tabcolsep) * \real{0.3333}}
  >{\raggedright\arraybackslash}p{(\columnwidth - 4\tabcolsep) * \real{0.3333}}@{}}
\toprule\noalign{}
\begin{minipage}[b]{\linewidth}\raggedright
Approach
\end{minipage} & \begin{minipage}[b]{\linewidth}\raggedright
Result reached
\end{minipage} & \begin{minipage}[b]{\linewidth}\raggedright
Why it does not finish RH
\end{minipage} \\
\midrule\noalign{}
\endhead
\bottomrule\noalign{}
\endlastfoot
Instantiate the final theorem with \(\chi\) & Actual endpoint type read & \nolinkurl{EveryCharacter} quantifies rational Dirichlet characters. \\
Pull back a Dirichlet character by the norm & Refuted, AM-K002(2) & The target is anti-invariant under complex conjugation and has odd order. \\
Express the Artin L-function by finitely many rational Dirichlet factors & Refuted in the Euler-factor/Artin sense, AM-K002(3) & Nontrivial commutator trace. This does not refute every conceivable deduction from A\_D. \\
Insert \(E\) as the existing cyclotomic/CM host & Refuted, AM-K002(1) & \(E\) cannot embed in a CM field, whereas \nolinkurl{IsChiField} requires a rational cyclotomic embedding. \\
Realize the cubic character geometrically & Completed explicit Hom-line realization, AM-K004 & The integral lattice, period functional, nonvanishing chain, and cost/window identities are not consequences of the geometric twist. \\
Generalize the analytic Landau stage & Completed reduction, AM-K005 & All source \nolinkurl{PoleData} hypotheses verified; the arithmetic window (K14) remains. \\
Use known generic CM automorphic source rows & Type checked & \nolinkurl{Row0b} is already CM-general, but the proved outer endpoint and host supply remain rational-character-specific. \\
Use a newer cubic-Hecke zero-free source & Consulted only & OpenAI (2026), Theorem 1.1, claims \(\Re s>7/8\). Its long proof was not audited here and is not an accepted premise stipulated by the user. The downloaded \nolinkurl{ComparatorChallenges/HeckeSevenEighths.lean} is a comparator statement containing \nolinkurl{sorry}, not a proof certificate. Even the stated half-plane would not prove RH. \\
Use Masullo's claimed all-Hecke theorem & Conditional only & No Masullo premise is made unconditional. The carried GRH\_NIGHT gate Q-MAS-GEOMETRIC-NARROW remains unchanged. \\
\end{longtable}

\textbf{Masullo locators as supplied, without a new draft audit.} \nolinkurl{INTAKE_CX_GRH.md} records the all-Hecke claim in §§1.2 and 1.2.3, PDF pp.~12--16; Theorem 13.3.1.2 on PDF p.~522 is blank in that draft. The same intake records the relevant gaps at §9.3.2, Lemmas 9.3.2.4--9.3.2.5 and Theorem 9.3.2.6, pp.~425--426; §11.3, p.~495; and §§12.3--12.4. \nolinkurl{INTAKE_GRH_NIGHT_part3.md} carries the stronger narrowed gate. Those records are cited with their existing check status; no new proof of those draft claims is supplied here.

\textbf{Next concrete attack.} Build the relative host for the explicit cubic conductor in (K4), with \(\chi\) represented by (K8), and prove that its arithmetic functional is exactly (K13). Then compute the actual lift costs and common integral structure. Success must yield (K14) for every \(\theta>1/2\); an estimate at a fixed \(\theta>1/2\) is a partial result only. The previously accepted DP64 leading budget and the twist's unchanged geometric height preclude treating a mere relabelling of the same CM data as a newly computed cost saving.

\subsection{10. References and source records}\label{references-and-source-records}

\begin{itemize}
\tightlist
\item
  Alpöge, L. (2026). \emph{Remark}. \url{https://alpo.ge/remark.pdf}. Timeline steps 2--20, pp.~3--4; Lean archive as specified in the prompt. Actual declarations and pinned copies listed above and in \nolinkurl{kummer_sources/manifest.json}.
\item
  Grothendieck, A., et al.~(1971/2003). \emph{Revêtements étales et groupe fondamental (SGA 1)}. Lecture Notes in Mathematics 224; recomposed edition. Exposé V, Théorème 4.1; Exposé XI §6, Proposition 6.1, Définition 6.2, Corollaires 6.3, 6.5. \url{https://arxiv.org/pdf/math/0206203}.
\item
  Harron, R. (2016, May 24). \emph{The shapes of pure cubic fields}. §2.1, Lemma 2.1, p.~5. \url{https://math.hawaii.edu/~rharron/research/Harron-PureCubicShapes.pdf}.
\item
  Hida, H. (2025, March 27). \emph{Non-vanishing of integrals of a mod p modular form}. Proof of Corollary 6.18, p.~45, induced-conductor formula only. \url{https://www.math.ucla.edu/~hida/EllrevF.pdf}. No use of its separate nonvanishing assertions is made.
\item
  Milne, J. S. (2020a). \emph{Class field theory} (v4.03). V §2 and Theorem 2.4(b), p.~156; V Theorem 3.27 and Remark 3.28, pp.~165--166; VIII §10, pp.~261--264. \url{https://www.jmilne.org/math/CourseNotes/CFT.pdf}.
\item
  Milne, J. S. (2020b). \emph{Complex multiplication} (v0.10). Proposition 1.4 and Corollary 1.5, p.~10; §5, p.~45. \url{https://www.jmilne.org/math/CourseNotes/CM.pdf}.
\item
  OpenAI. (2026, September 30). \emph{The quasi-Riemann hypothesis: A zero-free half-plane Re(s)\textgreater7/8}. Theorem 1.1, consulted but not used as an established premise in this report. \url{https://github.com/openai/math/tree/main/preprints/The-Quasi-Riemann-Hypothesis-September-30-2026}. Saved TeX SHA-256: \nolinkurl{42a5ee0febca59fd1def55cfd6c6808c322ef4237d7726303ea52f711deac6a3}.
\item
  Supplied programme sources: \nolinkurl{INTAKE_CX_GRH.md}; \nolinkurl{INTAKE_GRH_NIGHT_part3.md}; \nolinkurl{GRH_NIGHT.tex}, CX-ALG-024 and AG-LAMBDA; \nolinkurl{DP_MULT.tex}, DP6/DP10/DP32/DP34/DP42/DP64; \nolinkurl{ALPOGE_LEAN_MAP_20261120.md}, §§3--4. These are carried under the user's stated check status.
\end{itemize}

\textbf{Status for integration.} AM-K000--AM-K006 are proved in their stated scope, with whole-proof author readings recorded separately. The requested RH theorem remains open at Q3-KUMMER-WINDOW. No unconditional deterministic-factoring conclusion is asserted from this contribution. The \(p=r\) finite-field construction in AM-K000 is independently unconditional.

\textbf{Subsequent concrete progress.} The companion kummer\_cubic\_host.md proves AM-K007--AM-K011: explicit cubic ray-class characters at every allowed prime; finite-Euler-factor invariance of the normalized window; source-compatible splitting transport and its telescoping limit; and the actual CM tower \(\mathbf Q(\zeta_{3\cdot7^n})\) with an invariant coprime half-different, positive unimodular Hermitian lattice, principally polarized products of the CM elliptic twists, and the retained Artin Hom line. Adding the split prime \(13\) supplies every mathematical field of the source's AbelianUnramifiedExtensionConditions, including \(8\mid f\). The addendum proves why a relative tower ramified at inert \(2\) cannot preserve that invariant-half-different condition, and identifies the target's nonsplit and wild ramification that its displayed IsMu0/QDat families do not directly allow. The complete HostRaw/MainX replacement and the arithmetic window remain open.

\clearpage

\section{Cubic Kummer host: finite character, splitting transport, and an explicit CM tower}\label{cubic-kummer-host-finite-character-splitting-transport-and-an-explicit-cm-tower}

\textbf{Run:} ALP\_MACHINE, question 3. \textbf{Author:} GPT-6 Astra Pro; inherited effort not separately exposed; child /root/kummer\_extension. \textbf{Date:} 2026-11-21. \textbf{Companion:} kummer\_extension.md, AM-K000--AM-K006. This addendum proves smaller arithmetic and geometric subgates. It does not assert the Kummer RH window or an instance of the archive's complete HostRaw/MainX records.

The inputs already accepted in the main contribution retain their status. In particular the conductor table AM-K003, the elliptic Hom realization AM-K004, the field-independent analytic adapter AM-K005, and CX-ALG-024 are used without reopening the stipulated Dirichlet theorem. The new field called \(B_n\) below is a CM coefficient/endomorphism field; it is not the non-CM Kummer splitting field \(K(p^{1/3})\).

\subsection{1. AM-K007: an explicit finite ray-class host for every allowed cubic prime}\label{am-k007-an-explicit-finite-ray-class-host-for-every-allowed-cubic-prime}

Put \(K=\mathbf Q(\zeta)\), \(\zeta=\zeta_3\), \(\mathcal O=\mathbf Z[\zeta]\), and \(\lambda=1-\zeta\). The ring \(\mathcal O\) is norm-Euclidean, has class number one, and has units \(\mu_6\). We use these established elementary facts; the source is Murty and Esmonde (2005), §2.3, Exercise 2.3.2 and its solution, together with §9.1. The norm is \[
N(a+b\zeta)=a^2-ab+b^2.
\]

\subsubsection{\texorpdfstring{1.1. The case \(p\ne3\), necessarily \(p\equiv2\pmod3\)}{1.1. The case p\textbackslash ne3, necessarily p\textbackslash equiv2\textbackslash pmod3}}\label{the-case-pne3-necessarily-pequiv2pmod3}

The reduction map \[
\mu_6\longrightarrow(\mathcal O/(3))^\times
\] is an isomorphism. The target has order \(9-3=6\). The only unit congruent to \(1\) modulo \((3)=(\lambda^2)\) is \(1\): the nontrivial cube roots have \(\lambda\)-adic difference from \(1\) of valuation one, and reduction modulo \(\lambda\) distinguishes the two possible signs. Thus every \(a\) prime to \(3\) has a unique unit \(\epsilon(a)\) with \(\epsilon(a)\equiv a\pmod3\). Define \(a^\circ=a/\epsilon(a)\).

For the modulus \(\mathfrak m=(3p)\), the ray-class exact sequence and class number one give \[
\operatorname{Cl}_{(3p)}(K)
 \simeq
 \bigl((\mathcal O/(p))^\times\times
             (\mathcal O/(3))^\times\bigr)/\mu_6
 \simeq \mathbf F_{p^2}^{\,\times}.
\tag{H1}
\] The last map sends the ideal class of \((a)\), with \(a\) prime to \(3p\), to \(a^\circ\bmod p\). It is independent of the generator: if \(u\in\mu_6\), then \(\epsilon(ua)=u\epsilon(a)\). Every residue class can be represented using the Chinese remainder theorem. The kernel is exactly the principal ideals with a generator congruent to \(1\) modulo \(3p\). This proves the displayed isomorphism, including its explicit formula. The exact sequence being used is Milne (2020a), V, Theorem 1.7, pp.~150--151.

Identify the three cube roots in \(\mathbf F_{p^2}\) with \(\mu_3\subset\mathbf C\) by the reduction of the fixed \(\zeta\). Then \[
\eta_p((a))=
 \left((a^\circ\bmod p)^{(p^2-1)/3}\right)^{\!\uparrow}
\tag{H2}
\] is a nontrivial order-three ray-class character. The superscript means the inverse of the chosen reduction identification, not an arbitrary numerical logarithm.

The primitive Kummer character \(\chi_{p,3}\) has conductor dividing \((3p)\) by AM-K003 and has exact order three. A cyclic group has exactly two nontrivial characters of order three, which are inverses. Therefore \[
\chi_{p,3}=\eta_p\quad\hbox{or}\quad\eta_p^{-1}
\tag{H3}
\] as ray-class characters, with the possible inverse accounting for the chosen Artin/Frobenius convention. AM-K001 proves that the two resulting L-functions are identical. Thus (H2) identifies the desired L-function without a new cubic-reciprocity input.

This also recovers the conductor distinction. The character descends to modulus \((p)\) exactly when the order-three character of \(\mathbf F_{p^2}^\times\) is trivial on the image of \(\mu_6\). Its restriction to \(\mu_3\) is trivial exactly when \(9\mid p^2-1\), or \(p\equiv8\pmod9\) under the standing congruence. The order-two unit causes no condition for an order-three character. Hence the primitive modulus is \((p)\) in that case; otherwise AM-K003 gives \((p)\lambda^2=(3p)\). The argument includes \(p=2\); no injectivity of \(\mu_6\to\mathbf F_{p^2}^\times\) was assumed.

Complex conjugation acts on (H1) by \(z\mapsto z^p\). Since \(p\equiv-1\pmod3\), it acts on \(\eta_p\) by inversion, as required.

\subsubsection{\texorpdfstring{1.2. The wild case \(p=3\)}{1.2. The wild case p=3}}\label{the-wild-case-p3}

For \(\mathfrak m=(9)=\lambda^4\), normalize again by the unique unit modulo \(3\). This gives \[
\operatorname{Cl}_{(9)}(K)
 \simeq (1+3\mathcal O)/(1+9\mathcal O)
 \simeq (\mathcal O/3\mathcal O,+)
 \simeq (\mathbf Z/3\mathbf Z)^2.
\tag{H4}
\] The middle map is \(1+3u\mapsto u\bmod3\); multiplication becomes addition because the cross term is divisible by \(9\). It is bijective, and the unit normalization proves the first isomorphism.

Write \(u=a+b\zeta\) modulo \(3\). Conjugation sends \((a,b)\) to \((a-b,-b)\). A character of this additive group has the form \(\zeta^{sa+tb}\). It is anti-invariant precisely when \[
s(a-b)+t(-b)=-sa-tb\quad\text{for every }a,b,
\] which is equivalent to \(s=0\). Thus the two nontrivial anti-invariant characters are \[
\eta_3(1+3(a+b\zeta))=\zeta^b,\qquad \eta_3^{-1}.
\tag{H5}
\] The actual Kummer character has conductor \((9)\) and is anti-invariant, by AM-K002 and AM-K003. It must be one of these two. Again AM-K001 identifies their L-functions. This explicitly includes the wild conductor; it is not a tame specialization.

\subsubsection{1.3. Character separation and the ideal-indexed Dirichlet series}\label{character-separation-and-the-ideal-indexed-dirichlet-series}

Let \(H=\operatorname{Cl}_{\mathfrak m}(K)\) with the modulus just used. Its regular module over \(\mathbf Q(\zeta)\) has the exact projector \[
e_{\eta}=\frac1{|H|}\sum_{h\in H}\eta(h)^{-1}[h].
\tag{H6}
\] Translation by an ideal class \([\mathfrak a]\) acts on its image by \(\eta(\mathfrak a)\). The elementary identity \[
1+\eta(h)+\eta(h)^{-1}
   =3\,\mathbf 1_{\eta(h)=1}
\tag{H7}
\] gives the complete finite character separation. These statements follow directly by changing the index in (H6) and summing the three roots of unity.

A smaller integral realization uses the quotient \(H/\ker\eta\simeq C_3\): the rank-two augmentation representation has trace \(2\) on the identity and \(-1\) otherwise. Over \(\mathbf Q(\zeta)\) it splits as \(\eta\oplus\eta^{-1}\). The projector onto one line uses division by \(3\), so an integral construction must retain that denominator; (H6) is not asserted to be an integral idempotent at primes dividing \(|H|\).

For \(\Re s>1\), the ray-modulus L-function is exactly \[
L_{\mathfrak m}(s,\eta)
 =\frac16\sum_{(a,b)\in\mathbf Z^2\setminus\{(0,0)\}}
 \frac{\eta((a+b\zeta))}{(a^2-ab+b^2)^s},
\tag{H8}
\] where the numerator is zero when \((a+b\zeta)\) is not coprime to \(\mathfrak m\). Each nonzero integral ideal has precisely six generators, which proves the formula by absolute convergence. This is a concrete ideal-indexed host. It does not turn the series into a rational Dirichlet L-function.

\subsubsection{1.4. The harmless extra Euler factor is harmless for the actual window}\label{the-harmless-extra-euler-factor-is-harmless-for-the-actual-window}

When \(p\equiv8\pmod9\), the convenient modulus \((3p)\) omits the unramified Euler factor at \(\lambda\) of norm \(3\). Since \(\lambda\) is fixed by conjugation and the character is odd-order and anti-invariant, \(\chi(\lambda)=1\).

More generally, remove an unramified Euler factor of norm \(q>1\) and value \(u\), \(|u|=1\): \[
L_{\rm imp}(s)=(1-uq^{-s})L_{\rm prim}(s).
\] Use the same normalized ramp as AM-K005 for each function, including its own constant \(L'/L(1)\). With \(J=\lfloor\log_q x\rfloor\), direct subtraction gives \[
G_{\rm imp}(x)-G_{\rm prim}(x)
 =\log q\left(
     \sum_{k=1}^{J}u^k+
     \frac{x(u/q)^{J+1}}{1-u/q}
            \right).
\tag{H9}
\] Indeed the removed prime powers subtract \(\log q\sum_{k\le J}u^k(x/q^k-1)\), and the change in \(L'/L(1)\) adds \(x\log q\sum_{k\ge1}(u/q)^k\). Their finite sums cancel to (H9). Consequently the difference is \(O(1+\log x)\) for fixed \(q\).

For the present real-valued ramps, their AM-K005 windows \(W(\theta,\delta)\) are equivalent for every \(\theta>0\). To prove this with the actual quantifiers, start with a majorant for one ramp at exponent \(\eta/2\), and add \(C(1+\log x)\) to majorize the other. The added integral over \([Y,Y^{1+\delta}]\) is at most \[
C\int_Y^\infty (1+\log x)x^{-\theta-1}\,dx
=O_\theta(Y^{-\theta}(1+\log Y)).
\] Increase the prescribed lower threshold so that this and \(Y^{\eta/2}\) are each at most \(Y^\eta/2\). The same argument reverses the roles. Thus the explicit nonminimal modulus in (H1) introduces no new analytic gate.

\textbf{Closed subgate.} For all allowed cubic \(p\), the finite ray-class group, the actual character up to its L-equivalent inverse, its character projector, and the ideal-indexed L-series are explicit. No Kummer extension has been inserted as a CM field.

\subsection{2. AM-K008: the Kummer character is an admissible splitting twist}\label{am-k008-the-kummer-character-is-an-admissible-splitting-twist}

This statement works for every odd \(r\), not only \(3\). Put \(K^+=\mathbf Q(\zeta_r+\zeta_r^{-1})\), let \(c\) denote CM conjugation, and let \(\chi\) be the finite-order Kummer Hecke character. AM-K002 gives \(\chi^c=\chi^{-1}\).

For \(a\in\mathbb A_{K^+}^{\times}\), conjugation fixes \(a\), so \(\chi(a)^2=1\). Its order also divides the odd number \(r\), hence \[
\chi|_{\mathbb A_{K^+}^{\times}}=1.
\tag{H10}
\] At infinity it is trivial because a continuous finite-order character of each connected group \(\mathbf C^\times\) is trivial.

Let \(T=\operatorname{Res}_{K/K^+}^{1}\mathbf G_m\) and choose \(j\) with \(2j\equiv1\pmod r\), for example \(j=(r+1)/2\). On \(T(\mathbb A_{K^+})\) define \[
\nu(z)=\chi(z)^j.
\] It is unitary, continuous, finite-order, and trivial on \(T(K^+)\subset K^\times\). Moreover \[
\nu(x/cx)=\chi(x)^{2j}=\chi(x).
\tag{H11}
\] For \(r=3\), one can take \(j=2\). For a unitary group \(U(S)\) write \(\nu_S(g)=\nu(\det g)\).

\subsubsection{2.1. Exact transport of the oriented Weil record}\label{exact-transport-of-the-oriented-weil-record}

Suppose an actual oriented Weil datum over the CM extension \(K/K^+\) is given, with splitting characters \(\chi_S,\chi_T\) and action \(\omega\). Let \(a,b\in\mathbf Z\). Set \[
\chi_S'=\chi_S\chi^a,\qquad
\chi_T'=\chi_T\chi^b,\qquad
\omega'(g,h)=\nu_S(g)^{-b}\nu_T(h)^a\,\omega(g,h).
\tag{H12}
\] Then it satisfies the same oriented Weil-datum conditions, with the new splitting characters, on the same Heisenberg space and with the same additive character.

Here is a check of every field in the actual source record, not just an appeal to the name of the representation. In LocalNewformFunction.lean, the record IsWeilDatumO is at lines 53--75.

\begin{itemize}
\tightlist
\item
  \textbf{std:} the additive character is unchanged.
\item
  \textbf{hecS, hecT, resS, resT:} products remain Hecke characters; (H10) preserves the prescribed restriction to the real subfield.
\item
  \textbf{om\_pres:} a constant scalar multiple preserves each defining quasi-periodicity, level, continuity and archimedean smoothness condition of IsThetaFunO.
\item
  \textbf{om\_heis:} the scalar commutes with every linear Heisenberg operator, leaving the same symplectic action.
\item
  \textbf{om\_aut:} on rational unitary elements the determinants lie in \(T(K^+)\), where \(\nu=1\).
\item
  \textbf{om\_cont:} the new scalar is continuous; it multiplies the old continuous matrix coefficient.
\item
  \textbf{om\_unit:} the scalar has absolute value one, so every displayed squared-norm integral is identical.
\item
  \textbf{om\_split:} the source splitting factor, line 51, is \(\chi_T(x)^{-1}\chi_S(y)\). The new factor is this times \(\chi(x)^{-b}\chi(y)^a\). For a source determinant root, \(x/cx=\det g\) and \(y/cy=\det h\), so (H11) makes these exactly the scalar in (H12). The same positive real constant works.
\end{itemize}

For the last check, the source's split-place definition IsDetAt, TiedWeil.lean:77, also implies the determinant-root identity. At the selected idempotent \(e\) it reads \(x=e\det g+(1-e)\); the supported unitary element is the identity off \(e+ce\) and satisfies \(c(\det g)=(\det g)^{-1}\). Multiplying in the two orthogonal idempotent components gives \(x/cx=\det g\).

Multiplication by the scalar character also preserves the group-homomorphism property of \(\omega\), so the ambient type of the new action is correct. This completes the transport proof.

The related actual interface TiedWeil.lean:64,83--96 defines IsDetRoot and TwRule. Its four formulas explicitly permit a splitting-character change with trivial restriction to the real base, with the sign reversed on the reversed-rank edge. TwRule requires its character to belong to the supplied family X.Chr; membership of the new character in an existing family is not deduced from that record. Formula (H12) uses the oriented convention splitFacO, and therefore keeps its signs fixed. A reversal must use the source's corresponding formula rather than silently reuse a sign.

\subsubsection{2.2. Level and integral scope}\label{level-and-integral-scope}

The archimedean representation and archimedean norms are unchanged. At a finite place where \(\chi\) is unramified, \(\nu_S\) is trivial on every compact subgroup of \(U(S)\): compact determinant images have zero valuation, and an unramified character is trivial on local units. At each remaining place, replacing a compact level \(U_v\) by \(U_v\cap\ker\nu_S\) costs group index at most \(r\).

For \(r=3\), only the rational places \(p,3\) can contribute. Thus the product of these elementary local kernel indices is at most \(9\) for \(p\ne3\), and at most \(3\) for \(p=3\). This is solely the index needed to kill this one character on a given level. It is not the complete adelic covolume, a Hecke-orbit count, or the cost of the archive's arithmetic columns. Root-of-unity multiplication preserves an integral coefficient lattice when that lattice is a module over a ring containing those roots; projectors and saturation may introduce further denominators.

\textbf{Closed subgate.} The target character has exactly the relative restriction needed for the archive's splitting-change formulas. This is an actual transport theorem for a supplied Weil datum. It does not manufacture the required D7 datum or the global arithmetic theta comparison.

\subsection{3. AM-K009: coherent twists telescope and identify the remaining comparison}\label{am-k009-coherent-twists-telescope-and-identify-the-remaining-comparison}

Fix a chain of groups \(G_0,\ldots,G_m\) and kernels \(K_i(g_i,g_{i+1})\). State the convention explicitly: \[
(\Theta_i f)(h)=\int_{[G_i]}K_i(g,h)f(g)\,dg,
\] on any function space where these integrals exist. Let \(\nu_i\) be the determinant character from AM-K008 and \(t_i\in\mathbf Z\). Define \(M_i f=\nu_i^{t_i}f\) and \[
K_i'(g,h)=\nu_i(g)^{-t_i}\nu_{i+1}(h)^{t_{i+1}}K_i(g,h).
\] Then \[
\Theta_i'=M_{i+1}\Theta_iM_i^{-1},\qquad
\Theta_{m-1}'\cdots\Theta_0'
 =M_m(\Theta_{m-1}\cdots\Theta_0)M_0^{-1}.
\tag{H13}
\] Both formulas follow by substitution and cancellation of the adjacent multiplication operators. With a convention conjugating the input function, the corresponding powers are conjugated; the same cancellation remains. Each \(M_i\) is unitary on \(L^2\), and therefore norms and pairings of simultaneously transported vectors are unchanged.

Consequently, coherent splitting transport alone supplies no new prime-weighted arithmetic identity. It gives isomorphic analytic operators and the transported endpoints. To obtain the new explicit-formula functional one must identify a character-dependent arithmetic column, Hecke action or endpoint that carries (H2)/(H5), and verify its integral structure and its actual cost.

This distinction can be seen in the checked source, rather than inferred from broad terminology:

\begin{longtable}[]{@{}
  >{\raggedright\arraybackslash}p{(\columnwidth - 4\tabcolsep) * \real{0.3333}}
  >{\raggedright\arraybackslash}p{(\columnwidth - 4\tabcolsep) * \real{0.3333}}
  >{\raggedright\arraybackslash}p{(\columnwidth - 4\tabcolsep) * \real{0.3333}}@{}}
\toprule\noalign{}
\begin{minipage}[b]{\linewidth}\raggedright
Actual source interface
\end{minipage} & \begin{minipage}[b]{\linewidth}\raggedright
What it contains
\end{minipage} & \begin{minipage}[b]{\linewidth}\raggedright
What (H12) supplies
\end{minipage} \\
\midrule\noalign{}
\endhead
\bottomrule\noalign{}
\endlastfoot
LevChiEqOne.lean:23--64, TwDat and TwDat.Is & Level charts, stabilizers, arithmetic representatives, determinant roots, and the exact character relation to the common family & The permitted determinant character; it does not construct the charts or representatives \\
Same \url{file:73–77}, abelian₂ & A prescribed product of arithmetic scalar factors times jetL of the actual abelian theta section & No new equality with a prime sum follows from scalar transport \\
Same \url{file:87–101}, VacT and PinX & The moved-vacuum theta identity, its nonzero constant, fixed positive lattice pin, and the entire family of matching characters & Once supplied, the splitting factors transform by (H12); the arithmetic identity still needs proof for the new family \\
Same \url{file:103–115}, PinColT & Finite combination datum, rational datum, and a constant combining the level volume, arithmetic weights, vacuum constant and determinant-root evaluations & Individual phase factors are explicit; the full nonzero constant and covolume are not evaluated here \\
\nolinkurl{CharacterBoxes/Supply.lean:10,42} & FReq is a predicate on Scale, Frame and HostAx; GChi adds EAtH and VVle & ``Generic request'' is generic only within this HostAx domain, not over arbitrary CM fields \\
\end{longtable}

A precise remaining task is therefore the \textbf{twisted arithmetic comparison} for the explicit ray-class character and the actual CM lattice below: construct the arithmetic section and prove that its functional is AM-K005's \(G_\chi\), with all local multiplicities, conductor factors and integral denominators. That is narrower than leaving the character, finite host and local splitting laws unspecified. It still includes a substantive theorem; this report does not rename it as a proved window.

\subsection{4. AM-K010: an explicit CM supply tower with the induced cubic type}\label{am-k010-an-explicit-cm-supply-tower-with-the-induced-cubic-type}

The previous finite host is a ray-class group, whereas the geometric construction also requires a CM field, polarization and lattice. The following explicit tower supplies several precise HostDer-shaped ingredients while keeping the Kummer character in the descent of a CM abelian variety.

\subsubsection{4.1. A criterion for the invariant half-different}\label{a-criterion-for-the-invariant-half-different}

Let \(B/K\) be a finite Galois extension of a CM field \(K\), with \(B\) CM and conjugation \(c\) restricting to conjugation of \(K\). Write \(\mathfrak D=\mathfrak D_{B/K}\). There exists a Gal\((B/K)\)-stable integral ideal \(\mathfrak b\) with \[
\mathfrak D=\mathfrak b\,c(\mathfrak b),\qquad
\mathfrak b+c(\mathfrak b)=\mathcal O_B
\tag{H14}
\] if and only if \(B/K\) is unramified at every finite prime \(\mathfrak q\) of \(K\) fixed by \(c\).

For necessity, let \(\mathfrak q=c\mathfrak q\). The Galois group is transitive on primes above \(\mathfrak q\), so stability makes the valuation of \(\mathfrak b\) the same integer at all those primes. Conjugation permutes the same set, so \(c(\mathfrak b)\) has the same valuation there. Coprimality forces that integer to be zero, hence \(\mathfrak D\) has valuation zero at all these primes. A finite number-field extension is unramified at a prime exactly when its local different is trivial.

For sufficiency, partition the non-conjugation-fixed primes in the support of the relative discriminant into pairs \(\{\mathfrak q,c\mathfrak q\}\), and choose one member of every pair. Put into \(\mathfrak b\) all prime factors of \(\mathfrak D\) above the chosen members, with their full different exponents. Since the extension is Galois, \(\mathfrak D\) and this choice are Gal\((B/K)\)-stable. Conjugation preserves the relative different and supplies precisely the complementary factors. This proves (H14). No class-number-one assumption on \(B\) is required.

For \(K=\mathbf Q(\zeta_3)\), both the prime over \(2\) and the prime over \(3\) are conjugation-fixed. Thus an induced-type invariant half-different of this form excludes relative ramification at \(2\) or \(3\). In particular a literal relative extension ramified at \(2\), as in an added nontrivial real \(2\)-power cyclotomic tower, cannot have all these properties with this distinguished \(K\). This is a specific obstruction to retaining the archive's \(2\)-adic tower, not an obstruction to all CM host fields containing \(K\).

\subsubsection{4.2. The explicit tower}\label{the-explicit-tower}

For each \(n\ge1\) set \[
B_n=\mathbf Q(\zeta_{3\cdot7^n})
   =K(\zeta_{7^n}),\quad
d_n=6\cdot7^{n-1},\quad
a_n=7^{n-1}(6n-1).
\tag{H15}
\] Then \([B_n:K]=d_n\), and \(B_n\) is a Galois CM field of degree \(2d_n\). Its Galois group over \(\mathbf Q\) is \[
(\mathbf Z/3\mathbf Z)^\times
 \times(\mathbf Z/7^n\mathbf Z)^\times.
\] Complex conjugation is \((-1,-1)\). At \(3\), inertia is contained in the first factor; at \(7\), it is contained in the second; away from \(3,7\) it is trivial. Conjugation is in neither of the two nontrivial inertia groups. The inertia group for \(B_n/B_n^+\) is their intersection with \(\{1,c\}\), and is therefore trivial at every finite prime. Hence \[
B_n/B_n^+\ \text{is everywhere unramified at finite primes}.
\tag{H16}
\]

The prime \(7\) splits in \(K\): \[
(7)=(3+\zeta)(3+\zeta^2).
\] The two factors have norm \(7\) and are distinct. Each ramifies totally in \(B_n/K\), yielding primes \(P_n,cP_n\) of residue degree one and ramification degree \(d_n\). There is no other relative ramification.

The local ring is generated by \(\zeta_{7^n}\). The derivative of \[
\Phi_{7^n}(X)=\frac{X^{7^n}-1}{X^{7^{n-1}}-1}
\] at this root has valuation \[
n\,d_n-7^{n-1}=a_n,
\] because \(v_{P_n}(7)=d_n\) and \(v_{P_n}(\zeta_{7^n}^{\,7^{n-1}}-1)=7^{n-1}\). The monogenic different formula therefore gives \[
\mathfrak D_{B_n/K}=P_n^{a_n}(cP_n)^{a_n}.
\tag{H17}
\] Thus \(\mathfrak b_n=P_n^{a_n}\) is Gal\((B_n/K)\)-stable and satisfies (H14). This is a fully specified half-different, not an existence assertion for an unspecified field.

Discriminant transitivity and (H16) yield \[
|D_{B_n}|=3^{d_n}7^{2a_n},\qquad
D_{B_n^+}=3^{d_n/2}7^{a_n},\qquad
\operatorname{rd}(B_n)=\sqrt3\,7^{\,n-1/6}.
\tag{H18}
\] The roots of unity in \(B_n\) are exactly \(\mu_{6\cdot7^n}\). One elementary check is by ramification: any root order can involve only \(2,3,7\); a factor \(4\) would ramify at \(2\), a factor \(9\) would require ramification index at least \(6\) at \(3\) instead of \(2\), and a higher \(7\)-power would require larger ramification index at \(7\). The stated roots already lie in \(B_n\).

Every one of these roots has relative norm one to \(K\). Indeed \[
N_{B_n/K}(-1)=1,\quad
N_{B_n/K}(\zeta_3)=\zeta_3^{d_n}=1,\quad
N_{B_n/K}(\zeta_{7^n})=1.
\tag{H19}
\] The last equality is the constant term of the even-degree cyclotomic polynomial. These three roots generate the full group. This verifies the actual norm-of-roots condition, not the stronger and false assertion that \(B_n\) has only \(\pm1\) as roots.

The cyclotomic facts used here are Milne (2020c), Theorem 6.4, printed p.~99 (PDF p.~101), and Remark 6.6(b), printed p.~100 (PDF p.~102), together with the elementary local cyclotomic derivative computation displayed above. The prime splitting and all discriminant exponents have been shown explicitly.

\subsubsection{4.3. A positive self-dual Hermitian lattice}\label{a-positive-self-dual-hermitian-lattice}

Set \[
L_n=\mathfrak b_n^{-1},\qquad
h_n(x,y)=\operatorname{Tr}_{B_n/K}(x\,c(y)).
\tag{H20}
\] It is a free \(\mathcal O_K\)-module of rank \(d_n\), because \(\mathcal O_K\) is a PID. Its Hermitian dual is \[
L_n^\#=\mathfrak D_{B_n/K}^{-1}\,(cL_n)^{-1}
       =(\mathfrak b_n c\mathfrak b_n)^{-1}c\mathfrak b_n
       =\mathfrak b_n^{-1}=L_n.
\tag{H21}
\] The first equality follows from the definition of the inverse different as the trace-dual of \(\mathcal O_{B_n}\), extended to fractional ideals; it can also be checked after localization at each prime. Thus \(h_n\) is integral and unimodular on \(L_n\).

Choose the CM type \(\Phi_n\) consisting of the \(d_n\) embeddings extending a fixed embedding \(K\hookrightarrow\mathbf C\). Then \[
h_n(x,x)=\sum_{\sigma\in\Phi_n}|\sigma(x)|^2
\] under that embedding, so the form is positive definite. Choose an \(\mathcal O_K\)-basis and use the matrix convention \(h_n(x,y)=\bar y^{\,t}H_nx\) for coordinate columns. Then \(H_n\) is integral Hermitian and invertible over \(\mathcal O_K\). Its determinant is a positive rational unit, hence \(1\).

This exactly supplies the mathematical shape of HostDer.principal and HostDer.unimodular for this induced type, together with the unramified CM quotient and discriminant square. It supplies them as new mathematical objects, not as an instance of the existing parameterized HostRaw record.

\subsubsection{4.4. Actual CM abelian varieties and the cubic Hom line}\label{actual-cm-abelian-varieties-and-the-cubic-hom-line}

Use the elliptic curves from AM-K004: \[
A_0:y^2=x^3+1,\qquad A_p:y^2=x^3+p^2.
\] Both have \(\mathcal O_K\)-endomorphisms over \(K\). The chosen \(\mathcal O_K\)-basis of \(L_n\) realizes multiplication by \(\mathcal O_{B_n}\) as a subring of \(M_{d_n}(\mathcal O_K)\). Consequently it acts, over \(K\), on the actual products \[
\mathcal A_{0,n}=A_0^{d_n},\qquad
\mathcal A_{p,n}=A_p^{d_n}.
\] Give each product the polarization represented relative to its product principal polarization by the same positive integral Hermitian matrix \(H_n\). This is a principal polarization: positivity is the Riemann-form criterion, and unimodularity makes the polarization morphism an isomorphism. The entries are endomorphisms defined over \(K\), so the morphism is defined over \(K\). The adjoint of multiplication by \(b\in B_n\) for \(h_n\) is multiplication by \(c(b)\), proving that Rosati induces CM conjugation.

For completeness, the complex construction can be written directly as \[
\mathbf C^{\Phi_n}/\Phi_n(L_n).
\] Choose \(\xi=-1/\sqrt{-3}\) with the fixed embedding \(\sqrt{-3}\mapsto i\sqrt3\), so \(\operatorname{Im}\sigma(\xi)>0\) for \(\sigma\in\Phi_n\). The alternating form \(\operatorname{Tr}_{B_n/\mathbf Q}(\xi x c(y))\) is integral, positive in the Riemann convention and self-dual: the inverse different of \(K/\mathbf Q\) is generated by \(\xi\), and (H21) handles the relative trace. This produces the same polarized isomorphism class over \(\mathbf C\) as the Hermitian matrix construction, after scaling the common elliptic period. The Riemann-form and CM-pair framework is Milne (2020b), §2, Example 2.9, Example 2.11, and Proposition 3.17; Proposition 3.12 describes the induced CM isogeny type.

The coordinatewise isomorphism \[
g_p^{d_n}:\mathcal A_{p,n}\longrightarrow\mathcal A_{0,n}
\] over \(K(p^{1/3})\) is \(\mathcal O_{B_n}\)-linear because the scalar elliptic twist commutes with every matrix over \(\mathcal O_K\). It respects the displayed polarizations for the same reason. Its descent cocycle is the scalar \(\chi_{p,3}\) in \(\mathcal O_K^\times\subset\mathcal O_{B_n}^\times\).

On any one-dimensional \(B_n\)-eigensummand of \(H^1\) extending the chosen \(\zeta\)-eigensummand over \(K\), the ratio of the two Galois characters is therefore exactly \(\chi_{p,3}\), by the AM-K004 intertwining computation. Thus the explicit rank-one Artin Hom line survives in a principally polarized CM model with the lattice (H20). The splitting field is used only to trivialize the descent cocycle; it is not asserted to be a CM field.

The stable Faltings heights of the underlying products are \[
h_{\rm Fal}(\mathcal A_{0,n})=d_n h_{\rm Fal}(A_0),\qquad
h_{\rm Fal}(\mathcal A_{p,n})=d_n h_{\rm Fal}(A_0).
\tag{H22}
\] This uses additivity for products and the geometric isomorphism of \(A_p\) and \(A_0\). It is independent of the chosen principal polarization. It is not an evaluated theta-lift cost: derivatives, indices, normalized kernel constants and saturation still have to be computed. In particular, on the \(d_n\)-dimensional \(\zeta\)-eigenspace the determinant twist is \(\chi_{p,3}^{d_n}=1\) because \(3\mid d_n\), while each chosen \(B_n\)-eigensummand still has nontrivial character; a determinant or height alone does not recover that summand.

\subsubsection{4.5. Precisely which source requirements have and have not been supplied}\label{precisely-which-source-requirements-have-and-have-not-been-supplied}

\begin{longtable}[]{@{}
  >{\raggedright\arraybackslash}p{(\columnwidth - 4\tabcolsep) * \real{0.3333}}
  >{\raggedright\arraybackslash}p{(\columnwidth - 4\tabcolsep) * \real{0.3333}}
  >{\raggedright\arraybackslash}p{(\columnwidth - 4\tabcolsep) * \real{0.3333}}@{}}
\toprule\noalign{}
\begin{minipage}[b]{\linewidth}\raggedright
Source requirement
\end{minipage} & \begin{minipage}[b]{\linewidth}\raggedright
Supplied here
\end{minipage} & \begin{minipage}[b]{\linewidth}\raggedright
Remaining difference
\end{minipage} \\
\midrule\noalign{}
\endhead
\bottomrule\noalign{}
\endlastfoot
ConjEMemHostFIff.lean:89,95, CM field and CM type & \(B_n\) and the type induced from \(K\) & The original record has its own Frame and distinguished quadratic-field parameters \\
Same \url{file:135}, finite unramified CM quotient & (H16) & Proven for the new tower, not transported through HostRaw \\
Same \url{file:156}, invariant coprime half-different & (H17), \(\mathfrak b_n=P_n^{a_n}\) & Its tower prime is \(7\); retaining ramification at inert \(2\) violates §4.1 \\
Same \url{file:159}, unimodular trace lattice & (H20)--(H21) & Actual later arithmetic columns still need matching to this lattice \\
Same \url{file:163}, discriminant square & (H18) & The full source root-discriminant inequalities depend on its different parameters \\
Same \url{file:170}, norm of every root of unity is one & (H19) & Its separate stronger MuPM predicate fails here; no use of that predicate may be carried silently \\
Main CM abelian variety and cubic eigensummand & Actual \(\mathcal A_{0,n},\mathcal A_{p,n}\) with principal polarizations and the Hom line & No full MainX chain, terminal vector or normalized period has been constructed \\
ExistsPageN.lean:19,39 and Supply.lean:42 & No literal instance claimed & The original cyclotomic subgroup, \(2\)-adic PageN scale, character-specific EAtH and cohomology supply require a \(7\)-adic/relative replacement \\
\end{longtable}

The phrase ``generic CM theory'' therefore conceals two separate issues: several geometric objects now exist explicitly, but the original rational congruence/tower supply and its arithmetic prime functional are not consequences of that existence.

\subsection{5. AM-K011: the actual generic field record and its finite-place boundary}\label{am-k011-the-actual-generic-field-record-and-its-finite-place-boundary}

The source contains a geometric field interface weaker than HostRaw. It can be supplied completely at the level of mathematical data. This closes a named source input; the source's arithmetic character family still needs a separate extension.

\subsubsection{5.1. A degree-divisibility adjustment using a second split prime}\label{a-degree-divisibility-adjustment-using-a-second-split-prime}

Set \[
B_n^*=\mathbf Q(\zeta_{3\cdot7^n\cdot13}),\qquad
F_n^*=(B_n^*)^+,\qquad
f_n=[B_n^*:K]=72\cdot7^{n-1}.
\tag{H23}
\] The extra prime \(13\) splits in \(K\), since \((13)=(4+\zeta)(4+\zeta^2)\). Cyclotomic disjointness gives the displayed degree, so \(8\mid f_n\). The CM quotient \(B_n^*/F_n^*\) is unramified at all finite primes: inertia at \(3\), \(7\) or \(13\) is confined to that factor of the cyclotomic Galois group, whereas conjugation is \(-1\) in all three factors. Thus it lies in none of the inertia groups.

Relative ramification over \(K\) is confined to the two split rational primes \(7,13\). The half-different criterion in AM-K010 applies. More explicitly, choose the primes \((3+\zeta)\) and \((4+\zeta)\) of \(K\), and put \[
\mathfrak b_n^*
  =\prod_{P\mid(3+\zeta)}P^{a_n}
     \prod_{Q\mid(4+\zeta)}Q^{11}.
\tag{H24}
\] Adjoining the other prime-power roots is unramified locally, so the different exponent at \(7\) remains \(a_n\) and that at \(13\) is the tame exponent \(12-1=11\). This proves \(\mathfrak D_{B_n^*/K}=\mathfrak b_n^*c(\mathfrak b_n^*)\) with coprime, Gal\((B_n^*/K)\)-stable halves. Its inverse is therefore again a free positive unimodular Hermitian \(\mathcal O_K\)-lattice, now of rank \(f_n\). All the principally polarized product and Artin Hom constructions of §4.4 apply with this lattice.

The roots of unity are \(\mu_{6\cdot7^n\cdot13}\). Their relative norms are one: the norm of \(\zeta_3\) is one because \(3\mid f_n\); the norms of primitive \(7^n\)-th and \(13\)-th roots are powers of the constant terms of their even-degree cyclotomic polynomials; and \(f_n\) is even. Coprime-conductor discriminant multiplication gives \[
|D_{B_n^*}|=3^{12d_n}7^{24a_n}13^{22d_n},
\qquad
\operatorname{rd}(B_n^*)=\sqrt3\,7^{\,n-1/6}13^{11/12}.
\tag{H25}
\] Thus satisfying the degree condition does not require introducing relative ramification at inert \(2\).

\subsubsection{5.2. Every field of AbelianUnramifiedExtensionConditions}\label{every-field-of-abelianunramifiedextensionconditions}

The exact file is \nolinkurl{Formalism/LocalToGlobal/Columns/AbelianHostFieldConditions.lean.} Its Fields structure is at lines 19--42, Fields.Ok at lines 81--108, and AbelianUnramifiedExtensionConditions at lines 110--120. To avoid its letter K being confused with the fixed quadratic field, denote its coefficient field by \(C\).

Choose the source fields \[
X.F=F_n^*,\qquad X.E=B_n^*,\qquad C=X.K=B_n^*,
\] with their natural embeddings and conjugations. Take \(X.\mathrm{Tau}=\operatorname{Gal}(B_n^*/K)\), and let its embeddings be the corresponding automorphisms of \(B_n^*\). Then every Fields.Ok condition holds:

\begin{itemize}
\tightlist
\item
  \textbf{realF, rankE, star\_F, conjE, cplxE:} these are the CM-field identities for \(B_n^*/F_n^*\).
\item
  \textbf{emb\_star:} the cyclotomic automorphisms commute with conjugation.
\item
  \textbf{card\_Tau:} the cardinal is \(f_n=[F_n^*:\mathbf Q]\).
\item
  \textbf{emb\_inj:} distinct automorphisms are distinct embeddings.
\item
  \textbf{emb\_ne:} an embedding fixing \(K\) cannot equal a conjugated one, whose restriction to \(K\) is nontrivial.
\item
  \textbf{emb\_all:} every embedding \(B_n^*\to B_n^*\) is a Galois automorphism. Its restriction to \(K\) is either identity or conjugation, so it lies in the chosen set or its conjugate.
\item
  \textbf{cplxK, normalK, conjK:} the coefficient field \(C=B_n^*\) is a normal CM field, with its displayed complex embedding and conjugation.
\item
  \textbf{quad:} the source's Kq consists of elements whose images under all chosen embeddings agree. Since the identity embedding occurs, this is the fixed field of Gal\((B_n^*/K)\), hence precisely \(K\) and of rational degree two.
\end{itemize}

The remaining AbelianUnramifiedExtensionConditions fields are \textbf{eight\_dvd}, proved by (H23); \textbf{galE} and \textbf{abelE}, from cyclotomy; and \textbf{unram}, proved above. For the last field, the source ideal inequality is the ramification-index-one condition, exactly the finite unramified CM quotient already checked.

This constructs all mathematical data and proofs for that named record. It is not a Lean term or a build certificate. A larger normal cyclotomic coefficient field \(C\supset B_n^*\) may be used when finite character values require it: all embeddings of the normal field \(B_n^*\) still have their images in \(B_n^*\), so the same verification applies.

\subsubsection{5.3. Why the target is still outside the particular tame split-prime family}\label{why-the-target-is-still-outside-the-particular-tame-split-prime-family}

In the same file, IsMu0 at lines 314--320 requires a character to be unramified away from the two auxiliary rational primes \(g_i\), and tame everywhere. The preceding Primes.Ok record requires all the selected primes over \(g_i\) to have absolute norm \(g_i\) (lines 216--218), and QDat.Ok requires its primes to have prime absolute norm (line 347). These are substantive degree-one restrictions.

For \(p\ne3\) under the present hypothesis, \(p\equiv2\pmod3\) is inert in \(K\). Therefore every prime \(P\) over \(p\) in every number field containing \(K\) has \[
f(P/p)\ge2,\qquad NP=p^{f(P/p)},
\] so no such prime has prime absolute norm. AM-K003 says that \(\chi_{p,3}\) ramifies at the prime \((p)\) of \(K\). The extension \(B_n^*/K\) is unramified at \(p\), because \(p\) is neither \(7\) nor \(13\). The pulled-back character on \(B_n^*\) remains ramified: an unramified local extension leaves the inertia subgroup unchanged. Hence it cannot be placed solely in the displayed degree-one support.

At the prime above \(3\), the cases with conductor \(\lambda^2\) or \(\lambda^4\) are wild. Indeed tame inertia at residue characteristic \(3\) has no nontrivial quotient of order three, whereas these order-three characters are ramified. In (H23) the relative extension is unramified at \(3\), so the same obstruction persists after pullback. Thus all cases except \(p\equiv8\pmod9\) require a wild local character as well. In the exceptional congruence case the prime over \(3\) is unramified, but the inert-\(p\) degree restriction remains.

Consequently the general splitting transport AM-K008 is valid, while inserting its character into these particular IsMu0/QDat families is not automatic. A relative arithmetic adapter must allow the actual degree-two prime over \(p\), the wild conductor at \(3\) when present, and the resulting norm and lattice factors. This statement concerns these named families; it does not assert that no other source construction can represent the character.

\textbf{Closed source input and surviving gate.} The geometric field predicate AbelianUnramifiedExtensionConditions has an explicit supply with \(8\mid f\). The remaining work is the target's arithmetic character/column construction and its scale-uniform cost, with these precise nonsplit and wild local conditions retained.

\subsection{6. Resulting reduced programme}\label{resulting-reduced-programme}

For the cubic question, four concrete parts are now closed:

\begin{enumerate}
\def\labelenumi{\arabic{enumi}.}
\tightlist
\item
  the target character and all its finite conductor cases have a complete ray-class model;
\item
  the character is an allowed finite-order unitary splitting change, with exact source signs;
\item
  an explicit growing CM tower has an invariant coprime half-different, a positive unimodular trace lattice, the required norm-of-roots property, and a variant satisfying every mathematical field in AbelianUnramifiedExtensionConditions;
\item
  actual principally polarized products of the cubic CM elliptic twists carry the selected Artin Hom line.
\end{enumerate}

The remaining construction must place those objects in a relative arithmetic family, prove the exact equality between its arithmetic functional and AM-K005's normalized ideal-prime ramp, and compute its lift/period/lattice costs at all scales. The source's \(2\)-adic supply cannot simply be relabelled as the \(7\)-adic tower, and its displayed tame degree-one character family must be extended at the actual nonsplit and wild places. A successful estimate must give the window for every \(\theta>1/2\); all the identities here are compatible with the present gate Q3-KUMMER-WINDOW remaining open.

\subsection{References and source provenance}\label{references-and-source-provenance}

\begin{itemize}
\tightlist
\item
  Alpöge, L. (2026). \emph{Remark} and the stipulated Lean archive. \url{https://alpo.ge/remark.pdf}. Source files actually read: \nolinkurl{Formalism/LocalToGlobal/BesselHermitianTheta/TiedWeil.lean;} \nolinkurl{Formalism/LocalToGlobal/IndefiniteReal/LocalNewformFunction.lean;} \nolinkurl{Formalism/LocalToGlobal/Columns/AbelianHostFieldConditions.lean;} \nolinkurl{Development/AbelianVarietyTheory/Definitions/LatticeFrames/LevChiEqOne.lean;} \nolinkurl{Formalism/SharpRestriction/RamificationFrobenius/TotientGaloisClasses/ConjEMemHostFIff.lean;} its \nolinkurl{InertiaGaloisSymmetry/OfRawToHostRaw.lean;} \nolinkurl{Development/AutomorphicArithmetic/CharacterBoxes/Supply.lean} and ExistsPageN.lean.
\item
  Milne, J. S. (2020a). \emph{Class field theory} (v4.03). V, Theorem 1.7, pp.~150--151. \url{https://www.jmilne.org/math/CourseNotes/CFT.pdf}.
\item
  Milne, J. S. (2020b). \emph{Complex multiplication} (v0.10). §2, Examples 2.9 and 2.11; Propositions 3.12 and 3.17; Proposition 6.1. \url{https://www.jmilne.org/math/CourseNotes/CM.pdf}.
\item
  Milne, J. S. (2020c). \emph{Algebraic number theory} (v3.08). §6, Theorem 6.4, printed p.~99 (PDF p.~101); Remark 6.6(b)--(c), printed pp.~100--101 (PDF pp.~102--103). \url{https://www.jmilne.org/math/CourseNotes/ANT.pdf}.
\item
  Murty, M. R., \& Esmonde, J. (2005). \emph{Problems in algebraic number theory} (2nd ed.). Springer. §2.3, Exercise 2.3.2 and solution; §9.1. Author-hosted copy: \url{https://www.math.toronto.edu/~ila/2005_Book_ProblemsInAlgebraicNumberTheor.pdf}.
\item
  Companion contribution: AM-K001--AM-K005 in kummer\_extension.md. The earlier conductor, functional-equation, elliptic twist and analytic adapter proofs are incorporated by reference and not rederived here.
\end{itemize}

The archive was read from the shared pinned-source recovery, with its provenance files maintained by the parent, and from the child manifest for the earlier character files. No Lean compilation or newly formalized theorem is claimed by this addendum.

\clearpage

\section{AM-DO-017--020: literal moved-theta operators, a mixed-CM family, and exact split-place phases}\label{am-do-017020-literal-moved-theta-operators-a-mixed-cm-family-and-exact-split-place-phases}

Author: GPT-6 Astra Pro, ALP\_MACHINE \nolinkurl{/root/definite_oscillator}, 21 November 2026 UTC. The inherited reasoning effort is not separately exposed. This file gives full proofs of the new finite-operator assertions. AM-R-004, AM-DO-011--016, ALP-LIFT-07 and the source definitions are retained at their recorded status. Write \textbf{Σ\_tr} for AM-DO-015's normalized Mumford translate of F₁; it is distinct from the leading-coefficient-normalized hyperelliptic sigma function σ\_q used in the parallel calculation.

\subsection{1. The recovered literal source object}\label{the-recovered-literal-source-object}

The source definition sought in \nolinkurl{I.J.V.thMoved} has now been recovered and read. Its exact file is

\nolinkurl{Development/AbelianVarietyTheory/Definitions/AutomorphicMembership/IsDiscriminantTwelfthRootOfIdealLatticeIffDelta1.lean}.

The relevant locators are:

\begin{longtable}[]{@{}
  >{\raggedright\arraybackslash}p{(\columnwidth - 2\tabcolsep) * \real{0.5000}}
  >{\raggedright\arraybackslash}p{(\columnwidth - 2\tabcolsep) * \real{0.5000}}@{}}
\toprule\noalign{}
\begin{minipage}[b]{\linewidth}\raggedright
Source lines
\end{minipage} & \begin{minipage}[b]{\linewidth}\raggedright
Exact content
\end{minipage} \\
\midrule\noalign{}
\endhead
\bottomrule\noalign{}
\endlastfoot
18--24 & \nolinkurl{lamOf}, elliptic ideal lattice; the discriminant twelfth-root pin; an un-killed elliptic two-torsion point \\
36--51 & \nolinkurl{La}, \nolinkurl{Hf}, \nolinkurl{al}, \nolinkurl{U}, \nolinkurl{Iop} \\
110--118 & \nolinkurl{VacIn}: an elliptic \nolinkurl{PeriodPair}, its prescribed lattice, half-period, \nolinkurl{IsoKG}, its complex-linear extension, and \nolinkurl{KapIn} \\
130--138 & \nolinkurl{thMM}, \nolinkurl{thZero}, \nolinkurl{thHalf}, \nolinkurl{thMoved} \\
140--141 & \nolinkurl{eNat = chiFac · thMoved} \\
173--185 & literal conversion to the earlier \nolinkurl{Formalism.L2.Ker.Base} record \\
\end{longtable}

In mathematical notation its averaging operator is

\[
 I_{21}F=q_{21}^{-1/2}
       \sum_{x\in\Lambda_2/(\Lambda_1\cap\Lambda_2)}\alpha(x)U_xF,
 \qquad q_{21}=[\Lambda_2:\Lambda_1\cap\Lambda_2],
\tag{DO17.1}
\]

where the positive real square root is used and

\[
 U_xF(z)=e^{-\pi H(z,x)-\pi H(x,x)/2}F(z+x).
\tag{DO17.2}
\]

The source's \nolinkurl{La₁.addSubgroupOf La₂} is the intersection regarded as a subgroup of Λ₂; this matters for the index. Its sign is \texttt{al(x)=(-1)\^{}\{round((√d/2)H(x,x))\}}. In the admitted integral situations it is a sign semicharacter, so α⁻¹=α.

Writing Λ₀₀, Λ₁₀, Λ₁₁ for the standard, first-lattice-moved, and both-lattices-moved period lattices, the source literally uses

\[
 \mathrm{thMoved}_{L,L'}=I_{\Lambda_{10},\Lambda_{11}}
                       I_{\Lambda_{00},\Lambda_{10}}
                       \mathrm{thZero},
 \qquad\mathrm{thZero}=\mathrm{thMM}^{\boxtimes NN'}.
\tag{DO17.3}
\]

The exact seed is not left arbitrary: \nolinkurl{Formalism/LocalToGlobal/KernelLattices/LateralTheory.lean}, lines23--29, defines \nolinkurl{IsThM} by membership in the specified theta line and the fourth-tensor equation

\[
              \theta^{\boxtimes4}=\kappa\,\mathrm{kernel4}(L,\omega_0,\Lambda_M,g_C),
\tag{DO17.4}
\]

and takes \nolinkurl{thM} by choice from that predicate. The same file, lines128--178, repeats the complete base record and moved formulas. Lines230--238 separately name \nolinkurl{SignP} (comparison with the direct one-step average, up to sign) and \nolinkurl{OrderP} (equality of the two orders of moving the lattice factors). Thus phase and order compatibility are actual source obligations; norm preservation alone cannot replace them.

These additional modules were retrieved from the pinned mirror commit \nolinkurl{7cbf7a0c32fca1badf05931879a37c3f8a2b03b4}, with individual hashes in \nolinkurl{source_recovery_20261121T204914Z.json} and \nolinkurl{source_recovery_20261121T205055Z.json}. The intake does not supply original-author hashes for these paths. Retrieval and reading do not constitute a Lean build or a transitive proof audit.

\subsection{AM-DO-017. Exact normalized averaging between principal theta lines}\label{am-do-017.-exact-normalized-averaging-between-principal-theta-lines}

Let V be a complex vector space, H a positive Hermitian form linear in its first variable, and E=Im H. Let Λ₁,Λ₂ be commensurable full lattices selfdual for E:

\[
 \Lambda_i=\{v\in V:E(v,\Lambda_i)\subset\mathbf Z\}.
\tag{DO17.5}
\]

Let αᵢ:Λᵢ→U(1) satisfy

\[
 \alpha_i(x+y)=\alpha_i(x)\alpha_i(y)e^{\pi iE(x,y)},
\tag{DO17.6}
\]

and suppose their restrictions to Λ₀=Λ₁∩Λ₂ agree. Write T(Λᵢ,αᵢ) for the holomorphic functions satisfying U\_λF=αᵢ(λ)F. On this space use the probability-Haar norm

\[
 \|F\|_i^2=\frac1{\operatorname{vol}(D_i)}
     \int_{D_i}|F(z)|^2e^{-\pi H(z,z)}\,dz,
\tag{DO17.7}
\]

where Dᵢ is any fundamental domain for Λᵢ and dz is any fixed Lebesgue measure. This is independent of Dᵢ and of rescaling dz.

Then q={[}Λ₂:Λ₀{]}={[}Λ₁:Λ₀{]} is a positive finite integer and

\[
 \boxed{I_{21}F=q^{-1/2}\sum_{x\in\Lambda_2/\Lambda_0}
             \alpha_2(x)^{-1}U_xF}
\tag{DO17.8}
\]

is a well-defined nonzero isometry T(Λ₁,α₁)→T(Λ₂,α₂). In the source's sign-valued case this is exactly (DO17.1), including its square-root index. Both spaces are one-dimensional. No unspecified positive scalar occurs in (DO17.8).

\subsubsection{Proof}\label{proof-1}

Direct multiplication of the factors in (DO17.2) gives

\[
 U_xU_y=e^{-\pi iE(x,y)}U_{x+y},\qquad
 U_xU_y=e^{-2\pi iE(x,y)}U_yU_x.
\tag{DO17.9}
\]

Also

\[
 |U_xF(z)|^2e^{-\pi H(z,z)}
       =|F(z+x)|^2e^{-\pi H(z+x,z+x)}.
\tag{DO17.10}
\]

Thus translations are isometries whenever both sides are viewed on a common lattice on which their norms are periodic.

Because Λᵢ is selfdual, the alternating integral matrix of E on any Z basis of Λᵢ is unimodular. Its Pfaffian has absolute value1. Consequently both lattices have the same symplectic covolume, hence the same covolume for dz. Since Λ₀ has finite index in each, division of covolumes gives {[}Λ₁:Λ₀{]}={[}Λ₂:Λ₀{]}=q. A fundamental domain for Λ₀ is, up to null sets, the disjoint union of q translates of a fundamental domain for Λᵢ. For a Λᵢ-periodic norm density, dividing its integral by the domain volume therefore gives precisely (DO17.7). Thus the normalized norm on the common lattice equals the norm on either original lattice; there is no extra q in changing domains.

For x∈Λ₂, the function U\_xF lies in T(Λ₀,α₀), since E(x,Λ₀) is integral. The operators V\_x=α₂(x)⁻¹U\_x form an ordinary representation of the additive group Λ₂ by (DO17.6) and (DO17.9). If x∈Λ₀, then V\_xF=F. Consequently the sum (DO17.8) is independent of coset representatives and is fixed by all V\_y, y∈Λ₂. It is therefore in T(Λ₂,α₂).

The finite pairing

\[
 (\Lambda_1/\Lambda_0)\times(\Lambda_2/\Lambda_0)
       \longrightarrow U(1),\qquad
 (\lambda,x)\longmapsto e^{2\pi iE(\lambda,x)}
\tag{DO17.11}
\]

is well-defined and nondegenerate. Indeed, if x∈Λ₂ pairs trivially with every λ∈Λ₁, then selfduality of Λ₁ implies x∈Λ₁, hence x∈Λ₀. The same argument applies to the other variable. The two groups have order q, so this is a perfect finite pairing.

For distinct x,y modulo Λ₀, choose λ∈Λ₁ for which e\^{}\{2πiE(λ,x−y)\}≠1. The operator U\_λ is unitary on T(Λ₀,α₀), by (DO17.10). By (DO17.9), U\_xF and U\_yF are its eigenvectors with eigenvalues α₁(λ)e\^{}\{−2πiE(λ,x)\} and α₁(λ)e\^{}\{−2πiE(λ,y)\}. Their eigenvalues are distinct and have absolute value1. Taking their inner product and using unitarity shows that they are orthogonal. Each has norm \textbar\textbar F\textbar\textbar₁. Therefore

\[
 \|I_{21}F\|_2^2
     =q^{-1}\sum_{x\in\Lambda_2/\Lambda_0}\|U_xF\|_0^2
     =\|F\|_1^2.
\tag{DO17.12}
\]

This proves the isometry and nonvanishing.

For completeness, dimension one is the elementary principal-theta dimension statement. Choose a symplectic Z basis of Λᵢ, normalize its first half to the real standard basis, and write the second half as the columns of τ with Im τ positive. After a linear exponential gauge, the real-period relation gives a Fourier expansion indexed by Z\^{}g. The τ-period relation successively determines every Fourier coefficient from one coefficient, with the usual Gaussian factor exp(πinᵗτn) and the semicharacter shifts. Positivity gives normal convergence of that theta series. Hence the dimension is exactly one. Equivalently this is h⁰=1 for a principal polarization. ∎

The theta line on V/Λ₀ used in this proof has degree q: its symplectic covolume is q times that of a principal lattice. The q orthogonal vectors above are consistent with its q-dimensional space of sections. The argument did not confuse this degree with q², the order of its full theta kernel.

\subsubsection{Exact transport and exact residual phase}\label{exact-transport-and-exact-residual-phase}

If C:V→V′ is complex-linear with H′(Cx,Cy)=H(x,y), CΛᵢ=Λᵢ′ and αᵢ′(Cx)=αᵢ(x), then pullback by C commutes \textbf{exactly} with (DO17.8). Indeed it commutes with each U\_x, induces a bijection of the finite quotient sets, and preserves q. Thus any specified base equality with scalar c extends to the entire two-step family with that same scalar c.

For three such lattices, choose a nonzero F∈T(Λ₁,α₁). Since all three theta lines have dimension one, there is a unique phase

\[
 \mu_{321}=\frac{\langle I_{31}F,I_{32}I_{21}F\rangle_3}{\|F\|_1^2},
 \qquad |\mu_{321}|=1,
 \qquad I_{32}I_{21}=\mu_{321}I_{31}.
\tag{DO17.13}
\]

Our convention makes the inner product linear in its second argument. This phase is independent of F. It is explicitly a finite Gauss sum: expand the three averages, use (DO17.9), and use

\[
 \frac{\langle U_zF,U_tF\rangle}{\|F\|^2}
 =\begin{cases}
 e^{\pi iE(z,t)}\alpha_1(t-z),&t-z\in\Lambda_1,\\
 0,&t-z\notin\Lambda_1.
 \end{cases}
\tag{DO17.14}
\]

The zero case follows from the same perfect-pairing argument (using a common finite-index lattice for the finitely many vectors); the nonzero case follows from U\_\{−z\}U\_t=e\^{}\{πiE(z,t)\}U\_\{t−z\}. Equations(DO17.8), (DO17.9), and (DO17.14) are an exact finite prescription for μ, with all normalizations included. Neither its value1 nor its value±1 is automatic from nonvanishing. The special \nolinkurl{SignP} and \nolinkurl{OrderP} assertions in the original constructor must therefore be checked in a new constructor.

\subsection{AM-DO-018. A trace-one mixed-CM moved family with explicit finite inputs}\label{am-do-018.-a-trace-one-mixed-cm-moved-family-with-explicit-finite-inputs}

Retain the actual principally polarized J of AM-R-004. Let V=C³ in algebraic u coordinates, Λ=DΦ(O\_E), and H(u,v)=Σu\_i conjugate(v\_i)/G\_i. In the first-argument-linear convention, E\_H=Im H is the negative of the displayed Riemann intersection form in AM-R-004; changing this sign does not change selfduality. Put

\[
 w_0=-(\zeta+\zeta^2+\zeta^4),\qquad
 w_0+\bar w_0=1,
\]

acting through the actual CM action. On every w₀-stable E\_H-integral lattice define

\[
             \alpha_{\rm tr}(x)=e^{\pi iE_H(x,w_0x)}\in\{1,-1\}.
\tag{DO18.1}
\]

This is a semicharacter, and its formula agrees on intersections of lattices. It is the Appell--Humbert semicharacter of the rigidified trace-one Poincaré line L\_\{w₀\}.

\textbf{Proof.} The Rosati adjoint of w₀ is bar(w₀). Hence

\[
 E_H(x,w_0y)+E_H(y,w_0x)
       =E_H(x,(w_0-\bar w_0)y)
       =2E_H(x,w_0y)-E_H(x,y).
\]

For integral stable lattices, its exponential times πi equals exp(πiE\_H(x,y)). Expanding E\_H(x+y,w₀(x+y)) proves (DO17.6). For the line identification, the Poincaré line pulled back by the principal polarization is m\emph{L⊗p₁}L⁻¹⊗p₂*L⁻¹. Its Hermitian form on V×V is H(z,w′)+H(w,z′) and its semicharacter on (λ,μ) is exp(πiE\_H(λ,μ)), by the semicharacter identity for any representative L. Pullback by (1,w₀) has Hermitian form H, since w₀+bar(w₀)=1, and semicharacter (DO18.1). This proves the assertion with the canonical rigidifications. ∎

As an explicit sign check in the actual period basis, E\_R004(ζʲ,w₀ζʲ)=Tr(ξ bar(w₀))=−1 for each j=0,\ldots,5, with ξ=(ζ−ζ⁻¹)/7. Thus α\_tr takes value−1 on each adjacent-cycle generator γⱼ. AM-DO-012/ALP-LIFT-07's characteristic K∞=(101\textbar111) gives exactly these six signs. This checks the trace-one line against the analytic Σ\_tr convention rather than merely matching Chern classes.

\subsubsection{Matrix-lattice family}\label{matrix-lattice-family}

Fix positive integers a,b. For full O\_E-lattices N⊂E\^{}a and N′⊂E\^{}b that are selfdual for their standard Hermitian forms, let

\[
 M(N,N')=\operatorname{Hom}_{O_E}(N,N')\subset M_{b\times a}(E),
 \qquad\Lambda(N,N')=D\Phi(M(N,N')).
\tag{DO18.2}
\]

Here DΦ is applied entry by entry, and the Hermitian form is the sum of H on all ab entries. These lattices are full, commensurable, selfdual for the product E\_H, and stable under w₀. To check selfduality, localize at a rational prime and choose arbitrary module bases g,h of N,N′. Their Hermitian Gram matrices G\_N=g* g and G\_N′=h* h are integral and unimodular because the lattices are Hermitian selfdual. Writing X=hZg⁻¹ identifies the matrix pairing with the trace pairing whose Gram operator is Z↦G\_N′ Z G\_N⁻¹; this is an integral module automorphism. The scalar trace pairing is perfect on O\_E because ξO\_E is the codifferent, as proved in AM-R-004. Hence the matrix trace pairing is perfect on every local Hom lattice. Intersecting the local duals proves the global assertion. Equivalently the argument is projective-module duality Hom(N,N′)≃N*⊗N′. This includes every lattice pair obtained from the standard lattices by finite adelic unitary moves. No local classification of ramified or dyadic unimodular Hermitian lattices, and no global rational unitary basis, is assumed.

Let

\[
 F_{00}(u)=\prod_{i=1}^{b}\prod_{j=1}^{a}\Sigma_{\rm tr}(u_{ij}),
\]

and put

\[
 F_{N,N'}=I_{\Lambda(N,O_E^b),\Lambda(N,N')}
          I_{\Lambda(O_E^a,O_E^b),\Lambda(N,O_E^b)}F_{00},
\tag{DO18.3}
\]

using (DO17.8) and the common sign(DO18.1). This is an actual, nonzero holomorphic section of the corresponding principal trace-one line on V\^{}\{ab\}/Λ(N,N′). Every averaging coefficient and every finite quotient is specified. It has exactly the same probability-Haar L² norm as F₀₀.

Write

\[
 q_1=[\Lambda(N,O_E^b):\Lambda(N,O_E^b)\cap\Lambda(O_E^a,O_E^b)],\quad
 q_2=[\Lambda(N,N'):\Lambda(N,N')\cap\Lambda(N,O_E^b)].
\tag{DO18.4}
\]

The two q's are concrete integral lattice indices, computable by Smith normal form in the integral period basis. They are not arbitrary group volumes.

\subsubsection{Explicit rational Heisenberg intertwining}\label{explicit-rational-heisenberg-intertwining}

Use normalized complex coordinates z and the convention ψ∞(r)=exp(2πir), ψ\_f(r)=exp(−2πir) on Q. For t=A+τB with A,B∈Q³ define on finite Schwartz functions

\[
 (W_t\phi)(q_f)=\psi_f(-A^tq_f+A^tB/2)\,\phi(q_f-B).
\tag{DO18.5}
\]

Let Rφ(z)=e\^{}\{Q(z)\}Σ\_\{q∈Q³\}φ(q)e\^{}\{πiqᵗτq+2πiqᵗz\}. Then

\[
                     R(W_t\phi)=U_t(R\phi).
\tag{DO18.6}
\]

\textbf{Proof.} The multiplier from Q(z+t)−Q(z)−πH(z,t)−πH(t,t)/2 equals exp(2πiBᵗz+πiBᵗτB+πiAᵗB). In the theta sum for Rφ(z+t), replace q by q′−B. The remaining coefficient is φ(q′−B)exp(2πiAᵗq′−πiAᵗB), exactly(DO18.5) on rational points. Absolute convergence on compact sets permits this reindexing and all derivatives. The same calculation entry by entry proves the matrix version. ∎

For reference this map is injective on finite Schwartz functions: compact support gives a common denominator M with rational support in M⁻¹Z³; restricting z to the real cube of side M, Fourier uniqueness recovers every φ(q) times its nonzero Gaussian coefficient. The diagonal Q is dense in A\_f, so a locally constant function vanishing on Q is zero. The same holds in all matrix dimensions.

Let φ₀ be the tensor product of AM-DO-016's specified input

\[
 -i\,1_{\delta/2+\widehat{\mathbf Z}^3}(q_f)\,
             \psi_f(-q_f^t\varepsilon/2),
 \qquad\delta=(1,1,1),\ \varepsilon=(1,0,1).
\]

Replace each U\_t in (DO18.3) by W\_t, with the same α\_tr(t) and the same positive q\_i⁻¹/², to define the finite Schwartz function φ\_\{N,N′\}. Every t lies in the rational span of the standard period lattice, so (DO18.5) applies. The output is compactly supported and locally constant, with algebraic coefficients. AM-DO-016 and (DO18.6) give the \textbf{whole moved-family equality}

\[
 \boxed{R\phi_{N,N'}=\theta[0;0](0;\tau)^{ab}\,F_{N,N'}.}
\tag{DO18.7}
\]

The period scalar is the same for every N,N′ at these fixed ranks. This is more than a nonzero period or a single-fiber identity. It is an actual family of finite inputs and all-degree holomorphic theta equalities. It still does not identify φ\_\{N,N′\} with the prescribed split Weil translate of φ₀; that further assertion is addressed below.

\subsubsection{Normalization and explicit elementary costs}\label{normalization-and-explicit-elementary-costs}

Carry the exact norm from AM-R-004(R4.12), with κ=1, and the isometry of normalized Mumford translation. Then

\[
 \|F_{N,N'}\|^2
     =\det(2Y)^{-ab/2}|\theta[0;0](0;\tau)|^{-2ab}.
\tag{DO18.8}
\]

Thus the phase-preserving scalar that changes (DO18.7)'s reduced seed to norm1 is

\[
 c_{\rm unit}=\theta[0;0](0;\tau)^{ab}\det(2Y)^{ab/4},
 \qquad \|c_{\rm unit}F_{N,N'}\|=1.
\tag{DO18.9}
\]

This is an explicit local analytic normalization, not the terminal automorphic mass pairing in G1.

AM-R-006 now computes the algebraic scalar between the two actual seeds: Σ\_tr=cσ\_q, where the phase is fixed by (R6.3), c⁸=R\_S²/4⁸, v₂(c)=−2, v₇(c)=1/2, and the other finite valuations are zero. Consequently, if the same averages are applied to σ\_q\^{}\{⊠ab\}, then (DO18.7)'s common period is exactly (θ₀c)\^{}\{ab\}. Its squared norm is (DO18.8) divided by \textbar c\textbar\^{}\{2ab\}, and the phase-preserving norm-one scalar is (θ₀c)\^{}\{ab\}det(2Y)\^{}\{ab/4\}. These equalities follow from linearity of each specified average. They do not select an arbitrary algebraic root for c or absorb it silently into an unspecified period.

If normalized valuations satisfy v\_p(p)=1, the algebraic values of the constructed Schwartz input satisfy the elementary bound

\[
 v_p(\phi_{N,N'}(q_f))\ge-\tfrac12v_p(q_1q_2)
\tag{DO18.10}
\]

where a zero value has valuation +∞. Indeed the original nonzero coefficients, α\_tr signs and finite additive-character values are roots of unity; the only displayed denominators are the two inverse square roots of positive integers. The ultrametric inequality proves the claim, and cancellation may improve it. This is a statement about \textbf{finite Schwartz coefficients}, not an assertion that the corresponding reduced CM Taylor series is integral at a nonordinary prime.

At the fixed complex embedding, AM-DO-015 and triangle inequality give the completely explicit growth bound

\[
 |F_{N,N'}(u)|<49^{ab}\sqrt{q_1q_2}\,
      \exp\!\left(\frac\pi2\sum_{i,j,k}|u_{ij,k}|^2/G_k\right).
\tag{DO18.11}
\]

Each normalized U\_t preserves the Gaussian envelope, while a q-term average costs at most √q. The corresponding polydisk Cauchy bound follows exactly as in AM-R-004. This does not claim that √(q₁q₂) is optimal. L² norm preservation does not by itself eliminate this displayed pointwise cost for varying lattices.

\subsection{AM-DO-019. Exact boundary between this mixed family and the prescribed constructor}\label{am-do-019.-exact-boundary-between-this-mixed-family-and-the-prescribed-constructor}

There are two distinct questions.

\textbf{Unchanged elliptic \nolinkurl{VacIn}.} The recovered record fixes an elliptic period lattice and a complexification of an isomorphism over the quadratic CM subfield; its base fourfold kernel is built from that elliptic data. AM-R-003 proves that the primitive Φ=\{1,2,3\} Jacobian J has no nonzero homomorphism to or from an elliptic curve and no isogeny with an elliptic power. Therefore a full, line-preserving CM/Hodge transport from that unchanged elliptic constructor to J (or the corresponding powers) does not exist. This conclusion uses the actual category demanded by the record, not the false inference that absence of a canonical map excludes all maps. Arbitrary vector-space maps, selected coefficient maps, and a new mixed-CM constructor are not excluded.

\textbf{A new mixed-CM constructor.} Equations(DO18.1)--(DO18.11) give its trace-one theta seed, every moved lattice, every finite input, the exact average/covolume factors, and one common analytic period. The remaining prescribed-Weil equality is

\[
 \omega_{\psi,\chi}(g,h)\phi_0
       \stackrel{?}=\chi(x)^{-1}\chi(y)\,\phi_{N,N'},
 \quad x/\bar x=\det g,\quad y/\bar y=\det h.
\tag{DO19.1}
\]

The roots x,y are determinant roots, as in the actual source \nolinkurl{IsDetRoot}; substituting det(g),det(h) is generally wrong. The source also fixes ψ∞ with the opposite sign from (DO18.5); its polarization/character transport must be included before testing(DO19.1). In a common realized Heisenberg representation the two sides lie on the same one-dimensional lattice vacuum line, so their quotient is a phase once both are norm1. Its value must be computed from the chosen splitting on generators, or from the finite Gauss sum(DO17.13) together with the specified determinant-root character. Merely saying that both vectors are nonzero gives no equality.

The exact recovered source makes this distinction visible. \nolinkurl{MovedWeilVacuumEqVacuumOfMovedLattices.lean}, theorem \nolinkurl{vac_moved}, proves preservation of the moved lattice vacuum, theta membership and nonvanishing. It does \textbf{not} conclude the common c₀/determinant-root formula. \nolinkurl{LevChiEqOne.lean}, lines87--101, places that formula in \nolinkurl{VacT} and \nolinkurl{PinX.mv}. \nolinkurl{ThetaKernelSplitPlacesEqJetClassicalTheta.lean}, theorem \nolinkurl{thetaKernel_eq_jet_thC}, takes the equation \nolinkurl{xiN(om(g,h)vac)=κE} as an explicit hypothesis and then transports it to every polynomial jet. Thus the missing phase is exactly where the source says it is.

A failure to prove(DO19.1) is not a proof that it is false. Conversely, (DO18.7) does not prove it by construction, because the φ\_\{N,N′\} were defined through the two fixed finite averages, while the source fixes a group representation and its splitting independently. Nor do these analytic equalities prove the full integral row-lattice identities, the nonordinary formal-gauge comparison, the \nolinkurl{MainX} cohomological/topological bounds, or the terminal-mass equation from THETA3 G1. Those remain their stated separate gates.

\subsection{AM-DO-020. The phase is exactly the splitting character for relative split-place moves}\label{am-do-020.-the-phase-is-exactly-the-splitting-character-for-relative-split-place-moves}

The previously open phase in(DO19.1) can be resolved in a substantial precise scope: moves supported on finite places of F that split in E. This includes the primes above11 and29 for E=Q(ζ₇), and also the relative split dyadic places. This assertion is conditional on a Weil representation with the displayed covariance, unitarity and split-normalization axioms; it proves the required phase from those axioms, rather than replacing the prescribed representation by a new family of unrelated operators. It does not extend the old quadratic \nolinkurl{StdPair.xi} pin to the primitive type.

\subsubsection{Local statement}\label{local-statement}

Let k be a nonarchimedean field of characteristic zero, O its integers, and ψ a nontrivial additive character. Put E\_v=k×k with conjugation interchanging factors. Choose the split Schrödinger polarization of the rational symplectic Hom space. After identifying the momentum space with the ψ-dual of the position space X, write its form and Heisenberg action as

\[
 B((x,y),(x',y'))=\langle y,x'\rangle-\langle y',x\rangle,
 \qquad
 [\rho(x,y)\phi](t)=\psi(-\langle y,t\rangle+\langle y,x\rangle/2)
                           \phi(t-x).
\tag{DO20.1}
\]

Let L⊂X be a full O-lattice and L* its annihilator under ψ. The split O\_E-lattice is Λ\_L=L⊕L*. An integral trace-one endomorphism is w₀=(w,1−w), with w,1−w∈O. Its local semicharacter is

\[
 \alpha_L(x,y)=\psi\bigl((w-\tfrac12)\langle y,x\rangle\bigr),
       \qquad (x,y)\in L\oplus L^*.
\tag{DO20.2}
\]

Here the pairing is taken in the additive-character conductor convention, so its integral values and their O-multiples are killed by ψ. Formula(DO20.2) is the local factor of α\_tr=exp(πiE\_H(x,w₀x)): globally α\_tr=ψ\_f(B(x,w₀x)/2), using E\_H=−B and the product formula for the additive character. It follows immediately that

\[
 \alpha_L(x,y)^{-1}\rho(x,y)\phi(t)
                   =\psi(-\langle y,t\rangle)\phi(t-x)
       \quad\text{for }(x,y)\in\Lambda_L.
\tag{DO20.3}
\]

Indeed the extra scalar is ψ((1−w)〈y,x〉)=1. This cancellation is valid at residue characteristic two as well. Dropping the half-pairing term without the trace-one semicharacter would not justify that dyadic assertion.

The common eigenspace

\[
 \rho(\lambda)\phi=\alpha_L(\lambda)\phi
                     \quad(\lambda\in\Lambda_L)
\tag{DO20.4}
\]

is precisely the line spanned by1\_L. For y∈L\emph{, the modulation relation forces the support into (L})*=L; for x∈L, the translation relation makes the function constant there. Conversely1\_L satisfies the relations by(DO20.3). With any fixed Haar measure dt on X, its positive norm-one generator is

\[
                         v_L=\operatorname{vol}(L)^{-1/2}1_L.
\tag{DO20.5}
\]

For two position lattices L₀,L, the averaging operator(DO17.8), realized by ρ, satisfies the exact equality

\[
                             I_{L,L_0}v_{L_0}=v_L.
\tag{DO20.6}
\]

Our index notation in(DO20.6) means that the output lattice is L; equivalently I\_\{21\}v\_\{L₁\}=v\_\{L₂\}.

\subsubsection{Proof of the exact index and value in(DO20.6)}\label{proof-of-the-exact-index-and-value-indo20.6}

Let M=L∩L₀, N=L+L₀, u={[}L:M{]}, and v={[}L₀:M{]}. Duality gives

\[
 L^*\cap L_0^*=N^*,\quad
 [L^*:N^*]=[N:L]=[L_0:M]=v.
\]

Therefore

\[
 q=[\Lambda_L:\Lambda_L\cap\Lambda_{L_0}]=uv,
 \qquad \operatorname{vol}(L)/\operatorname{vol}(L_0)=u/v.
\tag{DO20.7}
\]

Use representatives x∈L/M and y∈L\emph{/N} in the finite average. By(DO20.3), its value on1\_\{L₀\} at t is

\[
 (uv)^{-1/2}\sum_{x\in L/M}\sum_{y\in L^*/N^*}
               \psi(-\langle y,t\rangle)1_{L_0}(t-x).
\tag{DO20.8}
\]

The possible support is N. For t∈N, finite character orthogonality makes the inner sum equal v when t∈L and zero otherwise. For t∈L there is exactly one class x∈L/M with t−x∈L₀. Hence(DO20.8) equals √(v/u)1\_L(t). Multiplying by vol(L₀)⁻¹/² gives(DO20.6). In particular its value at position0 is strictly positive. This proof also gives

\[
 I_{L_3,L_2}I_{L_2,L_1}=I_{L_3,L_1}
\tag{DO20.9}
\]

on the corresponding lattice-vacuum line. Thus the triangle phase is1 and the two orders of moving the source and target Hom-lattice factors agree in this split scope.

\subsubsection{Comparison with the independently prescribed split Weil action}\label{comparison-with-the-independently-prescribed-split-weil-action}

At a split place the unitary pair is GL\_a(k)×GL\_b(k). On its position Hom space, (g,h) acts by A:Z↦hZg⁻¹. Its determinant is

\[
                    \det_k A=(\det h)^a(\det g)^{-b}.
\tag{DO20.10}
\]

The explicit unitary Schrödinger change-of-variables operator

\[
 T_A\phi(t)=|\det A|_k^{-1/2}\phi(A^{-1}t)
\tag{DO20.11}
\]

has the requisite covariance, and sends v\_\{L₀\} to v\_\{AL₀\}. Its positive value at0 is vol(AL₀)⁻¹/².

Let ω(g,h) be the prescribed unitary Weil action and let its split-normalization character be

\[
             \chi_T(x)^{-1}\chi_S(y),\qquad
               x/\bar x=\det g,\quad y/\bar y=\det h.
\tag{DO20.12}
\]

Covariance transports(DO20.4), so ω(g,h)v\_\{L₀\}=λv\_\{AL₀\} for some λ; the elementary uniqueness of that line above suffices, without a further abstract irreducibility assertion. Unitarity gives \textbar λ\textbar=1. The functional ℓ(φ)=φ(0) is invariant under all momentum Heisenberg operators. It is therefore exactly a functional to which the source \nolinkurl{om_split} condition applies. That condition gives

\[
 \ell(\omega(g,h)v_{L_0})
       =r\,\chi_T(x)^{-1}\chi_S(y)\,\ell(v_{L_0}),\qquad r>0.
\tag{DO20.13}
\]

Both ℓ(v\_\{L₀\}) and ℓ(v\_\{AL₀\}) are positive. Comparing(DO20.13) with λv\_\{AL₀\}, and using \textbar λ\textbar=1, proves

\[
 \boxed{\omega(g,h)v_{L_0}
           =\chi_T(x)^{-1}\chi_S(y)\,v_{AL_0}.}
\tag{DO20.14}
\]

The positive scalar r is forced to be

\[
                    r=\bigl(\operatorname{vol}(L_0)/
                                  \operatorname{vol}(AL_0)\bigr)^{1/2}
                      =|\det A|_k^{-1/2}.
\tag{DO20.15}
\]

This computes both the phase and the positive normalizing factor. It does not silently discard the r in the source's split axiom.

\subsubsection{Global consequence for the mixed family}\label{global-consequence-for-the-mixed-family}

Apply this place by place to finite adelic g,h supported at relative split places. The finite quotients defining the global averages decompose into their primary components. The character and Heisenberg factors decompose in the same way; on the common lattice their weighted translation operators are an ordinary additive-group representation, so the local factors commute. The positive square-root normalization is the product of the positive local square roots. Consequently the two-step input φ\_\{N,N′\} of(DO18.7) is exactly the product of the positive local lattice-vacuum vectors at the moved places, with the original factors elsewhere. Equations(DO20.6) and(DO20.14) prove(DO19.1) in this scope, with its determinant-root character and the common scalar θ₀\^{}\{ab\} retained by(DO18.7).

For the even ranks of \nolinkurl{StdPair}, the restriction of χ\_S and χ\_T to A\_F\^{}× is η\_\{E/F\}\^{}\{even\}=1. Therefore changing determinant roots by fixed-field ideles does not change(DO20.12). For arbitrary ranks the statement uses the specified supported determinant roots; one must retain the splitting's stated root dependence.

There is an explicit way to align additive-character signs with the recovered source. Use its ψ\_F=(ψ\_Q⁻¹)∘Tr\_\{F/Q\}, whose real sign is−1, and set ξ₀=−ξ in \nolinkurl{symplO}, where ξ=(ζ−ζ⁻¹)/7. Then ψ\_F(Tr\_\{E/F\}(ξ₀h))=ψ\_Q(Tr\_\{E/Q\}(ξh)), the actual Heisenberg central character used above. This is a simultaneous form/character transport, so it leaves the realized representation unchanged. It changes the old constructor's fixed quadratic ξ₀ condition, which is already part of the mixed-CM constructor gate. It does not assert that the unchanged source \nolinkurl{StdPair.xi} is satisfied.

\subsubsection{A concrete rank-two to rank-four moved lattice}\label{a-concrete-rank-two-to-rank-four-moved-lattice}

Take a=2, b=4, π=2−ζ, and r=π/bar(π). The exact norm is

\[
 N_{E/\mathbf Q}(\pi)=\Phi_7(2)=127.
\]

The primes (π) and (bar π) are distinct: ζ is2 modulo(π), while its inverse is64, and2≠64 in F₁₂₇. Put

\[
              N=rO_E\oplus O_E,\qquad N'=O_E^4.
\tag{DO20.16}
\]

The lattice N is Hermitian selfdual since r bar(r)=1. Its Hom lattice has four entries r⁻¹O\_E and four entries O\_E. For one changed entry,

\[
 O_E\cap r^{-1}O_E=\bar\pi O_E,
 \qquad [O_E:\bar\pi O_E]=[r^{-1}O_E:\bar\pi O_E]=127.
\]

Hence the exact global indices in(DO18.4) are

\[
                         q_1=127^4,\qquad q_2=1.
\tag{DO20.17}
\]

One may choose the finite adelic unitary element g with components diag(r,1) at127 and identity elsewhere, h=1. Its support is relative split because127≡1 mod7. A determinant root is the finite idele whose127-components are π and whose other components are1. For trivial splitting characters at these even ranks, the character in(DO20.14) is1, so the actual split Weil translate equals the explicitly averaged input, with no residual phase. The finite average has127⁴ terms and exact coefficient127⁻², before signs and Heisenberg characters. Its elementary coefficient denominator bound is v₁₂₇≥−2. Its fixed-embedding Gaussian bound is49⁸·127². Its exact analytic period in(DO18.7) isθ₀⁸, and its norm-one comparison scalar isθ₀⁸ det(2Y)². These are concrete values for this moved kernel family, not a terminal mass or a proof that a selected rank-four row lies in a global theta-lift image.

\subsubsection{What fails to transport at the ramified prime7}\label{what-fails-to-transport-at-the-ramified-prime7}

The recovered source phase proof is more specific than a generic metaplectic uniqueness argument. \nolinkurl{.../KernelLattices/ColumnData/ColumnDataEquations/MuFinal.lean}, theorems \nolinkurl{muOneL_of} and \nolinkurl{muOneR_of}, uses \nolinkurl{J.negOneNorm D (4*e)} to find an integral l with

\[
                         l\bar l\equiv-1\pmod{4e}.
\tag{DO20.18}
\]

Together with an even-rank row/column symmetry this makes the source's triangle factor1. The source \nolinkurl{SplitPairOfGen.lean}, theorem \nolinkurl{splitPair_of_gen}, explicitly obtains \nolinkurl{IsCMUnram} from \nolinkurl{HostDer.unramEF}.

For E=Q(ζ₇), (DO20.18) has no solution when7\textbar e. Reduction at the unique prime above7 makes conjugation the identity on F₇, so the congruence would say l²=−1 in F₇. The squares are0,1,2,4 and do not contain6. This disproves that auxiliary congruence in the mixed example and prevents directly reusing that particular proof. It does not disprove the desired phase identity, which might be obtained with a different normalization or a more precise ramified computation. The present calculation closes split-place phases and leaves a sharply located nonsplit/ramified phase check, together with the independent global integral and terminal-mass gates.

\clearpage

\section{AM-DO-021: an exact ramified compact character and a rank-two to rank-four counterexample}\label{am-do-021-an-exact-ramified-compact-character-and-a-rank-two-to-rank-four-counterexample}

Author: GPT-6 Astra Pro, ALP\_MACHINE \nolinkurl{/root/definite_oscillator}; inherited reasoning effort not separately exposed; 21 November 2026 UTC.

This result uses the standard unitary Weil splitting attached to characters χ\_S,χ\_T in Kudla's normalization, with Pan's explicit comparison to the lattice model. It computes the missing scalar at the ramified prime 7. It does not infer that every representation satisfying only some of the source's listed axioms has already been identified with this standard splitting. The split-place proof AM-DO-020 did not require this additional identification.

\subsection{Statement}\label{statement-7}

Character labels use the source convention \nolinkurl{LocalNewformFunction.splitFacO}, whose source factor is \(\chi_T(x)^{-1}\) and target factor is \(\chi_S(y)\) on determinant roots. Thus conventional Harris--Kudla--Sweet tensor parameters are \((\chi_S,\chi_T^{-1})\); corrected AM-DO-023 records the same conversion at infinity. In particular every trivial-character assertion below is independent of that labeling choice.

Let E=Q(ζ₇), F=E⁺, and let v be the unique place of F above 7. Put K=F\_v, L=E\_v, O=O\_K, O\_L=O\_\{E\_v\}, and s=ζ₇−ζ₇⁻¹. Conjugation on L is written with a bar. The residue fields of K and L are both F₇, and s is an anti-invariant uniformizer of L. Let ξ=s/7.

For positive even integers a,b, take the standard Hermitian lattices O\_L\^{}a,O\_L\^{}b and their Hom lattice Λ=M\_\{b×a\}(O\_L). Use the actual rational symplectic form of AM-DO-018 and the corresponding additive character. Let φ₀ span the trace-one Heisenberg vacuum line of Λ. Let

\begin{verbatim}
K_a=U(I_a)(O),       K_b=U(I_b)(O).
\end{verbatim}

For t∈L¹=\{t:t bar(t)=1\}, define

\[
 \epsilon(t)=\begin{cases}1,&t\equiv1\pmod{s},\\-1,&t\equiv-1\pmod{s}.
 \end{cases}
\tag{DO21.1}
\]

Every norm-one t is a unit and has residue ±1, so this is a continuous quadratic character. Let χ\_S,χ\_T be the unitary splitting characters; at these even ranks their restrictions to K× are trivial. For determinant roots x/bar(x)=det g, y/bar(y)=det h, the exact scalar is

\[
 \boxed{\omega_{\psi,\chi}(g,h)\phi_0
   =\chi_T(x)^{-1}\chi_S(y)
      \epsilon(\det g)^{b/2}\epsilon(\det h)^{a/2}\phi_0,
       \quad g\in K_a,\ h\in K_b.}
\tag{DO21.2}
\]

In particular, take a=2,b=4, g=1 and h=diag(−1,1,1,1). Both lattices are unchanged. Both prescribed finite averages in AM-DO-018 are therefore the identity. Nevertheless, with s supported at the place v and y=−2s as a determinant root at v,

\[
               \omega(1,h)\phi_0=-\chi_S(y)\phi_0.
\tag{DO21.3}
\]

Thus the unchanged two-average formula (DO19.1), whose right side is +χ\_S(y)φ₀, is false for this actual standard splitting at these ranks. This is a specified nonzero-vector counterexample, not an argument from absence of a canonical map. Twisting a splitting character changes both sides by the same determinant-root factor, so it does not remove the residual minus sign.

For this fixed standard-lattice construction, agreement with the prescribed root factor on the \emph{entire compact pair} is equivalent to

\[
                            4\mid a\quad\text{and}\quad4\mid b.
\tag{DO21.4}
\]

Alternatively it holds on K\_a⁺×K\_b⁺, where K\_n⁺ consists of elements whose determinant has residue +1. Each K\_n⁺ has index 2. This changes the level or adds an orientation to the lattice data; it is not equality on the unchanged compact group. The new index is a 7-adic unit and an 11-adic unit. No conclusion about an iterated terminal mass or the net step-8 lift count follows without the other recorded arithmetic identities.

\subsection{Proof}\label{proof-2}

\subsubsection{1. Identification of the actual trace lattice with Pan's lattice}\label{identification-of-the-actual-trace-lattice-with-pans-lattice}

Use the additive character ψ\_K=ψ\_Q∘Tr\_\{K/Q₇\}, where ψ\_Q has conductor Z₇. The different of F/Q has v-adic exponent 2 (the discriminant is 7²), so ψ\_K has conductor exponent λ\_K=−2 in the convention that ψ\_K is trivial on p\_K\^{}\{λ\_K\}. The relative extension L/K is tamely ramified quadratic. The element s has L-valuation 1, and 7 has L-valuation 6; hence v\_L(−2ξ)=−5.

Pan (2001), §3.1, equation (3.1.a), uses

\[
 \langle\!\langle\ ,\ \rangle\!\rangle
       =\tfrac12\operatorname{Tr}_{L/K}(h\otimes\overline{h'}).
\tag{DO21.5}
\]

Take V=L\^{}b with h=I\_b and V′=L\^{}a with skew-Hermitian form h′=−2ξI\_a. Conjugating the source tensor factor identifies V⊗\emph{L V′ with the Hom realization used in AM-DO-009 and AM-DO-018. Then (DO21.5) is exactly Tr}\{L/K\}(ξ h\_Hom). The finite operators (DO18.5) instead have Heisenberg cocycle ψ\_f(E\_H/2), where E\_H is the negative of that trace form. Accordingly use ψ\_K\^{}\{-1\} with Pan's displayed trace form, or equivalently negate that form and keep ψ\_K. These two descriptions realize the same central character and operators. Lemmas 7.2--7.3 below give the same displayed αβ formulas for every nontrivial ψ\_K, so the computed determinant-root factor and residual real sign are unchanged.

Set the duality exponents of Pan §1.3, equation (1.3.a), to κ=0 and κ′=−5. Both standard O\_L lattices are selfdual for these conventions: for I\_b the duality ideal is O\_L, and for −2ξI\_a it is s\^{}\{-5\}O\_L. Moreover

\[
                   \kappa+\kappa'=-5=2\lambda_K-1,
\tag{DO21.6}
\]

which is exactly Pan's compatibility condition (3.1.b) in the ramified case. His good lattice (3.2.c) is therefore

\[
 B(L_0,L_0')=L_0^*\otimes L_0'\cap L_0\otimes L_0'^*
                  =O_L^{ab}=\Lambda.
\tag{DO21.7}
\]

It is selfdual for the trace symplectic form and ψ\_K. Since 2 is a unit at 7, the trace-one local character ψ\_K(B(x,w₀x)/2) is trivial on Λ: w₀ preserves Λ, B is integral in the additive-character convention, and its conductor ideal is stable under multiplication by 1/2. Thus φ₀ is exactly the ordinary selfdual-lattice vacuum in Pan's model.

For an element preserving Λ, the lattice operator M\_B is pullback by its inverse (Pan §2.3, equation (2.3.e), or the selfdual formula in §2.4). It fixes the vacuum supported on Λ. The actual split Weil operator is β\^{}Y times the Schrödinger operator, and Pan (3.3.a) compares the latter with α M\_B. Therefore its scalar on this vacuum is exactly αβ\^{}Y. This identifies the normalization being computed; one cannot replace that scalar by the value of the spherical symplectic splitting without this comparison.

\subsubsection{2. The integral Witt decomposition of the identity forms}\label{the-integral-witt-decomposition-of-the-identity-forms}

The two-dimensional Hermitian form I₂ over L/K is anisotropic. Indeed a nonzero isotropic vector would give N\_\{L/K\}(z)=−1 after dividing by one nonzero coordinate. Such a norm must have valuation zero, so z is a unit. Conjugation is the identity on the residue field, hence the residue of its norm is a square in F₇. The residue −1=6 is not a square. This is a contradiction.

The four-dimensional symmetric form I₄ over F₇ is the sum of two hyperbolic planes. For a direct elementary check, (1,1,1,2) is a nonzero isotropic vector. A vector pairing nontrivially with it can be adjusted by a multiple of it, using that 2 is invertible, to give a hyperbolic pair. The orthogonal binary complement has determinant −1 modulo squares, because det I₄=1 and a hyperbolic plane has determinant −1. A binary nondegenerate symmetric form is isotropic precisely when the negative of its determinant is a square: this follows immediately after orthogonal diagonalization by solving A x²+B y²=0. The complement is therefore another hyperbolic plane.

This mod-7 isometry lifts to O. To see this without assuming a lifting classification, suppose an invertible matrix C\_r realizes the isometry mod p\_K\^{}r. Correct C\_r by C\_r(1+p\_K\^{}r X). The first-order correction of the Gram matrix is XᵗJ+JX. For every symmetric error matrix S, choosing X=−J\^{}\{-1\}S/2 solves this linear equation. Inducting and taking the p\_K-adic limit gives an integral invertible isometry. Consequently I₄ is integrally isometric to H⊥H over O, and hence over O\_L. By taking orthogonal sums,

\[
 I_{4k}\simeq H^{\perp 2k},\qquad
 I_{4k+2}\simeq H^{\perp 2k}\perp I_2,
\tag{DO21.8}
\]

with the standard lattice decomposed as the corresponding integral orthogonal sum. Multiplying all forms by −2ξ preserves this decomposition and the lattice admissibility. These are precisely the decompositions needed for Pan's tensor-product identity (3.6.e).

\subsubsection{3. The compact scalar on a determinant section}\label{the-compact-scalar-on-a-determinant-section}

First let the moving group act on a single nondegenerate Hermitian or skew-Hermitian line, and write its element as t=x₀+s y₀∈L¹, t≠1. Pan Lemma 7.2 (printed p.207) states that, against a hyperbolic plane,

\[
                        \alpha(t)\beta^Y(t)=\chi((t-1)s).
\tag{DO21.9}
\]

Pan Lemma 7.3 (printed p.208) gives, against an anisotropic plane,

\[
 \alpha(t)\beta^Y(t)=
 \begin{cases}
   -\chi((t-1)s),&x_0\not\equiv1\pmod{\mathfrak p_K},\\
   +\chi((t-1)s),&x_0\equiv1\pmod{\mathfrak p_K}.
 \end{cases}
\tag{DO21.10}
\]

At a ramified quadratic place, a norm-one element has residue +1 or −1. Its residue agrees with that of x₀, so the extra factor in (DO21.10) is exactly ε(t). Also

\[
       \frac{(t-1)s}{\overline{(t-1)s}}
        =\frac{(t-1)s}{(t^{-1}-1)(-s)}=t.
\tag{DO21.11}
\]

Thus the argument of χ in these formulas is the actual determinant root, not the determinant t. Since χ\textbar K×=1, its value is independent of the chosen root. At t=1 the root may be chosen to be 1 and the formula follows from the identity operator.

Use (DO21.8) for the opposite even-dimensional space. Choose splitting characters on the hyperbolic and anisotropic summands with product χ; all but one may be trivial because every summand has even dimension. Pan (3.6.e) multiplies their αβ factors. Therefore the scalar for the one-dimensional moving subgroup is

\[
                   \chi(y)\epsilon(t)^{d/2},
       \qquad y/\bar y=t,
\tag{DO21.12}
\]

where d is the opposite even rank. For d divisible by 4 the anisotropic summand is absent; for d congruent to 2 modulo 4 it occurs exactly once. This proves (DO21.12), including its phase, for diagonal elements diag(t,1,\ldots,1) of the compact unitary group.

\subsubsection{4. Extension to the full compact group}\label{extension-to-the-full-compact-group}

Pan's Proposition 3.3 (printed pp.179--180) says that the scalar αβ\^{}Y on the compact special unitary group equals the character ζ defined there. When the opposite rank is even, ζ is trivial. Thus αβ\^{}Y factors through the determinant on the compact unitary group. The determinant map is surjective onto L¹, since t↦diag(t,1,\ldots,1) is a section and every norm-one element of L is integral. Formula (DO21.12) consequently determines its value on the entire compact group.

For the target group U(I\_b), d=a and this yields χ\_S(y)ε(det h)\^{}\{a/2\}. For the source group U(I\_a), the Hom action is contragredient: its determinant on the source factor is (det g)\^{}\{-1\}. The corresponding root is x\^{}\{-1\}, the splitting factor is χ\_T(x)\^{}\{-1\}, and ε((det g)\^{}\{-1\})=ε(det g). Hence its contribution is χ\_T(x)\^{}\{-1\}ε(det g)\^{}\{b/2\}. The two commuting unitary actions multiply. This proves (DO21.2).

For a=2,b=4 and h=diag(−1,1,1,1), ε(det h)=−1, proving (DO21.3). The trace-one seed is nonzero (AM-DO-015), so its finite vacuum is nonzero by the injectivity in AM-DO-018. The equality cannot hold with the positive sign. Both lattice indices are 1 since h is integral unitary. Finally the determinant-section elements with determinant −1 exist in each compact group; thus vanishing of both residual characters is equivalent to evenness of a/2 and b/2. This proves (DO21.4). The kernels K\_n⁺ have index exactly 2 by the same determinant section.

\subsection{Scope and next gate}\label{scope-and-next-gate}

The theorem separates two questions that were previously conflated. The mixed CM trace lattice and the finite averaging construction exist, and their relative split-place phases are exactly solved in AM-DO-020. The standard ramified compact phase is now also computed. For ranks 2→4 the literal lattice-indexed phase formula fails; at ranks divisible by 4, or after the stated compact restriction, this particular obstruction disappears. Neither alternative by itself proves the noncompact ramified move formula, the rational compatibility of the fixed two-step averages, an all-polynomial arithmetic \nolinkurl{MainX} identity, or a terminal pairing value. Those remain separate equations.

An arbitrary choice of a local character twist is not a way to change (DO21.3) while retaining the same determinant-root pins: the twist appears on both sides. A new constructor could add orientation/level information to its lattice objects or alter its arithmetic column transformation law. Such a construction must be specified and checked; this theorem does not exclude it, nor does it exclude actual nonzero unitary theta lifts (including the first torus lift being constructed in ALP-LIFT).

\subsection{Primary source and check status}\label{primary-source-and-check-status}

Pan, S.-Y. (2001). Splittings of the metaplectic covers of some reductive dual pairs. \emph{Pacific Journal of Mathematics, 199}(1), 163--226. \url{https://msp.org/pjm/2001/199-1/pjm-v199-n1-p10-p.pdf} . Exact locators used: §1.1 conductor convention; §1.3 equations(1.3.a--c); §2.3 equation(2.3.e) and §2.4; §3.1 equations(3.1.a--b); §3.2 equation(3.2.c); §3.3 equation(3.3.a) and Proposition; §3.6 equations(3.6.a,e); Lemmas7.2--7.3, printed pp.207--208. These statements and the displayed proofs of Lemmas7.2--7.3 were directly read. Proposition3.3 is applied as a published theorem, not advertised as independently re-proved here.

\clearpage

\section{ALP-LIFT-09: an actual first compact unitary theta lift from the primitive CM Jacobian}\label{alp-lift-09-an-actual-first-compact-unitary-theta-lift-from-the-primitive-cm-jacobian}

Run ALP\_MACHINE; \nolinkurl{/root/lift_accounting}, GPT-6 Astra Pro; inherited reasoning effort not separately exposed. This finding uses the actual period matrix and reduced theta of AM-R-004, the trace-one section of ALP-LIFT-07--08, and AM-DO-017--018's finite Heisenberg input. Those inputs retain their own proof/check records. The arithmetic class calculation below was also obtained independently by the root agent; it is recorded once here, as agreed. This is a genuine first theta lift, not a \nolinkurl{thetaG_exists} construction with arbitrary operators.

\subsection{The result and its boundary}\label{the-result-and-its-boundary}

Put \(E=\mathbf Q(\zeta_7)\), \(F=E^+\),

\[
 T=\operatorname{U}(1)(E/F),\qquad
 K_0=\{k\in(\widehat{\mathcal O}_E)^\times:k\bar k=1\}.
 \tag{V1}
\]

The following statements hold.

\begin{enumerate}
\def\labelenumi{\arabic{enumi}.}
\tightlist
\item
  \(T(\mathbb A_{F,f})=T(F)K_0\), and \(T(F)\cap K_0=\mu_{14}\). Consequently \([T]\simeq(T(F_\infty)\times K_0)/\mu_{14}\) as a compact group. For an open subgroup \(K\subset K_0\) with \(K\cap\mu_{14}=1\), its finite component count is exactly \(h_K=[K_0:K]/14\).
\item
  For the actual trace-one finite Schwartz vacuum of the Jacobian, the whole \(K_0\) preserves its one-dimensional eigenline. For any fixed honest global unitary Weil splitting, its action is an actual finite-image character \(\lambda_f\) of \(K_0\).
\item
  Let a homogeneous Fock jet have archimedean character \(\beta_\infty\), including the vacuum shift for that same splitting. If its actual theta value is nonzero, then \(\beta_\infty\lambda_f\) descends to one automorphic character \(\chi\) of \(T\). With probability Haar measure on \([T]\), its theta lift is exactly the original theta function on the target group. There is no finite averaging denominator.
\item
  This gives an explicit nonzero example for the totally definite target \(\operatorname U(2)\): take two copies of the actual Jacobian vacuum and take the \(u_1u_3\) derivative in each copy. The lift at the identity is \((\theta[0;0](0;\tau)c)^2\ne0\), where \(\Sigma_{\mathrm{tr}}=c\sigma_q\) and \(\partial_1\partial_3\sigma_q(0)=1\).
\item
  For an arbitrary finite Schwartz span fixed by a common \(K\), one fixed finite-coset automorphic input realizes every theta function of one fixed archimedean torus weight. Its exact squared \(L^2([T])\) norm is \(h_K\). This is the appropriate generalization when Heisenberg translations do not preserve the original lattice eigencharacter.
\end{enumerate}

The splitting, Schwartz vector, automorphic character and theta integral below are actual objects. However, this result does not identify a chosen character with an independently prescribed source character, instantiate the old elliptic \nolinkurl{VacIn}, construct an entire iterated tower, or prove its six \nolinkurl{AdmS} identities and terminal mass equation. The target automorphic \(L^2\) norm and the whole integral row lattice are not computed here. In particular, the displayed nonzero first lift is not a proof that Q1's requested lift count can be shortened.

\subsection{1. The norm-one torus has one maximal-integral finite class}\label{the-norm-one-torus-has-one-maximal-integral-finite-class}

We give the arithmetic proof, rather than assume a torus class-number formula. The basic inputs are Minkowski's ideal-class bound and Dirichlet's unit theorem (Milne, 2020, Theorems 4.3 and 5.1).

The cyclotomic discriminant is \(\lvert d_E\rvert=7^5\). Its Minkowski bound satisfies

\[
 \left(\frac4\pi\right)^3\frac{6!}{6^6}\sqrt{7^5}<5.
 \tag{V2}
\]

For an entirely rational check, use \(\pi>3\) and

\(\sqrt7<8/3\); the resulting upper bound is \(31360/6561<5\). The order of \(2\pmod7\) is three and the order of \(3\pmod7\) is six. Thus every prime ideal over 2 has norm 8 and every prime ideal over 3 has norm \(3^6\). There are no nontrivial integral ideals of norms 2, 3 or 4. Each ideal class has an integral representative of norm at most 4 by (V2), so

\[
                    \operatorname{Cl}(\mathcal O_E)=1.
 \tag{V3}
\]

For later reference, \(d_F=49\) gives Minkowski bound \(14/9<2\), so \(\operatorname{Cl}(\mathcal O_F)=1\) too; the latter is not needed to prove (V5).

Let \(a=\zeta+\zeta^{-1}\) and \(b=\zeta^2+\zeta^{-2}\). They are units, since their common minimal polynomial is

\(X^3+X^2-2X-1\). At the embeddings indexed by \(1,2,3\), their signs and the signs of \(-1\) are

\begin{longtable}[]{@{}ll@{}}
\toprule\noalign{}
Unit & Signs \\
\midrule\noalign{}
\endhead
\bottomrule\noalign{}
\endlastfoot
\(a\) & \(+,-,-\) \\
\(b\) & \(-,-,+\) \\
\(-1\) & \(-,-,-\) \\
\end{longtable}

These three sign vectors generate \(\{\pm1\}^3\): their negative-sign vectors \(011,110,111\) are independent over \(\mathbf F_2\). Dirichlet gives \(\#(\mathcal O_F^\times/(\mathcal O_F^\times)^2)=8\), because \(F\) is totally real of degree three and its roots of unity are \(\{\pm1\}\). The surjective signature map between these two groups of order 8 is an isomorphism. Therefore

\[
 \mathcal O_F^{\times,+}=(\mathcal O_F^\times)^2.
 \tag{V4}
\]

Take \(t\in T(\mathbb A_{F,f})\). Its fractional \(\mathcal O_E\)-ideal \(I(t)\) satisfies \(I(t)\overline{I(t)}=\mathcal O_E\). By (V3), write

\(I(t)=a_0\mathcal O_E\). Then \(a_0\bar a_0\) is a totally positive unit of \(F\). By (V4), it is \(b_0^2\) for a unit \(b_0\in\mathcal O_F^\times\). Thus \(\gamma=a_0/b_0\in T(F)\) has the same fractional ideal as

\(t\), and \(\gamma^{-1}t\in K_0\). This proves

\[
                 T(\mathbb A_{F,f})=T(F)K_0.
 \tag{V5}
\]

A rational element in \(K_0\) is an algebraic unit whose conjugates all have absolute value one. The kernel of the logarithmic unit map is the roots of unity, so it is in \(\mu(E)=\mu_{14}\). Conversely every such root lies in \(T(F)\cap K_0\). The assertion about \(\mu(E)\) also follows directly from the ramification of the cyclotomic field: an additional prime in the root's order, a factor 4, or a factor \(7^2\) would force an excluded ramified prime or a subfield of degree greater than six.

It follows that the natural map \(T(F_\infty)\times K_0\to[T]\) is onto, with kernel the diagonally embedded \(\mu_{14}\). It is an isomorphism after taking that quotient: the source is compact and the target Hausdorff. In particular, the automorphic quotient is compact without an appeal to reduction theory. If \(K\cap\mu_{14}=1\), then

\[
 T(F)\backslash T(\mathbb A_F)/(T(F_\infty)K)
       =K_0/(\mu_{14}K),\qquad
 h_K=[K_0:K]/14.
 \tag{V6}
\]

Such levels exist. For example, intersect any given open subgroup with the principal congruence subgroup at 3. Roots of order dividing 14 inject under reduction at 3, so this kills the intersection with \(\mu_{14}\). No unknown torus class representatives enter (V6).

\subsection{2. A completely specified honest global splitting exists here}\label{a-completely-specified-honest-global-splitting-exists-here}

Let \(V=E^n\) with its positive definite Hermitian form \(\sum x_i\bar y_i\), and let \(W=E\) be the skew-Hermitian line with coefficient \(-\xi\), where \(\xi=(\zeta-\zeta^{-1})/7\). On restriction to \(\mathbf Q\), the tensor alternating form is

\[
 \mathfrak b(x,y)=-\operatorname{Tr}_{E/\mathbf Q}
       \left(\xi\sum x_i\bar y_i\right).
 \tag{V7}
\]

This is precisely the form \(E_H=\operatorname{Im}H\) in AM-DO-018. Use the global additive character with \(\psi_\infty(r)=e^{2\pi ir}\) and \(\psi_f(r)=e^{-2\pi ir}\) on rational numbers, and over \(F\) compose it with \(\operatorname{Tr}_{F/\mathbf Q}\). Thus the finite Heisenberg formula below and the chosen tensor symplectic form have the same convention. This explicit sign choice is part of this construction; it is not an identification with an independently fixed source pin.

Here is an explicit choice of the required splitting characters. Put \(K_- =\mathbf Q(\sqrt{-7})\), \(\varpi=\sqrt{-7}\) with positive imaginary part, and let \(\epsilon_7:(\mathcal O_{K_-}/\varpi)^\times\to\{\pm1\}\) be the quadratic character of \(\mathbf F_7^\times\). The Minkowski bound \((2/\pi)\sqrt7<2\) proves that \(K_-\) has class number one. For an ideal \(\mathfrak a=(a)\) prime to \(\varpi\), set

\[
             u(\mathfrak a)=\epsilon_7(a)\,a/|a|.
 \tag{V8}
\]

This is independent of the generator because the only units are \(\pm1\) and \(\epsilon_7(-1)=-1\). It is multiplicative. Define a unitary idele-class character \(\mu_-\) by the unramified local ideal values (V8), and by

\[
 \mu_{-,\infty}(z)=|z|/z,\qquad
 \mu_{-,\varpi}(v)=\epsilon_7(v)\ (v\text{ a unit}),\qquad
 \mu_{-,\varpi}(\varpi)=i.
 \tag{V9}
\]

To verify rational triviality, write any \(a\in K_-^\times\) as \(a=\varpi^m b\) with \(b\) a \(\varpi\)-adic unit. The unramified finite factors multiply to \(\epsilon_7(b)b/|b|\); the ramified factor is \(i^m\epsilon_7(b)\); the infinite factor is \((i^mb/|b|)^{-1}\). Their product is one. This verifies the idele-class character, including its ramified normalization. Its restriction to \(\mathbb A_\mathbf Q^\times\) is \(\epsilon_{K_-/\mathbf Q}\): away from 7 the uniformizer value is the quadratic residue character, at 7 its value on units is \(\epsilon_7\) and its value on 7 is one, and at infinity it is the sign character. These are the local components of the quadratic character of discriminant \(-7\).

Since \(E=FK_-\), put

\[
                   \mu_E=\mu_-\circ N_{E/K_-}.
 \tag{V10}
\]

The norm compatibility of quadratic class-field characters gives \(\mu_E|_{\mathbb A_F^\times}=\epsilon_{E/F}\). Choose \(\chi_W=\mu_E\), and choose \(\chi_V=1\) when \(n\) is even and \(\chi_V=\mu_E\) when \(n\) is odd. They satisfy exactly

\[
 \chi_W|_{\mathbb A_F^\times}=\epsilon_{E/F}^{\dim W},
 \qquad
 \chi_V|_{\mathbb A_F^\times}=\epsilon_{E/F}^{\dim V}.
 \tag{V11}
\]

Fix the global unitary Weil splitting determined by these data. The construction and its rational compatibility are the standard global splitting theorem, in the explicit form stated by Borade et al.~(2025), §§3.1 and 4.1.1. We use only that construction and the definition of the theta integral from §4.1, not their root-number or central-\(L\)-value nonvanishing theorem. The two unitary actions commute and agree with the canonical rational splitting on rational points. Thus they give an honest representation of \(G(\mathbb A_F)\times T(\mathbb A_F)\), where \(G=\operatorname U(V)\). All scalar phases below refer to this fixed representation. In particular they are not phases that may be adjusted separately on different lattices.

The Schrödinger position space used in (V12) is the rational Lagrangian supplied by the integral period basis. It need not be \(F\)-stable. This causes no change in the representation: taking the trace of the central coordinate identifies the \(F\)-Heisenberg representation with the representation of the underlying rational symplectic space having the displayed character. A rational change of polarization transports both the Schwartz representation and the rational theta distribution. Thus the usual global unitary splitting acts on this specified rational Schrödinger model, and the theta sum in that model is indexed by \(\mathbf Q^{3n}\).

\subsection{3. The actual finite vacuum and its maximal-compact eigencharacter}\label{the-actual-finite-vacuum-and-its-maximal-compact-eigencharacter}

Use the actual period coordinates (u=Pz), \(Y=\operatorname{Im}\tau\) from AM-R-004. In each copy take \(\delta=(1,1,1)^t\), \(\varepsilon=(1,0,1)^t\), and

\[
 \phi_{0,f}(q)=
   -i\,1_{\delta/2+\widehat{\mathbf Z}^{\,3}}(q)
        \psi_f(-q^t\varepsilon/2).
 \tag{V12}
\]

This is the actual input of AM-DO-016/018, not a variable in a formal theta-family interface. For \(n\) copies use its tensor power. On the rational period vector \(t=A+\tau B\), its Heisenberg action is

\[
 (W_t\phi)(q)=\psi_f(-A^tq+A^tB/2)\phi(q-B).
 \tag{V13}
\]

Let \(\widehat\Lambda\) be the completion of the product period lattice identified with \(\widehat{\mathcal O}_E^n\), and put \(w_0=-(\zeta+\zeta^2+\zeta^4)\). Its finite trace-one Heisenberg eigencharacter (a semicharacter for the operators \(W_x\)) is

\[
 \alpha_f(x)=\psi_f\bigl(\mathfrak b(x,w_0x)/2\bigr),
       \qquad x\in\widehat\Lambda.
 \tag{V14}
\]

This is an adelic additive-character expression. The notation \(e^{\pi i\mathfrak b(x,w_0x)}\) is not used for arbitrary adelic \(x\). On the rational integral lattice (V14) equals the sign semicharacter of AM-DO-018, since its value is \(\pm1\). The eigenrelations (V13) for (V12) therefore extend from that dense rational lattice to all of \(\widehat\Lambda\), by continuity.

The finite eigenspace

\[
 \mathscr V_\Lambda=
 \{\phi:W_x\phi=\alpha_f(x)\phi
                         \text{ for all }x\in\widehat\Lambda\}
 \tag{V15}
\]

inside the finite Schwartz representation has dimension one. Here is a direct justification tied to the actual theta model. AM-DO-018's reading map \(R\) is injective on finite Schwartz functions and intertwines (V13) with normalized Mumford translation. It sends (V15) into the principal theta line on the actual complex torus, which has dimension one. Formula (V12) supplies a nonzero vector in (V15), proving equality of dimensions. For tensor powers the same reading map and the product principal theta line prove the identical assertion. This is also the usual selfdual-lattice vacuum uniqueness in the Heisenberg representation.

Every \(k\in K_0\) preserves \(\widehat\Lambda\), preserves \(\mathfrak b\), and commutes with \(w_0\). Consequently it preserves (V14). Heisenberg covariance of the \emph{fixed} honest Weil representation shows that it preserves (V15). There is therefore a unique scalar

\[
 \omega_f(k)\phi_{0,f}^{\otimes n}
       =\lambda_f(k)\phi_{0,f}^{\otimes n}.
 \tag{V16}
\]

The group law makes \(\lambda_f\) a character. Unitarity makes its values have absolute value one; smoothness makes its kernel open in \(K_0\), so its image is finite and consists of roots of unity. This proves that the scalar is an actual finite-order character, without guessing a Weil index. For an explicit evaluation prescription, the value of (V12) at \(q=\delta/2\) is \(i\); therefore

\[
 \lambda_f(k)=i^{-n}
   (\omega_f(k)\phi_{0,f}^{\otimes n})
                 (\delta/2,\ldots,\delta/2).
 \tag{V17}
\]

The same line argument applies to every selfdual \(\mathcal O_E\)-lattice vacuum of AM-DO-018, because scalar norm-one integral units preserve that lattice and its trace-one character. It does not assert that an arbitrary finite Heisenberg translate has the same character (V14).

\subsection{4. The reduced algebraic jet is an actual homogeneous Fock vector}\label{the-reduced-algebraic-jet-is-an-actual-homogeneous-fock-vector}

For one copy put

\[
 K(z,x)=\exp\left(\frac\pi2z^tY^{-1}z
                    +\pi i x^t\tau x+2\pi i x^tz\right),
 \qquad \varphi_{0,\infty}(x)=K(0,x).
 \tag{V18}
\]

The first factor is exactly the Chern reduction used throughout this run. Define the commuting creation operators

\[
 r_j=-\frac i2\sum_k(Y^{-1})_{jk}
           \left(\partial_{x_k}-2\pi i(\bar\tau x)_k\right).
 \tag{V19}
\]

A direct calculation gives

\[
 r_jK(z,x)=\left(2\pi ix_j+\pi(Y^{-1}z)_j\right)K(z,x)
                  =\partial_{z_j}K(z,x).
 \tag{V20}
\]

Replace the operators by their linear combinations dual to \(u=Pz\):

\(r_{u_i}=\sum_j(P^{-1})_{ji}r_j\). Thus for every multi-index \(\alpha\),

\[
 \varphi_{\alpha,\infty}(x)
    =\left.\partial_u^\alpha K(P^{-1}u,x)\right|_{u=0}
    =r_u^\alpha\varphi_{0,\infty}(x).
 \tag{V21}
\]

It is a polynomial times a positive Gaussian, hence Schwartz. The real torus \(T(F_\infty)\simeq(S^1)^3\) preserves the positive complex polarization. Its action preserves the vacuum line and acts diagonally on the three CM creation directions. One may check preservation directly: the differentiated real Heisenberg operator with parameter \((A,B)\) acts on the vacuum as \(-2\pi i(A+\tau B)^t x\). Its complex annihilator is therefore the kernel of \(A+\tau B\), which is preserved by the CM torus. Its unique Schwartz vacuum line is consequently preserved. Consequently a monomial of fixed CM multidegree is an eigenvector. For \(n\) copies and a multi-index \(\alpha=(\alpha_{ij})\), define its actual eigencharacter by

\[
 \omega_\infty(t)\varphi_{\alpha,\infty}
      =\beta_\alpha(t)\varphi_{\alpha,\infty}.
 \tag{V22}
\]

This fixes the sign and the vacuum character for the splitting in §2. The action on creation operators determines the relative weights, while the action on the vacuum supplies its scalar shift. No replacement of (V22) by an unshifted polynomial degree is made. The eigenspace assertion can equivalently be checked from (V19): Heisenberg covariance transports the holomorphic creation space by the three CM characters, and the simultaneous monomials have the corresponding characters multiplied by the vacuum character.

For clarity, the actual Gaussian norm is also explicit. With Lebesgue measure in the fixed rational position coordinates,

\[
 \|\varphi_{\alpha,\infty}\|_2^2
  =\det(2Y)^{-n/2}\,
       \prod_{i,j}\alpha_{ij}!\left(\frac\pi{G_i}\right)^{\alpha_{ij}}.
 \tag{V23}
\]

Indeed the vacuum norm is \(\det(2Y)^{-n/2}\). From (V19), \([r_{u_i}^{*},r_{u_j}]=\pi\delta_{ij}/G_i\), and \(r_{u_i}^{*}\varphi_{0,\infty}=0\). Moving annihilation operators past creation operators proves (V23) by induction. The finite vector (V12) has norm one, because its support has volume one and every nonzero value has absolute value one. Thus (V23) is also the global Schwartz norm for this input.

The theta series at the identity is

\[
 \vartheta_\alpha(1,1)
   =\sum_{q\in\mathbf Q^{3n}}
       \phi_{0,f}^{\otimes n}(q)\varphi_{\alpha,\infty}(q)
   =\left.\partial_u^\alpha
                  R(\phi_{0,f}^{\otimes n})(u)\right|_{u=0}.
 \tag{V24}
\]

This is an equality of an actual adelic theta sum and an actual reduced CM derivative. Differentiation is legitimate because the finite support lies in a bounded-denominator rational lattice and the positive Gaussian gives locally normally convergent series.

AM-DO-018 gives

\[
 R(\phi_{0,f}^{\otimes n})(u)
   =\theta[0;0](0;\tau)^n\prod_{j=1}^n\Sigma_{\mathrm{tr}}(u_j).
 \tag{V25}
\]

ALP-LIFT-08 gives \(\sigma_q(u)=u_1u_3-u_2^2+\text{higher terms}\), and AM-DO-015 identifies \(\Sigma_{\mathrm{tr}}=c\sigma_q\) with \(c\ne0\). Taking \(\alpha_{1j}=\alpha_{3j}=1\), \(\alpha_{2j}=0\) for every \(j\), gives the exact nonzero value

\[
             \vartheta_\alpha(1,1)
                  =(\theta[0;0](0;\tau)c)^n\ne0.
 \tag{V26}
\]

For this input the norm (V23) is \(\det(2Y)^{-n/2}(\pi^2/(G_1G_3))^n\). Neither (V26) nor that norm replaces the common period by an unrecorded scalar.

\subsection{5. One actual character gives the nonzero theta lift}\label{one-actual-character-gives-the-nonzero-theta-lift}

For the fixed splitting, write

\[
 \vartheta_\alpha(g,t)=
       \sum_q(\omega(g,t)\varphi_\alpha)(q),
 \qquad
 \varphi_\alpha=\phi_{0,f}^{\otimes n}
                           \otimes\varphi_{\alpha,\infty}.
 \tag{V27}
\]

The theta function is left invariant under rational points, by the rationally compatible global Weil construction. Its two group actions commute. Equations (V16) and (V22) therefore give

\[
 \vartheta_\alpha(g,t_\infty k)
       =\beta_\alpha(t_\infty)\lambda_f(k)
                            \vartheta_\alpha(g,1).
 \tag{V28}
\]

For \(\zeta\in\mu_{14}\), the diagonal element is rational. Evaluating (V28) at \(g=1\) and using the already established nonzero value (V26) proves

\[
                  \beta_\alpha(\zeta)\lambda_f(\zeta)=1.
 \tag{V29}
\]

This is precisely the finite-root-unit compatibility which is sometimes lost if only the polynomial degree is retained. Its proof uses the actual theta series and rational splitting; it is not an imposed compatibility condition assumed for a hypothetical lift.

By (V5)--(V6) and (V29), the formula

\[
 \chi_\alpha(\gamma t_\infty k)
        =\beta_\alpha(t_\infty)\lambda_f(k),
 \quad \gamma\in T(F),\ t_\infty\in T(F_\infty),\ k\in K_0,
 \tag{V30}
\]

is well-defined and is a continuous unitary automorphic character. All automorphic forms on this anisotropic torus are cuspidal, since a torus has no proper parabolic subgroups. Equation (V28), together with rational automorphy and (V5), now holds for every adelic \(t\) in the form

\[
            \vartheta_\alpha(g,t)
                   =\chi_\alpha(t)\vartheta_\alpha(g,1).
 \tag{V31}
\]

Choose Haar measure on \([T]\) of total mass one and use the theta-lift convention with a conjugated automorphic input. Then

\[
 \boxed{\Theta(\varphi_\alpha,\chi_\alpha)(g)
   =\int_{[T]}\vartheta_\alpha(g,t)\overline{\chi_\alpha(t)}\,dt
   =\vartheta_\alpha(g,1).}
 \tag{V32}
\]

There is no limiting procedure, convergence regularization, or unproved nonvanishing theorem in this calculation. The quotient is compact; (V31) makes the integral elementary. The source character has \(L^2\) norm one. Taking \(n=2\) proves the claimed actual nonzero \(U(1)\to U(2)\) lift. With a different fixed Haar measure, the right side is multiplied by its explicitly retained total volume.

The common-period-normalized identity value is exactly one by (V26). This identifies one actual nonzero matrix value. It is not an assertion that the whole automorphic function is primitive in the source's mass-weighted integral row module, and no valuation is assigned to a transcendental period itself.

\subsection{6. General finite inputs: an exact common first-image projector}\label{general-finite-inputs-an-exact-common-first-image-projector}

This extension explains exactly what remains for translated inputs. Let \(\mathcal S\) be a finite-dimensional span of finite Schwartz functions. Let \(K\subset K_0\) fix every vector in \(\mathcal S\), and shrink it at 3 so that \(K\cap\mu_{14}=1\). Such a common open subgroup exists by smoothness and finite dimensionality. Keep any fixed homogeneous archimedean weight \(\beta\). The statement applies to any finite-dimensional subspace of that weight, not only a single monomial. Write \(h=h_K\).

The identity finite component of \([T]/K\) is the image of \(T(F_\infty)K\), and it is an injective copy of \(T(F_\infty)\), because \(K\cap\mu_{14}=1\). Define

\[
 f_{K,\beta}(t)=
 \begin{cases}
 h\,\beta(t_\infty),&[t]=[t_\infty k],\ k\in K,\\
 0,&\text{on the other finite components}.
 \end{cases}
 \tag{V33}
\]

This is a well-defined smooth automorphic form of weight \(\beta\), with

\[
             \|f_{K,\beta}\|_2^2=h,
                 \qquad \|f_{K,\beta}\|_\infty=h.
 \tag{V34}
\]

Each component has measure \(1/h\). For any \(\phi_f\in\mathcal S\) and any archimedean vector of weight \(\beta\), the theta function on the identity component is \(\beta(t_\infty)\vartheta_\phi(g,1)\). Hence

\[
                \boxed{\Theta(\phi,f_{K,\beta})(g)
                           =\vartheta_\phi(g,1).}
 \tag{V35}
\]

The same input \(f_{K,\beta}\) works for the entire specified finite span. If the source automorphic input is required to have norm one, replacing it by \(h^{-1/2}f_{K,\beta}\) multiplies the output by \(h^{-1/2}\). Thus the exact added normalization cost is \(\frac12\log h\) in logarithmic norm, not an unspecified projection constant.

There is also an explicit finite character decomposition. Set \(A=K_0/K\), with its embedded subgroup \(\mu_{14}\). Let

\[
 \mathcal X_\beta=
   \{\eta\in\widehat A:
               \eta|_{\mu_{14}}=\beta^{-1}|_{\mu_{14}}\}.
 \tag{V36}
\]

Every character of a subgroup of a finite abelian group extends: one may extend successively over cyclic generators by choosing the necessary roots of unity. The extensions form a torsor for \(\widehat{A/\mu_{14}}\). Thus \(\#\mathcal X_\beta=h\), all its finite values are roots of unity, and each

\(\chi_\eta(t_\infty k)=\beta(t_\infty)\eta(k)\) descends to \([T]\). Finite character orthogonality gives

\[
                 f_{K,\beta}=\sum_{\eta\in\mathcal X_\beta}\chi_\eta.
 \tag{V37}
\]

No coefficient \(1/h\) occurs in (V37). For one component, choosing representatives \(a\in A/\mu_{14}\) gives the exact projection

\[
 \Theta(\phi,\chi_\eta)(g)
   =\frac1h\sum_{a\in A/\mu_{14}}
                   \overline{\eta(a)}\vartheta_\phi(g,a).
 \tag{V38}
\]

The summand is independent of representative because rational-root automorphy cancels the two characters. Formula (V38) displays the only averaging denominator; its \(p\)-adic contribution is at worst \(v_p(h)\) once the specified theta values have actually been placed in a common integral lattice. If (V35) is nonzero, at least one of the \(h\) irreducible-character lifts in (V37) is nonzero. This is a consequence for the explicitly specified finite family; it is not nonvanishing for a freely prescribed Hecke character.

\subsection{7. What this closes, and what it does not close}\label{what-this-closes-and-what-it-does-not-close}

The trace-one scalar torus argument supplies an actual nonzero first compact lift with the primitive mixed CM type. It also supplies a common actual first-image map for any selected finite Schwartz span at one weight, with an exact finite-index normalization. In that limited first-lift setting, existence of an automorphic input is no longer a merely formal interface.

For the base vacuum, the finite and archimedean input normalizations are (V12), (V23), and the norm-one character in (V30). The unknown target automorphic mass norm is a different quantity. The exact finite character (V16) is fixed by the global splitting; identifying it with a source-prescribed determinant-root character requires the corresponding local calculation. AM-DO-020 performs that calculation at relative split places, including 11, for the stated unramified splitting. No omitted ramified phase is declared equal to one here.

For DP59/DP65's permitted translated sections, applying §6 requires an actual common stabilizer \(K\), the exact archimedean weight, and the source's allowed automorphic input space. An unrestricted finite Schwartz span may be larger than the prescribed one. The finite-coset input in (V33) is generally a sum of characters, and its cost (V34) must be retained if the source requires unit norm. The accepted DP59/DP65 field-span and sharp-jet results are not re-proved or promoted here to a full iterated-image statement.

Nothing in (V32) or (V35) supplies the later branching map, all-degree \nolinkurl{ClsK.Core} with its specified phases, simultaneous terminal mass identities, the global \nolinkurl{MainX} laws, or a changed DP44 cost with all \(g_j,e_j,H\) identified. Those are the remaining gates for Q1's requested shorter tower.

\subsection{Primary references actually consulted}\label{primary-references-actually-consulted}

Borade, N., Franzel, J., Girsch, J., Yao, W., Yu, Q., \& Zelingher, E. (2025). \emph{On tori periods of Weil representations of unitary groups} (arXiv:2402.16808v2). \url{https://arxiv.org/html/2402.16808v2}. Exact used locators: §3.1, tensor alternating form and Heisenberg representation; §3.1.1, splitting characters; §4.1, rational global splitting, theta sum, and conjugated-input theta integral; §4.1.1, global character restrictions and tensor-product splitting. No nonvanishing theorem from that paper is used.

Milne, J. S. (2020). \emph{Algebraic number theory} (Version 3.08). \url{https://www.jmilne.org/math/CourseNotes/ANT.pdf}. Exact used locators: Theorem 4.3, printed pp.~70--71, Minkowski bound; Theorem 5.1, printed p.~85, unit rank; the logarithmic kernel discussion in §5, printed pp.~86--87.

Prasad, D. (2000). Theta correspondence for unitary groups. \emph{Pacific Journal of Mathematics, 194}(2), 427--438. \url{https://www.math.iitb.ac.in/~dprasad/pacific.pdf}. Used as an independent convention check: §1, printed pp.~427--428, the dependence on two splitting characters and the passage from metaplectic to honest unitary representations.

\subsection{Check status}\label{check-status}

The author read the complete saved proof, including the arithmetic class and character construction, finite trace-one eigenline, Chern/Fock derivative and Gaussian norm calculation, single-character integral, and general finite-coset projector. The exact algebra certificate in actual\_torus\_certificate.py passed. Actual check events and their scope are recorded in \nolinkurl{lift_accounting_CHECK_EVENTS.tsv.} No Lean compilation or formal proof is claimed.

\clearpage

\section{AM-R-007: unitary lattice classes have representatives supported at relative split places}\label{am-r-007-unitary-lattice-classes-have-representatives-supported-at-relative-split-places}

Author: \nolinkurl{/root}, GPT-6 Astra Pro; effort not externally exposed; ALP\_MACHINE, 21 November 2026. This is a concrete application of strong approximation to the moved-lattice comparison. It is a statement about double-coset representatives, not an assertion that the two-average theta family already has the required rational equivariance.

\subsection{Statement}\label{statement-8}

Let \(E/F\) be a quadratic extension of number fields, with its nontrivial involution denoted by a bar. Let \(V\) be a nondegenerate Hermitian space of dimension \(n\geq2\), let \(G=U(V)\), \(G^{\rm der}=SU(V)\), and \(T=\ker(N_{E/F}:\operatorname{Res}_{E/F}\mathbf G_m\to\mathbf G_m)\). Let \(\mathcal S\) be the set of finite places of \(F\) that split in \(E\), and choose one \(v_0\in\mathcal S\). For every compact open subgroup \(K\subset G(\mathbb A_{F,f})\), \[
 G(\mathbb A_{F,f})=G(F)\,G_{\mathcal S}^{\rm fin}\,K,
 \tag{R7.1}
\] where \(G_{\mathcal S}^{\rm fin}\) denotes adeles equal to the identity outside a finite subset of \(\mathcal S\). The finite subset is allowed to depend on the represented class.

For the actual field \(E=\mathbb Q(\zeta_7)\), \(F=E^+\), standard form \(V=E^n\) with Gram matrix \(I_n\), and its full standard lattice stabilizer \(K_0=\prod_vU(I_n,\mathcal O_{E,v})\), the stronger statement is \[
 G(\mathbb A_{F,f})=G(F)\,G(F_{v_{11}})\,K_0,
 \tag{R7.2}
\] where \(v_{11}\) is the unique place of \(F\) above 11. It splits in \(E/F\), since the subgroup generated by \(11\bmod7=4\) has order three and does not contain complex conjugation.

Thus every standard-genus Hermitian lattice class needed at maximal level has a representative obtained by moving only at 11. At a smaller level, (R7.1) supplies representatives using finitely many relative split places. Neither assertion bounds the denominators of the representative, the index of an intersection lattice, its stabilizer order, or a terminal mass.

\subsection{Proof of the split-place reduction for arbitrary compact level}\label{proof-of-the-split-place-reduction-for-arbitrary-compact-level}

Choose a diagonal Hermitian basis over \(E\); this is possible by Hermitian Gram--Schmidt in characteristic zero. It gives an explicit section of determinant, \[
 s:T\longrightarrow G,\qquad t\longmapsto
             \operatorname{diag}(t,1,\ldots,1).
 \tag{R7.3}
\] Indeed multiplying the first diagonal coordinate by a norm-one scalar preserves the diagonal form. This also proves surjectivity of determinant on all local and global point groups.

Replace \(K\), if necessary, by a product compact open subgroup \(K'=\prod_vK'_v\subset K\), standard at almost every place. Proving (R7.1) for \(K'\) proves it for \(K\), so write this product subgroup as \(K\) from now on. Put \(K_{T,v}=\det K_v\). This is an open subgroup of \(T(F_v)\): the local analytic section (R7.3), near the identity, shows that the image of an open identity neighborhood is open. At almost all places it contains all integral points of \(T\), and in the standard unramified model it equals them.

At a nonsplit finite place every norm-one element of \(E_v^\times\) is a unit: its valuation multiplied by the residue degree is the valuation of its norm, which is zero. Thus \(T(F_v)\) is compact, and \(K_{T,v}=T(F_v)\) at all but finitely many nonsplit places.

Given \(g\in G(\mathbb A_{F,f})\), set \(t=\det g\). There exists \(t_0\in T(F)\) such that \[
 t_0^{-1}t_v\in K_{T,v}
       \quad\text{at every nonsplit finite place }v.
 \tag{R7.4}
\] Only finitely many conditions need attention. For completeness the required weak approximation on this particular torus has a direct proof. Local Hilbert 90 writes each prescribed norm-one local element as \(a_v/\bar a_v\). Weak approximation in the number field \(E\) gives a nonzero \(a\in E\) sufficiently close to these finitely many \(a_v\). The continuous map \(a\mapsto a/\bar a\) then puts \(t_0=a/\bar a\) in the required open cosets. This uses only finite local approximation, not strong approximation for a torus with a finite omitted set.

At the split places, \(t_0^{-1}t_v\) belongs to \(K_{T,v}\) for all but finitely many \(v\), since it is an adele and \(K\) is standard almost everywhere. We may therefore write \[
 t=t_0\,t_{\mathcal S}\,k_T,
 \tag{R7.5}
\] where \(t_{\mathcal S}\) is supported at finitely many split places and \(k_T\in\prod_vK_{T,v}\). Choose \(k\in K\) with \(\det k=k_T\), choosing a lift at each place, and let \(a=s(t_{\mathcal S})\), a finite-support split-place adele. Then \[
 g_2=s(t_0)^{-1}gk^{-1},\qquad h=g_2a^{-1}
                     \in SU(V)(\mathbb A_{F,f}).
 \tag{R7.6}
\]

The group \(SU(V)\) is connected, simply connected, and absolutely almost simple for \(n\geq2\); over an algebraic closure it is \(SL_n\). At the chosen split place \(v_0\), it is \(SL_n(F_{v_0})\), hence noncompact. The strong approximation theorem for simply connected absolutely almost simple groups therefore applies with the omitted set consisting of all archimedean places together with \(v_0\). The exact input is Rapinchuk (2013), Theorem 2.3, printed p.~281, recording the Kneser--Platonov theorem in characteristic zero.

Apply it to the compact open subgroup \(K''=(aKa^{-1})\cap SU(V)(\mathbb A_{F,f})\). It gives \[
 h=\gamma u_0k'',\qquad
 \gamma\in SU(V)(F),\quad u_0\in SU(V)(F_{v_0}),
 \quad k''\in K''.
 \tag{R7.7}
\] To see the factorization directly from density, approximate \(h\) outside \(v_0\) by \(\gamma\) so that \(\gamma^{-1}h\in K''_v\) at each other finite place; put these components into \(k''\) and put its \(v_0\)-component equal to one. The remaining component is \(u_0\).

Since \(k''a=a(a^{-1}k''a)\), substitution in (R7.6) yields \[
 g=s(t_0)\gamma\,(u_0a)\,(a^{-1}k''a)k.
 \tag{R7.8}
\] The first factor is rational, the middle factor is supported at finitely many relative split places, and the last two factors lie in \(K\). This proves (R7.1). The order of factors in this argument matters: no nonrational adele was commuted past a rational matrix without conjugating the compact subgroup.

\subsection{Proof of the maximal-level statement for the cyclotomic field}\label{proof-of-the-maximal-level-statement-for-the-cyclotomic-field}

The torus calculation in ALP-LIFT-09, equations (V2)--(V6), proves \[
 T(\mathbb A_{F,f})=T(F)\,T(\widehat{\mathcal O}_F)
 \tag{R7.9}
\] for \(E=\mathbb Q(\zeta_7)\). Its ingredients are the class-number one calculation for \(E\) and the explicit full signatures of \(-1,\zeta+\zeta^{-1},\zeta^2+\zeta^{-2}\), which imply that every totally positive unit of \(F\) is a square. That proof is cited once under its own finding identifier; it is not counted again here.

For the standard lattice, diagonal norm-one units in (R7.3) show \(\det K_0=T(\widehat{\mathcal O}_F)\). Thus for each \(g\), equation (R7.9) permits \(t_{\mathcal S}=1\) in (R7.5). In the preceding proof \(a=1\), and (R7.8) has its only noncompact finite factor at \(v_{11}\). This is (R7.2). The unique prime of \(F\) above 11 follows because the image of 4 in \((\mathbb Z/7\mathbb Z)^\times/\{\pm1\}\) has order three; its two lifts to \(E\) exhibit the relative splitting.

\subsection{What the representative theorem reduces, and what it does not}\label{what-the-representative-theorem-reduces-and-what-it-does-not}

There is a useful exact conditional transport principle. Suppose two assignments on \(G(\mathbb A_{F,f})\) have the same rational left-equivariance and the same right-\(K\) equivariance, with specified invertible transport operators, and agree on all finite-support split-place representatives. Factorization (R7.1) then proves equality everywhere: apply the common right transport and common rational left transport to the representative equality. This short argument proves the principle, including its explicit equivariance hypotheses.

AM-DO-017 gives exact transport when all lattices, including the base lattice, are carried by a fixed complex-linear isometry. That does not yet give the rational left-equivariance needed for the fixed-base two-average family (DO18.3). Returning the moved base to the original base introduces its finite triangle phase. Therefore the present theorem, even combined with split-place phase equality, leaves the rational triangle/order law and compact vacuum character as real proof obligations. It cannot erase them by choosing representatives.

This theorem also has no quantitative height assertion. Strong approximation supplies existence in an open adelic set; it does not bound the matrices \(\gamma\), the supports or denominators of \(a\), the Smith indices \(q_1,q_2\), or the finite stabilizer factors. Those quantities enter AM-DO-018's explicit costs and AM-R-002's terminal content formula. A shorter lift chain still requires those actual data and all six linked identities.

In geometric terms the double cosets classify the indicated Hermitian lattice genus with level structure. The determinant sequence \(1\to SU\to U\to T\to1\) separates its simply connected derived part from its norm-one torus part. This is the reductive-group and torsor side of the arithmetic comparison; it does not identify integral theta sections, crystalline filtrations, or an equality of motivic L-functions. The exact SGA3 and Picard framework locators are maintained in the geometric dictionary.

\subsection{References}\label{references-1}

Rapinchuk, A. S. (2013). Strong approximation for algebraic groups. In E. Breuillard \& H. Oh (Eds.), \emph{Thin groups and superstrong approximation} (MSRI Publications, Vol. 61, pp.~269--298). Cambridge University Press. Exact theorem: Theorem 2.3, p.~281. Directly inspected primary author exposition: \url{https://library.slmath.org/books/Book61/files/70rapi.pdf} .

Platonov, V. P. (1969). The problem of strong approximation and the Kneser--Tits conjecture for algebraic groups. \emph{Mathematics of the USSR-Izvestiya, 3}(6), 1139--1147. \url{https://doi.org/10.1070/IM1969v003n06ABEH000838} . The theorem is used in the exact form quoted above; no reproof of strong approximation is claimed.

\clearpage

\section{AM-K017 --- integral logarithmic Hessians, stronger actual theta slopes, and exact mixed-jet recurrences}\label{am-k017-integral-logarithmic-hessians-stronger-actual-theta-slopes-and-exact-mixed-jet-recurrences}

Author: \nolinkurl{/root/kummer_extension}, GPT-6 Astra Pro, ALP\_MACHINE, 21 November 2026 UTC. This supplements AM-K012--016 in \nolinkurl{jacobian_reduced_theta.md}. It concerns the actual translated germ \(F\) and the actual trace-one germ \(\sigma_q\). Every reduction below is a reduction of Hurwitz derivatives, not of divided Taylor coefficients.

\subsection{1. Statement}\label{statement-9}

Retain \[
C:y^2=x^7-1,\quad E=\mathbb Q(\zeta_7),\quad
\omega_i=x^{i-1}dx/y,\quad D_i=\partial/\partial u_i,
\quad q(u)=u_1u_3-u_2^2.
\] Let \(W_{ij}=-\partial_i\partial_j\log\sigma_q\) be the actual algebraic Kleinian Hessian. Its explicit expressions, normalization, and integral Abel chart are AM-K012. For an odd prime \(p\ne7\), let \(q_i=p^{h_i}\) be the sufficient restricted-power nilpotence exponent in the table (K13.2).

\textbf{(a) Global constant lemma.} The rational function \[
\overline{D_i^{q_i-1}W_{ij}}
\tag{K18.1}
\] on the good special fibre is a constant. It is zero unless \[
iq_i+j\equiv0\pmod7.
\tag{K18.2}
\] In particular the pure constant vanishes whenever \(q_i\not\equiv-1\pmod7\). This proves the additional log-derivative hypothesis in (K13.7) geometrically, at all such primes. The three exact 11-adic checks in AM-K013 are consistent independent finite checks, no longer necessary premises for that vanishing.

\textbf{(b) Improved actual finite bounds.} At every prime above \(p\), for either \(F\) or \(\sigma_q\), \[
v_p(\partial_i^n f(0))\geq
\max\{0,c_i(p)n-s_p(n)/(p-1)\},
\tag{K18.3}
\] where the following constants are now proved:

\begin{longtable}[]{@{}llll@{}}
\toprule\noalign{}
\(p\bmod7\) & \(c_1(p)\) & \(c_2(p)\) & \(c_3(p)\) \\
\midrule\noalign{}
\endhead
\bottomrule\noalign{}
\endlastfoot
2 & \(1/(2p(p-1))\) & \(1/(2(p-1))\) & \(1/(2(p-1))\) \\
3 & \(1/((p-1)(p^3+1))\) & \(1/(2(p-1))\) & \(1/(2p(p-1))\) \\
4 & \(1/(2(p-1))\) & \(1/(2p(p-1))\) & \(1/(2(p-1))\) \\
5 & \(1/(2(p-1))\) & \(1/((p-1)(p^3+1))\) & \(1/(2p(p-1))\) \\
6 & \(1/(p^2-1)\) & \(1/(p^2-1)\) & \(1/(p^2-1)\) \\
\end{longtable}

For \(p\equiv1\pmod7\), this argument provides no positive constant. As before, these constants are lower bounds, not proved optimal finite slopes.

\textbf{(c) Full restricted-operator identities at 11.} In the three-variable Hurwitz algebra over \(\mathbb F_{11}[\zeta_7]\), the actual germs satisfy \[
\boxed{
\partial_1^{11}f=4u_3f,\qquad
\partial_2^{11}f=-\partial_1f,\qquad
\partial_3^{11}f=-u_2f,
\qquad f=F\ \hbox{or}\ \sigma_q.}
\tag{K18.4}
\] These are identities of the complete reduced Hurwitz series. In particular \[
\partial_1^{12}\partial_3^2\sigma_q(0)\equiv8\pmod{11}.
\tag{K18.5}
\] This actual source mixed derivative is a unit at both primes above 11. It disproves the naive additive extension of the three pure-direction constants to all mixed multiindices.

\textbf{(d) A failed next strengthening.} For \(\ell_n=\partial_i^n\log F(0)\), exact arithmetic gives \[
v_{11}(\ell_{2q_i})=1\qquad(i=1,2,3)
\tag{K18.6}
\] at both primes above 11. Thus an attempted improvement based on \(p^2\mid\ell_{2q_i}\) is false. The Bell-lemma equality case at that index is a real obstruction for this particular estimate, although cancellation in the full theta derivative could still give stronger bounds by a different argument.

\subsection{2. The integral extension lemma}\label{the-integral-extension-lemma}

We first give the general local-to-global argument, including its finite-place content.

Let \(R\) be the localization of \(\mathcal O_E\) at a prime above an odd \(p\ne7\), with residue field \(k\), and let \(\mathcal J/R\) be the smooth proper abelian scheme. Completion may be taken after the argument. All rational functions and line charts can thus be taken over the number field, where the complex comparison below applies. The total space is integral and normal, and its special fibre is geometrically integral. Let \(\mathcal L\) be the primitive integral theta line of ALP-LIFT-08, and \(s_q\) its primitive section at this prime. Write \(\Theta\) for its divisor. Its restriction to the special fibre is a nonzero theta divisor; in particular \(\Theta\) has no vertical divisor component.

Take an affine open \(\mathcal U\) on which \(\mathcal L\) is trivial, and write \(f_{\mathcal U}\) for the local equation of \(s_q\). This is a regular function. At the generic point of the special fibre it is a unit, since the section does not vanish on that entire fibre. The invariant derivatives \(D_i\) preserve regular functions on \(\mathcal U\), so \[
D_iD_j\log f_{\mathcal U}
=\frac{D_iD_jf_{\mathcal U}}{f_{\mathcal U}}
 -\frac{D_if_{\mathcal U}D_jf_{\mathcal U}}{f_{\mathcal U}^2}
\tag{K18.7}
\] is a rational function with no vertical pole.

The algebraic \(W_{ij}\) has no vertical pole either. Its pullback to the ordered-divisor chart \(\mathcal R\) of (K12.12) is explicitly integral by (K12.8)--(K12.10). This chart maps étale to \(\mathcal J\), and its reduction has a nonempty open image in the irreducible special fibre. Thus it detects the generic vertical valuation. An étale pullback cannot turn a negative vertical valuation into a nonnegative one. This proves the assertion without any integral Chern frame.

On the characteristic-zero generic fibre, define \[
\beta_{ij,\mathcal U}=W_{ij}+D_iD_j\log f_{\mathcal U}.
\tag{K18.8}
\] It is regular there. To justify the cancellation explicitly, in any analytic lift of a small open set write the actual reduced sigma as \(\sigma_q=f_{\mathcal U}g\), where \(g\) is the holomorphic, nonvanishing transition from the algebraic local frame to the reduced frame. Since \(W_{ij}=-D_iD_j\log\sigma_q\), the right side of (K18.8) is \(-D_iD_j\log g\), which is holomorphic. The right side is also rational over the generic algebraic fibre, hence has no horizontal divisor pole. Equivalently one may check each horizontal valuation after embedding into \(\mathbb C\).

By (K18.7) and the vertical integrality of \(W_{ij}\), (K18.8) has no vertical divisor pole. Every codimension-one point of the smooth total space is either horizontal or generic in the special fibre. A normal affine domain is the intersection of its height-one valuation rings. Therefore \[
\boxed{\beta_{ij,\mathcal U}\in\Gamma(\mathcal U,\mathcal O).}
\tag{K18.9}
\] This is an integral statement about the logarithmic \textbf{second derivative} of a local frame. It neither asserts an integral reduced frame nor asserts integral divided Taylor coefficients. In particular it is compatible with the supplied THETA3 C11 counterexample.

Now suppose \(\overline D_i^{q}=0\) on the special fibre, where \(q=p^h\). It also vanishes on its function field: the \(q\)-th iterate is a derivation, and the quotient rule extends the zero derivation from any dense affine open. Reducing (K18.8), \[
\overline D_i^{q-1}\overline W_{ij}
=\overline D_i^{q-1}\overline\beta_{ij,\mathcal U}
 -\overline D_i^{q}
       \left(\frac{\overline D_j\overline f_{\mathcal U}}
                    {\overline f_{\mathcal U}}\right).
\tag{K18.10}
\] The last term is zero and the first term is regular on \(\mathcal U_k\). The left side is therefore regular on every open of a cover of \(\mathcal J_k\), hence is a global regular function. Properness and geometric connectedness give \(H^0(\mathcal J_k,\mathcal O)=k\). This proves that it is a constant.

There is a slightly more general version that will be used at 11. If \(\overline D_i^{p}=a\overline D_k\), with \(a\in k\), then \[
\overline D_i^{p-1}\overline W_{ij}-a\overline W_{kj}
\tag{K18.11}
\] is a constant. Lift \(a\) to the DVR arbitrarily and apply (K18.9) to both Hessians. The logarithmic term is \((\overline D_i^{p}-a\overline D_k)
 (\overline D_j\overline f/\overline f)=0\), proving regularity exactly as above.

\subsection{3. CM characters and the improved slopes}\label{cm-characters-and-the-improved-slopes}

The source section has CM character 4: \[
\sigma_q(\zeta u_1,\zeta^2u_2,\zeta^3u_3)=\zeta^4\sigma_q(u).
\tag{K18.12}
\] This was established in AM-K014, and also follows from the invariant theta line and its leading quadratic. Twice differentiating gives \[
a^*W_{ij}=\zeta^{-i-j}W_{ij},\qquad
a^*(D_i^mW_{ij})=\zeta^{-i(m+1)-j}D_i^mW_{ij}.
\tag{K18.13}
\] At \(m=q_i-1\), a nonzero constant could have this character only if \(iq_i+j\equiv0\pmod7\), because \(\zeta\) still has order 7 in every residue field under consideration. This proves (K18.1)--(K18.2).

For the pure Hessian, \(i\ne0\pmod7\), so the exceptional condition is exactly \(q_i\equiv-1\pmod7\). Otherwise \[
D_i^{q_i-1}W_{ii}\in p\mathcal R.
\tag{K18.14}
\] Together with \(D_i^{q_i}\mathcal R\subset p\mathcal R\), this even improves the logarithmic estimate to \[
v_p(D_i^mW_{ii}|_a)\geq
\left\lfloor\frac{m+1}{q_i}\right\rfloor,
\qquad
v_p(\ell_n)\geq\left\lfloor\frac{n-1}{q_i}\right\rfloor\ (n\geq2).
\tag{K18.15}
\] For clarity, if \(m=Aq_i+B\), \(0\leq B<q_i\), the usual operator estimate gives \(p^A\). When \(B=q_i-1\), apply \((D_i^{q_i})^A\) to (K18.14) to obtain \(p^{A+1}\). This proves the displayed shifted floor bound; it is an iterated stable-ring statement.

In particular \(p\mid\ell_{q_i+1}\), the hypothesis of the improved Bell--Legendre bound (K13.7). The nonexceptional columns in the table of §1 use \(1/(2p^{h_i-1}(p-1))\). The exceptional columns retain the generic constant \(1/((p-1)(p^{h_i}+1))\) of (K13.6). These are exactly the listed entries. AM-K015 transfers both bounds to \(\sigma_q\) using the actual integral \([2]\)-ratio, so the conclusion concerns the source line and respects the source's odd level.

The exceptional possibility cannot be discarded merely by symmetry. The finite certificate computes \[
\partial_1^{28}\log F(0)\equiv2\pmod3,
\qquad
\partial_2^{126}\log F(0)\equiv1\pmod5.
\tag{K18.16}
\] Here the respective exponents are \(q_1=3^3=27\) and \(q_2=5^3=125\), both congruent to \(-1\pmod7\). These are exact nonzero constants in their whole cyclotomic residue algebras.

\subsection{4. Exact restricted-operator identities at 11}\label{exact-restricted-operator-identities-at-11}

\subsubsection{The algebraic restricted powers}\label{the-algebraic-restricted-powers}

The CM-label argument gives \(D_1^{11}=D_3^{11}=0\) on \(\mathcal J_k\), and \(D_2^{11}\) is a scalar multiple of \(D_1\). The scalar is \(-1\): the exact integral ODE (K12.7), evaluated at the branch divisor, gives \[
D_2^{11}y_j=-D_1y_j\quad(j=1,2,3),
\tag{K18.17}
\] and \(D_1y_1\) is a unit in the cyclotomic residue algebra. Its coefficient vector in \(1,\zeta,\ldots,\zeta^5\) is \((5,4,6,4,5,0)\); the certificate verifies its inverse. Since the two sides are invariant vector fields in the same one-dimensional labelled space, this one unit evaluation determines the scalar globally. Thus \[
\boxed{D_1^{11}=0,\qquad D_2^{11}=-D_1,
\qquad D_3^{11}=0.}
\tag{K18.18}
\]

\subsubsection{The logarithmic constants}\label{the-logarithmic-constants}

For \(i=1\), condition (K18.2) permits only \(j=3\); for \(i=3\), it permits only \(j=2\). The exact finite calculation, using the particularly simple functions \(W_{13}=x_1x_2x_3/4\) and \(W_{23}=-(x_1x_2+x_1x_3+x_2x_3)/4\), gives \[
-D_1^{10}W_{13}=4,\qquad
-D_3^{10}W_{32}=-1\pmod{11}.
\tag{K18.19}
\] The global constant lemma makes these values valid everywhere on the special fibre, not just at the evaluation divisor.

For the middle operator, (K18.11) shows that \(D_2^{10}W_{2j}+W_{1j}\) is a constant. Its CM character is \(-(22+j)\equiv-(1+j)\pmod7\), which is nonzero for every \(j\in\{1,2,3\}\). Therefore \[
D_2^{10}W_{2j}+W_{1j}=0\pmod{11}.
\tag{K18.20}
\]

Let \(H=\log F\), whose Hurwitz coefficients are integral and whose first and all odd-total-degree derivatives at the origin vanish. Differentiating \(\partial_1^{11}H\) once in direction \(j\) gives the constants (K18.19), so in the full three-variable Hurwitz algebra \[
\partial_1^{11}H=4u_3,\qquad
\partial_3^{11}H=-u_2.
\tag{K18.21}
\] There is no integration ambiguity: a series killed by all three divided-power derivations is constant, and these constants at zero vanish by parity. Similarly (K18.20) and parity give \[
\partial_2^{11}H+\partial_1H=0.
\tag{K18.22}
\]

In characteristic 11, every positive-augmentation element of the multivariate Hurwitz algebra has 11-th power zero. For a divided-power monomial with multiindex \(\alpha\ne0\), the valuation of the scalar in its 11-th power is \(\sum_j s_{11}(\alpha_j)>0\). Frobenius is additive, and each coefficient of an infinite product involves only finitely many terms, so the same assertion holds for a complete series. The Bell-polynomial identity modulo a prime is elementary here: in the polynomial \(B_p(H',\ldots,H^{(p)})\), every term except \((H')^p\) and \(H^{(p)}\) has coefficient divisible by \(p\). Indeed its coefficient has numerator \(p!\), whereas all factorials in its denominator have arguments less than \(p\). Thus \[
\partial_i^{11}F/F=(\partial_iH)^{11}+\partial_i^{11}H
=\partial_i^{11}H.
\tag{K18.23}
\] Equations (K18.21)--(K18.22) prove (K18.4) for \(F\).

\subsubsection{Transfer to the actual trace-one section}\label{transfer-to-the-actual-trace-one-section}

Use the integral algebraic comparison \(\sigma_q(2u)=r(u)F(2u)\) of AM-K015. On reductions of its coefficient functions, the three operators \(D_1^{11}\), \(D_2^{11}+D_1\), and \(D_3^{11}\) kill \(r\), by (K18.18). Each operator is a derivation modulo 11. Apply it to the product identity. The linear terms in (K18.4) scale under \(u\mapsto2u\), while \(2^{11}=2\) in characteristic 11. Dividing by the resulting unit factor and replacing \(2u\) by \(u\) proves exactly (K18.4) for \(\sigma_q\). No integral Chern frame was introduced.

The operator identities also hold for AM-R-006's actual moved source section \(\Sigma_{\rm tr}=c\sigma_q\), since \(c\) is constant; its being an 11-adic unit preserves every unit assertion below.

\subsection{5. Mixed derivatives, finite unit families, and their scope}\label{mixed-derivatives-finite-unit-families-and-their-scope}

Write \(S_\alpha=\partial^\alpha\sigma_q(0)\), and interpret a negative index as zero. Equation (K18.4) is equivalently the all-index recurrence \[
\begin{aligned}
S_{\alpha+11e_1}&=4\alpha_3S_{\alpha-e_3},\\
S_{\alpha+11e_2}&=-S_{\alpha+e_1},\\
S_{\alpha+11e_3}&=-\alpha_2S_{\alpha-e_2}
\end{aligned}\qquad(\bmod11).
\tag{K18.24}
\] The first and third equalities are just the Hurwitz Leibniz rule for multiplying by a coordinate. The second is a constant-coefficient operator equality. Each replacement lowers the total derivative order, so these relations reduce every derivative modulo 11 to a finite base cube with all three indices less than 11. They do not assert that every element of that base cube is known without further calculation.

Taking \(\alpha=e_1+2e_3\), and using \(S_{e_1+e_3}=1\), proves (K18.5). The tempting additive extension of the pure lower constants would instead give \[
v_{11}(S_{(12,0,2)})\geq
\frac{12}{20}+\frac2{20}
-\frac{s_{11}(12)+s_{11}(2)}{10}
=\frac3{10}>0,
\] which is false. Consequently pure-direction gains cannot be assigned additively to a selected mixed harmonic row without additional analysis.

There are explicit further unit witnesses. Since \(\partial_1u_3=0\), iteration of (K18.4), as an identity of the actual formal series, gives \[
\partial_1^{11k}\sigma_q=4^ku_3^k\sigma_q\pmod{11}.
\] Differentiating the leading nonzero jets yields \[
\begin{aligned}
S_{(11k+1,0,k+1)}&=4^k(k+1)! &&(0\leq k\leq9),\\
S_{(11k,2,k)}&=-2\,4^kk! &&(0\leq k\leq10),
\end{aligned}\qquad(\bmod11).
\tag{K18.25}
\] Every displayed value is a unit. In particular \[
S_{(100,0,10)}=8,\quad S_{(110,2,10)}=2,
\quad S_{(0,1100,10)}=8\pmod{11}.
\tag{K18.26}
\] The last equality uses \(\partial_2^{1100}\sigma_q=\partial_1^{100}\sigma_q\), obtained by iterating the middle identity in (K18.4). These are finite unit families and exact all-index mod-11 recurrences; they supply no asymptotic valuation upper bound.

For completeness, the same formal algebra gives \[
\partial_1^{121}\sigma_q=0,\qquad
\partial_3^{121}\sigma_q=0,\qquad
\partial_2^{1331}\sigma_q=0\pmod{11}.
\tag{K18.27}
\] Indeed \(u_j^{11}=11!\,u_j^{[11]}=0\) in the reduced Hurwitz algebra; iterating the first or third relation eleven times proves the first two equalities, and the middle relation then proves the last. Differentiating shows vanishing for every multiindex exceeding one of these bounds in its indicated component. The stronger pure-direction cutoffs 11,121,11 from AM-K015 are also immediate by setting the other variables to zero.

Neither (K18.24) nor (K18.27) is iterated as a divisibility statement modulo powers of 11. Such an iteration would require a stable integral module and a justified operator inclusion into \(11\) times that module. No such stronger inclusion is inferred from these formal congruences.

\subsection{6. The next proposed divisibility is false}\label{the-next-proposed-divisibility-is-false}

To test whether the improved Bell constant could be increased by forcing the equality-case logarithmic derivative to acquire an additional power of 11, the exact ODE was evaluated modulo \(121\). In the coefficient basis \(1,\zeta,\ldots,\zeta^5\), it gives

\begin{longtable}[]{@{}lrl@{}}
\toprule\noalign{}
Direction & \(q_i\) & \(\ell_{2q_i}/11\pmod{11}\) \\
\midrule\noalign{}
\endhead
\bottomrule\noalign{}
\endlastfoot
1 & 11 & \(3+3\zeta+8\zeta^3+5\zeta^4+8\zeta^5\) \\
2 & 121 & \(3+3\zeta+8\zeta^3+5\zeta^4+8\zeta^5\) \\
3 & 11 & \(3+3\zeta+3\zeta^2+3\zeta^3\) \\
\end{longtable}

Each displayed polynomial has gcd 1 with \(\Phi_7\) over \(\mathbb F_{11}\), so it is a unit at both primes above 11. This proves (K18.6). At the Bell index \(j=2q_i\), the lower bound in (K13.7) has \((p-1)v(\ell_j)+s_p(j)-1=10+2-1=11\), precisely the equality value \(j/(2p^{h_i-1})\). The proposed extra divisibility is therefore unavailable. This does not prove optimality of the final theta slope, because Bell terms can cancel.

\subsection{7. Reproducibility and source scope}\label{reproducibility-and-source-scope}

The exact finite certificate is \nolinkurl{jacobian_hessian_pcurvature_certificate.py}, with complete output in \nolinkurl{jacobian_hessian_pcurvature_certificate.json}. It checks the two constants (K18.19), the restricted scalar (K18.17), both exceptions in (K18.16), and all three failures in (K18.6). For modulus \(121\), it computes the initial cyclotomic inverses over \(\mathbb Q[Z]/\Phi_7\) and reduces them only after checking their denominator units; it does not treat \(\mathbb Z/121\mathbb Z\) as a field. It verifies the reduced inverses multiplicatively and records the underlying recurrence SHA256.

The stable integral Abel chart, actual Baker/Chern identification, and the classical Kleinian Hessian identities are fully sourced in AM-K012 and its bibliography. The actual primitive line/section and good Hodge basis are ALP-LIFT-05, ALP-LIFT-07, and ALP-LIFT-08. The normality, character, logarithmic-operator, and transfer arguments needed for the new statements are written in full here. The elementary fact that a normal noetherian domain is the intersection of its height-one valuation rings is the sole normality criterion used; no assertion of reduced-frame integrality is hidden in that step.

This closes a further actual-theta coefficient subgate and exhibits a concrete failure of a naive mixed extension. It does not identify a source-selected harmonic row, prove a matching sharp finite slope, supply the terminal every-\(\varepsilon\) window, or extend the accepted Dirichlet theorem to Kummer RH.

\clearpage

\section{AM-K018 --- exact source canonical gauge for the primitive Jacobian and its ordinary powers}\label{am-k018-exact-source-canonical-gauge-for-the-primitive-jacobian-and-its-ordinary-powers}

Author: \nolinkurl{/root/kummer_extension}, GPT-6 Astra Pro, ALP\_MACHINE, 21 November 2026 UTC. This is an exact comparison of formal gauges, not a claim that a reduced theta frame is integral. The accepted DP47, DP49, DP54, DP58, and DP59 theorems and the accepted source canonical-gauge existence theorem are used without rederiving their established conclusions.

\subsection{1. Statement and scope}\label{statement-and-scope}

Let \(J=\operatorname{Jac}(y^2=x^7-1)\), with its full cyclotomic action \(\mathcal O_E=\mathbb Z[\zeta_7]\), canonical principal polarization, and holomorphic CM type \(\Phi=\{1,2,3\}\). Put \[
w_0=-(a+a^2+a^4),\qquad w_0+w_0^\dagger=1,
\qquad L_0=(1,w_0)^*\mathcal P_\lambda.
\] Use the actual rigidified identification \(L_0=L_{2\infty}^{\rm rig}\) of ALP-LIFT-07 and the normalized section \(s_q\) of ALP-LIFT-08. AM-K012--014 identify its actual Chern-reduced germ as \(\sigma_q\).

At the prime 11, in the actual integral hyperbolic formal block and invariant product frame of DP47/DP49, the following comparison is exact: \[
\boxed{\sigma_q(\log z)=f_q(z)\,U(z)\,B(z)^{-1}.}
\tag{K19.1}
\] Here \(f_q\) is the algebraic section coefficient in the stated integral frame, \(U\) is precisely the normalized integral frame factor in DP47, and \(B\) is its actual primitive arbitrary-unit Poincaré splitting. There is no additional quadratic exponential. More generally every actual section \(s\) of \(L_0^{\otimes\kappa}\), represented in the same frame, has Chern-reduced germ \[
\boxed{\Theta_s(\log z)=f_s(z)U(z)^\kappa B(z)^{-\kappa}.}
\tag{K19.2}
\] The right side is exactly the normalized formal expression in DP54, DP58, and DP59. Equality holds as a complete characteristic-zero formal series, for all \(\kappa\); it is not restricted to small jet degree or small level.

Thus their accepted selected derivative statements apply to the actual Chern/Baker germs, with the same polynomial operator in the same invariant tangent directions. This includes the allowed theta translates in DP59 once their actual section and theta-group scalar are fixed. It does not identify an unspecified external automorphic row or its direction matrix.

The argument also applies to ordinary powers \(J^m\) with diagonal \(\mathcal O_E\)-action and their product polarization, because their holomorphic characters are repetitions of \(1,2,3\). It is \textbf{not} asserted for the whole DP65 family with conjugated \(E\)-actions: such a family may have opposite holomorphic characters across factors and admit nonzero balanced bilinear forms. Extra integral self-adjoint factor idempotents could remove those cross terms, but they have not been assumed here.

\subsection{2. Exact source interfaces actually inspected}\label{exact-source-interfaces-actually-inspected}

The following source files were recovered from the intake-pinned mirror and matched against the original-author archive SHA256 manifest. No Lean compilation or rederivation of the accepted formal theorems is claimed.

\begin{longtable}[]{@{}
  >{\raggedright\arraybackslash}p{(\columnwidth - 2\tabcolsep) * \real{0.5000}}
  >{\raggedright\arraybackslash}p{(\columnwidth - 2\tabcolsep) * \real{0.5000}}@{}}
\toprule\noalign{}
\begin{minipage}[b]{\linewidth}\raggedright
Source
\end{minipage} & \begin{minipage}[b]{\linewidth}\raggedright
Exact interface used
\end{minipage} \\
\midrule\noalign{}
\endhead
\bottomrule\noalign{}
\endlastfoot
\nolinkurl{ReciprocityTheory/ReductiveGroups/ArithmeticGeometry/BirationalGeometry/ToroidalCompactification/CocycleRightZeroThd.lean} & \nolinkurl{FB} at lines 86--158; every clause of \nolinkurl{FB.IsCanonGauge} at 174--183; \nolinkurl{hL} at 185; \nolinkurl{nabVal} at 187--189. \\
\nolinkurl{Development/AbelianVarietyTheory/Proofs/FormalGroupCharts/CanonGaugeSubIsBiadd.lean} & Complete theorem \nolinkurl{canonGauge_sub_isBiadd}, lines 16--60: differences are axis-normalized, additive in both arguments, and adjoint-balanced. \\
\nolinkurl{Development/AbelianVarietyTheory/Proofs/FormalGroupCharts/CanonGaugeIndependenceAll.lean} & \nolinkurl{gaugeExp_holds}, lines 19--21; \nolinkurl{canonGauge_independence_all}, lines 23--26. We do not replace its selected-operator hypotheses by unconditional gauge independence in a different CM type. \\
\nolinkurl{ReciprocityTheory/ReductiveGroups/ArithmeticGeometry/TaylorDiscriminants/WickVanishing/CanonB.lean} & Accepted theorem \nolinkurl{canonGauge : CanonGaugeSt}, line 334, supplies a canonical gauge for a finite free reduced integral algebra with involution; its opening construction and theorem signature were inspected. \\
\nolinkurl{ReciprocityTheory/ReductiveGroups/ArithmeticGeometry/TaylorDiscriminants/ConnectionKillsSeries.lean} & \nolinkurl{CanonGaugeSt}, lines 56--59, quantifies over the full finite free reduced algebra \(O\), with no level or jet-order bound; \nolinkurl{ConnectionIntegrality}, 77--83, is a separate supplied integral-connection statement, not integral reduced-frame coefficients. \\
\nolinkurl{Formalism/LocalToGlobal/ProjectiveLattices/ThGerm.lean} & Full \nolinkurl{IsFrame} clauses, 119--149; actual primitive diagonal frame \nolinkurl{eL}, 166--170; \nolinkurl{IsGerm}, 172--181; \nolinkurl{ChartAt}, 189--206; \nolinkurl{IsNab}, 208--217. In particular the source action algebra in \nolinkurl{ChartAt} is the full \nolinkurl{𝓞 E}, not merely the quadratic field generated by \(w_0\). \\
\end{longtable}

The gauge axioms, in ordinary notation and in the \textbf{same frame}, are \[
\begin{aligned}
h(0,Y)=h(X,0)&=0,\\
\mathrm{ph}_1&=\exp(h(X,Z)+h(Y,Z)-h(X+_FY,Z)),\\
\mathrm{ph}_2&=\exp(h(X,Y)+h(X,Z)-h(X,Y+_FZ)),\\
\mathrm{sy}&=\exp(h(Y,X)-h(X,Y)),\\
\mathrm{ad}(b)&=\exp(h([b]X,Y)-h(X,[b^\dagger]Y))
\quad(b\in\mathcal O_E).
\end{aligned}
\tag{K19.3}
\] The source normalized section is \(f\exp(\kappa h_L)\), where \(h_L(X)=h(X,[w_0]X)\). The signs in (K19.3), and hence in (K19.1), are fixed by these actual source clauses.

\subsection{3. A Chern-reduced biextension frame satisfies all those axioms}\label{a-chern-reduced-biextension-frame-satisfies-all-those-axioms}

Let \(V/\Lambda\) be the complex torus \(J(\mathbb C)\). In the first-argument-linear convention of AM-DO-011 and AM-R-004, its positive Hermitian form is \(H\). A reduced theta frame for a symmetric principal line has multiplier \[
j_L(\lambda,u)=\alpha(\lambda)
 \exp\left(\pi H(u,\lambda)+\frac\pi2H(\lambda,\lambda)\right),
\tag{K19.4}
\] where \(\alpha\) is its semicharacter and \(\operatorname{Im}H\) is integral on the lattice. Formula (K19.4) is the explicit multiplier already proved in AM-DO-011 from the theta series; equivalently it is Bannai--Kobayashi's Definition 1.8.

The symmetric polarization biextension is \(m^*L\otimes\operatorname{pr}_1^*L^{-1}
\otimes\operatorname{pr}_2^*L^{-1}\), with axial rigidifications. Its induced reduced frame on \(V\times V\) has multiplier \[
j_P((\lambda,\mu),(u,v))
=\exp\left(\pi\{H(u,\mu)+H(v,\lambda)+H(\lambda,\mu)\}\right).
\tag{K19.5}
\] Indeed the quotient of the three multipliers (K19.4) gives the two cross linear terms and \(\operatorname{Re}H(\lambda,\mu)\). The semicharacter quotient is \(\exp(\pi i\operatorname{Im}H(\lambda,\mu))\), giving the last term as written.

Both partial addition maps have scalar one in this reduced frame. This follows by adding the corresponding terms in (K19.5): the multiplier for \((u+u',v)\) is exactly the product of those for \((u,v)\) and \((u',v)\), with the shared lattice shift in the second variable, and likewise with the roles reversed. Its axial values are one.

The symmetry map also has scalar one. Swapping \((u,v)\) and \((\lambda,\mu)\) changes the exponent by \(\pi(H(\lambda,\mu)-H(\mu,\lambda))
=2\pi i\operatorname{Im}H(\lambda,\mu)\), whose exponential is one. The map is normalized on the axes.

Finally the polarization adjoint relation is \[
H(bu,v)=H(u,b^\dagger v)\quad(b\in E).
\tag{K19.6}
\] Inserting it in (K19.5) shows that the multipliers of \(P(bu,v)\) and \(P(u,b^\dagger v)\) coincide. Thus their canonical normalized adjunction has scalar one in the reduced frames. These maps realize the actual rigidified Poincaré addition, symmetry, and adjunction maps; uniqueness can also be seen by dividing two such normalized geometric maps, which gives a regular unit on the proper base and hence a constant fixed to one on the axes.

Choose an algebraic formal biextension frame \(e\), normalized on the axes, and let \(e_{\rm red}\) be the reduced frame just constructed. Write \[
e=\exp(h_C)e_{\rm red}.
\tag{K19.7}
\] Its formal logarithm has zero axial values. Computing the frame ratios for the four kinds of structure map gives exactly (K19.3). For example, the source \nolinkurl{IsFrame.ph1} compares its rigidified cubic generator with \(\mathrm{ph}_1 e(X+_FY,Z)e(X,Z)^{-1}e(Y,Z)^{-1}\), giving the negative addition coboundary displayed in (K19.3). Its \nolinkurl{sy} uses \(e(Y,X)/e(X,Y)\), and its \nolinkurl{ad} uses \(e([b]X,Y)/e(X,[b^\dagger]Y)\), giving the two remaining signs. Therefore \(h_C\) is a full-\(\mathcal O_E\) source canonical gauge, not only a candidate quadratic completion of a single section.

\subsection{4. Uniqueness for this full CM type and algebraic descent}\label{uniqueness-for-this-full-cm-type-and-algebraic-descent}

Let \(h,h'\) be two source canonical gauges in the \textbf{same formal frame}, for the full cyclotomic algebra. By the accepted source difference theorem, their difference is biadditive and adjoint-balanced. In normalized group logarithms \(U,V\), biadditivity forces a bilinear form \[
L(U,V)=\sum_{i,j=1}^3c_{ij}U_iV_j.
\tag{K19.8}
\] For completeness, additivity gives \(L(nU,V)=nL(U,V)\) for every positive integer \(n\); comparing homogeneous degrees over characteristic zero leaves first-variable degree one. Applying the same argument in the second variable gives (K19.8).

Take \(b=a\), the order-seven automorphism. Since \(a^\dagger=a^{-1}\), the balance equation gives \[
c_{ij}(\zeta^i-\zeta^{-j})=0.
\tag{K19.9}
\] For \(i,j\in\{1,2,3\}\), one has \(i+j\in\{2,3,4,5,6\}\), so every scalar in parentheses is nonzero. Thus \[
\boxed{h=h'.}
\tag{K19.10}
\] This is stronger than saying merely that an invariant quadratic diagonal vanishes. It does not require the algebraic frame itself to be invariant: the same-frame adjoint cocycles cancel in the difference.

The accepted source theorem \nolinkurl{CanonB.canonGauge} supplies a gauge over the characteristic-zero algebraic coefficient field for the actual algebraic formal data. Its hypotheses apply to \(\mathbb Z[\zeta_7]\), a reduced finite free \(\mathbb Z\)-algebra with conjugation. Compare that algebraic gauge, after complex base change, with \(h_C\). Equation (K19.10) proves equality and hence descent of \(h_C\)'s formal coefficients to the same algebraic field. This avoids treating a potentially nonconvergent \(p\)-adic formal frame as an analytic frame.

Now pass to the 11-adic field. If a chosen integral formal frame \(e'\) differs from an algebraic reference frame by a normalized formal unit \(J\), set \[
h_C'=h_C+\log J.
\tag{K19.11}
\] The logarithm is a characteristic-zero formal logarithm, with no assertion of integral coefficients. Frame substitution in (K19.3) shows that \(h_C'\) is canonical for the new cocycles. Thus uniqueness applies in the precise integral product frame of DP47. There is no illicit comparison of gauges belonging to different frame cocycles.

For \(J^m\) with its ordinary diagonal action, equation (K19.9) applies to every pair of variables, whatever their factor labels, and still forces zero. If different factors have conjugated \(E\)-actions, opposite labels can occur and the coefficient in (K19.9) can vanish. No uniqueness claim for that mixed situation is made.

\subsection{5. The trace-one diagonal and the DP47 primitive kernel}\label{the-trace-one-diagonal-and-the-dp47-primitive-kernel}

Pull the reduced biextension frame back along \(z\mapsto(z,w_0z)\). Its Hermitian form is \[
\begin{aligned}
H_P((z,w_0z),(w,w_0w))
&=H(z,w_0w)+H(w_0z,w)\\
&=H(z,(w_0+w_0^\dagger)w)=H(z,w).
\end{aligned}
\tag{K19.12}
\] The line is the actual \(L_0\), and its rigidified symmetric representative is \(L_{2\infty}^{\rm rig}\) by ALP-LIFT-07. Thus this pulled-back frame is precisely the normalized Chern-reduced frame used for \(\sigma_q\). There is no remaining linear or quadratic exponential: both reduced frames have the same multiplier, and the ratio of the normalized global line identifications is a holomorphic unit on the compact torus, hence the constant one.

This also checks the source primitive diagonal frame. Its \nolinkurl{ThGerm.eL} includes the rigidified primitive cube isomorphism \nolinkurl{wP}, not just the biextension frame. Under (K19.12), that actual isomorphism has constant one in reduced frames, fixed by \nolinkurl{IsPrim} and \nolinkurl{rig0}. Consequently the source coefficient \(f\exp(\kappa h_L)\) is exactly the Chern-reduced section coordinate.

DP49 supplies the actual polarized integral block \(G\times H\) at 11, with the \textbf{full cyclotomic action}, the principal Cartier duality, and the arbitrary-unit display models. In DP47 the invariant product gauge \(h_0\) is shown to satisfy the source axis, addition, symmetry, and adjunction identities for the commutative CM action preserving this decomposition. The full \(\mathcal O_E\)-action preserves the two integral idempotents; all its scalar endomorphisms on the one-dimensional formal factor commute with the Teichmüller action, and the transpose action is its actual adjoint on the other factor. Hence this is a full-\(\mathcal O_E\) canonical gauge, not merely a gauge for the smaller quadratic order generated by \(w_0\).

Compare \(h_0\) with \(h_C'\) in that same product frame. Equation (K19.10) gives \(h_0=h_C'\). The accepted DP47 identity \[
\exp h_{0,L}=U B^{-1}
\tag{K19.13}
\] therefore proves (K19.1)--(K19.2) with its exact sign and its actual integral unit \(U\). Both sides can be nonintegral as complete ordinary power series. This comparison does not require or prove an integral reduced theta frame.

The comparison contains no level parameter. It applies to every integer tensor power, and to the actual odd levels \(M=29^e\) of DP59 with \(\kappa=M^2\). At such a level the prescribed factor \(M^{-n}\) cancels exactly the change in an \(n\)-th derivative under \(u\mapsto Mu\), as already proved in AM-K015. There is no small-degree or small-\(M\) restriction concealed in the canonical-gauge argument.

\subsection{6. Fixed actual-base sharp slopes on two explicit mixed lines}\label{fixed-actual-base-sharp-slopes-on-two-explicit-mixed-lines}

The new gauge comparison permits a direct application of the accepted DP49 theorem to the single fixed section \(s_q\), with explicit lines. In the original \(u\)-coordinates take \[
f_{23}(t)=\sigma_q(0,t,t),\qquad
f_{13}(t)=\sigma_q(t,0,t).
\tag{K19.14}
\] Their leading terms are respectively \(-t^2\) and \(t^2\), by ALP-LIFT-08 or AM-K014. Thus the algebraic prefactor has nonzero restriction on these precise lines: multiplying by any normalized gauge exponential or integral frame unit cannot remove that leading term.

DP49's unit-arrow logarithmic coordinates \((X,Y_1,Y_2)\) have CM characters \((3,2,1)\). The two invariant bases are related by diagonal units at 11, since each is an integral eigenbasis with the same one-dimensional characters. Therefore the first line has unit \(G\)-component and unit first dual component, of restricted-power length two. The second has unit \(G\)-component, zero first dual component, and unit last dual component, of length one. All hypotheses of the fixed-line clauses of DP49 are now verified, including the nonzero prefactor restriction.

Consequently \[
\begin{aligned}
v_{11}(f_{23}^{(n)}(0))
 &\geq \frac n{1220}-\frac{s_{11}(n)}{10},
&\liminf_{n\to\infty}\frac{v_{11}(f_{23}^{(n)}(0))}{n}
 &=\frac1{1220},\\
v_{11}(f_{13}^{(n)}(0))
 &\geq \frac n{120}-\frac{s_{11}(n)}{10},
&\liminf_{n\to\infty}\frac{v_{11}(f_{13}^{(n)}(0))}{n}
 &=\frac1{120}.
\end{aligned}
\tag{K19.15}
\] Zeros are assigned valuation \(+\infty\) in the liminf. These are matching lower/asymptotic upper statements for actual derivatives of a single explicitly normalized section along two fixed mixed lines. The exact derivatives are respectively \((\partial_2+\partial_3)^n\sigma_q(0)\) and \((\partial_1+\partial_3)^n\sigma_q(0)\). They are distinct from individual mixed monomial jets and from the three pure-direction bounds of AM-K013--017.

As a consistency check, AM-K017's complete mod-11 identities give \(f_{13}^{(11)}=4t f_{13}\pmod{11}\); the shortest primitive mixed logarithmic obstruction indeed occurs in degree 12. The 23 direction has algebraic restricted-power length two, consistent with the first DP49 class in the displayed table. These congruences are not used to infer the liminf: the latter is the accepted DP49 conclusion after the now-verified gauge and nonvanishing premises.

For every admitted odd level, the normalized pullbacks of the same section retain both slopes by the exact \(M^{-n}\) cancellation. DP59's prescribed translated span likewise lies in this same actual gauge; its accepted ordinary and first-row harmonic jet witnesses therefore have their stated valuations as Chern-reduced derivatives of those actual sections. Their identification with a particular later automorphic row, or with the mixed-conjugate factors of DP65, remains a separate issue.

\subsection{7. What has been used and what remains}\label{what-has-been-used-and-what-remains}

The source records use full \(\mathcal O_E\)-adjunctions and their generic canonical gauge is supplied for all formal orders. The geometric full-CM structure on this Jacobian and the actual product kernel at 11 are the accepted DP49/DP47 inputs. The new comparison proves that their canonical gauge is the specific Chern/Baker gauge already used for the explicit periods, Gaussian bounds, determinant coefficients, and K017 recurrences. Thus these calculations now concern the same actual formal normalization.

This closes the gauge comparison for ordinary primitive-J powers and the DP59 section family. It does not prove that all those sections occur in a required iterated automorphic image; it does not compare an unknown selected direction matrix or minor; it does not remove all integral lattices, projection costs, or the terminal every-\(\varepsilon\) window. In particular the no-opposite-characters uniqueness proof must not be applied unchanged to the conjugated \(E\)-actions in DP65.

Primary literature for the reduced multiplier is Bannai, K., and Kobayashi, S. (2010), \emph{Algebraic theta functions and the p-adic interpolation of Eisenstein--Kronecker numbers}, Duke Mathematical Journal, 153(2), 229--295, Definition 1.8 in the author preprint \href{https://arxiv.org/pdf/math/0610163}{arXiv:math/0610163v4}; its concrete multiplier was independently computed in AM-DO-011. Exact accepted-result locators are the supplied \nolinkurl{research/DP47_integrated_fragment.tex}, \nolinkurl{DP49_integrated_fragment.tex}, \nolinkurl{DP54_integrated_fragment.tex}, \nolinkurl{DP58_integrated_fragment.tex}, and \nolinkurl{DP59_integrated_fragment.tex} inside \nolinkurl{RETURN_DP_MULT_checkpoint_06_20261121T0420Z}. The source retrieval records are \nolinkurl{../research/source_recovery_20261121T212734Z.json}, \nolinkurl{source_recovery_20261121T213026Z.json}, and \nolinkurl{source_recovery_20261121T213717Z.json}; each compared file matched its intake author hash.

\clearpage

\section{ALP-LIFT-10: the permitted half-torsion orbit in an actual first theta image}\label{alp-lift-10-the-permitted-half-torsion-orbit-in-an-actual-first-theta-image}

Run ALP\_MACHINE; /root/lift\_accounting, GPT-6 Astra Pro; inherited effort not separately exposed. This result uses accepted DP59 and DP62 for the algebraic permitted sections and their canonical normalizations. It does not re-prove their integral span or jet separation. It uses ALP-LIFT-09 for the actual global unitary splitting, first-lift integral, and torus class calculation, and AM-DO-020--021 for the fixed split and ramified compact phases.

\subsection{Statement}\label{statement-10}

Let \(J\), \(L\), \(s_q\), \(\sigma_q\), \(\Lambda\), and \(H\) have the normalizations of ALP-LIFT-08--09. Thus \(\sigma_q\) is the actual Chern-reduced germ of \(s_q\) with leading term \(u_1u_3-u_2^2\), and \(s_q\) is primitive at every odd prime other than7. Put \[
 \delta_7^2=-7,\qquad
 \alpha=(1+2\delta_7)^e,\qquad M=29^e,\qquad \kappa=M^2,
 \qquad e\ge1.
 \tag{O1}
\] Here \(\delta_7=1+2(\zeta+\zeta^2+\zeta^4)\), so the choice is exactly DP59's actual level. For every permitted root \(t'\in J[\alpha^2]\), let \(s_{t'}\) be the canonical symmetric theta translate of \([M]^*s_q\) from DP59.

Fix an even positive integer \(n\), and use the actual compact pair \(T=U(1)(E/F)\) and \(G=U(I_n)(E/F)\), with the honest splitting \(\chi_V=1,\chi_W=\mu_E\) of ALP-LIFT-09. Let \(P\) be any polynomial differential operator in the \(3n\) algebraic differential coordinates, homogeneous of total degree \(N\) and of a fixed degree at each of the three CM embeddings. All such polynomials have one fixed archimedean torus weight, denoted \(\beta\); it includes the actual vacuum character.

There is one explicitly defined automorphic input \(f_{M,\beta}\) on \(T\), independent of the allowed tuple \(\mathbf t'=(t'_1,\ldots,t'_n)\) and of \(P\) inside this weight, such that an actual theta lift has identity value \[
 \Theta(\phi_{\mathbf t',P},f_{M,\beta})(1)
  = C^n M^{-N}
       P(\partial_u)
         \left(\prod_{j=1}^n
                    \sigma_{s_{t'_j}}(u_j)\right)_{u=0},
 \qquad C=\theta[0;0](0;\tau)c.
 \tag{O2}
\] The notation \(\sigma_{s_{t'}}\) denotes the specified Chern-reduced germ of the algebraic section \(s_{t'}\), in the line of curvature \(\kappa H\). The constant \(c\ne0\) is fixed by \(\Sigma_{\rm tr}=c\sigma_q\); it is independent of \(M\) and the tuple. The finite Schwartz vector in (O2) is an explicitly specified Heisenberg translate of the fixed vacuum, with coefficients roots of unity and finite \(L^2\) norm one.

The automorphic input is supported on one finite component and has the exact norms \[
 \|f_{M,\beta}\|_2^2=h_M
       =56^2\,29^{3(e-1)},\qquad
 \|f_{M,\beta}\|_\infty=h_M.
 \tag{O3}
\] Consequently its normalization to norm one multiplies the right side of (O2) by \(h_M^{-1/2}\) and has logarithmic cost \[
 \log\sqrt{h_M}
      =\frac32\log M+\log56-\frac32\log29.
 \tag{O4}
\] Every prime divisor of \(h_M\) lies in \(\{2,7,29\}\), so both the projection and norm normalization have zero valuation cost at11, and at every prime outside this displayed exceptional set. At the DP59 levels \(M=O(N)\), the new logarithmic norm cost is \(O(\log N)\).

This is a simultaneous realization of the specified fixed-weight jet functionals by actual first theta lifts. It is not an identification of the entire section module with an automorphic representation, an assertion about one freely prescribed irreducible torus character, or a construction of the later definite lifts.

\subsection{1. Exact multiplication and translation of the actual reduced section}\label{exact-multiplication-and-translation-of-the-actual-reduced-section}

Write \[
 (U_x^H F)(v)
       =\exp\bigl(-\pi H(v,x)-\pi H(x,x)/2\bigr)F(v+x).
 \tag{O5}
\] This is the normalized Mumford translation used in AM-DO-017--018. The rigidified multiplication isomorphism \([M]^*L\simeq L^{M^2}\) carries the reduced section to \(\sigma_q(Mu)\) exactly. Indeed the pullback Hermitian form is \(M^2H\), its semicharacter is \(\alpha_{\rm tr}(M\lambda)=\alpha_{\rm tr}(\lambda)^{M^2}\), and the two normalized Appell--Humbert presentations have the same multiplier and the same origin rigidification. Their isomorphism is therefore the identity in these presentations; no extra quadratic exponential appears.

Choose a rational period representative \(r\in\kappa^{-1}\Lambda\) of \(t'\), and put \(x=Mr\). The equality of the two translated functions is immediate from (O5): \[
 U_r^{M^2H}\bigl(\sigma_q(Mu)\bigr)
                 =(U_x^H\sigma_q)(Mu).
 \tag{O6}
\] Both linear exponential terms are \(-\pi M^2H(u,r)\), both constant exponential terms are \(-\pi M^2H(r,r)/2\), and both arguments are \(Mu+Mr\).

The canonical symmetric lift in DP59 differs from \(U_r^{M^2H}\) by an explicitly determined sign only. To see this, let \(\iota\) be normalized inversion in the line presentation. Then \[
 \iota U_r^{M^2H}\iota^{-1}=U_{-r}^{M^2H}
                       =(U_r^{M^2H})^{-1}.
 \tag{O7}
\] Since the alternating self-pairing is zero, \((U_r^{M^2H})^\kappa=U_{\kappa r}^{M^2H}\), and the latter is the scalar \(\alpha_{\rm tr}(\kappa r)^\kappa
 =\alpha_{\rm tr}(\kappa r)\in\{\pm1\}\). Because \(\kappa\) is odd, multiplying \(U_r^{M^2H}\) by \[
                        \epsilon_r=\alpha_{\rm tr}(\kappa r)
 \tag{O8}
\] gives an operator whose \(\kappa\)-th power is one and which still satisfies (O7). The uniqueness of the symmetric finite-order lift, already proved in DP59 §4, identifies it with the permitted lift. In particular, \[
                 \sigma_{s_{t'}}(u)
                    =\epsilon_r(U_x^H\sigma_q)(Mu).
 \tag{O9}
\] Changing the representative gives the same section by that uniqueness. All phases introduced in (O8) are algebraic units.

The point represented by \(x\) lies in the smaller subgroup \[
                        [M]J[\alpha^2]=J[\alpha]\subset J[M].
 \tag{O10}
\] The inclusion follows from \(\alpha M=\alpha^2\bar\alpha\) and \(\alpha\bar\alpha=M\). Equality follows from the accepted DP59 kernel calculation: the kernel of \([M]\) on \(J[\alpha^2]\) is \(J[\alpha]\), so the image has the same order as \(J[\alpha]\). This is where the actual permitted half-orbit is used; replacing it by all of \(J[M]\) would lose the sharper level below.

Taking tensor products of (O9), and using homogeneity of \(P\), gives the all-degree identity \[
 M^{-N}P(\partial_u)\prod_j\sigma_{s_{t'_j}}(u_j)\big|_0
   =\left(\prod_j\epsilon_{r_j}\right)
            P(\partial_v)\prod_j(U_{x_j}^H\sigma_q)(v_j)\big|_0.
 \tag{O11}
\] The factor \(M^{-N}\) from the source therefore cancels the multiplication pullback exactly.

\subsection{2. The exact finite Schwartz family}\label{the-exact-finite-schwartz-family}

Let \(\phi_{0,f}\) be ALP-LIFT-09 (V12). For the selected representatives define \[
 \phi_{\mathbf t',f}
       =\bigotimes_{j=1}^n\bigl(\epsilon_{r_j}W_{x_j}\phi_{0,f}\bigr).
 \tag{O12}
\] The \(W_x\) are the explicit finite operators (V13). Their translates and modulation factors preserve compact support and local constancy. Their nonzero values have absolute value one and are roots of unity, because all period coefficients are rational. They are unitary, so (O12) has norm one.

The exact reading identity of AM-DO-018 and \(R\phi_{0,f}=C\sigma_q\) give \[
 R\phi_{\mathbf t',f}(v)
       =C^n\left(\prod_j\epsilon_{r_j}\right)
                      \prod_j(U_{x_j}^H\sigma_q)(v_j).
 \tag{O13}
\] There is one common period \(C^n\) for the whole family. Use the actual archimedean Schwartz vector \[
                   \phi_{P,\infty}
                         =P(r_u)\varphi_{0,\infty}^{\otimes n},
 \tag{O14}
\] with the creation operators of ALP-LIFT-09 (V19)--(V21). Its actual torus character is \(\beta\). Formula (V20) identifies the theta sum at the identity with the derivative of (O13). Thus the identity value on the right side of (O2) is already an actual theta series. It remains to realize it by one common automorphic input with its normalization controlled.

\subsection{3. The fixed compact vacuum character for even rank}\label{the-fixed-compact-vacuum-character-for-even-rank}

For the fixed splitting \(\chi_V=1\) and even \(n\), the scalar norm-one compact character of the base vacuum is \[
                  \lambda_f(k)=\epsilon_7(k_7)^{n/2},
 \tag{O15}
\] where \(\epsilon_7(t)\) is the residue of \(t\) in \(\{1,-1\}\) at the unique ramified prime over7. All other finite local factors are trivial.

Here is the precise justification of the unramified part, supplementing AM-DO-020--021's calculations. At every odd unramified quadratic place, the norm on integral units is surjective. For example the residue norm is surjective, and successive correction in the principal units uses the surjective residue trace. Choose a unit of norm \(-1\). It gives an integral primitive isotropic vector of \(I_2\); an integral pairing partner can be adjusted, using that 2 is a unit, to make an integral hyperbolic pair. Therefore \(I_n\) is an integral orthogonal sum of hyperbolic planes. Pan (2001), Lemma 6.2, printed p.197, gives scalar \(\alpha(t)\beta^Y(t)=\chi_V((t-1)\delta)=1\) against each such plane, with the identity element treated directly. His tensor identity (3.6.e) gives one for their sum. The selfdual trace-lattice normalization is the same lattice-model normalization identified in AM-DO-021 §1.

At relative split places, including2, AM-DO-020 gives the exact split determinant-root scalar. It is one on the compact torus when \(\chi_V=1\). The only remaining ramified place is7, and AM-DO-021 (DO21.12) gives exactly (O15). The simultaneous sign transport in ALP-LIFT-09 does not change this residual real sign. Thus the fixed compact subgroup \[
                  K_7^+=\ker\epsilon_7\subset K_{0,7}
 \tag{O16}
\] has index2 and fixes the vacuum for every even \(n\). When \(4\mid n\) this restriction is redundant, but retaining it gives one uniform exact level and cost for all even ranks.

\subsection{\texorpdfstring{4. The permitted half-orbit is fixed at level \(M\), not \(M^2\)}{4. The permitted half-orbit is fixed at level M, not M\^{}2}}\label{the-permitted-half-orbit-is-fixed-at-level-m-not-m2}

The prime29 splits completely in \(E/\mathbf Q\). There are three places \(v\) of \(F\) above29, each with \(F_v=\mathbf Q_{29}\) and \(E_v=F_v\times F_v\). Since the ideals \((\alpha)\) and \((\bar\alpha)\) are coprime and their product is \((M)\), the two valuations of \(\alpha\) in each relative pair are \((e,0)\) in one order.

The actual lattice decomposes at such a place into its two relative split isotropic summands, \[
                      \Lambda_v=X_v\oplus Y_v.
 \tag{O17}
\] Both are stable under \(w_0\) and the scalar norm-one torus. A point of \(J[\alpha]\) can be represented modulo \(\Lambda_v\) entirely in the summand on which \(\alpha\) has valuation \(e\), with its coordinate in \(29^{-e}X_v\) (or in \(29^{-e}Y_v\) in the other order).

Replacing a representative by an integral vector changes \(W_x\phi_{0,f}\) by a fixed root of unity: this follows from the Heisenberg multiplication law and the lattice vacuum eigenrelation. Such a scalar does not affect its stabilizer. It is therefore legitimate to use the pure isotropic representative at each place when computing the common stabilizer of the actual vectors (O12).

Take \[
 K_v(M)=\{t=(a,a^{-1})\in K_{0,v}:
                                  a\equiv1\pmod{29^e}\}.
 \tag{O18}
\] For \(x\in29^{-e}X_v\) and \(t\in K_v(M)\), put \(d=(t-1)x\). Then \(d\in X_v\subset\Lambda_v\). Both \(x\) and \(d\) lie in the same isotropic space, so \[
                  \mathfrak b(x,d)=0,\qquad
                  \mathfrak b(d,w_0d)=0.
 \tag{O19}
\] The lattice eigencharacter of \(d\) is consequently one, and the half-pairing in the multiplication law is zero. Thus \[
                         W_{tx}\phi_{0,f}=W_x\phi_{0,f}.
 \tag{O20}
\] The base compact scalar is one by (O15). Heisenberg covariance therefore proves that \(K_v(M)\) fixes every allowed translated vacuum exactly. The identical argument applies when the denominator lies in \(Y_v\).

At every finite prime not dividing29, the vector \(x\) is integral in the trace lattice, so \(W_x\phi_{0,f}\) is already a scalar multiple of the base vacuum. Every compact element fixing the base vacuum fixes that translate. Consequently the explicit product subgroup \[
 K_M=K_7^+\times\prod_{v\mid29}K_v(M)
             \times\prod_{v\nmid7\cdot29}K_{0,v}
 \tag{O21}
\] fixes every finite vector (O12), and their entire span, for every allowed tuple.

\subsection{5. Exact finite index and automorphic input}\label{exact-finite-index-and-automorphic-input}

The three factors at29 have indices \[
 [K_{0,v}:K_v(M)]
       =\#(\mathbf Z/29^e\mathbf Z)^\times
       =28\,29^{e-1}.
 \tag{O22}
\] The index at7 is2. Reduction at29 is injective on \(\mu_{14}\), since its order is prime to29, so \(K_M\cap\mu_{14}=1\). No extra auxiliary prime is required. ALP-LIFT-09 (V6) now gives \[
 h_M=\frac{2(28\,29^{e-1})^3}{14}
       =56^2\,29^{3(e-1)}.
 \tag{O23}
\] In particular \(v_{11}(h_M)=0\).

Define \(f_{M,\beta}\) by ALP-LIFT-09 (V33), with \(K=K_M\): it equals \(h_M\beta(t_\infty)\) on the identity finite component and zero on all other finite components. The definition is unambiguous because \(K_M\) kills the rational root-unit stabilizer. That component has measure \(1/h_M\). Consequently \[
 \Theta(\phi_{\mathbf t',f}\otimes\phi_{P,\infty},
                             f_{M,\beta})(g)
       =\vartheta_{\mathbf t',P}(g,1).
 \tag{O24}
\] This is a genuine automorphic theta integral for the fixed global splitting. Equations (O11), (O13), and (O14) give (O2) on setting \(g=1\). The elementary measure calculation in (V34) gives (O3), and taking its square root gives (O4). The same input works simultaneously for every allowed tuple and every \(P\) in the specified torus weight.

The input may be decomposed into \(h_M\) automorphic characters by (V36)--(V38). Its finite coefficients are \(h_M\) times roots of unity. Individual-character projections divide by \(h_M\); the simultaneous input (O24) has no such denominator. Even after normalization to norm one, \(h_M^{-1/2}\) is an 11-adic unit after a finite coefficient-field extension. Thus the algebraic value in (O2), after removing the retained common period \(C^n\), has exactly the same 11-adic valuation with the unit-norm input as it had before that normalization. This assertion is about the displayed value; it does not assert primitivity of every target row or compute the target mass norm.

For \(n=2\) this applies directly to pairs of DP59's ordinary allowed sections, including any homogeneous two-column differential chosen on ordinary powers of \(J\). Tensoring a fixed number of ordinary factors does not change the torus level index (O23).

\subsection{6. Extracting one norm-one character without losing the local valuation}\label{extracting-one-norm-one-character-without-losing-the-local-valuation}

The simultaneous input is useful for realizing a whole specified finite span. For an existence statement one may instead extract one of its irreducible-character components with norm one. The finite torus \(K_0\) preserves \(J[\alpha]\), since its scalar integral action commutes with \(\alpha\). Therefore, for every \(k\in K_0\), Heisenberg covariance gives \[
 \omega_f(k)\phi_{\mathbf t',f}
       =\lambda_f(k)\bigotimes_j
                    \bigl(\epsilon_{r_j}W_{kx_j}\phi_{0,f}\bigr).
 \tag{O26}
\] The tuple \((kx_j)_j\) still belongs to the permitted half-orbit. Choosing preimages under \([M]:J[\alpha^2]\to J[\alpha]\) and their rational period representatives identifies (O26) with an allowed tuple's vector, up to a root of unity. In particular, there is still the same common period \(C^n\).

Let \(f_{M,\beta}=\sum_\eta\chi_\eta\) be (V37). At the identity, each normalized character lift is an average \[
 C^{-n}\Theta(\phi_{\mathbf t',P},\chi_\eta)(1)
       =\frac1{h_M}\sum_{a\in K_0/(\mu_{14}K_M)}
                   \overline{\eta(a)}\,J_a,
 \tag{O27}
\] where every \(J_a\) is the normalized Chern jet of an allowed tuple, multiplied by a root of unity. Thus all these values are algebraic in a finite common coefficient field. The only displayed averaging denominator has valuations \[
 v_2(h_M)=6,\qquad v_7(h_M)=2,\qquad
 v_{29}(h_M)=3(e-1),\qquad v_p(h_M)=0\quad(p\nmid2\cdot7\cdot29).
 \tag{O28}
\] These are exact costs of the written projection formula; actual cancellation may improve them.

Since summing (O27) over all characters recovers the value in (O2), the ultrametric inequality has the following useful consequence. If a nonzero normalized Chern jet in (O2) has \(p\)-adic valuation \(V\), then at least one actual character lift in (O27) has valuation at most \(V\). Every such input character has probability-Haar \(L^2\) norm one. At11, and indeed at every \(p\nmid2\cdot7\cdot29\), (O27) additionally lies in the local integer-ring span of the displayed Chern-jet values. Neither assertion identifies in advance which character occurs.

There is an exact archimedean comparison as well. If the allowed normalized Chern jets in this fixed orbit are all bounded by \(B_P\), then each average in (O27) has absolute value at most \(B_P\): its \(h_M\) coefficients have modulus \(1/h_M\). A bound of this kind follows directly from the actual Gaussian estimate for \(\sigma_q\). If \[
 |\sigma_q(u)|\le A\exp\left(\frac\pi2\sum_i|u_i|^2/G_i\right),
 \tag{O29}
\] then (O5) preserves this same envelope exactly: expanding the quadratic form at \(u+x\) cancels its cross term and its \(H(x,x)\) term against the two real translation exponents. For \(P=\sum_\nu a_\nu D^\nu\) on \(n\) copies, polydisk Cauchy gives \[
 B_P=A^n\sum_\nu|a_\nu|\,\nu!
       \prod_{i,j:\nu_{ij}>0}
           \left(\frac{\pi e}{\nu_{ij}G_i}\right)^{\nu_{ij}/2}.
 \tag{O30}
\] The multiplication factor \(M^N\) has already canceled in (O11). This is a bound for these identity values; it is not an evaluation of the target automorphic \(L^2\) norm or of all source row-operator norms.

Thus two exact normalizations are available. One common input for the entire span has norm cost (O4). An extracted character has norm one, while its finite projection formula has the explicit denominator (O28). The latter has total logarithmic size \(\log h_M=3\log M+O(1)\), supported at the specified exceptional primes. Neither bookkeeping choice makes these costs disappear from a subsequent global comparison.

\subsection{7. Exact source-gauge transfer and explicit sharp first-image values}\label{exact-source-gauge-transfer-and-explicit-sharp-first-image-values}

\subsubsection{The source formal gauge is now matched on ordinary powers}\label{the-source-formal-gauge-is-now-matched-on-ordinary-powers}

Equations (O2), (O9), and (O11) identify the actual Chern-reduced germs of the permitted algebraic sections with an actual first theta integral. AM-K018 now proves the required comparison with the source's canonical formal gauge: \[
 \text{the source selected formal jet in its canonical gauge}
   =\text{the normalized Chern jet appearing in (O2)}.
 \tag{O25}
\] This is the same full-\(\mathcal O_E\) canonical gauge in the same integral biextension frame, including the DP47 normalized integral unit. The difference of two such gauges is adjoint-balanced bilinear in group logarithms, and no pair of ordinary holomorphic labels \(1,2,3\) is opposite modulo7. AM-K018 consequently identifies the Chern frame exactly with \(f_sU^\kappa B^{-\kappa}\), with no quadratic-exponential ambiguity. Its statement applies to all ordinary powers and every tensor level and permitted section here. I read that full proof and the consequential source clauses; the independent review is review\_kummer\_source\_gauge.md.

This comparison does not say that the reduced frame is integral. It identifies the nonintegral primitive-kernel factor with the precise source normalization whose derivative estimates were proved in DP54/DP58/DP59.

\subsubsection{One actual sharp harmonic first image}\label{one-actual-sharp-harmonic-first-image}

Take \(n=2\), and retain the first-row two-column harmonic operator \(P_{a,b}\) of DP51/DP58, with \[
 a=11^q,\qquad b=121\cdot11^q,\qquad
 N=a+b=122\cdot11^q,\qquad V_q=(11^q-1)/10 .
 \tag{O31}
\] Use the global algebraic eigenvectors dual to \(\omega_3\) and \(\omega_2\) for its two directions. Its supported monomials all have total degree \(a\) in the first direction and total degree \(b\) in the second, across the two columns. Hence it has one CM torus weight, even though it has several differential monomials. The DP49 unit-arrow eigenvectors differ from these actual algebraic CM eigenvectors by diagonal 11-adic units. Homogeneity in each label multiplies the entire selected harmonic value by a single unit, so the sharp valuation is unchanged. No algebraicity of those display-normalization comparison units is assumed or needed; the Fock polynomial itself uses the global algebraic directions just specified.

DP59, used as an accepted theorem, supplies \(M=29^e=O(N)\) and an allowed pair for which the source-normalized derivative has valuation exactly \(V_q\). Equation (O25) transfers that assertion to the actual Chern jet in (O2). Thus an actual \(U(1)\longrightarrow U(2)\) theta lift has \[
 v_{11}\!\left(C^{-2}
       \Theta(\phi_{\mathbf t',P_{a,b}},f_{M,\beta})(1)\right)
   =V_q .
 \tag{O32}
\] The common input has squared norm \(h_M\) as before; its normalization by \(\sqrt{h_M}\) is an 11-adic unit.

More strongly, at least one of the actual norm-one character inputs in (O27) has the same exact valuation \(V_q\). To check equality rather than only an upper bound, every orbital term in (O27) is a permitted pair, up to a root of unity. DP59's matching harmonic lower bound gives valuation at least \(V_q\) for each term. The denominator \(h_M\) is an 11-adic unit, so each projected value also has valuation at least \(V_q\). Their sum is (O32); at least one therefore has valuation exactly \(V_q\). This proves existence of an actual irreducible torus input, although it does not specify which character in advance.

The new conclusion is the first-image realization and its explicit projection cost. The sharp valuation and existence of the algebraic permitted pair themselves are the accepted DP59 result, not new assertions of this note.

\subsubsection{Ordinary mixed-line values require several torus weights}\label{ordinary-mixed-line-values-require-several-torus-weights}

For an ordinary DP58 line with one nonzero dual eigenaxis, use the global algebraic operator \(P=(D_3+D_j)^N\), with \(j=1\) or2. In the unit-arrow coordinates it becomes \((uD_X+vD_Y)^N\), with \(u,v\) units, exactly as admitted by DP58/DP59. Put \(N=(11^j+1)11^q\). It has \(N+1\) distinct CM multidegrees. It must not be inserted into the one-weight statement while keeping a single input \(f_{M,\beta}\).

Choose a permitted spectator section with unit value at the origin. Such a section exists by DP59's integral zero-jet surjectivity and integral permitted-orbit basis: if all its basis sections vanished modulo11 at the origin, no integral linear combination could have unit value. The Chern gauge has constant term one, so the same statement holds in the actual reduced frame. We work with this spectator in the second ordinary \(J\) factor, and the ordinary sharp section in the first, keeping target rank2.

Decompose \(P=\sum_{\ell=0}^N P_\ell\) into its distinct CM weights \(\beta_\ell\), including the binomial and direction coefficients. Put \[
 \phi_{P,\infty}=\sum_{\ell=0}^N\phi_{P_\ell,\infty},
 \qquad
 f_{M,P}=\sum_{\ell=0}^N f_{M,\beta_\ell}.
 \tag{O33}
\] Orthogonality of the different characters on the compact archimedean torus kills every cross term in the theta integral. The lift of this pair is exactly the sum of the individual-weight lifts, so at the identity it is the desired full ordinary derivative, multiplied by the spectator value and \(C^2\). Orthogonality also gives the exact input norm \[
 \|f_{M,P}\|_2^2=(N+1)h_M .
 \tag{O34}
\] Thus the normalized algebraic identity value has valuation \(V_q=(11^q-1)/10\), and the additional input logarithmic norm cost is \[
 \frac12\log(N+1)+\frac32\log M+
                  \log56-\frac32\log29=O(\log N).
 \tag{O35}
\] Here \(N+1\) is an 11-adic unit at every displayed degree: for \(q>0\) it is \(1\) modulo11, and for \(q=0\) it is \(2\) modulo11. Normalizing the whole input to norm one therefore does not change its 11-adic valuation. Other prime factors of \(N+1\) are real additional normalization denominators and are retained; the exceptional set (O28) alone applies only to the one-weight construction.

One may extract an individual automorphic character from (O33), obtaining a value of valuation at most \(V_q\). Equality for that selected weight does not follow from the ordinary-line lower bound, because the latter bounds the sum of its CM-weight components. No separate lower bound for arbitrary mixed monomial derivatives is asserted here.

\subsection{8. Remaining full-row and later-image comparisons}\label{remaining-full-row-and-later-image-comparisons}

\subsubsection{The full mixed-action row product and later lifts}\label{the-full-mixed-action-row-product-and-later-lifts}

The compact pair here uses ordinary powers \(J^n\). DP65's full-row product is instead \(J^P\times\bar J^{\,Q}\) with the conjugated \(E\)-action retained on the second block. Conjugating the scalar \(E\)-action changes the sign of the trace alternating form in that block. A single scalar-torus dual pair with this action is naturally realized with Hermitian signature \((P,Q)\), not the positive form \(I_{P+Q}\). The projection argument by itself does not remove this real-form distinction. Thus (O24) is not silently promoted to DP65's complete mixed-action harmonic image.

Even for the ordinary compact family, \(f_{M,\beta}\) is a specific finite sum of automorphic characters. Requiring one prescribed irreducible character, requiring an earlier selected image, or passing through a next definite theta lift imposes further representation-theoretic conditions. The exact first-image values and small additional cost in (O2)--(O4) do not prove survival through those conditions. No complete six-identity family or terminal mass pairing is constructed here.

There is a further global field distinction. The actual torus and first target here are over \(E/F=\mathbf Q(\zeta_7)/F\). AM-DO-022's all-finite even-rank family instead uses the DP57 host \(E'=\mathbf Q(\zeta_{28})\) acting on \(J^2\). These are not already a common tower. A restriction/extension-of-scalars kernel comparison or a new first-torus construction over that host would be required to join them.

\subsection{Exact accepted and primary locators}\label{exact-accepted-and-primary-locators}

DP\_MULT checkpoint06, DP59: DP\_MULT.tex, theorem at lines4815--4853, especially (TS1)--(TS4) and the factor \(M^{-n}\) in (TS5); research/theta\_translate\_span\_DP59.md, §§1,3,4,5,8. The subgroup, canonical symmetric finite-order lifts and integral span are carried accepted inputs.

DP\_MULT checkpoint06, DP62: DP\_MULT.tex, theorem at lines5321ff; research/theta\_translate\_norm\_integrality\_DP62.md. Its canonical section norm and good-prime integrality are carried accepted inputs, not identified with a target automorphic mass norm.

ALP-LIFT-09, actual\_torus\_first\_lift.md, (V1)--(V38); AM-DO-018, definite\_oscillator\_moved\_family.md, (DO18.5)--(DO18.7); AM-DO-020 split-place phase calculation; AM-DO-021, definite\_oscillator\_ramified\_phase.md, especially (DO21.12). AM-K018, jacobian\_source\_gauge\_match.md, sections2--5 and (K19.1)--(K19.13), supplies the exact ordinary-power source gauge comparison used in (O25). Its independent review is review\_kummer\_source\_gauge.md.

Pan, S.-Y. (2001). Splittings of the metaplectic covers of some reductive dual pairs. Pacific Journal of Mathematics, 199(1), 163--226. \url{https://msp.org/pjm/2001/199-1/pjm-v199-n1-p10-p.pdf}. Directly checked here: Lemma6.2, printed p.197, including its displayed proof and concluding scalar; the ramified compact specialization is cited under AM-DO-021's separately checked proof.

\subsection{Check status}\label{check-status-1}

The author read the entire saved proof, including (O1)--(O30), the fixed compact phases, the exact permitted half-orbit stabilizer, the finite index, both input normalizations, and the named remaining real-form and later-image gates. The gauge-transfer and sharp-first-image extension (O25), (O31)--(O35) received a separate complete reading after the whole-proof review of AM-K018. The exact finite algebra certificate \nolinkurl{permitted_orbit_certificate.py} passed. Events and their scopes are recorded in \nolinkurl{lift_accounting_CHECK_EVENTS.tsv.} No formal proof assistant execution is claimed.

\clearpage

\section{AM-DO-022: the full finite moved-vacuum identity on the repaired CM host}\label{am-do-022-the-full-finite-moved-vacuum-identity-on-the-repaired-cm-host}

Author: GPT-6 Astra Pro, ALP\_MACHINE \nolinkurl{/root/definite_oscillator}; inherited reasoning effort not separately exposed; 21 November 2026 UTC.

This finding carries DP57's field and integral-action assertions at their accepted status. It supplies a new finite Weil-phase comparison for the explicit mixed CM theta construction of AM-DO-017--020. The comparison is with the standard unitary splitting attached to the displayed characters, not with an arbitrary representation having only some of the source's axioms. The unchanged elliptic \nolinkurl{VacIn} is not an input to this proof.

\subsection{1. Statement with the exact product trace form}\label{statement-with-the-exact-product-trace-form}

Put

\[
 E=\mathbf Q(\zeta_7),\quad \xi=(\zeta_7-\zeta_7^{-1})/7,
 \quad E'=E(i)=\mathbf Q(\zeta_{28}),\quad F'=(E')^+,
 \quad A=J\times J.
 \tag{DO22.1}
\]

DP57 proves that \(E'/F'\) is unramified at every finite place, \(\mathcal O_{E'}=\mathcal O_E[i]\), and that the matrix \(I=\left(\begin{smallmatrix}0&-1\\1&0\end{smallmatrix}\right)\) gives the integral action of \(i\) on \(A\). Its product principal polarization has Rosati involution equal to conjugation on \(E'\). Its CM type is

\[
                         \Phi'=\{1,3,9,15,17,23\}\pmod{28}.
 \tag{DO22.2}
\]

The trace coefficient for this exact lattice and polarization is

\[
                                  \boxed{\xi'=\xi/2.}
 \tag{DO22.3}
\]

For positive even ranks \(a,b\), let \(S=E'^a\), \(T=E'^b\) have their identity Hermitian forms, and put

\[
 \Lambda_{00}=\operatorname{Hom}_{\mathcal O_{E'}}
                    (\mathcal O_{E'}^a,\mathcal O_{E'}^b).
 \tag{DO22.4}
\]

Its rational alternating form is \(E_H(X,Y)=-\operatorname{Tr}_{E'/\mathbf Q}
  (\xi'\operatorname{tr}(Y^*X))\), with the convention \(H\) linear in the first argument and \(E_H=\operatorname{Im}H\). Use \(\psi_\infty(t)=e^{2\pi it}\), and the product-compatible finite character \(\psi_f\). Let \(\omega_\chi\) be the standard unitary Weil representation in this realized Heisenberg model. Its two splitting characters satisfy

\[
 \chi_S|_{\mathbb A_{F'}^\times}=1,
 \qquad\chi_T|_{\mathbb A_{F'}^\times}=1,
 \tag{DO22.5}
\]

because \(a,b\) are even. In particular the globally trivial characters are permitted. The unitary Hom action is \(X\mapsto hXg^{-1}\). Character labels here use the recovered source convention \nolinkurl{LocalNewformFunction.splitFacO}: the source factor is the inverse of \(\chi_T\) evaluated at a determinant root. In the conventional Harris--Kudla--Sweet tensor labels used by Konno and Konno (2007), §3.3 and Lemma 3.1, this is the pair \((\xi,\xi')=(\chi_S,\chi_T^{-1})\), after the group and Heisenberg conventions are identified as in corrected AM-DO-023. This explicit labeling matters for nontrivial characters; it has no effect on the trivial-character comparison proved below.

Use the trace-one element \(w_0=-(\zeta_7+\zeta_7^2+\zeta_7^4)\), so \(w_0+\bar w_0=1\). Let \(\phi_{00}\) be the actual finite Schwartz input constructed in AM-DO-016 and AM-DO-018, with two copies of the Jacobian input for each of the \(ab\) entries. For finite adelic unitary elements \(g,h\), write

\[
 N=g\widehat{\mathcal O}_{E'}^a\cap E'^a,\qquad
 N'=h\widehat{\mathcal O}_{E'}^b\cap E'^b,
 \quad \Lambda_{N,N'}=\operatorname{Hom}_{\mathcal O_{E'}}(N,N').
 \tag{DO22.6}
\]

Let \(\phi_{N,N'}\) be the prescribed two normalized finite averages, first moving the source lattice and then the target lattice, exactly as in (DO18.3)--(DO18.7). Define the positive indices

\[
\begin{split}
 q_1&=[\Lambda_{N,0}:\Lambda_{N,0}\cap\Lambda_{00}],\\
 q_2&=[\Lambda_{N,N'}:\Lambda_{N,N'}\cap\Lambda_{N,0}].
\end{split}
 \tag{DO22.7}
\]

For determinant roots \(x/\bar x=\det g\), \(y/\bar y=\det h\), the exact equality at \textbf{all finite places} is

\[
 \boxed{\omega_\chi(g,h)\phi_{00}
            =\chi_T(x)^{-1}\chi_S(y)\phi_{N,N'}.}
 \tag{DO22.8}
\]

The reversed order of the two prescribed averages gives the same vector. The characters in (DO22.8) are evaluated at the roots, not at the determinants. Changing a root by an idele of \(F'\) leaves the factor unchanged by (DO22.5).

In particular, the ramified character of AM-DO-021 disappears after this specified repair. It is not discarded: no finite place of the new relative CM extension is ramified. At the places over 7 the extension is now unramified quadratic, and the argument below computes the phase there.

\subsection{2. Proof of the product lattice, polarization, and dyadic index}\label{proof-of-the-product-lattice-polarization-and-dyadic-index}

The \(\mathcal O_E\)-linear map

\[
          \mathcal O_E[i]\longrightarrow\mathcal O_E^2,
                       x+iy\longmapsto(x,y)
 \tag{DO22.9}
\]

is an isomorphism of free modules and intertwines multiplication by \(i\) with \(I\). For arbitrary \(x,y,x',y'\in E\), taking the relative trace gives

\[
\begin{split}
 &\operatorname{Tr}_{E'/\mathbf Q}
  \left[\frac{\xi}{2}(x+iy)(\bar x'-i\bar y')\right]\\
 &\hspace{15mm}=
      \operatorname{Tr}_{E/\mathbf Q}
             \left[\xi(x\bar x'+y\bar y')\right].
\end{split}
 \tag{DO22.10}
\]

Indeed the terms linear in \(i\) have relative trace zero and the remaining terms have relative trace twice themselves. AM-R-004 proves that the right side is the integral unimodular Riemann form of the product Jacobian. Thus (DO22.10) proves perfection of the actual integral trace pairing on \(\mathcal O_{E'}\), not merely a discriminant equality. Equivalently \(\xi'\mathcal O_{E'}=\mathfrak D_{E'/\mathbf Q}^{-1}\). The numerical check is

\[
 |N_{E'/\mathbf Q}(\xi')|=7^{-10}2^{-12}
                   =|d_{E'}|^{-1}.
 \tag{DO22.11}
\]

Here \(|N_E(\xi)|=7^{-5}\), the norm of \(\xi\) from \(E'\) is its square, and \([E':\mathbf Q]=12\). The cyclotomic discriminant is \(|d_{E'}|=7^{10}2^{12}\), consistent with (DO22.10).

In the actual differential coordinates of AM-R-004 write \(u=(u_1,u_2,u_3)\), \(v=(v_1,v_2,v_3)\), and

\[
                w_{i,+}=u_i+iv_i,\qquad w_{i,-}=u_i-iv_i.
 \tag{DO22.12}
\]

These are the two selected \(E'\)-eigencharacters over each old selected \(E\)-embedding. Since \(|u_i|^2+|v_i|^2=(|w_{i,+}|^2+|w_{i,-}|^2)/2\), the actual positive Hermitian form is

\[
 H_A=\sum_{i=1}^3\frac{|u_i|^2+|v_i|^2}{G_i}
          =\sum_{i=1}^3\frac{|w_{i,+}|^2+|w_{i,-}|^2}{2G_i}.
 \tag{DO22.13}
\]

The eigenbasis change on each pair of differentials is \(\left(\begin{smallmatrix}1&i\\1&-i\end{smallmatrix}\right)\), whose determinant is \(-2i\). Across the three pairs its determinant has normalized dyadic valuation 3. The eigenbasis is therefore not a primitive dyadic basis of the product Hodge lattice. In particular

\[
 \partial_{w_{i,+}}=\tfrac12(\partial_{u_i}-i\partial_{v_i}),
 \qquad
 \partial_{w_{i,-}}=\tfrac12(\partial_{u_i}+i\partial_{v_i}).
 \tag{DO22.14}
\]

A total degree \(D\) change of derivative coordinates has the explicit upper denominator bound \(2^D\); it need not attain that bound for a particular polynomial. At 11 this change is integral and invertible. No dyadic index or derivative denominator is removed by the field repair.

The trace-one semicharacter on the lattice in (DO22.9) is the product of the two old trace-one semicharacters, by (DO22.10) and because \(w_0\) acts diagonally. Thus its rigidified line is exactly \(L_{w_0}\boxtimes L_{w_0}\). Its reduced seed is \(\Sigma_{\rm tr}(u)\Sigma_{\rm tr}(v)\), or equivalently \(c^2\sigma_q(u)\sigma_q(v)\), with the exact \(c\) of AM-R-006.

\subsection{3. The local unramified comparison, including the splitting}\label{the-local-unramified-comparison-including-the-splitting}

We prove the part of (DO22.8) not supplied by AM-DO-020. Let \(k\) be a finite completion of \(F'\) of odd residue characteristic, and let \(L/k\) be its unramified quadratic extension in \(E'\). Let \(\mathcal O,\mathcal O_L\) denote their integers.

\subsubsection{3.1. Integral hyperbolic polarizations exist}\label{integral-hyperbolic-polarizations-exist}

The norm \(\mathcal O_L^\times\to\mathcal O^\times\) is surjective. One elementary verification starts with the surjective norm \(\mathbf F_{q^2}^\times\to\mathbf F_q^\times\). After a lift, correcting by \(1+\pi^r z\) changes its norm modulo \(\pi^{r+1}\) by \(1+\pi^r\operatorname{Tr}(z)\); the residue trace is surjective. Successive correction and completeness give the claim. Choose a unit \(z\) of norm \(-1\). Then \((z,1)\) is a primitive isotropic vector of the identity Hermitian lattice of rank 2. Unimodularity supplies a vector pairing to 1 with it. Subtracting one half of that vector's norm times the isotropic vector makes the second vector isotropic. Their pairing is a unit and their Gram determinant is a unit, so they form an integral basis. Hence

\[
       (\mathcal O_L^{2m},I_{2m})\simeq H^{\perp m}
                 \quad\text{integrally.}
 \tag{DO22.15}
\]

Scaling a Hermitian form by a fixed nonzero skew scalar gives the same integral isotropic direct-sum decomposition. Thus both even-rank factors in our dual pair possess polarizations with integral direct summands.

The trace-lattice conductor hypotheses used in Pan's comparison are satisfied. Let \(\lambda_k\) be the conductor exponent of \(\psi_k=\psi_{\mathbf Q_p}\circ\operatorname{Tr}_{k/\mathbf Q_p}\). Because \(L/k\) is unramified and \(\xi'\) generates the absolute codifferent, \(v_L(\xi')=\lambda_k\). Pan's half-trace convention is realized using the opposite skew form with coefficient \(-2\xi'\), and duality exponents \(0,\lambda_k\). Since 2 is a unit, their sum is exactly the unramified compatibility exponent in Pan (2001), §3.1, (3.1.b). His lattice (3.2.c) is the actual selfdual Hom lattice. Replacing both the symplectic form and the additive character by their negatives, if needed to pass between his Heisenberg convention and (DO18.5), transports the same representation. The scalar formulas used below are independent of this simultaneous sign transport.

At these odd places the trace-one semicharacter is trivial on every integral selfdual lattice under discussion: its defining half-pairing is integral in the additive-character conductor convention and 2 is a unit. Consequently the usual selfdual-lattice vacuum is the actual trace-one vacuum.

\subsubsection{3.2. Pure pullback agrees with the prescribed trivial-character splitting}\label{pure-pullback-agrees-with-the-prescribed-trivial-character-splitting}

First set \(\chi_S=\chi_T=1\). Polarize the target hyperbolic space as \(T=X_T\oplus Y_T\), with its integral lattice decomposed in the same way. The Hom symplectic space is polarized by

\[
 X=\operatorname{Hom}_L(S,X_T),\qquad
 Y=\operatorname{Hom}_L(S,Y_T).
 \tag{DO22.16}
\]

The whole source group preserves this polarization. In the Schrödinger model on \(X\), it has the natural unitary action

\[
                 T_g\phi(x)=\phi(xg).
 \tag{DO22.17}
\]

There is no absolute-determinant factor in (DO22.17): \(\det_k(X\mapsto Xg^{-1})
 =N_{L/k}(\det g)^{-b/2}=1\). The associated lift to the metaplectic circle extension is the natural Siegel-Levi splitting. The same assertion for a hyperbolic opposite unitary space is stated in Hanzer (2011), Corollary 2.4, printed p.128; its mixed-model form is equation (3.3), printed p.131.

We must identify (DO22.17) with the \textbf{fixed standard} splitting rather than choosing it independently. Here is that identification. On the standard compact source group, Pan (2001), Lemma 6.2, printed p.197, gives \(\alpha\beta^Y=1\) for each determinant-section element against a hyperbolic plane when its splitting character is trivial. Tensor multiplicativity, §3.6, equation (3.6.e), gives the same result against the whole even hyperbolic opposite space. Proposition 3.3 says that \(\alpha\beta^Y\) is trivial on the compact special unitary subgroup when the opposite dimension is even. The determinant map from the standard compact unitary group onto \(L^1\) is surjective: \(t\mapsto\operatorname{diag}(t,1,\ldots,1)\) is a section and every norm-one \(t\) is a unit. These statements show that the fixed standard splitting fixes the compact vacuum exactly, just as (DO22.17) does.

The ratio of two splittings of the same metaplectic extension is a continuous character of the source unitary group. This follows directly by multiplying their group-law identities; both cover the same symplectic action, so their ratio is central. The moving space is isotropic by (DO22.15), and for these local unitary groups

\[
                        [U(S),U(S)]=SU(S).
 \tag{DO22.18}
\]

Thus the ratio factors through the determinant. The abelianization assertion and the resulting uniqueness from the compact restriction are exactly the unitary case of the proof of Pan's Theorem 3.4, printed p.181. Its ramified exceptional determinant-index-two case is excluded here because \(L/k\) is unramified. Since the ratio is 1 on the compact subgroup and the latter surjects onto \(L^1\), it is 1 everywhere. This proves that (DO22.17) is the prescribed standard trivial-character action on the entire source group. The same proof, using a source hyperbolic polarization, identifies the target action.

This argument uses the published abelianization and splitting results with their hypotheses stated. It does not derive equality from scalar uniqueness alone: compact scalar calculation and determinant surjectivity are the two additional steps that fix the scalar.

\subsubsection{3.3. The one-factor average is exactly the Weil move}\label{the-one-factor-average-is-exactly-the-weil-move}

In (DO22.16) the original Hom lattice has the form \(P_0\oplus P_0^*\), where star denotes the \(\psi_k\)-annihilator. After an arbitrary source move \(g\), it has the form \(P_g\oplus P_g^*\), with \(P_g=\{xg^{-1}:x\in P_0\}\). This follows either from the perfect tensor pairing, or directly because the move is symplectic and preserves each of the two polarized subspaces. The positive normalized vacuum is

\[
                         v_P=\operatorname{vol}(P)^{-1/2}1_P.
 \tag{DO22.19}
\]

The proof of (DO20.6) is a calculation for \emph{arbitrary dual position lattices}. It does not require that the relative CM algebra be split once a polarization of this form is supplied. For clarity, with \(M=P\cap P_0\), \(u=[P:M]\), \(v=[P_0:M]\), the full symplectic index is \(q=uv\). The position and momentum finite sums send \(1_{P_0}\) to \(\sqrt{v/u}\,1_P\): the momentum sum is \(v\) on \(P\) and zero elsewhere in \(P+P_0\), and exactly one position class contributes there. Since \(\operatorname{vol}(P)/\operatorname{vol}(P_0)=u/v\), the normalized average sends \(v_{P_0}\) exactly to \(v_P\). At odd residue characteristic all half-pairings on either integral lattice are killed by \(\psi_k\), so this is precisely the prescribed trace-one average.

Equation (DO22.17) has the same effect, with the same positive value on the zero position. Therefore, in the original realized model as well,

\[
 I_{\Lambda_{N,0},\Lambda_{00}}\phi_{00}
                       =\omega_1(g,1)\phi_{00}.
 \tag{DO22.20}
\]

An arbitrary scalar used to identify the Schrödinger model cancels from this vector equality. Reversing source and target and using the other hyperbolic polarization proves

\[
 I_{\Lambda_{0,N'},\Lambda_{00}}\phi_{00}
                       =\omega_1(1,h)\phi_{00}.
 \tag{DO22.21}
\]

\subsubsection{3.4. Covariance fixes the two-factor phase and the order}\label{covariance-fixes-the-two-factor-phase-and-the-order}

Let \(A_g\) denote the symplectic Hom action of \(g\) and let \(A_h\) be that of \(h\). They commute, preserve the alternating form, and commute with \(w_0\). Hence they preserve the defining trace-one semicharacters. They send finite quotient representatives bijectively to the corresponding moved quotient representatives and preserve the indices. Heisenberg covariance therefore gives the exact operator identity on the appropriate vacuum line

\[
 I_{\Lambda_{N,N'},\Lambda_{N,0}}
  =\omega_1(g,1)\,
       I_{\Lambda_{0,N'},\Lambda_{00}}\,
            \omega_1(g,1)^{-1}.
 \tag{DO22.22}
\]

There is no scalar in (DO22.22): substitute the actual finite sums term by term, using \(\omega W_x\omega^{-1}=W_{A_gx}\), equality of the \(\alpha\)'s, and the unchanged positive square root of the index. Combining (DO22.20)--(DO22.22) yields

\[
\begin{split}
 \phi_{N,N'}
 &=I_{\Lambda_{N,N'},\Lambda_{N,0}}
                   I_{\Lambda_{N,0},\Lambda_{00}}\phi_{00}\\
 &=\omega_1(g,1)\omega_1(1,h)\phi_{00}
   =\omega_1(g,h)\phi_{00}.
\end{split}
 \tag{DO22.23}
\]

The reversed order gives \(\omega_1(1,h)\omega_1(g,1)\phi_{00}\), the same vector because the unitary splittings commute. This proves the required order equality without assuming that every triple of unrelated selfdual lattices has trivial metaplectic triangle phase.

\subsubsection{3.5. Restore the specified characters and determinant roots}\label{restore-the-specified-characters-and-determinant-roots}

At even ranks, a character \(\chi\) with \(\chi|_{k^\times}=1\) defines a character \(\nu_\chi\) of \(L^1\) by

\[
           \nu_\chi(t)=\chi(z),\qquad z/\bar z=t.
 \tag{DO22.24}
\]

Existence of the root is Hilbert 90; independence follows because two roots have quotient in \(k^\times\). Multiplicativity follows by multiplying roots. In particular Pan's Lemma 6.2 evaluates the compact splitting change at \((t-1)\delta\), whose quotient by its conjugate is \(t\), where \(\bar\delta=-\delta\).

The ratio between the \(\chi\)-splitting and the trivial-character splitting is a character of the moving group. The same compact calculation, followed by (DO22.18) and determinant surjectivity, shows that this character is exactly \(\nu_\chi\circ\det\) on the whole group. On the source member our convention uses the character \(\chi_T^{-1}\), equivalently the character on the contragredient source determinant. Thus its factor is \(\chi_T(x)^{-1}\). On the target it uses \(\chi_S\), giving \(\chi_S(y)\). Applying these factors to (DO22.23) proves the local assertion (DO22.8) with its exact root convention.

\subsection{4. The relative split places, including residue characteristic two}\label{the-relative-split-places-including-residue-characteristic-two}

AM-DO-020 proves (DO22.8) at every relative split place, using the prescribed split normalization, the positive indicator value, and the exact trace-one cancellation. Its proof includes residue characteristic two and does not replace the half-pairing by zero there.

All places of \(F'\) above 2 are relative split. One direct check is that the decomposition group of 2 in \(E/\mathbf Q\) is generated by \(2\pmod7\), namely \(\{1,2,4\}\); it excludes conjugation \(-1\pmod7\). Conjugation thus swaps the two primes of \(E\) over 2. Any prime of \(E'\) above one of them is swapped to a prime above the other. Consequently no prime above 2 is fixed by relative conjugation. For a quadratic extension this is exactly splitting of the prime in \(E'/F'\). This argument only identifies the local splitting type; it does not repeat DP57's proof of unramifiedness.

Every remaining finite place is either relative split, covered by AM-DO-020, or unramified quadratic of odd residue characteristic, covered by §3. The local finite averages factor as products of their primary components. Their indices multiply, their positive square roots multiply, and the character factors multiply. The original finite Schwartz vector factors as well. This proves the global finite identity (DO22.8) and its reversed-order counterpart. ∎

\subsection{5. Actual period, normalization, and arithmetic cost retained}\label{actual-period-normalization-and-arithmetic-cost-retained}

Write \(\theta_0=\theta[0;0](0;\tau)\). In the actual product coordinates of §2, apply the reading map of AM-DO-018 entry by entry. Let \(F_{N,N'}\) be the two-average theta germ with initial seed \((\Sigma_{\rm tr}\boxtimes\Sigma_{\rm tr})^{\boxtimes ab}\). The exact whole-family identity is

\[
\begin{split}
 R\phi_{N,N'}&=\theta_0^{2ab}F_{N,N'},\\
 R(\omega_\chi(g,h)\phi_{00})
   &=\chi_T(x)^{-1}\chi_S(y)\theta_0^{2ab}F_{N,N'}.
\end{split}
 \tag{DO22.25}
\]

The equality is of entire holomorphic functions in the common reduced gauge, so it holds after every fixed polynomial differential operator. It is stronger than a nonzero value at a single adelic point.

All analytic norms below use the probability Haar measure on the corresponding complex torus, as fixed and justified in AM-DO-017. They are not automatically the automorphic quotient measure in a terminal period. With that convention,

\[
 \|F_{N,N'}\|^2
          =\det(2Y)^{-ab}|\theta_0|^{-4ab},
 \qquad c_{\rm unit}=\theta_0^{2ab}\det(2Y)^{ab/2}.
 \tag{DO22.26}
\]

The finite Schwartz vector has norm one for the product rational position measure; both normalized averages are isometries on their vacuum lines. At a normalized nonarchimedean valuation \(v_p(p)=1\), its values satisfy

\[
           v_p(\phi_{N,N'}(q_f))\ge-\tfrac12v_p(q_1q_2).
 \tag{DO22.27}
\]

This follows from the \((q_1q_2)^{-1/2}\) coefficient, because the initial nonzero values and all translation/semicharacter values are roots of unity and the remaining finite sum is integral. It is an elementary bound, not an assertion that each coefficient attains it.

The fixed-embedding Gaussian envelope inherited from AM-R-004 and AM-DO-018 is

\[
 |F_{N,N'}(U)|
    <49^{2ab}\sqrt{q_1q_2}\,
                         e^{\pi H_{\rm Hom}(U,U)/2}.
 \tag{DO22.28}
\]

For the alternative seed \((\sigma_q\boxtimes\sigma_q)^{\boxtimes ab}\), replace the period \(\theta_0^{2ab}\) by \((\theta_0c)^{2ab}\), divide the squared norm by \(|c|^{4ab}\), and replace \(49^{2ab}\) by \((6125/16)^{2ab}\), using the proved AM-R-006 bound. The dyadic eigenbasis cost (DO22.14) remains additional if derivatives are expressed in \(E'\)-eigenforms.

To match the recovered source's opposite real additive-character sign, transport the form and character together: use \(\psi_{F'}=\psi_{\mathbf Q}^{-1}\circ\operatorname{Tr}_{F'/\mathbf Q}\) and the corresponding opposite trace coefficient. Their composition on the Heisenberg center is unchanged. This changes the old fixed quadratic coefficient pin; it is part of the new mixed constructor, not evidence that the old \nolinkurl{StdPair.xi} or elliptic \nolinkurl{VacIn} is satisfied.

\subsection{6. Precisely which gate this closes}\label{precisely-which-gate-this-closes}

The result supplies one actual common-period moved family, both prescribed finite averages with their exact covolume indices, its standard unitary Weil realization for every finite adelic move, and the equality of the two orders. It resolves the finite phase in (DO19.1) on this explicitly repaired host at all even ranks.

The recovered theorem \nolinkurl{thetaKernel_eq_jet_thC} in \nolinkurl{Development/AbelianVarietyTheory/Proofs/DecoratedLatticeProofs/ThetaKernelSplitPlacesEqJetClassicalTheta.lean} transports a vacuum identity to polynomial jets; its key hypothesis is an actual equality, rather than mere nonvanishing. Equation (DO22.25) provides the mathematical analogue of that hypothesis for this new mixed seed. No Lean build or formal instantiation is claimed.

AM-K018 subsequently proves the full-cyclotomic Chern/source formal gauge comparison for this underlying ordinary product of copies of \(J\), including the actual primitive kernel at 11. This comparison is used at its separately recorded proof-check status. It removes the previously listed gauge gate for these ordinary products; it does not assert that a reduced frame is integral or cover products with conjugated \(E\)-actions.

The following requirements are still separate: the old elliptic \nolinkurl{VacIn}/fourth-tensor seed pin; the actual integral row module with its specified coefficient type; the whole set of \nolinkurl{AdmS} covolume equations; membership of an arbitrarily selected row in the prescribed iterated theta image; and the terminal-mass equation in THETA3 Part II G1. Finite Weil equivariance is a concrete constructor result and cannot substitute for any of those arithmetic equations. AM-DO-026 records the more precise operator and class-kernel consequences.

\subsection{References and exact source use}\label{references-and-exact-source-use}

Pan, S.-Y. (2001). Splittings of the metaplectic covers of some reductive dual pairs. \emph{Pacific Journal of Mathematics, 199}(1), 163--226. \url{https://msp.org/pjm/2001/199-1/pjm-v199-n1-p10-p.pdf} Used locators: §3.1, (3.1.a)--(3.1.b), printed pp.177--178; §3.2, (3.2.c); Proposition 3.3, pp.179--180; Theorem 3.4 and especially its unitary abelianization/uniqueness argument, p.181; §3.6, (3.6.e), p.183; Lemma 6.2, p.197. The published abelianization and splitting theorems are used, not claimed as newly reproved here.

Hanzer, M. (2011). Rank one reducibility for unitary groups. \emph{Glasnik Matematički, 46}(66), 121--148. \url{https://web.math.pmf.unizg.hr/glasnik/46.1/46(1)-12.pdf} Used locators: Theorem 2.3 and Corollary 2.4, printed pp.127--128; equation (3.3), printed p.131. These give the natural polarized pullback action. Its identification with the fixed standard splitting is separately proved in §3.2 above.

Carried supplied result: \nolinkurl{DP_MULT.tex}, Proposition DP57, accepted unramified CM repair and endpoint realization, especially parts (a), (b), and (d), and the proof of the integral \(\mathcal O_{E'}\)-action. Local exact file: \nolinkurl{sources/RETURN_DP_MULT_checkpoint_06_20261121T0420Z/DP_MULT.tex}. Other carried results: AM-R-004, AM-R-006, AM-DO-017--020, ALP-LIFT-07--08, at their separately recorded proof-check status.

\clearpage

\section{AM-DO-024: the direct-average sign, an exact vacuum pairing, and adjoints}\label{am-do-024-the-direct-average-sign-an-exact-vacuum-pairing-and-adjoints}

Author: GPT-6 Astra Pro, ALP\_MACHINE \nolinkurl{/root/definite_oscillator}; inherited reasoning effort not separately exposed; 21 November 2026 UTC.

This result uses the repaired host and the \textbf{standard unitary splitting} proved in AM-DO-022. It computes a theta-line pairing with its measure normalization. It does not identify that pairing with the terminal automorphic mass in THETA3 G1.

\subsection{1. Statement with all normalizations}\label{statement-with-all-normalizations}

Put \(E'=\mathbf Q(\zeta_{28})\), \(F'=E'^{+}\), and use the principally polarized CM abelian variety \(A=J^2\), with integral trace coefficient \(\xi'=\xi/2\), constructed in DP57 and AM-DO-022. Fix positive even integers \(a,b\). Let \(S=E'^a\) and \(T=E'^b\) have positive identity Hermitian forms. The underlying Hom symplectic space, its character, the trace-one semicharacters, and its finite Schwartz vacuum \(\phi_{00}\) are exactly those of AM-DO-022; normalize \(\|\phi_{00}\|=1\).

For \(g\in U(S)(\mathbf A_{F',f})\) and \(h\in U(T)(\mathbf A_{F',f})\), write \(\Lambda_{00}\), \(\Lambda_{10}\), \(\Lambda_{11}\) for the standard, source-moved, and both-moved rational period lattices. Their actual action is \(X\mapsto hXg^{-1}\). All lattices carry the same trace-one rule \(\alpha(x)\), and its restrictions agree on intersections. Write

\[
 q_0=[\Lambda_{11}:\Lambda_{11}\cap\Lambda_{00}],\qquad
 I_{ji}=q_{ji}^{-1/2}
       \sum_{x\in\Lambda_j/(\Lambda_i\cap\Lambda_j)}
                         \alpha(x)^{-1}W_x,
 \tag{DO24.1}
\]

where \(q_{ji}=[\Lambda_j:\Lambda_i\cap\Lambda_j]\) and the square root is positive. On holomorphic theta functions replace \(W_x\) by the reduced translation \(U_x\) of AM-DO-017. For sign-valued \(\alpha\), the inverse in (DO24.1) equals \(\alpha\), as in the source definition.

Let \(\omega_1\) be the standard splitting with both unitary characters trivial. Then

\[
 \boxed{
 \omega_1(g,h)\phi_{00}
   =I_{11,00}\phi_{00}
   =I_{11,10}I_{10,00}\phi_{00}
   =I_{11,01}I_{01,00}\phi_{00}.}
 \tag{DO24.2}
\]

In particular the direct-average comparison has sign \(+1\). Its finite Schwartz pairing is

\[
 \boxed{\langle\phi_{00},\phi_{11}\rangle=q_0^{-1/2},
        \qquad \|\phi_{11}\|=1.}
 \tag{DO24.3}
\]

The inner product is conjugate-linear in the first argument. For general splitting characters \(\chi_S,\chi_T\), trivial on \(\mathbf A_{F'}^\times\), and determinant roots \(x/\bar x=\det g\), \(y/\bar y=\det h\), the formula is

\[
 \omega_\chi(g,h)\phi_{00}
       =\chi_T(x)^{-1}\chi_S(y)I_{11,00}\phi_{00}.
 \tag{DO24.4}
\]

The dependence is on these roots, not on evaluating the characters at \(\det g,\det h\) themselves. This is the root convention already fixed in AM-DO-009 and AM-DO-022.

Let \(F_{00}\) have initial seed \((\Sigma_{\mathrm{tr}}\boxtimes\Sigma_{\mathrm{tr}})^{\boxtimes ab}\) and put \(F_{11}=I_{11,00}F_{00}\). On the common complex torus quotient by \(\Lambda_{00}\cap\Lambda_{11}\), integrate its Hermitian norm density against \textbf{probability Haar measure}. With \(Y=\operatorname{Im}\tau\) and \(\theta_0=\theta[0;0](0;\tau)\),

\[
 \boxed{\langle F_{00},F_{11}\rangle
      =q_0^{-1/2}\det(2Y)^{-ab}|\theta_0|^{-4ab}.}
 \tag{DO24.5}
\]

Both norms squared are \(\det(2Y)^{-ab}|\theta_0|^{-4ab}\). For the alternative seed \((\sigma_q\boxtimes\sigma_q)^{\boxtimes ab}\), divide the right-hand side and both squared norms by \(|c|^{4ab}\), where the exact scalar \(\Sigma_{\mathrm{tr}}=c\sigma_q\) was computed in AM-R-006. The period connecting the finite Schwartz vacuum to that alternative seed is \((\theta_0c)^{2ab}\).

\subsection{2. Positivity of the actual local matrix coefficient}\label{positivity-of-the-actual-local-matrix-coefficient}

We first prove

\[
                 \langle\phi_{00},\omega_1(g,h)\phi_{00}\rangle>0
 \tag{DO24.6}
\]

at every finite place. This is a statement about the fixed standard splitting, rather than an arbitrary choice of a unit vector on each vacuum line.

\subsubsection{2.1. Relative split places}\label{relative-split-places}

At a relative split place, AM-DO-020 identifies both group actions with the normalized change of variables in a single Schrödinger position space. It includes residue characteristic two and retains the trace-one half-pairing there. In the model used in that proof, the initial and moved normalized vectors are, up to the same irrelevant unit scalar,

\[
 v_{P_0}=\operatorname{vol}(P_0)^{-1/2}\mathbf1_{P_0},\qquad
 v_P=\operatorname{vol}(P)^{-1/2}\mathbf1_P,
 \tag{DO24.7}
\]

for full position lattices. Their pairing is

\[
 \langle v_{P_0},v_P\rangle
   =\frac{\operatorname{vol}(P_0\cap P)}
          {\sqrt{\operatorname{vol}(P_0)\operatorname{vol}(P)}}>0.
 \tag{DO24.8}
\]

The quotient is independent of rescaling the position Haar measure.

\subsubsection{2.2. Odd unramified quadratic places}\label{odd-unramified-quadratic-places}

Let \(L/k\) be the local unramified quadratic extension. The positive identity lattices of even ranks \(a,b\) are integrally hyperbolic, by the explicit norm-unit and half-adjustment argument in AM-DO-022 §3.1. Choose those integral hyperbolic bases. Their lattice stabilizers \(K_S,K_T\) are hyperspecial maximal compact subgroups. The \textbf{standard trivial-character splitting fixes the initial vacuum} under \(K_S\times K_T\), by AM-DO-022 §3.2, including its Pan comparison and abelianization argument.

The Cartan decomposition for the split unramified unitary group gives

\[
 g=k_1a_\nu k_2,\qquad h=l_1b_\mu l_2,
 \tag{DO24.9}
\]

where \(k_i\in K_S\), \(l_i\in K_T\), and the middle terms are diagonal in the selected hyperbolic bases, with entries \(\varpi^{\nu_j},\varpi^{-\nu_j}\) and \(\varpi^{\mu_j},\varpi^{-\mu_j}\). The local hypotheses and these exact standard lattices/maximal split tori are recalled in Li, Rapoport, and Zhang (2024), §§1 and 3.4; §4.1 explicitly invokes the Cartan decomposition for \(K\backslash U(W_0)/K\). We use this published decomposition, without claiming a new proof of it.

Since the source and target splittings commute, (DO24.9) and compact invariance imply

\[
 \langle\phi_{00},\omega_1(g,h)\phi_{00}\rangle
       =\langle\phi_{00},\omega_1(a_\nu,b_\mu)\phi_{00}\rangle.
 \tag{DO24.10}
\]

Polarize the target as \(T=X_T\oplus Y_T\) in the same hyperbolic basis. The associated position space is \(\operatorname{Hom}_L(S,X_T)\). The whole source action preserves this polarization and is the pure normalized pullback identified with the standard splitting in AM-DO-022 §3.2. The target element \(b_\mu\) is in the Siegel Levi for this polarization. Its standard Weil operator is also a positive normalized pullback: in the standard splitting formula its Bruhat index is zero and its character factor is \(\chi_S=1\). The usual factor \(|\det A|^{1/2}\), or its reciprocal when writing the inverse change of variables, is positive. This is the polarized formula in Hanzer (2011), Theorem 2.3 and Corollary 2.4, together with the mixed model equation (3.3); the identification with the fixed splitting is the separate argument in AM-DO-022, not a new normalization choice.

The initial lattice splits as the position lattice and its exact dual because the trace coefficient is a codifferent generator. The initial vacuum is therefore a normalized indicator as in (DO24.7). The combined diagonal action sends it to another such normalized indicator. Its coefficient is (DO24.8), which is strictly positive. This proves (DO24.6) at an odd unramified quadratic place.

All relative dyadic places in this host are split, as proved in AM-DO-022 §4. The preceding two cases thus cover every finite place. The global matrix coefficient is a finite product of these positive local coefficients: outside finitely many places the move is compact and the coefficient is one. Consequently (DO24.6) holds globally.

\subsection{3. The positive coefficient fixes the direct-average phase}\label{the-positive-coefficient-fixes-the-direct-average-phase}

By the covariance of the actual Weil representation and the invariance of the trace-one rule under \(E'\)-linear isometries, \(\omega_1(g,h)\phi_{00}\) belongs to the vacuum line of \(\Lambda_{11}\). The vector \(I_{11,00}\phi_{00}\) belongs to that same one-dimensional line. Both have norm one: the former by unitarity and the latter by AM-DO-017. Thus they differ by a unique scalar \(\lambda\) of absolute value one.

The direct-average coefficient is exactly

\[
 \langle\phi_{00},I_{11,00}\phi_{00}\rangle=q_0^{-1/2}.
 \tag{DO24.11}
\]

Indeed choose zero as the representative of the zero coset in (DO24.1). That term contributes \(q_0^{-1/2}\). Every other term is orthogonal to \(\phi_{00}\): if \(x\notin\Lambda_{00}\), selfduality supplies \(\lambda_0\in\Lambda_{00}\) with a nontrivial commutator \(e^{2\pi i E_H(\lambda_0,x)}\). The two vectors are then eigenvectors of the same unitary Heisenberg operator with distinct eigenvalues. This is exactly the finite-pairing argument in AM-DO-017; the common semicharacter on the intersection ensures that there is no representative-dependent phase.

Both (DO24.6) and (DO24.11) are positive real numbers. Pairing \(\omega_1(g,h)\phi_{00}=\lambda I_{11,00}\phi_{00}\) with \(\phi_{00}\) therefore makes \(\lambda\) positive real. Since \(|\lambda|=1\), it follows that \(\lambda=1\). AM-DO-022 already proved that the actual Weil translate equals each two-step average. This proves (DO24.2) and (DO24.3). Its exact splitting-change formula proves (DO24.4).

The same finite orthogonality calculation on the common holomorphic theta space gives \(\langle F_{00},I_{11,00}F_{00}\rangle
 =q_0^{-1/2}\|F_{00}\|^2\). AM-DO-017 proves that the probability-Haar norms on the two principal quotients agree with their norms on the common finite cover. Inserting the explicit seed norm from (DO22.26) proves (DO24.5). The \(\sigma_q\) statement follows by multiplying both vectors by \(c^{-2ab}\), whose contribution to a Hermitian pairing is \(|c|^{-4ab}\). This completes the proof of all the stated pairing and sign assertions. \(\square\)

This argument applies to these actual dual-pair orbit lattices. It does not assert that the metaplectic triangle phase is one for three arbitrary selfdual lattices. Neither does it cover the original ramified \(\mathbf Q(\zeta_7)\) host with unchanged trivial characters: AM-DO-021 computes a nontrivial compact sign there.

\subsubsection{A concrete rank-two to rank-four value on the repaired host}\label{a-concrete-rank-two-to-rank-four-value-on-the-repaired-host}

Take \(a=2,b=4\), \(\pi=2-\zeta_7\), \(r=\pi/\bar\pi\), \(N=r\mathcal O_{E'}\oplus\mathcal O_{E'}\), and \(N'=\mathcal O_{E'}^4\). The two conjugate ideals generated by \(\pi\) and \(\bar\pi\) remain coprime after the integral extension from \(E\) to \(E'\). Their absolute norms are \(127^2\), because \(N_{E/\mathbf Q}(\pi)=127\) and \([E':E]=2\). For each of the four changed Hom entries,

\[
 \mathcal O_{E'}\cap r^{-1}\mathcal O_{E'}
       =\bar\pi\mathcal O_{E'},\qquad
 [\mathcal O_{E'}:\bar\pi\mathcal O_{E'}]
       =[r^{-1}\mathcal O_{E'}:\bar\pi\mathcal O_{E'}]=127^2.
 \tag{DO24.16}
\]

The remaining four entries are unchanged. Thus \(q_0=127^8\), and (DO24.3)--(DO24.5) become

\[
 \langle\phi_{00},\phi_{N,N'}\rangle=127^{-4},\qquad
 \langle F_{00},F_{N,N'}\rangle
        =127^{-4}\det(2Y)^{-8}|\theta_0|^{-32}.
 \tag{DO24.17}
\]

The common analytic period is \(\theta_0^{16}\), and the direct finite average has coefficient \(127^{-4}\) before its root-of-unity summands. This example retains the two Jacobian factors of the repaired host; omitting that doubling would give the different \(127^{-2}\) coefficient of AM-DO-020 over the original host.

\subsection{4. Exact adjoints of the finite operators}\label{exact-adjoints-of-the-finite-operators}

There is an additional operator identity that does not require the special host. Let \(\Lambda_1,\Lambda_2\), their common semicharacters, and their normalized theta lines satisfy AM-DO-017. Then

\[
                       I_{21}^{*}=I_{12},\qquad
                       I_{12}I_{21}=1,\qquad I_{21}I_{12}=1.
 \tag{DO24.12}
\]

The adjoints use the actual normalized Hermitian forms on the two lines, and the same identity holds in their finite Schwartz models.

\textbf{Proof.} Put \(\Gamma=\Lambda_1\cap\Lambda_2\), \(q=[\Lambda_i:\Gamma]\), and realize both principal theta lines inside the common space of sections for \(\Gamma\), with its probability-Haar norm. The operators \(\alpha_i(x)^{-1}U_x\), for \(x\in\Lambda_i/\Gamma\), form a unitary representation of that finite group on this common space. Its invariant subspace is exactly the principal theta line \(H_i\). Therefore

\[
 P_i=\frac1q\sum_{x\in\Lambda_i/\Gamma}\alpha_i(x)^{-1}U_x
 \tag{DO24.13}
\]

is its orthogonal projection. In particular \(P_i=P_i^*\). Equation (DO24.1) says \(I_{21}=\sqrt q\,P_2|_{H_1}\). For \(f\in H_1\), \(g\in H_2\),

\[
 \langle I_{21}f,g\rangle
 =\sqrt q\langle P_2f,g\rangle
 =\sqrt q\langle f,g\rangle
 =\sqrt q\langle f,P_1g\rangle
 =\langle f,I_{12}g\rangle.
 \tag{DO24.14}
\]

This proves the adjoint identity. Both operators are nonzero isometries between one-dimensional spaces by AM-DO-017, so they are unitary inverses. The finite Schwartz proof is identical, using the finite Heisenberg eigenrelations in place of the holomorphic periodicity. \(\square\)

For the two-step moved family this gives the \textbf{actual reversed operator}, with no unspecified constant:

\[
 (I_{11,10}I_{10,00})^*=I_{00,10}I_{10,11}.
 \tag{DO24.15}
\]

Together with (DO24.2), this is the finite-vacuum adjoint of the standard Weil action. A general determinant-character factor in (DO24.4) is conjugated, exactly as required by a unitary adjoint. No automorphic quotient volume has been replaced by a theta-line norm in this assertion.

\subsection{5. Literal source obligations supplied}\label{literal-source-obligations-supplied}

The recovered \nolinkurl{Formalism/LocalToGlobal/KernelLattices/LateralTheory.lean} defines \nolinkurl{SignP} at lines 230--233 and \nolinkurl{OrderP} at lines 235--238. Equation (DO24.2) proves their mathematical analogues for the new mixed-CM seed, with the stronger positive sign. The period and every polynomial jet remain those in (DO22.25), so the comparison uses one common seed scalar throughout.

The exact operator and pairing formulas are concrete inputs towards \nolinkurl{VacT}/\nolinkurl{ClsK.Core}. They do not instantiate the old elliptic \nolinkurl{VacIn}, do not choose the prescribed arithmetic row lattice, and do not identify the source \nolinkurl{Btop} terminal functional with (DO24.5). In particular a computed positive vacuum coefficient does not prove the terminal-mass equation or all six \nolinkurl{AdmS} predicates.

\subsection{References}\label{references-2}

Hanzer, M. (2011). Rank one reducibility for unitary groups. \emph{Glasnik Matematički, 46}(66), 121--148. \url{https://web.math.pmf.unizg.hr/glasnik/46.1/46(1)-12.pdf} Exact locators: Theorem 2.3 and Corollary 2.4, pp.127--128; equation (3.3), p.131. Used through the standard-splitting comparison proved in AM-DO-022.

Li, C., Rapoport, M., \& Zhang, W. (2024). Arithmetic fundamental lemma for the spherical Hecke algebra. \emph{Manuscripta Mathematica, 175}, 1--51. \url{https://doi.org/10.1007/s00229-024-01572-0} Exact locators: §1 (odd residue characteristic and unramified quadratic extension); §3.4 (split Hermitian form, standard selfdual lattice, hyperspecial subgroup, and diagonal maximal split torus); §4.1 (the Cartan decomposition on the corresponding double-coset space). No arithmetic fundamental lemma from this paper is used.

Carried and separately checked: AM-DO-017, AM-DO-020, AM-DO-022, AM-R-006, and the accepted DP57 host construction.

\clearpage

\section{AM-R-008: compact-determinant representatives and the exact cost of an automorphic determinant twist}\label{am-r-008-compact-determinant-representatives-and-the-exact-cost-of-an-automorphic-determinant-twist}

Author: \nolinkurl{/root}, GPT-6 Astra Pro; effort not externally exposed; ALP\_MACHINE, 21 November 2026 UTC.

This result refines AM-R-007 for explicit unitary levels whose determinant images equal the source torus level in ALP-LIFT-10. It does not identify those levels with the target stabilizer of the actual first lift. It determines the effect of a specified automorphic determinant character on the integral class-value module and its mass norm. It does not claim that changing a splitting character preserves the source's pinned coefficient data or the next theta kernel.

\subsection{1. Exact class representatives}\label{exact-class-representatives}

Let \(E=\mathbf Q(\zeta_7)\), (F=E\^{}+), (G=U(I\_n)), \(n\ge2\), and (T=U(1)(E/F)). Write \(K_0\subset G(\mathbb A_{F,f})\) for the full stabilizer of the standard Hermitian lattice and \(K_{T,0}=\det K_0\). For (M=29\^{}e), \(e\ge1\), define \(K\subset K_0\) by:

\begin{itemize}
\tightlist
\item
  at each of the three places over 29, \(g\equiv I_n\pmod{29^e}\);
\item
  at the ramified place over 7, the residue of \(\det g\) is (+1);
\item
  at every other finite place, \(g\in K_{0,v}\).
\end{itemize}

This is a normal subgroup of \(K_0\). Its determinant image is exactly the torus subgroup \(K_M\) of ALP-LIFT-10, equation (O21). Inclusion in that subgroup follows from the definitions. For the reverse inclusion, the section \(t\mapsto\operatorname{diag}(t,1,\ldots,1)\) belongs to the indicated group level whenever \(t\in K_M\).

Choose representatives \(a\in K_{T,0}\) for the finite quotient

\[
\mathcal A=K_{T,0}/(\mu_{14}K_M).
\tag{R8.1}
\]

Each representative can be chosen to be 1 outside 7 and 29, since the local factors of \(K_M\) equal (K\_\{T,0,v\}) elsewhere. Put \(s(a)=\operatorname{diag}(a,1,\ldots,1)\in K_0\). Then

\[
\boxed{
G(\mathbb A_{F,f})
 =\bigcup_{a\in\mathcal A}G(F)\,SU(I_n)(F_{v_{11}})\,s(a)K .}
\tag{R8.2}
\]

Thus every finite lattice class has a representative whose determinant is a compact norm-one idele supported at 7 and 29. Noncompact motion may be confined to the determinant-one group at 11. The set in (R8.2) need not be a disjoint union, and the theorem does not give a bound for the determinant-one matrix at 11.

\subsubsection{Proof}\label{proof-3}

ALP-LIFT-09 proves \(T(\mathbb A_{F,f})=T(F)K_{T,0}\), with \(T(F)\cap K_{T,0}=\mu_{14}\). For \(g\in G(\mathbb A_{F,f})\), write

\[
\det g=t_0\,a\,k_T,
\qquad t_0\in T(F),\quad a\in\mathcal A,\quad k_T\in K_M.
\tag{R8.3}
\]

The root of unity needed to choose \(a\) is absorbed in \(t_0\). Lift \(k_T\) to \(k\in K\) by the determinant section, and set

\[
h=s(t_0)^{-1}gk^{-1}s(a)^{-1}\in SU(I_n)(\mathbb A_{F,f}).
\tag{R8.4}
\]

The strong approximation argument of AM-R-007, with the omitted split place (v\_\{11\}), gives

\[
h=\delta u_{11}k_1,
\qquad \delta\in SU(I_n)(F),\quad
u_{11}\in SU(I_n)(F_{v_{11}}),\quad k_1\in K\cap SU(I_n)(\mathbb A_{F,f}).
\tag{R8.5}
\]

In the more general proof one uses the conjugate of \(K\) by (s\(a\)). Here it is the same subgroup because \(s(a)\in K_0\) and \(K\) is normal in \(K_0\). Substituting (R8.5) in (R8.4), with the factors in their actual order, gives

\[
g=s(t_0)\delta\,u_{11}s(a)\,
       \bigl(s(a)^{-1}k_1s(a)\bigr)k.
\tag{R8.6}
\]

The last factors lie in \(K\), proving (R8.2). The determinant of (u\_\{11\}s\(a\)) is \(a\), so it is compact. This proves every assertion about the representatives.

\subsubsection{Extension to any compatible target compact level}\label{extension-to-any-compatible-target-compact-level}

Let \(K\) now be any open compact subgroup of \(K_0\), and set \(D=\det K\). The group \(D\) is open in \(K_{T,0}\): the preimage of \(K\) under the continuous local determinant section is an open neighborhood of 1 contained in \(D\). Choose representatives \(a\) of \(K_{T,0}/(\mu_{14}D)\). The same conclusion (R8.2) holds with these representatives and this \(K\), without a normality assumption.

To prove it, write \(\det g=t_0ak_D\), and choose \(k\in K\) with \(\det k=k_D\). This lift exists by the definition of \(D\); it need not be the diagonal section. Set \(h\) as in (R8.4). Apply strong approximation to the actual conjugated compact \[
 L=s(a)Ks(a)^{-1}\cap SU(I_n)(\mathbb A_{F,f}).
 \tag{R8.16}
\] It gives \(h=\delta u_{11}\ell\), with \(\ell\in L\). The exact factor order is then \[
 g=s(t_0)\delta\,u_{11}s(a)
                \bigl(s(a)^{-1}\ell s(a)\bigr)k.
 \tag{R8.17}
\] The last two factors belong to \(K\). This proves the assertion. The representatives still have compact determinant; their support need not be confined to 7 and 29.

Every assertion of section 2 consequently holds with \(D\) in place of \(K_M\), provided \(\gamma_f|_D=1\). This supplies the exact twist isometry for any compatible actual target level. Refining a given level to make it compatible is a different operation; its index and norm conventions must still be retained. This extension was supplied and whole-proof checked by /root/lift\_accounting in agents/review\_root\_compact\_twist.md and read in full by /root.

\subsection{2. Exact integral and metric effect of a determinant character}\label{exact-integral-and-metric-effect-of-a-determinant-character}

Let \(\gamma=\gamma_\infty\gamma_f\) be a continuous unitary automorphic character of \(T(F)\backslash T(\mathbb A_F)\), and assume

\[
\gamma_f|_{K_M}=1.
\tag{R8.7}
\]

Fix a finite set (g\_i=u\_i s(a\_i)) of representatives of the class set used to present the automorphic space. Formula (R8.2) allows that choice for every class. No new class-number or mass formula is assumed in making this choice. Put

\[
\Gamma_i=G(F)\cap g_iKg_i^{-1}.
\tag{R8.8}
\]

For a positive definite unitary group these stabilizers are finite: the corresponding rational matrices preserve a fixed integral lattice, and their real embeddings lie in a compact group. A discrete lattice in the resulting finite-dimensional real matrix space meets a compact set in a finite set. This supplies the elementary stabilizer finiteness used below.

To avoid an implicit convention change, present coefficient-valued automorphic forms as functions satisfying

\[
f(\delta gk)=\rho(\delta_\infty)f(g),
\qquad \delta\in G(F),\quad k\in K.
\tag{R8.9}
\]

Define

\[
f^{\gamma}(g)=\gamma_f(\det g)f(g),\qquad
\rho^{\gamma}(z)=\gamma_\infty(\det z)^{-1}\rho(z).
\tag{R8.10}
\]

Then \(f^\gamma\) satisfies (R8.9) with \(\rho^\gamma\). Indeed \(\gamma_f(\det k)=1\) by (R8.7), while rational triviality of the global character gives

\[
\gamma_f(\det\delta)=\gamma_\infty(\det\delta_\infty)^{-1}.
\tag{R8.11}
\]

For \(\delta\in\Gamma_i\), its finite determinant belongs to \(K_M\), because conjugation does not change a determinant. Equations (R8.7)--(R8.11) therefore imply

\[
\gamma_\infty(\det\delta_\infty)=1,
\qquad V^{\Gamma_i}_{\rho}=V^{\Gamma_i}_{\rho^\gamma}.
\tag{R8.12}
\]

These are equal subspaces in the same underlying coefficient vector space, not just spaces with equal dimensions. On class values the map is the diagonal matrix

\[
\bigl(f(g_i)\bigr)_i\longmapsto
       \bigl(\gamma_f(a_i)f(g_i)\bigr)_i.
\tag{R8.13}
\]

Every \(\gamma_f(a_i)\) is a root of unity. The image of the compact group (K\_\{T,0\}) under a smooth character is finite: its kernel is open and hence has finite index. In particular the displayed scalar is an algebraic unit at every finite place and has absolute value one under every complex embedding.

Let \(\mathcal V\subset V\) be the specified coefficient lattice over a number ring containing these finitely many roots of unity. Use the same underlying lattice for the twisted coefficient space. Then (R8.12)--(R8.13) give an exact isomorphism

\[
\boxed{
\bigoplus_i\mathcal V^{\Gamma_i}_{\rho}
\ \xrightarrow{\ \sim\ }
\bigoplus_i\mathcal V^{\Gamma_i}_{\rho^\gamma}.}
\tag{R8.14}
\]

No denominator is introduced. All local contents of a vector are preserved, including those at 2, 7, 11 and 29. The same statement holds for a specified submodule transported by this actual map; it does not identify that transported submodule with an independently prescribed arithmetic row.

With any fixed Hermitian metric on \(V\), the character is unitary and the underlying metric may be used unchanged. Thus the map is also an exact isometry for the class mass norm

\[
\|f\|_{\rm mass}^2
   =\sum_i w_i\|f(g_i)\|^2,
\qquad w_i>0,
\tag{R8.15}
\]

including \(w_i=|\Gamma_i|^{-1}\) and any fixed common Haar normalization. It preserves the class-value supremum norm as well. There is no added logarithmic norm or denominator cost for this specified twist map in this presentation.

\subsection{3. Scope of the zero-cost statement}\label{scope-of-the-zero-cost-statement}

The subgroup K\_M in ALP-LIFT-10 is a source torus stabilizer. The target U(2) stabilizer, its index and its mass have not thereby been computed. The theorem here applies to the displayed U(n) level once a vector belongs to its class module; that membership is an additional hypothesis for any first-lift application. Independently of a finite class presentation, multiplication by a unitary determinant character preserves the scalar automorphic L2 norm. A refinement to a compatible target level must still retain the cost of that refinement.

The result removes a possible artificial height penalty coming from poorly normalized representatives of a determinant twist. It does not make a change of theta datum free in every other sense. The splitting characters, archimedean coefficient representation and source covariance pins have changed. To use the twist in a proposed lift chain, one must show that the new theta kernel has precisely these input/output characters, that the prescribed row submodules are the transported submodules in (R8.14), and that the subsequent pairing is nonzero. None of those assertions follows from an isometry of the ambient class-value modules.

There is also no new equality between the \(E\)-host first lifts and AM-DO-022's \(E'=\mathbf Q(\zeta_{28})\)-host kernels. The present theorem stays over \(E\). Applying it over \(E'\) would require the corresponding torus class calculation and compact determinant presentation; the old proof is not silently transplanted.

The geometric object is an associated coefficient bundle on the finite class presentation of a definite unitary quotient. Tensoring its representation by a rank-one character is an operation in the representation Grothendieck group, but (R8.13)--(R8.14) are actual integral maps with a metric, stronger than a virtual-class identity. The framework locators are SGA 3 I, 4.7.2--4.7.3 and SGA 6 VI, Theorem 3.3, as carried in the geometric dictionary. The arithmetic input is ALP-LIFT-09 (V2--V6), and the approximation theorem is the same Rapinchuk (2013), Theorem 2.3, p.~281, used in AM-R-007.

\clearpage

\section{AM-DO-023: exact honest-splitting weights and a specified second-edge vanishing}\label{am-do-023-exact-honest-splitting-weights-and-a-specified-second-edge-vanishing}

Author: GPT-6 Astra Pro, ALP\_MACHINE \nolinkurl{/root/definite_oscillator}; inherited reasoning effort not separately exposed; 21 November 2026 UTC.

The first theta lift in ALP-LIFT-09 is nonzero. The assertion below is that its displayed archimedean type does not survive one \textbf{specified} next positive-definite Hom theta edge with trivial splitting characters. It does not assert vanishing for every second edge, for an arbitrarily changed conjugation or polarization, for a changed character, or for a doubled Gram kernel. This calculation is over the original host \(E=\mathbf Q(\zeta_7)\), as is ALP-LIFT-09. AM-DO-022 instead uses \(E'=\mathbf Q(\zeta_{28})\); the two are not combined into one tower without an additional extension-of-field comparison.

\subsection{1. Conventions and the exact local formula}\label{conventions-and-the-exact-local-formula}

At a real CM place put \(S=\mathbf C^a\), \(T=\mathbf C^b\), with their positive identity Hermitian forms. Let \(\kappa>0\). The symplectic space is \(\operatorname{Hom}_{\mathbf C}(S,T)\), with form

\[
 b(X,Y)=-2\operatorname{Re}
          (i\kappa\operatorname{tr}(Y^*X))
       =2\kappa\operatorname{Im}\operatorname{tr}(Y^*X),
 \tag{DO23.1}
\]

and group action \(A(g,h)X=hXg^{-1}\). Use \(\psi(t)=e^{2\pi it}\) and the covariance

\[
 \omega(g,h)\rho(X)\omega(g,h)^{-1}=\rho(hXg^{-1}).
 \tag{DO23.2}
\]

Use the source convention for the honest unitary splitting: a change of the two character labels acts by \(\nu_{\Delta\chi_T}(\det g)^{-1}
  \nu_{\Delta\chi_S}(\det h)\), where \(\nu_\chi(z/\bar z)=\chi(z)\). This is the convention in \nolinkurl{LocalNewformFunction.splitFacO}, line 51, and AM-DO-022. Write its character labels as

\[
 \chi_S(z)=(z/\bar z)^{m/2},\quad m\equiv a\pmod2,
 \qquad
 \chi_T(z)=(z/\bar z)^{m'/2},\quad m'\equiv b\pmod2.
 \tag{DO23.3}
\]

Let \(\mathcal F=\operatorname{Sym}(T\otimes S^*)\) be the positive holomorphic Fock space, including its Gaussian vacuum. Then the \textbf{actual honest-splitting action} is

\[
 \boxed{\mathcal F\big|_{U(S)\times U(T)}
   =\det_S^{(-m'-b)/2}\boxtimes\det_T^{(m+a)/2}
                       \otimes\operatorname{Sym}(T\otimes S^*).}
 \tag{DO23.4}
\]

Here the integer powers are legitimate precisely because of the parity conditions in (DO23.3). The exponent \(m/2\) in the character on \(\mathbf C^\times\) must not be confused with a determinant power on \(U(b)\): the latter includes the additional \(a/2\) vacuum shift. For instance an odd-dimensional opposite space requires an honest splitting of the metaplectic circle extension, while the bare oscillator by itself only acts on the determinant double cover.

\subsubsection{Primary-source identification}\label{primary-source-identification}

Konno and Konno (2007), §3.1, printed pp.43--44, use a Hermitian space \(V\) linear in the second variable and a skew-Hermitian space \(W\) linear in the first variable, with symplectic form

\[
 \operatorname{Re}((v,v')\overline{\langle w,w'\rangle}).
 \tag{DO23.5}
\]

Take their \(V\) to be the positive target of rank \(b\), and their \(W\) to have coefficient \(-2i\kappa I_a\), hence signature \((p',q')=(0,a)\). Formula (DO23.5) becomes (DO23.1) entry by entry: \(\operatorname{Re}(2i\kappa\bar X Y)
 =2\kappa\operatorname{Im}(X\bar Y)\). Their matrix right action, equation (3.1), is \(X.\iota(h,g)=h^{-1}Xg\). Their Schrödinger representation satisfies covariance with the inverse of this right action. This can be checked directly from their §2.1 formula \(\omega(a)\phi(x)=|\det a|^{1/2}\phi(xa)\): conjugation sends translation by \(y\) to translation by \(ya^{-1}\). Thus their group arguments \((h,g)\), without inverting either group argument, give exactly (DO23.2).

There is a necessary character-label conversion. Konno and Konno's pair is \[
                       (\xi,\xi')=(\chi_S,\chi_T^{-1}).
\tag{DO23.5a}
\] Their source splitting changes by \(\nu_{\Delta\xi'}(\det g)\), whereas the source convention fixed above changes by its inverse. This can be read directly from the character factor \(\xi'(\det a)\) in the doubled source formula (3.10), printed p.48; the same positive determinant convention is stated for the zero-dimensional opposite member immediately above Lemma 3.1. Thus §4.1's source exponent is \(-m'\), not \(m'\). This is a conversion of splitting labels; no inversion of the group argument or change of the already verified Heisenberg action is made. Since \(d\psi=2\pi i\), one has \(d\psi\cdot i<0\). Proposition 4.4(ii), printed pp.68--69, gives on the positive target block the scalar \((m+q'-p')/2=(m+a)/2\), together with the standard creation action. On the negative source block it gives the scalar \((-m'+q-p)/2=(-m'-b)/2\), together with the dual creation action. Lemma 5.2(i), printed p.73, exponentiates these formulas to the compact groups. This proves (DO23.4) with the fixed standard Harris--Kudla--Sweet splitting, rather than merely determining it up to a determinant character.

As a direct sign check, in one normalized real coordinate the operators \(D_A=-2\pi ix\), \(D_B=-\partial_x\) implement the Heisenberg action used in (DO18.5). The rotation generator for multiplication by \(e^{i\theta}\) is \(L=-i\partial_x^2/(4\pi)+i\pi x^2\); it has vacuum eigenvalue \(i/2\) on \(e^{-\pi x^2}\). The creation operator has commutator \([L,r]=ir\). This confirms the positive target polynomial sign; the exact determinant-character shift is supplied by the preceding primary formula, not inferred from this sign check alone.

\subsection{2. The first lift's exact target types}\label{the-first-lifts-exact-target-types}

Use the actual Jacobian \(J\), its CM type \(\Phi=\{1,2,3\}\), and the splitting chosen in ALP-LIFT-09. At each of the three selected embeddings the element \(\xi=(\zeta_7-\zeta_7^{-1})/7\) is a positive multiple of \(i\). Thus its negative skew coefficient has the sign in (DO23.1) at every selected place. The character \(\chi_S=\mu_E\) has exponents

\[
                             (m_1,m_2,m_3)=(-1,-1,+1).
 \tag{DO23.6}
\]

Indeed \(\mu_E=\mu_-\circ N_{E/\mathbf Q(\sqrt{-7})}\), and \(\mu_{-,\infty}(z)=|z|/z\). The embeddings 1 and 2 restrict to the chosen imaginary-quadratic embedding, whereas 3 restricts to its conjugate; the latter reverses the exponent. This is also immediate from the quadratic residues \(\{1,2,4\}\pmod7\).

For the first edge \(a=1,b=2\), (DO23.4) gives the target vacuum characters

\[
                         1,\qquad1,\qquad\det_{U(2)}.
 \tag{DO23.7}
\]

ALP-LIFT-09's explicit nonzero vector applies \(\partial_1\partial_3\) in each of the two ordinary Jacobian factors. Therefore its total Fock degrees at the three selected CM places are \((2,0,2)\). In a rank-one source oscillator, the degree-\(d\) target polynomial space is \(\operatorname{Sym}^d\mathbf C^2\), an irreducible \(U(2)\)-module. The monomial with one creation in each of the two coordinates is a nonzero vector in \(\operatorname{Sym}^2\), even though it is not a highest-weight vector. Hence the target compact module containing the first lift is exactly

\[
  \operatorname{Sym}^2\mathbf C^2\ \boxtimes\ \mathbf1\
       \boxtimes\left(\det\otimes\operatorname{Sym}^2\mathbf C^2\right),
 \tag{DO23.8}
\]

with dominant highest weights

\[
                              (2,0),\qquad(0,0),\qquad(3,1).
 \tag{DO23.9}
\]

The actual theta map is equivariant for right translation on target automorphic functions: \(R_h\vartheta_\phi=\vartheta_{\omega(h)\phi}\). The source torus integral preserves this equivariance because the two unitary actions commute. It is nonzero on the displayed irreducible compact module, by the exact identity value in ALP-LIFT-09. Its kernel is an invariant submodule, hence is zero on that module. Thus (DO23.8) is a type of the actual nonzero automorphic first image, not merely a formal polynomial type before taking theta.

ALP-LIFT-09 labels its tensor source torus by direct multiplication \(t\). Identifying it with the Hom source can use \(g=t^{-1}\). This changes the source eigencharacter \(\beta\) by inversion. It does not change the target group, the opposite skew form, \(\mu_E\), or (DO23.7)--(DO23.9). The first-lift statement defined \(\beta\) by its actual action, so it is unaffected by this relabeling.

\subsection{3. Vanishing for the unchanged positive Hom second edge}\label{vanishing-for-the-unchanged-positive-hom-second-edge}

Now fix a second edge with source the positive \(U(2)\), target the positive \(U(4)\), the same positive-complex Hom convention (DO23.1)--(DO23.2), and \(\chi_S=\chi_T=1\). Formula (DO23.4) restricts its source Fock representation to

\[
           \det_{U(2)}^{-2}\otimes
                    \operatorname{Sym}(\mathbf C^4\otimes(\mathbf C^2)^*).
 \tag{DO23.10}
\]

Every irreducible source type is of the form \(\det^{-2}\otimes S_\lambda((\mathbf C^2)^*)\), where \(\lambda=(\lambda_1\ge\lambda_2\ge0)\). Its dominant highest weight is

\[
                         (-2-\lambda_2,-2-\lambda_1).
 \tag{DO23.11}
\]

For completeness, this follows by decomposing each homogeneous polynomial degree under the compact group, or by the polynomial Schur decomposition. Dualization reverses and negates the two dominant weights; twisting by \(\det^{-2}\) subtracts 2 from each. In particular both entries in (DO23.11) are at most \(-2\). None of the weights in (DO23.9) occurs; the second selected CM place alone already excludes the trivial \((0,0)\) type.

This proves that the actual first compact module (DO23.8) has zero local theta image under this specified second oscillator. It also proves the global theta integral is zero for every Schwartz input at that second edge, using the convention \(\int\vartheta(g,h)\overline{f(g)}\,dg\). One elementary global justification is to average first over the compact real source group. Its Haar average is the orthogonal projection onto matching compact types. The theta kernel's source types lie in (DO23.11), and \(f\)'s type is (DO23.9), so that projection is zero. The compact averaging can be inserted by right invariance of Haar measure. The definite automorphic quotient is compact, so no regularization or exchange of divergent integrals is involved.

The conclusion remains true for the contragredient of (DO23.8) at this unchanged edge: the second selected place is still trivial, which is absent from (DO23.10). Changing the source splitting character, changing the sign/polarization of the edge, or using a different coefficient module changes the displayed spectrum and must be analyzed separately. In particular a Gram composition with the adjoint theta operator is a different construction; it can detect the first image and is not ruled out by this calculation.

The recovered source uses the bilinear convention \(\int\vartheta(g,h)f(g)\,dg\) for \nolinkurl{liftT}, as stated in \nolinkurl{Development/AbelianVarietyTheory/Definitions/TwistedInnerProducts/WittSmallGroupEmbeddingMemPointwiseStabilizerIso.lean}, \nolinkurl{liftT_apply}, lines 160--167. It therefore tests the contragredient of the input type against the source oscillator. The preceding trivial-place argument also proves vanishing for this exact source convention at the specified unchanged edge. It must be used when comparing a general chain's nontrivial characters and coefficients; one cannot transfer the conjugated-input matching rule unchanged.

\subsection{4. A quantitative necessary shift, without asserting a global repair}\label{a-quantitative-necessary-shift-without-asserting-a-global-repair}

If the second edge keeps its Hom geometry but has target character \(\chi_T(z)=(z/\bar z)^{m'/2}\), in the source character convention, its source vacuum exponent is \(c=(-m'-4)/2\). Its source dominant weights are \((c-\lambda_2,c-\lambda_1)\). A given dominant weight \((r,s)\), \(r\ge s\), occurs in this polynomial model if and only if \(c\ge r\): choose \(\lambda=(c-s,c-r)\). Thus the types (DO23.9) require respectively

\[
                             m'_1\le-8,\qquad
                             m'_2\le-4,\qquad
                             m'_3\le-10.
 \tag{DO23.12}
\]

These are exact local occurrence conditions for this source module. They are not a construction of a global automorphic splitting character with those infinity types, and do not prove finite occurrence or global nonvanishing for the repaired second lift. A global character would also have to satisfy rational-unit and finite-level compatibility. The unchanged even-rank trivial characters have \(m'_i=0\) and fail all three inequalities.

\textbf{Correction record.} An earlier version identified Konno's \(\xi'\) with the source label \(\chi_T\), overlooking the inverse in \nolinkurl{splitFacO}. It consequently had \((m'-b)/2\) in (DO23.4) and the reversed conditions in (DO23.12). Equations (DO23.4), (DO23.5a), and (DO23.12) now use the source convention consistently. The actual first target types and the unchanged \(\chi_T=1\) second-edge vanishing were unaffected by that error.

\subsection{Primary reference}\label{primary-reference}

Konno, T., \& Konno, K. (2007). On doubling construction for real unitary dual pairs. \emph{Kyushu Journal of Mathematics, 61}(1), 35--82. \url{https://doi.org/10.2206/kyushujm.61.35} Primary PDF: \url{https://www.jstage.jst.go.jp/article/kyushujm/61/1/61_1_35/_pdf} Exact locators: §2.1, printed pp.37--38, Heisenberg and covariance convention; §3.1, pp.43--44, forms and equation (3.1); §3.3, pp.47--48, honest splittings; §4.1, p.49, character exponents; Proposition 4.4(ii), pp.68--69, exact Fock operators; Lemma 5.2(i), p.73, compact action. The paper's displayed last page is 82; an external metadata entry listing 85 is not used for the page range.

Carried first-lift construction: ALP-LIFT-09, especially (V7)--(V11), (V18)--(V26), and (V32), with the independent whole-proof check already recorded in \nolinkurl{definite_oscillator_CHECK_EVENTS.tsv}.

\clearpage

\section{ALP-LIFT-11: a fixed-conductor determinant repair of the specified second-edge type obstruction}\label{alp-lift-11-a-fixed-conductor-determinant-repair-of-the-specified-second-edge-type-obstruction}

Author: /root/lift\_accounting, GPT-6 Astra Pro, ALP\_MACHINE. 21 November 2026 UTC. This result stays over \(E=\mathbf Q(\zeta_7)\) and \(F=E^+\). It is not a common-host identification with the \(\mathbf Q(\zeta_{28})\) family of AM-DO-022.

AM-DO-023 supplies the primary-source-checked honest-splitting Fock formula and the specified second-edge vanishing theorem. They are used here as established results. The new construction is a global splitting character of uniformly bounded finite conductor, its exact effect on the actual first theta lift, and the resulting removal of that particular real-type obstruction. It does not assert nonvanishing of the next global lift.

\subsection{1. An explicit global norm-one character of any integral infinity type}\label{an-explicit-global-norm-one-character-of-any-integral-infinity-type}

Put \(T=U(1)(E/F)\) and \(K_{T,0}=(\widehat{\mathcal O}_E)^1\). ALP-LIFT-09 proves \[
 T(\mathbb A_{F,f})=T(F)K_{T,0},\qquad
 T(F)\cap K_{T,0}=\mu_{14}.
 \tag{R1}
\] In particular \[
 T(F)\backslash T(\mathbb A_F)
       \simeq(T_\infty\times K_{T,0})/\mu_{14},
 \tag{R2}
\] where the roots of unity are embedded diagonally. These arithmetic statements are used here without re-proving their Minkowski and unit-signature arguments.

Choose the three complex coordinates indexed by \(\Phi=\{1,2,3\}\), so that the rational root \(\zeta\) has archimedean coordinates \((\zeta,\zeta^2,\zeta^3)\). For any integer vector \(\mathbf c=(c_1,c_2,c_3)\), put \[
 \beta_{\mathbf c}(z_1,z_2,z_3)=
               z_1^{c_1}z_2^{c_2}z_3^{c_3},\qquad
 C_0=c_1+c_2+c_3,\quad S_0=c_1+2c_2+3c_3 .
 \tag{R3}
\] Here each \(z_i\) lies on the unit circle.

Choose one of the three split places \(v_*\) over29 and its first \(E\) factor so that \(\zeta\) reduces to \(16\) in \(\mathbf F_{29}^{\times}\). This is possible because16 has order7 modulo29. The element2 generates \(\mathbf F_{29}^{\times}\): \(2^{14}=-1\) and \(2^4=16\) has order7. The projection of \(K_{T,0,v_*}\) onto that first factor, followed by reduction, gives a surjection \[
 K_{T,0,v_*}\longrightarrow \mathbf F_{29}^{\times},
 \qquad \zeta\longmapsto 2^4,\quad -1\longmapsto2^{14}.
 \tag{R4}
\] It is injective on the rational subgroup \(\mu_{14}\).

Choose \(r\in\{0,\ldots,13\}\) with \[
                   r\equiv 7C_0+6S_0\pmod{14}.
 \tag{R5}
\] Define the finite compact character \(\rho_{\mathbf c}\) to be trivial at every place other than \(v_*\) and to satisfy \[
            \rho_{\mathbf c}(2)=\exp(2\pi i r/28)
 \tag{R6}
\] on the residue group at \(v_*\). On \(\zeta\) its value is \(\zeta^r\), while \(\beta_{\mathbf c}(\zeta)=\zeta^{S_0}\). Equation (R5) gives \(r\equiv-S_0\pmod7\). On \(-1\) the two values are \((-1)^r\) and \((-1)^{C_0}\), and \(r\equiv C_0\pmod2\). Thus \[
       \rho_{\mathbf c}|_{\mu_{14}}
                    =\beta_{\mathbf c}^{-1}|_{\mu_{14}}.
 \tag{R7}
\] The product \(\beta_{\mathbf c}\rho_{\mathbf c}\) therefore descends through (R2) to a continuous unitary automorphic character \(\gamma_{\mathbf c}\) of \(T\).

This construction is explicit on arbitrary ideles. Write \(t_f=\delta k\) with \(\delta\in T(F)\) and \(k\in K_{T,0}\), as in (R1), and set \[
 \gamma_{\mathbf c}(t_\infty,t_f)
      =\beta_{\mathbf c}(\delta_\infty^{-1}t_\infty)
                         \rho_{\mathbf c}(k).
 \tag{R8}
\] If \(\delta\) is changed by \(\epsilon\in\mu_{14}\), the compact factor is changed by \(\epsilon^{-1}\). Equation (R7) makes the product unchanged. Multiplicativity follows because the torus is abelian, and the same formula shows triviality on rational diagonal points.

Its infinity type is exactly \(\mathbf c\). Its finite restriction to the maximal compact has conductor dividing the first power of the fixed place \(v_*\), and all those compact values are roots of unity of order dividing28. In particular \[
 \gamma_{\mathbf c,f}|_{K_M}=1
                  \qquad\text{for every }M=29^e,\ e\ge1,
 \tag{R9}
\] where \(K_M\) is ALP-LIFT-10's source torus subgroup. The finite conductor bound is independent of the size of \(\mathbf c\); the unramified values at noncompact split places are not asserted to be roots of unity.

Every finite value is algebraic. In (R8), the finite character value is a root of unity, and \(\beta_{\mathbf c}(\delta_\infty)\) is a product of integer powers of conjugates of the algebraic number \(\delta\). For a fixed \(\mathbf c\) these values belong to a fixed finite coefficient field containing the normal closure of \(E\) and \(\mu_{28}\).

At a nonsplit finite place \(v\), every norm-one element of \(E_v^\times\) is a unit, since its valuation has norm valuation zero. Hence \(T(F_v)=K_{T,0,v}\). Our character is trivial on that entire group, not merely on a smaller unit subgroup: \[
                 \gamma_{\mathbf c,v}=1
                         \quad(v\text{ nonsplit finite}).
 \tag{R10}
\] This fact removes a potentially hidden ramified splitting factor in the next step.

\subsection{2. Exact change of the actual first splitting and theta integral}\label{exact-change-of-the-actual-first-splitting-and-theta-integral}

Let \[
 j:\mathbb A_E^\times\longrightarrow T(\mathbb A_F),
                         \qquad j(x)=x/\bar x .
 \tag{R11}
\] Keep the actual first pair of ALP-LIFT-09--10: the positive group \(G=U(I_2)\) and the skew line torus \(T\), with splitting characters \(\chi_V=1\) and \(\chi_W=\mu_E\). Replace only the latter by \[
              \chi_W'=\mu_E\,(\gamma_{\mathbf c}\circ j),
              \qquad\chi_V'=1 .
 \tag{R12}
\] The restriction of the multiplying factor to \(\mathbb A_F^\times\) is1, so the required restriction \(\chi_W'|_{\mathbb A_F^\times}=\epsilon_{E/F}\) is unchanged. Rational diagonal triviality follows from that of \(\gamma_{\mathbf c}\). Thus (R12) is an actual globally admissible splitting character, not just prescribed archimedean weights.

For the standard splitting used in these results, its actual Weil operators satisfy \[
 \boxed{\quad
 \omega'(g,t)=\gamma_{\mathbf c}(\det g)\,\omega(g,t).
 \quad}
 \tag{R13}
\] Here are all local checks of this identity. At every nonsplit finite place, (R10) says that \(\gamma_{\mathbf c,v}\circ j\) is literally the trivial character of \(E_v^\times\), so the local splitting has not changed. At every split finite place, the exact determinant root formula of AM-DO-020 evaluates the ratio on an element \(y\) with \(y/\bar y=\det g\); its value is precisely \(\gamma_{\mathbf c,v}(\det g)\). This includes possible unramified values at split places outside29. At each real place, the character \(\gamma_{\mathbf c}\circ j\) changes the exponent \(m_i\) of AM-DO-023 by \(2c_i\). Its exact Fock formula therefore changes the target determinant power by \(c_i\), giving the same scalar. The source action depends on the unchanged opposite character \(\chi_V\). These checks prove (R13) at every place and hence globally.

For any of the actual finite Schwartz inputs and Fock vectors already constructed, the theta series and first theta integral consequently satisfy \[
 \vartheta'_\phi(g,t)
       =\gamma_{\mathbf c}(\det g)\vartheta_\phi(g,t),
 \qquad
 \Theta'(\phi,f)(g)
       =\gamma_{\mathbf c}(\det g)\Theta(\phi,f)(g).
 \tag{R14}
\] The source torus input \(f\) is unchanged. In particular all identity values, including their retained common period and their 11-adic normalized valuations, are unchanged exactly. The twisted first lift is nonzero whenever the original one is nonzero. Since the character has modulus one, the scalar automorphic \(L^2\) norm and supremum norm of the entire first lift are unchanged as well.

\subsection{3. Removing the specified real-type obstruction in each integral convention}\label{removing-the-specified-real-type-obstruction-in-each-integral-convention}

For a first Fock vector of fixed CM multidegree \(\mathbf d=(d_1,d_2,d_3)\), put \[
                         \boldsymbol\nu=(0,0,1).
 \tag{R15}
\] AM-DO-023 proves that the actual first target compact types are \[
              \pi_i=\det^{\nu_i}\operatorname{Sym}^{d_i}
                               (\mathbf C^2),
 \qquad\operatorname{HW}(\pi_i)=(d_i+\nu_i,\nu_i).
 \tag{R16}
\] The specified next positive Hom edge \(U(2)\longrightarrow U(4)\) with both splitting characters trivial has source types \[
 \det^{-2}S_\lambda((\mathbf C^2)^*),\qquad
 \operatorname{HW}=(-2-\lambda_2,-2-\lambda_1),
 \quad\lambda_1\ge\lambda_2\ge0.
 \tag{R17}
\] In particular the untwisted types (R16) do not occur. This is the already proved scoped vanishing of AM-DO-023, not a universal no-go for every chain.

\subsubsection{The conjugated-input integral}\label{the-conjugated-input-integral}

First use the Hermitian theta integral \(\int\vartheta(g,h)\overline{f(g)}\,dg\), as in ALP-LIFT-09. Its source compact type must match that of \(f\) itself.

Choose in the global construction of section1 \[
                         c_i=-d_i-2-\nu_i .
 \tag{R18}
\] By (R13)--(R14), the actual twisted first image now has dominant weights \[
 (d_i+\nu_i+c_i,\nu_i+c_i)=(-2,-d_i-2).
 \tag{R19}
\] These are exactly the weights in (R17) for the one-row partition \(\lambda=(d_i,0)\). Thus the twisted first compact representation occurs in the second oscillator, with the same degree \(d_i\) at each place. The repair uses a genuine global splitting of finite conductor bounded independently of \(\mathbf d\).

For the explicit nonzero first vector of ALP-LIFT-09, \(\mathbf d=(2,0,2)\) and \(\boldsymbol\nu=(0,0,1)\). One takes \[
 \mathbf c=(-4,-2,-5),\qquad C_0=-11,\quad S_0=-23,
 \qquad r=9 .
 \tag{R20}
\] The residue character at29 is \(\rho(2)=\exp(2\pi i\,9/28)\), and the repaired first weights are \[
                       (-2,-4),\quad(-2,-2),\quad(-2,-4).
 \tag{R21}
\] All these occur in the specified second source model.

A single fixed repair also works for any finite prefix of degrees. Given bounds \(D_i\) and \(0\le d_i\le D_i\), use \(c_i=-D_i-2-\nu_i\). The resulting source weights are \((d_i-D_i-2,-D_i-2)\), which occur for \[
                    \lambda_i=(D_i,D_i-d_i).
 \tag{R22}
\] The next Fock degree is then \(2D_i-d_i\). This variant changes the needed polynomial degree and any corresponding archimedean coefficient norm; it is not a zero-cost identification of the next Fock vectors.

Conversely, for this conjugated-input composition with two fixed positive Hom edges and fixed splitting characters, the second source spectrum has a fixed upper bound on its highest-weight entries. Equation (R16) then excludes sufficiently large \(d_i\). The degree-dependent repair or a finite-prefix repair changes that datum. This observation concerns the displayed orientation and character choices; it does not apply without checking the alternating signs, contragredients and pins in a different proposed source chain.

\subsubsection{The source's literal bilinear integral has a fixed repair}\label{the-sources-literal-bilinear-integral-has-a-fixed-repair}

The recovered source liftT instead uses \[
                   \int_{[G]}\vartheta(g,h)f(g)\,dg .
 \tag{R23}
\] Its local matching condition pairs the oscillator source type with the contragredient of the input type. In this convention a stronger, degree-independent repair is available. Choose \[
              \mathbf c=\mathbf 2-\boldsymbol\nu=(2,2,1).
 \tag{R24}
\] Then the actual first image has weights \((d_i+2,2)\), whose contragredient weights are \[
                        (-2,-d_i-2).
 \tag{R25}
\] They are precisely (R17) for \(\lambda=(d_i,0)\), for every nonnegative degree \(d_i\) simultaneously. The next Fock degree stays \(d_i\). Thus neither a degree-dependent splitting nor a finite-prefix enlargement of the next polynomial degree is needed for the bilinear convention.

Here \(C_0=5\), \(S_0=9\), and (R5) gives \(r=5\). The fixed finite character is \(\rho(2)=\exp(2\pi i\,5/28)\), and \[
       (m_1,m_2,m_3)+2\mathbf c=(3,3,3).
 \tag{R26}
\] Consequently one actual fixed character \(\chi_W'=\mu_E(\gamma_{(2,2,1)}\circ j)\) repairs all of these real types in the literal bilinear composition. This is an independent convention-specific conclusion; it does not turn the two different theta integral conventions into the same map.

Equivalently, (R23) is the conjugated-input lift of \(\bar f\), whose compact type is the contragredient. This reformulation is useful when applying an inner-product theorem written with the conjugated-input convention.

\subsection{4. Integral class values and the actual cost scope}\label{integral-class-values-and-the-actual-cost-scope}

The identity values and scalar automorphic norms are preserved by (R14) without choosing representatives. An integral class-value claim requires a specified finite target level and its coefficient lattice. AM-R-008 proves the relevant precise statement: with compact-determinant representatives and a character trivial on \(\det K\), multiplication by \(\gamma_f(\det g)\) is a diagonal root-of-unity map. It preserves the entire local content and the mass norm, with the coefficient action twisted by \(\gamma_\infty^{-1}\) in that note's left-equivariant convention. This inverse is consistent with the positive right-translation twist in (R14); the two statements use their respective explicit function conventions.

For the explicit group level of AM-R-008, (R9) verifies its character hypothesis. There is consequently no additional denominator, content or logarithmic norm cost for the determinant twist on that class module. All its displayed finite class scalars have order dividing28.

The source torus stabilizer \(K_M\) in ALP-LIFT-10 must not be confused with a proved target \(U(2)\) stabilizer. This note does not assert that the original first vectors are invariant under the particular target group level of AM-R-008 without checking that target action. Every finite family of their Schwartz vectors has an open target stabilizer; it can be refined inside the kernel of \(\gamma_f\circ\det\). Applying a compatible compact-determinant class presentation then transports its actual class-value module by the same root-of-unity diagonal map. Its original stabilizer index, mass and coefficient lattice remain costs to be computed if needed.

Likewise this does not identify the transported module with an independently pinned arithmetic row. The first edge has ranks1 and2 and already lies outside the source's even-rank equal-character StdPair constructor. The later PinX.mv and PinColT.om interfaces require the same specified character on both members of an even-rank edge. Changing only \(\chi_W\) in (R12) is a valid actual first theta lift, but is not a verification of those different equal-character pins.

\subsection{5. What remains for a second actual lift}\label{what-remains-for-a-second-actual-lift}

The repaired first input has exactly the real representation required by the chosen second oscillator, and it retains the actual nonzero first value and first norm. Local occurrence at infinity is a necessary condition, not a proof that its global next theta integral is nonzero. One must still select actual finite second-kernel inputs, prove a nonzero global pairing, and retain the resulting mass and norm factors. A formal Gram picture or a hyperbolic Witt-tower persistence theorem does not establish this for a sequence of positive definite groups.

The \(\mathbf Q(\zeta_{28})\) host, the full DP65 conjugated-action rows, source tied direction matrices, and the terminal six-identity family remain separate comparisons. This construction removes a specific real-type obstruction without claiming the complete requested chain.

\subsection{Exact source locators}\label{exact-source-locators}

\begin{itemize}
\tightlist
\item
  ALP-LIFT-09, actual\_torus\_first\_lift.md, (V2)--(V11), torus class number and the original global splitting, and (V32)--(V38), actual first lift and projectors.
\item
  ALP-LIFT-10, permitted\_orbit\_first\_image.md, (O21), (O23)--(O28), (O31)--(O35), the exact fixed source levels, normalized values and input costs.
\item
  AM-DO-020, exact split-place determinant-root covariance; AM-DO-023, definite\_oscillator\_real\_second\_edge.md, (DO23.4)--(DO23.11), actual honest real target weights and the trivial-source-character case. Its corrected nontrivial source-character sign is not needed here.
\item
  Konno, T., \& Konno, K. (2007). On doubling construction for real unitary dual pairs. Kyushu Journal of Mathematics, 61(1), 35--82. Proposition4.4(ii), pp.68--69, and Lemma5.2(i), p.73. \url{https://www.jstage.jst.go.jp/article/kyushujm/61/1/61_1_35/_pdf}
\item
  AM-R-008, compact\_determinant\_twist\_accounting.md, (R8.2), (R8.10)--(R8.15), the actual integral class-value and mass isometry with its stated level hypothesis.
\item
  The actual source bilinear convention is \nolinkurl{WittSmallGroupEmbeddingMemPointwiseStabilizerIso.lean,} liftT\_apply, lines160--167, as directly checked in AM-DO-023/024's source audit.
\end{itemize}

\subsection{Check status}\label{check-status-2}

The exact finite-residue and weight calculations are included in \nolinkurl{fixed_conductor_type_repair_certificate.py.} Author and independent whole-proof checks, with their actual scope and time, are recorded separately in \nolinkurl{lift_accounting_CHECK_EVENTS.tsv.} No Lean execution is claimed.

\clearpage

\section{AM-K019 --- what the full theta differential module cannot supply}\label{am-k019-what-the-full-theta-differential-module-cannot-supply}

Author: \nolinkurl{/root/kummer_extension}, GPT-6 Astra Pro, ALP\_MACHINE, 21 November 2026 UTC. This is a new, elementary obstruction to one proposed iteration argument. It does not obstruct a selected harmonic functional, an evaluated jet ideal, or a quotient with separately controlled saturation. The accepted source connection-integrality theorem is used as a theorem.

\subsection{1. Exact statement}\label{exact-statement}

Let \(J=\operatorname{Jac}(y^2=x^7-1)\), let \(u_1,u_2,u_3\) be the invariant logarithmic coordinates of AM-K012, and let \(f\) be either the actual translated, value-one reduced theta function \(F\), or the actual trace-one reduced section \(\sigma_q\). On the complement of its zero divisor put \[
z_i=\partial_i\log f.
\tag{K20.1}
\] The following assertions hold.

\begin{enumerate}
\def\labelenumi{\arabic{enumi}.}
\tightlist
\item
  The six meromorphic functions \(u_1,u_2,u_3,z_1,z_2,z_3\) are algebraically independent over the field \(\mathbb C(J)\) of abelian functions. In particular, adjoining only the three logarithmic coordinate multipliers does not put \(f\) in a finite-dimensional vector space over \(\mathbb C(J)(u_1,u_2,u_3)\) preserved by ordinary invariant differentiation.
\item
  Adjoining the three \(z_i\) does give the expected polynomial rank-one connection module. In its natural integral polynomial lattice, however, every ordinary differentiated vector \(\partial_i^N f/f\) has a monic term \(z_i^N\). No positive restricted power of that connection maps the whole lattice into \(p\) times the lattice, and no lattice commensurable with it by a bounded power of \(p\) can have that property.
\item
  The source's actual integral connection on its auxiliary-variable module has an analogous obstruction: its creation term produces a unit coefficient of auxiliary degree \(N\) after \(N\) iterations. At prime-to-\(p\) tensor power and in a primitive tangent direction, a positive restricted power cannot map that whole integral module into \(p\) times itself, even after a bounded-denominator commensurable change of lattice.
\end{enumerate}

The distinction in assertion 2 is essential. The abstract polynomial lattice injects into the characteristic-zero evaluated Hurwitz series, but its reduction modulo \(p\) need not inject. In the present translated value-one example, each \(z_i\) has zero constant term, so \(z_i^p=0\) in the reduced Hurwitz ring. Positive \(p\)-adic bounds on evaluated derivatives therefore coexist with the polynomial-lattice obstruction.

\subsection{2. The lattice shifts span all six complex directions}\label{the-lattice-shifts-span-all-six-complex-directions}

Write \(J(\mathbb C)=V/\Lambda\), where \(V=\mathbb C^3\). Its principal polarization is represented by the nondegenerate positive Hermitian form \(H\), linear in the first variable. The actual reduced theta multiplier, checked in AM-K018, is \[
f(u+\lambda)=\alpha(\lambda)
 \exp\left(\pi H(u,\lambda)+\frac\pi2H(\lambda,\lambda)\right)f(u).
\tag{K20.2}
\] Consequently \[
z_i(u+\lambda)=z_i(u)+\pi H(e_i,\lambda).
\tag{K20.3}
\] If a translated representative changes the semicharacter or a constant factor, differentiation removes that change, so (K20.3) is identical for \(F\) and \(\sigma_q\).

Let \(C\) be the nonsingular complex matrix with \((\pi H(e_i,\lambda))_i=C\overline\lambda\). The shift subgroup of the six displayed functions is \[
\Gamma=\{(\lambda,C\overline\lambda):\lambda\in\Lambda\}
 \subset\mathbb C^6.
\tag{K20.4}
\] Choose a real lattice basis \(\lambda_1,\ldots,\lambda_6\). The six vectors \((\lambda_j,C\overline{\lambda_j})\) are complex-linearly independent. Indeed, a complex-linear functional \((a,b)\) vanishing on them vanishes on their real span. Since the \(\lambda_j\) span \(\mathbb C^3\) over \(\mathbb R\), this gives \[
a\lambda+bC\overline\lambda=0\qquad(\lambda\in\mathbb C^3).
\] Replacing \(\lambda\) by \(i\lambda\) and combining the two equations gives \(a\lambda=0\) and \(bC\overline\lambda=0\) for every \(\lambda\). Thus \(a=b=0\). The six-column square matrix has full complex rank.

It follows that \(\Gamma\) is Zariski dense in \(\mathbb C^6\): after an invertible complex linear transformation this is the familiar assertion that \(\mathbb Z^6\) is Zariski dense. The latter follows by induction on the number of variables, applying the fact that a nonzero univariate polynomial has only finitely many roots.

Suppose a polynomial relation among \(u,z\) has coefficients in \(\mathbb C(J)\). Its pullback is an identity of meromorphic functions on \(V\). Choose \(u_0\) avoiding the finitely many coefficient poles and the zero divisor of \(f\). Periodicity of the coefficient functions and of the theta divisor permits substitution of every \(u_0+\lambda\). The polynomial with coefficients evaluated at \(J(u_0)\) vanishes on \((u_0,z(u_0))+\Gamma\), so Zariski density makes every one of those coefficients zero. Varying \(u_0\) on a nonempty open set makes every original rational coefficient zero. This proves assertion 1's algebraic-independence claim.

The same independence holds after any algebraic extension of the abelian-function field, for example the finite ordered-divisor cover used in AM-K012. Algebraic extension does not change the transcendence degree of the composite field. Thus the explicit Abel-coordinate ring used for integral calculations causes no exception to the argument.

\subsection{\texorpdfstring{3. Monic derivatives exclude finite rank over abelian functions and \(u\)}{3. Monic derivatives exclude finite rank over abelian functions and u}}\label{monic-derivatives-exclude-finite-rank-over-abelian-functions-and-u}

Put \(K=\mathbb C(J)\). The Hessians \[
\partial_i z_j=-W_{ij}
\tag{K20.5}
\] are abelian functions. For \(F\) they are the translated Kleinian functions of AM-K012; for \(\sigma_q\) they are the corresponding unshifted ones. In either case \(D_iK\subset K\).

Fix \(i\). Recursively set \[
P_0(Z)=1,\qquad P_{N+1}(Z)=D_iP_N(Z)+ZP_N(Z),
\quad D_iZ=-W_{ii}.
\tag{K20.6}
\] Then \[
\frac{\partial_i^N f}{f}=P_N(z_i),\qquad
P_N(Z)=Z^N+\text{terms of degree below }N,
\quad P_N\in K[Z].
\tag{K20.7}
\] Both statements follow by induction. Differentiating a coefficient preserves \(K\); differentiating \(Z^r\) lowers its \(Z\)-degree; multiplication by \(Z\) produces the next unique leading term with coefficient one.

The \(P_N(z_i)\) are linearly independent over \(K(u_1,u_2,u_3)\), since \(z_i\) is transcendental over that field and their degrees are strictly increasing. A finite-dimensional vector space over that field, consisting of actual meromorphic germs and containing \(f\), cannot be preserved by \(\partial_i\): otherwise the infinite sequence \(f,\partial_i f,\partial_i^2 f,\ldots\) would be linearly dependent. This proves the stated finite-rank obstruction. It concerns ordinary differentiation on the actual reduced functions. It does not exclude a different connection on an auxiliary bundle, or assert that no rank-one connection exists after adjoining the \(z_i\).

For algebraic constant-field variants, this conclusion persists on extension from a number field to a characteristic-zero completion. To see the needed elementary descent, clear the finitely many rational denominators in a hypothetical relation. It becomes a finite linear relation, over the extended constant field, among coefficientwise number-field formal germs. Linear independence of any finite collection of such germs is witnessed by a nonzero finite coefficient minor; that minor remains nonzero after constant-field extension. Hence no such new relation appears. This statement is about finitely many rational-function coefficients and does not assert descent of an infinite analytic relation.

\subsection{4. The natural polynomial connection is integral but has no full-module gain}\label{the-natural-polynomial-connection-is-integral-but-has-no-full-module-gain}

Let \(\mathcal O\) be a good local coefficient DVR, with uniformizer \(p\), and use the explicit \(D\)-stable integral ordered-divisor ring \(R\) of AM-K012 near the translated non-effective characteristic. Thus \(D_iR\subset R\), \(W_{ij}\in R\), and all evaluated coefficients lie in the integral Hurwitz ring. A harmless unramified extension splits the branch points; no ramified change of uniformizer is being hidden.

Define \[
A=R[Z_1,Z_2,Z_3],\qquad
\delta_i|_R=D_i,\qquad \delta_i Z_j=-W_{ij}.
\tag{K20.8}
\] The derivations commute because the mixed third derivatives of the actual theta logarithm agree. On the free rank-one \(A\)-module \(M=Ae\) define \[
\nabla_i(ae)=(\delta_i a+Z_i a)e.
\tag{K20.9}
\] Its characteristic-zero realization is \(Z_i\mapsto z_i\), \(e\mapsto F\). The algebraic independence above makes this realization faithful over the fraction field. Formula (K20.7) gives \[
\nabla_i^Ne=P_N(Z_i)e,
\qquad [Z_i^N]P_N=1.
\tag{K20.10}
\] All coefficients are integral; one of them is a unit. Therefore \[
\nabla_i^Q M\not\subset pM\qquad\text{for every }Q\ge1.
\tag{K20.11}
\] This conclusion does not depend on the nilpotence length of the derivation \(D_i\) on \(R/pR\).

Here is the precise bounded-denominator strengthening. Suppose a lattice \(M'\) in \(M[1/p]\) satisfies \[
p^a e\in M',\qquad M'\subset p^{-C}M
\tag{K20.12}
\] for fixed nonnegative integers \(a,C\). These hypotheses hold for any lattice commensurable with \(M\) by a bounded power of \(p\). If, for some \(Q\ge1\), \(\nabla_i^Q M'\subset pM'\), then iteration gives \[
p^aP_{nQ}(Z_i)e=\nabla_i^{nQ}(p^ae)
 \in p^nM'\subset p^{n-C}M.
\tag{K20.13}
\] The coefficient of \(Z_i^{nQ}\) on the left is \(p^a\). For \(n>a+C\) this contradicts (K20.13). Thus no such commensurable lattice supplies the proposed uniform restricted-power divisibility.

\subsubsection{The evaluated Hurwitz image is a different lattice problem}\label{the-evaluated-hurwitz-image-is-a-different-lattice-problem}

For the translated \(F\), \(z_i(0)=0\). Reduction of a Hurwitz series with zero constant term has \(p\)-th power zero: the multinomial coefficients in a \(p\)-fold product cancel modulo \(p\), equivalently the only surviving Frobenius term is the \(p\)-th power of the constant coefficient. Hence \[
\operatorname{ev}(Z_i^p)=z_i^p=0
 \quad\text{in the characteristic-}p\text{ Hurwitz ring},
\tag{K20.14}
\] whereas \(Z_i^p\ne0\) in \(A/pA\). The reduction of evaluation has a nontrivial kernel even though characteristic-zero evaluation is injective. It also admits additional integral elements after saturation in its characteristic-zero image; for instance the image of \(Z_i^p/p\) is integral although the abstract polynomial is not in \(A\).

Consequently (K20.11) does \textbf{not} refute the actual congruences of AM-K017, nor the actual all-order lower bounds of AM-K013--015. It refutes the attempt to infer those bounds by a uniform divisibility operator on this unsaturated polynomial lattice, or on a lattice differing from it by a fixed power of \(p\). An iterative argument in the evaluated image must specify that image, its \(p\)-saturation, and the growth of the saturation cost. The isolated mod-\(p\) theta identities supply none of those bounds by themselves.

\subsection{5. The same issue is visible in the actual source connection}\label{the-same-issue-is-visible-in-the-actual-source-connection}

The source definition is \nolinkurl{ReciprocityTheory/ReductiveGroups/ArithmeticGeometry/TaylorDiscriminants/ConnectionKillsSeries.lean}, \nolinkurl{sgOf} at 18--19, \nolinkurl{sg1} at 21--22, \nolinkurl{nabla} at 31--35, \nolinkurl{Hset} at 41--43, and \nolinkurl{ConnectionIntegrality} at 77--83. The accepted theorem realizing the latter is \nolinkurl{WickVanishing/ExpConjConnectionIntegral.lean}, lines 40--82. These exact files were inspected; no rederivation of their established integral-connection assertion is made.

Use \(X\) for its formal-group variables and \(Z=(Z_1,\ldots,Z_d)\) for its auxiliary variables. Put \(s_i(X)=-\partial_{Y_i}h(X,0)\). The source operator is \[
\nabla_{u,\phi}
 =D_u^X+\sum_i(D_us_i+\phi_i)\partial_{Z_i}
  -\kappa\sum_i(Z_i-s_i)u_i.
\tag{K20.15}
\] Its admitted integral lattice of directions has \(\phi=\phi_0-Su\) with \(u,\phi_0\) integral and \(S=ds(0)\). Conjugating by the source factor \(\exp(\kappa h_L(X))\), which depends only on \(X\), gives an operator of the form \[
T=e^{-\kappa h_L}\nabla_{u,\phi}e^{\kappa h_L}
 =D_u^X+\sum_i A_i(X)\partial_{Z_i}
  -\kappa\sum_i u_iZ_i+b(X).
\tag{K20.16}
\] The accepted source theorem states that \(T\) preserves the integral formal-power-series module \(\mathcal N=\mathcal O[[X,Z]]\). Equivalently the coefficients in this particular first-order expression are integral; this also follows directly by applying preservation to \(1,X_j,Z_i\). For the argument below, preservation itself is sufficient.

Assume \(p\nmid\kappa\) and \(u\) is primitive, so some \(u_j\) is a unit. Starting at \(1\), \(T^N1\) is polynomial in \(Z\) of degree at most \(N\), with coefficients in \(\mathcal O[[X]]\). Every term of \(T\) except multiplication by \(-\kappa u\cdot Z\) preserves or lowers \(Z\)-degree. Induction therefore gives the exact highest homogeneous part \[
\operatorname{top}_Z(T^N1)=(-\kappa u\cdot Z)^N.
\tag{K20.17}
\] Its coefficient of \(Z_j^N\) is \((-\kappa u_j)^N\), a unit. Hence \[
T^Q\mathcal N\not\subset p\mathcal N
 \qquad(Q\ge1).
\tag{K20.18}
\] The same argument as (K20.12)--(K20.13), replacing \(e\) by \(1\), rules out a bounded-denominator commensurable lattice with such a uniform gain. The source's integral connection preserves an integral lattice; it does not make a positive power of each creation operator divisible by \(p\) on the whole lattice.

This is fully compatible with the source \nolinkurl{ValueSt}, lines 71--75: after the specified constant-coefficient evaluation, the relevant connection polynomial gives an ordinary invariant derivative. That evaluation annihilates the positive auxiliary powers responsible for (K20.17). Harmonic contractions and other selected quotients can likewise remove them. Any stronger gain must be proved for the actual selected evaluation or quotient, including its saturation and normalization, rather than deduced from full-module nilpotence that is false.

\subsection{6. Consequence for the current proof-transfer attempt}\label{consequence-for-the-current-proof-transfer-attempt}

AM-K018 identifies the full-\(\mathcal O_E\) source gauge with the actual Chern/Baker gauge on ordinary powers of \(J\). The obstruction above does not undo that identification. It specifies why the full mixed mod-11 congruences of AM-K017 cannot automatically be iterated through a finite-rank ordinary-differential module, and why even the actual source integral connection does not supply the proposed uniform restricted-power inclusion on its full auxiliary module.

The viable remaining arithmetic question is therefore more precise: determine the actual evaluated or harmonic-selected theta image, bound its \(p\)-saturation in the required degree, and compare that bound with the projection and normalization costs. The accepted DP54/58/59 examples now apply in the exact Chern gauge at 11; DP45/53 at the order-six Frobenius classes remain bare integral-function inputs until their actual primitive theta kernel or selected-value interface is supplied.

Primary source for (K20.2): Bannai, K., \& Kobayashi, S. (2010). Algebraic theta functions and the \(p\)-adic interpolation of Eisenstein--Kronecker numbers. \emph{Duke Mathematical Journal, 153}(2), 229--295, Definition 1.8 in \href{https://arxiv.org/pdf/math/0610163}{arXiv:math/0610163v4}. The actual multiplier and source compatibility were proved in AM-K018; the arguments in §§2--5 above are new elementary arguments, with the accepted source integral-connection theorem explicitly separated.

\clearpage

\section{AM-K020 --- two sharp pure directions of the actual trace-one theta section}\label{am-k020-two-sharp-pure-directions-of-the-actual-trace-one-theta-section}

Author: /root/kummer\_extension, GPT-6 Astra Pro, ALP\_MACHINE, 21 November 2026 UTC. The full actual source gauge was identified in AM-K018. The supplied DP41 bounded-multiplier lemma and DP45/53 bare cyclic-display bounds are used as accepted theorems. The new result concerns the single fixed actual section, not an arbitrary integral test function.

\subsection{1. Statement}\label{statement-11}

Let \(J=\operatorname{Jac}(y^2=x^7-1)\), and use the global invariant differential basis \(\omega_i=x^{i-1}dx/y\) and its logarithmic coordinates \(u_i\). Let \(\sigma_q\) be the actual reduced theta germ of the trace-one principal line \(L_{2\infty}\), normalized by \[
\sigma_q(u)=u_1u_3-u_2^2+\text{higher terms}.
\] Its complete algebraic normalization, actual Chern/source gauge, and integral content are AM-K014, AM-K018, and ALP-LIFT-08. Write \(v_p(p)=1\), with \(v_p(0)=+\infty\).

The following are exact sharp directional statements: \[
\boxed{\begin{aligned}
v_3\!\left(\partial_1^n\sigma_q(0)\right)
 &\ge \max\left\{0,\frac n{56}-\frac{s_3(n)}2\right\},
&
\liminf_{n\to\infty}
\frac{v_3(\partial_1^n\sigma_q(0))}{n}
 &=\frac1{56},\\[2mm]
v_5\!\left(\partial_2^n\sigma_q(0)\right)
 &\ge \max\left\{0,\frac n{504}-\frac{s_5(n)}4\right\},
&
\liminf_{n\to\infty}
\frac{v_5(\partial_2^n\sigma_q(0))}{n}
 &=\frac1{504}.
\end{aligned}}
\tag{K21.1}
\] Here \(s_p(n)\) is the base-\(p\) digit sum. These are derivatives along two precisely specified global CM eigenlines, on \(J\) itself. They require no conjugated factor and no comparison of mixed-action product gauges.

There are also explicit sharp witness sequences, not just a nonconstructive liminf. Put \[
n_3(r)=28\cdot3^r+18,\qquad n_5(r)=126\cdot5^r+30.
\] Then \[
\boxed{\begin{aligned}
v_3(\partial_1^{n_3(r)}\sigma_q(0))
 &=\frac{3^r+3-s_3(n_3(r))}{2}\quad(r\ge0),\\
v_5(\partial_2^{n_5(r)}\sigma_q(0))
 &=\frac{5^r-1}{4}\quad(r\ge0).
\end{aligned}}
\tag{K21.1a}
\] For \(r\ge3\), the first expression is \((3^r-1)/2\). Its values for \(r=0,1,2\) are \(0,1,5\); the last has one extra unit of valuation from a base-three carry. Every derivative in (K21.1a) is nonzero.

The same sharp slopes hold for the translated, value-one theta \(F\) in those directions. For \(\sigma_q\) they persist under every permitted pullback with the source normalization \(M^{-n}\), since its reduced pullback is \(\sigma_q(Mu)\) and the two factors cancel exactly. In particular the previously used odd levels \(M=29^e\) are prime to both 3 and 5.

The accepted DP45/53 bare integral-function slopes in these directions are respectively \(3/56\) and \(5/504\). Thus these actual normalized theta derivatives have slope equal to \(1/p\) times that bare slope at these two primes. This comparison is an actual matching lower/asymptotic upper result for one fixed section; it is not a conclusion about the total cost of a later automorphic row.

\subsection{2. An exact weighted-initial-form lemma}\label{an-exact-weighted-initial-form-lemma}

Let \(p\) be odd, \(q=p^h\), \(s=1/(p-1)\), and \[
H(t)=\sum_{n\ge2}\ell_n\frac{t^n}{n!}.
\] Assume \[
\ell_n=0\ (n\text{ odd}),\qquad
v(\ell_n)\ge \max\{0,\lfloor(n-2)/q\rfloor\},\qquad
v(\ell_{q+1})=0.
\tag{K21.2}
\] Put \[
c=\frac{s}{q+1},\qquad w_*=s-c,\qquad
I=[s-2c,s-c].
\tag{K21.3}
\] For a rational \(w\in I\), choose a finite ramified coefficient extension and \(\gamma\) with \(v(\gamma)=w\). Write \(\delta=w_*-w\in[0,c]\) and \[
\tau=s-(q+1)\delta\in[0,s].
\tag{K21.4}
\] Then \(H(\gamma T)\) has coefficient infimum exactly \(\tau\), and its entire initial form, after division by an element of valuation \(\tau\), is a nonzero scalar times \(T^{q+1}\). In particular that initial form has nonzero formal derivative in characteristic \(p\).

\textbf{Proof.} For even \(n\), Legendre's formula gives \[
v\!\left(\frac{\ell_n\gamma^n}{n!}\right)
=v(\ell_n)+s\,s_p(n)-(c+\delta)n.
\tag{K21.5}
\] An even positive integer has \(s_p(n)\ge2\). At \(n=q+1\), the digit sum is exactly two, and the unit assumption makes (K21.5) equal \(2s-(c+\delta)(q+1)=\tau\).

If \(2\le n<q+1\), the same expression is at least \(2s-(c+\delta)n\), exceeding \(\tau\) by at least \((c+\delta)(q+1-n)>0\).

If \(n>q+1\) is even, put \(k=\lfloor(n-2)/q\rfloor\ge1\). Then \(n\le(k+1)q+1\), so (K21.5) minus \(\tau\) is at least \[
k-(c+\delta)(n-q-1)
\ge k\{1-q(c+\delta)\}
\ge k\left(1-\frac{2qs}{q+1}\right)>0.
\tag{K21.6}
\] The last inequality holds because \(p\ge3\) implies \(2s\le1\). These estimates also bound all coefficients below and show that the higher-degree excess grows at least linearly in \(k\). Odd coefficients are zero. There is therefore one and only one initial monomial. Its degree is \(q+1\equiv1\pmod p\), proving the derivative assertion. ∎

Let \(\nu_w(G)=\inf_n(v([t^n]G)+wn)\). Suppose \(V(t)\ne0\) is any formal series for which \(V(\gamma T)\) is bounded for every rational \(w\in I\), and integral for \(w=w_*\). The accepted DP41 bounded-multiplier lemma applies to the full initial form just proved, not merely to the existence of a first pole. It gives \[
\nu_w(Ve^H)=-\infty\quad(w<w_*),
\qquad
0\le\nu_w(Ve^H)<\infty\quad(w\ge w_*).
\tag{K21.7}
\] For the first assertion take \(w\in I\), \(w<w_*\), so \(0\le\tau<s\), and use the nonzero-derivative initial form. The multiplier is only required to be nonzero and bounded; a zero at the origin is allowed. Monotonicity of \(\nu_w\) extends unboundedness to every smaller nonnegative weight. At \(w=w_*\), all coefficients of \(H(\gamma T)\) have valuation at least \(s\), so its exponential and its inverse are integral, with constant term one. This gives the asserted boundedness and finite infimum; larger weights follow by monotonicity.

If \(Ve^H=\sum a_nt^n\), (K21.7) implies \[
v(a_n)\ge-w_*n,\qquad
\liminf_{n\to\infty}\frac{v(a_n)}n=-w_*.
\tag{K21.8}
\] Indeed, a strictly larger liminf would make the coefficient infimum finite at some rational weight below \(w_*\), a contradiction. Multiplication by \(n!\) and \(s_p(n)/n\to0\) give \[
v(n!a_n)\ge cn-s\,s_p(n),\qquad
\liminf_{n\to\infty}\frac{v(n!a_n)}n=c.
\tag{K21.9}
\] All assertions also hold with \(H(2t)\) in place of \(H(t)\), since \(p\) is odd and multiplication of its \(n\)-th logarithmic derivative by \(2^n\) changes no valuation or initial-form nonvanishing.

\subsection{3. Apply the lemma to the actual translated theta}\label{apply-the-lemma-to-the-actual-translated-theta}

For the actual value-one \(F\), AM-K012 proves \[
\partial_i^2\log F=-W_{ii}(a+u)
\] in the explicit integral, invariant-derivation-stable Abel-coordinate ring. AM-K013 proves the all-order operator bound \(D_i^mR\subset p^{\lfloor m/q\rfloor}R\) at the indicated nilpotence exponent, and hence exactly the logarithmic bounds (K21.2), except for its final unit condition. This is an iterated assertion on a stable integral ring, not an iteration of one coefficient congruence.

The exceptional constants were computed exactly in the primary algebraic theta recurrence of AM-K017: \[
\begin{array}{c|c|c|c}
p&i&q&\ell_{q+1}\bmod p\\ \hline
3&1&27&2\\
5&2&125&1.
\end{array}
\tag{K21.10}
\] The computations are in \(\mathbb F_p[Z]/\Phi_7(Z)\), which is a field in both cases because \(p\) has order six modulo seven. The residues shown are nonzero scalars, so the logarithmic derivatives are units. The actual trace-one section is not being substituted for \(F\) in the definition of these logarithmic derivatives: \(F(0)=1\), whereas \(\sigma_q(0)=0\).

These exact values and their full algebraic ODE inputs are saved in:

\begin{verbatim}
work/agents/jacobian_hessian_pcurvature_certificate.py
work/agents/jacobian_hessian_pcurvature_certificate.json
work/agents/jacobian_theta_direction_jets.py
\end{verbatim}

The certificate regenerates the values using the branch divisor \((\zeta,\zeta^3,\zeta^5)\), the exact Kleinian Hessian, and division-free Hurwitz recurrences. Its recorded PASS is a finite arithmetic check. The all-order deduction here uses (K21.2) and the complete initial-form lemma above.

With \(V=1\), (K21.9) proves the sharp slopes \(1/56\) and \(1/504\) for the two indicated restrictions of \(F\).

\subsection{4. Transfer to the actual vanishing source section}\label{transfer-to-the-actual-vanishing-source-section}

AM-K015 proved an equality of actual reduced germs \[
\sigma_q(2u)=r(u)F(2u),
\tag{K21.11}
\] where \(r\) is a rational function on \(J\) regular and integral on a formal neighborhood of zero at every odd prime other than 7. This comes from the canonical rigidified \([2]\)-identification of the two theta lines and their actual integral sections; it is not an arbitrary integral formal prefactor. The equality is in the Chern gauge, so no additional quadratic exponential occurs. The value-one normalization of \(F\) is a unit at 3 and 5 by ALP-LIFT-07, and the \(q\)-normalized source section is primitive there by ALP-LIFT-08.

Write \(V_i(t)=r(te_i)\). Its restriction is nonzero in the two cases: \[
\sigma_q(t,0,0)=\frac{t^4}{12}+O(t^{18}),\qquad
\sigma_q(0,t,0)=-t^2+O(t^{16}).
\tag{K21.12}
\] These initial terms are exact coefficients of the actual section in AM-K014. Since \(F(2te_i)\) is a formal unit, (K21.11) implies \(V_i\ne0\). In fact \(V_i(0)=0\) in both cases. No unit quotient is assumed.

The remaining hypothesis of (K21.7) is boundedness of this actual restricted prefactor. The accepted DP45/53 theorem gives, in these globally labelled integral eigendirections, \[
c_i^{\mathrm{bare}}(p)=\frac{p}{(p-1)(p^3+1)}=pc,
\qquad q=p^3.
\tag{K21.13}
\] For every integral formal germ \(r\), its invariant Taylor coefficients therefore satisfy \[
v([t^n]r(te_i))\ge-(s-pc)n.
\tag{K21.14}
\] This is the accepted mixed bound with only direction \(i\) retained, followed by the exact factorial formula; any zero coefficients only improve it. The global invariant basis differs from DP45's primitive display basis by diagonal units, as explicitly checked in DP53.

Every \(w\in I\) satisfies \[
w\ge s-2c>s-pc,
\tag{K21.15}
\] because \(p\ge3\). Thus \(V_i(\gamma T)\) is integral and bounded throughout the whole initial-form interval. Its zero at the origin is harmless for DP41's nonzero bounded-multiplier lemma. Apply (K21.9) to \(V_i(t)\exp(H_i(2t))=\sigma_q(2te_i)\). Finally \(2^n\) is a unit, so its derivatives have exactly the same valuations as those of \(\sigma_q(te_i)\). This proves the matching liminf and the displayed digit-loss lower bound in (K21.1). The extra lower bound zero follows from the actual odd-prime Hurwitz integrality of \(\sigma_q\) in AM-K014.

\subsection{5. Explicit sharp subsequences of the fixed source section}\label{explicit-sharp-subsequences-of-the-fixed-source-section}

\subsubsection{5.1. Exact finite Sato determinant, without partition enumeration}\label{exact-finite-sato-determinant-without-partition-enumeration}

Two additional finite residues are \[
\partial_1^{18}\sigma_q(0)\equiv1\pmod3,\qquad
\partial_2^{30}\sigma_q(0)\equiv4\pmod5.
\tag{K21.17}
\] Here is a direct finite certificate and its truncation justification. Use AM-K014's actual Sato frame \[
v_j(t)=t^{\,2-n_j}A(t^{14})^{-\lfloor n_j/2\rfloor},
\quad A=1-TA^7,\quad(n_j)=(0,2,4,6,7,8,\ldots).
\] Write \(a_{m,j}=[t^m]v_j\). For a single KP time \(t_d=v\), put \(h_{dn}=v^n/n!\) and \(h_m=0\) when \(m<0\) or \(d\nmid m\). Then \[
\tau(v)=
\det\left(\sum_m a_{m,j}h_{m+i}(v)\right)_{1\le i,j\le L}
\pmod{v^{N+1}}
\qquad(L\ge\max\{3,dN\}).
\tag{K21.18}
\] To prove this, apply Cauchy--Binet, ordering the selected exponents \(m_1>\cdots>m_L\ge-L\). Set \(\lambda_i=m_i+i\). These form a partition, and the two determinants in each summand are respectively its stable Plücker coefficient and \(\det(h_{\lambda_j-j+i})=s_\lambda\), by Jacobi--Trudi. Every term of ordinary \(v\)-degree at most \(N\) has \(|\lambda|\le dN\), hence length at most \(L\). All its Plücker coefficients already stabilize at this \(L\): the added rows and columns have upper triangular identity tail, since \(v_j\) has leading exponent \(-j\) for \(j\ge4\) and only higher corrections. Thus no omitted partition can contribute. The sum over \(m\) is finite at each truncation, which also justifies Cauchy--Binet formally.

The exact integer coefficients needed for the matrix are \[
[T^k]A^{-m}
=\frac m k(-1)^{k-1}\binom{7k-m-1}{k-1}\quad(k\ge1),
\qquad [T^0]A^{-m}=1.
\tag{K21.19}
\] The generalized integer binomial is used if the upper entry is negative. Its displayed numerator is divisible by \(k\), as also follows from the integral implicit recursion for \(A\).

The script

\begin{verbatim}
work/agents/jacobian_sigma_hurwitz_determinant.py
\end{verbatim}

evaluates (K21.18) in the truncated Hurwitz ring. Its multiplication is \((ab)_n=\sum_{k=0}^n\binom nk a_kb_{n-k}\); it never inverts \(n!\) modulo \(p\). Complete unit pivoting leaves the actual two-dimensional constant-rank defect and an exact \(2\times2\) residual determinant. Finally the factor \(4\cdot2^{-n}\) converts the \(n\)-th derivative of \(\tau(v)=\sigma(v)\) into that of \(\sigma_q(u)=4\sigma(u/2)\).

For the two computations, \((d,N)=(5,18)\) and \((3,30)\), so \(L=90\) is sufficient. Each was executed again with \(L=104\); the complete retained derivative vectors agreed. The outputs are

\begin{verbatim}
work/agents/jacobian_sigma_hurwitz_p3_order18.json
work/agents/jacobian_sigma_hurwitz_p3_order18_tail14.json
work/agents/jacobian_sigma_hurwitz_p5_order30.json
work/agents/jacobian_sigma_hurwitz_p5_order30_tail14.json
\end{verbatim}

The nonzero residues through those cutoffs are respectively \(S_4=2,S_{18}=1\) modulo 3, and \(S_2=3,S_{16}=4,S_{30}=4\) modulo 5. Separate low-order executions reproduce \(S_4=2\), \(S_2=-2\), and \(S_6=1\) in their corresponding directions. These are exact finite computations, distinct from the all-order argument that follows. The one-command replay is

\begin{verbatim}
python work/agents/jacobian_sigma_sharp_certificate.py
\end{verbatim}

It regenerates both matrix sizes, verifies the entire retained vectors, and writes jacobian\_sigma\_sharp\_certificate.json with script hashes, exact critical weights, finite possible initial degrees, and witness digit sums. Its recorded status is PASS.

\subsubsection{5.2. The critical exponential coefficients}\label{the-critical-exponential-coefficients}

At \(w=w_*\), the scaled logarithm \(H(2\gamma T)\) has exactly one coefficient of valuation \(s\), in degree \(n_0=q+1\); every other coefficient has valuation strictly above \(s\). Write \(E(T)=\exp(H(2\gamma T))\). Its constant term is one, and \[
v([T^n]E)\ge s\quad(n>0),\qquad
v([T^n]E)=s
\ \Longleftrightarrow\ n=n_0p^r\ (r\ge0).
\tag{K21.20}
\] For the equivalence, factor the exponential formally into the exponentials of its separate logarithmic monomials. A nonzero coefficient from a logarithmic coefficient of valuation \(s+\epsilon\) has valuation at least \(s\,s_p(m)+m\epsilon>s\). The unique critical monomial contributes in degree \(mn_0\) with valuation \(ms-v(m!)=s\,s_p(m)\), which equals \(s\) exactly when \(m=p^r\). All other contributions at such a degree have larger valuation, so there is no cancellation. Each coefficient involves only finitely many factors, making this a formal argument rather than an analytic product assertion.

\subsubsection{5.3. The actual algebraic prefactor has a unique critical monomial}\label{the-actual-algebraic-prefactor-has-a-unique-critical-monomial}

Set \(\widetilde V(T)=r(\gamma T e_i)\) and \(G(T)=\sigma_q(2\gamma T e_i)=\widetilde V(T)E(T)\). By (K21.14), \[
v([T^j]\widetilde V)\ge\delta_0j,\qquad
\delta_0=(p-1)c>0.
\tag{K21.21}
\] Thus its coefficient infimum \(\nu\) is attained in finitely many degrees. Since \(E(0)=1\) and every positive coefficient of \(E\) has valuation at least \(s>0\), Gauss multiplicativity and reduction show that \(G\) and \(\widetilde V\) have the same infimum and the same entire initial form. In particular their initial degrees must satisfy the actual CM selection rule of \(\sigma_q\).

For \(p=3\), the exact unit \(S_{18}\) gives \[
\nu\le -v_3(18!)+18w_*=-8+\frac{27\cdot18}{56}
=\frac{19}{28}.
\] Any initial degree of \(\widetilde V\) is at most \((19/28)/\delta_0=19\). The source CM rule leaves only degrees 4 and 18. The value at 4 is \(-1+4w_*=13/14>19/28\), so \[
\nu=\frac{19}{28},\qquad
\operatorname{in}(\widetilde V)=aT^{18}\quad(a\ne0).
\tag{K21.22}
\]

For \(p=5\), the exact unit \(S_{30}\) gives \[
\nu\le -v_5(30!)+30w_*=-7+\frac{125\cdot30}{504}
=\frac{37}{84}.
\] Initial degrees are at most \((37/84)/\delta_0=111/2\), hence the selection rule leaves only 2, 16, 30, and 44. The value at degree 2 is \(2w_*=125/252>37/84\). Hurwitz integrality alone bounds the values at 16 and 44 below by respectively \(s\,s_5(16)-16c=1-16/504\) and \(s\,s_5(44)-44c=2-44/504\), both larger than \(37/84\). Therefore \[
\nu=\frac{37}{84},\qquad
\operatorname{in}(\widetilde V)=aT^{30}\quad(a\ne0).
\tag{K21.23}
\]

\subsubsection{5.4. Unique coefficient contributions give exact valuations}\label{unique-coefficient-contributions-give-exact-valuations}

In either case put \(m=18\) or 30, respectively, and \(n=n_0p^r+m\). The contribution \([T^m]\widetilde V\,[T^{n_0p^r}]E\) has valuation exactly \(\nu+s\). Every other contribution with a positive \(E\)-degree has larger valuation, by the uniqueness in (K21.22)--(K21.23) and (K21.20). The sole constant-\(E\) contribution is \([T^n]\widetilde V\); it also has larger valuation because \[
\delta_0(n_0+m)>\nu+s:
\qquad \frac{46}{28}>\frac{33}{28},\quad
\frac{156}{126}>\frac{29}{42}.
\tag{K21.24}
\] Thus for every \(r\ge0\), \[
v([T^{n_0p^r+m}]G)=\nu+s.
\tag{K21.25}
\] Undo the weight and the unit rescaling by 2, and multiply by \(n!\). Since \(v(S_m)=0\) and \(s_p(m)=2\) in both cases, the result is \[
v(S_n)=cn+\nu+s-s\,s_p(n)
=\frac{p^r+3-s_p(n)}{p-1}.
\tag{K21.26}
\] For \(p=5\), adding \(30=5+25\) to \(126\cdot5^r\) causes no base-five carry, so \(s_5(n)=4\) for every \(r\ge0\). For \(p=3\), adding \(18=2\cdot3^2\) causes precisely one carry when \(r=2\) and none for \(r=0,1\) or \(r\ge3\). This proves every assertion of (K21.1a), including the small-order exception. It also independently yields the two liminf upper bounds in (K21.1).

\subsection{6. Source scope, exact comparison, and the remaining general-prime condition}\label{source-scope-exact-comparison-and-the-remaining-general-prime-condition}

The resulting exact comparison is

\begin{longtable}[]{@{}
  >{\raggedright\arraybackslash}p{(\columnwidth - 6\tabcolsep) * \real{0.2000}}
  >{\raggedleft\arraybackslash}p{(\columnwidth - 6\tabcolsep) * \real{0.2667}}
  >{\raggedleft\arraybackslash}p{(\columnwidth - 6\tabcolsep) * \real{0.2667}}
  >{\raggedleft\arraybackslash}p{(\columnwidth - 6\tabcolsep) * \real{0.2667}}@{}}
\toprule\noalign{}
\begin{minipage}[b]{\linewidth}\raggedright
Prime and fixed global direction
\end{minipage} & \begin{minipage}[b]{\linewidth}\raggedleft
Actual bare DP45/53 slope
\end{minipage} & \begin{minipage}[b]{\linewidth}\raggedleft
Actual \(\sigma_q\) sharp slope
\end{minipage} & \begin{minipage}[b]{\linewidth}\raggedleft
Difference
\end{minipage} \\
\midrule\noalign{}
\endhead
\bottomrule\noalign{}
\endlastfoot
\(p=3\), \(\partial_1\) & \(3/56\) & \(1/56\) & \(1/28\) \\
\(p=5\), \(\partial_2\) & \(5/504\) & \(1/504\) & \(1/126\) \\
\end{longtable}

The operators are already specified global CM directions; no prime-dependent relabelling occurs in either statement. Each pure monomial is killed by conjugate-pair contractions whenever it is placed alone in a positive CM-labelled polynomial block. That elementary harmonic observation does not assert membership in an unspecified full row with its other degree, invariance, and HC conditions. Nor does the construction identify these particular pure derivative sequences with a required later theta-lift image.

The same proof gives a precise generalization criterion. For \(p\equiv3\pmod7\), take \(i=1\); for \(p\equiv5\pmod7\), take \(i=2\). In either case \(q=p^3\), the all-order logarithmic bound and the bare bound (K21.13) already hold for the actual Jacobian. If the single actual constant \[
\partial_i^{p^3+1}\log F(0)
\tag{K21.16}
\] is a \(p\)-adic unit, then the actual fixed source germ has the exact sharp slope \(1/((p-1)(p^3+1))\) in that direction. AM-K017 makes this a global Hessian restricted-power constant; it is independent of the chosen nonzero local theta frame. Its unit status for every prime in these two residue classes is not proved here. A unit claim cannot be inferred solely from the existence of a rank-one cyclic display without identifying the relevant universal-extension or primitive-kernel pairing.

Thus the new result is unconditional for the two explicitly certified primes 3 and 5, and gives a single checkable unit condition at every other prime in the two classes. It supplies actual directional upper sharpness without a conjugated-product gauge gap. It does not settle the full prime budget or the later selected-row cost.

\subsection{7. Exact accepted-result locators and verification boundary}\label{exact-accepted-result-locators-and-verification-boundary}

The accepted source archive is RETURN\_DP\_MULT\_checkpoint\_06\_20261121T0420Z. The used statements are:

\begin{itemize}
\tightlist
\item
  research/DP41\_integrated\_fragment.tex, proof paragraph ``Exponential obstruction, including arbitrary nonzero bounded multipliers,'' lines 131--147. Its lemma allows a complete ramified coefficient DVR, a logarithm with coefficient infimum in \([0,1/(p-1))\), a nonzero-derivative entire initial form, and any nonzero bounded multiplier. Its proof and hypotheses were read to verify applicability; the accepted lemma is not rederived here.
\item
  research/DP45\_integrated\_fragment.tex, the actual integral cyclic-display theorem and its primitive directions; research/DP53\_integrated\_fragment.tex, the full globally labelled spectrum and the diagonal-unit comparison of the fixed differential basis.
\item
  AM-K012--015, the actual Hessian, stable integral ring, {[}2{]}-ratio, and source sigma normalization; AM-K017, the complete finite certificate for (K21.10); AM-K018, the exact full-\(\mathcal O_E\) Chern/source gauge.
\end{itemize}

The new weighted inequalities (K21.5)--(K21.6), the multiplier verification (K21.11)--(K21.15), the exact finite determinant specialization and residues (K21.17)--(K21.19), and the actual explicit witness sequences (K21.20)--(K21.26) are the proof contributions of this note. The underlying genus-three theta identities are cited and source-audited in AM-K012--014; this note does not introduce a new analytic theta identity.

\clearpage

\section{AM-R-009: an actual nonzero first lift on the common unramified mixed-CM host}\label{am-r-009-an-actual-nonzero-first-lift-on-the-common-unramified-mixed-cm-host}

Author: /root, GPT-6 Astra Pro, ALP\_MACHINE, 21 November 2026 UTC. This proof constructs the first image on the same field and polarized abelian variety as AM-DO-022--024. It does not identify the rank-one starting edge with the source's prescribed even-rank arithmetic chain. The finite compact-vacuum assertion in section 3 has been checked against Pan (2001), section 6.2, including the rank-one source case. The source input has norm one with the specified probability Haar measure.

\subsection{1. The precise host, period, and eigen-coordinate}\label{the-precise-host-period-and-eigen-coordinate}

Use the accepted DP57 host, as normalized in AM-DO-022: \[
 E'=\mathbf Q(\zeta_{28})=\mathbf Q(\zeta_7,i),\qquad
 F'=E'^{+},\qquad A=J^2,\qquad \xi'=\xi/2.
 \tag{R9.1}
\] The action of \(i\) on the two ordinary \(J\) factors is \((u,v)\mapsto(-v,u)\). The trace form with coefficient \(\xi'\) is exactly the product principal polarization, with no isogeny or change to the underlying four Jacobian factors used below. Put \[
 w_j^+=u_j+iv_j,\qquad w_j^-=u_j-iv_j,\qquad
 s_A(u,v)=\sigma_q(u)\sigma_q(v),\qquad C=\theta_0c.
 \tag{R9.2}
\] All symbols \(\sigma_q,\theta_0,c\), and \(G_j\) retain their normalizations in AM-R-004--006. In particular \(\sigma_q(u)=u_1u_3-u_2^2+\) terms of degree greater than two, and the actual Bargmann theta realization of the original finite vacuum is \(C\sigma_q\).

Since \(\zeta_{28}=\zeta_7^2i^{-1}\), its eigenvalue on \(w_j^+\) has label \(8j-7\bmod28\), while its label on \(w_j^-\) is \(8j+7\bmod28\). Thus \[
 \Phi'=\{1,3,9,15,17,23\},\qquad w_2^+\text{ has label }9.
 \tag{R9.3}
\] The derivative in that coordinate is \[
 \partial_{w_2^+}=\tfrac12(\partial_{u_2}-i\partial_{v_2}).
 \tag{R9.4}
\] The first nonzero homogeneous term of \(s_A\) has degree four. Consequently, in its fourth derivative only the term with two derivatives on each factor contributes: \[
 \partial_{w_2^+}^{\,4}s_A(0)
   =2^{-4}\binom42(-i)^2(-2)(-2)=-\frac32.
 \tag{R9.5}
\] This is an exact coefficient calculation, independent of analytic tail estimates.

Take the target positive space \(V=E'^2\) and the source skew line \(W\) with coefficient \(-2\xi'\), in the half-trace Hom convention of AM-DO-022. On the underlying \(A^2=J^4\), let \[
 P=\partial_{w_{2,1}^+}^{\,4}\partial_{w_{2,2}^+}^{\,4}.
 \tag{R9.6}
\] The actual finite vacuum is the tensor of the four original finite Schwartz vacua; its norm is one. Its Bargmann realization is \[
 \mathcal R\phi_{00}=C^4s_A(z_1)s_A(z_2).
 \tag{R9.7}
\] Apply the actual creation operators corresponding to \(P\), as in ALP-LIFT-09, and denote the resulting global Schwartz vector by \(\phi_P\). Equations (R9.5)--(R9.7) give \[
 \boxed{\vartheta_{\phi_P}(1,1)=\frac94\,C^4\ne0.}
 \tag{R9.8}
\] Its degree vector is eight at CM label 9 and zero at all other labels. Hence this is a single source-torus isotypic Fock vector, not a sum whose nonzero components remain unidentified.

The period \(C^4\) has not been set equal to one. In the algebraic coordinate normalization the displayed coefficient is \(9/4\); it is an \(11\)-adic unit. The coordinate change (R9.4) has a dyadic denominator. No claim of integrality at 2 or at a changed bad-prime Hodge basis is inferred from the coefficient at 11.

\subsection{2. A globally admissible splitting, with its infinity labels}\label{a-globally-admissible-splitting-with-its-infinity-labels}

Let \(\mu_E\) be the explicit global character constructed in ALP-LIFT-09, so that \(\mu_E|_{\mathbf A_F^\times}=\epsilon_{E/F}\). Define \[
 \mu'=\mu_E\circ N_{E'/E},\qquad \chi_S=\mu',\qquad \chi_T=1.
 \tag{R9.9}
\] The fields satisfy \(E'=EF'\) and \(E\cap F'=F\). Restriction of the norm in (R9.9) to \(\mathbf A_{F'}^\times\) is \(N_{F'/F}\). Compatibility of the quadratic idele-class character with base change therefore gives \[
 \mu'|_{\mathbf A_{F'}^\times}
   =\epsilon_{E/F}\circ N_{F'/F}
   =\epsilon_{E'/F'}.
 \tag{R9.10}
\] This is the required restriction for a source skew line; the even-dimensional target requires the trivial restriction and has the chosen character 1. Use the same rationally compatible standard Weil representation as in the cited constructions.

At the ordered CM labels \(1,3,9,15,17,23\), the exponents \(\mu'_\tau(z)=(z/\bar z)^{m_\tau/2}\) are \[
 \mathbf m=(-1,+1,-1,-1,+1,-1),\qquad
 \boldsymbol\nu=(\mathbf m+\mathbf1)/2=(0,1,0,0,1,0).
 \tag{R9.11}
\] Indeed restriction to \(E\) gives the three labels \(1,3,2,1,3,2\); ALP-LIFT-09 and the corrected AM-DO-023 give exponent \(-1\) on labels 1 and 2 and \(+1\) on label 3. The target compact type of the actual first image is accordingly \[
 \pi_\infty=\mathop{\boxtimes}_{\tau\in\Phi'}
       \det^{\nu_\tau}\operatorname{Sym}^{d_\tau}\mathbf C^2,
 \quad d_9=8,\quad d_\tau=0\ (\tau\ne9).
 \tag{R9.12}
\] For the Hom source action \(X\mapsto Xt^{-1}\), its actual torus eigencharacter is \[
 \beta(t)=\prod_{\tau\in\Phi'}\tau(t)^{-1}\,\tau_9(t)^{-8}.
 \tag{R9.13}
\] Here the exponent \(-1\) at each place is the rank-two target's vacuum shift with \(\chi_T=1\). The character \(\chi_S\) changes the target action, not this source action. On the rational generator \(\zeta_{28}\), (R9.13) has exponent \[
 -\sum_{\tau\in\Phi'}\tau-8\cdot9=-68-72=-140\equiv0\pmod{28}.
 \tag{R9.14}
\] Thus \(\beta\) is trivial on \(\mu_{28}\).

\subsection{3. The torus quotient and the actual input}\label{the-torus-quotient-and-the-actual-input}

Put \(T'=U(1)(E'/F')\) and \(K_0=(\widehat{\mathcal O}_{E'})^1\). Masley (1976), Main Theorem, printed pp.~179 and 184, implies \(h(E')=1\): \(\varphi(28)=12\le24\), and 28 does not occur in the displayed list with class numbers 2 through 10. This published class-number result is an input, not a new calculation or a computational class-group certificate.

It implies the exact torus decomposition \[
 T'(\mathbf A_{F',f})=T'(F')K_0,\qquad
 T'(F')\cap K_0=\mu_{28}.
 \tag{R9.15}
\] Here is a short proof of the first assertion in the required norm-one form. For a norm-one finite idele \(t\), its fractional ideal \(I(t)\) satisfies \(I(t)\overline{I(t)}=\mathcal O_{E'}\). Every nonsplit prime therefore has exponent zero, and the exponents at a split conjugate pair are opposite. Choose a fractional ideal \(B\), supported on one prime in each such pair, so that \(I(t)=B/\bar B\). Class number one gives \(B=(b)\). The rational element \(b/\bar b\) has norm one and exactly the same ideal as \(t\). Their quotient belongs to \(K_0\). This proves the decomposition without assuming a statement about signatures of real units. For the intersection, a norm-one integral unit has every complex absolute value one; Kronecker's theorem makes it a root of unity, and the roots of unity of \(E'\) are precisely \(\mu_{28}\).

The source action preserves the actual one-dimensional finite vacuum eigenspace, because norm-one integral units preserve its trace lattice and semicharacter. Therefore it has an actual smooth unitary character \(\lambda_f:K_0\to\mathbf C^\times\). This is not an unspecified phase used to identify a moved kernel: it is the character of the fixed finite representation on its named one-dimensional line. Put \(K=\ker\lambda_f\) and \(\Gamma=\mu_{28}\cap K\). Smoothness and compactness imply that \(K\) has finite index. Equation (R9.15) gives \[
 h_K=\#T'(F')\backslash T'(\mathbf A_{F',f})/K
     =\frac{[K_0:K]}{[\mu_{28}:\Gamma]}.
 \tag{R9.16}
\] Rational invariance of the nonzero theta value (R9.8) gives \(\beta(\zeta)\lambda_f(\zeta)=1\) for every \(\zeta\in\mu_{28}\). Thus \(\beta|_\Gamma=1\); in fact (R9.14) makes \(\lambda_f|_{\mu_{28}}=1\).

Normalize Haar measure on the compact torus quotient \(T'(F')\backslash T'(\mathbf A_{F'})\) to have volume one. Its \(h_K\) finite components have equal volume \(1/h_K\). Define the actual input \(f_K\) to be \(h_K\beta(t_\infty)\) on the identity finite component and zero on the other components, using rational triviality to transport that description. The condition \(\beta|_\Gamma=1\) makes this definition unambiguous. Its exact norms are \[
 \|f_K\|_2^2=h_K,\qquad \|f_K\|_\infty=h_K.
 \tag{R9.17}
\] The source action on \(\phi_P\) is \(\beta\) at infinity and is trivial on \(K\) at finite places. The target action commutes with the source action. Integrating over the one component supporting \(f_K\), with the convention \[
 \Theta_P(g)=\int_{[T']}\vartheta_{\phi_P}(t,g)
                                      \overline{f_K(t)}\,dt,
 \tag{R9.18}
\] therefore gives the equality of actual automorphic functions \[
 \boxed{\Theta_P(g)=\vartheta_{\phi_P}(1,g),\qquad
        \Theta_P(1)=\frac94 C^4.}
 \tag{R9.19}
\] For norm-one source input \(f_K/\sqrt{h_K}\), the output is \(h_K^{-1/2}\Theta_P\). This proves a nonzero first lift on the common host without any unproved claim that \(\lambda_f\) is trivial.

\textbf{The actual compact character is trivial.} At every odd unramified quadratic place the target identity plane is integrally hyperbolic, as in AM-DO-022. Pan (2001), section 6.2, Lemma and equation (6.2.a), printed p.~197, treats exactly a one-dimensional moving unitary space and a two-dimensional isotropic opposite space. Its comparison has \(\alpha=1\) and \(\beta=\chi((t-1)\delta)\). The opposite character is \(\chi_T^{-1}=1\), so both factors are one. The lattice-model action fixes the actual selfdual-lattice vacuum, as specified in Pan section 3.3 and the normalization of AM-DO-022. There is no assumption that the moving rank is even; this primary lemma is itself the rank-one case. The other splitting character \(\chi_S=\mu'\) affects the target action only. At relative split places, including all dyadic places, AM-DO-020 gives the same compact invariance with the retained trace-one convention, since its source determinant-root factor is \(\chi_T(x)^{-1}=1\). The identity element acts as one as well. These cases cover every finite place of \(E'/F'\). Consequently \[
 \boxed{\lambda_f=1,\quad K=K_0,\quad h_K=1,\quad
        \|f_K\|_2=\|f_K\|_\infty=1.}
 \tag{R9.27}
\] Thus (R9.19) is already the first lift with norm-one source input. The general-component calculation above records exactly how a nontrivial finite character would have entered, without discarding such a cost by convention.

\subsection{4. Exact Fock normalization}\label{exact-fock-normalization}

In the \(w^\pm\) coordinates the positive Hermitian form on \(A\) is \[
 H_A=\sum_{j=1}^3
       \frac{|w_j^+|^2+|w_j^-|^2}{2G_j}.
 \tag{R9.20}
\] The creation operator for \(w_2^+\) is one half of the creation operator for \(u_2\) minus \(i\) times that for \(v_2\). Orthogonality of the two modes and the normalization in ALP-LIFT-09 give its one-creation squared norm \(\pi/(2G_2)\). The actual Gaussian used in (R9.7), before applying creation operators, has squared norm \(\det(2Y)^{-2}\), because there are four ordinary Jacobian factors. Hence the exact squared norm of the actual archimedean vector in (R9.6) is \[
 \|\phi_{P,\infty}\|^2=\det(2Y)^{-2}(4!)^2
                                  \left(\frac{\pi}{2G_2}\right)^8.
 \tag{R9.21}
\] The unit archimedean vector therefore has identity value \[
 C^4\det(2Y)\frac{3}{32}\left(\frac{2G_2}{\pi}\right)^4.
 \tag{R9.22}
\] Combining unit archimedean vector and norm-one source input adds the factor \(h_K^{-1/2}\). The determinant factor is necessary: an earlier draft used a norm relative to unit Gaussian and omitted it when normalizing the actual theta Gaussian; the present formulas correct that error. These are analytic normalizations; they do not replace the algebraic coefficient \(9/4\) or the common period \(C^4\).

\subsection{5. A fixed character repair for the literal bilinear second edge}\label{a-fixed-character-repair-for-the-literal-bilinear-second-edge}

This paragraph uses the corrected source convention of AM-DO-023 and the literal bilinear theta integral in ALP-LIFT-11's amended version. Let \[
 c_\tau=2-\nu_\tau,\qquad
 \mathbf c=(2,1,2,2,1,2).
 \tag{R9.23}
\] A global norm-one character with this infinity type can be constructed explicitly from (R9.15). Choose a split place over 29 with \(\zeta_{28}\mapsto2\pmod{29}\). The element 2 generates \(\mathbf F_{29}^\times\). The infinity character \(\beta_{\mathbf c}\) has exponent \[
 \sum_{\tau\in\Phi'}c_\tau\tau=116\equiv4\pmod{28}
 \tag{R9.24}
\] on \(\mu_{28}\). Let its finite compact character be trivial elsewhere and take the value \(\exp(2\pi i\,24/28)\) on 2 at this place. Its restriction to \(\mu_{28}\) is exactly \(\beta_{\mathbf c}^{-1}\), so their product descends through (R9.15) to a unitary automorphic character \(\gamma\). Equivalently, for \(t_f=\delta k\), define \(\gamma(t_\infty,t_f)
=\beta_{\mathbf c}(\delta_\infty^{-1}t_\infty)\rho(k)\). Independence of the choice of \(\delta\) follows from the displayed equality on roots of unity. Its finite conductor divides this fixed place over 29; its compact values have order dividing 7.

Replace \(\chi_S\) in (R9.9) by \(\chi_S'=\mu'(\gamma\circ j)\), where \(j(x)=x/\bar x\); leave \(\chi_T=1\). The restriction (R9.10) is unchanged. The actual splitting ratio is \(\gamma(\det g)\) on the target. To check every place, at nonsplit finite places the local norm-one torus is its maximal compact, and \(\gamma\) is trivial there; at split finite places use the determinant-root formula of AM-DO-020; at real places the corrected Fock formula changes the target power by \(c_\tau\). Thus \[
 \Theta'_P(g)=\gamma(\det g)\Theta_P(g)
 \tag{R9.25}
\] as actual automorphic functions. All identity values and scalar automorphic norms are preserved exactly. The first target types become \[
 \pi'_\tau=\det^2\operatorname{Sym}^{d_\tau}\mathbf C^2,
 \tag{R9.26}
\] and every infinity exponent of \(\chi_S'\) is 3.

For the next positive Hom edge \(U(2)\to U(4)\) with both splitting characters 1, the source Fock type of partition \((d_\tau,0)\) is \(\det^{-2}\operatorname{Sym}^{d_\tau}(\mathbf C^2)^*\). It is exactly the contragredient of (R9.26), as required by the literal bilinear integral \(\int\vartheta\,f\). The degrees have not increased. This is local occurrence and exact first-image preservation, not a proof of the global second integral's nonvanishing.

The second edge now lies on precisely the E-prime host of AM-DO-022--024, so its finite moved-vacuum formula, positive direct-average sign, common period, and probability-Haar theta-line pairing are available without an extension-of-field comparison. A prescribed global second image, its finite input, its automorphic mass and its selected coefficient-lattice costs remain to be proved.

\subsection{6. Geometric location and limits}\label{geometric-location-and-limits}

The abelian variety is the ordinary product \(J^4\), with the extra \(E'\)-action specified in (R9.1)--(R9.3); there is no substitution of the conjugate E-action product excluded from AM-K018. The value is an actual jet of the trace-one Picard line, transferred through its reduced Poincare trivialization to a global theta integral. The integral source input is an automorphic character of the norm-one torus; the repair is tensoring an actual target representation by a rank-one character. The exact FGA/SGA locators for these constructions are those in GROTHENDIECK\_DICTIONARY.md sections 2--5 (FGA 232 Theorem 3.1; FGA 236 Corollaries 1.7 and 2.3; SGA 7 I VII 2.7 and 2.9; SGA 3 I Proposition 4.7.3).

This closes the previously separate-host obstruction for these particular first and candidate second edges. It does not instantiate the old elliptic VacIn, the even-rank starting coefficient module, the full six arithmetic identities, or the DP44 gain/cost data.

\subsection{References and exact inputs}\label{references-and-exact-inputs}

Masley, J. M. (1976). Solution of small class number problems for cyclotomic fields. \emph{Compositio Mathematica, 33}(2), 179--186. Main Theorem, pp.~179 and 184; the table was visually checked on PDF pages 2 and 7 (one-based). \url{https://numdam.org/item/CM_1976__33_2_179_0.pdf}

AM-R-004--006 (actual periods, metric, trace-one scalar); AM-K018 (ordinary-product source gauge); AM-DO-020, AM-DO-022--024 (fixed finite conventions and E-prime host); corrected AM-DO-023 (DO23.3--DO23.9); ALP-LIFT-09 (V7--V11 and V18--V38); ALP-LIFT-11 (explicit torus-character extension and bilinear repair, used with the newly proved E-prime torus decomposition).

Pan, S.-Y. (2001). Splittings of the metaplectic covers of some reductive dual pairs. Pacific Journal of Mathematics, 199(1), 163--226. Section6.2, Lemma and equation(6.2.a), p197; section3.3 for the lattice normalization. \url{https://msp.org/pjm/2001/199-1/pjm-v199-n1-p10-p.pdf}

\clearpage

\section{AM-R-010: the permitted first-image family on the common E-prime host}\label{am-r-010-the-permitted-first-image-family-on-the-common-e-prime-host}

Author: /root, GPT-6 Astra Pro, ALP\_MACHINE, 21 November 2026 UTC. This extends the actual first-image construction, not an iterated nonvanishing or a cost computation for the full source chain.

\subsection{1. Statement and exact finite index}\label{statement-and-exact-finite-index}

Use \(E'=\mathbf Q(\zeta_{28})\), \(F'=E'^+\), \(A=J^2\), and the exact polarization and standard splitting of AM-R-009. Let the positive target have an even rank \(n\), and the source be \(T'=U(1)\). Thus the underlying theta variety is \(A^n=J^{2n}\), and its fixed vacuum period is \(C^{2n}\). For the level parameter \(M=29^e\), \(e\ge1\), use the same odd norm element \[
 \alpha=(1+2\sqrt{-7})^e,\qquad \alpha\bar\alpha=M
 \tag{R10.1}
\] as ALP-LIFT-10 and DP59, now acting on \(A\) through \(\mathcal O_E\subset\mathcal O_{E'}\). This \(M\) is an arithmetic level, not the number of theta lifts in DP44.

Let \(K_M'\subset K_0=(\widehat{\mathcal O}_{E'})^1\) be principal congruence modulo \(29^e\) at all six split places of \(F'\) over 29, and full maximal compact elsewhere. The common finite stabilizer, finite class count, and its prime contents are \[
 \boxed{\quad
 h_M'=\#T'(F')\backslash T'(\mathbf A_{F',f})/K_M'
       =28^5\,29^{6(e-1)}.
 \quad}
 \tag{R10.2}
\] In particular \[
 v_2(h_M')=10,\quad v_7(h_M')=5,\quad
 v_{29}(h_M')=6(e-1),\quad v_\ell(h_M')=0
 \quad(\ell\notin\{2,7,29\}).
 \tag{R10.3}
\] There is no ramified-7 index-two correction on this host. Every permitted translated theta germ of DP59 on these ordinary Jacobian factors, and every polynomial jet of total degree at most \(D\), is an actual first theta image. For a single fixed CM weight the common input has squared source norm \(h_M'\). For a finite set \(\mathcal B\) of weights, the single input realizing their entire direct sum has squared norm \[
                         |\mathcal B|\,h_M'.
 \tag{R10.4}
\] All polynomials of degree at most \(D\) require at most \[
                         |\mathcal B|\le\binom{D+6}{6}.
 \tag{R10.5}
\] This count concerns six total CM degrees, not the number of matrix entries in the target.

The corresponding unit-input logarithmic norm cost is exactly \[
 \frac12\log(h_M'|\mathcal B|)
 =3\log M+\frac52\log28-3\log29
                        +\frac12\log|\mathcal B|.
 \tag{R10.6}
\] For \(M\) polynomially bounded in \(D\), it is \(O(\log(D+2))\). This assertion concerns the torus input and its normalization. The actual archimedean Fock norm, period, target mass, and later lift costs are separate and retained.

\subsection{2. Why the indicated subgroup fixes the actual finite family}\label{why-the-indicated-subgroup-fixes-the-actual-finite-family}

The underlying abelian variety and trace form are the ordinary product of the ones in ALP-LIFT-10. Its identity \[
 U_r^{M^2H}\bigl(s(Mu)\bigr)=(U_x^H s)(Mu),
 \qquad x=Mr,
 \tag{R10.7}
\] including the canonical finite-order lift sign, applies to each factor and hence to their tensor product. This is the actual equality of theta-group translations used in that proof, not a new formal substitution for a selected theta germ. The normalized derivative \(M^{-d}\) cancels the factor \(M^d\) from a homogeneous derivative of degree \(d\).

The allowed half-orbit has local representatives supported at 29 with denominators dividing \(29^e\) in one of the two isotropic factors at each relative split place. The same property remains true after the extension \(E'/E\): 29 splits completely in \(E'\), and the valuation of \(\alpha\) at a selected prime is repeated at the two primes above it. Conjugation exchanges the selected and unselected factors. Thus the local allowed representatives are still in a single isotropic position lattice with denominator \(29^e\).

For \(t\equiv1\pmod{29^e}\), the displacement \((t-1)x\) is integral in that same isotropic factor. Its Heisenberg pairing with \(x\) is zero. At the odd prime 29 the trace-one character on that integral factor is one. Heisenberg covariance therefore gives \[
 \omega_f(t)W_x\phi_{00}
       =W_{tx}\omega_f(t)\phi_{00}=W_x\phi_{00}.
 \tag{R10.8}
\] The initial finite vacuum is fixed by the entire source maximal compact. For target rank two this is AM-R-009, equation (R9.27). For any even target rank, decompose its integrally hyperbolic lattice at an odd unramified place into orthogonal hyperbolic planes and apply the same Pan section6.2 comparison to each plane. Every source factor uses the opposite character 1, so their product is one. At relative split places, including dyadic places, the same assertion follows from AM-DO-020. The character of the odd source affects the target action and does not enter this calculation.

Away from 29 the permitted translation is integral, and its Heisenberg line is the same original line with its specified scalar. The compact source therefore fixes it. Tensoring over entries and places proves that \(K_M'\) fixes the entire allowed finite Schwartz span, not only an individual chosen vector.

At each of the six places over 29 the unit quotient has order \(28\,29^{e-1}\), so \[
 [K_0:K_M']=(28\,29^{e-1})^6.
 \tag{R10.9}
\] The reduction of \(\mu_{28}\) at any one such place is injective, because 2 has order 28 modulo 29. Hence \(K_M'\cap\mu_{28}=1\). AM-R-009 proves \(T'(\mathbf A_f)=T'(F')K_0\). Dividing (R10.9) by 28 gives (R10.2), and the prime valuations give (R10.3).

\subsection{3. One actual source input for every retained weight}\label{one-actual-source-input-for-every-retained-weight}

Let \(\beta\) be an actual archimedean source-torus weight occurring among the chosen creation vectors. Define \(f_\beta\) to equal \(h_M'\beta\) on the identity finite component of \([T']/K_M'\), and zero on the other components. Here the component is a copy of the full compact real torus, since its rational stabilizer is \(K_M'\cap\mu_{28}=1\). There is consequently no extra rational-root restriction on \(\beta\).

The components have probability-Haar mass \(1/h_M'\). If \(\phi_\beta\) is any allowed finite Schwartz combination with this Fock weight, compact invariance, commutation of the two unitary actions, and character orthogonality give \[
 \int_{[T']}\vartheta_{\phi_\beta}(t,g)
                    \overline{f_\beta(t)}\,dt
       =\vartheta_{\phi_\beta}(1,g),\qquad
 \|f_\beta\|_2^2=h_M'.
 \tag{R10.10}
\] This is an equality of actual automorphic functions for every target \(g\), with the fixed splitting and source measure of AM-R-009.

For distinct weights \(\beta,\beta'\), the analogous cross integral is zero by compact real character orthogonality, and \(f_\beta,f_{\beta'}\) are orthogonal. Set \(f_{\mathcal B}=\sum_{\beta\in\mathcal B}f_\beta\). For \(\phi=\sum_{\beta\in\mathcal B}\phi_\beta\), \[
 \Theta(\phi,f_{\mathcal B})(g)
      =\sum_{\beta\in\mathcal B}\vartheta_{\phi_\beta}(1,g),
 \quad
 \|f_{\mathcal B}\|_2^2=|\mathcal B|h_M',
 \quad
 \|f_{\mathcal B}\|_\infty=|\mathcal B|h_M'.
 \tag{R10.11}
\] The supremum equality follows at the identity real torus point, where all characters equal one. The input depends only on the finite weight set, not on the coefficients or chosen translates of \(\phi\).

In the six eigen-coordinates of \(A\), a degree-\(d\) monomial has a total-degree vector in \(\mathbf N^6\) of sum \(d\). Its source character is this vector, with sign reversed and with the same vacuum shift. Different vectors give distinct characters of the full real torus. Counting all vectors of total degree at most \(D\) proves (R10.5). A polynomial in the original ordinary-Jacobian coordinates has such a finite decomposition because \(u_j=(w_j^++w_j^-)/2\) and \(v_j=(w_j^+-w_j^-)/(2i)\). Those transformations have only dyadic denominators; they are integral and invertible at 11. The decomposition is orthogonal in the actual Fock Hermitian metric. It introduces no additional archimedean norm beyond the norm of the original vector, although the source input cost (R10.4) must be retained.

\subsection{4. What the construction transfers, and what it does not}\label{what-the-construction-transfers-and-what-it-does-not}

Equations (R10.7) and (R10.10) transfer every specified permitted jet value on the ordinary product \(J^{2n}\) to an actual first image on the E-prime host, with the exact period \(C^{2n}\). AM-K018 identifies its full-\(\mathcal O_E\) source gauge with the actual Chern gauge on those ordinary powers. The extra \(i\)-action mixes two ordinary copies; it does not replace a factor by the conjugate E-action. Thus the comparison excluded for the latter product is not used here.

For example, an already established DP59 harmonic witness on two ordinary Jacobian factors can be tensored with two permitted spectator sections having unit zero-jet, as in ALP-LIFT-10 section7. Its value on \(A^2=J^4\) and its normalized 11-adic valuation are then unchanged. Decomposing its polynomial into the six CM weights and using (R10.11) realizes that exact value in the actual first image. For a single chosen weight, the further finite-character averaging has denominator \(h_M'\), which is an 11-adic unit. One must not assume that an individual weight summand has the same valuation unless its lower bound is also proved; the exact combined value already follows from (R10.11).

The fixed determinant repair of AM-R-009 depends only on the vacuum shifts, not on these six degrees. It therefore gives local bilinear compatibility with the next positive \(U(2)\to U(4)\) edge for every such first image when \(n=2\). That next edge is over the same E-prime host as AM-DO-022--024. No scalar, finite index, or coordinate denominator is discarded in making this common-host statement.

The present result does not choose the next finite Schwartz vector, compute its integral content or norm, identify the first images with the source's prescribed starting module, or verify the six arithmetic row identities. In particular (R10.6) is not the full \(e_j\) of DP44, and the copied derivative valuation is not automatically the determinant gain \(g_j\). These are precise remaining comparisons, not consequences of first-image existence.

\subsection{5. Framework and exact inputs}\label{framework-and-exact-inputs}

The finite part is a theta-group representation on the section line of a Picard scheme, transported along the actual odd isogeny and its isotropic kernel. Its six weights are characters of the diagonalizable CM torus, and its norm formula uses a finite compact quotient with probability measure. The corresponding exact framework locators are FGA232 Theorem3.1; FGA236 Corollaries1.7 and2.3; SGA7 I VII2.7/2.9; SGA3 I Proposition4.7.3, with the qualifications in GROTHENDIECK\_DICTIONARY.md.

The exact carried inputs are DP59 (permitted half-orbit and zero-jet spectators), AM-K018 (ordinary-product gauge), ALP-LIFT-10 (O1--O28 and section7), AM-R-009 (R9.1--R9.27), and AM-DO-020/022 (finite split and unramified normalizations). Pan (2001), section6.2, Lemma and equation(6.2.a), p197, supplies the rank-one compact comparison used in each hyperbolic plane. No Lean execution is claimed.

\clearpage

\section{ALP-LIFT-12 --- an actual second definite lift of every repaired first-image vector}\label{alp-lift-12-an-actual-second-definite-lift-of-every-repaired-first-image-vector}

Author: /root/lift\_accounting, GPT-6 Astra Pro, ALP\_MACHINE, 21 November 2026 UTC.

This result proves nonvanishing for an actual second classical theta correspondence. Its second Schwartz function is allowed to vary through the full oscillator representation. It does not identify that function with a prescribed source moved kernel, and it supplies no second-output integral content or DP44 height bound.

\subsection{1. Statement and integral conventions}\label{statement-and-integral-conventions}

Set \[
 E=\mathbf Q(\zeta_7),\qquad F=E^+,\qquad
 T=U(1),\quad G=U(I_2),\quad H=U(I_4).
 \tag{S1}
\] The positive Hermitian forms and the corresponding skew forms use the Hom convention already fixed in ALP-LIFT-09 and AM-DO-023. In particular both groups in (S1) are compact at every real place. Use the actual rationally compatible first oscillator and the globally unitary splitting characters of ALP-LIFT-09.

Let \(F_1\ne0\) be any finite sum of the actual first theta vectors constructed there, allowing finite Schwartz translations and finitely many holomorphic Fock polynomials as in ALP-LIFT-10. Equivalently, one may start with a finite sum of first theta lifts of unitary automorphic characters of \(T\) using these polynomial Fock vectors. The source torus input may be the finite projector of ALP-LIFT-10: its finite character decomposition is included in this definition.

Use the one fixed global character of ALP-LIFT-11, \[
 \gamma=\gamma_{(2,2,1)},\qquad
 F_1^+(g)=\gamma(\det g)F_1(g).
 \tag{S2}
\] This is an actual first theta vector for the repaired first splitting. It has the same identity value and the same scalar automorphic norm. At a real place with original polynomial degree \(d_i\), its compact type has highest weight \[
 (d_i+2,2),\qquad
 (\pi_i^+)^\vee=\det^{-2}\operatorname{Sym}^{d_i}(\mathbf C^2)^*.
 \tag{S3}
\] These assertions are ALP-LIFT-11, (R23)--(R26), not new character or real-weight claims in this note.

For the second edge take the actual standard oscillator for \(G\times H\), with both splitting characters equal to \(1\). The assertion is:

\textbf{Theorem.} There is a genuine global Schwartz function \(\psi\) in this second oscillator representation for which \[
 \Theta^{\mathrm{bil}}_{2,\psi}(F_1^+)(h)
   :=\int_{G(F)\backslash G(\mathbf A_F)}
                \vartheta_{2,\psi}(g,h)F_1^+(g)\,dg
 \quad\text{is not identically zero.}
 \tag{S4}
\] This holds for each specific nonzero vector \(F_1^+\), including a nonzero sum across automorphic multiplicity spaces. A target translate of \(\psi\) can make the resulting value at \(h=1\) nonzero.

Formula (S4) is bilinear, as is the literal source liftT\_apply. Gan--Qiu--Takeda use the conjugated-input convention \(\int\vartheta(g,h)\overline{f(g)}\,dg\). We therefore apply their formula to \(f=\overline{F_1^+}\), in the contragredient representation. This change is necessary; applying their compact-type criterion directly to \(F_1^+\) would reverse the relevant source type.

\subsection{2. Temperedness of the actual first constituents}\label{temperedness-of-the-actual-first-constituents}

First fix one unitary torus character contributing to \(F_1\) and one irreducible automorphic constituent of its first image. Projection to this constituent is an equivariant map from the actual first oscillator. Localizing that map shows that each local constituent occurs as a quotient in the corresponding local theta correspondence. This is the usual local equivariance of the theta integral: right translation in one local group can be moved onto the same local Schwartz factor, while integration in the source torus supplies its local character. No assertion about a formal free thetaG parameter is involved.

At a nonsplit finite place, the source is a unitary character of the compact group \(U(1)\). Its one-dimensional Langlands parameter has bounded Weil image. Gan and Ichino, Theorem 4.4(i)(a), (ii)(b), printed p.~14, give the parameter of any nonzero almost equal rank lift: \[
 \phi_{\mathrm{first}}
       =(\phi_T\otimes\chi_V^{-1}\chi_W)\oplus\chi_W .
 \tag{S5}
\] Every character in (S5) is unitary. Hence this parameter is tempered, and so is its local representation. The two cases of that theorem cover both the presence and the absence of the distinguished character in \(\phi_T\); only an already nonzero first image is being considered. Theorem 4.4(iii) also says that the nonzero big theta lift of this tempered source is irreducible. The bounded-parameter/tempered-representation equivalence is the same paper's section 2.

At a split finite place the pair is of type \(\mathrm{GL}_1\times\mathrm{GL}_2\). Mínguez (2008), Theorem 1(2), printed p.~718, says that the small lift has parameter \[
                  \phi_T^\vee\oplus 1
 \tag{S6}
\] in the standard type-II normalization. Changing the honest splitting or the source convention only introduces unitary characters and, where appropriate, a contragredient. Thus (S6) is again bounded. The uniqueness assertion in Theorem 1(1) ensures that an irreducible quotient arising from the actual global first theta integral has this small-theta parameter. This argument includes split dyadic places.

The twist in (S2) is unitary at every place. Taking a contragredient also preserves temperedness. Consequently every finite local constituent of the conjugated input for (S4) is tempered. At real places the group is compact, so its irreducible unitary representations are tempered. We have proved: \[
 \text{Every irreducible constituent of }\overline{F_1^+}
       \text{ is tempered at every place.}
 \tag{S7}
\] For a sum of torus characters, apply the same argument to each summand before decomposing into irreducible target constituents. It does not require that the original projector be one character.

\subsection{3. Every finite local constituent occurs on the second edge}\label{every-finite-local-constituent-occurs-on-the-second-edge}

\subsubsection{3.1. The four-dimensional target form is hyperbolic locally}\label{the-four-dimensional-target-form-is-hyperbolic-locally}

At a nonsplit unramified finite place, the norm map on units of the unramified quadratic extension is surjective. In particular \(-1=N(z)\) for a unit \(z\), and \((1,z)\) is isotropic in \(I_2\). Thus \(I_2\) is hyperbolic and \(I_4\) is the sum of two hyperbolic planes.

For \(E/F\) in (S1), the only relative ramified finite place is above \(7\). At that place choose \(b\in\mathbf Q_7\) with \(b^2=-5\) and \(b\equiv3\pmod7\); Hensel's lemma applies because \(3^2+5\equiv0\pmod7\) and the derivative \(6\) is a unit. In \(I_4\) the two vectors \[
                 u=(1,2,b,0),\qquad v=(0,-b,2,1)
 \tag{S8}
\] are linearly independent and have zero self-pairing and zero mutual pairing. They span a two-dimensional isotropic subspace. Since \(I_4\) is nondegenerate, it is hyperbolic. This proves the required target assertion at every nonsplit finite place without an integral classification assertion. At a split place the second pair is simply \(\mathrm{GL}_2\times\mathrm{GL}_4\).

\subsubsection{3.2. A direct quotient at nonsplit places}\label{a-direct-quotient-at-nonsplit-places}

Let \(k=F_v\), \(L=E_v\) be a nonsplit finite extension, and write the hyperbolic target as \(X\oplus X^*\), with \(\dim_LX=2\). The second Schrödinger model is \[
                    \mathcal S(X\otimes_L W)
                       \simeq\mathcal S(M_2(L)),
 \tag{S9}
\] where \(W\) is the two-dimensional source space. The source group preserves this position space. Its action is right or inverse-right change of variables, depending on the matrix identification, times a smooth unitary character. The absolute determinant multiplier is one on \(U(W)\). Either right-action convention gives the following same quotient conclusion after inversion.

The orbit of the identity matrix is \(U(W)(k)\), and it is closed in \(M_2(L)\). Indeed it is cut out by the equation \(g^*Jg=J\) for a nonsingular matrix \(J\); that equation already forces \(g\) to be invertible. Restriction therefore gives an equivariant surjection \[
 \mathcal S(M_2(L))\longrightarrow C_c^\infty(U(W)(k)),
 \tag{S10}
\] with the same character twist.

Here is the surjectivity detail. If \(C\) is a compact open subset of the closed orbit, its complement in the orbit is closed in the ambient space. Each point of \(C\) has a compact open ambient neighborhood avoiding that complement. Finitely many such neighborhoods cover \(C\); their union intersects the orbit in exactly \(C\). A compactly supported locally constant function on the orbit has a finite partition into such sets. Extending their characteristic functions, and disjointifying the finitely many ambient neighborhoods, extends that function to an element of the ambient Schwartz space.

Every irreducible smooth representation \(\pi_v\) is a quotient of the right regular representation \(C_c^\infty(U(W)(k))\). For a nonzero smooth vector \(a\in\pi_v\), the map \[
 f\longmapsto\int_{U(W)(k)} f(g)\pi_v(g^{-1})a\,dg
 \tag{S11}
\] is equivariant for right translation. Taking the normalized indicator of a compact open subgroup fixing \(a\) shows that the map is nonzero. Irreducibility makes it surjective. The group is unimodular, so the substitution in this equivariance calculation introduces no modular character. A character twist of the regular representation has the same quotient property.

Equations (S10) and (S11) prove \[
                  \operatorname{Hom}_{G(k)}
                          (\omega_{2,v},\pi_v)\ne0
 \tag{S12}
\] for every irreducible smooth \(\pi_v\), in particular every constituent in (S7). This is a quotient argument, not an inference from a subquotient of an open-orbit submodule.

\subsubsection{3.3. Split places}\label{split-places}

At a split place use Mínguez, Theorem 1(2), with \(n=2,m=4\). It states directly that \(\operatorname{Hom}_{\mathrm{GL}_2}(\omega_{2,4},\pi_v)\ne0\) for every irreducible smooth \(\pi_v\), over a nonarchimedean field of any residue characteristic. Unitary splitting characters or a group automorphism do not change this all-representations assertion.

The nonsplit closed-orbit proof must not be copied verbatim here: the full-rank \(\mathrm{GL}_2\) orbit is not closed in the whole matrix space. Its regular quotient from an open submodule alone would not establish a quotient of the full oscillator. The cited type-II theorem closes precisely that point.

\subsection{4. The real local doubling factor is explicitly nonzero}\label{the-real-local-doubling-factor-is-explicitly-nonzero}

By (S3), each real component of the conjugated input has type \[
 \pi_i=\det^{-2}\operatorname{Sym}^{d_i}(\mathbf C^2)^*,
                     \qquad \dim\pi_i=d_i+1.
 \tag{S13}
\] AM-DO-023 and ALP-LIFT-11 show that this exact type occurs in the second positive Hom oscillator with trivial splitting characters. Choose an isometric \(U(2)\)-equivariant embedding \(\iota:\pi_i\hookrightarrow\omega_{2,i}\), a unit vector \(a\), and \(\phi=\iota(a)\). Such a vector is a polynomial times a Gaussian and is an actual Schwartz vector.

For the doubled oscillator, Gan--Qiu--Takeda, section 11.3, printed p.~52, give the equivariant map \(\delta\) normalized by evaluation at zero: \[
              \delta(\phi_1\otimes\overline{\phi_2})(0)
                         =\langle\phi_1,\phi_2\rangle .
 \tag{S14}
\] Their possible determinant character is trivial here because the second target splitting character is \(1\). Consequently its Siegel--Weil section at \(s_0=1\) satisfies \[
 \Phi_\delta(g,1;1)
       =\langle\omega_{2,i}(g)\phi,\phi\rangle
       =\langle\pi_i(g)a,a\rangle .
 \tag{S15}
\] Their local zeta integral in section 11.6, printed p.~54, uses the complex conjugate of the input matrix coefficient. The bars in these two formulas have been checked on rendered PDF pages; the plain-text extraction drops them. Thus Schur orthogonality gives the exact value \[
 \begin{split}
 Z_i(1,\Phi_\delta,a,a)
    &=\int_{U(2)}
          |\langle\pi_i(g)a,a\rangle|^2\,dg\\
    &=\frac{\operatorname{vol}(U(2))}{d_i+1}>0.
 \end{split}
 \tag{S16}
\] In the GQT normalization this yields \[
 Z_i^*(1,\Phi_\delta,a,a)
   =\frac{\operatorname{vol}(U(2))}
      {(d_i+1)L_i(3/2,\pi_i)}\ne0 .
 \tag{S17}
\] The standard local factor of a tempered representation is holomorphic and nonzero in \(\operatorname{Re}s>0\); the precise doubling-factor version is Lapid--Rallis, Proposition 5, also invoked in Gan--Ichino, section 6.7, printed p.~27. For these compact real groups the factors are the usual nonvanishing gamma factors.

This direct compact calculation is needed. GQT Theorem 11.7's general converse does not cover arbitrary definite real unitary groups at \(m=2n\). We are not applying that converse without its extra hypotheses; (S16) proves its required local nonvanishing in the present actual compact types.

\subsection{\texorpdfstring{5. The global \(L\)-value and the Rallis hypotheses}{5. The global L-value and the Rallis hypotheses}}\label{the-global-l-value-and-the-rallis-hypotheses}

Let \(\pi\) be an actual irreducible cuspidal subrepresentation of the space generated by \(\overline{F_1^+}\), with its actual embedding in automorphic functions. By (S7), its unramified Satake parameters have absolute value one. Hence the partial standard Euler product at \(s=3/2\) is absolutely convergent and nonzero: the logarithms are bounded by a fixed-degree multiple of \[
       \sum_{\mathfrak p}\sum_{r\ge1}
              \frac{N\mathfrak p^{-3r/2}}r<\infty.
 \tag{S18}
\] At a split place this can be viewed as the product over the two primes of \(E\). Enlarging the finite exceptional set has no effect on this conclusion. The finitely many remaining standard local factors are holomorphic and nonzero at \(3/2\) by the same tempered doubling-factor property used in (S17). Thus the completed factor in GQT is \[
                           L(3/2,\pi)\ne0 .
 \tag{S19}
\] There is no central-value nonvanishing conjecture in (S19). In particular one must not substitute \(L(1/2,\pi)\) for this evaluation point.

For clarity, one can avoid any ramified local-LLC factor identification in the last paragraph. GQT Theorem 11.4(ii) itself gives holomorphy of the completed factor once all local theta lifts are nonzero, which has been proved in sections 3--4. Its local factors have no zeros; its partial unramified product is nonzero by (S18). These facts also give (S19). The tempered local-factor property is a direct stronger check.

The target \(I_4\) is positive definite at every real place, so it is anisotropic over \(F\). Its global Witt index is \(r=0\), even though its finite local Witt indices are two. The parameters of GQT Theorem 11.4 are exactly \[
 \epsilon_0=0,\quad n=2,\quad m=4,\quad d(n)=2,\quad
 r=0,\quad s_{m,n}=(m-n)/2=1.
 \tag{S20}
\] Thus \(d(n)<m\le2d(n)\), \(r\le n\), and \(0<s_{m,n}\le d(n)/2\). The lower-tower vanishing assumption for \(j<r\) is empty. Both automorphic quotients are compact; all their automorphic forms are cuspidal because the groups have no proper rational parabolic subgroups.

At every finite place, (S12) and the split theorem give local theta occurrence. GQT Proposition 11.6(i), printed p.~55, then gives the nonzero normalized local zeta functional. At every real place the same conclusion follows from (S17). GQT Theorem 11.4(ii), printed p.~54, and (S19) therefore imply that the second global theta lift of this actual cuspidal embedding of \(\pi\) is nonzero.

With the paper's Haar and Petersson normalizations, its identity is \[
 \langle\Theta_{2,\psi_1}(f_1),
                    \Theta_{2,\psi_2}(f_2)\rangle
       =2\,L(3/2,\pi)
          \prod_v Z_v^*(1,\psi_{1,v}\otimes
                    \overline{\psi_{2,v}},f_{1,v},f_{2,v}).
 \tag{S21}
\] Replacing these measures by probability measures introduces the corresponding positive quotient-volume factors. No such factor has been silently assigned the value one in a height or norm budget. The archimedean factors (S16) are explicit; the finite test functions supplied by local nonvanishing have not been optimized or bounded arithmetically.

\subsection{6. Why every specified first vector survives, including multiplicities}\label{why-every-specified-first-vector-survives-including-multiplicities}

It remains to pass from constituent nonvanishing to the assertion about each specified \(F_1^+\). Let \(\mathcal A\) be the closed \(G(\mathbf A_F)\)-submodule generated by the finite collection of first vectors in question, after (S2). Define the common kernel \[
 \mathcal N=\{F\in\mathcal A:
          \Theta^{\mathrm{bil}}_{2,\psi}(F)=0
                         \text{ for every }\psi\}.
 \tag{S22}
\] Equivalently, intersect the kernels of the scalar functionals \[
 F\longmapsto\int_{[G]}
                    \vartheta_{2,\psi}(g,h)F(g)\,dg
 \quad(\psi,\ h\text{ arbitrary}).
\] They are continuous on \(L^2([G])\): the quotient is compact and a fixed theta kernel is a smooth square-integrable function there. Thus \(\mathcal N\) is closed. Changing variables in the integral moves right translation of \(F\) into the source action on \(\psi\), so \(\mathcal N\) is \(G(\mathbf A_F)\)-invariant.

The compact anisotropic automorphic quotient has discrete unitary spectrum, with finite multiplicities. A nonzero closed invariant subspace of such a Hilbert direct sum contains an irreducible automorphic subrepresentation. This statement allows a subspace inside a multiplicity space: it does not assume multiplicity one. If \(\mathcal N\ne0\), choose such an irreducible copy \(\pi^+\subset\mathcal N\). Its conjugate copy \(\overline{\pi^+}\) is an actual cuspidal embedding to which sections 2--5 apply.

The GQT identity applies to this particular embedding, with its induced Petersson form. Its factorization can always be obtained by scaling one local invariant pairing; that nonzero positive scale cannot change the conclusion. It gives some \(f\in\overline{\pi^+}\) and \(\psi\) for which \(\int\vartheta_{2,\psi}(g,h)\overline{f(g)}\,dg\) is nonzero. But \(\overline f\in\pi^+\subset\mathcal N\), a contradiction. Therefore \(\mathcal N=0\), proving (S4) for each nonzero specified first vector, not only for one favorable spectral projection.

Finally, if its output is nonzero at \(h_0\), replace \(\psi\) by its target translate \(\omega_2(h_0)\psi\). The value at the target identity becomes that previous value at \(h_0\). This does not compute the arithmetic height or denominators of the translated second Schwartz function.

\subsection{7. Precisely stated transfer to another CM host}\label{precisely-stated-transfer-to-another-cm-host}

The argument is structural over a CM extension \(K/K^+\). It applies with the same ranks if the following inputs have actually been supplied:

\begin{enumerate}
\def\labelenumi{\arabic{enumi}.}
\tightlist
\item
  An actual first \(U(1)\to U(I_2)\) theta vector from unitary characters and polynomial Fock vectors, with globally unitary splitting characters.
\item
  A globally admissible fixed character twist making each real first type \(\det^2\operatorname{Sym}^{d_\tau}\), with the same Hom and bilinear integral conventions.
\item
  The positive form \(I_4\) is hyperbolic at every nonsplit finite place.
\end{enumerate}

No torus class-number calculation enters the second-step proof once input 1 exists. If \(K/K^+\) is unramified at every finite nonsplit place, input 3 follows immediately from the norm argument in section 3.1. AM-R-009 supplies the intended inputs for \(K=\mathbf Q(\zeta_{28})\): its fixed twist has \(\mathbf c=(2,1,2,2,1,2)\), and its degree-eight vector at label9 has the actual first identity value \((9/4)C^4\). These inputs have now been independently whole-checked in review\_root\_unramified\_first\_lift.md, retaining the accepted DP57 and DO022 host facts. Thus there is the concrete corollary \[
 \exists\psi:\quad
 \Theta^{\mathrm{bil}}_{U(2)\to U(4),\psi}
       (\Theta'_P)\ne0,\qquad
 \Theta'_P(1)=\frac94C^4,
 \tag{S23}
\] over this same \(E'\), with the actual first source input of norm one from (R9.27). At infinity the degrees are \(d_9=8\) and zero at the other five places, so the unnormalized compact numerator in (S16), multiplied over all six places, is \(\prod_{\tau}\operatorname{vol}(U(2)_\tau)/9\). This does not normalize the finite second test functions. The field and lattice normalizations in this corollary are the stated \(E'\) data, not an identification of \(E'\) with the original \(E\) in (S1).

\subsection{8. Remaining gate and costs}\label{remaining-gate-and-costs}

This proves a genuine nonzero two-edge composition with an unrestricted actual second Schwartz input. For the sharp first coefficients of ALP-LIFT-10, it proves that the same nonzero first vector is not annihilated by every second kernel. It does not prove that its second value has the first coefficient's valuation.

The next arithmetic gate is to choose that second kernel inside the permitted finite moved family, prove its integral lattice and content, and bound the actual mass and archimedean norm of the nonzero pairing. The full common-character pins, stage row modules, harmonic branching maps and six terminal identities are also absent from this two-edge result. In particular the rank-one starting edge is still not the source's even-rank StdPair constructor. No shorter DP44 chain or net DP44 gain is claimed from (S4) or (S21).

\subsection{Primary references and exact locators}\label{primary-references-and-exact-locators}

\begin{itemize}
\tightlist
\item
  Gan, W. T., Qiu, Y., \& Takeda, S. (2014). The regularized Siegel--Weil formula (the second term identity) and the Rallis inner product formula. arXiv:1207.4709v3. Section11.2--11.3, printed pp.51--52, the conjugated-input integral and doubled Schwartz map; section11.6--11.7, printed p.54, normalized local factors and Theorem11.4; Proposition11.6(i), printed p.55; Theorem11.7, pp.56--57, the real-place caveat deliberately bypassed. \url{https://arxiv.org/pdf/1207.4709v3}
\item
  Gan, W. T., \& Ichino, A. (2015). The Gross--Prasad conjecture and local theta correspondence. arXiv:1409.6824v2. Section2, bounded parameters and tempered representations; Theorem4.4(i)(a), (ii)(b), (iii), printed p.14, the almost equal rank parameter and irreducibility; section6.7, printed p.27, the exact tempered local doubling-factor property of Lapid--Rallis Proposition5. \url{https://arxiv.org/pdf/1409.6824}
\item
  Mínguez, A. (2008). Correspondance de Howe explicite: paires duales de type II. Annales scientifiques de l'École Normale Supérieure, 41(5), 717--741. Theorem1(1)--(2), printed p.718; Corollary6.3 is its later proved version. The theorem includes every nonarchimedean residue characteristic. \url{https://www.numdam.org/item/ASENS_2008_4_41_5_717_0.pdf}
\item
  Lapid, E. M., \& Rallis, S. (2005). On the local factors of representations of classical groups. In Automorphic representations, L-functions and applications: Progress and prospects (pp.309--360). De Gruyter. Proposition5, used in exactly the form explicitly recorded in Gan--Ichino section6.7. The primary Gan--Ichino passage was directly read; no independent full reading of the Lapid--Rallis chapter is claimed. \url{https://doi.org/10.1515/9783110892703.309}
\end{itemize}

Internal dependencies: ALP-LIFT-09, (V22)--(V38); ALP-LIFT-10, actual finite moved family and source inputs; ALP-LIFT-11, (R23)--(R26); AM-DO-023, corrected real source action and the trivial-character second case. Check events are recorded separately, only after whole proof readings. This proof has not been formalized in Lean.

\clearpage

\section{AM-DO-025: the actual common-character compact chain and its exact degrees}\label{am-do-025-the-actual-common-character-compact-chain-and-its-exact-degrees}

Author: GPT-6 Astra Pro, ALP\_MACHINE \nolinkurl{/root/definite_oscillator}; inherited reasoning effort not separately exposed; 21 November 2026 UTC.

The selected increasing even-rank chain in the recovered source has an exact local compact-type match. The failure of the unchanged trivial-character second edge in AM-DO-023 is therefore not an obstruction to this selected chain. The computation below uses the source's actual character exponents, bilinear theta convention, and coefficient weights. It also computes the polynomial degree required at each edge. It does not prove global nonvanishing or the arithmetic row identities.

\subsection{1. Source data and conventions}\label{source-data-and-conventions}

Fix one selected complex embedding and an actual \nolinkurl{HCParam} for ranks

\[
 d=t+1,\qquad N_j=k+jd,\qquad 0\le j\le m.
 \tag{DO25.1}
\]

Here \(k,m\) are even, \(t\) is odd, and the other source rank conditions are retained. The real numbers \(e_j,b_j\) are integers, the \(u_i\) are half-integers, and

\[
 u_{k-1}>-b_0,\qquad
 e_j-b_j\ge\frac{d+2}{2},\qquad
 b_{j+1}-e_j\ge\frac{d+2}{2}.
 \tag{DO25.2}
\]

These are \nolinkurl{Params.lean}, \nolinkurl{HCParam}, lines 50--69, with the decreasing \(u_i\) forced by its interlacing inequalities. The actual character on edge \(j\to j+1\) has infinity exponent \(-e_j\): \nolinkurl{ChainParSucc.lean}, \nolinkurl{Com.Is.exp}, and \nolinkurl{DataTheta.lean}, \nolinkurl{ThetaX.chrOf_inf}. This means \(\mu_{j+1}(z)=(z/\bar z)^{-e_j}\).

Use \textbf{equal source labels} \(\chi_S=\chi_T=\mu_{j+1}\) on this definite edge, exactly as required by \nolinkurl{PinColT.om} and \nolinkurl{PinX.mv} in \nolinkurl{LevChiEqOne.lean}. In conventional Konno--Konno parameters the pair is \((\mu_{j+1},\mu_{j+1}^{-1})\). The inverse is necessary because the recovered \nolinkurl{splitFacO} is \(\chi_T(x)^{-1}\chi_S(y)\) on determinant roots. This conversion is checked and corrected explicitly in AM-DO-023.

For source rank \(a=N_j\) and target rank \(b=N_{j+1}\), the oscillator's holomorphic Fock module is consequently

\[
 \det_a^{e_j-b/2}\boxtimes\det_b^{-e_j+a/2}
       \otimes\operatorname{Sym}(\mathbf C^b\otimes(\mathbf C^a)^*).
 \tag{DO25.3}
\]

The source theta operator is bilinear:

\[
                  (T_\phi f)(h)=\int_{[U(a)]}
                          \vartheta_\phi(g,h)f(g)\,dg.
 \tag{DO25.4}
\]

This is literally \nolinkurl{liftT_apply} in \nolinkurl{Development/AbelianVarietyTheory/Definitions/TwistedInnerProducts/WittSmallGroupEmbeddingMemPointwiseStabilizerIso.lean}, lines 160--167, and \nolinkurl{KerDat.lift} in \nolinkurl{JointSmooth.lean}, lines 79--82. Thus the source oscillator must contain the \textbf{contragredient} of the input coefficient type. Confusing (DO25.4) with an integral against \(\bar f\) reverses this type test.

\subsection{2. Identification with the actual selected coefficient partitions}\label{identification-with-the-actual-selected-coefficient-partitions}

Write \(w_j=(w_{j,0},\ldots,w_{j,N_j-1})\) for the source's dominant integral weight \nolinkurl{Wts.lamTrue u j τ}. Write \(\nu_j\) for \nolinkurl{Wts.nuW u (j+1) τ}. The recovered source already proves the precise identities

\[
 \nu_{j,i}=w_{j,i}+e_j-\frac{N_{j+1}}2,
 \tag{DO25.5}
\]

and

\[
 w_{j+1}=
 \left(w_j-\frac d2\mathbf1_{N_j},\,
       \left(-e_j+\frac{N_j}2\right)\mathbf1_d\right).
 \tag{DO25.6}
\]

Exact locators are \nolinkurl{Formalism/LocalToGlobal/Columns/AbelianHostFieldConditions.lean}, \nolinkurl{Wts.lamN}, \nolinkurl{lamTrue}, and \nolinkurl{nuW}, lines 462--469; and \nolinkurl{Formalism/LocalToGlobal/Columns/BaseLemmas.lean}, \nolinkurl{lamTrue_old}, lines 496--510, \nolinkurl{lamTrue_new}, lines 512--530, and \nolinkurl{nuW_eq_true}, lines 542--552. These source identities are used as settled equalities, rather than replaced by freely chosen weights.

The \(\nu_j\) are actual nonnegative partitions, as proved in §3 below. The Schur summand indexed by \(\nu_j\) in (DO25.3) has source type

\[
 \det_a^{e_j-b/2}\otimes S_{\nu_j}((\mathbf C^a)^*),
 \tag{DO25.7}
\]

whose highest weight, in decreasing order, is

\[
 \left(e_j-b/2-\nu_{j,a-1},\ldots,
                    e_j-b/2-\nu_{j,0}\right)
               =(-w_{j,a-1},\ldots,-w_{j,0}).
 \tag{DO25.8}
\]

This is precisely the contragredient required by (DO25.4). The target type of the same Schur summand is

\[
 \det_b^{-e_j+a/2}\otimes S_{\nu_j}(\mathbf C^b).
 \tag{DO25.9}
\]

Its highest weight consists of \(-e_j+a/2+\nu_{j,i}\) for the first \(a\) entries and \(-e_j+a/2\) on the final \(d\) entries. Substitution of (DO25.5) gives exactly (DO25.6). Therefore every selected adjacent pair is a matching local compact theta type with the \textbf{same prescribed character on both members of the edge}.

The Schur summand occurs with multiplicity one under the compact product group. This is the polynomial Cauchy decomposition \(\operatorname{Sym}(T\otimes S^*)=
\bigoplus_\nu S_\nu(T)\otimes S_\nu(S)^*\). The primary oscillator reference is Konno and Konno (2007), §5.1, Fact 5.1, Lemma 5.2, the highest vectors (5.4)--(5.5), and Lemma 5.3, printed pp.72--74, with the character convention fixed in AM-DO-023. Since \(b>a\), every partition of length at most \(a\) appearing above is permitted. Its total homogeneous degree is \(|\nu_j|\). No determinant twist external to the selected source data has been introduced.

This is a representation-theoretic occurrence statement. An equivariant global theta map can still vanish on an occurring local type; the computation does not assert otherwise.

\subsection{3. Exact recurrence and gap consequences}\label{exact-recurrence-and-gap-consequences}

Define

\[
       \Delta_j=e_{j+1}-e_j-d\qquad(0\le j<m-1).
 \tag{DO25.10}
\]

Then the selected polynomial partitions satisfy the exact recurrence

\[
 \boxed{\nu_{j+1}
       =(\nu_j+\Delta_j\mathbf1_{N_j},\,
                         \Delta_j\mathbf1_d).}
 \tag{DO25.11}
\]

Moreover

\[
                  \nu_{0,k-1}\ge1,\qquad \Delta_j\ge2.
 \tag{DO25.12}
\]

In particular all displayed weights \(\nu_j\) are nonnegative partitions, with strictly positive final entry.

\textbf{Proof.} Apply (DO25.5) at \(j+1\), and insert (DO25.6). For an old entry the change from \(\nu_{j,i}\) is

\[
 -\frac d2+e_{j+1}-\frac{N_{j+2}}2-e_j
                                      +\frac{N_{j+1}}2
      =e_{j+1}-e_j-d=\Delta_j.
\]

For a new entry, its value is \(-e_j+N_j/2+e_{j+1}-N_{j+2}/2=\Delta_j\). This proves (DO25.11).

The two gaps in (DO25.2) add to \(e_{j+1}-e_j\ge d+2\), so \(\Delta_j\ge2\). At stage zero the source weight is \(w_{0,i}=u_i-\rho_k(i)\), where \(\rho_k(i)=(k-1)/2-i\). This is \nolinkurl{lamN_eq} and \nolinkurl{lamTrue_zero} in \nolinkurl{BaseLemmas.lean}, together with \nolinkurl{rhoVec} in \nolinkurl{DataCoeff.lean}, line 39. Hence

\[
 \nu_{0,k-1}=u_{k-1}+e_0-\frac{d+1}{2}.
 \tag{DO25.13}
\]

Because \(u_{k-1}\) is a half-integer, \(b_0\) is an integer, and \(u_{k-1}>-b_0\), one has \(u_{k-1}+b_0\ge1/2\). Adding \(e_0-b_0\ge(d+2)/2\) proves (DO25.13) is at least one. The starting weight is dominant, so all its entries are at least one. Recurrence (DO25.11) adds a nonnegative constant to every old entry and appends entries no greater than those old entries, proving dominance and positivity inductively. All entries are integers by the parity and integrality conditions. This proves (DO25.12) and the partition assertion. \(\square\)

\subsection{4. Exact degree evolution and its sharp rank-only lower bound}\label{exact-degree-evolution-and-its-sharp-rank-only-lower-bound}

Let \(D_j=|\nu_j|\), the homogeneous degree at edge \(j\to j+1\) for this embedding. Then

\[
 \boxed{D_0=\sum_{i=0}^{k-1}u_i
                   +k e_0-\frac{k(k+d)}2,\qquad
        D_{j+1}=D_j+N_{j+1}\Delta_j.}
 \tag{DO25.14}
\]

Therefore

\[
 \boxed{D_j\ge k(2j+1)+d\,j(j+1)}
 \tag{DO25.15}
\]

and across the \(m\) selected edges,

\[
\begin{split}
 \sum_{j=0}^{m-1}D_j
   &=mD_0+\sum_{r=0}^{m-2}(m-1-r)N_{r+1}\Delta_r,\\
 \sum_{j=0}^{m-1}D_j
   &\ge k m^2+\frac d3(m^3-m).
\end{split}
 \tag{DO25.16}
\]

\textbf{Proof.} Summing (DO25.5) at \(j=0\), the entries of \(\rho_k\) sum to zero, giving the first equality in (DO25.14). Summing (DO25.11) gives the second. Since \(D_0\ge k\), \(\Delta_r\ge2\), and \(\sum_{\ell=1}^{j}N_\ell=jk+d j(j+1)/2\), induction gives (DO25.15). Summing the exact recurrence counts the increment at index \(r\) exactly \(m-1-r\) times and gives the first line of (DO25.16). Finally \(\sum_{j=0}^{m-1}(2j+1)=m^2\) and \(\sum_{j=0}^{m-1}j(j+1)=(m^3-m)/3\), proving its second line. \(\square\)

The lower bounds are attained among the displayed \nolinkurl{HCParam} conditions. Indeed choose the \(u_i\) in decreasing consecutive half-integers, take \(u_{k-1}=-b_0+1/2\), set \(e_0-b_0=(d+2)/2\), and take both successive gaps in (DO25.2) equal to \((d+2)/2\). The interlacing integer \(v_i\) is then the unique midpoint between \(u_i\) and \(u_{i+1}\). These choices have \(\nu_0=\mathbf1_k\) and \(\Delta_j=2\), and every entry is of the required parity. The partitions become

\[
 \nu_j=((2j+1)^k,(2j)^d,(2j-2)^d,\ldots,2^d),
 \tag{DO25.17}
\]

where exponents denote repeated parts and the last list is empty for \(j=0\). The source's further \nolinkurl{Hyp3} noncompact-gap conditions may force larger spacings. The sharpness statement here concerns the stated \nolinkurl{HCParam} constraints, not a claimed sharp construction for all of \nolinkurl{Hyp3}.

The degree in (DO25.14) is the degree of the selected Schur summand, not an arbitrary direction used to illustrate a formal-group slope. It is a usable exact input to subsequent derivative, norm, and denominator estimates. It does not itself measure a net arithmetic gain: the finite family, integral row lattice, normalization periods, class masses, and terminal detector remain to be inserted. On the \(E'\) realization, the dyadic eigenbasis change from AM-DO-022 still has index \(2^3\) and degree-\(D\) denominator bound \(2^D\); no type calculation removes that basis cost.

\subsection{5. Exact source locators and retained gates}\label{exact-source-locators-and-retained-gates}

All paths below are relative to the recovered source root \nolinkurl{work/research/alpoge_source}.

\begin{longtable}[]{@{}
  >{\raggedright\arraybackslash}p{(\columnwidth - 2\tabcolsep) * \real{0.5000}}
  >{\raggedright\arraybackslash}p{(\columnwidth - 2\tabcolsep) * \real{0.5000}}@{}}
\toprule\noalign{}
\begin{minipage}[b]{\linewidth}\raggedright
Statement used
\end{minipage} & \begin{minipage}[b]{\linewidth}\raggedright
Exact source
\end{minipage} \\
\midrule\noalign{}
\endhead
\bottomrule\noalign{}
\endlastfoot
Ranks and gap hypotheses & \nolinkurl{Formalism/LocalRingCardinality/Params.lean}, \nolinkurl{Ranks.N}, lines 35--36; \nolinkurl{HCParam}, lines 50--69; \nolinkurl{sExp}, \nolinkurl{sExp'}, \nolinkurl{wG}, lines 71--75 \\
Compact Harish--Chandra parameter & \nolinkurl{Formalism/LocalRingCardinality/DataTheta.lean}, \nolinkurl{TPar.Xpar}, lines 545--548 \\
Character infinity type & \nolinkurl{DataTheta.lean}, \nolinkurl{ThetaX.chrOf_inf}; \nolinkurl{Development/AbelianVarietyTheory/Definitions/LatticeFrames/ChainParSucc.lean}, \nolinkurl{Com.Is.exp}, line 40 \\
Coefficient definitions & \nolinkurl{Formalism/LocalToGlobal/Columns/AbelianHostFieldConditions.lean}, \nolinkurl{Wts.lamTrue}, \nolinkurl{Wts.nuW}, lines 464--469 \\
Exact selected recurrence & \nolinkurl{Formalism/LocalToGlobal/Columns/BaseLemmas.lean}, \nolinkurl{lamTrue_old}, \nolinkurl{lamTrue_new}, \nolinkurl{nuW_eq_true}, lines 496--552 \\
Bilinear upward kernel & \nolinkurl{Development/AbelianVarietyTheory/Definitions/TwistedInnerProducts/WittSmallGroupEmbeddingMemPointwiseStabilizerIso.lean}, \nolinkurl{liftT_apply}, lines 160--167 \\
Equal characters on the definite edge & \nolinkurl{Development/AbelianVarietyTheory/Definitions/LatticeFrames/LevChiEqOne.lean}, \nolinkurl{PinX.mv}, lines 100--101, and \nolinkurl{PinColT.om}, line 110 \\
Integral polynomial convention & \nolinkurl{Formalism/LocalRingCardinality/DataCoeff.lean}, \nolinkurl{rhoVec}, \nolinkurl{hcToWeightQ}, \nolinkurl{DomWeight.dual}, and the Weyl-module definitions \\
\end{longtable}

No Lean compilation is claimed. The source equality bodies cited in §2 were directly read; the new assertions are the complete oscillator/type identification, explicit partition recurrence, and degree formulas proved above. Membership of the desired arithmetic row in the \textbf{global} iterated image, the six actual \nolinkurl{AdmS} statements, and the \nolinkurl{MainX} terminal equation remain distinct obligations.

Primary reference: Konno, T., \& Konno, K. (2007). On doubling construction for real unitary dual pairs. \emph{Kyushu Journal of Mathematics, 61}(1), 35--82. \url{https://doi.org/10.2206/kyushujm.61.35}

\clearpage

\section{AM-DO-026: actual mixed-kernel operators and the remaining source equations}\label{am-do-026-actual-mixed-kernel-operators-and-the-remaining-source-equations}

Author: GPT-6 Astra Pro, ALP\_MACHINE \nolinkurl{/root/definite_oscillator}; inherited reasoning effort not separately exposed; 21 November 2026 UTC.

This finding combines the proved mixed-CM construction with the literal recovered source interfaces. It proves an all-polynomial geometric kernel identity, its compact covariance, a concrete finite stabilizer, and the exact adjoint and class-mass formulas for the resulting actual operators. It then identifies the additional equations needed to obtain the source \nolinkurl{PinColT}, \nolinkurl{TieIo}, and terminal identities. These last equations are not inferred from one nonzero theta value.

The carried inputs are AM-DO-017--020 and AM-DO-022--025, AM-R-004--006, AM-K018, and the accepted DP57 host. Their conclusions are used at their recorded checked status. In particular the finite direct-average phase and pairing are already proved in AM-DO-024; they are not re-proved here.

\subsection{1. A geometric kernel identity for every polynomial and finite translate}\label{a-geometric-kernel-identity-for-every-polynomial-and-finite-translate}

Use \[
 E'=\mathbf Q(\zeta_{28}),\quad F'=E'^+,
 \quad A=J^2,\quad \xi'=\xi/2,
 \quad C=\theta_0c.
 \tag{DO26.1}
\] For positive even ranks \(a,b\), let \(S=E'^a,T=E'^b\) have the positive identity Hermitian forms. All Heisenberg signs, measures, and splitting labels are those of AM-DO-022. Let \(\Lambda=\operatorname{Hom}_{\mathcal O_{E'}}
(\mathcal O_{E'}^a,\mathcal O_{E'}^b)\), with its actual perfect trace pairing, and let \(\phi_0\) be its finite trace-one vacuum.

For moved lattices \(N,N'\), define \(F_{N,N'}\) \textbf{geometrically} by the two normalized finite Mumford averages from the seed \((\sigma_q\boxtimes\sigma_q)^{\boxtimes ab}\). Thus each average uses the actual finite quotient, its trace-one theta-group multiplier, and the positive factor \(q^{-1/2}\). This definition precedes the analytic theta comparison. For finite \(g,h\) moving the standard lattices to \(N,N'\), write \[
 A_{g,h}X=hXg^{-1},\qquad
 \epsilon_\chi(g,h)=\chi_T(x)^{-1}\chi_S(y),
 \quad x/\bar x=\det g,\quad y/\bar y=\det h.
 \tag{DO26.2}
\] The independence from the determinant roots is the even-rank restriction \(\chi_S|_{\mathbb A_{F'}^\times}
=\chi_T|_{\mathbb A_{F'}^\times}=1\).

For clarity, a finite adelic translation parameter does not itself define a complex tangent vector. If \(y\) is such a parameter and \(\Lambda'\) is the rational period lattice of \(F_{N,N'}\), choose \(r\in\Lambda'\otimes\mathbf Q\) with \(\delta=y-r\in\widehat\Lambda'\). Additive approximation supplies such \(r\). Define the actual adelic translation on this theta line by \[
 \mathcal U_y^{\Lambda'}F
   =\psi_f(-b(r,\delta)/2)\alpha_{\Lambda'}(\delta)\,U_rF.
 \tag{DO26.2a}
\] The finite semicharacter is the continuous extension of the actual trace-one lattice character. The Weyl law \(W_rW_\delta=\psi_f(b(r,\delta)/2)W_y\), together with \(W_\delta\phi_{N,N'}=\alpha_{\Lambda'}(\delta)\phi_{N,N'}\), proves that (DO26.2a) is independent of the chosen rational representative and that \(\mathcal R W_y\phi_{N,N'}=C^{2ab}
\mathcal U_y^{\Lambda'}F_{N,N'}\). Equivalently, independence follows from the reduced-translation group law and the rational product formula. All displayed central factors are retained.

Choose a finite collection of rational Heisenberg parameters \(x\) and coefficients \(d_x\), and set \[
 \phi_d=\sum_x d_xW_x\phi_0,\qquad
 F_{d;g,h}=\sum_x d_x\mathcal U_{A_{g,h}x}^{\Lambda'}F_{N,N'}.
 \tag{DO26.3}
\] Here \(W_x\) and \(\mathcal U_y^{\Lambda'}\) are the actual Schrödinger and reduced Mumford translations, with their central phase retained. In particular no finite adele has been silently treated as a complex translation vector. For a polynomial \(P\) in the actual invariant tangent coordinates, let \(r_P\) be its corresponding Chern creation operator. Then \[
 \boxed{\vartheta_{r_P\phi_d}(g,h)
   =C^{2ab}\epsilon_\chi(g,h)
                  \,[P(\partial)F_{d;g,h}]_{0}.}
 \tag{DO26.4}
\] This is one common period for the entire moved family and all polynomial degrees, not a separate scalar selected at each value.

\textbf{Proof.} Exact Heisenberg covariance gives \[
 \omega(g,h)\phi_d
 =\epsilon_\chi(g,h)\sum_xd_xW_{A_{g,h}x}\phi_{N,N'}.
\] AM-DO-022 gives the vacuum identity with period \(C^{2ab}\), and AM-DO-018 identifies the actual translations under the reading map. Therefore \(\mathcal R(\omega(g,h)\phi_d)
=C^{2ab}\epsilon_\chi(g,h)F_{d;g,h}\) as entire functions. The finite action commutes with every archimedean creation operator, because they act on different factors of the restricted tensor product. The explicit identity \(\mathcal Rr_P=P(\partial)\mathcal R\), checked in ALP-LIFT-09, then gives (DO26.4) by evaluation at zero. All sums in (DO26.3) are finite. The theta sum and its fixed-order derivatives converge normally by the Gaussian estimate already proved for the reading map. Hence no interchange of divergent expressions is used. \(\square\)

AM-K018 identifies the starting Chern-reduced formal gauge with the actual source primitive gauge for the underlying ordinary product of copies of \(J\), including the integral block at 11. The extra \(E'\)-action does not replace that product by a product with conjugated \(E\)-actions. Its complete generic gauge comparison also transports through the actual CM isogenies and torsion translations in these finite averages: these preserve the same holomorphic character set, so the same no-opposite-characters uniqueness argument applies. This does not supply an integral 11-adic frame after an arbitrary isogeny of 11-power degree. The accepted original integral frame transports directly only when that isogeny is prime to 11, or when an additional integral comparison is proved. No reduced theta frame is asserted integral. The product trace remains \(\xi/2\), and the eigenbasis change has dyadic determinant valuation 3 per \(A\), as in AM-DO-022.

For specified finite class representatives \(g_i,h_j\), define the geometric jet functionals \[
 \mathcal A_{ji}(P)=\epsilon_\chi(g_i,h_j)
                    [P(\partial)F_{d;g_i,h_j}]_0.
 \tag{DO26.5}
\] They are defined from actual line-bundle translations and averages, not by dividing an analytic theta value by a desired constant. For algebraic finite characters and coefficients they have algebraic values, by AM-K018 and algebraicity of the fixed torsion translations. For any chosen source level volume \(v_S\), (DO26.4) gives the concrete all-polynomial class-kernel equation \[
 v_S\vartheta_{r_P\phi_d}(g_i,h_j)
                   =v_SC^{2ab}\mathcal A_{ji}(P).
 \tag{DO26.6}
\] This is the exact mixed-family analogue of \nolinkurl{ClsK.Core}. Identification with its independently specified \nolinkurl{abelian₂}, frame scale, and chart characters is an additional comparison, detailed in section 5.

\subsection{2. The actual compact covariance and an explicit finite stabilizer}\label{the-actual-compact-covariance-and-an-explicit-finite-stabilizer}

For the selected common-character edge of AM-DO-025, put \(a=N_j,b=N_{j+1}\) and \(\chi_S=\chi_T=\mu_{j+1}\), with infinity exponent \(-e_j\). The corrected source splitting formula gives \[
 \omega_\infty(g,h)r_P\phi_\infty
 =\prod_\tau\det(g_\tau)^{e_{j,\tau}-b/2}
                    \det(h_\tau)^{-e_{j,\tau}+a/2}
      r_{\operatorname{mvZ}(g^{-1},h)P}\phi_\infty.
 \tag{DO26.7}
\] The polynomial action is expressed in the source's creation-label convention \nolinkurl{mvZ}; it is the tensor action on \(T\otimes S^*\). Equation (DO26.7) is AM-DO-023's honest-splitting formula applied to all creation polynomials, not merely to a lowest-weight line.

The exact recovered definitions give \[
 s_S=e_j-b/2=\texttt{sSharpC},\qquad
 s_D=e_j-a/2=\texttt{sDivC}.
 \tag{DO26.8}
\] Indeed \nolinkurl{Params.sExp'} and \nolinkurl{sExp} have these values at \(j+1\); their integer floors in \nolinkurl{Columns/Record.lean}, lines 1083--1087, are unchanged because \(e_j\) is integral and both ranks are even. Thus (DO26.7) is exactly the determinant and polynomial covariance required by \nolinkurl{ClsK.Real}, after matching the actual tangent frames. AM-DO-025 proves that the source partition \nolinkurl{nuW} selects the contragredient of \nolinkurl{lamTrue} on the source and the next \nolinkurl{lamTrue} on the target. The real covariance and the local type match are therefore established for this constructed compact family.

Here is also a definite finite-level assertion. Suppose every parameter in (DO26.3) belongs to \(M^{-1}\Lambda\), where \(M\) is a positive integer. Let \(K_S,K_T\) be the standard integral unitary compacts, refined by the principal congruence condition \(1\pmod{2M^2}\). For trivial characters, both groups fix \(\phi_d\) exactly: \[
                     \omega(K_S,K_T)\phi_d=\phi_d.
 \tag{DO26.9}
\] For nontrivial even-rank characters the same assertion holds after intersecting with the kernels of the two determinant characters \(\nu_{\chi_T}^{-1}\circ\det\) and \(\nu_{\chi_S}\circ\det\). These are open subgroups of the compacts, so the refinement is finite.

\textbf{Proof.} AM-DO-022 proves that the trivial-character standard compacts fix \(\phi_0\). The induced Hom action is congruent to the identity modulo \(2M^2\), including after inversion on the source. For \(x\in M^{-1}\Lambda\), write \(\delta=A_{g,h}x-x\in2M\widehat\Lambda\). Perfection and integrality of the actual trace pairing give \[
 \psi_f(b(x,\delta)/2)=1,\qquad
 \alpha_{\rm tr}(\delta)
      =\psi_f(b(\delta,w_0\delta)/2)=1.
\] For the first equality write \(x=\lambda/M,\delta=2M\eta\); the half-pairing is \(b(\lambda,\eta)\), which is integral. For the second it is \(2M^2b(\eta,w_0\eta)\), also integral. The exact Weyl group law therefore yields \(W_{x+\delta}\phi_0=W_x\phi_0\). Covariance fixes each summand in (DO26.3), proving (DO26.9). For the specified nontrivial characters the only remaining compact scalar is the displayed determinant factor from AM-DO-022; killing it gives the assertion. \(\square\)

This bound is deliberately uniform for the named finite span. It is not asserted to be the least level; ALP-LIFT-10 obtains a smaller one for its special isotropic half-orbit. Nor does (DO26.9) prove that an independently pinned source level already lies in this stabilizer. The index of the needed refinement remains an arithmetic cost.

\subsection{3. Exact adjoint and conjugation identities for the compact kernel}\label{exact-adjoint-and-conjugation-identities-for-the-compact-kernel}

Let \(X=[U(S)]\) and \(Y=[U(T)]\), with specified compatible quotient Haar measures. Both are compact because the groups are definite. Let \(K(x,y)\) be one of the actual polynomial theta kernels from section 1. Define the source-style \textbf{bilinear} operators \[
 (T_Kf)(y)=\int_XK(x,y)f(x)\,dx,
 \qquad
 (S_Kg)(x)=\int_YK(x,y)g(y)\,dy.
 \tag{DO26.10}
\] Use the inner product conjugate-linear in its first argument, and write \(C_Xf=\bar f\), \(C_Yg=\bar g\). Then the exact identities are \[
 \boxed{T_K^*=C_XS_KC_Y,\qquad
        C_YT_KC_X=T_{\bar K}.}
 \tag{DO26.11}
\] The second equation keeps the same source-to-target orientation; in the first, \(S_K\) reverses it. In particular \nolinkurl{liftS} with the same kernel is not the Hermitian adjoint of \nolinkurl{liftT} until these pointwise conjugations are inserted.

\textbf{Proof.} At fixed level the quotient is a finite union of compact real-group pieces. The finite support is fixed on each piece, and the positive Gaussian quadratic form has a uniform positive lower bound as the real group varies over its compact set. Its polynomial theta series therefore converges normally and gives a continuous bounded kernel. It is square integrable, and Fubini applies to the displayed operators, first for continuous functions and then to all \(L^2\) functions by the bounded integral-operator estimate. Directly, \[
 \begin{aligned}
 \langle T_Kf,g\rangle_Y
  &=\int_X\bar f(x)
       \left(\int_Y\overline{K(x,y)}g(y)\,dy\right)dx\\
  &=\langle f,C_XS_KC_Yg\rangle_X.
 \end{aligned}
\] Conjugating (DO26.10) term by term gives the second equality. This proves (DO26.11). \(\square\)

The corresponding bilinear pairing \(B_X(f_1,f_2)=\int_Xf_1f_2\) satisfies \[
                  B_Y(T_Kf,g)=B_X(f,S_Kg).
 \tag{DO26.12}
\] Consequently \[
             \langle f,T_K^*T_Kf\rangle_X=\|T_Kf\|_Y^2.
 \tag{DO26.13}
\] This is strictly positive for an actual nonzero first image. It constructs an actual detecting Gram operator; it does not identify that operator with a separately prescribed positive-rank next theta edge. Complex conjugating the theta series also gives a concrete kernel in the conjugate Weil representation, with the additive character and splitting scalars conjugated. This supplies the conjugations in (DO26.11) without guessing a source-target swapping character formula.

The scalar bookkeeping is exact. If \(K=c_0\mathcal K\), then \[
 T_K=c_0T_{\mathcal K},\qquad
 T_K^*=\bar c_0T_{\mathcal K}^*,\qquad
 T_K^*T_K=|c_0|^2T_{\mathcal K}^*T_{\mathcal K}.
 \tag{DO26.14}
\] Here \(c_0=C^{2ab}\), with any explicitly chosen creation-frame scaling retained. Neither the complex-torus probability measure of AM-DO-024 nor a Gaussian unit normalization silently substitutes for the automorphic quotient measure in these formulas.

For a compact subgroup embedding \(\iota:H\to G\), the actual restriction period \(B_{\rm top}(F,G)=\int_{[H]}\overline{F(\iota h)}G(h)\,dh\) obeys \[
 B_{\rm top}(C F,G)=\int_{[H]}F(\iota h)G(h)\,dh.
 \tag{DO26.15}
\] Thus the source's \nolinkurl{conjAF} in \nolinkurl{TieBm} has a concrete role: it converts its conjugate-linear first argument into the bilinear terminal pairing. Equation (DO26.15) does not specify the arithmetic terminal row or its scalar.

\subsection{4. Actual finite class masses and the operator matrix}\label{actual-finite-class-masses-and-the-operator-matrix}

Fix an open compact \(K_f\subset G(\mathbb A_f)\), and write \(g_i\) for the finitely many double-class representatives. Put \[
 \Gamma_i=G(F)\cap g_iK_fg_i^{-1},\qquad
       v_K=\operatorname{vol}(K_f),\qquad w_i=v_K/|\Gamma_i|.
 \tag{DO26.16}
\] The real compact group has Haar volume one. The rational intersection is finite: it is discrete in the compact product of the real group and the finite compact, hence finite. Quotient integration gives \[
 \int_{[G]}F(g)\,dg
       =v_K\sum_i\frac1{|\Gamma_i|}
                    \int_{G_\infty}F(g_ik)\,dk.
 \tag{DO26.17}
\] Indeed the piece is the diagonal quotient of \(G_\infty\times K_f\) by the finite group \(\Gamma_i\). Pushing forward Haar measure divides by its order; then integrating the right \(K_f\)-invariant function contributes \(v_K\). This proves (DO26.17) with the chosen compatible quotient measure.

If \(\rho\) is an irreducible unitary real coefficient type of dimension \(d_\rho\), and a unit functional \(\ell\) is used to realize its vectors as \(k\mapsto\ell(\rho(k)v)\), then \[
 \int_{G_\infty}|\ell(\rho(k)v)|^2\,dk
                          =\frac{\|v\|^2}{d_\rho}.
 \tag{DO26.18}
\] For completeness, the averaged positive rank-one operator \(\int\rho(k)^*\ell^*\ell\rho(k)\,dk\) commutes with \(\rho\), hence is scalar by irreducibility; its trace is one, so its scalar is \(1/d_\rho\). This proves (DO26.18). Multiplication of this realization by \(\sqrt{d_\rho}\) makes it isometric. Without that normalization, the factor in (DO26.18) must remain in every class metric.

Choose orthonormal real compact coefficient bases in each class piece, using that normalization and its actual invariant subspaces. Then the metric on the class coefficient vector is \(G_X=\operatorname{diag}(w_i I_i)\). Do the same on the target, with metric \(G_Y\). Expand the relevant finite-type kernel as \[
 K((i,k),(j,l))
       =\sum_{r,s}A_{ji,sr}\overline{u_{i,r}(k)}v_{j,s}(l).
 \tag{DO26.19}
\] Such an expansion is finite for a fixed polynomial input: its compact-group orbit is finite dimensional. Orthogonality and (DO26.17) now give the \textbf{actual operator matrices} \[
 \boxed{T=A G_X,\qquad
        T^{\dagger}=G_X^{-1}\bar T^{\,t}G_Y
                         =\bar A^{\,t}G_Y.}
 \tag{DO26.20}
\] The formulas hold blockwise on the invariant coefficient subspaces; empty blocks are omitted. They follow simply by integrating \(\bar u_{i,r}u_{i,r'}\) and then evaluating the defining Hermitian-adjoint equality. Thus the stabilizer denominators, level volume, real Schur factor, and conjugation are all explicit.

This construction proves a class-mass formula for the actual automorphic operator. An independently given integral polynomial or Weyl-module basis need not be this orthonormal basis. Changing to it requires its actual Gram matrix and lattice determinant; those quantities are not removed by (DO26.20). In particular an arbitrary choice of the source's \nolinkurl{rowAF} cannot be made isometric by assertion.

\subsection{5. Exact comparison with the recovered arithmetic pins}\label{exact-comparison-with-the-recovered-arithmetic-pins}

The definitions in \nolinkurl{LevChiEqOne.lean} give the following nonoptional equations. The symbols on their right sides are independently specified source data.

First, \nolinkurl{VacT} requires (DO26.4) at degree zero with the source \nolinkurl{xiS}, \nolinkurl{I.J.V.thMoved}, and determinant-root factor \(\mu(x)^{-1}\mu(y)\), uniformly in every finite move. AM-DO-022 proves that mathematical statement for the new mixed geometric family and the standard Weil representation. \nolinkurl{PinX.mv} additionally quantifies over every \nolinkurl{om} satisfying \nolinkurl{IsWeilDatumO}; identifying all such axiomatic data with this fixed standard realization requires its uniqueness statement. The old elliptic \nolinkurl{VacIn} is not the new seed merely because both are vacua.

Second, \nolinkurl{PinColT.cst} is, literally, \[
 c_{K,j}=v_{j,\alpha}\,w_{j,\alpha\gamma}\,c_{0,j}
       \mu_{j+1}(x_{j,\alpha})
       \mu_{j+1}(x_{j+1,\gamma})^{-1}
                 \quad\text{for every }\alpha,\gamma.
 \tag{DO26.21}
\] The volume is that of the source chart level, not a theta-line norm. If every factor is nonzero, independence of the chart is equivalent to all pairwise ratio equations \[
 \frac{v_{j,\alpha}w_{j,\alpha\gamma}}
      {v_{j,\alpha'}w_{j,\alpha'\gamma'}}
 \frac{\mu(x_{j,\alpha})}{\mu(x_{j,\alpha'})}
 \frac{\mu(x_{j+1,\gamma'})}{\mu(x_{j+1,\gamma})}=1.
 \tag{DO26.22}
\] This equivalence follows by dividing two instances of (DO26.21); conversely fix one chart to define \(c_{K,j}\) and use all ratios. It is a finite concrete compatibility test. The geometric period \(c_{0,j}=C^{2N_jN_{j+1}}\) is now known; the source chart factors and their arithmetic meanings cannot be replaced by freely chosen scalars just to force (DO26.22).

Third, \nolinkurl{ClsK.Core} requires for \textbf{every} polynomial \[
 v_{j,\alpha}\Theta_{j,\alpha\gamma}(h_a,h_{a'};P)
  =c_{K,j}\,\operatorname{abelian}_2
       \bigl(\operatorname{mvZ}(1,\operatorname{sc}_j I)
                            \operatorname{castP}(P)\bigr).
 \tag{DO26.23}
\] Our equation (DO26.6) supplies its analytic-to-geometric kernel comparison. The actual source functional is \(\rho_i(c)\rho_i'(c')\operatorname{jetL}
  (I.\mathrm{th}(\operatorname{rep}c,\operatorname{rep}c'))\), as defined at lines 73--77 of \nolinkurl{LevChiEqOne.lean}. Its frame scale, finite arithmetic datum, and these two character scalars must be matched to (DO26.5). The finite-level statement in section 2 and the real statement (DO26.7) supply corresponding \nolinkurl{Lev} and \nolinkurl{Real} identities for the newly constructed datum. They do not identify an independently fixed datum without this comparison.

Fourth, \nolinkurl{ClsK.Kval} says that the prescribed arithmetic row map \(K_j\), evaluated in its specified frames, is exactly the stabilizer-weighted geometric kernel sum, including the target factor \(\det(y')^{-s_D}\). Equation (DO26.20) proves the class integral formula for the actual operator. To deduce that particular \nolinkurl{Kval}, one must exhibit the same coefficient-lattice realization and its compact projection normalization. Local type occurrence from AM-DO-025 does not choose that realization.

The abstract \nolinkurl{ThetaX} ties are equally precise. With \(\operatorname{petY}_j=m_{Y,j}\operatorname{pet}_j\) and \(\operatorname{dnY}_j=m_{Y,j+1}\Theta_j\), they require \[
 \begin{aligned}
 m_{Y,j}&=\prod_v\operatorname{parIdx}(q_v,r,jt),\\
 \operatorname{petY}_j(\operatorname{rowAF}_j x,
                      \operatorname{rowAF}_j y)&=\langle x,y\rangle,\\
 \operatorname{petY}_j(\operatorname{dnY}_j
           \operatorname{rowAF}_{j+1}y,\operatorname{rowAF}_j x)
                &=\langle\operatorname{io}_j y,x\rangle.
 \end{aligned}
 \tag{DO26.24}
\] Equations (DO26.11) and (DO26.20) give the exact analytic adjoint and its actual matrix. They establish (DO26.24) only when the preassigned \nolinkurl{rowAF}, \nolinkurl{io}, metrics, and parahoric factors are identified with those matrices. Choosing a new abstract map equal to an adjoint would prove a different statement.

\subsection{6. Terminal mass and the literal two-step MainX equation}\label{terminal-mass-and-the-literal-two-step-mainx-equation}

The source terminal equation \nolinkurl{ClsT.tail} is \[
 v_T\sum_{\beta,a}\frac{1}{|\Gamma_{T,a}|}
    \ell_{\beta,a,G}\!\left(
      \operatorname{low}\bigl(\operatorname{rd}_{m,\beta}
                    \psi\circ\iota(h_a,1)\bigr)\right)
                 =c_P\,\operatorname{tail}(\psi,G).
 \tag{DO26.25}
\] It holds for every allowed arithmetic row \(\psi\) and every terminal coefficient vector \(G\), with the same \(c_P\). The actual finite theta-line overlap \(q_0^{-1/2}\), and the holomorphic overlap with period and metric factors in AM-DO-024, are completely determined. They are not (DO26.25): the latter contains a different quotient measure, an actual restriction embedding, coefficient projections, and a prescribed arithmetic terminal functional.

For fixed finite-dimensional row and terminal spaces, (DO26.25) is a finite matrix test. Choose their actual bases, let \(B\) be the matrix of the left side computed with its masses, and let \(T\) be the matrix of \nolinkurl{tail}. The requirement is \[
                              B=c_PT.
 \tag{DO26.26}
\] If \(T\ne0\), choose one nonzero entry \(T_{rs}\). Then there is such a scalar if and only if \(B_{ij}T_{rs}=B_{rs}T_{ij}\) for all entries. This condition also forces \(B_{ij}=0\) whenever \(T_{ij}=0\). The proof is division by the chosen entry, followed by bilinearity on the two bases. Thus one nonzero period fixes a candidate scalar but does not verify its uniformity. The source also requires \(c_P\ne0\) where its later nonvanishing arguments use it.

For the literal two-step case \(m=2\), \nolinkurl{MainX} uses the implication \[
 \begin{split}
 &\sum_i\alpha_i
   c_{V,0}^{(i)}c_{K,0}^{(i)}c_{K,1}^{(i)}c_P^{(i)}
   \operatorname{tail}^{(i)}
      (K_1^{(i)}K_0^{(i)}S_0^{(i)}\psi_i,F_i)=0\\
 &\quad\Longrightarrow\quad
 \sum_i\alpha_i c_{I,0}^{(i)}c_{I,1}^{(i)}
   \operatorname{per}_0^{(i)}
      (\psi_i,\operatorname{io}_0^{(i)}
                       \operatorname{ioT}_0^{(i)}F_i;y,v)=0
                         \quad\text{for all }y,v.
 \end{split}
 \tag{DO26.27}
\] The quantifier over all finite linear relations and common frames is essential. In the general even chain the factors are exactly \[
 c_{\rm Top}=\prod_i(c_{V,2i}c_{K,2i}c_{K,2i+1})c_P,
               \qquad c_{\rm Bot}=\prod_jc_{I,j}.
 \tag{DO26.28}
\] These are \nolinkurl{Tup.cTop}, \nolinkurl{cBot}, and \nolinkurl{MainX} in \nolinkurl{ValLastX.lean}. The new compact kernels prove substantive components that enter this system. They do not replace its Bessel/Jacobi two-step unfolding, its fixed row maps, or (DO26.25).

\subsection{7. The six actual analytic predicates and exact source scope}\label{the-six-actual-analytic-predicates-and-exact-source-scope}

The source's \nolinkurl{AdmS} is a conjunction of the following six predicates, not a name for the elementary compact kernel identities alone.

\begin{longtable}[]{@{}
  >{\raggedright\arraybackslash}p{(\columnwidth - 4\tabcolsep) * \real{0.3333}}
  >{\raggedright\arraybackslash}p{(\columnwidth - 4\tabcolsep) * \real{0.3333}}
  >{\raggedright\arraybackslash}p{(\columnwidth - 4\tabcolsep) * \real{0.3333}}@{}}
\toprule\noalign{}
\begin{minipage}[b]{\linewidth}\raggedright
Actual predicate
\end{minipage} & \begin{minipage}[b]{\linewidth}\raggedright
What the new construction proves
\end{minipage} & \begin{minipage}[b]{\linewidth}\raggedright
Obligation still present in the literal source
\end{minipage} \\
\midrule\noalign{}
\endhead
\bottomrule\noalign{}
\endlastfoot
\nolinkurl{RallisSiegelWeilFormula} & Actual compact kernels, exact prescribed local coefficient types, and separate concrete nonzero first lifts & For every \nolinkurl{Hyp3} parameter, a nonzero irreducible global chain in the specified types, with every upper vector a finite sum of prescribed theta lifts \\
\nolinkurl{LiftParameterConditions} & Actual compact restriction pairings and their conjugation convention & A nonzero terminal period against the specified noncompact discrete-series family for every allowed top irreducible representation \\
\nolinkurl{ThetaUnfoldingFreeData} & Exact compact Fubini/adjoint identities (DO26.11)--(DO26.12) & The source's Bessel/Jacobi period transfer with factorizable, real-finite downward Schwartz data and its universal nonzero-period implication \\
\nolinkurl{ThetaStablePeriodVanishing} & No cuspidality problem for the constructed anisotropic compact groups themselves & The stated lower-Witt-tower vanishing implication for the noncompact \nolinkurl{Y} chain \\
\nolinkurl{RealThetaLiftDiscreteSeries} & Exact compact real formula and the selected compact Schur correspondence & The specified noncompact discrete-series preservation and lower-signature vanishing statements \\
\nolinkurl{RealConstantNonvanishing} & Translation preserves the actual compact representation and its fixed local type & Stability of the source's particular noncompact \nolinkurl{IsDv} condition under all displayed finite and real actions \\
\end{longtable}

Consequently the correct progression is: an actual mixed theta seed; its exact all-finite moved family and direct pairing; its all-degree geometric class kernel; the selected real covariance; then the independently pinned arithmetic row and terminal equations. The first four now have explicit formulas. The remaining source comparisons are concrete algebraic and automorphic obligations, not evidence that no other map or construction can exist.

\subsubsection{Exact recovered-source locators}\label{exact-recovered-source-locators}

All paths below are relative to \nolinkurl{work/research/alpoge_source}.

\begin{longtable}[]{@{}
  >{\raggedright\arraybackslash}p{(\columnwidth - 2\tabcolsep) * \real{0.5000}}
  >{\raggedright\arraybackslash}p{(\columnwidth - 2\tabcolsep) * \real{0.5000}}@{}}
\toprule\noalign{}
\begin{minipage}[b]{\linewidth}\raggedright
Interface
\end{minipage} & \begin{minipage}[b]{\linewidth}\raggedright
Source locator
\end{minipage} \\
\midrule\noalign{}
\endhead
\bottomrule\noalign{}
\endlastfoot
Bilinear \nolinkurl{liftS}/\nolinkurl{liftT} & \nolinkurl{Development/AbelianVarietyTheory/Definitions/TwistedInnerProducts/WittSmallGroupEmbeddingMemPointwiseStabilizerIso.lean}, lines 157--167 \\
Finite-combination kernel and lift & \nolinkurl{Development/AbelianVarietyTheory/Definitions/LatticeFrames/JointSmooth.lean}, \nolinkurl{KerDat.pt}, \nolinkurl{lift}, lines 74--82 \\
Actual vacuum, character roots, chart constant & \nolinkurl{Development/AbelianVarietyTheory/Definitions/LatticeFrames/LevChiEqOne.lean}, \nolinkurl{abelian₂}, lines 73--77; \nolinkurl{VacT}, 87--93; \nolinkurl{PinX}, 95--101; \nolinkurl{PinColT}, 103--119 \\
Real, level, kernel and arithmetic row identities & \nolinkurl{Development/AbelianVarietyTheory/Definitions/LatticeFrames/OnFramesRdApply.lean}, \nolinkurl{ClsK.Real}, 93--102; \nolinkurl{Lev}, 103--107; \nolinkurl{Core}, 108--113; \nolinkurl{Kval}, 115--121 \\
Terminal class mass & Same file, \nolinkurl{ClsT.Is.top} and \nolinkurl{tail}, lines 143--150 \\
Exact determinant exponents & \nolinkurl{Formalism/LocalRingCardinality/Params.lean}, \nolinkurl{sExp}, \nolinkurl{sExp'}, lines 71--73; \nolinkurl{Formalism/LocalToGlobal/Columns/Record.lean}, \nolinkurl{sSharpC}, \nolinkurl{sDivC}, lines 1083--1087 \\
Conjugation, periods and row ties & \nolinkurl{Formalism/LocalRingCardinality/DataTheta.lean}, \nolinkurl{conjAF}, 645--649; \nolinkurl{Bper} and \nolinkurl{Btop}, 657--670; \nolinkurl{petY}, \nolinkurl{dnY}, \nolinkurl{TieMY}, \nolinkurl{TieRowIsom}, \nolinkurl{TieIo}, \nolinkurl{TieBm}, 730--752 \\
Two-step and terminal relation laws & \nolinkurl{Development/AbelianVarietyTheory/Definitions/LatticeFrames/ValLastX.lean}, \nolinkurl{ValX}, \nolinkurl{ValLastX}, \nolinkurl{KerX}, \nolinkurl{PerX}, \nolinkurl{LawsX}, \nolinkurl{LastX}, \nolinkurl{MainX}; \nolinkurl{Tup.cTop} and \nolinkurl{cBot} \\
Six analytic assumptions & \nolinkurl{Formalism/SharpRestriction/CuspSpanNormBounds/CayleySignature.lean}, lines 154--202 \\
Existing vacuum-to-polynomial comparison & \nolinkurl{Development/AbelianVarietyTheory/Proofs/DecoratedLatticeProofs/ThetaKernelSplitPlacesEqJetClassicalTheta.lean}, complete theorem \nolinkurl{thetaKernel_eq_jet_thC} \\
\end{longtable}

No Lean build or formal instantiation of a source record is claimed. The cited interfaces and theorem bodies were directly read. The new operator proofs above state exactly which representation, period, measure, and coefficient normalization they use.

\clearpage

\section{AM-K021 --- the cubic relative period and the exact CM-height boundary}\label{am-k021-the-cubic-relative-period-and-the-exact-cm-height-boundary}

\textbf{Status.} New proved reductions and scoped obstructions for question 3. The stipulated Alpöge theorem, its archive, and the earlier accepted DP statements retain their accepted status. This report does not establish the relative arithmetic window, a new RH theorem, or a new automorphic cost bound.

\textbf{Dependencies.} AM-K001--AM-K005 in kummer\_extension.md and AM-K007--AM-K011 in kummer\_cubic\_host.md supply the Artin identity, conductor table, completed functional equation, elliptic Hom realization, finite ray character, norm-one splitting twist, and explicit CM supply tower. They are used with those statements' existing proofs and scope. The new result strengthens the earlier observation that coherent splitting phases telescope: even a genuine cubic twist of the actual polarized CM variety leaves the stable geometric letter degrees unchanged. A separate finite-character argument excludes the cubic standard Artin character from every finite linear combination of CM-type height classes.

\subsection{1. The exact source scalar is geometric and quadratic}\label{the-exact-source-scalar-is-geometric-and-quadratic}

All paths in this section are relative to work/research/alpoge\_source. These are inspected source interfaces, not a claim to have compiled or independently reconstructed their full upstream proofs.

\begin{longtable}[]{@{}
  >{\raggedright\arraybackslash}p{(\columnwidth - 2\tabcolsep) * \real{0.5000}}
  >{\raggedright\arraybackslash}p{(\columnwidth - 2\tabcolsep) * \real{0.5000}}@{}}
\toprule\noalign{}
\begin{minipage}[b]{\linewidth}\raggedright
Source locator
\end{minipage} & \begin{minipage}[b]{\linewidth}\raggedright
Exact role
\end{minipage} \\
\midrule\noalign{}
\endhead
\bottomrule\noalign{}
\endlastfoot
\nolinkurl{Formalism/LocalRingCardinality/DataParams₂.lean,} lines 65--71 & chiK is the quadratic Dirichlet character modulo D.host.d; dlogL is the real part of its logarithmic derivative. \\
The same file, lines 126--127 & SizeData₂.logDerivativeAtOne is constrained by logDerivative\_at\_one\_eq to equal dlogL D 1. It is not a freely chosen real parameter. \\
The same file, line 224 & logOm is Euler's constant plus log(π / host.v) minus that logarithmic derivative. \\
The same file, lines 232--234 & MAIN contains the term −R.m · aT(t) · logOm for every host type t. \\
\nolinkurl{Formalism/LocalToGlobal/ProjectiveLattices/RosatiLatticeNorms/LamObj.lean,} definitions IsVecLine and LamObjGood & A letter line must carry the actual integral tangent lattice and the norm induced by the polarized complex uniformization. \\
\nolinkurl{Formalism/LocalToGlobal/ProjectiveLattices/RosatiLatticeNorms/LamPtClauses.lean,} lines 53--56 & LamDegSt requires every actual letter to have normalized degree logOm / 2. \\
\nolinkurl{Formalism/LocalToGlobal/ProjectiveLattices/RosatiLatticeNorms/CmFrameOfEq.lean,} lines 45--58 and 78--115 & The degree identity is transported to the frame letters and is a hypothesis of the construction. \\
\nolinkurl{Formalism/LocalToGlobal/CMDeterminantIdeals/StatementS.lean,} lines 63--73 & Step5 consumes those exact letter degrees in its determinant estimate. \\
\end{longtable}

In particular the actual source chain has two separate equalities: \[
 \log\Omega=
 \gamma+\log(\pi/v)
 -\frac{L'(1,\chi_K)}{L(1,\chi_K)},
 \qquad
 \widehat{\deg}_{\rm norm}(\ell)=\frac12\log\Omega.
 \tag{R1}
\] The first equality is a prescribed auxiliary quadratic analytic constant. The second is a fact about actual integral metrized lines. Renaming a splitting character does not change either equality.

The source IsVecLine is stronger than a choice of norm on an abstract one-dimensional vector space. At every finite place its fractional ideal is characterized by membership in the tangent lattice of a model. At every complex embedding its squared norm equals π times the polarized Hermitian length of the corresponding tangent vector. Thus an arbitrary metric rescaling chosen to force a desired new value would require a new comparison with those actual geometric conditions.

\textbf{Scope.} The original Dirichlet argument can use auxiliary quadratic hosts while proving conclusions for other Dirichlet characters. Consequently (R1) does not imply that every successful relative argument must replace the quadratic logarithmic derivative by the target cubic logarithmic derivative. It proves that this replacement is not an available instantiation of the unchanged record. The task could instead require a new character-weighted arithmetic comparison using the existing auxiliary heights.

\subsection{2. Cubic descent does not create a new stable CM letter degree}\label{cubic-descent-does-not-create-a-new-stable-cm-letter-degree}

Set \[
 K=\mathbf Q(\zeta),\quad \zeta=\zeta_3,\quad
 \alpha=p^{1/3}>0,\quad M=K(\alpha),
\] and retain the standing condition p=3 or p≡2 modulo 3. Consider the actual elliptic curves \[
 A_0:\ y^2=x^3+1,\qquad A_p:\ y^2=x^3+p^2.
\] Their geometric CM action is (x,y)↦(ζx,y).

\subsubsection{Proposition R2. Stable isometry under the cubic twist}\label{proposition-r2.-stable-isometry-under-the-cubic-twist}

After a finite extension where both curves have good reduction, the polarized CM isomorphism \[
 g_p:A_p\longrightarrow A_0,\qquad
 (x,y)\longmapsto (x/\alpha^2,y/p)
 \tag{R2}
\] identifies their integral tangent eigenlines, their polarized archimedean metrics, and their normalized Arakelov degrees. The same holds for each eigenline in the polarized CM products constructed in AM-K010--AM-K011 with the same fixed positive Hermitian form.

\textbf{Proof.} The displayed substitution gives \[
 (y/p)^2=(x/\alpha^2)^3+1
\] because α³=p.~It fixes the origin, commutes with the CM action, and hence preserves the canonical principal elliptic polarization. After a further finite extension both curves have abelian-scheme models. The generic-fibre isomorphism and its inverse extend uniquely between these models, so their maps on integral Lie algebras are inverse integral isomorphisms. At a complex embedding the polarized analytic uniformizations are isomorphic, and the induced Hermitian metric is carried by the same tangent map. The resulting integral metrized lines are isometric.

Normalized Arakelov degree is unchanged by a finite base extension: both finite and infinite contributions multiply by the extension degree before normalization. It therefore gives the same degree for the two stable eigenlines over their original field of definition.

For the products in AM-K010, apply g\_p on each elliptic factor. The scalar map commutes with the given matrix action of the larger CM coefficient field and with the fixed Hermitian form. It carries the corresponding polarization and each coefficient-field eigenline to the untwisted one. The preceding argument applies to every such line, not just to the determinant of the whole tangent space. ∎

The coordinate scaling is useful for detecting an incomplete accounting. For the algebraic differential ω=dx/y, \[
 g_p^*\omega_0=\alpha\,\omega_p,
 \qquad
 \int_{\gamma_p}\omega_p
  =\alpha^{-1}\int_{g_p\gamma_p}\omega_0.
 \tag{R3}
\] Thus the displayed complex period changes by p\^{}\{-1/3\}. This is a change of algebraic basis in an isometric stable line. Its finite principal-divisor contribution cancels its archimedean contribution by the product formula. Keeping only the −(log p)/3 in the period would not compute a normalized Arakelov degree.

This assertion concerns stable lines after good reduction and their actual polarized metrics. It does not identify unstable minimal models, remove conductor terms, or set the finite level costs at 2, 3 and p to zero. Those costs remain part of any relative arithmetic adapter.

\textbf{Consequence.} In the particular induced CM supply tower already constructed, a norm-one Kummer descent twist can preserve the nontrivial Hom character while leaving every degree used in (R1) unchanged. The nontrivial Galois character is real arithmetic information, but it is not a new stable CM-type height.

\subsection{3. A representation-level obstruction for all CM-type height classes}\label{a-representation-level-obstruction-for-all-cm-type-height-classes}

Write G=Gal(M/Q). With σ(α)=ζα and complex conjugation c fixing α, \[
 G\simeq S_3,\qquad c\sigma c^{-1}=\sigma^{-1}.
 \tag{R4}
\] Let ρ be the two-dimensional augmentation representation on the three roots of X³−p.~Its character is \[
\begin{array}{c|ccc}
 \text{cycle type}&1&(12)&(123)\\ \hline
 \chi_\rho&2&0&-1.
\end{array}
\tag{R5}
\] AM-K001 gives \[
 L_K(s,\chi)=L_{\mathbf Q}(s,\rho)
            =\zeta_{\mathbf Q(\alpha)}(s)/\zeta_{\mathbf Q}(s).
 \tag{R6}
\]

Let Q\^{}\{cm\} be the composite of all CM subfields of a fixed algebraic closure. Complex conjugation is central on every finite Galois CM subfield. Such subfields are closed under finite composition and Galois closure. Equivalently Q\^{}\{cm\} is fixed by the closed subgroup generated by all commutators with complex conjugation (Milne, 2020a, I, Proposition 1.4, Corollary 1.5 and Remark 1.6, p.~10).

\subsubsection{Proposition R7. No cubic component in a finite CM-type height class}\label{proposition-r7.-no-cubic-component-in-a-finite-cm-type-height-class}

Let B/Q be a finite Galois CM field, N=BM, and H=Gal(N/B). Inflate ρ to Gal(N/Q). For every complex class function A inflated from Gal(B/Q), \[
 \langle A,\chi_\rho\rangle_{\operatorname{Gal}(N/\mathbf Q)}=0.
 \tag{R7}
\] Consequently no finite complex linear combination of class functions factoring through Q\^{}\{cm\} equals χρ. In particular no finite combination of CM-type height Artin classes gives the target cubic Artin character.

\textbf{Proof.} B∩M is Galois over Q because B and M are Galois. The normal subfields of the S₃ extension M are Q, K and M. The last possibility would imply that complex conjugation is central on Gal(M/Q), because it is central on Gal(B/Q). This contradicts (R4). Hence B∩M is Q or K, and the restriction of H to M is respectively S₃ or its normal subgroup A₃.

The restriction of ρ to A₃ is the sum of its two nontrivial one-dimensional characters. It has no invariant vectors. Thus ρ\^{}H=0. Averaging ρ(h) over h∈H is the zero projector. For each g∈Gal(N/Q), \[
 \sum_{h\in H}\chi_\rho(gh)
 =\operatorname{Tr}\!\left(\rho(g)\sum_{h\in H}\rho(h)\right)=0.
\] A is constant on each H-coset. Summing the character inner product by those cosets proves (R7).

Any finite collection of locally constant functions on Gal(Q\^{}\{cm\}/Q) factors through one finite Galois CM field after enlarging their fields of definition. Their linear span is therefore orthogonal to χρ. Since the latter is an irreducible character with inner product one with itself, it cannot belong to this span. ∎

The standard CM-type class is AΦ\^{}0, the conjugacy average of the convolution Φ*Φ∨ on Gal(Q\^{}\{cm\}/Q). Its construction is exactly on this quotient (Yang, 2010, §3, equations (3.1)--(3.12); author manuscript pp.~9--10). Proposition R7 concerns that definition and its finite Artin support. It does not require a new case of the Colmez formula.

Thus the obstruction remains even if every relevant CM-type height realization is granted. A finite identity of Artin representations, or an equality of all corresponding Artin Euler factors, cannot turn those CM-type height classes into ρ. Indeed equality of Euler factors at all unramified primes gives equality of characters on a common finite Galois group by Frobenius density, contradicting (R7).

An equivalent short detection is \[
 [c,\sigma]=c\sigma c^{-1}\sigma^{-1}=\sigma^{-2}=\sigma,
\] which acts nontrivially under ρ. Every representation factoring through Q\^{}\{cm\} kills this commutator.

\subsubsection{Limits of the obstruction}\label{limits-of-the-obstruction}

\begin{enumerate}
\def\labelenumi{\arabic{enumi}.}
\tightlist
\item
  It is an obstruction to an Artin-character identity, not a proof that an individual real number L′(1,ρ)/L(1,ρ) cannot accidentally equal some CM height expression.
\item
  It does not exclude infinite approximation schemes or an arithmetic construction with non-CM descent data.
\item
  It does not exclude a relative comparison in which CM heights remain auxiliary and the cubic character is represented elsewhere.
\item
  It does not say that the cubic Artin line is absent from all tensor categories generated by CM abelian varieties. AM-K004 explicitly realizes it as a Hom line between two CM elliptic eigensummands. The failure is in the CM-type class function used for the stable height, which forgets this finite descent twist.
\end{enumerate}

For odd prime r the same argument applies to the augmentation representation of C\_r⋊(Z/rZ)\^{}×. Complex conjugation acts by inversion on C\_r, and {[}c,σ{]}=σ\^{}\{-2\} generates C\_r. The augmentation representation has no C\_r invariants. The cubic case has the additional explicit ray-character and rank-one lattice constructions used below.

\subsection{4. A concrete relative theta period for the target logarithmic derivative}\label{a-concrete-relative-theta-period-for-the-target-logarithmic-derivative}

This section gives an actual function and an exact integral formula. It is independent of any proposed CM-height identification.

The primitive conductor from AM-K003 has the following scalar ideal generator m: \[
\begin{array}{c|c|c}
 \text{case}&\mathfrak f=m\mathcal O_K&Q=3N\mathfrak f=|D_{\mathbf Q(\alpha)}|\\ \hline
 p\ne3,\ p\equiv8\pmod9&m=p&3p^2\\
 p\ne3,\ p\equiv2,5\pmod9&m=3p&27p^2\\
 p=3&m=9&243.
\end{array}
\tag{R8}
\] Let Λ=O\_K=Z{[}ζ{]}⊂C in the fixed complex embedding. Its Euclidean covolume is V=√3/2, and \textbar a+bζ\textbar²=a²−ab+b².

Define the periodic function \[
 w(a)=
 \begin{cases}
 \chi((a)),&(a,\mathfrak f)=1,\\
 0,&(a,\mathfrak f)\ne1.
 \end{cases}
 \qquad a\in\Lambda.
 \tag{R9}
\] Here w(0)=0. Because K has class number one, the actual character is already given explicitly by AM-K007 on the finite unit quotient. The definition is independent of an ideal generator, invariant under the six units, and periodic modulo mΛ. As a nontrivial multiplicative character on (O\_K/𝔣)\^{}×, extended by zero, it has mean zero.

Set \[
 \Theta_\chi(t)=\sum_{a\in\Lambda}
                  w(a)e^{-\pi t|a|^2},
 \qquad t>0,
 \tag{R10}
\] and \[
 I_\chi(s)=\int_0^\infty\Theta_\chi(t)t^{s-1}\,dt.
 \tag{R11}
\]

\subsubsection{Proposition R12. Convergence and exact Mellin identity}\label{proposition-r12.-convergence-and-exact-mellin-identity}

The integral in (R11), and the integrals obtained by inserting any fixed power of log t, converge absolutely locally uniformly in s∈C. Therefore Iχ is entire. Moreover \[
 I_\chi(s)=6\pi^{-s}\Gamma(s)L_K(s,\chi)
 \tag{R12}
\] by analytic continuation. The function Θχ is real-valued.

\textbf{Proof.} For t bounded below away from zero, (R10) is absolutely and locally uniformly convergent by the positive-definite norm on Λ. Because w(0)=0, it decays exponentially as t→∞.

For the other endpoint, let Λ* be the Euclidean dual lattice, with pairing \(\langle z,a\rangle=\operatorname{Re}(z\bar a)\). The finite characters of Λ/mΛ are indexed by \[
 \xi\in m^{-1}\Lambda^*/\Lambda^*.
\] With \[
 \widehat w(\xi)=m^{-2}
  \sum_{b\in\Lambda/m\Lambda}
      w(b)e^{-2\pi i\langle\xi,b\rangle},
\] finite Fourier inversion gives \[
 w(a)=\sum_\xi\widehat w(\xi)e^{2\pi i\langle\xi,a\rangle}.
\] Its constant coefficient is zero. Poisson summation for the two-dimensional Gaussian yields \[
 \Theta_\chi(t)=\frac1{Vt}
   \sum_{\xi\ne0}\widehat w(\xi)
      \sum_{\lambda^*\in\Lambda^*}
       \exp\!\left(-\frac{\pi|\lambda^*-\xi|^2}{t}\right).
 \tag{R13}
\] Each nonzero coset has positive distance from Λ*. There are finitely many cosets. Hence Θχ(t)=O(t\^{}\{-1\}e\^{}\{-c/t\}) for some explicit c\textgreater0. These two endpoint bounds remain integrable after multiplication by t\^{}\{Re(s)−1\}\textbar log t\textbar\^{}j, uniformly on every compact set of s. Differentiation under the integral and entireness follow.

For Re(s)\textgreater1, absolute convergence permits termwise integration. The Gaussian integral is Γ(s)(π\textbar a\textbar²)\^{}\{-s\}. Each nonzero integral ideal of O\_K has exactly six generators, all of the same norm and character value. Thus the sum is six times the Hecke ideal series and gives (R12) in this half-plane, then everywhere by continuation.

Finally \(\chi(\bar{\mathfrak a})=\overline{\chi(\mathfrak a)}\) by AM-K002, while norms and the conductor are invariant under conjugation. The substitution \(a\mapsto\bar a\) conjugates the sum without changing it, so Θχ(t) is real. No pointwise positivity is asserted. ∎

\subsubsection{Explicit tail bounds}\label{explicit-tail-bounds}

The preceding proof is effective. Since \textbar ŵ(ξ)\textbar≤1 and Λ*=(2i/√3)Λ, define \[
 B(a)=\left(1+\frac{2e^{-a}}{1-e^{-3a}}\right)^2-1
 \qquad(a>0).
\] The elementary inequalities n²≥1+3(n−1) for n≥1 and a²−ab+b²≥(a²+b²)/2 give \[
 |\Theta_\chi(t)|\le B(\pi t/2),
 \qquad
 |\Theta_\chi(t)|\le\frac1{Vt}
                    B\!\left(\frac{2\pi}{3m^2t}\right).
 \tag{R14}
\] For the second bound, the disjoint nonzero dual cosets in (R13) form a subset of \(m^{-1}\Lambda^*\setminus\{0\}\); summing the absolute values over that larger lattice proves the formula. These inequalities give certified truncation of the integral and its logarithmic moments. They are not a useful arithmetic-window estimate by themselves.

\subsubsection{Corollary R15. An exact replacement scalar as a logarithmic heat moment}\label{corollary-r15.-an-exact-replacement-scalar-as-a-logarithmic-heat-moment}

The denominator Iχ(1) is positive. One has \[
 \frac{L'_K(1,\chi)}{L_K(1,\chi)}
 =\log\pi+\gamma+
   \frac{\displaystyle\int_0^\infty
                       \Theta_\chi(t)\log t\,dt}
        {\displaystyle\int_0^\infty
                       \Theta_\chi(t)\,dt}.
 \tag{R15}
\] Consequently, for any v\textgreater0, the scalar having the same displayed analytic form as the source scalar is \[
\begin{aligned}
 \log\Omega_\chi(v)
 &:=\gamma+\log(\pi/v)
        -\frac{L'_K(1,\chi)}{L_K(1,\chi)}\\
 &=-\log v-
   \frac{\displaystyle\int_0^\infty
                       \Theta_\chi(t)\log t\,dt}
        {\displaystyle\int_0^\infty
                       \Theta_\chi(t)\,dt}.
\end{aligned}
\tag{R16}
\]

\textbf{Proof.} The quotient in (R6) and the positive residue of the Dedekind zeta function at one give L\_K(1,χ)\textgreater0. Equation (R12) gives Iχ(1)=6L\_K(1,χ)/π\textgreater0. Its logarithmic derivative at one, using Γ′(1)/Γ(1)=−γ, is \[
 I_\chi'(1)/I_\chi(1)
  =-\log\pi-\gamma+L_K'(1,\chi)/L_K(1,\chi).
\] Rearrangement proves both assertions. ∎

This is a finite ray-character weighted lattice theta function with a completely specified Gaussian and measure. It changes actual summands according to χ; it is not the result of conjugating an unchanged oscillator chain by scalar phases. The Fourier version (R13) also makes it a finite linear combination of shifted dual-lattice heat sums.

The construction supplies a precise relative period to compare against an arithmetic model. It does not supply that comparison. In particular this report does not identify (R16) with any line satisfying the unchanged IsVecLine and LamDegSt conditions in §1. An integral or spectral construction with this logarithmic moment would need its own finite-place lattice, determinant formula, and estimate.

\subsection{5. Why a leading regulator alone does not supply the source scalar}\label{why-a-leading-regulator-alone-does-not-supply-the-source-scalar}

Let F=Q(α) and Q=\textbar D\_F\textbar. AM-K003 supplies the exact completion \[
 \Lambda(s)=Q^{s/2}\Gamma_{\mathbf C}(s)L_K(s,\chi),
 \qquad
 \Gamma_{\mathbf C}(s)=2(2\pi)^{-s}\Gamma(s),
 \qquad
 \Lambda(s)=\Lambda(1-s).
 \tag{R17}
\] The value Λ(1) is nonzero and finite. Therefore L\_K(s,χ) has a simple zero at zero. Write \[
 L_K(s,\chi)=a s+b s^2+O(s^3),\qquad a\ne0.
 \tag{R18}
\]

\subsubsection{Proposition R19. The required invariant is a next-coefficient ratio}\label{proposition-r19.-the-required-invariant-is-a-next-coefficient-ratio}

The exact relation is \[
 \frac{L'_K(1,\chi)}{L_K(1,\chi)}
 =-\log Q+2\log(2\pi)+2\gamma
       -\frac{L_K''(0,\chi)}{2L_K'(0,\chi)}.
 \tag{R19}
\] If h\_F and Reg(F) denote the ordinary class number and regulator, then \[
 L_K(1,\chi)=\frac{2\pi h_F\operatorname{Reg}(F)}{\sqrt Q},
 \qquad
 L_K'(0,\chi)=h_F\operatorname{Reg}(F).
 \tag{R20}
\] Moreover, for the Laurent expansion \[
 \zeta_F(s)=\frac{R_F}{s-1}+C_F+O(s-1),
 \qquad \gamma_F=C_F/R_F,
\] one has \[
 L_K'(1,\chi)/L_K(1,\chi)=\gamma_F-\gamma.
 \tag{R21}
\]

\textbf{Proof.} The expansion Γ(s)=s\^{}\{-1\}−γ+O(s) gives \[
 \frac{\Lambda'(0)}{\Lambda(0)}
   =\tfrac12\log Q-\log(2\pi)-\gamma+b/a.
\] At one, \[
 \frac{\Lambda'(1)}{\Lambda(1)}
   =\tfrac12\log Q-\log(2\pi)-\gamma
      +\frac{L_K'(1,\chi)}{L_K(1,\chi)}.
\] The logarithmic derivative of (R17) makes these two expressions negatives. Since b/a=L\_K″(0,χ)/(2L\_K′(0,χ)), this proves (R19).

F has signature (1,1) and only two roots of unity. The analytic class number formula therefore gives R\_F=2πh\_F Reg(F)/√Q. This is L\_K(1,χ) by (R6). This specialization uses Milne (2020b), V, Theorem 2.4(a), p.~156. At zero, Λ(0)=2a. At one, Λ(1)=√Q L\_K(1,χ)/π. Equality gives the second assertion in (R20).

Finally divide the displayed Laurent expansion of ζ\_F by ζ\_Q(s)=(s−1)\^{}\{-1\}+γ+O(s−1). The constant term of the quotient is R\_F, and its linear coefficient is C\_F−γR\_F. Their quotient is (R21). ∎

Thus an explicit unit or regulator realization of the leading coefficient L\_K′(0,χ) is valuable, but it is not the invariant needed in (R1) with χ inserted. Formula (R19) requires the next coefficient relative to the leading one. Equivalently (R21) requires the Euler--Kronecker constant of the non-Galois cubic field, relative to that of Q.

For comparison with the theta construction, \[
 I_\chi(0)=6h_F\operatorname{Reg}(F),\qquad
 I_\chi(1)=12h_F\operatorname{Reg}(F)/\sqrt Q.
 \tag{R22}
\] Its logarithmic moment Iχ′(1)/Iχ(1), rather than either displayed value alone, is precisely the new scalar in (R16).

\subsection{6. The actual Artin motive and its ordinary period matrix}\label{the-actual-artin-motive-and-its-ordinary-period-matrix}

There is no need to posit a new variety merely to realize the representation.

Over K, let σ generate Gal(M/K). In the convention where Γσ acts by σ, the projector \[
 e_\chi=\frac13\sum_{j=0}^2\chi(\sigma)^{-j}\Gamma_{\sigma^j}
 \tag{R23}
\] on h⁰(Spec M), with coefficients containing μ₃, gives the rank-one χ component. Reversing the cohomological action convention exchanges χ and χ\^{}\{-1\}, whose L-functions coincide here.

Over Q, let π:Spec F→Spec Q. The degree-three unit and trace maps give the rational projector \[
 e_\rho=1-\frac13\pi^*\pi_*,
 \qquad
 h^0(\operatorname{Spec}F)=\mathbf1\oplus\mathcal M_\rho.
 \tag{R24}
\] Indeed π\_\emph{π\^{}}=3 on the trivial summand, so eρ²=eρ. Its Betti realization is the trace-zero part of the three-element permutation representation, hence is ρ. Its Artin L-function is exactly (R6).

For an explicit de Rham/Betti comparison, the trace-zero de Rham space is \[
 \ker(\operatorname{Tr}_{F/\mathbf Q}:F\to\mathbf Q)
   =\mathbf Q\alpha\oplus\mathbf Q\alpha^2.
\] Evaluating on the three embeddings gives the matrix \[
 \begin{pmatrix}
  \alpha&\alpha^2\\
  \zeta\alpha&\zeta^2\alpha^2\\
  \zeta^2\alpha&\zeta\alpha^2
 \end{pmatrix}.
 \tag{R25}
\] Its columns sum to zero, and the first two rows have determinant p(ζ²−ζ)≠0. This is a fully explicit algebraic comparison for the weight-zero Artin motive.

An ordinary Artin comparison matrix such as (R25) has algebraic entries. By itself it does not prescribe the logarithmic heat moment (R16), its integral metric, or its derivative-of-L determinant formula. The leading regulator statement (R20) supplies a different arithmetic object, involving units; (R19) explains exactly why that leading datum alone is still insufficient.

There is also the positive-dimensional realization from AM-K004: an eigensummand of Hom(h¹(A₀),h¹(A\_p)) over K carries χ or its inverse. Proposition R2 and (R25) explain how a nontrivial Artin descent can coexist with unchanged stable CM height. These are compatible facts, not contradictory notions of ``CM''.

\subsection{7. What this changes in the steps 2--20 audit}\label{what-this-changes-in-the-steps-220-audit}

The complete original step table remains in kummer\_extension.md, §7. The present proof resolves the following narrower questions.

\begin{longtable}[]{@{}
  >{\raggedright\arraybackslash}p{(\columnwidth - 4\tabcolsep) * \real{0.3333}}
  >{\raggedright\arraybackslash}p{(\columnwidth - 4\tabcolsep) * \real{0.3333}}
  >{\raggedright\arraybackslash}p{(\columnwidth - 4\tabcolsep) * \real{0.3333}}@{}}
\toprule\noalign{}
\begin{minipage}[b]{\linewidth}\raggedright
Part of the attempted transfer
\end{minipage} & \begin{minipage}[b]{\linewidth}\raggedright
New exact conclusion
\end{minipage} & \begin{minipage}[b]{\linewidth}\raggedright
Remaining obligation
\end{minipage} \\
\midrule\noalign{}
\endhead
\bottomrule\noalign{}
\endlastfoot
CM geometric supply and character descent & The explicit polarized cubic twists and induced CM tower exist; their actual stable tangent letters are isometric to the untwisted ones. & Compute descent-sensitive finite indices and the genuinely character-weighted arithmetic quantity. \\
The height term used in MAIN & The unchanged source record fixes a quadratic dlogL and actual degrees. The cubic S₃ character cannot be obtained by a finite CM-type height Artin-class identity. & A new comparison theorem, or an auxiliary-height argument that inserts the cubic character outside that scalar. \\
Ideal/ray-character analytic side & Θχ and its Mellin transform give the exact cubic L-function, including the wild p=3 case and primitive conductor. & Relate this actual relative period to the chosen arithmetic sections and lattice estimates. \\
Motivic realization & The rank-two Q-Artin projector and rank-one K-projector are explicit; a CM elliptic Hom realization is already available. & A period/metric/determinant construction representing the logarithmic moment, with all finite terms. \\
Leading special values & L(1), L′(0), and the ordinary unit regulator are exactly related. & Control L″(0)/(2L′(0)), or equivalently the relative Euler--Kronecker invariant, in the required arithmetic comparison. \\
Final analytic implication & AM-K005 already supplies the field-independent Landau/window adapter. & Prove its full family of arithmetic window bounds. None follows merely from the formulas above. \\
\end{longtable}

The relative Mellin construction is therefore a concrete alternative target, rather than an unspecified ``Hecke analogue of Chowla--Selberg''. Its input is the already explicit ray character, its output is exactly the required logarithmic derivative, and its normalization is fixed in (R10)--(R16). The open part is an arithmetic realization and estimate of this period that meets the actual source hypotheses.

\subsection{8. Sources and verification scope}\label{sources-and-verification-scope}

\begin{itemize}
\tightlist
\item
  Alpöge, L. (2026). Author Lean archive, source locations in §1. The local mirror is the parent-recovered source set. This report inspects the listed definitions and statements; it makes no new whole-archive compilation claim.
\item
  Milne, J. S. (2020a). \emph{Complex multiplication}. I, Proposition 1.4, Corollary 1.5 and Remark 1.6, p.~10. \url{https://www.jmilne.org/math/CourseNotes/CM.pdf}. Used only for the maximal CM quotient and central conjugation facts; the obstruction in Proposition R7 is proved here.
\item
  Milne, J. S. (2020b). \emph{Class field theory}. V, §2, especially Theorem 2.4(a)--(b), pp.~155--156. \url{https://www.jmilne.org/math/CourseNotes/CFT.pdf}. Used for the Hecke ideal convention, nonvanishing at one and the analytic class number formula.
\item
  Yang, T. (2010). The Chowla--Selberg formula and the Colmez conjecture. \emph{Canadian Journal of Mathematics, 62}(2), 456--472. \url{https://doi.org/10.4153/CJM-2010-028-x}. Author manuscript: \url{https://people.math.wisc.edu/~tonghaiyang/Colmez.pdf}. §3, equations (3.1)--(3.12), author pp.~9--10, defines the CM-type class functions on the maximal CM quotient. No unproved general Colmez formula is invoked in Proposition R7.
\item
  Internal proved prerequisites: AM-K001--AM-K005 and AM-K007--AM-K011, at the companion paths identified at the beginning. All new analytic identities in §§4--6 are proved above with explicit normalizations.
\end{itemize}

\textbf{Completed author checks.} The saved report was read in full and its arguments reconstructed: stable polarized tangent-line isometry; the common-Galois-group averaging in (R7); conductor and unit factors; both finite Fourier normalizations and Gaussian tail bounds; the Mellin logarithmic derivative; both gamma-factor expansions in (R19); the class-number and regulator constants; and the Artin projectors and trace-zero comparison matrix. The listed source interfaces and the cited primary definitions were inspected directly. No numerical or finite computational certificate was run for this report. The continuation in (R12) follows from its convergence proof, not from numerical agreement. The author check and source-use events are recorded in \nolinkurl{kummer_extension_CHECK_EVENTS.tsv.}

\clearpage

\section{AM-DO-027: exact common-chart phases and an integral unit normalization}\label{am-do-027-exact-common-chart-phases-and-an-integral-unit-normalization}

Author: GPT-6 Astra Pro, ALP\_MACHINE \nolinkurl{/root/definite_oscillator}; inherited reasoning effort not separately exposed; 21 November 2026 UTC.

This result closes the scalar compatibility in \nolinkurl{PinColT.cst} once the raw arithmetic finite datum has actually been realized. The normalization is permitted by the literal \nolinkurl{IsRatDatum} definition, not by changing the vacuum or inventing a new meaning for that predicate. On equal-volume charts with compact determinants it uses only roots of unity. It therefore has no additional finite denominator, content, or norm cost. An explicit example also shows why unramifiedness of the CM extension does not make these character phases disappear automatically.

\subsection{1. The precise common-character scalar equation}\label{the-precise-common-character-scalar-equation}

Fix one even-rank edge \(a\to b\) of the common \(E'=\mathbf Q(\zeta_{28})\) host, with the same actual splitting character \(\mu\) on both members, as required by \nolinkurl{PinColT.om}. Its restriction to \(\mathbb A_{F'}^\times\) is trivial. Write \[
 \nu_\mu(t)=\mu(x),\qquad x/\bar x=t.
 \tag{DO27.1}
\] Hilbert 90 and the displayed restriction make this a well-defined character of the norm-one torus. For source and target chart base elements \(k_\alpha,k'_\gamma\), define \[
 u_\alpha=\nu_{\mu,f}(\det k_\alpha)=\mu(x_\alpha),
 \qquad
 u'_\gamma=\nu_{\mu,f}(\det k'_\gamma)=\mu(x'_\gamma).
 \tag{DO27.2}
\] The roots are exactly the source's \nolinkurl{xr}, not the determinants themselves. Its \nolinkurl{TwDat.Is.kA} requires these root equations. Let \(v_\alpha\) be the actual source chart level volume. The literal condition is \[
        c_K=v_\alpha w_{\alpha\gamma}c_0
                         u_\alpha(u'_\gamma)^{-1}
                   \quad\text{for every }\alpha,\gamma,
 \qquad c_0=C^{2ab}.
 \tag{DO27.3}
\] This is \nolinkurl{LevChiEqOne.lean}, \nolinkurl{PinColT.cst}, lines113--114.

For unchanged raw scalars \(w_{\alpha\gamma}=1\), a common nonzero \(c_K\) exists if and only if

\begin{itemize}
\tightlist
\item
  all source volumes \(v_\alpha\) are equal;
\item
  \(u_\alpha\) is constant on the source chart set;
\item
  \(u'_\gamma\) is constant on the target chart set.
\end{itemize}

Here both chart sets are assumed nonempty. To prove necessity, fix a target chart and compare two source charts in (DO27.3). The absolute values give equality of their positive volumes, because \(|u_\alpha|=1\); the phases then give equality of the \(u_\alpha\). Fixing a source chart similarly proves the target assertion. Conversely those three conditions make the right side independent of the charts. This is a necessary and sufficient criterion, not merely a sufficient normalization rule.

Equivalently, the determinant conditions are \[
 \det(k_\alpha k_{\alpha_0}^{-1})\in\ker\nu_{\mu,f},
 \qquad
 \det(k'_\gamma(k'_{\gamma_0})^{-1})\in\ker\nu_{\mu,f}
                         \quad\text{for all }\alpha,\gamma.
 \tag{DO27.4}
\] If the determinant representatives are compact and the finite character is trivial on the maximal norm-one compact, these conditions hold automatically. Relative unramifiedness of \(E'/F'\) alone imposes no such triviality on the chosen character.

\subsection{2. Linearity is present in the actual source definitions}\label{linearity-is-present-in-the-actual-source-definitions}

The relevant definitions were read directly, rather than inferred from their names. In \nolinkurl{Formalism/LocalToGlobal/IndefiniteReal/LocalNewformFunction.lean}, lines238--261, the source has:

\begin{enumerate}
\def\labelenumi{\arabic{enumi}.}
\tightlist
\item
  \nolinkurl{IsSBHalf e Λ Φ}: support in the idempotent half, a common integral period, and a common denominator support condition.
\item
  \nolinkurl{IsPlaceModel … f₀ T}: \(T\) is explicitly a complex-linear map on functions, intertwines translations and modulations, and satisfies the unchanged vacuum normalization \(T(c\,\mathrm{halfInd})=f_0\).
\item
  \nolinkurl{IsThetaAt … f₀ Φ f}: \nolinkurl{AllSome} of those models with the assertion \(f=T\Phi\).
\end{enumerate}

\nolinkurl{IsRatDatum}, at lines25--29 of its namesake file, is the existence of such \(\Phi\), with its rational restriction equal to the specified arithmetic coefficient function \(\Xi\), and with that \nolinkurl{IsThetaAt} assertion. It follows that, for any complex scalar \(z\), \[
 \boxed{\operatorname{IsRatDatum}(\Xi,f_0,f)
       \Longrightarrow
       \operatorname{IsRatDatum}(z\Xi,f_0,zf).}
 \tag{DO27.5}
\] The level idempotent, lattice, vacuum, additive character, and \nolinkurl{IsPlaceModel} normalization are unchanged.

\textbf{Proof.} Use \(z\Phi\) as the new witness. Its support is contained in that of \(\Phi\), its same integral period works, and its same denominator works. Thus it remains \nolinkurl{IsSBHalf}. Its rational restriction is \(z\Xi\). Keep exactly the existing models \(T\); each is linear, so \(T(z\Phi)=zT\Phi=zf\). This proves the new assertion for every such model as well as for its existing witness, regardless of the \nolinkurl{AllSome} packaging. The condition \(T(c\,\mathrm{halfInd})=f_0\) is untouched. This proves (DO27.5). \(\square\)

For a nonzero \(z\), the support of the finite coefficient family is unchanged. \nolinkurl{JointSmooth.lean}, lines20--24, defines \[
 \operatorname{comb}(d,f_0)=\sum_xd_x\operatorname{heis}(x)f_0,
 \tag{DO27.6}
\] and \nolinkurl{IsFinComb} says that every support point has zero infinite part. Consequently \(\operatorname{comb}(zd,f_0)=z\operatorname{comb}(d,f_0)\), and \nolinkurl{IsFinComb} is retained. In particular this is a permitted rescaling of actual finite Schwartz inputs. Any fixed compact eigencharacter or real covariance of the input is also unchanged.

\subsection{3. Actual normalized finite inputs and the common constant}\label{actual-normalized-finite-inputs-and-the-common-constant}

Assume that for every chart pair there is an \textbf{actual raw datum} \[
 \operatorname{IsRatDatum}(\Xi_{\alpha\gamma},f_0,
                         \phi^0_{\alpha\gamma}),
 \qquad
 \phi^0_{\alpha\gamma}
            =\operatorname{comb}(d^0_{\alpha\gamma},f_0),
 \tag{DO27.7}
\] with the specified finite support condition. The arithmetic \(\Xi_{\alpha\gamma}\) in this assumption is the source's \nolinkurl{A.xi}; it has not been replaced by a freely chosen function. Its existence is an independent input to this proposition.

Fix any reference positive volume \(v_*\), and set \[
 w_{\alpha\gamma}
   =\frac{v_*}{v_\alpha}\,u_\alpha^{-1}u'_\gamma,
 \quad
 d_{\alpha\gamma}=w_{\alpha\gamma}d^0_{\alpha\gamma},
 \quad
 \phi_{\alpha\gamma}=w_{\alpha\gamma}
                                  \phi^0_{\alpha\gamma}.
 \tag{DO27.8}
\] Then (DO27.5)--(DO27.6) give exactly \nolinkurl{PinColT.dat}, with its required scalar multiple \(w_{\alpha\gamma}\Xi_{\alpha\gamma}\). Substitution in (DO27.3) gives, for every chart pair, \[
                          \boxed{c_K=v_*c_0.}
 \tag{DO27.9}
\] The common period is unchanged. No new choice of the source vacuum or of its unit norm is made. This proves the simultaneous \nolinkurl{dat} and \nolinkurl{cst} assertions from (DO27.7), not only a formal consistency of the numbers in \nolinkurl{cst}.

Suppose now that the source levels are conjugate to one fixed level inside the finite unitary group. A reductive local group is unimodular, so conjugation preserves Haar measure; their volumes are all the same \(v_*\). Suppose also that all the determinants in (DO27.2) belong to the maximal compact norm-one torus. The restriction of a continuous finite-adelic unitary character to that compact group has finite image. Indeed the compact is profinite, its character image is a compact profinite subgroup of the circle, and such a subgroup is finite. Therefore \[
                 w_{\alpha\gamma}=u_\alpha^{-1}u'_\gamma
                           \quad\text{is a root of unity.}
 \tag{DO27.10}
\] It follows exactly that \[
 \|\phi_{\alpha\gamma}\|_2
       =\|\phi^0_{\alpha\gamma}\|_2,
 \qquad
 v_{\mathfrak p}(w_{\alpha\gamma})=0
                              \quad\text{at every finite prime}.
 \tag{DO27.11}
\] Multiplication by this scalar is an invertible integral map on every coefficient lattice containing the character values. It preserves its content, all elementary divisors, and its saturation index. Every complex embedding also has absolute value one. These are equalities, not upper estimates.

There is no assertion that a later selected Schur projection has unit content. The statement concerns this particular scalar normalization of the existing chart inputs. Its root-of-unity change neither cures nor worsens a previously uncomputed lattice or projection denominator.

If the chart volumes differ, the exact remaining extra scalar is \(v_*/v_\alpha\). For commensurable compact opens it is rational: with \(C_\alpha=U_*\cap U_\alpha\), \[
 \frac{v_*}{v_\alpha}
        =\frac{[U_*:C_\alpha]}{[U_\alpha:C_\alpha]}.
 \tag{DO27.12}
\] The norm changes by that positive ratio and its local content changes by its valuation. Thus the character phase and the volume cost are separated by an exact formula. Neither is discarded when equal-volume hypotheses have not been verified.

On the present host compact-determinant representatives are available. R009 proves the norm-one finite torus has class number one. If \(g\) is any finite unitary representative, write \(\det g=\delta u\), with \(\delta\in T'(F')\) and \(u\) compact. Left multiplication by the rational unitary matrix \(\operatorname{diag}(\delta^{-1},1,\ldots,1)\) makes its determinant compact without changing its finite double class. This observation does not by itself prove that an independently pinned arithmetic chart has already made this representative choice; its frames must be transported along that rational action.

\subsection{4. An explicit nonconstant phase on unchanged lattices}\label{an-explicit-nonconstant-phase-on-unchanged-lattices}

The following example uses the same character on the two members of an even-rank edge and the same unramified CM host. It shows that the root-of-unity normalization in section 3 can be necessary.

Let \(a=6,b=20\). On the six real torus coordinates indexed by \(\Phi'=\{1,3,9,15,17,23\}\), prescribe the character \[
                         \beta(t)=\prod_{\tau\in\Phi'}t_\tau^{-8}.
 \tag{DO27.13}
\] At one split place \(v_*\) above 29 choose the first component so that \(\zeta_{28}\mapsto2\bmod29\), and define its residue character by \[
                         \rho(2)=\zeta_{28}^{12}.
 \tag{DO27.14}
\] Take the other finite compact components trivial. Since \(\sum\Phi'=68\), the value of (DO27.13) on the rational \(\zeta_{28}\) is \(\zeta_{28}^{-544}=\zeta_{28}^{16}\). Thus (DO27.14) is exactly its inverse on \(\mu_{28}\). The torus class-number-one decomposition from R009 constructs an actual unitary automorphic character \(\gamma\) with these infinity and finite compact restrictions. Put \[
                  \mu=\gamma\circ(z\mapsto z/\bar z).
 \tag{DO27.15}
\] Then \(\mu|_{\mathbb A_{F'}^\times}=1\), and \(\mu_\tau(z)=(z/\bar z)^{-8}\) at every selected place. This is an admissible same-character splitting for the two even-rank members. It also has an actual source \nolinkurl{HCParam} instance with admissible numerical \nolinkurl{Ranks} data: \(k=6,t=13,m=2,\mathrm{ord}=3\), \(u=(11/2,9/2,7/2,5/2,3/2,1/2)\), \(v=(5,4,3,2,1),b=(0,16),e=(8,24)\). The required interlacings hold, the final \(u\)-entry is positive, and both successive gaps equal \((t+3)/2=8\). The numerical rank conditions are also immediate: \(m+3=5\le6\), \(k+m-1+5=12\le13\), \(3m=6<13\), \(m\le2\mathrm{ord}=6\le k\), and \(m^2=4\le34\). No claim about further \nolinkurl{Hyp3} hypotheses is needed for this phase example.

Let the source charts use the common principal congruence level modulo 29 at all six places above 29 and the standard maximal compact elsewhere. Take two source base elements \(k_0=1\) and \(k_1\), where at \(v_*\) \[
                 k_1=\operatorname{diag}((2,2^{-1}),1,1,1,1,1)
                        \in U(6)(F'_{v_*}),
 \tag{DO27.16}
\] under the relative split identification, and at all other places \(k_1=1\). Both lie in the standard integral compact and hence preserve the identical selfdual lattice. Their source levels are equal, so their volumes are equal. Choose the target identity base chart. The determinant roots give \[
                 u_0=1,\qquad u_1=\zeta_{28}^{12}
                         =\zeta_7^{3},\qquad u'_0=1.
 \tag{DO27.17}
\] The two source level classes are distinct. If \(k_1=qu\), with \(q\in U(6)(F')\) and \(u\) in the stated principal level, then \(q\) is integral at every finite place. Its determinant is therefore a norm-one unit, hence a member of \(\mu_{28}\). At any other place over 29 that determinant reduces to 1. Reduction there is injective on \(\mu_{28}\), so \(\det q=1\). At \(v_*\) the congruence would instead give \(\det q=2\bmod29\), a contradiction.

For unscaled raw inputs, (DO27.3) asks simultaneously for \[
                         c_K=vc_0,
             \qquad      c_K=\zeta_7^{3}vc_0,
 \tag{DO27.18}
\] which is impossible since \(v>0,c_0\ne0\). This failure is a real character phase on actual different level classes, despite unchanged lattices. The underlying standard trivial-character vacuum is fixed, while the prescribed common-character source Weil action is \[
                    \omega_\mu(k_1,1)\phi_0
                           =\zeta_7^{-3}\phi_0.
 \tag{DO27.19}
\] Equation (DO27.19) is the exact inverse source determinant-root factor of AM-DO-022. The theta reading is injective and its entire vacuum germ is nonzero, so this is an actual kernel phase, not just a formal determinant computation. A sufficiently high polynomial jet detects it even if the origin value of the section vanishes.

The normalization of section 3 sets \(w_{00}=1\) and \(w_{10}=\zeta_7^{-3}\). It solves (DO27.18) in the literal \nolinkurl{PinColT.cst} convention, retaining the actual determinant-root factors in the kernel. Its scalar is already in \(E'\), is a unit at every finite prime, and has absolute value one at every complex embedding. The example is therefore a counterexample to automatic phase constancy, together with an exact permitted repair. It is not a counterexample to existence of \nolinkurl{PinColT}.

\subsection{5. Uniform character phases for the selected exponents}\label{uniform-character-phases-for-the-selected-exponents}

The same one-place construction works for arbitrary integer selected exponents \(-e_{j,\tau}\). Give the real torus that infinity type, and at the chosen place over 29 put \[
 \rho_j(2)=\zeta_{28}^{r_j},\qquad
            r_j\equiv\sum_{\tau\in\Phi'}\tau e_{j,\tau}
                                               \pmod{28}.
 \tag{DO27.20}
\] The inverse root-unit condition follows by evaluating at \(\zeta_{28}\), exactly as in section 4. Extend it using the same torus class decomposition and compose with \(z/\bar z\). This produces actual admissible equal-character edges with the selected infinity exponents. For compact-determinant chart representatives every scalar in (DO27.8) then lies in \(\mu_{28}\subset E'\), uniformly in the degrees and ranks. There is no new coefficient-field extension for these phases.

If a particular source instance has already fixed different characters \(\mu_j\), the construction in (DO27.20) does not replace them. Sections 1--3 apply to those characters as they stand; their actual compact values determine the unit scalars. Likewise, normalizing each edge does not establish the separate row-lattice, Schur-projection, or terminal-mass identities. It does establish that, once the raw finite datum is available, the common-chart character phases in \nolinkurl{PinColT.cst} can be removed within the expressly permitted scalar parameter, with the exact volume cost (DO27.12) and no additional unit cost.

\subsection{Exact source and carried-result locators}\label{exact-source-and-carried-result-locators}

\begin{itemize}
\tightlist
\item
  \nolinkurl{Formalism/LocalToGlobal/IndefiniteReal/LocalNewformFunction.lean}, \nolinkurl{IsSBHalf}, lines238--241; \nolinkurl{IsPlaceModel}, lines248--255; \nolinkurl{IsThetaAt}, lines257--261; \nolinkurl{splitFacO}, line51.
\item
  \nolinkurl{Development/AbelianVarietyTheory/Definitions/LatticeFrames/IsRatDatum.lean}, definition \nolinkurl{IsRatDatum}, lines25--29.
\item
  \nolinkurl{Development/AbelianVarietyTheory/Definitions/LatticeFrames/JointSmooth.lean}, \nolinkurl{comb} and \nolinkurl{IsFinComb}, lines20--24.
\item
  \nolinkurl{Development/AbelianVarietyTheory/Definitions/LatticeFrames/LevChiEqOne.lean}, \nolinkurl{TwDat.Is.kA}, lines52--53; \nolinkurl{PinColT.dat} and \nolinkurl{cst}, lines111--114.
\item
  AM-DO-022, exact standard finite splitting and determinant roots; AM-DO-025, selected infinity exponents; AM-DO-026, actual kernel normalization and chart equation; corrected AM-R-009, the common-host norm-one torus class decomposition and character extension. The published class-number input is carried through R009, not re-proved here.
\end{itemize}

No Lean compilation or complete \nolinkurl{PinColT} instantiation is claimed. The raw arithmetic datum (DO27.7) and the remaining fields of that record are explicitly retained as separate requirements.

\clearpage

\section{AM-R-011: the selected Schur norms, their row normalization, and a finite projection denominator}\label{am-r-011-the-selected-schur-norms-their-row-normalization-and-a-finite-projection-denominator}

Author: /root, GPT-6 Astra Pro, ALP\_MACHINE, 21 November 2026 UTC.

This result inserts the \textbf{actual selected partitions} of AM-DO-025 into the Fischer metric and the coefficient normalization used by the recovered source. It computes a concrete part of the archimedean cost, and a valid integral denominator for projection onto a fixed finite polynomial component. It does not identify that denominator with the least denominator of the source row map, or compute a full DP44 gain.

\subsection{1. The two partitions that must be distinguished}\label{the-two-partitions-that-must-be-distinguished}

At one selected CM embedding use the notation of AM-DO-025:

\[
 a=N_j,\quad b=N_{j+1}=a+d,\quad
 w_j=\operatorname{lamTrue}_j,\quad
 \nu_j=w_j+e_j-\frac b2\mathbf1_a.
 \tag{R11.1}
\]

The source coefficient representation uses the partition obtained by removing the final weight:

\[
 \eta_j=w_j-w_{j,a-1}\mathbf1_a,
 \qquad q_j=\nu_{j,a-1}.
 \tag{R11.2}
\]

Thus \(\eta_j\) is literally \nolinkurl{DomWeight.natW}, and \(w_{j,a-1}\) is \nolinkurl{detShift}. AM-DO-025 proves \(q_j\ge1\) and \(\nu_j\) is an integral partition. Consequently

\[
 \boxed{\nu_j=\eta_j+q_j\mathbf1_a,
 \qquad \eta_{j+1}=(\nu_j,0^d).}
 \tag{R11.3}
\]

For the second identity, the last block of the next weight has value \(-e_j+a/2\). Every old entry exceeds that value by \(\nu_{j,i}>0\), by the exact recurrence in AM-DO-025. Subtracting the last entry therefore leaves precisely the displayed partition. This proves (R11.3) from the actual weights, without choosing a different representation. The partition \(\eta_j\) describes the row coefficient metric; \(\nu_j\) describes the homogeneous kernel polynomial. They generally have different degrees.

For a partition \(\lambda\), put

\[
 H_\lambda=\prod_{(i,l)\in\lambda}
       (\lambda_i-l+\lambda_l'-i+1),\qquad H_\varnothing=1.
 \tag{R11.4}
\]

These are ordinary hook products, not the DP44 height constant \(H\). Trailing zero rows do not change them.

\subsection{2. The normalized Fischer calculation}\label{the-normalized-fischer-calculation}

On a complex \(b\)-by-\(a\) matrix variable \(Z\), use the Fischer inner product with distinct monomials orthogonal and

\[
 \|Z^\alpha\|_{\rm F}^2=\alpha!.
 \tag{R11.5}
\]

For \(\lambda_1\ge\cdots\ge\lambda_a\ge0\), let \(\lambda_{a+1}=0\), and define the primitive polynomial

\[
 p_\lambda(Z)=\prod_{r=1}^a
       \det(Z_{[1,r],[1,r]})^{\lambda_r-\lambda_{r+1}}.
 \tag{R11.6}
\]

Its leading diagonal monomial has coefficient one. In particular it is primitive in the ordinary integral monomial lattice. The exact norm is

\[
                         \boxed{\|p_\lambda\|_{\rm F}^2=H_\lambda.}
 \tag{R11.7}
\]

Here is a short derivation fixing the normalization, rather than a new proof of Schur--Weyl duality. Let \(D=|\lambda|\), and let \(f_\lambda\) be the dimension of the corresponding symmetric-group representation. Schur--Weyl duality and the hook formula give \(D!/f_\lambda=H_\lambda\). The exact references are Fulton and Harris (1991), Theorem 6.3(2)--(4), pp.77--78; Hook Length Formula 4.12, p.50; and the rectangular Cauchy decomposition in Exercise 6.11(b), p.80. These results are cited, not re-proved.

Choose the standard unitary metrics on the Schur modules, with highest vectors of norm one, compatible under the standard isometric inclusion \(\mathbf C^a\subset\mathbf C^b\). The degree-\(D\) Fischer reproducing kernel satisfies

\[
 \begin{aligned}
 \frac{(\operatorname{tr}(ZW^*))^D}{D!}
 &=\sum_{|\lambda|=D}\frac{f_\lambda}{D!}
                  \operatorname{tr}(S_\lambda(Z)S_\lambda(W)^*)\\
 &=\sum_{|\lambda|=D}\frac1{H_\lambda}
       \sum_{r,s}p_{rs}^\lambda(Z)\overline{p_{rs}^\lambda(W)}.
 \end{aligned}
 \tag{R11.8}
\]

The first equality is the trace of \((ZW^*)^{\otimes D}\) under the cited tensor decomposition. Functoriality and the unitary metrics give \(S_\lambda(W^*)=S_\lambda(W)^*\); this also follows from singular-value decomposition and the diagonal action on weights. The second equality expands the trace in orthonormal bases. The matrix coefficients \(p_{rs}^\lambda\) are linearly independent and exhaust the Cauchy summand. Comparison with the reproducing kernel therefore gives their common squared norm \(H_\lambda\) and their orthogonality. The highest-to-highest coefficient is (R11.6): both are highest vectors and evaluate on a diagonal matrix as \(\prod_i z_i^{\lambda_i}\), with coefficient one. This proves (R11.7), including its scalar.

\subsection{3. The determinant insertion has an exact scalar}\label{the-determinant-insertion-has-an-exact-scalar}

Let \(V_\eta\) be the irreducible compact coefficient module realized as the row orbit of \(p_\eta\) in square \(a\)-by-\(a\) polynomials. Multiplication by the full determinant gives

\[
 M_q:V_\eta\longrightarrow V_{\eta+q\mathbf1_a},
                         \qquad f\longmapsto(\det Z)^q f.
 \tag{R11.9}
\]

This is an integral polynomial map. On \(U(a)\), it intertwines the two representations after the unitary determinant twist. Hence its pulled-back Fischer form is a scalar multiple of the invariant form on \(V_\eta\). Evaluating on the highest vector and using (R11.7) proves, for every \(f\in V_\eta\),

\[
 \boxed{\|M_qf\|_{\rm F}^2
      =\frac{H_{\eta+q\mathbf1_a}}{H_\eta}\|f\|_{\rm F}^2.}
 \tag{R11.10}
\]

The ratio has the explicit integral product

\[
 \frac{H_{\eta+q\mathbf1_a}}{H_\eta}
       =\prod_{i=1}^{a}\frac{(\eta_i+q+a-i)!}{(\eta_i+a-i)!}
       =\prod_{i=1}^{a}\prod_{r=1}^{q}(\eta_i+a-i+r).
 \tag{R11.11}
\]

Indeed the factorial/Vandermonde expression for the hook product has the same Vandermonde denominator before and after the uniform shift by \(q\). Cancelling it gives (R11.11). This is also the hook-product expression directly used in the recovered source \nolinkurl{Divisibility/RecBridgeNum.lean}, theorem \nolinkurl{logHook_eq}, lines149--188. The displayed ratio is strictly positive.

The source defines its row metric as Fischer divided by the squared Fischer norm of its extremal vector: \nolinkurl{DataCoeff.lean}, \nolinkurl{normFock}, lines83--84, and \nolinkurl{HermitianForms/IsLocBasis.lean}, \nolinkurl{nfAll}, lines19--20. Thus on the actual row partition it is

\[
             \|f\|_{\mathrm{row},j}^2
                         =\|f\|_{\rm F}^2/H_{\eta_j}.
 \tag{R11.12}
\]

Combining (R11.10) and (R11.12), the selected determinant insertion has squared Fischer norm \(H_{\nu_j}\|f\|_{\mathrm{row},j}^2\). In normalized coefficient metrics it is an isometry. This last assertion concerns the coefficient map (R11.9); it is not a theta-operator isometry or an identity between automorphic measures.

In particular its hook ratios obey the exact cancellation

\[
 \boxed{\prod_{j=0}^{m-1}\frac{H_{\nu_j}}{H_{\eta_j}}
              =\frac{H_{\eta_m}}{H_{\eta_0}}.}
 \tag{R11.13}
\]

This follows from (R11.3). It prevents counting the same row-metric normalization repeatedly as an independent loss. No finite projection denominator, class mass, branch map, theta value, or DP44 term is asserted to cancel by (R11.13).

\subsection{4. Actual common-host Fock norm}\label{actual-common-host-fock-norm}

On \(E'=\mathbf Q(\zeta_{28})\), use the exact period and coordinates of AM-R-009. Each of the six coordinates is \(w_i^+=u_i+iv_i\) or \(w_i^-=u_i-iv_i\), and the actual creation operator is the one read by \(\partial/\partial w_i^\pm\). The monomial creation norm relative to the unit Gaussian is \(\alpha!(\pi/(2G_i))^{|\alpha|}\). This is the explicit orthogonal coordinate change in AM-R-009, retaining the factor \(1/2\) in each derivative. Let \(i(\tau)\in\{1,2,3\}\) name the underlying Jacobian coordinate of \(\tau\).

For a selected partition \(\nu_{j,\tau}\) at every real place, set \(D_{j,\tau}=|\nu_{j,\tau}|\). The tensor product of the primitive highest polynomials (R11.6), applied to the \textbf{actual} Gaussian, has exact squared archimedean norm

\[
 \boxed{\|\phi_j^{\mathrm{hw}}\|_\infty^2
  =\det(2Y)^{-ab}
       \prod_{\tau\in\Phi'}
       H_{\nu_{j,\tau}}
       \left(\frac{\pi}{2G_{i(\tau)}}\right)^{D_{j,\tau}}.}
 \tag{R11.14}
\]

There are \(ab\) copies of \(A=J^2\), so the squared vacuum norm is \(\det(2Y)^{-ab}\). At a fixed CM place all matrix entries have the same creation scale; homogeneity supplies its degree power and (R11.7) supplies its hook product. The CM coordinate blocks are orthogonal. These three facts prove (R11.14). Its square root is the norm, and its logarithm is

\[
 -\frac{ab}{2}\log\det(2Y)
   +\frac12\sum_\tau\left[
      \log H_{\nu_{j,\tau}}+
      D_{j,\tau}\log\frac{\pi}{2G_{i(\tau)}}\right].
 \tag{R11.15}
\]

The corresponding analytic-to-algebraic kernel still has period \(C^{2ab}\), by AM-DO-026. Unit normalization divides that kernel by the square root of (R11.14), and does not delete the period. Finite-vacuum norm one was used only for the unmoved standard vector; an additional finite linear combination retains its actual finite norm. A compact coefficient isometric embedding also retains the source Schur factor \(1/\dim V_{\eta_j}\) in Haar integration, as proved in AM-DO-026.

\subsection{5. A finite-degree integral projection bound}\label{a-finite-degree-integral-projection-bound}

Let \(L_D=\operatorname{Sym}^D(\mathbf Z^b\otimes(\mathbf Z^a)^*)\) be the ordinary integral polynomial lattice. Let \(E_\lambda\) be the rational projector onto its Cauchy component of shape \(\lambda\), and let \(L_\lambda=L_D\cap\operatorname{im}(E_\lambda)\). Then

\[
 \boxed{L_\lambda\subset E_\lambda L_D
                \subset H_\lambda^{-1}L_\lambda.}
 \tag{R11.16}
\]

To prove this without assuming integral semisimplicity, start on the tensor power \((\mathbf Z^b)^{\otimes D}\otimes((\mathbf Z^a)^*)^{\otimes D}\). Act on its first factor by the central symmetric-group idempotent

\[
 e_\lambda=\frac{f_\lambda}{D!}
             \sum_{\sigma\in S_D}\chi_\lambda(\sigma^{-1})\sigma
       =\frac1{H_\lambda}
             \sum_{\sigma\in S_D}\chi_\lambda(\sigma^{-1})\sigma.
 \tag{R11.17}
\]

The characters of the symmetric group are integers, and the central numerator commutes with the simultaneous permutation of both tensor factors. It therefore descends to the integral coinvariant quotient \(L_D\). Its denominator divides \(H_\lambda\). Over \(\mathbf Q\), Schur--Weyl duality identifies the descended operator with exactly the stated Cauchy projector. Hence \(H_\lambda E_\lambda L_D\subset L_\lambda\). Conversely it fixes every vector of \(L_\lambda\), proving the other inclusion in (R11.16).

At a prime \(p\), the possible saturation quotient is consequently killed by \(p^{v_p(H_\lambda)}\); its length is at most \(\operatorname{rank}(L_\lambda)v_p(H_\lambda)\). If every hook is a \(p\)-unit, the projection is integral and split. In particular \(p>\lambda_1+a-1\) suffices. This is an explicit upper bound for one projector, not a least-denominator formula. The source's particular Weyl/dual-Weyl lattice can admit a strictly better comparison. For example multiplication by \((\det Z)^q\) is already integral, so a large uniform shift does not by itself force the denominator bound in (R11.17) to be attained on that determinant-inserted submodule.

This finite-degree statement is consistent with AM-K019. It makes no assertion about a lattice simultaneously stable under all ordinary differential powers. Nor does a bound for a projector prove that the theta jet matrix on its image has a prescribed content or that a prescribed automorphic row belongs to that image.

\subsection{6. A concrete source-rank illustration}\label{a-concrete-source-rank-illustration}

Take the smallest displayed rank choice \(k=6,m=2,t=13,d=14\), and saturate only the \nolinkurl{HCParam} gaps as in AM-DO-025. Then

\[
 N_0=6,\quad N_1=20,\quad N_2=34,\qquad
 \nu_0=(1^6),\quad \nu_1=(3^6,2^{14}),
 \quad D_0=6,\quad D_1=46.
 \tag{R11.18}
\]

The row partitions are \(\eta_0=0\), \(\eta_1=(1^6,0^{14})\), and \(\eta_2=(3^6,2^{14},0^{14})\). Thus \(H_{\nu_0}=6!=720\), and the second determinant insertion ratio is

\[
 \frac{H_{\nu_1}}{720}
 =\prod_{i=1}^{20}\prod_{r=1}^{2}
                        (\eta_{1,i}+20-i+r).
 \tag{R11.19}
\]

All integers, prime valuations and the product identity for this example are reproduced by \nolinkurl{selected_schur_metric_certificate.py}. The small-polynomial checks there also expand (R11.6) in ordinary monomials and compare its exact Fischer norm to (R11.4); they are finite checks of normalization, not the proof of the general formula. The rank choice and gaps satisfy the conditions named here. No claim is made that they satisfy all the further noncompact \nolinkurl{Hyp3} conditions or already form a nonzero global chain.

\subsection{7. Framework and retained arithmetic gate}\label{framework-and-retained-arithmetic-gate}

The polynomial lattices and determinant maps are defined over \(\mathbf Z\); the rational projector in (R11.17) specifies exactly which primes need inverting for that chosen splitting. The associated Schur operations live in the representation Grothendieck ring and the special lambda-ring formalism of SGA6, Exposé V. The actual normalized metric and (R11.16) are stronger data than an equality of virtual classes. The representation ring is related to the flag construction by SGA6, Exposé VI, Theorem3.3, with the carried-source qualification in the geometric dictionary. The CM eigenspace splitting is the diagonalizable-group grading of SGA3, Exposé I, Proposition4.7.3. The period in (R11.14) remains the line-bundle metric of the Picard/theta construction, placed in FGA232 Theorem3.1 and SGA7 I VII2.7/2.9 in that dictionary.

For DP\_MULT Q3e these formulas supply the selected homogeneous degrees, an exact archimedean creation norm and its relation to the source's row metric. To produce actual \(g_j,e_j,H\), the next required data are still the allowed finite jet matrix on these partitions, its integral row realization, the actual nonzero iteration and the endpoint detector. The separate logarithmic quantities in (R11.15), (R11.16) and the class mass cannot simply be added as independent losses without checking the source maps: (R11.13) exhibits one exact normalization cancellation already.

\subsubsection{References and source locators}\label{references-and-source-locators}

\begin{itemize}
\tightlist
\item
  Fulton, W., \& Harris, J. (1991). \emph{Representation theory: A first course}. Springer. Hook Length Formula4.12, p.50; Theorem6.3(2)--(4), pp.77--78; Exercise6.11(b), p.80. The cited theorem statements and formulas were directly read.
\item
  Alpöge source, pinned commit recorded in the source audit: \nolinkurl{Formalism/LocalRingCardinality/DataCoeff.lean}, \nolinkurl{detShift}, \nolinkurl{natW}, \nolinkurl{extremal}, \nolinkurl{normFock} and \nolinkurl{normFockQ}, lines21--23, 55--56 and83--87; \nolinkurl{Formalism/LocalToGlobal/HermitianForms/IsLocBasis.lean}, \nolinkurl{nfAll}, lines19--20; \nolinkurl{Formalism/LocalToGlobal/Divisibility/RecBridgeNum.lean}, \nolinkurl{logHook_eq}, lines149--188. These files were directly read.
\item
  AM-DO-025, (DO25.5)--(DO25.17); AM-DO-026, (DO26.7)--(DO26.8) and (DO26.17)--(DO26.20); AM-R-009, corrected actual Gaussian norm; AM-K019, the full differential-module obstruction with its exact scope.
\end{itemize}

Whole-proof check events are recorded separately. No Lean execution or automatic proof verification is asserted.

\clearpage

\section{AM-K022 --- selected Schur jets in the actual permitted half-orbit}\label{am-k022-selected-schur-jets-in-the-actual-permitted-half-orbit}

Author: GPT-6 Astra Pro, ALP\_MACHINE /root/kummer\_extension, 21 November 2026 UTC. Inherited reasoning effort is not separately exposed.

This result gives an explicit arithmetic level and a single permitted finite Schwartz vector detecting any prescribed nonzero homogeneous Schur polynomial. It applies to the ordinary product of Jacobians on the common E′ host, including the actual first even-rank edge of AM-DO-025. It does not identify the second iterated image or compute its class-mass content.

The permitted half-orbit basis is \textbf{carried DP59}, not a new representation theorem. The new conclusions are the numerical interpolation threshold, its product specialization, the factorial-exact highest-symbol detection, the norm-one individual-translate conclusion, and the equality of all finite determinantal ideals after passing from jets to the permitted section basis.

\subsection{1. Fixed geometry and the permitted lattice}\label{fixed-geometry-and-the-permitted-lattice}

Work at 11 over an unramified complete DVR R large enough to contain the CM eigenvalues and the finite torsion and symmetric lifts in DP59. Normalize v₁₁(11)=1. Set \[
 J=\operatorname{Jac}(y^2=x^7-1),\qquad
 L=L_{2\infty},\qquad s=s_q,
\] where the actual q-normalized section is primitive at 11 by ALP-LIFT-08. Its reduced germ is σq by AM-K018. For an integer N≥1 write \[
 X=J^N,\qquad \mathcal L=L^{\boxtimes N},\qquad
 s_X=s^{\boxtimes N}.
 \tag{J1}
\] For the common host E′=Q(ζ₂₈), the positive Hom variety on an edge U(a)→U(b) is X=A\^{}\{ab\}=J\^{}\{2ab\}, with A=J², so N=2ab. Its holomorphic dimension is 6ab. The actual E′ eigenbasis differs from the two ordinary J bases by a matrix invertible at 11 (AM-DO-022). It may therefore be used for the integral tangent coordinates below.

Let \[
 \delta^2=-7,\qquad
 \alpha=(1+2\delta)^e,\qquad M=29^e,\qquad
 \alpha\bar\alpha=M,\qquad \kappa=M^2.
 \tag{J2}
\] The ideals (α) and \(\bar\alpha\) are coprime. On each J factor take the exact DP59 orbit \[
 s_{t'}=U_{t'}[M]^*s,\qquad t'\in J[\alpha^2],
\] with the canonical odd-order symmetric theta lift. Choose one representative of J{[}α²{]}/J{[}α{]}. DP59(a) gives an integral basis of \[
 W_M=[\alpha]^*H^0(J,L^M)
       \subset H^0(J,L^{M^2}),
 \qquad \operatorname{rank}_R W_M=M^3.
 \tag{J3}
\] The orbit is smaller than the full level-M² section module.

On X take the tensor basis \[
 \mathcal S_{\mathbf t'}
   =\boxtimes_{\ell=1}^{N}s_{t'_\ell},
 \qquad
 \mathcal W_{M,N}=W_M^{\otimes N}
   =[\alpha]^*H^0(X,\mathcal L^M).
 \tag{J4}
\] Its rank is M\^{}\{3N\}. Its actual rescaled Heisenberg parameters are \[
 \mathbf x=M\mathbf t'\in X[\alpha]\subset X[M].
 \tag{J5}
\] Thus all vectors used here belong to the allowed half-orbit. The construction never replaces it by the full level-M² theta group.

\subsection{2. An explicit integral interpolation threshold}\label{an-explicit-integral-interpolation-threshold}

\subsubsection{Proposition J6}\label{proposition-j6}

For every integer D≥0 and every M≥3D+2, restriction is surjective: \[
 H^0(J,L^M)\longrightarrow
 H^0\!\left(J,\ L^M\otimes\mathcal O_J/I_0^{D+1}\right).
 \tag{J6}
\] The analogous assertion holds on X=J\^{}N with its product line, with the same bound M≥3D+2, independent of N.

\textbf{Known geometric input.} On an abelian variety over a field, the third power of an ample line is very ample. Milne, \emph{Abelian Varieties}, I, §6, ``Ample divisors,'' printed p.~30 (PDF page 36 of AVc.pdf), states this in the chapter's arbitrary-characteristic setting. The higher-cohomology vanishing and integral H⁰ base change used below are the classical input already read and applied in DP59 §5, equations (TS11); its primary locator is Mumford (1966), §1(II), printed p.~289. No characteristic-zero vanishing theorem is substituted for the special-fibre assertion.

\textbf{Proof.} The special-fibre line L³ is very ample. Choose a section \(\bar h_0\) nonzero at zero, and sections \(\bar h_1,\bar h_2,\bar h_3\) vanishing there whose ratios \(\bar h_i/\bar h_0\) have linearly independent differentials. These choices exist because a very ample system has a surjective evaluation map and separates tangent vectors. They can be made over the residue field: these are surjective maps of finite-dimensional vector spaces, and an invertible minor can be chosen over that field.

Integral H⁰ base change lifts these sections to sections h\_i of L³ over R. Normalize h₀(0)=1 in a rigidified fibre. Subtract h\_i(0)h₀ from h\_i for i\textgreater0. The ratios \[
 t_i=h_i/h_0,\qquad 1\le i\le3,
\] are integral formal parameters at zero because their cotangent determinant is a unit. Integral linear combinations make their first differentials any specified integral basis, including the dual of the actual CM tangent basis. Only unit inverses are used.

For each integer m≥2 there is a section of L\^{}m whose value at zero is a unit. Here is the required special-fibre fact. The square L² is generated at every geometric point: choose a translate of a nonzero divisor of L so that neither it nor its opposite translate contains that point; the theorem of the square identifies their product line with L². The forbidden translation sets are two proper closed subsets. Their complement is nonempty over an algebraic closure. Evaluation surjectivity then descends to the residue field. L³ is generated because it is very ample. Every m≥2 is a sum of copies of 2 and 3, so products give the assertion for all m. Integral H⁰ base change lifts a section nonzero at zero to one with unit value over R.

Take m=M−3D≥2 and such a section η\_m of L\^{}m. The section η\_mh₀\^{}D trivializes L\^{}M along the D-th infinitesimal neighbourhood of zero. An arbitrary target jet in this frame is \[
 Q(t)=\sum_{|\beta|\le D}q_\beta t^\beta,\qquad q_\beta\in R.
\] The global section \[
 \eta_m\sum_{|\beta|\le D}
       q_\beta h_0^{D-|\beta|}
                h_1^{\beta_1}h_2^{\beta_2}h_3^{\beta_3}
 \tag{J7}
\] has exactly this jet. This proves (J6).

For the product, tensor the single-factor restriction maps through order D. A monomial of total degree at most D in all 3N parameters has degree at most D in each three-variable block. Products of the single-factor interpolating sections therefore span every such monomial. After multiplication by the product local frame unit, they span arbitrary total D-jets. The image is contained in H⁰(X,ℒ\^{}M) by the usual product-section map. This proves the product assertion without a dimension-dependent change of M. ∎

This is the numerical specialization r=3, b=2 of the integral homogenization in accepted DP56(a), not a new claim that an arbitrary ample line on an arbitrary proper variety has that threshold.

\subsubsection{Corollary J8. The same bound for the allowed half-orbit}\label{corollary-j8.-the-same-bound-for-the-allowed-half-orbit}

For M=29\^{}e≥3D+2, \[
 \mathcal W_{M,N}\longrightarrow
 H^0\!\left(X,\mathcal L^{M^2}
                    \otimes\mathcal O_X/I_0^{D+1}\right)
 \quad\text{is surjective over }R.
 \tag{J8}
\]

\textbf{Proof.} The étale isogeny {[}α{]} has degree prime to 11, hence induces an integral formal isomorphism at zero. The rigidified equality {[}α{]}*ℒ\^{}M≃ℒ\^{}\{M²\} is part of the DP59 biextension construction. Its comparison in integral local frames is multiplication by an integral unit. Pulling (J6) back along this formal isomorphism proves (J8). The tensor orbit basis in (J4) is the carried DP59 basis on each factor, so the surjection is from exactly the permitted lattice. ∎

One explicit choice is \[
 e(D)=\max\!\left(1,\left\lceil\log_{29}(3D+2)\right\rceil\right),
 \qquad M(D)=29^{e(D)},
 \qquad
 3D+2\le M(D)<29(3D+2).
 \tag{J9}
\] In particular κ=M²\textless29²(3D+2)². The threshold is for \textbf{ordinary algebraic jets modulo I₀\^{}\{D+1\}}, not for divided powers or derivative values normalized to one.

\subsection{3. Highest-symbol detection and its exact factorial}\label{highest-symbol-detection-and-its-exact-factorial}

Use any actual integral invariant tangent basis \(D_1,\ldots,D_{3N}\), and choose formal parameters t\_i with \(D_it_j(0)=\delta_{ij}\). In the characteristic-zero logarithmic coordinates u\_i dual to these derivations, \(t_i(u)=u_i+O(u^2)\).

For a section S of ℒ\^{}\{M²\}, let F\_S(u) be its actual Chern/source reduced germ in the common normalization of AM-K018. In an integral algebraic frame normalized at zero it has the form \[
 F_S(u)=f_S(t(u))\,G_M(u),\qquad G_M(0)=1.
 \tag{J10}
\] The gauge G\_M is a characteristic-zero series; this argument does not require it to be an integral ordinary power series. The equality uses the fixed source gauge, not a section-dependent gauge.

Let \[
 P(Z)=\sum_{|\beta|=D}c_\beta Z^\beta
          \in R[Z_1,\ldots,Z_{3N}]
\] be nonzero and homogeneous, and define the normalized functional \[
 \ell_{P,M}(S)=
       M^{-D}\,[P(\partial_u)F_S(u)]_{u=0}.
 \tag{J11}
\]

\subsubsection{Proposition J12}\label{proposition-j12}

For every \(\beta\in\operatorname{supp}P\), the image ideal of ℓP,M on the permitted lattice contains \(c_\beta\beta!\): \[
 c_\beta\beta!\in\ell_{P,M}(\mathcal W_{M,N}),
 \qquad
 \beta!=\prod_i\beta_i!,
 \qquad M\ge3D+2.
 \tag{J12}
\] Consequently at least one individual allowed tuple \(\mathbf t'\) satisfies \[
 \ell_{P,M}(\mathcal S_{\mathbf t'})\ne0,\qquad
 v_{11}\!\left(\ell_{P,M}(\mathcal S_{\mathbf t'})\right)
 \le \min_{\beta\in\operatorname{supp}P}
                     v_{11}(c_\beta\beta!).
 \tag{J13}
\]

\textbf{Proof.} By (J8), choose S with algebraic jet \[
 f_S(t)\equiv M^D t^\beta\pmod{(t)^{D+1}}.
 \tag{J14}
\] The scalar M\^{}D is a unit at 11. There are no terms below degree D. Substituting \(t(u)=u+O(u^2)\) and multiplying by \(G_M(u)=1+O(u)\) therefore gives \[
 F_S(u)\equiv M^D u^\beta\pmod{(u)^{D+1}}.
 \tag{J15}
\] All higher coefficients of the possibly nonintegral gauge are irrelevant to this equality. The homogeneous differential polynomial kills every degree-D monomial except the one matching its exponent: \[
 M^{-D}P(\partial_u)M^D u^\beta\big|_0
        =c_\beta\prod_i\beta_i!.
 \tag{J16}
\] This proves (J12).

Expand S in the integral basis (J4): \[
 S=\sum_{\mathbf t'}a_{\mathbf t'}\mathcal S_{\mathbf t'},
 \qquad a_{\mathbf t'}\in R.
 \tag{J17}
\] If all basis values had valuation greater than \(v_{11}(c_\beta\beta!)\), their integral linear combination could not equal the nonzero value in (J16). Choose β minimizing that valuation. At least one basis value obeys (J13). The tuple and its canonical theta lifts are global algebraic objects by DP59(c). The conclusion does not require the coefficients in (J17) themselves to be global algebraic numbers. ∎

The factorial in (J16) is essential. Interpolating the ordinary coefficient of t\^{}β to one produces derivative β!, not derivative one. An assertion of a unit derivative at arbitrary degree would divide by β! and would not be an integral interpolation statement.

The argument neither bounds nor uses the archimedean size of the coefficients in (J17). It uses them only in the local nonarchimedean valuation inequality to select one actual basis vector.

\subsection{4. The exact selected Schur polynomial}\label{the-exact-selected-schur-polynomial}

On the even-rank edge U(a)→U(b), b\textgreater a, let \(\nu_\tau\) be the actual nonnegative source partitions identified in AM-DO-025: \[
 \nu_{\tau,i}
   =(\lambda_{\mathrm{True},\tau})_i
                        +e_\tau-b/2,\qquad 1\le i\le a.
 \tag{J18}
\] There are six holomorphic labels τ on the E′ host. Put ντ,a+1=0 and write \(Z_\tau=(Z_{\tau,i,j})\) as an a-by-b matrix. Take the standard primitive highest bideterminant \[
 P_\nu(Z)=
 \prod_\tau\prod_{k=1}^{a}
       \det(Z_{\tau,\{1,\ldots,k\},\{1,\ldots,k\}})
                     ^{\nu_{\tau,k}-\nu_{\tau,k+1}}.
 \tag{J19}
\] Its homogeneous degree is exactly \[
 D=\sum_\tau|\nu_\tau|.
 \tag{J20}
\] It is the highest vector in the actual Cauchy summand of AM-DO-025. This polynomial already lies in the selected Schur space; applying a rational Schur projector to it is unnecessary.

The diagonal monomial of Pν has coefficient exactly one: \[
 \prod_{\tau,i}Z_{\tau,i,i}^{\nu_{\tau,i}}.
 \tag{J21}
\] Indeed every determinant must contribute its diagonal term to produce a monomial containing no off-diagonal variable. The exponent of its i-th diagonal variable is \(\sum_{k\ge i}(\nu_{\tau,k}-\nu_{\tau,k+1})=\nu_{\tau,i}\). This also proves that Pν is primitive in the polynomial coefficient lattice. It does not assert that every alternative Weyl or normalized Fock lattice has that same basis; their comparison is separate.

Applying (J13) with this diagonal monomial gives a single permitted tuple with \[
 a_{\nu,\mathbf t'}
 :=M^{-D}[P_\nu(\partial)F_{\mathcal S_{\mathbf t'}}]_0
 \ne0,
\qquad
 v_{11}(a_{\nu,\mathbf t'})
 \le
 \sum_{\tau,i}v_{11}(\nu_{\tau,i}!).
 \tag{J22}
\] Equivalently, \[
 v_{11}(a_{\nu,\mathbf t'})
 \le
 \frac{D-\sum_{\tau,i}s_{11}(\nu_{\tau,i})}{10}.
 \tag{J23}
\] This is an explicit upper bound for \textbf{one selected kernel coefficient on the allowed family}. It is not a sharp derivative slope and not a bound for an arbitrary later theta-image vector.

It works for every actual ν₀ satisfying the source recurrence and gap conditions. No fixed-tail or single-row restriction is imposed. Therefore it supplies the nonzero selected first-edge kernel needed for the parent's class-supported compact extraction on U(N₀)→U(N₁).

\subsection{5. The individual Schwartz vector has no translation-combination norm cost}\label{the-individual-schwartz-vector-has-no-translation-combination-norm-cost}

Let \(\phi_{0,f}\) be the actual finite vacuum on X=A\^{}\{ab\}, normalized to finite L² norm one as in AM-DO-022. For the tuple selected in (J22), use \[
 \phi_{\mathbf x,f}=W_{\mathbf x}\phi_{0,f},
 \qquad \mathbf x=M\mathbf t'\in X[\alpha].
 \tag{J24}
\] Every local Heisenberg translation is unitary for the fixed self-dual Schrödinger measure. The exact symmetric lift differs only by its fixed finite-order unit scalar. Hence \[
                 \|\phi_{\mathbf x,f}\|_2^2=1.
 \tag{J25}
\] In its permitted-orbit expansion this chosen vector has one coefficient equal to one and all others zero. Its coefficient maximum is one at every finite and complex embedding. Thus selecting this vector introduces no coefficient height or sum-of-orbit norm cost.

The level-normalized identity of ALP-LIFT-10 and AM-R-010 is \[
 M^{-D}P(\partial_u)
 U_{\mathbf t'}^{M^2H}\bigl(\sigma_q^{\boxtimes N}(Mu)\bigr)\big|_0
 =
 P(\partial_v)
 U_{\mathbf x}^{H}\sigma_q^{\boxtimes N}(v)\big|_0,
 \tag{J26}
\] with the specified canonical finite-order scalar retained. Accordingly AM-DO-026's actual kernel has the nonzero identity coefficient \[
           \vartheta_{r_{P_\nu}\phi_{\mathbf x,f}}(1,1)
                  =C^{2ab}a_{\nu,\mathbf t'}.
 \tag{J27}
\] Here \(C=\theta[0;0](0;\tau)c\), with the fixed algebraic scalar \(c\) of AM-R-006, is the one common period for the \(\sigma_q\) seed in AM-DO-022, §5. It is not a coefficient chosen for this tuple. The archimedean norm remains precisely the norm of the creation vector \(r_{P_\nu}\phi_{0,\infty}\); (J25) removes no Fock or period factor.

This statement supplies the finite vector and its norm. The source automorphic input, its stabilizer mass, its compact Schur-orthogonality factor, and its normalization are constructed separately in the parent's AM-R-012. None of those factors is set equal to one here.

\subsection{6. Exact finite saturation reduction for every chosen polynomial family}\label{exact-finite-saturation-reduction-for-every-chosen-polynomial-family}

The same interpolation gives a useful matrix statement that does not rely on a leading-symbol approximation.

Choose finitely many homogeneous polynomials P₁,\ldots,P\_q of degree at most D in a specified integral selected lattice, and include their individual source factors M\^{}\{-deg Pj\}. Choose an ordinary algebraic monomial basis of the total D-jet module in the fixed integral theta frame. Let \[
 J_D:R^{\,M^{3N}}\longrightarrow R^{\,\binom{3N+D}{3N}}
 \tag{J28}
\] be the matrix of the actual permitted section basis's ordinary jets. Let B be the q-by-\(\binom{3N+D}{3N}\) matrix of the corresponding normalized Chern/source derivative functionals on these jets. Its entries lie in Frac(R); no integral reduced frame is assumed. Let A be the q-by-M\^{}\{3N\} matrix of their actual allowed-tuple values. Then \[
                       A=BJ_D.
 \tag{J29}
\]

\subsubsection{Proposition J30. No extra finite-translation determinantal loss}\label{proposition-j30.-no-extra-finite-translation-determinantal-loss}

For M≥3D+2, the map J\_D is split surjective over R. For every k≥1 the fractional ideals generated by the k-by-k minors satisfy \[
                     I_k(A)=I_k(B).
 \tag{J30}
\] Equivalently, after adding zero columns, A and B have the same nonzero Smith factors over the DVR.

\textbf{Proof.} Surjectivity is (J8). The jet module is finite free, so choose an R-linear right inverse C with J\_DC=1. Thus B=AC. Cauchy--Binet applied to A=BJ\_D gives \(I_k(A)\subset I_k(B)\), since J\_D is integral. Applied to B=AC it gives the reverse inclusion, since C is integral. For fractional matrices first multiply by one common denominator; both k-minor ideals scale by its k-th power, so cancellation recovers (J30). Alternatively split the free domain of J\_D as its kernel plus a copy of its target; in adapted integral bases J\_D is \((1\ 0)\), and A is \((B\ 0)\). ∎

This equality is about the actual allowed half-orbit, not a larger arbitrary Schwartz space. It quantifies the absence of additional saturation loss in passing from the finite jet module to these section generators. The integral right inverse may have large archimedean entries, but no norm estimate for it is used in (J22) or (J25).

For the second automorphic lift, its class integral and the restriction to the already produced first image impose additional linear maps on this kernel data. Those maps can change the Smith factors, can cause cancellation, and can annihilate a vector. Proposition J30 does not assert that independently chosen jets at the identity can be prescribed at every class representative. The remaining second-stage problem is precisely the content of that \textbf{restricted class matrix}, with the source's actual selected coefficient lattice and compact contraction factors included.

The full differential-module obstruction AM-K019 is consistent with (J30). Here the degree-D jet module is finite, the level grows with D, and only the specified evaluated functionals are retained. No fixed-rank D-stable lattice, bounded saturation at all degrees, or iteration of a mod-11 coefficient congruence is asserted.

\subsection{7. Sources and checks}\label{sources-and-checks}

\begin{itemize}
\tightlist
\item
  DP\_MULT DP56(a), research/quantitative\_theta\_jet\_levels\_DP56.md, §§1--3: integral homogenization and its precise ordinary-jet convention. Carried; the explicit choice r=3,b=2 is justified in Proposition J6.
\item
  DP\_MULT DP59(a)--(c), research/theta\_translate\_span\_DP59.md, §§3--8: actual rigidified pullback, exact permitted half-orbit, integral basis, unramified normalization, and global algebraicity. Carried, not rederived.
\item
  Milne, J. S. (2008). \emph{Abelian varieties} (Version 2.0), I, §6, subsection ``Ample divisors,'' printed p.~30, paragraph immediately before Example 6.5. \url{https://www.jmilne.org/math/CourseNotes/AVc.pdf}. The precise used fact is that the third power of every ample line on an abelian variety is very ample; the chapter is over an arbitrary base field.
\item
  Mumford, D. (1966). On the equations defining abelian varieties. I. \emph{Inventiones Mathematicae, 1}, 287--354. §1(II), p.~289, as already audited in DP59 for the cohomology/base-change input. \url{https://pazuki.perso.math.cnrs.fr/index_fichiers/Mumford66.pdf}.
\item
  AM-DO-022 and AM-DO-025--026: actual E′ product polarization, unitary finite Heisenberg action, selected partitions, and all-polynomial kernel identity.
\item
  AM-K018 and ALP-LIFT-08: exact source gauge on ordinary J products and primitive actual principal section at 11.
\item
  ALP-LIFT-10 and AM-R-010: exact permitted parameter \(\mathbf x=M\mathbf t'\), retained symmetric lift scalar, and level-normalized derivative identity.
\end{itemize}

The result does not assume or bound an interpolating section's global coefficient height. The chosen individual allowed translate has coefficient vector a standard basis vector, so its finite Schwartz norm and translation-coefficient bound are exact.

\clearpage

\section{AM-R-012: an actual selected even-rank first lift with a single permitted finite translate}\label{am-r-012-an-actual-selected-even-rank-first-lift-with-a-single-permitted-finite-translate}

Author: /root, GPT-6 Astra Pro, ALP\_MACHINE, 21 November 2026 UTC.

The rank-one starting correspondence of AM-R-009--010 is not the starting edge in the recovered \nolinkurl{Ranks} record. The construction here uses the actual even ranks \(N_0\to N_1\), their selected compact coefficient types, and the same splitting character on both groups. It proves a controlled nonzero \textbf{first} theta lift with an exact source input and a single permitted finite Heisenberg translate. It does not yet identify its arithmetic row with the old elliptic \nolinkurl{StdPair}, or prove the later chain and six comparison identities.

The finite interpolation input is AM-K022, which specializes the accepted DP56/DP59 allowed half-orbit theorem to the actual selected polynomial. That proof and its check status are separate. In the form used here it supplies one permitted translate with finite Schwartz norm one and the nonzero jet bound (R12.5) below.

\subsection{1. Actual fields, types, polynomial and character}\label{actual-fields-types-polynomial-and-character}

Use \[
 E'=\mathbf Q(\zeta_{28}),\quad F'=E'^+,\quad A=J^2,
 \quad G=U(I_a),\quad H=U(I_b),\quad a=k,
 \quad b=k+d=N_1.
 \tag{R12.1}
\] All six selected complex places are those of AM-R-009, and the positive Hom oscillator, trace-one vacuum, and period are exactly those of AM-DO-022/026. Let \(w_0,w_1\) and \(\nu_{0,\tau}\) be the weights and partitions of an actual \nolinkurl{HCParam}, as in AM-DO-025. Write \[
 D=\sum_{\tau\in\Phi'}|\nu_{0,\tau}|,
 \quad P(Z)=\prod_{\tau\in\Phi'}p_{\nu_{0,\tau}}(Z_\tau).
 \tag{R12.2}
\] Here each highest polynomial is the primitive bideterminant of AM-R-011. The actual source compact type of the oscillator is \(\sigma=\rho_0^\vee\), and the target compact type is \(\rho_1\), where \(\rho_j\) has highest weight \(w_j\). This is precisely the equal-character identification proved in AM-DO-025, not the trivial-character rank-two special case.

There is an explicit globally admissible common character with the required exponents \(-e_{0,\tau}\). Indeed AM-R-009 proves \(T'(\mathbf A_f)=T'(F')K_{T,0}\), with intersection \(\mu_{28}\). On the real norm-one torus set \(\beta(t)=\prod_\tau t_\tau^{-e_{0,\tau}}\). At one split place over29 choose the factor where \(\zeta_{28}\mapsto2\), and define the compact residue character \[
 \varrho(2)=\exp\!\left(\frac{2\pi i}{28}
                         \sum_{\tau\in\Phi'}e_{0,\tau}\tau\right).
 \tag{R12.3}
\] It is trivial at every other compact place. Thus \(\varrho|_{\mu_{28}}=\beta^{-1}\), so \(\beta\varrho\) descends to a unitary automorphic norm-one character \(\gamma\). Put \(\mu(x)=\gamma(x/\bar x)\). Then \(\mu|_{\mathbf A_{F'}^\times}=1\), \(\mu_\infty(z)=(z/\bar z)^{-e_{0,\tau}}\), and the finite compact conductor divides29. Take \(\chi_S=\chi_T=\mu\) on the edge. At nonsplit finite places the norm-one group is compact and \(\gamma\) is trivial there; at other split places it is unramified. All values needed for the algebraic finite kernel lie in a finite algebraic coefficient extension. This is the explicit character construction of AM-R-009 applied to the present integral exponent vector.

Choose \[
 e=\max\{1,\lceil\log_{29}(3D+2)\rceil\},\qquad M=29^e.
 \tag{R12.4}
\] By AM-K022, after a fixed unramified coefficient extension at11, one can choose a \textbf{single} permitted product-translate \(x\) such that its normalized geometric jet \(a_x\) satisfies \[
 a_x\ne0,\qquad
 v_{11}(a_x)\le
       B_\nu:=\sum_{\tau\in\Phi'}\sum_{i=1}^{a}
                              v_{11}(\nu_{0,\tau,i}!).
 \tag{R12.5}
\] The integer in the bound is the derivative of the diagonal monomial in \(P\); it is not the hook product and not an unproved lower bound for every selected jet. The translate belongs to the actual DP59 half-orbit after its specified odd-level normalization. The corresponding finite Schwartz vector is \(\phi_{f,x}=W_x\phi_0\) with its actual central multiplier, and \(\|\phi_{f,x}\|=1\) by unitarity. No interpolation coefficient norm is hidden in this selection of one orbit vector.

Let \(\phi_x=r_P\phi_{f,x}\phi_\infty\). The actual kernel identity AM-DO-026, with the normalized level factors retained, gives \[
                  \vartheta_{\phi_x}(1,1)=C^{2ab}a_x\ne0.
 \tag{R12.6}
\]

\subsection{2. A concrete source level and its exact identity-class mass}\label{a-concrete-source-level-and-its-exact-identity-class-mass}

Let \(K_{a,0}\) be the standard full integral unitary compact. Refine it to \(K_a(M)\) by principal congruence modulo29\^{}e at all six split places over29, leaving all other places unchanged. Then \(K_a(M)\) fixes \(\phi_{f,x}\) exactly in the source action. The same assertion holds for the target subgroup \(K_b(M)\).

To see this, at29 every permitted parameter has denominator dividing \(M\) and lies in the selected isotropic split position space. The linear action \(X\mapsto hXg^{-1}\) preserves that position space. For \(g,h\equiv1\pmod M\), its displacement is integral in the same space. The pairing with the original parameter is zero, and the odd-place trace-one lattice character of that displacement is one. The exact Weyl law therefore fixes \(W_x\phi_0\). The compact determinant characters also equal one on this level by (R12.3). Away from29 the parameter is integral; the line is the original trace-one vacuum line, fixed by the standard compacts and the unramified splitting. These are exactly the split/unramified normalizations of AM-DO-022 and the isotropic half-orbit argument of AM-R-010, now applied entrywise to the matrix Hom space.

At the identity finite class the rational stabilizer is trivial: \[
                   \Gamma=G(F')\cap K_a(M)=\{1\}.
 \tag{R12.7}
\] First \(\Gamma\) is finite, because it preserves the integral lattice and lies in a compact group at infinity. Its image at one split place over29 is inside \(1+29M_a(\mathbf Z_{29})\), which is torsion free. For completeness, if a nonzero matrix \(B\) has minimum entry valuation \(s\ge1\), then \((1+B)^{29}-1\) has minimum valuation exactly \(s+1\): the term \(29B\) has that valuation, and all the other terms have strictly larger valuation. Raising to an integer prime to29 leaves the minimum valuation \(s\). A finite-order element in this group would therefore remain nonidentity after raising to its order, a contradiction. The local embedding of the rational matrix group is injective, proving (R12.7).

Normalize real Haar measure to have \(\operatorname{vol}(G_\infty)=1\), and finite Haar measure by \(\operatorname{vol}(K_{a,0})=1\). Use their compatible automorphic quotient measure; its total volume is not silently set to one. The exact index and the identity-class mass are \[
 \begin{aligned}
 I_a(M)&=[K_{a,0}:K_a(M)]
       =29^{6a^2(e-1)}|\mathrm{GL}_a(\mathbf F_{29})|^6\\
       &=M^{6a^2}\prod_{i=1}^{a}(1-29^{-i})^6,
       \qquad w=I_a(M)^{-1}.
 \end{aligned}
 \tag{R12.8}
\] The split reduction maps are surjective, with successive congruence quotients of cardinality \(29^{a^2}\). Counting ordered bases gives \(|\mathrm{GL}_a(\mathbf F_{29})|=29^{a^2}\prod_i(1-29^{-i})\). There are six places, and (R12.7) removes the stabilizer divisor at this class. This proves (R12.8) with the actual chosen measure.

\subsection{3. Compact matrix-coefficient extraction}\label{compact-matrix-coefficient-extraction}

Use the standard unitary metrics on the compact representations, with highest vector of norm one. Let \(u\in\sigma\) be the unit source vector corresponding to the source factor of the highest polynomial \(P\), and put \[
 d_\rho=\dim\rho_0=\dim\sigma,
 \qquad m(k)=\langle u,\sigma(k)u\rangle.
 \tag{R12.9}
\] The Hermitian convention is conjugate-linear in the first argument. At the identity finite class, the kernel can be written \[
             \vartheta_{\phi_x}(k,1)=L(\sigma(k)u),
             \qquad L(u)=C^{2ab}a_x,
 \tag{R12.10}
\] where the scalar relating the actual highest polynomial to unit Schur vectors is included in \(L\). This is the actual source equivariance of the theta kernel; its finite translate acts on a different tensor factor and does not change the compact type.

Define an automorphic function \(f_0\) supported on the identity finite class, whose value there is \(\overline{m(k)}\), and whose value on all other finite classes is zero. Equation (R12.7) makes this definition unambiguous. It is right \(K_a(M)\)-invariant and smooth at infinity, with right compact type \(\sigma^\vee=\rho_0\). This gives an actual automorphic input, rather than a free representation parameter. It is cuspidal because the definite rational group has no proper rational parabolic subgroup.

Schur orthogonality, or the average of the rank-one projection \(v\mapsto u\langle u,v\rangle\), gives \[
 \int_{G_\infty}\sigma(k)u\overline{m(k)}\,dk
                      =u/d_\rho,
 \qquad \int_{G_\infty}|m(k)|^2\,dk=1/d_\rho.
 \tag{R12.11}
\] Insert this into the literal \textbf{bilinear} theta integral, and use the mass \(w\) of the supporting class. One obtains the exact nonzero value and norm \[
 \boxed{\Theta^{\rm bil}_{\phi_x}(f_0)(1)
        =\frac{C^{2ab}a_x}{I_a(M)d_\rho},\qquad
        \|f_0\|_2^2=\frac1{I_a(M)d_\rho},\quad
        \|f_0\|_\infty=1.}
 \tag{R12.12}
\] The supremum is attained at the real identity. In the integral coefficient presentation this is a class vector with primitive highest coefficient: the highest-to-highest matrix coefficient is the primitive bideterminant with the integral determinant character attached, as in AM-R-011. Its value at the identity is one. No claim that all scalar values on the real compact group are algebraic integers is needed or intended.

Two useful alternative normalizations are consequently \[
 \begin{array}{c|c|c}
 \text{input}&\text{squared norm}&\text{output at 1}\\\hline
 I_a(M)d_\rho f_0&I_a(M)d_\rho&C^{2ab}a_x\\
 \sqrt{I_a(M)d_\rho}\,f_0&1&C^{2ab}a_x/\sqrt{I_a(M)d_\rho}.
 \end{array}
 \tag{R12.13}
\] Thus the source normalization cost for the extracting input is exactly \(\tfrac12\log(I_a(M)d_\rho)\). It has not been removed by using a nonprimitive scalar multiple of the row.

The compact dimension is explicitly \[
 d_\rho=\prod_{\tau\in\Phi'}\prod_{1\le i<l\le a}
      \frac{\eta_{0,\tau,i}-\eta_{0,\tau,l}+l-i}{l-i}.
 \tag{R12.14}
\] With fixed ranks and \(M<29(3D+2)\), the logarithmic source cost is \(O(\log(D+2))\), with the exact finite index given in (R12.8), rather than an unspecified level constant. The full archimedean Schwartz norm is (R11.14), and the finite Schwartz norm is one. Unit normalization of both input and Schwartz vector therefore divides (R12.6) by \(\sqrt{I_a(M)d_\rho}\,\|\phi_{x,\infty}\|\).

The normalized algebraic identity coefficient in the primitive row convention has the controlled upper valuation \[
 v_{11}\!\left(C^{-2ab}\Theta^{\rm bil}_{\phi_x}(f_0)(1)\right)
          \le B_\nu-v_{11}(I_a(M)d_\rho).
 \tag{R12.15}
\] This is a selected nonzero entry, not a lower content bound for every entry of an arithmetic row map. It is therefore not itself the gain \(g_0\) in DP44.

\subsection{4. A selected first coefficient with no positive 11-adic gain}\label{a-selected-first-coefficient-with-no-positive-11-adic-gain}

For the \nolinkurl{HCParam} example in AM-R-011, take the same parameters at all six CM places. Then \(a=6,b=20\), \(\nu_{0,\tau}=(1^6)\), \(\eta_{0,\tau}=0\), \(D=36\), \(d_\rho=1\), and the explicit choice (R12.4) is \(M=29^2=841\). The polynomial is the product of six six-by-six determinants, and \(B_\nu=0\).

Furthermore \[
 I_6(841)=29^{216}|\mathrm{GL}_6(\mathbf F_{29})|^6,
                         \qquad v_{11}(I_6(841))=0.
 \tag{R12.16}
\] Indeed the order of29 modulo11 is10: \(7^5=-1\) modulo11, while \(7\) and \(7^2\) are not one. None of \(29^i-1\), \(1\le i\le6\), is divisible by11. Equations (R12.5), (R12.12) and (R12.16) therefore give an actual permitted first lift for which \[
 v_{11}\!\left(C^{-240}\Theta^{\rm bil}_{\phi_x}(f_0)(1)\right)
                                  \le0.
 \tag{R12.17}
\] Thus a strictly positive uniform 11-adic content for all these allowed selected first coefficients would be false. The statement is an upper bound; equality zero would additionally require integrality of all the relevant reduced jet coefficients. Neither that extra integrality nor the further noncompact \nolinkurl{Hyp3} conditions are assumed in this example.

Its exact archimedean Schwartz norm is \[
 \|\phi_{x,\infty}\|^2
       =\det(2Y)^{-120}(6!)^6
                    \prod_{i=1}^{3}\left(\frac{\pi}{2G_i}\right)^{12}.
 \tag{R12.18}
\] The primitive scalar input has squared norm \(I_6(841)^{-1}\). All these quantities belong to the same actual first edge. They do not bound the gains of the later restricted image or the rebuilt-family DP44 difference.

\subsection{5. What has and has not been supplied}\label{what-has-and-has-not-been-supplied}

The construction supplies one actual first edge at the source's even ranks with its selected local types, a globally admissible equal character, one permitted finite input of norm one, a specific class-supported source function, and an exact nonzero value up to the explicitly chosen algebraic jet \(a_x\). The finite selection can be made deterministic by ordering the finite permitted orbit and choosing the first tuple with the valuation bound (R12.5). No search over an infinite Schwartz space is required for this first edge.

It does not prove that every source row survives, or that the resulting vector survives the next prescribed kernel. It does not replace the old elliptic \nolinkurl{HostRaw}/\nolinkurl{StdPair} by definition, identify the independently fixed arithmetic \nolinkurl{rowAF}, or verify the noncompact analytic assumptions in \nolinkurl{AdmS}. Those are still needed for the requested shorter chain and the DP44 comparison. The exact first value, norms and index above are available for that next step, without the previous rank-one starting mismatch.

Geometrically, the permitted translates are sections in the Picard/theta-group construction, and (R12.5) is finite-jet separation on the good integral model. The compact coefficient is a Schur functor over \(\mathbf Z\), with the rational class-function realization retaining its actual metric. The relevant locators are FGA232 Theorem3.1; FGA236 Corollaries1.7 and2.3; SGA7 I VII2.7/2.9; EGA IV4 section16 on principal parts; SGA3 I Proposition4.7.3 and SGA6 V/VI as qualified in the geometric dictionary. No new crystalline or motivic L-function implication is asserted by this first lift.

Exact carried inputs: DP56--59; AM-K022; AM-DO-022, AM-DO-025 and AM-DO-026; AM-R-009--011. The source bilinear integral is \nolinkurl{WittSmallGroupEmbeddingMemPointwiseStabilizerIso.lean}, \nolinkurl{liftT_apply}, lines160--167. The selected weights are \nolinkurl{Columns/BaseLemmas.lean}, \nolinkurl{lamTrue_old}, \nolinkurl{lamTrue_new}, and \nolinkurl{nuW_eq_true}, lines496--552. The source row normalization is \nolinkurl{DataCoeff.lean}, \nolinkurl{normFock}, lines83--84. No Lean build or formal source-record instantiation is claimed.

\clearpage

\section{ALP-LIFT-13 --- fixed finite first-image witnesses at 3 and 5}\label{alp-lift-13-fixed-finite-first-image-witnesses-at-3-and-5}

Author: /root/lift\_accounting, GPT-6 Astra Pro, ALP\_MACHINE, 21 November 2026 UTC.

This specializes the already constructed actual first theta correspondence to AM-K020's sharp derivatives of one fixed theta section. It does not reprove those derivative congruences. Unlike the general permitted-orbit projector, this untranslated specialization uses one unit-norm automorphic character and one fixed finite Schwartz vector.

\subsection{1. Actual vector and its exact coefficient}\label{actual-vector-and-its-exact-coefficient}

Retain \(E=\mathbf Q(\zeta_7)\), \(F=E^+\), the actual compact pair \(T=U(1)\), \(G=U(I_2)\), and the honest splitting of ALP-LIFT-09. Write \[
 a_{i,n}=\partial_i^n\sigma_q(0),\qquad
 C=\theta[0;0](0;\tau)c .
 \tag{F1}
\] The normalization of \(\sigma_q\) is the actual algebraic Chern/source gauge established in AM-K014/K018 and ALP-LIFT-08: \(\sigma_q(u)=u_1u_3-u_2^2+\text{higher terms}\). The global period \(C\) is retained.

On two ordinary Jacobian factors use the creation polynomial \[
       P_{i,n}=\partial_{i,1}^{\,n}
                         \partial_{1,2}\partial_{3,2}.
 \tag{F2}
\] The finite Schwartz vector is the same \(\phi_{0,f}^{\otimes2}\) for every \(i,n\). Its finite \(L^2\) norm is one. Let \(\phi_{i,n}\) have this finite part and the actual Fock creation vector corresponding to (F2). ALP-LIFT-09's exact Bargmann reading gives \[
 \vartheta_{\phi_{i,n}}(1,1)
   =C^2\,P_{i,n}\bigl(\sigma_q(u_1)\sigma_q(u_2)\bigr)_{0}
   =C^2a_{i,n}.
 \tag{F3}
\] Indeed the second factor contributes exactly \(\partial_1\partial_3\sigma_q(0)=1\). It is necessary to include this spectator derivative: an undifferentiated second trace-one section vanishes at the origin and would give zero.

The polynomial has one fixed multidegree across CM embeddings: its degree vector is \((n+1,0,1)\) if \(i=1\), and \((1,n,1)\) if \(i=2\). Thus it is one torus weight, including the actual vacuum shift, despite using two different eigenlines on the spectator factor.

\subsection{2. One actual source character, without a growing finite level}\label{one-actual-source-character-without-a-growing-finite-level}

Let \(\beta_{i,n}\) be that actual infinity torus character, and let \(\lambda_f\) be the character on the fixed finite vacuum line. ALP-LIFT-09 proves \[
 T(\mathbf A_{F,f})=T(F)K_0,\qquad
 T(F)\cap K_0=\mu_{14}.
 \tag{F4}
\] For every \(a_{i,n}\ne0\), rational invariance of the actual nonzero sum (F3) gives \[
           \beta_{i,n}(\zeta)\lambda_f(\zeta)=1
                         \quad(\zeta\in\mu_{14}).
 \tag{F5}
\] This is a consequence of the actual theta sum; no compatibility is assumed from a hypothetical lift.

Consequently \[
 \chi_{i,n}(\delta\,t_\infty k)
        =\beta_{i,n}(t_\infty)\lambda_f(k),
 \quad\delta\in T(F),\quad k\in K_0,
 \tag{F6}
\] defines a unitary automorphic character. Independence of the decomposition follows from (F5). The finite compact character is fixed as \(n\) varies. In the standard splitting its only nontrivial compact factor is the orientation character \(\epsilon\) at7, by AM-DO-020--021; the proof using the actual \(\lambda_f\) does not require this optional explicit description.

Normalize the compact torus quotient to volume one. The equality of actual automorphic functions is \[
 \Theta(\phi_{i,n},\chi_{i,n})(g)
       =\int_{[T]}\vartheta_{\phi_{i,n}}(g,t)
                                 \overline{\chi_{i,n}(t)}\,dt
       =\vartheta_{\phi_{i,n}}(g,1).
 \tag{F7}
\] It follows because the theta function transforms by \(\beta_{i,n}\lambda_f\), the group actions commute, and (F4) represents every finite component. Hence \[
 \|\chi_{i,n}\|_2=\|\chi_{i,n}\|_\infty=1,\qquad
 C^{-2}\Theta(\phi_{i,n},\chi_{i,n})(1)=a_{i,n}.
 \tag{F8}
\] There is no finite Fourier projector and no factor \(h_M^{\pm1/2}\) in (F7)--(F8). The algebraic normalized value in (F8) is what will be assigned a \(p\)-adic valuation. No \(p\)-adic valuation is assigned to the unnormalized complex period \(C^2\).

As a direct root-unit check, in the fixed Hom convention \(X\mapsto Xt^{-1}\), the source vacuum shift is \(-1\) at each CM place. Therefore the exponent of \(\beta_{i,n}\) on \(\zeta_7\) is \(-in-10\). Both witness sequences below have \(in\equiv4\bmod7\). Their \(n\) are even, so \(\beta_{i,n}(-1)=-1\), matching the finite orientation value \(\lambda_f(-1)=-1\). This also verifies (F5) explicitly for these sequences using the corrected AM-DO-023 source shift. Formula (F5) remains the coordinate-independent proof.

\subsection{3. Exact sharp actual first values}\label{exact-sharp-actual-first-values}

AM-K020, jacobian\_sharp\_small\_prime\_directions.md, equations (K21.1)--(K21.1a), supplies the nonzero derivatives \[
 n_3(r)=28\cdot3^r+18,\qquad
 n_5(r)=126\cdot5^r+30
 \tag{F9}
\] and their exact valuations. Substitution into (F8) gives the following actual first-image coefficients: \[
 \begin{split}
 v_3\!\left(C^{-2}
       \Theta(\phi_{1,n_3(r)},\chi_{1,n_3(r)})(1)\right)
   &=\frac{3^r+3-s_3(n_3(r))}{2},\\
 v_5\!\left(C^{-2}
       \Theta(\phi_{2,n_5(r)},\chi_{2,n_5(r)})(1)\right)
   &=\frac{5^r-1}{4}.
 \end{split}
 \tag{F10}
\] The new assertion in (F10) is their occurrence as these specified actual theta-lift values. The numerical derivative theorem itself belongs to AM-K020.

In particular the same first-image family, with fixed finite vacuum and unit-norm source character, has \[
 \liminf_{n\to\infty}
   \frac{v_3(C^{-2}\Theta(\phi_{1,n},\chi_{1,n})(1))}{n}
       =\frac1{56},\qquad
 \liminf_{n\to\infty}
   \frac{v_5(C^{-2}\Theta(\phi_{2,n},\chi_{2,n})(1))}{n}
       =\frac1{504},
 \tag{F11}
\] where the family is indexed by the nonzero coefficients; zero coefficients may be included with valuation \(+\infty\). For a zero coefficient not satisfying (F5), no extra character is being asserted by the definition of this family. The global lower bounds and the witness subsequences in AM-K020 give both directions of (F11). Replacing the denominator \(n\) by the actual total creation degree \(n+2\) leaves these limits unchanged.

At every fixed permitted level \(M=29^e\), take the untranslated section \(s_M=[M]^*s_q\) on each factor. ALP-LIFT-10, (O6)--(O11), gives exactly \[
 M^{-(n+2)}P_{i,n}
       \bigl(\sigma_q(Mu_1)\sigma_q(Mu_2)\bigr)_{0}
                         =a_{i,n}.
 \tag{F12}
\] The point \(t'=0\) is in the permitted half-orbit and its finite phase is1, so the finite vector in that identity is the unchanged \(\phi_{0,f}^{\otimes2}\). The level can be fixed once and for all; no interpolation condition \(M\asymp n\) is required for these particular actual rows.

\subsection{4. Costs that have and have not been determined}\label{costs-that-have-and-have-not-been-determined}

The source input in (F7) has norm one and the finite Schwartz norm is one. The exact archimedean Schwartz norm is not one. ALP-LIFT-09, (V23), gives \[
 \|\phi_{i,n,\infty}\|_2^2
   =\det(2Y)^{-1}\,
           n!\left(\frac{\pi}{G_i}\right)^n
                    \frac{\pi^2}{G_1G_3}.
 \tag{F13}
\] The two spectator creation operators lie in distinct orthogonal modes and have factorial1. The first factor contributes \(n!\); the two ordinary Gaussian vacua contribute \(\det(2Y)^{-1}\). Thus \[
 \log\|\phi_{i,n,\infty}\|_2
   =\frac12\log(n!)
      +\frac n2\log\frac{\pi}{G_i}
      -\frac12\log\det(2Y)
      +\log\pi-\frac12\log(G_1G_3).
 \tag{F14}
\] Equations (F13)--(F14), the actual period \(C^2\), and the change in infinity coefficient weight must remain in any comparison of globally normalized rows. A constant finite input cost is not a zero total DP44 cost.

ALP-LIFT-11's fixed bilinear determinant repair preserves all values in (F8) at the identity and all scalar automorphic norms. ALP-LIFT-12 then shows that each repaired nonzero first vector survives some actual \(U(2)\to U(4)\) theta kernel. That theorem does not preserve the exact valuations (F10) at the second output.

The selected later harmonic row, its integral lattice, and the source DP44 gain and error fields therefore remain separate. This result supplies actual fixed-finite first rows on ordinary \(J^2\); it does not identify a DP65 conjugated-action product with those ordinary factors.

\subsection{Dependencies and checks}\label{dependencies-and-checks}

AM-K020, equations(K21.1)--(K21.1a), is the exact new derivative input; it is cited rather than reproved here. ALP-LIFT-09, (V16), (V22)--(V32), provides the actual oscillator, period and class-one character calculation. ALP-LIFT-10, (O6)--(O12), provides the exact permitted pullback normalization. AM-DO-020--021 supplies the optional explicit compact character and corrected DO023 the source vacuum shift. AM-K018 supplies the ordinary-power source-gauge identification.

This specialization is a complete proof from those inputs. Whole-proof checks are recorded separately with their actual scope; no Lean execution or second-output saturation is claimed.

\clearpage

\section{AM-DO-028: the even-rank Hom/tensor and doubling-character dictionary}\label{am-do-028-the-even-rank-homtensor-and-doubling-character-dictionary}

Author: GPT-6 Astra Pro, ALP\_MACHINE \nolinkurl{/root/definite_oscillator}; inherited reasoning effort not separately exposed; 21 November 2026 UTC.

This result fixes the character ratio in the Rallis factor for the literal bilinear source theta operator. It uses the actual Hom action and the determinant-root convention of AM-DO-022--026. It applies to even-rank edges with splitting characters trivial on the base-field ideles. It does not claim a new general local theta theorem, a full global chain, or an arithmetic estimate for the Schwartz functions obtained by nonvanishing theorems.

\subsection{1. The tensor model includes a dual source space}\label{the-tensor-model-includes-a-dual-source-space}

Write \(E/F\) for the CM extension, including the common host \(E'=\mathbf Q(\zeta_{28})\) when that host is used. Let \(S=E^a,T=E^b\) have their identity Hermitian forms, with both \(a,b\) even. Use the actual trace-zero coefficient \(\xi\) of the symplectic form; on the common host this is \(\xi'=\xi_J/2\). The underlying symplectic action and form are \[
 X\longmapsto hXg^{-1},\qquad
 B(X,Y)=-\operatorname{Tr}_{E/F}
                   (\xi\operatorname{tr}(Y^*X)).
 \tag{DO28.1}
\] Composition with \(\operatorname{Tr}_{F/\mathbf Q}\) gives the rational form used previously. The additive character and self-dual measures are transported with this form.

Put \(V=T\) and let \(U=S^\vee\) have the skew-Hermitian form \(-\xi\) times the dual identity Hermitian form, linear in the first variable. The map \[
 V\otimes_EU\longrightarrow\operatorname{Hom}_E(S,T),
                  v\otimes\ell\longmapsto v\ell
 \tag{DO28.2}
\] is a symplectic isomorphism. Indeed, on pure tensors the product of the two forms is \(-\xi(v'^*v)(\ell'^*\ell)\), whose trace is exactly (DO28.1) after expanding the matrix trace. The source group identification is \[
 \alpha:U(S)\longrightarrow U(U),\qquad
                            \alpha(g)=g^{-t}.
 \tag{DO28.3}
\] Thus the tensor action \(h\otimes\alpha(g)\) is exactly the Hom action. This is an actual representation isomorphism, with the same Heisenberg form, rather than an identification of two abstract groups followed by an unspecified oscillator twist.

For a character \(\chi\) trivial on \(\mathbb A_F^\times\), write \[
 \nu_\chi(z/\bar z)=\chi(z).
 \tag{DO28.4}
\] The standard tensor splitting changes on \(U(U)\) by \(\nu_{\Delta\chi_V}(\det g_U)\) and on \(U(V)\) by \(\nu_{\Delta\chi_U}(\det h)\). This is the usual Kudla/Harris--Kudla--Sweet character-change convention. GQT equation (2.2), printed p.13, fixes the positive source sign already in the zero-dimensional opposite-member case. Gan (2022), section2.5, printed pp.14--15, states the two character changes explicitly with the Hilbert-90 inverse.

Since \(\det\alpha(g)=\det(g)^{-1}\), pulling back this splitting through (DO28.3) gives the source convention \[
 \nu_{\Delta\chi_T}(\det g)^{-1}
                         \nu_{\Delta\chi_S}(\det h)
 \tag{DO28.5}
\] when and only when the tensor labels are \[
                     \boxed{\chi_U=\chi_S,\qquad\chi_V=\chi_T.}
 \tag{DO28.6}
\] The determinant roots, rather than the determinants themselves, are evaluated in (DO28.4). Functorial transport of the fixed standard splitting through the symplectic isomorphism fixes the baseline as well as the character changes; no free character is added. In the common-host construction the baseline is also the explicitly checked trivial-character model of DO022.

This does not contradict the corrected Konno--Konno dictionary in DO023. Their matrix source argument there was kept as \(g\), not replaced by \(g^{-t}\); its character label is therefore \(\chi_T^{-1}\). Changing that label while omitting (DO28.3) changes the representation.

\subsection{2. The input is conjugated and transported}\label{the-input-is-conjugated-and-transported}

Let \(F\) be an actual automorphic vector in a unitary representation \(\sigma\) of \(U(S)(\mathbb A_F)\). The source operator is \[
        T_\phi F(h)=\int_{[U(S)]}\theta_\phi(g,h)F(g)\,dg.
 \tag{DO28.7}
\] GQT section11.2, printed pp.51--52, uses the integral against the conjugated input. The bar was checked in the rendered PDF, since text extraction omits it. Define the input in tensor coordinates by \[
         f(u)=\overline{F(\alpha^{-1}u)},\qquad
         \pi=\overline\sigma\circ\alpha^{-1}.
 \tag{DO28.8}
\] With Haar measure pushed forward under the group isomorphism, the GQT theta integral of \(f\) is exactly (DO28.7), for the transported actual Schwartz vector. There is no unspecified volume scalar. For Tamagawa measures the same conclusion follows from invariance under the rational group isomorphism.

At the level of standard base-change parameters, contragredience from complex conjugation and contragredience from (DO28.3) cancel. More precisely, at an unramified place the automorphism \(g\mapsto g^{-t}\) sends the Satake matrix to its inverse transpose, while \(\overline\sigma=\sigma^\vee\) for the unitary representation. Hence \[
                         \phi_\pi=\phi_\sigma.
 \tag{DO28.9}
\] The same statement holds whenever the usual local Langlands parameter is available and its contragredient functoriality is used. The unramified statement suffices to identify the partial global standard Euler product. At a compact real place the two operations likewise return the original highest weight.

This is why the actual selected input in DO025 retains its old parameter block under a same-character lift. Its Hom oscillator component is the dual of the input; the tensor source component is the input after the simultaneous change (DO28.3), (DO28.8).

\subsection{3. An exact doubling-integral twist identity}\label{an-exact-doubling-integral-twist-identity}

This part can be checked directly in GQT, without guessing the meaning of a printed multiplication sign in an \(L\)-factor. Let \(\chi\) be a unitary character of \(E^\times\backslash\mathbb A_E^\times\), trivial on \(\mathbb A_F^\times\). On the doubled unitary group \(G(U\oplus U^-)\), define \[
                         d_\chi=\nu_\chi\circ\det.
 \tag{DO28.10}
\] For its Siegel Levi element \(m(A)\), the unitary determinant is \(\det A/\overline{\det A}\). Therefore \[
                       d_\chi(m(A))=\chi(\det A).
 \tag{DO28.11}
\] Multiplication by \(d_\chi\) gives an exact isomorphism between the normalized induced representations \[
 I(s,1)\longrightarrow I(s,\chi),\qquad
                        \Phi\longmapsto\Phi_\chi=d_\chi\Phi.
 \tag{DO28.12}
\] The normalization by the parabolic modulus is unchanged, and the multiplier is independent of \(s\). Standard, holomorphic, and good sections are transported accordingly. For good sections defined using the standard intertwiner, this also follows directly: \(d_\chi\) is trivial on the unipotent radical, so the intertwiner changes only by its fixed nonzero value at the chosen Weyl representative. No zero or pole in \(s\) is introduced.

On the embedded pair \((g,1)\), the multiplier is \(\nu_\chi(\det g)\). In GQT section11.6, printed p.54, the local zeta integral is the integral of this section against the \textbf{conjugated} matrix coefficient of \(\pi\). That bar was also checked in the rendered PDF. Consequently \[
 \begin{split}
 Z_\chi(s,\Phi_\chi,f_1,f_2;\pi)
 &=\int\Phi(g,1)\nu_\chi(\det g)
                    \overline{\langle\pi(g)f_1,f_2\rangle}\,dg\\
 &=Z_1(s,\Phi,f_1,f_2;
                       \pi\otimes(\nu_\chi\circ\det)^{-1}).
 \end{split}
 \tag{DO28.13}
\] This holds first in the common region of absolute convergence, and elsewhere by meromorphic continuation. It is an equality of the actual local integrals with unchanged Haar measure. The same identity holds for the global doubling integral: the automorphic character \(d_\chi\) is trivial on rational points, so it factors through the Eisenstein sum, and its two group factors cancel against the corresponding input twists. Here there is a necessary distinction between \(\chi\) and \(\nu_\chi\): for \(t\bar t=1\), \(\chi(t)=\nu_\chi(t)^2\), since \(t/\bar t=t^2\). Thus the actual factor in GQT's global integral on printed p.53 is \[
 d_\chi(g_1)d_\chi(g_2)\chi(\det g_2)^{-1}
              =d_\chi(g_1)d_\chi(g_2)^{-1}.
 \tag{DO28.13a}
\] Its product with \(\overline{f_1(g_1)}f_2(g_2)\) is \(\overline{(f_1d_\chi^{-1})(g_1)}
       (f_2d_\chi^{-1})(g_2)\), exactly as required for the representation in (DO28.14). This also checks the determinant-root normalization in the global identity.

GQT defines its local standard factor from these doubling integrals and cites the normalized GCD statement of Yamana in section11.6. Thus, in the even-character case here, its factor is the untwisted standard factor of \[
                   \pi_0=\pi\otimes(\nu_\chi\circ\det)^{-1}.
 \tag{DO28.14}
\] At nonarchimedean places the equality of good-section ideals in (DO28.13) gives the same normalized reciprocal polynomial factor; its constant-term-one normalization removes a Laurent unit ambiguity. The archimedean standard gamma-factor interpretation gives the same twist when expressed through the standard parameter. In particular, at every unramified place its standard Euler parameter is \[
                              \phi_\pi\otimes\chi^{-1}.
 \tag{DO28.15}
\] No claim that an arbitrary nonstandard \(L\)-factor convention uses the same printed symbol is necessary. Equations (DO28.12)--(DO28.15) fix the convention by the actual integrals.

\subsection{4. The factor in the source's labels}\label{the-factor-in-the-sources-labels}

Apply section3 with \(\chi=\chi_V=\chi_T\), and use (DO28.9). The Rallis factor for (DO28.7), expressed in standard base-change notation, is \[
 \boxed{L_E(s,\phi_\sigma\otimes\chi_T^{-1})
        =L_E(s,\phi_\sigma^\vee\otimes\chi_T).}
 \tag{DO28.16}
\] The equality follows from conjugation of \(E/F\): the standard parameter is conjugate self-dual, and \(\chi_T^c=\chi_T^{-1}\) because its restriction to the base field is trivial. Conjugating all places of \(E\) permutes the Euler factors. At a split place of \(F\), this permutes its two \(E\)-places, so the assertion is for the paired local standard factor, not for a single arbitrarily chosen component in isolation.

There is a useful check directly in the source model. Replacing \(\chi_T\) by \(\chi_T\delta\) multiplies the source kernel by \(\nu_\delta(\det g)^{-1}\). Replacing the input by \(F(g)\nu_\delta(\det g)\) makes the integrand unchanged. The parameter in (DO28.16) is invariant under this simultaneous change. A proposed parameter \(\phi_\sigma^\vee\chi_T^{-1}\) would fail this exact identity and is therefore incorrect for the fixed source convention.

For an already established standard parameter block \(\mu_r\otimes S_d\) in an input to edge \(j\), the corresponding character in (DO28.16) is \[
                              \mu_r\mu_j^{-1},
 \tag{DO28.17}
\] or its inverse after the simultaneous conjugation used in the right side of (DO28.16). Here \(S_d\) denotes the actual length-\(d\) segment with whichever standard/Arthur parameter description has been proved for that local lift. This calculation does not supply that local theta theorem. When the segment has the usual unramified exponents \((d-1)/2,(d-3)/2,\ldots,-(d-1)/2\), evaluating at \(s=(d+1)/2\) gives the Hecke factors at \(1,2,\ldots,d\), with the character (DO28.17).

For the selected DO025 exponents, \(\mu_j/\mu_r\) has infinity exponent \(-e_j+e_r\). If \(j>r\), the proved strict increase \(e_j-e_r\ge(j-r)(d+2)>0\) shows that this is a nontrivial angular character. This verifies the sign and nontriviality needed when applying a separate Hecke \(L(1)\)-nonvanishing theorem. It is not a proof of that theorem or a quantitative lower bound for its value.

\subsection{5. Exact scope and sources}\label{exact-scope-and-sources}

The evenness restriction is essential to the simple global character \(d_\chi\) in (DO28.10): this argument assumes \(\chi|\mathbb A_F^\times=1\). It does not silently extend to an odd opposite rank, where a fixed reference character must first be divided out. It preserves the actual host, polarization, additive character, determinant roots, Haar measures, and bilinear input convention. A Rallis nonvanishing theorem still requires its local occurrence, global range, and cuspidality hypotheses. It provides no selected arithmetic row or bound on the later finite Schwartz input.

Primary locators used here:

\begin{itemize}
\tightlist
\item
  Gan, W. T., Qiu, Y., \& Takeda, S. (2014). \emph{The regularized Siegel--Weil formula (the second term identity) and the Rallis inner product formula}. Inventiones Mathematicae, 198, 739--831. The supplied arXiv version1207.4709, section2.8 and equation(2.2), printed pp.12--13; section11.2, pp.51--52; section11.5 and section11.6, pp.53--54; Theorem11.4, pp.54--55. The PDF was inspected for the bars omitted by its text extraction.
\item
  Gan, W. T. (2022). \emph{Automorphic forms and the theta correspondence}. Arizona Winter School lecture notes, section2.5, pp.14--15, character changes; sections2.6 and3.8, pp.15--16 and26, input/quotient convention; Theorem2.6, p.18, and Theorem3.3, p.27, the rank-one inverse-character factor. \url{https://swc-math.github.io/aws/2022/2022GanNotes.pdf}
\item
  AM-DO-022, DO023 as corrected, DO025, and DO026 for the actual Hom splitting and source integral. GQT's published results are used at their stated hypotheses, not re-proved.
\end{itemize}

The full source-to-doubling twist calculation above is new work in this run; its whole-proof check is recorded separately.

\clearpage

\section{ALP-LIFT-14 --- an actual selected even-rank chain with unrestricted later finite inputs}\label{alp-lift-14-an-actual-selected-even-rank-chain-with-unrestricted-later-finite-inputs}

Author: GPT-6 Astra Pro, ALP\_MACHINE /root/lift\_accounting, 21 November 2026 UTC. The inherited reasoning effort is not separately exposed.

This theorem starts from the actual even-rank first edge AM-R-012 and continues through the literal selected ranks and compact Schur types. It uses the common mixed CM host \(E'=\mathbf Q(\zeta_{28})\), the same character on both members of each edge, and the source's bilinear theta integral. Its later finite Schwartz vectors are unrestricted. Thus it proves actual nonzero iterated theta images in these geometric groups; it does not instantiate the six comparison identities, the old elliptic arithmetic row, or their finite contents and global sizes. In particular it does not prove a shorter chain or a DP44 improvement.

The accepted rank and weight recurrences are inputs. The new argument is the character construction, exact partial-standard-factor induction and boundary nonvanishing, combined with local occurrence and a multiplicity-safe global induction. Literature theta and Rallis theorems are cited at their actual hypotheses, not re-proved or counted as new findings.

\subsection{1. Exact statement and the actual common host}\label{exact-statement-and-the-actual-common-host}

Use the host, product polarization, additive character, trace-one vacuum and Hom oscillator of AM-DO-022: \[
 E'=\mathbf Q(\zeta_{28}),\quad F'=(E')^+,\quad
 A=J^2,\quad J=\operatorname{Jac}(y^2=x^7-1),
 \qquad G_j=U(I_{N_j}).
 \tag{E1}
\] Here the rational trace coefficient is \(\xi'=(\zeta_7-\zeta_7^{-1})/14\). The identity Hermitian spaces are positive at every real place. The extension \(E'/F'\) is unramified at every finite place, and all its dyadic places split. These actual host facts, including the integral vacuum normalization, are carried from DP57 and AM-DO-022.

Take an actual source Ranks record and HCParam for which the partitions of AM-DO-025 are defined. In particular \[
 d=t+1,\quad N_j=k+jd\quad(0\le j\le m),\qquad
 k,m,d\text{ even},\quad m\ge2,\quad d\ge k+m+5.
 \tag{E2}
\] Write \(w_{j,\tau}\) for the selected highest weights, and \(\rho_j=\boxtimes_{\tau\in\Phi'}\rho_{w_{j,\tau}}\), where \[
 \Phi'=\{1,3,9,15,17,23\}.
\] Use the edge indexing \(j=0,\ldots,m-1\). The common character on edge \(j\) is denoted \(\mu_j\); thus it is the character called \(\mu_{j+1}\) in the one-based edge notation of AM-DO-025. The exact carried recurrences are \[
 \begin{split}
 \nu_{j,\tau}
   &=w_{j,\tau}+e_{j,\tau}\mathbf1
                          -N_{j+1}\mathbf1/2,\\
 w_{j+1,\tau}
   &=(w_{j,\tau}-d/2,\,
                     (-e_{j,\tau}+N_j/2)^d),\\
 e_{j+1,\tau}-e_{j,\tau}&\ge d+2.
 \end{split}
 \tag{E3}
\] Each \(\nu_{j,\tau}\) is a nonnegative partition. The exact selected polynomial component of the oscillator on edge \(j\) is \[
                    \rho_j^\vee\boxtimes\rho_{j+1}.
 \tag{E4}
\] It has polynomial degrees \(|\nu_{j,\tau}|\). Neither its degree nor its partition is replaced by an unrestricted compact type.

\textbf{Theorem.} There are globally admissible common characters \(\mu_j\) with infinity exponents \(-e_{j,\tau}\), all with finite compact conductor dividing a fixed split place above29, and there are actual nonzero automorphic forms \(F_j\) of compact type \(\rho_j\), \(1\le j\le m\), such that \(F_1\) is the first lift of AM-R-012 and \[
 F_{j+1}(h)=
       \int_{[G_j]}\vartheta_{\phi_j}(g,h)F_j(g)\,dg
                        \ne0\qquad(1\le j<m).
 \tag{E5}
\] At infinity \(\phi_j\) lies in the exact component (E4); at finite places it can be chosen to be a finite linear combination of genuine Schwartz vectors. No bound for that combination's finite level, coefficients, or norm is claimed.

One may choose the induction so that the representation generated by \(F_j\) is spherical at all nonsplit finite places and at every split place outside29. More strongly, every nonzero vector in the resulting generated automorphic module survives some next finite Schwartz input in the selected real component. This last assertion concerns actual automorphic embeddings, including their multiplicities.

\subsection{2. Global characters and strict nontrivial ratios}\label{global-characters-and-strict-nontrivial-ratios}

Let \(T'\) be the norm-one torus. AM-R-009 proves, using the class number one of \(E'\), the actual factorization \[
 T'(\mathbb A_f)=T'(F')K_{T,0},\qquad
 T'(F')\cap K_{T,0}=\mu_{28}.
 \tag{E6}
\] Its ideal proof and the cited Masley class-number statement were independently checked in review\_root\_unramified\_first\_lift.md. We use that result, not a new class-number assertion.

For edge \(j\) put \(\beta_j(t)=\prod_{\tau\in\Phi'}t_\tau^{-e_{j,\tau}}\) on \(T'(\mathbb R)\). At one split29 factor choose the residue embedding \(\zeta_{28}\mapsto2\); 2 has order28 in \(\mathbf F_{29}^{\times}\). Define a character of that residue group by \[
 \varrho_j(2)=
 \exp\!\left(\frac{2\pi i}{28}
                    \sum_{\tau\in\Phi'}e_{j,\tau}\tau\right),
 \tag{E7}
\] and take it trivial on all other compact factors. Its restriction to \(\mu_{28}\) is \(\beta_j^{-1}\). Therefore (E6) gives a unitary automorphic torus character \(\gamma_j\) with infinity component \(\beta_j\) and compact component \(\varrho_j\): to define it on an adelic point, factor its finite part as \(qk\), move the rational \(q\) to the infinite component, and evaluate \(\beta_j\varrho_j\). Changing the factorization changes the two components by one element of \(\mu_{28}\), on which their product is one.

Set \[
 \mu_j(x)=\gamma_j(x/\bar x),\qquad
             \chi_{S,j}=\chi_{T,j}=\mu_j.
 \tag{E8}
\] This character is trivial on \(\mathbb A_{F'}^\times\), has \(\mu_{j,\tau}(z)=(z/\bar z)^{-e_{j,\tau}}\), and satisfies the even-rank splitting restrictions. At a nonsplit finite place the entire local norm-one torus is compact: its elements have valuation zero. Consequently \(\gamma_{j,v}\), and hence \(\mu_{j,v}\), is identically one there. At split places other than29 it is unramified. Its compact finite values are roots of unity of order dividing28.

For \(r<j\), define \[
                           \upsilon_{rj}=\mu_r\mu_j^{-1}.
 \tag{E9}
\] Its infinity exponent at every \(\tau\) is \(e_{j,\tau}-e_{r,\tau}>0\). Thus both \(\upsilon_{rj}\) and \(\upsilon_{rj}^2\) have nontrivial angular infinity type. In particular neither is a norm character. This strict condition will exclude the pole and the possible boundary-zero obstruction in section5. It follows from the actual recurrence (E3); it is not obtained by changing source exponents.

\subsection{3. Local theta occurrence and the unramified parameter recurrence}\label{local-theta-occurrence-and-the-unramified-parameter-recurrence}

Let \(a=N_j\), \(b=a+d\). At a split finite place, the groups are \(\mathrm{GL}_a\) and \(\mathrm{GL}_b\). Mínguez (2008), Theorem1(1)--(2), printed p.718, proves nonzero small theta for every irreducible input when \(a\le b\), including dyadic residual characteristic. We use its actual type-II representation, with the Fourier transform between the two matrix polarizations; equation(1.4), printed pp.721--722, fixes the half-modulus factors.

For precision, write the local common character as \(\mu=(\eta,\eta^{-1})\) on \(E'_v=F'_v\times F'_v\). The Hom oscillator changes by \[
                     \eta_T(\det g)^{-1}\eta_S(\det h).
 \tag{E10}
\] The literal bilinear integral pairs the source oscillator with the contragredient of the actual input \(\pi\). After removing (E10), the type-II source representation is \(\pi^\vee\otimes\eta_T\). Mínguez's target is its dual together with the standard rank-\(d\) trivial-representation segment, and restoring the target twist gives \[
 \phi_{\rm out}
       =(\phi_\pi\otimes\eta_S\eta_T^{-1})\,
                       \sqcup\,\eta_S[d].
 \tag{E11}
\] Thus the old block is unchanged when \(\eta_S=\eta_T\). This checks the sign by the actual Hom action, not by an unspecified theta-convention formula.

At nonsplit finite places the residue characteristic is odd and all \(\mu_{j,v}=1\). The even identity Hermitian lattices are integrally hyperbolic. Indeed the norm map on units of an unramified quadratic extension is surjective, so the binary identity form contains a primitive isotropic vector and splits a unimodular hyperbolic plane; iteration gives the assertion. The additive character used in the actual host may have conductor inherited from the different of \(F'/\mathbf Q\). This does not change the unramified comparison: write the actual pure-imaginary trace coefficient as \(c\delta\), with \(\delta\) a pure-imaginary unit and \(c\in F_v'^\times\). The absolute-codifferent condition and the trivial relative different imply that \(\psi_{F'}(c\,\cdot)\) has integral conductor. Rescaling the skew form and additive character together preserves their composite symplectic character, the lattice model and the unitary groups. This is the integral model already fixed in AM-DO-022.

Apply Liu (2021), Definition2.8 and Theorem6.2(1)--(2), printed pp.7 and21--22, with sign \(+\), source rank \(2r=a\), target rank \(2s=b\), and trivial splitting character. Proposition3.4, printed p.9, verifies its Assumption3.3 in odd residue characteristic. For every spherical source representation, the target small theta is the spherical representation stated there, and the characteristic vector of the selfdual lattice has nonzero image. No temperedness of the source is required. The extra standard eigenvalues are \(q_F^{\pm1},q_F^{\pm3},\ldots,q_F^{\pm(d-1)}\). The same powers appear in Liu (2011), Appendix, printed p.916, but we use the stronger 2021 theorem to ensure that every specialization survives; an annihilator identity alone would not prove that.

Write \(Q_w=\#(\mathcal O_{E'}/w)\). At an unramified place let \(\mathcal A_{j,w}\) be the standard Satake multiset of an actual constituent \(\pi_j\). We use the convention that an inducing character \(\alpha|\cdot|^u\) has eigenvalue \(\alpha(\varpi_w)Q_w^{-u}\). The two local calculations give exactly \[
 \mathcal A_{j+1,w}
   =\mathcal A_{j,w}\ \sqcup\
       \left\{\mu_{j,w}(\varpi_w)
              Q_w^{(d-1)/2-\ell}:0\le\ell<d\right\}.
 \tag{E12}
\] At a nonsplit place \(Q_w=q_F^2\), giving precisely the odd \(q_F\)-powers above; inverse paired Satake parameters absorb the contragredient convention. At a split place (E11) gives each chosen component, with its conjugate component included as well. This is a formula for unramified parameters. We make no claim that the same segment describes compact real Langlands parameters.

For each actual descendant irreducible constituent choose one ancestral initial constituent \(\pi_0\) in the first input. Iteration along that ancestry gives \[
 \mathcal A_{j,w}
  =\mathcal A_{0,w}\ \sqcup\
    \coprod_{0\le r<j}
     \left\{\mu_{r,w}(\varpi_w)
                  Q_w^{(d-1)/2-\ell}:0\le\ell<d\right\}.
 \tag{E13}
\] The starting class-supported vector can contain several initial constituents. Formula(E13) is not an assertion that one common \(\pi_0\) describes the whole generated automorphic module; the unitary bound in (E18) is uniform over all such ancestors.

\subsection{4. Exact Rallis factor in the literal source convention}\label{exact-rallis-factor-in-the-literal-source-convention}

AM-DO-028 is used here with its full Hom/tensor and conjugation dictionary. In tensor coordinates the source is \(S^\vee\) and its group map is \(\alpha(g)=g^{-t}\). GQT's conjugated-input convention is obtained from the actual bilinear input by \(f=\overline{F\circ\alpha^{-1}}\). The two contragrediences cancel on standard unramified parameters. Its tensor character is \(\chi_V=\chi_T\).

There is also a direct check of the twist. For \(\chi|_{\mathbb A_{F'}^\times}=1\), define \(\nu_\chi(z/\bar z)=\chi(z)\). On the doubled unitary group multiplication by \(d_\chi=\nu_\chi\circ\det\) identifies \(I(s,1)\) with \(I(s,\chi)\), because its restriction to the Siegel Levi \(A\) is \(\chi(\det A)\). The local doubling integral uses a conjugated matrix coefficient. Hence \[
 Z_\chi(s,d_\chi\Phi;\pi)
      =Z_1(s,\Phi;\pi\otimes(\nu_\chi\circ\det)^{-1}).
 \tag{E14}
\] The good-section ideals and normalized standard factors transform accordingly. In the global integral the factor on the second group is \(\chi(\det g_2)^{-1}\), and on the norm-one torus \(\chi(t)=\nu_\chi(t)^2\); thus the two doubled determinant factors cancel with the two input twists exactly. This verifies that no square root or sign is suppressed.

Accordingly the standard partial factor in source labels is \[
 \mathcal L_j^S(s)
   =L_{E'}^S(s,\operatorname{BC}(\pi_j)\otimes\mu_j^{-1}).
 \tag{E15}
\] Here BC notation denotes the standard unramified Euler parameters. A global automorphic base-change theorem is not being assumed. Choose a finite set \(S\) containing all ramified places and omit both \(E'\)-places over every omitted split \(F'\)-place. Conjugating \(E'/F'\) gives the equivalent factor \(L^S(s,\operatorname{BC}(\pi_j)^\vee\otimes\mu_j)\). The actual formula has ratios \(\mu_r\mu_j^{-1}\), not products \(\mu_r\mu_j\).

Substituting (E13) gives an identity of Euler products in a right half-plane, and therefore of their continuations: \[
 \mathcal L_j^S(s)=
 L_{E'}^S(s,\operatorname{BC}(\pi_0)\otimes\mu_j^{-1})
 \prod_{r<j}\prod_{\ell=0}^{d-1}
 L_{E'}^S(s-(d-1)/2+\ell,\upsilon_{rj}).
 \tag{E16}
\] The Rallis point for edge \(j\) is \(s_0=d/2\) in the doubled induction and \(s_*=(d+1)/2\) in the standard factor. Therefore \[
 \mathcal L_j^S(s_*)=
 L_{E'}^S(s_*,\operatorname{BC}(\pi_0)\otimes\mu_j^{-1})
                \prod_{r<j}\prod_{\ell=1}^{d}
                      L_{E'}^S(\ell,\upsilon_{rj}).
 \tag{E17}
\]

\subsection{5. The exact nonvanishing needed in (E17)}\label{the-exact-nonvanishing-needed-in-e17}

The initial representation need not be tempered. It is unitary because it occurs in the automorphic Hilbert space of a compact quotient. Its unramified spherical matrix coefficient has absolute value at most one. Macdonald's bounded-spherical-function criterion says that the logarithmic real Satake parameter lies in the convex hull of the Weyl orbit of \(\rho\). In the standard degree-\(k\) representation this gives \[
                         |\alpha_w|\le Q_w^{(k-1)/2}.
 \tag{E18}
\] For a split place this is the familiar \(\mathrm{GL}_k\) bound. For unramified \(U(k)\), the relative \(\rho\)-coordinates are \(k-1,k-3,\ldots,1\) in \(q_F\)-units and \(Q_w=q_F^2\), giving the same bound. The unitary center contributes no extra real exponent. Exact references are Macdonald (1971), ChapterIV, section4.7, Theorem4.7.1; Macdonald (1995), ChapterV, section3, Example, printed pp.299--300 supplies the split specialization.

By (E2), \[
 s_*-(k-1)/2=(d-k+2)/2\ge(m+7)/2>1.
 \tag{E19}
\] Thus the initial Euler product in (E17) is absolutely convergent and nonzero. Indeed its logarithmic series is bounded by a constant times \(k\sum_w Q_w^{-(d-k+2)/2}\), with all higher powers dominated by a convergent geometric series. Every individual factor is finite and nonzero. This argument does not assume a Ramanujan bound for \(\pi_0\).

For \(\ell\ge2\), the Hecke factors in (E17) are absolutely convergent and nonzero because \(\upsilon_{rj}\) is unitary. The boundary \(\ell=1\) requires a separate argument.

\textbf{Boundary lemma.} Let \(\upsilon\) be a unitary Hecke character with nontrivial angular infinity type, such that \(\upsilon^2\) also has nontrivial angular infinity type. Then \(L^S(1,\upsilon)\) is finite and nonzero.

\textbf{Proof.} Hecke--Tate continuation gives holomorphy at1 for both \(L^S(s,\upsilon)\) and \(L^S(s,\upsilon^2)\). No norm-character pole occurs because of their angular types. Omitting finitely many local factors preserves holomorphy there. For real \(\sigma>1\), expand the absolutely convergent Euler logarithms: \[
 \begin{split}
 3\log\zeta_{E'}^S(\sigma)
 +4\log|L^S(\sigma,\upsilon)|
 +\log|L^S(\sigma,\upsilon^2)|
 =\sum_{w\notin S}\sum_{n\ge1}
 \frac{3+4\operatorname{Re}(\upsilon(w)^n)
              +\operatorname{Re}(\upsilon(w)^{2n})}
                     {nQ_w^{n\sigma}}\ge0.
 \end{split}
 \tag{E20}
\] The inequality follows termwise from \(3+4\cos\theta+\cos2\theta=2(1+\cos\theta)^2\). Consequently \[
 \zeta_{E'}^S(\sigma)^3
 |L^S(\sigma,\upsilon)|^4|L^S(\sigma,\upsilon^2)|\ge1.
 \tag{E21}
\] If \(L^S(s,\upsilon)\) had a zero of order \(r\ge1\) at1, the left side would be \(O((\sigma-1)^{4r-3})\), since the zeta function has a simple pole and the last factor is bounded. It would tend to zero, a contradiction. This proves the lemma. The argument asserts neither absolute convergence nor an Euler-product limit at1. ∎

The continuation input is used only for nontrivial angular characters; see Prasad (2007), Lectures on Tate's thesis, section1, Theorem2, printed p.3. The elementary positivity proof above supplies the needed boundary nonvanishing. By (E9) it applies to every \(\upsilon_{rj}\). Thus \[
                   \mathcal L_j^S(s_*)\ne0
 \quad\text{and it is holomorphic at }s_*.
 \tag{E22}
\]

\subsection{6. Rallis hypotheses, including the real local test}\label{rallis-hypotheses-including-the-real-local-test}

For every later edge \(j\ge1\), (E2) gives \[
 a=N_j\ge k+d>d,\qquad
       a<b=a+d<2a,\qquad 0<s_0=d/2<a/2.
 \tag{E23}
\] The target identity Hermitian space is positive at every real place, hence anisotropic over \(F'\): an isotropic rational vector would remain isotropic under a real embedding, contradicting positivity. Its global Witt index is therefore \(r=0\). The groups have no proper rational parabolics, their adelic quotients are compact, and their automorphic representations are cuspidal. In GQT's notation \(\epsilon_0=0\), \(d(n)=n=a\), \(m=b\), and the hypothesis that lower lifts vanish for Witt indices less than \(r\) is vacuous. Although the numerical dimensions lie in the interval treated by the second-term theorem, the actual positive target with \(r=0\) is in its convergent case. GQT Proposition11.1(ii) and Theorem11.4 explicitly include this case.

At a nonsplit finite place, the spherical local occurrence was proved in section3. At a split finite place, including the possibly ramified representation at29 and all dyadic places, Mínguez gives occurrence for every irreducible input. GQT Proposition11.6(i) therefore gives a nonzero normalized local doubling test on the actual target's Siegel--Weil image at every finite place.

At a real place we do not invoke the incomplete converse discussed in GQT Theorem11.7. The actual selected Fock summand is (E4). Choose a unit vector in the target factor and an isometric copy of the source contragredient in that summand. The doubling isomorphism satisfies \(\delta(\phi_1\otimes\overline{\phi_2})(0)=\langle\phi_1,\phi_2\rangle\), with the determinant factor retained in (E14). Choose the source vector and this Fock vector so that the oscillator coefficient is the conjugate of the input coefficient in the GQT integral. Schur orthogonality then gives \[
 Z_\tau(s_0)=
 \int_{U(a)}|\langle\rho_j(g)v,v\rangle|^2\,dg
                 =\frac{\operatorname{vol}(U(a))}
                         {\dim\rho_{j,\tau}}>0.
 \tag{E24}
\] This is the same direct compact calculation proved in ALP-LIFT-12, now with an arbitrary selected Schur representation in place of a symmetric power. The standard real gamma factors are finite and nonzero at \(s_*>0\): a compact-group parameter is tempered, the character twist is angular and unitary, and its complex gamma arguments are \(s_*\) plus nonnegative half-integer angular sizes. Hence the normalized local test also is nonzero.

The construction uses the selected Schur component, not necessarily the identical highest polynomial at every intermediate evaluation. To pair a given source highest vector, a compact Weyl translate of the polynomial factor may be needed. This preserves its partition, total degree and unitary norm. Identifying one rigid highest-vector pin throughout the source's six identities is an additional comparison, not a consequence of (E24).

All local theta lifts are now nonzero. GQT Theorem11.4(ii) implies holomorphy of the completed standard factor at \(s_*\). Equation(E22) gives a nonzero partial factor. The finitely many missing nonarchimedean standard factors are reciprocal polynomials with constant term one and have no zeros; the real gamma factors have no zeros either. Holomorphy of the completed factor therefore rules out any pole among them at this point and implies \[
                             L(s_*,\pi_j\times\chi_V)\ne0.
 \tag{E25}
\] This deduction does not assume temperedness of the later representations.

The normalized unramified test is \(1/d_v(s_0,\chi_V)\). Its infinite product converges and is nonzero because \(s_0>0\), exactly as in GQT equation(11.3), printed p.54. Combining (E24), the finite local tests, and (E25), GQT Theorem11.4 gives a nonzero global theta lift on every actual irreducible cuspidal embedding satisfying the stated unramified parameters. The use is of the theorem's actual representation and measure, through (E14); it is not an inference from a formal Euler factor alone.

\subsection{7. Actual vectors, multiplicities and retention of the spherical hypothesis}\label{actual-vectors-multiplicities-and-retention-of-the-spherical-hypothesis}

The preceding paragraph is stronger than existence of one abstract automorphic representation with a nonzero lift. We give the vector argument because actual iterated images are the issue.

Let \(\mathcal H_j\) be the closed \(G_j(\mathbb A)\)-invariant automorphic subspace generated by the actual previous theta images with the prescribed real types. Every constituent in it has the unramified parameters (E13) for an ancestral initial constituent, by uniqueness of the local small theta quotient and the local calculations in section3. The actual starting vector has fixed finite level and real type, so it has only finitely many initial constituents. No common initial constituent is required: sections5--6 apply separately to each descendant, with the uniform bound (E18). Fix the real selected component (E4), while allowing all finite Schwartz vectors and all vectors of that component.

Let \(\mathcal N_j\subset\mathcal H_j\) be the joint kernel of all these next theta maps. The automorphic quotients are compact, so the maps defined by each smooth Schwartz kernel are continuous on the relevant Hilbert spaces. The joint kernel is closed. Equivariance of the Weil representation makes it invariant under \(G_j(\mathbb A)\): translating an input is absorbed by translating the finite Schwartz vector and its real selected component. If \(\mathcal N_j\ne0\), discrete spectral decomposition of the compact automorphic quotient supplies an actual irreducible cuspidal embedding contained in \(\mathcal N_j\). The Rallis argument of section6 applies to that embedding, with its inherited positive invariant inner product, and says that some selected theta map is nonzero on it. This contradicts its definition. Therefore \[
                              \mathcal N_j=0.
 \tag{E26}
\] This handles automorphic multiplicity copies rather than implicitly identifying a representation with one preferred embedding.

It follows that every nonzero current vector survives some next Schwartz vector. The latter may be a finite sum of decomposable vectors; a nonzero sum has a nonzero summand after any needed linear choices. Its output has type \(\rho_{j+1}\).

The next induction step needs spherical input at nonsplit places. This can be retained without asserting an unsupported fixed-vacuum global nonvanishing. Local uniqueness and Liu's theorem show that every output constituent is spherical at every nonsplit place; at split places outside29 the unramified type-II theorem gives the same conclusion. For a nonzero smooth output, only finitely many local compact invariance conditions need be restored. At one such place, if all translates of the vector had zero projection to the maximal-compact invariants, its cyclic subrepresentation would have no spherical constituent, a contradiction. Thus a finite translate followed by the compact averaging projector gives a nonzero vector invariant there. Averaging may be combined at the finite set of places. The operations commute with the other local actions, and equivariance absorbs them into the finite Schwartz input. All but finitely many components were already the vacuum, so there is no infinite operation.

The resulting finite input may be a Hecke combination of the original lattice vector. This argument does not say that it remains a single prescribed move, nor that its coefficients are units. It proves exactly the spherical condition required by section3 and no finite arithmetic estimate.

\subsection{8. Initialization and the resulting chain}\label{initialization-and-the-resulting-chain}

AM-R-012 supplies the actual starting vector, using the accepted DP59 half-orbit and the new interpolation result AM-K022. It does not use a rank-one precursor. For \(a=k\), \(b=k+d\), \(D=\sum_\tau|\nu_{0,\tau}|\), choose \[
 M=29^{\max(1,\lceil\log_{29}(3D+2)\rceil)}.
\] There is one permitted translate of finite Schwartz norm one with nonzero normalized selected jet \(a_x\), and \[
 v_{11}(a_x)\le
                  \sum_{\tau,i}v_{11}(\nu_{0,\tau,i}!).
 \tag{E27}
\] The derivative factorial is retained. The exact source input \(f_0\) is supported on the identity finite class at principal level \(M\); its compact component is the primitive highest matrix coefficient. With the actual quotient measure, \[
 \begin{split}
 I_a(M)&=M^{6a^2}\prod_{i=1}^a(1-29^{-i})^6,\\
 \|f_0\|_2^2&=(I_a(M)\dim\rho_0)^{-1},\\
 F_1(1)&=\frac{C^{2ab}a_x}{I_a(M)\dim\rho_0}\ne0.
 \end{split}
 \tag{E28}
\] The period \(C=\theta[0;0](0;\tau)c\) is the actual sigma-normalized period of AM-DO-022, not a newly chosen row scalar. The first vector is invariant under the standard compact outside29. Its generated constituents are therefore spherical there and satisfy (E12) for \(j=0\). Its source representation is unitary, so the crude initial bound (E18) applies.

Starting with this actual \(F_1\), apply (E26) successively for \(j=1,\ldots,m-1\), retaining spherical invariance as in section7. This produces the actual nonzero chain (E5). Since the finite Schwartz space has a basis over an algebraic coefficient extension containing the finite splitting values, a nonzero finite linear functional cannot vanish on all such basis vectors; the finite choices may be taken with algebraic coefficients. This observation gives no bound on their height or denominator.

The concrete HC-only example in AM-R-012 has \(6\to20\), \(D=36\), \(M=841\), \(\dim\rho_0=1\), and \(v_{11}(I_6(M))=0\). Its first algebraic identity coefficient has valuation at most zero. That example does not impose the separate noncompact Hyp3 conditions and is not evidence for a later gain. It is consistent with the present existence theorem, which transports nonzero vectors but does not transport their valuations.

\subsection{9. What remains for Q1}\label{what-remains-for-q1}

This result supplies an actual common-host even-rank chain and removes the earlier unrestricted-global-nonvanishing gate. It retains the exact source selected partitions, literal bilinear integral and equal-character requirement. It does not supply the source's full D7 constructor, its six AdmS identities, or an identification with the prescribed arithmetic rowAF.

For every later edge the missing arithmetic object is the restriction of the class-integrated theta matrix to the already produced input-image lattice, with its selected coefficient normalization and compact contraction included. AM-K022's finite jet matrix has the same determinantal ideals as the corresponding derivative-function matrix before that integration; it does not imply that the restricted class matrix is saturated. The finite local tests supplied by Rallis need not lie in the prescribed DP59/DP62 moved family. Their coefficients, common denominator, finite level, archimedean norm, and content can affect every \(g_j,e_j,H\) term in DP44.

Thus this theorem establishes actual analytic existence with exact local types. It provides no implication from individual derivative slopes, or from \(L(1)\ne0\), to a smaller required number of lifts. A quantitative lower bound for those \(L\)-values would itself be insufficient without the integral image and row comparisons.

\subsection{Primary sources and exact locators}\label{primary-sources-and-exact-locators}

\begin{itemize}
\tightlist
\item
  Gan, W. T., Qiu, Y., \& Takeda, S. (2014). The regularized Siegel--Weil formula (the second term identity) and the Rallis inner product formula. Inventiones Mathematicae, 198, 739--831. arXiv:1207.4709v3, sections11.2--11.8, printed pp.51--56; Proposition11.1(ii), Theorem11.4, Proposition11.6, equation(11.3). The conjugation bars on pp.52 and54 were checked in the rendered PDF. Local copy: lift\_refs/gqt2014.pdf.
\item
  Liu, Y. (2021). Theta correspondence for almost unramified representations of unitary groups. arXiv:2105.14223v1; published Journal of Number Theory, 230 (2022), 196--224. Definition2.8, Proposition3.4, Theorem6.2, printed pp.7,9,21--22. The sign \(+\), arbitrary spherical-parameter specialization is the one used here. Local copy: lift\_refs/liu2021\_unramified\_theta.pdf.
\item
  Liu, Y. (2011). Arithmetic theta lifting and L-derivatives for unitary groups, I. Algebra \& Number Theory, 5(7), 849--921. Appendix, printed pp.912--918, especially the map p.916 and PropositionA.5. This is an independent normalization check, not the sole proof of every spherical specialization. \url{https://msp.org/ant/2011/5-7/ant-v5-n7-p01-p.pdf}.
\item
  Mínguez, A. (2008). Correspondance de Howe explicite: paires duales de type II. Annales Scientifiques de l'École Normale Supérieure, 41(5), 717--741. Theorem1, p.718; equation(1.4), pp.721--722. \url{https://www.numdam.org/item/ASENS_2008_4_41_5_717_0.pdf}.
\item
  Macdonald, I. G. (1971). Spherical functions on a group of p-adic type. Ramanujan Institute Publications2. ChapterIV, section4.7, Theorem4.7.1, bounded-spherical-function criterion. Primary chapter transcribed by Arun Ram: \url{https://math.soimeme.org/~arunram/Resources/SFOAGOPTChapterIVCalculationOfTheSphericalFunctions.html}.
\item
  Macdonald, I. G. (1995). Symmetric functions and Hall polynomials (2nd ed.). Oxford University Press. ChapterV, section3, Example, printed pp.299--300. \url{https://math.berkeley.edu/~corteel/MATH249/macdonald.pdf}.
\item
  Prasad, D. (2007). Lectures on Tate's thesis. ICTP lecture notes, section1, Theorem2, printed p.3, used only for continuation and absence of a pole for a character with nontrivial angular infinity type. \url{https://www.math.iitb.ac.in/~dprasad/ictptate.pdf}.
\item
  AM-DO-022, AM-DO-025--026 and AM-DO-028 for the actual host, compact selected summands, source integral and doubling dictionary; AM-R-009 and AM-R-012 for the class-one character construction and actual first edge; AM-K022 for the exact permitted first jet. Their claims remain separately identified and are not counted again as new results.
\end{itemize}

\clearpage

\section{ALP-LIFT-15 --- one finite Heisenberg translate at every selected edge}\label{alp-lift-15-one-finite-heisenberg-translate-at-every-selected-edge}

Author: GPT-6 Astra Pro, ALP\_MACHINE /root/lift\_accounting, 21 November 2026 UTC. Inherited reasoning effort is not separately exposed.

This refines ALP-LIFT-14's unrestricted finite Schwartz existence. Every later nonzero lift can be obtained with \textbf{one actual finite Heisenberg translate of the normalized vacuum}, with finite \(L^2\) norm exactly one and coefficient vector consisting of a single one. The required finite translation denominator is not bounded. The full orbit here is larger than the prescribed DP59 half-orbit, and the resulting rational half-model function need not equal the source's raw Xi. Those restrictions remain genuine arithmetic obligations.

The proof below gives the complete finite basis calculation, its quantitative norm bound at a specified denominator, and its application to the actual selected chain. It does not invoke Stone--von Neumann irreducibility in place of a finite calculation.

\subsection{1. The actual finite lattice model}\label{the-actual-finite-lattice-model}

On an edge \(U(a)\to U(b)\) in the common \(E'=\mathbf Q(\zeta_{28})\) host, let \(W\) be the underlying rational symplectic space and let \[
             2g=\dim_{\mathbf Q}W=12ab,\qquad g=6ab.
 \tag{H1}
\] The actual lattice is the trace lattice on \(\operatorname{Hom}_{\mathcal O_{E'}}(\mathcal O_{E'}^a,\mathcal O_{E'}^b)\), with the alternating pairing and trace-one semicharacter of AM-DO-022. It is selfdual. Write \(\widehat\Lambda\subset W(\mathbb A_f)\) for its finite completion, and \(\phi_0\) for the actual normalized finite vacuum, so \(\|\phi_0\|_2=1\).

Let \(W_x\) denote the actual finite Weyl operator, with the source's additive character and selfdual measure. Fix the commutator convention \[
                 W_xW_y=\psi_f(B(x,y))W_yW_x.
 \tag{H2}
\] Changing both signs gives the same statements below. The actual vacuum character is denoted \(\alpha\): \[
                         W_\lambda\phi_0
                            =\alpha(\lambda)\phi_0
                 \quad(\lambda\in\widehat\Lambda).
 \tag{H3}
\] Its semicharacter law is exactly the one forced by the Weyl cocycle. The trace-one construction supplies a continuous sign character in the appropriate quadratic sense; its values on the lattice are roots of unity. This is the actual source vacuum, not an arbitrary integral frame for its Chern-reduced analytic germ.

For a positive integer \(N\), define the finite-dimensional eigenspace \[
 \mathcal S_N=
  \{f\in\mathcal S(W_{\rm pos}(\mathbb A_f)):
        W_{N\lambda}f=\alpha(N\lambda)f
                    \text{ for all }\lambda\in\widehat\Lambda\}.
 \tag{H4}
\] Here a rational polarization is chosen only to write the Schrödinger space. All the definitions through (H4) are invariant under a unitary change of model.

\subsection{2. Exact orthonormal basis}\label{exact-orthonormal-basis}

\textbf{Proposition.} Choose one representative for every element of \[
                     N^{-1}\widehat\Lambda/\widehat\Lambda.
\] Then \[
                 \{W_x\phi_0:
                    x\in N^{-1}\widehat\Lambda/\widehat\Lambda\}
 \tag{H5}
\] is an orthonormal basis of \(\mathcal S_N\), after choosing the indicated representatives. Changing a representative changes its vector by a unit scalar. In particular \[
                          \dim\mathcal S_N=N^{2g}.
 \tag{H6}
\] Every finite Schwartz vector belongs to some \(\mathcal S_N\).

\subsubsection{Proof: membership and orthogonality}\label{proof-membership-and-orthogonality}

If \(x\in N^{-1}\widehat\Lambda\) and \(\lambda\in\widehat\Lambda\), selfduality gives \(\psi_f(B(N\lambda,x))=1\). Equations(H2)--(H3) therefore put \(W_x\phi_0\) in \(\mathcal S_N\).

For distinct classes \(x,y\), selfduality of \(\widehat\Lambda\) gives \(\lambda\in\widehat\Lambda\) with \(\psi_f(B(\lambda,x-y))\ne1\). The two vectors \(W_x\phi_0,W_y\phi_0\) are eigenvectors of the unitary operator \(W_\lambda\) with different eigenvalues, by (H2)--(H3). Their inner product is therefore zero. Each has norm one because Weyl operators are unitary. The quotient has order \(N^{2g}\), since the lattice has rank \(2g\). It remains to check that \(\mathcal S_N\) has exactly this dimension and that the union exhausts the Schwartz space.

\subsubsection{Proof: an explicit Schrödinger calculation}\label{proof-an-explicit-schruxf6dinger-calculation}

Choose an integral symplectic basis of the unimodular lattice. Its existence follows by the usual symplectic Gram--Schmidt argument over \(\mathbf Z\): choose a primitive vector, use unimodularity to find a partner of pairing one, split their hyperbolic plane, and iterate on its integral orthogonal complement. In the associated Schrödinger model the reference vacuum is \(\mathbf1_{\widehat{\mathbf Z}^{\,g}}\), of norm one, and its lattice semicharacter is the Weyl half-pairing character \(\alpha_{\rm std}\). The ratio \(\alpha/\alpha_{\rm std}\) is a continuous character of \(\widehat\Lambda\), since the two semicharacters have the same cocycle. The perfect finite-adelic symplectic pairing identifies the character group of this compact lattice with \(W(\mathbb A_f)/\widehat\Lambda\). Thus a finite Weyl translate \(W_t\mathbf1_{\widehat{\mathbf Z}^{\,g}}\) has precisely the character \(\alpha\).

For the trace-one sign character, the ratio has order at most two, so \(t\) may be taken in \(\tfrac12\widehat\Lambda/\widehat\Lambda\); the general argument only needs a finite denominator. The eigenspace of the full lattice for \(\alpha_{\rm std}\) is visibly one-dimensional: invariance under integral translations makes a function constant on every integral coset, while invariance under all integral modulations forces its support into \(\widehat{\mathbf Z}^{\,g}\). Consequently the actual \(\phi_0\) differs from \(W_t\mathbf1_{\widehat{\mathbf Z}^{\,g}}\) by a scalar of absolute value one. This scalar is retained and does not affect the assertions.

Conjugation by \(W_t\) reduces (H4) to the corresponding eigenspace for the standard vacuum. For the pure position and pure momentum elements of \(N\widehat\Lambda\), its conditions say exactly \[
 f(u+Nz)=f(u)\quad(z\in\widehat{\mathbf Z}^{\,g}),\qquad
       \operatorname{supp}f\subset
                  N^{-1}\widehat{\mathbf Z}^{\,g}.
 \tag{H7}
\] The mixed elements impose no additional condition: their Weyl half-pairing factor is exactly \(\alpha_{\rm std}(N\lambda)\). Thus this space is the space of functions on \(N^{-1}\widehat{\mathbf Z}^{\,g}/
                                  N\widehat{\mathbf Z}^{\,g}\). Its dimension is \(N^{2g}\), proving (H6), and hence proving completeness of the orthonormal family (H5).

Every compactly supported locally constant function is invariant under a sufficiently small open additive subgroup and has support in a sufficiently large compact open subgroup. Both can be bounded by \(N\widehat{\mathbf Z}^{\,g}\) and \(N^{-1}\widehat{\mathbf Z}^{\,g}\) for one integer \(N\). Applying the fixed inverse translate \(W_t^{-1}\) first proves the same exhaustion for the actual character \(\alpha\). This completes the proof. ∎

This proof includes dyadic places. In particular it does not set the Weyl half-pairing character equal to one at2.

\subsection{3. Exact coefficients and a finite-window bound}\label{exact-coefficients-and-a-finite-window-bound}

For \(f\in\mathcal S_N\), using the convention that the inner product is linear in its second variable, write \[
 f=\sum_x c_xW_x\phi_0,\qquad
           c_x=\langle W_x\phi_0,f\rangle,\qquad
                    \sum_x|c_x|^2=\|f\|_2^2.
 \tag{H8}
\] Thus the conversion to actual finite translates has exactly the Parseval norm, not a bound with an unspecified lattice-volume factor.

Let \(\mathcal T\) be a linear map from this finite-dimensional Schwartz window to a normed space; an evaluated theta lift is an example. The triangle inequality and Cauchy--Schwarz give \[
 \max_x\|\mathcal T(W_x\phi_0)\|
 \ \ge\
 \frac{\|\mathcal T(f)\|}{\sum_x|c_x|}
 \ \ge\
 \frac{\|\mathcal T(f)\|}{N^g\|f\|_2}
 \quad\text{if }\mathcal T(f)\ne0.
 \tag{H9}
\] The first denominator is nonzero in this case. This is a quantitative selection bound once \(N\) is known. It does not provide a bound for the \(N\) supplied by a global nonvanishing theorem.

In particular a nonzero linear functional, or a nonzero theta map on a specified input, cannot vanish on every vector in (H5). One actual translate detects it. The detecting vector itself has finite norm one; its coefficient in the full-orbit presentation is one, not a square root of the orbit cardinality. The factor \(N^g\) in (H9) measures a possible loss in a numerical lower bound for the detected output, not a norm of the selected input.

For a local factor at a fixed rational prime \(q\), the same proof works with \(N=q^e\) and only that local lattice, with dimension \(q^{2ge}\). If the input is the vacuum at every other place, the detecting translation can be chosen supported at \(q\). This last assertion concerns only an input already having that support property; it does not prove that a global Rallis test has it.

\subsection{4. A single actual translate in the iterated chain}\label{a-single-actual-translate-in-the-iterated-chain}

Apply this proposition to a later edge of ALP-LIFT-14 and its actual nonzero current vector \(F_j\). That theorem gives a selected real Schur vector and a finite Schwartz vector \(f\) with \[
                    \Theta_{f,\phi_\infty}(F_j)\ne0.
\] If the supplied Schwartz function is a finite sum of tensor products with different real vectors, at least one summand has a nonzero output. Fix that selected real vector. Choose \(N\) with \(f\in\mathcal S_N\), and expand (H8). Linearity gives \[
 \Theta_{f,\phi_\infty}(F_j)
   =\sum_x c_x\Theta_{W_x\phi_0,\phi_\infty}(F_j).
\] At least one summand is nonzero. Replace the finite input by that single \(W_x\phi_0\). Its infinite part is unchanged, and the finite part is an actual Weyl translate with norm exactly one.

The generated target representation remains spherical outside29: this is a statement about its local irreducible constituents, which are the unique local theta quotients of the spherical source constituents, not a claim that the chosen output vector itself is fixed by the maximal compact. The occurrence and Euler-factor argument of ALP-LIFT-14 depend on the constituent's sphericity, not on every current vector being its spherical vector. Its joint-kernel argument applies to every nonzero current vector. Thus it can be iterated without an additional compact averaging that would turn the single translate back into a linear combination.

We obtain the stronger form of the chain: \[
 F_{j+1}
   =\Theta_{W_{x_j}\phi_{0,j},\,\phi_{\infty,j}}(F_j)\ne0,
 \qquad
          \|W_{x_j}\phi_{0,j}\|_{2,f}=1
                \quad(1\le j<m),
 \tag{H10}
\] where every \(\phi_{\infty,j}\) belongs to the exact selected component \(\rho_j^\vee\boxtimes\rho_{j+1}\), and every \(x_j\) has zero infinite component and finite support. No coefficient combination height or finite Schwartz norm growth occurs in this presentation. The unknown denominator and support of \(x_j\), and the output norms and contents, remain unknown.

At infinity the creation-vector norm is precisely the actual norm of \(\phi_{\infty,j}\) in AM-R-011. It is not made one by the finite normalization in (H10). Likewise the source input and automorphic Haar masses are not changed or omitted.

\subsection{5. Exact source interface and remaining restrictions}\label{exact-source-interface-and-remaining-restrictions}

The recovered source definitions give the following literal facts.

\begin{itemize}
\tightlist
\item
  \nolinkurl{Development/AbelianVarietyTheory/Definitions/LatticeFrames/JointSmooth.lean,} lines20--23: comb is a finite sum of heisO translates, and IsFinComb requires only that each translation parameter have zero infinite part.
\item
  \nolinkurl{Development/AbelianVarietyTheory/Definitions/LatticeFrames/IsRatDatum.lean,} lines21--28: the alternate IsFinComb predicate has the same content, whereas IsRatDatum separately requires an actual half-Schwartz function whose rational restriction is the prescribed \(\Xi\), followed by the specified theta transform.
\item
  \nolinkurl{Development/AbelianVarietyTheory/Definitions/LatticeFrames/LevChiEqOne.lean,} lines103--118: PinColT.dat requires both the finite-combination condition and the exact IsRatDatum tie to \(w\,A.\mathrm{xi}\). They are different conditions, and proving the first does not prove the second.
\end{itemize}

Consequently the single translate in (H10) satisfies the mathematical finite-combination condition, with the delta coefficient at \(x_j\). It does not, by that fact, satisfy PinColT.dat or identify the source raw Xi. No Lean record or build is claimed.

There is also a dimension obstruction to silently identifying the basis (H5) with the accepted DP59 half-orbit. For the actual \(g\)-dimensional Hom abelian variety, a level \(M\) window of the form (H4) has dimension \(M^{2g}\). The prescribed DP59 family \([\alpha]^*H^0(\mathcal L^M)\), with \(\alpha\bar\alpha=M\), has rank \(M^g\), and its rescaled Heisenberg parameters lie in one selected isotropic half-subgroup. For \(M>1\) this is a strictly smaller family. The single translate selected by the full-window argument can have both position and momentum components and need not lie in that half-subgroup.

Similarly, the actual raw Xi of \nolinkurl{Formalism/LocalToGlobal/Columns/FiniteSumNonzeroLargeDepth.lean,} lines930--950 and1116--1120, is assembled from the specified PhiAt, omPhi, E1at and local-volume factors, together with phiT. It is not an arbitrary Schwartz function. AM-DO-029 addresses its exact finite Fourier expansion separately. The argument here has not replaced that local data, nor proved that one of its prescribed finite rows detects the previous actual image.

The precise remaining quantitative task is therefore narrower: bound or prescribe the finite translation data while preserving a nonzero action on the already produced image, and identify the resulting class-integrated matrix with the fixed raw Xi and row lattice. Its denominators, class masses and output content still enter the DP44 accounting. A norm-one full-orbit input alone does not compute those quantities.

\subsection{Inputs and check scope}\label{inputs-and-check-scope}

AM-DO-022 supplies the actual selfdual trace lattice, vacuum, additive character and unitary Heisenberg action. ALP-LIFT-14 supplies the actual nonzero selected chain before this refinement. DP59's half-orbit rank is accepted background, cited only to distinguish the two families. The finite-dimensional basis, Parseval calculation, selection bound and single-translate chain refinement above are proved directly here.

No claim is made that the later translations can all be supported at29. Such a sharpening requires a fixed-vacuum nonvanishing test at every other place, including the actual dyadic model; general local theta occurrence alone does not identify that test vector.

\clearpage

\section{AM-R-013: two literal analytic source predicates on the mixed definite host}\label{am-r-013-two-literal-analytic-source-predicates-on-the-mixed-definite-host}

Author: /root, GPT-6 Astra Pro, ALP\_MACHINE, 21 November 2026 UTC.

This result extracts the quantified analytic content of the first two members of \nolinkurl{AdmS} from the actual mixed-host construction. The first is a corollary of ALP-LIFT-14 with an algebraic, rather than merely closed, theta image. The second uses the actual terminal interlacing weights and a finite rational translate to obtain a nonzero compact restriction. Its period has an exact positive value. Neither assertion identifies that period with the prescribed arithmetic \nolinkurl{perCol} or \nolinkurl{rowAF} lattice. In particular this is not a construction of the old elliptic \nolinkurl{D7}, a proof of all six \nolinkurl{AdmS} assertions, or a solution of \nolinkurl{MainX}'s terminal mass equation.

\subsection{1. The source predicates being matched}\label{the-source-predicates-being-matched}

The exact declarations are in \nolinkurl{Formalism/SharpRestriction/CuspSpanNormBounds/CayleySignature.lean}. \nolinkurl{RallisSiegelWeilFormula}, lines154--161, requires actual irreducible cuspidal subspaces \(\pi_j\) of the selected definite types, common edge characters with infinity exponents \(-e_j\), and, for every \(\phi\in\pi_{j+1}\), a finite identity \[
 \phi=\sum_{i=1}^{r_\phi}
       \Theta^{\mathrm{bil}}_{\Phi_i}(\phi_i),
 \qquad \phi_i\in\pi_j,
 \tag{R13.1}
\] with each \(\Phi_i\) factorizable and compact-finite. The predicate does not prescribe the arithmetic row lattice or a norm bound for these kernels.

\nolinkurl{LiftParameterConditions}, lines163--166, requires that every actual irreducible automorphic subspace of the selected last definite type have a nonzero terminal period against some cusp form of the prescribed last \(Y\)-type. The latter type is imposed on the conjugate form by \nolinkurl{IsDv}, lines123--125. For the ordinary compact restriction period the literal combination of conjugations is \[
 B_{\mathrm{top}}(\overline\phi,F)
       =\int_{[H]}\phi(\iota h)F(h)\,dh,
 \quad
 B_{\mathrm{top}}(f,F)
       :=\int_{[H]}\overline{f(\iota h)}F(h)\,dh.
 \tag{R13.2}
\] Here \(G=U(I_n)\), \(n=N_m\), and \(H=U(X')\) is the unitary group of the actual positive codimension-one subspace \(X'=z^\perp\), identified with the terminal \(Y_m\). The embedding \(\iota\) acts as the identity on the orthogonal line. At each real place it is conjugate to \(h\mapsto\operatorname{diag}(h,1)\); no global isometry \(X'\simeq I_{n-1}\) is assumed. This is the actual Hermitian restriction period, with its Haar measure retained. We do not use the auxiliary identity-action Hilbert-space model \nolinkurl{thetaG} to obtain either conclusion.

\subsection{2. An actual irreducible chain and finite factorizable spanning}\label{an-actual-irreducible-chain-and-finite-factorizable-spanning}

Use the common host \(E'=\mathbf Q(\zeta_{28})\), \(F'=E'^+\), the polarized \(J^2\), the same trace coefficient, and all the characters, groups and compact types of ALP-LIFT-14. For every fixed choice of the source rank and HC data there exist actual irreducible cuspidal subspaces \[
 \pi_j\subset\mathcal A(G_j),\qquad
        (\pi_j)_\infty\simeq\rho_j,
             \qquad 0\le j\le m,
 \tag{R13.3}
\] such that (R13.1) holds for every edge. All these are specified automorphic embeddings, so no multiplicity-one assertion for the global spectrum is required.

\textbf{Proof.} The first input of AM-R-012 has fixed finite level and compact real type. The automorphic space at that level and type is finite-dimensional: it is a finite direct sum, over the definite double-coset set, of invariant subspaces of the finite-dimensional coefficient representation. Its spectral decomposition therefore involves finitely many actual irreducible automorphic embeddings. Decompose this input accordingly. Its first theta image is nonzero, so at least one summand, belonging to an actual \(\pi_0\), has a nonzero image with the same selected first real component.

Let \(\mathcal T(\pi_0)\) be the algebraic span of these theta images, allowing all finite Schwartz functions and the selected real component. It is stable under the target finite and real actions by the exact equivariance of the Weil representation. All its automorphic constituents have the unique local small-theta target at every finite place, and the unique selected target compact type at infinity. Consequently they have one global irreducible isomorphism type. The automorphic multiplicity of that type is finite. The smooth algebraic isotypic part is a finite direct sum of irreducible modules, hence is semisimple. Its submodule \(\mathcal T(\pi_0)\) is semisimple as well and contains an actual irreducible submodule \(\pi_1\).

This argument concerns algebraic spans, not just a dense image in a Hilbert-space closure. By definition, every vector of \(\pi_1\) is a finite sum of theta images of vectors of \(\pi_0\). Every finite Schwartz function is a finite sum of restricted pure tensors. Every vector in the selected real component is a finite sum of tensors over the six archimedean places, and is compact-finite. Expanding these finite sums and using bilinearity makes each kernel factorizable without changing the image. This proves (R13.1) at the first edge.

ALP-LIFT-14, sections3--7, proves nonvanishing for every actual irreducible embedding with the resulting unramified parameters, including each possible multiplicity copy. Apply it to \(\pi_1\), then repeat the preceding algebraic-image argument. Induction gives (R13.3) and (R13.1) at every edge. The input and output are cuspidal because each positive definite rational group has no proper rational parabolic. The characters are exactly the globally constructed equal characters of LIFT14, so the required \nolinkurl{IsMu} condition holds. ∎

The proof uses only even \(k,m,d=t+1\), \(d\ge k+m+5\), the HC interlacing/parity, and \(e_j-b_j,b_{j+1}-e_j\ge(d+2)/2\). Thus it applies also when the source's auxiliary \nolinkurl{TPar.t} replaces the fixed \nolinkurl{Ranks.t}: \nolinkurl{TPar.t_ge}, \nolinkurl{Hyp3.Interlace}, \nolinkurl{TParity} and \nolinkurl{Centred} imply exactly these inequalities. The other \nolinkurl{Ranks} inequalities, such as \(3m<t\), are not used by the analytic nonvanishing proof. This proves the first predicate's mathematical quantifiers for that actual interpretation, whenever the surrounding geometric data have been supplied. It is not an assertion that a Lean \nolinkurl{ThetaX D} record has been instantiated.

\subsection{3. Exact terminal branching from the source HC lists}\label{exact-terminal-branching-from-the-source-hc-lists}

Fix one of the six complex places. Since \(m\) is even, the source terminal Harish--Chandra lists have no residual \(\eta\)-shift. From \nolinkurl{DataTheta.lean}, \nolinkurl{TPar.Xpar} and \nolinkurl{TPar.Dpar}, lines538--548, their decreasing rearrangements are \[
 \begin{split}
 x={}&(u_0,\ldots,u_{k-1},\,
          B_{t+1}(-e_0),\ldots,B_{t+1}(-e_{m-1})),\\
 y={}&(v_0,\ldots,v_{k-2},-b_0,\,
       B_t(-e_0),-b_1,B_t(-e_1),\ldots,
                    -b_{m-1},B_t(-e_{m-1})),\\
 B_l(c)={}&(c+(l-1)/2,\,c+(l-3)/2,\ldots,c-(l-1)/2).
 \end{split}
 \tag{R13.4}
\] In this display the terminal line \(y\) consists of its initial \(k-1\) entries followed by the \(m\) pairs \((-b_j,B_t(-e_j))\). It has length \(n-1\), whereas \(x\) has length \(n=k+m(t+1)\). All source signature labels are positive at \(j=m\).

The actual HC assumptions imply \[
 x_i>y_i>x_{i+1}\qquad(0\le i<n-1).
 \tag{R13.5}
\] Indeed, the initial entries alternate by the original \(u,v\) interlacing. Next \(u_{k-1}>-b_0\), while \(e_0-b_0\ge(t+3)/2\) places \(-b_0\) above the top of the first \(x\)-block. Within each block the \(t\) entries of \(B_t(-e_j)\) are exactly the midpoints between the \(t+1\) adjacent entries of \(B_{t+1}(-e_j)\). Between two blocks, \(b_{j+1}-e_j\ge(t+3)/2\) and \(e_{j+1}-b_{j+1}\ge(t+3)/2\) place \(-b_{j+1}\) strictly between their endpoint entries. These observations also justify the stated decreasing rearrangements and the absence of duplicates in the source finite sets.

Put \[
 \lambda_i=x_i-(n-1)/2+i,\qquad
 \beta_i=y_i-(n-2)/2+i.
 \tag{R13.6}
\] These are the dominant integral highest weights corresponding to the two HC lists. The \(x_i\) are half-integers and the \(y_i\) integers; hence each strict gap in (R13.5) is at least \(1/2\). Substituting (R13.6) gives \[
        \lambda_i\ge\beta_i\ge\lambda_{i+1}.
 \tag{R13.7}
\] The compact branching theorem therefore gives \[
 \dim\operatorname{Hom}_{U(n-1)}
      (V_\beta,\operatorname{Res}^{U(n)}_{U(n-1)}V_\lambda)=1.
 \tag{R13.8}
\] This is the ordinary branching theorem, cited rather than re-proved: Goodman and Wallach (2009), Theorem8.1.1, p.364. Restriction from the complex algebraic groups to their compact forms preserves these finite-dimensional representation multiplicities. Taking the external tensor product over the six places gives the actual terminal compact type \(\tau\) inside \(\rho_m|_{H_\infty}\), with multiplicity one. Its HC parameter is exactly \nolinkurl{Dpar m}, not a freely selected branching constituent.

\subsection{4. A compact restriction detector for every actual terminal embedding}\label{a-compact-restriction-detector-for-every-actual-terminal-embedding}

The next elementary argument supplies the global existence needed after the local branching statement. Local occurrence alone would not justify restriction of a fixed automorphic vector at a fixed finite point.

Let \(\pi\subset\mathcal A(G)\) be any nonzero irreducible automorphic embedding with \(\pi_\infty=\rho_m\). There are \(\phi\in\pi\) and a smooth cusp form \(F\) on \(H\) such that \[
 \overline F\text{ has real type }\tau,
 \qquad
 B_{\mathrm{top}}(\overline\phi,F)
       =\|\operatorname{Res}_H\phi\|_{L^2([H])}^{2}>0.
 \tag{R13.9}
\] Here \(\phi\) is chosen so that its restriction already has type \(\tau\); no projection denominator is suppressed.

\textbf{Proof.} First, \(G(F')\) is dense in \(G_\infty\). For the identity Hermitian form, its Lie algebra over \(F'\) is the vector space of matrices satisfying \(X^*=-X\). Archimedean weak approximation for this vector space is elementary: choose an \(F'\)-basis; the diagonal image of \(F'\) is dense in the product of its real completions, since a rational basis gives an invertible real embedding matrix and \(\mathbf Q^s\) is dense in \(\mathbf R^s\). Approximate each coordinate independently. The rational Cayley map \[
             X\longmapsto(1+X)(1-X)^{-1}
 \tag{R13.10}
\] is defined on every real skew-Hermitian matrix and maps onto the open dense subset of the compact unitary group having no eigenvalue \(-1\). Its inverse there is \((g-1)(g+1)^{-1}\). The same rational identities hold over \(F'\). Approximation through this chart proves density simultaneously at all real places. No bound on denominators of the rational approximant is asserted.

Choose an actual nonzero finite vector in \(\pi_f\), and let \(v\mapsto\phi_v\) be the resulting \(G_\infty\)-equivariant map from \(V_{\rho_m}\) to the automorphic embedding. Some finite point has nonzero evaluation on this map. A right finite translation allows us to take that point to be the identity. There is then a nonzero linear functional \(L\) on \(V_{\rho_m}\) with \[
                  \phi_v(k,1)=L(\rho_m(k)v).
 \tag{R13.11}
\] Let \(V_\tau\subset V_{\rho_m}\) be the nonzero branching subspace supplied by (R13.8). The span of \(\rho_m(G_\infty)V_\tau\) is a nonzero invariant subspace of the irreducible \(\rho_m\), so it is the whole space. Consequently there exist \(k\in G_\infty\) and \(v\in V_\tau\) with \(L(\rho_m(k)v)\ne0\). This inequality is open in \(k\). By the Cayley density just proved, choose \(q\in G(F')\) with \(L(\rho_m(q_\infty)v)\ne0\).

Right-translate \(\phi_v\) by \(q_f^{-1}\), and call the result \(\phi\). For \(h\in H_\infty\), left rational automorphy gives the exact identity \[
 \begin{split}
 \phi(\iota h,1)
 &=\phi_v(\iota h,q_f^{-1})
  =\phi_v(q_\infty\iota h,1)\\
 &=L\bigl(\rho_m(q_\infty)\rho_m(\iota h)v\bigr).
 \end{split}
 \tag{R13.12}
\] Thus the restriction is nonzero, already at the identity. Finite translation commutes with the real group action. Since \(v\in V_\tau\), the full restricted automorphic function \(\operatorname{Res}_H\phi\), at all its finite points, has real type \(\tau\).

Set \(F=\overline{\operatorname{Res}_H\phi}\). Restriction is well-defined on \([H]\) by rational automorphy. It is smooth and of some finite compact level. Both groups are anisotropic over \(F'\), so their quotients are compact and all their automorphic forms are cuspidal. Hence this \(F\) is permitted in the analytic terminal predicate, and its conjugate has precisely the required real type. Substitution into (R13.2) gives (R13.9). The integral is positive because a nonzero continuous function on the compact quotient is nonzero on an open set of positive Haar measure. ∎

\subsection{5. Retained class masses and the remaining arithmetic equation}\label{retained-class-masses-and-the-remaining-arithmetic-equation}

The value in (R13.9) is exact for the chosen measure. For any compact finite level \(K_H\) fixing the restricted function, choose representatives \(h_c\) for \(H(F')\backslash H(\mathbb A_f)/K_H\), and write \(\Gamma_c=H(F')\cap h_cK_Hh_c^{-1}\). Normalize real Haar measure to volume one and retain the finite volume \(\operatorname{vol}(K_H)\). Then the same value is \[
 \sum_c\frac{\operatorname{vol}(K_H)}{|\Gamma_c|}
       \int_{H_\infty}|\phi(\iota h,\iota h_c)|^2\,dh.
 \tag{R13.13}
\] This follows by the ordinary quotient-measure disintegration over the finite classes; it is the positive mass convention already used in AM-DO-026 and AM-R-012. It neither normalizes the total automorphic volume to one nor drops the stabilizers. Scaling \(F\) by its reciprocal would force a unit period over \(\mathbf C\), but AM-R-002 shows why this would still have an arithmetic content and norm cost.

Equations(R13.1) and(R13.9) supply the first two analytic predicates with their actual quantifiers and compact types. The other four predicates involve the noncompact \(Y_j\) descent, specified Fourier--Jacobi periods, the lower-tower vanishing, and real type stability. Those are not supplied by this proof. In addition, \nolinkurl{TieRowIsom}, \nolinkurl{TieIo}, and \nolinkurl{TieBm} in \nolinkurl{DataTheta.lean}, lines739--752, still impose separate arithmetic comparisons. In particular the explicit positive restriction period (R13.13) has not been proved equal to the source's prescribed terminal coefficient matrix. The finite translation \(q_f^{-1}\), the resulting level, and the integral coefficient content remain unquantified. These are precisely the limitations relevant to the requested DP44 lift-count comparison.

\subsection{6. Geometric interpretation and sources}\label{geometric-interpretation-and-sources}

The compact branching is a restriction of algebraic representations of the reductive unitary groups, which become \(\mathrm{GL}_n\) and \(\mathrm{GL}_{n-1}\) after splitting the CM extension. The integral weights define the usual Schur functors and determinant twists; their characters live in the representation \(\lambda\)-ring. No rational projector in this proof is asserted to preserve a chosen integral lattice. This distinction is the same one recorded in the SGA6/SGA3 entries of \nolinkurl{GROTHENDIECK_DICTIONARY.md}. The theta kernels still arise from the principal polarization's biextension and the CM Picard data already identified there (SGA7I, Exp.VII, Proposition2.9.3, pp.172--174; Exp.VIII, Theorem7.1, pp.299--300; FGA232, Theorem3.1, pp.148--149). The restriction norm is an archimedean realization. A metrized integral determinant-line comparison with it is additional work, not a formal consequence of functoriality of motives or of weights for sheaves over finite fields.

Primary new literature input: Goodman, R., \& Wallach, N. R. (2009). \emph{Symmetry, representations, and invariants}. Springer. Theorem8.1.1, p.364; DOI10.1007/978-0-387-79852-3. The theorem was directly inspected in the primary book PDF at \url{https://webhomes.maths.ed.ac.uk/~v1ranick/papers/goodwallx.pdf} (PDF page383). The compact branching theorem is cited as known.

Internal inputs, with their separate check histories: ALP-LIFT-14, sections3--8; AM-DO-025--026; AM-R-002, AM-R-012; AM-K022. Exact source declarations used above are \nolinkurl{CayleySignature.lean}, lines117--166, and \nolinkurl{DataTheta.lean}, \nolinkurl{TPar}, \nolinkurl{block}, \nolinkurl{Dpar}, \nolinkurl{Xpar}, \nolinkurl{ThetaX.IsStable}, \nolinkurl{typeX}, and \nolinkurl{TieBm}. No Lean execution or old-host instantiation is claimed.

\clearpage

\section{AM-DO-029: the actual raw Xi as a finite Heisenberg vector, its exact cost, and the 29-only mismatch}\label{am-do-029-the-actual-raw-xi-as-a-finite-heisenberg-vector-its-exact-cost-and-the-29-only-mismatch}

Author: GPT-6 Astra Pro, ALP\_MACHINE \nolinkurl{/root/definite_oscillator}; inherited reasoning effort is not separately exposed; 21 November 2026 UTC.

This finding concerns the \textbf{prescribed} \nolinkurl{xiAt} and the finite clauses of \nolinkurl{IsRatDatum}. It uses the carried THETA3-A12/A13/A15 finite-model results at their checked status, and the actual mixed host and normalization of AM-DO-022/026. It does not replace a prescribed array by an arbitrary Schwartz function. The new conclusions are its explicit finite-window realization on the mixed host, exact norm and denominator accounting, and a precise obstruction to identifying it with the previously constructed input supported only at 29.

\subsection{1. Literal source interface and local windows}\label{literal-source-interface-and-local-windows}

All source paths below are relative to \nolinkurl{work/research/alpoge_source}, pinned to commit \nolinkurl{7cbf7a0c32fca1badf05931879a37c3f8a2b03b4}. Reading these recovered definitions is not a Lean compilation or a transitive dependency audit.

\begin{longtable}[]{@{}
  >{\raggedright\arraybackslash}p{(\columnwidth - 4\tabcolsep) * \real{0.3333}}
  >{\raggedright\arraybackslash}p{(\columnwidth - 4\tabcolsep) * \real{0.3333}}
  >{\raggedright\arraybackslash}p{(\columnwidth - 4\tabcolsep) * \real{0.3333}}@{}}
\toprule\noalign{}
\begin{minipage}[b]{\linewidth}\raggedright
Source
\end{minipage} & \begin{minipage}[b]{\linewidth}\raggedright
Locator
\end{minipage} & \begin{minipage}[b]{\linewidth}\raggedright
Content used
\end{minipage} \\
\midrule\noalign{}
\endhead
\bottomrule\noalign{}
\endlastfoot
\nolinkurl{Development/AbelianVarietyTheory/Definitions/LatticeFrames/IsRatDatum.lean} & 19--28 & \nolinkurl{IsFinComb}; the three clauses of \nolinkurl{IsRatDatum} \\
\nolinkurl{Development/AbelianVarietyTheory/Definitions/LatticeFrames/LevChiEqOne.lean} & 68--71 & actual \nolinkurl{TwDat.xi}, obtained from \nolinkurl{xiAt} and the specified frames \\
\nolinkurl{Development/AbelianVarietyTheory/Definitions/AutomorphicMembership/XiAtPenultimateEqXiCpl1.lean} & 146--163 & the three branches \nolinkurl{XiEntry}, \nolinkurl{XiCpl1}, \nolinkurl{xiLast}; all contain \nolinkurl{phiT} \\
\nolinkurl{Formalism/LocalToGlobal/Columns/FiniteSumNonzeroLargeDepth.lean} & 646--648; 927--948; 981--986 & \nolinkurl{phiT}; the local \nolinkurl{PhiAt}, \nolinkurl{omPhi}, \nolinkurl{E1at}, \nolinkurl{XiAt}; \nolinkurl{joinCol}, \nolinkurl{PsiAt} \\
same & 1047--1065 & \nolinkurl{ChIdx.Ok.psiSupp}, \nolinkurl{psiInv}, \nolinkurl{xiSupp}, \nolinkurl{xiInv}, \nolinkurl{psiSuppX}, \nolinkurl{psiInvX} \\
same & 1108--1138 & \nolinkurl{Sall}, \nolinkurl{cAof}, \nolinkurl{cBof}, \nolinkurl{XiEntry}, \nolinkurl{XiCpl1} \\
\nolinkurl{Formalism/LocalToGlobal/IndefiniteReal/LocalNewformFunction.lean} & 28--37; 239--264; 270--279 & \nolinkurl{symplO}, \nolinkurl{heisO}; \nolinkurl{IsSBHalf}, \nolinkurl{IsPlaceModel}, \nolinkurl{IsThetaAt}; the actual \nolinkurl{localNewformFunction} \\
\end{longtable}

Fix one nonterminal edge \(0\leq i\leq m-2\), ranks

\[
 a=N(i),\qquad b=N(i+1),\qquad n=ab,
 \tag{DO29.1}
\]

and one allowed source and target decoration, local index, and frame. Write \(\Xi\) for the resulting scalar \nolinkurl{TwDat.xi}, before or after a specified nonzero chart scalar. The scalar is retained in every formula. Let \(\mathcal S=\mathrm{Sall}\) be its actual finite half-place set. The source's selected-place and prime hypotheses are retained; in particular this is a set of relative split places with no conjugate pair selected twice.

At a \(T\)-place \(P\), the factor is supported on

\[
 \operatorname{Hom}_{\mathcal O_{E,P}}(N_P,N'_P)
 \tag{DO29.2}
\]

and is constant modulo \(P\operatorname{Hom}(N_P,N'_P)\). In residue frames it is exactly

\[
 I_{v,u}(X)=1_{\{vX=u\}},\qquad u\ne0,\quad v\ne0.
 \tag{DO29.3}
\]

This includes matrix integrality. The assertion follows directly from the localized lattice condition in \nolinkurl{phiT} and the residue-linearity in \nolinkurl{Deco.Ok}; the exact incidence description and its rank-one Fourier span are carried THETA3-A13.

At a place \(P\in\mathrm{rowSG}\), let \(f_P,f'_P\) be the specified source and target frames and put

\[
 L_P^{\rm fr}=f'_P M_{b\times a}(\mathcal O_{E,P})f_P^{-1}.
 \tag{DO29.4}
\]

For an \nolinkurl{XiEntry} branch, the literal \nolinkurl{xiSupp} and \nolinkurl{xiInv} fields give support in \(q_P^{-A}L_P^{\rm fr}\) and period lattice \(q_P^B L_P^{\rm fr}\). For an \nolinkurl{XiCpl1} branch, the literal \nolinkurl{psiSuppX} and \nolinkurl{psiInvX} fields give the same bounds in the matrix variable, with the actual fixed representative of the extra column retained. Here \(q_P=\mathrm{qOf}(P)\), and the exponents mean scalar multiplication exactly as in the source. No assumption that an arbitrary function \nolinkurl{LocCst.psi} is continuous outside its prescribed depth is needed. These are already hypotheses of the actual \nolinkurl{ChIdx.Ok} record.

The first edge of every allowed chain is one of these branches, since the source has \(m\geq2\). The present assertion does \textbf{not} silently extend these nonterminal fields to \nolinkurl{xiLast}. That branch has the additional variable additive character in \nolinkurl{ftY}; its stated depth and terminal representative conditions must be used separately.

\subsection{2. The explicit extension and finite transform}\label{the-explicit-extension-and-finite-transform}

Use the actual mixed finite symplectic space from AM-DO-022:

\[
 E'=\mathbf Q(\zeta_{28}),\quad
 X=eW_f,\quad Y=\bar eW_f,\quad
 L=e\widehat\Lambda,
 \tag{DO29.5}
\]

where \(\Lambda=M_{b\times a}(\mathcal O_{E'})\), and \(e\) is the selected finite half idempotent for \(\mathcal S\). The mixed scalar trace coefficient is \(\xi'=\xi/2\); the product lattice and the dyadic basis index remain exactly those of AM-DO-022. Write \(B\) for the effective alternating pairing with the actual fixed additive character \(\psi_f\). Thus \(W_xW_y=\psi_f(B(x,y)/2)W_{x+y}\). This notation incorporates the source's standard-character sign, rather than changing it in a Fourier formula.

Let \(L^+\supseteq L\supseteq L^-\) be compact open lattices in \(X\) such that the local bounds of §1 imply

\[
 \operatorname{supp}\Phi\subseteq L^+,
 \qquad \Phi(t+v)=\Phi(t)\quad(v\in L^-).
 \tag{DO29.6}
\]

They can be chosen without an unspecified enlargement. At each \(T\)-place take the sum of the standard lattice and (DO29.2) for \(L_P^+\), and intersect the standard lattice with \(P\) times (DO29.2) for \(L_P^-\). At each other selected place take

\[
 L_P^+=L_P+q_P^{-A}L_P^{\rm fr},\qquad
 L_P^-=L_P\cap q_P^B L_P^{\rm fr}.
 \tag{DO29.7}
\]

Products give \(L^+,L^-\). Smith normal forms of the displayed local frame matrices give their exact elementary divisors and indices. Thus no frame index is suppressed. Define

\[
 R=[L^+:L],\qquad N=[L:L^-],\qquad
 C=L^+/L,\quad U=L/L^-.
 \tag{DO29.8}
\]

Choose representatives \(a\in L^+\) of \(C\), and representatives

\[
 y\in (L^-)^\perp/L^\perp\subset Y/L^\perp
 \tag{DO29.9}
\]

of the characters of \(U\), where annihilators use \(\psi_f\circ B\). This is a perfect finite pairing; in local coordinates it is the ordinary duality between a lattice quotient and its inverse-different dual. Its order is \(N\).

There is a unique locally constant function \(\Phi\) on \(X\), extended by zero off \(X\), satisfying

\[
 \Phi(e\gamma)=\Xi(\gamma)\quad
     (\gamma\in M_{b\times a}(E')).
 \tag{DO29.10}
\]

\textbf{Proof of existence and uniqueness.} At the finitely many selected places, use the factors of §1. The support and period statements make their product constant on every coset of \(L^-\) and zero off \(L^+\). Every required finite congruence class has a rational matrix representative by additive weak approximation, or equivalently the Chinese remainder theorem after clearing denominators. Its value is independent of the representative by the period statements. This defines the function and proves (DO29.10). Conversely rational matrices meet every open coset of \(L^-\), so no other locally constant extension with these periods is possible. Choose a positive integer \(k\) with \(kL\subseteq L^-\) and \(kL^+\subseteq L\); the local elementary divisors give one explicitly. Then (DO29.6) is exactly the support and period part of \nolinkurl{IsSBHalf}, and support on \(X=eW_f\) is its remaining clause. \(\square\)

The exact Fourier coefficients are

\[
 c_{a,y}=\frac1N\sum_{u\in U}
                 \Phi(a+u)\psi_f(-B(u,y)).
 \tag{DO29.11}
\]

Finite character orthogonality gives the literal identity

\[
 \Phi(t)=\sum_{a\in C,y\in\widehat U}
 c_{a,y}\,\psi_f(B(t-a,y))1_L(t-a).
 \tag{DO29.12}
\]

Let \(T\) be any normalized \nolinkurl{IsPlaceModel}, so that

\[
 T1_L=f_0,\quad T\Phi(\cdot+x)=W_xT\Phi,
 \quad T(\psi_f(B(\cdot,y))\Phi)=W_yT\Phi.
 \tag{DO29.13}
\]

Then the required finite vector is \textbf{exactly}

\[
 \boxed{f_\Xi=\sum_{a,y}c_{a,y}W_{-a}W_yf_0
   =\sum_{a,y}c_{a,y}\psi_f(-B(a,y)/2)W_{y-a}f_0.}
 \tag{DO29.14}
\]

Indeed apply modulation and then translation to \(1_L\) in (DO29.12), and use the displayed Weyl law. Every normalized model gives the same vector. This is precisely the carried A12 construction and the literal theorem \nolinkurl{thetaAt_eq_sum_heis_translates}, now with the coefficients of the \textbf{actual} array rather than freely chosen ones. The theorem \nolinkurl{exists_isThetaAt_of_vac}, lines124--132 of \nolinkurl{Development/AbelianVarietyTheory/Proofs/DivisorProofs/ExistsIsThetaAtOfVac.lean}, supplies existence in the original source wrapper; the actual mixed Schrödinger model supplies existence on the mixed host. Section5 states their precise finite comparison.

Equations (DO29.6), (DO29.10), and (DO29.14) verify all three mathematical clauses of \nolinkurl{IsRatDatum}. They also specify the exact finite list and coefficients for \nolinkurl{IsFinComb}; every parameter has zero infinite component. They do not assert that an unchanged elliptic \nolinkurl{D7} wrapper has been instantiated for the new CM type.

\subsection{3. Exact norm, denominators, and local newform examples}\label{exact-norm-denominators-and-local-newform-examples}

Normalize the finite Haar measure by \(\operatorname{vol}(L)=1\) and the finite vacuum by \(\|f_0\|=1\). Then

\[
 \boxed{\|f_\Xi\|^2
   =\sum_{a,y}|c_{a,y}|^2
   =\frac1N\sum_{a\in C,u\in U}|\Phi(a+u)|^2.}
 \tag{DO29.15}
\]

\textbf{Proof.} The functions \(\psi_f(B(t-a,y))1_L(t-a)\) have disjoint supports for different \(a\), and character orthogonality on \(L/L^-\) gives orthogonality for different \(y\) with the same \(a\). Each has norm1. The same orthogonality holds for their images \(W_{-a}W_yf_0\): lattice modulations distinguish different position cosets, lattice translations distinguish different characters, and these operators are unitary. The vacuum norm fixes the norm of every image at1. Consequently (DO29.14) is an orthonormal expansion. Finite Parseval gives the second equality. This proves the normalization even though the bare predicate \nolinkurl{IsPlaceModel} does not separately include an isometry axiom. \(\square\)

If \(H_\Xi=\max_{a,u}|\Phi(a+u)|\), then

\[
 \|f_\Xi\|\leq\sqrt R\,H_\Xi,
 \qquad \sum_{a,y}|c_{a,y}|
       \leq\sqrt{RN}\,\|f_\Xi\|
       \leq R\sqrt N\,H_\Xi.
 \tag{DO29.16}
\]

These follow respectively from the \(RN\) summands in (DO29.15) and Cauchy--Schwarz. They are bounds, whereas (DO29.15) is exact.

Let \(K_\Xi\) contain all actual table values and all character values in (DO29.11)--(DO29.14). The latter are roots of unity. If \(\delta\Phi(a+u)\) is integral over \(\mathcal O_{K_\Xi}\), then

\[
 \delta N c_{a,y}\in\mathcal O_{K_\Xi}.
 \tag{DO29.17}
\]

Equivalently, at a finite valuation \(v\),

\[
 v(c_{a,y})\geq
       \min_{a,u}v(\Phi(a+u))-v(N).
 \tag{DO29.18}
\]

The central phase in (DO29.14) is a unit and adds no valuation. The values of the actual array still contain its \nolinkurl{volN} factors, finite sums, specified characters and chart scalar; their valuation is \textbf{not} replaced by0. If the finite level and lattice changes are prime to11, the actual theta-group isomorphisms also preserve the integral11-adic theta lattice. Thus (DO29.17) is then a bound for the geometric coefficient realization after the common period has been removed. This assertion makes no claim at an auxiliary place of residue characteristic11.

For standard frames, at a degree-one selected \(T\)-prime the local indices are \(R_T=1,N_T=q^n\). At a selected degree-one row prime they are \(R_P=q^{An},N_P=q^{Bn}\). Thus, when all these frames are integral units, one may use

\[
 R=\prod_{P\in\mathrm{rowSG}}q_P^{An},\qquad
 N=T(u)^{[F':\mathbf Q]n}
           \prod_{P\in\mathrm{rowSG}}q_P^{Bn}.
 \tag{DO29.19}
\]

Repeated rational primes in the product are intentional: each selected half place contributes its matrix quotient. For other frames or residue degrees, (DO29.7)--(DO29.8), rather than this simplified formula, give the exact indices.

\subsubsection{3.1 The actual incidence factor}\label{the-actual-incidence-factor}

At a \(T\)-place with residue field \(\kappa\) of order \(q\), choose the residue additive character induced by the actual pairing. For nonzero \(u,v\) as in (DO29.3),

\[
 1_{M_{b\times a}(\mathcal O)}(t)I_{v,u}(\bar t)
   =q^{-a}\sum_{\eta\in\kappa^a}
       \psi_\kappa(-\eta\cdot u)
       \psi_\kappa(\eta\cdot v\bar t)1_L(t).
 \tag{DO29.20}
\]

This follows by summing a character of the \(a\)-dimensional additive group \(\kappa^a\): the sum is \(q^a\) when \(v\bar t=u\), and0 otherwise. Let \(y_{v,\eta}\in P^{-1}L^\perp/L^\perp\) be the exact dual vector for the character \(t\mapsto\psi_\kappa(\eta\cdot v\bar t)\), including the coefficient of the actual trace pairing. Therefore

\[
 f_{v,u}=q^{-a}\sum_{\eta\in\kappa^a}
                 \psi_\kappa(-\eta\cdot u)W_{y_{v,\eta}}f_0,
 \quad \|f_{v,u}\|^2=q^{-a},
 \quad \sum|c_\eta|=1.
 \tag{DO29.21}
\]

The map \(\eta\mapsto y_{v,\eta}\) is injective because \(v\ne0\). Thus the norm also follows directly from \(q^a\) orthogonal terms of coefficient absolute value \(q^{-a}\). Its denominator is exactly \(q^a\) in this expansion. The frequencies are the matrices \(v^t\eta\), of rank at most1; this is the source incidence restriction, not a full matrix-delta projector.

\subsubsection{\texorpdfstring{3.2 The literal \nolinkurl{localNewformFunction}}{3.2 The literal }}\label{the-literal-localnewformfunction}

Let the local field have residue order \(q\), uniformizer \(\varpi\), and let \(P\subset\mathrm{GL}_r(\mathcal O)\) be a compact open subgroup. Assume the specified finite characters are continuous and unitary, as in the actual compact type. The source function is

\[
 F\binom{A}{B}=
 \begin{cases}
 \lambda(A)\mu(\det A)^{-1},
       &A\in P,\ B\in\varpi M_{t\times r}(\mathcal O),\\
 0,&\text{otherwise}.
 \end{cases}
 \tag{DO29.22}
\]

This is the literal definition at the locator in §1; it includes both its unit block and its lower-block divisibility. Choose \(s\geq1\) such that \(K_r(\varpi^s)\subset P\) and the product character \(\lambda\mu^{-1}\circ\det\) is trivial there. Then (DO29.11) with \(R=1,N=q^{s(r+t)r}\) is an explicit finite Heisenberg expression for this function. More sharply its norm is

\[
 \boxed{\|F\|^2
  =q^{-tr}\,
     \frac{\prod_{j=1}^{r}(1-q^{-j})}
          {[\mathrm{GL}_r(\mathcal O):P]}.}
 \tag{DO29.23}
\]

To prove the period assertion, changing \(A\) by \(\varpi^sZ\) multiplies it on the right by an element of \(K_r(\varpi^s)\) whenever \(A\in\mathrm{GL}_r(\mathcal O)\). It therefore preserves both support and value; changes in \(B\) by \(\varpi^s\) preserve its divisibility. For the norm, its absolute value on its support is1. Additive matrix measure gives \(M_r(\mathcal O)\) volume1 and \(\mathrm{GL}_r(\mathcal O)\) volume \(\prod_j(1-q^{-j})\). Multiplication by any integral unit matrix preserves this additive measure, so \(P\) has exactly the displayed index fraction of this volume. The lower block has volume \(q^{-tr}\). Multiplication gives (DO29.23). This proves every factor and does not identify a group probability measure with additive matrix measure.

The actual \nolinkurl{PhiG}, \nolinkurl{PhiGS}, \nolinkurl{PhiQ}, \nolinkurl{E1at} and their specified frame operations are retained inside the table (DO29.11). In particular a finite auxiliary sum is not replaced by just the local function (DO29.22). Equation (DO29.15) computes its exact norm after that sum, including cancellation.

\subsection{4. A nonzero actual Xi cannot be a 29-only input}\label{a-nonzero-actual-xi-cannot-be-a-29-only-input}

Let \(\mathcal V_{29}\) be the span of all finite Heisenberg translates of the actual vacuum whose parameters are supported at places over29. This includes the actual permitted half-orbit used in AM-K022 and AM-R-012, and is deliberately larger than that half-orbit. At all other places every vector in \(\mathcal V_{29}\) has its fixed vacuum factor.

For an actual valid prime and decoration datum, suppose \(\Xi\not\equiv0\) is a nonterminal \nolinkurl{xiAt}. Then

\[
 \boxed{f_\Xi\notin\mathcal V_{29}.}
 \tag{DO29.24}
\]

Here equality is tested in one common finite model, enlarging the selected half set by the29 half places if necessary. The assertion is unchanged by a nonzero scalar multiplying \(\Xi\).

\textbf{Proof.} The prime axioms give three distinct rational primes \(T(u),g_0,g_1\). This is literal \nolinkurl{Primes.Ok.g_ne} and \nolinkurl{T_ne_g}, lines200--239 of \nolinkurl{Formalism/LocalToGlobal/Columns/AbelianHostFieldConditions.lean}. There is at least one selected place above each: their stated cardinality is \([F':\mathbf Q]>0\).

If \(T(u)\ne29\), choose a \(T\)-place \(P\nmid29\). Its factor (DO29.3) vanishes at0 because \(u\ne0\), and on an open neighborhood of0 by its period lattice. If \(T(u)=29\), choose a \(G\)-place \(P\nmid29\). Carried THETA3-A13 proves locally that nonvanishing of this factor forces the underlying matrix \(x:E'^a\to E'^b\) to be injective. In its \nolinkurl{XiEntry} proof, each enlarged summand has rank at most \(\operatorname{rank}x+c_i\) but requires a unit minor of size \(a+c_i\); in the penultimate branch \nolinkurl{joinCol} gives the same inequality with \(c_i=1\). Thus its local value at \(x=0\) is0, and the actual \nolinkurl{xiInv} or \nolinkurl{psiInvX} gives an open neighborhood on which it remains0. In all cases there is a mandatory factor at some \(P\nmid29\) that vanishes near0.

In a common half model, a vector in \(\mathcal V_{29}\) has function

\[
 G(t)=G_{29}(t_{29})
              \prod_{Q\in\mathcal S,\ Q\nmid29}1_{L_Q}(t_Q).
 \tag{DO29.25}
\]

This follows by applying only29-local operators to the tensor vacuum. The29 operators need not be pure position translations; their action is all contained in \(G_{29}\). If this function equalled the nonzero \(\Phi\) of (DO29.10), choose a point where the common value is nonzero. All its coordinates away from29 are in their \(L_Q\). Replace just its \(P\)-coordinate by0. Formula (DO29.25) leaves its nonzero value unchanged, whereas the mandatory source factor makes \(\Phi\) zero there. This is a contradiction. The same argument can be made entirely on rational samples: use a sufficiently small open coset at each selected place and additive weak approximation. The finite-model map is injective on the window by (DO29.15), so different functions cannot represent the same vector. \(\square\)

This obstruction is to the \textbf{29-only span}, not to the oscillator representation, not to arbitrary finite inputs, and not to a changed valid global tuple. Equation (DO29.14) gives the required nonzero vector with the other local factors included.

\subsection{5. What the mixed finite comparison proves}\label{what-the-mixed-finite-comparison-proves}

At every selected good relative split place away from2 and7, both the source trace coefficient and the mixed coefficient \(\xi'\) are local units. Express the two effective symplectic forms using the same fixed additive character. On \(X_P\oplus Y_P\) their pairing coefficients differ by a unit \(r_P\) of the real local field. The linear map

\[
 C_P(x+y)=x+r_Py,
 \qquad B_{\rm mixed}(C_Pz,C_Pw)=B_{\rm source}(z,w)
 \tag{DO29.26}
\]

is the explicit symplectic transport; \(r_P\) is by definition the ratio of the source coefficient to the mixed coefficient. It fixes the position variable and preserves both integral half lattices. It therefore fixes \(1_{L_P}\), transports every character in (DO29.11) to its actual dual vector, and preserves the finite measure normalization. It also commutes with \(X\mapsto hXg^{-1}\), since it is scalar on the dual half. The split-place determinant-root normalization from AM-DO-020 is retained. There is no extra unidentified phase in this normalized finite Heisenberg map.

Consequently the \textbf{same} raw table \(\Xi(\gamma)\) determines the mixed vector (DO29.14), with its exact finite norm and denominators. On the actual mixed theta family, AM-DO-026 gives its all-polynomial reading:

\[
 \vartheta_{r_P f_\Xi}(g,h)
 =C^{2ab}\epsilon_\chi(g,h)
   \left[P(\partial)\sum_{a,y}
     c_{a,y}\psi_f(-B(a,y)/2)
       \mathcal U_{A_{g,h}(y-a)}F_{N,N'}\right]_{0}.
 \tag{DO29.27}
\]

Every finite phase, chart scalar, and common period \(C^{2ab}\) is displayed or defined in the checked transported operator. This comparison constructs the finite part of a \textbf{new mixed datum}. It does not identify the old elliptic \nolinkurl{VacIn} seed with the primitive mixed-CM seed, and does not assert a new \nolinkurl{D7} instance without changing that wrapper. If a chosen auxiliary place is bad for (DO29.26), its actual lattice indices and dual pairing must be used instead; the abstract finite formula still holds but the unit-lattice simplification is unavailable.

\subsection{6. Consequences for the selected first edge and MainX}\label{consequences-for-the-selected-first-edge-and-mainx}

The identity \(\Xi(0)=0\), already settled in THETA3-A13, does \textbf{not} assert \(\vartheta_{r_Pf_\Xi}(1,1)=0\). The first is one value of a finite position function. The second is the weighted theta/jet sum (DO29.27). They are different linear functionals. In particular a nonintegral position translate from R012 itself can have finite value0 at the origin and a nonzero selected theta jet.

R012 proves nonvanishing of its selected jet on some member of the actual29 half-orbit. Equation (DO29.24) prevents using that statement to identify its finite vector directly with the nonzero prescribed \(f_\Xi\). Conversely (DO29.14) removes any merely existential finite-vector obstacle: the prescribed raw array is now an explicit admissible finite Heisenberg combination on the mixed host.

The remaining first-edge question is the concrete finite functional

\[
 \sum_{a,y}c_{a,y}\psi_f(-B(a,y)/2)
       [P(\partial)\mathcal U_{y-a}F_{00}]_0
             \stackrel{?}{\ne}0,
 \tag{DO29.28}
\]

or its actual matrix version over all allowed indices and decorations. A general nonzero theta vector or a nonzero term of this sum does not prove (DO29.28). The exact criterion in carried A15 remains applicable with the selected higher jet substituted for its first derivative. Its actual local incidence and rank restrictions remain inside the coefficients.

The next result treats a fixed nonzero auxiliary operator at primes disjoint from29 and11, acting on the \textbf{actual} K022 half-orbit. It will prove a quantitative selected-jet detector using simultaneous torsion-point interpolation. Even that result does not silently make a translated29 vacuum into the fixed \nolinkurl{Q}-factor at29 of an already prescribed \(\Xi\); that is a separate image or admissible-parameter equation.

\subsubsection{Check status}\label{check-status-3}

Author whole-proof review and independent review are recorded in \nolinkurl{definite_oscillator_CHECK_EVENTS.tsv} when actually performed. All new equations are finite-model or elementary lattice/measure arguments proved above. No Lean build or full MainX constructor is claimed.

\clearpage

\section{AM-K023 --- the actual selected jet matrix modulo 11}\label{am-k023-the-actual-selected-jet-matrix-modulo-11}

Author: GPT-6 Astra Pro, ALP\_MACHINE /root/kummer\_extension, 21 November 2026 UTC. Inherited reasoning effort is not separately exposed.

AM-K022 removes the finite half-orbit saturation loss from a specified finite jet matrix. This note computes the reduction of that jet matrix by a finite polynomial normal form. It also gives actual rectangular selected rows of unexpectedly large degree whose exact 11-content is zero. These are local selected kernel statements; they do not compute a later automorphic class integral.

\subsection{1. Every integral local section, not just the vanishing sigma section}\label{every-integral-local-section-not-just-the-vanishing-sigma-section}

Work over a finite unramified complete extension R of Z₁₁ containing the CM eigenvalues and the needed prime-to-11 torsion. Let k be its residue field. The ordinary Jacobian J, its principal trace-one line L=L₂∞, its q-normalized section, and the canonical source reduced gauge are those of AM-K018 and ALP-LIFT-08. Let F be the normalized nonvanishing even theta germ at D₁₃₅ from AM-K012. Its algebraic section is an integral unit at the origin, by ALP-LIFT-07, and F(0)=1. Its Hurwitz derivatives are integral.

Let κ be any positive integer prime to 11. Write 𝓜κ for the reduced germs of all integral formal local sections of L\^{}κ at zero. The following equality of integral modules is exact: \[
 \mathcal M_\kappa
   =\{\,h(u)F(u)^\kappa:
            h\text{ is an integral formal function on }\widehat J\,\}.
 \tag{NF1}
\] Here h(u) means substitution of the characteristic-zero group logarithm into an ordinary integral formal function. It does not mean that F is an integral ordinary power series in logarithmic coordinates.

\textbf{Proof.} The difference between L and the line of F is the rigidified two-torsion line \(\mathcal P_T\). Its pullback by {[}2{]} has its canonical rigidified trivialization. The {[}2{]}-pullback of F is a nonvanishing integral section near zero, so its κ-th power is an integral local frame for {[}2{]}*L\^{}κ under this trivialization. Hence any integral formal section S satisfies \[
 F_S(2u)=r(u)F(2u)^\kappa
\] for an arbitrary integral formal function r, and every such r occurs. The comparison has no exponential factor: this is the parallel rigidified trivialization of the pulled-back flat two-torsion line, exactly as in AM-K015. Since {[}2{]} is an integral formal automorphism at 11, replacing u by u/2 identifies these r with all integral formal functions h. This proves (NF1).

Every iterated invariant derivative of h has integral value at zero, because invariant vector fields preserve the integral formal coordinate ring. Thus h has integral Hurwitz coefficients. Products and inverses of Hurwitz series with integral coefficients and unit constant term are integral, as follows from the binomial product rule and the recursive coefficient equation for an inverse. Therefore every member of 𝓜κ is Hurwitz integral. This statement does not assert ordinary Taylor integrality. ∎

In particular the finite functional matrix B in AM-K022 can be taken integral at 11. Its presentation there as a matrix over Frac(R) was deliberately more general; changing its integral formal frame multiplies its jet columns by an integral invertible matrix and does not change its determinantal ideals.

\subsection{2. Universal restricted-operator relations}\label{universal-restricted-operator-relations}

For the ordinary coordinates u₁,u₂,u₃, AM-K017 proves both the algebraic restricted powers \[
 D_1^{11}=0,\qquad D_2^{11}=-D_1,\qquad D_3^{11}=0
\tag{NF2}
\] on the special fibre, and the full Hurwitz identities for F. The same identities on the completed local ring follow by continuity. For every f in 𝓜κ they imply \[
 \boxed{
 \partial_1^{11}f=4\kappa u_3f,\qquad
 \partial_2^{11}f=-\partial_1f,\qquad
 \partial_3^{11}f=-\kappa u_2f
 }\quad(\bmod11).
\tag{NF3}
\]

\textbf{Proof.} Write H=log F. AM-K017 gives \[
 \partial_1^{11}H=4u_3,\quad
 \partial_2^{11}H=-\partial_1H,\quad
 \partial_3^{11}H=-u_2.
\] The Hurwitz augmentation ideal has zero eleventh powers. Since the first derivatives of H have zero constant term, the prime Bell identity applied to κH gives (NF3) for F\^{}κ. For f=hF\^{}κ, the eleventh derivative of a product is a derivation modulo 11. The derivative of h is prescribed by (NF2). Thus the first and third formulas multiply by h; the middle one becomes \[
 (-D_1h)F^\kappa-h\,\partial_1F^\kappa
              =-\partial_1(hF^\kappa).
\] This proves (NF3) on the full module. In particular no division by the nonunit initial quadratic of sigma is used. ∎

Apply this to J² with ordinary coordinates u,v. Choose E′=Q(ζ₂₈) coordinates \[
 z_{i,+}=u_i+i v_i,\qquad z_{i,-}=u_i-i v_i.
\tag{NF4}
\] This is an integral change of coordinates at 11. The sign is the eigenvalue of the extra CM endomorphism \(I(u,v)=(-v,u)\). Its tangent derivations are \[
 \partial_{i,+}=(\partial_{u_i}-i\partial_{v_i})/2,\quad
 \partial_{i,-}=(\partial_{u_i}+i\partial_{v_i})/2.
\] Since \(i^{11}=-i\), the exact identities on every integral local section of the product line are \[
 \boxed{
 \partial_{1,\epsilon}^{11}f=2\kappa z_{3,\epsilon}f,\qquad
 \partial_{2,\epsilon}^{11}f=-\partial_{1,-\epsilon}f,\qquad
 \partial_{3,\epsilon}^{11}f=-\frac{\kappa}{2}z_{2,\epsilon}f.
 }\tag{NF5}
\] For example, the first right side is \((4\kappa/2)(u_3+i v_3)=2\kappa z_{3,+}\). The middle formula exchanges signs because the coefficient i is conjugated by Frobenius. The same formulas hold entry by entry on the actual Hom variety A\^{}\{ab\}=J\^{}\{2ab\}. A simultaneous unit normalization of coordinates gives the conjugate formulas, with unchanged matrix ranks and content.

\subsection{3. A finite normal form computes B modulo 11}\label{a-finite-normal-form-computes-b-modulo-11}

For each matrix entry introduce six commuting derivative symbols \(Z_{i,\epsilon}\). Let \(\Lambda_f(P)=[P(\partial)f]_0\). Equation (NF5) gives the following exact evaluation rules over k: \[
\begin{aligned}
 Z_{1,\epsilon}^{11}Q
     &\longmapsto 2\kappa\,\partial_{Z_{3,\epsilon}}Q,\\
 Z_{2,\epsilon}^{11}Q
     &\longmapsto -Z_{1,-\epsilon}Q,\\
 Z_{3,\epsilon}^{11}Q
     &\longmapsto -(\kappa/2)\partial_{Z_{2,\epsilon}}Q.
\end{aligned}\tag{NF6}
\] For instance \(Q(\partial)(z_jf)|_0
=(\partial_{Z_j}Q)(\partial)f|_0\). This is the ordinary Leibniz rule, including its exact integer multiplicities.

Apply (NF6) linearly until every individual exponent is less than 11. Each nonzero replacement lowers total degree, by 12, 10, or 12 respectively. Thus the process terminates. Let Tκ(P) denote its output, in the vector space \[
 \mathcal C=\operatorname{span}_k
   \{Z^\beta:0\leq\beta_j<11\text{ for every individual variable}\}.
\tag{NF7}
\] The procedure is confluent, and the following matrix statement is exact.

\subsubsection{Proposition NF8}\label{proposition-nf8}

For any finite chosen family P₁,\ldots,Pq of degrees at most D, the rank modulo 11 of its integral source jet matrix B equals the rank of the coefficient vectors \[
                    T_\kappa(P_1),\ldots,T_\kappa(P_q).
\tag{NF8}
\] The normalization factors \(M^{-\deg P_j}\) in AM-K022 are units and do not change this rank. For one polynomial P, its derivative image ideal has valuation zero if and only if \(T_\kappa(P)\ne0\). If its normal form is zero, the image is contained in 11R. The latter assertion alone gives no exact higher content.

\textbf{Proof.} All replacements preserve the evaluation on every member of 𝓜κ, by (NF6). The evaluation functionals associated with the cube monomials of degree at most D are linearly independent over k. Indeed choose integral algebraic formal parameters with the prescribed tangent basis. The degree-ordered matrix evaluating \(Z^\gamma\), \(0\leq\gamma_j<11\), on \(t^\beta F^\kappa\) has zero entry if \(|\gamma|<|\beta|\); in equal degree its entries are zero except at \(\gamma=\beta\), where they are \(\beta!\). These diagonal entries are units because every βj\textless11. Thus this square triangular submatrix is invertible.

Any two terminal reductions of P have identical functionals on all integral prefactors. Their difference lies in the cube and is therefore zero, proving confluence. The same independence proves (NF8). All the cube degrees encountered are at most the original degree, so the argument uses precisely the finite D-jet matrix. ∎

For several selected rows this computes exactly the number of Smith factors that are 11-adic units. It does not compute positive valuations of the remaining factors, and it does not iterate a congruence modulo 11 as a divisibility statement modulo higher powers.

Combining (NF8) with AM-K022, J30, gives the identical rank and unit-content criterion for the actual permitted half-orbit tuple matrix. Thus this calculation concerns actual allowed kernels, not an enlarged arbitrary Schwartz span.

\subsection{4. Closed calculation for rectangular selected rows}\label{closed-calculation-for-rectangular-selected-rows}

Take source rank a and target rank b\textgreater a. For each of the six labels use the first a columns, and write \(X_{i,\epsilon}\) for that a-by-a matrix of derivative symbols. Define \[
 P_{\mathbf n}=
 \prod_{\epsilon=\pm}\prod_{i=1}^3
                 \det(X_{i,\epsilon})^{n_{i,\epsilon}}.
\tag{NF9}
\] It is the primitive highest bideterminant for the rectangular partitions \((n_{i,\epsilon}^a)\), with determinant twists restored as in AM-DO-025.

Put \[
 C_a(n)=n(n+1)\cdots(n+a-1).
\tag{NF10}
\] The classical ordinary Cayley identity, in its integral polynomial specialization, is \[
 \det(\partial_X)\det(X)^n=C_a(n)\det(X)^{n-1}
 \quad(n\ge1),
\tag{NF11}
\] with zero on a constant. The primary reference is Caracciolo, Sokal, and Sportiello (2013), Theorem 2.1, equation (2.1); the directly consulted author preprint arXiv:1105.6270v2 gives this on printed p.~7. The short argument below checks the normalization used here; the identity is not a new result.

For completeness, over Q the left side is a relative invariant of the left GLₐ action with character det\^{}\{n−1\}, so on the dense invertible matrices it is a constant times det\^{}\{n−1\}. At X=I, the derivative indexed by a permutation π has value \(\operatorname{sgn}(\pi)n^{c(\pi)}\): each disjoint cycle can be assigned independently to one of the n determinant factors in \(\det(I+T)^n\), and its sign is the cycle sign. The signs in det(∂) cancel these signs. The resulting sum is \(\sum_\pi n^{c(\pi)}\), equal to \(n(n+1)\cdots(n+a-1)\) by inserting the last symbol either as a new cycle or at one of the a−1 positions in the old cycles. This proves (NF11) over Q and hence as an integral polynomial identity.

Since \(\det(X)^{11}=\det(X^{[11]})\) modulo 11, rules (NF6) and (NF11) give the following exact recurrences for the evaluated polynomial:

\begin{longtable}[]{@{}
  >{\raggedright\arraybackslash}p{(\columnwidth - 4\tabcolsep) * \real{0.3333}}
  >{\raggedright\arraybackslash}p{(\columnwidth - 4\tabcolsep) * \real{0.3333}}
  >{\raggedright\arraybackslash}p{(\columnwidth - 4\tabcolsep) * \real{0.3333}}@{}}
\toprule\noalign{}
\begin{minipage}[b]{\linewidth}\raggedright
Condition
\end{minipage} & \begin{minipage}[b]{\linewidth}\raggedright
Scalar factor
\end{minipage} & \begin{minipage}[b]{\linewidth}\raggedright
Exponent change
\end{minipage} \\
\midrule\noalign{}
\endhead
\bottomrule\noalign{}
\endlastfoot
\(n_{1,\epsilon}\ge11\) & \((2\kappa)^a C_a(n_{3,\epsilon})\) & \(n_{1,\epsilon}\mapsto n_{1,\epsilon}-11,\ n_{3,\epsilon}\mapsto n_{3,\epsilon}-1\) \\
\(n_{2,\epsilon}\ge11\) & \((-1)^a\) & \(n_{2,\epsilon}\mapsto n_{2,\epsilon}-11,\ n_{1,-\epsilon}\mapsto n_{1,-\epsilon}+1\) \\
\(n_{3,\epsilon}\ge11\) & \((-\kappa/2)^a C_a(n_{2,\epsilon})\) & \(n_{3,\epsilon}\mapsto n_{3,\epsilon}-11,\ n_{2,\epsilon}\mapsto n_{2,\epsilon}-1\) \\
\end{longtable}

A factor Cₐ(0) means that the result is zero. These are finite reduction identities for functionals, not identities of the original symbol polynomials before quotienting by (NF6).

\subsubsection{Corollary NF12. Exact unit-content examples beyond the factorial threshold}\label{corollary-nf12.-exact-unit-content-examples-beyond-the-factorial-threshold}

Suppose all six exponents equal n=11q+r, where \[
 1\le q\le r,\qquad r+q\le10,\qquad a+r\le11.
\tag{NF12}
\] Then the exact 11-content of this selected row is zero. More explicitly its nonzero normal form is \[
\begin{split}
 T_\kappa(P_{\mathbf n})={}&
 \kappa^{4aq}
 \left(\prod_{\ell=0}^{q-1} C_a(r-\ell)\right)^4\\
 &{}\cdot
 \prod_{\epsilon=\pm}
    \det(X_{1,\epsilon})^{r+q}
    \det(X_{2,\epsilon})^{r-q}
    \det(X_{3,\epsilon})^{r-q}.
\end{split}\tag{NF13}
\]

\textbf{Proof.} Reduce the two first-label exponents q times. Each third-label exponent becomes \(10q+r\). Reduce these q times; each middle-label exponent becomes \(10q+r\), and the third exponents become r−q. Reduce the middle exponents q times; they become r−q and each first-label exponent becomes r+q. The scalar factors in the table give (NF13); the signs and powers of 2 cancel. Every factor in every \(C_a(r-\ell)\) has an integer representative between 1 and 10, by (NF12), so the displayed scalar is nonzero modulo 11.

Each remaining determinant power has exponent at most 10. Therefore each individual symbol exponent is less than 11, and its polynomial lies in the cube. Its diagonal monomial has coefficient one, so the result is nonzero. Proposition NF8 proves zero content. ∎

Examples include n=12 for any even a≤10. For a=6, q=r=5 gives n=60 and total degree \(6an=2160\), yet the selected row has a unit value on some actual permitted tuple once M≥6482. The least permitted power is M=29³=24389, and its finite Schwartz vector still has norm one. The simple highest-symbol factorial bound in AM-K022 would only have given \(6a\,v_{11}(60!)=180\), so the normal form improves that finite bound to the exact value zero.

These rectangles have an explicit source HC specialization at the first edge. For any already admissible rank record with k=a and d=t+1, put H=(d+2)/2 and, at all six embeddings, \[
 b_j=j(d+2),\qquad e_j=H+j(d+2),\qquad
 u_i=a+n-\tfrac32-i,\qquad v_i=a+n-2-i.
\tag{NF14}
\] Every interlacing difference is 1/2, the two e,b gaps equal H, and the last u is n−1/2\textgreater−b₀=0. The integer/half-integer conditions hold. Since \(w_{0,i}=u_i-\rho_a(i)\), equation DO25.5 gives exactly \(\nu_{0,i}=n\). Thus these are actual selected first-edge partitions, not freely relabelled Schur weights. For example the HC data a=6,d=14,n=60 give the selected edge 6→20. This checks the HC interface; any separate ambient arithmetic field or global size conditions retain their own status.

There is also a basic unbounded-rank family: if all rectangular exponents are at most 10, the highest polynomial already lies in the cube, for every a. Its row has zero content regardless of its total degree, and (NF14) realizes these first-edge HC weights whenever an admissible rank record is supplied. Neither family by itself controls every later stage of one fixed growing chain.

\subsection{5. Scope and sources}\label{scope-and-sources}

The novel step is (NF1)--(NF8): the actual K017 relations extend to every integral local section, giving a rigorous finite normal form for B rather than a congruence for only one theta function. NF12--NF13 are the new selected-row application inside that finite quotient; NF11 is the classical Cayley identity with its integral normalization explicitly checked.

Caracciolo, S., Sokal, A. D., \& Sportiello, A. (2013). Algebraic/combinatorial proofs of Cayley-type identities for derivatives of determinants and pfaffians. Advances in Applied Mathematics, 50(4), 474--594. Primary author preprint: \url{https://arxiv.org/pdf/1105.6270}, version 2 (31 December 2012), Theorem 2.1 and equation (2.1), printed p.~7. Only the ordinary determinant identity and its normalization are used here.

Carried inputs are AM-K012, AM-K015, AM-K017, AM-K018, ALP-LIFT-07--08, and AM-K022; their exact primary/source locators are in those reports. The actual E′ coordinates and selected Cauchy decomposition are AM-DO-022 and AM-DO-025, especially DO25.3 and DO25.5--9. No extra claim about reduced ordinary Taylor integrality, bounded differential-module saturation, or a later automorphic class matrix is used.

The literal HC check (NF14) uses \nolinkurl{Formalism/LocalRingCardinality/Params.lean,} lines 50--69, and the actual selected-weight identity DO25.5. The finite independent arithmetic certificate is \nolinkurl{selected_jet_normal_form_certificate.py,} with output in selected\_jet\_normal\_form\_certificate.json. It checks the integer determinant-two Cayley identity for n=1,\ldots,40; 250 one-entry full normal forms; the complete six-matrix determinant-two expansions at n=12 and n=24 (4096 and 531441 nonzero input monomials modulo 11); a vanishing example at n=23; and the rank-six n=60 level arithmetic. It does not infer any infinite or higher-content conclusion from those finite checks.

The full finite jet matrix is integral and its rank modulo 11 is now computable from finite polynomial operations. Positive Smith valuations, control as every selected weight grows, and the effect of restricting the next class matrix to a previously produced theta image remain separate problems.

\clearpage

\section{AM-K024 --- a scalable sharp upper bound in the actual ordinary E′ selected rows}\label{am-k024-a-scalable-sharp-upper-bound-in-the-actual-ordinary-e-selected-rows}

Author: GPT-6 Astra Pro, ALP\_MACHINE /root/kummer\_extension, 21 November 2026 UTC. Inherited reasoning effort is not separately exposed.

This is a fixed-rank, varying-HC-family statement for the actual ordinary host A=J² with E′=Q(ζ₂₈). It applies to the selected compact Cauchy summands identified in AM-DO-025. It uses the exact primitive-kernel coefficient from accepted DP58. It does not introduce conjugate-action J factors, and it does not identify the output of a later automorphic class integral.

The result gives both a bound for the finite jet functional matrix B and a single allowed finite Schwartz input realizing that bound. It excludes a uniform local slope larger than 1/1220 across the indicated source-selected family. There is no net global cost conclusion.

\subsection{1. Fixed selected edge and exact statement}\label{fixed-selected-edge-and-exact-statement}

Fix an admissible source rank record and a stage j with \[
 a=N_j,\qquad b=N_{j+1}=a+d,\qquad d=t+1\ge2.
\tag{SC1}
\] Fix word-independent source HC data at each of the six holomorphic embeddings of E′. Let \(\nu_\tau^*\) be the resulting actual selected partitions on this edge. They are positive antitone partitions of length a by AM-DO-025. Put \[
 D_0=\sum_\tau|\nu_\tau^*|,\qquad
 \mathcal F_0=\prod_{\tau}\prod_{i=1}^{a}\nu_{\tau,i}^*!,\qquad
 c_0=v_{11}(\mathcal F_0).
\tag{SC2}
\] All these data are fixed as r varies.

Label the six embeddings by \((i,\epsilon)\), \(i=1,2,3\), \(\epsilon=\pm\), using \(z_{i,\pm}=u_i\pm i v_i\) on A=J². These are the actual E′ eigencharacters, with the underlying J character i. Set \[
 A_r=11^r,\qquad B_r=121\,11^r,\qquad
 T_r=A_r+B_r+D_0=122\,11^r+D_0.
\tag{SC3}
\]

\subsubsection{Theorem SC4}\label{theorem-sc4}

For each r≥0 there are integers \(0\le h_r\le A_r\), \(0\le k_r\le B_r\) and actual source HC data whose selected partitions on this edge are \[
\begin{array}{ll}
 \nu_{3,+}=\nu_{3,+}^*+h_r e_1,&
 \nu_{3,-}=\nu_{3,-}^*+(A_r-h_r)e_1,\\
 \nu_{2,+}=\nu_{2,+}^*+k_r e_1,&
 \nu_{2,-}=\nu_{2,-}^*+(B_r-k_r)e_1,\\
 \nu_{1,+}=\nu_{1,+}^*,&
 \nu_{1,-}=\nu_{1,-}^*.
\end{array}\tag{SC4}
\] There is a primitive integral polynomial \(P_r\) in the exact source-and-target Cauchy summand for these selected types whose degree is Tᵣ and whose normalized source jet functional has image-ideal valuation at most \[
                         \frac{11^r-1}{10}+c_0.
\tag{SC5}
\]

Choose the least M=29\^{}e≥3Tᵣ+2 and κ=M². A single tuple of the actual DP59 half-orbit then satisfies \[
 a_r=M^{-T_r}[P_r(\partial)F_{\mathcal S_{\mathbf t'}}]_0\ne0,
 \qquad
 v_{11}(a_r)\le\frac{11^r-1}{10}+c_0.
\tag{SC6}
\] Its finite Schwartz vector has norm one; no interpolating coefficient height is used. Its actual identity theta-kernel value is \(C^{2ab}a_r\), with the one fixed period \(C=\theta_0c\).

Consequently, for any γ\textgreater1/1220 and fixed constant C₀, the lower bound \[
              v_{11}(a)\ge\gamma T-C_0
\tag{SC7}
\] cannot hold for every allowed selected local coefficient in this fixed-rank HC family. This assertion concerns the union of the explicitly displayed selected types. It does not assert a bound for every vector in one fixed type or one fixed later theta image.

\subsection{2. The source HC change is literal}\label{the-source-hc-change-is-literal}

At each embedding, increase only the initial entry u₀ of the fixed source HC record by an arbitrary nonnegative integer ητ. Leave its other u entries, every v entry, and every e,b entry fixed.

In \nolinkurl{Formalism/LocalRingCardinality/Params.lean,} lines 50--69, the only interlacing condition involving u₀ is \(u_0>v_0\). It remains true after this increase. The last-u condition is unaffected because k≥2; the half-integer condition remains true; the e,b gaps and integer conditions are unchanged. Thus this is an actual HC record. Choosing the same change for every word preserves the source's word-compatibility equations.

The initial source weight is \(w_{0,i}=u_i-\rho_k(i)\). Increasing u₀ by ητ increases its first entry by ητ and changes no other entry. The exact recurrence DO25.6 copies this change into the first entry of every subsequent source weight. Formula DO25.5 therefore gives \[
                     \nu_{j,\tau}=\nu_{j,\tau}^*+\eta_\tau e_1.
\tag{SC8}
\] Taking the four η's in (SC4) gives precisely those selected partitions. Their dominance is preserved by increasing the first entry.

All e entries remain unchanged, so the actual common source-and-target character on every edge is unchanged. This is stronger than freely choosing six unrelated partitions. It does not, however, preserve a size window that independently bounds u₀. Such a window must grow with the indicated HC changes.

\subsubsection{Full analytic Hyp3 is preserved when the fixed baseline satisfies it}\label{full-analytic-hyp3-is-preserved-when-the-fixed-baseline-satisfies-it}

Suppose additionally that the fixed baseline TPar satisfies the source's stronger Hyp3 in CayleySignature.lean, lines 110--118, including the condition \(1801\le r+t\). Then every u₀ modification used above satisfies that same Hyp3.

To check every clause, the 1801 condition and the parity and eta clauses are unchanged. Interlace was checked above. DataTheta.lean, TPar.Dpar, lines 538--544, depends only on v,b,e,t,eta, so every discrete parameter set and therefore every NCGaps inequality is literally unchanged. In TPar.Centred, lines 550--554, only the inequalities involving u₀ can change. Interlacing implies \(u_i\ge u_{\mathrm{last}}>-b_0\), and \(e_j\ge e_0>b_0\): the e,b entries alternate strictly increasingly. Thus \(u_i+e_j>0\). Increasing u₀ consequently increases its absolute value \(|u_0+e_j|\); its required lower bound is preserved. All other Centred inequalities are unchanged.

This proves full analytic Hyp3 preservation, as well as preservation of IsMu because e is fixed. Hence SC4--SC20 apply within a genuinely Hyp3-admissible analytic parameter family whenever one such baseline is supplied. They still do not preserve separately fixed SizeData₂ upper windows or identify an iterated arithmetic row.

\subsection{3. A primitive source-highest polynomial with one spare column}\label{a-primitive-source-highest-polynomial-with-one-spare-column}

Use an a-by-b symbol matrix \(Z_\tau\) for each label. Reserve column 1. On columns 2,\ldots,a+1 define the residual prefix-minor polynomial \[
 Q_0(Z)=
 \prod_\tau\prod_{\ell=1}^{a}
 \det\bigl(Z_{\tau,\{1,\ldots,\ell\},
                         \{2,\ldots,\ell+1\}}\bigr)^{
                           \nu_{\tau,\ell}^*-\nu_{\tau,\ell+1}^*},
 \qquad \nu_{\tau,a+1}^*=0.
\tag{SC9}
\] These columns exist because b≥a+1. The same columns can be reused for the six labels: they are independent sets of variables.

For \(0\le h\le A_r\), \(0\le k\le B_r\), put \[
\begin{split}
 P_{r,h,k}={}&
 Z_{3,+,1,1}^{h}
 Z_{3,-,1,1}^{A_r-h}
 Z_{2,+,1,1}^{k}
 Z_{2,-,1,1}^{B_r-k}\,Q_0(Z).
\end{split}\tag{SC10}
\] Every factor in Q₀ is invariant when an earlier row is added to a later row. A monomial using only row 1 retains this invariance. The row degrees are exactly (SC4), with h,k in place of hᵣ,kᵣ. The product of the residual diagonal terms and the displayed first-column monomial has coefficient one. Hence Pᵣ,h,k is primitive in the polynomial coefficient lattice.

\subsubsection{Exact common source-and-target Schur label}\label{exact-common-source-and-target-schur-label}

The row-invariance assertion places this polynomial in the source-highest subspace of its indicated row weight. It also determines the target representation. Indeed, for each embedding the actual Cauchy decomposition is \[
 \operatorname{Sym}(T\otimes S^*)
       =\bigoplus_\mu S_\mu(T)\otimes S_\mu(S)^*.
\tag{SC11}
\] In a source irreducible \(S_\mu(S)^*\), the specified maximal-unipotent invariant line has the source highest weight corresponding to μ, and no other such invariant weight. Since the partitions label pairwise distinct source highest weights, the invariant subspace of the row weight ν in (SC11) occurs only in its μ=ν summand. Thus \[
 P_{r,h,k}\in
 \bigotimes_\tau
      S_{\nu_\tau}(T_\tau)\otimes
              (\text{source highest line in }S_{\nu_\tau}(S_\tau)^*).
\tag{SC12}
\] Restoring the fixed determinant twists in DO25.3 gives exactly the contragredient selected source type and the selected target type DO25.7--9.

The target vector need not be its highest vector, but it lies wholly in that same irreducible target representation. It does not split among different target Schur labels. No rational projector or projection denominator is introduced: the integral polynomial itself is already in the required Cauchy summand. A demand for one separately fixed target-highest coefficient would be an additional restriction, not part of (SC12).

This argument is over characteristic zero, where the source comparison lives. Its integral assertion is the explicit polynomial with coefficients in R. Any further comparison with a differently normalized abstract Weyl pairing retains the normalization factors studied in the parent's selected-Schur metric report; none is set to one here.

\subsection{4. The actual primitive-kernel input and the exact filler factorial}\label{the-actual-primitive-kernel-input-and-the-exact-filler-factorial}

The Hom variety is \(A^{ab}=J^{2ab}\) with its ordinary J factors. On the first ordinary J factor at row 1, column 1, let X₃ and X₂ be the differential symbols for the two primitive integral directions of actual E-labels 3 and 2. They satisfy \[
             X_3=Z_{3,+,1,1}+Z_{3,-,1,1},\qquad
             X_2=Z_{2,+,1,1}+Z_{2,-,1,1}.
\tag{SC13}
\] This follows from the exact coordinate change z±=u±iv. Any unit changes to the DP49 unit-arrow display coordinates preserve the valuations below.

Accepted DP58, equations ES11--12 and the mixed monomial evaluation recorded in DP65 EH16--17, gives the actual unit-parameter primitive kernel in these ordinary directions: \[
 v_{11}\!\left(
 \partial_{X_3}^{A_r}\partial_{X_2}^{B_r}
                   F_{-\kappa}(0)\right)
                         =\frac{11^r-1}{10}.
\tag{SC14}
\] The parameter −κ is a unit for every allowed M. AM-K018 matches the full-O\_E source gauge on these ordinary J factors. DP54 and DP59 give the integral unit algebraic prefactor that cancels the local frame unit, so (SC14) is an actual permitted integral-local-section functional value. The pure primitive kernel is not substituted for the fixed sigma section without this exact prefactor comparison.

On the filler columns choose the ordinary integral algebraic monomial whose linear logarithmic term is \[
              \prod_{\tau,i}z_{\tau,i,i+1}^{\nu_{\tau,i}^*}.
\tag{SC15}
\] It has degree D₀. The residual highest polynomial Q₀ has matching monomial coefficient one; every other supported monomial has a different exponent. Since its degree is D₀, higher terms of the coordinate changes or of every reduced frame cannot contribute. Therefore the filler evaluation is exactly \[
                    Q_0(\partial)(\text{filler germ})|_0
                                      =\mathcal F_0.
\tag{SC16}
\] No factorial is inverted. The underlying change between the ordinary J coordinates and the E′ coordinates has unit determinant at 11. Thus the prefactor in (SC15) is integral in the actual formal coordinate lattice.

Take the product of the first-column primitive-kernel prefactor, these filler monomials, and normalized unit prefactors on unused entries. It is an actual integral formal local section of \(\mathcal L^\kappa\). Factorization between the first column and filler columns gives \[
 v_{11}\!\left(
 M^{-T_r}
 [X_3^{A_r}X_2^{B_r}Q_0(\partial)F_S]_0
 \right)
                         =\frac{11^r-1}{10}+c_0.
\tag{SC17}
\] The displayed \(M^{-T_r}\) is a unit. This is a construction in the finite jet module B; globalizing it is the separate integral interpolation step of AM-K022.

\subsection{5. Weight extraction without a projection denominator}\label{weight-extraction-without-a-projection-denominator}

Expanding (SC13) gives the exact identity of differential polynomials \[
 X_3^{A_r}X_2^{B_r}Q_0
    =\sum_{h=0}^{A_r}\sum_{k=0}^{B_r}
          \binom{A_r}{h}\binom{B_r}{k}P_{r,h,k}.
\tag{SC18}
\] Every coefficient is an integer. Apply this equality to the single integral formal section just constructed. If every nonzero component value had valuation greater than the right side of (SC17), their integral linear combination would have greater valuation too. Consequently at least one pair (hᵣ,kᵣ) has \[
 v_{11}\!\left(
  M^{-T_r}[P_{r,h_r,k_r}(\partial)F_S]_0
          \right)
                         \le\frac{11^r-1}{10}+c_0.
\tag{SC19}
\] Terms whose binomial coefficient is divisible by 11 cause no difficulty: the argument only uses integrality, and does not divide by a binomial coefficient. This proves the image-ideal bound (SC5) for an actual selected row. It does not assert equality or a matching lower bound for that row.

Now use AM-K022 with D=Tᵣ. Its allowed half-orbit span realizes every integral algebraic Tᵣ-jet for M≥3Tᵣ+2. Equivalently, J30 identifies its derivative image ideal with that of the finite jet module. Expanding one realizing section in its integral allowed basis and applying the nonarchimedean valuation inequality chooses an individual tuple obeying (SC6). The canonical tuple consists of global algebraic points and sections by carried DP59(c), even when the interpolating coefficients were complete-local.

The finite Schwartz vector is a single Heisenberg translate of the fixed norm-one vacuum. Hence its norm is one, and its translation coefficient vector is a standard basis vector. The actual all-polynomial identity DO26.4 gives the fixed period \(C^{2ab}\), with no degree-dependent replacement.

Finally \[
 \frac{(11^r-1)/10+c_0}{122\,11^r+D_0}
                              \longrightarrow\frac1{1220}.
\tag{SC20}
\] This proves the failure of every larger uniform slope in (SC7).

\subsection{6. What this closes and what it leaves}\label{what-this-closes-and-what-it-leaves}

The old DP65 construction used both ordinary and conjugate-action J factors in a generalized harmonic module. The present proof uses only ordinary J powers in the actual E′ Hom variety and proves membership in the precise compact source-and-target selected summand. The full-O\_E gauge comparison from AM-K018 therefore applies directly. The first-row growth is implemented inside the actual source HC record; the common edge character is unchanged.

The bound is for actual selected \textbf{kernel coefficients} in the permitted half-orbit. At the first edge, the parent's class-supported construction AM-R-012 can use such a nonzero coefficient, retaining its source mass, compact norm, selected-polynomial norm and common period. At a later edge, the input is already constrained to a previous automorphic theta image. The restricted class matrix may cancel these coefficients or avoid these witnesses. Neither (SC19) nor the exact finite saturation theorem identifies that image or preserves this value after integration.

Likewise, varying the first HC gaps grows the source size window. The archimedean Schur norm, source class mass, denominator normalizations and the global height identity remain separate costs. The local upper slope cannot by itself be promoted to a statement about the net error or gain of the full Alpöge machine.

\subsection{7. Exact input locators}\label{exact-input-locators}

\begin{itemize}
\tightlist
\item
  Accepted DP58, research/explicit\_sharp\_jets\_linear\_levels\_DP58.md, §3, ES11--12 and the mixed monomial evaluation; accepted DP65, research/embedding\_labelled\_harmonic\_witness\_DP65.md, §5, EH16--17, for the explicitly recorded evaluation of that same carried primitive kernel.
\item
  Accepted DP54, research/global\_theta\_finite\_jets\_DP54.md, and DP59, research/theta\_translate\_span\_DP59.md, §§6--8: the integral unit frame prefactor, its finite-jet realization and the actual allowed translate lattice.
\item
  AM-K018: full-O\_E source gauge match on ordinary J products; AM-K022: explicit M≥3D+2 threshold, factorial convention, permitted individual tuples, and equality of finite determinantal ideals.
\item
  AM-DO-022: ordinary E′ host and prime-to-11 coordinate change. AM-DO-025, DO25.3--9: exact determinant twists, Cauchy decomposition and selected source/target labels.
\item
  \nolinkurl{Formalism/LocalRingCardinality/Params.lean,} HCParam, lines 50--69: literal u₀ variation. \nolinkurl{Formalism/LocalToGlobal/Columns/AbelianHostFieldConditions.lean,} lines 462--469, and Columns/BaseLemmas.lean, lines 496--552: actual selected-weight identities, carried as in AM-DO-025.
\item
  \nolinkurl{Formalism/LocalRingCardinality/DataTheta.lean,} TPar.Interlace, Dpar and Centred, lines 529--554; \nolinkurl{Formalism/SharpRestriction/CuspSpanNormBounds/CayleySignature.lean,} NCGaps and the stronger Hyp3, lines 110--118: exact preservation of the full baseline analytic admissibility conditions.
\end{itemize}

All new reductions are proved above. No Lean execution, full later-chain realization, or finite computational certificate is claimed by this report.

\clearpage

\section{AM-K025 --- an adjacent-gap factorial bound for the actual highest bideterminant}\label{am-k025-an-adjacent-gap-factorial-bound-for-the-actual-highest-bideterminant}

Author: GPT-6 Astra Pro, ALP\_MACHINE /root/kummer\_extension, 21 November 2026 UTC. Inherited reasoning effort is not separately exposed.

\textbf{Scope.} The general bound below applies to every selected partition in the ordinary E′ construction. Its zero-content growing-chain example satisfies the literal HC equations only; it does not assert the additional Hyp3, NCGaps, or global SizeData₂ inequalities. Those stronger hypotheses may force much larger adjacent gaps. The result is a local kernel image-ideal bound before restriction to a previous automorphic theta image.

\subsection{1. The bound uses adjacent differences, not the row lengths}\label{the-bound-uses-adjacent-differences-not-the-row-lengths}

Fix a source rank a, target rank b≥a, and the six actual E′ holomorphic labels τ. Let \(\nu_\tau\) be nonnegative antitone partitions of length a. Put \[
 \delta_{\tau,n}=\nu_{\tau,n}-\nu_{\tau,n+1}\ge0,
 \qquad \nu_{\tau,a+1}=0,
\tag{AG1}
\] and use the actual primitive highest bideterminant \[
 P_\nu=\prod_\tau\prod_{n=1}^a
                  \det(Z_{\tau,[1,n],[1,n]})^{\delta_{\tau,n}},
 \qquad D=\sum_\tau|\nu_\tau|.
\tag{AG2}
\] Both its source and target vectors are highest vectors in the exact selected Cauchy summand. No projection onto either compact type is required.

For M=29\^{}e≥3D+2, let \(I_\nu\) be the image ideal of the normalized derivative functional \(M^{-D}P_\nu(\partial)\) on the actual integral formal section lattice, equivalently on the actual permitted DP59 half-orbit by AM-K022 J30. Then \[
 \boxed{
 0\le v_{11}(I_\nu)
       \le \sum_{\tau}\sum_{n=1}^a
                                n\,v_{11}(\delta_{\tau,n}!).
 }\tag{AG3}
\] Equivalently, the upper bound is \[
 \frac1{10}\sum_{\tau,n}
          n\bigl(\delta_{\tau,n}-s_{11}(\delta_{\tau,n})\bigr).
\tag{AG4}
\] If every adjacent difference is at most 10, then \(I_\nu=R\): its exact 11-content is zero, and an individual allowed tuple has unit value.

\textbf{Proof.} For one fixed label and one prefix size n, use the anti-diagonal determinant term \[
       (-1)^{n(n-1)/2}\prod_{i=1}^n Z_{i,n+1-i}.
\tag{AG5}
\] All these choices are simultaneously uniquely leading terms. To see this, choose a constant C\textgreater ab and assign positive variable weights \(w_{ij}=C-ij\). A term indexed by a permutation π has weight \(nC-\sum_i i\pi(i)\). Its maximum is uniquely achieved by the reversing permutation: if i\textless j but π(i)\textless π(j), swapping those two images decreases the sum by \((j-i)(\pi(j)-\pi(i))>0\). Repeated swaps give the unique reversal. Thus the product of the anti-diagonal terms is the uniquely leading term of (AG2) and has coefficient ±1.

Different prefix sizes use disjoint anti-diagonals, because their variables satisfy different equations i+j=n+1. Consequently the leading monomial has exponent exactly \(\delta_{\tau,n}\) on each of the n variables of that anti-diagonal. Its exponent factorial is therefore \[
                   \beta!=\prod_{\tau,n}
                                      (\delta_{\tau,n}!)^n.
\tag{AG6}
\] Apply the highest-symbol interpolation AM-K022 J12 to this monomial. It gives \(\pm\beta!\in I_\nu\), with no factorial inverse or reduced-frame integrality assumption. This proves the upper bound. AM-K023 NF1 proves Hurwitz integrality of the entire actual integral local section module at 11, so \(I_\nu\subset R\), giving the lower bound. If every δ≤10, (AG6) is a unit. The integral allowed basis then contains an individual tuple with unit evaluation. ∎

This can be substantially smaller than the diagonal-monomial bound \(\sum_{\tau,i}v_{11}(\nu_{\tau,i}!)\). The improvement is an exact monomial calculation inside the same highest bideterminant, with the same selected source and target types.

\subsection{2. Literal minimal HC chain: exact unit content at every stage}\label{literal-minimal-hc-chain-exact-unit-content-at-every-stage}

This section is \textbf{HC-only}. It does not claim admissibility under the stronger 1801 or NCGaps hypotheses.

Fix an already admissible rank record with \(k\) even, \(d=t+1\) even, and put H=(d+2)/2. At every holomorphic embedding use \[
 b_j=j(d+2),\qquad e_j=H+j(d+2),\qquad
 u_i=k-\tfrac12-i,\qquad v_i=k-1-i.
\tag{AG7}
\] The two interlacing gaps are 1/2; the last u is 1/2\textgreater−b₀; and both e,b gaps equal H. All required integer or half-integer conditions hold. Formula DO25.5 gives \(\nu_0=(1^k)\), while the literal recurrence DO25.11 has \(\Delta_j=e_{j+1}-e_j-d=2\). Therefore \[
 \nu_j=((2j+1)^k,(2j)^d,(2j-2)^d,\ldots,2^d)
\tag{AG8}
\] at each label. Its nonzero adjacent differences are \[
       \delta_k=1,\qquad
       \delta_{k+hd}=2\quad(1\le h\le j).
\tag{AG9}
\] Thus every selected highest bideterminant in this chain has exact local content zero at 11.

For source rank a=k+jd and target rank b=a+d, its total six-label degree and an exact interpolated value are \[
 D_j=6\{k(2j+1)+d\,j(j+1)\},
\tag{AG10}
\] \[
 M^{-D_j}P_{\nu_j}(\partial)F_S(0)
       =2^{\,6kj+3d\,j(j+1)}
\tag{AG11}
\] for a suitable integral formal section S; the six identical determinant signs multiply to +1. The level can be the least power 29 at least \(3D_j+2\). An individual permitted tuple has unit value, although its value need not be the specific power of 2 in (AG11). Its finite Schwartz norm is one.

The nonzero degrees grow quadratically with j. This rules out a positive local lower bound depending only on total selected degree across all HC records. It leaves open the more restrictive family forced by the global machine, and it does not imply that these unit coefficients survive any later class integral.

\subsection{3. The universal finite cutoff clarifies this scope}\label{the-universal-finite-cutoff-clarifies-this-scope}

AM-K023 NF5 implies, for every integral formal section on each A entry, \[
 \partial_{1,\pm}^{121}f=0,\qquad
 \partial_{3,\pm}^{121}f=0,\qquad
 \partial_{2,\pm}^{1331}f=0
                       \quad(\bmod11).
\tag{AG12}
\] Indeed iterating the first or third eleventh-power identity eleven times gives a multiple of \(z^{11}=0\) in the Hurwitz algebra; its coordinate multiplier commutes with that derivative. The middle identity then reduces its 1331st power to a 121st first-label power.

Consequently a nonzero evaluated derivative modulo 11 on A\^{}\{ab\} has, in each entry, exponents at most 120,120 on the first labels, 1330,1330 on the middle labels, and 120,120 on the third labels. Their sum is 3140 per entry. Therefore \[
             D>3140ab
 \quad\Longrightarrow\quad
 B(P)\subset11R
\tag{AG13}
\] for every homogeneous polynomial P of degree D, including every selected row. This is a uniform first divisibility statement for the full integral local section module. It cannot be iterated after dividing by 11 without a further stable-lattice argument.

There is no conflict with (AG8): its degree is far below this dimensional cutoff for the indicated ranks. More generally the adjacent-gap criterion and the cutoff give finite, independently checkable tests that any proposed admissible HC specialization must pass.

\subsection{4. Sources and status}\label{sources-and-status}

The finite interpolation and permitted span are AM-K022, using carried DP59. The integral local source gauge and Hurwitz module are AM-K018 and AM-K023. The exact selected partitions and determinant twists are AM-DO-025. The HC clauses are \nolinkurl{Formalism/LocalRingCardinality/Params.lean,} lines 50--69.

All new claims are proved above by the leading-monomial argument and the already established operator identities. No new computational certificate or Lean execution is asserted. The general bound (AG3) remains valid for the larger source gaps; only the special unit family is restricted to the stated HC scope.

\clearpage

\section{ALP-LIFT-16 --- fixed spherical tests, 29-supported single translates, and exact level masses}\label{alp-lift-16-fixed-spherical-tests-29-supported-single-translates-and-exact-level-masses}

Author: GPT-6 Astra Pro, ALP\_MACHINE /root/lift\_accounting, 21 November 2026 UTC. Inherited reasoning effort is not separately exposed.

This result sharpens ALP-LIFT-15 in the actual common \(E'=\mathbf Q(\zeta_{28})\) host. The later finite inputs can all be single norm-one Heisenberg translates supported only at29. The proof includes an exact split-place spherical-vacuum calculation valid in residue characteristic2. The translation depths remain unbounded, and the prescribed raw Xi and prior-image arithmetic lattice are not identified.

The additional quantitative output is an explicit sufficient principal level and its class mass. Its 11-adic valuation is independent of the unknown 29-depth. This does not compute a later gain: the arithmetic entries and their cancellations are separate.

\subsection{1. The split-place vacuum coefficient, with the actual measure}\label{the-split-place-vacuum-coefficient-with-the-actual-measure}

Let \(F\) be any nonarchimedean local field of characteristic zero, with residue cardinality \(q\), uniformizer \(\varpi\), and ring \(\mathcal O\). Put \[
 G=\operatorname{GL}_a(F),\quad K=\operatorname{GL}_a(\mathcal O),
                    \qquad b>a.
 \tag{S1}
\] Normalize Haar measure on \(G\) by \(\operatorname{vol}(K)=1\). At a split place of the common unitary host, AM-DO-020/022 identifies the actual normalized finite vacuum with \(\phi_0=\mathbf1_{M_{b,a}(\mathcal O)}\) after the unitary measure change determined by its lattice. Its \(L^2\) norm is one. The source action, without its determinant-root character, is \[
       \omega(g)\phi(X)=|\det g|^{b/2}\phi(Xg).
 \tag{S2}
\] For a general lattice measure, the normalized characteristic vector is \(\operatorname{vol}(L)^{-1/2}\mathbf1_L\); its matrix coefficient is unchanged by the corresponding unitary identification. Thus no additive-lattice volume is lost in the normalization used here.

Write \(a_\lambda=\operatorname{diag}(\varpi^{\lambda_1},\ldots,
\varpi^{\lambda_a})\), with \(\lambda_1\ge\cdots\ge\lambda_a\). A direct intersection-of-lattices calculation gives \[
 A_b(g):=\langle\omega(g)\phi_0,\phi_0\rangle
       =q^{-b\sum_i|\lambda_i|/2}
                  \quad(g\in Ka_\lambda K).
 \tag{S3}
\] Indeed the \(i\)-th column lies in \(\mathcal O^b\cap\varpi^{-\lambda_i}\mathcal O^b\), of volume \(q^{-b\max(-\lambda_i,0)}\), while the determinant factor is \(q^{-b\sum_i\lambda_i/2}\). The formula holds at2 as well as at odd residue characteristic. It uses the actual split Schrödinger action, not an assumed integral diagonal/antidiagonal basis for a doubled space.

\subsection{2. A matrix-disintegration proof of the exact spherical transform}\label{a-matrix-disintegration-proof-of-the-exact-spherical-transform}

Let \(\pi\) be an irreducible unramified unitary representation of \(G\), let \(m_\pi(g)\) be its normalized spherical matrix coefficient, and let \(\alpha_1,\ldots,\alpha_a\) be its standard normalized Satake parameters. Define \[
 C_{a,b}(q)=\prod_{i=0}^{a-1}(1-q^{i-b}),\qquad
                           s_*=(b-a+1)/2.
 \tag{S4}
\] We will prove, in the convergence range specified below, \[
 \boxed{\int_G A_b(g)m_\pi(g)\,dg
        =C_{a,b}(q)L(s_*,\pi)L(s_*,\pi^\vee).}
 \tag{S5}
\] For a unitary representation the two \(L\)-values at a real point are complex conjugates. Thus the right side is strictly positive whenever the factors are finite and nonzero.

\subsubsection{2.1 Exact full-rank matrix disintegration}\label{exact-full-rank-matrix-disintegration}

For an integer \(B\ge a\), put \[
             f_B(h)=\mathbf1_{M_a(\mathcal O)}(h)
                                           |\det h|^{B/2}.
\] Let \(X_0=\binom{I_a}{0}\in M_{B,a}(F)\). Every full-rank matrix can be written \(kX_0h\), with \(k\in\operatorname{GL}_B(\mathcal O)\), \(h\in G\). To see this, intersect its column space with \(\mathcal O^B\). This intersection is a saturated rank-\(a\) direct summand over the DVR, so an integral basis extends to a basis of \(\mathcal O^B\); its change-of-basis matrix is the required \(k\).

With probability Haar measure \(dk\) on \(\operatorname{GL}_B(\mathcal O)\), matrix Lebesgue measure satisfies \[
 \int_{M_{B,a}(F)}F(X)\,dX
  =C_{a,B}(q)\int_{\operatorname{GL}_B(\mathcal O)}
        \int_G F(kX_0h)|\det h|^B\,dh\,dk.
 \tag{S6}
\] The rank-deficient set has measure zero, as follows by induction and Fubini from the vanishing of a nonzero matrix minor. Both sides on the full-rank orbit are left compact-invariant and transform under right multiplication by \(h_0\in G\) with the same Jacobian \(|\det h_0|^{-B}\). Uniqueness of a relatively invariant measure on this transitive homogeneous space makes them proportional. The proportionality constant is fixed by the open set of matrices whose reduction has rank \(a\): in the right-hand coordinates this is exactly \(h\in K\), so the integral without the constant is one. Its additive measure is \[
 \frac{\prod_{i=0}^{a-1}(q^B-q^i)}{q^{Ba}}
                 =\prod_{i=0}^{a-1}(1-q^{i-B}).
\] This proves (S6), including its constant. The homogeneous-space measure assertion can equivalently be checked in a local minor chart; the same Jacobian and normalization give (S6).

Using (S6) on the product of two characteristic functions gives \[
       A_B(g)=C_{a,B}(q)\int_G f_B(hg)f_B(h)\,dh
             =C_{a,B}(q)(\check f_B*f_B)(g),
              \qquad\check f_B(h)=f_B(h^{-1}).
 \tag{S7}
\] The inverse-argument operator \(|\det g|^{-B/2}\phi(Xg^{-1})=\omega(g^{-1})\phi(X)\) has the same vacuum coefficient, since \(A_B(g^{-1})=A_B(g)\) by (S3). It is not the left action in(S2).

\subsubsection{2.2 The one-sided spherical zeta integral}\label{the-one-sided-spherical-zeta-integral}

With the normalization in (S1), the elementary unramified Godement--Jacquet calculation is \[
 \int_G\mathbf1_{M_a(\mathcal O)}(h)|\det h|^z
                     m_\pi(h)\,dh
       =\prod_{i=1}^a
            (1-\alpha_i q^{-z+(a-1)/2})^{-1}.
 \tag{S8}
\] Replacing \(m_\pi(h)\) by \(m_\pi(h^{-1})\) replaces the \(\alpha_i\) by their inverses. The product of the two formulas is therefore independent of this convention.

For an explicit normalization check, use normalized induction from unramified characters \(\chi_i\), with \(\chi_i(\varpi)=\alpha_i\), and its spherical vector \(v^\circ\) satisfying \(v^\circ(nak)=\delta_B(a)^{1/2}\prod_i\chi_i(a_i)\). Averaging \(v^\circ(kh)\) over \(K\) gives the normalized spherical function. Since the matrix indicator and determinant power are left \(K\)-invariant, that averaging disappears from the integral. In the Iwasawa integral \(h=nak\), integrality requires \(a_i\in\mathcal O\) and \(n_{ij}\in a_j^{-1}\mathcal O\). Integrating these upper off-diagonal entries contributes \(\prod_j|a_j|^{-(j-1)}\). The remaining exponent of each \(|a_i|\) is \[
          z-\frac{a+1-2i}{2}-(i-1)=z-(a-1)/2.
\] The \(a\) geometric series give (S8). This calculation is initially in a half-plane of absolute convergence. The normalized spherical function is determined by its Satake character, so it applies also at reducible principal-series parameters, by the usual spherical-function identity or its continuous specialization. It is not restricted to generic irreducible representations.

A primary independent locator for precisely this normalization is Macdonald (1995), ChapterV, section4, equations(4.4)--(4.5), printed p.301, with the convergence statement immediately following; his section3 records the inverse convention for the spherical transform. Thus (S8) is a known formula used with its constant checked here.

For \(B\) sufficiently large, \(f_B\in L^1(G)\); the spherical transform is multiplicative on \(K\)-bi-invariant convolution. Equations(S7)--(S8) give \[
 \int_G A_B(g)m_\pi(g)\,dg
  =C_{a,B}(q)
     L((B-a+1)/2,\pi)L((B-a+1)/2,\pi^\vee).
 \tag{S9}
\]

\subsubsection{2.3 Extension and absolute convergence at the actual rank}\label{extension-and-absolute-convergence-at-the-actual-rank}

Put \(r(\lambda)=\sum_i|\lambda_i|\), and introduce an independent complex variable \(T\). The Cartan sum \[
 J_\pi(T)=\sum_{\lambda_1\ge\cdots\ge\lambda_a}
   \operatorname{vol}(Ka_\lambda K)m_\pi(a_\lambda)
                                                    T^{r(\lambda)}
 \tag{S10}
\] is a power series convergent near zero. For all sufficiently large integers \(B\), (S9) identifies its value at \(T=q^{-B/2}\) with \[
 R_\pi(T)=
 \frac{\prod_{i=0}^{a-1}(1-q^iT^2)}
      {\prod_{i=1}^{a}
       (1-q^{(a-1)/2}T\alpha_i)
       (1-q^{(a-1)/2}T\alpha_i^{-1})}.
 \tag{S11}
\] These values accumulate at zero, where both functions are analytic. The identity theorem therefore proves equality of their power series near zero, and hence equality wherever the series converges and the rational function has no pole. This step does not assume that an identity at infinitely many large real dimensions alone determines a meromorphic function of the dimension.

For completeness, suppose \[
                    |\alpha_i|,\ |\alpha_i^{-1}|\le q^s
                        \qquad(1\le i\le a).
 \tag{S12}
\] Macdonald's explicit spherical formula and double-coset volume formula give constants \(C,A\), depending on this representation, such that \[
 \begin{split}
 |m_\pi(a_\lambda)|
  &\le C(1+r(\lambda))^A q^{-\langle\rho,\lambda\rangle}
                         q^{s r(\lambda)},\\
 \operatorname{vol}(Ka_\lambda K)
  &\le C q^{2\langle\rho,\lambda\rangle},
 \qquad
 \rho=((a-1)/2,(a-3)/2,\ldots,(1-a)/2).
 \end{split}
 \tag{S13}
\] These estimates follow directly from Macdonald (1995), ChapterV, (2.9) and (3.4), printed pp.298--299. For distinct Satake parameters the finite Weyl sum gives the first bound without the polynomial factor. If parameters coincide, the apparent denominators cancel; taking the finite-order limits differentiates monomials \(\alpha^{w\lambda}\), adding only a polynomial in the integer entries of \(\lambda\). There are finitely many Weyl terms and finitely many possible coincidences, so a fixed \(A\) suffices. This explains the bound also at repeated parameters. The same formula holds after a central shift of \(\lambda\); the displayed bound uses both inverse and direct parameters.

Since \(\langle\rho,\lambda\rangle\le(a-1)r(\lambda)/2\), (S13) proves absolute convergence of (S10) when \[
                       |T|q^{(a-1)/2+s}<1.
 \tag{S14}
\] Only polynomially many dominant integer tuples have a given \(r(\lambda)\), so the geometric decay suffices.

In the actual selected chain let \(a=N_j\), \(b=a+d\), \(j\ge1\). Formula(E13) of ALP-LIFT-14 and the initial bound(E18) imply \[
                |\alpha_i|,\ |\alpha_i^{-1}|
                        \le q^{(d-1)/2}.
 \tag{S15}
\] Indeed the initial rank is \(k<d\), all inserted blocks have length \(d\), and the finite character values on uniformizers have absolute value one. Unitarity gives the same bound for the inverse multiset. A further unramified unitary determinant twist preserves the bound.

At \(T=q^{-b/2}\), (S14) holds with \[
                |T|q^{(a-1)/2+(d-1)/2}=q^{-1}.
 \tag{S16}
\] Thus the actual integral in (S5) is absolutely convergent. Every denominator in (S11) has its nonconstant term of modulus at most \(q^{-1}<1\); the numerator is positive since \(b>a-1\). Substitution gives (S5), with a strictly positive value for unitary \(\pi\). This proof includes residue characteristic2.

\subsection{3. The actual split Rallis test, including the character ratio}\label{the-actual-split-rallis-test-including-the-character-ratio}

At a split place write the source common character as \(\mu=(\eta,\eta^{-1})\). Outside29 it is unramified and unitary. The actual Hom source operator contributes \(\eta(\det g)^{-1}\), by AM-DO-020/022. AM-DO-028 transports GQT's conjugated-input convention and dual-source group back to the literal bilinear integral. Its local matrix coefficient is therefore that of \[
                         \widetilde\pi=\pi\otimes\eta^{-1}.
\] The doubled evaluation identity \(\delta(\phi_0\otimes\overline{\phi_0})(0)
 =\|\phi_0\|^2=1\) shows that the section at \((g,1)\) is exactly the actual vacuum coefficient (S3) with this determinant-root factor. No choice of integral doubled coordinates is needed for this identity.

Consequently the actual local zeta value is \[
 Z_v(s_0;\phi_0,\phi_0,v^\circ,v^\circ)
  =C_{a,b}(q)
       L(s_*,\pi\eta^{-1})L(s_*,\pi^\vee\eta),
 \qquad s_0=d/2.
 \tag{S17}
\] These two factors are exactly the split local standard factor in the source convention fixed by AM-DO-028. The normalized value is \[
                            Z_v^*(s_0)=C_{a,b}(q)>0.
 \tag{S18}
\] If the chosen local group Haar measure gives \(K\) volume \(v_K\) instead of one, both (S17) and (S18) are multiplied by \(v_K\). This measure factor is retained. The formula is valid at the dyadic split places of the actual host; it closes the test-vector point explicitly left open at the end of ALP-LIFT-15.

\subsection{4. Fixed tests at all other places outside29}\label{fixed-tests-at-all-other-places-outside29}

At an odd nonsplit finite place, every stage character is trivial and the actual model is the integral unramified model, as checked in ALP-LIFT-14 section3. Liu (2021), Lemma6.1, identifies the doubled lattice characteristic vector with his unramified standard section. His Proposition5.6 with sign \(+\) gives \[
 Z_v(s_0)=\frac{L_v(s_*,\pi)}{b_a(s_0)},\qquad
 b_a(s)=\prod_{i=1}^a(1-(-1)^i q^{-2s-i})^{-1}.
 \tag{S19}
\] The odd residue characteristic makes the diagonal/antidiagonal lattice change integral, so the doubled product of the actual unit lattice vacua is this lattice characteristic vector. The normalization at the identity is one on both sides. The additive-character rescaling needed at a ramified base-field place was already carried out in LIFT14 without changing the composite symplectic character or normalized vacuum. Therefore the normalized local value is exactly \[
             b_a(s_0)^{-1}
             =\prod_{i=1}^a(1-(-1)^i q^{-d-i})>0,
 \tag{S20}
\] again multiplied by \(v_K\) for a different group Haar normalization.

Thus at every finite place outside29 the \textbf{actual fixed vacuum} and the unique spherical source line give a nonzero normalized local test. All exceptional dyadic places are split and were treated directly in section3. Almost everywhere these are the standard factors in GQT equation(11.3), so their global normalized product converges and is nonzero. At the real places the selected Schur test remains the exact positive \(\operatorname{vol}(U(a))/\dim\rho_{j,\tau}\) of LIFT14(E24). At29 arbitrary finite Schwartz functions remain available, and GQT Proposition11.6(i) supplies the local tests there.

\subsection{5. Restricting the global nonzero family to29}\label{restricting-the-global-nonzero-family-to29}

Let \(\pi\) be any actual cuspidal embedding arising in LIFT14. At each finite place outside29, \(\pi_v^{K_v}\) is one-dimensional. The Rallis argument, with the fixed tests of sections3--4, therefore gives a nonzero lift using spherical source vectors outside29 and the vacuum there, while varying the Schwartz function only at29 and in the selected real component.

This yields every-vector nonvanishing on the away-from29 spherical space, not only one representation-level occurrence. Here are the multiplicity and cancellation details.

For one global irreducible isomorphism type, its spherical space away from29 is \[
   \mathcal M_\pi\otimes\pi_{29}\otimes\rho_j
             \otimes\bigotimes_{v\nmid29,\ v<\infty}\pi_v^{K_v},
 \tag{S21}
\] where \(\mathcal M_\pi\) is its finite-dimensional actual automorphic multiplicity space. The remaining \(\pi_{29}\otimes\rho_j\) is irreducible for \(G(F'_{29})\times G_\infty\), with the product over the six29 places understood. The joint kernel of the theta family with all other finite inputs fixed is invariant under this group. In a finite sum of copies of an irreducible smooth module, a submodule is of the form \(\mathcal M'\otimes
(\pi_{29}\otimes\rho_j)\). If it were nonzero, a nonzero line in \(\mathcal M'\) would define an actual global cuspidal embedding on which all the fixed-away-from29 theta maps vanish. The Rallis argument of the preceding paragraph applies to that embedding, giving a contradiction. Hence this joint kernel is zero.

Different global source isomorphism types cannot cancel their theta images. All target constituents of one source have the unique local small-theta target. At a nonsplit place outside29 the source is spherical and its Satake multiset is recovered by subtracting the fixed inserted block in (E12). At a split place the reciprocal uniqueness in Mínguez's Howe correspondence gives the same injectivity, including at29. Thus two inequivalent global sources, whose real types are already fixed, give inequivalent global targets. Their automorphic isotypic subspaces are orthogonal. A nonzero image of one source component therefore cannot be canceled by the others.

Every smooth current vector has a finite spectral decomposition at its finite level and real type. The preceding two paragraphs show that every nonzero current vector fixed by the standard compacts outside29 survives a test varying only at29. The output is again fixed by those target compacts because the actual vacuum is fixed by them. Hence this property iterates from the first vector of AM-R-012.

Now apply the local version of ALP-LIFT-15 to the29-component. For some finite \(e_j\ge1\), it lies in the depth-\(29^{e_j}\) window. One basis vector detects the nonzero map. We obtain \[
 F_{j+1}=
  \Theta_{W_{x_j}\phi_{0,j},\,\phi_{\infty,j}}(F_j)\ne0,
 \quad
 x_j\in29^{-e_j}\Lambda_{29,j}/\Lambda_{29,j},
 \quad
 x_{j,v}=0\ (v\nmid29),
 \quad
 \|W_{x_j}\phi_{0,j}\|_{2,f}=1.
 \tag{S22}
\] Its infinite translation component is zero. The real vector stays in the exact selected Schur component. This is a full Heisenberg orbit, not the prescribed isotropic half-orbit.

\subsection{6. An explicit sufficient level and its exact index}\label{an-explicit-sufficient-level-and-its-exact-index}

Put \(M=29^e\), and let \(x\in M^{-1}\Lambda_{29}\). If \(g\) and \(h\) are principal-congruence elements modulo \(M^2\) in the source and target integral unitary groups, then \[
                 (g,h)x-x\in M\Lambda_{29}.
 \tag{S23}
\] This follows from the Hom action \(X\mapsto hXg^{-1}\). All four position/momentum matrix blocks satisfy the same integral congruence. Let the displacement be \(\delta\). The Weyl phase \(B(x,\delta)/2\) is integral because \(x\in M^{-1}\Lambda_{29}\), \(\delta\in M\Lambda_{29}\), the pairing is integral on the lattice, and29 is odd. The trace-one character of \(\delta\) is one at29. Thus \(W_{x+\delta}\phi_0=W_x\phi_0\) exactly. The common splitting character has conductor dividing29 and is one on the same principal subgroups. Covariance therefore shows that the whole translated finite vacuum is fixed by these source and target levels.

For rank \(n\), let \(K_n(M^2)\) be principal congruence modulo \(29^{2e}\) at all six split29 places and the standard compact elsewhere. Its exact index is \[
 \begin{split}
 I_n(M^2)&=[K_{n,0}:K_n(M^2)]\\
 &=29^{6n^2(2e-1)}
                  |\operatorname{GL}_n(\mathbf F_{29})|^6\\
 &=29^{12en^2}\prod_{i=1}^n(1-29^{-i})^6.
 \end{split}
 \tag{S24}
\] Therefore (S22) has an explicit sufficient target level of this form. The exponent \(e_j\) has not been bounded in terms of the selected degree.

Every rational class stabilizer at \(K_n(M^2)\) is trivial. Indeed it is finite, because it preserves the corresponding finite lattice and is compact at infinity. At one29 place it embeds into a conjugate of \(1+29M_n(\mathbf Z_{29})\), which is torsion free by the binomial valuation argument in AM-R-012. A conjugate has the same property. Consequently, with \(\operatorname{vol}(G_\infty)=1\) and \(\operatorname{vol}(K_{n,0})=1\), \textbf{every} finite class has mass exactly \(I_n(M^2)^{-1}\), not merely the identity class.

\subsection{7. The 11-adic and archimedean mass accounting}\label{the-11-adic-and-archimedean-mass-accounting}

The order of29 modulo11 is10, and \[
                           v_{11}(29^{10}-1)=1.
\] If \(10\nmid i\), \(v_{11}(29^i-1)=0\). If \(i=10\ell\), the elementary lifting-the-exponent identity gives \(v_{11}(29^i-1)=1+v_{11}(\ell)\). Hence, with \(r=\lfloor n/10\rfloor\), \[
 \boxed{v_{11}(I_n(M^2))
          =6\sum_{\ell=1}^{r}(1+v_{11}(\ell))
          =6\bigl(r+v_{11}(r!)\bigr).}
 \tag{S25}
\] It is independent of \(e\). The residue check can be made without a large integer: \(29^{10}\equiv89\pmod{121}\), so its difference from one has valuation exactly one. The standard lifting argument follows by expanding \((1+11u)^\ell-1\), first for \(\ell\) prime to11 and then iterating the 11th power. This proves the valuation formula.

For example the value in (S25) is0 at rank6,12 at rank20, 18 at rank34, and24 at rank48. Thus the unknown translation depth does not introduce an increasing 11-adic denominator through these class masses. This assertion concerns these Haar factors only.

The archimedean size of the index does depend on depth: \[
 \frac12\log I_n(M^2)
       =6en^2\log29+
                     3\sum_{i=1}^n\log(1-29^{-i}).
 \tag{S26}
\] For clarity, this is a class-coordinate normalization, not a claim that the already produced input can be replaced by an arbitrary class delta. If a scalar automorphic function of compact type \(\rho\) is fixed by this finite level, choose unitary bases and write its class restrictions as \(F_c(k)=\operatorname{tr}(A_c\rho(k))\). Since all stabilizers are trivial, Schur orthogonality gives \[
 \|F\|_2^2=
  \frac{1}{I_n(M^2)\dim\rho}
                  \sum_c\|A_c\|_{\rm HS}^2.
 \tag{S27}
\] Consequently each matrix entry of \(A_c\), in the chosen unitary bases, is at most \(\sqrt{I_n(M^2)\dim\rho}\,\|F\|_2\). Comparison with the source's primitive Schur coefficient lattice remains the separate normalization already studied in AM-R-011. Formula(S27) does not identify the arithmetic row.

\subsection{8. Remaining gate and exact scope}\label{remaining-gate-and-exact-scope}

The new construction controls the finite support and norm: all later inputs are one actual29-supported Heisenberg translate, with coefficient one and finite \(L^2\) norm one. It supplies an explicit sufficient principal level and exact class masses. The later selected real norms, the input-vector norms, and the output content are not set to one.

The depths \(e_j\) are not bounded. More fundamentally, these full-orbit inputs need not equal the source's specified raw Xi or lie in its DP59 half-orbit. AM-DO-029 records the actual finite expansion of that raw function and its separate local support conditions. Neither (S22) nor the mass formula proves that the prescribed row kernel acts nontrivially on the previous image lattice, or that its integral image is saturated.

No later \(g_j\), source \(e_j\)-error, complete \(H\), six comparison identities, or DP44 shortening is deduced from this result. The value of (S25) is useful concrete accounting data; it is not itself a row gain.

\subsection{Primary locators and carried inputs}\label{primary-locators-and-carried-inputs}

\begin{itemize}
\tightlist
\item
  Macdonald, I. G. (1995). Symmetric functions and Hall polynomials (2nd ed.). Oxford University Press. ChapterV, (2.9), p.298; (3.4), p.299; (4.4)--(4.5), p.301, including its convergence statement. \url{https://math.berkeley.edu/~corteel/MATH249/macdonald.pdf}. The formula was read directly; the matrix-disintegration normalization and the application at the actual rank are given in full above.
\item
  Liu, Y. (2021). Theta correspondence for almost unramified representations of unitary groups. arXiv:2105.14223v1. Notation5.5 and equation(5.1), p.15; Proposition5.6, p.16, sign \(+\); Lemma6.1, p.20; Theorem6.2, pp.21--22. These are the actual lattice-vector and normalized local integral statements used in section4.
\item
  Gan, W. T., Qiu, Y., \& Takeda, S. (2014). The regularized Siegel--Weil formula (the second term identity) and the Rallis inner product formula. arXiv:1207.4709v3, sections11.3--11.8, especially the doubled evaluation on p.52, equation(11.3), Theorem11.4 and Proposition11.6.
\item
  Mínguez, A. (2008). Correspondance de Howe explicite: paires duales de type II. Annales Scientifiques de l'École Normale Supérieure, 41(5), 717--741. Theorem1, p.718, including the multiplicity-one/reciprocal uniqueness and all-residue-characteristic assertion.
\item
  Actual host and phases: AM-DO-020/022; exact source-to-doubling convention: AM-DO-028. Actual nonzero chain and Satake blocks: ALP-LIFT-14. Full finite Heisenberg basis: ALP-LIFT-15. First input and torsion-free principal-level proof: AM-R-012. These retain their separate check histories.
\end{itemize}

\clearpage

\section{AM-DO-030: auxiliary finite kernels and simultaneous selected jets in the actual permitted half-orbit}\label{am-do-030-auxiliary-finite-kernels-and-simultaneous-selected-jets-in-the-actual-permitted-half-orbit}

Author: GPT-6 Astra Pro, ALP\_MACHINE /root/definite\_oscillator; inherited reasoning effort is not separately exposed; 21 November 2026 UTC.

The carried inputs are DP59/DP62, AM-K018/022, AM-DO-022/026/029, and AM-R-012. This finding keeps a \textbf{fixed nonzero auxiliary finite Heisenberg operator}, proves simultaneous integral jet interpolation on the actual DP59 half-orbit, and selects one permitted translate that remains nonzero after that operator. Its finite norm, coefficient valuation and an explicit sufficient level are retained. A prescribed local Q-factor at29 is not absorbed into an operator supported away from29.

\subsection{1. Actual variety, orbit and auxiliary vector}\label{actual-variety-orbit-and-auxiliary-vector}

For the even-rank positive edge \(U(a)\to U(b)\), use

\[
 X=A^{ab}=J^{2ab},\qquad g=\dim X=6ab,\qquad
 \mathcal L=L_{2\infty}^{\boxtimes 2ab}.
 \tag{DO30.1}
\]

Its ordinary product CM action, principal trace-one line, common vacuum period \(C^{2ab}\), actual invariant differential coordinates, and finite vacuum \(\phi_0\) are the ones already checked in the carried inputs. In particular the finite vacuum has norm1; its archimedean Gaussian norm is separate. No conjugated \(J\)-factor is substituted for an ordinary factor.

Let

\[
 \alpha=(1+2\sqrt{-7})^e,\qquad M=29^e,\qquad
 \alpha\bar\alpha=M,\qquad e\ge1.
 \tag{DO30.2}
\]

The actual allowed section lattice is

\[
 \mathcal W_M=[\alpha]^*H^0(X,\mathcal L^M)
          \subset H^0(X,\mathcal L^{M^2}).
 \tag{DO30.3}
\]

It has the \textbf{actual} integral DP59 product basis \(\mathcal S_\lambda\), indexed by the prescribed quotient of \(X[\alpha^2]\). Its rescaled parameters are

\[
 x_\lambda=M t'_\lambda\in X[\alpha].
 \tag{DO30.4}
\]

The corresponding finite vectors are \(\phi_\lambda=\epsilon_\lambda W_{x_\lambda}\phi_0\), with their specified finite-order scalars \(\epsilon_\lambda\). These scalars are units at11 and have complex absolute value1. AM-K022 identifies this basis and its exact level-normalized reading; it is not replaced by all sections of \(\mathcal L^{M^2}\).

Fix an odd integer \(d\) with \(\gcd(d,11\cdot29)=1\). Let \(\mathcal Q\) be finite places supported on the rational primes dividing \(d\). Choose finite Heisenberg parameters \(y\), supported on \(\mathcal Q\), with \(dy\in\widehat\Lambda\), where \(\Lambda\) is the actual perfect trace lattice. Consider

\[
 \mathcal O=\sum_y b_yW_y,\qquad
 \phi_{\rm aux}=\mathcal O\phi_0\ne0.
 \tag{DO30.5}
\]

The \(b_y\) are algebraic, and the operator is fixed before choosing the29 orbit. This includes the finite incidence and other auxiliary-place combinations in AM-DO-029 whenever all their parameters are supported away from29 and11. It does not assume that a general such combination is an idempotent projector.

Combine terms that differ by an integral auxiliary lattice vector, using the actual Weyl phase and lattice semicharacter. On the vacuum, and on the whole29-only span, this produces

\[
 \mathcal O=\sum_{j=1}^{s} b_jW_{y_j},\qquad b_j\ne0,
 \tag{DO30.6}
\]

where the \(y_j\) have distinct classes modulo \(\widehat\Lambda\). In (DO30.6) equality is an equality of operators \textbf{on that span}; it is not asserted on every Schwartz vector. Indeed an integral auxiliary Heisenberg operator acts on the fixed auxiliary vacuum by its actual scalar, and commutes with every29-local operator.

The classes define distinct points \(t_j\in X[d]\). The correspondence follows by rational approximation modulo the integral lattice, or directly from the finite quotient \(d^{-1}\Lambda/\Lambda\). Distinct classes remain distinct after reduction at11, since \(X[d]\) is finite étale on the good11-adic model and becomes constant after an unramified extension. We may further extend the coefficient DVR to contain the \(b_j\); this extension need not be unramified. The valuation is always normalized by \(v_{11}(11)=1\).

The finite vectors \(W_{y_j}\phi_0\) are orthonormal. This is the selfdual-lattice character argument used in AM-DO-017/029: a nonintegral difference is detected by a lattice Heisenberg eigencharacter. Consequently

\[
 S_{\rm aux}^2:=\|\phi_{\rm aux}\|^2
                    =\sum_{j=1}^{s}|b_j|^2>0.
 \tag{DO30.7}
\]

Since the auxiliary and29 factors are disjoint,

\[
 \|\mathcal O\phi_\lambda\|^2=S_{\rm aux}^2
 \quad\text{for every allowed }\lambda.
 \tag{DO30.8}
\]

More generally \(\mathcal O\) restricts to a scaled isometry from the29 span, with scale \(S_{\rm aux}\). This follows by writing that span as its variable29 factor tensored with the fixed vacuum elsewhere. There is no interpolation-coefficient norm in (DO30.8).

\subsection{2. A common torsion representative and its exact phase}\label{a-common-torsion-representative-and-its-exact-phase}

A rational representative of an auxiliary finite parameter need not act by a fixed scalar correction on every29-translated vacuum. We choose a representative that does.

Because \(M\) is invertible modulo \(d\), put

\[
 z_j=[M]^{-1}t_j\in X[d].
 \tag{DO30.9}
\]

Choose a lift \(z_j\in d^{-1}\Lambda\), and put \(r_j=M z_j\). Then \(r_j\) represents the finite class \(y_j\), and its component at29 lies in \(M\widehat\Lambda_{29}\). Thus

\[
 \delta_j=y_j-r_j\in\widehat\Lambda,\qquad
 (\delta_j)_{29}\in M\widehat\Lambda_{29}.
 \tag{DO30.10}
\]

Every allowed \(x_\lambda\) is supported at29 with denominator dividing \(M\). Perfectness of the integral trace pairing gives \(\psi_f(B(\delta_j,x_\lambda))=1\). At all other places the second parameter is0. Therefore \(W_{\delta_j}\) acts on the whole permitted orbit by the same scalar as on \(\phi_0\). Applying the Weyl law gives the explicit root of unity

\[
 \varepsilon_{j,M}
 =\psi_f(-B(r_j,\delta_j)/2)\alpha_\Lambda(\delta_j).
 \tag{DO30.11}
\]

In the convention of AM-DO-026, the reading identity is

\[
 W_{y_j}\phi_\lambda
       \ \longleftrightarrow\
 \varepsilon_{j,M}\,U_{r_j}^{H}F_\lambda.
 \tag{DO30.12}
\]

All phases in (DO30.11) are part of the actual finite reading. They are independent of \(\lambda\), and are11-adic units. Here \(F_\lambda\) is the reduced principal-level function whose pullback by \([M]\) is \(\mathcal S_\lambda\).

The reduced translation formula gives the exact identity

\[
 (U_{r_j}^{H}F_\lambda)(Mu)
       =U_{z_j}^{M^2H}\bigl(F_\lambda(Mu)\bigr).
 \tag{DO30.13}
\]

Indeed both sides have argument \(Mu+Mz_j\), linear exponential term \(-\pi M^2H(u,z_j)\), and constant term \(-\pi M^2H(z_j,z_j)/2\). Thus the auxiliary operator, after the common pullback, evaluates the section at the \textbf{distinct} points \(z_j\), with the fixed phases (DO30.11). After \([\alpha]^*\), the points on the base of (DO30.3) are

\[
 q_j=[\alpha]z_j=[\bar\alpha]^{-1}t_j.
 \tag{DO30.14}
\]

They remain distinct, including modulo11, because \(\alpha\) is invertible on \(X[d]\).

\subsubsection{The translated line and the11-adic fibre scalar}\label{the-translated-line-and-the11-adic-fibre-scalar}

Usually \(z_j\notin K(\mathcal L^{M^2})\), since its order is prime to \(M\). Thus \(U_{z_j}^{M^2H}\) is a comparison with a \textbf{translated line}, differing by the corresponding torsion line in \(\operatorname{Pic}^0(X)\). It is not an automorphism in the theta group of \(\mathcal L^{M^2}\).

This causes no loss of an integral fibre unit at11. Pull back by \([d]\). The reduced identity now uses the point \(z_j/d\), whose order divides \(d^2\), on \(\mathcal L^{d^2M^2}\). Since \(d^2\mid d^2M^2\), that point belongs to its actual theta kernel. DP62's relative theta group and odd symmetric normalization apply: the finite-order lift is integral, and its multiplier in any primitive integral fibre frames is a unit. The normalized analytic reduced lift differs by the fixed finite-order scalar, also a unit. The pullback is finite étale at11, and its map on the formal neighbourhood of the zero section is an isomorphism after choosing the zero preimage. The fibre at zero of the pulled-back line is the original fibre. Consequently the constant in the given translated-line formal comparison is this same unit. This is a local fibre assertion, not a global trivialization of the nontrivial torsion line.

The argument asserts a unit \textbf{fibre comparison}, not a unit value of the translated theta section. It does not divide by a theta-null. AM-K018 gives the same normalized canonical gauge on the ordinary \(J\)-product; its positive character CM type is unchanged by a torsion translation. The higher coefficients of this gauge need not be integral ordinary power-series coefficients, and no such claim is used below.

\subsection{3. Simultaneous integral jets with an explicit threshold}\label{simultaneous-integral-jets-with-an-explicit-threshold}

Let \(R\) be a complete coefficient DVR over11 on which the points \(q_1,\ldots,q_s\), the abelian scheme and the line are defined. Their specializations are distinct. For \(D\ge0\), set

\[
 M_0(s,D)=3(s-1)(D+1)+3D+2=3s(D+1)-1.
 \tag{DO30.15}
\]

Then for every \(M\ge M_0(s,D)\), restriction is surjective:

\[
 H^0(X,\mathcal L^M)\longrightarrow
 \bigoplus_{j=1}^{s}
 H^0\!\left(X,\mathcal L^M\otimes
                     \mathcal O_X/I_{q_j}^{D+1}\right).
 \tag{DO30.16}
\]

This is a simultaneous-point extension of AM-K022's single-point argument. The geometric inputs remain the same carried facts: \(\mathcal L^3\) is very ample on an abelian variety, every \(\mathcal L^r\) with \(r\ge2\) is generated at each point, and positive powers have the integral \(H^0\) base change used in DP59/K022. Exact primary locators are given in §7.

\textbf{Proof.} Fix \(j\). Choose a section \(h_{j0}\) of \(\mathcal L^3\) whose value at \(q_j\) is a unit, and sections \(h_{j1},\ldots,h_{jg}\) vanishing there whose ratios \(h_{j\ell}/h_{j0}\) form integral formal parameters. Very ampleness on the special fibre gives the required independent cotangent vectors. Integral base change lifts them; subtracting their values times \(h_{j0}\) makes their vanishing exact. Only unit inverses are needed.

For each \(k\ne j\), choose \(k_{jk}\) in \(H^0(X,\mathcal L^3)\) with unit value at \(q_j\) and value0 at \(q_k\). Very ampleness separates the two special-fibre points. Lift such a section, and subtract its value at \(q_k\) times a section normalized to1 there. That subtracted scalar is in the maximal ideal, so the value at \(q_j\) remains a unit. Put

\[
 H_j=\prod_{k\ne j}k_{jk}^{D+1}.
 \tag{DO30.17}
\]

It is a section of \(\mathcal L^{3(s-1)(D+1)}\), is a unit in the fibre at \(q_j\), and vanishes to order at least \(D+1\) at every other \(q_k\). Put

\[
 r=M-3(s-1)(D+1)-3D\ge2.
 \tag{DO30.18}
\]

Choose \(\eta_j\in H^0(X,\mathcal L^r)\) with unit value at \(q_j\). The section \(H_j\eta_jh_{j0}^{D}\) is a local frame of \(\mathcal L^M\) through the required jet order at \(q_j\). Express any desired jet in this frame as \(\sum_{|\beta|\le D}c_\beta t_j^\beta\), with \(c_\beta\in R\). The global section

\[
 H_j\eta_j
   \sum_{|\beta|\le D}c_\beta
       h_{j0}^{D-|\beta|}
            \prod_{\ell=1}^{g}h_{j\ell}^{\beta_\ell}
 \tag{DO30.19}
\]

has that jet at \(q_j\) and zero \(D\)-jet at every other point. Summing these sections over \(j\) realizes all desired jets. This proves (DO30.16) over \(R\), without division by factorials or interpolation determinants of nonunit valuation. \(\square\)

Pulling back by the actual étale isogeny \([\alpha]\) and its DP59 rigidified line isomorphism gives

\[
 \mathcal W_M\twoheadrightarrow
   \bigoplus_{j=1}^{s}
       \operatorname{Jet}^{D}_{z_j}(\mathcal L^{M^2}).
 \tag{DO30.20}
\]

The formal tangent maps are invertible over \(R\): \(\alpha\bar\alpha=M\) and \(M\) is a unit at11. All line comparison constants are units. Thus this is surjectivity from the actual permitted lattice, not from a larger section space.

One explicit permitted level is

\[
 e(s,D)=\max\{1,\lceil\log_{29}(3s(D+1)-1)\rceil\},
 \quad M=29^{e(s,D)},\quad
 M_0(s,D)\le M<29M_0(s,D).
 \tag{DO30.21}
\]

The threshold recovers \(3D+2\) when \(s=1\).

\subsection{4. The selected functional and exact factorial valuation}\label{the-selected-functional-and-exact-factorial-valuation}

Let \(P=\sum_{|\beta|=D}c_\beta Z^\beta\) be a nonzero homogeneous polynomial in the actual integral CM invariant tangent basis. Define

\[
 B_P=\min_{\beta\in\operatorname{supp}P}
                              v_{11}(c_\beta\beta!),
 \qquad \beta!=\prod_\ell\beta_\ell!.
 \tag{DO30.22}
\]

For the selected primitive bideterminant \(P_\nu\) of AM-K022, its diagonal coefficient is1, so

\[
 B_P\le B_\nu=\sum_{\tau,i}v_{11}(\nu_{\tau,i}!).
 \tag{DO30.23}
\]

For \(S\in\mathcal W_M\), form the exact auxiliary functional

\[
 \ell_{\mathcal O,P,M}(S)
 =M^{-D}\left[P(\partial_u)
     \sum_{j=1}^{s}b_j\varepsilon_{j,M}
             U_{z_j}^{M^2H}F_S(u)\right]_{u=0},
 \tag{DO30.24}
\]

where the translated functions are read in their actual translated-line formal charts, as in §2. All constants are the fixed ones in the reading map and the primitive fibre comparisons, not section-dependent gauge choices.

If \(M\) obeys (DO30.21), at least one \textbf{individual} permitted basis section satisfies

\[
 \boxed{\ell_{\mathcal O,P,M}(\mathcal S_\lambda)\ne0,
 \qquad
 v_{11}\!\left(\ell_{\mathcal O,P,M}(\mathcal S_\lambda)\right)
       \le \min_j v_{11}(b_j)+B_P.}
 \tag{DO30.25}
\]

\textbf{Proof.} Choose \(j\) with \(v_{11}(b_j)\) minimal and a monomial \(\beta\) attaining \(B_P\). By (DO30.20), prescribe zero \(D\)-jets at all \(z_k\), \(k\ne j\). At \(z_j\), choose integral formal parameters with first derivatives the given invariant basis. Prescribe the algebraic coefficient jet \(M^D t^\beta\) in a primitive line frame. The scalar \(M^D\) is an11-adic unit.

The translated source/Chern gauge has a unit constant \(u_{j,M}\). Its higher coefficients are immaterial: the chosen algebraic jet has no term of degree below \(D\). Consequently its degree-\(D\) reduced term is precisely \(u_{j,M}M^D u^\beta\). All other points make zero contribution. Thus for this integral section \(S\),

\[
 \ell_{\mathcal O,P,M}(S)
    =b_j\varepsilon_{j,M}u_{j,M}c_\beta\beta!.
 \tag{DO30.26}
\]

The two displayed comparison scalars are units, so this value has valuation \(\min_jv_{11}(b_j)+B_P\) and is nonzero. Expand \(S\) in the \textbf{integral} basis \(\{\mathcal S_\lambda\}\) of \(\mathcal W_M\). If every basis value had strictly larger valuation, their integral linear combination could not equal (DO30.26). This proves (DO30.25). No archimedean norm of that interpolating linear combination is used; the final selected input is one basis vector. \(\square\)

This proof also explains why the factorial is unavoidable: an ordinary coefficient prescribed to1 has derivative \(\beta!\). It does not prove a unit derivative at arbitrary degree. It also does not require integral coefficients for the full canonical reduced gauge.

The all-polynomial reading AM-DO-026 and the exact level cancellation (DO30.13) give

\[
 \vartheta_{r_P\mathcal O\phi_\lambda}(1,1)
   =C^{2ab}\,\ell_{\mathcal O,P,M}(\mathcal S_\lambda).
 \tag{DO30.27}
\]

This is a nonzero \textbf{actual} selected kernel value, with the one common period retained. The polynomial still lies in the actual selected compact type, since every auxiliary operator is finite and commutes with the archimedean creation action.

\subsection{5. Exact cost and one explicit source compact level}\label{exact-cost-and-one-explicit-source-compact-level}

The chosen finite Schwartz vector is \(\mathcal O\phi_\lambda\). Its coefficients in its reduced auxiliary orbit are the \(b_j\) times unit phases; its exact norm is \(S_{\rm aux}\) from (DO30.7). There is no dependence on \(M\) in that finite norm. The full Schwartz norm is

\[
 \|r_P\phi_{0,\infty}\|\,S_{\rm aux};
 \tag{DO30.28}
\]

the period, highest-vector hook factor and Gaussian covolume remain those of AM-R-011/012. Normalizing the finite vector to norm1 would divide its value by \(S_{\rm aux}\); that extra scalar is not silently made integral.

For the actual \(T\)-incidence product of AM-DO-029, with \(a\)-dimensional source residue space at each selected \(T\)-place, the reduced frequency set and coefficients are

\[
 s=\prod_{P\in v_T}q_P^a,\qquad
 |b_j|=s^{-1},\qquad
 S_{\rm aux}^2=s^{-1}.
 \tag{DO30.29}
\]

The distinct frequencies are the rank-one matrices \(v_P^t\eta_P\). If no \(q_P\) is11, then \(v_{11}(b_j)=0\); thus (DO30.25) retains the same factorial bound \(B_P\). The sufficient level is now \(M\ge3s(D+1)-1\). For one rational prime \(T\) split into six selected half places, \(s=T^{6a}\), not \(T^{6ab}\). Other prescribed \(G/Q\) factors are included by their exact reduced coefficients from (DO29.11), with their actual \(s\), norm and minimum valuation.

Here is an explicit finite stabilizer for the new first lift. Suppose the auxiliary places are odd relative split places, the original compact vacuum is fixed there, and each \(y_j\) has denominator at most \(\varpi_v^{c_v}\) at \(v\), \(c_v\ge1\). Take principal congruence

\[
 K_{{\rm aux},a,v}
   =1+\varpi_v^{2c_v}M_a(\mathcal O_v)
       \ \subset\mathrm{GL}_a(\mathcal O_v),
 \tag{DO30.30}
\]

and similarly on the target. Intersect with the kernel of any separately specified vacuum determinant character if it is ramified there. For the R012 character, which has compact conductor only at29, that extra intersection is unnecessary.

To check fixation, a matrix in (DO30.30) changes \(y_j\) by \(\delta\in\varpi_v^{c_v}\widehat\Lambda_v\). Both \(B(y_j,\delta)/2\) and the trace-one lattice exponent for \(\delta\) are integral at this odd place. Their character values are1. Thus the Weyl covariance formula fixes each \(W_{y_j}\phi_0\), with its actual central phase, not merely its line. At29 use the existing principal congruence of exponent \(e\); at all other places use the original vacuum compact subgroup.

In the unramified split situation its source index is explicitly

\[
 I_{{\rm aux},a}
    =\prod_{v\in\mathcal Q}
      q_v^{a^2(2c_v-1)}
                    |\mathrm{GL}_a(\kappa_v)|.
 \tag{DO30.31}
\]

The exponent comes from the \(2c_v-1\) successive additive matrix kernels. It is the index of the sufficient stabilizer, not an assertion that it is minimal. There can be11-factors in \(|\mathrm{GL}_a(\kappa_v)|\), even when the residue characteristic is not11.

Substitution into the checked class-supported construction R012 gives source identity-class mass \[
 \frac1{I_a(M)I_{{\rm aux},a}},
\] primitive compact-matrix-coefficient input norm squared \[
 \frac1{I_a(M)I_{{\rm aux},a}d_\rho},
\] and actual lifted value

\[
 \frac{C^{2ab}\,
       \ell_{\mathcal O,P,M}(\mathcal S_\lambda)}
          {I_a(M)I_{{\rm aux},a}d_\rho}.
 \tag{DO30.32}
\]

This uses the same trivial rational stabilizer supplied by the29 congruence, and the same compact Schur orthogonality. It does not change the selected real weights or the common splitting. These are source mass factors; they are distinct from the finite Schwartz norm (DO30.8).

For a fixed auxiliary operator, \(s,d\), \(I_{{\rm aux},a}\) and \(S_{\rm aux}\) are fixed and \(\log M=O(\log(D+2))\). If they vary with \(D\), the actual bound is

\[
 \log M<\log29+\log(3s(D)(D+1)-1),
 \tag{DO30.33}
\]

with the displayed auxiliary index and coefficient data retained. In particular one must not call the overhead logarithmic in \(D\) while silently allowing an exponentially growing \(s(D)\).

\subsection{6. Exact remaining raw-array equation}\label{exact-remaining-raw-array-equation}

The theorem now proves nonvanishing after every fixed auxiliary operator \(\mathcal O\) supported away from29 and11 for which \(\mathcal O\phi_0\ne0\). It does so on the actual permitted29 half-orbit, with a single selected basis vector. It is stronger than nonvanishing of some unrestricted finite Schwartz input.

It does \textbf{not} identify that new input with a fixed complete xiAt. Its finite half-function has the form

\[
 \Phi_{\rm aux}(t_{\rm aux})
                \Phi_{\lambda,29}(t_{29}),
 \tag{DO30.34}
\]

where \(\Phi_{\lambda,29}\) is the actual29-translated vacuum. If the prescribed array also contains a fixed local Q-factor at29, the required equation is still \[
 \Phi_{\lambda,29}
       \stackrel{?}{=}\Phi_{Q,29},
\] or membership of the corresponding family in the actual source's allowed local-index span, with all its frame and character rules. Neither K022's interpolation nor (DO30.25) proves that equation.

Applying a nonzero operator supported at29 cannot simply be included in this proof: the disjoint-factor tensor argument and the invertibility of \(\alpha\) on its torsion points are then unavailable. The fixed Q29 image, and the resulting prescribed row/operator identity, therefore remain precise open gates. Later restricted image saturation, Kval, terminal mass and the MainX inequalities are also not consequences of (DO30.32).

\subsubsection{6.1 Literal reference-character and prime-set audit}\label{literal-reference-character-and-prime-set-audit}

The two source symbols called mu0 have different types and must be distinguished. In Development/AbelianVarietyTheory/Definitions/AutomorphicMembership/ \nolinkurl{TrueColumnChoicesHeckeCharReferenceEqOne.lean:}

\begin{longtable}[]{@{}
  >{\raggedright\arraybackslash}p{(\columnwidth - 2\tabcolsep) * \real{0.5000}}
  >{\raggedright\arraybackslash}p{(\columnwidth - 2\tabcolsep) * \real{0.5000}}@{}}
\toprule\noalign{}
\begin{minipage}[b]{\linewidth}\raggedright
Locator
\end{minipage} & \begin{minipage}[b]{\linewidth}\raggedright
Actual object or condition
\end{minipage} \\
\midrule\noalign{}
\endhead
\bottomrule\noalign{}
\endlastfoot
175--190 & the index \(I.\mathrm{mu0}\), with \(\mathrm{tup}(I.\mathrm{mu0})=1\) \\
202--205 & \(I.\mathrm{glob.mu0}=I.C.\mathrm{mu0}\), the family of reference Hecke characters \\
211--215 & \(I.Q(a)=\mathrm{QE}(\mathrm{sel},\mathrm{tup}(a))\) \\
222--226 & the actual tying of the reference characters to Chars \\
273--280 & the proved identity \(I.Q(I.\mathrm{mu0})=\varnothing\) \\
295--299 & the reference twisting character \(I.\xi(I.\mathrm{mu0})=1\) \\
\end{longtable}

Thus the Q-set records the \textbf{finite-order twisting tuple}; it is not the conductor of the reference splitting character. In particular a Q29 factor is not forced merely by saying that a chosen reference character has conductor29.

However the reference character cannot have arbitrary ramification. In \nolinkurl{Formalism/LocalToGlobal/Columns/AbelianHostFieldConditions.lean,} IsMu0, lines314--323, requires it to be unramified away from the two G-primes, tame everywhere, and to have the specified residue character at each G-place. Chars.Tied, lines325--334, imposes these conditions on every reference character. Chars.Ok.chi\_om, lines302--305, prescribes \[
 \chi_{j,P}=\omega_G^{\,2(j+1)}
\] for one faithful residue character at each rational G-prime, and chi\_inj requires the characters to be distinct as \(j\) varies. Consequently the R012 reference character, \textbf{when its29 compact character is nontrivial}, is incompatible with keeping both G-primes away from29. Replacing it by a source-tied reference character is an additional explicit construction, with these residue conditions retained. If29 is one of the G-primes, its G-factor rather than a Q-factor remains in the raw array. If the particular29 compact character is trivial, this ramification obstruction does not apply; the other source tying conditions still do. Moreover, the later good-data clause g\_big in that same file, line488, requires \(2fk<g_i-1\). On the actual mixed host \(f=6\), and the allowed ranks have \(k\ge6\), so \(g_i>73\). Therefore29 cannot be one of these G-primes. For a reference character whose29 compact component is nontrivial, relabelling29 as G does not repair the full source compatibility. This is a further literal prime condition; the analytic characters and theta chain do not claim to satisfy it.

Even when the reference tuple is used and every G/T auxiliary prime avoids29, the new input in (DO30.27) is not automatically the same raw datum. The literal TwDat.Is.idem in \nolinkurl{Definitions/LatticeFrames/LevChiEqOne.lean,} line62, requires the predicate \[
 \operatorname{IsIdemSum}
   \bigl(A.e,\operatorname{Sall}(I.\mathrm{sel},I.\mathrm{tup}(a),u)\bigr).
\] The definition in \nolinkurl{Formalism/LocalToGlobal/IndefiniteReal/DatumVec.lean,} lines78--80, is exactly \[
 \bigl(\exists (e_P)_P,\;
   (\forall P\in S,\operatorname{IsIdemAt}(e_P,P))
   \ \land\ e=\sum_{P\in S}e_P\bigr)
       \ \land\ e\,\overline e=0.
 \tag{DO30.34a}
\] There is no function named IdemSum in this assertion. The support conclusion follows from the literal predicate, not from its name. In \nolinkurl{Formalism/LocalToGlobal/HeckeThetaElements/IsHeckeElt.lean,} lines68--70, IsIdemAt is IsSplitIdem together with \[
 x\in P\quad\Longleftrightarrow\quad
       (e_P/x)_f\notin\widehat{\mathcal O}_E
 \qquad(0\ne x\in\mathcal O_E).
 \tag{DO30.34b}
\] Here IsSplitIdem includes \(e_P^2=e_P\), \(e_P\overline{e_P}=0\), and \(e_P+\overline{e_P}\) equal to one finite base-place idempotent; see \nolinkurl{BesselHermitianTheta/TiedWeil.lean,} lines68--72. Its infinite component is zero, as also proved in EqZeroOrEqOfPlaceIdem.lean, lines134--140.

For completeness, the finite support implication has a short direct proof. Identify the finite tensor adelic algebra with the product of the completions of \(E\). Idempotence says that each component of \(e_P\) is0 or1. For any prime \(Q\ne P\), choose a nonzero \(x\in Q\setminus P\), possible because distinct maximal ideals are comaximal. If the \(Q\)-component of \(e_P\) were1, then the \(Q\)-component of \(e_P/x\) would be nonintegral. This would contradict (DO30.34b), since \(x\notin P\). Thus every such component is0. The \(P\)-component must be1: otherwise all finite components vanish, and any nonzero \(x\in P\) contradicts (DO30.34b). Consequently \(e_P\) is exactly the finite idempotent supported at \(P\). Formula (DO30.34a) therefore gives a zero29 component whenever Sall has no place above29. The normalized place model then sends its half-indicator to the \textbf{fixed} vacuum. Thus the fixed array gives \(\mathcal O\phi_0\), while (DO30.27) gives \(\mathcal O W_{x_\lambda}\phi_0\). The extra29 factor remains an actual change of finite input. One cannot add a29 component to \(e\) and still claim the same literal idem clause.

There is also a precise obstruction to absorbing a nonzero29 position shift by the available linear frame/evaluation actions. At a relative split29 place, the positive Hom polarization identifies every unitary-pair action with \(t\mapsto h^{-1}tg\), times its nonzero determinant-root and Jacobian scalar. Hence every linear translate of the vacuum has position function a nonzero scalar times the indicator of a lattice \textbf{containing0}. A pure position translate by \(x\notin L\) has function \(1_L(t+x)\), whose support does \textbf{not} contain0. They cannot be equal. This proves the claimed obstruction for all such linear source/target evaluations at29, not only integral frames. It does not rule out equality of an isolated theta value, a finite linear combination of group translates, or a genuinely enlarged datum with new affine parameters. It also does not assert that the selected tuple in (DO30.25) is nontrivial: if the zero tuple already works, it is an actual raw witness. The theorem does not force that special outcome.

Accordingly the empty-Q choice removes one possible misidentification of the ramification set, but does not eliminate the raw-array matching equation. The finite linear source model fixes the origin; its frame changes cannot be used as undeclared affine Heisenberg translations.

\subsection{7. Exact sources and check scope}\label{exact-sources-and-check-scope}

\begin{itemize}
\tightlist
\item
  AM-K022, selected\_schur\_permitted\_jets.md, J1--J5: actual permitted lattice and its integral basis; J6--J17: the carried single-point interpolation and highest-symbol factorial argument.
\item
  DP59, research/theta\_translate\_span\_DP59.md, §§3--8; DP62, research/theta\_translate\_norm\_integrality\_DP62.md, TN2--TN5 and §§2--4: actual rigidified multiplication, odd symmetric theta lifts and integral finite-order normalization. These supplied results are used as carried inputs, not reproved in full.
\item
  Milne, J. S. (2008). \emph{Abelian varieties} (Version2.0), I, §6, ``Ample divisors,'' printed p.30, paragraph before Example6.5. \url{https://www.jmilne.org/math/CourseNotes/AVc.pdf}. The arbitrary-characteristic very-ampleness input is the one already read in AM-K022.
\item
  Mumford, D. (1966). On the equations defining abelian varieties. I. \emph{Inventiones Mathematicae, 1}, 287--354, §1(II), p.289. \url{https://pazuki.perso.math.cnrs.fr/index_fichiers/Mumford66.pdf}. The positive-power cohomology/base-change input is carried from DP59/K022.
\item
  Polishchuk, A. \emph{Analogue of Weil representation for abelian schemes}, author preprint arXiv:alg-geom/9712021, §1, pp.4--5, definition preceding Lemma1.1, Proposition1.2 and Corollary1.3. \url{https://arxiv.org/pdf/alg-geom/9712021}. The precise relative theta-group input is the one already checked in DP62.
\item
  AM-K018, jacobian\_source\_gauge\_match.md, §§3--5: exact canonical reduced gauge on the ordinary product; no integral reduced power-series gauge is assumed.
\item
  AM-DO-026, (DO26.2a)--(DO26.4): actual adelic-to-Mumford phase and all-polynomial reading. AM-DO-029, (DO29.11)--(DO29.23): exact prescribed finite coefficients and norms.
\item
  AM-R-012, §§2--3: the actual compact selected first input, stabilizer mass, and Schur-orthogonality factor.
\end{itemize}

Proof-check events are recorded only when actually performed in the agent TSV. No Lean compilation or complete source MainX instance is asserted.

\clearpage

\section{AM-K026 --- a full-Hyp3 family, exact selected norm correction, and the remaining cost interval}\label{am-k026-a-full-hyp3-family-exact-selected-norm-correction-and-the-remaining-cost-interval}

Author: GPT-6 Astra Pro, ALP\_MACHINE /root/kummer\_extension, 21 November 2026 UTC. Inherited reasoning effort is not separately exposed.

This result goes beyond the HC-only unit examples. It supplies an explicit family satisfying every numerical clause of the source's stronger analytic Hyp3, computes the adjacent-gap content bound and the exact hook/mass factors, and proves that the K024 upper slope survives the selected target metric normalization. It still does not identify the second restricted class-image lattice or the full D7/SizeData₂ arithmetic adapter.

The classical/source inputs are retained as theorems. In particular the noncompact-gap spacing lemma is used at its exact recovered interface, and the hook norms are the parent's AM-R-011.

\subsection{1. A concrete full-Hyp3 baseline}\label{a-concrete-full-hyp3-baseline}

Fix an admissible source rank record with k,m even, t odd, r=k+m−1, d=t+1, and \(1801\le r+t\). Put \[
 S=r+mt,\qquad Q=S^3+2,\qquad G=2Q+d=2S^3+t+5.
\tag{HC1}
\] At every CM embedding define \[
 u_i=k-\tfrac12-i,\quad v_i=k-1-i,\quad
 b_j=jG,\quad e_j=(j+\tfrac12)G.
\tag{HC2}
\] The integer G is even. Set eta to the required \(-\operatorname{infTy}(\eta_T)\).

\subsubsection{Proposition HC3}\label{proposition-hc3}

These parameters satisfy the literal HCParam record and the stronger Hyp3 of CayleySignature.lean, lines 110--118, including the 1801 condition, Centred and NCGaps. The actual selected compact partitions are \[
 \nu_j=
 Q\,((2j+1)^k,(2j)^d,(2j-2)^d,\ldots,2^d)
\tag{HC3}
\] at every one of the six labels.

\textbf{Proof.} All u values are half-integral and v,b,e are integral. The two u,v interlacing gaps equal 1/2; the last u equals 1/2\textgreater−b₀=0. The b,e entries alternate strictly increasingly, and both HC gaps equal G/2, which exceeds \((t+3)/2\). Thus both the HC record and TPar.Interlace and TParity hold.

For Centred, u and v are positive, and every e is at least G/2. Thus the required u+e and v+e absolute values exceed \((t+3)/2\). For any i,j, \(|e_j-b_i|=|j+1/2-i|G\ge G/2\). If j′≠j and \(x\in\operatorname{block}(-e_{j'},t)\), then \[
 |x+e_j|\ge |j-j'|G-(t-1)/2
                   \ge G-(t-1)/2>(t+3)/2.
\] This checks every Centred clause in DataTheta.lean, lines 550--554.

For NCGaps apply the accepted source theorem noncompact\_gaps\_of\_spacing in \nolinkurl{ExistsTowerOfLocalSetting.lean,} lines 99--103. Its required lower spacing is \(2(r+mt)^3+t+3\), and G exceeds that by 2. Its other hypotheses are \(-b_0<v_i\), pairwise b spacing at least G, and \(e_j<b_{j+1}\); all were just checked. This gives the exact NCGaps predicate for every j≤m. The 1801 and eta clauses hold by construction.

Finally \(w_{0,i}=u_i-\rho_k(i)=k/2\). Hence DO25.5 gives \(\nu_0=e_0-d/2=Q\) in every row. The exact increment is \(\Delta_j=e_{j+1}-e_j-d=G-d=2Q\). Recurrence DO25.11 gives (HC3). ∎

This supplies the source's full analytic parameter conditions. It does not assert that these numbers satisfy every separately prescribed global upper window, or that an unchanged old induced-CM D7 record has been instantiated on the ordinary E′ host.

\subsection{2. Compulsory first divisibility and an explicit upper bound}\label{compulsory-first-divisibility-and-an-explicit-upper-bound}

Write \(a=k+jd,\ b=a+d\), and \[
 R_j=kj+d\,j(j+1)/2.
\] The total degree and the anti-diagonal upper bound from AM-K025 are \[
 D_j=6Q\{k(2j+1)+d\,j(j+1)\},
\tag{HC4}
\] \[
 U_j=6k\,v_{11}(Q!)+6R_j\,v_{11}((2Q)!).
\tag{HC5}
\] The nonzero prefix differences are Q at k and 2Q at k+d,\ldots,k+jd. In particular \[
 U_j=\frac{D_j}{10}
       -\frac6{10}\{k\,s_{11}(Q)+R_j\,s_{11}(2Q)\}.
\tag{HC6}
\]

For this family every stage satisfies the universal mod-11 cutoff: \[
                              D_j>3140ab.
\tag{HC7}
\] Indeed every part is at least Q, so \(D_j\ge6aQ\). The target rank b is at most \(N_m=k+m(t+1)=S+1\). Since S≥r+t≥1801, the elementary inequality \(6(S^3+2)>3140(S+1)\) proves (HC7).

There is also an exact minimum, rather than just this coarse comparison: \[
 \frac{D_j}{ab}
   =\frac{6Q}{d}\left(1-\frac{k^2}{ab}\right),\qquad
 \min_{0\le j<m}\frac{D_j}{ab}
                         =\frac{6Q}{k+d}.
\tag{HC7a}
\] To verify the identity, expand \(ab=k^2+kd(2j+1)+d^2j(j+1)\). The product ab strictly increases with j, so the displayed ratio strictly increases. Its minimum is at the first edge.

Consequently the exact selected highest-row image ideal satisfies \[
                      \boxed{1\le v_{11}(I_{\nu_j})\le U_j.}
\tag{HC8}
\] The lower bound is only one factor of 11. Neither the large degree nor the cutoff justifies iterating it after dividing the module by 11. The width of (HC8) is an unresolved quantitative gap, not an exact gain calculation.

Increasing the u₀ entries as in AM-K024 preserves this same full Hyp3 and only increases D at the fixed edge. Thus every member of that family remains above the cutoff, while K024 still gives selected coefficients with upper asymptotic rate at most 1/1220. These two facts are compatible: compulsory first divisibility does not determine an asymptotic slope.

\subsection{3. A fully explicit two-edge instance}\label{a-fully-explicit-two-edge-instance}

Take \[
 k=6,\quad m=2,\quad t=1795,\quad d=1796,\quad r=7,
\quad \mathrm{ord}=3.
\tag{HC9}
\] All the source rank inequalities hold, and r+t=1802≥1801. The values are \[
 S=3597,\quad Q=46\,539\,457\,175,\quad
 G=93\,078\,916\,146.
\tag{HC10}
\] Thus \[
 (b_0,b_1)=(0,93\,078\,916\,146),\qquad
 (e_0,e_1)=(46\,539\,458\,073,139\,618\,374\,219).
\] The u entries are \(11/2,9/2,7/2,5/2,3/2,1/2\); the v entries are 5,4,3,2,1.

The exact numerical consequences are:

\begin{longtable}[]{@{}lrr@{}}
\toprule\noalign{}
Quantity & Stage 0 & Stage 1 \\
\midrule\noalign{}
\endhead
\bottomrule\noalign{}
\endlastfoot
Compact edge a→b & 6→1802 & 1802→3598 \\
Selected partition at each label & \(Q^6\) & \((3Q)^6,(2Q)^{1796}\) \\
Total six-label degree D & 1,675,420,458,300 & 1,008,044,642,410,500 \\
D/(3140ab), approximately & 49,350.1104664 & 49,514.7022745 \\
Proven lower bound for \(v_{11}(I_\nu)\) & 1 & 1 \\
Adjacent-gap upper bound U & 167,542,045,704 & 100,804,464,197,676 \\
Least e with \(29^e\ge3D+2\) & 9 & 11 \\
Permitted level M & 14,507,145,975,869 & 12,200,509,765,705,829 \\
\end{longtable}

The certified factorial values are \[
 v_{11}(Q!)=4\,653\,945\,714,\qquad
 v_{11}((2Q)!)=9\,307\,891\,431.
\tag{HC11}
\] Every integer and every source rank, Interlace, Centred and NCGaps check in this example is reproduced by \nolinkurl{hyp3_selected_cost_certificate.py.} The program compares all opposite-sign discrete parameters for j=0,1,2; the minimum nonvacuous gap is G, against the required \(S^3=46\,539\,457\,173\). Common eta shifts cancel from these differences.

\subsection{4. Exact hook and class-mass factors for that instance}\label{exact-hook-and-class-mass-factors-for-that-instance}

Let Hν be the hook product of one label. The source coefficient partition is \(\eta_0=0\) and \(\eta_j=(\nu_{j-1},0^d)\) for j≥1. Thus the actual uniform insertion exponent is Q at stage 0 and 2Q thereafter. AM-R-011 gives the exact recurrence \[
 H_{\nu_j}
   =H_{\eta_j}
       \prod_{i=1}^{a}
        \frac{(\eta_{j,i}+q_j+a-i)!}{(\eta_{j,i}+a-i)!},
 \quad q_0=Q,\quad q_j=2Q\ (j\ge1).
\tag{HC12}
\] The ratio is the actual coefficient norm insertion. Its product over stages telescopes to \(H_{\nu_{m-1}}/H_{\eta_0}\); these ratios must not also be counted as independent uncancelled row-metric losses.

For the example:

\begin{longtable}[]{@{}
  >{\raggedright\arraybackslash}p{(\columnwidth - 4\tabcolsep) * \real{0.2727}}
  >{\raggedleft\arraybackslash}p{(\columnwidth - 4\tabcolsep) * \real{0.3636}}
  >{\raggedleft\arraybackslash}p{(\columnwidth - 4\tabcolsep) * \real{0.3636}}@{}}
\toprule\noalign{}
\begin{minipage}[b]{\linewidth}\raggedright
Exact 11-adic factor
\end{minipage} & \begin{minipage}[b]{\linewidth}\raggedleft
Stage 0
\end{minipage} & \begin{minipage}[b]{\linewidth}\raggedleft
Stage 1
\end{minipage} \\
\midrule\noalign{}
\endhead
\bottomrule\noalign{}
\endlastfoot
\(v_{11}(H_\nu)\), one label & 27,923,674,284 & 16,800,744,033,680 \\
\(v_{11}(H_\nu/H_\eta)\), one label & 27,923,674,284 & 16,772,820,359,396 \\
\(v_{11}(I_a(M))\), R012 class index & 0 & 1182 \\
\(v_{11}(\dim\rho)\), six-label source coefficient & 0 & 0 \\
\end{longtable}

Here the mass convention is the exact AM-R-012 convention \[
 I_a(M)=M^{6a^2}\prod_{\ell=1}^{a}(1-29^{-\ell})^6.
\tag{HC13}
\] At 11 its valuation is independent of the exponent e in M=29\^{}e, because M is an 11-unit. The source dimension is computed by the Weyl dimension formula for η; at stage 1 it is the six-label product for \((Q^6,0^{1796})\). This is an integer dimension, not a denominator inserted into a polynomial projection.

For scale, high-precision numerical evaluations give \[
 \log H_{\nu_0}\approx6.5798135173845434\cdot10^{12},\qquad
 \log H_{\nu_1}\approx4.0756480430066282\cdot10^{15}.
\tag{HC14}
\] These logarithms are numerical evaluations of the exact factorial products, not standalone interval certificates. All exact valuations are computed with integers.

The actual primitive highest creation norm remains, by AM-R-011, \[
 \log\|\phi_j^{\rm hw}\|_\infty
  =-\frac{ab}{2}\log\det(2Y)
       +3\log H_{\nu_j}
       +|\nu_j|\sum_{i=1}^3\log\frac{\pi}{2G_i}.
\tag{HC15}
\] The algebraic-to-analytic kernel period is still \(C^{2ab}\). Subtracting its logarithmic absolute value gives an additional explicit term \(-2ab\log|C|\); that period is not removed by (HC12). Source class extraction also retains (HC13) and its compact dimension exactly as in R012.

These formulas exhibit concrete archimedean and finite factors on a full-Hyp3 family. The interval (HC8) is too wide to identify the actual local gain, and no sum of the separate factors is labelled a total DP44 error.

\subsection{5. Exact metric correction for the scalable K024 selected vector}\label{exact-metric-correction-for-the-scalable-k024-selected-vector}

The target vector in K024 was shown to lie in the exact selected Cauchy summand, but it need not be its highest vector. Its normalization can be computed exactly, with no projector bound.

Fix one residual partition \(\lambda\) of length a. On the spare column let \[
       P_h=Z_{1,1}^{h}\,Q_\lambda(Z_{\text{columns }2,\ldots,a+1}).
\tag{HC16}
\] The variables in its two factors are disjoint. The Fischer monomials are orthogonal, so \[
                    \|P_h\|_{\rm F}^2=h!\,H_\lambda.
\tag{HC17}
\] Its selected source partition is \(\lambda+h e_1\). In the Cauchy realization with unit source highest vector and the target compact metric whose highest vector has norm one, AM-R-011 therefore gives the exact target squared norm \[
                  \theta_\lambda(h)
                  =\frac{h!H_\lambda}{H_{\lambda+h e_1}}.
\tag{HC18}
\] No independent hook penalty is added to (HC17): (HC18) is the comparison between the two ways of expressing that same norm.

Put \(L=\lambda_1+a-1\) and \(c_j=\lambda_1-\lambda_j+j-1\), j=2,\ldots,a. The hook factorial/Vandermonde formula from R011 gives \[
 \boxed{\theta_\lambda(h)=
  \frac1{\binom{h+L}{L}}
       \prod_{j=2}^{a}\frac{h+c_j}{c_j}.}
\tag{HC19}
\] The c\_j are distinct integers in \{1,\ldots,L\}. Cancelling them from the rising factorial gives the equivalent manifestly positive formula \[
 \theta_\lambda(h)
   =\prod_{s\in\{1,\ldots,L\}\setminus\{c_2,\ldots,c_a\}}
                         \frac{s}{h+s}.
\tag{HC20}
\] There are exactly λ₁ factors. In particular \[
        0<\theta_\lambda(h)\le1,\qquad
        |\log\theta_\lambda(h)|=O_\lambda(\log(h+1)).
\tag{HC21}
\]

The corresponding source compact dimension obeys \[
 \frac{\dim S_{\lambda+h e_1}(\mathbf C^a)}
      {\dim S_\lambda(\mathbf C^a)}
                         =\prod_{j=2}^{a}\frac{h+c_j}{c_j}.
\tag{HC22}
\] Removing a determinant twist does not change that dimension, so it is precisely the dimension ratio for the actual row coefficient η.

For the finite valuation, Kummer's carry formula, or directly Legendre's formula, gives \[
 v_{11}\binom{h+L}{L}\le\lfloor\log_{11}(h+L)\rfloor
 \quad(h+L\ge1).
\] Thus (HC19) implies \[
 v_{11}(\theta_\lambda(h))
 \ge-\lfloor\log_{11}(h+L)\rfloor-\sum_{j=2}^{a}v_{11}(c_j).
\tag{HC23}
\] The analogous upper bound and (HC22) give absolute valuation bounds of order log(h+1), with the residual rank and partition fixed.

Apply this at the four growing labels of K024, with \((h,A_r-h,k,B_r-k)\), and write Θ for the product of their θ factors. Scaling the target vector to norm one multiplies its polynomial and its kernel coefficient by Θ\^{}\{-1/2\}. This is an explicit algebraic scalar; it is not silently treated as an integral lattice unit. Combining (HC23) with K024 gives \[
 v_{11}(a_r/\sqrt{\Theta})
       \le\frac{11^r-1}{10}+c_0+O(\log T_r).
\tag{HC24}
\] The upper slope bound remains 1/1220. With the primitive polynomial kept instead, its exact actual creation norm is \[
\begin{split}
 \|\phi_{r,h,k}\|_\infty^2={}&
  \det(2Y)^{-ab}
  h!(A_r-h)!k!(B_r-k)!
       \prod_\tau H_{\nu_\tau^*}\\
 &{}\cdot
  \left(\frac{\pi}{2G_3}\right)^{A_r}
  \left(\frac{\pi}{2G_2}\right)^{B_r}
  \prod_\tau
   \left(\frac{\pi}{2G_{i(\tau)}}\right)^{|\nu_\tau^*|}.
\end{split}\tag{HC25}
\] Dividing this expression by Θ gives exactly the R011 norm with the new selected hook products. These are alternative normalizations of the same vector, not additive costs.

For fixed ranks, the AM-R-012 source class index has \(\log I_a(M)=O(\log T_r)\) at M=O(Tᵣ), and its 11-valuation is constant in e. Equation (HC22) gives \(\log\dim\rho=O(\log T_r)\) and \(v_{11}(\dim\rho)=O(\log T_r)\). Therefore the exact class-supported first extraction and the target metric comparison change this fixed-rank valuation upper bound by at most O(log Tᵣ). The creation norm (HC25), the fixed period, and all other primes remain present.

For the explicit first-edge baseline \(\lambda=(Q^6)\), this simplifies to \[
 \dim S_{(h,0^5)}(\mathbf C^6)=\binom{h+5}{5},\qquad
 \theta_\lambda(h)=
            \frac{\binom{h+5}{5}}{\binom{h+Q+5}{Q+5}}.
\tag{HC26}
\] These formulas make the full-Hyp3 normalized first family completely explicit. They do not control an arbitrary irreducible component after further global decomposition or the second restricted class-image lattice.

\subsection{6. What is still needed for actual g/e}\label{what-is-still-needed-for-actual-ge}

The actual finite local section matrix B is now bounded above on selected families, and its first forced divisibility on a fully admissible analytic family is proved. The permitted finite half-orbit adds no saturation loss by K022. The selected compact metric and first source class extraction have explicit factors.

What remains is not a choice of polynomial normalization: it is the integral image of the already produced source vector under the next selected class matrix, including its fixed branch maps and mass pairing. A larger local gain might hold on that restricted image even when it fails on the full permitted selected kernel family. Conversely no existence theorem for a later nonzero theta representation supplies its valuation content or a controlled integral projection. The current results neither prove such a gain nor rule it out.

\subsubsection{Source locators and certificate}\label{source-locators-and-certificate}

\begin{itemize}
\tightlist
\item
  \nolinkurl{Formalism/LocalRingCardinality/Params.lean,} Ranks and HCParam, lines 7--69.
\item
  \nolinkurl{Formalism/LocalRingCardinality/DataTheta.lean,} TPar.Interlace, Dpar and Centred, lines 529--554.
\item
  \nolinkurl{Formalism/SharpRestriction/CuspSpanNormBounds/CayleySignature.lean,} NCGaps and stronger Hyp3, lines 110--118.
\item
  \nolinkurl{Formalism/SharpRestriction/CuspSpanNormBounds/UnimodularSchurMeasures/ExistsTowerOfLocalSetting.lean,} noncompact\_gaps\_of\_spacing, lines 99--103; the separate integer constructor begins at line 181 and its Hyp3 theorem at line 199. The proof here uses the spacing theorem with the displayed constant-midpoint family; it does not claim that this is the different terminal midpoint choice made by that constructor.
\item
  AM-R-011, R11.7--17: exact Fischer hook norm, determinant insertion and normalized row metric. AM-R-012: exact class-supported extraction and mass. AM-K022--025: finite jets, common source gauge, normal forms, adjacent-gap upper bound and the scalable witness.
\item
  \nolinkurl{hyp3_selected_cost_certificate.py} and its JSON output: exact rank/Hyp3 comparisons, factorial and hook valuations, source dimension valuation, 29-level and index valuation; numerical logarithms are explicitly marked as such.
\end{itemize}

\clearpage

\section{AM-R-014: genuine Witt-lower cuspidality and real-type stability in the source predicates}\label{am-r-014-genuine-witt-lower-cuspidality-and-real-type-stability-in-the-source-predicates}

Author: /root, GPT-6 Astra Pro, ALP\_MACHINE, 21 November 2026 UTC.

This note verifies the mathematical contents of the fourth and sixth analytic predicates in \nolinkurl{CayleySignature.AdmS} for genuine unitary theta integrals. The fourth is an application of the established theta tower property. The sixth is a property of actual smooth isotypic subspaces. Neither is asserted as a new tower theorem. The precise application and the distinction between full smooth and compact-finite vectors are recorded because the source's quantified interfaces require them.

Together with AM-R-013 this accounts for four of the six analytic interfaces, subject to the stated genuine geometric interpretation. The Fourier--Jacobi/free-data assertion and the real discrete-series up/down assertion are not proved here. No complete arithmetic \nolinkurl{D7} or \nolinkurl{ThetaX D} record is claimed.

\subsection{1. The two literal interfaces}\label{the-two-literal-interfaces}

\nolinkurl{Formalism/SharpRestriction/CuspSpanNormBounds/CayleySignature.lean}, lines177--184, defines \nolinkurl{ThetaStablePeriodVanishing} as follows. Given a stable subspace of cusp forms on the source \(Y_{j+1}\), assume that all its theta lifts to the specified lower spaces \(Y_{j,i}^{\mathrm{low}}\), for \[
 1\le i\le q_j,
 \qquad q_j=\left\lfloor\frac{m-j}{2}\right\rfloor,
 \tag{R14.1}
\] vanish for every Schwartz input. Its lift to \(Y_j\) must then be cuspidal. This is a statement about constant terms; it does not assert that the surviving lift is nonzero.

Lines196--199 define \nolinkurl{RealConstantNonvanishing}. Despite its name, the literal conclusion here is stability of \nolinkurl{IsDv} under every finite adelic action and each actual real unitary group action. The definition of \nolinkurl{IsDv}, lines123--125, is \[
       \overline F\in\operatorname{typ}(Y_j,D_j).
 \tag{R14.2}
\] There is no lower bound on a real integral in this particular declaration. A different integral-nonvanishing assertion must not be substituted for it.

The abstract record only requires \nolinkurl{YsLow_zero} at lower index zero (\nolinkurl{DataTheta.lean}, lines639--640). To apply the tower theorem we choose the positive lower indices to be the actual Witt-lower Hermitian spaces, as described next. The auxiliary proof model \nolinkurl{SpanPowerMembership/ThetaFamilyAdmissibleThree.lean}, lines67--70, uses a zero-dimensional space at every positive lower index; that auxiliary choice is not used in the present mathematical application.

\subsection{2. Actual spaces and all rational parabolics}\label{actual-spaces-and-all-rational-parabolics}

Keep the common CM host \(E'/F'\) and even integers \(k,m\), with odd \(t\), and put \[
 n_j=k+m-1+jt,\qquad p_j=n_j-q_j.
 \tag{R14.3}
\] A concrete choice of the required rational Hermitian space is \[
 Y_j=(E')^{n_j},\qquad
       h_j=I_{p_j}\oplus(-I_{q_j}).
 \tag{R14.4}
\] Its standard integral lattice is unimodular. At every real place its signature is exactly \((p_j,q_j)\). Each positive/negative coordinate pair is a rational hyperbolic plane: if its orthogonal vectors are \(a,b\), of norms \(1,-1\), then \((a+b)/2,a-b\) are isotropic and have mutual Hermitian pairing one. Consequently the rational Witt index is at least \(q_j\). It is at most \(q_j\) by any real signature. Thus it equals \(q_j\).

For \(0\le i\le q_j\) take \[
 Y_{j,i}^{\mathrm{low}}
    =I_{p_j-i}\oplus(-I_{q_j-i}),
 \qquad
 Y_j\simeq\mathbb H^i\perp Y_{j,i}^{\mathrm{low}}.
 \tag{R14.5}
\] These are actual rational spaces in one Witt tower, of dimensions \(n_j-2i\). An equivalent construction works for an arbitrary Hermitian \(Y_j\) with the indicated real signature: take its actual Witt decomposition, whose rational index \(q\) is at most \(q_j\), and impose the lower-lift hypothesis at least for \(1\le i\le q\). No rational hyperbolic plane is inferred from signature alone in that more general case.

Every maximal proper rational parabolic of the concrete \(U(Y_j)\) is conjugate to the stabilizer of a rational isotropic subspace of dimension \(i\), for some \(1\le i\le q_j\). Witt's extension theorem gives this conjugacy for each dimension. Its Levi has the factors \(\operatorname{Res}_{E'/F'}\mathrm{GL}_i\) and \(U(Y_{j,i}^{\mathrm{low}})\). Thus the list in (R14.1) contains all proper maximal constant-term tests, not just a chosen archimedean subgroup. For \(q_j=0\) there is no such rational parabolic and cuspidality is automatic.

This choice also admits the required positive embeddings at every nonterminal stage, because \[
 p_j-N_j
   =m-1-j-\left\lfloor\frac{m-j}{2}\right\rfloor
   =\left\lceil\frac{m-j}{2}\right\rceil-1\ge0
       \quad(j<m).
 \tag{R14.6}
\] At the terminal stage \(Y_m=I_{N_m-1}\) embeds as the standard codimension-one positive subspace of \(X_m=I_{N_m}\). These statements supply this concrete analytic geometry. They do not supply every other field, auxiliary-prime, level, lattice and arithmetic condition in the full source \nolinkurl{Spaces} record.

\subsection{3. Character parity is constant along each lower tower}\label{character-parity-is-constant-along-each-lower-tower}

Use the genuine global Weil representation for each pair, with its compatible splitting characters. More precisely, take the genuine \nolinkurl{Chr} domain to be the group of unitary Hecke characters with integral angular exponents of \(E'\) whose restriction to \(\mathbb A_{F'}^\times\) belongs to \(\{1,\epsilon_{E'/F'}\}\), with the chosen odd splitting character \(\eta_T\) in that group. This restriction is part of the mathematical character domain. The literal \nolinkurl{IsMu} condition specifies infinity types alone; it would not force trivial finite base-field restriction in the group of all Hecke characters. Within the stated domain, its integral infinity types force the restriction to be one: the alternative quadratic restriction has the nontrivial sign character at every real place, whereas an integral exponent of \(z/\bar z\) is trivial on the real units. In the source notation the characters on the two members of the descending edge are \[
 \chi_{j+1}=\operatorname{car}(\mu,j+1,j+1),\qquad
 \chi_j=\operatorname{car}(\mu,j+1,j).
 \tag{R14.7}
\] Here the even-rank character \(\mu\) restricts trivially to \(\mathbb A_{F'}^\times\), and the usual odd-rank character \(\eta_T\) restricts to the quadratic character \(\epsilon_{E'/F'}\). The source \nolinkurl{car} inserts \(\eta_T^{-1}\) at an even index and makes no insertion at an odd index. Since \(k+m-1\) and \(t\) are odd, \[
 n_j\equiv j+1\pmod2,
 \qquad
 \chi_j|_{\mathbb A_{F'}^\times}
          =\epsilon_{E'/F'}^{\,n_j}.
 \tag{R14.8}
\] Removing \(i\) hyperbolic planes changes the dimension by \(2i\). The same character is therefore a valid splitting character at every lower target in (R14.5). We use these compatible restrictions, with the fixed additive character, in applying the tower formula. This argument does not identify characters of different parity or apply an even-rank doubling twist formula to an odd-rank group.

\subsection{4. Application of the unitary constant-term formula}\label{application-of-the-unitary-constant-term-formula}

The primary input is Wu's unitary tower formula: Proposition3.1, equations(3.2)--(3.4), author preprint pp.9--12. In a mixed Schwartz model, restriction to the smaller unitary Levi identifies the maximal constant term of a cuspidal theta lift with the lower theta lift of the restricted Schwartz function. The general Levi factor adds the explicitly stated character and modulus. Only vanishing is needed below. The symmetric direction is included in the proposition.

Let \(F\) be any cusp form in the source subspace of the literal predicate, and \(\Phi\) any Schwartz input. Fix a maximal parabolic indexed by \(i\) and a target point \(h\). Right equivariance replaces \(\Phi\) by \(\omega(h)\Phi\), so it suffices to test the constant term at the identity in that mixed model. The tower formula is \[
 \operatorname{CT}_{Q_i}\Theta_{Y_j}(\Phi,F)\big|_{U(Y_{j,i}^{\mathrm{low}})}
       =\Theta_{Y_{j,i}^{\mathrm{low}}}(r_i\Phi,F),
 \tag{R14.9}
\] where \(r_i\) is the Schwartz restriction at zero in the extra isotropic coordinates. The compatible Haar measures give this displayed normalization; changing them multiplies the two sides by the corresponding nonzero measure constant and has no effect on the following implication.

The lower-lift hypothesis quantifies over every Schwartz function, so the right side of (R14.9) is zero, including after the above right translation. The constant term therefore vanishes everywhere on that maximal parabolic. This holds for every \(i\), and hence for every proper rational parabolic: the constant term along a smaller parabolic is an iterated constant term through a containing maximal one. This proves cuspidality of the lift.

Although Wu states the proposition using an irreducible cuspidal representation, its pointwise identity applies to every cusp form. In the cited proof, equations(3.21)--(3.22) use only the vanishing of its proper constant terms; equations(3.23)--(3.25) give the surviving zero orbit. No projection onto an irreducible constituent is used here. In particular an arbitrary stable, possibly nonclosed source subspace in the literal predicate causes no semisimplicity gap. Absolute convergence for the cuspidal input is part of the cited formula.

Thus \nolinkurl{ThetaStablePeriodVanishing} holds in this genuine Witt-lower interpretation. Once compatible splitting characters are fixed, the constant-term argument uses no further \nolinkurl{Hyp3} inequalities and needs no nonzero lower image. Character parity was checked above using the stated parity conditions. Thus it implies the corresponding restricted source quantifiers. It supplies no bound for the norm or denominators of the descended vector.

\subsection{5. Full smooth isotypic spaces and the sixth predicate}\label{full-smooth-isotypic-spaces-and-the-sixth-predicate}

Use the usual smooth automorphic space of moderate growth for the ambient real-group action, with its cusp subspace of rapidly decreasing smooth cusp forms and its usual Schwartz seminorms. In particular the cuspidal vectors under discussion are in finite sums of smooth globalizations of actual cuspidal automorphic embeddings. This is also the cusp domain used in section4: arbitrary smooth functions satisfying only a formal zero-constant-term condition are not being substituted into an absolutely convergent theta integral. The genuine theta maps and constant-term maps are taken in these standard smooth spaces. For a parameter collection \(D\), let \(\operatorname{typ}(V,D)\) be the algebraic sum of the smooth real isotypic subspaces with the specified real representation types. If no such representation exists, this space is zero. At an indefinite real place a discrete-series type here means its full smooth globalization, rather than only its compact-finite vectors. Compact-finiteness remains a condition on the test kernels at the places where the source explicitly asks for it.

Each real-group operator preserves every such isotypic subspace. Every finite adelic right translation commutes with the real action. If \(T:V_\sigma^\infty\longrightarrow\mathcal A\) is an equivariant realization of one of the prescribed real types, then \(R(g_f)T\) is another such realization (or zero). It follows that \[
 R(g_f)\operatorname{typ}(V,D)
       \subset\operatorname{typ}(V,D),\qquad
 R(h_\infty)\operatorname{typ}(V,D)
       \subset\operatorname{typ}(V,D).
 \tag{R14.10}
\] Pointwise complex conjugation commutes with right translation of functions by a fixed real or finite group element. Consequently \[
 \overline F\in\operatorname{typ}(V,D)
 \ \Longrightarrow\ 
 \overline{R(g)F}=R(g)\overline F
                    \in\operatorname{typ}(V,D).
 \tag{R14.11}
\] This is exactly the finite and real stability in \nolinkurl{RealConstantNonvanishing}, for every parameter collection and every stage. No numerical nonzero constant is inferred from the name of the declaration.

For a noncompact real group, it would be incorrect to make this argument with only the compact-finite vectors as an ambient space carrying every real translation. Such a subspace need not be stable under that action. The smooth globalization in this interpretation resolves that actual domain issue. On the definite side of AM-R-013, the selected real representation is finite-dimensional, so that issue does not change its algebraic finite-sum identities.

\subsection{6. Framework translation and exact remaining gates}\label{framework-translation-and-exact-remaining-gates}

The parabolic and Levi groups are actual reductive group schemes over the number field. Their character and representation gradings belong to the SGA3 framework, in particular ExposeI, Proposition4.7.3 and Corollary4.7.3.1, already inspected in the geometric dictionary. The passage to a Schur representation class belongs to the carried SGA6 ExposeVI, Theorem3.3. Those geometric statements do not prove the automorphic constant-term formula; Wu's analytic theorem is the separate input used above. The polarized kernel remains an actual theta-group/biextension object in SGA7I, ExposeVII, sections2.7 and2.9, pp.169--174, with its normalizations retained as in AM-DO-026/029.

The two remaining analytic assertions are \nolinkurl{ThetaUnfoldingFreeData} (lines168--175), with its actual Fourier--Jacobi period and factorizable free-data witness, and \nolinkurl{RealThetaLiftDiscreteSeries} (lines186--194), with its precise real up/down parameter and all smaller-target vanishing. Completing these is not equivalent to the still separate arithmetic row ties, raw-Xi selected nonvanishing, terminal mass or complete cost array.

\subsection{References}\label{references-3}

Wu, C. (2013). Irreducibility of theta lifting for unitary groups. \emph{Journal of Number Theory, 133}(10), 3296--3318. Author preprint posted 2014: \url{https://arxiv.org/pdf/1409.0770} . Proposition3.1 and equations(3.2)--(3.4), printed p.9; proof, especially equations(3.21)--(3.25), pp.11--12. The primary statement and these proof passages were directly inspected in this run.

Gan, W. T., Qiu, Y., \& Takeda, S. (2014). The regularized Siegel--Weil formula (the second term identity) and the Rallis inner product formula. \emph{Inventiones Mathematicae, 198}, 739--831. \url{https://arxiv.org/pdf/1207.4709} . Section11.2, author preprint pp.51--52, records the first-occurrence cuspidality consequence for the full classical-group setting, including unitary groups. The more explicit unitary formula used here is Wu's Proposition3.1.

The exact source declarations and mathematical uses are listed above. This verification is not a new proof of the cited tower property and includes no Lean execution.

\clearpage

\section{AM-K027 --- the exact 11-adic slope in an actual full-Hyp3 selected family}\label{am-k027-the-exact-11-adic-slope-in-an-actual-full-hyp3-selected-family}

Author: GPT-6 Astra Pro, ALP\_MACHINE \nolinkurl{/root/kummer_extension}, 21 November 2026 UTC. Inherited reasoning effort is not separately exposed.

The accepted DP51 common-weight theorem supplies the missing lower bound for AM-K024. The full CM gauge comparison AM-K018 now makes that theorem applicable to the actual ordinary host \(A=J^2\), including every integral local section and the exact DP59 half-orbit. The resulting fixed-rank, varying-weight selected family has \textbf{exact asymptotic normalized derivative slope \(1/1220\)}. The statement remains a selected kernel/first-extraction result. It does not identify the later restricted automorphic image or resolve the literal raw-Xi matching equation.

\subsection{1. The common-weight bound in the actual source gauge}\label{the-common-weight-bound-in-the-actual-source-gauge}

Let R be a finite unramified complete extension of \(\mathbb Z_{11}\) containing the needed CM eigenvalues and torsion. Normalize \(v(11)=1\). Set \[
 J=\operatorname{Jac}(y^2=x^7-1),\qquad
 L=L_{2\infty},\qquad X=J^\ell,
\] with its product principal line \(\mathcal L\), where \(\ell\) is fixed. Use the actual Chern/source reduced gauge of AM-K018, and let \(\kappa\) be a positive integer prime to 11. A section is integral here if its algebraic coefficient in an integral formal frame of \(\mathcal L^\kappa\) is integral. No ordinary Taylor integrality of the reduced frame is assumed.

Put \[
                   w=\frac{121}{1220},\qquad
                   c=\frac1{1220},\qquad s=\frac1{10}.
\tag{EX1}
\]

\subsubsection{Proposition EX2. Full actual local-module estimate}\label{proposition-ex2.-full-actual-local-module-estimate}

For every integral formal local section S and its actual reduced germ \(F_S\), in the primitive ordinary J logarithmic coordinates or any integral invertible linear change of them, \[
       v([z^\alpha]F_S)+w|\alpha|\ge0.
\tag{EX2}
\] Consequently every integral homogeneous differential polynomial \(P=\sum_{|\alpha|=D}p_\alpha Z^\alpha\) satisfies \[
 v\bigl(P(\partial)F_S(0)\bigr)
 \ge \frac D{1220}
       -\frac1{10}\max_{\alpha\in\operatorname{supp}P}
                                 \sum_m s_{11}(\alpha_m).
\tag{EX3}
\] Zero values have valuation \(+\infty\). For D≥1 the weaker uniform estimate is \[
 v\bigl(P(\partial)F_S(0)\bigr)
 \ge \frac D{1220}
          -3\ell\bigl(1+\lfloor\log_{11}D\rfloor\bigr).
\tag{EX4}
\] The constants in (EX2)--(EX4) are independent of \(\kappa\) as long as \(11\nmid\kappa\), and apply to the allowed levels \(\kappa=M^2\), \(M=29^e\).

\textbf{Proof and exact accepted interface.} DP49 identifies each actual formal J with its polarized arbitrary-unit singleton/complement block of height three. AM-K018, (K19.2) and (K19.13), identifies the source/Chern reduced section with \[
              F_S=f_S\,\prod_{b=1}^{\ell}U_b^\kappa B_b^{-\kappa}
\tag{EX5}
\] in the actual integral product frame. Here \(f_S\) is an arbitrary integral formal function on the product, and the \(U_b\) are integral units. These exact identifications are carried results; in particular there is no additional rational quadratic exponential.

Use accepted DP51 §6, (HA17)--(HA21), with p=11, h=3 and primitive parameter \(-\kappa\). Its common-weight substitution makes every full primitive kernel \(B_b^{-\kappa}\) integral at weight w, and makes the inverse logarithm maps integral at that same common weight. The latter statement applies to a joint prefactor as well as a separated product of prefactors: composing the integral series \(f_S\prod_bU_b^\kappa\) with the product of those zero-constant integral inverse maps is integral. Each resulting formal coefficient uses only finitely many terms of the prefactor. Thus no separated-factor hypothesis is needed for (EX5). This gives (EX2) directly from the accepted common-weight theorem.

If \(z=Cz'\) with \(C\in\operatorname{GL}_{3\ell}(R)\), the same common scale is applied to every coordinate. Substitution by C preserves integral coefficients at that scale; its inverse gives the converse. Hence (EX2) applies to the exact \(E'=\mathbb Q(\zeta_{28})\) coordinates \(z_{i,\pm}=u_i\pm i v_i\): both this matrix and its inverse are integral at 11. It also applies after the unit diagonal comparisons with the display eigenbasis.

Finally \(\partial^\alpha F_S(0)=\alpha![z^\alpha]F_S\). The Legendre identity \[
 v(\alpha!)=\frac{|\alpha|-\sum_ms_{11}(\alpha_m)}{10}
\] and \(1/10-w=1/1220\) prove (EX3) by the nonarchimedean triangle inequality. Bounding each digit sum by \(10(1+\lfloor\log_{11}D\rfloor)\) proves (EX4). These factorial manipulations are the same normalization as accepted DP51; no new display theorem is asserted. ∎

AM-K023 separately proves Hurwitz integrality of this entire actual local module. Thus the left side of (EX3) is also nonnegative, and an integer ceiling may be taken over the present unramified coefficient ring. The bound is a genuine full weighted coefficient estimate. It is not obtained by iterating the mod-11 normal form or by assuming a finitely generated differential-stable lattice, which AM-K019 excludes in the relevant unrestricted setting.

\subsection{2. A sharper row-degree error for the selected module}\label{a-sharper-row-degree-error-for-the-selected-module}

On the actual edge \(U(a)\longrightarrow U(b)\), the Hom variety is \(J^{2ab}\), so \(3\ell=6ab\). Let the selected row degrees be the six partitions \(\nu_\tau\), and put \[
 D=\sum_\tau|\nu_\tau|,\qquad
 d_{11}(n)=\begin{cases}0&n=0,\\1+\lfloor\log_{11}n\rfloor&n>0.\end{cases}
\] Every monomial in a polynomial of these row degrees satisfies \(\sum_{j=1}^b\alpha_{\tau,i,j}=\nu_{\tau,i}\). Since every entry in that row is at most \(\nu_{\tau,i}\), (EX3) gives \[
 \boxed{
 v\bigl(P(\partial)F_S(0)\bigr)
 \ge \frac D{1220}
                 -b\sum_{\tau,i}d_{11}(\nu_{\tau,i}).}
\tag{EX6}
\] This holds for every integral polynomial in the specified row-weight space, in particular for the fixed source-highest part of the actual selected Cauchy summand used here. The entire abstract tensor product of the source and target Schur modules also has other source row weights; (EX6), with these same individual row degrees, is not an assertion about all those other weight vectors. It introduces no abstract Schur projector. For the primitive highest prefix-minor polynomial using only the first a columns, b in (EX6) may be replaced by a: \[
 v\bigl(P_{\boldsymbol\nu}(\partial)F_S(0)\bigr)
 \ge \frac D{1220}
                 -a\sum_{\tau,i}d_{11}(\nu_{\tau,i}).
\tag{EX7}
\] These are statements in the explicit polynomial coefficient lattice. Changing to a different rationally normalized Weyl lattice requires its actual scaling factor.

For every \(M\ge3D+2\), AM-K022 identifies the image ideals of the formal jet matrix and the permitted half-orbit tuple matrix. The normalizing factor \(M^{-D}\) is an 11-unit. Consequently (EX3), (EX6), and (EX7) are also bounds for their actual allowed normalized kernel coefficients. The fixed period \(C^{2ab}\) remains outside these algebraic coefficient estimates.

\subsection{3. Exact asymptotic slope for AM-K024}\label{exact-asymptotic-slope-for-am-k024}

Fix one edge and its fixed residual partitions \(\nu_\tau^*\), as in AM-K024. In particular a,b and \[
 D_0=\sum_\tau|\nu_\tau^*|,\qquad
 c_0=v\left(\prod_{\tau,i}\nu_{\tau,i}^*!\right)
\] do not vary. Let \(Q_0\) be the residual prefix-minor polynomial on columns 2,\ldots,a+1. Define the finite constant \[
 K_0=\frac1{10}\max_{\beta\in\operatorname{supp}Q_0}
                                     \sum_m s_{11}(\beta_m).
\tag{EX8}
\] It can be bounded explicitly by \(a\sum_{\tau,i}d_{11}(\nu_{\tau,i}^*)\). Set \[
 A_r=11^r,\quad B_r=121\,11^r,\quad
 T_r=122\,11^r+D_0.
\tag{EX9}
\] AM-K024 supplies \(h_r\in[0,A_r]\), \(k_r\in[0,B_r]\) and the primitive integral polynomial \[
 P_r=Z_{3,+,1,1}^{h_r}Z_{3,-,1,1}^{A_r-h_r}
     Z_{2,+,1,1}^{k_r}Z_{2,-,1,1}^{B_r-k_r}Q_0.
\tag{EX10}
\] It already lies in the \textbf{same exact source-and-target Schur summand} with the selected partitions prescribed by (SC4). It is not merely a source-highest vector awaiting an unspecified target projection.

Choose \(M_r=29^{e_r}\) least with \(M_r\ge3T_r+2\). Let \(\mathfrak I_r\) be the derivative image ideal for \(P_r\) on the integral formal \(T_r\)-jet module of \(\mathcal L^{M_r^2}\); equivalently, by AM-K022 J30, it is the ideal generated by the actual allowed half-orbit tuple values. AM-K024 proves it is nonzero and \[
          v(\mathfrak I_r)\le \frac{11^r-1}{10}+c_0.
\tag{EX11}
\]

The four growing variables in (EX10) are disjoint from all variables of \(Q_0\). Therefore the digit term in (EX3) is exactly bounded by \[
 \frac{s_{11}(h_r)+s_{11}(A_r-h_r)
       +s_{11}(k_r)+s_{11}(B_r-k_r)}{10}+K_0
 \le4d_{11}(T_r)+K_0.
\] Combining this lower bound with (EX11) proves \[
 \boxed{
 \frac{T_r}{1220}-4d_{11}(T_r)-K_0
 \ \le\ v(\mathfrak I_r)
 \ \le\ \frac{T_r}{1220}
                 +c_0-\frac1{10}-\frac{D_0}{1220}.}
\tag{EX12}
\] The same two bounds hold for an individual allowed tuple chosen to attain the image ideal. Such a tuple exists because the half-orbit basis is finite and generates that ideal; taking its least-valued nonzero entry attains the minimum. Its finite Schwartz vector is one translate of the fixed norm-one vacuum, with no interpolation coefficient height included.

It follows in particular that \[
 \boxed{
 \lim_{r\to\infty}\frac{v(\mathfrak I_r)}{T_r}
 =\lim_{r\to\infty}\frac{v(a_r)}{T_r}
 =\frac1{1220}.}
\tag{EX13}
\] This is a limit for the constructed sequence, stronger than a mere liminf upper bound. For the indices r covered by (EX13), the tuple, the selected first-row increments, and the level may vary; the ranks, residual partitions, host, and common edge character remain fixed.

By AM-K024 §2 and AM-K026 §1, choose the fixed residual partitions from the explicit full-Hyp3 baseline. Every subsequent nonnegative \(u_0\) increment preserves \textbf{all} clauses of that stronger analytic Hyp3, including Centred, NCGaps and the 1801 condition. Thus (EX13) holds inside that full analytic admissibility family. It does not preserve separately fixed SizeData₂ upper windows, nor does it by itself supply the complete old D7 arithmetic host.

\subsection{4. Exact selected-metric and first-extraction normalization}\label{exact-selected-metric-and-first-extraction-normalization}

For a fixed residual partition \(\lambda\) and first-row increment h, AM-K026 proves that the target vector represented by the primitive polynomial has squared norm \[
 \theta_\lambda(h)=\frac{h!H_\lambda}{H_{\lambda+h e_1}}
  =\prod_{s\in S_\lambda}\frac{s}{h+s},
\tag{EX14}
\] where \(S_\lambda\subset\{1,\ldots,\lambda_1+a-1\}\) is a fixed set of \(\lambda_1\) positive integers. Let \(\Theta_r\) be the product of the four such factors for the increments in (EX10). The product formula gives both valuation bounds \[
 -|S_\lambda|\log_{11}(h+\lambda_1+a)
 \le v(\theta_\lambda(h))
 \le\sum_{s\in S_\lambda}v(s).
\tag{EX15}
\] Thus \(v(\Theta_r)=O(\log(T_r+1))\), with a fixed constant depending on the residual partitions. Normalizing the target vector multiplies its algebraic value by \(\Theta_r^{-1/2}\), an explicitly retained algebraic scalar. Equations (EX12)--(EX15) imply \[
             \lim_{r\to\infty}
       \frac{v(a_r/\sqrt{\Theta_r})}{T_r}=\frac1{1220}.
\tag{EX16}
\] There is no claim that this square-root normalization is integral.

In the concrete compact class-supported first extraction of AM-R-012, the further unit-input factor is \(1/\sqrt{I_a(M_r)\dim\rho_r}\). For fixed a,b, \[
 I_a(M)=M^{6a^2}\prod_{i=1}^a(1-29^{-i})^6
\] has 11-valuation independent of the exponent e, and the dimension ratios under first-row growth are products of fixed finitely many linear factors in the growing increments. Hence \(v(\dim\rho_r)=O(\log(T_r+1))\). Retaining the common period, that specific first value has the form \[
 C^{2ab}\,
       \frac{a_r}{\sqrt{\Theta_r I_a(M_r)\dim\rho_r}},
\tag{EX17}
\] with the exact mass convention of AM-R-012. Its algebraic coefficient after removal of the fixed period again has slope \(1/1220\). This conclusion is about that explicitly constructed first extraction. The raw source-array matching and any later restriction to a previously produced automorphic image are separate gates, not consequences of the normalization calculation.

The full archimedean creation norm is still the exact expression in AM-K026 (HC25), or equivalently the corresponding product of new hooks after dividing by \(\Theta_r\). Its factorial growth has not been replaced by an \(O(\log T_r)\) assertion. Only the specified target comparison and fixed-rank class/dimension factors have logarithmic logarithm or valuation corrections. These are different quantities and must not be conflated in a net gain/error computation.

\subsection{5. Stronger compulsory bounds for the concrete full-Hyp3 baseline}\label{stronger-compulsory-bounds-for-the-concrete-full-hyp3-baseline}

The bounds above also sharpen AM-K026's first-divisibility conclusion without using any iteration of a congruence. In its exact example \[
 k=6,\quad m=2,\quad t=1795,\quad d=1796,
 \qquad Q=46\,539\,457\,175,
\] the two edge partitions are \[
 \nu_0=(Q^6),\qquad
 \nu_1=((3Q)^6,(2Q)^{1796})
\] at each of the six labels. Every positive row part has 11-adic digit length 11. Applying (EX6) and (EX7), then taking integer ceilings, gives:

\begin{longtable}[]{@{}
  >{\raggedright\arraybackslash}p{(\columnwidth - 8\tabcolsep) * \real{0.1579}}
  >{\raggedleft\arraybackslash}p{(\columnwidth - 8\tabcolsep) * \real{0.2105}}
  >{\raggedleft\arraybackslash}p{(\columnwidth - 8\tabcolsep) * \real{0.2105}}
  >{\raggedleft\arraybackslash}p{(\columnwidth - 8\tabcolsep) * \real{0.2105}}
  >{\raggedleft\arraybackslash}p{(\columnwidth - 8\tabcolsep) * \real{0.2105}}@{}}
\toprule\noalign{}
\begin{minipage}[b]{\linewidth}\raggedright
Edge
\end{minipage} & \begin{minipage}[b]{\linewidth}\raggedleft
Degree D
\end{minipage} & \begin{minipage}[b]{\linewidth}\raggedleft
Lower bound for every integral selected polynomial in the fixed source-highest row space
\end{minipage} & \begin{minipage}[b]{\linewidth}\raggedleft
Lower bound for the primitive highest prefix-minor coefficient
\end{minipage} & \begin{minipage}[b]{\linewidth}\raggedleft
Adjacent-gap upper bound for that highest coefficient
\end{minipage} \\
\midrule\noalign{}
\endhead
\bottomrule\noalign{}
\endlastfoot
6 → 1802 & 1,675,420,458,300 & 1,372,581,866 & 1,373,293,082 & 167,542,045,704 \\
1802 → 3598 & 1,008,044,642,410,500 & 825,838,183,001 & 826,051,784,873 & 100,804,464,197,676 \\
\end{longtable}

The upper bounds are AM-K026's exact adjacent-gap factorial bounds, hence apply to the image ideals and to some individual permitted tuple. The lower bounds hold for every integral section in the corresponding local module and every indicated polynomial in the stated fixed row space. The middle two columns do not assert that every element of a rationally rescaled abstract selected lattice is integral in this polynomial coefficient lattice.

These intervals are genuine local coefficient bounds, but they still leave a wide range for the exact positive Smith valuation of this particular fixed baseline. The asymptotically exact conclusion (EX13) instead concerns the explicitly varied first-row family. Neither statement computes the selected class matrix restricted to a previous theta image.

\subsection{6. Sources, new content, and check scope}\label{sources-new-content-and-check-scope}

The display and common-weight estimates are accepted results, not newly derived results:

\begin{itemize}
\tightlist
\item
  Supplied DP51, \nolinkurl{research/harmonic_asymptotic_rate_DP51.md}, §6, (HA17)--(HA21): entire polynomial module, exact common weight and factorial digit error. The full statement and proof were read for this application.
\item
  Supplied DP46, \nolinkurl{research/unit_arrow_kernel_invariance_DP46.md}, §§4--6, especially (UI16)--(UI20): arbitrary unit arrows, integral inverse-map profile and invariant integral frames.
\item
  Supplied DP59, \nolinkurl{research/theta_translate_span_DP59.md}, §§6--8, especially the lower-bound argument in §7: actual integral half-orbit basis, formal prefactor and normalized derivative convention.
\item
  AM-K018 (K19.2), (K19.13): exact actual source/Chern gauge on ordinary J powers; AM-K022 J30: equality of all finite determinantal ideals; AM-K024 (SC4)--(SC20): the exact selected witness, full-Hyp3 preservation and its fixed filler; AM-K026 (HC18)--(HC25): exact target norm and mass comparisons.
\item
  AM-DO-025: actual source-and-target Cauchy labels; AM-R-011: selected metric; AM-R-012: concrete class-supported first extraction. Their source and primary bibliographic locators are retained in those reports.
\end{itemize}

The new deductions are the transfer of the accepted full common-weight estimate to every actual ordinary-product local section, its row-support version (EX6)--(EX7), the matching selected-family sandwich and limit (EX12)--(EX13), and the retained normalized limit (EX16)--(EX17). The finite arithmetic table has a separate exact certificate. No Lean build, new display proof, full source instance, positive Smith-factor computation, or later-image valuation theorem is claimed.

\clearpage

\section{AM-DO-031: the literal terminal Xi, its finite realization, exact norms, and an explicit nonzero raw value}\label{am-do-031-the-literal-terminal-xi-its-finite-realization-exact-norms-and-an-explicit-nonzero-raw-value}

Author: GPT-6 Astra Pro, ALP\_MACHINE /root/definite\_oscillator; inherited reasoning effort is not separately exposed; 21 November 2026 UTC.

The finite realization in AM-DO-029 covered the nonterminal branches. This finding treats the actual remaining branch xiLast. It gives explicit support and period lattices, the corresponding finite Heisenberg expansion, its exact norm, and a rational matrix at which the raw array has an explicitly known nonzero value. None of these statements identifies that value with a theta value, a selected automorphic matrix coefficient, or the MainX terminal mass.

\subsection{1. Source and notation}\label{source-and-notation}

All recovered source files are relative to work/research/alpoge\_source, commit 7cbf7a0c32fca1badf05931879a37c3f8a2b03b4. The claims below are a mathematical reading of those definitions, not a claim to have compiled their Lean dependency graph.

\begin{longtable}[]{@{}
  >{\raggedright\arraybackslash}p{(\columnwidth - 4\tabcolsep) * \real{0.3333}}
  >{\raggedright\arraybackslash}p{(\columnwidth - 4\tabcolsep) * \real{0.3333}}
  >{\raggedright\arraybackslash}p{(\columnwidth - 4\tabcolsep) * \real{0.3333}}@{}}
\toprule\noalign{}
\begin{minipage}[b]{\linewidth}\raggedright
Source
\end{minipage} & \begin{minipage}[b]{\linewidth}\raggedright
Locator
\end{minipage} & \begin{minipage}[b]{\linewidth}\raggedright
Content
\end{minipage} \\
\midrule\noalign{}
\endhead
\bottomrule\noalign{}
\endlastfoot
\nolinkurl{Formalism/LocalRingCardinality/Params.lean} & 6--17, 32--40 & actual rank restrictions; r, N, rS \\
\nolinkurl{Formalism/LocalToGlobal/Columns/FiniteSumNonzeroLargeDepth.lean} & 345--347, 1103--1104 & Q-frame condition and frsF \\
same & 711--715, 780--806 & literal PhiQ, PhiG, PhiGS \\
same & 817, 838--852, 1039--1042 & qOf; limited-depth additive characters; depth hypotheses \\
same & 1123--1131, 1142--1156 & ybSet, repY, PhiTop, ftY \\
\nolinkurl{Development/AbelianVarietyTheory/Definitions/GlobalTwisting/IntegralGLSuccEmbeddingMemParahoricTimesPrincipa.lean} & 391--393, 420--441 & locTop, entB\_top, terminal rank identity \\
same & 444--509, 531--547 & G-support, congruence invariance, representative denominator \\
\nolinkurl{Development/AbelianVarietyTheory/Definitions/AutomorphicMembership/XiAtPenultimateEqXiCpl1.lean} & 124--130, 146--163 & admissible terminal representative; xiOver, xiLast, xiAt \\
\nolinkurl{Formalism/LocalToGlobal/Columns/LevGSym.lean} & 11--13, 23--28, 108--112 & symIdx and the exact symmetric unipotent support \\
\nolinkurl{Development/AbelianVarietyTheory/Definitions/AutomorphicMembership/TrueColumnChoicesHeckeCharReferenceEqOne.lean} & 242 & actual sharp=true condition \\
\nolinkurl{Development/AbelianVarietyTheory/Definitions/LatticeFrames/LevChiEqOne.lean} & 54--55, 68--71 & standard base lattices and the prescribed framed raw array \\
\end{longtable}

Put

\[
 h=m-1,\qquad r=k+m-1,\qquad
 a=N(m-1)=r+ht,\qquad
 b=N(m)=a+t+1.
 \tag{DO31.1}
\]

The actual restrictions include even \(m\ge2\), even \(k\ge m+3\), and \(t\ge r+5\). In particular every matrix block used below exists. Set \(r_S=m/2-1\). Then \(0\le r_S\) and \(k+r_S<r\).

Write \(\mathsf M=\mathrm{cst.M}\ge1\); this symbol is unrelated to the auxiliary level \(29^e\) of AM-DO-030. At a selected finite prime \(P\), write

\[
 q=q_P=|\mathcal O_E/P|,\qquad
 d_P=v_P(q).
 \tag{DO31.2}
\]

Thus \(|q|_P=q^{-d_P}\). The rational integer q appearing in the source is not replaced silently by a uniformizer. For the prescribed G-primes, Primes.Ok.vG\_card, vG\_norm, and piG\_cop give \(2[E^+:\mathbf Q]\) distinct degree-one primes after adjoining the conjugate set; the degree formula then forces ramification index1. Thus \(d_P=1\). These axioms are in \nolinkurl{AbelianHostFieldConditions.lean,} lines207--218. At a general Q-prime the formulas below retain \(d_P\). The common degree-one case is \(d_P=1\).

Fix actual allowed terminal decorations \(p,p'\) and \(Y\in\mathrm{ybSet}(p.N)\). In their prescribed local frames, write the matrix variable as

\[
 X=\begin{pmatrix}Z\\ \eta\end{pmatrix},\qquad
 Z\in M_{(b-1)\times a}(E_P),\quad
 \eta\in M_{1\times a}(E_P),
 \quad
 y'_P=f_P^{-1}\operatorname{repY}(Y).
 \tag{DO31.3}
\]

The top rank identity and entB\_top say exactly that the terminal PhiTop reads the ordinary upper matrix Z. There are no remaining auxiliary-column powers of q hidden in this assertion: \(cOf(m-1)=0\). Any piT scaling of a terminal target decoration is already in its specified frame. Applying an additional scaling would change the literal array.

With \(\epsilon_{\mathrm{Sg}}=-1\), its local factor is

\[
 F_P(X)=\Phi_{\mathrm{Top},P}(Z)\,
 \begin{cases}
  1_{\mathcal O_P^a}(\eta)\,\psi_P(-\eta y'_P),
       &P\in G,\\[2mm]
  q^{-\mathsf M a}
  1_{q^{-\mathsf M}\mathcal O_P^a}(\eta)\,
       \psi_P(-\eta y'_P),&P\in Q.
 \end{cases}
 \tag{DO31.4}
\]

Here G denotes the G branch and Q denotes the complement of G in rowSG; for the source's disjoint prime data these are its usual G and Q sets. The prescribed separation of the T-places from rowSG is retained. The complete raw xiLast is the product of these factors and the T-incidence factors.

\subsection{2. Representative bounds and finite windows}\label{representative-bounds-and-finite-windows}

\subsubsection{2.1 The actual representative has the needed denominator}\label{the-actual-representative-has-the-needed-denominator}

In the prescribed frames,

\[
 y'_P\in
 \begin{cases}
 q^{-1}\mathcal O_P^a,&P\in G,\\
 q^{-\mathsf M}\mathcal O_P^a,&P\in Q.
 \end{cases}
 \tag{DO31.5}
\]

The G assertion is precisely repY\_int at the source locator above. Here is also the Q argument, which does not need an unproved continuity assertion for psi.

By ybSet, there is a vector y with \(y_P\in q^{-\mathsf M}N_P\), and every member of Y differs from y by a vector in \(q^{\mathsf M}N_P\). Since \(\mathsf M\ge1\), \(q^{\mathsf M}N_P\subset q^{-\mathsf M}N_P\). The source's epsilon representative belongs to Y: PtAt.repY\_mem proves this for the actual terminal point; the same conclusion follows directly because the displayed coset contains y. Hence \(\operatorname{repY}(Y)_P\in q^{-\mathsf M}N_P\). DecoF's frame condition identifies \(N_P\) with \(\mathcal O_P^a\), proving the Q case.

\subsubsection{2.2 No global additive-character hypothesis is introduced}\label{no-global-additive-character-hypothesis-is-introduced}

LocCst.Ok only prescribes additivity on \(q^{-D_P}\mathcal O_P\), and value1 on \(\mathcal O_P\). These conditions suffice. At G, the arguments in (DO31.4) belong to \(q^{-1}\mathcal O_P\), and \(D_P\ge1\). At Q they belong to \(q^{-2\mathsf M}\mathcal O_P\), and ChIdx.Ok.depthQ gives \(D_P\ge2\mathsf M\).

The localized rational module modulo its integral submodule has the same finite quotient as its completion. Additivity and psi\_int make psi constant on those integral cosets, so it extends uniquely to this finite local quotient. On either such finite additive quotient, psi is nonzero: additivity applied to x and -x gives \(\psi(x)\psi(-x)=\psi(0)=1\). The quotient is finite, so every value is a root of unity. Thus it has complex absolute value1 and is a unit at every finite place of its coefficient field.

\subsubsection{2.3 Support and period lattices}\label{support-and-period-lattices}

In the chosen frames let \(L=M_{b\times a}(\mathcal O_P)\). The following windows are valid:

\begin{longtable}[]{@{}
  >{\raggedright\arraybackslash}p{(\columnwidth - 4\tabcolsep) * \real{0.3333}}
  >{\raggedright\arraybackslash}p{(\columnwidth - 4\tabcolsep) * \real{0.3333}}
  >{\raggedright\arraybackslash}p{(\columnwidth - 4\tabcolsep) * \real{0.3333}}@{}}
\toprule\noalign{}
\begin{minipage}[b]{\linewidth}\raggedright
Place
\end{minipage} & \begin{minipage}[b]{\linewidth}\raggedright
Support lattice \(L_P^+\)
\end{minipage} & \begin{minipage}[b]{\linewidth}\raggedright
Period lattice \(L_P^-\)
\end{minipage} \\
\midrule\noalign{}
\endhead
\bottomrule\noalign{}
\endlastfoot
T & \(L\) & \(P L\) \\
G & \(L\) & \(qL\) \\
Q & \(M_{(b-1)\times a}(\mathcal O_P)\oplus q^{-\mathsf M}\mathcal O_P^a\) & \(qM_{(b-1)\times a}(\mathcal O_P)\oplus q^{\mathsf M}\mathcal O_P^a\) \\
\end{longtable}

At T this is the carried exact incidence function \(1_{\{vX=u\}}\), including integral support. At G the source proves the support and congruence statements directly. They can also be checked in (DO31.4): PhiTop depends only on the reduction of Z, and changing \(\eta\) by q times an integral row changes \(\eta y'_P\) by an integral element.

At Q, PhiQ requires every entry of Z to be integral, requires the entries in its bottom-left \((b-1-r)\)-by-r block to lie in q times the local ring, and evaluates a residue character on the top r-by-r determinant. Adding q times an integral matrix preserves all these conditions and the determinant residue. Changing \(\eta\) by \(q^{\mathsf M}\) times an integral row changes its pairing with \(y'_P\) by an integral element. Both old and new character arguments remain in the permitted depth. This proves invariance.

All three support lattices contain L and all three period lattices lie in L. Therefore the indices in these frames are

\[
\begin{array}{c|c|c}
 &R_P=[L_P^+:L]&N_P=[L:L_P^-]\\ \hline
 T&1&q^{ab}\\
 G&1&q^{ab}\\
 Q&q^{d_P\mathsf M a}
   &q^{d_P((b-1)a+\mathsf M a)}.
\end{array}
\tag{DO31.6}
\]

For an arbitrary rational choice of base lattice, transport these exact lattices by \(X\mapsto f'_P Xf_P^{-1}\), and use the sum/intersection and Smith-normal-form construction of AM-DO-029. Formula (DO31.6) is not a license to discard the determinant of a nonintegral frame. For the literal standard TwDat base points, both source and target lattices are standard by TwDat.Is.base; the frames at their places are integral isomorphisms, so (DO31.6) applies without a volume factor.

\subsection{3. Exact finite Heisenberg construction}\label{exact-finite-heisenberg-construction}

Extend the local product by the vacuum half-indicator outside its selected places. The resulting finite Schwartz function \(\Phi_{\mathrm{last}}\) agrees with the literal rational xiLast on its source half-model. Its finite support and local constancy follow from the windows just proved, rather than from an assumption about psi beyond its prescribed depth.

Let \(L^+\supset L\supset L^-\) be the actual product windows, \(R=[L^+:L]\), and \(N=[L:L^-]\). For coset representatives \(a_0\in L^+/L\), and \(y\in(L^-)^\perp/L^\perp\), set

\[
 c_{a_0,y}
 =\frac1N\sum_{u\in L/L^-}
       \Phi_{\mathrm{last}}(a_0+u)\,
       \psi_f(-B(u,y)).
 \tag{DO31.7}
\]

Fourier inversion gives

\[
 \Phi_{\mathrm{last}}(t)
 =\sum_{a_0,y}c_{a_0,y}\,
      \psi_f(B(t-a_0,y))1_L(t-a_0).
 \tag{DO31.8}
\]

Consequently the literal normalized place model sends it to

\[
 f_{\mathrm{last}}
 =\sum_{a_0,y}c_{a_0,y}W_{-a_0}W_y\phi_0
 =\sum_{a_0,y}c_{a_0,y}
       \psi_f(-B(a_0,y)/2)\,W_{y-a_0}\phi_0.
 \tag{DO31.9}
\]

These are the same actual signs and central half-phase as AM-DO-029, with the terminal support and period lattices substituted. They are not a choice of an unspecified intertwining scalar.

For completeness, (DO31.8) is proved by fixing a support coset \(a_0+L\) and applying character orthogonality to the finite group \(L/L^-\). Summing \(\psi_f(B(u-u',y))\) over its full character group gives N if \(u=u'\bmod L^-\), and0 otherwise. The vectors \(W_{-a_0}W_y\phi_0\) in (DO31.9) are orthonormal: different a\_0 have disjoint position supports, and for a fixed a\_0 the same finite-character orthogonality applies. Thus

\[
 \|f_{\mathrm{last}}\|^2
 =\sum_{a_0,y}|c_{a_0,y}|^2
 =\frac1N\sum_{t\in L^+/L^-}
                  |\Phi_{\mathrm{last}}(t)|^2.
 \tag{DO31.10}
\]

If \(\delta\Phi_{\mathrm{last}}\) is integral in a coefficient DVR, then \(\delta N c_{a_0,y}\) is integral. Also

\[
 \sum_{a_0,y}|c_{a_0,y}|
 \le\sqrt{RN}\,\|f_{\mathrm{last}}\|.
 \tag{DO31.11}
\]

All finite Heisenberg parameters in (DO31.9) have zero infinite component, so this is an actual IsFinComb-type finite construction on the mixed host. It solves the finite function-equation part of a new mixed datum. It does not manufacture the old elliptic VacIn wrapper or the remaining arithmetic tying records.

\subsection{4. Exact local norms}\label{exact-local-norms}

Use additive Haar measure giving the standard position lattice volume1, equivalently the normalized finite vacuum measure already fixed in AM-DO-017/029.

\subsubsection{4.1 Incidence places}\label{incidence-places}

At T the nonzero residue row v defines a surjective map \(M_{b\times a}(\mathbf F_q)\to\mathbf F_q^a\), \(X\mapsto vX\). Each fibre has \(q^{(b-1)a}\) points. Hence

\[
 \|F_T\|^2=q^{-a}.
 \tag{DO31.12}
\]

\subsubsection{4.2 Q places, with the actual scalar q retained}\label{q-places-with-the-actual-scalar-q-retained}

The top r-by-r block must have determinant a unit; the bottom-left \((b-1-r)\)-by-r block lies in q times the local ring. All other upper entries are unrestricted integral entries. Their support has measure

\[
 q^{-d_P r(b-1-r)}
       \prod_{j=1}^{r}(1-q^{-j}).
\]

The last row has measure \(q^{d_P\mathsf M a}\), and its scalar amplitude in (DO31.4) has square \(q^{-2\mathsf M a}\). Every character value on support has absolute value1. Therefore

\[
 \boxed{\;
 \|F_Q\|^2
 =q^{(d_P-2)\mathsf M a-d_P r(b-1-r)}
       \prod_{j=1}^{r}(1-q^{-j}).
 \;}
 \tag{DO31.13}
\]

In the degree-one, unramified case this specializes to \(q^{-\mathsf M a-r(b-1-r)}
\prod_{j=1}^{r}(1-q^{-j})\). The scalar \(q^{-\mathsf M a}\) in ftY was never replaced by a Haar-normalization guess.

\subsubsection{4.3 G places in the actual sharp case}\label{g-places-in-the-actual-sharp-case}

The upper matrix Z has \(b-1=a+t\) rows. The bottom t rows vanish in the residue field. Its upper a-by-a matrix has block sizes

\[
 (r,\underbrace{t,\ldots,t}_{h\text{ times}}).
\]

The first diagonal block is upper unipotent in the order given by symIdx. In addition the k entries with row \(p<k\) and column \(k+r_S\) are zero. These are k distinct entries that would otherwise be free: symIdx(p)=p+r\_S is strictly less than \(k+r_S\), and that column is fixed by symIdx. Thus the first block contributes \(q^{r(r-1)/2-k}\) choices.

All blocks strictly below the diagonal vanish. The h diagonal t-by-t blocks must be invertible. Indeed the formula contains the inverse of the unit-character evaluation of the determinant of the whole upper a-by-a block; if it is not a unit that value, and hence the local factor, is zero. Since the first block is unipotent and the matrix is block upper triangular, unit determinant is equivalent to invertibility of each of those t-by-t blocks.

The blocks from the first block to the t-blocks contribute \(q^{rht}\), and the blocks above the diagonal among the t-blocks contribute \(q^{t^2h(h-1)/2}\). Thus the number of upper residue matrices on which PhiTop is nonzero is

\[
 q^{\,r(r-1)/2-k+rht+t^2h(h-1)/2}
       |\mathrm{GL}_t(\mathbf F_q)|^h.
 \tag{DO31.14}
\]

This counts exactly the displayed PhiGS conditions. In particular the apparently extra inequalities with a row \(p<r\) in the later block condition place no restriction on the upper rectangular blocks: the natural-number subtraction gives \((p-r)/t=0\), so its strict inequality is false.

The last row has integral support of volume1. The character values have absolute value1. Dividing (DO31.14) by \(q^{(b-1)a}\), and using \(|\mathrm{GL}_t(\mathbf F_q)|
=q^{t^2}\prod_{j=1}^{t}(1-q^{-j})\), gives

\[
 \boxed{\;
 \|F_G\|^2
 =q^{-C_G}
       \prod_{j=1}^{t}(1-q^{-j})^{m-1},\qquad
 C_G=\frac{r(r+1)}2+k+mrt+
                 \frac{m(m-1)}2t^2.
 \;}
 \tag{DO31.15}
\]

The actual source has sharp=true. If one instead uses its separate unsharp PhiG branch, there are no k extra zero entries and the same formula holds with \(C_G-k\). This is a comparison of the two literal branches, not a change to the prescribed data.

For the smallest numerical ranks previously used, \((m,k,t)=(2,6,13)\), one has \((r,a,b)=(7,20,34)\) and \(C_G=385\). Thus the actual G norm square is \(q^{-385}\prod_{j=1}^{13}(1-q^{-j})\). At a degree-one Q-place the norm square is \(q^{-20\mathsf M-182}\prod_{j=1}^{7}(1-q^{-j})\). These finite factors can be very small even though the finite function is nonzero.

\subsubsection{4.4 Global norm and chart scalar}\label{global-norm-and-chart-scalar}

For standard TwDat lattices, the exact finite norm is the product of (DO31.12), (DO31.13), and (DO31.15), over the actual selected places. For a specified raw chart scalar w it is multiplied by \(|w|^2\). If the chart uses different base lattices, its actual additive Jacobian or the equivalent Smith-normal-form formula (DO31.10) must be retained.

The archimedean selected-polynomial norm and its hook and Gaussian factors are separate, as in AM-R-011. A finite norm calculation cannot cancel them by renaming the vacuum.

\subsection{5. An exact nonzero value of the prescribed terminal array}\label{an-exact-nonzero-value-of-the-prescribed-terminal-array}

At every G-place choose the framed local matrix

\[
 X_G=\begin{pmatrix}I_a\\0\end{pmatrix}
 \in M_{b\times a}(E_P).
\]

All diagonal characters have value1, all additive arguments in PhiGS or PhiG are0, and its last row is0. Hence \(F_G(X_G)=1\). At every Q-place choose the same matrix. Its upper top-r determinant is1, its lower-left block is0, and its last row is0. Hence \(F_Q(X_Q)=q^{-\mathsf M a}\). These conclusions hold for every allowed representative \(y'_P\).

At T solve \(vX=u\). This is possible since \(v\ne0\): choose a column vector w with \(vw=1\) and take \(X=wu\). If a full-rank representative is desired, extend w and a basis of \(\ker v\) to construct one, using \(b>a\). The incidence factor is1.

Take sufficiently small local cosets around these matrices, using the periods from Section2 and the actual rational frames. Additive approximation, or a Chinese-remainder argument on a common fractional lattice, supplies a rational matrix \(\gamma\in M_{b\times a}(E)\) in all these cosets. It may simultaneously be chosen integral outside the selected finite places. Concretely, clear the finitely many denominators, impose the finitely many integral congruences by CRT, and cancel the denominator at the remaining primes by additional congruences. Distinct selected prime ideals give comaximal congruence ideals.

The exact local periodicity now proves

\[
 \boxed{\quad
 \Xi_{\mathrm{last}}(\gamma)
       =\prod_{P\in Q}q_P^{-\mathsf M a}\ne0.
 \quad}
 \tag{DO31.16}
\]

For the reference twisting tuple, the source proves \(Q=\varnothing\). Thus this exact raw value is1. With the chart scalar w it is multiplied by w.

The existence assertion is effective once the actual frames, representative and prime ideals are specified: the proof is finite CRT on explicit lattice cosets. It does not assert a height bound for the smallest such gamma; no such bound is needed for (DO31.16), and one is not hidden in the finite norm calculation.

If all these auxiliary residue characteristics are different from11, the unscaled terminal finite vector is primitive over the11-adic coefficient ring. Indeed every table value in (DO31.7) is either0 or a root of unity times \(\prod_Qq_P^{-\mathsf M a}\), hence integral; every denominator N is an11-adic unit. Therefore every Fourier coefficient is integral. At least one is a unit: otherwise Fourier inversion would make every table value nonunit, contradicting (DO31.16). This is primitivity of the finite Schwartz coefficient vector. It is not primitivity of the selected theta image.

\subsection{6. What this closes, and what it leaves open}\label{what-this-closes-and-what-it-leaves-open}

The terminal finite construction is now as explicit as the nonterminal construction in AM-DO-029. Its support/period data, Heisenberg coefficients, phase, norm, denominator and a nonzero raw value are all fixed by formulas. In particular there is no remaining existential obstacle to realizing the prescribed terminal scalar function in the mixed finite oscillator model.

After applying the actual selected archimedean creation polynomial P, its theta reading is the same checked finite sum as in AM-DO-026/029:

\[
 \vartheta_{r_P f_{\mathrm{last}}}(1,1)
 =C^{2ab}\sum_{a_0,y}c_{a_0,y}
    \psi_f(-B(a_0,y)/2)
       [P(\partial)\mathcal U_{y-a_0}F_{00}]_0.
 \tag{DO31.17}
\]

Nonvanishing of (DO31.16) does not prove nonvanishing of (DO31.17). Its terms can cancel, or the vector can lie in the kernel of the specified automorphic projection. The prescribed previous image and terminal arithmetic detector must still be paired through the actual source operator. The remaining MainX requirements include that operator/kernel identity, the integral row normalizations, and the exact terminal mass and tail inequalities. This finding does not substitute an arbitrary nonzero period for those equations.

\subsection{Review status}\label{review-status}

Author whole-proof check: PASS, 2026-11-21T23:13:49.764065+00:00. The complete saved finding and its displayed source formulas were read. The G support count, actual q-scalar at Q, CRT value, Fourier signs and finite primitivity scope were checked. Independent review is separate; no Lean build is claimed.

\clearpage

\section{AM-K028 --- the literal unramified etaT does not exist for the ordinary J² host type}\label{am-k028-the-literal-unramified-etat-does-not-exist-for-the-ordinary-juxb2-host-type}

Author: GPT-6 Astra Pro, ALP\_MACHINE \nolinkurl{/root/kummer_extension}, 21 November 2026 UTC. Inherited reasoning effort is not separately exposed.

For the actual positive host \[
 E'=\mathbb Q(\zeta_{28}),\qquad A=J^2,
 \qquad\Phi'=\{1,3,9,15,17,23\}\pmod{28},
\tag{ET1}
\] there is \textbf{no Hecke character satisfying the source's literal \nolinkurl{IsEtaT}}, when its chosen embeddings are this actual CM type. The obstruction already follows from its \nolinkurl{unr} and \nolinkurl{inf} clauses and from the defining triviality of a Hecke character on principal ideles. It is independent of all theta-lift estimates, the old \nolinkurl{IsTen} issue, and the separate raw-Xi issue.

This statement is about the specified field/type pair. It is not a claim that such a character is impossible for every CM type, every larger host, or every infinity exponent.

\subsection{1. Exact source conventions}\label{exact-source-conventions}

The definition in \nolinkurl{Formalism/LocalToGlobal/Columns/AbelianHostFieldConditions.lean}, lines 308--312, is:

\begin{verbatim}
structure IsEtaT (X : Fields) (σ₀ : X.K →+* ℂ) (χ : HChar X.E) : Prop where
  unitary : χ.IsUnitary
  unr : ∀ v, χ.UnrAt v
  res : χ.ResEta X.F
  inf : χ.HasInf (X.embC σ₀) (fun _ => 1)
\end{verbatim}

The exact imported definitions were recovered and read in \nolinkurl{Formalism/LocalToGlobal/UnitaryNorms/IsUnitaryNormOf.lean}:

\begin{longtable}[]{@{}
  >{\raggedright\arraybackslash}p{(\columnwidth - 2\tabcolsep) * \real{0.5000}}
  >{\raggedright\arraybackslash}p{(\columnwidth - 2\tabcolsep) * \real{0.5000}}@{}}
\toprule\noalign{}
\begin{minipage}[b]{\linewidth}\raggedright
Lines
\end{minipage} & \begin{minipage}[b]{\linewidth}\raggedright
Literal definition or consequence
\end{minipage} \\
\midrule\noalign{}
\endhead
\bottomrule\noalign{}
\endlastfoot
11--19 & \nolinkurl{HChar L} is a continuous homomorphism on the ideles, trivial on \nolinkurl{diagUnits L k} for every \(k\in L^\times\). \\
37 & \nolinkurl{infDiag} is the product of the infinite local embeddings of a principal element. \\
59--65 & \nolinkurl{loc}, \nolinkurl{locK}, and \nolinkurl{inf} are restrictions of that same global character. \\
71 & \nolinkurl{UnrAt v} means the local character is one on every local unit. \\
90--91 & \nolinkurl{HasInf φ n} prescribes the angular exponent n itself. \\
93--94 & \nolinkurl{HasInfQ φ a} instead requires an integer n with \(n=2a\). \\
194--204 & \nolinkurl{TrivOn} and \nolinkurl{ResEta} impose separate restrictions to the real subfield. \\
\end{longtable}

In particular the actual infinity convention is \[
 \chi_\infty(k)=
   \prod_{\tau\in\Phi'}
       \left(\frac{\tau(k)}{|\tau(k)|}\right)^{n_\tau}.
\tag{ET2}
\] The \nolinkurl{IsEtaT} exponent is exactly \(n_\tau=1\), with no factor two and no inverse. The nearby \nolinkurl{IsMu0.inf} uses \nolinkurl{HasInfQ} and has a different convention. Confusing the two would change this test.

\nolinkurl{Chars.Tied.etaT}, in \nolinkurl{AbelianHostFieldConditions.lean}, lines 325--333, actually requires this \nolinkurl{IsEtaT}; it is not an optional later specialization. The separate \nolinkurl{res} and \nolinkurl{unitary} clauses can only restrict the possibilities further.

\subsection{2. The global-unit obstruction}\label{the-global-unit-obstruction}

\subsubsection{Lemma ET3}\label{lemma-et3}

Let L be a number field, let \(\chi:\mathbb A_L^\times/L^\times\to\mathbb C^\times\) be a continuous Hecke character, and assume every finite local component is unramified. Then \[
                    \chi_\infty(\varepsilon)=1
             \qquad(\varepsilon\in\mathcal O_L^\times).
\tag{ET3}
\] Consequently, if \(\chi\) has the angular infinity prescription (ET2), then necessarily \[
       \prod_\tau
       \left(\frac{\tau(\varepsilon)}{|\tau(\varepsilon)|}\right)^{n_\tau}
                           =1
       \qquad(\varepsilon\in\mathcal O_L^\times).
\tag{ET4}
\]

\textbf{Proof.} Put \(U_f=\prod_{v<\infty}\mathcal O_{L_v}^\times\), embedded in the ideles with infinite component one. The direct sum of the local unit groups, namely the subgroup whose entries are one outside a finite set, is dense in this product. The unramified hypotheses kill every factor and hence this dense subgroup. Continuity then gives \(\chi|_{U_f}=1\). For a global unit \(\varepsilon\), its principal idele is the product of its infinite idele and a member of \(U_f\). The global character is one on that principal idele by definition. Dividing by the finite value one proves (ET3), and (ET2) gives (ET4). ∎

The density argument is included because the finite components of a principal root of unity need not be one at all but finitely many primes. One must not replace that idele by an unjustified finite product. Equivalently, the usual Hecke product formula gives (ET3), but no unproved \nolinkurl{ProductRule} field is needed here.

\subsubsection{Theorem ET5}\label{theorem-et5}

For the field and actual CM type (ET1), no character satisfies \nolinkurl{IsEtaT}.

\textbf{Proof.} Apply (ET4) to the global unit \(\varepsilon=\zeta_{28}\). Every archimedean absolute value is one. The actual six embeddings give \[
 \prod_{a\in\Phi'}\zeta_{28}^a
       =\zeta_{28}^{1+3+9+15+17+23}
       =\zeta_{28}^{68}=\zeta_{28}^{12}.
\tag{ET5}
\] This has order \(28/\gcd(28,12)=7\), so it is not one. Lemma ET3 contradicts the simultaneous \nolinkurl{unr} and \nolinkurl{inf} clauses. ∎

There is a useful check using the original primitive CM field. The type \(\Phi'\) restricts to two copies of \(\{1,2,3\}\) modulo 7. Therefore the unit \(\zeta_7\) has prescribed infinity value \(\zeta_7^{2(1+2+3)}=\zeta_7^{12}=\zeta_7^5\ne1\). This is the same order-seven obstruction. No ambiguity in the choice of a primitive 28th root affects it.

Changing the coefficient embedding or applying a Galois conjugation to the whole field/type pair multiplies the exponent 12 by a unit modulo 28, so the nontrivial value still has order seven. Complex conjugation merely inverts it. A different choice of integral frame or a norm-one theta translate does not affect a global infinity-type obstruction.

\subsection{3. What twisting or taking powers can and cannot change}\label{what-twisting-or-taking-powers-can-and-cannot-change}

Every everywhere-finite-unramified character with \textbf{uniform integer angular exponent n} on this same type must satisfy \[
                     28\mid68n,
             \quad\text{equivalently}\quad 7\mid n.
\tag{ET6}
\] Thus the literal exponent one is excluded. Taking a seventh power can kill this particular root-unit value, but changes the infinity exponent to seven and therefore does not produce the prescribed \nolinkurl{IsEtaT.inf}. Taking a root to restore exponent one would restore the same contradiction if it were still unramified everywhere finite.

Similarly, multiplying by an everywhere-unramified finite-order character cannot compensate (ET5): that character has trivial value on all finite global-unit components and trivial continuous angular infinity type. Any Hecke character with the prescribed angular type one must have some finite ramification. Allowing such ramification is a change to the literal \nolinkurl{unr} clause, with additional local and global data to control; it is not a solution of the same source record.

These statements concern a necessary unit condition. They do not assert that the exponent-seven character exists with every remaining source restriction merely because (ET6) is then satisfied. Other units and the exact restriction \nolinkurl{ResEta} must also be checked.

\subsection{4. Ordinary J powers and a type-preserving extension}\label{ordinary-j-powers-and-a-type-preserving-extension}

\subsubsection{Proposition ET7. Compatible full CM actions on ordinary J powers}\label{proposition-et7.-compatible-full-cm-actions-on-ordinary-j-powers}

Suppose a full CM field L of degree \(6s\) acts on \(J^s\), contains the specified central cyclotomic field \(E=\mathbb Q(\zeta_7)\), and its chosen CM type \(\Psi\) is the actual holomorphic type of that ordinary product. Any Hecke character of L that is unramified at every finite place and has uniform angular exponent one on \(\Psi\) necessarily satisfies \[
                                7\mid s.
\tag{ET7}
\]

\textbf{Proof.} On \(H^0(J^s,\Omega^1)\), each of the three E-characters \(1,2,3\) occurs with multiplicity s. On restricting the full CM type \(\Psi\) to E, its embeddings therefore have exactly those multiplicities. The prescribed angular value on the global unit \(\zeta_7\in E\subset L\) is \[
 \prod_{\tau\in\Psi}\tau(\zeta_7)
                    =\zeta_7^{s(1+2+3)}=\zeta_7^{6s}.
\] Lemma ET3 says this value is one. Since 6 is invertible modulo 7, one obtains \(7\mid s\). ∎

This applies to compatible maximal CM subfields in the standard matrix endomorphism algebra of an ordinary J power. It does not apply unchanged to products with conjugated E-actions, to a type unrelated to the actual Hodge representation, or to a CM field that does not contain the specified central E. For the host (ET1), s=2, recovering the contradiction above. Passing merely from J² to another ordinary multiplicity prime to seven cannot repair it.

\subsubsection{Totally real extension of the current host}\label{totally-real-extension-of-the-current-host}

Let \(F'=(E')^+\), let \(\widetilde F/F'\) be totally real of degree e, and set \(\widetilde E=E'\widetilde F\). Since every real element of E′ lies in F′, one has \([\widetilde E:E']=e\). Give \(\widetilde E\) the CM type induced from \(\Phi'\), consisting of every lift of each chosen embedding. The old root \(\zeta_{28}\) then has prescribed angular value \[
                  \zeta_{28}^{68e}.
\tag{ET8}
\] Hence any literal unramified type-one character on this induced pair necessarily has \[
                               7\mid e.
\tag{ET9}
\] For degrees prime to seven the obstruction persists, including extensions of degrees that are powers of two. This is also Proposition ET7 with the resulting ordinary multiplicity \(s=2e\). Degree seven kills this \textbf{particular} obstruction without changing the exponent one on the enlarged CM type.

There is an elementary explicit degree-seven candidate extension: let B be the unique degree-seven subfield of \(\mathbb Q(\zeta_{29})\). Its fixing subgroup in the cyclic group of order 28 has order four and contains complex conjugation, so B is totally real. Since \([F':\mathbb Q]=6\) is coprime to 7, \(F'\cap B=\mathbb Q\); therefore \(\widetilde F=F'B\) has degree seven over F′. For this candidate, (ET8) is one.

Nothing here proves that all unit angular values on \(\widetilde E\) are trivial, that an unramified character with \nolinkurl{ResEta} exists, or that the enlarged CM lattice satisfies the required source frame and polarization conditions. New units of the extension, its class group, the finite restrictions, and the enlarged arithmetic host remain exact tasks. This is a narrowly identified possible repair, not a completed one.

Nor is the obstruction universal among CM types on E′ itself. For example the type induced from the quadratic residue choice modulo seven, \[
                   \{1,9,11,15,23,25\}\pmod{28},
\] has exponent sum 84, so this root-unit test is trivial for it. That is a different type from the ordinary primitive-J² type in (ET1); passing to it does not preserve the geometric construction already checked here. No existence conclusion is inferred from its passing this one test.

\subsection{5. Consequence for the current machine adapter}\label{consequence-for-the-current-machine-adapter}

The positive ordinary host, its exact theta germ and canonical gauge, its selected compact Cauchy shapes, and its permitted finite-jet arithmetic are unaffected as independently stated results. But on this unchanged field/type pair, the literal \nolinkurl{Chars.Tied} record cannot be completed because its required \nolinkurl{etaT} does not exist. This is stronger than merely leaving its construction open. It does not question the accepted source theorem, whose hypotheses concern its own arithmetic host.

Accordingly, AM-K027's exact selected-family slope and the class-supported constructions remain results in the actual indicated positive local/global model, with the full source arithmetic adapter still absent. A valid eventual adaptation must change or enlarge the host/type or explicitly replace the unramified auxiliary-character interface and account for that change. It cannot fill this field with a conventionally normalized character on the same pair.

\subsection{6. Provenance and checks}\label{provenance-and-checks}

The source definition of \nolinkurl{IsEtaT} was already present in the intake-pinned recovery. The additional imported file \nolinkurl{IsUnitaryNormOf.lean} was retrieved in this run from the exact pinned mirror commit \nolinkurl{7cbf7a0c32fca1badf05931879a37c3f8a2b03b4}:

\href{https://raw.githubusercontent.com/CocoSgt/riemann-hypothesis-lean/7cbf7a0c32fca1badf05931879a37c3f8a2b03b4/src/Formalism/LocalToGlobal/UnitaryNorms/IsUnitaryNormOf.lean}{\nolinkurl{https://raw.githubusercontent.com/CocoSgt/riemann-hypothesis-lean/7cbf7a0c32fca1badf05931879a37c3f8a2b03b4/src/Formalism/LocalToGlobal/UnitaryNorms/IsUnitaryNormOf.lean}}

Its SHA256 is \nolinkurl{503df2de36607142871ffe2399c86c903953535658077b7e06b4f6ce593d6739}. It is \textbf{absent from the supplied selected author-file manifest}, so this run does not claim a new byte comparison of that member with the author's original archive. The precise retrieval and that null comparison status are recorded in \nolinkurl{work/research/source_recovery_20261121T231352Z.json}. The infinity convention is directly visible in the recovered source, and the mathematical obstruction is proved above from that displayed convention.

The actual host/type input is accepted DP57 and AM-DO-022 (DO22.1)--(DO22.3). The only new mathematical input is the elementary global-unit argument and its explicitly scoped consequences. The exact residue sums and orders have a separate arithmetic certificate. No Lean compilation, unramified-character existence theorem, or full enlarged-host construction is claimed.

\clearpage

\section{ALP-LIFT-17 --- prescribed column types exclude the spherical chain}\label{alp-lift-17-prescribed-column-types-exclude-the-spherical-chain}

Author: /root/lift\_accounting, GPT-6 Astra Pro, ALP\_MACHINE, 21 November2026UTC.

This result identifies an actual local representation obstruction to placing the chain of ALP-LIFT-14--16 in the source's explicit column modules. It is stronger than the observation that a generic-unipotent projector kills a spherical vector: at every noninitial, nonterminal stage, the specified determinant-block character is absent from the entire unramified local representation. From stage2 onward, even a tame determinant twist of an unramified representation cannot have the prescribed type.

The scope is the literal ColV/colLev\_j family. The later ChIdx.levG\_j is a separate variable family and is not identified with colLev\_j by the source record. No annihilation theorem for every raw Xi entry, or no-go theorem for a different ramified theta chain, is asserted.

\subsection{1. A finite flag lemma, with full proof}\label{a-finite-flag-lemma-with-full-proof}

Let \(\kappa\) be a finite field, \(G=\operatorname{GL}_n(\kappa)\), and let \(B\) be its standard upper triangular subgroup. Fix a block decomposition \[
 n=r+n_1+\cdots+n_s,\qquad n_i\ge1,\quad s\ge1.
 \tag{F1}
\] Let \(P=LN\) be the standard parabolic with Levi \[
 L=\operatorname{GL}_r(\kappa)\times
                       \prod_{i=1}^s\operatorname{GL}_{n_i}(\kappa).
\] For any subgroup \(H_0\subseteq\operatorname{GL}_r(\kappa)\), put \[
 H=(H_0\times\prod_{i=1}^s\operatorname{GL}_{n_i}(\kappa))N.
\] Let \(\lambda_0:H_0\to\mathbf C^\times\) and \(\chi_i:\kappa^\times\to\mathbf C^\times\) be characters, and define \[
 \lambda(h_0,h_1,\ldots,h_s;n)
            =\lambda_0(h_0)\prod_{i=1}^s\chi_i(\det h_i).
 \tag{F2}
\] Thus the character is trivial on the parabolic radical \(N\). If at least one \(\chi_i\) is nontrivial, then \[
 \operatorname{Hom}_H
       (\lambda,\mathbf C[G/B])=0.
 \tag{F3}
\]

Here \(\mathbf C[G/B]\) is the permutation representation on complete flags. To prove(F3), use the finite Bruhat decomposition \(G=PWB\), with permutation matrices \(W\). Since \(H\) contains \(N\) and every Levi factor except the first, every \(H\)-orbit on \(G/B\) has a representative \[
                           \ell_0 wB,
 \tag{F4}
\] where \(\ell_0\) acts only on the first \(r\)-block. The equality \(P=H\operatorname{GL}_r(\kappa)\) follows directly by writing a parabolic matrix as a block diagonal matrix times \(N\) and conjugating \(N\) past the first block.

Choose an index \(i\) with \(\chi_i\ne1\). Let \(D_i\) be the diagonal torus in the \(i\)-th determinant block, with identity in every other block. It is contained in \(H\), commutes with \(\ell_0\), and is conjugated by \(w\) into the full diagonal torus of \(B\). Consequently \[
           D_i\subset H\cap \ell_0wBw^{-1}\ell_0^{-1}.
 \tag{F5}
\] The determinant map \(D_i\to\kappa^\times\) is surjective, so \(\lambda|_{D_i}\ne1\). A \(\lambda\)-eigenfunction on the orbit of(F4) must vanish there: its value at the representative is both unchanged by the stabilizer and multiplied by a nontrivial value of \(\lambda\). It then vanishes on the entire orbit. This holds on every orbit, proving(F3).

The same conclusion holds with every character inverted. No assumption about \(\lambda_0\), genericity, or the order of the finite field is used.

For later use, write the character projector as \[
            e_\lambda=|H|^{-1}\sum_{h\in H}\lambda(h)^{-1}h.
 \tag{F6}
\] Equation(F3) says that this operator is zero on the full flag permutation module. It is therefore zero on every submodule and quotient of that module. If \(Q\) is any parabolic subgroup, the permutation module \(\mathbf C[G/Q]\) embeds in \(\mathbf C[G/B]\), after conjugating \(Q\) to contain \(B\), by pullback of functions along \(G/B\to G/Q\). Hence(F6) is also zero on every quotient of every parabolic permutation module.

This proof uses only the finite Gaussian-elimination/Bruhat decomposition and the explicit torus in(F5). It does not invoke a classification of finite unipotent representations.

\subsection{2. Every unramified local representation is excluded}\label{every-unramified-local-representation-is-excluded}

Let \(F_v\) be a nonarchimedean field, with ring \(\mathcal O\), uniformizer \(\varpi\), residue field \(\kappa\), and put \[
 K=\operatorname{GL}_n(\mathcal O),\qquad
 K_1=1+\varpi M_n(\mathcal O).
\] Let \(\pi\) be an irreducible smooth complex representation with a nonzero \(K\)-fixed vector \(v\). Its entire space is generated by the translates of \(v\). We prove directly that \(\pi^{K_1}\), as a \(K/K_1\)-module, is spanned by quotients of the parabolic permutation modules in section1.

Let \(e_{K_1}\) be averaging with probability Haar measure. Smith normal form gives \[
       \operatorname{GL}_n(F_v)=
         \bigcup_{\mathbf a}K\operatorname{diag}
                (\varpi^{a_1},\ldots,\varpi^{a_n})K.
\] Thus \(\pi^{K_1}\) is spanned by the \(K\)-translates of \[
 w_{\mathbf a}
       =e_{K_1}\pi(\operatorname{diag}(\varpi^{a_i}))v.
 \tag{F7}
\] This follows by expressing any vector as a finite linear combination of translates of \(v\), then averaging. No topological completion of the representation is being taken.

Write \(a=\operatorname{diag}(\varpi^{a_i})\). The subgroup \(K\cap aKa^{-1}\) fixes \(\pi(a)v\). Because \(K_1\) is normal in \(K\), its averaging operator commutes with \(K\), so that subgroup also fixes \(w_{\mathbf a}\). Its reduction modulo \(\varpi\) is a parabolic subgroup: entry \((i,j)\) must have valuation at least \(\max(0,a_i-a_j)\), and the resulting block zero conditions are exactly the flag stabilizer for the equal-exponent blocks. Every matrix in that parabolic has a lift satisfying these conditions, by taking the required off-block entries to be zero. Therefore the \(K/K_1\)-orbit of \(w_{\mathbf a}\) is a quotient of \[
                   \mathbf C[\operatorname{GL}_n(\kappa)/Q_{\mathbf a}]
 \tag{F8}
\] for a parabolic \(Q_{\mathbf a}\).

Inflate \(H,\lambda\) from section1 to their full inverse image \(\widetilde H\subset K\). Since \(K_1\subset\widetilde H\) and the inflated character is trivial on \(K_1\), every \(\lambda\)-eigenvector belongs to \(\pi^{K_1}\). By(F6)--(F8), its character projector vanishes on every spanning submodule. We obtain \[
 \boxed{\pi^{\widetilde H,\lambda}=0
          \quad\hbox{whenever some }\chi_i\ne1.}
 \tag{F9}
\] This excludes the type from the entire representation, including all its nonspherical vectors and all its translates. It does not assume that \(\pi\) is tempered or generic.

The proof also works for any module algebraically generated by \(K\)-fixed vectors. In particular it applies constituent by constituent, with actual automorphic multiplicities, to the away-from29 spherical module in ALP-LIFT-16.

\subsection{3. Tame determinant twists do not fix distinct blocks}\label{tame-determinant-twists-do-not-fix-distinct-blocks}

Let \(\eta:F_v^\times\to\mathbf C^\times\) be a tame character, so it is trivial on \(1+\varpi\mathcal O\). Write \(\overline\eta\) for its residue character. For \(\pi=\pi_{\mathrm{sph}}\otimes(\eta\circ\det)\), a vector of \(\widetilde H\)-type \(\lambda\) would give in \(\pi_{\mathrm{sph}}\) a vector of character \[
       \lambda\cdot(\overline\eta\circ\det)^{-1}.
 \tag{F10}
\] If \(H_0\) is upper unitriangular, its determinant is one; the determinant-block characters in(F10) are \[
                           \chi_i\overline\eta^{-1}.
\] By(F9) they must all be trivial. Thus a necessary condition is \[
                 \chi_1=\cdots=\chi_s=\overline\eta.
 \tag{F11}
\] In particular, two distinct prescribed determinant-block characters exclude every tame determinant twist of every unramified representation. This statement is deliberately restricted to tame twists; no claim about arbitrary wild twists is inferred from the finite reduction argument.

\subsection{4. Exact source identification of the finite type}\label{exact-source-identification-of-the-finite-type}

All paths are relative to the recovered source work/research/alpoge\_source, pinned to commit 7cbf7a0c32fca1badf05931879a37c3f8a2b03b4. The following were read directly.

\begin{longtable}[]{@{}
  >{\raggedright\arraybackslash}p{(\columnwidth - 4\tabcolsep) * \real{0.3333}}
  >{\raggedright\arraybackslash}p{(\columnwidth - 4\tabcolsep) * \real{0.3333}}
  >{\raggedright\arraybackslash}p{(\columnwidth - 4\tabcolsep) * \real{0.3333}}@{}}
\toprule\noalign{}
\begin{minipage}[b]{\linewidth}\raggedright
Source
\end{minipage} & \begin{minipage}[b]{\linewidth}\raggedright
Exact locator
\end{minipage} & \begin{minipage}[b]{\linewidth}\raggedright
Role
\end{minipage} \\
\midrule\noalign{}
\endhead
\bottomrule\noalign{}
\endlastfoot
\nolinkurl{Formalism/LocalToGlobal/Columns/GenericUnipotentCharacter.lean} & 155--180 & uniUT and its block upper triangular description \\
same & 232--256 & supChar, detChar, lamGen, extChi \\
same & 437--452 & HU, HX, SX and lamX \\
\nolinkurl{Formalism/LocalToGlobal/Columns/ColForms.lean} & 92--105 & colLev and ColV \\
\nolinkurl{Formalism/LocalToGlobal/Columns/ColFormsLemmas.lean} & 137--145 & colLev\_of\_lt and the terminal branch \\
\nolinkurl{Formalism/LocalToGlobal/Columns/AbelianHostFieldConditions.lean} & 273--306 & Chars, chi\_inj, chi\_om \\
same & 205--224;488 & split degree-one G-place conditions, distinct primes, and g\_big \\
same & 315--323 & IsMu0 finite restrictions and tameness \\
\nolinkurl{Formalism/LocalToGlobal/UnitaryNorms/IsUnitaryNormOf.lean} & 66--75 & LevelLe, UnrAt, TameAt \\
\nolinkurl{Formalism/LocalRingCardinality/Params.lean} & 5--20;30--40 & ranks and \(N_j=k+j(t+1)\) \\
\nolinkurl{Formalism/LocalToGlobal/Columns/FiniteSumNonzeroLargeDepth.lean} & 1017--1023 & only the initial levG is tied to colLev by zero\_G \\
\end{longtable}

For \(1\le j<m\), put \(n=N_j=k+j(t+1)\) and \(r_j=k+j\). The literal source has \[
 \begin{split}
 H_{X,j}&=\operatorname{uniUT}
                (n,\ r_j,\ [t,\ldots,t]),\\
 n&=r_j+jt,\\
 \lambda_{X,j}
       &=\psi_{S_X}\prod_{i=0}^{j-1}\chi_i(\det h_i).
 \end{split}
 \tag{F12}
\] The \(j\) determinant blocks have size \(t\), and the first block is \(U_{r_j}\). The character \(\psi_{S_X}\) is supported on the displayed selected simple-root entries of that first unitriangular block. It is trivial on the interblock radical. Thus(F12) is exactly the situation of(F1)--(F2), with \(H_0=U_{r_j}\). The omission of one simple root in \(S_X\) at some stages does not affect the argument, since \(\lambda_0\) in the finite lemma was arbitrary.

There is no additional unused trailing block: the identity \(n=r_j+jt\) is exact. The terminal declaration at \(j=m\) uses HXm and an additional fixed-column condition, so the theorem is not extended to that branch without a separate calculation.

After choosing the coefficient-field embedding into \(\mathbf C\), the characters have the same orders as in the source field. The automorphic convention may require \(\lambda_{X,j}\) or its inverse; both are excluded by the same proof.

\subsection{5. Every prescribed determinant character is nontrivial}\label{every-prescribed-determinant-character-is-nontrivial}

Let \(P\) be a selected G-place over the rational auxiliary prime \(q=g_\epsilon\), \(\epsilon\in\{0,1\}\). The source has residue field of cardinality \(q\), where \(q\) is prime. The prime-field map therefore identifies its residue units with \(\mathbf F_q^\times\).

The chi\_om condition specifies a faithful character \(\omega:\mathbf F_q^\times\to K^\times\) such that \[
                        \chi_i=\omega^{\,2(i+1)}
                         \qquad(0\le i<m).
 \tag{F13}
\] This is the exact source character, not a new choice. Its order is \[
                \operatorname{ord}(\chi_i)
                   =\frac{q-1}{\gcd(q-1,\,2(i+1))}.
 \tag{F14}
\] The source size and rank inequalities give \[
       q-1>2fk\ge2k>2m\ge2(i+1)>0.
 \tag{F15}
\] Here \(f\ge1\), and \(k\ge m+3\). Hence every \(\chi_i\) is nontrivial, and the characters for distinct \(i<m\) are distinct. The latter also appears as the separate chi\_inj field. This proves the nontriviality hypothesis of(F9) at every G-place for every \(1\le j<m\).

The conclusion for the literal column type is therefore \[
 \boxed{\pi^{\widetilde H_{X,j},\lambda_{X,j}^{\pm1}}=0
   \quad(1\le j<m)}
 \tag{F16}
\] for every unramified local representation at such a place. For \(2\le j<m\), the same vanishing holds for every tame determinant twist of an unramified representation, by the distinctness and(F11).

For the actual host \(E'=\mathbf Q(\zeta_{28})\), \(f=6\) and the admissible even rank has \(k\ge6\). Thus \(q-1>72\), so both prescribed G primes are different from29. Every local constituent of the chain in LIFT14--16 is unramified at those places. Consequently its nonzero stage-\(j\) module cannot itself lie in the explicit source ColV with this finite type, for \(1\le j<m\). Deepening the29 level, choosing a nonspherical vector in the same local representation, or taking any further finite group translate within that representation cannot remove(F16).

This is a local representation obstruction, not an inequality between gains.

\subsection{6. The initial stage and the limits of the obstruction}\label{the-initial-stage-and-the-limits-of-the-obstruction}

At \(j=0\), the source has \(H_{X,0}=U_k\) and its full nontrivial generic-unipotent character, because \(m\ge2\). There are no determinant blocks. A \(K\)-fixed vector is fixed by \(U_k\); its nontrivial character projector is therefore zero. A determinant twist does not change this calculation on \(U_k\).

This initial-vector statement does not imply that every vector of an unramified representation is killed. The finite flag permutation module can have a generic-unipotent component, and section1 gives no exclusion when there are no nontrivial determinant blocks. A source-type projection or different vector in the same representation must therefore be tested separately at \(j=0\).

The standard initial class-supported vector of AM-R-012 is spherical at the auxiliary G-places, so it is not itself a literal initial generic-type vector there. This is distinct from the stronger entire-representation obstruction at \(j\ge1\).

The source's ChIdx record requires only zero\_G: \[
                            \mathrm{levG}_0=\mathrm{colLev}_0.
 \tag{F17}
\] Its later levG\_j and conjugating indices are additional data. Therefore(F16) must not be promoted to a theorem that every permitted raw PhiG/PhiGS or Xi kernel annihilates every vector in the existing chain. That stronger assertion would require its actual right-source covariance and the identification of its later finite level. The direct raw-Xi identity and auxiliary finite interpolation in AM-DO-029/030 keep that question separate.

\subsection{7. Exact finite conductor forced by the source character pin}\label{exact-finite-conductor-forced-by-the-source-character-pin}

The existing source IsMu0 condition already requires a different finite character family from the conductor29 family in LIFT14. At each selected G-place its unit restriction is \[
                         \mu_i|_{\mathcal O_{E',P}^{\times}}
                                  =\chi_i^{-1}.
 \tag{F18}
\] It is tame there and unramified away from the two G primes. This was also audited in AM-DO-030; it is a literal source condition, not a new character-existence theorem.

The following conductor accounting follows from(F13)--(F18). The source atg condition, applied to integral representatives of each nonzero residue class, gives the displayed unit restriction; tameness makes it a character of precisely that finite quotient. At every selected G-place the character is therefore nontrivial on units but trivial on principal units, so its conductor exponent is exactly one. At the conjugate place the restriction is the inverse, because \(\mu_i\) is trivial on the diagonal base-field ideles; it is also nontrivial and has conductor exponent one.

For completeness, the degree-one prime, norm, cardinality and coprimality conditions really do account for every place above each \(g_\epsilon\). The two principal ideals generated by \(\pi_G\) and its conjugate are coprime, have product \((g_\epsilon)\), and have equal norms. Each norm is therefore \(g_\epsilon^f\). Each selected half already contains \(f\) distinct prime ideals of norm \(g_\epsilon\), so all their exponents are one and no additional prime factor is possible. Conjugation gives the other half.

Thus any genuine realization of this source character pin on the actual degree12 host has \[
 \mathfrak f(\mu_i)=(g_0g_1)\mathcal O_{E'},\qquad
                  N_{E'/\mathbf Q}\mathfrak f(\mu_i)=(g_0g_1)^{12}.
 \tag{F19}
\] The equality is conditional on the genuine source IsMu0 realization; it does not assert that the global unit compatibility needed for that realization has been solved. It records the exact finite conductor required if it is solved.

This is useful normalization data, but it is not by itself a lower bound for the full H term. A unitary determinant character can act isometrically in a suitable class presentation, and the changed kernels, levels, local tests and arithmetic content must still be compared. AM-R-008's integral isometry has its own hypotheses and does not identify the altered row matrices.

\subsection{8. What a repair would have to change}\label{what-a-repair-would-have-to-change}

A compatible new chain must admit the prescribed nontrivial local determinant blocks at the G-places. The current spherical constituents cannot acquire them merely by selecting vectors or deepening level elsewhere. From stage2 onward, one tame scalar determinant twist cannot replace the distinct blocks either.

A natural remaining route is to use globally admissible stage characters ramified as in(F18) and local theta constituents with their correspondingly ramified inserted blocks. Type-II local theta occurrence itself remains available at these split places. That observation does not prove the required local type, the exact raw-Xi action on the previous image, or global nonvanishing with the prescribed finite kernel. It also does not solve the independently audited odd-character etaT pin.

The full new result is(F16), its tame-twist strengthening, and the exact conditional conductor(F19). It neither disproves a different actual mixed theta tower nor supplies the missing \(g_j,e_j,H\) arrays, terminal row, six arithmetic identities, or a shorter DP44 chain.

\clearpage

\section{AM-K029 --- a degree-seven repair of the literal unramified etaT character}\label{am-k029-a-degree-seven-repair-of-the-literal-unramified-etat-character}

Author: GPT-6 Astra Pro, ALP\_MACHINE \nolinkurl{/root/kummer_extension}, 21 November 2026 UTC. Inherited reasoning effort is not separately exposed.

The root-unit obstruction of AM-K028 can be repaired on an explicit larger field while keeping the ordinary primitive-J Hodge type. The following construction proves the full character interface, including the base-field restriction \nolinkurl{ResEta}; it does not merely remove one necessary unit test.

Write \[
 E_0=\mathbb Q(\zeta_{28}),\quad F_0=E_0^+,
 \quad\Phi_0=\{1,3,9,15,17,23\},
\] and let B be the unique degree-seven subfield of \(\mathbb Q(\zeta_{29})\). Put \[
                         K=E_0B,\qquad F=F_0B.
\tag{ER1}
\] Give K the induced CM type \(\Psi\), containing all seven lifts of each member of \(\Phi_0\).

\subsubsection{Theorem ER2}\label{theorem-er2}

The extension K/F is CM and unramified at every finite place, \([K:\mathbb Q]=84\), and \([F:\mathbb Q]=42\). There exists a continuous unitary Hecke character \[
                   \eta:\mathbb A_K^\times/K^\times\longrightarrow S^1
\] that is unramified at every finite place, has infinity component \[
        \eta_\infty(z)=\prod_{\tau\in\Psi}\frac{z_\tau}{|z_\tau|},
\tag{ER2}
\] and restricts to the quadratic character \(\eta_{K/F}\) on \(\mathbb A_F^\times\). Thus it satisfies the literal \nolinkurl{IsEtaT} convention checked in AM-K028. Its normalized finite values \[
                 \eta(\varpi_P)\,(NP)^{-1/2}
\tag{ER3}
\] all lie in a single finite algebraic coefficient field, which may be enlarged to contain its complex conjugates. Therefore the separate \nolinkurl{etaVal} field can also be supplied after that coefficient-field enlargement.

No class-number calculation for the degree-84 or degree-42 fields is needed for this existence proof. No bound on the degree, discriminant, or heights of the auxiliary coefficient field is asserted. The full polarized host, original \nolinkurl{IsTen}, all source size predicates, and later selected-image bounds are separate conditions.

\subsection{1. The field and its roots of unity}\label{the-field-and-its-roots-of-unity}

The group \((\mathbb Z/29\mathbb Z)^\times\) is cyclic of order 28. Its subgroup of order four contains complex conjugation and fixes B; thus B is totally real, cyclic of degree seven, and ramified only at 29, with ramification index seven there. The fields E₀ and B have coprime degrees 12 and 7, so are linearly disjoint. Hence K has degree 84, and its conjugation fixes exactly F, of degree 42.

Accepted DP57 proves that E₀/F₀ is unramified at every finite place. This property is preserved by base change. More concretely, locally a finite étale extension of a DVR stays étale after base change, and its finite field components give the completions of K/F. Hence K/F is everywhere finite-unramified.

The rational ramification indices in K are \[
                     e_2=2,\qquad e_7=6,\qquad e_{29}=7,
\tag{ER4}
\] and all other finite rational primes are unramified. These follow from the standard cyclotomic indices in \(\mathbb Q(\zeta_{28})\), the totally ramified degree-seven extension at 29, and their disjoint ramification sets.

It follows that \[
                                  \mu(K)=\mu_{28}.
\tag{ER5}
\] Indeed a primitive \(\ell^a\)-th root in K would force the cyclotomic ramification index \(\varphi(\ell^a)\) to divide the corresponding index in (ER4). No odd prime other than 7 can occur: at 29 the required index 28 does not divide 7, and all remaining odd primes are unramified. A 49th root would require index 42 at 7, and an eighth root would require index four at 2. Thus the root order divides 28; E₀ already contains all these roots. The same argument gives \(\mu(E_0)=\mu_{28}\).

For a CM field L and a unit u, \(u/\bar u\) is a root of unity. Here is the elementary fact used: all conjugates of this algebraic integer unit have absolute value one. The characteristic polynomials of its positive powers have bounded integer coefficients, so only finitely many such powers can occur; two powers are equal, proving finite order. This also gives the usual CM-unit index bound below without a unit-group classification theorem.

\subsection{2. The entire unit phase image in E₀}\label{the-entire-unit-phase-image-in-eux2080}

Define \[
 a_0(x)=\prod_{\phi\in\Phi_0}\frac{\phi(x)}{|\phi(x)|}
                       \quad(x\in E_0^\times).
\tag{ER6}
\] We prove the exact identity \[
                             a_0(\mathcal O_{E_0}^\times)=\mu_7.
\tag{ER7}
\]

First every real unit \(\varepsilon\in\mathcal O_{F_0}^\times\) has phase one. The canonical quadratic Artin character of the known extension E₀/F₀ is unramified at every finite place and has local sign at every real place. On a principal real unit, its finite value is one, so global reciprocity gives \(\operatorname{sgn}N_{F_0/\mathbb Q}\varepsilon=1\). This sign is exactly \(a_0(\varepsilon)\). This uses the quadratic character on \textbf{F₀}, which exists by class field theory; it does not presume the obstructed \nolinkurl{etaT} character on E₀.

Next put \(u_0=1-\zeta_{28}\). It is a unit because \[
 \Phi_{28}(X)=X^{12}-X^{10}+X^8-X^6+X^4-X^2+1,
                      \qquad \Phi_{28}(1)=1.
\] Its conjugation quotient is \[
                       \frac{u_0}{\bar u_0}=-\zeta_{28}=\zeta_{28}^{15}.
\tag{ER8}
\] For an arbitrary unit u, write \(u/\bar u=\zeta_{28}^a\). Choose \(\epsilon\in\{0,1\}\) with \(a-15\epsilon\) even and choose j with \(2j\equiv a-15\epsilon\pmod{28}\). Then \(u/(\zeta_{28}^j u_0^\epsilon)\) is fixed by conjugation and is an integral unit. Thus \[
 \mathcal O_{E_0}^\times
   =\mu_{28}\,\mathcal O_{F_0}^\times\,\{1,u_0\}.
\tag{ER9}
\] This proves exactly the needed Hasse-index-two description; no missing generator is assumed.

On \(\mu_{28}\), the phase is \(a_0(\zeta_{28}^j)=\zeta_{28}^{68j}\), whose image has order seven. For \(0<a<28\), \[
 \frac{1-e^{2\pi ia/28}}{|1-e^{2\pi ia/28}|}
                             =-i e^{\pi ia/28}.
\] Using the six integers in \(\Phi_0\), this gives \[
                    a_0(u_0)=(-i)^6\zeta_{56}^{68}=\zeta_7^5.
\tag{ER10}
\] Equations (ER9)--(ER10), and the real-unit result, prove (ER7).

\subsection{3. All new units of K satisfy the required phase condition}\label{all-new-units-of-k-satisfy-the-required-phase-condition}

Let \[
 a(x)=\prod_{\tau\in\Psi}\frac{\tau(x)}{|\tau(x)|}
                   =a_0(N_{K/E_0}x).
\tag{ER11}
\] The equality follows by multiplying over all lifts of each embedding; the absolute values multiply in the same way.

For a unit U of K, put \(v=N_{K/E_0}U\), an E₀-unit. By (ER5), write \(U/\bar U=\zeta\in\mu_{28}\). Since this root lies in E₀, \[
                       \frac v{\bar v}=\zeta^7.
\] Consequently \[
 a(U)^2=a_0(v)^2
       =\prod_{\phi\in\Phi_0}\phi(v/\bar v)
       =a_0(\zeta^7)=1.
\tag{ER12}
\] On the other hand, (ER7) gives \(a(U)=a_0(v)\in\mu_7\). The only element common to \(\mu_2\) and \(\mu_7\) is one. Therefore \[
                            a(\mathcal O_K^\times)=\{1\}.
\tag{ER13}
\] This treats every new unit in the degree-seven extension, not only the roots inherited from E₀.

\subsection{4. Extending a unit-compatible angular type to an unramified character}\label{extending-a-unit-compatible-angular-type-to-an-unramified-character}

We use the following elementary ideal construction. Let L be a CM field, let b be a unitary angular character of \((L\otimes\mathbb R)^\times\), and suppose \(b(\mathcal O_L^\times)=1\). On the principal fractional ideals define \[
                         c((x))=b(x)^{-1}.
\tag{ER14}
\] It is well-defined by the unit condition. It extends to a homomorphism \(c:I_L\to S^1\) on all fractional ideals. To see this without an abstract extension theorem, write the finite class group as a direct sum of cyclic groups, choose representative ideals \(\mathfrak a_j\) of their generators and relations \(\mathfrak a_j^{h_j}=(x_j)\), and choose roots \(c(\mathfrak a_j)^{h_j}=b(x_j)^{-1}\) in \(S^1\). Every ideal has a principal remainder after choosing the canonical ranges of the generator exponents. The displayed root equations make its value independent of carries and make multiplication exact.

For a finite idele \(z_f\), let \(I(z_f)=\prod_PP^{v_P(z_P)}\). Then \[
                         \chi(z)=b(z_\infty)c(I(z_f))
\tag{ER15}
\] is continuous, unitary, unramified at all finite places, and trivial on L× by (ER14). The ideal-valued map is locally constant, and b is continuous, so continuity has not been omitted. A character with trivial infinity component and unramified finite components similarly factors through the ordinary finite ideal class group.

By (ER7), this constructs an unramified unitary character \(\chi_7\) on E₀ with angular type \(a_0^7\). By (ER13), it constructs a character \(\chi_0\) on K with angular type a. The latter already has all desired clauses except the exact restriction to F.

\subsection{5. Repairing the restriction without a large class-number calculation}\label{repairing-the-restriction-without-a-large-class-number-calculation}

The only numerical class-number input is \[
                                  h(F_0)=1.
\tag{ER16}
\] This is explicitly tabulated in Masley (1978), Table 3, printed p.316: conductor 28, degree 6, class number 1. I inspected that primary table image. The field is the maximal real subfield of E₀, whose Galois group is cyclic of order six. No conjectural class-number bound is used.

Let \(\eta_0=\eta_{E_0/F_0}\). The restriction of \(\chi_7\) to F₀ has sign at each real place, since seven is odd, and is unramified at all finite places. Hence \((\chi_7|_{\mathbb A_{F_0}^\times})/\eta_0\) has trivial infinity component and factors through \(\mathrm{Cl}(F_0)\). By (ER16), \[
                        \chi_7|_{\mathbb A_{F_0}^\times}=\eta_0.
\tag{ER17}
\]

Put \(\eta_1=\eta_{K/F}\) and \[
 \beta=\chi_7\circ N_{K/E_0},\qquad
 \delta=\chi_0^7\beta^{-1},\qquad
 \tau=(\chi_0|_{\mathbb A_F^\times})\eta_1^{-1}.
\tag{ER18}
\] The norm pulls local units to local units. Thus β is unramified everywhere finite and its infinity component is \(a^7\). It follows that δ is an ordinary class character of K. Likewise τ is an ordinary class character of F, because the infinity component of \(\chi_0|_F\) is the product of signs, exactly that of \(\eta_1\).

The quadratic base-change identity gives \[
                \beta|_{\mathbb A_F^\times}
             =\eta_0\circ N_{F/F_0}=\eta_1.
\tag{ER19}
\] It can also be checked locally here. At a finite prime of F₀ where E₀ splits, both sides are trivial. At an inert finite prime E₀/F₀ is unramified quadratic; after a residue extension of degree f it splits precisely when f is even. The unramified value of \(\eta_0\) on the norm of a uniformizer is \((-1)^f\), the same as the new quadratic character. Both sides are trivial on units. At infinity the norm and sign give the same product of local signs. This proves (ER19), including primes where F/F₀ itself is ramified.

Let \(j:\mathrm{Cl}(F)\to\mathrm{Cl}(K)\) be extension of ideals. From (ER18)--(ER19), \[
                                  \delta\circ j=\tau^7.
\tag{ER20}
\] The kernel of j has exponent dividing two: applying the ideal norm to an extended ideal gives its square. If x lies in this kernel, then (ER20) gives \(\tau(x)^7=1\), while \(x^2=1\) gives \(\tau(x)^2=1\). Therefore \(\tau(x)=1\). Thus \(\tau^{-1}\) is a well-defined character on \(\operatorname{im}j\).

Every character of a subgroup of a finite abelian group extends to the whole group: adjoining one generator at a time, choose a root of the specified value of its least positive power lying in the previous subgroup. All these values remain roots of unity. Apply this to obtain a class character \(\lambda\) of K with \[
                                   \lambda\circ j=\tau^{-1}.
\tag{ER21}
\] Finally set \[
                                      \eta=\chi_0\lambda.
\tag{ER22}
\] It is unitary, unramified at every finite place, has infinity component a, and its restriction to F is \(\tau\eta_1\tau^{-1}=\eta_1\). This proves all four literal \nolinkurl{IsEtaT} clauses. The argument has handled the capitulation compatibility explicitly; it has not inferred \nolinkurl{ResEta} merely from the unit condition.

\subsection{6. A single finite coefficient field for etaVal}\label{a-single-finite-coefficient-field-for-etaval}

The finite values required in the source are normalized by \((NP)^{-1/2}\). Their common algebraicity can be built into the preceding construction.

Define the algebraic type norm \[
                T(x)=\prod_{\phi\in\Phi_0}\phi(N_{K/E_0}x)\in E_0^\times.
\tag{ER23}
\] Since E₀ is Galois, every factor and the product lie in E₀. For the chosen complex embedding, \[
              a(x)=\frac{T(x)}{\sqrt{N_{K/\mathbb Q}x}},
\tag{ER24}
\] where the norm is positive because K is CM. In particular \(T(U)=1\) on every unit by (ER13).

The principal-ideal rule \(\psi((x))=T(x)^{-1}\) is therefore well-defined. Extend it to all ideals by the cyclic-class construction of §4, choosing only the finitely many roots \[
                         \psi(\mathfrak a_j)^{h_j}=T(x_j)^{-1}.
\tag{ER25}
\] All ψ-values then lie in one finite extension of E₀. Every chosen root has complex absolute value \((N\mathfrak a_j)^{-1/2}\), since \(\mathfrak a_j^{h_j}=(x_j)\); the same absolute-value identity holds on all ideals. Thus \(c(I)=\psi(I)\sqrt{NI}\) is the unitary ideal character used to construct \(\chi_0\).

The correcting class character λ has finite order. Adjoining its finitely many root-of-unity values, one obtains a fixed finite coefficient field containing \[
                        \eta(\varpi_P)(NP)^{-1/2}
                                      =\psi(P)\lambda(P)
\tag{ER26}
\] for every finite prime P. Closing that field under conjugation is still a finite extension. This supplies the existence part of the source \nolinkurl{etaVal} condition, with its exact \nolinkurl{valTw(1/2)} normalization. No bound for the necessary class orders or this coefficient field is part of the theorem.

\subsection{7. Persistence under further real extensions and remaining host work}\label{persistence-under-further-real-extensions-and-remaining-host-work}

The repaired character persists under any further totally real extension \(F_1/F\). Put \(K_1=KF_1\), use the induced CM type, and pull η back by \(N_{K_1/K}\). Its infinity exponent remains one at every induced embedding, its finite components remain unramified, and the same local quadratic base-change check proves the exact restriction \(\eta_{K_1/F_1}\). Thus the character obstruction need not be solved again when a further real extension is required for other numerical host conditions.

This construction concerns the ordinary type of a possible \(J^{14}\) host. It proves the auxiliary character and its normalized value field, with no assertion here that the full integral polarized lattice is the standard product lattice or that all old source real, rank, tensor, and size records are instantiated. Those additional comparisons must retain their own finite and archimedean costs. The degree-seven step repairs one actual previously impossible field of the source adapter; it does not complete the GRH extension or the MainX window.

\subsection{8. Primary references and precise use}\label{primary-references-and-precise-use}

\begin{itemize}
\tightlist
\item
  Masley, J. M. (1978). Class numbers of real cyclic number fields with small conductor. \emph{Compositio Mathematica, 37}(3), 297--319. Table 3, printed p.316, conductor28/degree6 row: \(h(F_0)=1\). Primary PDF: \url{https://numdam.org/item/CM_1978__37_3_297_0.pdf}. The table image was inspected directly, rather than relying on search snippets or a secondary class-number list.
\item
  Milne, J. S. (2020). \emph{Class field theory} (Version4.03, 6 August 2020). Chapter V, Proposition4.12, pp.177; Proposition5.2 and Theorem5.3, pp.178--179: the quadratic Artin character, its local restrictions, principal-idele triviality, and norms. \url{https://www.jmilne.org/math/CourseNotes/CFT.pdf}. Only these standard class-field inputs are used; the unit generation, character-extension, and capitulation correction are proved above.
\item
  Accepted DP57 and AM-DO-022: the original ordinary J² field/type and finite-unramified CM extension E₀/F₀. AM-K028: exact source \nolinkurl{HasInf}, \nolinkurl{UnrAt}, \nolinkurl{ResEta}, \nolinkurl{IsEtaT}, and \nolinkurl{etaVal} conventions and their provenance.
\end{itemize}

The exact finite arithmetic checks are separate from the proof of character existence. No Lean compilation, new analytic L-function theorem, large-field class-number computation, or effective coefficient-height estimate is claimed.

\clearpage

\section{AM-K030 --- an explicit principal product host for the degree-seven character repair}\label{am-k030-an-explicit-principal-product-host-for-the-degree-seven-character-repair}

Author: GPT-6 Astra Pro, ALP\_MACHINE \nolinkurl{/root/kummer_extension}, 21 November 2026 UTC. Inherited reasoning effort is not separately exposed.

AM-K029 supplies the literal unramified \nolinkurl{IsEtaT} character on an induced-type degree-84 CM field. This finding supplies its actual integral principally polarized ordinary-Jacobian host. The identification is an isomorphism of principally polarized complex abelian varieties, with a full integral CM action. It is stronger than an isogeny or a rational trace-form comparison.

Keep \[
 E_0=\mathbb Q(\zeta_{28}),\quad F_0=E_0^+,\quad
 \Phi_0=\{1,3,9,15,17,23\},\quad A=J^2,
\] where \(J:y^2=x^7-1\) has its actual principal polarization. Let B be the unique degree-seven subfield of \(\mathbb Q(\zeta_{29})\), let \(K=E_0B\), and let \(F=F_0B\). The induced type \(\Psi\) consists of all seven lifts of every embedding in \(\Phi_0\).

Accepted DP57 and AM-DO-022 identify A with the type-\(\Phi_0\) CM torus whose lattice is \(\mathcal O_{E_0}\) and whose exact principal trace coefficient is \[
 \xi_0=\frac{\zeta_7-\zeta_7^{-1}}{14}.
\tag{PH1}
\] All sign conventions below use that same alternating form. Nothing changes its source/target Hom convention.

\subsection{1. An explicit selfdual normal basis}\label{an-explicit-selfdual-normal-basis}

Set \(\zeta=\zeta_{29}\), \[
 H=\{1,12,17,28\}\subset(\mathbb Z/29\mathbb Z)^\times,
 \qquad C_i=2^iH\quad(i\in\mathbb Z/7\mathbb Z),
\] and define the Gaussian periods \[
             \eta_i=\sum_{h\in H}\zeta^{2^ih}.
\tag{PH2}
\] The integer 2 has order 28 modulo 29, and H is its subgroup of order four. In particular H contains \(-1\), so the periods and B are real. Let \(\sigma\) be the generator of \(\operatorname{Gal}(B/\mathbb Q)\) induced by \(\zeta\mapsto\zeta^2\). Then \(\sigma\eta_i=\eta_{i+1}\).

The following seven elements are completely explicit: \[
 \begin{split}
  a&=(-8,-4,-2,-5,-2,-4,-4),\\
  e_0&=\frac1{29}\sum_{i=0}^6a_i\eta_i,\qquad
  e_j=\sigma^j(e_0)\quad(0\le j\le6).
 \end{split}
\tag{PH3}
\]

\subsubsection{Theorem PH4}\label{theorem-ph4}

Let \(\mathfrak p\) be the unique prime of B above 29. Then \[
 \mathfrak D_{B/\mathbb Q}=\mathfrak p^6,
 \qquad \mathcal A_B:=\mathfrak p^{-3},
 \qquad
 \mathcal A_B=\bigoplus_{j=0}^6\mathbb Z e_j,
 \qquad \operatorname{Tr}_{B/\mathbb Q}(e_je_k)=\delta_{jk}.
\tag{PH4}
\] Every \(e_j\) has \(\mathfrak p\)-valuation exactly \(-3\), and \(\operatorname{Tr}_{B/\mathbb Q}e_j=1\).

\textbf{Proof of the ramification and trace identities.} The prime 29 is totally ramified in \(\mathbb Q(\zeta)\), and hence is totally ramified in its degree-seven subfield B. There are no other ramified primes. This extension is tame because \(29\nmid7\); the different exponent is \(7-1=6\). These are the standard cyclotomic and tame-different facts; precise references appear below.

The sum of the seven periods is \(-1\). Therefore \(\operatorname{Tr}\eta_i=-1\). To compute \(\operatorname{Tr}(\eta_i\eta_j)\), sum the products over the seven simultaneous translates of the two cosets. The coefficient of \(\zeta^t\) is constant for nonzero residues t, since multiplication by any nonzero residue permutes all terms. If \(i=j\), there are 28 terms with exponent zero: for each of the seven cosets, there are four ordered pairs \((x,-x)\). The remaining 84 terms distribute three to every nonzero residue. If \(i\ne j\), there is no exponent-zero term and all 112 terms distribute four to every nonzero residue. Hence \[
                  \operatorname{Tr}(\eta_i\eta_j)=29\delta_{ij}-4.
\tag{PH5}
\] In particular the periods are linearly independent, their trace determinant is \(29^6\), and, by the different calculation, they are an integral basis of \(\mathcal O_B\).

For the vector a in (PH3), direct integer multiplication gives \[
 \sum_i a_i=-29,\qquad
 \sum_i a_i a_{i+k}=\begin{cases}145&k=0,\\116&k\ne0.\end{cases}
\tag{PH6}
\] One may check the shorter equivalent computation with \(b=a+4(1,\ldots,1)=(-4,0,2,-1,2,0,0)\): its sum is \(-1\), and its cyclic correlations are 25 at zero and \(-4\) elsewhere. Combining (PH5) and (PH6) yields \[
 \operatorname{Tr}(e_0e_k)
 =\frac{29\sum_i a_i a_{i-k}-4(\sum_i a_i)^2}{29^2}
 =\delta_{0k}.
\tag{PH7}
\] The trace-one assertion follows from \(\operatorname{Tr}\eta_i=-1\). Galois invariance gives all entries in (PH4).

\textbf{Proof of membership in the square-root inverse different.} Write \[
             P(X)=\sum_i a_i\sum_{h\in H}X^{(2^ih\bmod29)},
\] using representatives in \(\{1,\ldots,28\}\). The exact polynomial identity needed for the local check is \[
\begin{split}
 P(1+T)\equiv T^{16}\bigl(&1-8T-13T^2+2T^3+7T^4-14T^5
 +14T^6-5T^7\\
 &-12T^8+6T^9-4T^{10}+4T^{11}-8T^{12}\bigr)\pmod{29}.
\end{split}
\tag{PH8}
\] This is a 29-term binomial computation; the accompanying certificate verifies every coefficient, not just the first nonzero one. In the completion of \(\mathbb Q(\zeta)\) at 29, \(t=\zeta-1\) has valuation one and 29 has valuation 28. Equation (PH8) shows that \(P(\zeta)\) has valuation exactly 16: all terms of lower degree have valuation at least 28, its degree-16 coefficient is a unit, and all higher-degree terms have strictly larger valuation. The extension \(\mathbb Q(\zeta)/B\) has ramification degree four, so \[
                    v_{\mathfrak p}(P(\zeta))=4,
                    \qquad v_{\mathfrak p}(e_0)=4-7=-3.
\tag{PH9}
\] At every other finite prime, e₀ is integral because its only rational denominator is 29. The same statements hold for all its conjugates. Thus \(L:=\sum\mathbb Z e_j\subset\mathcal A_B\).

The trace-dual of a fractional ideal I is \(I^*=\mathfrak D_{B/\mathbb Q}^{-1}I^{-1}\). Consequently \(\mathcal A_B^*=\mathcal A_B\). Equation (PH7) says \(L^*=L\). Taking duals of \(L\subset\mathcal A_B\) gives \(\mathcal A_B\subset L\), proving equality. This finishes the proof. No classification of positive unimodular lattices, approximate lattice search, or unproved integral-normal-basis assertion is needed. ∎

\subsection{2. A positive generator and an actual integral CM action}\label{a-positive-generator-and-an-actual-integral-cm-action}

Define \[
 \pi_B=N_{\mathbb Q(\zeta)/B}(1-\zeta)
       =\prod_{h\in H}(1-\zeta^h)
       =4-2\eta_0+\eta_4.
\tag{PH10}
\] The last identity is obtained by multiplication and the relation \(\sum_i\eta_i=-1\); the certificate checks it exactly in the cyclotomic algebra.

The factors with exponents \(h\) and \(-h\) multiply to \(|1-\zeta^h|^2>0\). This remains true at every real embedding of B. Thus \(\pi_B\) is totally positive. Norm transitivity and \(N_{\mathbb Q(\zeta)/\mathbb Q}(1-\zeta)=29\) give \[
              N_{B/\mathbb Q}\pi_B=29,
              \quad \mathfrak p=(\pi_B),
              \quad \mathcal A_B=\pi_B^{-3}\mathcal O_B.
\tag{PH11}
\]

Multiplication by any \(\beta\in\mathcal O_B\) on the basis e gives an integral matrix: \(\mathcal A_B\) is an \(\mathcal O_B\)-ideal. Its entries are explicitly \[
                 (M_\beta)_{jk}=\operatorname{Tr}(\beta e_k e_j),
\tag{PH12}
\] and are symmetric. The matrices give a faithful ring embedding \(\mathcal O_B\hookrightarrow M_7(\mathbb Z)\), with every B-element self-adjoint for the standard Euclidean form. As an optional directly usable generator check, multiplication by \(\eta_0\) is \[
 M_{\eta_0}=\begin{pmatrix}
 -2&0&-1&0&0&0&-1\\
 0&0&0&-1&0&0&1\\
 -1&0&0&1&0&1&-1\\
 0&-1&1&1&0&1&0\\
 0&0&0&0&0&-1&0\\
 0&0&1&1&-1&0&1\\
 -1&1&-1&0&0&1&0
 \end{pmatrix}.
\tag{PH13}
\] The conjugate-period matrices are obtained by conjugating this matrix by the cyclic permutation \(e_j\mapsto e_{j+1}\). Since the periods form an integral basis, these matrices specify the full \(\mathcal O_B\)-action. The characteristic polynomial is \[
 X^7+X^6-12X^5-7X^4+28X^3+14X^2-9X+1.
\tag{PH14}
\] Neither the characteristic polynomial nor a numerical eigenvalue computation is being used to infer the integral-lattice theorem; it was proved above.

\subsection{3. The actual product principal polarization}\label{the-actual-product-principal-polarization}

The discriminants of E₀ and B are coprime, and the degrees 12 and 7 are coprime. The integral-compositum theorem therefore gives \[
            \mathcal O_K=\mathcal O_{E_0}\otimes_{\mathbb Z}\mathcal O_B,
 \qquad |d_K|=|d_{E_0}|^7 |d_B|^{12}
                  =2^{84}7^{70}29^{72}.
\tag{PH15}
\] Consider the full lattice \[
       \Lambda=\mathcal O_{E_0}\otimes_{\mathbb Z}\mathcal A_B
               =\pi_B^{-3}\mathcal O_K\subset K.
\tag{PH16}
\] It is stable under \(\mathcal O_K\). On pure tensors define the alternating form by the product of the old principal alternating form and the real trace pairing of B. Equivalently, in K it is the trace form with coefficient \(\xi_0\). For arbitrary \(x_j,y_j\in E_0\), the exact identity is \[
 \operatorname{Tr}_{K/\mathbb Q}
  \left(\xi_0\Bigl(\sum_jx_je_j\Bigr)
                    \overline{\Bigl(\sum_ky_ke_k\Bigr)}\right)
       =\sum_{j=0}^6\operatorname{Tr}_{E_0/\mathbb Q}
                           (\xi_0x_j\bar y_j).
\tag{PH17}
\] It follows directly from (PH4), because conjugation fixes B. Thus the tensor alternating form is integral and unimodular on \(\Lambda\). Its positivity follows from the positivity of the old Hermitian form and of the real trace form. With the induced complex structure of \(\Psi\), the map \[
         (x_0,\ldots,x_6)\longmapsto\sum_jx_je_j
\tag{PH18}
\] is an isomorphism of polarized complex tori \[
      (A^7,\lambda_A^{\boxtimes7})
           \simeq (K\otimes\mathbb R/\Lambda,\lambda_\Lambda)
           \simeq (J^{14},\lambda_J^{\boxtimes14}).
\tag{PH19}
\] For the complex-linearity assertion, B acts through its regular real representation on the seven copies; E₀ supplies the same complex structure on each copy. Diagonalizing that real representation gives exactly all seven lifts of every old selected embedding. This is the stated induced type, not a change to conjugate Jacobian factors.

The commuting integral actions of \(\mathcal O_{E_0}\) and \(\mathcal O_B\), together with (PH15), supply the full integral \(\mathcal O_K\)-action on the actual product. Rosati acts as conjugation on E₀ by DP57 and as the identity on B by (PH12); it is therefore complex conjugation on K. The field degree is \(84=2\dim J^{14}\), so this is a full CM action.

If the source uses the standard lattice \(\mathcal O_K\), transport (PH16) by multiplication by \(\pi_B^{-3}\). The exact coefficient then becomes \[
             \boxed{\xi_K=\xi_0\pi_B^{-6}.}
\tag{PH20}
\] Since \(\pi_B\) is totally positive, this has the same selected positivity signs. Moreover \[
        \xi_K\mathcal O_K=\mathfrak D_{K/\mathbb Q}^{-1},
       \qquad |N_{K/\mathbb Q}\xi_K|
                      =2^{-84}7^{-70}29^{-72}.
\tag{PH21}
\] The ideal equality also follows from the established perfect trace pairing, so it is stronger than a mere equality of absolute norms. The factor \(\pi_B^{-6}\) is part of the new host data. Keeping the old \(\xi_0\) while silently replacing \(\Lambda\) by \(\mathcal O_K\) would give the wrong principal lattice.

\subsection{4. The theta line and the CM coordinate comparison}\label{the-theta-line-and-the-cm-coordinate-comparison}

Use the same trace-one element \[
                      w_0=-(\zeta_7+\zeta_7^2+\zeta_7^4)
                      \quad(w_0+\bar w_0=1).
\] It acts diagonally in (PH18). Equation (PH17) splits its actual trace-one semicharacter into the product of the old semicharacters. Thus the rigidified symmetric principal line is exactly the external product of the fourteen ordinary Jacobian trace-one lines. With the normalization of AM-R-006 and AM-LIFT-08, transport the already fixed product section through (PH18). Its actual reduced germ in product coordinates is \[
        \prod_{r=1}^{14}\Sigma_{\mathrm{tr}}(u^{(r)})
            =c^{14}\prod_{r=1}^{14}\sigma_q(u^{(r)}).
\tag{PH22}
\] Here \(\sigma_q=u_1u_3-u_2^2+\cdots\), and c is the already fixed algebraic scalar. In particular its order of vanishing at the origin is 28. No theta-null unit is asserted. This is the line at the invariant characteristic \(2\infty\), not the distinct non-effective translated characteristic.

Let \(\tau_0,\ldots,\tau_6\) be the seven real embeddings of B and put \[
                        U_{kj}=\tau_k(e_j).
\tag{PH23}
\] Equation (PH4) says \(U^tU=I_7\). Hence U is a real orthogonal matrix, \(U^{-1}=U^t\), and \(\det U=\pm1\). Write the seven A-copy coordinates as \((u_{i,j},v_{i,j})\), \(i=1,2,3\), \(j=0,\ldots,6\), and define \[
 w_{i,\pm,j}=u_{i,j}\pm i v_{i,j},\qquad
 z_{i,\pm,k}=\sum_{j=0}^6U_{kj}w_{i,\pm,j}.
\tag{PH24}
\] These are full K-eigenvectors for the actual induced CM action. The Hermitian metric is exactly \[
 H=\sum_{i,j}\frac{|u_{i,j}|^2+|v_{i,j}|^2}{G_i}
   =\sum_{i,\pm,k}\frac{|z_{i,\pm,k}|^2}{2G_i}.
\tag{PH25}
\] Thus the B-diagonalization itself has no archimedean operator-norm loss.

At every prime away from 29, all entries of U and \(U^{-1}\) are integral. At 11, (PH24) is consequently an integral invertible change of the primitive Hodge eigenbasis, including its old \(\pm i\) blocks, since 2 is a unit. This proves a precise integral comparison at 11, not merely an algebraic splitting of de Rham cohomology. AM-K018 then applies to this ordinary \(J^{14}\) product: the central \(\zeta_7\)-action has no invariant holomorphic quadratic form, and its full-CM Chern/reduced-gauge comparison is unchanged. No claim is made for the conjugate-factor products outside that lemma's scope.

At 29, every entry of U has normalized valuation \(-3/7\), by (PH9); its inverse has the same bound. A homogeneous derivative of degree D transported by this matrix therefore has coefficient valuations at least \(-3D/7\). This is an upper bound on the possible denominator, not a claim that every derivative attains it. The old \(\pm i\) coordinate change retains its dyadic denominator bound \(2^D\). The additional standard-lattice rescaling in (PH20), if used to choose logarithmic coordinates, must also be included explicitly; it is a diagonal multiplication by the displayed conjugates of \(\pi_B^{\pm3}\). All these diagonal factors are units at 11.

This section identifies the line, product section, metric, and local coordinate change. It does not assert that the old selected-Schur upper-bound witnesses retain their old individual CM weight after the seven-copy change. The common joint integral lower estimate can be transported through an integral invertible 11-coordinate change; a new selected upper witness needs its own weight calculation.

\subsection{5. Finite-level costs that change under the repair}\label{finite-level-costs-that-change-under-the-repair}

The relative CM extension K/F remains unramified at every finite place by AM-K029. At 29 it is split, because 29 splits completely in E₀. Nevertheless the absolute real field F has six primes above 29, each with residue field \(\mathbb F_{29}\) and absolute ramification index seven. Rational 29 is therefore \textbf{not} a uniformizer of these local real fields.

For example, for the standard split rank-a compact group at all six places, a principal congruence subgroup defined by the rational ideal \((M)\), with \(M=29^e\) and \(e\ge1\), has the exact index \[
     I_a^{\mathrm{new}}(M)
        =M^{42a^2}\prod_{i=1}^a(1-29^{-i})^6.
\tag{PH26}
\] Indeed the local residue length is \(7e\), so each local index is \(29^{7ea^2}\prod_{i=1}^a(1-29^{-i})\), by reduction to \(\operatorname{GL}_a(\mathbb F_{29})\) and the successive additive matrix kernels. Multiplying six places gives (PH26). This statement concerns that precisely defined principal subgroup; it is not an automatic identification of every source level symbol with rational M.

The old 29 half-orbit construction still has to be instantiated with the new actual lattice and its exact \(\alpha\bar\alpha=M\) data. The good-prime jet-surjectivity argument of AM-K022 applies at 11 once those data are specified, but the former absolute-uniformizer and degree-six mass formulas cannot simply be copied. Terminal \nolinkurl{qOf} depths must likewise retain their actual local valuations.

Combining this theorem with AM-K029 gives an actual product principal CM host and its literal auxiliary character. It does not settle the independent original \nolinkurl{IsTen}, full real/rank/size records, prescribed finite Xi equality, or the later selected-image arithmetic window. The degree-seven repair is one concrete part of that adapter, with its new finite costs exposed.

\subsection{6. Sources, novelty, and certificate scope}\label{sources-novelty-and-certificate-scope}

\begin{itemize}
\tightlist
\item
  Milne, J. S. (2020). \emph{Algebraic number theory} (Version 3.08, 19 July 2020). Proposition 6.2 and Theorem 6.4, printed pp.96--99, for prime cyclotomic ramification; Lemma 6.5, printed p.100, and Remark 6.6(c), printed p.101, for the integral compositum and its discriminant. Primary author PDF: \url{https://www.jmilne.org/math/CourseNotes/ANT.pdf}. The exact compositum statements were read in the source.
\item
  Erez, B. (1991). A survey of recent work on the square root of the inverse different. \emph{Astérisque, 198--200}, 133--152. Section 1, equations (1.1)--(1.2), printed p.134, for the tame different and trace-dual ideal; Theorem 2.9, printed p.139, records the general selfdual-normal-basis theorem in this tame abelian setting. Primary author paper: \url{https://smf.emath.fr/system/files/filepdf/AST_1991__198-199-200__133_0.pdf}. The general theorem is background only: (PH3)--(PH9) verify the particular basis directly.
\item
  Accepted DP57; AM-DO-022, §§1--2; AM-R-004/006 and AM-LIFT-08: the original Jacobian, product lattice and exact theta normalization. AM-K018: the full-CM gauge comparison for ordinary powers. AM-K029: the enlarged field, induced type and literal unramified-character construction.
\end{itemize}

The general square-root-inverse-different and selfdual-normal-basis facts are classical, not claimed as new. The concrete contributions here are the displayed basis and its exact local certificate for conductor 29, the integral product-polarization transport, the exact coefficient (PH20), and the finite coordinate/cost comparison for this host. The certificate checks rational/cyclotomic algebra, matrix identities, all coefficients in (PH8), and the displayed norm/index formulas. No Lean build, unbounded analytic derivative computation, or full source-host instance is claimed.

\clearpage

\section{AM-K032 --- an actual product principal host with degree divisible by eight}\label{am-k032-an-actual-product-principal-host-with-degree-divisible-by-eight}

Author: GPT-6 Astra Pro, ALP\_MACHINE /root/kummer\_extension, 22 November 2026 UTC; inherited reasoning effort not separately exposed.

The degree-seven construction of AM-K029/030 repairs the odd character but has real degree42. This finding supplies an explicit further real extension of degree4, together with an exact polarized product identification. The resulting host is the actual principally polarized variety \(J^{56}\), with full integral CM action, real degree168, and \(8\mid168\). Its odd character, all uniform even stage characters, and the permitted alpha half-orbit persist. The primitive-type common-field obstruction remains: the source's defined Fields.Kq is still \(\mathbb Q\). Consequently this is not a complete D7 constructor or a MainX estimate.

The new ingredient is elementary weighted trace linear algebra. No classification of positive unimodular lattices and no uncomputed class group is needed.

\subsection{1. Two exact rank-two weighted trace forms}\label{two-exact-rank-two-weighted-trace-forms}

For \(a\in\{1,3\}\), put \[
 d_a=a^2+4,\qquad R_a=\mathbb Q(\sqrt{d_a}),\qquad
 u_a=\frac{a+\sqrt{d_a}}2,
 \qquad \delta_a=\frac{u_a^{-1}}{\sqrt{d_a}}.
\tag{QH1}
\] Thus \((d_1,d_3)=(5,13)\). Each \(d_a\) is a squarefree discriminant congruent to1 modulo4. The ring of integers is \[
                 \mathcal O_{R_a}=\mathbb Z[u_a],
           \qquad u_a^2-au_a-1=0,
           \qquad u_a^{-1}=u_a-a.
\tag{QH2}
\] For a=3, the generator differs by1 from the usual \((1+\sqrt{13})/2\), so it gives the same integral ring. The discriminant of the displayed basis is \(d_a\).

\subsubsection{Proposition QH3}\label{proposition-qh3}

The element \(\delta_a\) is totally positive, generates the inverse different of \(R_a/\mathbb Q\), and satisfies \[
 \operatorname{Tr}_{R_a/\mathbb Q}(\delta_a)=1,\qquad
 \operatorname{Tr}_{R_a/\mathbb Q}(\delta_a u_a)=0,\qquad
 \operatorname{Tr}_{R_a/\mathbb Q}(\delta_a u_a^2)=1.
\tag{QH3}
\] Thus the weighted trace form \((x,y)\mapsto\operatorname{Tr}(\delta_a xy)\) has matrix \(I_2\) in the \textbf{integral} basis \((1,u_a)\). Multiplication by \(u_a\) in that basis is the symmetric integral matrix \[
                         M_a=\begin{pmatrix}0&1\\1&a\end{pmatrix}.
\tag{QH4}
\]

\textbf{Proof.} Write \(s=\sqrt{d_a}>a>0\) in the positive embedding. Then \[
 \delta_{a,+}=\frac{s-a}{2s}>0,\qquad
 \delta_{a,-}=\frac{s+a}{2s}>0.
\tag{QH5}
\] Their sum is1. Since \(\delta_a u_a=1/\sqrt{d_a}\), its trace is0. Since \(\delta_a u_a^2=u_a/\sqrt{d_a}\), its trace is1. This proves all three entries of the Gram matrix. In particular the integral lattice is selfdual for the weighted trace form, so \(\delta_a\mathcal O_{R_a}=\mathfrak D_{R_a/\mathbb Q}^{-1}\) by the definition of the inverse different. This also follows from the derivative \(2u_a-a=\sqrt{d_a}\) and the unit \(u_a^{-1}\), but the Gram computation already proves it directly. The multiplication matrix follows from \(u_a^2=1+au_a\). Finally \[
                    N_{R_a/\mathbb Q}(\delta_a)=1/d_a.
\tag{QH6}
\] All statements are exact. ∎

Let \(R=\mathbb Q(\sqrt5,\sqrt{13})\). Coprimality of the two discriminants implies \[
 \mathcal O_R=\mathcal O_{R_1}\otimes_{\mathbb Z}\mathcal O_{R_3},
 \qquad [R:\mathbb Q]=4,\qquad d_R=5^2 13^2.
\tag{QH7}
\] In the basis \[
                     1,\ u_1,\ u_3,\ u_1u_3,
\tag{QH8}
\] the form \(\operatorname{Tr}_{R/\mathbb Q}(\delta_1\delta_3xy)\) is \(I_4\). The two generating actions are commuting symmetric integral Kronecker matrices. These facts follow by multiplying the two trace identities, and the certificate checks the four-by-four matrices explicitly.

\subsection{2. Tensoring the actual CM lattice}\label{tensoring-the-actual-cm-lattice}

Retain the notation and the \textbf{actual} polarized product construction of AM-K030: \[
 \begin{split}
 E_0&=\mathbb Q(\zeta_{28}),\qquad
 B\subset\mathbb Q(\zeta_{29}),\quad[B:\mathbb Q]=7,\\
 K&=E_0B,\quad F=K^+,\quad [K:\mathbb Q]=84,\\
 \mathcal A_B&=\pi_B^{-3}\mathcal O_B,
 \qquad \pi_B=N_{\mathbb Q(\zeta_{29})/B}(1-\zeta_{29}),\\
 \Lambda&=\mathcal O_{E_0}\otimes\mathcal A_B
               =\pi_B^{-3}\mathcal O_K,\qquad
 \xi_0=\frac{\zeta_7-\zeta_7^{-1}}{14}.
 \end{split}
\tag{QH9}
\] The carried lattice \((\Lambda,\operatorname{Tr}_{K/\mathbb Q}(\xi_0x\bar y))\) is exactly the product principal CM lattice of \(J^{14}\), by the explicit trace-orthonormal basis of \(\mathcal A_B\). Its type is induced from \(\Phi_0=\{1,3,9,15,17,23\}\) and contains every real B-lift.

Set \[
                  K^*=KR,
              \qquad F^*=FR.
\tag{QH10}
\] The discriminants of K and R are coprime, and K is unramified at5 and13. Thus the fields are linearly disjoint and \[
 \begin{gathered}
 [K^*:\mathbb Q]=336,\qquad [F^*:\mathbb Q]=168,\\
 \mathcal O_{K^*}=\mathcal O_K\otimes\mathcal O_{R_1}
                               \otimes\mathcal O_{R_3}.
 \end{gathered}
\tag{QH11}
\] For completeness, any nontrivial subfield of the biquadratic R contains one of \(\mathbb Q(\sqrt5),\mathbb Q(\sqrt{13}),\mathbb Q(\sqrt{65})\), all ramified at5 or13; hence K cannot intersect R nontrivially. The integral tensor assertion is the coprime-discriminant lemma, applied twice.

Define \[
 \Lambda^*=\Lambda\otimes\mathcal O_{R_1}\otimes\mathcal O_{R_3}
                   =\pi_B^{-3}\mathcal O_{K^*},
 \qquad
 E^*(x,y)=\operatorname{Tr}_{K^*/\mathbb Q}
                       (\xi_0\delta_1\delta_3x\bar y).
\tag{QH12}
\] Use every extension of each old CM embedding to K*, so the new type has four lifts of each of the42 old embeddings.

\subsubsection{Theorem QH13}\label{theorem-qh13}

The CM lattice (QH12) is an actual product principal polarization: \[
                    (A^*,L^*)\simeq
                       (J,L_{2\infty})^{56}
\tag{QH13}
\] as principally polarized abelian varieties. It has a full action of \(\mathcal O_{K^*}\), with Rosati involution equal to complex conjugation. The CM extension \(K^*/F^*\) is unramified at every finite prime. On the standard lattice \(\mathcal O_{K^*}\), the coefficient of its alternating trace form is \[
                    \boxed{\xi^*=\xi_0\pi_B^{-6}\delta_1\delta_3.}
\tag{QH14}
\]

\textbf{Proof.} Decompose an element of \(\Lambda^*\) in the four-element real basis (QH8), with coefficients in \(\Lambda\). Taking the trace first over the two quadratic extensions and applying (QH3) kills the mixed terms and leaves the sum of four copies of the old alternating trace form. Consequently the lattice is the orthogonal direct sum of four old principal symplectic lattices. Since all four real lifts of each old holomorphic embedding are chosen, the corresponding complex structure is the product complex structure. This proves (QH13), not merely an isogeny statement. Positivity is also immediate either from this product identification or from the total positivity of both \(\delta_a\).

The two new integral actions are the matrices (QH4), tensored with identity matrices; they commute with each other and with the old \(\mathcal O_K\)-action. Their tensor ring is the full integer ring in (QH11). They are selfadjoint because their matrices are symmetric in the product principal lattice. The old Rosati involution is conjugation; together these facts give conjugation on the full K* action. The CM type is the asserted induced type.

The extension K/F is finite-unramified by AM-K029/030. Finite étale quadratic algebras remain étale under base change. Applying this to F\emph{/F proves the finite-unramified assertion. The quadratic base change is still a field: F} is totally real, whereas a generator of the quadratic CM extension K/F cannot lie in a totally real field.

Finally, write \(x=\pi_B^{-3}x'\) and \(y=\pi_B^{-3}y'\) in (QH12). The element \(\pi_B\) is real, so the coefficient becomes (QH14). This multiplication transports the entire lattice and form; it is not an assertion that the coordinate factor is a unit at29. ∎

As an independent norm/discriminant check, \[
 \begin{split}
 |d_{K^*}|&=|d_K|^4d_R^{84}
      =2^{336}7^{280}29^{288}5^{168}13^{168},\\
 N_{K^*/\mathbb Q}(\xi^*\mathcal O_{K^*})
      &=|d_{K^*}|^{-1}.
 \end{split}
\tag{QH15}
\] The latter is the norm of the inverse different ideal. Equivalently, the principal trace-dual lattice calculation already shows that \(\xi^*\mathcal O_{K^*}=\mathfrak D_{K^*/\mathbb Q}^{-1}\).

\subsection{3. Actual line, gauge, eigenbasis and finite coordinate costs}\label{actual-line-gauge-eigenbasis-and-finite-coordinate-costs}

The trace-one element \[
                 w_0=-(\zeta_7+\zeta_7^2+\zeta_7^4),
                      \qquad w_0+\bar w_0=1,
\tag{QH16}
\] acts diagonally on all56 ordinary J factors. Under the polarized product identification, its symmetric trace-one line is the external product of \(L_{2\infty}\). One way to fix the symmetric representative is to use its invariance under \(\zeta_7\): the endomorphism \(\zeta_7-1\) is invertible on the two-torsion, so the principal polarization has a unique symmetric representative fixed by this action. This is the same argument used for the actual source line in ALP-LIFT-07/08 and K030.

Transport the already fixed product section, rather than introducing a new unconstrained section. With the notation of AM-R-006 and K030 its reduced germ is \[
                    c^{56}\prod_{\ell=1}^{56}\sigma_q(u_\ell),
\tag{QH17}
\] and its origin multiplicity is112. Here \(\sigma_q\) is the actual q-normalized principal germ and c is the fixed derivative-Thomae scalar. The transported principal section is primitive at11 by the carried finite normalization; no theta-null unit at the origin is claimed, since the section vanishes there.

To compare the new CM coordinates, first retain K030's trace-orthonormal B coordinates z, whose Hermitian form has coefficient \(1/(2G_i)\) on each old E₀ label \((i,\pm)\) and B embedding. For one new quadratic extension, the two ordinary product coordinates \((z_0,z_1)\) become \[
 \begin{pmatrix}z_+\\z_-\end{pmatrix}
      =V_a\begin{pmatrix}z_0\\z_1\end{pmatrix},
 \quad
 V_a=\begin{pmatrix}1&u_{a,+}\\1&u_{a,-}\end{pmatrix},
 \quad \det V_a=-\sqrt{d_a}.
\tag{QH18}
\] This is the embedding matrix of (1,u\_a), and \[
     V_a^{\mathsf t}
           \operatorname{diag}(\delta_{a,+},\delta_{a,-})V_a=I_2.
\tag{QH19}
\] In particular it is a \textbf{weighted} identity. The unrescaled V\_a is not asserted to be orthogonal. After both extensions the Hermitian form is exactly \[
       H=\sum_{i,\pm,k,\epsilon,\epsilon'}
          \frac{\delta_{1,\epsilon}\delta_{3,\epsilon'}}{2G_i}
                      |z_{i,\pm,k,\epsilon,\epsilon'}|^2.
\tag{QH20}
\] The positive square-root rescaling by \(\sqrt{\delta_{1,\epsilon}\delta_{3,\epsilon'}}\) is an archimedean isometry to the product metric. Formula (QH20), rather than an untracked scalar bound, is the appropriate input for the selected Fock norm.

Both V\_a and V\_a\^{}\{-1\} are integral at11, since the u\_a are algebraic units and their determinants are11-units. Combined with K030's B comparison and the old \(\pm i\) comparison, the full K* eigenbasis is therefore an integral invertible change of the primitive ordinary J tangent basis at11. The central \(\zeta_7\) still has holomorphic labels1,2,3. Their pairwise sums never vanish modulo7, including repeated factors. Thus AM-K018's uniqueness proof for the full-CM canonical/source gauge applies to this ordinary product. This supplies the same gauge comparison for (QH17); no extension to a product with conjugated J factors is asserted.

The additional finite coordinate costs are explicit:

\begin{longtable}[]{@{}
  >{\raggedright\arraybackslash}p{(\columnwidth - 2\tabcolsep) * \real{0.5000}}
  >{\raggedright\arraybackslash}p{(\columnwidth - 2\tabcolsep) * \real{0.5000}}@{}}
\toprule\noalign{}
\begin{minipage}[b]{\linewidth}\raggedright
Prime
\end{minipage} & \begin{minipage}[b]{\linewidth}\raggedright
Additional effect from \(R/\mathbb Q\)
\end{minipage} \\
\midrule\noalign{}
\endhead
\bottomrule\noalign{}
\endlastfoot
11 & V\_a and its inverse are integral invertible; all \(\delta_a\) are units \\
29 & No new denominator from V\_a; the B and standard-lattice \(\pi_B^{\pm3}\) factors of K030/DO034 remain \\
5 & Entries of \(V_1^{-1}\) have valuation \(-1/2\), normalized by \(v_5(5)=1\); the other V is invertible integral \\
13 & Entries of \(V_3^{-1}\) have valuation \(-1/2\); the other V is invertible integral \\
Other primes & Both new embedding matrices and their inverses are integral invertible \\
\end{longtable}

Indeed \[
 V_a^{-1}=\frac1{u_{a,-}-u_{a,+}}
              \begin{pmatrix}u_{a,-}&-u_{a,+}\\-1&1\end{pmatrix}.
\tag{QH21}
\] At the discriminant prime every numerator entry is a unit and the denominator has valuation1/2. Thus transforming a homogeneous degree-D polynomial from new eigen-derivatives to the ordinary product derivatives has additional coefficient denominator at most D/2 at5 and at13. This is an upper bound for the coordinate change, not a claim that an actual selected theta value attains it. The old dyadic and29 costs are retained. The explicit metric (QH20) and these local changes must be included when extending DO034's selected polynomial calculation; they are not removed by the11-unit comparison.

\subsection{4. Exact ranks, conductor and explicit full-congruence G primes}\label{exact-ranks-conductor-and-explicit-full-congruence-g-primes}

The smallest real extension degree h making \(8\mid42h\) is4. The present construction realizes it. The exact source f\_le\_k therefore requires \[
       \boxed{f=168,\qquad \operatorname{ord}=84,\qquad k\ge168.}
\tag{QH22}
\] The abstract Frame can again use its parameter e=0 and distinguished element0. Its group may be \(C_{21}\times C_2\times C_2\), of order84. Its quotient to HGrp(0) is trivial and satisfies the literal Frame clauses. In fact \(\operatorname{Gal}(K^*/\mathbb Q)\simeq C_2^4\times C_{21}\), compatible with an abstract \(G_0\times A\) decomposition. This observation alone does not construct the source's full field/type record.

For m=2 and k=168, take \[
 \begin{gathered}
 t=1633,\quad r=169,\quad d=1634,\quad r+t=1802,\\
                    (N_0,N_1,N_2)=(168,1802,3436).
 \end{gathered}
\tag{QH23}
\] All source rank inequalities hold. Substituting these ranks in K031's same numerical parameter formula gives \[
 \begin{gathered}
 S=3435,\quad S^3=40\,530\,337\,875,\quad
 C=S^3+k+t+11=40\,530\,339\,687,\\
 u_i=(337-2i)C+\tfrac12,\quad
 v_i=(336-2i)C,\quad
 (b_0,b_1)=(-2C,2C),\quad(e_0,e_1)=(0,4C).
 \end{gathered}
\tag{QH24}
\] The identical proof verifies every numerical parity, interlace, centred and NCGaps clause of Hyp3 with \(\eta_{\rm num}=-1/2\). The certificate separately checks the actual new finite lists. The source selected recurrence yields \[
 \begin{split}
 \nu_{0,i}&=(337-2i)C+i-984\quad(0\le i<168),\\
 \min_i\nu_{0,i}&=3C-817>0,\qquad
 \Delta=4C-1634,\qquad
 \nu_1=(\nu_0+\Delta,\Delta^{1634}).
 \end{split}
\tag{QH25}
\] The actual new selected polynomial must use all168 CM blocks. The highest bideterminant construction in DO034 is defined separately at each CM place and therefore gives a polynomial in these168 blocks by the same formula; its diagonal coefficient remains one. DO034's saved first-extraction theorem is stated for42 places. Extending its global formula requires the present quadratic metric data and new local indices, rather than silently retaining its numerical exponents. Equations (QH20)--(QH21) supply those additional metric/coordinate data; the separate first-lift analysis retains responsibility for that extension.

The rational conductor is exactly \[
                         N^*=812\cdot5\cdot13=52\,780.
\tag{QH26}
\] Containment in \(\mathbb Q(\zeta_{52780})\) follows from the three cyclotomic containments. Conversely, containment in any cyclotomic field forces divisibility by812 by K031, and ramification of the two new quadratic subfields forces divisibility by5 and13. This is exactly the source's minimum-divisibility IsCond criterion.

The following explicit primes satisfy the full new congruence: \[
 \begin{array}{c|r|l}
 i&g_i^*&\pi_i^*\\\hline
 0&22\,285\,932\,761&52\,781+52\,780\sqrt{-7}\\
 1&64\,072\,175\,441&211\,121+52\,780\sqrt{-7}.
 \end{array}
\tag{QH27}
\] The finite certificate checks primality by direct trial division through the integer square roots149284 and253124, respectively. It also checks \[
 \pi_i^*\bar\pi_i^*=g_i^*,\qquad
 \pi_i^*\equiv1\pmod{52780\mathcal O_{\mathbb Q(\sqrt{-7})}},\qquad
 g_i^*\equiv1\pmod{52780},\qquad
                 g_i^*-1>2\cdot168^2=56\,448.
\tag{QH28}
\] The conjugate ideals are coprime because \(\gcd(g_i^*,2\operatorname{Re}\pi_i^*)=1\). The primes split completely in K*, and each chosen half contains168 degree-one primes. Thus the actual norm, full PiGCong, residue cardinality and G-size clauses are met in the same external quadratic subfield as K031. The literal piG\_Kq clause is still not met.

For faithful residue characters one may use primitive roots3 and6 at the two primes. The certificate verifies their order from the complete factorizations of \(g_i^*-1\). The common residue-root field can be chosen using \[
 L^*=\operatorname{lcm}(28,g_0^*-1,g_1^*-1)
                  =13\,526\,981\,751\,068\,240.
\tag{QH29}
\] This explicit coefficient extension is large; no claim that its degree or arithmetic height is negligible is made.

\subsection{5. Actual characters and the unchanged permitted family}\label{actual-characters-and-the-unchanged-permitted-family}

Let \(\eta\) be K029's actual odd character on K. Define \[
                    \eta^*=\eta\circ N_{K^*/K}.
\tag{QH30}
\] By the norm-compatibility proof in K029, this character is unitary, finite-unramified, has angular exponent1 at each new CM embedding, and restricts to \(\eta_{K^*/F^*}\) on \(\mathbb A_{F^*}^\times\). Its normalized etaVal values lie in one finite coefficient field: at a prime P of residue degree f over p, the normalized value is the f-th power of the old normalized value. A norm of a local uniformizer has valuation f; its unit factor is killed by the unramified character, so no choice of uniformizer affects this assertion.

For any integers e\_j choose the LIFT18 finite-order zero-infinity characters \(\nu_j\) of E₀ at the primes (QH27), and set \[
         \boxed{\mu_j^*=(\nu_j\circ N_{K^*/E_0})(\eta^*)^{-2e_j}.}
\tag{QH31}
\] The full proof is HC16--21 of K031 and applies verbatim: every new G completion is the same degree-one \(\mathbb Q_{g_i^*}\), so its local norm is the identity; the even eta power changes infinity weight while changing no compact unit character. Hence all exact uniform IsMu0 clauses hold, without a divisibility condition on e\_j and without a large-field class-number assumption. A single finite field contains all muVal and etaVal values. For j=0,1 the exact conductor is \[
     \mathfrak f(\mu_j^*)=g_0^*g_1^*\mathcal O_{K^*},
       \qquad N\mathfrak f(\mu_j^*)=(g_0^*g_1^*)^{336}.
\tag{QH32}
\] Qualitative T-prime choices use K031's ray-class construction with modulus \(52780(g_0^*g_1^*)^2\); no quantitative upper window is inferred.

The same central element \[
           \alpha=(1+2\sqrt{-7})^e,\qquad M=29^e,
                    \alpha\bar\alpha=M
\tag{QH33}
\] acts diagonally on all56 J factors and is congruent to1 modulo2. Its conjugate ideal remains coprime. Therefore on an edge of ranks a,b the permitted product DP59 basis is exactly the h=4 instance of K031: \[
 \begin{gathered}
 X=(A^*)^{ab}=J^{56ab},\quad \dim X=168ab,\\
 \mathcal W_M=[\alpha]^*H^0(X,\mathcal L^M),
       \quad\operatorname{rank}\mathcal W_M=M^{168ab},\\
 \mathcal L=(L^*)^{\boxtimes ab},\qquad
 t'\in X[\alpha^2]/X[\alpha],\qquad Mt'\in X[\alpha].
 \end{gathered}
\tag{QH34}
\] The ordinary D-jet threshold remains \(M\ge3D+2\), and each individually selected finite translate has norm1. The auxiliary multijet threshold in DO030 is unchanged under its stated coprimality assumptions. The selected highest-symbol factorial bound is K022's bound in the new integral eigenbasis. It does not transfer an old six-place sharp directional witness to one new full-CM selected weight.

At29 both real quadratics split, as witnessed by \[
                       11^2\equiv5\pmod{29},\qquad
                       10^2\equiv13\pmod{29}.
\tag{QH35}
\] Thus F* has24 primes over29, each with ramification index7 and residue field \(\mathbb F_{29}\), and all are split in K\emph{/F}. The rational principal-level index in a rank-a split unitary factor is consequently \[
        \boxed{I^*(M,a)=M^{168a^2}
                       \prod_{i=1}^a(1-29^{-i})^{24}.}
\tag{QH36}
\] This follows from \(|\operatorname{GL}_a(\mathcal O_v/29^e)|
 =29^{7ea^2}\prod_{i=1}^a(1-29^{-i})\) at each of24 places. It is the rational principal-level factor, not a substitution for every source level or the prescribed ramified G-column index. The latter uses the new primes,168 selected places at each G, and the separate LIFT19/DO034 formulas.

\subsection{6. The remaining constructor obstruction}\label{the-remaining-constructor-obstruction}

The source's eight\_dvd condition is now fulfilled by the actual host. The source's common-field requirement is not. The induced type includes every R- and B-lift of each E₀ embedding. If every selected embedding agrees on x, agreement across all lifts first forces \(x\in E_0\); the automorphisms3 and17 in \(\Phi_0\) then force \(x\in\mathbb Q\), exactly as in K031(HC28). Hence \[
                       \mathrm{Fields.Kq}=\mathbb Q.
\tag{QH37}
\] This restates the already identified primitive-type obstruction on the new host. Neither the larger field nor the odd character changes it. In particular the external \(\mathbb Q(\sqrt{-7})\) containing the explicit prime generators cannot be relabeled as the source-defined common quadratic field.

The completed pieces are therefore an actual principal CM product host with8-divisible real degree, a literal odd Hecke character, all uniform even character conditions, explicit rank/prime/congruence data, and the unchanged permitted half-orbit. The source's full field/type record, exact later selected image, prescribed raw-Xi chain and quantitative MainX window remain separate. This finding makes no conclusion about GRH or a Kummer extension of the accepted theorem.

\subsection{7. Primary and exact input locators}\label{primary-and-exact-input-locators}

\begin{itemize}
\tightlist
\item
  Milne, J. S. (2020). \emph{Algebraic number theory} (Version3.08), Example2.41, printed p.39, for quadratic integer rings and discriminants; Lemma6.5, printed p.100, and Remark6.6(c), printed p.101, for the coprime-discriminant integral tensor and discriminant formula. Primary author PDF: \url{https://www.jmilne.org/math/CourseNotes/ANT.pdf}. The weighted trace identity itself is proved in full in (QH3)--(QH6).
\item
  AM-K029/030: the actual finite-unramified degree-seven CM host, odd character, trace-orthonormal inverse-different basis and full principal product identification. AM-K031 HC1--28: source numerical records, character norm pullback, permitted isogeny and exact common-field scope.
\item
  \nolinkurl{Formalism/LocalRingCardinality/Params.lean:6–22;} \nolinkurl{Formalism/LocalToGlobal/Columns/AbelianHostFieldConditions.lean:59–76,108–125,206–238,310–331,486–494,560–579;} \nolinkurl{Formalism/LocalToGlobal/UnitaryNorms/IsUnitaryNormOf.lean:66–94:} the exact unmodified field, character, degree and rank clauses. Paths are relative to work/research/alpoge\_source.
\item
  AM-K018 and ALP-LIFT-07/08: the actual principal line, finite normalization and full-CM reduced-gauge uniqueness for ordinary J powers. AM-R-006: fixed scalar c.~Carried DP59, AM-K022 and AM-DO-030/034: actual allowed half-orbit, finite jets and new-block selected construction.
\end{itemize}

The exact certificate \nolinkurl{degree_four_principal_host_repair_certificate.py} and its JSON output check both weighted trace matrices, integral symmetric actions and tensor commutation, new degree/discriminant arithmetic, the actual new rank/gap/partition lists, direct primality and full prime congruences, primitive-root witnesses and the29 splitting/index identities. They are finite arithmetic checks. The polarized variety, character and canonical-gauge proofs are the mathematical arguments above and their explicitly stated carried inputs; no Lean build or analytic window check is claimed.

\clearpage

\section{ALP-LIFT-18 --- exact literal-G character compatibility and an explicit admissible family}\label{alp-lift-18-exact-literal-g-character-compatibility-and-an-explicit-admissible-family}

Author: /root/lift\_accounting, GPT-6 Astra Pro, ALP\_MACHINE, 21 November2026UTC.

This result solves the even stage-character compatibility problem for the actual mixed host, with the source's specified G-residue characters. It gives a necessary and sufficient congruence, constructs the characters when it holds, and supplies numerical Harish--Chandra data satisfying that congruence and the displayed source gap conditions. The next local type and prescribed-kernel nonvanishing problems are not solved by the character construction.

The norm-one torus class-number statement of AM-R-009 is used as a checked input. The exact residue characters and conductor conventions are the ones audited in ALP-LIFT-17. The independently obstructed unramified odd character etaT is not supplied here.

\subsection{1. Host, the actual G half, and the root norm}\label{host-the-actual-g-half-and-the-root-norm}

Put \[
 E'=\mathbf Q(\zeta_{28}),\qquad F'=E'^+,\qquad
 K_q=\mathbf Q(\sqrt{-7}),\qquad
 \Phi'=\{1,3,9,15,17,23\}.
 \tag{G1}
\] The actual primitive CM type on \(J^2\) is \(\Phi'\), with \(\sigma_a(\zeta_{28})=\zeta_{28}^a\). The G half in the source is defined instead by a prime generator \(\pi_G\in K_q\), with \(\pi_G\bar\pi_G=q\), where \(q\equiv1\pmod{28}\) is prime. Its selected places are the six places of \(E'\) above the chosen prime of \(K_q\), namely the support of \(\bar\pi_G\). They give one member of every conjugate pair over \(q\).

The subgroup \(\operatorname{Gal}(E'/K_q)\) consists of those units modulo28 whose reduction modulo7 is a quadratic residue: \[
                    \Psi'=\{1,9,11,15,23,25\}.
\] For an elementary check of this quadratic subfield, put \(u=\zeta_7+\zeta_7^2+\zeta_7^4\) and \(v=\zeta_7^3+\zeta_7^5+\zeta_7^6\). Then \(u+v=-1\) and direct multiplication gives \(uv=2\), so \((u-v)^2=-7\). A residue-class automorphism fixes \(u-v\) exactly when it preserves \(\{1,2,4\}\), which gives the displayed subgroup. Directly, \[
                \sum_{a\in\Psi'}a=84,\qquad
       N_{E'/K_q}(\zeta_{28})=\zeta_{28}^{84}=1.
 \tag{G2}
\] The other coset gives the same norm-one assertion. It is important that this norm uses the induced G half, rather than the primitive CM type \(\Phi'\), whose exponent sum is68.

At the selected degree-one G places, all residue fields are canonically \(\mathbf F_q\) through their prime fields. For any global root \(z\in\mu_{28}\), (G2) gives \[
               \prod_{P\mid\bar\pi_G}(z\bmod P)=1
                                      \quad\hbox{in }\mathbf F_q^\times.
 \tag{G3}
\] This is simply the reduction of its relative norm.

\subsection{2. Prescribed compact characters and the exact congruence}\label{prescribed-compact-characters-and-the-exact-congruence}

Let \(g_0,g_1\) be two such G primes, with the source's selected halves. For each choose a faithful residue character \(\omega_\epsilon:\mathbf F_{g_\epsilon}^\times
 \to\mathbf C^\times\). At stage \(i\), the literal source condition is \[
        \chi_{i,P}=\omega_\epsilon^{\,2(i+1)}
                 \qquad(P\mid\bar\pi_{G,\epsilon}),
 \tag{G4}
\] and the desired Hecke character \(\mu_i\) is required to be unitary, trivial on \(\mathbb A_{F'}^\times\), tame, unramified outside \(g_0g_1\), and to satisfy \[
         \mu_{i,P}|_{\mathcal O_{E',P}^\times}=\chi_{i,P}^{-1}.
 \tag{G5}
\] These are the IsMu0 conditions, before the independent odd-character and arithmetic-kernel pins.

Let \(T'=\ker(N_{E'/F'}:\operatorname{Res}_{E'/F'}\mathbf G_m
\to\mathbf G_m)\), and \(K_0=(\widehat{\mathcal O}_{E'})^1\). At a relative split finite place, projection to its selected E-prime identifies the compact norm-one group with the local unit group. Define a compact character \(\rho_i\) of \(K_0\) by \[
       \rho_i(t)=
          \prod_{\epsilon=0}^1
          \prod_{P\mid\bar\pi_{G,\epsilon}}
                            \chi_{i,P}(t_P\bmod P)^{-1},
 \tag{G6}
\] and take it to be one at every other finite place. Equations(G3)--(G4) imply \[
                         \rho_i|_{\mu_{28}}=1.
 \tag{G7}
\]

For any integer vector \(\mathbf e_i=(e_{i,a})_{a\in\Phi'}\), let \[
          \beta_i(t_\infty)=
                    \prod_{a\in\Phi'}\sigma_a(t_\infty)^{-e_{i,a}}.
 \tag{G8}
\] This is a unitary character of \(T'_\infty\simeq U(1)^6\). AM-R-009 proves \[
 T'(\mathbb A_{F',f})=T'(F')K_0,\qquad
                          T'(F')\cap K_0=\mu_{28}.
 \tag{G9}
\] Therefore \(\beta_i\rho_i\) descends to an automorphic character of \(T'(\mathbb A_{F'})\) if and only if its restriction to the diagonal roots is one. By(G7), the exact condition is \[
 \boxed{\sum_{a\in\Phi'}a e_{i,a}\equiv0\pmod{28}.}
 \tag{G10}
\] For a uniform vector \(e_{i,a}=e_i\), this is equivalent to \[
                                7\mid e_i,
 \tag{G11}
\] since \(\sum\Phi'=68\) and \(\gcd(68,28)=4\).

There is no further norm-one global-unit condition hidden here: the full intersection in(G9) is precisely \(\mu_{28}\). This uses the actual class-number-one torus decomposition, rather than testing only one necessary root relation in an arbitrary CM field.

\subsection{3. Explicit construction and necessity for the actual Hecke character}\label{explicit-construction-and-necessity-for-the-actual-hecke-character}

Assume(G10). Given a finite norm-one idele \(t_f\), write \(t_f=\delta k\), with \(\delta\in T'(F')\), \(k\in K_0\), using(G9), and set \[
 \gamma_i(t_\infty,t_f)
       =\beta_i(\delta_\infty^{-1}t_\infty)\rho_i(k).
 \tag{G12}
\] If another decomposition is used, the two rational factors differ by an element of \(T'(F')\cap K_0=\mu_{28}\). Equations(G7),(G10) make(G12) independent of that choice. The formula is multiplicative, continuous, unitary, and trivial on the diagonally embedded \(T'(F')\). It is therefore an actual automorphic character.

Put \(j(x)=x/\bar x\) on ideles and define \[
                              \mu_i(x)=\gamma_i(j(x)).
 \tag{G13}
\] It is trivial on the base-field ideles. At infinity its restriction is \[
           \prod_{a\in\Phi'}
               \left(\frac{\sigma_a(x)}{|\sigma_a(x)|}\right)^{-2e_{i,a}},
 \tag{G14}
\] which is exactly HasInfQ with exponent vector \(-\mathbf e_i\): the source definition sets the integral angular exponent to twice the HasInfQ entry. At a selected G-prime, an idele supported in that E-component maps under \(j\) to the norm-one pair \((x,x^{-1})\). Equations(G6),(G13) therefore give precisely(G5), with its inverse sign. At other split places it is unramified.

At a nonsplit finite place the extension \(E'/F'\) is unramified, and every norm-one element is a unit. Moreover every norm-one unit is \(a/\bar a\) for a local unit \(a\): start with local Hilbert90, then multiply its witness by a base-field uniformizer power to make its valuation zero, using relative ramification index one. Since \(\rho_i=1\) there, \(\gamma_i\) is the trivial local torus character, and \(\mu_i\) is trivial on the entire local multiplicative group. All required outside-G unramified conditions follow. The G characters factor through residue units, so tameness is exact. Their conductor is the one in LIFT17(F19).

Conversely, suppose a genuine \(\mu_i\) satisfies these exact finite restrictions and HasInfQ\((-\mathbf e_i)\). The idele map \(j:\mathbb A_{E'}^\times\to T'(\mathbb A_{F'})\) is surjective by local Hilbert90, with kernel \(\mathbb A_{F'}^\times\); the global rational version is also surjective. Hence \(\mu_i\) descends to an automorphic \(\gamma_i\) on \(T'\). At each relative split finite place its restriction to \(K_0\) is forced by the unit character. At a nonsplit place the preceding integral Hilbert90 argument shows that an unramified \(\mu_i\) descends to a character trivial on the whole compact torus. Thus \(\gamma_i|_{K_0}=\rho_i\), and its infinity character is(G8). Rational triviality on \(\mu_{28}\) forces(G10).

This proves necessity as well as sufficiency. The conclusion is about the exact IsMu0 finite and infinity conditions, with its uniformizer values defined by the constructed character. It does not allow independently prescribing incompatible muVal entries.

\subsection{4. Concrete primes and a fixed coefficient field}\label{concrete-primes-and-a-fixed-coefficient-field}

For the actual minimal numerical choice \(k=6,m=2\), take \[
 \begin{array}{c|c|c}
 q&\pi_G&\pi_G\bar\pi_G\\ \hline
 113&1+4\sqrt{-7}&113\\
 281&13+4\sqrt{-7}&281 .
 \end{array}
 \tag{G15}
\] Both rational primes are \(1\bmod28\), so they split completely in \(E'\). The elements are algebraic integers in \(K_q\). Their conjugate ideals are coprime: a common prime would divide both their norm \(q\) and their sum \(2\) or26, which is impossible. Consequently they define the required six-element selected halves, with residue fields \(\mathbf F_q\). Their lower size bound is \[
            g_\epsilon-1>2\cdot6\cdot6=72.
\] This verifies the displayed G-prime conditions used here. No additional primitive-field D7 or etaT condition is inferred.

The faithful residue characters can take values in \[
                              K=\mathbf Q(\zeta_{560}),
 \tag{G16}
\] since \(560=\operatorname{lcm}(28,112,280)\). This is a fixed degree192 cyclotomic field, independent of all the stage weights. Formula(G12) also has algebraic finite values in this field: the rational norm-one factor \(\delta\) belongs to \(E'\), all its selected complex conjugates are Galois conjugates in \(E'\), and its absolute values are one. Thus the integer-power product \(\beta_i(\delta_\infty)\) belongs to \(E'\); the remaining factor is a root of unity in the displayed cyclotomic field.

This identifies actual finite character values, including values on uniformizers, by a formula. It does not assert that all those algebraic values are units at11. On compact units their values are roots of unity and hence units at every finite place. Values on noncompact determinant representatives require the separate normalization in the class presentation.

For these two stages the exact conductor is \[
       \mathfrak f(\mu_i)=31753\,\mathcal O_{E'},\qquad
                   N_{E'/\mathbf Q}\mathfrak f(\mu_i)=31753^{12}.
 \tag{G17}
\] This conductor is independent of the size of \(\mathbf e_i\). Its norm alone is not a complete H or a theta-row error term.

\subsection{5. Nonempty source gap locus with the required character congruence}\label{nonempty-source-gap-locus-with-the-required-character-congruence}

The compatibility condition does not force the numerical Harish--Chandra gap locus to be empty. Here is a fully explicit uniform example satisfying the numerical clauses of Hyp3. Keep \[
       k=6,\quad m=2,\quad r=7,\quad t=1795,\quad d=t+1=1796.
 \tag{G18}
\] The source rank inequalities hold: \(k,m\) are even, \(t\) is odd, \(m+3\le k\), \(r+5\le t\), \(3m<t\), and \(1801\le r+t=1802\). With the actual degree-six CM type count, take its corresponding order parameter3, so \(m\le6\le k\). The remaining displayed quadratic rank bound is immediate.

Choose any positive multiple \(C\) of7 satisfying \[
                  C>(r+mt)^3+k+t+10.
 \tag{G19}
\] At every one of the six embeddings define, with zero-based indices, \[
 \begin{aligned}
 u_i&=(2(k-i)+1)C+\tfrac12
                                &&(0\le i<k),\\
 v_i&=2(k-i)C                    &&(0\le i<k-1),\\
 b_j&=(4j-2)C,\qquad e_j=4jC     &&(0\le j<m),\\
 \eta&=-\tfrac12 .
 \end{aligned}
 \tag{G20}
\] The \(u_i\) are half-integral, the other parameters are integral, and \[
 u_i>v_i>u_{i+1},\qquad
 u_{k-1}=3C+\tfrac12>-b_0=2C,\qquad
 b_0<e_0<b_1<e_1.
\] Every required centred distance is at least \(\min(2C,4C-(t-1)/2)\), except for the positive \(u+e\) and \(v+e\) distances, which are still larger. Thus each exceeds \((t+3)/2\).

For NCGaps, the only false-sign element at stage0 is \(-b_1=-2C\). The other \(b\)-element is \(2C\), and all \(v_i\) are at least \(4C\), so every opposite-sign difference is at least \(4C>(r+mt)^3\). At stages1 and2 every element in the literal Dpar list has true sign, so the cross-sign condition is empty. The common \(\eta\)-shift at odd stages does not alter any difference. These statements directly check the source definition in DataTheta.lean538--555 and \nolinkurl{CayleySignature.lean106–118.}

The splitting weights are \(e_0=0,e_1=4C\), both divisible by7. Hence(G10) holds at both stages, and sections2--3 supply actual global even characters with exactly the prescribed G residue restrictions and no29 ramification. Their infinity exponents are the required \(-e_j\). The selected first partition is nonnegative: its smallest entry is \(3C-898>0\). The degree and other selected real data are therefore genuine finite quantities, although very large.

For a concrete integer, one may take \[
 C=46539458987,\qquad
 (e_0,e_1)=(0,186157835948),\qquad
 (b_0,b_1)=(-93078917974,93078917974).
 \tag{G21}
\] Here \((r+mt)^3=46539457173\), and(G19) holds. The symbols \(e_j\) in(G20)--(G21) are the Harish--Chandra splitting weights, not the DP44 error array.

The numerical \(\eta=-1/2\) matches an odd splitting character with infinity exponent \(+1/2\). It does not supply the original all-finite-unramified etaT pin, which AM-K028 separately obstructs on this host. Thus this paragraph verifies the stated numerical gap and character conditions, not a complete D7 instance.

\subsection{6. What is now fixed and what remains}\label{what-is-now-fixed-and-what-remains}

The even global characters are now explicit on a nonempty actual parameter locus, with the literal G conductor, exact residue characters, fixed coefficient field, and exact unit compatibility. The conductor29 character choice of LIFT14 was convenient for its spherical chain, but it is not needed for existence of these different stage characters.

At a split G place, changing to these characters changes the local theta parameters: the inserted rank-d block carries the stage character with its nontrivial residue restriction. This is exactly the kind of local change that the obstruction in LIFT17 requires. It is not obtained by choosing another vector in the old spherical constituent or by one uniform tame determinant twist.

To proceed to a chain with the prescribed column types, one must identify the ramified local constituents and show that the specified generic/determinant-block type and raw Xi map survive on the previous image. Their levels, projection masses and integral content remain to be calculated. Neither the existence of the characters nor their fixed finite conductor proves that local type occurrence, global prescribed-kernel nonvanishing, an integral gain, or a shorter DP44 chain.

\subsection{Source and input locators}\label{source-and-input-locators}

\begin{itemize}
\tightlist
\item
  AM-R-009, first\_lift\_unramified\_host.md, equationR9.15 and its norm-one ideal proof: actual class-number-one torus decomposition and the complete root-unit intersection. Masley(1976) remains the checked class-number input.
\item
  ALP-LIFT-17, sections4--7: exact source residue characters, G-prime conditions, size bounds, tameness and conductor.
\item
  \nolinkurl{Formalism/LocalToGlobal/UnitaryNorms/IsUnitaryNormOf.lean,} lines66--94: LevelLe, UnrAt, TameAt, HasInf and HasInfQ.
\item
  \nolinkurl{Formalism/LocalToGlobal/Columns/AbelianHostFieldConditions.lean,} lines205--224,273--306,315--323 and488: G halves, chi\_om, IsMu0 and g\_big.
\item
  \nolinkurl{Formalism/LocalRingCardinality/DataTheta.lean,} lines501--555, and \nolinkurl{Formalism/SharpRestriction/CuspSpanNormBounds/CayleySignature.lean,} lines106--125: the exact numerical parameter and gap conditions.
\item
  AM-K028: the separate odd-character obstruction is carried, not used as a substitute for the even-character construction.
\end{itemize}

All recovered source paths are relative to work/research/alpoge\_source, pinned to commit 7cbf7a0c32fca1badf05931879a37c3f8a2b03b4. No Lean build or complete source constructor is claimed.

\clearpage

\section{AM-DO-032: the actual terminal Y-family, exact Fourier Gram, and its projection denominator}\label{am-do-032-the-actual-terminal-y-family-exact-fourier-gram-and-its-projection-denominator}

Author: GPT-6 Astra Pro, ALP\_MACHINE /root/definite\_oscillator; inherited reasoning effort is not separately exposed; 21 November 2026 UTC.

This is a further calculation on the literal xiLast of AM-DO-031. It varies only the Y-parameter already present in the recovered source. It identifies its finite Fourier family, the exact Gram matrix, and the denominator of its orthogonal projector. Its arithmetic denominator is for this specified finite family; it does not identify the global MainX terminal detector lattice.

\subsection{1. The literal finite parameter set}\label{the-literal-finite-parameter-set}

Retain the notation and all source hypotheses of AM-DO-031: \[
 a=N(m-1),\quad b=N(m),\quad r=k+m-1,\quad
 \mathsf M=\mathrm{cst.M},\quad
 q_P=|\mathcal O_E/P|,\quad d_P=v_P(q_P).
 \tag{DO32.1}
\] In particular the T-places and rowSG are disjoint, and G denotes the G branch while Q denotes its complement inside rowSG. For the actual standard TwDat base points, the local frame lattices are standard.

The source definitions are:

\begin{longtable}[]{@{}
  >{\raggedright\arraybackslash}p{(\columnwidth - 2\tabcolsep) * \real{0.5000}}
  >{\raggedright\arraybackslash}p{(\columnwidth - 2\tabcolsep) * \real{0.5000}}@{}}
\toprule\noalign{}
\begin{minipage}[b]{\linewidth}\raggedright
Source
\end{minipage} & \begin{minipage}[b]{\linewidth}\raggedright
Locator
\end{minipage} \\
\midrule\noalign{}
\endhead
\bottomrule\noalign{}
\endlastfoot
\nolinkurl{Formalism/LocalToGlobal/Columns/FiniteSumNonzeroLargeDepth.lean} & ybSet and repY,1123--1131; ftY,1150--1156 \\
\nolinkurl{Development/AbelianVarietyTheory/Definitions/AutomorphicMembership/XiAtPenultimateEqXiCpl1.lean} & actual terminal point and representative,124--130; xiOver/xiLast,146--157 \\
\nolinkurl{Development/AbelianVarietyTheory/Definitions/GlobalTwisting/IntegralGLSuccEmbeddingMemParahoricTimesPrincipa.lean} & locTop,391--393; representative bound,531--547 \\
\end{longtable}

For fixed source decoration, ybSet is the set of cosets of the product-local period module in the product-local support module. Therefore it is canonically the finite set \[
 \mathcal Y=
 \prod_{P\in G}(q_P^{-1}\mathcal O_P^a/\mathcal O_P^a)
 \ \times
 \prod_{P\in Q}
     (q_P^{-\mathsf M}\mathcal O_P^a/
                      q_P^{\mathsf M}\mathcal O_P^a).
 \tag{DO32.2}
\] This statement concerns cosets, not the particular epsilon representative chosen by repY. To prove it, send a rational coset to its local residue classes. Injectivity is exactly the definition of its period module as the intersection of the local conditions. Surjectivity follows by CRT after clearing the finitely many denominators, as in AM-DO-031 Section5. The prescribed frames identify the localized source lattice with each standard local lattice. Thus every tuple in (DO32.2) is an actual allowed Y.

Changing repY within its coset changes its local component by \(\mathcal O_P^a\) at G and by \(q_P^{\mathsf M}\mathcal O_P^a\) at Q. On the supports of ftY, the resulting character argument is integral. LocCst.Ok.psi\_int therefore proves that the local finite function is unchanged. No representative phase is suppressed.

\subsection{2. The exact local Fourier calculation}\label{the-exact-local-fourier-calculation}

At a G-place, \(d_P=1\), and the source limited-depth character induces a nontrivial additive character on \(q^{-1}\mathcal O_P/\mathcal O_P\). Since q is the rational prime and the residue degree is1, its pairing with \(\mathcal O_P/q\mathcal O_P\) is perfect. The last-row functions are exactly \[
 \varphi_y(\eta)=1_{\mathcal O_P^a}(\eta)
                       \psi_P(-\eta y),
 \qquad y\in q^{-1}\mathcal O_P^a/\mathcal O_P^a.
 \tag{DO32.3}
\] With additive measure giving \(\mathcal O_P^a\) volume1, they are an orthonormal basis of the functions supported there and constant modulo \(q\mathcal O_P^a\): \[
 \langle\varphi_y,\varphi_{y'}\rangle
          =1_{\{y=y'\}}.
 \tag{DO32.4}
\] Indeed the integral of a nontrivial character on a compact additive group is0: translation by an element at which it is not1 multiplies that integral by a nontrivial scalar. For a trivial character the integral is the group's volume1.

At a Q-place the source character is only required on \(q^{-D_P}\mathcal O_P\), where \(D_P\ge2\mathsf M\). Its annihilator of the local ring within this domain is an \(\mathcal O_P\)-module: \[
 \mathfrak a_{\psi,P}
 =\{z\in q^{-D_P}\mathcal O_P:
                  \psi_P(z\mathcal O_P)=1\}
 =P^{c_P},
 \qquad -d_P<c_P\le0.
 \tag{DO32.5}
\] Here and below P denotes its maximal ideal in the completion. To justify every part of (DO32.5), the set is closed under addition and multiplication by \(\mathcal O_P\), contains \(\mathcal O_P\), and lies in a fixed fractional ideal. Hence it is a fractional ideal \(P^{c_P}\). The source condition \(\psi_P(q^{-1})\ne1\) excludes \(q^{-1}\) from it, so \(c_P>-d_P\), while its inclusion of the ring gives \(c_P\le0\). The same conductor ideal is therefore already determined on the smaller argument range needed for ftY. This reasoning uses only the finite additive quotient and does not extend the source character arbitrarily.

The actual Q functions, including the source scalar, are \[
 \varphi_y(\eta)=q^{-\mathsf M a}
 1_{q^{-\mathsf M}\mathcal O_P^a}(\eta)
                      \psi_P(-\eta y),
 \qquad
 y\in H_P:=
 q^{-\mathsf M}\mathcal O_P^a/
                       q^{\mathsf M}\mathcal O_P^a.
 \tag{DO32.6}
\] The pairing radical is \[
 R_P=
 q^{\mathsf M}P^{c_P}\mathcal O_P^a/
                       q^{\mathsf M}\mathcal O_P^a,
 \qquad |R_P|=q^{-c_Pa}.
 \tag{DO32.7}
\] The notational product \(P^{c_P}\mathcal O_P^a\) means componentwise multiplication by that fractional ideal. The bounds in (DO32.5) imply that its numerator is inside the support module. Two allowed y's give exactly the same function if their difference is in \(R_P\). Otherwise their inner product is0. More precisely, \[
 \boxed{\quad
 \langle\varphi_y,\varphi_{y'}\rangle
 =q^{(d_P-2)\mathsf M a}
                1_{\{y-y'\in R_P\}}.
 \quad}
 \tag{DO32.8}
\]

To prove (DO32.8), the character integral is nonzero precisely when \((y-y')q^{-\mathsf M}\mathcal O_P\subset P^{c_P}\), coordinate by coordinate. This is exactly (DO32.7). Its nonzero value before the amplitude is the volume \(q^{d_P\mathsf M a}\); the squared amplitude is \(q^{-2\mathsf M a}\). Their product gives the displayed factor. If the character is trivial on the full support, the two functions agree pointwise, not just up to a scalar, since their value at \(\eta=0\) is the same nonzero number.

After removing the duplicate parameters, the number of Q functions is \[
 |H_P/R_P|=q^{(2d_P\mathsf M+c_P)a}.
 \tag{DO32.9}
\] They span exactly the functions supported on \(q^{-\mathsf M}\mathcal O_P^a\) and constant modulo \(q^{\mathsf M}P^{c_P}\mathcal O_P^a\). This follows from finite Fourier inversion for the resulting perfect quotient pairing.

In the degree-one case \(d_P=1\), (DO32.5) forces \(c_P=0\). There are no duplicates, the common norm square is \(q^{-\mathsf M a}\), and the matrix of values \[
 q^{-\mathsf M a}\psi_P(-\eta y),
 \qquad \eta,y\in H_P,
 \tag{DO32.10}
\] is unitary for counting measure: \(|H_P|=q^{2\mathsf M a}\). This explains the literal scalar in ftY without replacing it by a guessed Haar normalization.

\subsection{3. Gram matrix of the prescribed terminal family}\label{gram-matrix-of-the-prescribed-terminal-family}

Let \(f_Y\) be the actual finite Heisenberg realization of xiLast for Y, using AM-DO-031. All upper-matrix and T-incidence factors are fixed as Y varies. The local positions at distinct selected places are independent. Consequently the preceding last-row inner products tensor with the fixed nonzero upper and T factors.

Let \[
 \overline{\mathcal Y}=
 \prod_{P\in G}
        (q_P^{-1}\mathcal O_P^a/\mathcal O_P^a)
 \times\prod_{P\in Q} H_P/R_P,
 \qquad n_Y=|\overline{\mathcal Y}|.
 \tag{DO32.11}
\] Choose one representative for each reduced class. The remaining choices give identical finite vectors. Then \[
 \boxed{\quad
 \langle f_Y,f_{Y'}\rangle_f
                =\kappa\,1_{\{Y=Y'\}},
 \qquad
 \det\operatorname{Gram}(f_Y)=\kappa^{n_Y},
 \quad}
 \tag{DO32.12}
\] where the nonzero rational number \(\kappa\) is exactly the unscaled finite norm square computed in AM-DO-031: \[
\begin{split}
 \kappa={}&
 \prod_{P\in T}q_P^{-a}
 \prod_{P\in G}
    \left[
     q_P^{-C_G}
       \prod_{i=1}^{t}(1-q_P^{-i})^{m-1}
    \right]\\
 &{}\cdot
 \prod_{P\in Q}
    \left[
     q_P^{(d_P-2)\mathsf M a-d_P r(b-1-r)}
       \prod_{i=1}^{r}(1-q_P^{-i})
    \right],\\
 C_G={}&r(r+1)/2+k+mrt+m(m-1)t^2/2.
\end{split}
\tag{DO32.13}
\] The scalar is unchanged by the actual Y-representative. A prescribed common chart scalar w multiplies the Gram by \(w\overline w\); it must be retained. Nonstandard lattice frames retain their actual Jacobian as in AM-DO-031.

Thus this is a source-defined orthogonal family, not a new arbitrary choice of finite inputs. It is a basis of its actual finite span. The identity (DO32.12) is a finite oscillator pairing, not an automorphic Petersson pairing.

\subsection{4. Integral saturation versus the exact pairing denominator}\label{integral-saturation-versus-the-exact-pairing-denominator}

Assume all selected residue characteristics differ from11, and work over a coefficient valuation ring containing the finite character values. Conjugate coefficient lattices are paired when the chosen 11-place is not fixed by the coefficient involution. Every root of unity remains an11-adic unit.

The lattice generated by the distinct \(f_Y\) is saturated in the finite function module on its ambient support/period window, intersected with its characteristic-zero span. Here is a direct argument. Finite Fourier inversion uses only orders of finite p-groups with \(p\ne11\); its denominators are units. The invariant submodule under the radicals in (DO32.7) is split by the corresponding averaging idempotent, whose denominator is again a unit. Finally the fixed upper/T factor has a unit-valued coordinate: use the local identity and incidence matrices of AM-DO-031 Section5. Evaluation at that coordinate splits tensor multiplication by that fixed factor, after multiplying by its unit inverse. This proves saturation.

Nevertheless the pairing on this saturated lattice need not be unimodular. From (DO32.13), \[
 \boxed{\quad
 v_{11}(\kappa)
 =\sum_{P\in G}(m-1)
                  \sum_{i=1}^{t}v_{11}(q_P^i-1)
   +\sum_{P\in Q}\sum_{i=1}^{r}
                         v_{11}(q_P^i-1).
 \quad}
 \tag{DO32.14}
\] All powers of q are units at11, so this formula holds for the general-Q calculation as well. For \(o_P=\operatorname{ord}_{11}(q_P)\mid10\) and \(h_L=\lfloor L/o_P\rfloor\), the exact elementary lifting-the-exponent formula gives \[
 \sum_{i=1}^{L}v_{11}(q_P^i-1)
   =h_L\,v_{11}(q_P^{o_P}-1)+v_{11}(h_L!).
 \tag{DO32.15}
\] Indeed the terms vanish unless \(o_P\mid i\); for \(i=o_Pj\) their valuation is \(v_{11}(q_P^{o_P}-1)+v_{11}(j)\).

In the actual rank range, \(t\ge13\). Thus every G-place away from11 contributes at least \(m-1>0\) to (DO32.14), because \(o_P\le10\). This is compatible with the coefficient primitivity of every \(f_Y\): coefficient primitivity and a unimodular restricted pairing are different conditions.

Apply the already proved detector calculation AM-R-002 to this specified finite lattice. For an individual basis vector \(f_Y\), its pairing ideal against this same lattice is exactly \(\kappa R\). Its dual vector within this span is \(\kappa^{-1}f_Y\), with exact denominator valuation \(v_{11}(\kappa)\) and complex norm \(\kappa^{-1/2}\). With the Hermitian pairing linear in its first argument, the orthogonal projector onto the whole span is \[
 \Pi_{\mathrm{term}}
 =\kappa^{-1}\sum_{Y\in\overline{\mathcal Y}}
                  f_Y\langle\,\cdot\,,f_Y\rangle_f.
 \tag{DO32.16}
\] It has that exact denominator on the ambient finite function lattice, with its volume normalizations retained. To see that the upper bound is attained, factor the space into the fixed upper/T vector and its full last-row Fourier subspace. A unit-valued coordinate of the fixed factor gives a unit numerator in its rank-one projector. The last-row projection has unit denominators at11. Its diagonal at an allowed last-row position is a p-power unit. Hence no smaller11-power can clear the displayed projector.

This is a restriction on a detector required to remain in this terminal finite span or on its orthogonal projector. It is not a denominator lower bound for every detector in the entire finite oscillator. Since the span is saturated, an integral functional on it can extend to the ambient coefficient module; that extended functional need not be represented by an integral vector in this same orthogonal subspace. Nor is the global MainX detector lattice already identified with this finite span.

\subsection{5. Exact effect of applying a theta operator}\label{exact-effect-of-applying-a-theta-operator}

Let T denote any fixed actual selected theta operator out of this finite input space, with the other archimedean and automorphic inputs fixed and their actual measures retained. Use the source Hermitian convention, linear in the first argument. Equations (DO32.12) and (DO32.16) give the genuine adjoint identity \[
 T\Pi_{\mathrm{term}}T^\dagger
 =\kappa^{-1}
    \sum_{Y\in\overline{\mathcal Y}}
          (Tf_Y)\langle\,\cdot\,,Tf_Y\rangle.
 \tag{DO32.17}
\] Restrict the target to the finite-dimensional image of the finite input window, so every linear scalar detector has a Riesz adjoint. For any such detector \(\ell\), this says \[
 \sum_Y|\ell(Tf_Y)|^2
 =\kappa\,
   \|\Pi_{\mathrm{term}}T^\dagger\ell^\dagger\|^2.
 \tag{DO32.18}
\] This is exact positivity and finite Plancherel, including its scalar. It proves nonvanishing of one actual allowed Y precisely when the displayed projected adjoint is nonzero.

The new calculation fixes the projection, its rank, Gram and arithmetic denominator. It does not assert that a prescribed prior theta image has nonzero projection to this finite span. A zero projected adjoint is still possible; the source's local types and arithmetic row must decide it. Likewise the actual terminal mass equation uses class stabilizers, branching maps and the specified row pairing. Substituting the finite Gram scalar for that global identity would be unjustified.

The calculation therefore removes a concrete finite Fourier and normalization ambiguity, and supplies a computable factor for a future matching proof. It does not replace the remaining selected-image or terminal-mass equation by the mere fact that xiLast is a nonzero function.

\subsection{Review status}\label{review-status-1}

Author whole-proof review: PASS, 2026-11-21T23:28:15.931357+00:00. The complete saved proof was read; source-linear-first Hermitian rank-one operator notation was corrected during that review. Exact quotient/radical, Gram, saturation, projector denominator and theta-adjoint scope checked. Independent review remains separate; no Lean build claimed.

\clearpage

\section{AM-DO-033: an actual finite Steinberg first input, with exact extraction and costs}\label{am-do-033-an-actual-finite-steinberg-first-input-with-exact-extraction-and-costs}

Author: GPT-6 Astra Pro, ALP\_MACHINE /root/definite\_oscillator, 21 November2026UTC; inherited reasoning effort not separately exposed.

This is a focused extension of AM-DO-030 and AM-R-012. It supplies a nonzero selected first lift whose source input has the actual initial generic-unipotent type at the auxiliary G places. It uses a specific finite Steinberg vector, one actual permitted29 half-translate, and a compact matrix coefficient at each G place. It computes the finite Schwartz norm and the source mass exactly. The new finite vector is not identified with the prescribed raw PhiG array. No later restricted-image assertion is included.

The global construction uses the already checked common host E'=Q(zeta28), A=J², its trace-one line, and the selected even edge U(a)→U(b), with a=k and b=k+t+1. A valid common splitting character with the required infinity type is assumed fixed. It is unramified off G and29, tame at G, and has at most the residue-level compact conductor at29 used in R012. It may have the literal tame G restrictions constructed in LIFT18, in which case it is unramified at29. Replacing the splitting while keeping its infinity exponents does not change the identity kernel value or the Heisenberg reading; its finite compact action is explicitly accounted for below. The local argument works over every finite residue field; the global extension to a different host additionally requires that host's actual mixed-theta/DO030 comparison, not just an abstract extension of its number field.

\subsection{1. A primitive generic vector with a known metric}\label{a-primitive-generic-vector-with-a-known-metric}

Fix one odd relative split auxiliary place with residue field κ=F\_q, q prime to11·29. Put

\[
 \bar G=\mathrm{GL}_a(\kappa),\quad B=TU,\quad
 n_U=a(a-1)/2,\quad d_\tau=q^{n_U}.
 \tag{DO33.1}
\]

Choose the precise nondegenerate character ψ\_U of U required by the initial column convention. For the source this is the character λ\_X,0; it is nontrivial on every simple-root subgroup. Let \(\mathcal F=\mathbf C[\bar G/B]\), with its delta basis orthonormal. The standard Steinberg submodule τ is irreducible over C, has dimension d\_tau, and has a one-dimensional B-fixed space. These are the classical results of Steinberg(1957), Theorems1--2, pp.348--350, and Geck(2015), Theorem2.1 and Corollary2.3, pp.3--4. They are used as published results.

Let w0 be the longest permutation, and define

\[
 v=\sum_{u\in U}\psi_U(u)^{-1}[uw_0B],\qquad
                         u_0=q^{-n_U/2}v.
 \tag{DO33.2}
\]

Then

\[
 v\in\tau,\quad\tau(u)v=\psi_U(u)v,
 \quad\|v\|^2=q^{n_U},\quad\|u_0\|=1.
 \tag{DO33.3}
\]

Here is the explicit membership and normalization check. The Steinberg vector in the flag model is \(e=\sum_w(-1)^{\ell(w)}[wB]\). Apply \(\sum_{u\in U}\psi_U(u)^{-1}u\) to it. For w≠w0 there is a simple root subgroup in U∩wBw\^{}−1 on which ψ\_U is nontrivial: otherwise w\^{}−1 would send every simple root to a negative root, forcing w=w0. Summing over that subgroup cancels its entire contribution. For w0 its U-stabilizer is trivial, so the result is \((-1)^{\ell(w_0)}v\). This proves membership. Left translation reindexes the sum and gives the asserted character. The q\^{}\{n\_U\} flags in the longest cell are distinct, so the displayed norm is the ordinary flag norm. No orthogonality is being asserted for the different classical basis \(\{ue\}\).

Set, with Hermitian inner product conjugate-linear in its first argument,

\[
 m(g)=\langle u_0,\tau(g)u_0\rangle.
 \tag{DO33.4}
\]

Then m(1)=1, \textbar m(g)\textbar≤1, and every m(g) belongs to q\^{}−n\_U times the ring generated by the values of ψ\_U. This follows because g permutes the orthonormal flag basis in (DO33.2). In particular m(g) is integral at11, and m(1) is a unit there. Also m(gu)=ψ\_U(u)m(g), giving the exact right-generic character needed below.

\subsection{2. An actual residue-matrix Schwartz vector}\label{an-actual-residue-matrix-schwartz-vector}

Use the DO020/022 unitary split position model \(M_{b,a}(F_v)\), with measure giving its integral lattice volume1. Let η be the character on F\_v\^{}× such that the common split character is \((\eta,\eta^{-1})\) in the selected E half. It is tame at this place. The actual source action is

\[
 \omega(g)f(X)=\eta(\det g)^{-1}|\det g|^{b/2}f(Xg).
 \tag{DO33.5}
\]

Indeed DO20.11 is the inverse pullback by the Hom map \(X\mapsto Xg^{-1}\). On the integral compact the absolute-value factor is1 and η depends only on the residue determinant. The inverse-argument operator gives the same spherical vacuum coefficient but is not used as the source action here.

Put \(X_0=\begin{bmatrix}I_a\\0\end{bmatrix}\). The orbit \(X_0\bar G\subset M_{b,a}(\kappa)\) is free. For v'∈τ define a locally constant function, supported on its inverse image in \(M_{b,a}(\mathcal O_v)\), by

\[
 \mathcal E(v')(X)=
 \begin{cases}
 \eta(\det t)\langle u_0,\tau(t)v'\rangle,
                    &X\bmod\mathfrak p=X_0t,\ t\in\bar G,\\
 0,&\text{otherwise}.
 \end{cases}
 \tag{DO33.6}
\]

Freeness makes this unambiguous. The actual source action satisfies

\[
                  \omega(g)\mathcal E(v')
                        =\mathcal E(\tau(\bar g)v')
                      \qquad(g\in\mathrm{GL}_a(\mathcal O_v)).
 \tag{DO33.7}
\]

To check it, evaluate at X=X0t modulo the maximal ideal. The scalar η(det g)\^{}−1 in(DO33.5) cancels the new factor η(det g) in η(det(tg)); the matrix coefficient becomes \(\langle u_0,\tau(t)\tau(g)v'\rangle\). Off the orbit, both sides vanish. Thus the determinant correction in(DO33.6) is fixed by the prescribed splitting; it is not an arbitrary type twist. In the literal source situation η=χ\_0\^{}−1.

Each residue cell has additive measure q\^{}−ab. Finite Schur orthogonality therefore gives the exact identity

\[
 \|\mathcal E(v')\|^2=\gamma_q\|v'\|^2,
 \qquad
 \gamma_q=q^{-ab}\frac{|\mathrm{GL}_a(\kappa)|}{d_\tau}.
 \tag{DO33.8}
\]

In particular the raw choice \(f_q=\mathcal E(u_0)\) has \(f_q(X_0)=1\), norm squared γ\_q, and a pure source τ-type. Its values are algebraic and11-integral. Its unit-norm version \(f_q^\circ=\gamma_q^{-1/2}f_q\) is an isometric τ vector. This normalization is algebraic after adjoining the positive square root; its possible11-adic denominator is retained in section5.

\subsection{3. The exact auxiliary Heisenberg expansion}\label{the-exact-auxiliary-heisenberg-expansion}

The finite model is the normalized actual vacuum model of DO029, not an unrelated finite permutation representation. Here support is L, period is \(\mathfrak pL\), and the finite Fourier quotient has size q\^{}ab. Consequently

\[
 f_q=\sum_{y\in(\mathfrak pL)^\perp/L^\perp}b_{q,y}W_y\phi_{0,v},
 \qquad
 b_{q,y}=q^{-ab}\sum_{X\in M_{b,a}(\kappa)}
                              f_q(X)\psi_v(-B(X,y)).
 \tag{DO33.9}
\]

All actual Fourier and trace-one phases are included in this formula, as in DO29.11--13. Its parameters have denominator at most the rational auxiliary prime when the place is of degree one. They are supported at this auxiliary place. The nonzero coefficients number at most q\^{}ab, and

\[
 \sum_y|b_{q,y}|^2=\gamma_q,
 \qquad \min_{b_{q,y}\ne0}v_{11}(b_{q,y})=0.
 \tag{DO33.10}
\]

The first equality is Parseval with the vacuum-normalized measure. For the second, every coefficient is11-integral, since q is a unit and all character values are roots of unity. If every coefficient had positive valuation then Fourier inversion at X0 could not give f\_q(X0)=1. This also proves that the auxiliary operator acts nontrivially on the vacuum, the actual hypothesis needed in DO030.

Tensor this construction over all selected G places, and keep the vacuum at the other auxiliary places. Write

\[
 \tau_G=\bigotimes_v\tau_v,\quad
 d_G=\prod_vq_v^{a(a-1)/2},\quad
 \Gamma_G=\prod_v\gamma_{q_v},\quad
 s_G=\#\{\text{nonzero tensor Fourier coefficients}\}.
 \tag{DO33.11}
\]

The raw finite auxiliary vector has norm squared Γ\_G, its nonzero coefficients have minimum11-valuation0, and \(s_G\le\prod_vq_v^{ab}\). All these are finite, explicit quantities independent of the selected polynomial degree D. The corresponding unit auxiliary vector has norm1 exactly; its coefficient minimum is \(-\tfrac12v_{11}(\Gamma_G)\).

For the actual degree12 host, six selected relative split places occur at each of g0,g1. Thus

\[
 \begin{split}
 d_G&=(g_0g_1)^{3a(a-1)},\\
 \Gamma_G&=(g_0g_1)^{6(-ab+a(a+1)/2)}
       \prod_{\epsilon=0}^1\prod_{i=1}^a(1-g_\epsilon^{-i})^6,\\
 s_G&\le(g_0g_1)^{6ab}.
 \end{split}
 \tag{DO33.12}
\]

For a host with f selected places above each completely split G prime, the same local product gives \(d_G=(g_0g_1)^{f a(a-1)/2}\), \(\Gamma_G=(g_0g_1)^{f(-ab+a(a+1)/2)}
\prod_{\epsilon,i}(1-g_\epsilon^{-i})^f\), and \(s_G\le(g_0g_1)^{fab}\). These are local counting identities, conditional globally on that host's actual theta-family construction.

\subsection{4. A nonzero selected first kernel and a genuine source input}\label{a-nonzero-selected-first-kernel-and-a-genuine-source-input}

Fix the actual primitive selected polynomial P of degree D at infinity, as in R012, with factorial content bound B\_P from DO30.22. Apply DO030 to the raw auxiliary vector above. Choose

\[
 e\ge1,\quad M=29^e\ge3s_G(D+1)-1.
 \tag{DO33.13}
\]

Then one individual permitted DP59 half-orbit translate λ has

\[
 a_\lambda\ne0,\quad v_{11}(a_\lambda)\le B_P,\qquad
 \vartheta_{r_P\mathcal O_G\phi_\lambda}(1,1)
                                      =C^{2ab}a_\lambda.
 \tag{DO33.14}
\]

This conclusion uses simultaneous jets on the actual allowed section lattice, with the fibre units and the common period already established in DO030. It is not inferred merely from nonvanishing of the finite Schwartz vector. The29 operation commutes with every G-local compact action, so the resulting kernel still has the pure finite source type τ\_G and the selected real source type \(\sigma=\rho_0^\vee\).

At29 take the same level K\_a(M) as R012. At each G place retain the full integral compact, and let the source function vary there according to its finite matrix coefficient. At every other place retain the original integral compact. Let K* denote this product. Its volume is I\_a(M)\^{}−1, with

\[
 I_a(M)=M^{6a^2}\prod_{i=1}^a(1-29^{-i})^6.
 \tag{DO33.15}
\]

The rational stabilizer at the identity class is still trivial: it is finite by definiteness and contained in a torsion-free principal29 congruence. No extra principal congruence at G is necessary to define a matrix-coefficient function on this full compact class.

Let u\_inf be the unit real source vector used in R012, let u\_G=⊗u0, and put

\[
 m_\infty(k)=\langle u_\infty,\sigma(k)u_\infty\rangle,
 \qquad m_G(k_G)=\langle u_G,\tau_G(k_G)u_G\rangle.
 \tag{DO33.16}
\]

On the identity finite K* class define \(F(k_\infty,k_f)=\overline{m_\infty(k_\infty)m_G(k_G)}\), and set F=0 on the other classes. The trivial stabilizer makes this an unambiguous automorphic smooth function. It is cuspidal because the rational source group is definite. Its source real type is ρ0 and its finite G-type is τ\_G\^{}∨.

Two applications of compact Schur orthogonality, retaining the full compact probability measures at G, give

\[
 \boxed{
 \Theta^{\rm bil}_{r_P\mathcal O_G\phi_\lambda}(F)(1)
       =\frac{C^{2ab}a_\lambda}{I_a(M)d_\rho d_G}\ne0,
 \qquad
 \|F\|_2^2=\frac1{I_a(M)d_\rho d_G},\quad\|F\|_\infty=1.}
 \tag{DO33.17}
\]

To verify the extraction rather than just its norm, write the kernel on the supporting class as \(L(\sigma(k_\infty)u_\infty\otimes\tau_G(k_G)u_G)\), where L at the original tensor is the value in(DO33.14). Averaging each factor against its conjugated matrix coefficient returns that tensor divided by d\_rho d\_G. The support volume is I\_a(M)\^{}−1. This proves the first identity, and averaging the squares proves the second.

For every initial source unipotent u at a G place, (DO33.4) implies F(gu)=ψ\_U(u)\^{}−1F(g). This is the literal adelic initial column convention: \nolinkurl{Columns/RecLevel.lean}, lines50--52, defines \nolinkurl{levChr=lam ∘ levRed} and \nolinkurl{levTyp=levChr⁻¹}. Thus the sign agrees with the source's initial generic-unipotent type. The present conclusion does not assert that F meets every other decoration or column constraint of a full \nolinkurl{ColV} datum.

The nonzero theta image in(DO33.17) has at least one irreducible automorphic source constituent with this finite Steinberg K-type. Such a constituent has a nonzero Iwahori- fixed vector because τ\^{}B is one-dimensional. The projection to an irreducible automorphic constituent commutes with the compact actions; a nonzero compact irreducible subspace remains injective in any nonzero projected copy. Hence this is a statement about an actual source constituent, including its automorphic occurrence, not just a local abstract type.

\subsection{5. Norm-one formulas and arithmetic costs}\label{norm-one-formulas-and-arithmetic-costs}

Let \(S_G=\sqrt{\Gamma_G}>0\), and use the same positive square root at the chosen complex embedding. The raw finite Schwartz vector has norm S\_G, whereas the source input F has the norm in(DO33.17). With unit-normalized finite input and source input, the exact value is

\[
 \Theta^{\rm bil}_{r_P\,S_G^{-1}\mathcal O_G\phi_\lambda}
      (\sqrt{I_a(M)d_\rho d_G}\,F)(1)
       =\frac{C^{2ab}a_\lambda}
                    {S_G\sqrt{I_a(M)d_\rho d_G}}.
 \tag{DO33.18}
\]

The real Schwartz norm \(\|r_P\phi_{0,\infty}\|\), including its Gaussian covolume and hook factor, remains R011/R012's exact quantity. If that factor is also made unit, (DO33.18) is divided by that norm.

At11, since every q\_v is a unit,

\[
 v_{11}(d_G)=0,\qquad
 v_{11}(\Gamma_G)=
      \sum_v\sum_{i=1}^a v_{11}(q_v^i-1).
 \tag{DO33.19}
\]

Thus making the finite input unit norm has an exact local scalar cost \(\tfrac12v_{11}(\Gamma_G)\). It is not automatically a unit merely because the G primes differ from11. For the primitive raw input, however, the selected coefficient bound is

\[
 v_{11}\!\left(C^{-2ab}\Theta^{\rm bil}(F)(1)\right)
       \le B_P-v_{11}(I_a(M)d_\rho).
 \tag{DO33.20}
\]

The finite Steinberg dimension introduces no11 valuation. Its archimedean source normalization cost is the additional \(\tfrac12\log d_G\), and the finite input normalization has the explicit S\_G factor. For fixed ranks and fixed G primes all these are fixed numbers. Equation(DO33.13) gives \(\log M=O(\log(D+2))\), with the potentially large explicit coefficient s\_G retained. If the auxiliary primes or ranks vary, their displayed factors must also vary in the cost.

This proves an initial auxiliary-type construction with the chosen ramified splitting characters when those are used. It supplies neither a lower content bound for every row nor the next prescribed raw-Xi map. It also leaves the odd global etaT problem, later integral image saturation, Kval and the MainX terminal mass outside its scope.

\subsection{Sources and proof status}\label{sources-and-proof-status}

The new steps are the explicit generic flag vector metric, its determinant-corrected free-orbit embedding, its actual finite Fourier coefficients, and the finite compact extension of R012's class extraction. The mixed-family and jet comparison are the checked DO020/022/026/029/030 and R012 inputs. LIFT18 supplies the exact even character construction when its root-unit congruence holds.

Steinberg, R. (1957). Prime power representations of finite linear groups II. \emph{Canadian Journal of Mathematics,9},347--351. \url{https://doi.org/10.4153/CJM-1957-041-1} . Theorems1--2, pp.348--350, read directly in the publisher PDF.

Geck, M. (2015). On the ℓ-modular composition factors of the Steinberg representation. \url{https://arxiv.org/pdf/1504.04157} . Theorem2.1, Corollary2.3, pp.3--4, read directly; only their characteristic-zero and integral flag-model assertions are used here.

Source sign: recovered \nolinkurl{Formalism/LocalToGlobal/Columns/RecLevel.lean}, lines50--52; initial finite type \nolinkurl{GenericUnipotentCharacter.lean}437--452 and \nolinkurl{ColForms.lean}92--105, as audited in LIFT17. No Lean compilation or full literal raw-datum instantiation is claimed.

An exact finite diagnostic is saved in \nolinkurl{definite_oscillator_steinberg_certificate.py} and its JSON output. It checks GL2(F3), b=3, with a nontrivial determinant splitting: the flag vector has norm squared3, the free-orbit Schwartz vector has norm squared16/729, all source generic covariance identities hold, and Fourier inversion and Parseval agree exactly in Q(zeta3). It also checks that for a=6, v11(\textbar GL6(F113)\textbar)=1 while the corresponding valuations at281 and337 are0. The finite diagnostic is a normalization check, not a substitute for the general proof above.

\clearpage

\section{AM-R-015: the selected real discrete-series descent and its exact first occurrence}\label{am-r-015-the-selected-real-discrete-series-descent-and-its-exact-first-occurrence}

Author: /root, GPT-6 Astra Pro, ALP\_MACHINE, 21 November2026UTC.

This finding computes the real correspondence on the actual signed Harish--Chandra lists in the source. The insertion block has length t, its compact sign is positive, and its exact first-occurrence invariant is t. Every smaller target obtained by removing a positive number of hyperbolic planes is therefore zero. The analytic consequence is the fifth source predicate in the genuine oriented-kernel interpretation specified below. This is an application of established real theta theorems, with a new explicit check of the selected parameters and the source's conjugation convention. It supplies no integral row tie or arithmetic terminal value.

\subsection{1. The oriented vertical kernel}\label{the-oriented-vertical-kernel}

The direction of the integral must be distinguished from the direction of the Hom model. For the edge from the smaller space Y\_j to the larger space Y\_(j+1), write its oriented theta kernel as \(\theta_j(g,h;\Phi)\), with g on the smaller group. The actual downward operator is

\[
 \mathcal D_j(\Phi,F)(g)
       =\int_{[U(Y_{j+1})]}\theta_j(g,h;\Phi)F(h)\,dh.
 \tag{R15.1}
\]

It is bilinear and has no conjugation on the kernel. Its Hilbert-space adjoint upward operator uses the conjugated kernel. This is exactly \nolinkurl{IndefiniteReal/OrthogonalComplementCharacterisation.lean}, liftInt/IsIo, lines41--55, and IsIoUp, lines78--88; \nolinkurl{LinksChar.lean}, lines283--287, ties ioY/ioX to that convention. \nolinkurl{ChainParSucc.lean}, line132, is the adjoint upward formula. The abstract downward th has reversed group order: \nolinkurl{BesselHermitianTheta/TiedWeil.lean}, lines131--142, specifies the swapped characters and transpose transport without an additional complex conjugation.

Consequently one must not replace (R15.1) by a newly chosen standard Hom kernel from the larger space to the smaller space. That replacement would alter the representation and its character shifts.

The small-group Hom action is contragredient. Transport it by \(g\mapsto g^{-t}\) to the ordinary tensor dual pair, and transport the form and additive character together. In the standard real tensor convention of Ichino section4.1, the additive character is \(x\mapsto e^{-2\pi i x}\). The compact weight formula is his Lemma7.10, including the cross-character shifts. In Hom coordinates the corresponding joint representation is

\[
       \overline{\sigma_j}\boxtimes\sigma_{j+1}
 \tag{R15.2}
\]

when the tensor correspondence pairs \(\sigma_j\) with \(\sigma_{j+1}\). The transpose-inverse on the small factor changes its discrete-series parameter to the contragredient parameter; this can also be read on its compact Cartan, where every weight is negated, followed by the compact Weyl reordering. This statement applies to the odd and indefinite ranks as well: character changes are ratios with trivial real restriction, while the fixed odd baseline is pinned by the standard real compact-weight formula. It does not extrapolate the even-only determinant-root formula of AM-DO-028 to an unspecified odd splitting.

Explicitly, for signatures \((p,q),(r,s)\), the small Hom factor has compact vacuum weights \[
 \bigl(-(r-s)/2-c_{\rm big}\bigr)^{[p]};
 \bigl((r-s)/2-c_{\rm big}\bigr)^{[q]},
\] whereas the large factor has weights \[
 \bigl((p-q)/2+c_{\rm small}\bigr)^{[r]};
 \bigl(-(p-q)/2+c_{\rm small}\bigr)^{[s]}.
\] These are Lemma7.10's two vacuum lists after negating the small compact weights. Their integrality follows from the rank and character parities. They fix the odd baseline as well as all integer character twists. This identifies the standard genuine splitting chosen here; it does not claim that a partial list of abstract source axioms alone uniquely determines a real splitting.

Thus an input F with \(\overline F\) of type \(\sigma_{j+1}\) can have downward type only \(\overline{\sigma_j}\). The parameter calculation proving the indicated correspondence follows.

\subsection{2. Cross-character centering of the actual lists}\label{cross-character-centering-of-the-actual-lists}

Fix a selected complex embedding. Retain the source hypotheses Interlace, TParity and Centred, and write

\[
 n_j=k+m-1+jt,\qquad q_j=\lfloor(m-j)/2\rfloor,
 \qquad p_j=n_j-q_j,\qquad N_j=k+j(t+1).
 \tag{R15.3}
\]

The source odd character has half-angular weight \(h=\pm1/2\), so the parameter \(\eta=-h\) has exactly the two values permitted by TPar.eta\_half. Let \(\chi_i=\operatorname{car}(\mu,j+1,i)\) for this edge. Then

\[
 c_i=\operatorname{infTy}\chi_i
          =-e_j+1_{i\ \mathrm{even}}\eta.
 \tag{R15.4}
\]

If \(D_i^0\) denotes the unshifted signed list in Dpar i, the source definition gives \(D_i=D_i^0+1_{i\ \mathrm{odd}}\eta\). Hence the two cross-character shifts cancel exactly:

\[
       D_j-c_{j+1}=D_j^0+e_j,
       \qquad D_{j+1}-c_j=D_{j+1}^0+e_j.
 \tag{R15.5}
\]

These are the actual half-angular shifts m0/2 and n0/2 in the real theta convention. They are not arbitrary simultaneous translations. The rank parities and the odd half-integral h give the ordinary Harish--Chandra lattice \(\mathbb Z+(n_i-1)/2\) on each group.

Sort each of the following lists in decreasing order:

\[
 \begin{split}
 A_j={}&\{v_l+e_j:0\le l<k-1\}
       \cup\{e_j-b_i:0\le i\le j\}\\
      &\quad\cup\bigcup_{i<j}\bigl(B_t(-e_i)+e_j\bigr),\\
 B_j={}&\{e_j-b_i:i>j,\ i-j\ \mathrm{even}\},\\
 E_j={}&\{e_j-b_i:i>j,\ i-j\ \mathrm{odd}\},\\
 B_t(0)={}&((t-1)/2,(t-3)/2,\ldots,-(t-1)/2).
 \end{split}
 \tag{R15.6}
\]

The letter E\_j here is a list, not a number field. Interlace orders the centers e\_i and the b\_i. Every member of A\_j is positive: the v terms lie above -b\_0, the b\_i terms have i\textless=j, and each earlier block lies entirely above the current center by Centred. Every member of B\_j and E\_j is negative. More precisely Centred gives

\[
   a\ge(t+3)/2\quad(a\in A_j),\qquad
   b,e\le-(t+3)/2\quad(b\in B_j,e\in E_j).
 \tag{R15.7}
\]

All entries are integers. They are distinct: interlacing separates the v and b entries, while Centred separates each inserted block from every other block and from these entries. The centered signed lists are exactly

\[
 \begin{array}{c|c|c}
 &\text{positive signature slot}&\text{negative signature slot}\\\hline
 D_j-c_{j+1}&A_j\cup B_j&E_j\\
 D_{j+1}-c_j&A_j\cup B_t(0)\cup E_j&B_j.
 \end{array}
 \tag{R15.8}
\]

Indeed the future b-signs flip when j increases by one, the old positive entries retain their signs, and the new block is entirely in the positive signature slot. This also follows directly from \nolinkurl{DataTheta.lean}, Dpar, lines538--543.

The exact list lengths are

\[
 |A_j|=N_j,\quad |B_j|=q_{j+1},\quad |E_j|=q_j,
 \tag{R15.9}
\]

so (R15.8) has precisely the source signatures, not merely the correct total ranks:

\[
 (p_j,q_j)=(N_j+q_{j+1},q_j),\qquad
 (p_{j+1},q_{j+1})=(N_j+q_j+t,q_{j+1}).
 \tag{R15.10}
\]

\subsection{3. Nonzero upward correspondence and the exact downward type}\label{nonzero-upward-correspondence-and-the-exact-downward-type}

Apply the sufficiently separated discrete-series case of real theta lifting: Ichino (2022), Lemma8.2, printed p.35, with its explicit regular-case proof citing Li (1990). The applicable separation is \((t+1)/2\); (R15.7) is stronger. The four lists in that result are \((A_j,B_j,\varnothing,E_j)\), and the two extra signature dimensions are t and0, respectively, by (R15.10).

The resulting cohomological induction has compact Levi block U(t) and one-dimensional remaining Levi factors. In the normalized induction parameter its inserted block is zero repeated t times. The compact Levi rho adds the distinct sequence B\_t(0); all remaining entries remain the entries in (R15.8). Equivalently, refine the compact Levi to a compact Cartan and use induction in stages. The inequalities (R15.7) keep this refinement in the good range and the final parameter regular. Thus the representation is the discrete series with exactly the second signed list of (R15.8), after restoring the shift c\_j. It is \(\sigma_{j+1}\), with parameter D\_(j+1).

Therefore, in the tensor convention with these exact characters,

\[
             \theta_{p_{j+1},q_{j+1}}(\sigma_j)
                         =\sigma_{j+1}\ne0.
 \tag{R15.11}
\]

Howe duality then gives the reverse small lift

\[
             \theta_{p_j,q_j}(\sigma_{j+1})=\sigma_j.
 \tag{R15.12}
\]

Only the regular case of Lemma8.2 is used. The general theorem in Ichino's introduction invokes a global Arthur multiplicity formula; that general conditional input is not being used here. Its regular case is the earlier local Li theorem, explicitly isolated in Lemma8.2. No new proof of Li's theorem or of Howe duality is claimed.

\subsection{4. Exact first occurrence and all lower-target vanishing}\label{exact-first-occurrence-and-all-lower-target-vanishing}

For the reverse correspondence, center the larger parameter by the smaller-space character c\_j. By (R15.8) it contains the contiguous central block B\_t(0), all in the positive signature slot. This block is maximal: its next possible endpoints would be \(\pm(t+1)/2\), which are absent by (R15.7). Its length is therefore exactly t in Atobe's Definition1.6, not merely at least t.

The remaining positive entries are A\_j; the remaining negative entries in the positive slot are E\_j; the negative-slot entries are B\_j, all negative. Consequently Atobe's two base indices are

\[
     k_{\sigma_{j+1}}=t,\qquad
     r_{\sigma_{j+1}}=|A_j|+|B_j|=p_j,\qquad
     s_{\sigma_{j+1}}=|E_j|=q_j.
 \tag{R15.13}
\]

Atobe (2020), Theorem1.7(2), last sentence, determines this balanced Witt tower unconditionally for discrete series. Its permitted shifts are precisely the nonnegative integers. The smaller targets in the source have shift -i. Hence

\[
 \boxed{\quad
 \theta_{p_j-i,q_j-i}(\sigma_{j+1})=0
       \qquad(1\le i\le q_j).
 \quad}
 \tag{R15.14}
\]

The theorem concerns the big-theta nonvanishing criterion as well, so these are zero intertwining spaces, not merely zeros of a chosen Schwartz vector. The discrete-series clause is explicitly declared unconditional on printed p.5 of Atobe's published paper. The later correction to its separate k=-1 tempered case is irrelevant because here k=t\textgreater=1.

For the additive-character convention, Ichino section6.2, printed p.20, explicitly identifies Atobe's convention with his section4.1 choice, using the compact-weight comparison of Lemma7.10. His Theorem6.1(ii) gives the same balanced-tower conclusion here: its exceptional negative-shift possibility requires both next central endpoints, whereas (R15.7) excludes both. Thus a change of the additive-character sign is not being silently absorbed into the signature labels.

\subsection{5. Consequence for the genuine fifth source predicate}\label{consequence-for-the-genuine-fifth-source-predicate}

Use the genuine smooth cusp domain of AM-R-014 and the downward orientation (R15.1). If \(\overline F\) has the prescribed real type D\_(j+1), then F has the contragredient type. Equation(R15.2), Howe duality and (R15.12) imply that every nonzero cuspidal downward image has contragredient real type D\_j. Thus its complex conjugate has type D\_j, which is exactly IsDv at stage j.

One can justify the assertion constituent by constituent in the cuspidal unitary spectrum: every irreducible real constituent of the image is a quotient of the corresponding big theta module, whose unique irreducible quotient is fixed by (R15.12). The actual cusp inputs lie in finite sums of smooth automorphic embeddings, as specified in AM-R-014; no decomposition of arbitrary smooth functions is assumed. For compact-finite test data the assertion holds on the Harish--Chandra modules and then their smooth globalizations.

For a lower target, (R15.14) makes every such local intertwining space zero at each chosen real place. It therefore kills the global theta integral for every Schwartz input. Equivalently, first use the zero Harish--Chandra theta module on dense compact vectors and extend by continuity of the cuspidal theta integral in the usual Schwartz/smooth-cusp topologies. The measures and normalizing nonzero scalars do not affect this vanishing.

These are exactly the two conclusions of \nolinkurl{CayleySignature.RealThetaLiftDiscreteSeries}, lines186--194. Together with AM-R-013/014, they verify five of the six analytic predicates in the explicitly stated genuine model. The remaining analytic predicate is ThetaUnfoldingFreeData, with its precise Fourier--Jacobi period and factorizable witness. A full arithmetic ThetaX/D7 instance, the odd global character pin on a repaired host, the raw-Xi selected action and the terminal mass remain separate.

\subsection{6. Geometric framework and source status}\label{geometric-framework-and-source-status}

The signed weights, central character shifts and compact-Levi Schur classes live in SGA3 I Proposition4.7.3/Corollary4.7.3.1 and the carried SGA6 VI Theorem3.3 representation/flag framework. The oriented kernel remains an actual polarized Hom/theta-group line, in the SGA7 I VII2.7/2.9 setting, pp.169--174. Neither a virtual class nor an isocrystal slope proves the analytic real correspondence; the specific local theorems cited above supply that independent input.

The source lists, parity and centering are DataTheta.lean538--555 and \nolinkurl{CayleySignature.lean104–125.} The rank/signature formulas are Params.lean and the definitions sigYD/rankYD in the recovered source. No Lean execution is claimed.

\subsection{References}\label{references-4}

Atobe, H. (2020). On the non-vanishing of theta liftings of tempered representations of U(p,q). \emph{Advances in Mathematics, 363},106984. \url{https://doi.org/10.1016/j.aim.2020.106984} . Author repository copy: \url{https://eprints.lib.hokudai.ac.jp/repo/huscap/all/84508/Adv.\%20math\%20363_106984.pdf} . Definition1.6 and Theorem1.7, printed pp.4--5; the unconditional discrete-series scope is explicit immediately after that theorem.

Ichino, A. (2022). Theta lifting for tempered representations of real unitary groups. \emph{Advances in Mathematics, 398},108188. \url{https://doi.org/10.1016/j.aim.2022.108188} . Author preprint: \url{https://arxiv.org/pdf/2008.06174} . Section4.1 p.10; Lemma7.10 pp.30--31; regular-case Lemma8.2 p.35 and its proof. The general conditional theorem is distinguished from this local regular case. Theorem6.1(ii) and section6.2, printed pp.19--20, fix the balanced tower and the additive-character convention.

Li, J.-S. (1990). Theta lifting for unitary representations with nonzero cohomology. \emph{Duke Mathematical Journal, 61}(3),913--937. \url{https://doi.org/10.1215/S0012-7094-90-06135-6} . The precise regular case used here is read in Ichino's Lemma8.2; direct retrieval of the original publisher text was blocked by a service security page. No direct whole-paper reading of Li is claimed.

Review status: author whole-proof reading PASS, 2026-11-21T23:40:17.491120+00:00. Independent whole-proof review by /root/definite\_oscillator also PASS, recorded separately in definite\_oscillator\_review\_r015.md. The author reread the complete revised proof, including the explicit compact vacuum shifts, exact source down/adjoint distinction, regular local input, first-occurrence invariant and smooth-cusp scope. No arithmetic constructor or Lean build is certified.

\clearpage

\section{AM-R-016: the final two free-data transfers, and all six analytic predicates at lift count two}\label{am-r-016-the-final-two-free-data-transfers-and-all-six-analytic-predicates-at-lift-count-two}

Author: /root, GPT-6 Astra Pro, ALP\_MACHINE, 21 November2026UTC.

This finding proves the third analytic predicate for the final two edges of the concrete genuine tower of R014. Consequently, when m=2, it supplies the entire \nolinkurl{ThetaUnfoldingFreeData} predicate. Together with R013--015 it gives all six mathematical analytic predicates in that interpretation. The conclusion concerns freely chosen Schwartz data. It does not supply the source's arithmetic host constructor, its prescribed raw Xi maps, an integral row identification, or the cost difference of DP44.

The point is that these two edges involve only ordinary compact restriction and rank-zero Fourier--Jacobi periods. There is no positive-rank unipotent unfolding in this case. We retain the conjugations, the odd splitting character and the Haar measures.

\subsection{1. Spaces, characters and oriented kernels}\label{spaces-characters-and-oriented-kernels}

Let E/F be the CM extension used by the genuine analytic construction, with its fixed trace-zero coefficient xi and standard additive character. Let eta be its chosen unitary odd splitting character, with restriction epsilon\_(E/F) to the base-field ideles. Let mu be any even splitting character, trivial on the base-field ideles. The finite conductor of eta is allowed in this analytic assertion. On the old host its everywhere-unramified pin is obstructed by K028; K029 supplies that pin on its larger field.

Use the standard rational spaces of R014, with k,m even and t odd. Put d=t+1. At the final three indices write

\[
 A=X_{m-2}=I_a,\quad B=X_{m-1}=I_b,\quad
 C=X_m=I_c,\qquad b=a+d,\quad c=b+d.
 \tag{R16.1}
\]

Let L be the last positive coordinate line of C and put H=L\^{}perp. The corresponding vertical spaces and embeddings are exactly

\[
 Y_{m-2}=A\perp(-L),\qquad Y_{m-1}=B,\qquad Y_m=H.
 \tag{R16.2}
\]

Here L and -L have Gram matrices 1 and -1. Thus their dimensions and signatures agree with R014: the first space in(R16.2) has signature(a,1), the second is positive of rank b, and the last is positive of rank c-1. The positive embeddings are the first summand A, the identity of B, and H into C, respectively. All the quotients {[}U(A){]}, {[}U(B){]} and {[}U(H){]} are compact.

Write theta\_(S,T)\^{}(alpha,beta)(g,h;Phi) for the oriented Hom theta kernel, where alpha is the splitting character attached to S and beta the character attached to T. Its Hom symplectic form and action are

\[
 \mathcal B_{S,T}(X,Y)
   =-\operatorname{Tr}_{E/F}\bigl(\xi\operatorname{tr}(Y^*X)\bigr),
 \qquad X\longmapsto hXg^{-1}.
 \tag{R16.3}
\]

The adjoint is taken for the indicated Hermitian forms. We use the actual compatible global Weil splittings and self-dual additive measures. Orthogonal sums have product measures and the tensor product theta distribution. Every automorphic quotient carries the same fixed Haar measure at each occurrence below; its total volume is not set to one unless already so normalized.

The downward operator uses this smaller-first kernel without conjugation:

\[
 D_{S,T}(\Phi,F)(g)
       =\int_{[U(T)]}\theta_{S,T}(g,h;\Phi)F(h)\,dh.
 \tag{R16.4}
\]

This is R015's convention, and agrees with the actual downward operator after the source's rank-dependent transpose transport. We do not replace it by a Hermitian adjoint.

For edge j=m-2, which is even, \nolinkurl{car} prescribes characters (mu eta\^{}(-1),mu) on (Y\_(m-2),B). The horizontal pair(A,B) has characters(mu,mu). For edge j=m-1, which is odd, it prescribes(mu,mu eta\^{}(-1)) on(B,H), while(B,C) has(mu,mu). Here mu is the character for the particular edge; the two edge characters need not coincide.

\subsection{2. The one common rank-zero Fourier--Jacobi kernel}\label{the-one-common-rank-zero-fourierjacobi-kernel}

On the pair(B,-L) choose a Weil representation with character eta\^{}(-1) attached to -L. Its other splitting character may be any character trivial on the base field. Define

\[
 \vartheta_-(g;s)=\theta_{B,-L}(g,1;s),\qquad
 \mathcal J_B(\phi,F;s)
       =\int_{[U(B)]}\phi(g)F(g)
                         \overline{\vartheta_-(g;s)}\,dg.
 \tag{R16.5}
\]

The restriction g-\textgreater theta\_(B,-L)(g,1;s) is independent of that other splitting character after the standard character-change identification: changing the character attached to B multiplies the action of U(-L), and that argument is the identity. Therefore the same space of kernels(R16.5) can be used on both adjacent edges even when their mu characters differ.

This is precisely the rank-zero Fourier--Jacobi period: its unipotent group is trivial, its rational quotient has volume1, and its auxiliary theta factor is conjugated. The source formula is \nolinkurl{LastSeesawPeriodNonvanishing.lean}, theorem \nolinkurl{per_zero}, lines58--94, together with \nolinkurl{fjSrcFun} in \nolinkurl{WittSmallGroupEmbeddingMemPointwiseStabilizerIso.lean}, lines228--233. In the source theta-function model, evaluation of the transformed theta function at zero is the displayed theta kernel; it is not evaluation of a bare Schrödinger function.

We require two elementary compatible realizations of this same auxiliary kernel. First, conjugating the Hom Weil representation for(B,L) changes L to -L and inverts both splitting characters. For its character pair(mu,eta), it gives a conjugate-linear Schwartz isomorphism C\_L with

\[
 \theta_{B,L}^{\mu,\eta}(g,1;q)
          =\overline{\vartheta_-(g;C_Lq)}.
 \tag{R16.6}
\]

Second, Hermitian adjoint is anti-symplectic between Hom(-L,B) and Hom(B,-L). Indeed tr((Y\^{}\emph{)\^{}}X\^{}\emph{)=overline(tr(Y\^{}}X)), and bar(xi)=-xi; substitution in(R16.3) changes the sign of the form. Adjoint also interchanges the two group actions. Composing its Schwartz transport with complex conjugation therefore gives a conjugate-linear isomorphism C\_ad satisfying

\[
 \theta_{-L,B}^{\eta^{-1},\mu}(1,g;q)
             =\overline{\vartheta_-(g;C_{\rm ad}q)}.
 \tag{R16.7}
\]

Under this conjugate-adjoint exchange the splitting characters are exchanged, giving(mu,eta\^{}(-1)); they are not both inverted again. The interchange of source and target determinant actions cancels those inversions. At a split place the source scalar is chi\_T(det g)\^{}(-1) and the target scalar is chi\_S(det h), which verifies this cancellation directly. At infinity the auxiliary U(B) vacuum weight is 1/2+h when eta has half-angular weight h; the two unbarred line kernels in(R16.6)--(R16.7) both have weight -1/2-h. This also checks the odd baseline.

For completeness, the normalization of(R16.6)--(R16.7) has no free scalar. Negating one Hermitian form negates the Heisenberg form; complex conjugation changes psi to psi\^{}(-1), so the two changes cancel. Adjoint carries rational position points to rational position points when the polarizations are transported with it. Hence conjugation commutes with the rational theta sum. The self-dual measures are transported at the same time. The linear, quadratic and Fourier generators of the Weil representation respect these maps; the Fourier identity is the usual conjugation of its integral kernel. Thus the equality includes the theta distribution itself, not just an abstract representation isomorphism.

These are also the exact general theta identities recorded in \nolinkurl{IsWeilDatumConjNegHermitianSpaces.lean}, \nolinkurl{weil_conjTarget}, and \nolinkurl{ConjWeilDatumViaAdjointExchangedPairStatement.lean}. They are used as accepted source identities within their scope. The Hom computation above explains their identical normalization for the genuine E/F model without asserting that the arithmetic D7 record has been constructed.

\subsection{3. Terminal transfer: a positive line and a rank-zero period}\label{terminal-transfer-a-positive-line-and-a-rank-zero-period}

Fix a smooth cusp form phi on U(B) and a smooth cusp form F on U(H). The source's terminal bilinear period is

\[
 \mathcal P_H(f,F)=\int_{[U(H)]}f(\iota h)F(h)\,dh
                 =B_{\rm top}(\overline f,F).
 \tag{R16.8}
\]

Orthogonal factorization C=H perp L, with the character identity (mu eta\^{}(-1))*eta=mu, gives the exact theta identity

\[
 \theta_{B,C}^{\mu,\mu}(g,(h,1);\Phi\boxtimes q)
  =\theta_{B,H}^{\mu,\mu\eta^{-1}}(g,h;\Phi)
                    \overline{\vartheta_-(g;C_Lq)}.
 \tag{R16.9}
\]

To check factorization, the Hom space is the symplectic direct sum Hom(B,H) plus Hom(B,L). Its chosen position lattice over F is the corresponding direct sum; the rational theta sum of a product is the product of the two absolutely convergent sums. The tensor Weil action restricts to the displayed product. The target character multiplication above is exactly its splitting requirement. This is the standard orthogonal-sum theta tensor rule, with no index or volume factor.

Fubini on the two compact automorphic quotients now gives

\[
 \boxed{
 \mathcal P_H\bigl(\Theta_{B,C}(\phi;\Phi\boxtimes q),F\bigr)
       =\mathcal J_B\bigl(\phi,D_{B,H}(\Phi,F);C_Lq\bigr).}
 \tag{R16.10}
\]

This equality retains the two fixed Haar measures on both sides. The vertical kernel is unbarred; only the auxiliary rank-zero kernel is conjugated.

Suppose the left terminal period is nonzero for an arbitrary Schwartz input Phi'. Then it is nonzero for some input Phi boxtimes q with Phi factorizable and real compact-finite, and q likewise. Here is the needed free-data argument. The left side is a continuous linear functional of Phi' on the Schwartz space. At finite places, compactly supported locally constant functions on a direct sum are finite sums of products. At real places, finite sums of product Hermite functions are dense in the Schwartz topology. For the positive spaces B,H,L, these polynomial-Gaussian functions are finite under the relevant compact unitary groups. Finite sums of global restricted pure tensors of such functions are therefore dense. If the functional vanished on every pure product it would vanish on this dense span and hence on Phi', a contradiction. Choose one product on which it is nonzero. Equation(R16.10) then gives the required nonzero lower Fourier--Jacobi period.

The continuity assertion can be seen directly without an interchange of a merely formal series. At fixed finite support and period, the adelic quotients have finitely many charts with compact real parameters. The rational theta series and all fixed-order derivatives are uniformly bounded on those charts by a Schwartz seminorm and a convergent lattice majorant. Multiplying by the fixed smooth forms and integrating over the compact quotients is continuous. The finite-place Schwartz space is the union of these fixed support/period spaces.

\subsection{4. Penultimate transfer: ordinary Bessel restriction}\label{penultimate-transfer-ordinary-bessel-restriction}

Write V=A perp(-L), as in(R16.2), and fix a horizontal factorizable compact-finite input Phi' for(A,B). The source decomposition and the character identity mu*eta\^{}(-1)=mu eta\^{}(-1) give

\[
 \theta_{V,B}^{\mu\eta^{-1},\mu}((g,1),h;
                                      \Phi'\boxtimes q)
   =\theta_{A,B}^{\mu,\mu}(g,h;\Phi')
                         \overline{\vartheta_-(h;C_{\rm ad}q)}.
 \tag{R16.11}
\]

Let phi be a smooth cusp form on U(A), F one on U(B), and define the ordinary codimension-one Bessel period by

\[
 \mathcal B_A(\phi,G)=\int_{[U(A)]}\phi(g)G((g,1))\,dg.
 \tag{R16.12}
\]

The Bessel unipotent rank is zero. Thus there is no nontrivial unipotent character or hidden unipotent volume in(R16.12). Fubini gives the exact identity

\[
 \boxed{
 \mathcal B_A\bigl(\phi,D_{V,B}(\Phi'\boxtimes q,F)\bigr)
   =\mathcal J_B\bigl(\Theta_{A,B}(\phi;\Phi'),F;
                                                    C_{\rm ad}q\bigr).}
 \tag{R16.13}
\]

If the right side is nonzero for some auxiliary datum s, it is nonzero for some factorizable compact-finite s. The same continuity argument applies, now to the auxiliary Schwartz space alone. Polynomial-Gaussian functions for the pair(B,-L) are finite under U(B) times U(1), and their span is dense. Choose one global pure tensor with nonzero period and take q=C\_ad\^{}(-1)s. The induced rational adjoint transport preserves pure tensors and compact finiteness. Then Phi' boxtimes q is factorizable and finite under K\_V times K\_B. At infinity K\_V=U(A) times U(1), and the diagonal U(B) action on the tensor product of two finite modules is finite. Equation(R16.13) is the required lower nonzero Bessel period.

No cuspidality claim for the noncompact group U(V) has been inserted into Fubini: only its restriction to U(A) is integrated. For the selected input type, R015 gives vanishing at all smaller targets and R014 then gives cuspidality of this actual downward lift. Thus the constructed form belongs to the genuine cusp domain when that is required by the source. Independently, the literal free-data predicate quantifies over every cusp form G equal to this lift; either reading uses the same identity(R16.13).

\subsection{5. Quantifiers, scope and exact implications}\label{quantifiers-scope-and-exact-implications}

For j=m-1, a nonzero terminal \nolinkurl{PerNZ} is(R16.8) and the existential free datum is supplied by(R16.10) and section3. For j=m-2, upper \nolinkurl{PerNZ} is exactly(R16.5), and(R16.13) supplies the lower ordinary Bessel period. The \nolinkurl{pos} choices are the actual standard positive embeddings in(R16.2). The functions in the source \nolinkurl{bB} have the conjugation \nolinkurl{Bj(conjAF(phi),F)}; with the ordinary Hermitian first slot this is precisely the bilinear phi*F convention used here. More explicitly, in this genuine interpretation the rank-zero \nolinkurl{Pos} datum at the identity embedding B into B is the unique empty isotropic flag on its zero orthogonal complement. Its \nolinkurl{FJd} is the full Schwartz space for the fixed auxiliary pair (B,-L), with the cross-character eta\^{}(-1). Thus the existential upper \nolinkurl{PerNZ} ranges over exactly the kernels in(R16.5), rather than over an unexamined additional class of positions. At the lower codimension-one Bessel step the isotropic flag is again empty; the nonisotropic coordinate in the complementary line does not occur in the rank-zero unipotent character, so every such coordinate choice gives the same ordinary restriction. At j=m-1 the source uses the auxiliary conjugated theta factor, and at j=m-2 the period is independent of any auxiliary FJ datum because its Bessel unipotent rank is zero.

The conclusion is uniform in the two edge characters mu, the chosen forms and the nonzero upper period. It does not assume a special nonzero matrix coefficient or a selected automorphic multiplicity copy. It therefore proves the literal mathematical quantifiers of the third predicate at these final two indices for the genuine model specified here.

When m=2 the set of indices j\textless m is exactly\{0,1\}. Consequently R013--016 cover all six analytic predicates for that model, provided the stated genuine characters and theta spaces are supplied. This closes the analytic free-data gate for lift count2. It does not identify these freely chosen kernels with the raw arithmetic PhiY/PhiG/Xi family. In particular, the density argument provides existence and no bound on the selected finite level, coefficient denominators or norms.

For m\textgreater2 the earlier free-data edges still involve genuine positive-rank Bessel/Fourier--Jacobi unfolding. The accepted source UpBF/UpFB statements also require their local test-space package. Those additional edges are not proved in this note.

\subsection{6. Exact locators and check status}\label{exact-locators-and-check-status}

\begin{itemize}
\tightlist
\item
  Alpöge, L. (2026). \emph{Remark}, step8, with the stipulated pinned archive at commit7cbf7a0c32fca1badf05931879a37c3f8a2b03b4. \nolinkurl{Formalism/SharpRestriction/CuspSpanNormBounds/CayleySignature.lean}, \nolinkurl{car}114--115, \nolinkurl{bB}131--133, \nolinkurl{PerNZ}150--152, \nolinkurl{ThetaUnfoldingFreeData}168--175.
\item
  The same archive, \nolinkurl{Development/AbelianVarietyTheory/Proofs/TwistedOperators/LastSeesawPeriodNonvanishing.lean}, \nolinkurl{per_zero}58--94; and \nolinkurl{Definitions/TwistedInnerProducts/WittSmallGroupEmbeddingMemPointwiseStabilizerIso.lean}, \nolinkurl{fjSrcFun}228--233.
\item
  The same archive, \nolinkurl{Proofs/DecoratedLatticeProofs/ThetaTensorRuleTargetOrthogonalSum.lean}, \nolinkurl{tensorTarget}75ff.; \nolinkurl{ThetaTensorRuleOrthogonalSumSourceStatement.lean}; \nolinkurl{IsWeilDatumConjNegHermitianSpaces.lean}, \nolinkurl{weil_conjTarget}71--136; \nolinkurl{ConjWeilDatumViaAdjointExchangedPairStatement.lean}.
\item
  ALP\_MACHINE R013 section1, R014 sections2--4, R015 sections1--2: the genuine spaces, actual bilinear conventions, character parity and the already checked remaining analytic predicates.
\end{itemize}

Whole-proof review status: author complete rereading PASS, 2026-11-21T23:56:58.217024+00:00. Independent whole-proof review by /root/definite\_oscillator PASS on22November2026UTC, including the saved rank-zero Pos/FJd clarification; see agents/definite\_oscillator\_review\_r016.md and CHECK\_EVENTS.tsv. No Lean build is requested or performed. The new mathematical work is the exact final-two-edge specialization, the compatible single auxiliary kernel, the compact Fubini identities and the factorizable compact-finite witness argument.

\clearpage

\section{ALP-LIFT-19 --- ramified theta constituents with the literal nonterminal column types}\label{alp-lift-19-ramified-theta-constituents-with-the-literal-nonterminal-column-types}

Author: /root/lift\_accounting, GPT-6 Astra Pro, ALP\_MACHINE. Written 21 November2026UTC.

This theorem repairs the local representation obstruction of ALP-LIFT-17. At a degree-one split G-place, start with an irreducible representation containing the finite Steinberg type, and use the distinct tame stage characters required by IsMu0. Every later local theta representation then contains the literal adelic column type, with the correct inverse character. An explicit finite induced vector is primitive in its natural coefficient lattice and has a computable positive Hermitian norm.

The theorem is local and independent of the degree of the global CM field. It therefore applies to the old degree12 host or to an enlarged host whenever the same degree-one G-place and stage-character hypotheses hold. The global character construction of LIFT18 is used only for its stated old host. No class-number-one assertion, G-prime example, or mass formula is silently transported to the degree84 repaired host.

The finite Steinberg representation, Casselman's Iwahori theorem, Zelevinsky's irreducibility criterion, and Mínguez's local theta theorem are literature inputs. The new work here is their exact assembly with the source columns, the explicit Mackey component and its primitive vector/norm, and the consequent removal of this specific local type obstruction. The literal raw PhiG/Xi kernel and the integral lattice of the iterated global image remain separate.

\subsection{1. Local data and the source sign}\label{local-data-and-the-source-sign}

Let F be a nonarchimedean local field of characteristic zero with residue field \(\kappa=\mathbf F_q\). For the literal degree-one G places, q is the rational G prime. Put \[
 K_n=\operatorname{GL}_n(\mathcal O_F),\qquad
 K_n^1=1+\varpi M_n(\mathcal O_F).
 \tag{T1}
\] Fix integers \(k\ge1,t\ge1,m\ge2\), set \(d=t+1\), and write \[
 n_j=k+jd=(k+j)+jt,\qquad r_j=k+j.
 \tag{T2}
\] Let \(\eta_0,\ldots,\eta_{m-1}:F^\times\to\mathbf C^\times\) be unitary tame characters whose residue characters \(\bar\eta_i:\kappa^\times\to\mathbf C^\times\) are pairwise distinct and nontrivial.

At a selected source G-prime, IsMu0 and chi\_om say exactly \[
 \eta_i|_{\mathcal O_F^\times}
       =\chi_i^{-1},\qquad
 \chi_i=\omega^{\,2(i+1)},
 \tag{T3}
\] where \(\omega\) is faithful. The source inequality \(q-1>2fk>2m\) ensures the required nontriviality and distinctness. This is the character condition already audited in LIFT17; LIFT18 constructs actual global even characters with it on its specified weight locus over the degree12 host.

The adelic column convention has an inverse. In the pinned source, Columns/RecLevel.lean50--52 defines \[
 \operatorname{levChr}(\lambda)=\lambda\circ\operatorname{levRed},
 \qquad
 \operatorname{levTyp}(\lambda)=\operatorname{levChr}(\lambda)^{-1}.
 \tag{T4}
\] Consequently the desired adelic right-translation character at stage \(j<m\) is \(\lambda_{X,j}^{-1}\), whose determinant-block factors are the \(\bar\eta_i\) in(T3).

Choose a nontrivial additive character \(\psi:\kappa\to\mathbf C^\times\). For \(0\le j<m\), define \[
 \begin{split}
 H_j&=\bigl(U_{r_j}(\kappa)\times
                 \operatorname{GL}_t(\kappa)^j\bigr)N_{P_j},\\
 P_j&=P_{(r_j,t,\ldots,t)}\subseteq\operatorname{GL}_{n_j}(\kappa),\\
 \lambda_j(h_0,h_1,\ldots,h_j;n)
   &=\psi_{S_j}(h_0)^{-1}
                    \prod_{i=0}^{j-1}\bar\eta_i(\det h_{i+1}),\\
 S_j&=\{0\le a<k+j-1:\ a+1\ne k\text{ or }j+1=m\}.
 \end{split}
 \tag{T5}
\] Here \(U_r\) is upper unitriangular, and \(\psi_S\) applies \(\psi\) to the sum of the indicated simple-root entries. This is exactly HX, SX and the inverse of lamX in \nolinkurl{GenericUnipotentCharacter.lean437–452.} The character is trivial on \(N_{P_j}\). For \(j=0\), it is the full generic character of \(U_k\). For \(1\le j<m-1\), it omits only the root joining the first k coordinates to the last j coordinates in \(U_{k+j}\). For \(j=m-1\), it is fully generic in \(U_{k+j}\). The terminal HXm case \(j=m\) is not included.

\subsection{2. The exact local representation theorem}\label{the-exact-local-representation-theorem}

Let \(\pi_0\) be an irreducible admissible representation of \(\operatorname{GL}_k(F)\) such that \[
 \operatorname{St}_k(\kappa)\hookrightarrow\pi_0^{K_k^1}
 \quad\hbox{as }\operatorname{GL}_k(\kappa)\hbox{-modules}.
 \tag{T6}
\] For \(j\ge0\), form normalized parabolic induction \[
 \Pi_j=
 \pi_0\times(\eta_0\circ\det_{\operatorname{GL}_d})
       \times\cdots\times
                (\eta_{j-1}\circ\det_{\operatorname{GL}_d}).
 \tag{T7}
\]

\textbf{Theorem.} For every \(0\le j\le m\), the representation \(\Pi_j\) is irreducible. For each \(0\le j<m\), \(\Pi_{j+1}\) is the actual local theta descendant of \(\Pi_j\) under the literal bilinear Hom convention with the same splitting character \(\eta_j\) on both members of edge \(n_j\to n_{j+1}\). For \(0\le j<m\), its finite column space is nonzero: \[
                   \Pi_j^{\widetilde H_j,\lambda_j}\ne0,
 \tag{T8}
\] where \(\widetilde H_j\) is the full inverse image of \(H_j\) in \(K_{n_j}\). Thus the desired finite type occurs in the entire actual ramified theta representation. This is the opposite local conclusion from LIFT17's unramified chain.

\subsubsection{2.1 Cuspidal support and irreducibility}\label{cuspidal-support-and-irreducibility}

The finite Steinberg representation has a one-dimensional Borel-fixed subspace. For example, in its integral basis \(\{u e:u\in U_k\}\), the vector \(\sum_u u e\) is fixed by \(U_k\), and diagonal matrices permute its summands because they fix e and normalize \(U_k\). Thus(T6) gives a nonzero Iwahori-fixed vector in \(\pi_0\).

Casselman(1980), Proposition2.6, printed p.396, states that every irreducible admissible representation with a nonzero Iwahori-fixed vector embeds in an unramified principal series. Hence the cuspidal support of \(\pi_0\) consists entirely of unramified characters of \(F^\times\). This statement applies to \(\operatorname{GL}_k\), including its center; Casselman's paper is formulated for connected reductive groups.

The cuspidal support of the one-dimensional representation \(\eta_i\circ\det_{\operatorname{GL}_d}\) is \[
 \bigl\{\eta_i|\cdot|^{(1-d)/2},
             \eta_i|\cdot|^{(3-d)/2},\ldots,
                         \eta_i|\cdot|^{(d-1)/2}\bigr\}.
 \tag{T9}
\] Its residue character is uniformly \(\bar\eta_i\). The distinctness in(T3) therefore puts different factors of(T7) on disjoint cuspidal lines; none meets the unramified support of \(\pi_0\). No two segments from different factors can be linked, since an unramified absolute-value twist does not change the residue character.

Zelevinsky(1980), Proposition8.5, printed p.196, now says that the product(T7) is irreducible. Its Proposition8.6, printed p.197, records the equivalent factorization by distinct cuspidal lines. This argument does not require \(\pi_0\) to be tempered or generic. If \(\pi_0\) is unitary, so is(T7), because every inducing factor is unitary and normalized parabolic induction preserves unitarity.

\subsubsection{2.2 Identification with the actual theta descendant}\label{identification-with-the-actual-theta-descendant}

At a relative split place, write the common unitary character as \((\eta_j,\eta_j^{-1})\) on \(F^\times\times F^\times\). The actual Hom oscillator carries the scalar \[
             \eta_j(\det g)^{-1}\eta_j(\det h).
 \tag{T10}
\] The source lift is bilinear, \(\int\theta(g,h)f(g)\,dg\), so the untwisted type-II source representation is \(\Pi_j^\vee\otimes\eta_j\). Mínguez(2008), Theorem1(2), printed p.718, gives nonzero local theta for every irreducible input in the smaller rank. Its target Langlands quotient has the contragredient old parameters and the parameters of the rank-d trivial representation. Restoring(T10) therefore leaves the old \(\Pi_j\) block unchanged and inserts \(\eta_j\circ\det_{\operatorname{GL}_d}\).

This identifies the target with(T7) for \(j+1\): the induced product is irreducible by section2.1, and its Langlands quotient is precisely the one in Mínguez's theorem. Equivalently its standard parameter is the direct sum of the old parameter and(T9). This use of the Langlands quotient is not an assertion that an arbitrary reducible parabolic induction equals small theta.

The same convention was checked against the actual Hom action in LIFT14(E10--E11) and AM-DO-028. The finite right-translation character in(T4) is retained; no Hermitian-input convention is substituted.

It remains to prove(T8). We do so constructively.

\subsection{3. A finite Mackey component adapted to the source columns}\label{a-finite-mackey-component-adapted-to-the-source-columns}

For this section write \(G=\operatorname{GL}_{n_j}(\kappa)\) and \[
 Q=P_{(k,d,\ldots,d)},\qquad
 \mathcal M_j=
 \operatorname{Ind}_{Q}^{G}
       \left(\operatorname{St}_k\boxtimes
                  \mathop{\boxtimes}\displaylimits_{i=0}^{j-1}\bar\eta_i\circ\det_d\right).
 \tag{T11}
\] Finite induction is unnormalized; its function model consists of functions \(f:G\to V_{\operatorname{St}}\) with \(f(qg)=\tau(q)f(g)\), acted on by right translation.

Give the Q blocks the coordinate order \[
        A,\ L_0,T_0,\ L_1,T_1,\ \ldots,L_{j-1},T_{j-1},
 \qquad \dim A=k,\quad\dim L_i=1,\quad\dim T_i=t.
\] Let w be the permutation sending the P-coordinate order \[
            A,L_0,\ldots,L_{j-1},\ T_0,\ldots,T_{j-1}
\] to this Q-coordinate order. Put \(H=w^{-1}Qw\cap P_j\). The image of H in the Levi \(L_{P_j}=\operatorname{GL}_{r_j}\times\operatorname{GL}_t^j\) is exactly \[
              R=P_{(k,1,\ldots,1)}
                       \times\operatorname{GL}_t^j.
 \tag{T12}
\] This can be checked entry by entry. Within the first P-block, the Q-upper-triangular conditions impose the parabolic blocks \(k,1,\ldots,1\); every matrix with those conditions is allowed. The T-block diagonal matrices are unrestricted. Thus every element of R has a lift in H.

The representation \(\tau\) in(T11) is trivial on \(w(H\cap N_{P_j})w^{-1}\). Indeed off-diagonal Q blocks act trivially by inflation. Within \(L_i\oplus T_i\), the remaining intersection has matrices \(\left(\begin{smallmatrix}1&*\\0&I_t\end{smallmatrix}\right)\), on which a determinant character is trivial. On the image R, the action is \[
 \operatorname{St}_k\boxtimes
     (\bar\eta_0,\ldots,\bar\eta_{j-1})
       \boxtimes
          \mathop{\boxtimes}\displaylimits_{i=0}^{j-1}\bar\eta_i\circ\det_t.
 \tag{T13}
\]

It follows that \(\mathcal M_j^{N_{P_j}}\) contains the component \[
 \mathcal V_j\boxtimes
       \mathop{\boxtimes}\displaylimits_{i=0}^{j-1}\bar\eta_i\circ\det_t,\qquad
 \mathcal V_j=
 \operatorname{Ind}_{P_{(k,1,\ldots,1)}}^{\operatorname{GL}_{r_j}}
        (\operatorname{St}_k\boxtimes\bar\eta_0\boxtimes
                                  \cdots\boxtimes\bar\eta_{j-1}).
 \tag{T14}
\] Here is an explicit proof, so a quotient in a Mackey filtration is not mistaken for a submodule. Take functions in(T11) supported on the finite double coset QwP\_j and invariant under right \(N_{P_j}\). Their restriction to wP\_j is the induced-function model from H to P\_j. The trivial action of \(H\cap N_{P_j}\) just proved makes such a function descend to \(L_{P_j}\). Its left transformation law is exactly(T13) on R. Conversely every induced function with that law inflates to P\_j and extends uniquely to QwP\_j with the Q-equivariance of(T11), and by zero elsewhere. These constructions are inverse and \(L_{P_j}\)-equivariant. They give the actual subspace(T14).

This is the Mackey component whose intersection-size matrix has entries k in the first-first position, one from each new Q block in the first P block, t in its corresponding determinant block, and zero elsewhere. No new free trailing block is used.

\subsection{4. Generic and degenerate generic vectors in the first block}\label{generic-and-degenerate-generic-vectors-in-the-first-block}

We use the following elementary finite induction observation. Let \(R'=P_{(a_1,\ldots,a_s)}\subseteq\operatorname{GL}_r(\kappa)\), and suppose each representation \(\sigma_i\) of \(\operatorname{GL}_{a_i}(\kappa)\) has a vector for the chosen generic upper-unitriangular character. Then \(\operatorname{Ind}_{R'}^{\operatorname{GL}_r}
(\boxtimes_i\sigma_i)\) also has such a vector.

To prove it, take the permutation which reverses the order of the blocks while preserving the internal order in each block. Its conjugate parabolic meets \(U_r\) in the product of the within-block upper-unitriangular groups, with the blocks in reversed order. Choose the tensor of the given generic vectors. Define an induced function on the corresponding \(R'wU_r\) double coset by right multiplication with the desired generic character, and zero elsewhere. The transformation law on the stabilizer agrees with the tensor vector by construction, so the function is well-defined and nonzero. It is a generic eigenvector under right \(U_r\). This construction uses no averaging denominator. It works with \(\psi^{-1}\) just as with \(\psi\).

The finite Steinberg representation is regular on \(U_k\) (Steinberg1957, Theorem2(i), printed p.349), so it contains the required \(\psi^{-1}\)-eigenvector. Applying the preceding observation to(T14) proves the fully generic case \(j=m-1\).

For \(1\le j<m-1\), the source character omits the root at k. Transitivity of finite induction rewrites \(\mathcal V_j\) as \[
 \operatorname{Ind}_{P_{(k,j)}}^{\operatorname{GL}_{k+j}}
  \left(\operatorname{St}_k\boxtimes
      \operatorname{Ind}_{B_j}^{\operatorname{GL}_j}
                (\bar\eta_0\boxtimes\cdots\boxtimes\bar\eta_{j-1})\right).
 \tag{T15}
\] The subspace supported on \(P_{(k,j)}\) and fixed by its radical contains exactly the displayed Levi representation. Its first factor has the generic \(\psi^{-1}\) vector; the second has the generic \(\psi^{-1}\) vector by the preceding open-orbit construction with all block sizes one. Their tensor is therefore an eigenvector for the character of \(U_{k+j}\) which is generic inside the k and j blocks and trivial on the joining radical. That is precisely \(\psi_{S_j}^{-1}\). Together with(T14), this proves \[
                         \mathcal M_j^{H_j,\lambda_j}\ne0.
 \tag{T16}
\] For \(j=0\), the same assertion is simply the generic Steinberg vector. The proof treats both source branches exactly.

\subsection{5. Passage to the actual p-adic representation}\label{passage-to-the-actual-p-adic-representation}

By Iwasawa decomposition, the compact function model of(T7) restricts to induction from \(Q(F)\cap K_{n_j}\). Taking \(K_{n_j}^1\)-invariants gives the exact finite induced model \[
 \Pi_j^{K_{n_j}^1}\simeq
 \operatorname{Ind}_{Q(\kappa)}^{\operatorname{GL}_{n_j}(\kappa)}
       \left(\pi_0^{K_k^1}\boxtimes
                  \mathop{\boxtimes}\displaylimits_{i=0}^{j-1}\bar\eta_i\circ\det_d\right).
 \tag{T17}
\] For clarity, right \(K_{n_j}^1\)-invariance means that a function on K factors through its finite quotient. Its value is fixed by \((Q\cap K_{n_j}^1)\), whose Levi projection is \(K_k^1\times(K_d^1)^j\). These are exactly the coefficient invariants in(T17). Conversely such a finite induced function inflates to K and extends by Iwasawa; modulus characters are one on the compact intersection. This proves the isomorphism rather than merely an inclusion.

By(T6), the finite module(T11) embeds in(T17). Equation(T16) proves(T8), completing the theorem.

\subsection{6. A primitive finite vector and its exact norm}\label{a-primitive-finite-vector-and-its-exact-norm}

There is a useful quantitative statement before any global automorphic integration. Let R be a ring of integers enlarged to contain the values of \(\psi\) and the \(\bar\eta_i\). Use the Steinberg lattice \[
        \operatorname{St}_{k,R}
          =\bigoplus_{u\in U_k}R\,u e
          \ \subset R[\operatorname{GL}_k(\kappa)/B_k].
 \tag{T18}
\] The basis and its integral realization are Steinberg's Theorem1; the same statement over arbitrary coefficient rings is recorded in Geck(2015), Theorem2.1. Give the flag permutation module its Hermitian metric with flags orthonormal.

The generic vector \(v_\psi=\sum_{u\in U_k}\psi_U(u)\,u e\), with the sign chosen to give right type \(\psi_U^{-1}\), is primitive in(T18). Its coefficients in that basis are roots of unity. Moreover \[
                     \|v_\psi\|^2=|U_k|=q^{k(k-1)/2}.
 \tag{T19}
\] To see the norm directly, expand e as the alternating sum of Weyl flags. For every Weyl element except the longest, its \(U_k\)-stabilizer contains a simple-root subgroup on which \(\psi_U\) is nontrivial, so the sum cancels. Only the longest Bruhat cell remains; its \(|U_k|\) flags are distinct and have coefficients of absolute value one. This proves(T19). It does not assert that the whole basis \(\{u e\}\) is orthonormal.

Give(T11) its normalized compact-induction metric \[
                    \|f\|^2=\frac1{|G|}\sum_{g\in G}\|f(g)\|^2.
\] The constructions in sections3--4 can be made over R with no division. At one selected coset their value is \(v_\psi\); at every nonzero coset it has the same norm, because the Steinberg action is unitary and all character values have modulus1. Hence they produce a primitive vector \(v_j\) in the natural R-valued induced lattice with exact type \(\lambda_j\).

Its support size in \(Q\backslash G\) is \[
 q^{\,d j(j-1)/2+\epsilon_j k j},\qquad
 \epsilon_j=
 \begin{cases}
  1,&j=m-1,\\
  0,&0\le j<m-1.
 \end{cases}
 \tag{T20}
\] For \(j=0\) this means one coefficient. For \(j\ge1\), the fiber of \(H\backslash P_j\to R\backslash L_{P_j}\) has \[
        [N_{P_j}:H\cap N_{P_j}]=q^{t j(j-1)/2}.
\] Indeed precisely the entries from \(L_\ell\) to \(T_i\) with \(\ell>i\) are allowed in the P radical and forbidden by Q. There are \(t\,j(j-1)/2\) such entries; all others are unrestricted.

In the fully generic first-block construction the open U orbit has size \[
 |U_{k+j}|/|U_k|=q^{kj+j(j-1)/2}.
\] In the omitted-root construction(T15), the support lies in \(P_{(k,j)}\) and its flag orbit is the longest cell for \(\operatorname{GL}_j\), of size \(q^{j(j-1)/2}\). Combining these counts proves(T20).

Write the Gaussian multinomial index as \[
 \mathcal J_j(q)=[G:Q]
 =\frac{\prod_{s=1}^{k+jd}(q^s-1)}
        {\left(\prod_{s=1}^k(q^s-1)\right)
         \left(\prod_{s=1}^d(q^s-1)\right)^j}.
 \tag{T21}
\] Every left Q coset has \(|Q|\) elements, so(T19)--(T21) give \[
 \boxed{\|v_j\|^2=
       \frac{q^{\,k(k-1)/2+d j(j-1)/2+\epsilon_j k j}}
                       {\mathcal J_j(q)}.}
 \tag{T22}
\] This is a norm of an explicit primitive finite vector, before mapping into a chosen automorphic image lattice. The embedding(T6) may be scaled to be an isometry on its finite Steinberg copy, in which case the same norm holds in the compact model(T17). Such an isometric scaling does not by itself identify the arithmetic lattice inherited from a global coefficient row.

For \(11\nmid q\), let \(h=\operatorname{ord}_{11}(q)\) and \(c=v_{11}(q^h-1)\). Put \[
             D_q(n)=c\lfloor n/h\rfloor
                         +v_{11}(\lfloor n/h\rfloor!).
\] The elementary lifting-the-exponent formula gives \[
 v_{11}\mathcal J_j(q)
       =D_q(k+jd)-D_q(k)-jD_q(d),\qquad
 v_{11}(\|v_j\|^2)=-v_{11}\mathcal J_j(q).
 \tag{T23}
\] Indeed only s divisible by h contribute to \(\prod_{s=1}^n(q^s-1)\), and their valuations are \(c+v_{11}(s/h)\). This displays an actual finite normalization factor. It is not an assertion that the global theta row has that content or that this factor is a DP44 gain.

\subsection{7. Global use and remaining gates}\label{global-use-and-remaining-gates}

For this section assume the rank gaps, selected real types and angular weight inequalities of LIFT14, in addition to the local hypotheses above. These global hypotheses are stronger than the bare local inequalities \(k\ge1,t\ge1\).

The local theorem applies to any nonzero global first lift whose initial local representations at all G places contain the finite Steinberg type, and whose common stage characters have the literal unit restrictions(T3). Such a first input is constructed on the old degree12 host by the finite-type extension of AM-DO-030: use a full-rank residue-matrix orbit to place the actual Steinberg compact type in the auxiliary finite Schwartz vector, including the determinant character in(T10); simultaneous multijet interpolation supplies a permitted29 half-translate with nonzero selected kernel coefficient; compact Schur orthogonality then extracts an actual nonzero first lift. The exact extraction and normalizations belong to AM-DO-033, rather than being inferred from finite type occurrence alone.

For later edges the analytic induction of LIFT14 carries over after enlarging its omitted finite set S to include the G primes. All nonsplit local components and their fixed spherical tests are unchanged. Every G place is split, so Mínguez supplies local occurrence there. Outside S the same standard-parameter recurrence holds, and the partial Rallis factor remains \[
 L_{E}^{S}\!\left(\frac{d+1}{2},
             \operatorname{BC}(\pi_0)\mu_j^{-1}\right)
       \prod_{r<j}\prod_{\ell=1}^{d}L_E^{S}(\ell,\mu_r\mu_j^{-1}).
 \tag{T24}
\] Here BC denotes the unramified standard parameters, as in LIFT14, not a new global base-change theorem. The initial factor converges absolutely by the same crude unitary bound and source rank gap. Each ratio has nontrivial angular infinity type, so the \(\ell=1\) factors are finite and nonzero by the checked boundary argument in LIFT14; all other factors are in their absolutely convergent range.

Thus the same Rallis proof gives nonzero later lifts with unrestricted finite Schwartz vectors, with the selected real Schur component retained. Its joint-kernel proof continues to apply to each actual automorphic multiplicity copy. At the G places every constituent has the local form(T7), for its own initial ancestor, and hence contains(T8).

One can consequently choose nonzero stage vectors of the exact nonterminal finite column type. To justify this selection, take the product of the compact character projectors at the finitely many G places. If it killed every finite translate of a nonzero cyclic automorphic submodule, that submodule could have no constituent containing the corresponding local types, contrary to(T8). Therefore some finite translate followed by the projector is nonzero. These operations commute with the selected real type and are absorbed into the preceding finite Schwartz input by Weil equivariance. The output remains in the actual generated previous theta image.

This assertion does not impose a predetermined raw PhiG, PhiGS or Xi vector. A local type can occur while a particular local intertwiner vanishes on the previous image; an integral projector can also change content. Nor does it establish the terminal HXm condition, the odd etaT pin on the old host, the full arithmetic D7 constructor, or the six comparison identities. The degree84 host needs its own global even stage characters and first finite-type extraction hypotheses before this conditional transfer can be invoked there.

The new local primitive vector(T22) and the first extraction normalizations do not compute the later global image lattice. The missing quantitative object remains the class-integrated selected theta matrix restricted to that prior image, with the source's fixed row normalization. Its saturation, coefficient denominators, levels and archimedean size determine the actual \(g_j,e_j,H\) required by DP44.

\subsection{8. Finite diagnostic and scope of verification}\label{finite-diagnostic-and-scope-of-verification}

The exact certificate \nolinkurl{ramified_column_type_certificate.py} checks the first nonterminal model \(k=j=t=1\) over both \(\mathbf F_3\) and \(\mathbf F_5\), with a nontrivial quadratic determinant character. It verifies every matrix-line covariance identity for the constructed vector, its support and the norm in(T22), and the character-projector multiplicity. The generic branch has multiplicity1; the omitted-root branch has multiplicity2. The latter corrected an initial auxiliary test expectation of multiplicity1. No uniqueness claim is used or made in the theorem.

The same certificate verifies the forbidden-entry root count in 100 small block configurations and(T23) against exact integer Gaussian multinomials in1344 cases. These finite checks support the sign and normalization; the general local and global proofs are the arguments above. No Lean compilation, raw-Xi realization or global integral saturation check is claimed.

\subsection{Primary sources and exact source locators}\label{primary-sources-and-exact-source-locators}

\begin{itemize}
\tightlist
\item
  Casselman, W. (1980). The unramified principal series of p-adic groups. I. The spherical function. Compositio Mathematica, 40(3), 387--406. Propositions2.4 and2.6, printed pp.395--396. \url{https://numdam.org/item/CM_1980__40_3_387_0.pdf}.
\item
  Geck, M. (2015). On the l-modular composition factors of the Steinberg representation. arXiv:1504.04157v2. Theorem2.1, printed p.3, and Corollary2.3, printed pp.4--5. This records the integral Steinberg basis and proves the Borel-fixed line. \url{https://arxiv.org/pdf/1504.04157}.
\item
  Mínguez, A. (2008). Correspondance de Howe explicite: paires duales de type II. Annales Scientifiques de l'École Normale Supérieure, 41(5), 717--741. Theorem1, printed p.718, and equation(1.4), pp.721--722. \url{https://www.numdam.org/item/ASENS_2008_4_41_5_717_0.pdf}.
\item
  Steinberg, R. (1957). Prime power representations of finite linear groups II. Canadian Journal of Mathematics, 9, 347--351. Theorems1--2 and the integral-basis corollary, pp.348--350. \url{https://doi.org/10.4153/CJM-1957-041-1}. Original PDF: \url{https://www.cambridge.org/core/services/aop-cambridge-core/content/view/1A82E966D7A0832A4FD80689E5F67CAE/S0008414X00044709a.pdf/prime-power-representations-of-finite-linear-groups-ii.pdf}.
\item
  Zelevinsky, A. V. (1980). Induced representations of reductive p-adic groups. II. On irreducible representations of GL(n). Annales Scientifiques de l'École Normale Supérieure, 13(2), 165--210. Proposition8.5, p.196; Proposition8.6, p.197. \url{https://www.numdam.org/article/ASENS_1980_4_13_2_165_0.pdf}.
\end{itemize}

Pinned source root: work/research/alpoge\_source, commit 7cbf7a0c32fca1badf05931879a37c3f8a2b03b4.

\begin{itemize}
\tightlist
\item
  \nolinkurl{Formalism/LocalToGlobal/Columns/RecLevel.lean50–52:} exact inverse adelic type.
\item
  \nolinkurl{Formalism/LocalToGlobal/Columns/GenericUnipotentCharacter.lean437–452:} HX, SX, lamX, including omitted-root branch.
\item
  \nolinkurl{Formalism/LocalToGlobal/Columns/ColForms.lean92–105:} literal colLev and ColV.
\item
  \nolinkurl{Formalism/LocalToGlobal/Columns/AbelianHostFieldConditions.lean273–306,315–323,488:} chi\_om, distinctness, IsMu0 and g\_big.
\item
  AM-DO-028: actual Hom/bilinear splitting dictionary.
\item
  AM-DO-030/033: fixed auxiliary multijet and finite-type first extraction.
\item
  ALP-LIFT-14: later selected global analytic induction, including its every-vector and multiplicity scope.
\end{itemize}

\clearpage

\section{ALP-LIFT-20: the actual raw first G factor and its exact Steinberg component}\label{alp-lift-20-the-actual-raw-first-g-factor-and-its-exact-steinberg-component}

Author: GPT-6 Astra Pro, ALP\_MACHINE \nolinkurl{/root/lift_accounting}. Date: 21--22 November 2026 UTC. Inherited reasoning effort is not exposed.

This result concerns the actual source functions \nolinkurl{PhiG}, \nolinkurl{PhiGS}, \nolinkurl{PsiAt}, and \nolinkurl{XiCpl1} when \nolinkurl{m=2}. It identifies a particular allowed first coupled G factor exactly, rather than replacing that factor by a freely chosen Schwartz function. Its Steinberg component is nonzero, with a fully computed scalar and norm. Compact integration against the initial Steinberg input extracts this component while leaving the prescribed raw factor in the kernel.

The local theorem is independent of the global CM field and holds at every degree-one relative split G place satisfying the source coordinate conventions. The global corollary uses a separately established actual mixed theta family, a selected polynomial, and a valid common stage character. It does not instantiate the whole \nolinkurl{ChIdx.Ok}, identify other raw T/Q factors, handle \nolinkurl{PhiTop}, or prove a later restricted-image statement.

\subsection{1. Exact source specialization}\label{exact-source-specialization}

All source locations below refer to the recovered pinned commit \nolinkurl{7cbf7a0c32fca1badf05931879a37c3f8a2b03b4}. The principal file is \nolinkurl{Formalism/LocalToGlobal/Columns/FiniteSumNonzeroLargeDepth.lean}, abbreviated FS. No Lean execution is claimed.

Take \nolinkurl{m=2}, put \nolinkurl{a=k}, \nolinkurl{b=k+t+1}, and specialize the first coupled kernel to \nolinkurl{m2=0}. The actual rank definitions give

\[
 r=a+1,\quad r_S=0,\quad N_0=a,\quad N_1=b,
 \quad c_0=1,\quad c_1=0.
 \tag{RG1}
\]

Here \nolinkurl{cOf R j = R.m-1-j} (FS815), and \nolinkurl{rS=m/2-1} is in \nolinkurl{Formalism/LocalRingCardinality/Params.lean}. The formula is valid for every admissible source rank; no small-rank numerical test is being substituted for those rank constraints.

Fix a degree-one G place with uniformizer the rational prime \nolinkurl{q}, residue field \nolinkurl{kappa=F_q}, and integral coordinate lattice. Write \nolinkurl{O} for the local valuation ring. Fix the exact additive character \nolinkurl{psi_0} in the source, and let

\[
 G_q=\operatorname{GL}_a(\kappa),\qquad
 U_q=\{u\in G_q:u\text{ is upper unitriangular}\},
 \qquad
 \psi_U(u)=\psi_0\!\left(\sum_{i=0}^{a-2}u_{i,i+1}\right).
 \tag{RG2}
\]

The character is the actual initial nondegenerate column character. Its inverse is the adelic source-input character: \nolinkurl{Columns/RecLevel.lean}50--52 defines \nolinkurl{levTyp=levChr^(-1)}.

In \nolinkurl{PsiAt} the second frame is literally one (FS984--986). The first frame is also one at this stage: \nolinkurl{ChIdx.Ok.zero_h} states \nolinkurl{L.h 0 b P=1} for every initial index (FS1019--1023). Thus this specialization does not impose a new choice of initial frame beyond the source's own condition. In arbitrary global lattice frames, the following calculation is made after applying the exact \nolinkurl{inFr} and target inverse frame in \nolinkurl{XiCpl1}.

Choose the allowed target residue class

\[
                       y\equiv q^{-1}e_a\pmod{O^b}.
 \tag{RG3}
\]

The subscript is zero-based: \nolinkurl{e_a} is the first coordinate immediately after the first \nolinkurl{a} target coordinates. This class belongs to \nolinkurl{q^(-1)O^b/O^b}, exactly the G component in \nolinkurl{ybSet} (FS1125--1133). The assertion is independent of the representative chosen by \nolinkurl{repY}.

Let \nolinkurl{X} be a \nolinkurl{b by a} matrix. The definition \nolinkurl{joinCol X y} is the \nolinkurl{b by (a+1)} matrix \nolinkurl{[X,y]} here (FS981--983). Both index permutations \nolinkurl{fIdx} in \nolinkurl{entB} are the identity in (RG1). The source last-column scaling exponent is one, all its other exponents are zero, and every target exponent is zero. Hence

\[
             \operatorname{entB}(\operatorname{joinCol}(X,y))
                           =[X,qy].
 \tag{RG4}
\]

This follows directly from FS684--694; it is not a change of Schwartz normalization. Since both group arguments of \nolinkurl{omPhi} are one, all determinant, valuation, and splitting factors in \nolinkurl{Loc.om} are one (FS553--561, 940--945).

\subsubsection{The raw-factor identity}\label{the-raw-factor-identity}

Put \nolinkurl{X_0=[I_a;0]} and define

\[
 f_U(X)=
 \begin{cases}
   \psi_U(u),&X\in M_{b,a}(O),\quad\bar X=X_0u,\ u\in U_q,\\
   0,&\text{otherwise}.
 \end{cases}
 \tag{RG5}
\]

Then, for either value of the source \nolinkurl{sharp} flag,

\[
 \boxed{\ \operatorname{PsiAt}(0,P,1,X,y)=f_U(X)\ }
                  \qquad\text{whenever (RG3) holds.}
 \tag{RG6}
\]

Here is the complete predicate check. In both \nolinkurl{PhiG} and \nolinkurl{PhiGS}, the internal variable called \nolinkurl{a} is now \nolinkurl{r=a+1}, since the stage parameter is \nolinkurl{j=1}. The first integrality condition requires all entries in \nolinkurl{[X,qy]} to be integral. The condition on rows at least \nolinkurl{a+1} makes their first \nolinkurl{a+1} entries zero modulo \nolinkurl{q}. In the first \nolinkurl{a+1} rows and columns, the first \nolinkurl{a+1} diagonal entries are one and the entries below the diagonal are zero. The last column is already \nolinkurl{e_a} modulo \nolinkurl{q} by (RG3). Consequently the first \nolinkurl{a} columns have upper unitriangular top \nolinkurl{a by a} block and every subsequent row is zero. This is precisely the support in (RG5).

For \nolinkurl{PhiGS}, \nolinkurl{symIdx a 0} is the identity (\nolinkurl{Columns/LevGSym.lean}11--13). Its extra zero condition in column \nolinkurl{a} is already imposed by \nolinkurl{qy=e_a} modulo \nolinkurl{q}. Its other off-diagonal conditions are exactly the same lower zeros just used. There are no determinant-block conditions, since their column range starts at \nolinkurl{r} and ends strictly before \nolinkurl{r}.

On the support, the leading \nolinkurl{(a+1) by (a+1)} matrix is \nolinkurl{diag(u,1)} modulo \nolinkurl{q}, so its determinant is one. The inverse \nolinkurl{chi_0} determinant factor therefore equals one, even though the common splitting character itself can be ramified. In \nolinkurl{lamGvS}, the second additive sum has length \nolinkurl{r-a-1=0}. The first sum is exactly the exponent in (RG2). In \nolinkurl{lamGv}, the only additional simple-root entry is the entry in row \nolinkurl{a-1}, column \nolinkurl{a}; it is zero by (RG3). Thus both functions give \nolinkurl{psi_U(u)}. Off the stated support one of the displayed source predicates fails. This proves (RG6) using the literal definitions FS780--810 and FS927--931.

Adding an integral vector to \nolinkurl{y} changes \nolinkurl{qy} by a multiple of \nolinkurl{q}, so the same proof also establishes representative independence directly; no unproved \nolinkurl{ChIdx.Ok.psiInv} is used to obtain (RG6).

At finitely many G places the classes (RG3) can be imposed simultaneously, including after the actual target frames. Indeed the additive Chinese remainder theorem for the localized fractional lattice gives an element of the global target vector space with those finitely many prescribed classes and with the required integral classes at the other finite places in \nolinkurl{rowSG}. Its class is an element of the source \nolinkurl{ybSet}; every representative of that class has the same local factors. This is only an existence statement for that allowed class. No archimedean height bound for a chosen rational representative is asserted.

\subsection{2. Finite Fourier coefficients and exact raw norm}\label{finite-fourier-coefficients-and-exact-raw-norm}

Give \texttt{M\_\{b,a\}(O)} additive volume one, as in the actual split Schrödinger model. Each residue cell has volume \nolinkurl{q^(-ab)}. Set

\[
            n_U=a(a-1)/2,\qquad d_q=q^{n_U}=|U_q|.
 \tag{RG7}
\]

Every nonzero value of \nolinkurl{f_U} is a root of unity, its value at \nolinkurl{X_0} is one, and

\[
                         \|f_U\|^2=q^{-ab}d_q.
 \tag{RG8}
\]

This is its exact Schwartz norm, not a normalized finite group-counting norm. Its finite value lattice is primitive at every prime at which the character values are units. In particular it is primitive at11 when \nolinkurl{q} is prime to11.

There is an especially simple actual Heisenberg expansion. Let \texttt{V=M\_\{b,a\}(kappa)} and let \nolinkurl{W} be the linear subspace of matrices whose only possible nonzero entries are in the strictly upper triangular part of their top \nolinkurl{a by a} block. Then \nolinkurl{dim W=n_U}, the support of \nolinkurl{f_U} is the affine space \nolinkurl{X_0+W}, and its character there is \nolinkurl{ell(w)=sum_i w_(i,i+1)}. Use the perfect finite pairing \nolinkurl{B} and actual Fourier phase of DO29.11--13; identify \nolinkurl{ell} with the corresponding character of \nolinkurl{W}. For the actual finite vacuum \nolinkurl{phi_0} one has

\[
 f_U=\sum_y b_y W_y\phi_0,\qquad
 b_y=
 \begin{cases}
 q^{n_U-ab}\psi_0(-B(X_0,y)),&B(w,y)=\ell(w)\ \forall w\in W,\\
 0,&\text{otherwise}.
 \end{cases}
 \tag{RG9}
\]

If the finite pairing is expressed as characters rather than scalar-valued forms, the condition means equality of the two characters on \nolinkurl{W}; this is invariant under the coordinate choice. To prove the formula, insert \nolinkurl{X=X_0+w} into the finite Fourier formula \nolinkurl{q^(-ab) sum_X f_U(X) psi_0(-B(X,y))}. The sum over \nolinkurl{W} is \nolinkurl{q^(n_U)} when its character is trivial and zero otherwise. Perfectness of the pairing gives exactly \nolinkurl{q^(ab-n_U)} nonzero coefficients. Therefore

\[
 \#\{y:b_y\ne0\}=q^{ab-n_U},\qquad
 \sum_y|b_y|=1,\qquad
 \sum_y|b_y|^2=q^{n_U-ab}.
 \tag{RG10}
\]

These are the actual phase-correct Heisenberg operators, not just abstract coordinate translations. They are unitary, so the operator in (RG9) has norm at most one. Its exact value-lattice and Fourier-coefficient lattices should not be confused: the latter have explicit \nolinkurl{q}-power denominators \nolinkurl{q^(ab-n_U)}. These denominators are units at11 but must be retained at the auxiliary prime and in any global coefficient-height accounting.

\subsection{3. Exact Steinberg projection in the actual split action}\label{exact-steinberg-projection-in-the-actual-split-action}

Let \nolinkurl{eta} be the local selected-half common splitting character. In the literal source character pin, \nolinkurl{bar eta=chi_0^(-1)}. The genuine left source action is

\[
 \omega(g)f(X)=\eta(\det g)^{-1}
                    |\det g|^{b/2}f(Xg).
 \tag{RG11}
\]

This is the corrected DO33.5/LIFT16(S2) convention. On the integral compact it factors on the present functions through \nolinkurl{G_q} with its displayed determinant correction. In particular \nolinkurl{f_U} has right-source \nolinkurl{psi_U} type under \nolinkurl{U_q}; the determinant factor is trivial there.

Use the actual finite Steinberg representation \nolinkurl{St_q} and the unit generic vector \nolinkurl{u_0} constructed in DO33.2--4. Its dimension is \nolinkurl{d_q}; the \nolinkurl{psi_U} eigenspace is one dimensional. These facts are the published finite Steinberg theorems (Steinberg, 1957, Theorems1--2; Geck, 2015, Theorem2.1) and the already checked normalization DO033. With the Hermitian form conjugate-linear in its first argument put

\[
 m(g)=\langle u_0,\operatorname{St}_q(g)u_0\rangle,
 \qquad
 f_{\rm St}(X_0g)=\eta(\det g)m(g),
 \tag{RG12}
\]

and extend \nolinkurl{f_St} to residue cells, zero off the full-rank orbit \nolinkurl{X_0G_q}. This is exactly DO33.6 with \nolinkurl{v'=u_0}. It has pure source \nolinkurl{St_q} type, value one at \nolinkurl{X_0}, and

\[
                 \|f_{\rm St}\|^2=
                          q^{-ab}|G_q|/d_q.
 \tag{RG13}
\]

Let \nolinkurl{P_St} be the orthogonal projection onto the full \nolinkurl{St_q} isotypic subspace for the actual action (RG11), including all its multiplicity copies in the residue Schwartz space. Then

\[
 \boxed{\ P_{\rm St}f_U=\alpha_q f_{\rm St},\qquad
                  \alpha_q=\frac{d_q^2}{|G_q|}>0.\ }
 \tag{RG14}
\]

Here is a proof that controls the whole isotypic projection, not only its nonzero inner product. On the free orbit \nolinkurl{X_0G_q}, tensoring functions by \nolinkurl{eta(det g)} identifies the actual right-source action with the ordinary right regular representation. It leaves \nolinkurl{f_U} unchanged, since every point in its support has determinant one. The Steinberg matrix coefficients span the full Steinberg isotypic space in this regular representation. For any \nolinkurl{v,w in St_q}, their inner products with \nolinkurl{f_U} are

\[
 \begin{split}
 &q^{-ab}\sum_{u\in U_q}
       \overline{\langle w,\operatorname{St}_q(u)v\rangle}
                                         \psi_U(u)\\
 &\hspace{12mm}=q^{-ab}d_q\,
       \overline{\langle w,u_0\rangle\langle u_0,v\rangle}.
 \end{split}
 \tag{RG15}
\]

The equality is the finite \nolinkurl{psi_U} projector and the one-dimensionality of its image. Thus the projection is a multiple of the single coefficient \nolinkurl{m}, including in the multiplicity direction. By finite Schur orthogonality the squared norm of \nolinkurl{m} on \nolinkurl{G_q} is \nolinkurl{|G_q|/d_q}. Its inner product with \nolinkurl{f_U} is \nolinkurl{d_q}. The scalar is therefore exactly \nolinkurl{d_q^2/|G_q|}, proving (RG14). All other residue orbits are invariant under the source compact and have zero restriction of \nolinkurl{f_U}, so they do not add any component to this projection.

Consequently

\[
 \|P_{\rm St}f_U\|^2
       =q^{-ab}\frac{d_q^3}{|G_q|},\qquad
 \frac{\|P_{\rm St}f_U\|^2}{\|f_U\|^2}
       =\frac{d_q^2}{|G_q|}=\alpha_q.
 \tag{RG16}
\]

The same calculation gives the more precise compact filter actually used in a theta lift:

\[
 \boxed{\ \int_{\operatorname{GL}_a(O)}
       \overline{m(\bar g)}\,\omega(g)f_U\,dg
              =\frac{\alpha_q}{d_q}f_{\rm St}.\ }
 \tag{RG17}
\]

The compact has measure one. Finite Schur orthogonality makes this filter \nolinkurl{d_q^(-1)} times the projection onto \nolinkurl{u_0} on each Steinberg copy and zero on other irreducible types. Equation(RG14) then proves (RG17). This also verifies the sign: the kernel has \nolinkurl{psi_U} covariance and the source matrix coefficient is its complex conjugate, with covariance \nolinkurl{psi_U^(-1)} as required by the literal bilinear lift.

There is no claim that \nolinkurl{P_St} is an integral projector. For a fixed place above11 with \nolinkurl{q} prime to11,

\[
 v_{11}(\alpha_q)=-v_{11}(|G_q|)
             =-\sum_{i=1}^{a}v_{11}(q^i-1).
 \tag{RG18}
\]

The explicit denominators are part of the compact averaging cost. The raw factor (RG5) itself remains primitive.

\subsection{4. A first global lift retaining the raw G factors}\label{a-first-global-lift-retaining-the-raw-g-factors}

This section is a corollary of the local theorem and the already established actual-family/jet/extraction inputs, not an assertion of new global theta-family existence. Use DO033 on the actual \nolinkurl{J^2} host, or its separately verified degree-seven host version DO034 once that version and its rank/character hypotheses are in place. Keep a selected primitive polynomial \nolinkurl{P}, its actual period scalar \nolinkurl{mathcal P}, and its content upper bound \nolinkurl{B_P}. In DO033, \nolinkurl{mathcal P=C^(2ab)}; on a different host its own period formula must be used. Fix any further nonzero auxiliary finite factor for which the actual simultaneous-jet theorem has been established. It is retained in all its norms and coefficient bounds. For clarity, the displayed formulas below take that additional factor to be the normalized vacuum; its known scalar factors can be multiplied back.

At each selected G place take the exact raw factor (RG5), arising from the allowed class (RG3). Write

\[
 d_G=\prod_vd_{q_v},\quad
 A_G=\prod_v\alpha_{q_v},\quad
 \Gamma_U=\prod_vq_v^{-ab}d_{q_v},\quad
 \Gamma_{\rm St}=\prod_vq_v^{-ab}|G_{q_v}|/d_{q_v}.
 \tag{RG19}
\]

Then \nolinkurl{Gamma_U=A_G Gamma_St}. The exact raw G Schwartz norm is \nolinkurl{sqrt(Gamma_U)}, and its tensor Heisenberg expansion has coefficient absolute sum one and \nolinkurl{prod_v q_v^(ab-n_U)} nonzero terms. The Steinberg projection is \nolinkurl{A_G} times the tensor of the DO033 vectors.

Apply the actual simultaneous-jet theorem to that tensor of the nonzero \nolinkurl{f_St} vectors, not merely to an abstract finite representation. With its Fourier support size \nolinkurl{s_St}, one can take a single permitted29 half-translate at

\[
                   M=29^e\ge3s_{\rm St}(D+1)-1,
 \qquad s_{\rm St}\le\prod_vq_v^{ab}.
 \tag{RG20}
\]

It gives an actual selected coefficient \nolinkurl{a_lambda} with \nolinkurl{a_lambda != 0} and \nolinkurl{v_11(a_lambda)<=B_P}, after the same period removal and coordinate hypotheses as in DO033/034. The29 operator and the real selected differential commute with the G-compact action because they act at other places. Thus the same individual permitted translate works before the compact filtering of the raw factor; no replacement of \nolinkurl{f_U} by \nolinkurl{f_St} in the prescribed kernel is necessary.

Let \nolinkurl{I=vol(K*)^(-1)} be the exact source compact index for the selected29 level, with full integral compacts at G, and let \nolinkurl{d_rho} be the selected real source dimension. The identity class has trivial rational stabilizer at the levels in DO033/034. Use their actual class-supported source function \nolinkurl{F=conj(m_infty product_v m_v)}, zero on other classes. It has initial literal \nolinkurl{psi_U^(-1)} type at every G place. Successive application of (RG17) and real Schur orthogonality gives the exact raw-kernel value

\[
 \boxed{\ \Theta_{\rm raw}^{\rm bil}(F)(1)
       =\frac{A_G\,\mathcal P\,a_\lambda}{I d_\rho d_G}
       \ne0,\qquad
 \|F\|_2^2=\frac1{I d_\rho d_G},\quad\|F\|_\infty=1.\ }
 \tag{RG21}
\]

The proof performs the compact integral on the raw theta kernel, replacing its G factor by the right-hand side of (RG17) inside that integral. Absolute convergence of the definite theta integral and the finite compact average justify this operation. It is exactly the source's bilinear convention, not a silently conjugated-input theta map.

After making the raw finite input and source input unit norm, the value becomes

\[
 \Theta_{\rm raw,unit}^{\rm bil}(F_{\rm unit})(1)
       =\frac{A_G\,\mathcal P\,a_\lambda}
                    {\sqrt{\Gamma_U}\sqrt{I d_\rho d_G}}.
 \tag{RG22}
\]

If the real Schwartz vector is also made unit, divide by its exact norm from R011/DO034. Every factor is displayed; in particular the Steinberg averaging factor is not omitted. For \nolinkurl{q_v} prime to11,

\[
 \begin{split}
 v_{11}(\Gamma_U)&=v_{11}(d_G)=0,\\
 v_{11}(A_G)&=-\sum_v\sum_{i=1}^{a}v_{11}(q_v^i-1),\\
 v_{11}\!\left(\mathcal P^{-1}
                 \Theta_{\rm raw}^{\rm bil}(F)(1)\right)
   &\le B_P-v_{11}(I d_\rho)
                  -\sum_v\sum_{i=1}^{a}v_{11}(q_v^i-1).
 \end{split}
 \tag{RG23}
\]

This is an upper content witness for one actual first coefficient. It is neither a lower divisibility theorem for all rows nor a comparison with the full global DP44 cost. It also does not say that the original raw theta kernel has only Steinberg type.

For \nolinkurl{f} selected places above each of two completely split rational G primes \nolinkurl{g_0,g_1}, the local factors specialize to

\[
 \begin{split}
 d_G&=(g_0g_1)^{f a(a-1)/2},\\
 \Gamma_U&=(g_0g_1)^{f(-ab+a(a-1)/2)},\\
 A_G&=(g_0g_1)^{-fa}
       \prod_{\epsilon=0}^{1}\prod_{i=1}^{a}
                                     (1-g_\epsilon^{-i})^{-f}.
 \end{split}
 \tag{RG24}
\]

These identities use \nolinkurl{|GL_a(F_q)|=q^(a^2) product_i(1-q^(-i))} and are field-general local counts. The source index \nolinkurl{I} must still be the correct one for the actual global host: old \nolinkurl{J^2} and repaired \nolinkurl{J^14} have different rational29 ramification and therefore different exponents in that index.

\subsection{5. What this closes and what remains}\label{what-this-closes-and-what-remains-1}

The local equality (RG6) closes the raw initial G-factor identification for the displayed allowed coupled class when \nolinkurl{m=2}. Equations(RG14) and(RG17) prove nonzero interaction with the actual initial Steinberg input and compute its exact cost. The conditional global corollary proves that this raw G factor can be retained in the actual first lift once the cited actual-family and selected-jet hypotheses are satisfied.

The statement does not identify a complete \nolinkurl{XiCpl1} unless its other actual T/Q factors and their finite family are also supplied. It does not prove nonvanishing of \nolinkurl{XiAt} at later stages for \nolinkurl{m>2}, identify the terminal \nolinkurl{PhiTop}, construct all six \nolinkurl{AdmS} ties, or compute the full DP44 changed-data cost. Those are separate exact source gates. No isotypic occurrence is being substituted for preservation of a prescribed later arithmetic image lattice.

\subsection{6. Sources and checks}\label{sources-and-checks-1}

Accepted inputs: DO033's generic vector, actual split action, phase-correct Fourier model, simultaneous jets, and compact class extraction; DO029/030 for the actual Heisenberg and auxiliary multijet comparison; R012/R011 for the selected real polynomial and norm. Their proved scope is retained. DO034 is needed separately before the new host is used in the global corollary.

Steinberg, R. (1957). Prime power representations of finite linear groups II. \emph{Canadian Journal of Mathematics, 9}, 347--351. \url{https://doi.org/10.4153/CJM-1957-041-1} . Theorems1--2, pp.348--350: published Steinberg construction, irreducibility, and regular restriction to the upper unipotent. Read in the publisher PDF during this run.

Geck, M. (2015). \emph{On the ell-modular composition factors of the Steinberg representation}. \url{https://arxiv.org/pdf/1504.04157} . Theorem2.1 and Corollary2.3, pp.3--5: the integral flag model and characteristic-zero properties used in DO033.

Author whole-proof reading was completed on 21 November 2026, including all source predicates, the full isotypic multiplicity argument, the compact filter, the conditional global extraction, and every displayed norm/valuation factor. The exact diagnostic passed 729 residue-matrix predicate comparisons, 48 full Steinberg projections, 48 matrix-coefficient filters, and the full 729-point Fourier transform. Its raw norm is \nolinkurl{1/243}, projection scalar \nolinkurl{3/16}, and projected norm is \nolinkurl{1/1296}. The separate TSV records the check time and exact scope.

The exact diagnostic \nolinkurl{raw_first_g_steinberg_certificate.py} checks both raw predicates and the full finite Steinberg projection/filter in \nolinkurl{GL_2(F_3)}, with a nontrivial determinant splitting and all residue matrices. It also checks the explicit raw Fourier expansion and all displayed finite norms. This diagnostic is not a proof of the general rank statement or a verification of a full Lean datum.

\clearpage

\section{AM-K031 --- exact enlarged-host ranks, uniform stage characters, and the unchanged half-orbit isogeny}\label{am-k031-exact-enlarged-host-ranks-uniform-stage-characters-and-the-unchanged-half-orbit-isogeny}

Author: GPT-6 Astra Pro, ALP\_MACHINE \nolinkurl{/root/kummer_extension}, 21 November 2026 UTC; inherited reasoning effort not separately exposed.

This finding supplies concrete numerical rank/gap data, compatible split primes and actual Hecke characters on the CM host of AM-K029/030. It also proves that the permitted DP59 half-orbit survives the real-field enlargement. It is deliberately not a \nolinkurl{Dat.Ok} or D7 instance: the unmodified source still requires \(8\mid f\) and a common quadratic CM subfield for the selected type. The latter is the already audited primitive-type obstruction of AM-R-003, not a new general impossibility theorem.

Throughout the first three sections, \[
 E_0=\mathbb Q(\zeta_{28}),\quad F_0=E_0^+,\quad
 K=E_0B,\quad F=F_0B,\quad [B:\mathbb Q]=7,
\] where \(B\subset\mathbb Q(\zeta_{29})\), and the actual type contains every lift of \(\Phi_0=\{1,3,9,15,17,23\}\). Thus \([F:\mathbb Q]=f=42\), and the actual principal abelian variety is \(J^{14}\). Let \(\eta\) be the literal \nolinkurl{IsEtaT} character constructed in AM-K029.

\subsection{1. The source rank constraints and an explicit gap locus}\label{the-source-rank-constraints-and-an-explicit-gap-locus}

The exact record \nolinkurl{Formalism/LocalRingCardinality/Params.lean:6–22} imposes \[
 \begin{gathered}
 k,m\text{ even},\qquad t\text{ odd},\qquad 2\le m,
 \qquad m+3\le k,\\
 k+m+4\le t,\qquad 3m<t,\qquad
 m\le f=2\operatorname{ord}\le k,\qquad
 m^2\le k+m(t+1).
 \end{gathered}
\tag{HC1}
\] The real-place type count is tied to this f by \nolinkurl{AbelianHostFieldConditions.lean:567}. In particular the old \(k=6\) cannot remain the starting rank on the degree-84 host: \[
                         \operatorname{ord}=21,\qquad k\ge42.
\tag{HC2}
\] The inherited \nolinkurl{Frame} record itself causes no numerical obstruction to order21. Take its additive group to be \(\mathbb Z/21\mathbb Z\), its cyclotomic parameter to be zero, its projection to the trivial group \(HGrp(0)=(\mathbb Z/4\mathbb Z)^\times/\langle-1\rangle\), and its distinguished element to be zero. The projection and cyclic-kernel clauses in \nolinkurl{ZsmulGeneratorEqZeroOfProjection.lean:16–24} then hold. This is a mathematical record construction, not a claim of Lean compilation or of its compatibility with every later host record.

For the shortest allowed length \(m=2\) and smallest starting rank \(k=42\), the rank clauses alone give the least odd t as49. The numerical Hyp3 condition \(1801\le r+t\), where \(r=k+m-1\), is stronger. A convenient simultaneous choice is \[
 \boxed{f=k=42,\quad m=2,\quad r=43,\quad t=1759,
         \quad d=t+1=1760.}
\tag{HC3}
\] It gives \[
       n_V=r+t=1802,
       \qquad (N_0,N_1,N_2)=(42,1802,3562).
\tag{HC4}
\] The actual G-size clause is \nolinkurl{AbelianHostFieldConditions.lean:488}, namely \[
                 2fk<g_i-1.
\tag{HC5}
\] Here its right lower threshold is3528, not the value obtained by retaining \(k=6\).

The following uniform parameters verify all numerical parity, interlacing, centred-distance and opposite-sign gap inequalities occurring in Hyp3. Put \[
 \begin{split}
 S&=r+mt=3561,\qquad S^3=45\,156\,047\,481,\\
 C&=S^3+k+t+11=45\,156\,049\,293.
 \end{split}
\tag{HC6}
\] At every one of the 42 embeddings use zero-based indices and set \[
 \begin{aligned}
 u_i&=(2(k-i)+1)C+\tfrac12 &&(0\le i<k),\\
 v_i&=2(k-i)C &&(0\le i<k-1),\\
 b_j&=(4j-2)C,\qquad e_j=4jC &&(j=0,1),\\
 \eta_{\mathrm{num}}&=-\tfrac12.
 \end{aligned}
\tag{HC7}
\] In particular \[
 (e_0,e_1)=(0,180\,624\,197\,172),\qquad
 (b_0,b_1)=(-90\,312\,098\,586,90\,312\,098\,586).
\tag{HC8}
\]

\textbf{Verification.} The u's are half-integral, the other displayed weights are integral, and \(u_i>v_i>u_{i+1}\), while \(u_{k-1}=3C+1/2>-b_0=2C\). Both consecutive b/e gaps are \(2C\), exceeding \((t+3)/2=881\).

For the exact \nolinkurl{Centred} predicate, \(u_i+e_j\ge3C+1/2\), \(v_i+e_j\ge4C\), and \(|e_j-b_i|\ge2C\). The only cross-block term has centre difference \(4C\) and a block half-width \((t-1)/2=879\), so its absolute value is at least \(4C-879>881\). These are all four kinds of quantifiers in \nolinkurl{DataTheta.lean:550–555}.

For \nolinkurl{NCGaps}, the literal boolean in \nolinkurl{Dpar} is \(i\le j\) or \(i-j\) even. At stage0 the only false-sign entry is \(-b_1=-2C\); the other b-entry is \(2C\), and all v-entries are at least \(4C\). Thus every opposite-sign distance is at least \(4C>S^3\). At stages1 and2 every entry has true sign, so no opposite-sign comparison remains. The common odd-stage shift by \(\eta_{\mathrm{num}}\) changes no distance. This is the full finite sign analysis for \(m=2\), not a gap check only for adjacent entries.

The source's extra \nolinkurl{e_even} condition at line578 also holds: \(\sum_\tau e_j=42e_j\) is even. In the half-angular convention of \nolinkurl{HasInfQ}, the constructed odd character has infinity type \(+1/2\), so the required numerical sign is precisely \(\eta_{\mathrm{num}}=-1/2\). This checks the numerical equality; an actual \nolinkurl{ThetaX.infTy} realization remains part of the analytic adapter.

For comparison with AM-DO-025, the first selected partition at each place is \[
           \nu_{0,i}=(85-2i)C+i-921\quad(0\le i<42),
\tag{HC9}
\] whose smallest entry is \(3C-880>0\). The next partition is \[
         \nu_1=(\nu_0+\Delta,\Delta^{1760}),
         \qquad \Delta=4C-1760>0.
\tag{HC10}
\] These follow by substitution into the already checked selected-weight recurrence, not by asserting that an old six-place polynomial remains a single 42-place weight. A new selected polynomial is constructed in the new CM blocks separately.

\subsection{2. Exact conductor and split G primes, including PiGCong}\label{exact-conductor-and-split-g-primes-including-pigcong}

The abelian field K has rational conductor \[
                              N=28\cdot29=812.
\tag{HC11}
\] It lies in \(\mathbb Q(\zeta_{812})\). Conversely, a cyclotomic field containing K contains i and \(\zeta_7\), forcing respectively \(4\mid N'\) and \(7\mid N'\); ramification at29 forces \(29\mid N'\). Hence \(812\mid N'\). This verifies the minimum-divisibility formulation of \nolinkurl{IsCond} in \nolinkurl{RealCyclotomicTwoPower.lean:54}, including every conjugate embedding since K is normal over \(\mathbb Q\).

Write \(\delta=\sqrt{-7}\in E_0\), and use the following \textbf{explicit} primes and generators in the external quadratic subfield \(K_q=\mathbb Q(\delta)\): \[
 \begin{array}{c|c|c}
 i&g_i&\pi_i\\\hline
 0&4\,615\,409&1+812\delta\\
 1&5\,276\,377&813+812\delta.
 \end{array}
\tag{HC12}
\] Both rational integers are prime. The accompanying certificate checks primality by trial division through the integer square root, as well as \[
 \pi_i\bar\pi_i=g_i,\qquad
 \pi_i\equiv1\pmod{812\mathcal O_{K_q}},\qquad
                   g_i\equiv1\pmod{812}.
\tag{HC13}
\] Their conjugate ideals are coprime: a common prime would divide both the norm \(g_i\) and the sum \(2\operatorname{Re}\pi_i\), whereas the corresponding rational gcd is one. Coprimality persists after extension to K.

Thus both primes split completely in K. The support of \(\bar\pi_i\) consists of42 distinct primes, each with residue field canonically \(\mathbb F_{g_i}\) through its prime field. Equations (HC12)--(HC13) supply the norms, coprimality, cardinalities, residue norms, distinctness, full \nolinkurl{PiGCong}, and strict G-size inequalities. The value \(812\) in \nolinkurl{PiGCong} is not replaced by28. The explicitly retained exception is \nolinkurl{piG_Kq}: the source's \textbf{defined} common field \nolinkurl{Fields.Kq} is \(\mathbb Q\) for this type, rather than the external \(K_q=\mathbb Q(\sqrt{-7})\) used here. Section6 explains that literal failure.

At each i choose a faithful character \(\omega_i:\mathbb F_{g_i}^{\times}\to\mu_{g_i-1}\). For example 3 and5 are primitive roots for \(g_0\) and \(g_1\), respectively; the certificate checks their orders using the complete factorization of \(g_i-1\). Use the literal residue characters \[
             \chi_{j,P}=\omega_i^{2(j+1)}
                       \qquad(P\mid\bar\pi_i).
\tag{HC14}
\] They satisfy the equality-on-prime-field and \nolinkurl{chi_om} clauses. For \(j=0,1\) they are nontrivial and distinct. Their values lie in \[
        \mathbb Q(\zeta_L),\qquad
        L=\operatorname{lcm}(28,g_0-1,g_1-1)=14\,995\,460\,592.
\tag{HC15}
\] This is a fixed coefficient extension; its large degree is not discarded in any later height budget.

\subsection{3. Actual global stage characters for every uniform integral weight}\label{actual-global-stage-characters-for-every-uniform-integral-weight}

\subsubsection{Theorem HC16}\label{theorem-hc16}

For any integers \(e_j\), there are continuous unitary Hecke characters \(\mu_j\) of K which are trivial on \(\mathbb A_F^\times\), tame at every finite place, unramified away from \(g_0g_1\), have infinity component \[
       (\mu_j)_\infty(z)
             =\prod_{\tau\in\Psi}
                    \left(\frac{z_\tau}{|z_\tau|}\right)^{-2e_j},
\tag{HC16}
\] and satisfy \[
       \mu_{j,P}|_{\mathcal O_{K,P}^\times}=\chi_{j,P}^{-1}
                  \qquad(P\mid\bar\pi_i).
\tag{HC17}
\] These are the exact \nolinkurl{IsMu0} character conditions for a uniform splitting weight. The values required as \nolinkurl{muVal} lie in a single finite coefficient field, after adjoining the host embeddings and normal closure if required. There is no new class-number-one assumption on K or F.

\textbf{Construction and proof.} Apply the checked ALP-LIFT-18 construction on E₀ to these new split primes, with \textbf{zero} infinity weight at every embedding. Its root compatibility is then automatic. Concretely, on the compact norm-one torus its character is the product of \(\omega_i^{-2(j+1)}\) at the six selected E₀-primes. The product on global roots is one because \(N_{E_0/K_q}(\mu_{28})=1\), as proved in LIFT18(G2)--(G3). The checked class-one torus decomposition of AM-R-009 makes this an actual automorphic torus character. Composing with \(x\mapsto x/\bar x\) gives a finite-order Hecke character \(\nu_j\) of E₀, trivial on \(\mathbb A_{F_0}^\times\), with zero infinity type and the required inverse residue characters. Its values lie in the roots of unity in (HC15). This is an application of the existing character theorem; its class-number proof is not rederived here.

Pull this actual character back to K by the idele norm: \[
              \widetilde\nu_j=\nu_j\circ N_{K/E_0}.
\tag{HC18}
\] At every chosen G prime the extension is completely split, so the completion map is \(K_P=(E_0)_{P_0}=\mathbb Q_{g_i}\). On an idele supported in one such component, the local norm is the identity. Thus (HC18) has the required residue character at \textbf{each} of the seven lifts of every old selected prime; it does not raise a single local residue character to the seventh power. Norms take local units to units at every other place, so the pullback is unramified there, including the new ramification place29. The norm of an F-idele lies in \(\mathbb A_{F_0}^\times\), hence the pullback is trivial on \(\mathbb A_F^\times\).

Now use the actual odd character of AM-K029: \[
                   \boxed{\mu_j=\widetilde\nu_j\,\eta^{-2e_j}.}
\tag{HC19}
\] The restriction of η to \(\mathbb A_F^\times\) is the quadratic character \(\eta_{K/F}\). Its even power in (HC19) is therefore trivial on that group. Since η is everywhere finite-unramified, it changes no compact residue restriction or tame conductor. Its angular exponent is one, so its \(-2e_j\) power supplies exactly (HC16), which is the literal \nolinkurl{HasInfQ(-e_j)} convention. All the claimed properties follow.

For the two stages in (HC3), the finite conductor is exactly \[
          \mathfrak f(\mu_j)=g_0g_1\mathcal O_K,
          \qquad N_{K/\mathbb Q}\mathfrak f(\mu_j)=(g_0g_1)^{84}.
\tag{HC20}
\] Both halves have conductor exponent one because the selected residue characters are nontrivial and base-field triviality forces the inverse character at the conjugate place. Unramified η cannot cancel a nontrivial residue character.

To check the common algebraic field, let \(a_P=\eta(\varpi_P)(NP)^{-1/2}\), the \nolinkurl{etaVal} supplied in AM-K029. At a finite uniformizer, \[
        \mu_j(\varpi_P)
          =\widetilde\nu_j(\varpi_P)\,a_P^{-2e_j}(NP)^{-e_j}.
\tag{HC21}
\] The first factor is in the fixed finite group of roots from (HC15), the second lies in the fixed \nolinkurl{etaVal} field, and the last is rational. Thus one coefficient field works for every integer choice of the uniform weights. Taking its normal closure and adjoining K preserves finiteness. No effective small degree, discriminant or uniformizer-height bound is inferred. ∎

This is stronger than merely pulling back LIFT18's nonzero infinity weights, which would retain its sufficient \(7\mid e_j\) restriction. Here all integer uniform weights work. For example the explicit \(C\) in (HC6) is \(5\bmod7\), so \(e_1=4C\equiv6\bmod7\); it is covered by (HC19) and not by the old uniform congruence.

The construction is stable under further totally real extension: pull η back by norm, and pull \(\nu_j\) directly from E₀. At G-primes chosen completely split in the enlarged CM field, each local norm remains the identity. Every integer uniform infinity weight again works. This gives a parameterized character construction for later real-field repairs without reusing an unjustified large-field torus class-number-one assertion.

\subsection{4. The fixed alpha and the actual permitted DP59 half-orbit}\label{the-fixed-alpha-and-the-actual-permitted-dp59-half-orbit}

The alpha used in AM-K022/DO030 is in \(\mathbb Q(\sqrt{-7})\), not in \(\mathbb Q(i)\): \[
       \alpha=(1+2\delta)^e,\quad M=29^e,\quad e\ge1,
                   \qquad \alpha\bar\alpha=M.
\tag{HC22}
\] The alternative \((5+2i)^e\) has the same norm identity, but changing to it is unnecessary. We retain (HC22), which acts diagonally on every ordinary J factor.

The ideals of \(1+2\delta\) and its conjugate are coprime in every extension of E₀. A common prime would divide their sum2 and their product29, which is impossible. Thus the ideals of α and \(\bar\alpha\) remain coprime in K and in every further real extension. At rational29, their valuations in K are \(7e\) on the two disjoint six-prime halves, rather than e. This ramification does not change the element identity or the coprimality statement.

The element \(1+2\delta\) is congruent to one modulo2. It therefore acts trivially on the relevant two-torsion characteristic, and the exact odd symmetric pullback normalization carried in DP59 applies to the same invariant theta line. Under the K030 polarized product isomorphism, α is still the same diagonal endomorphism on all14 J factors: it lies in the central old E₀-action and commutes with the real trace-lattice transport.

More generally suppose a further real-field construction realizes the induced host as \(J^{14h}\) with its product line, h a positive integer. On an edge of ranks a,b put \[
       X=J^{14hab},\quad g_X=42hab,\quad
       \mathcal L=L_{2\infty}^{\boxtimes14hab}.
\tag{HC23}
\] Tensor the \textbf{carried} single-J DP59 basis, using precisely (HC22). It gives an integral basis of \[
       \mathcal W_M=[\alpha]^*H^0(X,\mathcal L^M)
                         \subset H^0(X,\mathcal L^{M^2}),
       \qquad \operatorname{rank}\mathcal W_M=M^{42hab}.
\tag{HC24}
\] Its parameters are the prescribed representatives of \(X[\alpha^2]/X[\alpha]\), with actual rescaled translation \[
                         x=Mt'\in X[\alpha].
\tag{HC25}
\] Indeed \(\alpha Mt'=\alpha^2\bar\alpha t'=0\), and coprimality makes \(\bar\alpha\) invertible on α-primary torsion. This remains the permitted half-orbit, not the full level-\(M^2\) theta group.

At11, α is étale and all new CM-coordinate changes under consideration are integral invertible by K030 or by the corresponding later host proof. Consequently the K022 ordinary-jet surjectivity, equality of all determinantal image ideals, and one-translate norm-one conclusion apply with \[
               M\ge3D+2,
       \qquad e(D)=\max\{1,\lceil\log_{29}(3D+2)\rceil\}.
\tag{HC26}
\] The threshold is independent of h,a,b. Each individually selected finite Heisenberg translate has finite Schwartz norm one; its coefficient vector is one standard basis vector. No new coefficient-combination height or norm is introduced by the real extension. The Fock norm, source representation dimension, class mass, actual period and finite29 coordinate changes are separate and remain in the cost.

The auxiliary multijet statement of DO030 also persists with its same threshold \(3s(D+1)-1\), provided its auxiliary torsion/order and Schwartz support remain prime to29 and11. Its proof uses the same half-orbit, integral selfdual trace lattice, and the isogeny's invertibility on that auxiliary torsion. Rational29 now has the new local ramification depth, so its explicit class-index factor is PH26 of K030, not the old degree-six formula. None of this identifies the DP59 translate with the separate literal source raw-Xi input.

\subsection{5. Auxiliary T-prime existence and size scope}\label{auxiliary-t-prime-existence-and-size-scope}

For any finite collection of words and any specified finite lower size threshold, the norm/congruence part of the T-prime interface can be supplied in the same external quadratic subfield. Let \[
                      L_T=812(g_0g_1)^2.
\tag{HC27}
\] Take the ray class field of \(\mathbb Q(\sqrt{-7})\) for the ideal modulus \(L_T\mathcal O_{K_q}\), and its normal closure over \(\mathbb Q\). There are infinitely many rational primes splitting completely in this finite normal extension, by the unconditional Chebotarev theorem. At such a prime q, a chosen prime ideal of \(K_q\) lies in the trivial ray class, hence has a generator \(\pi_T\equiv1\pmod{L_T}\). Since the prime has degree one and \(K_q\) is imaginary quadratic, \(\pi_T\bar\pi_T=q\). Its conjugate ideals are distinct and coprime. Moreover \(q\equiv1\pmod{812}\), so it splits completely in K and supplies the required42-element half with residue norm q. Choosing distinct such q outside a finite set gives the requested finite family.

For any fixed finite weight data one can choose a real size parameter a with \(2fN_m+1<4a\) and every source spread at most \(2a\), then choose these T-primes larger than \(4a\). This proves qualitative nonemptiness of the corresponding norm, congruence and lower-size clauses. It supplies no effective upper bound for T and no assertion that this choice meets the quantitative MainX window. The same literal \nolinkurl{piT_Kq} obstruction remains as for G, since the source's common field is not this external quadratic subfield.

\subsection{6. Literal source failures and exact scope}\label{literal-source-failures-and-exact-scope}

The following distinctions are necessary.

\begin{longtable}[]{@{}
  >{\raggedright\arraybackslash}p{(\columnwidth - 2\tabcolsep) * \real{0.5000}}
  >{\raggedright\arraybackslash}p{(\columnwidth - 2\tabcolsep) * \real{0.5000}}@{}}
\toprule\noalign{}
\begin{minipage}[b]{\linewidth}\raggedright
Interface
\end{minipage} & \begin{minipage}[b]{\linewidth}\raggedright
Result on this host
\end{minipage} \\
\midrule\noalign{}
\endhead
\bottomrule\noalign{}
\endlastfoot
Actual principal CM variety and full integral action & K030 gives \(J^{14}\), degree84, induced type, exact coefficient and lattice \\
Literal \nolinkurl{IsEtaT} and normalized \nolinkurl{etaVal} & K029 proves all character clauses and a finite common coefficient field \\
Numerical \nolinkurl{Ranks}, with f tied to degree & (HC1)--(HC4) supply \(f=k=42,m=2,t=1759\) \\
Numerical parity, interlace, centred and NCGaps inequalities & (HC6)--(HC10) verify them; this alone is not a \nolinkurl{ThetaX} instance \\
Exact uniform even stage-character conditions & (HC16)--(HC21) prove them for every integer uniform weight \\
G norms, coprimality, residue cardinalities, full PiGCong and lower bound & (HC12)--(HC14) give explicit values in an external quadratic subfield \\
DP59 half-orbit and integral 11-jets & (HC22)--(HC26) retain the actual allowed family with parameterized dimensions \\
\nolinkurl{eight_dvd : 8 ∣ f} & Fails for \(f=42\); a further real extension of degree divisible by4 is necessary numerically \\
\nolinkurl{Fields.Ok.quad}, \nolinkurl{Primes.Ok.piG_Kq}, \nolinkurl{piT_Kq} & Fail for the unchanged primitive type: its defined common field is \(\mathbb Q\) \\
Full D7, prescribed raw Xi, later selected image and quantitative cost window & Not supplied by these constructions \\
\end{longtable}

For the common-field assertion, the induced type contains all B-lifts of each E₀ embedding, so an element on which every selected embedding agrees must first lie in E₀. The selected E₀ embeddings contain the identity, 3 and17. The subgroup generated by3 is \(\{1,3,9,27,25,19\}\), and17 is outside it; together they generate all12 automorphisms of E₀. The common element is consequently rational. Thus \[
                  \mathrm{Fields.Kq}=\mathbb Q.
\tag{HC28}
\] The source defines this field in \nolinkurl{AbelianHostFieldConditions.lean:60–76}, requires degree2 at lines108 and157, and requires the prime generators to belong to it at lines210 and226. This is a precise restatement of the prior primitive-type obstruction. It cannot be repaired by relabeling the external \(\mathbb Q(\sqrt{-7})\) as \nolinkurl{Fields.Kq}.

For an induced type under any further totally real extension, the same common-field obstruction persists. Agreement across all lifts first forces descent to E₀, and the preceding group argument still gives \(\mathbb Q\). In contrast, the numerical condition \(8\mid42h\) is equivalent to \(4\mid h\). At its first possible real extension degree4, f becomes168 and the rank constraint becomes \(k\ge168\), with G-size threshold at least \(2\cdot168^2=56\,448\). A construction of that larger polarized host must be supplied separately; the formal divisibility check alone would not construct it.

\subsection{7. Sources and verification scope}\label{sources-and-verification-scope-1}

\begin{itemize}
\tightlist
\item
  Exact source: \nolinkurl{Formalism/LocalRingCardinality/Params.lean:6–22,28–40}; \nolinkurl{Formalism/LocalToGlobal/Columns/AbelianHostFieldConditions.lean:60–76,108–125,206–238,293–331,486–494,560–579}; \nolinkurl{Formalism/LocalToGlobal/UnitaryNorms/RealCyclotomicTwoPower.lean:54}; \nolinkurl{Formalism/LocalRingCardinality/DataTheta.lean:501–555}; \nolinkurl{Formalism/SharpRestriction/CuspSpanNormBounds/CayleySignature.lean:106–125}.
\item
  Inherited Frame: \nolinkurl{Formalism/SharpRestriction/RamificationFrobenius/TotientGaloisClasses/ZsmulGeneratorEqZeroOfProjection.lean:16–24} and \nolinkurl{Formalism/FrobeniusOrbits/Statement.lean:27–29}. Additional recovered members are pinned to commit \nolinkurl{7cbf7a0c32fca1badf05931879a37c3f8a2b03b4}; their retrieval manifests explicitly record when no intake-author hash exists. No original-byte match is inferred from that absence.
\item
  AM-R-009 and ALP-LIFT-18(G2)--(G14): the checked E₀ norm-one torus and its actual character construction, used with zero infinity weight. AM-K029/030: the repaired odd character, full product host and exact finite comparisons. AM-DO-025: selected partition recurrence. Carried DP59 and AM-K022/DO030: the permitted basis and ordinary/multijet interfaces.
\item
  Milne, J. S. (2020). \emph{Class field theory} (Version4.03). Chapter V, Theorem3.6, printed p.158, and Theorem3.23, printed p.164, for ray class fields and unconditional Chebotarev. Primary author PDF: \url{https://www.jmilne.org/math/CourseNotes/CFT.pdf}. These are used only for the qualitative T-prime choice, not an effective size estimate.
\end{itemize}

The finite certificate checks every numerical rank clause, all displayed gap comparisons including every actual opposite-sign pair, the prime norms/congruences/primality and primitive-root orders, the common-field finite group computation, and the alpha norm/coprimality arithmetic. It does not infer analytic nonvanishing, a selected-image bound, a large-field class number, a character theorem, or a Lean instance from those finite checks.

The independently replayable files are \nolinkurl{enlarged_host_ranks_characters_alpha_certificate.py} and its JSON output. The certificate uses integer/Fraction arithmetic and direct trial division, with no probable-prime substitution. In particular it checks the literal natural subtraction in the Dpar sign rule and reconstructs the selected partition from the source weight recurrence before comparing it with (HC9)--(HC10).

\clearpage

\section{AM-DO-034: a new selected first lift on the degree-seven principal host}\label{am-do-034-a-new-selected-first-lift-on-the-degree-seven-principal-host}

Author: GPT-6 Astra Pro, ALP\_MACHINE \nolinkurl{/root/definite_oscillator}, 21 November 2026 UTC. Inherited reasoning effort is not separately exposed.

This extends AM-DO-030/033 to the actual principally polarized \nolinkurl{J^14} host of AM-K030. The selected polynomial is constructed afresh in the forty-two CM matrix blocks. The proof does not transport the old six-block upper witness through a change of individual CM weights. It retains the standard-lattice rescaling, the ramified rational-29 level, the finite Steinberg norm, and the exact class-supported extraction.

The inputs are the proved geometry and integral action of AM-K030, the finite splitting proof of AM-DO-022, the local selected-type calculation of AM-DO-025, and the actual DP59 half-orbit theorem in the checked form AM-K022. The construction is conditional only on the displayed stage character and rank/weight data being provided. AM-K031 is supplying concrete such even characters and auxiliary primes. None of this asserts the unchanged full \nolinkurl{Dat.Ok}, the original \nolinkurl{Fields.Kq} wrapper, or every numerical source-rank condition on this new host.

\subsection{1. Actual host and two coordinate conventions}\label{actual-host-and-two-coordinate-conventions}

Keep the notation of AM-K030:

\[
 E_0=\mathbf Q(\zeta_{28}),\qquad
 B\subset\mathbf Q(\zeta_{29}),\quad[B:\mathbf Q]=7,
 \qquad K=E_0B,\quad F=K^+,\quad [F:\mathbf Q]=42.
 \tag{DO34.1}
\]

Let \nolinkurl{Psi} be the induced CM type, with forty-two members. Let \nolinkurl{pi_B} be K030's totally positive generator of the prime of B above29 and \nolinkurl{e_0,...,e_6} its explicit orthonormal trace basis of \nolinkurl{pi_B^(-3) O_B}. The actual principally polarized variety is

\[
 \mathcal A=(J^2)^7=J^{14},\qquad
 \Lambda=\pi_B^{-3}\mathcal O_K,
 \quad\xi_0=\frac{\zeta_7-\zeta_7^{-1}}{14}.
 \tag{DO34.2}
\]

In the standard \nolinkurl{O_K} lattice convention the same polarization has coefficient

\[
 \boxed{\xi_K=\xi_0\pi_B^{-6}.}
 \tag{DO34.3}
\]

This is an equality of the actual integral trace form, not only of its discriminant. It identifies the standard line with the product trace-one line. Its reduced seed in ordinary product coordinates is \nolinkurl{c^14 product sigma_q}, as in K030(PH22).

Index the CM embeddings by \nolinkurl{tau=(i,epsilon,l)}, with \nolinkurl{i=1,2,3}, \nolinkurl{epsilon=+,-}, and \nolinkurl{l=0,...,6}. For every matrix entry of a Hom space, write the fourteen ordinary-J coordinates as \nolinkurl{(u_(i,j),v_(i,j))}, and put

\[
 w_{i,\pm,j}=u_{i,j}\pm i v_{i,j},\qquad
 z_{i,\pm,l}=\sum_jU_{lj}w_{i,\pm,j},
 \quad U_{lj}=\tau_l(e_j),\quad U^tU=1.
 \tag{DO34.4}
\]

These are K030's product-eigen coordinates on the lattice \nolinkurl{Lambda}. To use the standard \nolinkurl{O_K} lattice, define instead

\[
       \zeta_\tau=\tau(\pi_B)^3z_\tau,
       \qquad
       \partial_{\zeta_\tau}
                 =\tau(\pi_B)^{-3}\partial_{z_\tau}.
 \tag{DO34.5}
\]

Indeed the map from the standard lattice to \nolinkurl{Lambda} is multiplication by \nolinkurl{pi_B^(-3)}, so its product-eigen coordinate is \nolinkurl{z_tau=tau(pi_B)^(-3) zeta_tau}. This proves (DO34.5) and explains why (DO34.3) and (DO34.5) must be used together. It is not legitimate to keep both the old coefficient \nolinkurl{xi_0} and the new standard lattice.

All the changes in (DO34.4)--(DO34.5) and their inverses are integral at11. In particular the standard \nolinkurl{zeta} derivations give an integral invariant tangent basis after a coefficient extension. The \nolinkurl{U} matrix is real orthogonal, so it introduces no real operator-norm loss. The old \nolinkurl{plus/minus i} change and the positive diagonal factors in (DO34.5) still have to be included in norms and other local valuations.

\subsection{2. A new highest polynomial in the actual forty-two blocks}\label{a-new-highest-polynomial-in-the-actual-forty-two-blocks}

Fix even ranks \nolinkurl{a=k}, \nolinkurl{b=k+t+1}, with \nolinkurl{b>a}, and actual source weight data of the form used in AM-DO-025 at all \nolinkurl{tau in Psi}. Let \nolinkurl{e_tau} be the exponent parameter and let \nolinkurl{w_(0,tau)} be the actual dominant initial coefficient weight. Define

\[
 \nu_{\tau,i}=w_{0,\tau,i}+e_\tau-b/2,
 \quad 1\le i\le a,
 \qquad d_\tau=|\nu_\tau|,
 \qquad D=\sum_{\tau\in\Psi}d_\tau.
 \tag{DO34.6}
\]

AM-DO-025 proves that these are nonnegative partitions under the displayed source HC gap conditions. That calculation is local at one real place and does not depend on the number of CM embeddings. It therefore applies to each of the forty-two places here.

For \nolinkurl{Z_tau} the \nolinkurl{a by b} matrix of formal standard-coordinate creation variables, set \nolinkurl{nu_(tau,a+1)=0} and define

\[
 P_\nu(Z)=
 \prod_{\tau\in\Psi}\prod_{r=1}^{a}
 \det\bigl(Z_{\tau,[1,r],[1,r]}\bigr)
                   ^{\nu_{\tau,r}-\nu_{\tau,r+1}}.
 \tag{DO34.7}
\]

This is a new polynomial, defined directly in the new CM eigenspaces. It has total degree D and separate degree \nolinkurl{d_tau} in each block. Its diagonal monomial is

\[
             \prod_{\tau,i}Z_{\tau,i,i}^{\nu_{\tau,i}}
 \tag{DO34.8}
\]

with coefficient exactly one. Only the diagonal term of every minor can produce a monomial with no off-diagonal variable, and summing its exponents telescopes to \nolinkurl{nu_(tau,i)}. Thus \nolinkurl{P_nu} is primitive in the standard CM polynomial coefficient lattice, in particular at11.

For the common splitting character with infinity component \nolinkurl{mu_tau(z)=(z/bar z)^(-e_tau)}, the actual compact Fock module at \nolinkurl{tau} is

\[
 \det_a^{e_\tau-b/2}\boxtimes
 \det_b^{-e_\tau+a/2}\otimes
       \operatorname{Sym}(\mathbf C^b\otimes(\mathbf C^a)^*).
 \tag{DO34.9}
\]

The bideterminant (DO34.7) belongs already to its Cauchy summand indexed by \nolinkurl{nu_tau}. Its source type is exactly \nolinkurl{rho_(0,tau)^vee} and its target type is the source's \nolinkurl{rho_(1,tau)}, by DO25.5--9. Taking the tensor product over \nolinkurl{Psi} gives source type \nolinkurl{sigma=rho_0^vee} and target type \nolinkurl{rho_1}. No Schur projector, averaging denominator, or claim that the transported old polynomial still has one CM weight is needed. This resolves the first-edge weight-projection gate for a newly constructed selected polynomial.

The standard-coordinate polynomial and the analogous product-eigen polynomial satisfy the exact differential identity

\[
 P_\nu(\partial_\zeta)
       =c_\nu P_\nu(\partial_z),\qquad
 c_\nu=\prod_{\tau\in\Psi}
                      \tau(\pi_B)^{-3d_\tau}.
 \tag{DO34.10}
\]

Separate homogeneity proves this identity. In particular it is an11-adic unit scalar, but it is not discarded at29 or at infinity. With the valuation normalized by \nolinkurl{v_29(29)=1},

\[
            v_{29}(c_\nu)=-3D/7.
 \tag{DO34.11}
\]

Each entry of \nolinkurl{U} also has valuation \nolinkurl{-3/7} at29. Expressing the full standard derivative polynomial in ordinary product derivatives therefore gives the coefficient bound \nolinkurl{-6D/7} at29. At2 it retains the old upper denominator bound \nolinkurl{2^D}. These are upper bounds on denominators, not assertions that every coefficient attains them. At11 the coordinate changes are invertible and none of these denominators appears.

\subsection{3. The actual allowed half-orbit survives the field repair}\label{the-actual-allowed-half-orbit-survives-the-field-repair}

On the edge use the actual variety and line

\[
 X=\mathcal A^{ab}=J^{14ab},\qquad
 g=\dim X=42ab,\qquad
 \mathcal L=L_{2\infty}^{\boxtimes14ab}.
 \tag{DO34.12}
\]

Choose

\[
 \alpha=(1+2\sqrt{-7})^e\in E_0\subset K,
 \quad M=29^e,\quad e\ge1,
 \qquad\alpha\bar\alpha=M.
 \tag{DO34.13}
\]

Multiplication by alpha commutes with the B-action and acts diagonally on the fourteen ordinary factors. Its actual Rosati adjoint is bar alpha. Hence its polarization norm is exactly M on the new product. The ideals \nolinkurl{(alpha)} and \nolinkurl{(bar alpha)} remain coprime after extension to K: they lie on the two different relative split halves at29. At a selected prime of K above29 their nonzero valuations are \nolinkurl{7e}, since the absolute ramification degree there is seven.

Take the exact DP59 tensor product basis on the fourteen ordinary factors, and then on all \nolinkurl{ab} entries. It is an integral basis of

\[
 \mathcal W_M=[\alpha]^*H^0(X,\mathcal L^M)
         \subset H^0(X,\mathcal L^{M^2}),
 \qquad \operatorname{rank}\mathcal W_M=M^{42ab}.
 \tag{DO34.14}
\]

The prescribed quotient representatives give parameters

\[
           t'_\lambda\in X[\alpha^2],\qquad
           x_\lambda=M t'_\lambda\in X[\alpha]\subset X[M].
 \tag{DO34.15}
\]

They retain their actual odd-order symmetric lift scalars. They are exactly the permitted product orbit, not the whole theta group of \nolinkurl{L^(M^2)}. All statements follow by the actual diagonal product identification, so they do not require a new conjectural CM theta-span theorem. At11, \nolinkurl{[alpha]} is étale and is a formal isomorphism because its degree is a power of29. The integral principal line and section are the ones in K030, so the K018 reduced/Chern gauge applies to this ordinary product.

The actual finite vacuum is the transport of the product vacuum through the exact symplectic lattice isometry in K030(PH17--20). Its finite norm is one. At every nonsplit odd place, K/F is unramified and the opposite even identity Hermitian lattice is integrally hyperbolic. The proof in DO022 §3 therefore applies with the new local field and the coefficient \nolinkurl{xi_K}, which generates the actual absolute codifferent. Its Pan conductor condition uses that exact valuation and does not require absolute unramifiedness. All dyadic places are relatively split, because splitting of E0/F0 at2 persists under base change. DO020 covers them. Consequently the standard splitting and the two normalized moved-lattice averages agree, with their actual determinant-root characters, exactly as in DO22.8. This is a transport of the proved local hypotheses and proof, not an inference from equality of vacuum lines alone.

The same Heisenberg reading therefore has the fixed product period \nolinkurl{C^(14ab)}, where C is the already fixed ordinary-J period of DO022/R006. All finite Weyl and symmetric-lift scalars remain in the reading. The equality holds as a reduced germ and hence after the actual polynomial \nolinkurl{P_nu(partial_zeta)}.

\subsection{4. Simultaneous jets and one selected nonzero coefficient}\label{simultaneous-jets-and-one-selected-nonzero-coefficient}

Fix an auxiliary finite vector

\[
 \phi_{\rm aux}=\mathcal O\phi_0\ne0,
 \qquad\mathcal O=\sum_{r=1}^{s}b_rW_{y_r},
 \tag{DO34.16}
\]

supported at odd rational primes away from11 and29. Reduce the parameters modulo the actual selfdual lattice, combine their exact phases, and retain distinct classes with nonzero coefficients as in DO030. Let a fixed odd integer \nolinkurl{d} prime to11·29 kill these classes. Then they give s distinct points of \nolinkurl{X[d]}.

Every step of DO030 is valid on (DO34.12), for the following specific reasons. First, the new lattice is perfect by K030, so the same rational representative \nolinkurl{r_j=M z_j}, where \nolinkurl{M z_j=t_j}, has the fixed phase on the whole orbit. This uses the rational inclusions \nolinkurl{delta_(j,29) in M Lambda_29} and \nolinkurl{x_lambda in M^(-1) Lambda_29}; it does not substitute29 for a uniformizer. Second, the prime-to11 torsion points remain distinct on the good model. Third, the translated-line comparison is still made by the actual \nolinkurl{[d]} pullback and its integral symmetric theta lift, so its fibre scalar is an11-unit. No translated theta-null is divided out. Finally, the integral simultaneous-jet proof depends on \nolinkurl{L^3} being very ample and on the actual product line and base change, not on the number of factors.

Thus the same explicit threshold gives the surjection

\[
 M\ge3s(D+1)-1
 \quad\Longrightarrow\quad
 \mathcal W_M\twoheadrightarrow
       \bigoplus_{j=1}^{s}
           \operatorname{Jet}^D_{z_j}(\mathcal L^{M^2}).
 \tag{DO34.17}
\]

For completeness, sections of \nolinkurl{L^3} that vanish at all but one of the s specializations supply the factor of degree \nolinkurl{3(s-1)(D+1)}. At the selected point choose g formal parameters as ratios of sections of \nolinkurl{L^3}, and homogenize an arbitrary D-jet with degree \nolinkurl{3D}. The remaining exponent is at least2 and supplies a section nonzero at that point. The construction uses only unit inverses, so it is integral. Pullback by the actual alpha isogeny gives (DO34.17).

Take the standard CM tangent basis in (DO34.5). Its diagonal monomial (DO34.8) has coefficient one, and an ordinary jet with that coefficient evaluates under \nolinkurl{P_nu} to the factorial

\[
       \prod_{\tau,i}\nu_{\tau,i}!,\qquad
       B_\nu=\sum_{\tau,i}v_{11}(\nu_{\tau,i}!).
 \tag{DO34.18}
\]

Prescribe this degree-D monomial at an auxiliary point whose \nolinkurl{b_r} has minimum valuation, and prescribe zero D-jets at all other auxiliary points. Multiply the monomial by the unit \nolinkurl{M^D} to retain the source's normalized level reading. There are no lower terms. Therefore the possibly nonintegral higher coefficients of the Chern gauge cannot affect the degree-D derivative. The resulting functional has valuation \nolinkurl{min_r v_11(b_r)+B_nu}. Expanding the interpolating section in the integral actual basis (DO34.14) implies that at least one individual basis section has valuation no larger.

Consequently one permitted lambda satisfies

\[
 a_\lambda\ne0,\qquad
 v_{11}(a_\lambda)\le\min_rv_{11}(b_r)+B_\nu,
 \qquad
 \vartheta_{r_{P_\nu}\mathcal O\phi_\lambda}(1,1)
                   =C^{14ab}a_\lambda.
 \tag{DO34.19}
\]

This proves selected first nonvanishing in the forty-two CM blocks. It does not assert the old fixed-row sharp slope, which concerned a different explicitly selected directional witness. The present factorial upper bound is its own fully proved statement.

\subsection{5. Finite Steinberg input on the new host}\label{finite-steinberg-input-on-the-new-host}

Suppose \nolinkurl{g_0,g_1} are distinct odd rational primes, different from11 and29, that split completely in K. Fix the forty-two selected relative split places above each. Let the common character mu be unramified away from these G places and29, tame at G, and trivial on the principal rational29 level used below. Its infinity exponent must be the one in (DO34.6). The concrete character supplied in K031 is unramified at29 and meets this last condition.

At each G place use exactly the determinant-corrected free residue-matrix construction of DO033. Its genuine source action is \nolinkurl{eta(det g)^(-1)|det g|^(b/2) f(Xg)}. It contains the pure finite Steinberg type, has raw value one at \nolinkurl{[I_a;0]}, and has squared norm

\[
          \gamma_q=q^{-ab}\frac{|\mathrm{GL}_a(\mathbf F_q)|}
                                         {q^{a(a-1)/2}}.
 \tag{DO34.20}
\]

Its Fourier expansion is an actual auxiliary Heisenberg combination with at most \nolinkurl{q^(ab)} terms, minimum coefficient valuation zero at11, and squared coefficient norm gamma\_q. These facts are the local theorem of DO033 and require no new field extension argument once the local fields are actually degree one.

Tensoring over the eighty-four selected G places gives

\[
 \begin{split}
 d_G&=(g_0g_1)^{21a(a-1)},\\
 \Gamma_G&=(g_0g_1)^{42(-ab+a(a+1)/2)}
       \prod_{\epsilon=0}^{1}\prod_{i=1}^{a}
                                     (1-g_\epsilon^{-i})^{42},\\
 s_G&\le(g_0g_1)^{42ab},\qquad
 \min_rv_{11}(b_r)=0,\qquad
 \|\mathcal O_G\phi_\lambda\|_f^2=\Gamma_G.
 \end{split}
 \tag{DO34.21}
\]

Apply (DO34.19) with this actual vector and choose

\[
 e=\max\{1,\lceil\log_{29}(3s_G(D+1)-1)\rceil\},
                  \qquad M=29^e.
 \tag{DO34.22}
\]

The result is one permitted half-translate, still carrying the selected real source type and the pure finite Steinberg source type, with nonzero coefficient and upper valuation \nolinkurl{B_nu}. Its finite norm is exactly as in (DO34.21); no norm of an interpolating linear combination is being used.

\subsection{6. The new level, rational stabilizer, and exact extraction}\label{the-new-level-rational-stabilizer-and-exact-extraction}

At29 use principal congruence modulo the rational ideal \nolinkurl{(M)=(29^e)} at all six real-field primes over29. At each G place keep the full integral compact, and let the source function vary there by its finite matrix coefficient. Keep the full integral compact elsewhere. Denote the resulting product by \nolinkurl{K*}.

The principal29 subgroup fixes every permitted finite translate exactly. Its parameter lies in the selected relative split isotropic position space, with denominator dividing M. Acting by \nolinkurl{g,h=1 mod M} displaces it by an integral vector in that same space. The symplectic pairing and the odd-place trace-one phase on that displacement are zero. This proves the exact Weyl fixation, even though the valuation of M at that local prime is \nolinkurl{7e}. The compact character is trivial there by hypothesis.

Normalize real Haar measure to volume one, and give the standard full finite compact volume one. K030(PH26) gives

\[
 I_a^{K}(M)=[K_{a,0}:K^*]
       =M^{42a^2}\prod_{i=1}^{a}(1-29^{-i})^6,
                  \qquad\operatorname{vol}K^*=(I_a^K(M))^{-1}.
 \tag{DO34.23}
\]

The notation here records the index caused by the29 refinement; the G compact factors have not been shrunk. There are six residue fields, each equal to \nolinkurl{F_29}, and each congruence length is \nolinkurl{7e}. Replacing the exponent6 on the residue product by42 would be incorrect.

The rational stabilizer \nolinkurl{G(F) intersect K*} is trivial. It is finite because it preserves the global integral lattice and is compact at all real places. At a place over29 it lies in \nolinkurl{1+29 M_a(O_(F_v))}. If a matrix B has minimum uniformizer valuation \nolinkurl{s>=7}, the unique lowest term of \nolinkurl{(1+B)^29-1} is \nolinkurl{29 B}, of valuation \nolinkurl{s+7}: all middle terms have valuations at least \nolinkurl{7+2s}, and \nolinkurl{B^29} has valuation \nolinkurl{29s>s+7}. Raising to an integer prime to29 preserves the minimum valuation. Thus this congruence subgroup is torsion free. The rational matrix embedding is injective, proving the assertion.

Let \nolinkurl{u_infty} be the unit source highest vector and \nolinkurl{u_G} the product of DO033's unit generic Steinberg vectors. Set \nolinkurl{d_rho=dim rho_0}. On the identity finite \nolinkurl{K*} class define

\[
 F(k_\infty,k_f)=
 \overline{\langle u_\infty,\sigma(k_\infty)u_\infty\rangle
             \langle u_G,\tau_G(k_G)u_G\rangle},
 \tag{DO34.24}
\]

and set it to zero on all other classes. Triviality of the rational stabilizer makes this a well-defined automorphic function. It is cuspidal because the rational unitary group is definite. Its real type is \nolinkurl{rho_0}, and its G compact type is the contragredient finite Steinberg type, with the literal initial inverse generic character.

On that class the actual kernel is \nolinkurl{L(sigma(k_infty)u_infty tensor tau_G(k_G)u_G)} with value \nolinkurl{C^(14ab) a_lambda} at the original tensor. Applying compact Schur orthogonality to both unit vectors returns the original tensor divided by \nolinkurl{d_rho d_G}. The class volume is the one in (DO34.23). Therefore

\[
 \boxed{
 \Theta^{\rm bil}_{r_{P_\nu}\mathcal O_G\phi_\lambda}(F)(1)
       =\frac{C^{14ab}a_\lambda}{I_a^K(M)d_\rho d_G}\ne0,
 \quad
 \|F\|_2^2=\frac1{I_a^K(M)d_\rho d_G},\quad
 \|F\|_\infty=1.}
 \tag{DO34.25}
\]

This is an actual first automorphic lift on the new host. It has an actual irreducible source constituent containing the finite Steinberg type. Finite-dimensional compact projection to a nonzero constituent is injective on that irreducible compact copy. Its Borel-fixed line therefore gives an Iwahori-fixed vector in the constituent. This supplies the initial hypothesis of the local theorem ALP-LIFT-19 on this host, whenever its stage characters and global rank-gap hypotheses are also supplied.

No class-number-one statement about K or F is needed: the input is supported on a single actual class, whose mass was computed using its trivial stabilizer.

\subsection{7. Norms and costs that remain after the repair}\label{norms-and-costs-that-remain-after-the-repair}

Let \nolinkurl{h(nu)} be the hook product of a partition. In the standard-coordinate convention (DO34.5), the exact real Schwartz norm is

\[
 \|r_{P_\nu}\phi_{0,\infty}\|^2
 =\det(2Y)^{-7ab}\,
   |c_\nu|^2
   \prod_{\tau\in\Psi}
        \left[ h(\nu_\tau)
          \left(\frac{\pi}{2G_{i(\tau)}}\right)^{d_\tau}\right].
 \tag{DO34.26}
\]

The vacuum exponent is \nolinkurl{14ab/2=7ab}. The Fischer norm of each primitive highest bideterminant is its hook product, by AM-R-011. Orthogonality of U and the metric in K030(PH25) give the creation scale \nolinkurl{pi/(2G_i)}. Equation (DO34.10) supplies the exact remaining scalar. This derives (DO34.26) in the actual product measure; it does not normalize the real Gaussian to one.

Unit-normalizing the finite input and the source input changes (DO34.25) to

\[
 \frac{C^{14ab}a_\lambda}
        {\sqrt{\Gamma_G}\sqrt{I_a^K(M)d_\rho d_G}}.
 \tag{DO34.27}
\]

If the real Schwartz vector is also made unit, divide this value by the square root of (DO34.26). Thus the standard-lattice scalar, all Gaussian/hook factors, and the auxiliary finite normalization remain visible.

At11 one has

\[
 \begin{split}
 v_{11}(d_G)&=0,\\
 v_{11}(I_a^K(M))&=6\sum_{i=1}^{a}v_{11}(29^i-1),\\
 v_{11}(\Gamma_G)&=42\sum_{\epsilon=0}^{1}
                         \sum_{i=1}^{a}v_{11}(g_\epsilon^i-1),\\
 v_{11}\!\left(C^{-14ab}\Theta^{\rm bil}(F)(1)\right)
     &\le B_\nu-v_{11}(I_a^K(M)d_\rho).
 \end{split}
 \tag{DO34.28}
\]

For the candidate primes \nolinkurl{4,615,409} and \nolinkurl{5,276,377} supplied by K031 and rank \nolinkurl{a=42}, all three integers \nolinkurl{29,g_0,g_1} have order10 modulo11, and \nolinkurl{v_11(q^10-1)=1}. Thus the sums over \nolinkurl{1<=i<=42} equal4, giving the exact fixed factors

\[
          v_{11}(I_{42}^{K}(M))=24,
          \qquad v_{11}(\Gamma_G)=336.
 \tag{DO34.29}
\]

This last calculation is independent of e. Unit finite normalization has scalar valuation cost168 in this example. It does not invalidate the raw minimum coefficient valuation zero used in (DO34.19).

For fixed ranks, fixed G primes, and variable selected weights, (DO34.22) has \nolinkurl{log M=O(log(D+2))}; the class index and compact dimension consequently give an \nolinkurl{O(log(D+2))} source extraction cost. The coordinate factor in (DO34.10), its29 denominators, and the real hook/Fock norm have their own degree dependence and are not included in that logarithmic assertion. In particular the possible \nolinkurl{6D/7} ordinary-coordinate 29 denominator is not a constant-level cost.

\subsection{8. Scope and exact inputs}\label{scope-and-exact-inputs}

This proves an actual selected first kernel, with a single permitted29 half-translate, a specified finite Steinberg auxiliary vector, and an exact class-supported source input on the new principal host. It closes the first CM-weight and actual-half-orbit transfer gap. It does not identify this auxiliary vector with the literal raw \nolinkurl{PhiG} or \nolinkurl{Xi}, prove the terminal restricted Gram nonvanishing, or compute the integral lattice of a later iterated theta image. It also does not turn the enlarged field into the unchanged original host record when unrelated field/rank wrappers fail.

The local selected Schur identity and splitting proof are used as settled earlier findings. The novelty here is their exact simultaneous realization on K030's standard principal lattice and allowed orbit, with the forty-two-block polynomial and the changed costs.

Exact locators:

\begin{itemize}
\tightlist
\item
  AM-K030(PH3--9, PH17--26): explicit trace basis, actual integral product polarization, standard coefficient, K-eigen coordinates, and rational29 congruence index.
\item
  AM-DO-022 §§3--4: standard finite splitting comparison, Pan conductor and compact phase, arbitrary odd local base field, determinant roots, and split dyadic case.
\item
  AM-DO-025(DO25.3--9): literal common-character compact Cauchy summand and source/target coefficient match; source \nolinkurl{Columns/BaseLemmas.lean} lines496--552.
\item
  AM-K022(J3--5, J6--16, J19--27): actual DP59 orbit, integral jet threshold, factorial detection and source-normalized differential reading.
\item
  AM-DO-030 §§2--4: fixed auxiliary phases, translated-line unit comparison and simultaneous integral jets.
\item
  AM-DO-033 §§1--5: primitive finite Steinberg vector, actual determinant-corrected residue model, Fourier norm and compact source extraction.
\item
  AM-R-011: exact Fischer coefficient formula, highest bideterminant hook norm, and creation/Gaussian scaling.
\item
  Primary geometric facts are the already read Milne, \emph{Abelian Varieties}, I§6, p.30, and Mumford(1966), \emph{On the equations defining abelian varieties I}, §1(II), p.289, as documented in K022/DO030.
\end{itemize}

\clearpage

\section{AM-R-017: explicit auxiliary T-primes and a first source class at maximal 29-level}\label{am-r-017-explicit-auxiliary-t-primes-and-a-first-source-class-at-maximal-29-level}

Author: /root, GPT-6 Astra Pro, ALP\_MACHINE, 22 November2026UTC. Effort is not externally exposed.

This result makes two previously qualitative choices concrete. It supplies one explicitly certified auxiliary T-prime for each of the K031 and K032 hosts. It also constructs a nonzero first automorphic input with the G generic type and nontrivial T dilation characters, without imposing a principal level at29 to remove rational stabilizers. Its exact class mass is computed. The theta lift of this input with the complete prescribed raw kernel is not proved nonzero here: averaging a translated29 kernel over the maximal29 compact remains a separate issue. The unchanged source \nolinkurl{Fields.Kq} obstruction also remains.

\subsection{1. Two explicit norm primes with the full source congruence}\label{two-explicit-norm-primes-with-the-full-source-congruence}

Use the already checked fields of K031 and K032, and write \(\delta=\sqrt{-7}\) in their common external quadratic subfield. The following data use the exact G-primes of those findings:

\begin{longtable}[]{@{}lrrrrrr@{}}
\toprule\noalign{}
Product host & f & k & Rational conductor c & g0 & g1 & n \\
\midrule\noalign{}
\endhead
\bottomrule\noalign{}
\endlastfoot
J14 & 42 & 42 & 812 & 4615409 & 5276377 & 3 \\
J56 & 168 & 168 & 52780 & 22285932761 & 64072175441 & 16 \\
\end{longtable}

In each row set \[
 L=c(g_0g_1)^2,\qquad A=1+nL,\qquad
 \pi_T=A+L\delta,\qquad T=A^2+7L^2.
 \tag{R17.1}
\] Thus every number in the construction is specified without a prime-search oracle. The exact prime integers, also recorded in the certificate, are:

\begin{itemize}
\tightlist
\item
  J14: \nolinkurl{3710360310942963702130514092646914581678165356491757211050633}.
\item
  J56: \nolinkurl{3045759889845354427724717877501093821489629367456155582288640401018208099619883018305672886530161}.
\end{itemize}

Both T are prime. Here is the complete certificate method, including the factored divisors required to make the primality conclusion deterministic. For J14 use \(F=586L\), and for J56 use \(F=44L\). Their factorizations are \[
\begin{aligned}
 F_{14}&=2^3\cdot7\cdot29\cdot293\cdot g_0^2g_1^2,\\
 F_{56}&=2^4\cdot5\cdot7\cdot11\cdot13\cdot29\cdot g_0^2g_1^2.
\end{aligned}
\tag{R17.2}
\] Each listed prime factor is verified by trial division through its integer square root. The replay verifies \(F\mid T-1\), \(F^2>T\), and, for every prime \(r\mid F\), \[
 z_r^{T-1}\equiv1\pmod T,\qquad
 \gcd(z_r^{(T-1)/r}-1,T)=1,
 \quad z_2=3,\quad z_r=2\ (r\ne2).
 \tag{R17.3}
\] These are precisely Pocklington's factored-divisor criterion. The exact modular residues and gcds are in \nolinkurl{checks/explicit_t_prime_and_stabilizer_certificate.json}. Candidate discovery used a separate primality screen, but that screen is not an input to the saved replay or to this primality proof. The replay uses only integer arithmetic, modular powers, gcds, and trial division on the small factors in(R17.2).

The norms and congruences required in the source follow exactly: \[
 \pi_T\bar\pi_T=T,\qquad \pi_T\equiv1\pmod{c(g_0g_1)^2},
 \qquad T\equiv1\pmod c.
 \tag{R17.4}
\] The certificate checks \(\gcd(T,2A)=1\). A prime dividing both \(\pi_T\) and its conjugate would divide T and their sum \(2A\), which is impossible. The two ideals therefore are coprime, also after extending to either host. Since the host is contained in the c-th cyclotomic field, T splits completely. Exactly f primes divide \(\pi_T\), all to exponent one and with residue field \(\mathbf F_T\): at each conjugate pair one must divide the product T, and coprimality forces exactly one to divide \(\pi_T\). This verifies the norm, congruence, coprimality, half-cardinality and residue-norm clauses for this explicit choice.

The first T has61 digits and the second97. This supplies exact finite values, not an asymptotic upper bound for a changing-weight family. For example, retaining the K031/K032 constants C, define \[
 A_{\rm size}=2(2k+1)C+2+2fN_2.
 \tag{R17.5}
\] The two values are respectively \nolinkurl{7676528679020} and \nolinkurl{27317450103536}. The replay verifies \(2fN_2+1<4A_{\rm size}<T\). These are the \nolinkurl{a_gt} and \nolinkurl{T_gt} inequalities of \nolinkurl{Sizes} for this size parameter. The independently specified \nolinkurl{spread_le} inequality must still be verified for whichever full collection of source words is used; it is not inferred merely from finding a large T. No quantitative MainX window is claimed by(R17.5).

The field-membership clause \nolinkurl{piT_Kq} is deliberately not asserted. K031 proved that the source's defined common field is Q for this induced primitive type. The present \(\pi_T\) lies in the actual external \(\mathbf Q(\sqrt{-7})\), as do the previously constructed G generators.

\subsection{2. The exact rational stabilizer of a standard definite lattice}\label{the-exact-rational-stabilizer-of-a-standard-definite-lattice}

Let E be either host, let F be its maximal real subfield, and let \(X=(E^a,I_a)\). Write \[
 K_0=\prod_{v<\infty}U(I_a,\mathcal O_{F_v}),\qquad
 \Gamma=U(X)(F)\cap K_0.
 \tag{R17.6}
\] At split places this notation means the standard \(\mathrm{GL}_a\) compact; equivalently K0 stabilizes \(\mathcal O_E^a\) at every finite place. Then \[
                  \Gamma=\mu(E)^a\rtimes S_a.
 \tag{R17.7}
\]

Indeed every entry x of a matrix in Gamma is integral, and at every complex embedding \(|\tau(x)|\le1\), by the unitary row norm. If x is nonzero, its algebraic norm is a nonzero integer, so the product of these absolute values is at least1. Every factor must therefore equal1. Such an integral x is a root of unity: all powers of x lie in the bounded set of the Minkowski lattice \(\mathcal O_E\) with these absolute-value bounds; that set is finite, so two powers coincide. The row norm now forces exactly one nonzero entry in each row. Orthogonality forces their column positions to be distinct. This proves containment in the displayed monomial group, and the reverse containment is immediate.

For both fields \(\mu(E)=\mu_{28}\). K031 has Galois group \(C_2^2\times C_{21}\), and K032 has \(C_2^4\times C_{21}\), each of exponent42. Their conductors are c from the table. A cyclotomic subfield arising from an additional root of unity must have conductor dividing c. The possible new factors29, and also5 or13 in the second case, would require a quotient Galois group of exponent divisible by4, which these groups do not possess. The2-part of c is4 and its7-part is7. Since i and \(\zeta_7\) already belong to both fields, exactly \(\mu_{28}\) remains.

Choose one of the T places and identify its residue field with \(\kappa=\mathbf F_T\). Reduction embeds \(\mu_{28}\) in \(\kappa^\times\), since T is prime to28. Suppose \(b\in\kappa^\times\) has order greater than a in the quotient \(\kappa^\times/\mu_{28}\), and take \[
 u=(1,b,b^2,\ldots,b^{a-1}),\quad
 P_u=\{g\in\mathrm{GL}_a(\kappa):ug=s(g)u\},\quad
 H_u=\ker(s:P_u\to\kappa^\times).
 \tag{R17.8}
\] Then the exact intersections with the rational group are \[
 \boxed{\quad\Gamma\cap\widetilde P_u=\mu_{28}I_a,
               \qquad\Gamma\cap\widetilde H_u=\{1\}.\quad}
 \tag{R17.9}
\] Tildes denote reduction preimages inside K0 at that place, with no extra restriction at the other places.

To prove(R17.9), pass to \(\kappa^\times/\mu_{28}\). A monomial g in \(P_u\) must carry the set of coordinate classes \(\{1,[b],\ldots,[b]^{a-1}\}\) to itself by multiplication by \([s(g)]\). The class \([s(g)]\) belongs to the cyclic subgroup generated by \([b]\), because it carries one member of that set to another. A proper consecutive interval of length a in a cyclic group of order greater than a has trivial translation stabilizer: its cyclic consecutive gaps have a unique gap greater than1, which any stabilizing translation must fix. The same claim is immediate for a=1. Thus \([s(g)]=1\). All coordinate classes are distinct, so the permutation of g is the identity. Every diagonal entry equals s(g), and g is scalar. Fixing u then forces that scalar to be1. Conversely every scalar root lies in \(P_u\), as required.

For both numerical rows we may take b=2. The certificate checks \[
                     2^{28j}\not\equiv1\pmod T
                         \quad(1\le j\le a),
 \tag{R17.10}
\] with a=42 or168. Thus no primitive-root search or complete factorization of T-1 is required. Formula(R17.8) is an explicit residue functional. The argument remains valid if the other chosen T places use the same functional, and if the G places are kept at their full maximal compact.

\subsection{3. A nonzero class-supported input and its exact mass}\label{a-nonzero-class-supported-input-and-its-exact-mass}

Let \(\rho_\infty\) be any of the specified irreducible compact types on \(U(a)^f\), let \(d_\rho=\dim\rho_\infty\), and choose a unit vector \(v_\rho\). At each G place use the finite Steinberg representation and its unit generic vector from DO033/LIFT20. Denote their dimensions by \(d_{G,P}\), and put \(d_G=\prod_{P\in G}d_{G,P}\). Their central characters are trivial, because they occur in the finite flag module \(\operatorname{Ind}_B^{\mathrm{GL}_a(\kappa)}1\).

Let H be K0 with its component at each of the f chosen T places replaced by \(\widetilde P_u\); use(R17.8) at every such place. With each maximal compact having volume1, \[
 d_T=[\mathrm{GL}_a(\mathbf F_T):P_u]
       =\frac{T^a-1}{T-1},\quad
 \operatorname{vol}(H)=d_T^{-f},\quad
       \Gamma_H=U(X)(F)\cap H=\mu_{28}I_a.
 \tag{R17.11}
\] The index is the number of projective nonzero row vectors: the full general linear group is transitive on them and Pu is the stabilizer. Equation(R17.9) gives the last assertion; no congruence at29 has entered.

Write \(\omega_\infty\) for the character of \(\mu_{28}\) on \(\rho_\infty\). Choose multiplicative residue characters \(\nu_P:\mathbf F_T^\times\to\mathbf C^\times\), all nontrivial, such that \[
       \omega_\infty(\zeta)\prod_{P\mid\pi_T}\nu_P(\bar\zeta_P)=1
                       \qquad(\zeta\in\mu_{28}).
 \tag{R17.12}
\] This choice always exists. Restriction from the character group of the cyclic group \(\mathbf F_T^\times\) to that of its subgroup \(\mu_{28}\) is surjective. Each restriction has \((T-1)/28>1\) extensions, so it has a nontrivial extension even if the prescribed restriction is trivial. Prescribe trivial restrictions at all but one place and the required inverse character at the remaining place. This also explains the finite coefficient extension needed for the \(\nu_P\): roots of order T-1 must be included. Their coefficient-field cost is not discarded.

For \(h\in H\) and \(k\in U(a)^f\), define \[
 m(k,h)=\langle v_\rho,\rho_\infty(k)v_\rho\rangle
   \prod_{P\in G}\langle v_P,\operatorname{St}_P(\bar h_P)v_P\rangle
   \prod_{P\mid\pi_T}\nu_P(s(\bar h_P)).
 \tag{R17.13}
\] Use the conjugate-linear-first Hermitian convention in this display. Define \(F\) on \([U(X)]\) to be \(\overline{m(k,h)}\) on \(U(X)(F)U(a)^fH\), and zero off this class. This is well-defined: two presentations differ by an element of \(\Gamma_H\), and(R17.12) says exactly that the resulting product multiplier is1. F is smooth, has the required dual real compact type and inverse G-generic character, and is invariant under \(\widetilde H_u\) at T. At29 it is invariant under the full maximal compact. It is nonzero, since F(1)=1. It is an automorphic cusp form because U(X) is anisotropic over F.

The exact scalar-dilation augmentation at each T place is zero: right multiplication by an element of Pu with s-coordinate t multiplies F by \(\nu_P(t)^{-1}\), and \(\sum_{t\in\mathbf F_T^\times}\nu_P(t)^{-1}=0\). This is the local augmentation condition used for the initial ramified column. Reading it as the full source \nolinkurl{RamT} record additionally uses its specified decoration/dilation dictionary; the scalar identity alone does not instantiate that independent record.

Normalize Haar measure on the compact real group to1, as in DO033. The quotient of the support class has measure \(d_T^{-f}/28\). Schur orthogonality at the real group and the independent G compacts therefore gives \[
 \boxed{\quad \|F\|^2
       =\frac{1}{28\,d_T^f\,d_\rho\,d_G},
       \qquad F(1)=1,\qquad\|F\|_\infty\le1.\quad}
 \tag{R17.14}
\] Every T character has modulus1 on its support. The factor28 comes from the actual rational stabilizer and is not absorbed into Haar measure. The finite G scalar denominators and the coefficient metric of \(v_\rho\) remain those of DO033/R011; this statement does not replace their primitive integral models by an arbitrarily rescaled unit vector.

For the explicit ranks the new projective-level index is an11-adic unit: \[
                         v_{11}(d_T)=0.
 \tag{R17.15}
\] For J14, T is2 modulo11 and a=42, so neither numerator nor denominator of dT is divisible by11. For J56, T is1 modulo11 and a=168, so \(v_{11}(T^a-1)-v_{11}(T-1)=v_{11}(168)=0\). The certificate checks these exact valuations. The analytic norm cost \(d_T^f\), despite its11-adic unit valuation, is still large and fully retained in(R17.14).

\subsection{4. What this removes, and what it leaves}\label{what-this-removes-and-what-it-leaves}

The old first-input construction made the rational stabilizer trivial by a principal level at29. Equations(R17.8)--(R17.14) supply a different actual input with the required first G generic behavior and T augmentation while keeping maximal level at29. Thus rational-stabilizer compatibility, by itself, is no longer an obstruction to a source input at that level. The two numerical T choices also replace the qualitative prime choice for these fixed hosts by exact certified integers.

There is still no nonvanishing assertion for the raw kernel applied to this F. If a previously chosen29 Weyl translate is inserted into the theta kernel, the source integral at maximal29 level averages that translate over the compact. Its resulting selected coefficient can vanish. Neither the explicit T-prime nor(R17.14) proves that projection nonzero. The previous-image integral lattice, the terminal row map, the global MainX mass, and the changed-family DP44 cost remain separate. Likewise(R17.4) does not repair \nolinkurl{Fields.Kq}, and(R17.5) does not prove the complete quantitative size window for an unbounded weight family.

\subsection{5. Exact sources and checks}\label{exact-sources-and-checks}

\begin{itemize}
\tightlist
\item
  K031, (HC3)--(HC16) and section5; K032, sections4--6: the two actual product hosts, G-primes, conductors, ranks and constants. These are used with their recorded whole-proof checks.
\item
  Alpöge's stipulated pinned source at commit \nolinkurl{7cbf7a0c32fca1badf05931879a37c3f8a2b03b4}: \nolinkurl{Formalism/LocalToGlobal/Columns/AbelianHostFieldConditions.lean}, \nolinkurl{Primes.Ok}206--238 and \nolinkurl{Sizes}486--494; these are the literal norm, residue, congruence, common-field and size conditions distinguished above. \nolinkurl{LevChiEqOne.lean}, \nolinkurl{RamT}85--87, is the separate global decoration/dilation condition.
\item
  Pocklington, H. C. (1914--1916). The determination of the prime or composite nature of large numbers by Fermat's theorem. \emph{Proceedings of the Cambridge Philosophical Society, 18},29--30. The precise factored-divisor form used here is also stated as Theorem1 in the FLINT official documentation, \emph{Primality proving: APRCL and ECPP}, \url{https://flintlib.org/doc/algorithms/number-theory/primality-proving.html}. The latter statement was consulted directly; this is not a claim that the inaccessible original PDF was read in full.
\item
  DO033 sections1 and4, DO034 section6, R011: matrix-coefficient normalization, actual compact masses and the separate coefficient metric. The new class calculation retains their conventions.
\end{itemize}

The exact replay is \nolinkurl{scripts/explicit_t_prime_and_stabilizer_certificate.py}. Its finite checks certify(R17.1)--(R17.5),(R17.10),(R17.15). The proofs of(R17.7)--(R17.14) are the mathematical arguments above, not consequences of a numerical experiment. Whole-proof review is recorded separately.

\clearpage

\section{ALP-LIFT-21: the raw terminal G factor survives the actual previous theta constituent}\label{alp-lift-21-the-raw-terminal-g-factor-survives-the-actual-previous-theta-constituent}

Author: GPT-6 Astra Pro, ALP\_MACHINE /root/lift\_accounting. Date: 22 November 2026 UTC. Inherited reasoning effort is not exposed. The rank-filtration reduction in section4 was suggested independently by /root while the explicit full-rank intertwiner was being derived.

This theorem concerns the actual terminal local G factor of XiCpl2 for m=2, with the same allowed coupled class used in LIFT20. It proves nonvanishing in the actual small theta quotient of each LIFT19 first constituent. It proves more than occurrence of a compatible compact type: the fixed raw terminal Schwartz vector is injective on the entire prescribed first-column type. In a specified Haar-normalized induced model its norm and its effect on a natural integral coefficient lattice are computed.

The normalization of that local induced model is explicit. Its comparison with a globally pinned automorphic coefficient model is not supplied by uniqueness of the local theta correspondence. No assertion about the final MainX row, all six AdmS identities, or the full DP44 cost follows merely from the local theorem.

\subsection{1. Ranks, actual raw factor, and its compact groups}\label{ranks-actual-raw-factor-and-its-compact-groups}

Fix a degree-one relative split G place of residue characteristic different from11. Write \(F_v\) for its local field, \(O\) for its ring of integers, and \(\kappa=\mathbf F_q\) for its residue field, with the rational prime \(q\) a uniformizer. Normalize multiplicative Haar measure by \(\operatorname{vol}(\mathrm{GL}_n(O))=1\), and additive matrix measure by \(\operatorname{vol}(M_{u,v}(O))=1\).

Take the literal source ranks with \(m=2\). Put

\[
 d=t+1,\quad r=k+1,\quad b=k+d=r+t,\quad N=b+d=k+2d.
 \tag{TG1}
\]

Thus the terminal edge is \(U(b)\to U(N)\). Every actual Ranks datum has enough coordinates for the following blocks. The argument uses \(t\ge1\), so in particular applies to the larger values required by the source.

The source locators are relative to the pinned recovered tree work/research/alpoge\_source, commit 7cbf7a0c32fca1badf05931879a37c3f8a2b03b4:

\begin{itemize}
\tightlist
\item
  \nolinkurl{Formalism/LocalRingCardinality/Params.lean,} definitions of r, N, rS, and cOf in the column file.
\item
  \nolinkurl{Formalism/LocalToGlobal/Columns/FiniteSumNonzeroLargeDepth.lean} (FS), PhiG/PhiGS780--810, PhiTop1141--1147, ftY1150--1155, and XiCpl2 1160--1167.
\item
  \nolinkurl{Development/AbelianVarietyTheory/Definitions/GlobalTwisting/IntegralGLSuccEmbeddingMemParahoricTimesPrincipa.lean,} entB\_top420--441 and the terminal rank identity.
\item
  Checked DO031, equations(DO31.1)--(DO31.5), identifies those exact raw functions and their actual terminal frames.
\end{itemize}

No Lean execution is claimed. We use the definitions and the proved finite model in DO031 with its stated scope.

Write a terminal matrix in its actual local frames as

\[
 X=\begin{pmatrix}Z\\ \ell\end{pmatrix}\in M_{N,b}(F_v),
 \qquad Z\in M_{N-1,b}(F_v),\quad \ell\in M_{1,b}(F_v).
 \tag{TG2}
\]

The terminal top function has no remaining auxiliary-column power: cOf(m-1)=0, so entB\_top makes it read \(Z\) itself. The last-row factor is the actual integral character function in DO31.4. In the convention of that equation it is

\[
             1_{O^b}(\ell)\psi_v(-\ell y),\qquad
                y\equiv q^{-1}e_k\pmod{O^b}.
 \tag{TG3}
\]

The class of \(y\) is exactly LIFT20(RG3), now in the rank-b middle space. If the last row is instead written before the source's piT unit rescaling, a fixed unit multiplies the argument in(TG3). Every proof below keeps the resulting root of unity, and all support, norm, and content assertions are unchanged. We use DO031's already-rescaled convention.

Let \(\eta_0,\eta_1\) be the actual tame selected-half characters on the first and second edges. Their residue characters are nontrivial and distinct, and the literal pin is \(\bar\eta_i=\chi_i^{-1}\). Denote the finite groups

\[
 \begin{aligned}
 H&=\left\{
   \begin{pmatrix}u&A\\0&h\end{pmatrix}:
       u\in U_r(\kappa),\ A\in M_{r,t}(\kappa),\
       h\in\mathrm{GL}_t(\kappa)\right\},\\
 H^\sharp&=\{g\in H:g e_k=e_k\}.
 \end{aligned}
 \tag{TG4}
\]

The index \(k\) is zero-based, so it is the last of the first \(r=k+1\) coordinates. Thus \(H^\sharp\) has \(u=\operatorname{diag}(u_k,1)\), with \(u_k\in U_k\), while \(A,h\) are unchanged. Write \(J,J^\sharp\) for their inverse images in \(\mathrm{GL}_b(O)\). Their cardinalities and Haar volumes are

\[
 \begin{split}
 h_q&=|H|=q^{r(r-1)/2+rt}|\mathrm{GL}_t(\kappa)|,\\
 h_q^\sharp&=q^{-k}h_q,\qquad
 v_J=\frac{h_q}{|\mathrm{GL}_b(\kappa)|},\qquad
 v_{J^\sharp}=q^{-k}v_J.
 \end{split}
 \tag{TG5}
\]

Put

\[
 \lambda(g)=
   \psi_0\!\left(-\sum_{i=0}^{r-2}u_{i,i+1}\right)
                         \bar\eta_0(\det h)\qquad(g\in H).
 \tag{TG6}
\]

This is exactly the inverse adelic first-column character from LIFT19(T1)--(T5), in the \(j=1=m-1\) fully generic branch. The source RecLevel.lean50--52 gives the inverse convention. We use the same notation for its inflation to \(J\) and its restriction to \(J^\sharp\).

\subsubsection{Exact raw support and phase}\label{exact-raw-support-and-phase}

Let \(\Phi,\Phi^\sharp\) denote the actual nonsharp and sharp terminal local factors, including(TG3). They vanish unless \(X\) is integral and, modulo \(q\), has the form

\[
 \bar X=\begin{pmatrix}B\\0_{(d-1)\times b}\\w\end{pmatrix},
 \qquad
 B=\begin{pmatrix}u&A\\0&h\end{pmatrix}\in H
       \text{ or }H^\sharp,\qquad w\in\kappa^b.
 \tag{TG7}
\]

On that support,

\[
 \begin{split}
 \Phi(X)&=
  \psi_0\!\left(\sum_{i=0}^{r-2}u_{i,i+1}-w_k\right)
                       (\bar\eta_1\bar\eta_0^{-1})(\det h),\\
 \Phi^\sharp(X)&=
  \psi_0\!\left(\sum_{i=0}^{k-2}u_{i,i+1}-w_k\right)
                       (\bar\eta_1\bar\eta_0^{-1})(\det h).
 \end{split}
 \tag{TG8}
\]

An empty sum is zero. For the source ranks the sums are nonempty; this convention also gives the small diagnostic.

We verify the raw predicates. In PhiTop the internal variable a=r+(m-1)t is \(b\), the full source width. Its integrality condition makes every entry of \(Z\) integral. The subsequent row condition makes the rows below its first \(b\) rows zero modulo \(q\). The first \(r\) columns have unitriangular leading block and zeros below. The remaining \(t\) columns give an arbitrary upper-right rectangle and a square bottom-right block. If its determinant is zero, unitV and its inverse give zero, so the function vanishes. Otherwise that block is \(h\) in(TG7).

For rS=0, symIdx k 0 is the identity. The additional PhiGS condition fixes column \(k\) to \(e_k\), precisely \(H^\sharp\). The determinant term is \(\chi_0(\det h)\chi_1(\det B)^{-1}\), with \(\det B=\det h\), giving the factor in(TG8). The additive term is the full \(U_r\) character for PhiG; in PhiGS it is the \(U_k\) character because its second root sum has length r-k-1=0. Finally(TG3) contributes \(\psi_0(-w_k)\). Changing \(y\) by an integral vector does not affect it. This proves(TG7)--(TG8) from the actual raw definitions.

Both functions are therefore supported in the full-rank open subset of \(M_{N,b}(F_v)\): the leading \(b\) by \(b\) minor is a unit. Their support is compact in that open set. They are fixed by the target principal congruence \(K_N^1\). Indeed its determinant character is trivial and it changes integral matrix entries only modulo \(q\).

Their exact additive Schwartz norms are

\[
 \Gamma_{\rm raw}=q^{-Nb+b}h_q,\qquad
 \Gamma_{\rm raw}^\sharp=q^{-Nb+b}h_q^\sharp
                              =q^{-k}\Gamma_{\rm raw}.
 \tag{TG9}
\]

The support has \(h_q q^b\) or \(h_q^\sharp q^b\) residue matrices, each of volume \(q^{-Nb}\), and every nonzero value has modulus one.

\subsection{2. The actual preceding constituent}\label{the-actual-preceding-constituent}

Let \(\pi_0\) be an irreducible actual initial local constituent with a finite Steinberg type as in LIFT19. It has Iwahori-fixed vectors and unramified cuspidal support, by the Casselman result already used there. The actual first and second local theta constituents in the literal bilinear convention are

\[
 \pi=\pi_0\times(\eta_0\circ\det_d),\qquad
 \pi^+=\pi\times(\eta_1\circ\det_d).
 \tag{TG10}
\]

The product denotes normalized parabolic induction. Both are irreducible: the added residue characters are distinct from the initial unramified line and from each other. This is LIFT19's application of Zelevinsky's disjoint-line theorem and Mínguez's actual type-II theta correspondence. These published results are used rather than claimed as new.

The source here is the actual constituent \(\pi\) of a first theta image, not an arbitrary representation chosen only by its archimedean weight. Its prescribed compact-type space

\[
                          V_\lambda=\pi^{J,\lambda}
 \tag{TG11}
\]

is nonzero by LIFT19(T16)--(T17). Every nonzero vector in this space works below. It is \(K_b^1\)-fixed. Restricting its character to \(J^\sharp\) gives the same vectors for the literal sharp kernel.

\subsection{3. The full-rank Haar-normalized intertwiner}\label{the-full-rank-haar-normalized-intertwiner}

The actual common-character split oscillator acts by

\[
 \omega(g,h)\phi(X)=
 \eta_1(\det g)^{-1}\eta_1(\det h)
 |\det g|^{N/2}|\det h|^{-b/2}\phi(h^{-1}Xg).
 \tag{TG12}
\]

This is the genuine left-action convention of DO033.5 and corrected LIFT16(S2). The source variable acts by \(X\mapsto Xg\). The source bilinear lift takes diagonal coinvariants of \(\omega\otimes\pi\).

Write \(\Omega\) for the full-rank open set and \(X_0=[I_b;0]\). For \(\phi\in S(\Omega)\), \(v\in\pi\), define

\[
 {\cal I}(\phi,v)=
   \int_{\mathrm{GL}_b(F_v)}
       \phi(X_0g)\eta_1(\det g)^{-1}
                    |\det g|^{N/2}\pi(g)v\,dg.
 \tag{TG13}
\]

This is an ordinary compact integral, with no meromorphic regularization. The image of \(g\mapsto X_0g\) is closed in \(\Omega\); the inverse image of the compact support of \(\phi\) is compact in \(\mathrm{GL}_b\). The smooth integrand takes values in a finite-dimensional span on that compact set, so its integral lies in \(\pi\).

Let \(P=P_{(b,d)}\) be the upper parabolic of \(\mathrm{GL}_N\). Set

\[
             {\cal T}(\phi,v)(h)
                      ={\cal I}(\omega(1,h)\phi,v).
 \tag{TG14}
\]

Then T factors through the source diagonal coinvariants and takes values in the actual normalized induction \(\pi^+\). The complete covariance checks are as follows.

First \({\cal I}(\omega(g_0,1)\phi,\pi(g_0)v)={\cal I}(\phi,v)\). Substitute \(g'=g g_0\) in(TG13). The two eta factors combine to \(\eta_1(\det g')^{-1}\), the absolute-value factors combine identically, and \(\pi(g)\pi(g_0)=\pi(g')\). Unimodularity preserves Haar.

Next let \(p=[A\ C;0\ D]\in P\), so \(p^{-1}X_0=X_0A^{-1}\). Substitute \(g=A g'\) in(TG13), using(TG12), to obtain exactly

\[
 {\cal I}(\omega(1,p)\phi,v)=
 |\det A|^{d/2}|\det D|^{-b/2}
          \eta_1(\det D)\pi(A){\cal I}(\phi,v).
 \tag{TG15}
\]

The unipotent C disappears because it fixes the embedded source space. The absolute-value factor is \(\delta_P(p)^{1/2}\). Therefore T(phi,v)(p h) has precisely the normalized induction transformation law. Target translation by \(h_0\) sends the section to \(h\mapsto T(\phi,v)(h h_0)\). Smoothness follows from a small compact fixing \(\phi\), and \(P\backslash\mathrm{GL}_N\) is compact.

The full-rank coinvariant identification is

\[
 (S(\Omega)\otimes\pi)_{\mathrm{GL}_b}
       \simeq
 \operatorname{Ind}_{P}^{\mathrm{GL}_N}
                  (\pi\boxtimes\eta_1\det_d),
 \tag{TG16}
\]

with map T and the displayed normalization. Here is an elementary verification in addition to the covariance calculation. Full-rank matrices form a locally trivial principal GL\_b bundle over the p-adic Grassmannian. That compact, totally disconnected space admits a finite clopen refinement of trivializing charts for each locally constant function in question. On a chart, multiply the fibre function by \(\eta_1(\det g)^{-1}|\det g|^{N/2}\). The source action becomes ordinary right translation on \(C_c^\infty(\mathrm{GL}_b)\), together with \(\pi\) on its coefficient vector. Its diagonal coinvariant map is \(f\otimes v\mapsto\int f(g)\pi(g)v\,dg\), an isomorphism onto \(\pi\).

To verify that assertion, a vector fixed by a compact open K is represented by \(\operatorname{vol}(K)^{-1}1_K\otimes v\). Changing to a smaller such K leaves its class unchanged by finite coset decomposition and the diagonal relations. Conversely, partition the compact support of f into cosets of a small compact open. The same diagonal relations and finite averaging reduce its class to that of its displayed integral. This proves both injectivity and surjectivity. Finite partitions of unity on the clopen charts glue the fibre maps, with transition law(TG15). This proves(TG16).

\subsection{4. Identification with the actual small theta quotient}\label{identification-with-the-actual-small-theta-quotient}

The open-set map must be related to the full oscillator. We use the published rank filtration in Mínguez(2008), section2, pp.722--724. The boundary quotients are the rank j strata with \(j<b\). Corollary2.3 describes their possible irreducible source quotients: after removing oscillator determinant normalization, such a quotient is a quotient of an induced representation with a nonzero trivial GL\_(b-j) factor. Modulus characters only introduce unramified twists of that factor. Its cuspidal support therefore contains an unramified line.

For(TG12) the relevant source quotient is \(\pi^\vee\). Removing \(\eta_1^{-1}\) replaces its cuspidal inertia lines by those of \(\pi^\vee\otimes\eta_1\), up to unramified powers and the inverse-transpose automorphism used in the alternate rectangular-matrix convention. Those lines are

\[
                         \bar\eta_1,\qquad
                         \bar\eta_1\bar\eta_0^{-1}.
 \tag{TG17}
\]

Both are nontrivial. Simultaneously inverting them under the alternate convention changes nothing. Uniqueness of cuspidal support therefore excludes the source quotient from every boundary stratum. No pole-free regularization is assumed.

There is an actual nonzero oscillator map to \(\pi^\vee\otimes\pi^+\) by LIFT19 and Mínguez's theta theorem. Its restriction to \(S(\Omega)\) cannot vanish: otherwise a nonzero map from a boundary stratum would follow by the finite rank filtration, contradicting the preceding argument. By(TG16) and irreducibility of \(\pi^+\), Schur's lemma makes the target-equivariant Hom space from the full-rank coinvariants to \(\pi^+\) one-dimensional. Thus the actual small-theta map restricts to a nonzero scalar multiple of T.

The coinvariant and tensor-Hom formulations agree here without a completed-tensor ambiguity. A compact open source subgroup fixing a given Schwartz vector makes the dependence on the input factor through its invariants. Admissibility of \(\pi\) makes that space finite dimensional, so the corresponding map has coefficients in the smooth contragredient tensor the target.

Every vector detected by T is therefore nonzero in the actual small theta quotient. One may normalize that local quotient map so its restriction equals(TG14). The scalar comparing this choice with a globally pinned automorphic realization is not computed by uniqueness. All norm and lattice statements retain that distinction.

\subsection{5. Exact survival of the prescribed column}\label{exact-survival-of-the-prescribed-column}

For \(v\in V_\lambda\), evaluate T at the identity. The raw factor on \(X_0g\) is supported on \(g\in J\); its last row is zero. The determinant terms give

\[
               \Phi(X_0g)\eta_1(\det g)^{-1}
                                      =\lambda(g)^{-1}.
 \tag{TG18}
\]

The determinant has absolute value one and \(\pi(g)v=\lambda(g)v\), so the integrand in(TG13) is exactly v. On \(J^\sharp\) the restricted character gives the identical calculation. Hence

\[
 \boxed{\quad
 {\cal T}(\Phi,v)(1)=v_Jv,\qquad
 {\cal T}(\Phi^\sharp,v)(1)=v_{J^\sharp}v
                   \quad(v\in V_\lambda).\quad}
 \tag{TG19}
\]

Both fixed raw kernels are injective on the entire type space. Every nonzero prior-image vector in this prescribed local column survives. This does not assert survival for arbitrary vectors outside the column. Without the type condition, identity evaluation is exactly the compact character projector: T(Phi,v)(1)=v\_J e\_(J,lambda)v, and similarly for the sharp subgroup. Thus(TG19) concerns a fixed raw function, not a Schwartz choice depending on v.

\subsection{6. Exact compact-model norms}\label{exact-compact-model-norms}

The actual automorphic constituent \(\pi\) is unitary. Give \(\pi^+\) its normalized induction metric on \(K_N=\mathrm{GL}_N(O)\), retaining the chosen unitary metric on \(\pi\). Put

\[
 {\cal J}_{b,d}(q)=[\mathrm{GL}_N(\kappa):P_{(b,d)}(\kappa)]
     =\frac{\prod_{i=1}^{N}(q^i-1)}
        {\prod_{i=1}^{b}(q^i-1)\prod_{i=1}^{d}(q^i-1)}.
 \tag{TG20}
\]

The sections are target \(K_N^1\)-fixed. Their finite parabolic quotient points have measure \(1/{\cal J}_{b,d}\). A nonzero value can occur only when the residue source b-plane is a graph

\[
 L_z=\left\{\begin{pmatrix}x\\0\\zx\end{pmatrix}:x\in\kappa^b\right\},
                    \qquad z\in\kappa^{1\times b},
 \tag{TG21}
\]

where the zero has d-1 coordinates. The raw support requires the plane to lie in the space spanned by the first b and last coordinates and project invertibly onto the first b, giving exactly these graphs. The reduction of a compact frame has full rank, so a contributing g is integral and invertible modulo q; no other finite plane can contribute. For a graph choose the integral representative \(h_z=I_N-e_Nz\), with z extended by zeros on the last d coordinates. Then \(h_z^{-1}X_0=[I_b;0;z]\).

For \(v\in V_\lambda\), all determinant and representation factors cancel as in(TG18). The finite averages are

\[
 \begin{split}
 {\cal T}(\Phi,v)(h_z)
  &=v_J\left(|H|^{-1}\sum_{g\in H}\psi_0(-zge_k)\right)v,\\
 {\cal T}(\Phi^\sharp,v)(h_z)
  &=v_{J^\sharp}\psi_0(-z_k)v.
 \end{split}
 \tag{TG22}
\]

For the sharp subgroup \(ge_k=e_k\). For H, the first k entries of \(ge_k\) vary independently over \(\kappa\), its k-th entry is one, and its subsequent entries are zero. Additive orthogonality therefore gives

\[
 {\cal T}(\Phi,v)(h_z)=
 \begin{cases}
    v_J\psi_0(-z_k)v,&z_0=\cdots=z_{k-1}=0,\\
    0,&\text{otherwise}.
 \end{cases}
 \tag{TG23}
\]

There are exactly \(q^{b-k}\) nonzero finite quotient values in the nonsharp case and \(q^b\) in the sharp case. Hence for every \(v\in V_\lambda\),

\[
 \boxed{
 \begin{split}
 \|{\cal T}(\Phi,v)\|^2
    &=\frac{q^{b-k}v_J^2}{{\cal J}_{b,d}(q)}\|v\|^2,\\
 \|{\cal T}(\Phi^\sharp,v)\|^2
    &=\frac{q^b v_{J^\sharp}^2}{{\cal J}_{b,d}(q)}\|v\|^2.
 \end{split}}
 \tag{TG24}
\]

Their squared constants have ratio \(q^{-k}\), the same as the raw Schwartz norms(TG9). After unit-normalizing the finite raw input, the two maps have the same squared scaling constant:

\[
 \boxed{
 {\cal C}_q=
 \frac{q^{Nb-k}h_q}{|\mathrm{GL}_b(\kappa)|^2{\cal J}_{b,d}(q)}
 =\frac{q^{Nb}h_q^\sharp}
        {|\mathrm{GL}_b(\kappa)|^2{\cal J}_{b,d}(q)}.}
 \tag{TG25}
\]

These are local induced-model metric identities. The real selected Schwartz norm, global mass, and scalar comparing this model to an automorphic row are separate. The finite sharp modification alone does not improve this unit-input scaling. For the primitive LIFT19(T18)--(T20) vector, its natural metric is \(\|v\|^2=q^{k(k+1)/2}/{\cal J}_{k,d}(q)\); it may be substituted into(TG24) directly.

\subsection{7. Exact content in the natural induced lattices}\label{exact-content-in-the-natural-induced-lattices}

Let R be the valuation ring at a place above11 of a coefficient extension containing the additive and tame character values. Use the natural finite models supplied by LIFT19:

\[
 \begin{split}
 L_1&=\operatorname{Ind}_{P_{(k,d)}(\kappa)}
                           ^{\mathrm{GL}_b(\kappa)}
       (\mathrm{St}_{k,R}\boxtimes\bar\eta_0\det_d),\\
 L_2&=\operatorname{Ind}_{P_{(b,d)}(\kappa)}
                           ^{\mathrm{GL}_N(\kappa)}
       (L_1\boxtimes\bar\eta_1\det_d).
 \end{split}
 \tag{TG26}
\]

Fix one finite Steinberg embedding into \(\pi_0^{K_k^1}\). These are the selected finite submodules of the actual compact models obtained from that embedding, with specified integral inducing lattices. They need not exhaust all principal-congruence invariants of either constituent. Their comparison with a global arithmetic row lattice is extra data. Put \(L_\lambda=L_1\cap V_\lambda\). It is saturated in \(L_1\): its defining equations are the kernels of R-linear maps \(\pi(h)-\lambda(h)\) into a torsion-free module.

Divide T by \(v_J\) or \(v_{J^\sharp}\). Equations(TG22)--(TG23) show that every finite quotient value is zero or a root of unity times v. The identity value is exactly v; parabolic transition maps preserve the integral lattices. Consequently

\[
 v_J^{-1}{\cal T}(\Phi,-),\quad
 v_{J^\sharp}^{-1}{\cal T}(\Phi^\sharp,-):
                         L_\lambda\longrightarrow L_2
 \tag{TG27}
\]

are primitive saturated injections. To prove saturation, suppose a uniformizer divides the image in \(L_2\). Identity evaluation shows it divides v in \(L_1\); saturation of \(L_\lambda\) puts the divided vector in the domain. Torsion-freeness then identifies its image with the divided image. This does not assume existence of an integral H-projector.

If \(c_R(v)\) denotes the content valuation in the indicated free R-module, then

\[
 \begin{split}
 c_R({\cal T}(\Phi,v))&=c_R(v)+v_{11}(v_J),\\
 c_R({\cal T}(\Phi^\sharp,v))&=c_R(v)+v_{11}(v_{J^\sharp}),\\
 v_{11}(v_J)=v_{11}(v_{J^\sharp})
              &=-\sum_{i=t+1}^{b}v_{11}(q^i-1).
 \end{split}
 \tag{TG28}
\]

Here \(v_{11}(11)=1\), and q is an11-unit. The matrix-coefficient normalization and natural induced lattices are fixed in this statement. It cannot be identified with the global DP44 gain \(g_j\) without the missing row-lattice and scalar comparison.

\subsection{8. Global implication and remaining exact gates}\label{global-implication-and-remaining-exact-gates}

The local result applies at every degree-one selected G place of a repaired host with distinct nontrivial literal tame characters. It is independent of the number of real embeddings. It therefore applies to the actual first constituents from K031/DO034, and to later repaired hosts only after their actual first-family and character hypotheses are supplied.

The fixed terminal raw G factor cannot annihilate a nonzero previous-image vector in \(V_\lambda\). Finite tensor products of these column spaces satisfy the product versions of(TG19), (TG24), and(TG28). Thus this particular G-place test is no longer an unspecified local nonvanishing condition on the m=2 route.

This does not put a particular global first vector into every required column by an integral operator of known height. LIFT19 supplies occurrence and an explicit natural vector; the comparison with the exact global previous-image lattice remains necessary. Nor does this argument treat the fixed T/Q terminal factors, the permitted29 family at the second edge, the final global row, or the normalization of all local quotient maps by one global theta period. These are separate gates before a full DP44 changed-data comparison or a shorter source-compatible chain is claimed.

\subsection{9. Sources and proof status}\label{sources-and-proof-status-1}

Mínguez, A. (2008). Correspondance de Howe explicite: Paires duales de type II. Annales scientifiques de l'École Normale Supérieure, 41(5), 717--741. \url{https://www.numdam.org/item/ASENS_2008_4_41_5_717_0/} . Theorem1, p.718; action conventions in section1.4, pp.721--722; exact rank filtration, Corollary2.3, and Theorem2.4 with its restriction argument, pp.722--724. Read directly in the primary PDF. The theorem is used as a published input.

Other published representation inputs are those already checked in LIFT19: Casselman(1980), Proposition2.6; Zelevinsky(1980), Propositions8.5--8.6; and the finite Steinberg theorem. They are not counted again as new results.

The new calculations are the raw terminal compact support and character, the Haar-normalized open-rank integral in the actual source convention, its value on the whole prescribed column, the graph-plane norm, and the saturated natural-lattice comparison. The author whole-proof reading has been completed, including the raw source predicates, all source and target covariance signs, full-rank coinvariants, the boundary restriction argument, every graph value, both norm normalizations, and saturation without an integral isotypic projector.

The exact diagnostic \nolinkurl{raw_terminal_g_certificate.py} passed two finite cases: q=5,k=1,t=1 with distinct nontrivial residue characters, and q=3,k=2,t=1 for the higher-k raw geometry. It checks1,958 raw determinant/input phase cancellations and412 full graph averages, the support counts, both Haar volumes, and equal unit-input metric scaling. The second case does not assert the distinct character hypothesis, and neither small case is an admissible full source Ranks datum. The general quotient and lattice proofs are independent of these finite checks. Exact check events are in \nolinkurl{lift_accounting_CHECK_EVENTS.tsv.}

\clearpage

\section{AM-DO-035: all first coupled raw factors, their compact filters, and the remaining 29-place gate}\label{am-do-035-all-first-coupled-raw-factors-their-compact-filters-and-the-remaining-29-place-gate}

Author: GPT-6 Astra Pro, ALP\_MACHINE \nolinkurl{/root/definite_oscillator}. 22 November 2026 UTC. Inherited reasoning effort is not separately exposed.

This completes the local G/T/Q compatibility calculation for the first coupled kernel when \nolinkurl{m=2}. It uses the actual \nolinkurl{ybSet} class, the actual common splitting character, and an initial source input satisfying the prescribed G, T and Q conditions. In particular, the Q determinant character cancels; it is not an obstruction to \nolinkurl{zero_Q}. A nontrivial T diamond character gives the initial ramification condition and an explicit nonzero compact filter.

Combined with ALP-LIFT-20 and AM-DO-030/034, this gives a nonzero first lift with \textbf{all raw G/T/Q factors retained}, after one allowed auxiliary 29-place Heisenberg translate and a source level refinement there. It does not identify that translated kernel or refined input with the literal unmodified \nolinkurl{XiCpl1}/\nolinkurl{MfV}. Section 8 states the remaining equation and proves the precise limitation of absorbing it by lattice moves.

The local identities hold for every source-admissible T prime and twisting tuple. The global corollary assumes fixed such data, with all its auxiliary primes away11 and29. This file does not certify an explicit T prime or instantiate the complete \nolinkurl{Primes.Ok} record.

\subsection{1. Source locations and conventions}\label{source-locations-and-conventions}

All source paths below are in the recovered original commit \nolinkurl{7cbf7a0c32fca1badf05931879a37c3f8a2b03b4}. Abbreviate \nolinkurl{Formalism/LocalToGlobal/Columns/FiniteSumNonzeroLargeDepth.lean} by FS. These are source readings, not a Lean execution.

\begin{longtable}[]{@{}
  >{\raggedright\arraybackslash}p{(\columnwidth - 2\tabcolsep) * \real{0.5000}}
  >{\raggedright\arraybackslash}p{(\columnwidth - 2\tabcolsep) * \real{0.5000}}@{}}
\toprule\noalign{}
\begin{minipage}[b]{\linewidth}\raggedright
Source locator
\end{minipage} & \begin{minipage}[b]{\linewidth}\raggedright
Clause used
\end{minipage} \\
\midrule\noalign{}
\endhead
\bottomrule\noalign{}
\endlastfoot
FS594--598, 600--608, 636--641 & \nolinkurl{homAt}, the two scaled modules, and the actual \nolinkurl{kerOf} sum \\
FS646--655 & \nolinkurl{phiT} and the chosen lattice/residue frames \nolinkurl{liftFr} \\
FS684--694, 714--719 & \nolinkurl{entB} and the raw Q factor \nolinkurl{PhiQ} \\
FS780--810, 981--986 & raw G/GS factors, \nolinkurl{joinCol}, and \nolinkurl{PsiAt} \\
FS999--1010 & actual uniformizer and residue tying of the common character \\
FS1019--1023 & initial index, \nolinkurl{hal=1}, prescribed G type and full trivial Q level \\
FS1038--1061 & depths, support and representative-period conditions \\
FS1103--1138 & actual frames, \nolinkurl{Sall}, \nolinkurl{cAof}, \nolinkurl{cBof}, \nolinkurl{ybSet}, \nolinkurl{repY}, and \nolinkurl{XiCpl1} \\
FS1246--1259 & inverse source and direct target determinant characters \\
\nolinkurl{Columns/Adelic.lean}139--145 & mirabolic source level at T \\
\nolinkurl{Columns/RecLevel.lean}50--52 & the adelic source character is the inverse column character \\
\nolinkurl{Columns/ScaledRamifiedLattice.lean}81--94, 160--169 & diamond operators, their sum \nolinkurl{NmG}, and \nolinkurl{MramV} \\
\nolinkurl{Development/AbelianVarietyTheory/Definitions/LatticeFrames/LevChiEqOne.lean}61, 100, 114--117 & actual common label \nolinkurl{mu0_j * xi(a)}, raw datum and scalar tie \\
\nolinkurl{Columns/AbelianHostFieldConditions.lean}314--334 & reference characters unramified away G, and their actual tying \\
\nolinkurl{Formalism/LocalToGlobal/Divisibility/Tuples.lean}, \nolinkurl{GoodPl}, \nolinkurl{tupGrp}, \nolinkurl{Qprime}, \nolinkurl{resChar} & twist ramification, tame residue characters and degree-one Q places \\
\end{longtable}

Put

\[
 m=2,\qquad a=k,\qquad b=k+t+1,\qquad m_2=0,
 \qquad r=a+1,\quad r_S=0,\quad c_0=1,\quad c_1=0.
 \tag{DO35.1}
\]

Take the base decorated lattices at the first two ranks. Their underlying lattices are the standard integral lattices. Their T functionals are nonzero residue covectors; the base choice is the last-coordinate functional. Their Q frames are the identity. At G the selected \nolinkurl{liftFr} need not literally be the identity, but preserves the integral lattice and induces the prescribed identity residue frame. Therefore all residue formulas below remain valid in those exact frames.

At a relative split degree-one place write \nolinkurl{O} for the local ring, \nolinkurl{q} for its rational uniformizer and residue cardinality, and \nolinkurl{kappa=F_q}. Additive matrix Haar gives \nolinkurl{M_(b,a)(O)} volume one. The genuine source action, with the actual common splitting character having local component \nolinkurl{eta}, is

\[
 [\omega(g)f](X)
   =\eta(\det g)^{-1}|\det g|^{b/2}f(Xg).
 \tag{DO35.2}
\]

This is the split formula checked in AM-DO-020/024 and ALP-LIFT-20. On a compact subgroup the absolute-value factor is one. The formula uses a genuine left representation; replacing \nolinkurl{Xg} by \nolinkurl{Xg^(-1)} without changing the other conventions would reverse the calculation below.

Let \nolinkurl{xi} denote the finite-order twisting tuple and \nolinkurl{mu_0} the first reference character. The actual character in (DO35.2) is

\[
                          \eta=\mu_0\xi.
 \tag{DO35.3}
\]

This is exactly \nolinkurl{TwDat.Is.chr} and \nolinkurl{PinX.mv} in the cited source. At G the tuple is unramified and the reference character is prescribed. At T both factors are unramified. At Q the reference factor is unramified, whereas \nolinkurl{xi} has its nontrivial tame residue character. The Q set belongs to the tuple, not to the reference character.

\subsection{2. One admissible global target class for all G and Q places}\label{one-admissible-global-target-class-for-all-g-and-q-places}

Write \nolinkurl{m_c=cst.M} to distinguish this fixed positive integer from the auxiliary isogeny degree \nolinkurl{M=29^e}. At every selected G place prescribe

\[
             y_P\equiv q^{-1}e_a\pmod{O^b};
 \tag{DO35.4}
\]

at every Q place prescribe

\[
             y_P\equiv e_a\pmod{q^{m_c}O^b}.
 \tag{DO35.5}
\]

Here \nolinkurl{e_a} is zero-indexed: it is the first coordinate immediately following the first \nolinkurl{a} target coordinates. These prescriptions are made in the target frames in FS1103--1106. They are compatible with the base frames described above.

\textbf{Claim.} A single rational target vector \nolinkurl{y} defines an actual class \nolinkurl{Y} in \nolinkurl{ybSet} satisfying every prescription (DO35.4)--(DO35.5).

\textbf{Proof.} The selected G and Q prime ideals are pairwise distinct; the Q set is disjoint from G and T by the tuple's good-place condition. Multiply by \nolinkurl{D_aux=g_0 g_1}. At each G prime the required congruence on \nolinkurl{D_aux y} is integral modulo that prime; at each Q prime it is an integral congruence modulo its \nolinkurl{m_c}-th power, since \nolinkurl{D_aux} is a unit there. The Chinese remainder theorem in \nolinkurl{O_K}, coordinate by coordinate, gives an integral vector satisfying all these congruences. Divide it by \nolinkurl{D_aux}. This gives a rational vector with the stated local conditions. In the actual selected G frames the same proof uses their integral coordinate maps; their reductions are the specified residue frames. In particular the congruence \nolinkurl{q y=e_a mod q} is unchanged by a frame congruent to the identity.

The vector is in the module defining \nolinkurl{ybSet}: its G valuation is at least \nolinkurl{-1}, and its Q valuation is at least zero, hence at least \nolinkurl{-m_c}. Set

\[
 Y=y+\operatorname{locLat}
       \bigl(N_1;\ e_G=0,\ e_Q=m_c\bigr).
 \tag{DO35.6}
\]

This is the literal set in FS1128--1130. Every representative of \nolinkurl{Y}, including \nolinkurl{repY Y}, has the same prescribed residues. At G ALP-LIFT-20 proves the exact representative invariance; at Q it follows from the determinant and divisibility conditions below (or from the actual \nolinkurl{psiInv} clause). No independent local \nolinkurl{y} is being substituted for the one global class. ∎

T imposes no condition on \nolinkurl{y}: its factor is \nolinkurl{phiT(X)}. Its distinct prime support makes its incidence conditions compatible with this same G/Q class. For example, one may choose a rational \nolinkurl{X} with \nolinkurl{X=X_0 mod G}, the Q block residues in Section 3, and the T last row condition in Section 4, again by the Chinese remainder theorem. Thus the combined finite raw function is nonzero before any theta or compact averaging assertion is made.

\subsection{3. The exact Q factor and cancellation of its compact character}\label{the-exact-q-factor-and-cancellation-of-its-compact-character}

Let

\[
 X=\begin{pmatrix}A\\ r\\ B\end{pmatrix},
 \quad A\in M_a(K_P),\ r\in M_{1,a}(K_P),\
 B\in M_{t,a}(K_P),
 \qquad \chi=\operatorname{resChar}(\xi,P).
 \tag{DO35.7}
\]

At Q the exponents in \nolinkurl{entB} are all zero. Both group arguments in \nolinkurl{PsiAt} are one, by \nolinkurl{zero_h} and its definition. For the class (DO35.5), the top \nolinkurl{(a+1)} square block of \nolinkurl{[X,y]} has last column \nolinkurl{e_a} modulo \nolinkurl{q}, and the lower \nolinkurl{t} rows of that column vanish modulo \nolinkurl{q}. Consequently the actual source definition gives

\[
 f_Q(X)=
 \chi(\det\bar A)\,
 1_{\{A\in\operatorname{GL}_a(O),\ r\in O^a,
                   \ B\in qM_{t,a}(O)\}}.
 \tag{DO35.8}
\]

The determinant character is extended by zero on nonunits, exactly as \nolinkurl{unitV} in FS. This proves (DO35.8) also when the top block is singular. A change \nolinkurl{y -> y+q^(m_c)z}, with integral \nolinkurl{z} and \nolinkurl{m_c>=1}, changes neither reduction nor lower-row divisibility.

For every \nolinkurl{g in GL_a(O)} the support is preserved and

\[
                     f_Q(Xg)=\chi(\det\bar g)f_Q(X).
 \tag{DO35.9}
\]

But the actual label (DO35.3) has exactly the residue character \nolinkurl{chi} at Q. Substituting (DO35.9) into (DO35.2) proves

\[
                  \boxed{\omega(g)f_Q=f_Q\quad(g\in\operatorname{GL}_a(O)).}
 \tag{DO35.10}
\]

Thus the raw Q factor is fixed by the full source compact. This is precisely the initial condition \nolinkurl{zero_Q=LevQFam.top}, whose definition is \nolinkurl{H=top, lam=1} (FS409). Using only the reference label \nolinkurl{mu_0} in (DO35.2) would incorrectly leave an extra character. The inverse source residue character \nolinkurl{rhoS=chi^(-1)} in FS1246--1248 confirms the same cancellation independently.

Its exact additive norm is

\[
 \|f_Q\|^2
  =q^{-ta}\prod_{i=1}^{a}(1-q^{-i})=:\Gamma_Q.
 \tag{DO35.11}
\]

Indeed \nolinkurl{A} has invertible reduction with the displayed probability, \nolinkurl{r} is arbitrary integral, and \nolinkurl{B} has \nolinkurl{ta} entries divisible by \nolinkurl{q}. Every nonzero value has modulus one. The value at \nolinkurl{[I_a;0;0]} is one.

The function is supported on integral matrices and periodic modulo \nolinkurl{q}. Its finite Fourier expansion in the Heisenberg vacuum basis is therefore an actual finite combination at P. All denominators divide a power of \nolinkurl{q}; at a coefficient place over11 with \nolinkurl{q!=11}, at least one Fourier coefficient is a unit. This last assertion does not require a Gauss-sum estimate: Fourier inversion over a finite group of order prime to11 is an integral automorphism, and (DO35.8) has the unit value one. The free row \nolinkurl{r} forces the corresponding dual row to vanish, so its Fourier support has size at most \nolinkurl{q^(a(b-1))}. Parseval gives exactly (DO35.11).

\subsection{4. The T factor, source level and initial ramification condition}\label{the-t-factor-source-level-and-initial-ramification-condition}

Let \nolinkurl{u in kappa^a} and \nolinkurl{v in kappa^b} be the nonzero source and target residue covectors. In the base decorations they are the last-coordinate covectors. The exact factor \nolinkurl{phiT} is

\[
 f_{T,u}(X)
  =1_{M_{b,a}(O)}(X)\,1_{\{v\bar X=u\}}.
 \tag{DO35.12}
\]

Its norm is \nolinkurl{q^(-a)} because these are \nolinkurl{a} independent residue conditions. Put

\[
 H_u=\{g\in\operatorname{GL}_a(\kappa):ug=u\},\qquad
 P_u=\{g:ug=s(g)u\text{ for some }s(g)\in\kappa^\times\}.
 \tag{DO35.13}
\]

The quotient \nolinkurl{P_u/H_u} is \nolinkurl{kappa^times}, and

\[
 d_T=[\operatorname{GL}_a(\kappa):P_u]
       =\frac{q^a-1}{q-1}.
 \tag{DO35.14}
\]

Choose any nontrivial unitary residue character \nolinkurl{nu} of \nolinkurl{kappa^times}. The rational T prime is odd in the source, so such a character exists. Work over a coefficient extension containing its values. Define

\[
 f_{T,\nu}=\sum_{s\in\kappa^\times}\nu(s)^{-1}f_{T,su}.
 \tag{DO35.15}
\]

The supports are disjoint, so this is nonzero and

\[
                    \|f_{T,\nu}\|^2=(q-1)q^{-a}.
 \tag{DO35.16}
\]

At T the character (DO35.3) is unramified. Equation (DO35.2) therefore gives, for a lift of \nolinkurl{g in P_u},

\[
 \omega(g)f_{T,u}=f_{T,s(g)^{-1}u}.
 \tag{DO35.17}
\]

With maximal-compact Haar of volume one, extend \nolinkurl{nu(s(g))} by zero from the inverse image of \nolinkurl{P_u}. Equal-size fibres of \nolinkurl{s:P_u -> kappa^times} prove the exact compact filter

\[
 \boxed{\int_{K_T}1_{P_u}(g)\nu(s(g))\,
                  \omega(g)f_{T,u}\,dg
       =\frac{f_{T,\nu}}{d_T(q-1)}
       =\frac{f_{T,\nu}}{q^a-1}.}
 \tag{DO35.18}
\]

The scalar function in this integral has absolute value one on \nolinkurl{P_u}, zero outside, norm squared \nolinkurl{1/d_T}, and is right \nolinkurl{H_u} invariant. This matches the literal T mirabolic level in \nolinkurl{Adelic.lean}139--145. Right multiplication by a lift with \nolinkurl{s(g)=c} multiplies it by \nolinkurl{nu(c)}. Since \nolinkurl{nu} is nontrivial,

\[
                    \sum_{c\in\kappa^\times}\nu(c)=0.
 \tag{DO35.19}
\]

Thus its T diamond sum vanishes. This is exactly the defining local condition in \nolinkurl{NmG} and \nolinkurl{ColRamV}, not a cuspidality claim. The source diamond operation multiplies the decorated covector by \nolinkurl{c} (\nolinkurl{StdLatUnimodular.lean}113--116); changing the identification of right translation with that operation replaces \nolinkurl{nu} by its inverse, and leaves (DO35.19) unchanged. Its scalar action on the algebraic coefficient is the identity in \nolinkurl{diamondOperatorOnForms'}.

For completeness, the finite Fourier coefficients can also be computed exactly. Fix a nontrivial residue additive character \nolinkurl{psi}. The coefficient indexed by \nolinkurl{ell in kappa^a} on the rank-one matrix frequency \nolinkurl{v^t ell} is

\[
 b_\ell=q^{-a}\sum_{s\ne0}\nu(s)^{-1}\psi(-s\,\ell u^t).
 \tag{DO35.20}
\]

It vanishes when \nolinkurl{ell u^t=0}; otherwise it is a nonzero Gauss sum divided by \nolinkurl{q^a}, of complex modulus \nolinkurl{q^(1/2-a)}. To see nonvanishing and the norm without importing a formula, put \nolinkurl{G(c)=sum_(s!=0) nu(s)^(-1) psi(cs)}. For \nolinkurl{c!=0}, substitution \nolinkurl{s -> c^(-1)s} shows all \nolinkurl{|G(c)|} agree. Orthogonality of \nolinkurl{psi} gives \nolinkurl{sum_c |G(c)|^2=q(q-1)} while \nolinkurl{G(0)=0}, hence \nolinkurl{|G(c)|^2=q}. The product of the Gauss sum and its conjugate is \nolinkurl{q}, so these coefficients are units at every place over11 when \nolinkurl{q!=11}. Their number is exactly

\[
                      s_T=(q-1)q^{a-1}.
 \tag{DO35.21}
\]

Parseval recovers (DO35.16). Their total complex absolute value is \nolinkurl{(q-1)/sqrt(q)}. No division by \nolinkurl{q-1} has been made in the auxiliary vector (DO35.15); the division belongs to the actual compact filter (DO35.18) and must be charged there.

\subsection{5. The complete auxiliary vector and exact fixed costs}\label{the-complete-auxiliary-vector-and-exact-fixed-costs}

At each G place use the Steinberg vector \nolinkurl{f_St} of AM-DO-033, with its determinant factor for the actual reference label. ALP-LIFT-20 proves, for the actual raw factor \nolinkurl{f_U},

\[
 \int_{K_G}\overline{m_G(g)}\omega(g)f_U\,dg
          =\frac{\alpha_G}{d_G}f_{St},
 \quad d_G=q^{a(a-1)/2},\quad
 \alpha_G=\frac{d_G^2}{|\operatorname{GL}_a(\kappa)|}.
 \tag{DO35.22}
\]

Here \nolinkurl{m_G} is the matrix coefficient of a unit generic vector in the finite Steinberg representation. Its inverse-unipotent source type is the literal initial \nolinkurl{colLev_0} character. The raw factor remains in the kernel on the left of (DO35.22).

Form the actual finite auxiliary vector

\[
 \phi_{\rm aux}
   =\bigotimes_{P\in G}f_{St,P}
       \otimes\bigotimes_{P\in T}f_{T,\nu_P}
       \otimes\bigotimes_{P\in Q}f_{Q,P}
       \otimes\bigotimes_{P\notin G\cup T\cup Q}\phi_{0,P}.
 \tag{DO35.23}
\]

Its finite norm squared is the exact product

\[
 \Gamma_{\rm aux}
  =\prod_{P\in G}\left(q_P^{-ab}\frac{|\operatorname{GL}_a(\kappa_P)|}{d_{G,P}}\right)
    \prod_{P\in T}(q_P-1)q_P^{-a}
    \prod_{P\in Q}\Gamma_{Q,P}.
 \tag{DO35.24}
\]

It is nonzero. Each factor is an actual finite Heisenberg combination of the local vacuum; their product is one such combination with pairwise distinct translated classes after like terms are combined. The previous arguments and DO033 show that the coefficient array is primitive at11. A valid bound for its number of terms is

\[
 s\le
  \prod_{P\in G}q_P^{ab}
  \prod_{P\in T}(q_P-1)q_P^{a-1}
  \prod_{P\in Q}q_P^{a(b-1)}.
 \tag{DO35.25}
\]

Every auxiliary rational prime here is assumed different from 11 and29. This is the exact hypothesis needed for the integral prime-to-11 translated-line comparisons in DO030. The empty Q set, supplied by the reference tuple, is explicitly allowed.

The raw finite vector with the unprojected G and T factors has norm squared

\[
 \Gamma_{\rm raw}
  =\prod_{P\in G}q_P^{-ab}d_{G,P}
    \prod_{P\in T}q_P^{-a}
    \prod_{P\in Q}\Gamma_{Q,P}.
 \tag{DO35.26}
\]

Unlike the selected G vector, the raw G norm has no GL-order denominator. If Q is empty, its finite norm has zero11 valuation. For nonempty Q the explicit factor (DO35.11) retains its possible 11 valuation. One must not equate primitive coefficients with unit Hilbert norm.

\subsection{6. Global compact extraction on the actual degree-seven host}\label{global-compact-extraction-on-the-actual-degree-seven-host}

Use the genuine host and selected highest polynomial of AM-DO-034, with the actual rank and stage characters supplied in AM-K031. In particular \nolinkurl{f=42}, \nolinkurl{a=42}, \nolinkurl{b=1802} is its \nolinkurl{m=2,t=1759} example. Keep the standard-lattice coefficient \nolinkurl{xi_K=xi_0 pi_B^(-6)} and its actual eigen-coordinate scalar. Let \nolinkurl{D} be the total selected degree and \nolinkurl{B_nu} the factorial valuation bound in DO034. Let

\[
 M=29^e\ge3s(D+1)-1.
 \tag{DO35.27}
\]

The simultaneous jet theorem DO030, in the proved host form DO034, applied to the nonzero primitive vector (DO35.23), gives one actual permitted half-orbit translate \nolinkurl{W_(x_lambda)} with

\[
 \theta_{P_\nu\otimes W_{x_\lambda}\phi_{\rm aux}}(1,1)
      =\mathcal P\,a_\lambda\ne0,\qquad
 v_{11}(a_\lambda)\le B_\nu,
 \quad \mathcal P=C^{14ab}.
 \tag{DO35.28}
\]

All gauges, the ordinary product line and the fixed root-of-unity translation phases are those already proved in DO030/034. The factor \nolinkurl{C} is the same analytic period, not a newly normalized value of the raw T or Q function.

Define \nolinkurl{K^e} to be the maximal source compact outside29 and the rational principal congruence subgroup of level \nolinkurl{29^e} at the six F places above29. Its exact index in the maximal compact is

\[
 I_a^K(M)=M^{42a^2}\prod_{i=1}^{a}(1-29^{-i})^6.
 \tag{DO35.29}
\]

The ramification index is seven and there are six residue places; these must not be replaced by forty-two unramified residue factors. DO034 proves that the finite rational stabilizer of this compact is trivial. The congruence fixes every relevant half-orbit vector; the other auxiliary places are disjoint from29.

On the identity finite class define the actual automorphic source function by

\[
 F(g_\infty,k)=
   \overline{m_\infty(g_\infty)}
   \prod_{P\in G}\overline{m_{G,P}(k_P)}
   \prod_{P\in T}1_{P_u}(k_P)\nu_P(s(k_P)),
 \quad k\in K^e,
 \tag{DO35.30}
\]

and set it to zero off that class. Here \nolinkurl{m_infty} is the unit highest source matrix coefficient used in DO034, and \nolinkurl{d_rho} is the dimension of that compact real representation. Triviality of the rational stabilizer proves this definition is well defined on the automorphic quotient, rather than only on a local chart. It is smooth at finite places, compact-finite at infinity, and square-integrable because the source unitary quotient is compact.

At G it has the initial inverse generic character; at T it has the literal mirabolic invariance and zero diamond sum from Section 4; at Q it is fixed by the full compact. These statements are about this actual global function. They do not follow merely from nonzero local factors. Its only extra source level is the displayed refinement at29. With maximal finite compact volumes one and compact real volume one, direct orthogonality gives

\[
 \|F\|^2
   =\frac1{I_a^K(M)d_\rho d_{G,\rm all}d_{T,\rm all}},
 \quad \|F\|_\infty=1,
 \tag{DO35.31}
\]

where \nolinkurl{d_(G,all)=product_G d_G} and \nolinkurl{d_(T,all)=product_T d_T}. No class-number-one assertion is used.

Now keep the \textbf{raw} G, T and Q factors in the theta kernel, and apply the same permitted29 translate to its vacuum there. The compact integrations (DO35.10), (DO35.18), (DO35.22) commute, and the real integration is the selected matrix-coefficient projection of DO034. Thus the literal bilinear lift has value

\[
 \boxed{
 \Theta_{\rm raw,G/T/Q}^{\rm bil}(F)(1)
   =\frac{A_G\,\mathcal P\,a_\lambda}
      {I_a^K(M)d_\rho d_{G,\rm all}
          \prod_{P\in T}(q_P^a-1)}\ne0,
 \qquad A_G=\prod_{P\in G}\alpha_{G,P}.}
 \tag{DO35.32}
\]

This is an equality for actual global automorphic and Schwartz inputs. It retains the raw local functions and the common character (DO35.3). The local G denominator and the T averaging denominator are not suppressed in a claim of nonzero pairing.

If the source and raw finite Schwartz vector are separately normalized to norm one, divide (DO35.32) by

\[
 \sqrt{\Gamma_{\rm raw}}\,
 \bigl(I_a^K(M)d_\rho d_{G,\rm all}d_{T,\rm all}\bigr)^{-1/2}.
 \tag{DO35.33}
\]

If the real Schwartz factor is normalized as well, divide by its exact norm in DO034, including \nolinkurl{det(2Y)^(-7ab)}, the hook factors, the eigen-coordinate metrics and \nolinkurl{|c_nu|^2}. No new real norm formula is being inferred from compact projection.

For fixed ranks and fixed G/T/Q data, (DO35.27) can be chosen with \nolinkurl{M<29(3s(D+1)-1)} (except the harmless \nolinkurl{e>=1} endpoint). Consequently the logarithm of (DO35.29) is \nolinkurl{O(log(D+1))}. All G/T/Q factors in (DO35.31)--(DO35.33) are fixed constants in that assertion. This does not remove the separate linear-degree coordinate denominator, Fock and factorial costs in DO034. If the T prime or the Q tuple varies with D, the displayed products, rather than this fixed-data simplification, apply.

\subsection{\texorpdfstring{7. What is and is not coupled to \nolinkurl{cst.M}, \nolinkurl{cAof} and \nolinkurl{cBof}}{7. What is and is not coupled to ,  and }}\label{what-is-and-is-not-coupled-to-cst.m-caof-and-cbof}

The first raw function at \nolinkurl{m2=0} contains no \nolinkurl{E1at} sum. Its G and Q local formulas are \nolinkurl{PsiAt} with both group arguments one. The integer \nolinkurl{m_c=cst.M} only controls the admitted Q target class and its period here. The choices (DO35.4)--(DO35.5) work for every \nolinkurl{m_c>=1}. The first-factor support is integral in X and periodic modulo q, so the first-factor clauses are compatible with

\[
 A\ge m_c,\quad B\ge m_c,\quad B>0,\quad
 D_G=1,\quad D_Q\ge\max(2m_c,3).
 \tag{DO35.34}
\]

For example \nolinkurl{m_c=A=B=1,D_Q=3} satisfies these displayed bounds. This is not an assertion that all group, character and later-row clauses of \nolinkurl{ChIdx.Ok} have been instantiated.

The exact integers entering the finite \nolinkurl{kerOf} expression are

\[
 c_A=\left(\prod_{P\in G\cup Q}q_P\right)^A,
 \qquad
 c_B=T(u)\left(\prod_{P\in G\cup Q}q_P\right)^B.
 \tag{DO35.35}
\]

Their exponents count the selected prime ideals; in particular a rational fully split G prime occurs f times. These integers are units at29 when all actual auxiliary places avoid29. Increasing \nolinkurl{cst.M}, A or B cannot insert a29 component into \nolinkurl{Sall} or into the actual half-position idempotent.

One must also avoid the converse false inference that every rational representative x in \nolinkurl{kerOf} is integral at29. \nolinkurl{homAt} in FS594 imposes conditions only at \nolinkurl{Sall}, so representatives may have denominators elsewhere. The finite quotient and the congruence/additivity identities are defined using that same localized module. Changing such a representative does not enlarge the actual idempotent support or change \nolinkurl{IsRatDatum}. The obstruction in the next section is about that actual finite input, not about the incidental denominator of a representative chosen in the displayed finite sum.

\subsection{8. The remaining literal datum and source-level equations}\label{the-remaining-literal-datum-and-source-level-equations}

For the degree-seven host, every29 place is absolutely ramified of index seven. A source Q place must be \nolinkurl{GoodPl}, which includes degree one and rational-uniformizer valuation one. Thus29 is not an eligible Q place. The source G primes are the large degree-one primes of K031, and the T conditions likewise require degree-one selected places and their explicit congruences. Hence the actual G/T/Q \nolinkurl{Sall} used here avoids29.

The literal idempotent predicate is \nolinkurl{IsIdemSum(A.e,Sall)} (LevChiEqOne62), and AM-DO-030 proved from \nolinkurl{IsIdemAt} that it has zero29 component. The finite input fixed by \nolinkurl{IsRatDatum} is therefore the raw G/T/Q vector with the standard vacuum at29. Equation (DO35.32) instead uses \nolinkurl{W_(x_lambda)} there, and its source function has the refined level \nolinkurl{K^e}. Neither difference is removed by Section 3's Q character cancellation.

There is a precise limitation on repair by actual lattice moves. At a relative split29 place, all source and target unitary actions in the Hom model preserve the position polarization. They act by an invertible linear change \nolinkurl{X -> h^(-1)Xg}, with a nonzero determinant-root scalar and positive Jacobian factor. A moved vacuum is consequently a nonzero scalar times the indicator of a linear lattice containing zero. The actual \nolinkurl{thMoved} family comes from those source and target lattice moves; its root variables in \nolinkurl{VacT} are determinant roots, not affine translation parameters. A nontrivial permitted pure position translate has indicator \nolinkurl{1_L(X+x)} with \nolinkurl{x notin L}; its support does not contain zero. Thus

\[
 \boxed{W_x\phi_0\ne c\,\omega(g,h)\phi_0
 \quad\text{for }x\notin L,\ c\ne0,}
 \tag{DO35.36}
\]

for every individual source/target lattice or frame move at that place. The statement includes nonintegral linear moves. Changing the selected class, the row frame, or a determinant root does not change this origin-support proof.

It does not exclude an equality of a single global theta value, a finite linear combination of group translates, nonvanishing after a compact projection, or a nonzero witness already having \nolinkurl{x_lambda=0}. These are different possible routes, and none is disproved by (DO35.36).

The remaining first-row assertion at maximal29 level can be written concretely. Let \nolinkurl{P_(29,0)} be compact averaging on the source29 factor, and let \nolinkurl{L_(G/T/Q,Pnu)} be the actual theta functional after the already computed G/T/Q and real filters. One needs, for an actual permitted global initial input at that level, a nonzero value of

\[
 L_{G/T/Q,P_\nu}\bigl(P_{29,0}\phi_{0,29}\bigr)
       =L_{G/T/Q,P_\nu}(\phi_{0,29}).
 \tag{DO35.37}
\]

The proved theorem supplies a nonzero value with \nolinkurl{W_(x_lambda) phi_(0,29)} and a smaller compact. A construction which removes the source stabilizer by a generic T decoration would allow a maximal29 class-supported input, but its integral against that translated kernel contains \nolinkurl{P_(29,0) W_(x_lambda) phi_(0,29)}. Nonzero unaveraged value does not imply nonzero averaged value. This is the exact projection gate; it is not settled by the finite nonvanishing of (DO35.23).

Likewise moving the auxiliary isogeny to another split ray prime does not alone solve (DO35.37): if that prime is outside \nolinkurl{Sall} the same issue recurs, while if it lies in Q its raw factor is the prescribed (DO35.8), not a freely selectable translate. The completed local identities provide the actual tests to use in a future local or global nonvanishing argument.

Even after such an analytic assertion is proved, identifying the source input with a specified integral \nolinkurl{MfL}/\nolinkurl{NlamC} row, or proving the literal \nolinkurl{PinCol}/\nolinkurl{MainX} mass equations, requires those actual lattice and scalar ties. The compact values and norms in this finding do not instantiate the full arithmetic records.

\subsection{9. Proof-check status}\label{proof-check-status}

Author whole-proof review and the separate finite exact diagnostic are recorded in \nolinkurl{definite_oscillator_CHECK_EVENTS.tsv} after the completed file has been read. No computational check is used as a substitute for the proofs above. No Lean build is claimed.

\clearpage

\section{AM-K033 --- explicit sharp selected content in the enlarged CM eigenbasis}\label{am-k033-explicit-sharp-selected-content-in-the-enlarged-cm-eigenbasis}

Author: GPT-6 Astra Pro, ALP\_MACHINE /root/kummer\_extension, 22 November 2026 UTC; inherited reasoning effort not separately exposed.

This finding identifies an explicit pair of CM labels on the actual \(J^{14}\) host, and on its weighted \(J^{56}\) extension, that retains the exact sharp primitive-kernel derivative. A fixed filler places this pair inside the full source-and-target selected Schur summand. The resulting image ideals in the actual permitted DP59 half-orbit have asymptotic slope \(1/1220\), with a bounded additive error before compact metric normalization.

The proof computes the full critical logarithmic form after the seven-copy eigenbasis change. It does not infer selected sharpness merely from an invertible change of the ambient tangent lattice. The accepted display and exponential theorems are used at their exact interfaces; their proofs are not rederived. The conclusion concerns the selected kernel lattice. A previously produced automorphic image or a separately imposed auxiliary finite G-vector can still impose a further restriction.

\subsection{1. Actual hosts, coordinates, and accepted primitive input}\label{actual-hosts-coordinates-and-accepted-primitive-input}

There are two cases:

\begin{longtable}[]{@{}
  >{\raggedleft\arraybackslash}p{(\columnwidth - 6\tabcolsep) * \real{0.2857}}
  >{\raggedright\arraybackslash}p{(\columnwidth - 6\tabcolsep) * \real{0.2143}}
  >{\raggedright\arraybackslash}p{(\columnwidth - 6\tabcolsep) * \real{0.2143}}
  >{\raggedleft\arraybackslash}p{(\columnwidth - 6\tabcolsep) * \real{0.2857}}@{}}
\toprule\noalign{}
\begin{minipage}[b]{\linewidth}\raggedleft
Parameter h
\end{minipage} & \begin{minipage}[b]{\linewidth}\raggedright
Host field
\end{minipage} & \begin{minipage}[b]{\linewidth}\raggedright
Principal variety
\end{minipage} & \begin{minipage}[b]{\linewidth}\raggedleft
CM type count f
\end{minipage} \\
\midrule\noalign{}
\endhead
\bottomrule\noalign{}
\endlastfoot
1 & \(K=E_0B\), \(E_0=\mathbb Q(\zeta_{28})\), \(B\subset\mathbb Q(\zeta_{29})\), \([B:\mathbb Q]=7\) & \(J^{14}\) & 42 \\
4 & \(K^*=K(\sqrt5,\sqrt{13})\) & \(J^{56}\) & 168 \\
\end{longtable}

The integral actions, principal product identifications and standard CM lattice coefficients are AM-K030 and AM-K032. Work at11 over a finite unramified complete coefficient DVR R large enough for the CM eigenvalues and torsion being used, with \(v(11)=1\). Each level can use a larger finite unramified extension; none of the bounds depends on its degree. The coefficient-field extension needed for an unrelated global etaVal is not being asserted unramified.

Use the K030 basis \(e_0,\ldots,e_6\) of \(\pi_B^{-3}\mathcal O_B\), with \[
 e_j=\sigma^j e_0,\qquad
 \sigma=(\zeta_{29}\mapsto\zeta_{29}^{\,2})|_B,\qquad
 U_{lj}=\sigma^l(e_j),\qquad U^{\mathsf t}U=UU^{\mathsf t}=I_7.
\tag{CS1}
\] Here indices l,j lie in \(\mathbb Z/7\mathbb Z\), and \(\pi_l=\sigma^l(\pi_B)\). All entries of U and its inverse and all \(\pi_l^{\pm1}\) are integral at11.

Let \(\mathcal R\) be a singleton when h=1. When h=4, let \(\mathcal R\) be the four embeddings of \(\mathbb Q(\sqrt5,\sqrt{13})\); for \(\rho=(\epsilon,\epsilon')\) put \[
 \delta_\rho=\delta_{1,\epsilon}\delta_{3,\epsilon'},
 \qquad
 u_a=\frac{a+\sqrt{a^2+4}}2,\quad
 \delta_a=\frac{u_a^{-1}}{\sqrt{a^2+4}},\quad a=1,3.
\tag{CS2}
\] For h=1 set \(\delta_\rho=1\). Every \(\delta_\rho\) is an11-unit.

The standard CM logarithmic coordinates will be denoted \[
            \zeta_{i,\varepsilon,l,\rho},
       \qquad i=1,2,3,\quad \varepsilon\in\{+,-\},
                \quad l\in\mathbb Z/7\mathbb Z,\quad\rho\in\mathcal R.
\tag{CS3}
\] For h=1, they are exactly \[
 \zeta_{i,\pm,l}=\pi_l^3\sum_{j=0}^{6}U_{lj}
                                  (u_{i,j}\pm i v_{i,j}).
\tag{CS4}
\] For h=4, apply in addition the two embedding matrices \[
 V_a=\begin{pmatrix}1&u_{a,+}\\1&u_{a,-}\end{pmatrix},
 \qquad
 V_a^{\mathsf t}\operatorname{diag}(\delta_{a,+},\delta_{a,-})V_a=I_2,
\tag{CS5}
\] on the four ordinary real tensor factors, as in K032(QH18--21). The same final \(\pi_l^3\) rescaling gives the standard lattice coordinates. All changes and their inverses are integral at11. They preserve the ordinary J label i.

For one ordinary J, the accepted DP49 decomposition and K018's actual source-gauge comparison give an integral normalized algebraic frame in which the reduced coefficient of a local section of level \(\kappa\) is \[
                   f\,U_0^\kappa\exp(-\kappa\beta).
\tag{CS6}
\] Here f is an arbitrary integral formal function, \(U_0\) is an integral unit, and \(\beta\) is the actual arbitrary-unit primitive logarithm. This formula is used on identical copies with the identical choice of ordinary tangent basis \(u_i\). Thus every copy has the same critical unit below. Choosing \(f=U_0^{-\kappa}\) is an actual integral formal-section choice. It does not assert that the fixed sigma section itself equals the pure kernel.

Put \[
                  s=\frac1{10},\qquad w=\frac{121}{1220},
                         c=s-w=\frac1{1220}.
\tag{CS7}
\] Choose a finite ramified scaling extension with \(v(\varpi)=w\), and put \[
                      \mu=\frac{\varpi^{122}}{11^{12}},
                           \qquad v(\mu)=s.
\tag{CS8}
\] Only this auxiliary Gauss-weight argument uses the ramified extension. The derivative values themselves remain in the original unramified coefficient ring.

\subsubsection{Exact accepted whole-series input}\label{exact-accepted-whole-series-input}

DP46(UI16--19), with h=3 and j=1, and the common-weight substitution of DP51 §6 give \[
 \mu^{-1}\beta(\varpi u_1,\varpi u_2,\varpi u_3)
        \text{ integral},\qquad
 \overline{\mu^{-1}\beta(\varpi u_1,\varpi u_2,\varpi u_3)}
                         =\theta\,u_3u_2^{121},
                   \quad \theta\ne0.
\tag{CS9}
\] This is the reduction of the \textbf{entire} scaled logarithm, not its first ordinary-degree term. To spell out the application: in DP46's contracted dual inverse, raising the later input \(y_2\) to the common weight w makes that input vanish in the reduction, since its profile weight is strictly smaller than w. The reduced inverse has \(Y_1=y_1\) and \(Y_2=-\bar a_1^{-1}y_1^{11}\); substitution into the complete minimum \(\bar a_2^{-1}S Y_2^{11}\) gives (CS9). The diagonal unit change from those display coordinates to the fixed ordinary CM eigenforms changes only the nonzero unit \(\theta\). Invariant integral frame changes are strictly above this critical form by DP46 §5. This is a direct use of the accepted full inverse-map and initial-form statements.

For an integral unit parameter \(\lambda\), the right side of (CS9) is multiplied by \(\bar\lambda\). In the actual tensor-level application \(\lambda=-\kappa\), where \(\kappa=M^2\), \(M=29^e\); hence it is always a unit.

\subsection{2. Frobenius computes the new critical pairs}\label{frobenius-computes-the-new-critical-pairs}

Bars denote residue-field reduction at11. For the unramified constants this is reduction modulo11; for the auxiliary ramified scaling extension it means reduction modulo its maximal ideal. On B, arithmetic Frobenius is the class of11 in \((\mathbb Z/29\mathbb Z)^\times/H\), where \(H=\{1,12,17,28\}\). One has \(11\in2^4H\) and \(121\equiv5\in2H\pmod{29}\). Therefore \[
                   \overline{U_{lj}}^{\,121}
                            =\overline{U_{l+1,j}},\qquad
                   \bar\pi_l^{\,121}=\bar\pi_{l+1}.
\tag{CS10}
\] This follows for the denominators29 as well because they are11-units. Equivalently, if P shifts the row l to l+1, then \(\bar U^{(121)}=P\bar U\).

On the two quadratic real fields, Frobenius squared is the identity, since their Galois groups have exponent2 and11 is unramified. Consequently \[
                   \bar V_a^{(121)}=\bar V_a.
\tag{CS11}
\] Also \(i^{121}=i\) modulo11.

Let \(\beta_{\rm prod}\) be the sum of the primitive logarithms on all14h ordinary J factors at one fixed matrix entry of a Hom variety. Since the principal polarization and the chosen invariant frame are products, this is the logarithm of the actual product primitive kernel in (CS6).

\subsubsection{Theorem CS12. Full critical form in standard CM coordinates}\label{theorem-cs12.-full-critical-form-in-standard-cm-coordinates}

For every unit \(\lambda\), in the coordinates (CS3), \[
 \boxed{\ 
 \overline{\mu^{-1}\lambda
       \beta_{\rm prod}(\varpi\boldsymbol\zeta)}
 =
 \frac{\bar\lambda\theta}{2}
 \sum_{\varepsilon,l,\rho}
       \bar\delta_\rho\,\bar\pi_{l+1}^{-6}\,
       \zeta_{3,\varepsilon,l+1,\rho}\,
                \zeta_{2,-\varepsilon,l,\rho}^{121}.\ }
\tag{CS12}
\] The left side denotes the same actual logarithm expressed in the new coordinates before the common scaling. Every displayed coefficient is a nonzero residue. All scaled coefficients before reduction are integral.

\textbf{Proof.} Start with (CS9) on the ordinary factors and sum them. On one pair of ordinary J copies write \(w_{i,\pm}=u_i\pm i v_i\). In characteristic11, since \(121\) is a Frobenius power and fixes i, \[
 u_3u_2^{121}+v_3v_2^{121}
  =\frac12\left(w_{3,+}w_{2,-}^{121}
                          +w_{3,-}w_{2,+}^{121}\right).
\tag{CS13}
\] For each sign the B-coordinate change is z=Uw, with inverse \(w=U^{\mathsf t}z\). Thus the coefficient pairing between the new B labels k,l is \[
 \sum_j\bar U_{kj}\bar U_{lj}^{121}
       =\sum_j\bar U_{kj}\bar U_{l+1,j}
       =\begin{cases}1&k=l+1,\\0&k\ne l+1.\end{cases}
\tag{CS14}
\] This uses the actual trace orthogonality, rather than a complex-unitary approximation or a generic invertibility assertion.

In the quadratic extension, the inverse-coordinate pairing is \[
  (\bar V_a^{-1})^{\mathsf t}
                 (\bar V_a^{-1})^{(121)}
   =\bar V_a^{-\mathsf t}\bar V_a^{-1}
   =\operatorname{diag}(\bar\delta_{a,+},\bar\delta_{a,-}),
\tag{CS15}
\] by (CS5) and (CS11). Taking the tensor product for the two quadratics introduces exactly \(\bar\delta_\rho\) and forces the same real embedding \(\rho\) on the paired labels. For h=1 this step is absent.

Finally \(z_{i,\varepsilon,l,\rho}=\pi_l^{-3}\zeta_{i,\varepsilon,l,\rho}\). On a nonzero pair the two factors contribute \[
              \bar\pi_{l+1}^{-3}
                        (\bar\pi_l^{-3})^{121}
                              =\bar\pi_{l+1}^{-6}.
\tag{CS16}
\] Equations (CS13)--(CS16) prove the displayed form.

All coordinate substitutions and inverses are integral at11, and the same scale \(\varpi\) is used on every input. Substitution therefore commutes with reduction and preserves integrality of the entire scaled logarithm in (CS9). At a fixed ordinary degree there are only finitely many coefficients involved. This verifies the whole-series assertion and rules out an unexamined logarithmic term at the same critical valuation. ∎

There are fourteen nonzero pairs when h=1 and fifty-six when h=4. In particular, fix l and \(\rho\), and set \[
                \tau_3=(3,+,l+1,\rho),\qquad
                \tau_2=(2,-,l,\rho).
\tag{CS17}
\] On the plane where only \(X=\zeta_{\tau_3}\) and \(Y=\zeta_{\tau_2}\) are nonzero, (CS12) is exactly one nonzero monomial \(\gamma XY^{121}\). Thus the critical-exponential lemma DP58(ES6--12) applies directly to this explicit new-CM plane.

For \[
                  A_r=11^r,\qquad B_r=121\,11^r,\qquad r\ge0,
\tag{CS18}
\] it gives, for the actual product primitive kernel \(F_\lambda=\exp(\lambda\beta_{\rm prod})\), \[
 \begin{split}
 v\bigl([\zeta_{\tau_3}^{A_r}\zeta_{\tau_2}^{B_r}]F_\lambda\bigr)
                  &=-\frac{121\,11^r-1}{10},\\
 \boxed{\ 
 v\!\left(
       \partial_{\tau_3}^{A_r}\partial_{\tau_2}^{B_r}
                      F_\lambda(0)\right)
                     =\frac{11^r-1}{10}.\ }
 \end{split}
\tag{CS19}
\] The first equality uses the full critical monomial, uniformly for every unit \(\lambda\). The second adds the two exact Legendre factorial valuations; both exponents are powers of11. Every derivative in (CS19) is nonzero.

This calculation locates the sharp primitive kernel in a particular pair of the enlarged CM eigenspaces. It does not enlarge the accepted singleton/complement display theorem to an arbitrary higher-height CM word.

\subsection{3. An explicit actual selected Schur family}\label{an-explicit-actual-selected-schur-family}

Fix an admissible source rank record and an edge with \[
                    a=N_j,\qquad b=N_{j+1}=a+d,
                         \qquad d=t+1\ge2.
\tag{CS20}
\] Choose fixed word-independent HC data on all f embeddings and let \(\nu_\tau^*\) be its actual selected positive partitions of length a, as in the literal identities checked in DO025. The numerical full-Hyp3 baselines of K031 and K032 are concrete choices on the two hosts. Put \[
 D_0=\sum_\tau|\nu_\tau^*|,\qquad
 \mathcal F_0=\prod_{\tau,i}\nu_{\tau,i}^*!,\qquad
                        c_0=v(\mathcal F_0).
\tag{CS21}
\] Ranks, the edge, the baseline, both fields and all stage characters remain fixed as r varies.

Reserve column1 of each a-by-b creation-symbol matrix \(Z_\tau\). Define on columns2 through a+1 \[
 Q_0(Z)=\prod_\tau\prod_{q=1}^{a}
    \det Z_{\tau,[1,q],[2,q+1]}^{
                                  \nu_{\tau,q}^*-\nu_{\tau,q+1}^*},
          \qquad\nu_{\tau,a+1}^*=0.
\tag{CS22}
\] Here \([2,q+1]\) means the consecutive q columns \(2,\ldots,q+1\). Each such column exists because b≥a+1. Set \[
 \boxed{\ 
 P_r=Z_{\tau_3,1,1}^{A_r}
             Z_{\tau_2,1,1}^{B_r}\,Q_0(Z),\qquad
 T_r=A_r+B_r+D_0=122\,11^r+D_0.\ }
\tag{CS23}
\] This polynomial is explicitly prescribed; no unknown distribution among the seven copies is selected.

Increase the first HC entry \(u_0\) by \(A_r\) at \(\tau_3\), by \(B_r\) at \(\tau_2\), and by zero at every other place. Make the same modifications for every word. The selected partitions on the fixed edge become exactly \[
 \nu_{\tau_3}=\nu_{\tau_3}^*+A_r e_1,\qquad
 \nu_{\tau_2}=\nu_{\tau_2}^*+B_r e_1,\qquad
                       \nu_\tau=\nu_\tau^*\text{ otherwise}.
\tag{CS24}
\] All e entries are unchanged. Thus the actual uniform characters supplied in K031/K032 stay the same, as do their compact G restrictions and conductors. The word row condition and the common first/last-u equations in Wts.Ok are preserved: the row data v,b,e do not change, and the first-u modification is common to every word.

Every HC inequality survives an increase of the largest u. The full numerical Hyp3 survives as well when the baseline satisfies it. The Dpar and NCGaps lists contain no u; the parity, eta and1801 clauses are unchanged. Interlacing and \(u_{\rm last}>-b_0\), together with \(e_j\ge e_0>b_0\), imply \(u_i+e_j>0\). Increasing \(u_0\) increases the absolute value in the only changing Centred inequalities. This checks all clauses, not only HC interlacing. Independently fixed upper SizeData windows are not preserved.

\subsubsection{Proposition CS25. Exact source-and-target summand}\label{proposition-cs25.-exact-source-and-target-summand}

The polynomial \(P_r\) is primitive integral and already belongs to the exact compact Cauchy summand indexed by the partitions (CS24), with the source and target determinant twists of DO025.

\textbf{Proof.} Prefix minors are invariant under adding an earlier row to a later row. Multiplication by a first-row variable retains this invariance. The row degrees telescope to (CS24). The product of the minor diagonal monomials and the two first-column powers has coefficient one, and no other determinant choices produce the same all-diagonal filler monomial. Thus \(P_r\) is primitive in the actual polynomial coefficient lattice.

Over characteristic zero, the Cauchy decomposition at each CM place is \[
       \operatorname{Sym}(T_\tau\otimes S_\tau^*)
          =\bigoplus_\lambda S_\lambda(T_\tau)\otimes S_\lambda(S_\tau)^*.
\tag{CS25}
\] In the source representation of a particular summand, its maximal-unipotent invariant line has precisely the source highest weight labelled by that \(\lambda\), and no other such invariant weight. Hence the source-invariant row weight \(\nu_\tau\) of \(P_r\) selects only \(\lambda=\nu_\tau\). The target vector can be non-highest, but it lies entirely in \(S_{\nu_\tau}(T_\tau)\), rather than being a sum of different target Schur types. Tensoring over the CM places and restoring DO025's fixed determinant twists proves the claim. No rational Schur projector or unspecified lattice denominator occurs. ∎

\subsection{4. Exact formal value and the permitted image ideal}\label{exact-formal-value-and-the-permitted-image-ideal}

On the Hom variety \[
            X=(J^{14h})^{ab}=J^{14hab},
                   \qquad \mathcal L=L_{2\infty}^{\boxtimes14hab},
\tag{CS26}
\] take the actual source/Chern reduced gauge, which is K018's full-CM gauge by K030/K032. Fix any allowed \(M=29^e\) and \(\kappa=M^2\).

At row1, column1 choose the integral formal prefactor in (CS6) that cancels the product of the normalized algebraic frame units. This produces the actual product primitive kernel with parameter \(-\kappa\). On the filler entries choose an integral algebraic monomial whose leading logarithmic term in the standard CM basis is \[
                    \prod_{\tau,i}
                         \zeta_{\tau,i,i+1}^{\,\nu_{\tau,i}^*}.
\tag{CS27}
\] Such an integral monomial exists because the new CM tangent basis is integral invertible at11: choose integral formal parameters with precisely those linear differentials. Its degree is D₀. On unused entries use normalized integral frame prefactors, giving reduced value1 at zero.

The first matrix entry and all filler entries are different product factors. The degree-D₀ filler derivative cannot see higher terms of its algebraic coordinates or of its reduced frame. The matching diagonal monomial has coefficient one in Q₀, and all other exponent vectors differ. Consequently its evaluation is \textbf{exactly} \(\mathcal F_0\); no factorial is inverted. Combining with (CS19) gives an actual integral formal section \(S_r\) such that \[
 \boxed{\ 
 v\!\left(
 M^{-T_r} P_r(\partial)F_{S_r}(0)\right)
                =\frac{11^r-1}{10}+c_0.\ }
\tag{CS28}
\] The unit \(M^{-T_r}\) is the source's normalized level reading. This is a value on an integral formal-section lattice; it is not yet a value of the fixed untranslated sigma section.

Let \(\mathfrak I_r\) be the image ideal of this normalized functional on the full integral formal \(T_r\)-jet module of \(\mathcal L^{M^2}\). Choose \[
      M_r=29^{e_r},\qquad
      e_r=\max\{1,\lceil\log_{29}(3T_r+2)\rceil\}.
\tag{CS29}
\] K031/K032 and the exact DP59/K022 half-orbit theorem give \[
 \mathcal W_{M_r}=[\alpha]^*H^0(X,\mathcal L^{M_r}),
 \quad\alpha=(1+2\sqrt{-7})^{e_r},\quad
 \operatorname{rank}\mathcal W_{M_r}=M_r^{fab},
\tag{CS30}
\] with its actual integral translated basis indexed by \(X[\alpha^2]/X[\alpha]\), and rescaled translations \(M_rt'\in X[\alpha]\). Its map to the integral \(T_r\)-jets is surjective. Therefore \(\mathfrak I_r\) is exactly the ideal generated by its actual allowed tuple values. From (CS28), \[
                           v(\mathfrak I_r)
                          \le\frac{11^r-1}{10}+c_0.
\tag{CS31}
\] This ideal is nonzero. A tuple attaining the least nonzero valuation exists because the basis is finite. It is one actual algebraic allowed tuple, although the formal interpolating coefficients may have been complete-local. Its finite Schwartz vector has norm one. The interpolation coefficient vector is not charged as an actual chosen input.

\subsection{5. A matching lower bound with constant error}\label{a-matching-lower-bound-with-constant-error}

Let \[
 K_0=\frac1{10}
       \max_{\beta\in\operatorname{supp}Q_0}
                                 \sum_m s_{11}(\beta_m).
\tag{CS32}
\] This is a finite fixed constant, independent of r and M. One computable bound is \[
                 K_0\le a\sum_{\tau,i}d_{11}(\nu_{\tau,i}^*),
 \qquad
 d_{11}(n)=
 \begin{cases}0&n=0,\\1+\lfloor\log_{11}n\rfloor&n>0.\end{cases}
\tag{CS33}
\] The accepted DP51 full common-weight bound, transferred to the actual integral local module in K027(EX2--5), applies to every section here. The new standard CM change is integral invertible, so that full bound applies in these coordinates as well.

Every monomial of \(P_r\) has the two fixed growing exponents Aᵣ and Bᵣ in columns disjoint from Q₀. Both have base11 digit sum1. The factorial digit loss is therefore exactly bounded by \(1/5+K_0\), rather than by a number growing with r. Thus \[
 \boxed{\ 
 \frac{T_r}{1220}-\frac15-K_0
 \ \le\ v(\mathfrak I_r)\
 \le\frac{T_r}{1220}
                  +c_0-\frac1{10}-\frac{D_0}{1220}.\ }
\tag{CS34}
\] The same lower bound holds for every integral section; the individual allowed tuple attaining \(\mathfrak I_r\) satisfies both bounds. All derivative values are Hurwitz-integral by K023, so the lower bound may also be replaced by its nonnegative integer ceiling.

In particular \[
             \boxed{\ 
             \lim_{r\to\infty}\frac{v(\mathfrak I_r)}{T_r}
              =\lim_{r\to\infty}\frac{v(a_r)}{T_r}
                                      =\frac1{1220}.\ }
\tag{CS35}
\] More precisely, \(v(\mathfrak I_r)-T_r/1220\) stays in the fixed interval displayed in (CS34). This proves sharpness in an \textbf{explicitly prescribed} full-CM selected family on either enlarged host. It is stronger than an ambient uniform-slope obstruction, and it has no hidden CM weight projection.

For the empty filler the same argument even gives \(v(\mathfrak I_r)=(11^r-1)/10\): the lower bound is less than this integer by only1/10. That empty-filler observation is not an assertion of the full source HC/Hyp3 conditions. The nonempty actual baseline filler is retained in (CS21)--(CS35).

For the concrete first-edge baselines already constructed in K031/K032, the constants are as follows. The last column concerns only comparison of the growing degree \(122\cdot11^r\) with D₀; it does not replace the large fixed upper error c₀.

\begin{longtable}[]{@{}
  >{\raggedright\arraybackslash}p{(\columnwidth - 8\tabcolsep) * \real{0.1579}}
  >{\raggedleft\arraybackslash}p{(\columnwidth - 8\tabcolsep) * \real{0.2105}}
  >{\raggedleft\arraybackslash}p{(\columnwidth - 8\tabcolsep) * \real{0.2105}}
  >{\raggedleft\arraybackslash}p{(\columnwidth - 8\tabcolsep) * \real{0.2105}}
  >{\raggedleft\arraybackslash}p{(\columnwidth - 8\tabcolsep) * \real{0.2105}}@{}}
\toprule\noalign{}
\begin{minipage}[b]{\linewidth}\raggedright
Actual host and fixed edge
\end{minipage} & \begin{minipage}[b]{\linewidth}\raggedleft
D₀
\end{minipage} & \begin{minipage}[b]{\linewidth}\raggedleft
c₀
\end{minipage} & \begin{minipage}[b]{\linewidth}\raggedleft
Computable bound for K₀
\end{minipage} & \begin{minipage}[b]{\linewidth}\raggedleft
First r with \(122\cdot11^r\ge D_0\)
\end{minipage} \\
\midrule\noalign{}
\endhead
\bottomrule\noalign{}
\endlastfoot
\(J^{14}\), 42→1802 & 3,504,831,920,337,006 & 350,483,192,023,230 & 899,640 & 13 \\
\(J^{56}\), 168→1802 & 194,467,812,219,985,248 & 19,446,781,221,824,712 & 60,484,032 & 15 \\
\end{longtable}

These are exact finite sums of the actual selected baseline parts and their factorial valuations, reproduced in the certificate. Thus the bounded-error asymptotic theorem retains a potentially substantial finite baseline cost. It does not itself certify a useful fixed numerical MainX window.

\subsection{6. Metric and first-image scope}\label{metric-and-first-image-scope}

The primitive target vector in the Cauchy summand has the exact metric correction already proved in K026: \[
 \Theta_r=
 \frac{A_r!\,H_{\nu_{\tau_3}^*}}
               {H_{\nu_{\tau_3}^*+A_re_1}}\,
 \frac{B_r!\,H_{\nu_{\tau_2}^*}}
               {H_{\nu_{\tau_2}^*+B_re_1}}.
\tag{CS36}
\] Here H is the hook product. Each quotient is a fixed finite product \(s/(n+s)\), so its real logarithm and its11-adic valuation have size \(O(\log(T_r+1))\), with the residual partitions fixed. Scaling the target vector to norm one multiplies the coefficient by \(\Theta_r^{-1/2}\); it is not treated as an integral unit. Hence its valuation slope is still \(1/1220\).

All real creation factors remain explicit. In standard coordinates put \[
        h_\tau=\frac{\delta_\rho\,\pi_l^{-6}}{2G_i},
           \qquad \tau=(i,\varepsilon,l,\rho).
\tag{CS37}
\] At the chosen complex realization these numbers are positive. The actual squared archimedean Schwartz norm of the primitive vector is \[
 \begin{split}
 \|r_{P_r}\phi_{0,\infty}\|^2
   ={}&\det(2Y)^{-7hab}\,
          A_r!B_r!\prod_\tau H_{\nu_\tau^*}\\
      &{}\cdot(\pi h_{\tau_3})^{A_r}
                (\pi h_{\tau_2})^{B_r}
                 \prod_\tau(\pi h_\tau)^{|\nu_\tau^*|}.
 \end{split}
\tag{CS38}
\] The two first-column powers and all filler variables are disjoint, so their Fischer norms multiply. The hook normalization is the exact parent R011/K026 normalization, and K030/K032 give the metric coefficients. This factorial archimedean cost is not an \(O(\log T_r)\) term. Dividing (CS38) by \(\Theta_r\) is an alternative target normalization of the same vector, not an additional independent loss.

For the rational29 principal class refinement, the exact index is \[
             I_a(M)=M^{fa^2}
                         \prod_{q=1}^{a}(1-29^{-q})^{6h}.
\tag{CS39}
\] Its11-adic valuation is independent of the exponent in M. If a class-supported extraction with the chosen finite compact type is supplied, its compact dimension factor varies by the two fixed-rank Weyl ratios \[
 \prod_{j=2}^{a}
 \frac{A_r+\nu_{\tau_3,1}^*-\nu_{\tau_3,j}^*+j-1}
                  {\nu_{\tau_3,1}^*-\nu_{\tau_3,j}^*+j-1}
 \quad\text{and the analogous factor with }B_r,\tau_2.
\tag{CS40}
\] These have11-valuation \(O(\log(T_r+1))\). Thus these specified metric/class corrections alone preserve the local slope. The fixed product period is \(C^{14hab}\) in the conventions of DO034, and remains separate.

There are two important restrictions on this conclusion. First, the value \(a_r\) in (CS35) comes from the unmodified allowed translated vacuum and the selected polynomial. DO034's finite Steinberg auxiliary vector, or LIFT20's literal raw G factor, defines a different functional. A separate simultaneous-germ argument is required before assigning that restricted functional the same sharp upper slope. The factorial nonvanishing theorem for an arbitrary auxiliary vector does not establish this sharper assertion. Second, the next automorphic stage acts on a previous theta image, which is a further integral sublattice. Neither (CS34) nor a nonzero theta-representation theorem determines its saturation or prevents cancellation after the class integral.

The new theorem therefore closes the enlarged-CM selected-kernel sharpness gap. It does not close the later-image arithmetic or the complete DP44 gain/error inequality. The common-type field obstruction Fields.Kq=\(\mathbb Q\) remains, including on the8-divisible host.

\subsection{7. Exact inputs and finite certificate}\label{exact-inputs-and-finite-certificate}

\begin{itemize}
\tightlist
\item
  Accepted DP46, unit\_arrow\_kernel\_invariance\_DP46.md, UI16--19 and §5: full scaled inverse maps, critical logarithmic initial form, and invariant-frame comparison. Accepted DP51, harmonic\_asymptotic\_rate\_DP51.md, §6 HA17--21: the common weight on every variable and full local-module bound. Accepted DP58, explicit\_sharp\_jets\_linear\_levels\_DP58.md, ES6--12: the critical-exponential coefficient theorem. These statements were read at their full interfaces; no new display proof is claimed.
\item
  AM-K018 K19.2/K19.13: exact source/Chern gauge on the ordinary product. K030 PH17--26 and K032 QH13--21: actual product principal lattices, full CM coordinates and integral11 comparisons. K031 HC22--26 and K032 QH33--36: the actual central alpha and permitted half-orbit.
\item
  AM-K022 J6--16/J30: ordinary jet surjectivity and equality of the permitted image ideals. AM-K024 SC9--19: filler and exact Cauchy-label argument. AM-K026 HC18--25 and parent R011: target hook comparison and Fischer normalization. AM-K027 EX2--7: accepted common-weight transfer to arbitrary joint integral local prefactors.
\item
  \nolinkurl{Formalism/LocalToGlobal/Columns/AbelianHostFieldConditions.lean:442–469,} Wts.Ok and selected weight definitions; \nolinkurl{Columns/BaseLemmas.lean:496–552,} selected recurrence; \nolinkurl{LocalRingCardinality/DataTheta.lean:529–555} and \nolinkurl{SharpRestriction/CuspSpanNormBounds/CayleySignature.lean:110–118,} exact Hyp3 clauses. Paths are relative to work/research/alpoge\_source. The old unmodified Fields.Kq condition is not changed.
\end{itemize}

The finite certificate \nolinkurl{selected_cm_frobenius_sharp_content_certificate.py} and its JSON output verify the actual seven-copy Frobenius permutation and all49 trace-pair entries over the cyclotomic residue algebra, the plus/minus pairing, both weighted quadratic pairings, the standard-lattice unit factors, the explicit support pairs, and exact Legendre and baseline arithmetic. They do not infer an infinite coefficient theorem or a selected automorphic image from finite samples. Those assertions, where made, are proved by the cited accepted whole-series theorems and the arguments above. No Lean execution is claimed.

\clearpage

\section{ALP-LIFT-22: the two prescribed sharp G kernels compose on their actual middle level}\label{alp-lift-22-the-two-prescribed-sharp-g-kernels-compose-on-their-actual-middle-level}

Author: GPT-6 Astra Pro, ALP\_MACHINE /root/lift\_accounting. Date: 22 November 2026 UTC. Inherited reasoning effort is not exposed.

This result closes a further local gate for the literal m=2 sharp coupled construction. The raw first G kernel does not require an intermediate generic projector, an arbitrary translate, or a new finite Schwartz choice before the second edge. Its actual image already has the middle level forced by ChIdx.Ok. The prescribed sharp terminal G kernel is injective on that image. The coupled Y-variable has only one possible G residue class, so no local sum over Y can cancel this component.

The result is local at G, with a fixed Haar normalization of the actual small-theta quotient models. It leaves the globally pinned quotient scalars, the other places, global Y-summation, and the exact arithmetic automorphic row lattice separate.

\subsection{1. The literal middle source level}\label{the-literal-middle-source-level}

Use LIFT21's notation \[
 m=2,\quad d=t+1,\quad r=k+1,\quad b=k+d,\quad N=b+d.
 \tag{CG1}
\] The place is degree one and relative split, with residue \(\kappa=\mathbf F_q\), \(q\ne11\), and the two actual tame characters \(\eta_0,\eta_1\) have distinct nontrivial residues. All source references below use the recovered commit 7cbf7a0c32fca1badf05931879a37c3f8a2b03b4. Let FS abbreviate \nolinkurl{Formalism/LocalToGlobal/Columns/FiniteSumNonzeroLargeDepth.lean.}

For m2=0 the coupled middle index is one. The source forces \[
             c_1=0,\qquad h_1=1.
 \tag{CG2}
\] The first is the literal cOf definition; the second is ChIdx.Ok.last\_h (FS1025--1026). Thus the following calculation does not choose a new intermediate frame.

When c=0, every Ufl in FS468 is the whole space and each extra coordinate vector eN(n,n+i) is zero. Therefore the conditions defining R1 are automatic: \[
  R1=\mathrm{GL}_b,\qquad N1=\{1\},\qquad n1Arg=0,
  \qquad \operatorname{prCat}(1,P,1)=\operatorname{Jat}(1,P).
 \tag{CG3}
\] The inclusion inX is the identity, proving the last equality from Cbe and prX. These are the exact definitions FS470--541, not a hypothesis about a new model.

For the actual sharp flag, Jat is JGsym. Since rS=0, symIdx is the identity. At c=0 the lattice flag in flagS is the usual residue prefix flag, except at index k: bd5(k,0,k)=0, so its lattice is \(qO^b\). The conditions in flagU therefore give the following exact residue subgroup: \[
 H^\sharp=
 \left\{\begin{pmatrix}u&A\\0&h\end{pmatrix}:
       u=\operatorname{diag}(u_k,1),\ u_k\in U_k(\kappa),
       A\in M_{k+1,t}(\kappa),\ h\in\mathrm{GL}_t(\kappa)
 \right\}.
 \tag{CG4}
\]

Here are the individual flag checks. For s\textless k, the condition on the s-th basis vector makes its diagonal entry one and its lower entries zero. At s=k it fixes the \emph{entire} column to e\_k because the modified lattice is qO\^{}b. Preservation of the prefix flag at r=k+1 gives the upper parabolic block shape. The next listed index is r+t=b and adds no condition; later indices also add none. Preservation of the ordinary lattice gives invertibility. Conversely every lift of(CG4) satisfies these conditions. Thus the local subgroup is the full inverse image \(J^\sharp\) of \(H^\sharp\). The flag definitions are FS738--766, with LevGSym.lean11--13 and its entrywise membership lemma108--141.

The grpG and chrG clauses of ChIdx.Ok (FS1072--1083) now force the actual middle levG group to be \(H^\sharp\). The n1Arg part of chiInv is one by(CG3); lamRow on this group is \[
        \psi_0\!\left(\sum_{i=0}^{k-2}(u_k)_{i,i+1}\right)
                                          \chi_0(\det h).
 \tag{CG5}
\] Hence the actual adelic middle type, with the inverse convention of RecLevel.lean50--52, is \[
 \lambda^\sharp=
        \psi_{U_k}^{-1}\,\bar\eta_0(\det h).
 \tag{CG6}
\] All character values are the actual residue values; the determinant inverse has not been dropped.

This middle type is not the generic ColV1 type. The source distinction is explicit: ColForms.lean105--114 defines MfV=ColV0, and FS1408--1417 imports that initial space by startC. The actual coupled kernel uses midGSD=(chF \ldots{} 1).mem (FS1220) and ChKer.cp on that space (FS1236--1237). In chF the group is FrmG of L.levG and L.levQ (FS1090--1102). It therefore uses the type(CG6), not a substituted ColV1. No assertion that the whole ChIdx.Ok record has been constructed is made.

\subsection{2. The coupled G residue class is forced}\label{the-coupled-g-residue-class-is-forced}

The allowed G classes of Y are \(q^{-1}O^b/O^b\). Before choosing any class, the first sharp raw kernel reads the matrix \([X,qy]\), by LIFT20(RG4). Its PhiGS condition fixes the last source column, which is column k in this specialization, to e\_k in every target row. Therefore \[
 \operatorname{PsiAt}_{\rm sharp}(X,y)\ne0
       \quad\Longrightarrow\quad qy=e_k\pmod q.
 \tag{CG7}
\] The diagonal condition fixes the k-th entry; the special sharp condition fixes all entries above it; the remaining lower and external-row conditions fix all entries below it. These include all b target coordinates.

Thus exactly one G residue class can contribute: \[
                         y=q^{-1}e_k\pmod{O^b}.
 \tag{CG8}
\] For that class the first kernel is exactly LIFT20's \[
 f_U(X)=
 \begin{cases}
    \psi_{U_k}(u),&X\in M_{b,k}(O),\
                       \bar X=[I_k;0]u,\ u\in U_k(\kappa),\\
    0,&\text{otherwise}.
 \end{cases}
 \tag{CG9}
\] A representative change by an integral vector does not affect either first or terminal G factor. The same allowed class is used on both coupled edges. A local coupled sum over Y therefore has only this G term. Distinct global Y classes can still differ at other places; this statement does not remove their global summation.

\subsection{3. The first raw output is an identity-parabolic section}\label{the-first-raw-output-is-an-identity-parabolic-section}

Let \(\pi_0\) be an actual initial constituent with unramified cuspidal support and a finite Steinberg type as in LIFT19/LIFT20. Put \[
 V_0=\pi_0^{J_0,\psi_{U_k}^{-1}},\qquad
 J_0=\operatorname{red}^{-1}U_k(\kappa),\qquad
 a_q=\operatorname{vol}(J_0)
                =\frac{|U_k(\kappa)|}{|\mathrm{GL}_k(\kappa)|}.
 \tag{CG10}
\] This is precisely the initial MfV local type. It contains the explicit nonzero Steinberg generic vector.

The full-rank integral of LIFT21(TG13)--(TG16), now on the edge k to b with character eta0, identifies the actual first small theta quotient with \[
                      \pi_1=
          \operatorname{Ind}_{Q_{(k,d)}}^{\mathrm{GL}_b}
                 (\pi_0\boxtimes\eta_0\det_d).
 \tag{CG11}
\] Its boundary is excluded exactly as in LIFT21: after removing eta0, every source inertia line is the nontrivial residue of eta0 or its inverse. The induction is irreducible by LIFT19. Thus the compact integral is a fixed Haar normalization of the actual quotient, with the same explicit global-normalization caveat.

Define \(s(v)\) in the compact induced model to be supported on the identity parabolic cell \(Q(O)K_b^1\), with \[
 s(v)\!\left(\begin{pmatrix}A&C\\0&D\end{pmatrix}\right)
                 =\pi_0(A)\eta_0(\det D)v
                         \qquad\text{on }Q(O).
 \tag{CG12}
\] Here v is \(K_k^1\)-fixed. The expression is well-defined on the compact cell and extends to normalized induction by Iwasawa. Its identity value is v.

The first fixed raw kernel has the exact image \[
                  {\cal T}_0(f_U,v)=a_q\,s(v)
                                      \qquad(v\in V_0).
 \tag{CG13}
\] To prove this, its residue columns span only the standard k-plane. The compact target section therefore vanishes on every other parabolic quotient point. At identity the source integral is supported on J0; the source character is trivial on its determinant, and \(\psi_{U_k}(u)\pi_0(u)v=v\). It equals \(a_qv\). Target parabolic covariance determines its values on the only remaining cell, yielding (CG13). No compact projector onto a different target type was applied.

The image s(v) has precisely(CG6). Indeed \(H^\sharp\subset Q_{(k,d)}(\kappa)\), and its Levi entries in that parabolic are u\_k and a block with determinant det h. Right translation of the identity-cell section thus acts by \(\psi_{U_k}^{-1}\bar\eta_0(\det h)\). Its support is preserved since this subgroup is contained in the parabolic. Therefore \[
                     s(V_0)\subset
                      \pi_1^{J^\sharp,\lambda^\sharp}.
 \tag{CG14}
\]

For precision, s(v) is invariant under the joining root \(1+cE_{k-1,k}\), which lies in the Q radical. A projector requiring a nontrivial generic character on that root would kill it. This is an obstruction only to imposing that different compact type on this particular output. The actual sharp coupled middle level is(CG6), so that projector is not part of the present source construction.

\subsection{4. The actual terminal raw kernel preserves this exact image}\label{the-actual-terminal-raw-kernel-preserves-this-exact-image}

Let \(\Phi^\sharp_y\) be the actual terminal PhiTop/ftY function with the same class(CG8). The sharp proof of LIFT21(TG19), (TG22), and(TG24) uses only the restricted character of \(J^\sharp\), not an extension to the larger generic group J. It consequently applies to every vector in \(\pi_1^{J^\sharp,\lambda^\sharp}\). Explicitly its identity integral is \[
          {\cal T}_1(\Phi_y^\sharp,w)(1)
                            =v_{J^\sharp}w,
     \qquad w\in\pi_1^{J^\sharp,\lambda^\sharp},
 \tag{CG15}
\] since on Jsharp the raw character times eta1's source determinant inverse is lambda-sharp inverse. On every graph-plane representative h\_z it is \(v_{J^\sharp}\psi_0(-z_k)w\), and it is zero on the other finite parabolic cells. The same boundary-rank proof identifies it with the actual second small theta quotient.

Combining(CG13)--(CG15) proves the exact composed identity \[
 \boxed{\quad
  \bigl[{\cal T}_1(\Phi_y^\sharp,
                 {\cal T}_0(f_U,v))(1)\bigr](1)
                         =a_q v_{J^\sharp}v
                          \qquad(v\in V_0).\quad}
 \tag{CG16}
\] There are two identity evaluations, in the nested induced models. Since both scalar volumes are nonzero, the two \emph{fixed raw} G kernels compose injectively on the whole initial generic type V0. No local Hecke rotation, new Schwartz input, or intermediate projection is required.

The residue class(CG8) is forced by the first raw kernel, so this is also the only possible local sharp coupled G term. This assertion does not replace the global coupled sum, its mass factors, or its other-place kernels.

\subsection{5. Exact norms and natural integral content}\label{exact-norms-and-natural-integral-content}

Put \(n_U=k(k-1)/2\) and use the Gaussian parabolic indices \({\cal J}_{u,v}(q)\) defined in LIFT21. The first identity-cell insertion satisfies \[
 \|s(v)\|^2={\cal J}_{k,d}(q)^{-1}\|v\|^2,\qquad
 \|f_U\|^2=q^{-bk}|U_k(\kappa)|.
 \tag{CG17}
\] The raw second factor has \(\|\Phi^\sharp_y\|^2=q^{-Nb+b}|H^\sharp|\). The exact composite induced-model norm is \[
 \boxed{
 \|{\cal T}_1(\Phi^\sharp_y,{\cal T}_0(f_U,v))\|^2
   =\frac{q^b(a_qv_{J^\sharp})^2}
                {{\cal J}_{k,d}(q){\cal J}_{b,d}(q)}
                                                \|v\|^2.}
 \tag{CG18}
\] The first norm is a single cell of the compact quotient; the second is the q\^{}b graph cells of LIFT21. This proves the formula for every vector in V0, not only a norm-one test vector.

If both finite raw kernels are unit-normalized, the squared scaling constant on V0 becomes \[
 \boxed{
 \frac{q^{bk}|U_k|}
          {|\mathrm{GL}_k(\kappa)|^2{\cal J}_{k,d}(q)}
 \cdot
 \frac{q^{Nb}|H^\sharp|}
          {|\mathrm{GL}_b(\kappa)|^2{\cal J}_{b,d}(q)}.}
 \tag{CG19}
\] No real selected-polynomial norm, global class mass, or automorphic quotient scalar is absorbed into this formula.

Fix an integral finite Steinberg lattice inside the actual \(\pi_0^{K_k^1}\), and use its natural finite induced lattices on both subsequent representations. Its initial generic intersection is saturated. The normalized first insertion \(a_q^{-1}{\cal T}_0(f_U,-)=s\) is integral, and identity evaluation returns the input. The normalized second map has only zero or root-of-unity multiples of that input among its finite graph values. Consequently \[
 (a_qv_{J^\sharp})^{-1}
        {\cal T}_1(\Phi_y^\sharp,{\cal T}_0(f_U,-))
 \tag{CG20}
\] is a primitive saturated injection of the initial generic lattice into the second natural induced lattice. The saturation proof is the identity-evaluation argument of LIFT21, applied twice; it uses no integral isotypic averaging.

With the valuation normalized by \(v_{11}(11)=1\), the exact scalar content change in these lattices is \[
 \begin{split}
 v_{11}(a_qv_{J^\sharp})
   &=-\sum_{i=1}^{k}v_{11}(q^i-1)
     -\sum_{i=t+1}^{b}v_{11}(q^i-1),\\
 a_qv_{J^\sharp}
   &=\frac{q^{k(k-1)+(k+1)t}
                       |\mathrm{GL}_t(\kappa)|}
           {|\mathrm{GL}_k(\kappa)|
                       |\mathrm{GL}_b(\kappa)|}.
 \end{split}
 \tag{CG21}
\] Thus this local raw two-edge map introduces exactly its displayed Haar denominator and no hidden additional content on the natural induced generic lattice. The comparison with the global arithmetic row lattice and its period normalization is still necessary before identifying this with a DP44 gain.

\subsection{6. Scope}\label{scope}

This theorem uses the actual sharp middle level forced by the source's coupled ChIdx structure, not the separate generic ColV1 level. It closes the two fixed G-edge interaction, including the local Y-variable, for the m=2 route. It applies place by place to every degree-one G prime of a repaired host once its actual first constituent and literal stage characters are supplied.

It does not construct the complete ChIdx.Ok record, prove global nonvanishing of the sum over all remaining Y data, calibrate the local quotient maps by one global pinned theta period, or establish the row-lattice comparison at11. The actual T/Q factors and the other places must be handled with their own tests. The original quadratic-type wrapper and other global constructor conditions are also outside this local calculation.

\subsection{7. Proof checks and references}\label{proof-checks-and-references}

The actual first raw factor is LIFT20. The full-rank quotient construction, published boundary theorem, and sharp terminal graph values are LIFT21. Their hypotheses and normalization limitations are retained. The new source readings are R1 at c=0, last\_h, grpG/chrG, the precise middle chF definition, and the distinction between that space and ColV1.

The finite diagnostic \nolinkurl{composed_raw_sharp_g_certificate.py} checks the forced class, the modified middle flag, first-output character, and the exact norm/content factors. The author has completed the whole-proof reading, including the exact middle source definitions, the unique G residue class, both actual quotient identifications, the identity-parabolic insertion, the larger sharp eigenspace, all composed norms, and the saturated lattice argument. The separate TSV records its time.

The finite diagnostic passed206 first-class tests and1,958 middle flag/character tests. It also checks that the different joining-root generic projector kills the particular identity section, while the actual sharp type fixes it. The displayed Haar content formula gives local composite valuations -10 at each of the two K031 G primes for k42,t1759; for the K032 prime pair and k168,t1633 it gives -73 and-184. These are local model scalars, not global DP44 gains. Small finite geometry cases are diagnostic, not full admissible Ranks examples.

\clearpage

\section{AM-K034 --- complete prime-to-11 torsion gauges and sharp selected auxiliary content}\label{am-k034-complete-prime-to-11-torsion-gauges-and-sharp-selected-auxiliary-content}

Author: GPT-6 Astra Pro, ALP\_MACHINE /root/kummer\_extension, 22 November 2026 UTC; inherited reasoning effort not separately exposed.

This finding proves the full formal-germ comparison that was deliberately absent from AM-DO-030. A prime-to-11 torsion translation has an integral horizontal difference-line frame, and this is the difference frame of the actual reduced theta translation. Consequently a fixed nonzero auxiliary finite Heisenberg operator preserves the complete selected local image ideal, up to the coefficient content of that operator. Combined with AM-K033 and simultaneous interpolation in the actual DP59 half-orbit, this transfers the sharp selected slope \(1/1220\) to the fixed finite Steinberg family of AM-DO-034. No integrality of a Chern-reduced ample-line frame, no theta-null unit, and no saturation of a later automorphic image is asserted.

\subsection{1. The two frames and their actual connections}\label{the-two-frames-and-their-actual-connections}

Let R be a complete DVR of mixed characteristic \((0,11)\), with \(v(11)=1\). The actual data below may first be defined over a number field and then completed at11. Finite coefficient extensions are allowed; no extension is claimed unramified when its extra scalar coefficients require ramification. Let \(X/R\) be an abelian scheme, \(\mathcal L\) a symmetric rigidified relatively ample line, and let its complex realization carry the translation-invariant positive Hermitian form H. In our application \(X\) is an ordinary power of J and \(\mathcal L\) is its product principal line.

Fix an integer \(\kappa\), a positive odd integer d prime to11, and a section \(t\in X[d](R)\). Denote translation by \(T_t\). The rigidified degree-zero difference line is

\[
 \mathcal Q_t=
 T_t^*(\mathcal L^\kappa)\otimes\mathcal L^{-\kappa}
       \otimes p^*(t^*\mathcal L^\kappa)^{-1},
 \qquad p:X\longrightarrow\operatorname{Spec}R.
\tag{TG1}
\]

The displayed rigidification uses the given one on \(0^*\mathcal L\). The polarization homomorphism and the rigidified theorem of the square identify \(\mathcal Q_t\) with the value at \(\kappa t\) of the Poincare family. Therefore it has a rigidification-preserving isomorphism

\[
                 \vartheta_t:\mathcal Q_t^{\otimes d}
                                     \overset{\sim}{\longrightarrow}\mathcal O_X.
\tag{TG2}
\]

This isomorphism is unique after its origin value has been fixed. Indeed two such isomorphisms differ by a global unit on the proper geometrically connected abelian scheme, hence by a unit of R, and the rigidification fixes it. We use that uniqueness, rather than an unspecified line-translation isomorphism, below.

\subsubsection{Lemma TG3. Integral torsion horizontal frame}\label{lemma-tg3.-integral-torsion-horizontal-frame}

There is a unique integrable connection on \(\mathcal Q_t\) for which (TG2) is horizontal. Its origin-normalized horizontal frame on the formal completion \(\widehat X_0\) is an integral frame. If q is an integral formal frame with origin value1 and

\[
                  a=\vartheta_t(q^{\otimes d})\in
                    1+\mathfrak m_0R[[z_1,\ldots,z_g]],
\tag{TG3}
\]

where \(\mathfrak m_0=(z_1,\ldots,z_g)\), then that frame is exactly

\[
                  q_{\mathrm{hor}}=a^{-1/d}q.
\tag{TG4}
\]

Here \(a^{-1/d}\) means the unique formal d-th root with constant1. It and its inverse belong to \(R[[z]]^\times\).

\textbf{Proof.} On a local integral frame q define \[
                        \nabla q=\frac1d\,d\log(a)\,q.
\tag{TG5}
\] This glues under a frame change \(q'=bq\), since \(a'=b^d a\). Because d is a unit and a is a local unit, this is an integral algebraic connection; it is integrable by the displayed formula. Any connection making (TG2) horizontal must have (TG5), proving uniqueness.

The recursive coefficient equation for \(b^d=a^{-1}\), \(b(0)=1\), has at each positive total degree the new coefficient multiplied by d and all other terms already determined in R. Since d is a unit, it has a unique solution in \(1+\mathfrak m_0R[[z]]\). This argument works in several variables and in ramified R. Differentiating \(b^d a=1\) gives \(\nabla(bq)=0\), proving (TG4). No division by a theta section or by one of its values occurs. ∎

\subsubsection{Lemma TG6. This is the connection in the actual reduced translation}\label{lemma-tg6.-this-is-the-connection-in-the-actual-reduced-translation}

Write the reduced translation of a coefficient function F as

\[
 U_t^{\kappa H}F(u)=
   \exp\!\left(-\pi\kappa H(u,t)
                 -\frac{\pi\kappa}{2}H(t,t)\right)F(u+t).
\tag{TG6}
\]

After normalizing the difference-line fibre at0, the difference frame between this reduced translated frame and the original reduced frame is exactly the horizontal frame of Lemma TG3. In particular its complete formal ratio to an integral difference-line frame is the unit in (TG4). This is equality of normalized connections and full germs.

\textbf{Proof.} Use the first-argument-linear convention for H. The original reduced coefficient multiplier is \[
 j_\kappa(\lambda,u)=\alpha(\lambda)^\kappa
       \exp\!\left(\pi\kappa H(u,\lambda)
                      +\frac{\pi\kappa}{2}H(\lambda,\lambda)\right).
\tag{TG7}
\] A substitution into (TG6) shows that its multiplier is \[
 j_\kappa(\lambda,u)\,\chi_t(\lambda),\qquad
 \chi_t(\lambda)=
       \exp\!\left(\pi\kappa\{H(t,\lambda)-H(\lambda,t)\}\right).
\tag{TG8}
\] The second factor is constant in u and has absolute value1. Since \(dt\in\Lambda\) and \(\operatorname{Im}H\) is integral on \(\Lambda\), one has \(\chi_t(\lambda)^d=1\).

The reduced frame of the ample line has squared norm \(\exp(-\pi\kappa H(u,u))\). In the translated frame specified by (TG6), its metric is the same function: the frame multiplier is \(\exp(\pi\kappa H(u,t)+\pi\kappa H(t,t)/2)\), whose absolute square cancels the cross and constant terms in \(\exp(-\pi\kappa H(u+t,u+t))\). The quotient frame of the rigidified line (TG1) consequently has constant norm1. It is horizontal for its flat Chern connection.

Moreover its d-th tensor power is a global holomorphic nowhere-vanishing frame of \(\mathcal Q_t^d\), since \(\chi_t^d=1\). Compare it with (TG2). Their quotient is a global holomorphic unit on the compact complex torus, hence a constant; the origin rigidification makes that constant1. Thus the actual reduced quotient frame is horizontal for exactly the connection making (TG2) horizontal. By the uniqueness in Lemma TG3 it is that connection, including its normalized formal frame.

For completeness, this comparison does not require a possibly nonconvergent integral formal frame to be evaluated over the complex numbers. Compare the algebraic torsion connection and an algebraic reference local frame after complex base change. The equality just proved descends to their finite algebraic jets. Complete at11 and then change to any integral formal frame; (TG3)--(TG5) give the resulting complete formal equality. ∎

The formal root is integral in \textbf{integral algebraic parameters}. Its expression in invariant logarithmic coordinates need not have integral divided Taylor coefficients. Similarly the original ample-line reduced frame is not claimed integral. Lemma TG6 concerns only the prime-to-11 torsion difference line and does not contradict the THETA3 C11/C19 counterexamples.

\subsubsection{Actual fibre normalization}\label{actual-fibre-normalization}

The origin normalization in Lemma TG6 must be compared with the one in the actual finite Heisenberg reading. AM-DO-030, section2, supplies this comparison: pull back by an integer prime to11 killing the torsion, use the relative symmetric theta-group lift of accepted DP62, and descend through the resulting formal etale isomorphism. Its finite-order scalar is an11-unit. One may use \(2d\) instead of d throughout this auxiliary comparison; both are11-units, and this also matches the \(2d\)-divisibility in the primary Mumford construction. It changes no level or source parity condition.

Thus the actual scalar is an11-unit, rather than an unrecorded nonzero scalar. The complete connection comparison in Lemma TG6 is the additional argument beyond DO030's fibre statement. There is no assumption that \(t\in K(\mathcal L^\kappa)\), and no claim that translation by t acts globally as a theta-group automorphism of \(\mathcal L^\kappa\).

The precise primary comparison is Bannai and Kobayashi (2010), author preprint arXiv:math/0610163v4, section1.3, formulas (14) and the definition of \(U^M_{v_0}\), PDF pp.14--15; Proposition2.3, PDF pp.21--22, with the abelian-scheme observation following Corollary2.4; Lemma2.7 and Proposition2.8, PDF pp.24--25. The latter identify this normalized reduced translation with Mumford's algebraic translation. No ordinary-reduction integrality theorem from that paper is used here.

\subsection{2. Equality of the complete translated coefficient modules}\label{equality-of-the-complete-translated-coefficient-modules}

Now take the actual ordinary J-product, with the full source gauge already identified in AM-K018. Write its accepted primitive logarithm as \(\beta_{\mathrm{prod}}\), and let \(U_0\) denote the integral normalized frame factor. In an actual integral formal frame, the complete reduced coefficient module of \(\mathcal L^\kappa\) at0 is

\[
 \mathscr C_\kappa=
 U_0^\kappa\exp(-\kappa\beta_{\mathrm{prod}})
                    R[[z_1,\ldots,z_g]].
\tag{TG9}
\]

All series are compared through the same invariant tangent coordinates. The coordinates z inside \(R[[z]]\) are algebraic formal parameters; the logarithmic variables used for differentiation are distinguished from them.

\subsubsection{Proposition TG10. Translated modules agree}\label{proposition-tg10.-translated-modules-agree}

For each prime-to-11 torsion point t, take every integral formal section near t, translate it to0 by the \textbf{actual} reduced operator (TG6), and include the finite Heisenberg fibre scalar fixed above. The resulting complete coefficient module is exactly \(\mathscr C_\kappa\).

\textbf{Proof.} Translation identifies the formal neighbourhood of t with \(\widehat X_0\) over R. Choose integral frames of \(\mathcal L^\kappa\) near0 and t, and a primitive integral vector in \(t^*\mathcal L^\kappa\). Their difference is an integral frame q of (TG1), up to an integral unit. By Lemma TG6 its reduced difference frame is \(a^{-1/d}q\); by the actual fibre comparison its remaining constant is an11-unit. Hence the translated reduced coefficient of an arbitrary integral local section is \[
 c_t\,h_t(z)\,
        U_0^\kappa\exp(-\kappa\beta_{\mathrm{prod}})\,f(z),
       \qquad c_t\in R^\times,\quad h_t\in R[[z]]^\times,
       \quad f\in R[[z]].
\tag{TG11}
\] Every f occurs. Multiplication by \(c_th_t\) is an automorphism of the complete formal function module and of every finite jet quotient. This proves equality in (TG9), not only an inclusion or equality of their leading terms. The signs of \(h_t\) depend on which integral frame is called q; either sign is an integral unit, and the normalized connection comparison has already fixed the actual reduced operator. ∎

This application uses the actual Chern/source-gauge equality of K018 on the \textbf{untranslated} ordinary product, and the explicit reduced translation (TG6) on the translated line. It does not invoke CM invariance at the translated point, which is generally false. It does not replace the integral connection in Lemma TG3 with an arbitrary holomorphic connection on a torsion line.

\subsection{3. Exact auxiliary selected image ideal}\label{exact-auxiliary-selected-image-ideal}

Use the enlarged hosts of AM-K033: h=1 or4, \(f=42h\), with actual principal variety \(A=J^{14h}\). Fix a source edge a to b, the fixed residual selected partitions, the explicit polynomials \(P_r\), and degrees

\[
                  T_r=122\,11^r+D_0
\tag{TG12}
\]

from K033(CS20--24). The full numerical Hyp3 preservation, exact source-and-target Cauchy summand, and growing first-row weights remain exactly those proved there.

Fix a nonzero auxiliary finite operator supported away from11 and29, in DO030's combined form

\[
 \mathcal O=\sum_{j=1}^{s}b_j W_{y_j},\qquad
 b_j\ne0,\qquad
 \mathcal O\phi_0\ne0,
 \qquad t_j\in X[d]\ \text{distinct},
 \quad X=A^{ab}.
\tag{TG13}
\]

Here d is odd and prime to \(11\cdot29\). The distinct classes are the classes after DO030's actual lattice phases have been combined. Set

\[
 \eta=\min_j v(b_j),\qquad
                    \mathfrak b=(b_1,\ldots,b_s)\subset\operatorname{Frac}(R).
\tag{TG14}
\]

The coefficients and operator are fixed as r varies. Put \(M=29^e\), \(\kappa=M^2\), and \(\alpha=(1+2\sqrt{-7})^e\). Choose

\[
                   M\ge 3s(T_r+1)-1 .
\tag{TG15}
\]

The exact DP59 lattice is \[
 \mathcal W_M=[\alpha]^*H^0(X,\mathcal L^M)
                 \subset H^0(X,\mathcal L^{M^2}),
\tag{TG16}
\] with its actual permitted integral basis \(\mathcal S_\lambda\). No full level-\(M^2\) orbit is substituted for it.

Let \(\ell_r=M^{-T_r}P_r(\partial)\vert_0\). Define the finite fractional ideal \[
                 \mathfrak J_{r,\kappa}=\ell_r(\mathscr C_\kappa)
\tag{TG17}
\] using full integral \(T_r\)-jets. Define \(\mathfrak I^{\rm aux}_{r,M}\) to be the ideal generated by the actual allowed auxiliary tuple values \[
 a^{\rm aux}_{r,\lambda}
 =M^{-T_r}P_r(\partial)
       \left(\sum_j b_j\varepsilon_{j,M}
                   U_{z_j}^{M^2H}F_{\mathcal S_\lambda}\right)(0),
       \qquad z_j=[M]^{-1}t_j .
\tag{TG18}
\] The \(\varepsilon_{j,M}\) are exactly the11-unit finite reading phases of DO030(11--13). Pullback notation in (TG18) uses the section-level functions, as in DO030(13); it is not a second multiplication by M.

\subsubsection{Theorem TG19. Exact content multiplication}\label{theorem-tg19.-exact-content-multiplication}

\[
 \boxed{\quad
           \mathfrak I^{\rm aux}_{r,M}
                 =\mathfrak b\,\mathfrak J_{r,\kappa},
       \qquad
           v(\mathfrak I^{\rm aux}_{r,M})
                 =\eta+v(\mathfrak J_{r,\kappa}).
       \quad}
\tag{TG19}
\]

\textbf{Proof.} By Proposition TG10 the full integral \(T_r\)-jet coefficient module at each translated point \(z_j\), read at0 with its actual reduced operator, is exactly \(\mathscr C_\kappa\). Translation keeps the invariant tangent basis; the multiplication maps and alpha are etale at11, so their maps on formal parameters are integral isomorphisms. Thus the image ideal of the j-th summand on its independent jet module is \(b_j\varepsilon_{j,M}\mathfrak J_{r,\kappa}=b_j\mathfrak J_{r,\kappa}\).

The points \(q_j=\alpha z_j=\bar\alpha^{-1}t_j\) on the base of (TG16) remain distinct modulo11. DO030's simultaneous integral jet theorem, valid on these ordinary products by DO034 and the host constructions, states that (TG15) makes the map \[
  \mathcal W_M\longrightarrow
          \bigoplus_{j=1}^{s}
             J^{T_r}_{z_j}(\mathcal L^{M^2})
\tag{TG20}
\] surjective. It is obtained from the actual \(H^0(\mathcal L^M)\) by alpha pullback and the explicit \(3s(T_r+1)-1\) interpolation bound. The map is over R; formal complete-local interpolation coefficients are sufficient for equality of image ideals.

On the independent direct sum, the image of the sum functional is the sum of the individual image ideals. Surjectivity therefore gives \[
 \mathfrak I^{\rm aux}_{r,M}
       =\sum_j b_j\mathfrak J_{r,\kappa}
       =(b_1,\ldots,b_s)\mathfrak J_{r,\kappa}.
\] This proof also shows explicitly why potential cancellation among different auxiliary terms cannot force extra content: prescribe a desired jet at one point and zero jets at all the other points. It is not an assertion that the individual global theta functions have disjoint support. Since R is a DVR, valuation of the product is the sum, proving (TG19). ∎

The finite list of basis values generates this ideal. Its minimum valuation is attained by one permitted tuple, whose actual finite vector is \(\mathcal O\phi_\lambda\). Selecting that tuple does not require charging the norm of an interpolating coefficient vector. The vector has the exact fixed norm \[
                \|\mathcal O\phi_\lambda\|_f^2
                               =\sum_j|b_j|^2,
\tag{TG21}
\] by DO030's disjoint-place scaled isometry.

\subsection{4. Sharp selected bounds after the auxiliary restriction}\label{sharp-selected-bounds-after-the-auxiliary-restriction}

K033's proof of the formal ideal bounds is uniform for every unit \(\kappa=M^2\), not just its originally chosen least interpolation level. Therefore (CS34) and (TG19) give

\[
 \boxed{\quad
 \eta+\frac{T_r}{1220}-\frac15-K_0
 \ \le\ v(\mathfrak I^{\rm aux}_{r,M})\
 \le\eta+\frac{T_r}{1220}
                  +c_0-\frac1{10}-\frac{D_0}{1220}.
 \quad}
\tag{TG22}
\]

Every basis value satisfies the lower bound, and a permitted nonzero basis value attains the ideal valuation and hence satisfies both bounds. In particular

\[
 \boxed{\qquad
 \lim_{r\longrightarrow\infty}
      \frac{v(a^{\rm aux}_{r,\lambda(r)})}{T_r}
              =\frac1{1220}.
 \qquad}
\tag{TG23}
\]

The bounded additive error includes \(\eta\) and the full fixed baseline constants from K033. The conclusion concerns the \textbf{same explicit selected family} \(P_r\) in the full enlarged CM basis. The auxiliary operator is fixed, actual, and nonzero. There is no enlargement of the harmonic subspace, no appeal to a unit in the ambient unselected space, and no restriction to the degree-\(T_r\) leading algebraic term.

The fixed auxiliary coefficients may be fractional; (TG22) retains their exact minimum valuation. If they are all integral and one is a unit, then \(\eta=0\). If a scalar normalization of the finite vector is imposed, its valuation is added explicitly. These statements use the Hurwitz derivative functional \(\ell_r\), rather than the divided Taylor coefficient. The factorial shifts are already those in K033.

\subsection{5. The actual finite Steinberg family and its costs}\label{the-actual-finite-steinberg-family-and-its-costs}

For the h=1 host, AM-DO-034 section5 supplies the finite Steinberg operator \(\mathcal O_G\), with fixed ranks a,b and the actual stage characters of K031. Its coefficients have minimum11-valuation0. It satisfies

\[
 \begin{aligned}
 d_G&=(g_0g_1)^{21a(a-1)},\\
 \Gamma_G:=\|\mathcal O_G\phi_\lambda\|_f^2
   &=(g_0g_1)^{42(-ab+a(a+1)/2)}
         \prod_{\epsilon=0}^{1}
           \prod_{i=1}^{a}(1-g_\epsilon^{-i})^{42},\\
 s_G&\le(g_0g_1)^{42ab}.
 \end{aligned}
\tag{TG24}
\]

Thus (TG22)--(TG23) apply to this actual prescribed finite compact type, with \(\eta=0\). The first-row modifications in K033 leave the e entries, finite stage characters, ranks, G-primes, and the operator itself unchanged. This fixed-operator hypothesis is essential for the asymptotic statement.

The proof of TG19 works without change on the h=4 host. When the analogous local finite Steinberg construction is made with its actual168-place record and G-primes, its parameters are the corresponding \(f=168\) local products. This local-product substitution does not silently claim that every global first-image assertion in the h=1 DO034 file has already been reconstructed for h=4.

There is a simple completely explicit sufficient level for either local product. Put \(Q=g_0g_1\), choose the smallest positive integer L with \(Q\le29^L\), and define

\[
 e_r=Lfab+\left\lceil\log_{29}\bigl(3(T_r+1)\bigr)\right\rceil,
                         \qquad M_r=29^{e_r}.
\tag{TG25}
\]

Here the ceiling means the least integer k with \(29^k\ge3(T_r+1)\); it can be computed using integer arithmetic. Then \(M_r\ge3s_G(T_r+1)\), proving (TG15). This bound intentionally keeps the auxiliary orbit size:

\begin{longtable}[]{@{}
  >{\raggedright\arraybackslash}p{(\columnwidth - 10\tabcolsep) * \real{0.1304}}
  >{\raggedleft\arraybackslash}p{(\columnwidth - 10\tabcolsep) * \real{0.1739}}
  >{\raggedleft\arraybackslash}p{(\columnwidth - 10\tabcolsep) * \real{0.1739}}
  >{\raggedleft\arraybackslash}p{(\columnwidth - 10\tabcolsep) * \real{0.1739}}
  >{\raggedleft\arraybackslash}p{(\columnwidth - 10\tabcolsep) * \real{0.1739}}
  >{\raggedleft\arraybackslash}p{(\columnwidth - 10\tabcolsep) * \real{0.1739}}@{}}
\toprule\noalign{}
\begin{minipage}[b]{\linewidth}\raggedright
Host baseline
\end{minipage} & \begin{minipage}[b]{\linewidth}\raggedleft
a
\end{minipage} & \begin{minipage}[b]{\linewidth}\raggedleft
b
\end{minipage} & \begin{minipage}[b]{\linewidth}\raggedleft
\(Q=g_0g_1\)
\end{minipage} & \begin{minipage}[b]{\linewidth}\raggedleft
L
\end{minipage} & \begin{minipage}[b]{\linewidth}\raggedleft
Fixed term \(Lfab\) in \(e_r\)
\end{minipage} \\
\midrule\noalign{}
\endhead
\bottomrule\noalign{}
\endlastfoot
K031, \(J^{14}\) & 42 & 1802 & \(4615409\cdot5276377\) & 10 & 31,787,280 \\
K032, \(J^{56}\), analogous local product & 168 & 1802 & \(22285932761\cdot64072175441\) & 15 & 762,894,720 \\
\end{longtable}

The large fixed exponent is a sufficient interpolation cost, not an assertion that the least useful level is this large. It is \(O(1)+O(\log T_r)\) for fixed source ranks and G-primes, but it must not disappear from a concrete MainX window.

For these two exact local products the certificate gives respectively \[
 (v_{11}I_a,\ v_{11}\Gamma_G,\ v_{11}d_G)
       =(24,336,0)\quad\hbox{and}\quad(408,21336,0).
\tag{TG25a}
\] The second tuple uses the24 actual real-field29 places of K032, rather than replacing their residue exponent by168. These are fixed finite normalization costs; the h=4 tuple by itself is not a new global extraction theorem.

On the h=1 host, the actual class-supported first extraction of DO034(25) now has a selected tuple with (TG22), and

\[
 \Theta_{\mathcal O_G\phi_\lambda,r_{P_r}\phi_\infty}(F_r)(1)
       =\frac{C^{14ab}a^{\rm aux}_{r,\lambda}}
                    {I_a(M)\,d_{\rho,r}\,d_G},\qquad
 \|F_r\|_2^2=\frac1{I_a(M)d_{\rho,r}d_G}.
\tag{TG26}
\]

The target vector is the primitive vector of K033(CS25), with its exact hook metric (CS36); the source highest vector used in the compact matrix coefficient is the matching source factor. The Cauchy argument already puts this possibly non-highest target vector in the single selected target summand. Compact extraction does not require it to be a target-highest vector.

Here \[
 I_a(M)=M^{42a^2}\prod_{i=1}^{a}(1-29^{-i})^6.
\tag{TG27}
\] Its11-valuation is independent of M; \(v_{11}(d_G)=0\). The changing compact dimension \(d_{\rho,r}\) has valuation \(O(\log(T_r+1))\), by the exact fixed-rank Weyl ratios in K033(CS40). Thus, after removing the displayed fixed period and retaining the stated algebraic normalization, the actual first coefficient in (TG26) has valuation slope \(1/1220\). Unit source and finite Schwartz normalization replaces the denominator by \(\sqrt{\Gamma_G}\sqrt{I_a(M)d_{\rho,r}d_G}\), exactly as in DO034(27); the fixed \(\Gamma_G\) and these compact corrections also preserve that slope.

This does not discard the archimedean Schwartz norm. It remains exactly K033(CS38), including \(A_r!B_r!\), the residual hook products, and the positive CM metric factors. Those real factorial costs are not \(O(\log T_r)\), and normalizing the full archimedean Schwartz vector changes the real cost in the complete gain/error accounting. The complete DP44 inequality is not inferred from the local slope.

The theorem applies to any other \textbf{specified nonzero fixed} auxiliary vector that has an actual finite Heisenberg expansion of the form (TG13). In particular it can be used with a literal raw G factor once its fixed expansion and the relevant compact extraction have been supplied. It does not identify a raw factor with the finite Steinberg projection by definition, and it does not assert that a later theta image contains this entire selected lattice.

\subsection{6. Scope, checks, and references}\label{scope-checks-and-references}

The classical finite-torsion flat-frame argument in Lemmas TG3/TG6 is written in full to fix the normalized connection, the integral root, and the actual reduced multiplier. The new application is the exact complete-module comparison TG10, the actual permitted auxiliary ideal identity TG19, and its sharp selected consequence TG22 on the enlarged CM family.

Carried accepted inputs are DP47/49/51/58/59/62 at the interfaces already specified in K018, K022, and K033. Exact local/source identifications are K018(K19.2/K19.13), K033(CS9/CS12/CS19/CS25/CS28/CS34), DO030(9--20), and DO034(12--29). No Lean execution is claimed. The simultaneous jet theorem and the accepted display bounds are used, not rederived.

Primary literature: Bannai, K., \& Kobayashi, S. (2010). Algebraic theta functions and the p-adic interpolation of Eisenstein--Kronecker numbers. \emph{Duke Mathematical Journal, 153}(2), 229--295. Author preprint: \href{https://arxiv.org/pdf/math/0610163}{arXiv:math/0610163v4}, exact locators given in section1. Those primary passages concern the reduced multiplier and algebraic Mumford translation; no positive-characteristic ordinary hypothesis from later sections is imported.

The companion exact certificate checks formal torsion-root and connection signs in finite rational jets, including the excluded d=11 case, and computes the displayed sufficient level constants and the actual fixed finite11-costs. These finite calculations check the indicated identities and arithmetic only. The infinite full-germ and image-ideal assertions are the proofs above, not inferences from a finite series sample.

\clearpage

\section{AM-R-018: complete numerical Sizes for the explicit fixed-prime selected families}\label{am-r-018-complete-numerical-sizes-for-the-explicit-fixed-prime-selected-families}

Author: /root, GPT-6 Astra Pro, ALP\_MACHINE, 22 November2026UTC. Reasoning effort is not externally exposed.

This result checks the remaining spread condition left explicit in R017. It gives all four clauses of the literal numerical \nolinkurl{Sizes} record on the K031/K032 baselines and on finite, explicitly bounded portions of the K033 growing selected family. It does not construct \nolinkurl{Dat.Ok}, the global \nolinkurl{SizeData₂}, or a net lift saving. In particular the common-type-field obstruction remains. The fixed T-primes cannot support an unbounded growing-weight family under this record.

\subsection{1. Literal source spread, including the rank term}\label{literal-source-spread-including-the-rank-term}

Use the actual m=2 data of K031/K032. Thus

\[
 d=t+1,\quad N_j=k+jd,\quad
 u_i=(2(k-i)+1)C+\tfrac12,\quad
 e_0=0,\quad e_1=4C.
 \tag{R18.1}
\]

Here k is42 or168, d is1760 or1634, and C is respectively 45156049293 or40530339687. The selected dominant weights are the actual \nolinkurl{lamTrue}, with their source determinant shifts. Write these weights as \(w_j\). The checked source recurrence is

\[
 \begin{aligned}
 w_{0,i}&=(2(k-i)+1)C+1-k/2+i,\\
 w_{1}&=(w_0-d/2,\ (k/2)^d),\\
 w_{2}&=(w_1-d/2,\ (-4C+(k+d)/2)^d).
 \end{aligned}
 \tag{R18.2}
\]

The notation \(x^d\) in a weight list means d repeated entries. These equations follow directly from \nolinkurl{lamTrue_old} and \nolinkurl{lamTrue_new} in \nolinkurl{Columns/BaseLemmas.lean}496--530, with \nolinkurl{hcToWeightQ} at stage zero. They are the already checked recurrence, used here to compute its extrema.

The source does \textbf{not} define spread as first minus last alone. \nolinkurl{LocalRingCardinality/DataCoeff.lean}34--35 defines it, for positive rank N, by

\[
              \operatorname{spread}(w)=w_0-w_{N-1}+N.
 \tag{R18.3}
\]

All rank terms in the following calculation are retained. The three exact baseline spreads are

\[
 \begin{aligned}
 S_0&=2(k-1)C+1,\\
 S_1&=(2k+1)C+1+d/2,\\
 S_2&=(2k+5)C+1+d/2.
 \end{aligned}
 \tag{R18.4}
\]

\textbf{Proof.} The entries of \(w_0\) decrease by \(2C-1>0\). Its last entry is \(3C+k/2\). In \(w_1\) the old last entry is \(3C+k/2-d/2>k/2\), since \(3C>d/2\), so the new constant block is the minimum. In \(w_2\) the smallest old entry is \(k/2-d/2\), which exceeds the new last entry \(-4C+(k+d)/2\) by \(4C-d>0\). At both steps the first entry remains the shifted old first entry. These observations also verify dominance at the block joins. Substituting these first and last entries and the ranks \(k,k+d,k+2d\) into (R18.3) gives(R18.4). Finally \(S_1-S_0=3C+d/2>0\) and \(S_2-S_1=4C>0\). Thus \(S_2\) is the largest baseline spread. ∎

If \(u_0\) at one embedding is increased by any integer \(n\ge0\), every weight in(R18.2) changes only in its first entry, which increases by n.~Dominance persists, and each literal spread is exactly \(S_j+n\). The same extrema hold for any interlaced changes of the middle u entries that retain the first and last entries and the same row data. Consequently the calculation applies simultaneously to any finite source word collection with those common extremes.

\subsection{2. The growing K033 family and an exact size parameter}\label{the-growing-k033-family-and-an-exact-size-parameter}

At K033's two distinguished embeddings the first u entry is increased by

\[
                   A_r=11^r,\qquad B_r=121\,11^r
                         \quad(r\ge0),
 \tag{R18.5}
\]

respectively. Other embeddings have no increase. The fixed row data, finite characters, G-primes, ranks and source interlacing remain those already checked in K033. The largest literal spread over all stages and embeddings is therefore exactly

\[
                        S(r)=S_2+B_r.
 \tag{R18.6}
\]

Choose the integral real size parameter

\[
                 a(r)=\left\lceil\frac{S_2+B_r}{2}\right\rceil.
 \tag{R18.7}
\]

This a is the \nolinkurl{Sizes} real parameter, not the source rank, the lift count, or the ordinary differential degree. The exact record \nolinkurl{AbelianHostFieldConditions.lean}486--494 is

\[
 \begin{gathered}
 2fk<g_i-1\ (i=0,1),\qquad 4a<T,\\
             2fN_2+1<4a,\qquad
 \operatorname{spread}(w_{j,\tau})\le2a\quad(0\le j\le2).
 \end{gathered}
 \tag{R18.8}
\]

The first inequalities were checked in K031/K032. The last follows exactly from(R18.6)--(R18.7). The middle lower bound holds already at r=0 by the exact values below and continues to hold because a(r) increases. Thus the only remaining condition on r is \(4a(r)<T\), with T the corresponding explicit prime of R017.

\begin{longtable}[]{@{}
  >{\raggedright\arraybackslash}p{(\columnwidth - 8\tabcolsep) * \real{0.1579}}
  >{\raggedleft\arraybackslash}p{(\columnwidth - 8\tabcolsep) * \real{0.2105}}
  >{\raggedleft\arraybackslash}p{(\columnwidth - 8\tabcolsep) * \real{0.2105}}
  >{\raggedleft\arraybackslash}p{(\columnwidth - 8\tabcolsep) * \real{0.2105}}
  >{\raggedleft\arraybackslash}p{(\columnwidth - 8\tabcolsep) * \real{0.2105}}@{}}
\toprule\noalign{}
\begin{minipage}[b]{\linewidth}\raggedright
Host
\end{minipage} & \begin{minipage}[b]{\linewidth}\raggedleft
\(S_0\)
\end{minipage} & \begin{minipage}[b]{\linewidth}\raggedleft
\(S_1\)
\end{minipage} & \begin{minipage}[b]{\linewidth}\raggedleft
\(S_2\)
\end{minipage} & \begin{minipage}[b]{\linewidth}\raggedleft
Exact permitted r interval
\end{minipage} \\
\midrule\noalign{}
\endhead
\bottomrule\noalign{}
\endlastfoot
\(J^{14}\) & 3,702,796,042,027 & 3,838,264,190,786 & 4,018,888,387,958 & \(0\le r\le55\) \\
\(J^{56}\) & 13,537,133,455,459 & 13,658,724,475,337 & 13,820,845,834,085 & \(0\le r\le90\) \\
\end{longtable}

For complete reproducibility, these endpoints are certified by the integer inequalities

\[
 \begin{split}
 121\,11^{R_*}&\le
       2\left\lfloor\frac{T-1}{4}\right\rfloor-S_2
                         <121\,11^{R_*+1},\\
 R_*&=55\text{ for }J^{14},\qquad
 R_*=90\text{ for }J^{56}.
 \end{split}
 \tag{R18.9}
\]

\textbf{Proof of the endpoints.} For positive integers S and T, \(4\lceil S/2\rceil<T\) is equivalent to \(S\le2\lfloor(T-1)/4\rfloor\). Substitute \(S=S_2+121\,11^r\), and use strict increase in r. The two exact boundary comparisons in(R18.9), all four clauses of(R18.8) throughout the stated intervals, and the failed next endpoint are verified by integer arithmetic in the saved certificate. The primality certificates for T are carried from R017; they are not replaced by a numerical primality screen here. ∎

At the first excluded r, even allowing a real size parameter cannot restore this particular fixed-T family. Indeed its spread condition forces \(a\ge S(r)/2\), and the certificate also verifies \(2S(r)\ge T\) there. The strict inequality \(4a<T\) is then impossible. This excludes only that fixed T and growing-weight choice. It does not exclude choosing a new T-prime as the weights grow.

The K033 growing degree first exceeds its fixed baseline degree at r=13 on \(J^{14}\) and r=15 on \(J^{56}\). Accordingly the nonempty ranges13--55 and15--90 also meet the entire numerical(R18.8). The large fixed valuation errors and all finite and real normalization costs in K033/K034 are unchanged; baseline domination is not a MainX gain/error inequality.

\subsection{3. Scope and exact sources}\label{scope-and-exact-sources}

R017 now has the missing \nolinkurl{spread_le} check for the stated baselines and these growing selected weights. The result supplies a complete numerical \nolinkurl{Sizes} witness, conditional only on the already specified weight/prime data. It does not make the primitive CM type satisfy \nolinkurl{Fields.Ok.quad}, identify the integral row and previous-image lattices, or compute the global DP44 budget. K033/K034's unbounded sharp-slope assertion concerns their selected section modules; its use with this fixed T is restricted to the explicit finite intervals just proved.

Exact stipulated source locators, pinned to commit \nolinkurl{7cbf7a0c32fca1badf05931879a37c3f8a2b03b4}, are \nolinkurl{Formalism/LocalRingCardinality/DataCoeff.lean}34--39, \nolinkurl{Formalism/LocalToGlobal/Columns/BaseLemmas.lean}467--530, and \nolinkurl{Formalism/LocalToGlobal/Columns/AbelianHostFieldConditions.lean} 443--494. K031 HC3--10, K032 QH rank/weight data, K033 CS20--35 and R017 R17.1--5 are carried at their recorded checked scope. The script is \nolinkurl{scripts/selected_fixed_prime_sizes_certificate.py}; its complete arithmetic output is in \nolinkurl{checks/selected_fixed_prime_sizes_certificate.json}. This is a mathematical source/parameter verification, not a Lean build. Whole-proof readings are recorded separately from the finite certificate.

\clearpage

\section{ALP-LIFT-23: exact Fourier contraction of the prescribed coupled Q family}\label{alp-lift-23-exact-fourier-contraction-of-the-prescribed-coupled-q-family}

Author: GPT-6 Astra Pro, ALP\_MACHINE /root/lift\_accounting. Date: 22 November 2026 UTC. Inherited reasoning effort is not exposed.

This finding evaluates the entire literal local Y sum in the two coupled finite kernels for m=2 at a degree-one Q place. The auxiliary depth disappears. The resulting finite function is nonzero, has an explicit Gauss-sum scalar, and has primitive content at every coefficient place over11 when q is different from11. This is an identity of the specified finite Schwartz functions. It is not a claim that the second automorphic theta integral is injective on the preceding image.

\subsection{1. Literal source data and conventions}\label{literal-source-data-and-conventions}

Use the pinned source tree work/research/alpoge\_source at commit 7cbf7a0c32fca1badf05931879a37c3f8a2b03b4. Put \[
 m=2,\quad r=k+1,\quad d=t+1,\quad b=k+d=r+t,\quad N=b+d.
 \tag{Q1}
\] The first edge has ranks k,b, and the second edge has ranks b,N. Let Q be a relative split place whose local field F is the rational q-adic field, so q is a rational prime and a uniformizer. Write O for its ring of integers. This explicit unramified degree-one hypothesis is retained; residue degree one alone would not imply it.

This hypothesis holds at every literal tuple Q place. As independently audited by /root/definite\_oscillator, Tuples.lean64--67 defines DegOne by prime absolute norm q together with q not belonging to the square of the prime ideal, and GoodPl adds relative splitting. The first condition gives residue degree one; the second gives ramification index one, since the local valuation of q is that index. The split CM completion is therefore the rational q-adic field. The tuple condition in Tuples.lean112--114 kills inertia outside GoodPl, while Qprime128 is defined by failure to kill inertia. Thus Qprime is contained in GoodPl. The larger host's ramified29 place is excluded from this Q set. These exact source clauses were read independently here in \nolinkurl{Formalism/LocalToGlobal/Divisibility/Tuples.lean.}

Take the actual standard frames and the forced initial and terminal frame choices of ChIdx.Ok.zero\_h and last\_h. The limited-depth character is trivial on O and nontrivial on q\^{}\{-1\}O. Its induced nontrivial residue additive character is \[
             \psi_0(z)=\psi(q^{-1}\widetilde z).
 \tag{Q2}
\] Only the finite argument range prescribed by LocCst.Ok is used.

Let chi be the actual tame residue character of the tuple at Q. We first assume chi is nontrivial, and extend it by zero at0. The trivial-character version is given in section5. Let \(\mathsf M=\mathrm{cst.M}\ge1\), with depth at least \(2\mathsf M\). The actual local parameter set, by the already checked DO032, is \[
       Y_Q=q^{-\mathsf M}O^b/q^{\mathsf M}O^b.
 \tag{Q3}
\] Representatives may be changed within these cosets because both literal factors are unchanged by that change.

The exact source locators are:

\begin{itemize}
\tightlist
\item
  \nolinkurl{Formalism/LocalToGlobal/Columns/FiniteSumNonzeroLargeDepth.lean} (FS): entB684--694; unitV669; PhiQ711--715; omPhi939--941; joinCol982--983 and PsiAt985--986; zero\_h1020, last\_h1025; XiCpl11134--1138; PhiTop1142--1147 and ftY1150--1156; XiCpl21160--1167.
\item
  The same file: ybSet1125--1129, and the actual coupled Y action and CplKer in1216--1237.
\item
  Accepted and checked DO031.3--DO031.5 and DO032.2--DO032.10 give the literal terminal last-row normalization and finite quotient, including the source scalar \(q^{-\mathsf M b}\).
\item
  The Q middle group is the reduction preimage of \(P_{r,t}\), by Jat909--913, last\_h, and grpQ1062--1066. It does not generally fix an individual chosen class \(y=e_k\). The present calculation retains the whole prescribed family.
\end{itemize}

The actual terminal last row in these frames is denoted by \(\eta\in F^{1\times b}\); any piT unit already absorbed by the source frame remains absorbed here. No further coordinate rescaling is introduced.

\subsection{2. The two literal functions}\label{the-two-literal-functions}

Write the first matrix as \[
                 X=\begin{pmatrix}A\\C\end{pmatrix},
 \qquad A\in M_{r,k}(F),\quad C\in M_{t,k}(F),
 \qquad y=\begin{pmatrix}z\\w\end{pmatrix}.
 \tag{Q4}
\] Because the Q exponent in lamYat is zero and the two relevant frames are the identity, entB of joinCol is precisely the ordinary matrix \([X,y]\), with the added column in position k. Thus the first raw factor is exactly \[
 F_X(y)=
 \mathbf1_{M_{r,k}(O)}(A)\mathbf1_{qM_{t,k}(O)}(C)
 \mathbf1_{O^r}(z)\mathbf1_{qO^t}(w)
       \chi\!\left(\det[\bar A,\bar z]\right).
 \tag{Q5}
\] In particular it is zero when \(\bar A\) has rank smaller than k. This includes all support conditions, including the lower auxiliary column condition; it is not just the specialization \(y=e_k\).

Write the second matrix as \(\binom Z\eta\), with \(Z\in M_{N-1,b}(F)\). Its terminal factor is \[
 H(Z)\,q^{-\mathsf M b}
 \mathbf1_{q^{-\mathsf M}O^b}(\eta)\psi(-\eta y),
 \qquad H(Z)=\Phi_{\mathrm{Top},Q}(Z).
 \tag{Q6}
\] Here H has all of the literal PhiQ conditions: Z is integral, its bottom-left \((N-1-r)\)-by-r block is divisible by q, and its leading r-by-r determinant is a unit, on which chi is evaluated.

The actual local contraction over its entire allowed Y set is therefore \[
 {\cal C}_{\mathsf M}(X,Z,\eta)
   =H(Z)q^{-\mathsf M b}
     \mathbf1_{q^{-\mathsf M}O^b}(\eta)
        \sum_{y\in Y_Q}F_X(y)\psi(-\eta y).
 \tag{Q7}
\] The source writes a finite sum, not a Haar integral over Y. The scalar in (Q7) is exactly its ftY scalar.

\subsection{3. Closed formula and independence of depth}\label{closed-formula-and-independence-of-depth}

For a matrix \(\bar A\) of rank k define its nonzero cofactor row \(\ell_A\in\kappa^{1\times r}\) by \[
                  \ell_A z=\det[\bar A,z].
 \tag{Q8}
\] Put \[
                    G(\chi,\psi_0)=
                       \sum_{s\in\kappa^\times}\chi(s)\psi_0(s).
 \tag{Q9}
\] If the support conditions on A,C in (Q5) hold and \(\bar A\) has rank k, then the exact answer is \[
 \boxed{
 {\cal C}_{\mathsf M}(X,Z,\eta)=
 \begin{cases}
 H(Z)\,q^{-d}G(\chi,\psi_0)\chi(-c)^{-1},
 &\eta\in q^{-1}O^b,\quad
   \overline{q\eta_{\mathrm{top}}}=c\ell_A,\quad c\in\kappa^\times,\\
 0,&\text{otherwise}.
 \end{cases}}
 \tag{Q10}
\] It is also zero when those A,C support or rank conditions fail. There is no condition on \(\overline{q\eta_{\mathrm{low}}}\). In particular the entire answer is independent of \(\mathsf M\ge1\).

Proof. Because of (Q5), the contributing representatives in (Q3) are exactly \[
       z\in O^r/q^{\mathsf M}O^r,\qquad
       w\in qO^t/q^{\mathsf M}O^t.
\] The determinant character depends only on \(\bar z\). The sum over all higher lifts in both blocks is zero unless \(\eta\in q^{-1}O^b\). When this condition holds, those sums contribute \(q^{(\mathsf M-1)b}\). The lower residue block is zero. Multiplication by the source scalar \(q^{-\mathsf M b}\) therefore leaves \[
       H(Z)q^{-b}\sum_{z\in\kappa^r}
            \chi(\ell_Az)\psi_0(-az),
       \qquad a=\overline{q\eta_{\mathrm{top}}}.
 \tag{Q11}
\] This argument also works for \(\mathsf M=1\): there are no higher lifts, and the support of ftY is already \(q^{-1}O^b\).

If \(a\) is not a scalar multiple of \(\ell_A\), there is a vector \(z_0\in\ker\ell_A\) with \(az_0\ne0\). Translation by a suitable multiple of \(z_0\) preserves the determinant factor and multiplies the character sum by a scalar different from1. The sum is zero. If \(a=c\ell_A\), every fibre of \(\ell_A\) has \(q^k\) elements, so the sum equals \[
           q^k\sum_{s\in\kappa^\times}\chi(s)\psi_0(-cs).
 \tag{Q12}
\] For \(c=0\), it vanishes because chi is nontrivial. For \(c\ne0\), substitute \(u=-cs\); the answer becomes \(q^k\chi(-c)^{-1}G(\chi,\psi_0)\). Since \(b-k=d\), (Q10) follows. Every manipulation takes place inside the specified finite additive quotient, so no unrestricted extension of the source psi has been assumed.

A concrete nonzero value is obtained with \(A=\binom{I_k}{0}\), C=0, \(\eta=q^{-1}e_k^\vee\), and Z having identity leading b-by-b block and all remaining rows zero. Then \(\ell_A=e_k^\vee\), H(Z)=1, and \[
      {\cal C}_{\mathsf M}(X,Z,\eta)
                       =q^{-d}G(\chi,\psi_0)\chi(-1)^{-1}\ne0.
 \tag{Q13}
\]

\subsection{4. Exact norms, content, and height of this scalar}\label{exact-norms-content-and-height-of-this-scalar}

All measures in this section give integral additive matrix lattices volume1. The standard finite Gauss identity is proved here in the needed normalization. For \(c\in\kappa\) put \(G_c=\sum_{s\ne0}\chi(s)\psi_0(cs)\). Then \(G_0=0\), and substitution shows all \(|G_c|\), \(c\ne0\), are equal. Additive orthogonality gives \(\sum_c|G_c|^2=q(q-1)\). Hence \[
                       |G(\chi,\psi_0)|^2=q.
 \tag{Q14}
\] This proof applies after every complex embedding of the finite coefficient field, since both induced characters remain nontrivial.

For fixed supported X of full residue rank, the nonzero eta set in (Q10) has measure \((q-1)q^t\): there are \(q-1\) top fractional classes and \(q^t\) arbitrary lower fractional classes, each followed by the integral lattice of volume1. Thus \[
 \int_F\cdots\int_F
   \left|{\cal C}_{\mathsf M}(X,Z,\eta)\right|^2\,d\eta
       =|H(Z)|^2(q-1)q^{-d}.
 \tag{Q15}
\] This also equals \(|H(Z)|^2\int|F_X(y)|^2dy\), since the first support has lower volume \(q^{-t}\) and the determinant is nonzero on the proportion \((q-1)/q\) of its top residue classes. Thus the depth-independent contraction is exactly the additive Fourier transform in this variable, including its normalization.

For the finite counting measure on (Q3), the number of nonzero first values is \[
                 (q-1)q^{\mathsf M b-d}.
 \tag{Q16}
\] The number of nonzero transformed values, modulo \(q^{\mathsf M}O^b\), is \((q-1)q^{t+\mathsf M b}\). Their squared amplitude \(q^{1-2d}\) gives the same number (Q16), as required by the actual unitary finite Fourier matrix.

The upper terminal factor has norm \[
 \Gamma_{\mathrm{top},Q}
  =q^{-r(N-1-r)}\prod_{i=1}^r(1-q^{-i}).
 \tag{Q17}
\] The probability that an r-by-k residue matrix has rank k is \(\prod_{i=0}^{k-1}(1-q^{i-r})\). Integrating also in X and Z gives the exact joint Schwartz norm \[
 \boxed{
 \|{\cal C}_{\mathsf M}\|_2^2
   =q^{-tk}(q-1)q^{-d}
       \prod_{i=0}^{k-1}(1-q^{i-r})\,
       \Gamma_{\mathrm{top},Q}.}
 \tag{Q18}
\] This is the norm of the contracted raw array in its product matrix variables; it is not the norm of an automorphic composite operator.

Both G and its complex conjugate are algebraic integers, and (Q14) is the algebraic identity \(G\overline G=q\). Consequently G is a unit at every coefficient place above11 if \(q\ne11\). Every nonzero value in (Q10) is then an11-adic unit. Its minimum valuation is exactly zero, by (Q13). No division by \(q-1\), a finite general linear group order, or a class mass has been introduced.

The magnitude of the scalar in (Q10) is \(q^{1/2-d}\). Its absolute logarithmic Weil height is exactly \[
                  h\!\left(q^{-d}G(\chi,\psi_0)\right)
                               =(d-\tfrac12)\log q.
 \tag{Q19}
\] Indeed its inverse is \(q^{d-1}\overline G\), an algebraic integer, and every complex conjugate of that inverse has absolute value \(q^{d-1/2}>1\). Computing the height of this algebraic integer and using \(h(x^{-1})=h(x)\) proves (Q19). This is the height of the specified finite scalar. It must not be renamed the source global H.

\subsection{5. Trivial residue character and exact scope}\label{trivial-residue-character-and-exact-scope}

If chi is trivial, still extended by zero at0, (Q11) remains valid. The value is zero unless \(a=c\ell_A\), and its two remaining values are \[
 {\cal C}_{\mathsf M}=
 \begin{cases}
 H(Z)q^{-d}(q-1),&c=0,\\
 -H(Z)q^{-d},&c\ne0.
 \end{cases}
 \tag{Q20}
\] The norm (Q15) remains the same. The main primitive-value conclusion remains true for \(q\ne11\), because the \(c\ne0\) value is a unit, although the \(c=0\) value can have positive valuation if11 divides \(q-1\). The nontrivial-character height formula(Q19) is not asserted for the trivial case.

The full allowed global Y set factors into these local finite quotients and the G quotients by the already checked CRT statement DO032.2. Therefore a finite product of the literal Q contractions is exactly the corresponding partial global Y sum. At sharp G, LIFT22 shows that the first factor leaves just its single forced local class. These statements remove cancellation in the raw finite array and retain the exact Q Fourier scalar and its norms.

They do not show that this raw function survives the next automorphic integral. At Q the two stages have the same tuple residue character, so the all-ramified boundary exclusion used in LIFT21 at G is not available. The fixed-Q-class first theta function is also not generally in the full prescribed middle parabolic type by itself: that covariance belongs to the coupled family. An argument for the second image must retain this family or supply the exact local quotient comparison. No claim about a globally pinned scalar, the iterated integral row lattice, the complete D7 data, or a DP44 gain is made.

\subsection{6. Checks}\label{checks}

The exact finite replay is \nolinkurl{raw_coupled_q_fourier_certificate.py.} It compares the literal finite Y sum with (Q10)/(Q20), including the nontrivial M=2 higher-depth cancellation, in cyclotomic integer coordinates. It also checks the support counts, Parseval factors, the cofactor orientation, and the unit normalization. Small finite examples are diagnostics of the identity, not full Ranks data. Whole-proof reading and replay events are recorded separately.

\clearpage

\section{AM-DO-036: the unchanged first raw kernel survives at maximal 29-level}\label{am-do-036-the-unchanged-first-raw-kernel-survives-at-maximal-29-level}

Author: GPT-6 Astra Pro, ALP\_MACHINE \nolinkurl{/root/definite_oscillator}. 22 November 2026 UTC. Inherited reasoning effort is not separately exposed.

The first raw T factor has a uniform nonzero local test on every unitary source vector with its prescribed mirabolic invariance when \nolinkurl{b>=4a}. The proof separates its compact term from an explicitly bounded nonunit-determinant tail. The raw G and Q tests are exact compact projectors. Together with the checked spherical tests elsewhere and AM-R-017's actual class-supported input, these facts prove nonvanishing with the \textbf{unchanged full first G/T/Q finite kernel and the maximal vacuum at29}. No auxiliary29 Heisenberg translate is inserted.

The statement uses the new mixed theta realization and its actual stage characters. It does not instantiate the unchanged elliptic \nolinkurl{VacIn}, the source's \nolinkurl{Fields.Kq} record, every \nolinkurl{ChIdx.Ok} clause, an integral prior image lattice, or the MainX terminal-mass comparison. A positive target quotient mass remains in the exact norm identity.

\subsection{1. Exact local convention and the full-rank period operator}\label{exact-local-convention-and-the-full-rank-period-operator}

Let F be a nonarchimedean characteristic-zero local field, O its ring, q its residue cardinality, and

\[
 G=\operatorname{GL}_a(F),\quad K=\operatorname{GL}_a(O),\quad
 K_b=\operatorname{GL}_b(O),\quad b\ge4a,\quad z=b/2.
 \tag{DO36.1}
\]

Normalize group Haar by \nolinkurl{vol(K)=vol(K_b)=1} and matrix Haar by \nolinkurl{vol(M_(b,a)(O))=1}. Let eta be the actual unitary local splitting component. The genuine Hom Weil action is

\[
 [\omega(g)f](X)=\eta(\det g)^{-1}|\det g|^{b/2}f(Xg).
 \tag{DO36.2}
\]

This is the source formula of AM-DO-020/024 and ALP-LIFT-20, and the common global character is exactly \nolinkurl{mu_0 xi}, as checked in DO035 from \nolinkurl{LevChiEqOne.lean}61,100. In particular it is unramified at T, has residue chi at Q, and is tame with the prescribed residue character at G.

Let pi be any unitary representation of G on a Hilbert space V. Use conjugate-linear-first inner products. Put \nolinkurl{X_0=[I_a;0]} and define the Bochner integral

\[
 B_f(v;k)=\int_G f(kX_0h)|\det h|^{b/2}
                   \eta(\det h)^{-1}\pi(h)v\,dh,
                       \qquad k\in K_b.
 \tag{DO36.3}
\]

For the raw functions below \nolinkurl{|f|<=1_(M_(b,a)(O))}. Since \nolinkurl{kX_0} has an integral left inverse, \nolinkurl{f(kX_0h)!=0} implies \nolinkurl{h in M_a(O)}. The convergence and all interchanges follow from the exact bound proved in Section 2.

The exponent is \textbf{b/2}, with no unrecorded modular-character shift. Directly, for \nolinkurl{g in G}, simultaneous replacement of \nolinkurl{f,v} by \nolinkurl{omega(g)f,pi(g)v} gives the integral with variable \nolinkurl{hg}. The determinant powers cancel after \nolinkurl{t=hg}, and

\[
 \eta(\det g)^{-1}\eta(\det(tg^{-1}))^{-1}
                         =\eta(\det t)^{-1}.
\]

Hence

\[
                 B_{\omega(g)f}(\pi(g)v;k)=B_f(v;k).
 \tag{DO36.4}
\]

This is the simultaneous coinvariant convention for the literal bilinear theta lift. It checks both the character sign and the half-determinant exponent without importing a different local induced-model normalization.

Write

\[
 C_{a,b}(q)=\prod_{i=0}^{a-1}(1-q^{i-b}).
 \tag{DO36.5}
\]

The exact full-rank disintegration proved in LIFT16(S6) is

\[
 \int_{M_{b,a}(F)}H(X)dX
   =C_{a,b}(q)\int_{K_b}\int_G
                    H(kX_0h)|\det h|^b\,dh\,dk.
 \tag{DO36.6}
\]

For clarity, every full-rank X has such a presentation because its column space intersected with \nolinkurl{O^b} is a saturated direct summand. The two measures have the same right Jacobian; their constant is the probability that a residue \nolinkurl{b by a} matrix has rank a, which is precisely (DO36.5). The rank-deficient set has additive measure zero. These facts also give a direct proof of (DO36.6) in the present convention.

Substitution into the squared norm of (DO36.3), with the second integration variable written \nolinkurl{t=hg}, now proves the exact identity

\[
 \boxed{\quad
 Z_f(v):=\int_G\langle f,\omega(g)f\rangle
                    \langle v,\pi(g)v\rangle\,dg
       =C_{a,b}(q)\int_{K_b}\|B_f(v;k)\|^2dk\ge0.
 \quad}
 \tag{DO36.7}
\]

Indeed the two eta factors on the squared-norm side become \nolinkurl{eta(det h) eta(det t)^(-1)=eta(det g)^(-1)}; the two determinant powers become \nolinkurl{|det h|^b |det g|^(b/2)}, and the vector inner product is \nolinkurl{<v,pi(g)v>}. Equation (DO36.6) gives the other matrix coefficient. With the opposite convention for writing both matrix coefficients, simultaneous \nolinkurl{g -> g^(-1)} gives the same real value. This is the actual local doubling test after the Hom/tensor and bilinear/conjugated-input transport of DO028.

\subsection{2. Exact integrability and the uniform tail inequality}\label{exact-integrability-and-the-uniform-tail-inequality}

For \nolinkurl{z>a-1}, define

\[
 D_a(z;q)=\int_{G\cap M_a(O)}|\det h|^zdh
          =\prod_{i=0}^{a-1}(1-q^{i-z})^{-1},
 \qquad E_a(z;q)=D_a(z;q)-1.
 \tag{DO36.8}
\]

Here is a self-contained normalization check. For additive matrix measure, let \nolinkurl{J_a(s)=int_(M_a(O)) |det X|^s dX}. The first column has exact common valuation n with measure \nolinkurl{q^(-an)(1-q^(-a))}. A primitive first column is carried to \nolinkurl{e_1} by an integral invertible matrix. The remaining columns then contribute a free first row and an \nolinkurl{(a-1)} square matrix. Therefore

\[
 J_a(s)=\frac{1-q^{-a}}{1-q^{-a-s}}J_{a-1}(s),\qquad J_0(s)=1.
\]

The multiplicative Haar measure is \nolinkurl{dh=prod_(i=1)^a(1-q^(-i))^(-1)|det h|^(-a)dX}, because K has the displayed additive measure. Substitute \nolinkurl{s=z-a} to obtain (DO36.8). The convergence condition is \nolinkurl{s>-1}, equivalently \nolinkurl{z>a-1}. In particular, \nolinkurl{E_a} is exactly the integral over integral matrices with nonunit determinant, not a heuristic count of Cartan double cosets.

For \nolinkurl{x_i=q^(i-z)} put

\[
 S_a(z;q)=\sum_{i=0}^{a-1}x_i
                =q^{-z}\frac{q^a-1}{q-1}.
\]

When \nolinkurl{z>=2a},

\[
 q^a S_a(z;q)\le\frac{1-q^{-a}}{q-1}<1.
\]

The elementary inequality \nolinkurl{prod(1-x_i)>=1-sum x_i} consequently gives

\[
 \boxed{\quad
 0<E_a(z;q)\le\frac{S_a(z;q)}{1-S_a(z;q)}
                 <\frac1{q^a-1}\qquad(z\ge2a).
 \quad}
 \tag{DO36.9}
\]

This includes q=2. It is strictly stronger than needed at the actual T primes of R017.

The bound \nolinkurl{||B_f(v;k)||<=D_a(z;q)||v||} proves the absolute convergence in Section 1. For the Fubini step in (DO36.7), take absolute values before integrating h,t,k; the result is bounded by \nolinkurl{C_(a,b) D_a(z;q)^2 ||v||^2}. Thus the positive identity is proved by absolutely convergent integrals, with no meromorphic continuation or local theta-occurrence theorem substituted for this fixed-vector calculation.

\subsection{3. The untouched T incidence factor has a positive test}\label{the-untouched-t-incidence-factor-has-a-positive-test}

Use the actual source and target residue covectors \nolinkurl{u!=0,v_0!=0} and the raw T function from DO035:

\[
       f_T(X)=1_{M_{b,a}(O)}(X)1_{\{v_0\bar X=u\}}.
 \tag{DO36.10}
\]

Its prescribed source compact is

\[
 J_T=\{g\in K:u\bar g=u\},\qquad
                     \operatorname{vol}(J_T)=\frac1{q^a-1}.
 \tag{DO36.11}
\]

Choose a primitive integral matrix X* of rank a with \nolinkurl{v_0 X*=u} modulo q. Such a choice exists because b\textgreater a: after an integral target-coordinate change make \nolinkurl{v_0} the last-coordinate functional, use \nolinkurl{I_a} in the first a rows, u in the last row and zero elsewhere. Extend its primitive columns to \nolinkurl{k* in K_b}, so \nolinkurl{k* X_0=X*}.

Let \nolinkurl{v in V^(J_T)}. The compact part of (DO36.3) is exactly

\[
 \int_K f_T(X*h)\eta(\det h)^{-1}\pi(h)v\,dh
                   =\frac{v}{q^a-1}.
 \tag{DO36.12}
\]

Indeed eta is unramified at T, the condition is precisely \nolinkurl{h in J_T}, and pi fixes v there. The remaining integral has norm at most \nolinkurl{E_a(b/2;q)||v||}, by unitarity of pi and eta. Define

\[
 \Delta_T(a,b;q)=\frac1{q^a-1}-E_a(b/2;q)>0.
 \tag{DO36.13}
\]

Then \nolinkurl{||B_(f_T)(v;k*)||>=Delta_T||v||}. The same equality of the integrands holds for every k congruent to k* modulo q: the raw function depends only on reduction inside integral matrices, and every such primitive matrix has an integral left inverse. This open subset has measure \nolinkurl{|GL_b(F_q)|^(-1)}. Equation (DO36.7) proves

\[
 \boxed{\quad
 Z_{f_T}(v)\ge
 \frac{C_{a,b}(q)}{|\operatorname{GL}_b(\mathbf F_q)|}
       \Delta_T(a,b;q)^2\|v\|^2>0\quad(v\ne0).
 \quad}
 \tag{DO36.14}
\]

This uses the raw incidence function, not the character-filtered auxiliary vector of DO035. A nontrivial T diamond character of the source input is fully compatible with its \nolinkurl{J_T} invariance; the lemma makes no claim that a T coefficient is spherical.

\subsection{4. The raw G and Q tests are exact compact projectors}\label{the-raw-g-and-q-tests-are-exact-compact-projectors}

At a G place let U be the upper unitriangular subgroup of \nolinkurl{GL_a(F_q)}, let psi\_U be the exact source generic character, and let \nolinkurl{J_G} be its full reduction preimage in K. ALP-LIFT-20 identified the actual first raw G/GS function, for the admitted class \nolinkurl{y=q^(-1)e_a mod O^b}, as

\[
 f_G(X)=\psi_U(\bar h)
       \quad\text{if }X\in M_{b,a}(O),\ \bar X=X_0\bar h,
                        \ \bar h\in U,
 \tag{DO36.15}
\]

and zero otherwise. If \nolinkurl{pi(j)v=psi_U(bar j)^(-1)v} for j in J\_G, then

\[
 B_{f_G}(v;1)
      =\frac{|U|}{|\operatorname{GL}_a(\mathbf F_q)|}\,v.
 \tag{DO36.16}
\]

To check this, \nolinkurl{f_G(X_0h)!=0} forces \nolinkurl{h in J_G}; its determinant has residue one, so the tame eta factor is one. The two generic characters cancel. There is no nonunit-determinant tail, since every residue matrix in U is invertible. The same equation holds on \nolinkurl{K_b(1)}, giving

\[
 Z_{f_G}(v)\ge
 \frac{C_{a,b}(q)}{|\operatorname{GL}_b(\mathbf F_q)|}
   \left(\frac{|U|}{|\operatorname{GL}_a(\mathbf F_q)|}\right)^2
                      \|v\|^2.
 \tag{DO36.17}
\]

This works on the whole actual inverse-generic source space; it is not limited to one chosen Steinberg multiplicity line.

At Q use (DO35.8), the same admissible global \nolinkurl{ybSet} class, and the actual label \nolinkurl{eta=mu_0 xi}. For every K-fixed vector,

\[
                      B_{f_Q}(v;1)=v.
 \tag{DO36.18}
\]

Here nonzero support forces \nolinkurl{h in K}, and the factor \nolinkurl{chi(det h)} in the raw function cancels \nolinkurl{eta(det h)^(-1)}. The same open \nolinkurl{K_b(1)} argument gives

\[
             Z_{f_Q}(v)\ge
              \frac{C_{a,b}(q)}{|\operatorname{GL}_b(\mathbf F_q)|}
                                       \|v\|^2.
 \tag{DO36.19}
\]

Thus all the first raw factors have a strictly positive fixed test on their actual initial source spaces. Finite nonzero Schwartz functions alone would not imply these statements; the projector and tail calculations are essential.

\subsection{5. Maximal vacua at every other finite place}\label{maximal-vacua-at-every-other-finite-place}

At a split place outside G/T/Q, eta is unramified and \nolinkurl{f=phi_0=1_(M_(b,a)(O))}. For every K-fixed v, (DO36.3) has compact term v and tail at most \nolinkurl{E_a(b/2;q)||v||}, independently of k. Hence

\[
 Z_{\phi_0}(v)\ge
       C_{a,b}(q)(1-E_a(b/2;q))^2\|v\|^2.
 \tag{DO36.20}
\]

In particular this applies at all the relative split29 places, including their absolute ramification index seven, and at the dyadic split places. The trace coefficient and additive character are the actual ones; the unitary identification with this Schrodinger model is DO022/034. Changing the absolute base-field ramification does not change the local residue-field calculation.

At an odd nonsplit place the extension is unramified, the even identity lattices are integrally hyperbolic, and the stage character is trivial: it is unramified and trivial on the base field. The actual fixed vacuum calculation LIFT16(S19--20), from Liu's unramified doubling section, gives on a unit spherical vector

\[
 Z_v= L_v(s_*,\pi_v)\prod_{i=1}^a(1-(-1)^i q_v^{-d-i}),
 \quad d=b-a,\quad s_*=(d+1)/2.
 \tag{DO36.21}
\]

The additive-character/codifferent rescaling and the actual lattice section were verified there and in DO022; no new choice of local section is made here. Macdonald's bounded spherical parameter criterion, in the form checked in LIFT14(E18), gives \nolinkurl{|alpha|<=Q^((a-1)/2)} with \nolinkurl{Q=q_v^2}. Consequently, putting \nolinkurl{x_v=q_v^(-(b-2a+2))}, a uniform positive lower bound is

\[
 Z_v\ge (1+x_v)^{-a}
                   \prod_{i=1}^a(1-q_v^{-d-i})>0.
 \tag{DO36.22}
\]

The local L-value is positive in this region. The unitary spherical parameter is stable under conjugate-inverse, and the unramified unitary-group standard parameter is stable under inverse. Thus its eigenvalue multiset is conjugation-stable. Nonreal conjugate factors have positive product, and every real factor \nolinkurl{1-alpha Q^(-s_*)} is positive because its absolute correction is less than one. The lower bound follows from \nolinkurl{|1-alpha Q^(-s_*)|<=1+x_v}. This also proves finiteness without a temperedness assumption.

The bounds (DO36.20), (DO36.22) differ from one by a summable quantity as the finite places vary. In the split case \nolinkurl{E_a=O_a(q^(-(b/2-a+1)))}, with exponent at least a+1\textgreater=2; in the nonsplit case the exponent is at least2a+2. Their infinite product is therefore strictly positive. This proves the needed fixed-vacuum assertion at29 as part of the same global family, without any later change of finite input.

\subsection{6. A direct convergent global norm identity in the correct range}\label{a-direct-convergent-global-norm-identity-in-the-correct-range}

Let E/F be either actual mixed host of K031 or K032, and \nolinkurl{G=U(I_a)}, \nolinkurl{H=U(I_b)}, with the positive real signatures, the exact common stage character, and the actual CM vacuum. The fields and trace coefficients are the ones proved in K030/K032 and DO034; this is not an extension of the old elliptic constructor. Both groups are anisotropic over F. For the two concrete first edges,

\[
 (a,b)=(42,1802)\quad\text{or}\quad(168,1802),
                              \qquad b\ge4a.
 \tag{DO36.23}
\]

Take the complete first raw finite vector from DO035 with the single admitted global Y class and the source/target T decorations of R017. At29 and every other place outside G/T/Q keep the maximal vacuum. At infinity take the actual selected highest polynomial P\_nu from DO034, with its source vector w and its true positive norm N\_infty. Choose R017's real matrix coefficient from this same w, so the input lies on the paired contragredient line at each real place. This is a choice of the matched vector in the known selected summand, not an inference that any source vector pairs with a rigid highest polynomial.

The source function F is now the \textbf{unchanged actual function} of R017, with the inverse character-name convention explained in the independent review of that finding. It is nonzero, maximal at29, has the initial inverse G character, the T mirabolic invariance and zero diamond sum, and the full trivial Q level. Its exact norm for maximal finite compact volumes one and real compact volume one is

\[
                    \|F\|^2=\frac1{28d_T^f d_\rho d_G}.
 \tag{DO36.24}
\]

\subsubsection{Correct literature interface}\label{correct-literature-interface}

GQT Theorem11.4 assumes \nolinkurl{a<b<=2a}; it must \textbf{not} be cited directly for (DO36.23). Instead use GQT Theorem7.1, printed p.33: for the anisotropic target, its Siegel--Weil formula holds in every dimension, with constant \nolinkurl{c_(b,a)=1} when \nolinkurl{s_0=(b-a)/2>0}. The doubling construction of section11.3, printed p.52, and the unfolding in sections11.5--11.6, printed pp.53--54, are the same absolutely convergent identities in this range. The character/conjugation transport is exactly DO028. No continuation through a pole is needed.

There is an especially direct unfolding here. The source Hermitian space is anisotropic. Every rational maximal isotropic subspace of its doubled space \nolinkurl{S plus (-S)} has zero intersection with either summand. Its two projections are therefore isomorphisms, so it is the graph of a rational isometry. Hence the rational double cosets in the doubling unfolding have just the open graph orbit, with stabilizer the diagonal G. Unfolding the Eisenstein sum is consequently the ordinary integral of a matrix coefficient against the right regular automorphic coefficient; there are no other rational orbits to dispose of.

Let dg be the chosen source Haar measure, let dh be the chosen target Haar measure, and put \texttt{V\_H=vol\_dh(H(F)\textbackslash{}\textbackslash{}H(A))}. Use the source inner product with dg. The resulting identity is

\[
 \boxed{\quad
 \|\Theta_\phi^{\rm bil}F\|_{dh}^{2}
       =V_H\int_{G(\mathbb A)}
            \langle\phi,\omega(g)\phi\rangle
            \langle F,R(g)F\rangle_{dg}\,dg.
 \quad}
 \tag{DO36.25}
\]

Here \nolinkurl{Theta^bil} is the integral against F, as in the literal source. For Tamagawa target measure \nolinkurl{V_H=2}; this is the factor \nolinkurl{[E:F]} in the doubled identity. Indeed GQT7.1 uses the normalized target theta integral and states \nolinkurl{tau(H)/[E:F]=1} for the present unitary group. Thus the unnormalized target integral is \nolinkurl{[E:F]} times the Siegel--Weil Eisenstein value. Scaling dh scales both sides of (DO36.25) by the same constant. Scaling dg scales the theta lift by that constant and scales the integral on the right twice, once through dg and once through the source inner product. This checks (DO36.25) for the measures used in R017 without assigning the target quotient mass the value one.

All interchanges are legitimate. The source and target quotients are compact, so their theta integrals are uniformly convergent on compact representatives. The remaining adelic matrix-coefficient integral is absolutely convergent: at each bad finite place Section 2 bounds it by \nolinkurl{C_(a,b)D_a(b/2)^2}; at almost every split place its absolute integral is exactly that same expression for the maximal vacuum, which is \nolinkurl{1+O_a(q^(-(b/2-a+1)))}. At nonsplit places the hyperspherical vacuum coefficient is nonnegative. Indeed compact invariance reduces it to a Cartan diagonal in the Levi of an integral hyperbolic polarization; the trivial splitting character gives the positive dilation action on the lattice indicator. The unitary change from that polarization preserves the selfdual lattice vacuum. Apply the same fixed-vacuum formula (DO36.21) to the trivial spherical source representation: it computes its absolute integral, and the parameter bound gives \nolinkurl{1+O_a(q^(-(b-2a+2)))}. This is a summable majorant. Infinity is compact. This is the ordinary convergent Siegel--Weil/unfolding regime, stronger than the continuation required by the second-term theorem.

\subsection{7. Strict positivity for the fixed global input, including multiplicities}\label{strict-positivity-for-the-fixed-global-input-including-multiplicities}

For a local factor define the bounded convolution operator

\[
 Q_v=\int_{G(F_v)}\langle\phi_v,\omega(g)\phi_v\rangle
                                               R(g)\,dg.
 \tag{DO36.26}
\]

At every split finite place (DO36.7) proves positivity of this operator on any unitary representation, including the actual automorphic Hilbert space. At G it preserves the inverse-generic subspace; at T it preserves the \nolinkurl{J_T}-fixed subspace; at Q and all spherical places it preserves the maximal-compact fixed subspace. These preservation assertions follow from the corresponding left/right compact covariance of its coefficient. At G/T/Q its restriction has the strictly positive uniform lower bounds (DO36.17), (DO36.14), (DO36.19). The nonsplit spherical operators are positive and have (DO36.22) on their fixed spaces. One may obtain the latter statement by decomposing the unitary spherical representation; the bound is uniform in its bounded spherical parameter.

At infinity compact Schur orthogonality gives, on the specified source contragredient line,

\[
                         Q_\infty=\frac{N_\infty^2}{d_\rho}.
 \tag{DO36.27}
\]

The actual F lies in that line in every real representation copy, by its matrix-coefficient construction in R017. It lies in the required fixed/type subspace at every finite place. Operators at different places commute. Finite products of these positive operators are bounded below on this subspace by the products of their displayed lower bounds. At the unramified places the deviations from one are uniformly summable, so the products converge on this subspace and their lower-bound product is still positive. Equivalently, decompose the finite-level, fixed-real-type automorphic space into its finitely many irreducible constituents, retaining their Hilbert multiplicity spaces; the same uniform bounds apply to every copy. The convolution operators act trivially on those multiplicity spaces. No multiplicity-one assumption and no choice of a favored irreducible projection is needed.

Let \nolinkurl{ell_v} denote (DO36.20) or (DO36.22) at the remaining spherical places and let \nolinkurl{c_G,c_T,c_Q} be the lower-bound constants in (DO36.17),(DO36.14),(DO36.19). Then

\[
 \begin{split}
 \|\Theta_\phi^{\rm bil}F\|_{dh}^{2}
 &\ge V_H\frac{N_\infty^2}{d_\rho}
       \prod_{P\in G}c_{G,P}
       \prod_{P\in T}c_{T,P}
       \prod_{P\in Q}c_{Q,P}
       \prod_{v\notin G\cup T\cup Q}\ell_v\ \|F\|^2\\
 &>0.
 \end{split}
 \tag{DO36.28}
\]

This proves that the exact nonzero R017 input survives the complete fixed first raw kernel at maximal29 level. The target value at the identity need not be nonzero; the proved assertion is positive global norm and hence a nonzero actual target automorphic form. It is stronger than a nonzero unrestricted local theta lift and addresses the particular finite input which the raw first coupled factor prescribes.

\subsection{8. Costs, source scope and range audit}\label{costs-source-scope-and-range-audit}

The finite raw Schwartz norm is exactly DO035(26). The real norm \nolinkurl{N_infty^2} is the actual DO034 Fischer/Gaussian/hook norm, including the pi\_B coordinate scalar and, on K032, its weighted quadratic metric. All finite projector and tail constants in (DO36.28) are explicit functions of the fixed ranks and residue cardinalities. They can be very small; this proof does not label them negligible. There is no auxiliary29 principal-level cost and no29 translation depth in this theorem. The T characters and their finite coefficient field are exactly those retained by R017.

The norm formula retains \nolinkurl{V_H}, the actual target quotient mass. Its positivity is enough for the present analytic nonvanishing. It is not an identification with the source's terminal mass or a solution of the MainX scalar equation. The construction is a genuine mixed-theta first image with the actual raw local G/T/Q factors and source compact types; the unchanged \nolinkurl{Fields.Kq}, the integral row comparisons and all later restricted-image contents remain separate.

Finally, the range distinction in Section 6 does not correct an earlier checked theorem. LIFT12 uses a=2,b=4, the allowed upper endpoint of GQT11.4. LIFT14(E23) explicitly applies that theorem only to later j\textgreater=1, where a=N\_j\textgreater d and \nolinkurl{a<b=a+d<2a}; its initial edge was constructed separately. LIFT16 continues those later edges, and LIFT19 has the same later-edge hypotheses after DO033's initial construction. The new first edge is the case requiring the separate GQT7.1 convergent-range derivation above.

\subsection{9. Primary references and proof checks}\label{primary-references-and-proof-checks}

\begin{itemize}
\tightlist
\item
  Gan, W. T., Qiu, Y., \& Takeda, S. (2014). The regularized Siegel--Weil formula (the second term identity) and the Rallis inner product formula. \emph{Inventiones Mathematicae, 198}, 739--831. arXiv:1207.4709v3. Theorem7.1, printed p.33; doubling section11.3, p.52; sections11.5--11.6, pp.53--54; the actual range restriction of Theorem11.4, p.54. \url{https://arxiv.org/pdf/1207.4709}. These portions were read directly in the primary PDF/text; its omitted conjugation bars use the already checked rendered-PDF convention in DO028.
\item
  Macdonald's bounded spherical criterion and Liu's actual unramified vacuum test are used in the exact checked form LIFT14(E18) and LIFT16(S19--20), with their primary locators. No new local theta classification is being asserted.
\item
  Literal raw factors, characters and target class: DO035 Sections1--4 and its source locators; LIFT20's raw first G equality; R017's exact source class, central compatibility and mass; DO034's actual new-host polynomial and metric.
\end{itemize}

The author whole-proof review and independent local review are recorded separately after the complete text is read. No Lean build is claimed.

\clearpage

\section{ALP-LIFT-24: the fixed second T kernel preserves the actual first projective-line type}\label{alp-lift-24-the-fixed-second-t-kernel-preserves-the-actual-first-projective-line-type}

Author: GPT-6 Astra Pro, ALP\_MACHINE /root/lift\_accounting. Date: 22 November 2026 UTC. Inherited reasoning effort is not exposed.

This finding replaces the first-edge determinant-tail estimate by an exact local period calculation suitable for the second edge. The raw T incidence function has a nonzero image on the actual first theta constituent. Its scalar in a specified induced model is a standard local L-factor divided by q\^{}n−1. The proof retains the nontrivial T dilation character of R017, uses the entire Schwartz space, and does not infer quotient survival from a formally nonzero full-rank integral.

The local normalization below is explicit. Its comparison with globally pinned arithmetic rows remains separate. In particular the L-factor can have nontrivial valuation at11; it is not discarded as an integral unit.

\subsection{1. Local groups and the actual T type}\label{local-groups-and-the-actual-t-type}

Let F be a characteristic-zero nonarchimedean local field with ring O, uniformizer pi, and residue field kappa of cardinality q. Assume its residue characteristic is larger than all ranks used below; this is satisfied by either explicit rational T prime of R017. Normalize vol(GL\_n(O))=1 in every rank. Write \[
 K_n=\mathrm{GL}_n(O),\quad
 J_n=\{g\in K_n:e_n^\vee\bar g=s(g)e_n^\vee,
                           s(g)\in\kappa^\times\},\quad
 H_n=\ker s.
 \tag{MT1}
\] Here e\_n denotes the last coordinate, with one-based indexing in this file. Thus J\_n is the reduction preimage of the parabolic of block sizes(n−1,1), and \[
 d_n=[K_n:J_n]=\frac{q^n-1}{q-1},\qquad
                  [K_n:H_n]=q^n-1.
 \tag{MT2}
\] Let nu be a nontrivial character of kappa\^{}times. For a smooth representation pi\_n define \[
 V_n(\nu)=\{v:\pi_n(j)v=\nu(s(j))v\quad(j\in J_n)\}.
 \tag{MT3}
\] For the R017 source, nu is the inverse of its named character nu\_P. All formulas use the actual representation character(MT3), so there is no implicit exchange of the two signs.

The literal T function for n→m is \[
 \phi_{u,v}(X)=\mathbf1_{M_{m,n}(O)}(X)
                              \mathbf1_{\{v\bar X=u\}},
 \tag{MT4}
\] where u and v are the prescribed nonzero residue covectors. This is FS646--649, phiT, as checked in DO035. The stage character eta is unramified at T. The actual Hom action is \[
 \omega_\eta(g,h)\phi(X)=
 \eta(\det g)^{-1}\eta(\det h)
 |\det g|^{m/2}|\det h|^{-n/2}\phi(h^{-1}Xg).
 \tag{MT5}
\] The source lift uses the bilinear coinvariant convention. These are the same signs as LIFT20--22 and DO036.

Integral source and target coordinate changes take u,v to the last coordinate covectors. They preserve the two Haar normalizations and the local Schwartz norm, and their eta values are1. We therefore use these coordinates locally. This does not assert an uncharged global change of arithmetic basis.

\subsection{2. Which representations contain this type}\label{which-representations-contain-this-type}

We use the following precise consequence of the principal-series type theorem and unitary classification.

\textbf{Local type lemma.} If an irreducible smooth unitary representation pi\_n has \(V_n(\nu)\ne0\), then \[
          \pi_n=\sigma\times\xi,\qquad
 \dim V_n(\nu)=1,
 \tag{MT6}
\] where sigma is an irreducible unitary spherical representation of GL\_(n−1)(F), xi is a unitary character of F\^{}times with residue nu, and the induction is normalized. Conversely this product has the displayed one-dimensional type.

Here and subsequently rank zero has its usual trivial interpretation. All actual source ranks are larger than one.

Proof. Restrict(MT3) to the upper Iwahori. Its torus character is \((1,\ldots,1,\nu)\), and it is trivial on the two relevant unipotent subgroups. Roche(1998), Theorem7.7, identifies this tame Iwahori character as the type for the principal-series inertial component. Thus pi\_n is a subquotient of a principal series with n−1 unramified characters and exactly one character with residue nu. The explicit residue-characteristic hypothesis in Roche, p.379, is satisfied here.

Separate the two inertial lines. Zelevinsky's classification and its disjoint-support irreducibility, already used in LIFT19, give \(\pi_n=\sigma\times\xi\), with sigma in the unramified principal series component and xi the unique character on the nu line. We explain why sigma is spherical and why the type has dimension1; being in that component alone would not suffice.

On the K\_n principal-congruence invariants of the full principal series, the compact representation is the finite principal series induced from the characters(1,\ldots,1,nu). Its J\_n,nu eigenspace has dimension one. Indeed the Bruhat double cosets of a Borel and P\_(n−1,1) are represented by the n possible positions of the distinguished last coordinate. If that coordinate does not match the nu coordinate, the diagonal torus in the intersection has independent unit entries at those two positions. The two characters then disagree, so the corresponding intertwiner is zero. In the matching coset they agree, and the space of scalar intertwiners is one-dimensional. This is the ordinary finite Mackey calculation; all compact invariants are exact in characteristic zero.

The unramified GL\_(n−1) principal series has its unique spherical irreducible subquotient sigma\_sph. The product sigma\_sph×xi is irreducible, by the distinct inertial supports, and has a nonzero J\_n,nu vector. Because the entire principal series has only the one such line, it is its only irreducible subquotient containing that line, with multiplicity one. Hence sigma=sigma\_sph and(MT6) follows.

Finally Tadic(1986), TheoremD, says every unitary irreducible is a product of the stated unitary Speh and complementary blocks, with unique factors, and such products are irreducible. Each block is contained in a single inertial line. Since the nu line in pi\_n has multiplicity one, its block is the unitary character xi. Removing that block leaves a unitary product, which is sigma. This proves the unitary assertion as well as the stated irreducibility.

A convenient realization of the line is the normalized induced model on P\_(n−1,1). Choose a unit spherical vector v\_sigma and let v\_0 have compact restriction \[
 v_0(k)=
 \begin{cases}
       \nu(s(k))v_\sigma,&k\in J_n,\\
       0,&k\notin J_n.
 \end{cases}
 \tag{MT7}
\] Then \(v_0(1)=v_\sigma\) and \[
                         \|v_0\|^2=d_n^{-1}.
 \tag{MT8}
\] The defining transformation under P∩K and the right J\_n character both follow directly. The lower principal unipotent congruence subgroup is in J\_n and acts trivially.

\subsection{3. The whole-Schwartz period at the actual point}\label{the-whole-schwartz-period-at-the-actual-point}

For this section pi\_n may be any irreducible smooth representation, and eta any unitary character. Put \[
 m=n+d,\qquad z_0=\frac{d+1}{2},\qquad
 X_0=\begin{pmatrix}I_n\\0\end{pmatrix}.
 \tag{MT9}
\] For Re(z) sufficiently large define \[
 I_z(\phi,w)=\int_{\mathrm{GL}_n(F)}
    \phi(X_0g)|\det g|^{z+(n-1)/2}
                      \eta(\det g)^{-1}\pi_n(g)w\,dg.
 \tag{MT10}
\] For unitary pi\_n the ordinary vector integral converges in a right half-plane, since phi is compactly supported and the determinant integral is bounded there. For general pi\_n, one can define its coordinates by smooth matrix coefficients in the usual GJ half-plane.

Suppose \[
                        L(z,\pi_n\otimes\eta^{-1})
                   \text{ is finite at }z=z_0.
 \tag{MT11}
\] Then(MT10) has a well-defined holomorphic vector value at z0 for every Schwartz phi on the entire matrix space, including functions that meet the rank-deficient locus.

Here are the details. The restriction g→phi(X0g) extends to a Schwartz function on M\_n(F). Each smooth dual matrix coefficient of(MT10) is a Godement--Jacquet integral. The published GJ theorem, in the all-irreducibles version recalled in Jiang and Luo(2023), Theorem2.4(1)--(2), pp.313--314, makes it \[
                    L(z,\pi_n\eta^{-1})\,
                              P(q^z,q^{-z})
 \tag{MT12}
\] for a Laurent polynomial P. Its normalization of Haar can change the constant polynomial, but cannot change this holomorphy claim.

The restriction of phi is invariant under left multiplication by a sufficiently small compact open subgroup J of GL\_n.~Take J also in the kernel of eta. Changing variables shows that I\_z is J-fixed wherever the vector integral converges. Admissibility makes \(V_{\pi_n}^J\) finite-dimensional. Choose smooth dual coordinates on that space; (MT12) continues each coordinate and proves holomorphy at z0. This constructs an actual smooth vector, not just an element of the algebraic bidual. It is independent of the choice of J and the coordinates by the identity theorem.

Define \[
              T_z(\phi,w)(h)=I_z(\omega_\eta(1,h)\phi,w).
 \tag{MT13}
\] The exact source covariance is \[
 I_z(\omega_\eta(g_0,1)\phi,\pi_n(g_0)w)
           =|\det g_0|^{z_0-z}I_z(\phi,w).
 \tag{MT14}
\] For \(p=\begin{pmatrix}A&B\\0&D\end{pmatrix}\in P_{n,d}(F)\), the exact target covariance is \[
 I_z(\omega_\eta(1,p)\phi,w)=
   \eta(\det D)|\det A|^{z-1/2}|\det D|^{-n/2}
                       \pi_n(A)I_z(\phi,w).
 \tag{MT15}
\] Both follow by changing g to gg0 or Ag in the convergent region; all eta factors and determinant exponents are displayed in(MT5). They are identities of rational functions in q\^{}z, so remain true at every common holomorphic value.

At z=z0, (MT14) is the literal simultaneous source coinvariant identity, and (MT15) is exactly normalized parabolic induction. Thus \[
  T_{z_0}:\ 
   \bigl(\mathcal S(M_{m,n}(F))\otimes\pi_n\bigr)_{\mathrm{GL}_n(F)}
       \longrightarrow
     \operatorname{Ind}_{P_{n,d}}^{\mathrm{GL}_m}
                    (\pi_n\boxtimes\eta\circ\det_d)
 \tag{MT16}
\] is a genuine target-equivariant map. It is smooth because a compact open subgroup fixing phi on the target fixes its induced function.

If pi\_n and eta are unitary, the target in(MT16) is irreducible by Tadic's unitary induction theorem. Any nonzero map(MT16) is therefore the actual small theta quotient map, up to a single nonzero scalar, by Mínguez(2008), Theorem1. This last statement compares nonvanishing with the actual theta quotient. It does not identify that scalar with a globally pinned arithmetic one.

This construction is useful precisely when an ordinary full-rank Bochner integral no longer converges. It avoids assuming that a map defined only on the open rank stratum automatically extends across its boundary.

\subsection{4. Exact value of the untouched incidence function}\label{exact-value-of-the-untouched-incidence-function}

Now impose(MT6), and take eta unramified. Choose a primitive matrix \(X_*\in M_{m,n}(O)\) with \[
                   vX_*=u,\qquad u=e_n^\vee.
 \tag{MT17}
\] Its columns extend to an integral basis, so \(X_*=k_*X_0\) with \(k_*\in K_m\). The corresponding value of(MT13), at \(h=k_*^{-1}\), is the period obtained by replacing X0 with X*.

Put \[
 F_u(g)=\mathbf1_{M_n(O)}(g)\mathbf1_{\{u\bar g=u\}},\qquad
 F_\nu(g)=\sum_{c\in\kappa^\times}\nu(c)^{-1}
                 \mathbf1_{M_n(O)}(g)\mathbf1_{\{u\bar g=cu\}}.
 \tag{MT18}
\] The restriction of the raw function(MT4) to X*g is exactly F\_u. For w in the line(MT3), right multiplication by an integral diagonal matrix with last entry c shows \[
                I_z(F_\nu,w)=(q-1)I_z(F_u,w).
 \tag{MT19}
\] This equation is initially in the convergent half-plane and then continues. The unit determinant has absolute value1 and its eta value is1; the character in F\_nu cancels the character of w. Moreover F\_nu is bi-J\_n equivariant with character nu\^{}\{-1\}, so its period operator has image in the one-dimensional line(MT3).

We compute its scalar by evaluating the induced vector at1. Absorb the unramified eta into the factors in(MT6), and write \(\sigma'=\sigma\eta^{-1}\), \(\xi'=\xi\eta^{-1}\). The residue of xi' is still nu. Every matrix in the support of F\_nu has a unique block form \[
 g=\begin{pmatrix}A&B\\C&e\end{pmatrix},\quad
 A\in M_{n-1}(O),\ B\in O^{n-1},
 C\in\pi O^{1\times(n-1)},\ e\in O^\times.
\] It is invertible exactly when \(A'=A-Be^{-1}C\) is invertible, and \[
 g=
 \begin{pmatrix}I&Be^{-1}\\0&1\end{pmatrix}
 \begin{pmatrix}A'&0\\0&e\end{pmatrix}
 \begin{pmatrix}I&0\\e^{-1}C&1\end{pmatrix}.
 \tag{MT20}
\] The last factor is in the principal lower congruence subgroup. In the induced model(MT7), \[
 v_0(g)=|\det A'|^{1/2}\sigma'(A')\nu(\bar e)v_\sigma.
 \tag{MT21}
\] The nu factor is canceled by F\_nu. Write \(c_j=\prod_{i=1}^j(1-q^{-i})\). Multiplicative Haar on GL\_j is \(c_j^{-1}|\det g|^{-j}\) times additive matrix Haar. The change \(A\mapsto A'\) has additive Jacobian1. Integrating B,C,e contributes \(q^{-(n-1)}(1-q^{-1})\). Converting the remaining additive measure in A' to normalized GL\_(n−1) Haar therefore contributes the exact constant \[
 \frac{c_{n-1}}{c_n}\,q^{-(n-1)}(1-q^{-1})
                       =\frac{q-1}{q^n-1}=d_n^{-1}.
 \tag{MT22}
\] The remaining determinant exponent is \[
        z+(n-1)/2+1/2-n+(n-1)=z+(n-2)/2.
\] This is exactly the rank n−1 standard GJ exponent. The spherical calculation LIFT16(S8), with vol(K\_(n−1))=1, now gives \[
                  I_z(F_\nu,v_0)=d_n^{-1}
                                   L(z,\sigma')v_0.
 \tag{MT23}
\] One can evaluate against the spherical functional to use(S8): the integral operator is bi-K\_(n−1) and hence acts on its one-dimensional spherical line. Thus no genericity or Whittaker model is required, even if sigma is a non-generic spherical quotient.

Combining(MT19) and(MT23), at every holomorphic value z0, \[
 \boxed{
   T_{z_0}(\phi_{u,v},v_0)(k_*^{-1})
      =\frac{L((m-n+1)/2,\sigma\eta^{-1})}{q^n-1}\,v_0.}
 \tag{MT24}
\] If v0 is replaced by its unit normalization, the same scalar applies. Formula(MT8) still records the norm of the displayed identity-normalized compact section.

For nonzero L and(MT11), this proves that the raw incidence function survives the actual small theta quotient on every nonzero vector of the entire line(MT3), with no change of finite Schwartz input and no supplementary compact projector.

\subsection{5. Application to the actual first image}\label{application-to-the-actual-first-image}

Use the source ranks \[
       a=k,\quad d=t+1>a,\quad n=a+d,\quad m=a+2d.
 \tag{MT25}
\] For K031 these are(42,1760,1802,3562); for K032 they are (168,1634,1802,3436). Let pi\_0 be any actual local constituent of R017's initial function. Its nontrivial T projective-line character gives \[
                         \pi_0=\sigma_0\times\xi,
 \tag{MT26}
\] where sigma0 is unitary spherical of rank a−1 and xi is unitary with the prescribed nontrivial residue. Let eta0,eta1 be the actual unramified unitary splitting characters at this T place.

The actual first local theta constituent is \[
 \pi_1=\pi_0\times(\eta_0\circ\det_d)
       =\bigl(\sigma_0\times\eta_0\circ\det_d\bigr)\times\xi.
 \tag{MT27}
\] Mínguez's parameter theorem gives this identification, and Tadic makes the displayed unitary induction irreducible. The rank n−1 factor \(\sigma_1=\sigma_0\times\eta_0\circ\det_d\) is spherical: normalized induction of its spherical vectors has a nonzero K-fixed vector, and the induction is irreducible. It is also unitary.

The fixed raw first incidence function transports its projective-line character to the actual first output. Explicitly, if \(j\in J_n\) has \(v\bar j=c v\), its target action on(MT4) is the incidence function with source covector cu. For a lift \(g_c\in J_a\) with \(u\bar g_c=cu\), that function is \(\omega(g_c^{-1},1)\phi_{u,v}\). The source coinvariant identity then gives \[
 [\omega(1,j)\phi_{u,v},w]
       =[\phi_{u,v},\pi_0(g_c)w]
       =\nu(c)[\phi_{u,v},w].
 \tag{MT28}
\] Thus every actual first output lies in \(V_{\pi_1}(\nu)\). No change of the first output by a Hecke operator is used. DO036 proves nonvanishing of the complete global first raw kernel; locally its T component is nonzero by the same first-edge test or by(MT24) at ranks a,n.~The line in(MT28) is consequently the one to which the second-edge theorem applies.

It remains to verify all holomorphy and nonzero conditions for the second edge. The unramified Satake parameters of sigma1 are the parameters alpha\_i of sigma0 and the d parameters \[
              \eta_0(\pi)\,q^{(d-1)/2-j},
                                \qquad j=0,\ldots,d-1.
 \tag{MT29}
\] Since sigma0 is unitary spherical of rank a−1, the checked bounded spherical-parameter criterion gives \(|\alpha_i|\le q^{(a-2)/2}\). At \(z_0=(d+1)/2\), with eta1 unitary, \[
 |\alpha_i\eta_1(\pi)^{-1}q^{-z_0}|
       \le q^{-(d-a+3)/2}<q^{-1},\qquad
 |\eta_0\eta_1^{-1}(\pi)q^{(d-1)/2-j-z_0}|
       =q^{-1-j}.
 \tag{MT30}
\] The ramified xi eta1\^{}\{-1\} factor has local L-function1. Therefore \(L(z,\pi_1\eta_1^{-1})=L(z,\sigma_1\eta_1^{-1})\) is finite and nonzero at z0. This proves(MT11) for all Schwartz vectors on the second edge, not only for the test being evaluated.

Consequently, for every nonzero actual first vector w in its prescribed T line, the unchanged raw second T factor has \[
 \boxed{
   T^{(1)}(\phi_T,w)(k_*^{-1})=
       \frac{L((d+1)/2,\sigma_1\eta_1^{-1})}{q^{a+d}-1}\,w
                                  \ne0.}
 \tag{MT31}
\] This is an actual preceding-image statement in the local small theta quotient. It does not merely show that the two compact types occur abstractly somewhere in a representation.

\subsection{6. Explicit size bounds and remaining normalization gates}\label{explicit-size-bounds-and-remaining-normalization-gates}

The scalar in(MT31), denoted lambda\_T, satisfies \[
 \frac{(1+q^{-1})^{-(n-1)}}{q^n-1}
       \le|\lambda_T|
       \le\frac{(1-q^{-1})^{-(n-1)}}{q^n-1}.
 \tag{MT32}
\] This follows by applying \(1-|z|\le|1-z|\le1+|z|\) to the n−1 Euler factors in(MT30). It is a uniform estimate over all initial constituents of the fixed source level and type. The large T prime does not make the basic Haar denominator \(q^n-1\) negligible.

If the parameters and the chosen local normalization are algebraic, the exact11-adic scalar valuation is \[
 v_{11}(\lambda_T)=
 v_{11}\!\left(L((d+1)/2,\sigma_1\eta_1^{-1})\right)
                              -v_{11}(q^n-1).
 \tag{MT33}
\] The first term is retained. No claim that the L-value is an11-unit is implied by its complex nonvanishing or its unitary normalization.

For the exact T primes of R017 the finite replay gives v11(q\^{}a−1)=v11(q\^{}n−1)=0 on K031, and both valuations equal1 on K032. These are only the displayed Haar denominators; the separate L-factor term in(MT33) is still retained.

The theorem settles the fixed local T test on the actual first line. Its normalization is the GJ/normalized-induction convention (MT10), Haar vol(K)=1, and compact section(MT7). Comparison with a globally pinned theta map can multiply it by a nonzero scalar. The global coupled Q and G tests have their separately recorded normalizations. Their assembly into the prescribed arithmetic row lattice, the target mass, all six tied source data, and DP44 costs require those comparisons. The analytic GJ theorem is not an integral-saturation theorem.

The scalar(MT24) is not justified by citing the generic newform test alone: sigma1 is usually non-generic because of its one-dimensional GL\_d block. The direct Schur-complement calculation and the all-irreducibles GJ theorem are essential here.

\subsection{7. Primary sources and check scope}\label{primary-sources-and-check-scope}

\begin{itemize}
\tightlist
\item
  Roche, A. (1998). Types and Hecke algebras for principal series representations of split reductive p-adic groups. Annales scientifiques de l'École Normale Supérieure,31(3),361--413. Residual-characteristic restriction p.379; tame type and Theorem7.7/Remark7.8, pp.391--392; Theorem9.4, p.404. \url{https://www.numdam.org/article/ASENS_1998_4_31_3_361_0.pdf} The theorem is used as a published input; the finite one-line calculation in section2 explains the stronger spherical conclusion.
\item
  Tadic, M. (1986). Classification of unitary representations in irreducible representations of general linear group (non-Archimedean case). Annales scientifiques de l'École Normale Supérieure,19(3),335--382. TheoremD and its block definitions, pp.338--339. The characteristic-zero appendix removes the historical dependency discussed immediately after that theorem. \url{https://www.numdam.org/article/ASENS_1986_4_19_3_335_0.pdf}
\item
  Jiang, D., \& Luo, Z. (2023). Certain Fourier operators and their associated Poisson summation formulae on GL1. Pacific Journal of Mathematics,326(2),301--372. Section2B, equation(2-10), Theorem2.4(1)--(2), pp.313--314. \url{https://msp.org/pjm/2023/326-2/pjm-v326-n2-p05-p.pdf} This is the directly read all-irreducibles restatement of Godement--Jacquet, ChapterI, Theorem3.3; the original monograph is not claimed to have been read in full. Haar constants are computed independently in(MT22) and LIFT16(S8).
\item
  Mínguez, A. (2008). Correspondance de Howe explicite: paires duales de type II. Annales scientifiques de l'École Normale Supérieure,41(5),717--741. Theorem1, printed p.718; actual oscillator normalization(1.4), pp.721--722. \url{https://www.numdam.org/item/ASENS_2008_4_41_5_717_0.pdf}
\item
  LIFT16(S8), LIFT14(E18) for the independently checked spherical Haar normalization and bounded Satake criterion; DO035--036 and R017 for the literal incidence function, actual first input, and its nontrivial T character; the source FS646--649 and actual character definition in LevChiEqOne.lean61,100.
\end{itemize}

The finite diagnostic \nolinkurl{raw_t_projective_period_certificate.py} checks the exact compact character average, block Haar constant and two-edge Satake inequalities. Such computations check the normalization and finite algebra; the GJ continuation, type theorem, and local theta identification are the proofs above and their published inputs. Whole-proof and independent reviews are recorded separately. No Lean build is claimed.

\clearpage

\section{AM-K035 --- an explicit mixed CM period and source Fock-frame adapter}\label{am-k035-an-explicit-mixed-cm-period-and-source-fock-frame-adapter}

Author: GPT-6 Astra Pro, ALP\_MACHINE /root/kummer\_extension, 22 November 2026 UTC; inherited reasoning effort is not separately exposed.

This finding supplies an actual equation downstream of the failed quadratic-field constructor. It identifies the period vector of the principally polarized hosts \(J^{14}\) and \(J^{56}\), computes their source-style embedding-lattice coordinate generators, and calculates the exact change from the source's unit-normalized creation operators to algebraic reduced-theta derivatives. The resulting all-polynomial equation has the literal shape of ClsK.Core, including its per-place RealD.sc. It does not inhabit the old Fields.Kq or elliptic VacIn records, and it does not identify an independently specified later arithmetic row map by definition.

Carried inputs are the accepted DP57 and canonical-gauge existence theorems, AM-R-004 and AM-R-006, ALP-LIFT-09, AM-K018, AM-K030/032--034, and AM-DO-017/018/022/026/027. Their proved conclusions are used at their recorded scope. The old common-quadratic obstruction is settled by AM-R-003 and is not proved again here.

\subsection{1. The exact source equations being replaced}\label{the-exact-source-equations-being-replaced}

The uniform analytic restriction occurs before the generic class-kernel equations:

\begin{itemize}
\tightlist
\item
  \nolinkurl{Development/AbelianVarietyTheory/Definitions/LatticeFrames/SplitPart.lean,} lines22--24 and38--49, uses one complex period \(\delta\), one real metric scalar sc, and the hypotheses \[
  \tau(\xi_0)=i/\sqrt d,\qquad \mathrm{sc}\,|\delta|^2=2/\sqrt d
  \quad\text{for every }\tau.
  \tag{PF1}
  \]
\item
  \nolinkurl{Development/AbelianVarietyTheory/Definitions/AutomorphicMembership/IsDiscriminantTwelfthRootOfIdealLatticeIffDelta1.lean,} lines53--67, defines the algebraic embedding lattice mmM, its coordinate-line intersection ellT, and IsDgen. This coordinate generator is logically distinct from the elliptic period in its subsequent VacIn.
\item
  \nolinkurl{Formalism/LocalToGlobal/IndefiniteReal/LocalNewformFunction.lean,} lines124--141, defines crea and the actual normalizing factor \[
  \mathrm{fockScale}(\tau)=2\sqrt{2\pi|\tau(\xi_0)|},
  \qquad a_\tau^+=\mathrm{fockScale}(\tau)^{-1}\mathrm{crea}_\tau.
  \tag{PF2}
  \] DatumVec.lean, lines130--134, uses these operators in pureTensorO. They are not the unnormalized operators \(r_\tau\) of ALP-LIFT-09.
\item
  \nolinkurl{Development/AbelianVarietyTheory/Definitions/LatticeFrames/OnFramesRdApply.lean,} lines25--39 and108--121, already allows a nonzero complex scalar RealD.sc(j,τ) separately at every place. Its Core equation is \[
  v_j\,\vartheta_{\mathrm{pureTensorO}(P)}(g_c,h_{c'})
   =c_{K,j}\,\mathrm{abelian}_2(c,c')
      \bigl(\mathrm{mvZ}(1,\operatorname{diag}(\mathrm{sc}_{j,\tau}))P\bigr).
  \tag{PF3}
  \] The following Kval clause is a separate equality for the prescribed arithmetic map \(t.K_j\).
\end{itemize}

Thus a common scalar as in (PF1) is not required by the displayed downstream equation (PF3). An exact varying period and metric profile can be tested there without asserting a new instance of the old constructor.

\subsection{2. Actual periods and metric on the two product hosts}\label{actual-periods-and-metric-on-the-two-product-hosts}

Write \(E=\mathbb Q(\zeta_7)\), \(E_0=E(i)\), and retain the degree-seven real field \(B\subset\mathbb Q(\zeta_{29})\), positive generator \(\pi_B\), and trace-orthonormal basis \(e_0,\ldots,e_6\) from AM-K030. There are two cases: \[
\begin{array}{c|c|c|c}
h&R&K_h&A_h\\ \hline
1&\mathbb Q&E_0B&J^{14}\\
4&\mathbb Q(\sqrt5,\sqrt{13})&E_0BR&J^{56}.
\end{array}
\tag{PF4}
\] The full CM type is indexed by \[
\tau=(i,\varepsilon,l,\rho),\quad
i\in\{1,2,3\},\quad \varepsilon\in\{+1,-1\},\quad
l\in\mathbb Z/7\mathbb Z,
\tag{PF5}
\] where \(\rho\) runs over the real embeddings of \(R\), omitted if \(h=1\). The restrictions are \(\zeta_7\mapsto\zeta_7^i\), \(i_{\mathbb C}\mapsto\varepsilon i_{\mathbb C}\), and the l-th B-embedding. Put \[
\begin{split}
\pi_l&=\sigma^l(\pi_B),\\
\delta_\rho&=1\quad(h=1),\qquad
\delta_\rho=\rho(\delta_1\delta_3)\quad(h=4),\\
\xi_h&=\frac{\zeta_7-\zeta_7^{-1}}{14}
                 \pi_B^{-6}\delta_1\delta_3
       \quad\text{with the last factor omitted if }h=1.
\end{split}
\tag{PF6}
\] All \(\pi_l,\delta_\rho\) are positive. Define the actual periods \[
\Omega_i=-\frac{2i_{\mathbb C}}7
          B(i/7,1/2)(\zeta_7^i-1),\qquad
\Omega_\tau=\Omega_i,
\tag{PF7}
\] and retain the actual Hodge Gram coefficients of ALP-LIFT-03 and AM-R-004: \[
G_i=\frac{2\pi}{7}
 \frac{\Gamma(i/7)\Gamma(1/2-i/7)}
      {\Gamma(1-i/7)\Gamma(1/2+i/7)}>0.
\tag{PF8}
\]

\subsubsection{Theorem PF9: an actual period-vector replacement}\label{theorem-pf9-an-actual-period-vector-replacement}

There are explicit algebraic invariant coordinates \(\zeta_\tau\) on the universal cover of \(A_h\) in which \[
\boxed{\Lambda_h=\operatorname{diag}(\Omega_\tau)
                       \Psi(\mathcal O_{K_h}),\qquad
H_h(\zeta,\zeta)=\sum_\tau h_\tau|\zeta_\tau|^2,}
\tag{PF9}
\] where \[
h_\tau=\frac{\delta_\rho\pi_l^{-6}}{2G_i},\qquad
q_\tau=-2i_{\mathbb C}\tau(\xi_h)>0,\qquad
\boxed{h_\tau|\Omega_\tau|^2=q_\tau.}
\tag{PF10}
\] The same formulas apply entry by entry on the Hom torus. Its period lattice for actual moved modules \(N,N'\) is \(\operatorname{diag}(\Omega_\tau)\Psi(\operatorname{Hom}_{\mathcal O_{K_h}}(N,N'))\).

\textbf{Proof.} On a pair of ordinary Jacobian factors use \[
w_{i,\varepsilon,j}=u_{i,j}+\varepsilon i_{\mathbb C}v_{i,j}.
\] Its period lattice is \(\Omega_i\Psi(\mathcal O_{E_0})\), directly from AM-R-004 and the integral identification \(\mathcal O_{E_0}=\mathcal O_E[i]\). For the seven B-factors put \[
z_{i,\varepsilon,l}=\sum_{j=0}^6\sigma^l(e_j)w_{i,\varepsilon,j},
\qquad
\zeta_{i,\varepsilon,l}=\pi_l^3z_{i,\varepsilon,l}.
\tag{PF11}
\] The first transformation identifies the lattice with \(\Omega_i\Psi(\pi_B^{-3}\mathcal O_{K_1})\), because the \(e_j\) form an actual integral basis of that fractional ideal. The second gives the lattice in (PF9). Orthogonality of the first matrix and the displayed rescaling give the metric when \(h=1\).

For \(h=4\), apply the embedding matrices of the integral bases \((1,u_1)\), \((1,u_3)\) in AM-K032. Their exact identities \[
V_a^{\mathsf T}\operatorname{diag}(\delta_{a,+},\delta_{a,-})V_a=I_2
\] give the factor \(\delta_\rho\), and their integral tensor product gives \(\mathcal O_{K_4}\). No isogeny or unspecified period factor occurs.

For the final normalization put \(x=i/7\). The beta-gamma and reflection formulas give \[
\begin{split}
\frac{|\Omega_i|^2}{G_i}
&=\frac87
\frac{\Gamma(x)\Gamma(1-x)}
     {\Gamma(1/2+x)\Gamma(1/2-x)}
\sin^2(\pi x)\\
&=\frac87\sin(\pi x)\cos(\pi x)
=\frac47\sin(2\pi i/7).
\end{split}
\tag{PF12}
\] Since \(\tau(\xi_h)=i_{\mathbb C}\sin(2\pi i/7)\delta_\rho\pi_l^{-6}/7\), (PF10) follows. The Hom assertion is the same scalar embedding entry by entry. The alternating-form convention is \[
\operatorname{Im}H_h(X,Y)
 =-\operatorname{Tr}_{K_h/\mathbb Q}
       \bigl(\xi_h\operatorname{tr}(Y^*X)\bigr),
\tag{PF13}
\] exactly as in AM-DO-022. ∎

\subsection{3. Exact coordinate generators for the embedding lattice}\label{exact-coordinate-generators-for-the-embedding-lattice}

Let \(M=K_h\), or any number field containing it, and \(R_M=\mathcal O_M\). The fields \(K_h\) are Galois, so all selected embeddings land in M. Define \[
\mathcal M_\Psi
   =\operatorname{span}_{R_M}
          \{(\tau(a))_{\tau\in\Psi}:a\in\mathcal O_{K_h}\}
       \subset M^{42h},
\qquad
\ell_\tau=\{c\in M:ce_\tau\in\mathcal M_\Psi\}.
\tag{PF14}
\] This is the algebraic shape of mmM and ellT for the actual standard ideal \(\mathcal O_{K_h}\); it contains no complex period. Put \[
d_i=\prod_{\substack{j\in\{1,2,3\}\\j\ne i}}
                       (\zeta_7^i-\zeta_7^j),
\] \[
\boxed{cD_\tau
   =d_i(2\varepsilon i_{\mathbb C})\pi_l^6
      \begin{cases}
      1,&h=1,\\
      \rho(\sqrt5)\rho(\sqrt{13}),&h=4.
      \end{cases}}
\tag{PF15}
\]

\subsubsection{Theorem PF16: the exact source-style generator}\label{theorem-pf16-the-exact-source-style-generator}

\[
\boxed{\ell_\tau=cD_\tau R_M.}
\tag{PF16}
\] For normalized extensions of rational valuations, \[
\begin{array}{c|ccccc}
p&2&7&29&5,13&\text{all others}\\ \hline
v_p(cD_\tau),\ h=1&1&1/3&6/7&0&0\\
v_p(cD_\tau),\ h=4&1&1/3&6/7&1/2&0.
\end{array}
\tag{PF17}
\] In particular all \(cD_\tau\) are units at11.

\textbf{Partial E-factor.} Let \(r_i=\zeta_7^i\), \(i=1,2,3\). The image after \(R_M\)-span is polynomial evaluation at these three roots. Division by the monic \(F(T)=\prod_{j=1}^3(T-r_j)\) reduces it to degree at most two. If this polynomial vanishes at the two roots other than \(r_i\), successive monic division shows that it is a scalar over \(R_M\) times \(\prod_{j\ne i}(T-r_j)\). Thus its possible values at \(r_i\), with the others zero, form precisely \(F'(r_i)R_M=d_iR_M\). This calculation uses the selected three embeddings; the full E-different would give the wrong ideal.

\textbf{Fully selected factors.} For a number field L with all embeddings in M, the dual of its full embedding lattice for the coordinate dot product is the embedding lattice of \(\mathfrak D_{L/\mathbb Q}^{-1}\), after \(R_M\)-span. This follows by writing a basis and its trace-dual basis. If \(\mathfrak D_L=(\Delta_L)\), then \[
ce_\sigma\in\operatorname{span}_{R_M}\operatorname{Emb}(\mathcal O_L)
\ \Longleftrightarrow\
c\,\sigma(\Delta_L^{-1}\mathcal O_L)\subset R_M
\ \Longleftrightarrow\
c\in\sigma(\Delta_L)R_M.
\tag{PF18}
\] The factors are therefore \(2\varepsilon i_{\mathbb C}\) for \(\mathbb Q(i)\), \(\pi_l^6\) for B, and \(\rho(\sqrt5),\rho(\sqrt{13})\) for the real quadratics. The B-different is the exact result of AM-K030; no integral power basis is assumed.

\textbf{Tensor-product step.} The coprime discriminants give the integral tensor decompositions established in AM-K030/032. Their embedding matrices are Kronecker products. For full-rank lattices \(L_1\subset M^m,L_2\subset M^n\) over a Dedekind domain, the coordinate-line ideal of \(L_1\otimes L_2\) at \((i,j)\) is the product of the two coordinate-line ideals. To prove this, localize at a prime and use basis matrices A,B over the DVR. Then \[
(A\otimes B)^{-1}(e_i\otimes e_j)
     =(A^{-1}e_i)\otimes(B^{-1}e_j).
\] The smallest entry valuation is the sum of the two minima, so the valuation threshold for multiplication by c to become integral is the sum of the two thresholds. Equality at every localization proves the ideal equality and hence (PF16).

Each nonzero difference of seventh roots has normalized 7-valuation \(1/6\) and no other finite support; \(\pi_B\) has 29-valuation \(1/7\) and no other support. The remaining different generators give the other entries in (PF17). ∎

This is an exact solution of the embedding-coordinate generator problem with \(M=K_h\). It does not assert that these are primitive Hodge bases at bad primes or that the reduced frame is integral. ALP-LIFT-08's section contents and K033's reduced-gauge denominators remain. At11 the change to jetC(cD,−) multiplies each fixed multidegree by a unit.

\subsection{4. The literal source creation factor, including its sign}\label{the-literal-source-creation-factor-including-its-sign}

The source convention in \nolinkurl{BesselHermitianTheta/TiedWeil.lean59–62} is \(\psi_\infty(t)=\exp(-2\pi it)\). In one CM coordinate \(w=x+iy\), \[
b_{\rm src}(w,v)=\operatorname{Tr}_{\mathbb C/\mathbb R}
       (\tau(\xi_h)w\bar v)
       =-q_\tau\operatorname{Im}(w\bar v).
\] Thus the literal heisO in \nolinkurl{LocalNewformFunction.lean29–35} is \[
T_v f(w)=e^{\pi iq_\tau\operatorname{Im}(w\bar v)}f(w+v).
\tag{PF19}
\] Its differentiated operators are \[
D_1=\partial_x+\pi iq_\tau y,\quad
D_i=\partial_y-\pi iq_\tau x,\quad
\mathrm{annih}=2\partial_{\bar w}+\pi q_\tau w,\quad
\mathrm{crea}=2\partial_w-\pi q_\tau\bar w.
\tag{PF20}
\] In particular a Fock-vacuum theta function is a Gaussian times a holomorphic function, with the positive Gaussian sign shown here.

For the complete Hom space define the mixed reading \[
(\mathcal R_\Omega f)(u)
 =\exp\!\left(\frac\pi2\sum_{\tau,r,s}h_\tau|u_{\tau,r,s}|^2\right)
       f\bigl((-\Omega_\tau^{-1}u_{\tau,r,s})_{\tau,r,s}\bigr).
\tag{PF21}
\] The finite component is held fixed. This is the actual varying-period counterpart of xiN.

\subsubsection{Theorem PF22: exact source polynomial normalization}\label{theorem-pf22-exact-source-polynomial-normalization}

Let \(F=\mathcal R_\Omega f\) be holomorphic and P any complex polynomial in the Hom creation labels. With identity source/target orthonormal frames, use the literal pureTensorO of (PF2). Then \[
\boxed{[\mathrm{pureTensorO}(f,P)](0)
 =\left[P(\lambda_\tau\partial_{\tau,r,s})F\right](0),\qquad
\lambda_\tau=-\frac{\Omega_\tau}{\sqrt{\pi q_\tau}}.}
\tag{PF22}
\] The square root is positive real. Every \(\lambda_\tau\ne0\), and \[
\boxed{|\lambda_\tau|^2\,\pi h_\tau=1.}
\tag{PF23}
\]

\textbf{Proof.} Conjugating crea by the Gaussian gives \[
e^{\pi q|w|^2/2}\,\mathrm{crea}\,
 e^{-\pi q|w|^2/2}=2\partial_w-2\pi q\bar w.
\tag{PF24}
\] The summands commute since \(\partial_w\bar w=0\). In its n-th power every term containing \(\bar w\) vanishes at zero. The remaining term is \(2^n\partial_w^n\), and \(u=-\Omega_\tau w\) contributes \((-\Omega_\tau)^n\). Dividing by \(\mathrm{fockScale}^n=(2\sqrt{\pi q_\tau})^n\) gives the displayed coefficient. Distinct matrix coordinates commute, proving every monomial and then every polynomial. Equation(PF23) follows from(PF10). ∎

There is no hidden parity change in(PF21). In the normalized coordinates of DO-018 let \(z=A+\tau_{\rm per}B\), \(Y=\operatorname{Im}\tau_{\rm per}\). The archimedean Schrödinger action is \[
\rho_\infty(A,B)\varphi(x)
 =e^{-2\pi iA^{\mathsf T}x+\pi iA^{\mathsf T}B}\varphi(x-B).
\] On \(\varphi_\infty(x)=e^{\pi i x^{\mathsf T}\tau_{\rm per}x}\), theta evaluation gives \[
\begin{split}
\vartheta_{\rho_\infty(A,B)\varphi_\infty\otimes\varphi_f}
 &=e^{\pi iB^{\mathsf T}z}
      \sum_q\varphi_f(q)e^{\pi iq^{\mathsf T}\tau_{\rm per}q-2\pi iq^{\mathsf T}z}\\
 &=e^{-\pi H(z,z)/2}\,R\varphi_f(-z).
\end{split}
\tag{PF25}
\] The second equality is the exact identity \[
-\frac\pi2 z^{\mathsf T}Y^{-1}\bar z
+\frac\pi2 z^{\mathsf T}Y^{-1}z=\pi iB^{\mathsf T}z.
\] Thus(PF21) returns precisely the DO-018/026 reduced function \(R\varphi_f\), with its original finite translations and phases. It applies to every finite Schwartz input. The simultaneous transport \(b_{\rm src}=-\operatorname{Im}H\), \(\psi_{\rm src}=\exp(-2\pi i\,\cdot)\) gives the same realized Heisenberg multiplier. No relabeling of finite translations is needed.

The actual algebraic-coordinate creation operator \(r_\tau\) of ALP-LIFT-09 has commutator/norm coefficient \(\pi h_\tau\); the source operator is \(\lambda_\tau r_\tau\), with coefficient1 by(PF23). Passing from K033/K034 to source pureTensorO must retain these complex period factors. No finite valuation is assigned to them.

\subsection{5. The all-polynomial arithmetic frame equation}\label{the-all-polynomial-arithmetic-frame-equation}

For a reduced germ F set \[
\mathcal J_{cD,F}(P)
 =\sum_m P_m\left(\prod_{\tau,r,s}cD_\tau^{m_{\tau,r,s}}\right)
                         \partial^m F(0).
\tag{PF26}
\] This is exactly jetL(cD,F): jetC uses ordinary derivatives, not Taylor coefficients divided by \(m!\). Let \(\mathcal D_s\) multiply every variable at τ by \(s_\tau\); the literal mvZ(1,diag(s)) does this. Monomial evaluation gives \[
\boxed{\mathcal J_{cD,F}(\mathcal D_{\lambda/cD}P)
                   =P(\lambda_\tau\partial_\tau)F(0).}
\tag{PF27}
\]

Fix actual finite class representatives \(g_c,h_{c'}\), a finite Schwartz datum \(\varphi_d\), and the independently constructed finite Mumford averages \(F_{d;c,c'}\) of DO-026 on \(A_h^{ab}\). They are defined geometrically before the analytic comparison. Let \(\epsilon_\chi(c,c')\) be the actual determinant-root splitting factor. The seed is the product of \(14hab\) copies of \(\sigma_q\), and its fixed comparison period is \(C^{14hab}\), \(C=\theta_0c\), with the phase fixed in AM-R-006. The product and finite Mumford intertwining retain this same period for every moved lattice and every polynomial. Consequently \[
\boxed{
\vartheta_{\mathrm{pureTensorO}(\varphi_d,P)}(g_c,h_{c'})
 =C^{14hab}\epsilon_\chi(c,c')\,
      \mathcal J_{cD,F_{d;c,c'}}
               (\mathcal D_{\lambda/cD}P).
}
\tag{PF28}
\] The finite Weil-phase premise is exactly the fixed unramified-CM realization of DO-022 and its product-host transport. It is not an assertion about arbitrarily chosen nonzero vacuum vectors. All finite phases and normalized-average indices of DO-026 remain in F.

The required varying source frame is explicit: \[
\boxed{\mathrm{sc}_{j,\tau}
       =-\frac{\Omega_\tau}{cD_\tau\sqrt{\pi q_\tau}}.}
\tag{PF29}
\] For a stage it is independent of degree and finite chart. Take \[
\mathrm{abelian}^{\rm geom}_2(c,c')
  =\epsilon_\chi(c,c')\mathcal J_{cD,F_{d;c,c'}}.
\] With \(c_K=v_jC^{14hab}\), equation(PF28) is(PF3) for this actual geometric family. If the chart characters have already been separated as in LevChiEqOne.abelian₂, use those characters and the common root-of-unity normalization of DO-027; its scalar criterion is unchanged. This is not a definition by division of an analytic value by a desired period: the sections and finite averages determine the functional before the comparison.

The numerical period/frame comparison left in DO-026§5 is therefore explicit for these kernels. The varying profile(PF10) replaces the uniform quadratic scale, with the source unit-Fock normalization preserved.

\subsection{6. One field for all jet orders in a fixed finite chart family}\label{one-field-for-all-jet-orders-in-a-fixed-finite-chart-family}

\subsubsection{Proposition PF30}\label{proposition-pf30}

Fix finitely many actual moved module pairs, finite translations, and character/Schwartz coefficients in(PF28). There is one number field \(M_1\), independent of derivative order, containing all values \[
\mathcal J_{cD,F_{d;c,c'}}(P),\qquad P\in M_1[Z_{\tau,r,s}].
\tag{PF30}
\]

\textbf{Proof.} The product ppav, full endomorphism action, rigidified trace-one line and normalized algebraic section are defined over a number field. Each commensurable CM lattice in the finite list is algebraically isogenous to that product: clearing denominators gives an isogeny with finite torsion kernel, whose algebraic quotient realizes the analytic torus. After a finite extension the kernels, rigidified polarizations and finitely many line isomorphisms are defined. Enlarge once for the finitely many torsion points, finite coefficients, roots of unity and square roots of averaging indices.

Choose algebraic formal frames and invariant tangent bases over this field. Their formal group laws, addition and adjunction cocycles have coefficients there. The full diagonal \(\mathcal O_E\)-action on every isogenous ordinary-Jacobian product retains holomorphic labels1,2,3 with repetitions, without opposite pairs. The accepted existence theorem and same-frame uniqueness in K018 identify the algebraic canonical gauge with the reduced Chern gauge over that same field in all orders. The rigidified translation gauges have the same algebraic descent; K034 gives the stronger integral statement at11 for prime-to11 torsion.

Hence the finitely many reduced section germs are complete series over one number field, rather than requiring a different extension for each coefficient. Invariant logarithm substitutions also have coefficients there in characteristic zero. Adjoin the finitely many explicit \(cD_\tau\) if necessary. All formal derivatives and finite sums in(PF26) remain in this field. ∎

No effective discriminant or degree bound is asserted for \(M_1\). This is not yet literal Amb.JetsRat in ThDistC.lean239--243: that statement quantifies over \textbf{every} selfdual \(N,N'\), all admitted \(c,x\), with one preassigned X.K. Passing from a finite class family to that uniform statement needs its prescribed algebraic transport and character descent. Adjoining an unrelated field for each lattice would not prove it.

\subsection{7. The next exact arithmetic equations and cost scope}\label{the-next-exact-arithmetic-equations-and-cost-scope}

The calculation separates three obligations.

\begin{enumerate}
\def\labelenumi{\arabic{enumi}.}
\tightlist
\item
  \textbf{Solved for the geometric family:} the periods, embedding-line generators, literal source creation factors, and all-polynomial equation(PF28). They do not need a common imaginary quadratic subfield.
\item
  \textbf{Still required for the prescribed arithmetic row:} identify the geometric \(\epsilon_\chi\mathcal J_{cD,F}\) with the independently fixed row functional and map. The exact next equation is ClsK.Kval, \nolinkurl{OnFramesRdApply.lean115–121:} \[
  \operatorname{val}_{j+1}(t.K_j\psi,c')
   =\sum_c|\operatorname{Stab}(c)|^{-1}
      \mathrm{abelian}_2(c,c')(\mathrm{subY}(F_c))
           \det(y')^{-s_D},
  \tag{PF31}
  \] with the actual selected polynomial, source row, target row and determinant factors. DO-026 computes this class integral for its actual operator, but local type alone does not identify an externally specified \(t.K_j\) or the integral saturation of a later image.
\item
  \textbf{Still required for a uniform arithmetic constructor:} Amb.JetsRat with its fixed descent field and the remaining constructor hypotheses or separately proved replacements. The old false Fields.Kq record remains false; PF30 is a finite-family result.
\end{enumerate}

At11 all cD factors are units, so K033/K034's selected arithmetic content bounds hold unchanged for jetC(cD,−). Their matching slope1/1220 is before multiplication by the complex factors(PF29). The Fock norm is preserved exactly by(PF23), and the analytic common period is retained. No full DP44 cost estimate, later-image saturation, terminal identity, or every-\(\varepsilon\) window follows from this adapter.

\subsection{Sources and verification}\label{sources-and-verification}

The primary period source is Tadokoro, Y. (2014), \emph{The period matrix of the hyperelliptic curve \(w^2=z^{2g+1}-1\)}, \href{https://arxiv.org/abs/1211.6910v5}{arXiv:1211.6910v5}, Propositions3.2 and4.1 and Theorem4.4, as instantiated and checked in AM-R-004. Reduced theta multipliers and normalized translation are as in Bannai, K., \& Kobayashi, S. (2010), \emph{Algebraic theta functions and p-adic interpolation of Eisenstein--Kronecker numbers}, Duke Mathematical Journal, 153(2),229--295, Definition1.8 and§1.3, \href{https://arxiv.org/pdf/math/0610163}{author preprint}. Both bibliographic titles were checked against their primary arXiv records during this review. The creation computation is directly from the recovered literal definitions above; it is not an asserted extension of a source theorem beyond its hypotheses.

The companion certificate checks exact cyclotomic interpolation/norms, finite scalar/jet normalizations, and the source-sign Gaussian identity as polynomial identities. It is a finite form check, not a substitute for the all-order proofs, a transitive audit of accepted theorems, or a Lean build. Actual hashes and check events are recorded separately.

\clearpage

\section{AM-DO-037: the prescribed coupled Q family survives the actual second local theta quotient}\label{am-do-037-the-prescribed-coupled-q-family-survives-the-actual-second-local-theta-quotient}

Author: GPT-6 Astra Pro, ALP\_MACHINE /root/definite\_oscillator. 22 November 2026 UTC. Inherited reasoning effort is not separately exposed.

The literal two-edge Q family for m=2 has a nonzero image under the actual successive local theta quotients. The initial representation is any irreducible unitary spherical representation of GL\_k at that Q place. A specified unipotent target evaluation, followed by evaluation in the first induced model, is an explicit nonzero scalar times the initial spherical vector. The proof retains the entire prescribed Y sum and both raw kernels. It does not replace either kernel by an unrestricted Schwartz input.

The actual source definition of DegOne forces the local field at every tuple Q place to be Q\_q. Thus this hypothesis is realized on both the K031 and K032 hosts, even though those hosts have absolutely ramified places elsewhere. The exact local quotient scalar below is normalized by explicit Haar integrals. Identification with the separately pinned global arithmetic row and terminal mass is not asserted.

\subsection{1. The literal Q-place hypothesis includes absolute unramifiedness}\label{the-literal-q-place-hypothesis-includes-absolute-unramifiedness}

All source locators refer to the pinned source tree at commit 7cbf7a0c32fca1badf05931879a37c3f8a2b03b4. In \nolinkurl{Formalism/LocalToGlobal/Divisibility/Tuples.lean:}

\begin{longtable}[]{@{}
  >{\raggedright\arraybackslash}p{(\columnwidth - 2\tabcolsep) * \real{0.5000}}
  >{\raggedright\arraybackslash}p{(\columnwidth - 2\tabcolsep) * \real{0.5000}}@{}}
\toprule\noalign{}
\begin{minipage}[b]{\linewidth}\raggedright
Lines
\end{minipage} & \begin{minipage}[b]{\linewidth}\raggedright
Literal condition
\end{minipage} \\
\midrule\noalign{}
\endhead
\bottomrule\noalign{}
\endlastfoot
58--60 & SplitAt: the extension of the real-field prime is a product of two distinct CM primes \\
64--65 & DegOne(v): its absolute norm is prime, and that rational prime is not in v squared \\
67 & GoodPl is SplitAt, DegOne, and avoidance of the prescribed rational set S \\
112--114 & A tuple kills local inertia at every place which is not GoodPl \\
128 & Qprime is the set where that inertia character is nontrivial \\
149--161 & IsResChar and resChar give its tame residue character \\
198--200 & IsTupleMap excludes the coefficient-place rational prime from the corresponding Q set \\
\end{longtable}

Let v belong to Qprime of a tuple. The last condition in tupGrp implies GoodPl(v): otherwise the inertia character would be trivial. Write q=Norm(v). The first DegOne clause says q is a rational prime, so the residue degree is one. In the localization O\_(F,v), the valuation of the rational prime q is the absolute ramification index e. The second clause q notin v\^{}2 says e=1. Hence

\[
                 [F_v:\mathbf Q_q]=ef=1.
 \tag{DO37.1}
\]

Relative splitting gives E\_P=F\_v for either selected CM prime P. Thus the precise local hypothesis used in LIFT23 is automatic. In particular the absolutely ramified29 places of K031/K032 cannot occur in Q; no replacement of q by a ramified uniformizer is needed for this literal tuple family.

The tuple is trivial on base-field ideles and kills wild inertia. At a split place its two residue characters are therefore inverses. Since the place belongs to Q, both are nontrivial. Consequently the selected residue character chi used below is nontrivial. If the tuple is evaluated at a coefficient place over11 through IsTupleMap, then q is different from11 as well. For a bare tuple without that additional condition the local theorem below still holds at q=11, but its Gauss sum is not claimed an11-unit.

This argument checks an actual definition. It does not use the false implication that residue degree one alone gives F\_v=Q\_q.

\subsection{2. A whole-Schwartz normalized quotient map}\label{a-whole-schwartz-normalized-quotient-map}

This section records the precise published interface needed to pass from an explicit integral to the actual small theta quotient. Let n\textless m, delta=m-n, G\_n=GL\_n(F), K\_n=GL\_n(O), and

\[
 \kappa_n(q)=\prod_{i=1}^{n}(1-q^{-i}),\qquad
 \operatorname{vol}(K_n)=1,\qquad
 \operatorname{vol}(M_n(O))=1.
 \tag{DO37.2}
\]

Thus multiplicative matrix Haar is dg=kappa\_n(q)\^{}(-1)\textbar det g\textbar\^{}(-n)d\_add g. Let eta be unitary and pi an irreducible unitary representation of G\_n.~Use the actual split Hom action

\[
 [\omega_\eta(g,h)\phi](X)=
 \eta(\det h)\eta(\det g)^{-1}
 |\det h|^{-n/2}|\det g|^{m/2}\phi(h^{-1}Xg).
 \tag{DO37.3}
\]

For E\_n={[}I\_n;0{]} and Re(z) sufficiently large, put

\[
 I_z(\phi,v)=\int_{G_n}
       \phi(E_ng)|\det g|^{z+(n-1)/2}
                 \eta(\det g)^{-1}\pi(g)v\,dg.
 \tag{DO37.4}
\]

This is defined on the whole Schwartz space M\_(m,n), not just on its full-rank open stratum. Restriction to E\_n M\_n(F) is a Schwartz function on M\_n(F). Pairing (DO37.4) with any smooth dual vector is therefore a Godement--Jacquet zeta integral. Jiang and Luo (2023), Theorem2.4(1)--(2), printed pp.313--314, states that every such integral is L(z,pi eta\^{}(-1)) times a Laurent polynomial in q\^{}z.

The vector-valued assertion follows without a weak-versus-strong holomorphy assumption. For fixed phi, the restricted Schwartz function on M\_n(F) is invariant under left multiplication by a sufficiently small open compact J contained in ker(eta det). Changing variables shows I\_z(phi,v) lies in pi\^{}J for every z where it converges. Admissibility makes pi\^{}J finite dimensional. Apply the scalar theorem to finitely many smooth dual vectors separating this space. Thus I\_z is a meromorphic vector with no pole wherever L(z,pi eta\^{}(-1)) is finite.

Direct changes of variables, first g to gg\_0 and then g to At, give

\[
 I_z(\omega(g_0,1)\phi,\pi(g_0)v)
       =|\det g_0|^{(\delta+1)/2-z}I_z(\phi,v),
 \tag{DO37.5}
\]

and, for p=block(A,B;0,D) in the upper parabolic P\_(n,delta),

\[
 I_z(\omega(1,p)\phi,v)=
 \eta(\det D)|\det A|^{z-1/2}|\det D|^{-n/2}
                       \pi(A)I_z(\phi,v).
 \tag{DO37.6}
\]

At z\_0=(delta+1)/2 these are simultaneous source coinvariance and precisely normalized parabolic covariance. Consequently

\[
 T_{\eta,\pi}(\phi\otimes v)(h)
                  = I_{z_0}(\omega(1,h)\phi,v)
 \tag{DO37.7}
\]

is, when L is finite there, an actual map from the whole oscillator tensor pi to

\[
           \pi\times(\eta\circ\det_\delta)
      =\operatorname{Ind}_{P_{n,\delta}}^{G_m}
                    (\pi\boxtimes\eta\circ\det_\delta),
 \tag{DO37.8}
\]

where induction is normalized. It intertwines the target with right translation. Holomorphy in the induced compact model is also finite dimensional for a fixed input: phi has a fixed open compact target stabilizer, so only finitely many compact cosets and the fixed finite-dimensional pi\^{}J are needed.

Tadic (1986), TheoremD(i)--(ii), printed pp.338--339, implies that normalized induction of irreducible unitary representations of general linear groups is irreducible and unitary. Indeed express the factors as products of the elementary unitary factors in part(ii), concatenate them, and apply part(i). Hence (DO37.8) is irreducible. Once (DO37.7) is nonzero, Minguez (2008), Theorem1(1), printed p.718, identifies it with the actual irreducible local theta quotient, with multiplicity one.

The source is paired bilinearly. Thus the simultaneous coinvariant uses pi, while the source irreducible quotient of the oscillator alone is pi contragredient; this is the same Hom convention as DO028. The old parameter is therefore pi in (DO37.8). No source contragredient has been dropped.

Only the holomorphic-multiple assertion of Jiang--Luo is used. Its displayed Gaussian convention is not used to fix Haar constants; all kappa factors below follow from (DO37.2). This avoids an implicit conversion between additive and maximal-compact probability Haar measures.

\subsection{3. The representations and the two raw kernels}\label{the-representations-and-the-two-raw-kernels}

At the actual Q place put

\[
 a=k,\quad r=a+1,\quad d=t+1,\quad b=a+d=r+t,\quad N=b+d.
 \tag{DO37.9}
\]

Let pi\_0 be an irreducible unitary spherical representation of GL\_a(F), and take v in pi\_0\^{}(K\_a). Write eta\_0,eta\_1 for the two actual common stage labels. Their residue characters at Q are the same chi, since each is the unramified reference label times the same tuple character. Set

\[
                  \vartheta=\eta_0\eta_1^{-1}.
 \tag{DO37.10}
\]

This is an unramified unitary character. In particular \textbar vartheta(q)\textbar=1, with \textbar q\textbar=q\^{}(-1).

For the first quotient, L(z,pi\_0 eta\_0\^{}(-1))=1: the parameter of pi\_0 is unramified, whereas twisting by the nontrivial inertia character chi\^{}(-1) leaves no inertia invariants. The map (DO37.7) is holomorphic at z=(d+1)/2 and has target

\[
                  \pi_1=\pi_0\times(\eta_0\circ\det_d).
 \tag{DO37.11}
\]

It is nonzero, for example on a Schwartz function supported near E\_a K\_a with the matching eta\_0 character. Thus it is the actual first local theta quotient by Section2.

For the second quotient,

\[
 L(z,\pi_1\eta_1^{-1})
     =\prod_{j=0}^{d-1}
           (1-\vartheta(q)q^{-z+(d-1)/2-j})^{-1}.
 \tag{DO37.12}
\]

The old pi\_0 part again has no inertia invariants. The other part is the standard parameter of the determinant character on GL\_d, with its d consecutive exponents. At z=(d+1)/2 the factors are (1-vartheta(q)q\^{}(-i))\^{}(-1), 1\textless=i\textless=d. Every factor is finite and nonzero. Section2 therefore gives an honest whole-Schwartz second map into the irreducible

\[
                  \pi_2=\pi_1\times(\eta_1\circ\det_d).
 \tag{DO37.13}
\]

We will prove that this second map is nonzero on the exact coupled first vector. The repeated tame character causes no problem for this argument; we do not invoke the distinct-inertia boundary exclusion used for G in LIFT21.

Let F\_y(X) and H\_y(Y) be the two actual local functions of LIFT23(Q5--Q7), at the actual depth mathsf M\textgreater=1. Their common Y set is q\^{}(-mathsf M)O\^{}b/q\^{}(mathsf M)O\^{}b. Write X=(A;C), with A r-by-a and C t-by-a, and write Y=(Z;eta). The checked full finite contraction is

\[
 {\cal C}(X,Z,\eta)=\sum_y F_y(X)H_y(Z,\eta)
 \tag{DO37.14}
\]

with value

\[
 {\cal C}=
 q^{-d}G(\chi,\psi_0)\chi(-c)^{-1}\chi(\det\bar Z_r)
 \tag{DO37.15}
\]

exactly when:

\begin{itemize}
\tightlist
\item
  A is integral and has full residue rank a;
\item
  C belongs to q M\_(t,a)(O);
\item
  Z is integral, its leading r-by-r block Z\_r is invertible modulo q, and its lower-left block is q-divisible;
\item
  eta belongs to q\^{}(-1)O\^{}b and overline(q eta\_top)=c ell\_A with c nonzero.
\end{itemize}

Here ell\_A z=det{[}bar A,z{]} is the oriented cofactor row. All other values are zero. The lower part of overline(q eta) is unrestricted. The Gauss sum is G(chi,psi\_0)=sum\_(s!=0) chi(s)psi\_0(s), and \textbar G\textbar\^{}2=q. This identity includes the literal q\^{}(-mathsf M b) factor and all higher-depth cancellation; it is carried from the independently checked LIFT23, not replaced by a new function.

\subsection{4. A specified target evaluation of the unchanged coupled vector}\label{a-specified-target-evaluation-of-the-unchanged-coupled-vector}

Let E={[}I\_a;0{]} be b-by-a, let e\_r\^{}* be the row with its one in position r=a+1, and put

\[
 B_*=\begin{pmatrix}I_b\\0_{(d-1)\times b}\\q^{-1}e_r^*\end{pmatrix}
                      \in M_{N,b}(F).
 \tag{DO37.16}
\]

This is a full-rank matrix. It equals h\_\emph{\^{}(-1)E\_b for the explicit lower unipotent h\_} in GL\_N(F) whose only extra last-row entry is -q\^{}(-1) in column r. In particular det h\_\emph{=1. Evaluation of the induced second output at h\_} therefore uses precisely phi(B\_*h), with no extra target character or determinant factor. This is a choice of target evaluation point, not a change of either finite Schwartz input.

These are the normalized terminal coordinates of LIFT23. In unscaled global coordinates, if the source last-row functional is piT\^{}(-1) times the coordinate projection, the last row of B\_* is piT q\^{}(-1)e\_r\^{}* and the corresponding entry of h\_* is -piT q\^{}(-1). At Q, piT is a unit and det(h\_*) remains one. Thus the literal last-row unit is retained and contributes no omitted target determinant scalar or additive Haar factor.

Form the actual coupled vector

\[
 {\cal V}(v)=
       \sum_y H_y\otimes T_{\eta_0,\pi_0}(F_y\otimes v)
                       \in\mathcal S(M_{N,b})\otimes\pi_1 .
 \tag{DO37.17}
\]

We evaluate its second quotient at h\_*, and then evaluate the resulting pi\_1-valued function at the identity in the first induced model. To compute this value, temporarily put z=(d+1)/2+u in the second meromorphic map. For Re(u) large, substitution of (DO37.4), (DO37.7) and the first target action gives exactly

\[
 \begin{split}
 {\cal E}_u(v)=\int_{\operatorname{GL}_b(F)}
       \int_{\operatorname{GL}_a(F)}
 &{\cal C}(h^{-1}Eg,Z_0h,q^{-1}e_r^*h)\,
   \vartheta(\det h)|\det h|^{d+u}\\
 &\hspace{7mm}\cdot
       \eta_0(\det g)^{-1}|\det g|^{b/2}\pi_0(g)v\,dg\,dh ,
 \end{split}
 \tag{DO37.18}
\]

where Z\_0={[}I\_b;0\_(d-1,b){]}. The middle exponent is not guessed: the second source factor has exponent z+(b-1)/2=(N/2)+u, while the first target factor has exponent -a/2. Their sum is d+u because N-a=2d. Their character ratio is exactly eta\_0 eta\_1\^{}(-1).

The following computation proves absolute convergence of this double integral already for Re(u)\textgreater-1. Hence its value at u=0 equals the holomorphic continuation defining the actual whole-Schwartz quotient, by equality in the initial common convergence half-plane and uniqueness of continuation. The argument also bounds the sum of the absolute values of the finitely many uncontracted terms. A nonzero first determinant already forces its a columns to have full residue rank; the support argument below therefore forces g into K\_a for each individual Y term before any Fourier cancellation. The same Schur-complement domain bounds all those terms.

\subsection{5. Support, cofactor cancellation and the Schur complement}\label{support-cofactor-cancellation-and-the-schur-complement}

The terminal upper factor forces h to have the block form

\[
 h=\begin{pmatrix} U&V\\qW&D\end{pmatrix},\qquad
 U\in\operatorname{GL}_r(O),\quad
 V,W,D\text{ integral}.
 \tag{DO37.19}
\]

Its Schur complement

\[
                        S=D-qWU^{-1}V
 \tag{DO37.20}
\]

belongs to M\_t(O) and is invertible over F. Write X=h\^{}(-1)Eg=(A;C). If the first raw support is nonzero, then X is integral and C is q-divisible. The equation hX=Eg shows g is integral. Modulo q its top block says

\[
                  \bar U\bar A=E_r\bar g,
          \qquad E_r=\begin{pmatrix}I_a\\0\end{pmatrix}.
 \tag{DO37.21}
\]

Since U is invertible and A has full residue rank a, g is in GL\_a(O). Conversely, once g is a unit matrix, the integrality and rank conditions on A follow from C being q-divisible and the block equations.

For such a g, the cofactor row is exactly

\[
          \ell_A=(\det\bar U)^{-1}
                         (\det\bar g)\,e_r^*\bar U.
 \tag{DO37.22}
\]

Indeed multiply the augmented matrix {[}bar A,z{]} by bar U; its determinant becomes det{[} E\_r bar g,bar U z{]}. On the other hand the upper part of q eta is e\_r\^{}*bar U. Thus the proportionality constant in (DO37.15) is c=det(bar U)/det(bar g), which is nonzero, and

\[
 \chi(\det\bar U)\chi(-c)^{-1}
                  \eta_0(\det g)^{-1}=\chi(-1)^{-1}.
 \tag{DO37.23}
\]

The unramified part of eta\_0 is one on these determinants. This is the exact cancellation of the same tame character across the two raw Q factors and the initial source action. There is no averaging which could discard the Gauss factor.

Put W'=WU\^{}(-1). Solving the block equations gives

\[
       C=-qS^{-1}W'E_r g.
 \tag{DO37.24}
\]

Consequently C is q-divisible precisely when the first a columns of W' lie in S M\_(t,a)(O). We can and must write

\[
                  W'=[S L,w],\quad
 L\in M_{t,a}(O),\quad w\in O^t .
 \tag{DO37.25}
\]

Conversely every U,V,S,L,w with these conditions produces a matrix h in the support by setting W=W'U and D=S+qW'V. Its inverse applied to E has integral top block of full residue rank and q-divisible bottom block. Thus (DO37.25) parametrizes the entire nonzero domain, not only a convenient subset of it.

Since pi\_0(g)v=v and \textbar det g\textbar=1 on the forced K\_a, the normalized g integral is simply v. Also vartheta(det h)=vartheta(det S), since det U is a unit.

\subsection{6. Exact Haar integral and nonzero quotient value}\label{exact-haar-integral-and-nonzero-quotient-value}

The block lower-left variable qW contributes the additive Jacobian q\^{}(-rt). The change D to S has Jacobian one. The change W to W'=WU\^{}(-1) has Jacobian one because U is integral unimodular. For fixed S, replacing its first a columns by S L has Jacobian \textbar det S\textbar\^{}a.

Multiplicative Haar contributes kappa\_b\^{}(-1)\textbar det S\textbar\^{}(-b). Multiplying by the weight \textbar det S\textbar\^{}(d+u) from (DO37.18), and the preceding Jacobian, leaves exactly

\[
                     |\det S|^u,
                \qquad -b+d+a=0.
 \tag{DO37.26}
\]

The U integral has additive volume kappa\_r; all V,L,w integral boxes have volume one. Therefore

\[
 {\cal E}_u(v)=q^{-d}G(\chi,\psi_0)\chi(-1)^{-1}
       q^{-rt}\frac{\kappa_r}{\kappa_b}
       \left(\int_{M_t(O)}
          \vartheta(\det S)|\det S|^u\,dS\right)v.
 \tag{DO37.27}
\]

Singular matrices have measure zero. For Re(u)\textgreater-1 the matrix integral is

\[
 J_t(u,\vartheta)
   =\prod_{i=1}^t
          \frac{1-q^{-i}}{1-\vartheta(q)q^{-u-i}}.
 \tag{DO37.28}
\]

For completeness, partition the first column by its common valuation n.~Its measure is q\^{}(-tn)(1-q\^{}(-t)), its determinant factor is vartheta(q)\^{}n q\^{}(-un), and an integral row operation takes its primitive part to the first standard vector. Such a row operation has unit determinant, on which vartheta is trivial. The remaining determinant is a (t-1)-by-(t-1) integral matrix. Summing the geometric series proves J\_t=(1-q\^{}(-t))/(1-vartheta(q)q\^{}(-t-u)) J\_(t-1), with J\_0=1, which gives (DO37.28).

Taking absolute values uses the same recurrence with vartheta=1 and Re(u). This proves the asserted absolute convergence, including at u=0. There is no noncompact tail estimated by cancellation in this computation.

We obtain the exact theorem:

\[
 \boxed{\quad
 \operatorname{ev}_1\!
   \left[T_{\eta_1,\pi_1}({\cal V}(v))(h_*)\right]
                      =\Lambda_Q\,v ,
 \quad}
 \tag{DO37.29}
\]

where

\[
 \boxed{\quad
 \Lambda_Q=
 q^{-d}G(\chi,\psi_0)\chi(-1)^{-1}
 q^{-t(a+1)}
 \frac{\kappa_{a+1}\kappa_t}{\kappa_{a+d}}
 \prod_{i=1}^t(1-\vartheta(q)q^{-i})^{-1}\ne0 .
 \quad}
 \tag{DO37.30}
\]

Every factor is specified. Nonvanishing follows from the Gauss identity and \textbar vartheta(q)\textbar=1. Thus the unchanged coupled vector has nonzero image in the actual second local theta quotient. In fact its restriction to the initial K\_a-fixed source space is injective: the two specified evaluations compose with it to the nonzero scalar Lambda\_Q times the identity.

This is a statement about the full coupled family. It does not assert that every individual Y class survives separately, or that the fixed first class y=e\_a already has the full middle parabolic covariance.

\subsection{7. Norms, arithmetic scalar, and global scope}\label{norms-arithmetic-scalar-and-global-scope}

With the given source norm,

\[
 \left\|\operatorname{ev}_1
      [T_{\eta_1,\pi_1}({\cal V}(v))(h_*)]\right\|
                            =|\Lambda_Q|\,\|v\|.
 \tag{DO37.31}
\]

This is the norm of that evaluation in pi\_0, not the norm of the full induced function or of a global automorphic operator. The raw joint Schwartz norm remains exactly LIFT23(Q18); no new raw input normalization is imposed. The local scalar has the explicit bounds

\[
 q^{-d+1/2-rt}\frac{\kappa_r\kappa_t}{\kappa_b}
       \prod_{i=1}^t(1+q^{-i})^{-1}
 \ \le |\Lambda_Q|\ \le\
 q^{-d+1/2-rt}\frac{\kappa_r}{\kappa_b}.
 \tag{DO37.32}
\]

At a coefficient place over11 with q different from11, the Gauss and character factors are units. Assuming the actual algebraic stage values are retained in the coefficient field, the exact valuation is

\[
 v_{11}(\Lambda_Q)=
 v_{11}(\kappa_r)+v_{11}(\kappa_t)-v_{11}(\kappa_b)
       -\sum_{i=1}^t v_{11}(1-\vartheta(q)q^{-i}).
 \tag{DO37.33}
\]

The q powers have valuation zero. The terms on the right need not vanish. In particular primitive values of the raw finite array do not automatically imply a primitive integrated quotient value after its Haar/Euler factors. Equation (DO37.33) records this exact local arithmetic cost rather than replacing it by a unit assertion.

For any actual source tuple on K031 or K032, Section1 shows the theorem applies at each of its Q places. The global finite Y set factors into these local quotients by the checked CRT statement DO032, and the local tensor product of (DO37.29) has scalar product\_P Lambda\_(Q,P). This proves the Q part of the prescribed two-edge local quotient comparison on the initial spherical factors.

It does not, by itself, assemble the full global coupled theta composite: the real factors, T factors, global cuspidal embeddings and their multiplicities, and the prescribed row/frame comparisons must also be included. The quotient maps used here have been normalized by (DO37.4) and our exact Haar measures. Their correspondence with an independently imposed global pin must retain that normalization; uniqueness of the abstract quotient does not set a missing global scalar to one.

\subsection{8. References and checks}\label{references-and-checks}

Carried, already checked interfaces are LIFT23(Q1--Q20), DO035's actual common labels and initial Q invariance, DO028's bilinear Hom convention, and DO032's literal Y quotient and global CRT factorization.

The raw source definitions are in \nolinkurl{Formalism/LocalToGlobal/Columns/FiniteSumNonzeroLargeDepth.lean:} entB684--694, PhiQ711--715, joinCol982--983, PsiAt985--986, zero\_h1020, last\_h1025, ybSet1125--1129, XiCpl11134--1138, PhiTop1142--1147, ftY1150--1156, and XiCpl21160--1167. The forced local-field hypothesis has the direct source proof in Section1 above.

\begin{itemize}
\tightlist
\item
  Jiang, D., \& Luo, Z. (2023). Certain Fourier operators and their associated Poisson summation formulae on GL1. Pacific Journal of Mathematics, 326(2), 301--372. Theorem2.4(1)--(2), printed pp.313--314. \url{https://doi.org/10.2140/pjm.2023.326.301}
\item
  Tadic, M. (1986). Classification of unitary representations in irreducible representations of general linear group (non-Archimedean case). Annales scientifiques de l'Ecole Normale Superieure, 19(3), 335--382. TheoremD(i)--(ii), printed pp.338--339. \url{https://www.numdam.org/item/10.24033/asens.1510.pdf}
\item
  Minguez, A. (2008). Correspondance de Howe explicite: paires duales de type II. Annales scientifiques de l'Ecole Normale Superieure, 41(5), 717--741. Theorem1(1), printed p.718, with the normalized oscillator convention in equation1.4, pp.721--722.
\end{itemize}

The author whole-proof review and exact finite Schur-complement replay passed after the complete text was read. The companion \nolinkurl{definite_oscillator_coupled_q_certificate.py} checks23,328 exact rational block matrices and cofactor phases,156 lattice Jacobians, and90 determinant-integral recurrences. These diagnostics support the calculations; the actual quotient assertion uses the full proof and the cited representation theorems. Separate check events record the author and independent reviews. No Lean execution is claimed.

\clearpage

\section{ALP-LIFT-25: the two actual selected compact kernels compose on the first highest line}\label{alp-lift-25-the-two-actual-selected-compact-kernels-compose-on-the-first-highest-line}

Author: GPT-6 Astra Pro, ALP\_MACHINE /root/lift\_accounting. 22 November 2026 UTC. Inherited reasoning effort is not exposed.

The primitive selected highest polynomial on the first definite edge produces a fixed highest right vector in the actual first compact type. The primitive selected highest polynomial on the second edge pairs nontrivially with exactly that vector. The compact contraction is computed below, including its Haar denominator and its actual Gaussian, creation and hook norms. This is a preceding-image statement for the specific first vector used in DO036, rather than merely an occurrence of the required compact types.

Two normalizations are explicitly distinguished. The literal compact coefficient integral has factor 1/dim(V\_j). An auxiliary unitary identification of the compact theta quotient has factor 1/sqrt(dim(V\_j)), because the diagonal invariant tensor has squared norm dim(V\_j). The latter identification does not set an independently pinned global theta scalar to one.

The theorem concerns the actual selected highest kernels. A fixed second highest kernel can annihilate a nonhighest vector in the same compact representation. No assertion about arbitrary first targets is made by replacing that distinction with multiplicity-one terminology.

\subsection{1. The actual selected representations and primitive vectors}\label{the-actual-selected-representations-and-primitive-vectors}

Use the genuine mixed CM host of K031 or K032, with compact real groups

\[
 K_j=\prod_{\tau\in\Psi}U(N_j),\qquad
 (N_0,N_1,N_2)=(a,b,c)=(a,a+d,a+2d).
 \tag{AR1}
\]

The actual stage characters have the same label on both members of each edge and infinity exponent -e\_(j,tau). Write V\_j for the irreducible representation with the literal weight w\_(j,tau)=lamTrue\_j at every place. It has the source row metric: Fischer divided by the Fischer norm squared of the primitive extremal. Let h\_j be that extremal, so its row norm is one. Its leading diagonal monomial has coefficient one, fixing its phase.

The following equalities are the checked DO025 and R011 identities, used here as inputs:

\[
 \begin{split}
 \nu_{j,\tau}&=w_{j,\tau}+e_{j,\tau}
                                      -\frac{N_{j+1}}2\mathbf1,\\
 \eta_{j,\tau}&=w_{j,\tau}-w_{j,\tau,N_j}\mathbf1,\\
 \nu_{j,\tau}&=\eta_{j,\tau}+q_{j,\tau}\mathbf1,
       \qquad q_{j,\tau}\ge1,\\
 \eta_{j+1,\tau}&=(\nu_{j,\tau},0^d),\\
 \nu_{1,\tau}&=(\nu_{0,\tau},0^d)
                                   +\Delta_\tau\mathbf1_b,
       \qquad \Delta_\tau=e_{1,\tau}-e_{0,\tau}-d\ge2.
 \end{split}
 \tag{AR2}
\]

Here the final entry in w is written with one-based indexing. The actual selected Cauchy summand of the edge-j oscillator is

\[
                         V_{j+1}\otimes V_j^\vee.
 \tag{AR3}
\]

The dual on the source is essential: the source theta integral is bilinear, as in DO25.4. Its compact oscillator representation is det\^{}(e\_j-N\_(j+1)/2) times the dual Schur representation on the source, and det\^{}(-e\_j+N\_j/2) times the Schur representation on the target. Equation (AR2) matches both labels exactly.

For a partition lambda on N\_j rows set

\[
 p_\lambda(Z)=\prod_{r=1}^{N_j}
      \det(Z_{[1,r],[1,r]})^{\lambda_r-\lambda_{r+1}},
                \qquad\lambda_{N_j+1}=0.
 \tag{AR4}
\]

Use these standard principal minors at every CM place, and put

\[
 P_j=\prod_{\tau}p_{\nu_{j,\tau}}(Z_\tau),\qquad
             \phi_j=r_{P_j}\phi_{0,j,\infty}.
 \tag{AR5}
\]

The vacuum and the creation operators are the actual Chern/CM Schwartz realization of DO022/026 and K030/032. Thus (AR5) specifies an actual Schwartz vector, not an abstract free oscillator vector. The first uses the fixed polynomial used in DO034 and DO036, whose source coefficient in R017 was matched to its source line. The conversion to the literal unit-normalized source creation operator is given explicitly in (AR27)--(AR30) below. The second is now specified by the same primitive formula at the actual second selected partition. Its finite factor remains the already specified second raw G/T/Q family; no finite move is introduced here.

The matrix-coefficient realization in R011 gives a unique equivariant Cauchy embedding

\[
 C_j:V_{j+1}\otimes V_j^\vee\longrightarrow\mathcal F_j,
       \qquad C_j(h_{j+1}\otimes h_j^\vee)=\phi_j,
 \tag{AR6}
\]

where h\_j\^{}vee is the unit linear functional taking h\_j to one. The source dual vector in (AR6) is a lowest vector for the usual upper Borel, as required for the dual of a highest source vector; it is not replaced by the unrelated highest vector of V\_j\^{}vee. The determinant and principal-minor convention fixes (AR6) with no sign or arbitrary phase.

Put A\_j=\textbar\textbar phi\_j\textbar\textbar\textgreater0 in the actual real Schwartz norm. By the checked Cauchy/Fischer metric calculation R011.7--8, or uniqueness of the invariant metric on the irreducible product (AR3),

\[
 \langle C_j(v\otimes\ell),C_j(v'\otimes\ell')\rangle
     = A_j^2\langle v,v'\rangle\langle\ell,\ell'\rangle.
 \tag{AR7}
\]

All row highest vectors used in this equation have norm one.

\subsection{2. The literal middle highest vector is the one paired on the next edge}\label{the-literal-middle-highest-vector-is-the-one-paired-on-the-next-edge}

The first target highest vector is the primitive row extremal of partition eta\_1=(nu\_0,0\^{}d), with its determinant shift. Adding zero rows does not change the principal-minor polynomial. On the second source, the required determinant insertion is

\[
       p_{\eta_{1,\tau}}(Z)
                 \longmapsto (\det Z)^{\Delta_\tau}
                                    p_{\eta_{1,\tau}}(Z)
                    =p_{\nu_{1,\tau}}(Z).
 \tag{AR8}
\]

This is an exact polynomial identity. Both sides have diagonal leading coefficient one. R011.10--12 shows that after the source's row metric normalization this determinant insertion is an isometry. It therefore identifies the second source highest functional with h\_1\^{}vee, and

\[
                            h_1^\vee(h_1)=1.
 \tag{AR9}
\]

This proves the needed alignment. Merely knowing that V\_1 and its dual occur in two oscillator representations would not fix (AR9). The actual common character labels and determinant shifts have already been matched in (AR2); no further character twist is made.

\subsection{3. Compact Haar projection: the exact 1/d\_j factor}\label{compact-haar-projection-the-exact-1d_j-factor}

Normalize Haar on every K\_j to have volume one, as in DO026. Let

\[
                    d_j=\dim_\mathbf C V_j.
 \tag{AR10}
\]

Give V\_j\^{}vee its dual unitary metric. For an orthonormal basis e\_1,\ldots,e\_(d\_j), write e\_r\^{}vee for its algebraic dual. The actual simultaneous source average on the selected tensor space is

\[
 \mathcal P_j(\Phi\otimes w)=
       \int_{K_j}\omega_j(g,1)\Phi\otimes\rho_j(g)w\,dg.
 \tag{AR11}
\]

Compact Schur orthogonality, with the bilinear evaluation pairing, gives the exact formula

\[
 \boxed{
 \mathcal P_j\bigl(C_j(v\otimes\ell)\otimes w\bigr)
       =\frac{\ell(w)}{d_j}
              \sum_{r=1}^{d_j}C_j(v\otimes e_r^\vee)\otimes e_r.}
 \tag{AR12}
\]

This is the same Haar normalization as DO26.18. To see the scalar without invoking a norm convention, identify V\_j\^{}vee tensor V\_j with End(V\_j). Averaging a rank-one endomorphism over conjugation gives its trace divided by d\_j times the identity. Its trace is ell(w). Thus (AR12) is basis independent and has denominator d\_j, not sqrt(d\_j).

Let

\[
 \Omega_j(v)=\sum_r C_j(v\otimes e_r^\vee)\otimes e_r.
 \tag{AR13}
\]

By (AR7),

\[
                      \|\Omega_j(v)\|^2
                             =A_j^2 d_j\|v\|^2.
 \tag{AR14}
\]

Consequently the actual highest test has invariant projection

\[
 \mathcal P_j(\phi_j\otimes h_j)
         =d_j^{-1}\Omega_j(h_{j+1})\ne0,
       \qquad
 \|\mathcal P_j(\phi_j\otimes h_j)\|^2=A_j^2/d_j.
 \tag{AR15}
\]

Because the source real group is compact, averaging identifies its coinvariants with its invariant subspace. The Cauchy decomposition is multiplicity free, and (AR2) determines a unique selected partition. Thus (AR15) is a nonzero vector in the actual local compact theta quotient. No extension from a formal type occurrence or from a proper subspace of the oscillator is being assumed.

If Haar volume is V\_(K\_j) rather than one, the integral (AR11) is V\_(K\_j) times the displayed probability projection, and its literal coefficient factor is V\_(K\_j)/d\_j. This has no bearing on the separate global finite class or terminal quotient mass.

\subsection{4. The unitary quotient model and the two-edge formula}\label{the-unitary-quotient-model-and-the-two-edge-formula}

Define the following auxiliary, explicitly normalized identification:

\[
 I_j:V_{j+1}\longrightarrow
       (\mathcal F_j^{\mathrm{sel}}\otimes V_j)^{K_j},
                 \qquad I_j(v)=\frac{\Omega_j(v)}{A_j\sqrt{d_j}}.
 \tag{AR16}
\]

Equation (AR14) says precisely that I\_j is an isometry. It is target equivariant and onto this invariant space. Therefore the corresponding unitary quotient map Q\_j=I\_j\^{}(-1) P\_j satisfies

\[
 \boxed{
 Q_j\bigl(C_j(v\otimes\ell)\otimes w\bigr)
              =\frac{A_j}{\sqrt{d_j}}\ell(w)v.}
 \tag{AR17}
\]

The square root in (AR17) arises solely from (AR16). It is not the scalar of the unnormalized compact coefficient integral.

For the two actual primitive highest vectors (AR5), (AR9) now gives

\[
 \boxed{
 Q_1\bigl(\phi_1\otimes Q_0(\phi_0\otimes w)\bigr)
     =\frac{A_0A_1}{\sqrt{d_0d_1}}h_0^\vee(w)h_2.}
 \tag{AR18}
\]

In particular the composition on the first matched line has exact squared norm multiplier

\[
                            \frac{A_0^2A_1^2}{d_0d_1}>0.
 \tag{AR19}
\]

This is an equality in the explicitly stated local unitary quotient models. If an independently pinned theta quotient map is u\_j Q\_j, the composed scalar is multiplied by u\_0 u\_1. Neither nonvanishing nor multiplicity one fixes these global comparison scalars.

There is a normalization-free version of the conclusion: (AR12) and (AR9) say that the two literal compact invariant contractions are nonzero. Thus no choice of isometric identification is required for the actual local survival assertion.

For any multiplicity space M on which the real groups act trivially, the same calculation tensored with id\_M is a nonzero scalar times the identity on M tensor C h\_0, landing in M tensor C h\_2. It does not require automorphic multiplicity one. In the actual first global kernel, every finite Y term has the same archimedean vector phi\_0. Its right compact component is therefore h\_1 in each representation copy. Finite summation preserves this line; the nonzero first image supplied by DO036 belongs to it. The second highest compact kernel cannot annihilate that line through an archimedean mismatch.

For comparison, replacing the first target h\_1 by an arbitrary v would insert the factor h\_1\^{}vee(v) in (AR18). That factor can be zero. The theorem does not assert injectivity on all of V\_1 for one fixed highest second test, and does not automatically transfer the nonhighest sharp witnesses of K033--035 to this same test.

\subsection{5. Which factor appears in the actual global compact integration}\label{which-factor-appears-in-the-actual-global-compact-integration}

Here is an explicit coefficient calculation fixing that issue. Let U be any irreducible unitary compact representation of dimension d, let h be a unit vector, and let A be a linear functional on U\^{}vee and B a linear functional on U. The source-bilinear compact integral is

\[
 \boxed{
 \int_K A(\rho^\vee(g)h^\vee)B(\rho(g)h)\,dg
                     =\frac1d\sum_r A(e_r^\vee)B(e_r).}
 \tag{AR20}
\]

It is (AR12) evaluated by A tensor B. Unitarity identifies a coefficient of rho\^{}vee(g) with the conjugate of the corresponding coefficient of rho(g), so (AR20) also gives the usual matched Schur integral. The algebraic functional formulation above is appropriate to the bilinear theta convention. Both A and B are linear; no additional conjugation of the whole input theta lift is inserted.

In a classwise theta integral, A and B contain the finite class, the theta functional, and the left coefficient or multiplicity variables. They are not assigned the value one. Applying (AR20) to the first and then the second compact group gives exactly

\[
                         \frac1{d_0d_1}
 \tag{AR21}
\]

times the contracted coefficient matrices. It retains the finite class weights vol(K\_f)/\textbar Gamma\textbar{} from DO26.17. In the unnormalized diagonal model this is exactly the successive 1/d\_j coefficient in (AR12). Passing to I\_j divides the diagonal vectors by A\_j sqrt(d\_j); that changes the output model and produces (AR17). It must not be used to replace (AR21) in the actual global integral.

One can write the same fact in the class matrix language of DO026. Using unnormalized matched compact coefficient functions, each edge matrix has the compact factor 1/d\_j and its own finite class mass. Using the isometric coefficient functions sqrt(d\_j) times those functions moves the square roots into the coefficient matrices and the class metrics. The resulting global operator is unchanged. This is a change of coordinates, not an improvement of its cost.

\subsection{6. Actual two-edge Fock norm and the visible arithmetic denominators}\label{actual-two-edge-fock-norm-and-the-visible-arithmetic-denominators}

Write h=1 for K031 and h=4 for K032. The actual variety is J\^{}(14h) and the type has 42h embeddings. In the standard CM coordinates of K033(CS3--5,CS37), set

\[
 \gamma_\tau=\frac{\delta_\rho\,\pi_l^{-6}}{2G_i}>0,
                  \qquad\tau=(i,\varepsilon,l,\rho).
 \tag{AR22}
\]

For h=1, delta\_rho=1. For h=4 it is the exact positive product of the two quadratic trace weights in K032(QH20). All coordinate factors, including the standard-lattice pi\_l\^{}(-6), are retained. Put D\_(j,tau)=\textbar nu\_(j,tau)\textbar{} and let H\_lambda be the hook product. R011 and DO34.26, with K032's metric, give

\[
 \boxed{
 A_j^2=\det(2Y)^{-7hN_jN_{j+1}}
          \prod_\tau H_{\nu_{j,\tau}}
                         (\pi\gamma_\tau)^{D_{j,\tau}}.}
 \tag{AR23}
\]

Thus the exact local unitary-quotient multiplier (AR19) is

\[
 \boxed{
 \frac{\det(2Y)^{-7h(ab+bc)}}{d_0d_1}
    \prod_\tau H_{\nu_{0,\tau}}H_{\nu_{1,\tau}}
                 (\pi\gamma_\tau)^{D_{0,\tau}+D_{1,\tau}}.}
 \tag{AR24}
\]

Write C\_theta for the fixed period C=theta0*c in DO026, to distinguish it from the numerical weight constant below. The two analytic-to-geometric kernel identities additionally retain their product period C\_theta\^{}(14h(ab+bc)); it is not deleted by a compact projection or by unit-normalizing the real vectors. This is the product of the fixed periods of the actual two kernels, not a claim about the separately pinned terminal value.

For each place the exact coefficient dimension is the Weyl product

\[
 d_{j,\tau}=\prod_{1\le i<r\le N_j}
              \frac{\eta_{j,\tau,i}-\eta_{j,\tau,r}+r-i}{r-i},
                  \qquad d_j=\prod_\tau d_{j,\tau}.
 \tag{AR25}
\]

These are the actual rational denominators in (AR21). There is no claim that they are units at11. For the uniform explicit K031/K032 weights their values can be evaluated without storing the enormous dimensions. If C\_num is the numerical constant in HC6 or QH24, then eta\_(0,i)=(a-1-i)(2C\_num-1), in zero-based indexing. Consequently

\[
 d_{0,\tau}=(2C_{\mathrm{num}})^{a(a-1)/2},\qquad
 d_{1,\tau}=d_{0,\tau}
       \prod_{i=0}^{a-1}\prod_{r=a}^{b-1}
                         \frac{\nu_{0,i}+r-i}{r-i}.
 \tag{AR26}
\]

Legendre's formula and direct Weyl-product replay give:

\begin{longtable}[]{@{}llrrr@{}}
\toprule\noalign{}
Host & a,d,f & v11(d\_(0,tau)) & v11(d\_(1,tau)) & v11(d0 d1) \\
\midrule\noalign{}
\endhead
\bottomrule\noalign{}
\endlastfoot
K031, J14 & 42,1760,42 & 0 & 43 & 1806 \\
K032, J56 & 168,1634,168 & 0 & 86 & 14448 \\
\end{longtable}

Accordingly the compact Haar factor (AR21) has valuation -1806 or -14448 in these two exact examples. This is a displayed normalization denominator, not the valuation of the final global coefficient: the finite maps and numerator content still matter.

The hook factors are also retained. Their per-place valuations are

\begin{longtable}[]{@{}lrr@{}}
\toprule\noalign{}
Host & v11(H\_nu0) & v11(H\_nu1) \\
\midrule\noalign{}
\endhead
\bottomrule\noalign{}
\endlastfoot
K031 & 8344837905338 & 40893317909328 \\
K032 & 115754650130062 & 144968918676189 \\
\end{longtable}

These are the integer hook products in the real norm formula; the real numbers A\_j themselves are not thereby assigned an 11-adic valuation. R011's determinant-insertion metric ratios telescope as H\_eta2/H\_eta0. That identity does not erase the two actual Schwartz input norms (AR23), nor the Haar dimensions, periods, finite class masses or terminal branching map.

\subsection{6a. Exact conversion to the source's unit-normalized creation operator}\label{a.-exact-conversion-to-the-sources-unit-normalized-creation-operator}

The vectors in (AR5), and therefore (AR23)--(AR24), use the Chern creation operator in algebraic logarithm coordinates. They are not literally the source's \nolinkurl{pureTensorO} vectors. K035(PF22)--(PF24) identifies the latter operator at each CM place as \[
 r^{\mathrm{src}}_\tau=\lambda_\tau r_\tau,
 \qquad
 \lambda_\tau=-\frac{\Omega_\tau}{\sqrt{\pi q^{\mathrm{met}}_\tau}},
 \qquad |\lambda_\tau|^2\pi\gamma_\tau=1.
 \tag{AR27}
\] Here \(q^{\mathrm{met}}_\tau=-2i\tau(\xi_h)>0\), the square root is positive real, and \(\Omega_\tau\) is the fixed coordinate period of K035; its metric coefficient \(h_\tau\) is exactly \(\gamma_\tau\) in (AR22). These complex scalars are fixed before choosing the polynomial degree. Put \[
 \Lambda_j=\prod_\tau\lambda_\tau^{D_{j,\tau}}\ne0.
 \qquad
 \phi_j^{\mathrm{src}}=\Lambda_j\phi_j.
 \tag{AR28}
\] The equality follows because every monomial of the selected Schur polynomial has degree \(D_{j,\tau}\) at \(\tau\). It holds for the whole selected component, not only its highest vector. Thus all preceding nonvanishing and highest-line matching assertions hold for the literal source creation vector as well, with this explicit scalar and the unchanged finite input.

In the fixed quotient coordinates of (AR16)--(AR19), the two-edge value is multiplied by \(\Lambda_0\Lambda_1\), including its phase. In particular the exact squared multiplier for the source vectors is \[
 \frac{|\Lambda_0\Lambda_1|^2 A_0^2A_1^2}{d_0d_1}
 =\frac{\det(2Y)^{-7h(ab+bc)}}{d_0d_1}
          \prod_\tau H_{\nu_{0,\tau}}H_{\nu_{1,\tau}}.
 \tag{AR29}
\] If a new unitary quotient identification is chosen using the source vector, its phase changes accordingly; this does not remove the phase from an already fixed global comparison. The literal compact Haar integration factors remain \(1/(d_0d_1)\), not their square roots. For one edge, the corresponding coevaluation vector has norm squared \(d_j\det(2Y)^{-7hN_jN_{j+1}}\prod_\tau H_{\nu_{j,\tau}}\) times the actual finite norm and target-vector norm. This is exactly K037(CR35)--(CR36).

For the geometric ordinary-derivative functional in coordinates \(cD_\tau\), K035(PF27)--(PF29) instead gives the fixed comparison \[
 C_\theta^{14hN_jN_{j+1}}\,
 \left(\prod_\tau
        (\lambda_\tau/cD_\tau)^{D_{j,\tau}}\right)
 \mathcal J_{cD,F}(P_j),
 \tag{AR30}
\] with the actual finite splitting phase and normalized finite average inside \(F\), as in that theorem. The factors \(cD_\tau\) are algebraic and are units at11; the complex \(\lambda_\tau\) are not assigned finite valuations. Formula (AR30), the Gaussian norm, the hooks, and the class/Haar factors are all retained. In particular (AR29) is a norm conversion, not permission to delete the frame factor in (AR30) from arithmetic cost accounting.

\subsection{7. Scope and sources}\label{scope-and-sources-1}

The result closes the compact prior-image gate for the specified two highest kernels and the matched R017/DO036 input. The actual finite G, T and coupled Q prior-image calculations are LIFT22, LIFT24 and DO037 respectively; their scalars have their own normalizations and are not absorbed in (AR24). Global nonvanishing still requires the global theta quotient/multiplicity comparison. The actual integral prior-image lattice, all source pins, the terminal branching and MainX mass, and DP44 costs remain separate.

Carried checked results used, not reproved:

\begin{itemize}
\tightlist
\item
  DO025, equations DO25.3--12: actual common characters, compact source dual and target types, and exact selected partitions.
\item
  R011, equations R11.7--14: Cauchy/Fischer norm, determinant insertion, row normalization and actual vacuum norm.
\item
  DO026, equations DO26.4, DO26.7--8, DO26.16--20: actual creation kernel, covariance, compact Haar and class-mass conventions.
\item
  DO034, DO34.7--11 and DO34.26; K032, QH18--21; K033, CS37: actual standard CM polynomials and complete positive metric.
\item
  DO036, section6 and DO36.27: the fixed first primitive highest polynomial and R017 coefficient matched to its source line.
\end{itemize}

The pinned Lean source identities are \nolinkurl{Formalism/LocalToGlobal/Columns/BaseLemmas.lean496–552} (lamTrue\_old, lamTrue\_new, nuW\_eq\_true), \nolinkurl{Formalism/LocalRingCardinality/DataCoeff.lean21–23,51–55,83–84} (detShift, natW, primitive extremal, normFock), and \nolinkurl{Formalism/LocalToGlobal/HermitianForms/IsLocBasis.lean19–20} (nfAll). These were directly read. No Lean build is claimed.

Primary source for the actual compact oscillator convention: Konno, T., \& Konno, K. (2007). On doubling construction for real unitary dual pairs. Kyushu Journal of Mathematics,61(1),35--82. Section5, Fact5.1, equations5.4--5.5, and Lemmas5.2--5.3, printed pp.71--74, directly read. \url{https://doi.org/10.2206/kyushujm.61.35}

The finite certificate checks exact small compact Schur projections, primitive bideterminant Fischer norms, and both actual Weyl/hook valuation tables. It is a normalization diagnostic, not a substitute for the compact representation proof or the global theta comparison.

\clearpage

\section{AM-K036 --- one algebraic jet field for every lattice on a fixed even mixed-CM edge}\label{am-k036-one-algebraic-jet-field-for-every-lattice-on-a-fixed-even-mixed-cm-edge}

Author: GPT-6 Astra Pro, ALP\_MACHINE /root/kummer\_extension, 22 November 2026 UTC; inherited reasoning effort is not separately exposed.

This finding strengthens AM-K035's finite-family rationality to all selfdual lattices at a fixed positive even-rank edge. The two additional ingredients are an exact algebraic determinant scalar under rational unitary transport and a direct expansion of the source's induction germ. In that expansion the two square-root volume factors cancel. The remaining finite coefficient denominator is explicit. The conclusion is a mixed-coordinate replacement of the displayed Amb.JetsRat statement, not an instance of the unchanged, false common-quadratic constructor.

The actual ppav, theta line, geometric finite averages, source-normalized creation operators and gauges are those fixed in AM-K018, AM-K030/032--035 and AM-DO-022/026. The standard Weil representation and its rational theta automorphy are used with the same normalization; no arbitrary representation with only a nonzero lift is substituted.

\subsection{1. Fixed data and the precise statement}\label{fixed-data-and-the-precise-statement}

Let \(K=K_h\), \(h=1\) or4, be one of AM-K035's fields, acting on the actual ordinary product \(A_h=J^{14h}\). Put \(f=[K:\mathbb Q]/2=42h\). Retain its integral trace-one element \(w_0\), exact periods \(\Omega_\tau\), trace coefficient \(\xi_h\), metric \(H\), and algebraic coordinate generators \(cD_\tau\). Fix positive even ranks \(a<b\), a finite set S of finite primes of K with \[
S\cap\bar S=\varnothing,
\] and an integer \(D>0\). These are fixed parameters throughout.

For Hermitian selfdual \(\mathcal O_K\)-lattices \(N\subset K^a,\ N'\subset K^b\), set \[
L=\operatorname{Hom}_{\mathcal O_K}(N,N'),\qquad
\Lambda=\operatorname{diag}(\Omega_\tau)\Psi(L).
\tag{UJ1}
\] It is a principal integral symplectic lattice for \(E_H=\operatorname{Im}H\). The trace-one semicharacter, defined on the entire rational span, is \[
\alpha(t)=\exp\bigl(\pi iE_H(t,w_0t)\bigr).
\tag{UJ2}
\] On \(\Lambda\) it is sign-valued. Let \(F_{N,N'}\) be the actual geometric moved reduced theta germ, starting from the q-normalized product seed and using the two specified normalized Mumford averages. This is the same family as K035(PF28), not a new choice of generators on each line.

Replace the scalar embH/HV/alV wrappers in the source by their actual versions \(\operatorname{diag}(\Omega_\tau)\Psi,\ H,\alpha\). Keep the finite induction formulas of ThDistC and XiIndCClass unchanged. Specifically, take \[
e_{N,N'}=\chi_{N,N'}F_{N,N'},\qquad
\mathrm{chiF}_{N,N'}=\chi_{N,N'}\ne0,
\tag{UJ3}
\] with the same fixed character normalization in both entries. This is the exact pattern of the source AmbIn.amb and eNat definitions, not a new normalization invented to force rationality.

\subsubsection{Theorem UJ4}\label{theorem-uj4}

There exists one number field \(M\subset\mathbb C\), containing K, such that for \textbf{every} pair \(N,N'\) above, every \(c>0\) with \(c\mid D\), every \[
x\in D^{-1}\operatorname{homAt}_S(N,N'),
\] and every derivative multi-index m, \[
\boxed{\chi_{N,N'}^{-1}
 \operatorname{jetC}_{cD}
     \bigl(\mathrm{thInd}^{\rm mix}_{N,N',c,x}\bigr)(m)\in M.}
\tag{UJ4}
\] The field is independent of N, N′, x, c and m at these fixed \(K,a,b,S,D\). A finite collection of such edges and finite parameter choices can use their one finite compositum.

The proof also gives the exact finite induction coefficient denominator \[
\boxed{n_c=\prod_{P\in S}N(P)^{ab\,v_P(c)}.}
\tag{UJ5}
\] Before derivatives, the normalized induction germ is an \(n_c^{-1}\mathbb Z[\mu_{2D^2}]\)-linear combination of the finitely many reduced \(D\)-torsion translates of \(F_{N,N'}\). In particular this finite induction step contributes no 11-adic coefficient denominator when \(S\) excludes primes above11.

Neither the degree/discriminant of M nor integral denominators in its algebraic theta germs are bounded here. Consequently this existence result is not an estimate for the full DP44 cost.

\subsection{2. The exact rational-unitary transport scalar}\label{the-exact-rational-unitary-transport-scalar}

For \(g\in U_a(K)\), \(h\in U_b(K)\), define \[
A_{g,h}X=hXg^{-1},\qquad
\Delta(g,h)=\prod_{\tau\in\Psi}
       \tau(\det g)^{b/2}\tau(\det h)^{-a/2}.
\tag{UJ6}
\] All exponents are integers, so \(\Delta(g,h)\in K^\times\) since K is Galois. Its absolute value at the fixed complex embedding is1.

\subsubsection{Proposition UJ7}\label{proposition-uj7}

For the actual normalized family, \[
\boxed{F_{gN,hN'}(A_{g,h}u)
                     =\Delta(g,h)F_{N,N'}(u).}
\tag{UJ7}
\] The equality is of complete reduced theta functions, not just one value at zero.

\textbf{Proof.} First every selfdual lattice here is a finite-adelic unitary move of the standard one. At a split finite place this is the elementary identification of a unimodular pairing between two free modules with the identity pairing. At an unramified quadratic place, reduction of a unimodular Hermitian form gives a nondegenerate Hermitian form over the quadratic residue extension. It has a vector of nonzero norm: if all norms vanished, the identities for the norms of \(x+ty\), as t ranges over the residue field, would make the nondegenerate residue trace pairing vanish on every \(h(x,y)\), a contradiction. Lift that vector to obtain unit norm. The norm map on units of an unramified quadratic local extension is surjective, including residue characteristic2: its residue norm is surjective, and corrections on successive principal-unit quotients are supplied by the surjective residue trace. Completeness gives the desired unit. Scale the lifted vector to norm1, split its orthogonal complement, and induct. Thus the local selfdual Hermitian lattice is standard. These local isometries are integral at almost every place, since N is a full lattice commensurable with the standard one; they form a finite-adelic isometry. This argument uses finite-unramifiedness and does not assert global class number1.

Use the standard unitary Weil representation with trivial splitting characters, permitted because a,b are even. AM-DO-022's finite identity gives, first from the standard lattice and then by the representation law, \[
\varphi_{gN,hN'}=\omega_f(g,h)\varphi_{N,N'}.
\tag{UJ8}
\] Its hypotheses hold on K: the relative CM extension is finite-unramified on both hosts, and the actual trace-one integral lattice and its positive averages are the ones already fixed. Thus (UJ8) includes the required phase; it is not inferred from vacuum-line dimension.

The compact archimedean action on the vacuum has the exact character from DO-026(26.7), at infinity exponent zero: \[
\omega_\infty(g,h)\varphi_\infty
   =\beta(g,h)\varphi_\infty,\qquad
\beta(g,h)=\prod_\tau
       \tau(\det g)^{-b/2}\tau(\det h)^{a/2}
              =\Delta(g,h)^{-1}.
\tag{UJ9}
\] For any archimedean Heisenberg parameter w, rational theta automorphy and Heisenberg covariance give \[
\begin{split}
\Theta\bigl(\rho_\infty(w)\omega_f(g,h)\Phi\bigr)
 &=\Theta\bigl(\omega_\infty(g,h)^{-1}\rho_\infty(w)\Phi\bigr)\\
 &=\beta(g,h)^{-1}
       \Theta\bigl(\rho_\infty(A_{g,h}^{-1}w)\Phi\bigr),
\end{split}
\tag{UJ10}
\] where \(\Phi=\varphi_\infty\otimes\varphi_{N,N'}\). In the first line one inserts \(\omega_\infty(g,h)\omega_f(g,h)\), whose rational action fixes the theta functional. K035(PF25) identifies this entire-function calculation with the actual reduced reading, including its sign.

The map \(A_{g,h}\) preserves H, and it commutes with the period matrix \(\operatorname{diag}(\Omega_\tau)\), because each \(\Omega_\tau\) is scalar on its matrix block. The Gaussian factors therefore cancel on both sides of(UJ10). The same common period \(C^{14hab}\) occurs for both finite inputs, so it cancels too. This proves(UJ7).

If one instead uses the actual common stage character with infinity exponent \(-e_\tau\), its finite determinant-root factor for rational g,h is \(\prod_\tau\tau(\det g)^{-e_\tau}\tau(\det h)^{e_\tau}\), by the global character identity. Multiplication by the archimedean character \(\prod_\tau\tau(\det g)^{e_\tau-b/2}
\tau(\det h)^{-e_\tau+a/2}\) again gives \(\Delta^{-1}\). Thus the geometric scalar does not depend on this permitted splitting-character choice. ∎

The same isometry transports reduced translations exactly: \[
(U_{A_{g,h}t}F_{gN,hN'})(A_{g,h}u)
             =\Delta(g,h)(U_tF_{N,N'})(u).
\tag{UJ11}
\] This follows by inserting(UJ7) into the explicit multiplier \(U_tF(u)=e^{-\pi H(u,t)-\pi H(t,t)/2}F(u+t)\). No extra quadratic exponential or period appears. Both its scalar and its linear coordinate matrix are algebraic over K.

\subsection{3. Finitely many unitary classes at each fixed rank}\label{finitely-many-unitary-classes-at-each-fixed-rank}

Here is a full reduction of the required finiteness to the standard Minkowski second theorem. It avoids assuming that every selfdual lattice is globally standard.

For a Hermitian selfdual \(\mathcal O_K\)-lattice \(N\subset K^n\), the underlying positive integral symmetric form is \[
Q_N(x,y)=\operatorname{Tr}_{K/\mathbb Q}
                       \sum_jx_j\bar y_j.
\tag{UJ12}
\] Its trace-dual lattice is \(\mathfrak D_K^{-1}N\). Consequently its determinant is \[
\det Q_N=|\operatorname{disc}K|^n,
\qquad \operatorname{rank}_{\mathbb Z}N=n[K:\mathbb Q].
\tag{UJ13}
\] The dual assertion follows locally from Hermitian selfduality and then globally by intersection; the index is the norm of the different to the n-th power.

For clarity, positive integral Z-lattices of fixed rank d and fixed determinant \(\mathcal D\) have finitely many isometry classes. Indeed each nonzero vector has squared norm at least1. Let \(\lambda_1\le\cdots\le\lambda_d\) be the successive Euclidean minima. Minkowski's second theorem gives \[
\prod_i\lambda_i\le
        \frac{2^d\sqrt{\mathcal D}}{\operatorname{vol}(B_d)}.
\tag{UJ14}
\] Since every \(\lambda_i\ge1\), all minima are bounded by the right side. Choose independent vectors of bounded norm attaining them. Their integral Gram entries are bounded by Cauchy--Schwarz, so only finitely many Gram matrices occur. Their generated sublattice has bounded determinant and hence bounded index in the original lattice. An overlattice of index at most I lies in \((I!)^{-1}\) times that sublattice. There are finitely many such subgroups modulo the sublattice, proving the finiteness assertion. This argument does not incorrectly assume that successive-minimum vectors form a basis.

For each of the finitely many underlying trace lattices, multiplication by every fixed generator \(\eta\) of the finitely generated \(\mathbb Z\)-algebra \(\mathcal O_K\) has operator norm at most \[
\max_{\sigma:K\hookrightarrow\mathbb C}|\sigma(\eta)|
\] for \(Q_N\). In a fixed Z-basis there are therefore only finitely many integral matrices for each generator, hence finitely many possible \(\mathcal O_K\)-actions. An isometry of trace lattices commuting with this action extends to a K-linear isometry of the original Hermitian spaces: for every \(\eta\in K\), equality of the traces \(\operatorname{Tr}(\eta h(x,y))\) determines \(h(x,y)\) by the nondegeneracy of the field trace. It is thus an element of \(U_n(K)\).

It follows that the Hermitian selfdual lattices have finitely many \(U_n(K)\)-classes. Apply this separately at ranks a and b and choose representatives \[
N_1,\ldots,N_r,\qquad N'_1,\ldots,N'_s.
\tag{UJ15}
\] Every pair is \(gN_i,hN'_j\) for rational unitary g,h. Equations(UJ7) and(UJ11) are the exact scalar control needed to use this finite list.

Minkowski's inequality is classical, not a new claim here. A primary modern proof and exact statement are Henk(2002), Theorem1.3 and§3, pp.2 and5--7, cited below. All deductions from it in this section are given explicitly.

\subsection{4. The actual induction formula and volume cancellation}\label{the-actual-induction-formula-and-volume-cancellation}

This section uses the source's actual order of quotient representatives. The definition IsReps(\(\Lambda_1,\Lambda_2,S_0\)) means that \(S_0\) lies in \(\Lambda_2\) and represents \(\Lambda_2/(\Lambda_1\cap\Lambda_2)\), not the reversed quotient.

Let \(\mathcal A_c,\mathcal B_c\) denote the images in the complex period space of halfAO and halfBO. Their local rational lattices are \[
\begin{array}{c|ccc}
&P\in S& P\in\bar S& P\notin S\cup\bar S\\ \hline
\mathcal A_c&c^{-1}L_P&cL_P&L_P\\
\mathcal B_c&cL_P&c^{-1}L_P&L_P.
\end{array}
\tag{UJ16}
\] Thus \[
D\Lambda\subset\mathcal A_c,\mathcal B_c
                  \subset D^{-1}\Lambda\qquad(c\mid D).
\tag{UJ17}
\] The lattices \(\mathcal A_c,\mathcal B_c\) are selfdual for \(E_H\); the paired exponents at conjugate primes give this immediately from the selfduality of \(\Lambda\). Their restrictions of the integral trace-one element \(w_0\) are stable, so \(\alpha\) is their actual semicharacter.

Choose the source representatives \[
S_A\subset\Lambda\cap\mathcal A_c
\quad\text{for }(\Lambda\cap\mathcal A_c)/
                        (\mathcal B_c\cap\Lambda\cap\mathcal A_c),
\] \[
S_B\subset\mathcal B_c
\quad\text{for }\mathcal B_c/(\Lambda\cap\mathcal B_c).
\tag{UJ18}
\] Local calculation from(UJ16) gives \[
|S_A|=|S_B|=n_c
       =\prod_{P\in S}N(P)^{ab\,v_P(c)}.
\tag{UJ19}
\] For \(S_B\) the product initially runs over \(\bar S\); its equality follows because c is a rational integer, so conjugate primes have equal norms and valuations.

The finite representative sets in(UJ18) exist by the displayed local indices. The lift predicate is also inhabited for every x under consideration. In fact take the full fractional lattice \(\mathcal A_D\), whose localization at \(P\in S\) is \(D^{-1}L_P\). The Chinese remainder theorem for this finitely generated projective lattice supplies \(s_0\in\mathcal A_D\) with \(s_0-x\in cL_P\) for every \(P\in S\): this is reduction modulo \((cD)\mathcal A_D\) at the finitely many relevant primes. It satisfies the lift predicate with auxiliary integer \(c'=D\). Consequently the source's choice operator does choose an element satisfying that predicate, rather than an arbitrary value of an uninhabited choice.

Let \(s_0\) be that actual source liftW, allowing any of its permitted choices. Its two conditions are: \[
s_0\in\mathcal A_{c'}\text{ for some }c'>0,\qquad
s_0-x\in c\,\operatorname{homAt}_S(N,N').
\tag{UJ20}
\] Here and below x and \(s_0\) are sent to the period space by the fixed embedding when used in U or α. If \(x\in D^{-1}\operatorname{homAt}_S\), then at every \(P\in S\), \[
s_0\in D^{-1}L_P+cL_P\subset D^{-1}L_P.
\] Outside S, the first condition in(UJ20) makes \(s_0\) integral: at conjugate primes it is in \(c'L_P\), and elsewhere in \(L_P\). Therefore \[
\boxed{s_0\in D^{-1}\Lambda.}
\tag{UJ21}
\] This bound holds for every choice permitted by liftW, without a bound on \(c'\).

The literal definitions are \[
\begin{split}
\widehat\Phi(y)&=v^{-1/2}\frac v{|S_A|}
               \sum_{t\in S_A}\Phi(t)\operatorname{brk}(y,t)^{-1},\\
\xi_S(e)&=\sum_{y\in S_B}\widehat\Phi(y)\alpha(y)U_y e,\\
\xi_{\rm Ind}(e,s_0)&=\alpha(s_0)U_{-s_0}\xi_S(e),\\
\mathrm{thInd}(e)&=v^{-1/2}\xi_{\rm Ind}(e,s_0),
\end{split}
\tag{UJ22}
\] where \(v=\mathrm{volC}>0\) and \(\Phi=1_{\mathcal B_c}\). Positivity of v follows directly from its finite product of ratios of nonzero local lattice indices; the localized full lattices have finite-index intersection in both directions. The two occurrences of \(v^{-1/2}\), together with v, cancel exactly. Furthermore only one \(t\in S_A\) belongs to \(\mathcal B_c\): it is the representative of the zero class. For this t and every \(y\in S_B\subset\mathcal B_c\), selfduality makes \(E_H(t,y)\) integral, so \(\operatorname{brk}(y,t)=1\).

Using(UJ3) and linearity in the seed gives the exact simplified formula \[
\boxed{
\chi_{N,N'}^{-1}\mathrm{thInd}^{\rm mix}_{N,N',c,x}
 =\frac{\alpha(s_0)}{n_c}
       \sum_{y\in S_B}\alpha(y)U_{-s_0}U_yF_{N,N'}.
}
\tag{UJ23}
\] In particular no square root of a varying lattice-volume remains. This cancellation does not even require the parity hypothesis; parity is used for the rational determinant powers in(UJ6).

The source character cancellation in(UJ23) is literal. \nolinkurl{TrueColumnChoicesHeckeCharReferenceEqOne.lean143–152} sets e to eNat and chiF to chiFac; \nolinkurl{IsDiscriminantTwelfthRootOfIdealLatticeIffDelta1.lean140–141} defines eNat as chiFac times thMoved. For the actual characters these factors are nonzero. If a formally chosen factor were zero, both the seed and all these induction jets would be zero; that case would give algebraicity trivially rather than justify division by zero.

\subsection{5. A bounded torsion reduction, with the coefficient denominator retained}\label{a-bounded-torsion-reduction-with-the-coefficient-denominator-retained}

For \(s,t\in D^{-1}\Lambda\), integrality of \(E_H\) and of \(w_0\) gives \[
\alpha(s)\in\mu_{2D^2},\qquad
\exp(\pi iE_H(s,t))\in\mu_{2D^2}.
\tag{UJ24}
\] The reduced Heisenberg law \[
U_sU_t=e^{-\pi iE_H(s,t)}U_{s+t}
\] therefore changes a product of these translations only by a root in that fixed finite group. For \(\lambda\in\Lambda\), \[
U_{t+\lambda}F_{N,N'}
   =e^{\pi iE_H(t,\lambda)}\alpha(\lambda)\,
                              U_tF_{N,N'}.
\tag{UJ25}
\] This scalar lies in \(\mu_{2D}\). By(UJ17),(UJ21), every \(y-s_0\) in(UJ23) lies in \(D^{-1}\Lambda\). Reduction modulo \(\Lambda\) now proves \[
\boxed{
\chi^{-1}\mathrm{thInd}^{\rm mix}_{N,N',c,x}
\in n_c^{-1}\mathbb Z[\mu_{2D^2}]
       \operatorname{-span}
       \{U_tF_{N,N'}:t\in D^{-1}\Lambda/\Lambda\}.
}
\tag{UJ26}
\] The last set has exactly \(D^{2fab}\) torsion points, with arbitrary representatives allowed by(UJ25). Neither x nor an unbounded auxiliary lift denominator enlarges it. This proves the asserted explicit coefficient denominator and the 11-unit consequence in the theorem. It does not assert integrality of the individual reduced torsion germs.

\subsection{6. Completion of the uniform field proof}\label{completion-of-the-uniform-field-proof}

For each of the finitely many representative pairs in(UJ15), take every \(D\)-torsion translate in(UJ26). K035(PF30), using the actual algebraic CM gauges of K018, supplies one number field for the complete formal germs in this finite collection. Enlarge their finite compositum once to contain K, every \(cD_\tau\), and \(\mu_{2D^2}\); call the resulting field M.

For an arbitrary pair \(N=gN_i,N'=hN'_j\), the algebraic linear map \(A_{g,h}\) identifies \[
D^{-1}\Lambda_{i,j}/\Lambda_{i,j}
   \ \longrightarrow\ D^{-1}\Lambda_{N,N'}/\Lambda_{N,N'}.
\] Equations(UJ7),(UJ11) express every translated germ on the right as an algebraic scalar \(\Delta(g,h)\in K\) times a representative translated germ composed with \(A_{g,h}^{-1}\). Its matrix in the fixed \(\zeta_\tau\) coordinates has coefficients in K, because the period vector is scalar on each τ-block. Consequently every such complete series has coefficients in M. Equation(UJ26) then places every normalized induction germ in M{[}{[}u{]}{]}. Its derivatives of every order, and the weights \(cD_\tau^m\), remain in M. This proves(UJ4). ∎

The source Amb.JetsRat parameters a,b,S,D are fixed outside its quantifiers; it does not ask for one field for all ranks simultaneously. Thus the theorem has exactly its all-lattice/all-order quantifier pattern for any fixed even edge after replacing the impossible scalar/quadratic geometry by the explicit mixed geometry. If the prescribed coefficient field is smaller than M, a separate descent to that smaller field would still be necessary. The theorem permits choosing the coefficient field once; it does not assert that every preassigned field already contains the values.

\subsection{7. What this closes, and what remains}\label{what-this-closes-and-what-remains-2}

Together K035 and K036 now give an explicit varying-period Core identity, exact algebraic coordinate generators, and a single algebraic field for every admitted induction jet on each fixed even edge. The apparently arbitrary normalization chiFac cancels from the relevant arithmetic jet, and the source induction's volume factors cancel before differentiation. These are equations for the actual geometric kernel family and the literal finite induction operation.

The following are still separate:

\begin{itemize}
\tightlist
\item
  The false Fields.Kq/elliptic IsTen wrapper is not inhabited by this result.
\item
  The selected arithmetic row and subsequent image must satisfy the actual Kval and terminal comparison equations, rather than being chosen to make them true.
\item
  The field M is not given with a useful degree, discriminant, or height bound. Rational transport matrices and representative theta sections have not been bounded uniformly in the parameters.
\item
  The exact coefficient denominator \(n_c\) controls only the finite induction combination in(UJ26). It does not control reduced-gauge denominators, selected Schur saturation, all nonarchimedean row contents, or the full e/H or DP44 cost.
\end{itemize}

For the actual S that excludes11, the finite induction coefficient ring in(UJ26) is 11-integral. This is compatible with, and does not improve by itself, the sharp selected slope1/1220 in K033/K034. That slope concerns the complete reduced-gauge module and selected derivative operator; it cannot be converted into a stronger later-row bound solely by the present algebraicity theorem.

\subsection{Source locators and verification scope}\label{source-locators-and-verification-scope}

Direct source locators:

\begin{itemize}
\tightlist
\item
  ThDistC.lean72--106: Amb, IsLift, liftW, thInd and jetK;239--243: the fixed-rank Amb.JetsRat quantifiers.
\item
  XiIndCClass.lean17--34: brk, hatS, xiS and xiInd;171--197: halfAO/halfBO;216--237: period and half-period groups and representative order.
\item
  \nolinkurl{TorsionIntegrality/QuadraticSemicharacters/IsSemiChar.lean60–61:} actual IsReps orientation.
\item
  \nolinkurl{Columns/FiniteSumNonzeroLargeDepth.lean594–605:} homAt, divLat, mulLat;617: IsSdLat.
\item
  \nolinkurl{TrueColumnChoicesHeckeCharReferenceEqOne.lean143–152} and \nolinkurl{IsDiscriminantTwelfthRootOfIdealLatticeIffDelta1.lean140–141:} exactly matching eNat/chiFac normalization.
\item
  \nolinkurl{LocalNewformFunction.lean59–78:} actual global rational automorphy and Heisenberg covariance in IsWeilDatumO; the realized character calculation is DO-026(26.7).
\end{itemize}

Primary literature for the only additional standard theorem: Henk, M. (2002). \emph{Successive minima and lattice points}, Theorem1.3 and§3. \href{https://arxiv.org/pdf/math/0204158}{Author manuscript, arXiv:math/0204158}. The paper states the Minkowski bound on p.2 and gives its proof on pp.5--7. The finiteness application and all transport/induction reductions above are spelled out, rather than cited as an unspecified arithmetic compactness argument.

The companion certificate is a finite exact normalization and local-exponent check. It does not enumerate the enormous lattice-class set, construct M explicitly, prove an asymptotic cost bound, or execute Lean. Author and independent proof checks are recorded with their actual scope in the check-event TSVs.

\clearpage

\section{AM-R-019: exact normalization of the source terminal branching polynomial}\label{am-r-019-exact-normalization-of-the-source-terminal-branching-polynomial}

Author: /root, GPT-6 Astra Pro, ALP\_MACHINE, 22 November2026UTC. Reasoning effort is not externally exposed.

This finding computes the literal kapBr in the prescribed terminal operator. The classical Gelfand--Tsetlin norm theorem is an input; the normalization bridge to the source's product of minors is proved here. The resulting factorial product gives exact real and finite-place costs without enumerating the polynomial or its spectral projector. For the concrete m=2 hosts it reduces to at most (k+1)(k+2)/2 nontrivial factors, and the integer cBr has a finite factorial formula. It does not prove that the global preceding image has a nonzero terminal restriction or identify its integral lattice.

\subsection{1. Literal source normalization}\label{literal-source-normalization}

Let N=n+1 with n\textgreater=1, and let dominant integral weights lambda=(lambda\_1,\ldots, lambda\_N), mu=(mu\_1,\ldots,mu\_n) interlace:

\[
 a_i=\mu_i-\lambda_{i+1}\ge0,\qquad
 b_i=\lambda_i-\mu_i\ge0\quad(1\le i\le n).
 \tag{R19.1}
\]

For the source's SameShift condition assume mu\_n=lambda\_N. Subtract their common determinant shift from every entry when realizing their polynomial modules. All differences below are unchanged. Thus the proof may use lambda\_N=mu\_n=0.

Let X be an N-by-N matrix. Write Delta\_i for the determinant of its first i rows and columns, and Z\_i for the determinant with rows1,\ldots,i−1,N and columns1,\ldots,i, in that order. The actual source vectors and number are

\[
 p_\lambda=\prod_{i=1}^{n}\Delta_i^{\lambda_i-\lambda_{i+1}},
 \qquad \xi_{\lambda,\mu}=\prod_{i=1}^{n}\Delta_i^{a_i}Z_i^{b_i},
 \qquad \kappa(\lambda,\mu)
    =\frac{\|\xi_{\lambda,\mu}\|_F^2}{\|p_\lambda\|_F^2}.
 \tag{R19.2}
\]

Here monomials are orthogonal and \textbar\textbar X\^{}alpha\textbar\textbar{}\emph{F\^{}2=alpha!. These are exactly xiBr and kapBr in \nolinkurl{Columns/FiniteSumNonzeroLargeDepth.lean119–133,} with aBr/bBr41--44. The accepted source interlacing theorem says that xi is the nonzero GL\_n-highest vector of weight mu in the lambda representation and coefLast(xi)=p\_mu under SameShift. For the intermediate lowering steps we also need its elementary polynomial realization without SameShift. Every Delta\_i and Z\_i is killed by E\_pq for p\textless q\textless=n: either row q is absent, or row p is already present, producing a repeated row. The row weights of xi telescope to mu\_1,\ldots,mu\_n and sum(lambda)−sum(mu). Its column raising operators also kill it, and its column weight is lambda because a\_i+b\_i=lambda\_i− lambda}(i+1). The Cauchy decomposition therefore places it in the row representation of highest weight lambda for every interlacing mu. Evaluation at the matrix X\_* of Section3 gives one, proving it is nonzero. This supplies the intermediate realization directly; it does not extend the source SameShift interface silently.

\subsection{2. Exact factorial formula}\label{exact-factorial-formula}

\subsubsection{Theorem R19.3}\label{theorem-r19.3}

\[
\boxed{
\kappa(\lambda,\mu)=
\prod_{1\le i\le j<N}
 \frac{(\lambda_i-\mu_j+j-i)!}{(\mu_i-\mu_j+j-i)!}
\prod_{1\le i<j\le N}
 \frac{(\mu_i-\lambda_j+j-i-1)!}
      {(\lambda_i-\lambda_j+j-i-1)!}.}
\tag{R19.3}
\]

Every factorial argument is a nonnegative integer. A more useful form, with x\_ij=mu\_i−mu\_j+j−i, is

\[
\boxed{
\kappa(\lambda,\mu)=
 \prod_{1\le i\le j\le n}
 \frac{(x_{ij}+b_i)!(x_{ij}+a_j)!}
      {x_{ij}!(x_{ij}+a_j+b_i)!}.}
\tag{R19.4}
\]

In particular 0\textless kappa\textless=1. A factor is one if a\_j=0 or b\_i=0. Formula(R19.4), with Legendre's factorial valuation, computes v\_p(kappa) exactly; no cancellation or integrality guess is needed.

\textbf{Proof of the last claims from(R19.3).} In the second product put j=r+1 and group it with the term (i,r) in the first. The four arguments become x+b\_i, x+a\_r, x, x+a\_r+b\_i. The ratio is \(\prod_{s=1}^{a_r}(x+s)/(x+b_i+s)\), hence is positive and at most one. Interlacing and dominance give x\textgreater=0 and all remaining arguments nonnegative. It remains to prove the precise normalization in(R19.3).

\subsection{3. Lowering-operator normalization: a complete minor calculation}\label{lowering-operator-normalization-a-complete-minor-calculation}

Use the row operators

\[
 E_{pq}=\sum_{c=1}^{N}X_{pc}\frac{\partial}{\partial X_{qc}},
 \qquad h_i=E_{ii}-i+1.
\]

They satisfy E\_pq\^{}*=E\_qp for the Fock metric. Molev's lowering operator, equation(2.10), is

\[
 z_{Ni}=\sum_{i<i_1<\cdots<i_s<N}
 E_{i_1i}E_{i_2i_1}\cdots E_{Ni_s}
       \prod_{j\notin\{i_1,\ldots,i_s\}\atop i<j<N}(h_i-h_j).
 \tag{R19.5}
\]

The empty path term is E\_Ni times all displayed diagonal factors. Molev's Lemma2.5 and multiplicity-free branching put its image of xi on the mu−epsilon\_i highest line whenever a\_i\textgreater0. We claim

\[
\boxed{
 z_{Ni}\xi_{\lambda,\mu}
   =\left(\prod_{j=i+1}^{N}
       (\mu_i-\lambda_j+j-i-1)\right)
        \xi_{\lambda,\mu-\epsilon_i}.}
\tag{R19.6}
\]

To determine the scalar, evaluate at X\_* whose first n rows are the corresponding coordinate rows and whose last row has every entry one. Each Delta\_l and Z\_l for l\textless=n is one, so both relevant xi vectors evaluate to one. Fix a path p\_0=i\textless p\_1\textless\ldots\textless p\_s\textless N. The corresponding differential product in(R19.5) is the coefficient of t\_0\ldots t\_s in xi evaluated after the simultaneous row replacements

\[
 X_{p_r,*}\mapsto e_{p_r}^*+t_r e_{p_{r+1}}^*\quad(r<s),
 \qquad
 X_{p_s,*}\mapsto e_{p_s}^*+t_s(1,\ldots,1).
 \tag{R19.7}
\]

Indeed iterated differentiation represents the operator product by left multiplication with the elementary row-update matrices in reverse order. Their pairwise cross-products in this order vanish, so the simultaneous replacements are exact. Every differentiation is first order in its own variable; coefficient and derivative agree.

Put B\_0=1 and B\_r=1−t\_(r−1)B\_(r−1). Direct triangular determinant expansion of(R19.7) gives

\[
 \Delta_l=\begin{cases}1,&l<p_s,\\1+t_sB_s,&l\ge p_s,\end{cases}
 \qquad
 Z_{p_r}=B_r\ (1\le r\le s),\qquad Z_l=1\ (l\notin\{p_1,\ldots,p_s\}).
 \tag{R19.8}
\]

For example the upper triangular path matrix has inverse-column entries 1,−t\_(s−1),t\_(s−2)t\_(s−1),\ldots, which sum to B\_s; the rank-one determinant formula gives the first identity. Replacing row p\_r by the last row gives the same inverse-column sum truncated at r, proving the second. Rows not on the path do not connect to the missing column, so their determinants remain one.

Set A\_j=sum\_(l=j)\^{}n a\_l. Substitution in xi is therefore

\[
 (1+t_sB_s)^{A_{p_s}}\prod_{r=1}^sB_r^{b_{p_r}}.
\]

First extract t\_s; it contributes A\_(p\_s) and one extra factor B\_s. Then extract t\_(s−1),\ldots,t\_0 successively using B\_r=1−t\_(r−1)B\_(r−1). The result is exactly

\[
 (-1)^s A_{p_s}\prod_{r=1}^s(b_{p_r}+1).
 \tag{R19.9}
\]

This also covers the empty path, with value A\_i. The rightmost diagonal factors in(R19.5) act first on xi, with eigenvalues M\_j=mu\_i−mu\_j+j−i. Write c\_j=b\_j+1 and M\_i=0. Summing paths according to their largest vertex, and then using A\_l=sum\_(r\textgreater=l)a\_r, gives the scalar

\[
\begin{aligned}
S_i&=A_i\prod_{j>i}M_j
 -\sum_{l>i}c_lA_l
        \prod_{i<j<l}(M_j-c_j)\prod_{j>l}M_j\\
&=\sum_{r=i}^{n}a_r
        \prod_{i<j\le r}(M_j-c_j)\prod_{j>r}M_j.
\end{aligned}
\tag{R19.10}
\]

Both equalities are finite product expansions; the second is the telescoping identity for replacing each M\_j by M\_j−c\_j in order. Now M\_j−c\_j=M\_(j−1)+a\_(j−1). Thus the last expression is

\[
 \prod_{r=i}^{n}(M_r+a_r)-\prod_{r=i}^{n}M_r
 =a_i\prod_{r=i+1}^{n}(\mu_i-\lambda_{r+1}+r-i),
\]

since M\_i=0. This is the coefficient in(R19.6). If a\_i=0, branching makes the image zero and the same scalar calculation has value zero. This completes the normalization proof.

\subsection{4. Application of the classical norm theorem}\label{application-of-the-classical-norm-theorem}

Starting at p\_lambda, apply the lowering operators in Molev's Lemma2.6/Theorem2.7 order, so the largest i acts first. Each intermediate weight interlaces; at coordinate i, decrease from lambda\_i to mu\_i. Formula(R19.6) depends on that coordinate and lambda, independently of the previously decreased coordinates. The total scalar is therefore

\[
 C_{\lambda,\mu}
 =\prod_{1\le i<j\le N}
       \frac{(\lambda_i-\lambda_j+j-i-1)!}
            {(\mu_i-\lambda_j+j-i-1)!},
 \qquad \xi_{\mu}^{\mathrm{GT}}=C_{\lambda,\mu}\xi_{\lambda,\mu}.
 \tag{R19.11}
\]

The Fock metric divided by \textbar\textbar p\_lambda\textbar\textbar\_F\^{}2 is the contravariant metric with unit highest vector used in Molev's Proposition2.4. Choose the lower rows of the GT pattern to be the initial segments of mu. The factors in that proposition from levels below N are one. Its top-level factor is

\[
 \|\xi_{\mu}^{\mathrm{GT}}\|^2
  =\left[\prod_{1\le i\le j<N}
       \frac{(\lambda_i-\mu_j+j-i)!}{(\mu_i-\mu_j+j-i)!}\right]
       C_{\lambda,\mu}.
\]

Dividing by C\_(lambda,mu)\^{}2 proves(R19.3). This invokes the published norm theorem; it does not reprove it or identify a basis by its name alone. The direct minor calculation supplies the needed normalization to the actual source vector. ∎

\subsection{5. The exact integer spectral factor}\label{the-exact-integer-spectral-factor}

Set d0=sum\_i a\_i, L\_i=lambda\_(i+1), U\_i=min(lambda\_i,lambda\_(i+1)+d0), and R\_ij=mu\_j+i−j. Then L\_i\textless=mu\_i\textless=U\_i. The literal source integer is

\[
 c_{\rm Br}=\prod_{i=1}^{n}\prod_{v=L_i\atop v\ne\mu_i}^{U_i}
                          \prod_{j=1}^{n}(v-R_{ij}).
 \tag{R19.12}
\]

For j\textless i, R\_ij\textgreater U\_i; for j\textgreater i, R\_ij\textless L\_i, by interlacing. Consequently it is a nonzero integer, and its exact signed factorial evaluation is

\[
 c_{\rm Br}=\prod_i\left[(-1)^{\mu_i-L_i}(\mu_i-L_i)!(U_i-\mu_i)!
   \prod_{j<i}\frac{(-1)^{U_i-L_i}(R_{ij}-L_i)!}
          {(R_{ij}-U_i-1)!(R_{ij}-\mu_i)}
   \prod_{j>i}\frac{(U_i-R_{ij})!}
          {(L_i-R_{ij}-1)!(\mu_i-R_{ij})}\right].
 \tag{R19.13}
\]

\textbf{Proof.} The j=i product has its single zero omitted and is the displayed two factorials with mu\_i−L\_i negative factors. For j\textless i, all factors are negative and one is omitted, leaving U\_i−L\_i of them; the full consecutive product of their absolute values is divided by the omitted R\_ij−mu\_i. For j\textgreater i all factors are positive and the same argument gives the last expression. Their product is exactly(R19.12). If L\_i=U\_i=mu\_i, every factor is one. No enormous interval needs to be enumerated. ∎

\subsection{6. The actual terminal operator norm and its scalar valuations}\label{the-actual-terminal-operator-norm-and-its-scalar-valuations}

The accepted source identities InterlaceAll.IoHw, CoefXi, BrTranspose and BeOnEq give

\[
 \operatorname{brStd}(\xi_{\lambda,\mu})
       =c_{\rm Br}\kappa(\lambda,\mu)p_\mu,
 \qquad \|\operatorname{brStd}\|=|c_{\rm Br}|\sqrt\kappa.
 \tag{R19.14}
\]

For the tensor over CM embeddings and the literal t\_tau=tau(star piT), write nBr\_tau=sum lambda\_tau−sum mu\_tau. Then

\[
\boxed{
\|\operatorname{brAt}\|
 =\prod_\tau |t_\tau|^{-nBr_\tau}
                  |c_{{\rm Br},\tau}|\sqrt{\kappa_\tau}.}
\tag{R19.15}
\]

For R017's actual norm element, every \textbar t\_tau\textbar=sqrt(T). Thus its logarithm is exactly

\[
 -\tfrac12\log T\sum_\tau nBr_\tau
   +\sum_\tau\log|c_{{\rm Br},\tau}|
   +\tfrac12\sum_\tau\log\kappa_\tau.
 \tag{R19.16}
\]

At11, t\_tau is a unit. The scalar on the primitive xi-to-p\_mu line therefore has valuation

\[
 v_{11}(c_{{\rm Br},\tau})+v_{11}(\kappa_\tau),
 \tag{R19.17}
\]

exactly calculable from(R19.4),(R19.13) and v\_p(r!)=sum\_(s\textgreater=1)floor(r/p\^{}s). This is a scalar valuation on that specified line; it is not the Smith content of the full projected integral row lattice. The latter also depends on the actual inclusion matrix and is retained as an open obligation.

\subsection{7. Concrete selected m=2 weights on J14 and J56}\label{concrete-selected-m2-weights-on-j14-and-j56}

Use K031/K032 and R018. Let k be42 or168, d be1760 or1634, C be45156049293 or40530339687, and N=k+2d. The terminal column and row weights, before removing their common determinant shift, are

\[
\begin{aligned}
 \lambda_i&=(2(k-i)+3)C+i-k/2-d &&(1\le i\le k),\\
 \lambda_{k+1},\ldots,\lambda_{k+d}&=k/2-d/2,\\
 \lambda_{k+d+1},\ldots,\lambda_N&=-4C+(k+d)/2;\\
 \mu_i&=2(k-i+1)C+i-k/2-d &&(1\le i<k),\\
 \mu_k&=2C+k/2-d,\\
 \mu_{k+1},\ldots,\mu_{k+d-1}&=k/2-d/2,\\
 \mu_{k+d}&=-2C+k/2,\\
 \mu_{k+d+1},\ldots,\mu_{N-1}&=-4C+(k+d)/2.
\end{aligned}
\tag{R19.18}
\]

These are lamTrue\_old/new and sgTrue\_eq, with RowPar.sgHC, not newly chosen interlacing weights. In particular the common shift is −4C+(k+d)/2. Their nonzero a,b entries are

\[
\begin{array}{c|cc}
 i&a_i&b_i\\ \hline
 1,\ldots,k-1&C-1&C\\
 k&2C-d/2&C\\
 k+d&2C-d/2&2C-d/2.
\end{array}
\tag{R19.19}
\]

All others vanish. Hence the nontrivial factors of(R19.4) number at most(k+1)(k+2)/2, and

\[
 d0=(k+3)C-(k-1)-d,\qquad
 nBr=(k-2)C+k/2.
 \tag{R19.20}
\]

Increasing the first u entry by an integer A\textgreater=0 changes only lambda\_1, adding A; the row mu, every a\_i and d0 stay fixed, while b\_1 and nBr increase by A. Put A\_\emph{=max(0,d0−(lambda\_1−lambda\_2)). Then cBr is independent of A for A\textgreater=A\_}: only U\_1 varies, and it has reached lambda\_2+d0 by that point. Moreover

\[
\boxed{
 \frac{\kappa(\lambda+A\epsilon_1,\mu)}{\kappa(\lambda,\mu)}
  =\prod_{j=1}^{n}
    \frac{(\lambda_1-\mu_j+j)_{a_j}}
         {(\lambda_1+A-\mu_j+j)_{a_j}},}
\tag{R19.21}
\]

where(x)\_r=x(x+1)\ldots(x+r−1), with(x)\_0=1. This follows by grouping the i=1 factors in(R19.4); every other factor is unchanged. In particular kappa is a positive rational function of A with exact decay degree d0. The real term in (R19.16) is therefore explicit even on K033's growing family: after the cBr threshold its extra logarithm is −(A/2)log T−(d0/2)log A+O(1) at the growing embedding. The constant depends only on the fixed host and fixed row. The source size restriction from R018 still limits A at a fixed T; this asymptotic formula is not an assertion that fixed-T weights remain admissible as A tends to infinity.

\subsubsection{Exact11-adic values on the actual weight family}\label{exact11-adic-values-on-the-actual-weight-family}

The saved certificate evaluates the following integers by Legendre sums and checks the actual row weights against the independent literal RowPar.sgHC formula:

\begin{longtable}[]{@{}lrr@{}}
\toprule\noalign{}
Host & d0 & A\_* \\
\midrule\noalign{}
\endhead
\bottomrule\noalign{}
\endlastfoot
J14 & 2,032,022,216,384 & 1,941,710,117,799 \\
J56 & 6,930,688,084,676 & 6,849,627,405,303 \\
\end{longtable}

\begin{longtable}[]{@{}lrrr@{}}
\toprule\noalign{}
Host & v11(kappa) & v11(cBr), baseline & v11(cBr), saturated \\
\midrule\noalign{}
\endhead
\bottomrule\noalign{}
\endlastfoot
J14 & −118 & 1,431,126,153,655,645 & 2,122,569,126,603,571 \\
J56 & −3423 & 4,747,460,542,799,389 & 7,100,307,556,522,475 \\
\end{longtable}

These are per embedding. The kappa valuation is unchanged after increasing lambda\_1 by A=11\^{}r or121\emph{11\^{}r for any r\textgreater=13. Indeed every positive integer factor in the numerator rising factorials of(R19.21) is at most S\_2−2, where S\_2 is the exact terminal spread in R018. In both hosts S\_2−2\textless11\^{}13. Adding 11\^{}r or121}11\^{}r therefore preserves the valuation of each factor. Both A choices also exceed A\_* at r\textgreater=13, so the displayed cBr valuation is the saturated one. This is an all-r proof, not an extrapolation from the two sampled endpoints in the certificate.

For the actual family with two growing CM embeddings and all other embeddings fixed, the valuation of the full highest-line branching scalar in(R19.17) is consequently

\begin{longtable}[]{@{}lr@{}}
\toprule\noalign{}
Host & v11 of the complete specified highest-line scalar, r\textgreater=13 \\
\midrule\noalign{}
\endhead
\bottomrule\noalign{}
\endlastfoot
J14 & 61,490,184,399,427,986 \\
J56 & 802,279,065,217,168,460 \\
\end{longtable}

At the two growing embeddings use the saturated cBr value; at the other f−2 embeddings use its baseline value. Add all kappa valuations. Since piT is an11-unit there is no extra T contribution at this prime. The sign of cBr is negative in both the baseline and saturated per-place cases, as checked by(R19.13). These large exact numbers are scalar normalizations of the prescribed branch on its primitive highest line. They do not give the valuation of an unknown global terminal vector or the full integral projector matrix. R018's fixed-T admissible ranges still apply when these formulas are used in that numerical family.

\subsection{8. Costs, framework, sources and check scope}\label{costs-framework-sources-and-check-scope}

This removes the unevaluated branching-polynomial norm from the terminal norm budget. It also replaces the spectral interval product by an exact finite factorial formula. It does not replace the actual global terminal vector by xi, identify a detecting row at zero cost, or assign this factor twice in DP44: the source's separate T-term already carries its own prescribed period power. R002's content/dual-norm formulas must still be applied to the actual image and row lattice.

The polynomial representation and its restriction live in the integral representation/flag setting of SGA3 I4.7.3/4.7.3.1 and SGA6 VI3.3 recorded in the dictionary. The factorial scalar and Fock norm refine this representation class to a metrized map. The polarized theta line is still the Picard/biextension object of FGA232 Theorem3.1 and SGA7 I VII2.7/2.9, VIII7.4(b); those representability results do not determine the present metric normalization. The formal/crystalline slopes of cote4 and cote54 scans5--7 are unchanged. This is a DP\_MULT Q3e/THETA3 G1 normalization result, with no consequence asserting Kummer RH.

Exact stipulated source locators, commit 7cbf7a0c32fca1badf05931879a37c3f8a2b03b4:

\begin{itemize}
\tightlist
\item
  \nolinkurl{Columns/FiniteSumNonzeroLargeDepth.lean41–55,119–151:} aBr,bBr,rSet,cBr,xiBr,kapBr,brStd,brAt.
\item
  \nolinkurl{KernelLattices/InterlaceAll.lean54–71} and125--147: BrTranspose,IoHw,CoefXi,beOnVal,BeOnEq. Their validity is accepted as instructed; no Lean build is performed.
\item
  \nolinkurl{Columns/AbelianHostFieldConditions.lean408–421,464–472:} RowPar.sgHC,sgTrue,lamTrue; \nolinkurl{Columns/BaseLemmas.lean496–530,} 575--582: the actual terminal weight formulas.
\end{itemize}

Molev, A. I. (2002). \emph{Gelfand--Tsetlin bases for classical Lie algebras} (arXiv:math/0211289v2). \url{https://arxiv.org/abs/math/0211289} . Proposition2.4, equation(2.10), Lemmas2.5--2.6 and Theorem2.7, printed pp10--11 (PDF pp11--12). These exact primary loci were read in this run. The norm theorem is cited; the source-specific minor normalization is proved in Sections3--4 above.

Finite exact diagnostics in branch\_norm\_diagnostic.py compare 250 complete Fock norms and251 lowering actions for ranks2--5; the combined path replay checks2550 paths. They support the full proof and are not substituted for it. The companion arithmetic certificate records exact actual-weight valuations without expanding huge factorials. Whole-proof reviews are logged separately after the entire finished proof is read.

\clearpage

\section{ALP-LIFT-26: the full prescribed sharp two-edge coupled lift is nonzero}\label{alp-lift-26-the-full-prescribed-sharp-two-edge-coupled-lift-is-nonzero}

Author: GPT-6 Astra Pro, ALP\_MACHINE /root/lift\_accounting. 22 November 2026 UTC. Inherited reasoning effort is not exposed.

This result assembles the checked actual local quotient maps into the fixed global m=2 coupled operator on either repaired mixed CM host. The initial automorphic function is the actual R017 function. The finite kernels are the unchanged sharp G/T/Q kernels, with maximal vacua at the remaining places. The real kernels are the two primitive selected highest polynomials of LIFT25. The full prescribed Y quotient is summed. No later Schwartz replacement or Hecke translation is made.

The distinction between this statement and unrestricted theta nonvanishing is essential. Rallis is first used to prove that the global map on every automorphic multiplicity copy is nonzero. Local Howe uniqueness then identifies the kernel of that global map with the product of the actual local small-theta quotient kernels. The explicit local calculations detect the particular coupled tensor. We do not infer fixed-vector survival from GQT Proposition11.6 alone.

The resulting nonzero analytic automorphic vector is not yet the source's arithmetic row or terminal period. Its global Rallis and class-mass normalization, integral content, and projection to the prescribed terminal branching component remain explicit obligations.

\subsection{1. Actual data and precise statement}\label{actual-data-and-precise-statement}

Let E/F be either K031's degree84 CM host with its principally polarized J14, or K032's degree336 CM host with its principally polarized J56. Write

\[
 (a,d,b,c)=(42,1760,1802,3562)
       \quad\hbox{or}\quad(168,1634,1802,3436),
 \qquad G_j=U(I_{a+jd}),\quad j=0,1,2.
 \tag{GS1}
\]

All three Hermitian spaces are positive at every real embedding. The relative CM extension is unramified at finite places; its dyadic places split. The actual selfdual trace lattices, additive character, vacuum, creation operators and splitting are those of K030/032 and DO022/034, as used in DO036. These geometric and character results are carried inputs, not reasserted constructions.

Denote the actual common splitting character on edge j by chi\_j. In the source notation it is mu\_j times the fixed tuple character xi. Both members of that edge have this same character. In particular

\[
 \upsilon=\chi_0\chi_1^{-1}=\mu_0\mu_1^{-1},\qquad
 \chi_j|_{\mathbb A_F^\times}=1.
 \tag{GS2}
\]

The fixed tuple character cancels in the ratio. At the two rational G primes and all their selected places, the two residue characters are distinct and nontrivial. At T both chi\_j are unramified. At Q they have the same specified nontrivial residue. These are the literal conditions checked in LIFT18/19, K031/032, DO035, and LIFT24. At infinity the exponents are -e\_j, with e\_1-e\_0=4C\textgreater0 for the concrete uniform data. Hence upsilon and upsilon squared have nontrivial angular infinity type.

Let F\_0 be R017's actual nonzero class-supported automorphic function on G\_0. It has the matched selected highest real line, finite Steinberg generic type at G, prescribed nontrivial projective-line type at T, and maximal compact invariance at Q and all other finite places. Its exact norm and source level are the ones already computed in R017 and DO036; no new source input is substituted.

At infinity use phi\_(j,infty) from LIFT25(AR5) on both edges.

These are the Chern-coordinate creation vectors. For the literal source \nolinkurl{pureTensorO} vectors one must multiply edge j by the fixed nonzero scalar \(\Lambda_j=\prod_\tau\lambda_\tau^{D_{j,\tau}}\) of LIFT25(AR27)--(AR28), as identified in K035(PF22)--(PF23). All statements below about the Chern-normalized operator therefore hold for the literal source operator after multiplication by \(\Lambda_0\Lambda_1\). Its squared output norm gains \(|\Lambda_0\Lambda_1|^2\); the exact source coevaluation norms are K037(CR35), or LIFT25(AR29). The coefficient comparison retains \(C_\theta^{14h(ab+bc)}\prod_\tau(\lambda_\tau/cD_\tau)
 ^{D_{0,\tau}+D_{1,\tau}}\), as in LIFT25(AR30). This conversion does not change the finite kernels or prior-image nonvanishing, and its complex period factors are not assigned finite valuations. At finite places use the literal sharp XiCpl1 and XiCpl2 factors, including ftY's exact power q\^{}(-M b), all character phases, and the actual piT in the terminal row. Let the complete coupled tensor be

\[
 \Psi=\sum_{Y\in\mathcal Y}\phi_{1,Y}\otimes\phi_{0,Y},
 \tag{GS3}
\]

where the sum is expressed in standard local frames. Section2 proves that this is exactly the literal moving-lattice Y sum, not a new choice of representatives. Define the actual bilinear coupled lift

\[
 \mathcal B F_0
   =\sum_{Y\in\mathcal Y}
       \Theta_{1,\phi_{1,Y}}^{\rm bil}
          \bigl(\Theta_{0,\phi_{0,Y}}^{\rm bil}F_0\bigr),
 \qquad
 \Theta_{j,\phi}^{\rm bil}F(h)
    =\int_{[G_j]}\theta_{j,\phi}(g,h)F(g)\,dg.
 \tag{GS4}
\]

\textbf{Theorem.} For either host and its stated actual data,

\[
                         \mathcal B F_0\ne0.
 \tag{GS5}
\]

The result holds on every nonzero vector in the same finite-level, matched-highest input subspace whose G representation lies in the finite Steinberg isotypic subspace. It includes all automorphic multiplicity copies and their finite linear combinations.

This is a theorem about the complete fixed coupled operator. In particular at least one admitted matched Y pair has a nonzero double lift. It does not say that DO036's previously selected single Y\_0 must survive the second edge: at Q the proof uses the full Y sum. The nonzero function on G\_2 does not by itself give a nonzero literal brC/tilQ row period. The source ChKer.cp2 subsequently evaluates at tilQ(d), and BtOf then applies PtilFF. The Q graph used in(GS22) is not in that embedded terminal subgroup; it cannot establish that restricted endpoint detection. This remaining terminal gate is distinct from(GS5).

\subsection{2. Exact factorization of the literal Y family}\label{exact-factorization-of-the-literal-y-family}

The source definitions are \nolinkurl{Columns/FiniteSumNonzeroLargeDepth.lean}1120--1169: locLat, ybSet, repY, XiCpl1, PhiTop, ftY, and XiCpl2. The Y support and period modules are intersections of the prescribed local modules. DO032 proves their CRT identification

\[
 \mathcal Y\simeq
 \prod_{P\in G}(q_P^{-1}O_P^b/O_P^b)
 \times\prod_{P\in Q}(q_P^{-M}O_P^b/q_P^M O_P^b).
 \tag{GS6}
\]

For completeness, localization sends a global coset injectively to its local cosets because the global period module is precisely the intersection of those local conditions. Clearing the finitely many denominators and applying CRT gives surjectivity. The prescribed middle frame identifies its localized lattice with O\_P\^{}b. This argument applies to every middle lattice, not just the identity class.

There is no archimedean chamber condition on repY. A change of repY inside its coset changes y\_P by O\_P\^{}b at G and q\_P\^{}M O\_P\^{}b at Q. The first raw factor is periodic for exactly those changes; on the terminal support, ftY's change has integral additive-character argument. LocCst.Ok.psi\_int therefore makes it one. These periodicity checks are DO031 Section2 and DO032 Section1, with the first sharp G check in LIFT22(CG7--9) and the full Q check in LIFT23/DO037. No phase depending on the global representative remains.

At a middle adelic frame h, use that frame in both raw factors. The global coset Y maps by (GS6) to the local classes of h\_P\^{}\{-1\}y. Those classes run through the entire fixed standard quotient exactly once. Therefore the finite sum of the products of the two raw local functions is exactly

\[
 \Psi_f=\bigotimes_{v<\infty}\Psi_v,\qquad
 \Psi_v=
 \begin{cases}
  \displaystyle\sum_{y\in\mathcal Y_v}
                 \phi_{1,v,y}\otimes\phi_{0,v,y},&v\in G\cup Q,\\
  \phi_{1,v}\otimes\phi_{0,v},&v\notin G\cup Q.
 \end{cases}
 \tag{GS7}
\]

All factors in the source product use the selected place over v; the opposite split component is already encoded in the honest unitary Weil representation. It is not an additional independent Y variable. At infinity there is the fixed tensor phi\_(1,infty) tensor phi\_(0,infty). Equation(GS7) is the finite distributive law after CRT, with the exact source scalars inside the factors.

At G the first sharp factor forces the sole class y=q\^{}(-1)e\_a; LIFT22(CG7--9) proves this in every target coordinate. At Q all classes in (GS6) remain. The complete local Q sum is the one in DO037, including its Fourier/Gauss scalar.

\subsubsection{Why this is one fixed Schwartz tensor under the diagonal action}\label{why-this-is-one-fixed-schwartz-tensor-under-the-diagonal-action}

The following pointwise operator identity is stronger than a bare CRT count. Every place at which Y varies is relative split. Identify its unitary group with GL over F\_v using the selected E-place, as in DO020 and LIFT24. Every prescribed local frame is then an actual GL matrix and hence an element of this split unitary group; its opposite E-component is its inverse adjoint. Write g,h,t for the initial, middle and terminal frames at this place. The genuine Hom Weil action on the two Schwartz spaces is

\[
\begin{aligned}
 [\omega_0(g,h)f](X)
   &=\kappa_0(g,h)f(h^{-1}Xg),\\
 \kappa_0(g,h)
   &=\eta_0(\det g)^{-1}\eta_0(\det h)
                         |\det g|^{b/2}|\det h|^{-a/2},\\
 [\omega_1(h,t)H](Z)
   &=\kappa_1(h,t)H(t^{-1}Zh),\\
 \kappa_1(h,t)
   &=\eta_1(\det h)^{-1}\eta_1(\det t)
                         |\det h|^{c/2}|\det t|^{-b/2}.
\end{aligned}
 \tag{GS7a}
\]

These are the honest split formulas with the prescribed character phase, not only symplectic covariance up to a scalar. For the fixed standard functions f\_(0,z),f\_(1,z) define Psi\_v by(GS7). Linearity of the two actual Weil operators gives exactly

\[
\begin{aligned}
 &[(\omega_1(h,t)\otimes\omega_0(g,h))\Psi_v](Z,X)\\
 &\quad=\kappa_1(h,t)\kappa_0(g,h)
    \sum_{z\in\mathcal Y_v}
      f_{1,z}(t^{-1}Zh)f_{0,z}(h^{-1}Xg).
\end{aligned}
 \tag{GS7b}
\]

For the literal middle lattice N\_h, its local support and period modules are hL\_+ and hL\_-, where L\_+/L\_- is the standard quotient in(GS6). The source frsF formula in FS1103--1104 uses precisely that same frame h in both XiCpl1 and XiCpl2. Thus its raw coordinate sum at (g,h,t) is the sum displayed after the two kappa factors in (GS7b): the local coordinate h\^{}(-1)repY runs through z in L\_+/L\_-, and the previously proved periodicity removes the difference between representatives. This comparison holds for arbitrary X and Z before taking theta sums. Multiplying by the actual frame factors(GS7a), as required by the genuine Weil reading of those raw arrays, gives equality with the diagonal-middle action on the \textbf{same fixed Psi\_v} at every h, not insertion of a new h-dependent projector.

In particular the middle character and modulus in that product are

\[
       (\eta_0\eta_1^{-1})(\det h)|\det h|^{(c-a)/2}
             =(\eta_0\eta_1^{-1})(\det h)|\det h|^d.
 \tag{GS7c}
\]

They are independent of z and agree with the exact middle factor in DO037's double full-rank period. They have not been replaced by one. The initial and terminal factors in(GS7a) likewise remain present.

At places outside G union Q there is no Y insertion at all: the fixed tensors transform by their genuine Weil operators. Infinity also has the fixed selected tensor. Taking the product of(GS7b) and these unchanged factors proves the global fixed-tensor identity. Finally summing over the two rational matrix variables gives the actual product of theta kernels. This proves that its global composition is evaluated through q\_1(1 tensor q\_0), as used in (GS19) and(GS27), rather than through an unproved moving-family operator. The fixed common period/creation factors in the actual mixed realization multiply both descriptions equally and are independent of Y; they remain in the chosen phi\_j and A\_j.

The first and middle quotients are compact. Theta series for these Schwartz vectors converge uniformly on compact representatives, and the Y sum is finite. Thus the double integral and finite sum in (GS4) may be interchanged without a regularization or rearrangement. The formula is the actual composition of two global theta maps evaluated on the tensor(GS3), with the actual quotient Haar measures.

\subsection{3. The standard factor for every actual first constituent}\label{the-standard-factor-for-every-actual-first-constituent}

Fix an irreducible automorphic constituent pi\_0 of the input space. It is unitary. DO036 proves nonzero first lift on every embedding of pi\_0 that occurs here, using the unchanged first raw kernel and its strict positive norm bound. Its real type is the literal selected one, and its first descendant is denoted pi\_1.

At an unramified place w of E outside a finite set S, local theta gives exactly

\[
 \mathcal A_{1,w}=\mathcal A_{0,w}\sqcup
 \{\chi_{0,w}(\varpi_w)Q_w^{(d-1)/2-l}:0\le l<d\}.
 \tag{GS8}
\]

Here Q\_w is the residue cardinality in E and A is the standard Satake multiset. At relative split places this is Mínguez's type-II formula with the actual Hom factor eta\_T(det g)\^{}(-1)eta\_S(det h). At nonsplit places it is Liu's unramified theorem in the plus tower, with the paired eigenvalues in Q\_w=q\_v\^{}2 units. These are the checked LIFT14(E10--13) formulas. Absolute ramification of the host is allowed; the assertion is only made outside S when identifying Euler parameters.

DO028 fixes the bilinear-input and Hom-to-tensor conventions. In particular the standard factor on edge1 uses chi\_1 inverse. Set

\[
 s_0=d/2,\qquad s_*=(d+1)/2.
 \tag{GS9}
\]

Substitution of (GS8), initially in a right half-plane, gives

\[
 \mathcal L_1^S(s_*)=
 L_E^S(s_*,\operatorname{BC}_{\rm unr}(\pi_0)\chi_1^{-1})
          \prod_{l=1}^{d}L_E^S(l,\upsilon).
 \tag{GS10}
\]

The notation BC\_unr means only the actual standard unramified Euler parameters. No global automorphic base-change representation or Arthur parameter is assumed. The equality at s\_* is obtained by continuing the identity of Euler products. In particular we do not infer the first constituent's real Langlands parameter from (GS8).

The crude unitary spherical bound used and checked in LIFT14 is

\[
                  |\alpha_w|\le Q_w^{(a-1)/2}.
 \tag{GS11}
\]

It applies to every initial constituent, without temperedness. The real part available in its first factor in (GS10) is

\[
 s_*-(a-1)/2=(d-a+2)/2=860\quad\hbox{or}\quad734.
 \tag{GS12}
\]

That Euler product is therefore absolutely convergent and nonzero. The factors at l\textgreater=2 are also absolutely convergent and nonzero. For l=1, upsilon and its square have nontrivial angular type by (GS2). The already proved LIFT14 boundary lemma applies over this field as well: Hecke--Tate continuation gives holomorphy at1, and

\[
 3+4\cos\theta+\cos2\theta=2(1+\cos\theta)^2\ge0
 \tag{GS13}
\]

in the Euler logarithms excludes a zero at1. Consequently (GS10) is holomorphic and nonzero at the stated point. This is not a claim of absolute Euler-product convergence for the l=1 factor, nor a quantitative lower bound for that value.

The finite G ratio is also nontrivial, which independently excludes the trivial Hecke character. The angular condition and its square are retained because they are the hypotheses of the exact boundary proof being used. The fixed stage characters are not varied.

\subsection{4. Rallis gives a nonzero global second map on every embedding}\label{rallis-gives-a-nonzero-global-second-map-on-every-embedding}

For the second edge n=b=1802 and m=c=b+d, we have

\[
                         n<m<2n,\qquad r=0.
 \tag{GS14}
\]

The target is positive at a real place and hence anisotropic over F; the same holds at all real places here. Thus its global Witt index is r=0 and the lower-tower vanishing hypothesis is empty. All automorphic representations of these anisotropic groups are cuspidal.

This is the \textbf{anisotropic convergent case within Theorem11.4's dimension window}, not the geometric second-term range. GQT Proposition11.1(ii), printed p.53, and Theorem11.4, printed pp.54--55, apply exactly in this case.

At a split finite place the actual small theta lift is nonzero by Mínguez, including the ramified G/T/Q inputs. At a nonsplit finite place the input is spherical and the relative extension is unramified; Liu's Theorem6.2 applies to the identity Hermitian lattices. The absolute-codifferent rescaling of the skew form and additive character is the one already checked in LIFT14 and DO022, so it preserves the actual lattice vector and the composite symplectic character. No classification for a split place is substituted at a nonsplit place. The actual basic vector at every remaining place is nonzero in the quotient by LIFT16; this includes places over29 with their actual absolute ramification. For K031/K032,29 is in fact relatively split, since it splits in Q(zeta28), and adjoining the real fields preserves that splitting. Dyadic places are also relatively split. The odd nonsplit places use Liu's theorem over their actual local fields.

At a real place, LIFT25 gives the exact selected summand and a nonzero compact contraction on the first highest line. More generally the matching source/target vectors give the local doubling test

\[
 \int_{U(b)}|\langle\rho_1(g)v,v\rangle|^2\,dg
               =\operatorname{vol}(U(b))/\dim\rho_1>0.
 \tag{GS15}
\]

The real standard gamma factors are finite and nonzero at s\_*\textgreater0: the compact unitary representation has bounded parameter and the character twist is angular. This supplies the real local test directly; it does not invoke the unresolved general real converse in GQT11.6.

At finite places GQT Proposition11.6(i) gives existence of a nonzero normalized test on the actual Siegel--Weil image. Therefore all local theta lifts and all required local tests are nonzero. GQT11.4(ii) gives holomorphy of the completed standard factor at s\_*. Its partial factor(GS10) is nonzero. The finitely many missing standard local factors have no zeros, and the real gamma factors have no zeros; holomorphy of the completed factor therefore excludes a pole in any missing factor and makes its value nonzero.

GQT11.4 now proves that the global second theta map is nonzero on \textbf{every actual irreducible embedding} of pi\_1. The Euler product and local tests depend on the representation, while its positive invariant inner product is the one inherited from that embedding. Replacing an embedding by a nonzero vector in its multiplicity space changes the norm by its positive squared length and does not change the argument.

At this point only existence for an irreducible embedding has been proved. The following factorization step is what makes that statement useful for the fixed raw vector.

\subsection{5. Global factorization and the automorphic multiplicity map}\label{global-factorization-and-the-automorphic-multiplicity-map}

We state the precise algebraic argument to avoid identifying a representation with a preferred automorphic copy.

Let pi be an irreducible unitary cuspidal representation in a source automorphic space, M\_pi its Hilbert multiplicity space, and let S be the algebraic Schwartz space, with K-finite vectors at infinity. Assume all local small theta lifts sigma\_v are nonzero and that the global theta map is nonzero on every embedding belonging to M\_pi. Fix nonzero actual local quotient maps in the bilinear convention

\[
 q_v:S_v\otimes\pi_v\longrightarrow\sigma_v.
 \tag{GS16}
\]

At almost every place normalize q\_v to send the basic tensor to the unit spherical vector. Let q be their restricted tensor product. Then the actual global theta map has the form

\[
 \Theta_\pi=A_\pi\otimes q:
 M_\pi\otimes(S\otimes\pi)
            \longrightarrow M_\sigma\otimes\sigma,
 \qquad \sigma=\bigotimes'_v\sigma_v,
 \tag{GS17}
\]

where A\_pi is \textbf{injective}.

Here is a proof. The target automorphic Hilbert space is a discrete unitary sum of irreducibles; fixed finite level and real type give finite-dimensional subspaces. Projecting the global theta map to an irreducible target constituent gives at each place a nonzero local balanced Hom. Local Howe duality forces that constituent to have local representation sigma\_v and makes the local Hom one-dimensional. At finite nonsplit places we use Gan--Takeda, Theorem1.2; at split places Mínguez; at compact real places this is the actual multiplicity-free Cauchy decomposition. Thus all target projections have the same global representation sigma.

The tensor product of the local one-dimensional Hom spaces is also one-dimensional. To see this without an infinite tensor ambiguity, fix the basic vectors outside a finite set. On that finite set, successive local Hom factorization makes every coefficient map a multiple of the product q\_v. Enlarging the set preserves the same multiple because the basic values were normalized to one. Every algebraic input is in one such finite tensor subspace. This gives (GS17), with all remaining coefficients contained in A\_pi.

If A\_pi killed a nonzero m in M\_pi, the embedding v -\textgreater{} m tensor v would have zero theta lift for every Schwartz vector. This contradicts the assumed nonvanishing on that actual embedding. Hence A\_pi is injective. No multiplicity-one assertion for the automorphic spectrum is used.

There is also a useful metric refinement. Give the source and target representation factors their chosen invariant unitary metrics and give the multiplicity spaces the actual quotient metrics. The global doubling identity, applied to any unit multiplicity vector, has the same representation matrix coefficient and the same local factors. Polarization in the multiplicity vectors therefore gives

\[
                    A_\pi^*A_\pi=c_\pi\,1_{M_\pi},
                           \qquad c_\pi>0.
 \tag{GS18}
\]

This also follows by applying the norm identity to every nonzero linear combination of two embeddings. The scalar is retained; it is not set to one by a choice of local quotient model.

The argument only uses irreducible quotients of the local big theta module. It does not assume that the big theta module is itself irreducible or that an arbitrary invariant form on it is positive. Automorphic unitarity provides the semisimple target projections.

For our first edge, DO036 supplies the hypothesis for every pi\_0 embedding with the initial prescribed type. For the second edge, Section4 supplies it for every embedding of its actual descendant. We obtain injective maps A\_(0,pi0) and A\_(1,pi1).

Finally, distinct initial irreducible representations cannot acquire the same second global descendant: at every place the local Howe correspondence is injective on its nonzero domain. Applying that injectivity on the second and then first edge would identify all local components of the two initial representations, hence the initial restricted tensor products themselves. Their global multiplicity copies are already handled by A\_1 A\_0.

\subsection{6. The actual local coupled tensors are injective on the input types}\label{the-actual-local-coupled-tensors-are-injective-on-the-input-types}

For a constituent pi\_0, let W\_(0,v) be the part of its local space occupied by R017's initial type. The projection of F\_0 to pi\_0 stays in the tensor product of these subspaces: spectral projections commute with all local actions and compact projectors.

At G this is the finite Steinberg isotypic subspace intersected with the prescribed inverse-generic subspace. At T it is the nontrivial projective-line character space. At Q and at all other finite places it is the maximal compact fixed line. At infinity it is the selected highest line. For every place define

\[
 C_v(w)=q_{1,v}\bigl((1\otimes q_{0,v})(\Psi_v\otimes w)\bigr).
 \tag{GS19}
\]

For a sum in Psi\_v this means the corresponding finite sum of compositions, with the literal coefficients. The checked local results give the following exact detection identities.

\subsubsection{G places}\label{g-places}

LIFT22(CG16) gives, in the nested actual induced quotient models,

\[
                   [C_v(w)(1)](1)=a_q v_{J^\sharp}\,w,
 \qquad a_q=|U_a(\mathbb F_q)|/|GL_a(\mathbb F_q)|.
 \tag{GS20}
\]

The second scalar is the exact Haar volume of the actual sharp middle group. It is nonzero. Only the forced Y class contributes. This identity holds on the whole initial generic type, so it remains injective on the finite Steinberg subspace. In particular no erroneous full generic projector on the joining root is introduced.

\subsubsection{T places}\label{t-places}

LIFT24 proves that the actual first output of the raw T factor lies on the required one-dimensional projective-line type of pi\_1. DO036 proves it is nonzero for the initial type. For every vector w\_1 on that line, the fixed second raw incidence factor satisfies

\[
 q_{1,v}(\phi_{1,v}\otimes w_1)(k_*^{-1})
  =\frac{L((d+1)/2,\sigma_{1,v}\eta_{1,v}^{-1})}{q^b-1}\,w_1.
 \tag{GS21}
\]

Here pi\_1=sigma\_1 times xi in LIFT24's exact local induction and eta\_j is the local component of chi\_j. The displayed L-factor is finite and nonzero. Thus C\_v is injective on the initial T line. The first local scalar is not replaced by one or by its compact leading term; its nonzero GJ/tail value remains part of C\_v.

\subsubsection{Q places}\label{q-places}

DO037 proves for the \textbf{complete local Y sum} that evaluation at the actual unit-transported target graph and then at the first identity has value

\[
             \ell_v(C_v(w))=\Lambda_{Q,v}w,
 \tag{GS22}
\]

where, with r=a+1, t=d-1, theta=eta\_0 eta\_1 inverse and kappa\_n=product\_(i=1)\^{}n(1-q\^{}(-i)),

\[
 \Lambda_{Q,v}=q^{-d}G(\chi,\psi_0)\chi(-1)^{-1}
       q^{-t(a+1)}\frac{\kappa_{a+1}\kappa_t}{\kappa_b}
        \prod_{i=1}^{t}(1-\theta(\varpi)q^{-i})^{-1}\ne0.
 \tag{GS23}
\]

The source input here is its full maximal compact fixed line, exactly as in R017. The graph has determinant one after the actual piT-unit transport. DO037 retains the cofactor sign, Gauss sum, all Haar constants and the initial compact restriction. Its theorem is a whole-Schwartz small-quotient identity, not merely a finite-array Fourier calculation. It proves injectivity of C\_v on that line.

\subsubsection{Remaining finite places}\label{remaining-finite-places}

Both vectors are the actual selfdual lattice vacua. The fixed-vacuum theorem LIFT16 gives nonzero image of the spherical line on each edge, with the actual additive character and local fields. Thus C\_v is nonzero on a one-dimensional space. We may normalize its quotient models so that the basic vectors have image one; the inverse normalization is then retained in A\_j. This does not change the actual global operator. In particular the absolute ramification at29 does not introduce a principal29 level or a new finite translation.

\subsubsection{Real places}\label{real-places}

LIFT25 proves that the first highest output is exactly the source line paired by the second primitive highest polynomial. In its specified auxiliary unitary quotient models,

\[
 C_\infty(h_0)
       =\frac{A_0A_1}{\sqrt{d_0d_1}}h_2\ne0,
 \tag{GS24}
\]

where A\_j is the actual selected Schwartz norm and d\_j the exact dimension of the selected source representation, taking the tensor over real places. The literal two compact coefficient integrations instead carry 1/(d\_0 d\_1), as proved there. These are different descriptions of the same map in different output metrics; neither is silently substituted for the other. The actual Gaussian, creation, hook and period factors of LIFT25 remain present.

Equations(GS20--24) give nonzero left inverses up to their displayed scalars on every W\_(0,v). Therefore their finite tensor product is injective. Outside finitely many places these are normalized basic lines. Consequently the restricted tensor map

\[
                    C_{\pi_0}=\bigotimes'_v C_v
 \tag{GS25}
\]

is injective on the actual initial type space W\_(0,pi0). There is no estimate asserting that the left inverses have norm one or are integral. Their scalar and evaluation costs are not discarded.

\subsection{7. Proof of fixed global nonvanishing}\label{proof-of-fixed-global-nonvanishing}

Decompose the actual finite-level, fixed-real-type input into its finitely many irreducible isotypic components,

\[
                  F_0=\sum_{\pi_0}F_{0,\pi_0},\qquad
 F_{0,\pi_0}\in M_{0,\pi_0}\otimes W_{0,\pi_0}.
 \tag{GS26}
\]

For each component use the two factorizations(GS17). The complete CRT tensor identity(GS7), not a choice of a nonzero summand, then gives the actual coupled output component

\[
 (\mathcal BF_0)_{\pi_2}
   =\bigl((A_{1,\pi_1}A_{0,\pi_0})\otimes C_{\pi_0}\bigr)
                                                    F_{0,\pi_0}.
 \tag{GS27}
\]

Both multiplicity maps are injective, and C\_pi0 is injective on the indicated type space. Thus a nonzero input component has a nonzero output. Distinct pi\_0 components have inequivalent pi\_2 descendants by local Howe injectivity and are orthogonal in the target automorphic space. They cannot cancel. Since F\_0 is nonzero, at least one component is nonzero, proving(GS5).

This proof permits arbitrary finite linear combinations in the initial type space and arbitrary multiplicity vectors. It does not replace a vector by a favorable irreducible projection before declaring the original sum nonzero. It also does not claim that a fixed highest second kernel preserves nonhighest real vectors; those vectors are outside the stated input space and LIFT25 explicitly records that obstruction.

\subsection{8. Exact metric identity and costs still retained}\label{exact-metric-identity-and-costs-still-retained}

With the chosen invariant unitary quotient models, write

\[
 A_{j,\pi_j}^*A_{j,\pi_j}
                =c_{j,\pi_j}\,1,\qquad c_{j,\pi_j}>0.
 \tag{GS28}
\]

Equation(GS27) and orthogonality give the exact norm identity

\[
 \boxed{\quad
 \|\mathcal BF_0\|^2=
 \sum_{\pi_0}c_{0,\pi_0}c_{1,\pi_1}
   \|(1\otimes C_{\pi_0})F_{0,\pi_0}\|^2>0.\quad}
 \tag{GS29}
\]

The constants are actual global theta norms, not arbitrary pins. The first is determined by DO036's convergent norm identity with its actual quotient mass. For the second, in Tamagawa measures, GQT gives the following exact expression. Choose decomposable local tensors x\_v=phi\_v tensor v\_v with q\_v(x\_v) nonzero, normalize q\_v/basic metrics as above, and set

\[
 b_v=\frac{Z_v^*(s_0,\delta(\phi_v\otimes\overline{\phi_v}),
                                     v_v,v_v)}
                  {\|q_v(x_v)\|^2}.
 \tag{GS30}
\]

Then

\[
                  c_{1,\pi_1}
      =2L(s_*,\pi_1\times\chi_1)\prod_v b_v>0.
 \tag{GS31}
\]

The factor on the right is expressed in the actual GQT convention through DO028, hence has the standard unramified factor(GS10). In(GS30) the local input and oscillator are transported by g -\textgreater{} g\^{}(-t) and conjugation exactly as in DO028. Thus its printed GQT matrix-coefficient bar is retained; the notation does not replace the actual bilinear input by an unrelated Hermitian lift. At almost all places b\_v=1/d\_v(s\_0,chi\_1); the product of these factors converges and is nonzero by GQT(11.3). At the finitely many remaining places b\_v is the actual normalized doubling form relative to the specified local quotient metric. Equation(GS17) and the Rallis identity imply that this ratio is independent of a nonzero detecting tensor. One can see this by fixing all other local tensors with nonzero quotient and polarizing in x\_v: the global norm factors through the local irreducible unitary quotient. This also proves that the normalized local form is a nonzero scalar multiple of its positive quotient metric in this globalizable case. Its scalar is not assumed positive separately; the complete real quantity in(GS31) is positive. This is the needed radical statement, obtained after global factorization rather than assumed from local occurrence.

For source and target measures dg\_j=h\_j dg\_j\^{}Tam with h\_j\textgreater0, the relative multiplicity metric constant on edge j changes by h\_j h\_(j+1). Indeed the theta integral changes by h\_j, its target squared norm by h\_(j+1), and its source multiplicity squared norm by h\_j. Thus the formula transports to the actual probability-real, maximal-finite-compact measures without setting any quotient mass or global comparison scalar to one.

Equations(GS20--24),(GS29--31) retain all finite Gauss/projector volumes, T local factors, real dimension/hook/Gaussian factors and global L-values. They do not evaluate the finite number of initial automorphic constituents or the constants c\_j numerically. The nonzero global L(1,upsilon) factor is not assigned a uniform lower bound. No 11-adic valuation is attached to an unspecified real norm.

Most importantly, a nonzero complex automorphic output does not identify its integral image lattice. The actual algebraic global quotient scalars, the saturation and content of the prior-image matrix, the selected terminal branching projection, the TopForm row and the MainX terminal-mass equation remain separate. The unchanged elliptic Fields.Kq/VacIn constructor is still not a constructor for the primitive mixed host. Accordingly(GS5) is not a completed DP44 comparison or a shorter-chain theorem.

\subsection{9. Sources and review status}\label{sources-and-review-status}

Published inputs used at their stated hypotheses:

\begin{itemize}
\tightlist
\item
  Gan, W. T., Qiu, Y., \& Takeda, S. (2014). The regularized Siegel--Weil formula (the second term identity) and the Rallis inner product formula. \emph{Inventiones Mathematicae, 198}, 739--831. \url{https://arxiv.org/abs/1207.4709} . Sections11.2--11.8, especially Proposition11.1(ii), Theorem11.4, Proposition11.6 and equation (11.3), printed pp.51--56. Local primary copy: \nolinkurl{work/agents/lift_refs/gqt2014.pdf} and \nolinkurl{.txt}.
\item
  Gan, W. T., \& Takeda, S. (2016). A proof of the Howe duality conjecture. \emph{Journal of the American Mathematical Society, 29}, 473--493. \url{https://arxiv.org/abs/1407.1995} . Theorem1.2 and the exact definition of the maximal semisimple quotient, pp.1--3 of arXiv v4, were read directly. This covers all residual characteristics for nonsplit unitary pairs. The split case is expressly distinguished there and supplied by Mínguez.
\item
  Mínguez, A. (2008). Correspondance de Howe explicite: paires duales de type II. \emph{Annales Scientifiques de l'École Normale Supérieure, 41}(5), 717--741. Theorem1 and equation(1.4), carried from the direct readings in LIFT14/19/21/24.
\item
  Liu, Y. (2022). Theta correspondence for almost unramified representations of unitary groups. \emph{Journal of Number Theory, 230}, 196--224. arXiv:2105.14223v1, Theorem6.2 and Proposition3.4, used in the exact spherical plus-tower scope already checked in LIFT14/16.
\item
  DO028 for the complete Hom/tensor, conjugation and standard factor dictionary; DO032 for the literal CRT/representative identity; DO036 for the first actual global map; DO037, LIFT22/24/25 for the exact actual local coupled quotients; R017 and K031/032 for the initial function and actual host.
\end{itemize}

The global multiplicity argument, the fixed coupled assembly and the exact norm decomposition are proved above. The literature theorems and carried local findings are not counted again as new results. No Lean build or automated proof of the global theorem is claimed. Whole-proof review events are recorded only after the finished text is read in full.

\clearpage

\section{AM-K037 --- an actual mixed CM arithmetic class-row adapter}\label{am-k037-an-actual-mixed-cm-arithmetic-class-row-adapter}

Author: GPT-6 Astra Pro, ALP\_MACHINE /root/kummer\_extension. 22 November 2026 UTC; inherited reasoning effort is not separately exposed.

This finding constructs the arithmetic maps of the actual mixed family from its geometric reduced-theta kernels. It then proves their exact comparison with a canonical finite sum of genuine theta lifts. The coefficient spaces, their Weyl lattices, the class stabilizers, and the geometric kernels are fixed before this comparison. No lattice is fitted to analytic values.

In particular, the separate ``identify an externally specified new-host Kval map'' obligation in K035 is unnecessary for this constructed replacement family. The map below has the literal subD/subY formula of Kval. It need not be an instance of the false old Fields.Kq constructor. The remaining quantitative question concerns the image of this actual map on its chosen lattice, its successive images, and its norms.

The accepted coefficient-module results are used as theorems. In addition to the already checked K035--036 and DO025--027, the proof uses the source's Weyl GL-stability and normalized Weyl/dual-Weyl pairing. An elementary localization argument below extends that integral stability from generic invertible matrices to all matrices. This supplies a stronger statement than mere characteristic-zero membership: the selected universal substitution itself preserves the Weyl lattice.

\subsection{1. Fixed data and the exact source interfaces}\label{fixed-data-and-the-exact-source-interfaces}

Use either actual polarized host of K030/K032: \[
 (K_h,A_h)=(\mathbb Q(\zeta_{28})B,J^{14})\quad(h=1),
 \qquad
 (\mathbb Q(\zeta_{28})B(\sqrt5,\sqrt{13}),J^{56})\quad(h=4).
 \tag{CR1}
\] Keep its actual CM type, trace pairing, source additive-character convention, theta line, finite vacuum and all chosen finite operators. For one edge let \(a=N_j<b=N_{j+1}\) be even and use the same prescribed character \(\mu_{j+1}\) on the two members of the pair. At each CM label \(\tau\), put \[
\begin{aligned}
 \lambda_\tau&=\operatorname{lamTrue}_{j,\tau},&
 \lambda'_\tau&=\operatorname{lamTrue}_{j+1,\tau},\\
 s_{S,\tau}&=e_{j,\tau}-b/2,&
 s_{D,\tau}&=e_{j,\tau}-a/2,\\
 \eta_\tau&=\lambda_\tau-\lambda_{\tau,a-1}\mathbf1_a,&
 q_\tau&=\lambda_{\tau,a-1}+s_{S,\tau},\\
 \nu_\tau&=\eta_\tau+q_\tau\mathbf1_a.&
\end{aligned}
\tag{CR2}
\] The source weight identities, as already checked in DO025 and R011, give \[
 q_\tau\ge1,\qquad
 \lambda'_{\tau,b-1}=-s_{D,\tau},\qquad
 \eta'_\tau=(\nu_\tau,0^{b-a}).
 \tag{CR3}
\] Write \(D_\tau=|\nu_\tau|\). These are the actual selected partitions; they are not substituted by more convenient highest weights.

The source locators used in the construction are as follows.

\begin{longtable}[]{@{}
  >{\raggedright\arraybackslash}p{(\columnwidth - 2\tabcolsep) * \real{0.5000}}
  >{\raggedright\arraybackslash}p{(\columnwidth - 2\tabcolsep) * \real{0.5000}}@{}}
\toprule\noalign{}
\begin{minipage}[b]{\linewidth}\raggedright
Recovered source file
\end{minipage} & \begin{minipage}[b]{\linewidth}\raggedright
Exact use
\end{minipage} \\
\midrule\noalign{}
\endhead
\bottomrule\noalign{}
\endlastfoot
\nolinkurl{Formalism/GradedWeylDivisibility/StatementA.lean,} 26--38, 68--80 & Column equations for dualWeyl, rowDiv, and the integral Weyl lattice generated from hw \\
\nolinkurl{Formalism/LocalRingCardinality/DataCoeff.lean,} 53--84, 115--143 & nablaSq, delSq, square/round actions, normalized Fischer metric and the Weyl/dual-Weyl pairing \\
\nolinkurl{Formalism/LocalRingCardinality/RankBoundMonotoneInDimension.lean,} 902--909 & Proven factSqStable, invoking the accepted WeylGLStable theorem over every commutative base ring \\
\nolinkurl{Formalism/LocalToGlobal/Divisibility/CoefMod.lean,} 60--80, 178--226, 309--330 & VV, its polynomial action, the fixed standard/transferred lattices, rhoV, rhoVr and delAtr \\
\nolinkurl{Formalism/LocalToGlobal/Columns/CoefModB.lean,} 55--90 & Literal symmetric bilinear pairBK and its rhoVr/rhoV invariance \\
\nolinkurl{Formalism/LocalToGlobal/KernelLattices/ThDistC.lean,} 19--43, 61--65, 137--148 & subD, jetLin, jetMap, retrVV, JetInVV, ordinary derivatives and the distribution map \\
\nolinkurl{Development/AbelianVarietyTheory/Definitions/LatticeFrames/EvalAtOneExtremalMul.lean,} 17--38, 90--105 & Vn, evalAt, mvZ, subY, jetL and subY\_eq\_subD \\
\nolinkurl{Development/AbelianVarietyTheory/Definitions/LatticeFrames/OnFramesRdApply.lean,} 96--121 & Real, Lev, Core and the exact Kval equation \\
\nolinkurl{Formalism/LocalToGlobal/Columns/Record.lean,} 1285--1300 & The source integral column uses delAtr, not a replaced dual-Weyl lattice \\
\end{longtable}

Take one number field \(M\), stable under complex conjugation, containing the host, the fixed finite character values, all fixed chart constants, and the all-order geometric jet field supplied by K036 for these ranks and allowed denominators. This is a finite extension chosen once. The choice has no effective degree or discriminant estimate here.

For every actual finite class pair \((c,d)\), the independently defined geometric kernel is the finite Mumford average and translation sum of DO026, with its actual character and central phases. Its normalized functional will be denoted \[
 J_{dc}(P)=\epsilon_{dc}\,
             [P(cD_\tau\partial_{\tau,rs})F_{dc}]_0\in M .
 \tag{CR4}
\] The coefficients, phases and translations in \(F_{dc}\) are fixed by the chosen finite Schwartz datum, before any analytic theta value is considered. Any already specified chart factors of DO027 are included in \(\epsilon_{dc}\). Formula (CR4) uses ordinary derivatives. No division by a multiindex factorial is inserted.

The finite input and row levels are required to have their actual covariance \[
 \omega_f(k,k')\phi_{j,f}
       =\chi_j(k)^{-1}\chi_{j+1}(k')\phi_{j,f}.
 \tag{CR4a}
\] This is not inferred from weights at infinity. DO026 supplies explicit finite stabilizers for every fixed geometric finite span, with the necessary determinant-character kernels; its finite refinement index is retained. When a more specifically prescribed raw finite family and level are used, their checked actual covariance is the input to the same theorem. Nothing here identifies an arbitrary pinned level with that stabilizer. Several chart blocks may have different fixed finite vectors, each satisfying its displayed source and target covariance.

The exact source-creation comparison proved in K035 is \[
 \vartheta_{\operatorname{pureTensorO}(\phi_f,P)}(g_c,h_d)
   =C^{14hab}J_{dc}\!\left(P(s_\tau Z_{\tau,rs})\right),
 \quad
 s_\tau=-\frac{\Omega_\tau}
                   {cD_\tau\sqrt{\pi q^{\rm met}_\tau}},
 \quad q^{\rm met}_\tau=-2i\,\tau(\xi_h)>0.
 \tag{CR5}
\] The notation \(q^{\rm met}_\tau\) distinguishes this positive metric coefficient from the integer \(q_\tau\) in (CR2). The common period \(C=\theta_0c\), the periods \(\Omega_\tau\), and the algebraic coordinate generators \(cD_\tau\) are exactly those of K035. In particular the scalar in (CR5) is not selected separately at different derivative orders.

\subsection{2. The universal substitution, including its integral Weyl image}\label{the-universal-substitution-including-its-integral-weyl-image}

For a commutative ring \(R\), let \(\nabla_\eta(R)\) and \(\Delta_\eta(R)\) be the source's dual-Weyl and Weyl polynomial modules. Their highest polynomial is the product of leading minors \[
 p_\eta(Y)=\prod_{r=1}^a
             \det(Y_{[1,r],[1,r]})^{\eta_r-\eta_{r+1}},
 \qquad \eta_{a+1}=0.
 \tag{CR6}
\] Over a characteristic-zero field these two modules have the same span. Their integral lattices are retained separately.

Use \(a\)-by-\(a\) variables \(U\), \(b\)-by-\(b\) variables \(Y\) and \(b\)-by-\(a\) variables \(Z\), at each \(\tau\). Define \[
 {\cal Q}(P)(Y,Z)=
 \left[
     \left(\prod_\tau\det(U_\tau)^{q_\tau}\right)P(U)
 \right]_{U_\tau=Z_\tau^{\mathsf T}Y_{\tau,[1,b],[1,a]}}.
 \tag{CR7}
\] This is exactly subD after the source determinant insertion. Its evaluation at \(Y=y\) is exactly subY(y) applied to that insertion. All coefficients of the substitution are integers.

\subsubsection{Theorem CR8: the actual selected substitution has integral image}\label{theorem-cr8-the-actual-selected-substitution-has-integral-image}

For every commutative ring \(R\), \[
\begin{split}
 {\cal Q}\left(\bigotimes_\tau\nabla_{\eta_\tau}(R)\right)
   &\subset
       \left(\bigotimes_\tau\nabla_{\eta'_\tau}(R)\right)[Z],\\
 {\cal Q}\left(\bigotimes_\tau\Delta_{\eta_\tau}(R)\right)
   &\subset
       \left(\bigotimes_\tau\Delta_{\eta'_\tau}(R)\right)[Z].
\end{split}
\tag{CR8}
\] Here membership means that every \(Z\)-coefficient belongs to the displayed module in the \(Y\)-variables. The second assertion retains the source's chosen Weyl lattice; it is stronger than JetInVV over a field.

\textbf{Proof of the first assertion.} It is enough to work at one label. Multiplication by \(\det(U)^q\) changes the column weight from \(\eta\) to \(\nu=\eta+q\mathbf1\), and preserves invariance under adding an earlier column to a later one. Replacing \(U\) by \(Z^{\mathsf T}Y_{[,1:a]}\) preserves these same column identities in \(Y\). The columns beyond \(a\) do not occur. Thus the polynomial has column weight \((\nu,0)\) and satisfies every defining upper-column invariance of \(\nabla_{\eta'}(R[Z])\). Comparing \(Z\)-coefficients proves the assertion, directly over \(R\).

\textbf{Integral extension of row stability.} First work over \(\mathbb Z\) and fix a target Weyl lattice \(\Delta\) in its finite-degree ambient polynomial module \(A\). Let \(G\) be a generic \(b\)-by-\(b\) matrix. The accepted factSqStable theorem, with final weight zero, implies \[
                  \operatorname{rowAct}(G)\Delta
       \subset\Delta\otimes_{\mathbb Z}\mathbb Z[G,\det(G)^{-1}].
 \tag{CR9}
\] The harmless transpose in sqAct is removed by taking its group argument to be \(G^{\mathsf T}\). The variable determinant has a leading monomial with coefficient \(1\) or \(-1\), for any monomial order. Consequently multiplication by \(\det(G)\) is injective on \((A/\Delta)[G]\), even when \(A/\Delta\) has torsion: the leading coefficient of a nonzero polynomial in this module is multiplied by a unit, and cannot disappear.

The left side of (CR9) already has coefficients in \(A[G]\). Its class in \((A/\Delta)[G]\) vanishes after inverting \(\det(G)\), so some power of \(\det(G)\) kills it. The preceding injectivity forces that class to be zero before localization. Thus \[
                 \operatorname{rowAct}(G)\Delta\subset\Delta[G].
 \tag{CR10}
\] Polynomial specialization proves stability under every row matrix, not only invertible ones. This argument uses flat base change for the polynomial and localized polynomial rings; it does not assume that the integral Weyl lattice is saturated in its ambient module.

\textbf{Proof of the Weyl assertion.} Multiplication by \(\det(U)^q\) takes \(p_\eta\) exactly to \(p_\nu\). It commutes with every off-diagonal row divided power, because an elementary row addition does not change the determinant. The diagonal row divided powers act on a row-homogeneous polynomial of row weight \(\mu\) by the integer \(\binom{\mu_i}{r}\). Every divided-power word applied to the highest polynomial remains row homogeneous (or becomes zero). Diagonal factors in such a word can therefore be removed at the cost of integer scalar factors. The defining divided-power generators show \[
                  \det(U)^q\Delta_\eta(\mathbb Z)
                                  \subset\Delta_\nu(\mathbb Z).
 \tag{CR11}
\] Embedding the first \(a\) rows and columns into the target square matrix takes \(p_\nu\) to \(p_{(\nu,0)}\). The same row divided-power words on those first \(a\) rows show that it takes the entire module in (CR11) into the target Weyl lattice.

Now apply (CR10), and specialize the first \(a\) rows of its generic matrix \(G\) to \(Z^{\mathsf T}\); the remaining rows can be zero. The embedded polynomial does not use those remaining rows of \(GY\). Its image is exactly (CR7). Therefore all its coefficients belong to the target Weyl lattice over \(\mathbb Z\). All Weyl generators are defined by integral polynomial formulas. Base change sends their integral span to the defining span over any \(R\). This proves the result over \(R\), and tensoring the separate labels proves (CR8). \(\square\)

For any \(M\)-linear polynomial functional \(J\), define \[
                    \Phi_J(P)(Y)=J_Z({\cal Q}(P)(Y,Z)).
 \tag{CR12}
\] Theorem CR8 proves the exact JetInVV premise of ThDistC for these weights. Its retrVV projection is the identity on this image, whatever complementary subspace was used to define that auxiliary projection. This is a proved membership statement, not a choice of projection which fits the desired output.

In standard integral coordinates, if all degree-profile values \(J(Z^\alpha)\) belong to a fractional ideal \(\mathfrak a\), then \[
                       \Phi_J(\Delta_\eta)
                                   \subset\mathfrak a\Delta_{\eta'}.
 \tag{CR13}
\] There is no hook-product or Weyl/dual-Weyl index denominator in (CR13). Such indices remain relevant to other projections and duality operations; the assertion concerns this particular universal substitution and this particular input Weyl lattice.

\subsection{3. The canonical Cauchy map and its exact metric normalization}\label{the-canonical-cauchy-map-and-its-exact-metric-normalization}

Let \(V_j\) denote the source polynomial space VV with action \(\rho_j^\vee=\rhoVr\), and let \(W_j\) be the same polynomial space with action \(\rho_j=\rhoV\). Use the source's nondegenerate symmetric bilinear form \[
 B_j(P,v)=\operatorname{pairBK}_{\lambda_j}(P,v)
         =\frac{\operatorname{Fischer}_{\rm bil}(P,v)}
                        {\prod_\tau H_{\eta_{j,\tau}}}.
 \tag{CR14}
\] It identifies \(W_j\) with the algebraic dual of \(V_j\), and \[
 B_j(\rho_j^\vee(g)P,\rho_j(g)v)=B_j(P,v).
 \tag{CR15}
\] This is exactly the source convention; no Hermitian conjugation is inserted into the arithmetic map. Complex conjugation gives the positive compact Hermitian form when norms are taken.

For an invertible frame \(y\), let \(v_y\in W_{j+1}\) be uniquely defined by \[
 B_{j+1}(Q,v_y)=
          \prod_\tau\det(y_\tau)^{-s_{D,\tau}}Q(y).
 \tag{CR16}
\] Such evaluation vectors span \(W_{j+1}\), because a polynomial which vanishes on all invertible matrices is zero. The target membership in CR8 therefore defines a unique bilinear map \[
 \iota_j:V_j\otimes W_{j+1}\longrightarrow M[Z],\qquad
 \iota_j(P,v_y)=
        \prod_\tau\det(y_\tau)^{-s_{D,\tau}}{\cal Q}(P)(y,Z).
 \tag{CR17}
\] Equivalently, contract the target polynomial coefficients of \({\cal Q}(P)\) with \(v\) using \(B_{j+1}\). This definition fixes its normalization without choosing a projector or an arbitrary scalar in a multiplicity-one assertion.

Under the actual source oscillator convention, \[
 \omega_\infty(g,h)\operatorname{pureTensorO}(P)
  =\prod_\tau\det(g_\tau)^{s_{S,\tau}}
                \det(h_\tau)^{-s_{D,\tau}}\,
       \operatorname{pureTensorO}
            (P(h_\tau^{\mathsf T}Z_\tau g_\tau^{-\mathsf T})).
 \tag{CR18}
\] The map \(\iota_j\), followed by this creation map, intertwines \(\rho_j^\vee\boxtimes\rho_{j+1}\) with the actual selected oscillator summand.

\textbf{Verification of the two actions.} On the source, substituting \(Zg^{-\mathsf T}\) into (CR7) changes its matrix argument to \(g^{-1}Z^{\mathsf T}Y_{[,1:a]}\). The determinant insertion contributes \(\det(g)^{-q_\tau}\); multiplication by the vacuum character \(\det(g)^{s_{S,\tau}}\) leaves \(\det(g)^{-\lambda_{\tau,a-1}}\), precisely rhoVr. On the target, substituting \(h^{\mathsf T}Z\) changes \(y\) to \(hy\). The factor \(\det(h)^{-s_D}\) combines with \(\det(y)^{-s_D}\) to give \(\det(hy)^{-s_D}\). Equation (CR15) gives \(\rho_{j+1}(h)v_y=v_{hy}\). These identities prove (CR18) on the spanning vectors \(v_y\), and hence on the whole map.

The canonical map also has a fixed norm. Normalize the compact metrics on \(V_j,W_{j+1}\) by (CR14), so the extremal polynomial has norm one. With the ordinary Fischer norm on creation-label polynomials, \[
 \boxed{\|\iota_j(P,v)\|_{\rm F}^2
       =\left(\prod_\tau H_{\nu_\tau}\right)
                              \|P\|_j^2\|v\|_{j+1}^2}
 \tag{CR19}
\] for pure tensors; the corresponding tensor-product metric identity holds on their span.

To check the scalar, take extremal source \(P=p_\eta\) and the target frame \(y=1\). Its vector \(v_1\) is exactly \(p_{\eta'}\). Indeed evaluation at \(1\) vanishes on every row-weight space except the highest row weight \(\eta'\), whose line is spanned by \(p_{\eta'}\); that polynomial evaluates to one. Row weights are orthogonal in the Fischer form, and its normalized squared norm is one. This proves \(v_1=p_{\eta'}\) with no scalar. Formula (CR17) now gives the primitive leading-minor polynomial \(p_\nu(Z)\), with coefficient one on its leading diagonal monomial. Its squared Fischer norm is \(H_\nu\), by the checked R011 calculation. The selected Cauchy component is irreducible and has multiplicity one (DO025); invariance of the metrics then proves (CR19). This is a normalization check, using the accepted Cauchy decomposition rather than re-proving it.

The integral duality is also fixed: pairBK pairs the standard dual-Weyl lattice perfectly with the standard Weyl lattice in the opposite action. Thus CR8 and (CR17) take a source integral Weyl vector and a target integral dual-Weyl test vector to an integral polynomial. This statement names the two integral lattices explicitly; their characteristic-zero spans coincide. No use of a complex orthonormal basis below changes either arithmetic lattice.

\subsection{4. Arithmetic classes, unchanged integral lattices, and Kval}\label{arithmetic-classes-unchanged-integral-lattices-and-kval}

Fix the actual finite row levels and finite characters. For clarity first use one chart at each stage, with level \(U_j\), character \(\chi_j\), finite class set \[
 {\cal C}_j=G_j(F)\backslash G_j(\mathbb A_f)/U_j,
 \qquad
 \Gamma_c=G_j(F)\cap g_cU_jg_c^{-1}.
 \tag{CR20}
\] The groups are totally definite. Their fixed-rank class sets are finite by K036; adding the specified finite decorations refines them by finite sets. Stabilizers are finite. These are the actual arithmetic and adelic classes, with their actual lattices and decorations, not a class-number-one replacement.

For \(\gamma\in\Gamma_c\), put \(\kappa_\gamma=g_c^{-1}\gamma_fg_c\in U_j\). The arithmetic class-row space is \[
 {\cal H}_j(M)=
 \left\{(P_c)\in\bigoplus_{c\in{\cal C}_j}V_j:
     \rho_j^\vee(\gamma_\infty)P_c
                =\chi_j(\kappa_\gamma)P_c
                  \text{ for every }\gamma\in\Gamma_c\right\}.
 \tag{CR21}
\] This sign is fixed by the scalar automorphic realization in (CR26) below, whose right \(U_j\)-character is \(\chi_j\).

Let \(L_{j,c}\) be the literal delAtr lattice for the lattice Lsh attached to \(c\), including its source determinant twist and any fixed finite coefficient decoration. Define \[
 {\cal L}_j={\cal H}_j(M)\cap\bigoplus_c L_{j,c}.
 \tag{CR22}
\] This is the standard integral coefficient module used here. Its definition precedes the theta map. No averaging projector is used to redefine it. It is a full lattice in its \(M\)-space: the ambient lattices are full, the invariance equations are linear over \(M\), and clearing denominators in a basis of their solution space supplies a spanning set in (CR22). Locally over a DVR it is the saturated intersection with the specified subspace.

Define the actual arithmetic map by the finite formula \[
 \boxed{({\cal K}_jP)_d
           =\sum_{c\in{\cal C}_j}\frac1{|\Gamma_c|}
                              \Phi_{J_{dc}}(P_c).}
 \tag{CR23}
\] Everything on the right was fixed geometrically or algebraically in §§1--3. It is an \(M\)-linear map, without fitting to analytic values. Each summand belongs to the prescribed target polynomial space by CR8. The target stabilizer equations follow either directly from the rational and finite-level covariance of the same geometric kernels, or from the independent genuine-theta comparison proved in §5. Thus it maps \({\cal H}_j\) to \({\cal H}_{j+1}\).

For every frame \(y'\), formula (CR23) gives \[
\boxed{\begin{split}
 &\operatorname{val}_{j+1}({\cal K}_jP,d;y')\\
 &=\sum_c|\Gamma_c|^{-1}
    J_{dc}\!\left(
        \operatorname{subY}(y')
        \left[\prod_\tau\det(U_\tau)^{q_\tau}P_c(U)\right]\right)
                       \prod_\tau\det(y'_\tau)^{-s_{D,\tau}}.
\end{split}}
\tag{CR24}
\] The source frame value is \[
 \operatorname{val}_j(P,c;y)
   =\left[\prod_\tau\det(y_\tau)^{q_\tau}P_c(y)\right]
                         \prod_\tau\det(y_\tau)^{-s_{S,\tau}},
 \tag{CR25}
\] because \(q_\tau-s_{S,\tau}=\lambda_{\tau,a-1}\). Thus (CR24) is exactly the displayed source Kval formula, with \(n=q\), \(F_c=\det^qP_c\) and abelian₂ equal to the already fixed \(J_{dc}\). The equality holds for all rows and all frames, not only for one nonzero coefficient. It uses ordinary polynomial evaluation and the actual same-field jets.

For several finite chart labels \(\alpha,\gamma\), replace the source sum by \(\sum_\alpha\sum_{c\in{\cal C}_{j,\alpha}}\) and use the predetermined block kernel \(J_{\gamma d,\alpha c}\). The same proof applies blockwise. Equal source chart volumes give the common scalar in §5. If unequal volumes are used, their exact ratios must be included by the permitted finite input normalization of DO027, or retained as explicit factors in the analytic formula. They cannot be omitted.

\subsection{5. Comparison with a canonical sum of genuine theta lifts}\label{comparison-with-a-canonical-sum-of-genuine-theta-lifts}

Let \(v_j=\operatorname{vol}(U_j)\), with infinite compact Haar of volume one. For a row \(P\) and \(v\in W_j\), define the actual scalar automorphic function by \[
 f_{P,v}(g_ck_\infty k_f)
       =\chi_j(k_f)B_j(P_c,\rho_j(k_\infty)v).
 \tag{CR26}
\] Equation (CR21) makes this well defined: replacing the representative by \(\gamma\in\Gamma_c\) multiplies the finite factor by \(\chi_j(\kappa_\gamma)\) and the bilinear coefficient by its inverse. Conversely these are precisely its stabilizer conditions. Hence this is the canonical coefficient realization of the actual row, with its prescribed right level character.

Take any basis \(v_i\) of \(W_j\) and its pairBK-dual basis \(p_i\) of \(V_j\), so \(B_j(p_i,v_l)=\delta_{il}\). For \(v\in W_{j+1}\), define \[
 {\cal T}_j(P;v)(h)
  =\sum_i\int_{[G_j]}
       f_{P,v_i}(g)\,
       \vartheta_{\operatorname{pureTensorO}
                  (\phi_{j,f},\,\iota_j(p_i,v))}(g,h)\,dg .
 \tag{CR27}
\] This is a finite sum of genuine source-style bilinear theta lifts, with the already fixed actual finite Schwartz function. It is basis independent because \(\sum_i v_i\otimes p_i\) is the canonical coevaluation tensor. It has not been defined through \({\cal K}_j\). For chart blocks use their already fixed finite Schwartz functions and sum the corresponding integrals.

\subsubsection{Theorem CR28: exact analytification of the constructed arithmetic map}\label{theorem-cr28-exact-analytification-of-the-constructed-arithmetic-map}

For every row \(P\in{\cal H}_j(M)\), every \(v\in W_{j+1}\), and every target adele \(h\), \[
 \boxed{{\cal T}_j(P;v)(h)
      =v_j C^{14hab}
              \left(\prod_\tau s_\tau^{D_\tau}\right)
                         f_{{\cal K}_jP,v}(h).}
 \tag{CR28}
\] This is the mixed arithmetic adapter. Its per-place scale is the literal RealD.sc of K035, and its class map has the exact Kval formula (CR24).

\textbf{Proof.} Write \(g=g_ck\), \(h=h_dl\) on compact real pieces. The finite source level character of the kernel is \(\chi_j^{-1}\), so it cancels (CR26); the target finite character is retained. The real covariance (CR18) sends its archimedean vector to \(\iota_j(\rho_j^\vee(k)p_i,\rho_{j+1}(l)v)\). Consequently the sum over \(i\) in the integrand is contracted by \[
 \sum_i\rho_j(k)v_i\otimes\rho_j^\vee(k)p_i
                         =\sum_i v_i\otimes p_i.
 \tag{CR29}
\] Pairing its first factor with \(P_c\) gives exactly \(P_c\). Thus the sum of integrands is independent of \(k\), and equals the theta functional at \((g_c,h_d)\) applied to \(\iota_j(P_c,\rho_{j+1}(l)v)\).

This polynomial is homogeneous of degree \(D_\tau\) at every label. Formula (CR5) therefore contributes the single common factor \(C^{14hab}\prod_\tau s_\tau^{D_\tau}\). Definitions (CR12), (CR14) and (CR17) identify the remaining value with \[
 B_{j+1}(\Phi_{J_{dc}}P_c,\rho_{j+1}(l)v).
 \tag{CR30}
\] Finally the actual quotient integration formula of DO026 is \[
 \int_{[G_j]} F(g)\,dg
  =v_j\sum_c|\Gamma_c|^{-1}
                                  \int_{G_{j,\infty}}F(g_ck)\,dk .
 \tag{CR31}
\] Its infinite integral is one after (CR29). Substitution of (CR23) gives (CR28). Normal convergence of the fixed-degree Gaussian theta kernels on the compact quotient justifies the integrals and their finite sums, as in DO026. The scalar in (CR28) is nonzero. Genuine theta automorphy also proves the target equivariance asserted after (CR23), without changing its arithmetic definition. \(\square\)

There is no unrecorded factor \(1/\dim W_j\) in (CR28). It cancels because (CR27) uses the complete coevaluation tensor. A single fixed matrix-coefficient realization instead has the Schur factor \(1/\dim W_j\), exactly as in DO026. The complete-type construction has its own cost, computed next.

\subsection{6. Exact norm and lattice bookkeeping}\label{exact-norm-and-lattice-bookkeeping}

Put \(d_j=\dim_\mathbb C W_j\). With an orthonormal compact basis \(v_i\), its pairBK-dual \(p_i\) is also orthonormal under the corresponding dual metric. For every row \(P\), \[
 \boxed{\sum_i\|f_{P,v_i}\|_{L^2([G_j])}^2
          =v_j\sum_c|\Gamma_c|^{-1}\|P_c\|_j^2.}
 \tag{CR32}
\] Indeed the sum of squared matrix coefficients at each \(k\) is \(\|P_c\|^2\), by unitarity; then use (CR31). For one unit \(v\), irreducibility instead gives \[
 \|f_{P,v}\|_{L^2}^2
       =\frac{v_j}{d_j}\sum_c|\Gamma_c|^{-1}\|P_c\|_j^2 .
 \tag{CR33}
\] Equations (CR32)--(CR33) distinguish the Hilbert-Schmidt metric of the complete coefficient realization from a single coefficient embedding.

The norm of the actual archimedean coevaluation vector \[
 {\cal E}_j(v)=
       \sum_i v_i\otimes
          \operatorname{pureTensorO}
                    (\phi_{j,f},\iota_j(p_i,v))
 \tag{CR34}
\] is, with the literal unit-normalized creation operators, \[
 \boxed{\|{\cal E}_j(v)\|^2
     =d_j\left(\prod_\tau H_{\nu_\tau}\right)
             \|\phi_{j,f}\|_f^2
             \|\phi_{\infty,0}\|_\infty^2
             \|v\|_{j+1}^2.}
 \tag{CR35}
\] This follows from CR19 and orthonormality of the \(v_i\). The actual product-host vacuum has \[
             \|\phi_{\infty,0}\|_\infty^2
                              =\det(2Y_J)^{-7hab}.
 \tag{CR36}
\] There are \(14hab\) ordinary Jacobian factors, each with squared Gaussian norm \(\det(2Y_J)^{-1/2}\). K035's change to the literal normalized creation operator preserves the Fischer normalization. The finite norm in (CR35), all hook products, the source dimension, the level volume, and the period in (CR28) remain actual factors. The arithmetic class-row lattice is still (CR22), not the orthonormal lattice used only to compute this norm.

There is also an exact local formulation of the denominator question. At a finite place \(w\) of \(M\), choose actual local bases \(B_{c,\tau}\), \(B'_{d,\tau}\) of the source and target lattices Lsh. The source definitions give \[
 L_{j,c,w}=\rho_j^\vee(B_c)\Delta_{\eta,\mathcal O_w},
 \qquad
 L_{j+1,d,w}=\rho_{j+1}^\vee(B'_d)
                                      \Delta_{\eta',\mathcal O_w}.
 \tag{CR37}
\] All determinant twists in these equalities are retained. The polynomial covariance gives exactly \[
 (\rho_{j+1}^\vee(B'))^{-1}\Phi_J\rho_j^\vee(B)
                         =\Phi_{J^{B,B'}},
 \tag{CR38}
\] where \[
 J^{B,B'}(P)=
 \prod_\tau\det(B_\tau)^{s_{S,\tau}}
                \det(B'_\tau)^{-s_{D,\tau}}\,
 J\!\left(P((B'_\tau)^{\mathsf T}
                                  Z_\tau B_\tau^{-\mathsf T})\right).
 \tag{CR39}
\] For example, (CR39) follows by writing \(\rho^\vee(B)P(U)=\det(B)^{-\lambda_{\rm last}}
P(B^{-1}U)\), absorbing \(q-\lambda_{\rm last}=s_S\) into the determinant insertion, and evaluating the target at \(B'Y\). This also verifies every transpose and sign.

Let \(\mathfrak j_{dc,w}\) be the fractional ideal generated by the finite collection \[
       J_{dc}^{B_c,B'_d}(Z^\alpha),
       \qquad |\alpha_\tau|=D_\tau\quad\text{for every }\tau .
 \tag{CR40}
\] Theorem CR8 gives the concrete bound \[
 \Phi_{J_{dc}}(L_{j,c,w})
                         \subset\mathfrak j_{dc,w}L_{j+1,d,w}.
 \tag{CR41}
\] After the actual stabilizer contraction, \[
 {\cal K}_j({\cal L}_{j,w})
     \subset \pi_w^{\,b_{j,w}}{\cal L}_{j+1,w},
 \qquad
 b_{j,w}=\min_{c,d}
       \left(v_w(\mathfrak j_{dc,w})-v_w(|\Gamma_c|)\right).
 \tag{CR42}
\] Zero ideals have valuation \(+\infty\). The saturated intersection property in (CR22) justifies passing from the ambient coefficient lattices to the equivariant row lattice in (CR42). Empty spaces are treated in the usual way.

This is an inclusion, not an equality for the sharp content. Cancellation in the class sum can increase valuations. Restricted earlier images can increase them further. The actual sharp invariant is obtained from the matrix of \({\cal K}_j:{\cal L}_{j,w}\to{\cal H}_{j+1}(M_w)\) in local integral bases and, for later stages, from its restriction to the actual preceding image. No class projection, arbitrary choice of image lattice, or ambient nonzero vector proves that invariant. Formula (CR42) isolates it using the actual map rather than leaving Kval itself unspecified.

\subsection{7. Composition, terminal compatibility, and the remaining gate}\label{composition-terminal-compatibility-and-the-remaining-gate}

Repeating the construction along the prescribed finite even-rank chain defines actual arithmetic maps \({\cal K}_j\), all over one finite compositum field, on the unchanged class coefficient lattices. At every edge their analytic complete-type theta realizations satisfy (CR28). Iteration therefore retains the explicit product \[
 \prod_j
 \left[v_j C^{14hN_jN_{j+1}}
                  \prod_\tau s_{j,\tau}^{|\nu_{j,\tau}|}\right].
 \tag{CR43}
\] The final arithmetic column is in precisely the rhoVr polynomial convention used by the literal branching operator brC and bilinear pairBK. Its ordinary terminal restriction can be paired with the rhoV terminal row as in DO038. There is no additional transpose scalar. DO038 keeps the exact branch normalization and terminal stabilizer mass; those are not redefined in this construction.

There is a further finite-assembly distinction. A single ordinary ColV row does not automatically become the source's coupled evGSD row by pulling back along tilC: the latter has a prescribed, possibly nontrivial Q character. The actual finite Y/Q expansion can be assembled from the blockwise maps (CR23), with its fixed coefficients. If its actual finite-array covariance is verified, linearity of CR24 and CR28 proves the same arithmetic/analytic comparison for that coupled row, and DO038 supplies its terminal mass identity. No such covariance is inferred from an arbitrary nonzero ordinary theta output. The literal source colC has codomain evGSD.EqFun in \nolinkurl{FiniteSumNonzeroLargeDepth.lean1410–1418,} and PerInv1451--1454 expresses its cancellation, but transferring that particular coupled assembly requires its actual finite kernel hypotheses. The basic theorem here is the ordinary class-row adapter and its blockwise extension, not an unsupported assertion that all mixed terminal arrays already satisfy those hypotheses.

This result supplies an actual replacement arithmetic family and proves its geometric Kval equation and analytification. It does not prove that the old common-quadratic Fields record exists. It also does not infer nonvanishing of every later restricted image, optimal valuation gains, the total DP44 height cost, or an every-epsilon window from the existence of these maps. Those are now questions about the concrete maps (CR23), lattices (CR22), and normalizations (CR28), (CR35), and the specified terminal pairing.

In particular K033/K034's sharp selected first-kernel content cannot automatically be substituted for the content of a later class-contracted image. Conversely an additional ``unknown externally frozen Kval map'' is not a reason to withhold this actual construction: its maps have already been fixed geometrically and their comparison proved.

\subsection{Verification and primary references}\label{verification-and-primary-references}

The companion exact certificate tests the universal substitution on small nontrivial shapes, constructs the corresponding integral Weyl lattices from divided-power generators, checks actual coefficient membership, verifies the matrix covariance (CR39), and checks coevaluation and mass normalizations. These finite tests do not replace the general proofs or assert a Lean build. Author and independent review events are recorded separately.

The representation inputs are the accepted recovered source theorems at the locators above and the primary references already checked for DO025/R011: Fulton, W., \& Harris, J. (1991). Representation theory: A first course. Springer, Theorem 6.3, pp.77--78, Hook Length Formula 4.12, p.50, and Exercise 6.11(b), p.80; and Konno and Konno (2007), section 5.1, Fact 5.1, Lemma 5.2, formulas (5.4)--(5.5), and Lemma 5.3, pp.72--74. The geometric reduced-theta input is the actual algebraic construction and all-order field comparison proved in K035--036, using Bannai and Kobayashi's primary reduced-theta conventions. No general RH extension or unstated extension of an accepted coefficient theorem is used.

\clearpage

\section{AM-DO-038 --- the literal terminal period, its top realization, and exact stabilizer cancellation}\label{am-do-038-the-literal-terminal-period-its-top-realization-and-exact-stabilizer-cancellation}

Author: GPT-6 Astra Pro, ALP\_MACHINE \nolinkurl{/root/definite_oscillator}. 22 November 2026 UTC. Inherited reasoning effort is not separately exposed.

This finding constructs the top automorphic realization from the fixed arithmetic terminal pairing and proves its terminal class-mass equation. The arithmetic tail is the source's \nolinkurl{PtilFF}; it is not defined by dividing an analytic period by a desired scalar. The final arithmetic column can be supplied by the independent geometric mixed-CM operators and class-row construction. The identity applies uniformly to every such column, including zero columns.

The branching operator remains the literal \nolinkurl{brC}, including \nolinkurl{brConst}, \nolinkurl{kapBr}, and the power of the conjugate T-uniformizer. Its closed factorial evaluation is separate. No terminal nonvanishing is inferred from nonvanishing of the last theta lift alone.

\subsection{1. Exact source data and coefficient convention}\label{exact-source-data-and-coefficient-convention}

\begin{longtable}[]{@{}
  >{\raggedright\arraybackslash}p{(\columnwidth - 2\tabcolsep) * \real{0.5000}}
  >{\raggedright\arraybackslash}p{(\columnwidth - 2\tabcolsep) * \real{0.5000}}@{}}
\toprule\noalign{}
\begin{minipage}[b]{\linewidth}\raggedright
Source
\end{minipage} & \begin{minipage}[b]{\linewidth}\raggedright
Content used
\end{minipage} \\
\midrule\noalign{}
\endhead
\bottomrule\noalign{}
\endlastfoot
\nolinkurl{Columns/FiniteSumNonzeroLargeDepth.lean}, 127--152, 201--218 & \nolinkurl{xiBr}, \nolinkurl{kapBr}, \nolinkurl{brConst}, \nolinkurl{brStd}, \nolinkurl{brAt}, \nolinkurl{brC}, \nolinkurl{tilQ}, \nolinkurl{PtilFF} \\
Same file, 220--245, 1410--1418, 1451--1454 & \nolinkurl{PtilC}, the actual \nolinkurl{evRho}/\nolinkurl{evGSD}, coupled \nolinkurl{colC}, and \nolinkurl{PerInv} \\
\nolinkurl{Columns/CoefModB.lean}, 55--66 & Bilinear Fischer \nolinkurl{pairBK}, divided by the extremal Fischer norm \\
\nolinkurl{Columns/ColForms.lean}, 22--37, 53--74; \nolinkurl{CoefCol.lean}, 59--75, 97--99; \nolinkurl{RecLevel.lean}, 50--52 & Inverse level characters, \nolinkurl{rhoVr} columns, \nolinkurl{rhoV} detectors, and the inverse character of the adelic reading \\
\nolinkurl{Columns/RowForms.lean}, 221--229, 253--269; \nolinkurl{GenericUnipotentCharacter.lean}, 407--429, 526--557, 613 & Literal terminal residue groups and their characters \\
\nolinkurl{Columns/RhoEqOfSmulEq.lean}, 134--146, 534--536 & \nolinkurl{csOf} uses stabilizer cardinal inverses \\
\nolinkurl{Columns/Record.lean}, 1558--1610, 1901--1913 & \nolinkurl{lowC_ev}, \nolinkurl{periodC}, antilinear \nolinkurl{jMCL}, \nolinkurl{recPer}, and lowering on arithmetic rows \\
\nolinkurl{KernelLattices/InterlaceAll.lean}, 74--84 & \nolinkurl{lowerK}, \nolinkurl{periodK} \\
\nolinkurl{KernelLattices/KernelLocalFactors/InterlaceDifferenceShiftInvariant.lean}, 334--345 & \nolinkurl{PtilFF_eq}, \nolinkurl{PtilC_eq} \\
\nolinkurl{HermitianForms/LastCoordinateUniformizerVectorSpansOrthogonalLi.lean}, 58--77, 84--104 & \nolinkurl{ClsPTied} and \nolinkurl{ValPTied} class, stabilizer, coefficient, and metric transports \\
\nolinkurl{HermitianForms/FormsTied.lean}, 44--46 & \nolinkurl{PerTie}: \nolinkurl{period(w,G)=ipP(jM G,w)}, with \nolinkurl{jM} antilinear \\
Same file, 23--31; \nolinkurl{Columns/StdLatUnimodular.lean}, 165--178; \nolinkurl{Columns/TiedXpU.lean}, 172--199 & Exact decorated insertion and its rational/finite frame transport \\
\nolinkurl{LatticeFrames/IsRatDatum.lean}, 139--147 & \nolinkurl{TopForm}, including its stabilizer multiplier \\
\nolinkurl{LatticeFrames/TopFormula.lean}, 27--38 & \nolinkurl{IsJ}, \nolinkurl{TopFormula} \\
\nolinkurl{LatticeFrames/StabOrd.lean}, 24--37 & \nolinkurl{volLev}, \nolinkurl{stabOrd}, \nolinkurl{IsReps} \\
\nolinkurl{LatticeFrames/OnFramesRdApply.lean}, 137--150 & \nolinkurl{ClsT.Is.reps}, \nolinkurl{inv}, \nolinkurl{top}, \nolinkurl{tail} \\
\end{longtable}

The newly retrieved HermitianForms definitions and their mirror commit, URLs, and SHA256 hashes are under \nolinkurl{definite_oscillator_source_refs/}. They are not represented as having an unavailable original-author intake hash. Their actual formulas are checked below; there is no Lean execution.

Let the final column rank be N, let n=N−1, and let V\_lambda and V\_sigma be the source polynomial coefficient spaces over C. Fix the source interlacing weights and the exact operator

\[
 B=\operatorname{brC}
 =\left(\prod_\tau
   \tau(\overline{\pi_T})^{-nBr(\lambda_\tau,\sigma_\tau)}\right)
          \operatorname{brStd}:V_\lambda\longrightarrow V_\sigma.
 \tag{DO38.1}
\]

The source's \nolinkurl{lowC_ev} extends precisely this B to frame functions. The source's \nolinkurl{PtilFF_eq} identifies its finite period with \nolinkurl{periodK} after B. These source results are cited at their existing status, not re-proved here. The same file gives \nolinkurl{cBr_ne_zero} at1364, the nonzero interlacing branch at1599--1600 and1828--1831, and \nolinkurl{chain_brC_ne} at1861--1867.

Write

\[
 [P,Q]_\sigma=\operatorname{pairBK}_\sigma(P,Q),\qquad
 \langle P,Q\rangle_\sigma=[P,\overline Q]_\sigma.
 \tag{DO38.2}
\]

The Hermitian form is linear in its first argument. Its normalization is exactly \nolinkurl{nfAll}: ordinary monomial factorials divided by the extremal Fischer norm. Positivity and the source metric comparison are the checked R011 conclusions. In a nonorthonormal algebraic basis the actual Gram matrix remains in (DO38.2). The final column uses the contragredient \nolinkurl{rhoVr} model; the terminal detector uses \nolinkurl{rhoV}, paired by the bilinear form. Complex conjugation identifies their compact Hermitian realizations.

The frame weight \nolinkurl{sgT} is the literal shifted weight \nolinkurl{sgTrue}: \nolinkurl{TiedXpV.lean}, 216--217, proves this equality, and 136--142, 163--167, 201--206 prove that its natural partition and normalized Fischer form equal those used by the unshifted row polynomial. Thus the determinant factor in the frame function is retained even though it does not alter (DO38.2). The actual \nolinkurl{brC} is equivariant under the block embedding with last coordinate one; this is \nolinkurl{branchingProjection_act} in \nolinkurl{InterlaceDifferenceShiftInvariant.lean}, 3146--3158, together with \nolinkurl{PairBInv} in \nolinkurl{CoefModB.lean}, 86--89. The fixed source weights satisfy the stated same-shift hypothesis. No unshifted replacement of the branch is made.

The bilinear invariance can also be checked without treating that predicate declaration as a theorem: the ordinary Fischer pairing is \(P(\partial)Q|_0\). Replacing P by its inverse-transpose linear substitution and Q by the matching direct substitution cancels the two matrices by the chain rule. The determinant twists are reciprocal, and the common extremal norm is a scalar. This proves the displayed \nolinkurl{rhoVr}/\nolinkurl{rhoV} invariance for the actual polynomial spaces.

\subsection{2. Actual classes and equality of their stabilizers}\label{actual-classes-and-equality-of-their-stabilizers}

Let T be the totally definite rank-n unitary group on the actual mixed CM host. Fix the finite row level U and its finite-order unitary character chi. Infinite compact Haar has volume one. Set

\[
 v_U=\operatorname{vol}(U),\quad
 \mathcal C=T(F)\backslash T(\mathbb A_f)/U,\quad
 \Gamma_c=T(F)\cap h_cUh_c^{-1}.
 \tag{DO38.3}
\]

The fixed-rank lattice-class proof of K036 section3 applies at this odd rank as well; only its later determinant-power argument needs even ranks. Finite decorations refine each lattice class by a finite set. The local standardness argument there identifies the actual selfdual lattice genus with adelic unitary moves. A local frame supplies h\_c; its residue frame supplies the stipulated split-place decoration. Changing the frame by U or the rational basis gives exactly the two equivalences in (DO38.3). No global class-number-one assertion is used.

The literal \nolinkurl{ClsPTied.y_bij} expresses this class bijection, and its \nolinkurl{stabI_inj} and \nolinkurl{stabI_range} identify the arithmetic stabilizer with Gamma\_c.~To verify the adelic side explicitly, a rational gamma is counted by \nolinkurl{stabOrd(U,h_c)} exactly when

\[
                       h_c^{-1}\gamma_fh_c\in U.
 \tag{DO38.4}
\]

In \nolinkurl{gamma h_c=h_c k kappa}, the infinite element is forced to be gamma\_infinity and the finite element to be (DO38.4); conversely these choices give the equality. The source range condition is precisely (DO38.4). Thus

\[
 \boxed{\operatorname{stabOrd}(U,h_c)=|\Gamma_c|
       =\operatorname{Nat.card}(\operatorname{stabilizer}_2(c)).}
 \tag{DO38.5}
\]

Also \nolinkurl{volLev=v_U}, by its definition as the measure of the product of the infinite compact group and U. DO026(26.17) then gives the actual class integration coefficient v\_U/\textbar Gamma\_c\textbar. This is equality of the specific stabilizers, not an equality inferred from two total masses.

For the constructed decorated module this identification is direct: an arithmetic stabilizer has a rational isometry preserving the lattice and a residue-frame element in the chosen subgroup. The residue-frame map is surjective, so its residue element is uniquely determined by that rational isometry. Conversely (DO38.4) supplies that element by reduction. These are inverse maps and prove (DO38.5) for the constructed family, rather than assuming an uninhabited \nolinkurl{ClsPTied} record on the new host.

\subsection{3. The arithmetic pairing before its automorphic realization}\label{the-arithmetic-pairing-before-its-automorphic-realization}

Let M be the terminal coefficient-row module and retain its fixed arithmetic pairing

\[
 \mathscr P(W,G)=\sum_{c\in\mathcal C}|\Gamma_c|^{-1}
                                      [W_c,G_c]_\sigma.
 \tag{DO38.6}
\]

This is \nolinkurl{periodK}. The source \nolinkurl{periodC} is this pairing in the chosen row basis; its matrix is not replaced by an identity matrix. Let H\_G,c be the Hermitian Riesz vector of \nolinkurl{W_c -> [W_c,G_c]_sigma}. In the identity frame and real polynomial model,

\[
                         H_{G,c}=\overline{G_c}.
 \tag{DO38.7}
\]

This holds termwise in the Fischer formula, even in a nonorthonormal primitive basis. The polynomial of the source \nolinkurl{eW(jM G)_c} is exactly this Riesz vector. In another orthonormal frame the actual Gram transport gives the equivalent Riesz equation, without a new scalar.

Indeed \nolinkurl{ValPTied.ipP(z,z')} is the sum of \nolinkurl{nfAll(Q,P)}, where P belongs to z and Q to z'. Together with \nolinkurl{PerTie} it gives

\[
 \mathscr P(W,G)=\sum_c|\Gamma_c|^{-1}
                               \langle W_c,H_{G,c}\rangle_\sigma.
 \tag{DO38.8}
\]

The order and the antilinearity of jM make this linear in both W and G. The fixed arithmetic tail is therefore

\[
 \boxed{\mathscr T(\psi,G)=\operatorname{PtilFF}(\psi,G)
   =\sum_c|\Gamma_c|^{-1}
                       [B\psi(d_c),G_c]_\sigma.}
 \tag{DO38.9}
\]

Here, exactly as in the source, \nolinkurl{PtilFF} takes an already restricted function \(\psi:\operatorname{RowPt}\to V_\lambda\). If f is a full final column, put \(f^\downarrow(d)=f(\operatorname{tilC}d)\). Then \nolinkurl{PtilC(f,G)=PtilFF(f^downarrow,G)}. The source \nolinkurl{tilC} is \nolinkurl{tilQ} followed by forgetting the additional Q decorations. A final coupled column \nolinkurl{colC} instead lands directly in \nolinkurl{evGSD.EqFun}; it is already such a restricted function. These two domains are not identified by notation. In either case B and the arithmetic period are fixed before the automorphic realization.

Write \(\lambda_U=\operatorname{rowTy}.\mathrm{chr}\). The actual source actions are

\[
 \psi(gd)=\lambda_U(g_f)\rho^\vee_\lambda(\operatorname{ext}g_\infty)
                  \psi(d),\qquad
 G(gd)=\lambda_U(g_f)^{-1}\rho_\sigma(g_\infty)G(d).
 \tag{DO38.9a}
\]

The first is \nolinkurl{evRho}; the second is \nolinkurl{formRho} for \nolinkurl{rowTy.inv}. Equivariance of B and invariance of \nolinkurl{pairBK} prove that every summand in (DO38.9) is independent of the representative. They also prove that \(W=B\psi\) lies in the opposite coefficient space. This is the literal \nolinkurl{PerInv} equation, with the two finite characters canceled, not a new character convention. In the source adelic reading \nolinkurl{levTyp=(levChr)^-1}; hence both the lowered column and its Riesz family have reading character \(\chi=\lambda_U^{-1}\), and the scalar top automorphic form will have character \(\overline\chi\).

For a full column f, (DO38.9a) requires the actual block level inclusion and character restriction (\nolinkurl{LevExt}, \nolinkurl{Record.lean}, 1501--1519). For the source m=2 terminal groups it holds when \(P.H\le U_r\) and the displayed characters restrict as prescribed; in particular \nolinkurl{rPsym} has this inclusion by \nolinkurl{GenericUnipotentCharacter.lean}, 526--557. For the coupled column it is part of its actual \nolinkurl{evGSD.EqFun} codomain. Mere nonvanishing of a finite Schwartz factor does not establish this equivariance.

\subsubsection{3.1 The decorated insertion is the actual group restriction}\label{the-decorated-insertion-is-the-actual-group-restriction}

This transport can be constructed on the new mixed host. Choose the base final decoration to be the literal insertion of the base row decoration. The source formula \nolinkurl{Deco.til} has underlying lattice \(N\oplus O_Ee_N\), T-functional \(x\mapsto\overline{\pi_T}x_N\), and G/Q residue frame with last coordinate \(x_N/\pi_T\). Thus \(\operatorname{diag}(h_c,1)\) sends the base decoration to \(\operatorname{tilQ}(d_c)\) at every place, preserving these exact last-coordinate normalizations. Here \(\pi_T\) is a unit at the G/Q places, as in the literal source disjointness assumptions.

Let a be the chosen final decorated class containing this insertion. There are a rational unitary \(\gamma_c\) and a finite final-level element \(\kappa_c\) such that

\[
 \gamma_{c,f} h_a=\operatorname{diag}(h_c,1)\kappa_c
                         \quad\text{in }T_N(\mathbb A_f).
 \tag{DO38.9b}
\]

This is the defining double-coset equality, so it does not require class number one or a new existence theorem. The rational decorated move and its residue component are exactly the arithmetic move used to pass from the chosen final representative to \nolinkurl{tilQ(d_c)}. Substitution in \nolinkurl{rdS} gives the actual restricted frame function

\[
 \Phi_{f,c}(y)=\chi_N(\kappa_c)^{-1}
                            f_a(\gamma_{c,\infty}^{-1}y).
 \tag{DO38.9c}
\]

The arithmetic equivariance gives precisely the same expression for the reading of \(f^\downarrow(d_c)\). This is \nolinkurl{EvzTie}, including its inverse finite character, derived from (DO38.9b). The embeddings fix the line generated by \(z_0=\pi_Te_N\), since fixing that line vector is equivalent to fixing \(e_N\). Changing representatives modifies both sides by the same rational and finite actions, already canceled in (DO38.9a). The factor \(\overline\pi_T^{-nBr}\) in B is still present: neither (DO38.9b) nor (DO38.9c) discards it.

Thus \nolinkurl{evzC}, \nolinkurl{lowC_ev}, and \nolinkurl{recPer_lower_emb} apply to this actual canonical geometric reading of an admitted full column. For a coupled final row, the same terminal mass calculation applies once its fixed decorated kernel expansion has this reading and the equivariance (DO38.9a). The source \nolinkurl{colC} has exactly that equivariant codomain, but transplanting a standard local Y-sum to moving middle classes needs its actual covariance and transport proof. Neither a nonzero quotient of the standard local sum nor pointwise CRT alone supplies that proof. An unrelated class map or an arbitrary finite rephasing would not satisfy (DO38.9c).

\subsection{4. A canonical actual top automorphic realization}\label{a-canonical-actual-top-automorphic-realization}

Let V be the compact coefficient representation in the source frame convention, with \texttt{rho(k)Phi(y)=Phi(k\^{}\{-1\}y)}. Use (DO38.2) and let J identify the top compact representation with V. One may take the actual polynomial representation and J the identity; an equivalent compact DSModel is transported by that same intertwiner. This fixes \nolinkurl{IsJ}, rather than freely rescaling each row or each class.

The admitted coefficient families satisfy the literal \nolinkurl{FiC} condition

\[
 \rho(\gamma_\infty)H_c
   =\chi(h_c^{-1}\gamma_fh_c)H_c
                  \quad(\gamma\in\Gamma_c).
 \tag{DO38.10}
\]

For g=gamma h\_c k kappa define

\[
 H(g)=\chi(\kappa)\rho(k^{-1})H_c,
 \qquad e_T(G,v)(g)=\langle Jv,H_G(g)\rangle_\sigma.
 \tag{DO38.11}
\]

The first equation is exactly the \nolinkurl{rdP}/\nolinkurl{rdS} reading, whose frame argument is \texttt{gamma\_infinity\^{}\{-1\}\ g\_infinity\ y}. Any two decompositions of g differ by an element of Gamma\_c.~Substituting (DO38.10) shows that they give the same value. The resulting function is left T(F)-invariant, smooth, and finite under the infinite compact group. It transforms by conjugate(chi) on the right under U. Since T is anisotropic, its quotient is compact and has no proper rational parabolic constant terms; the function is cuspidal.

The formula is linear in G because H\_G is antilinear and occupies the second Hermitian argument. It is linear in v. Unitarity gives

\[
                       e_T(G,v)(gk)=e_T(G,\pi(k)v)(g).
 \tag{DO38.12}
\]

At a finite representative it gives

\[
 \boxed{e_T(G,v)(h_c)
   =\langle Jv,H_{G,c}\rangle_\sigma
   =|\Gamma_c|\mathscr P(\operatorname{single}_c(Jv),G).}
 \tag{DO38.13}
\]

This is the literal \nolinkurl{TopForm}/\nolinkurl{TopFormula}, with its forced stabilizer multiplier. The single-class family here belongs to the ambient space of coefficient families; it need not itself satisfy the stabilizer constraint. The source period is defined on that ambient space. Equivalently one may first average Jv over the finite stabilizer with its prescribed character: pairing with H\_G,c gives the same value, since H\_G,c has exactly that character. For the full coefficient-row module, (DO38.11) gives the full top multiplicity space at this level: evaluating any intertwining automorphic embedding at each h\_c gives a functional on V, whose Riesz vector satisfies (DO38.10). Riesz duality recovers the unique coefficient family. A separately restricted arithmetic subspace need not be onto that full multiplicity space.

The actual norm, for the displayed J, is

\[
 \boxed{\|e_T(G,v)\|^2
   =\frac{v_U}{\dim V}\,\|Jv\|^2
                       \sum_c|\Gamma_c|^{-1}\|H_{G,c}\|^2.}
 \tag{DO38.14}
\]

This is the checked class and compact-Schur formula of DO026(26.17)--(26.18). Thus this realization is not declared isometric to an independently normalized source row. Its level volume and Schur dimension remain explicit; they are not class-dependent constants in (DO38.13).

\subsubsection{4.1 Which side the terminal branch acts on}\label{which-side-the-terminal-branch-acts-on}

For clarity, let \(P_c\) be the finite-class coefficient of a full column, and write its complete frame function as \(\Phi_{P_c}(y)=\operatorname{fnOf}(P_c)(y)\), including the source determinant shift. At a finite class representative and an infinite compact element k, the literal reading is

\[
       \operatorname{rd}(h_c k,y)=\Phi_{P_c}(k y),\qquad
       \operatorname{lowC}(\Phi_{P_c})
                         =\operatorname{ev}_\sigma(BP_c).
 \tag{DO38.14a}
\]

The second equality acts on the whole frame function before a value of y is chosen. It therefore branches the class coefficient P\_c. A single highest-weight Schwartz kernel instead fixes one right compact seed in the scalar automorphic output. Choosing the identity frame can make that particular scalar function restrict only to the top H-type. This does not imply that P\_c is a highest polynomial: in a matrix coefficient, the class coefficient and the fixed right seed occupy the two different factors. Right translations vary the seed while retaining the class coefficient.

For example, in the elementary compact representation Sym\^{}2(C\^{}2), let the right seed be e\_1\^{}2 and the class functional be the coefficient of e\_2\^{}2. The scalar matrix coefficient on diag(U(1),1) is identically zero, although that class functional has a nonzero restriction to the e\_2\^{}2 branch; a right rotation of the seed detects it. This example only illustrates the distinction. No such new rotation is inserted in the fixed arithmetic operator. Its actual test is BP\_c after the literal decorated transport (DO38.9c).

In the source \nolinkurl{PerX}, the proportionality factor a(y,v) is permitted to vanish for particular frames or compact vectors. Consequently a vanishing period at a single highest seed neither proves nor disproves the arithmetic tail. A complete-type Kval realization is what recovers P\_c and makes the test in (DO38.14a) meaningful.

\subsection{5. The exact terminal class-mass equation}\label{the-exact-terminal-class-mass-equation}

Choose the canonical terminal finite trace \nolinkurl{delta_1} and one terminal chart. Define the class functional from the fixed period by

\[
 \ell_{c,G}(\Phi)=|\Gamma_c|
              \mathscr P(\operatorname{single}_c\Phi,G).
 \tag{DO38.15}
\]

For an arbitrary frame function, (DO38.15) uses the literal \nolinkurl{periodC} extension through \nolinkurl{Ev.extend}; it agrees with the fixed bilinear polynomial pairing on V by \nolinkurl{periodC_eq_ipCl}. Thus it is linear on the full domain required by \nolinkurl{ClsT.ell}; on V it equals \nolinkurl{<Phi,H_G,c>_sigma}. Equations (DO38.12)--(DO38.13) give \nolinkurl{ClsT.Is.top}. The opposite final-column and top finite characters cancel in their product, giving \nolinkurl{ClsT.Is.inv}; section2 gives \nolinkurl{reps}.

Let Phi\_psi,c be the final column restriction at h\_c, read by the stipulated decorated insertion and frame map. The literal \nolinkurl{lowC} transport says, for every coefficient row,

\[
 \operatorname{lowC}(\Phi_{\psi,c})
      =\operatorname{ev}_\sigma(B\psi(d_c)).
 \tag{DO38.16}
\]

This is the source \nolinkurl{lowC_ev}/\nolinkurl{recPer_lower_emb} identity on the entire polynomial space. In the new family it is implemented by those same decorated insertions and readings. It would not follow for an unrelated class map merely from a geometric theta identity.

Substituting (DO38.5), (DO38.15), and (DO38.16) into the literal terminal class-mass sum gives

\[
 \begin{split}
 v_U\sum_c|\Gamma_c|^{-1}
             \ell_{c,G}(\operatorname{lowC}\Phi_{\psi,c})
 &=v_U\sum_c|\Gamma_c|^{-1}
                     [B\psi(d_c),G_c]_\sigma\\
 &=v_U\operatorname{PtilFF}(\psi,G).
 \end{split}
 \tag{DO38.17}
\]

Therefore, uniformly in both rows and in every class,

\[
                         \boxed{c_P=v_U=\operatorname{volLev}(U)>0.}
 \tag{DO38.18}
\]

The coefficient Gram, hooks, \nolinkurl{kapBr}, and T-uniformizer powers remain inside the fixed bilinear pairing and B on both sides. None has been normalized away. For maximal finite compact volumes one,

\[
                 c_P=\prod_{v\text{ carrying a row level}}
                                         [K_v:U_v]^{-1}.
 \tag{DO38.19}
\]

At a split residue level H\_v in GL\_n(k\_v), its factor is \nolinkurl{|H_v|/|GL_n(k_v)|}, since reduction is surjective and its principal kernel lies in U\_v. The value is thus an explicit positive rational index product, without enumerating individual global class stabilizers. A different Haar choice rescales both sides by that same factor.

\subsubsection{5.1 Literal G/Q volumes and their denominators}\label{literal-gq-volumes-and-their-denominators}

Put \(r=k+m-1\) and \(n=r+mt=N_m-1\). The top row has T-parameter one and no T-decoration. Its G-level is the full preimage of \(H_G=\operatorname{HPar}(n;r,t,\ldots,t;P.H)\), and its Q-level is the full preimage of \(H_Q=P_{r,mt}\). These are the groups in \nolinkurl{rowTyG} and \nolinkurl{rowTyQ}, not the kernels of their characters. Their orders are

\[
 \begin{split}
 |H_G(\mathbf F_q)|
  &=|P.H(\mathbf F_q)|\,|\operatorname{GL}_t(\mathbf F_q)|^m
                  q^{rmt+\binom m2t^2},\\
 |H_Q(\mathbf F_q)|
  &=|\operatorname{GL}_r(\mathbf F_q)|
    |\operatorname{GL}_{mt}(\mathbf F_q)|q^{rmt}.
 \end{split}
 \tag{DO38.19a}
\]

To prove the first formula, choose the leading diagonal block in P.H, each of the m remaining diagonal blocks in GL\_t, and all upper block entries freely. Their dimensions are rmt and binom(m,2)t\^{}2. All lower blocks are zero, and block triangularity is exactly invertibility of the displayed diagonal blocks. The second formula is the same argument with two unrestricted diagonal blocks. Consequently

\[
 \boxed{c_P=
 \prod_{v\in G}
   \frac{|P.H(\mathbf F_{q_v})|\,|GL_t(\mathbf F_{q_v})|^m
             q_v^{rmt+\binom m2t^2}}{|GL_n(\mathbf F_{q_v})|}
 \prod_{v\in Q}
   \frac{|GL_r(\mathbf F_{q_v})|\,|GL_{mt}(\mathbf F_{q_v})|
             q_v^{rmt}}{|GL_n(\mathbf F_{q_v})|}.}
 \tag{DO38.19b}
\]

Here every actual real-field place is counted once; the selection of one CM prime above it is only the split identification. If an additional finite refinement is imposed, multiply by its index inverse. The displayed canonical choice has maximal compact at 29; no principal 29 refinement is introduced to make the class formulas work. The exact logarithmic denominator cost of c\_P is \(\sum_v\log[K_v:U_v]\), with possible cancellation when combined with other independently specified constants left visible.

For q prime to 11, let e=ord\_11(q) and c=v\_11(q\^{}e−1). Directly from \(|GL_j(\mathbf F_q)|=q^{j(j-1)/2}\prod_{i=1}^j(q^i-1)\) and the odd-prime lifting-of-exponent identity,

\[
 v_{11}|GL_j(\mathbf F_q)|
   =\lfloor j/e\rfloor c+v_{11}(\lfloor j/e\rfloor!).
 \tag{DO38.19c}
\]

The identity follows because 11 divides q\^{}i−1 exactly when e divides i, and then v\_11(q\^{}(e s)−1)=c+v\_11(s), obtained by expanding successive 11th powers and then a power prime to 11. For the source \nolinkurl{rPsym}, P.H is a subgroup of the finite upper unipotent group, so its order is a power of q. For m=2, r=k+1 and rS=0 it consists of U\_k with the last column fixed, and its exact order is q\^{}(k(k−1)/2).

\begin{longtable}[]{@{}
  >{\raggedright\arraybackslash}p{(\columnwidth - 10\tabcolsep) * \real{0.1364}}
  >{\raggedright\arraybackslash}p{(\columnwidth - 10\tabcolsep) * \real{0.1364}}
  >{\raggedleft\arraybackslash}p{(\columnwidth - 10\tabcolsep) * \real{0.1818}}
  >{\raggedleft\arraybackslash}p{(\columnwidth - 10\tabcolsep) * \real{0.1818}}
  >{\raggedleft\arraybackslash}p{(\columnwidth - 10\tabcolsep) * \real{0.1818}}
  >{\raggedleft\arraybackslash}p{(\columnwidth - 10\tabcolsep) * \real{0.1818}}@{}}
\toprule\noalign{}
\begin{minipage}[b]{\linewidth}\raggedright
Host
\end{minipage} & \begin{minipage}[b]{\linewidth}\raggedright
k, t, n, number f of selected places per G prime
\end{minipage} & \begin{minipage}[b]{\linewidth}\raggedleft
G prime q
\end{minipage} & \begin{minipage}[b]{\linewidth}\raggedleft
v11(
\end{minipage} & \begin{minipage}[b]{\linewidth}\raggedleft
GL\_n
\end{minipage} & \begin{minipage}[b]{\linewidth}\raggedleft
)
\end{minipage} \\
\midrule\noalign{}
\endhead
\bottomrule\noalign{}
\endlastfoot
K031 & 42, 1759, 3561, 42 & 4615409 & 390 & 191 & −8 \\
K031 & same & 5276377 & 390 & 191 & −8 \\
K032 & 168, 1633, 3435, 168 & 22285932761 & 754 & 357 & −40 \\
K032 & same & 64072175441 & 1888 & 896 & −96 \\
\end{longtable}

Thus the exact G contribution is v11(c\_P,G)=−672 on K031 and −22848 on K032 for this row group. The Q factors are the inverse Gaussian binomial indices in (DO38.19b) and must still be included at their actual primes. These numbers concern the terminal volume only; they are not the cost of the branch, the theta kernel, or an optimal integral terminal detector.

\subsubsection{5.2 The full terminal period, including its real scalar}\label{the-full-terminal-period-including-its-real-scalar}

There is also an exact real normalization, useful for distinguishing the arithmetic branch from a single highest seed. Let B* be the Hermitian adjoint of the fixed B for the two source normalized Fischer metrics, and put \(d_\sigma=\dim V_\sigma\). The compact branch has multiplicity one. Since B is a nonzero intertwiner, Schur's lemma and positivity give

\[
 BB^*=\beta_B I,\qquad
 \Pi_\sigma=\beta_B^{-1}B^*B,\qquad
 \beta_B=\|B\|^2>0.
 \tag{DO38.19d}
\]

The scalar is explicit in root R019(19.15):

\[
 \boxed{\beta_B=\prod_\tau
  |\tau(\overline\pi_T)|^{-2nBr_\tau}
                    |c_{{\rm Br},\tau}|^2\kappa_\tau.}
 \tag{DO38.19e}
\]

For a source frame y let \(E_y(P)=\Phi_P(y)\). This is a fixed linear evaluation functional on the complete coefficient space, with its determinant factor included. Then the actual period for an admitted full-column reading is

\[
 \boxed{\int_{[T]} e_T(G,v)(h)\,
                  \operatorname{rd}(f,y)(\iota h)\,dh
       =v_U\,\frac{E_y(B^*Jv)}{\beta_Bd_\sigma}
                         \operatorname{PtilC}(f,G).}
 \tag{DO38.19f}
\]

For an actually assembled coupled reading, replace PtilC(f,G) by PtilFF(psi,G) under precisely the reading condition of section3.1. Thus the source PerX scalar is \(a(y,v)=E_y(B^*Jv)/(\beta_Bd_\sigma)\); c\_P remains v\_U.

\textbf{Proof.} Work on one finite class after its fixed decorated transport, and write P for that class coefficient and \(\ell_G(W)=[W,G_c]_\sigma\). By (DO38.11)--(DO38.12), the two compact factors are

\[
   \ell_G(\rho_\sigma(k)Jv),\qquad
              E_y(\rho_\lambda(\iota k)^{-1}P).
 \tag{DO38.19g}
\]

The integral kills the H-types other than sigma. On its single sigma summand substitute \(\Pi_\sigma P=B^*BP/\beta_B\). B* intertwines, so the second factor becomes \(\beta_B^{-1}(E_y\circ B^*)(\rho_\sigma(k)^{-1}BP)\). The bilinear compact matrix-coefficient identity is

\[
 \int_{T_\infty}\ell_1(\rho(k)v)\ell_2(\rho(k)^{-1}w)\,dk
             =d_\sigma^{-1}\ell_1(w)\ell_2(v).
 \tag{DO38.19h}
\]

It follows by averaging the rank-one endomorphism \(v\otimes\ell_2\): its average commutes with the irreducible representation, and its trace is \(\ell_2(v)\). Applying the average to w and then \(\ell_1\) proves (DO38.19h). Thus the one-class integral is \([BP,G_c]_\sigma E_y(B^*Jv)/(\beta_Bd_\sigma)\). Multiplication by the exact class measure v\_U/\textbar Gamma\_c\textbar{} and summation gives (DO38.19f). No complex conjugation of the bilinear period is inserted. The conjugation is already in the adjoint and Riesz vectors.

If v is nonzero then B\emph{Jv is a nonzero polynomial frame function, so some frame y has nonzero E\_y(B}Jv). Even unitary frames suffice: their Zariski density follows by complexifying the Lie algebra, or equivalently by the fact that the compact group's matrix coefficients separate its irreducible vectors. This proves that the proportionality factor is not identically zero, while allowing it to vanish at an individual highest frame. It does not prove that the particular arithmetic coefficient family has BP\_c nonzero.

There is a quantitative archimedean normalization. Let \(d_\lambda=\dim V_\lambda\), take \(\|Jv\|=1\), and use the standard unitary frame orbit. Evaluation at the identity frame is the unit highest-coefficient functional: all other weight spaces vanish there, because the column weights are fixed and only diagonal monomials survive; the normalized highest polynomial has value and norm one. Hence every E\_y on this compact orbit has norm one. Compact Schur orthogonality on V\_lambda gives

\[
 \int_{U(N)_\infty}|E_y(B^*Jv)|^2\,dy
        =\frac{\beta_B}{d_\lambda}.
 \tag{DO38.19i}
\]

Some unitary frame therefore satisfies \(|a(y,v)|\ge1/(d_\sigma\sqrt{\beta_Bd_\lambda})\). Rescaling v by a(y,v)\^{}−1 makes the factor a exactly one, with norm at most \(d_\sigma\sqrt{\beta_Bd_\lambda}\) for the canonical isometric J. The resulting actual period is exactly v\_U PtilC(f,G). This is an archimedean bound and a normalization of a freely chosen compact test vector. It neither changes the fixed arithmetic branch nor supplies a height bound for an algebraic frame or an integral detector.

\subsection{6. Exact nonzero values and retained finite-projector scope}\label{exact-nonzero-values-and-retained-finite-projector-scope}

Suppose the lowered family

\[
                         W_c=B\psi(d_c)
 \tag{DO38.20}
\]

is nonzero and belongs to the admitted opposite coefficient space. The full canonical terminal module contains its Riesz dual, given in the identity frame by G\_c=conjugate(W\_c). For that actual row,

\[
 \boxed{\operatorname{PtilFF}(\psi,G)
           =\sum_c|\Gamma_c|^{-1}\|W_c\|_\sigma^2>0.}
 \tag{DO38.21}
\]

The terminal class-mass value is c\_P times this number. If W is algebraic over a conjugation-stable number field, G is algebraic over that same field. Its integral denominator still depends on its actual coefficient lattice. A smaller prescribed detector space need not contain G. Nonzero final theta output does not imply (DO38.20) nonzero: the fixed branch and fixed decorated insertion can kill it. That exact coefficient test remains; generic branching or arbitrary rational rotation is not substituted for it.

For a nonidentity prescribed finite trace, retain its arithmetic coefficient operator. If e preserves the admitted row module and its exact action is E, then \nolinkurl{trc(e)e_T(G,v)=e_T(EG,v)}; the same proof replaces G by EG in (DO38.17). This follows by applying the actual class reading to every finite translation, including its character and rational frame factors. Several terminal charts give the corresponding sum of coefficient actions. Equality with an independently fixed, unmodified \nolinkurl{PtilFF} requires that additional operator equality. No such nonidentity-projector equality is asserted here.

\subsection{7. Consequence for the canonical two-step family}\label{consequence-for-the-canonical-two-step-family}

For identity terminal trace the literal m=2 top factor becomes

\[
                  c_{\mathrm{Top}}=c_{V,0}c_{K,0}c_{K,1}\,v_U.
 \tag{DO38.22}
\]

The kernel periods and creation frames come from the independent DO026/K035/K036 family and its canonical class-row operator. Equation (DO38.17) supplies the terminal finite mass tie for its admitted ordinary columns read by the literal decorated insertion and brC, and for any coupled final row with the actual reading and covariance specified in section3.1. The top realization has been constructed from the fixed arithmetic pairing, not left as an arbitrary external normalization. An ordinary output whose pullback fails the Q character condition is not silently treated as such a coupled row.

To invoke MainX one must verify its actual \nolinkurl{LawsX} identities for this same family, including the geometric arithmetic column maps and any coupled finite assembly. MainX itself is a vanishing implication and does not require a separately proved nonzero lowered family. Its contrapositive would give a nonzero prescribed tail from a nonzero bottom period of the stipulated descended top input. Alternatively, (DO38.20)--(DO38.21) give a direct terminal detector if that lowered family is known to be nonzero. The integral-image and cost estimates remain additional arithmetic questions in either route. The old common-quadratic field wrapper is not inhabited. This is a mathematical adapter for the specified mixed geometry and coefficient rows, with its finite trace explicitly fixed.

\subsection{Verification}\label{verification}

The companion finite certificate checks unequal stabilizer masses, nonidentity Gram matrices, the bilinear/Hermitian Riesz convention, a nontrivial small branching polynomial, the exact real Schur scalar with a nonidentity source metric, and finite-index normalization. It passes 144 unequal-mass identities, 432 Riesz identities, 144 positive full-detector tests, 64 real Schur tests, six enumerated finite-group cases, and 814 direct-versus-LTE valuation checks, as well as the literal branching and zero-seed examples. It does not enumerate the actual large-rank class set or establish nonvanishing of (DO38.20) for the selected theta image. Author and independent whole-proof events are recorded separately. No Lean is run.

\clearpage

\section{AM-R-020 --- the whole integral terminal branch and a uniform content bound}\label{am-r-020-the-whole-integral-terminal-branch-and-a-uniform-content-bound}

Author: GPT-6 Astra Pro, ALP\_MACHINE /root. 22 November 2026 UTC. Reasoning effort is not externally exposed.

R019 evaluated the literal terminal branch on its distinguished polynomial highest line and computed its exact complex norm. That polynomial line is not, in general, the complete integral image of the branch. This finding bounds the branch on the actual standard Weyl lattice, without identifying it with that line or assuming that its image is a scalar multiple of the target Weyl lattice. The bound is uniform in the growing first weight of the two concrete hosts. It concerns the whole branch operator; its restriction to an earlier theta image is explicitly excluded.

\subsection{1. Fixed conventions and accepted inputs}\label{fixed-conventions-and-accepted-inputs}

First remove the common final weight from the interlacing pair, so that we have partitions lambda of length N=n+1 and mu of length n with lambda\_N=mu\_n=0. The general same-shift pair is obtained by the common determinant twist; it has exactly these polynomial lattices. The literal T-uniformizer exponent is retained separately below and is not changed by this notational reduction.

Write Delta\_lambda and Nabla\_lambda for the integral Weyl and dual-Weyl polynomial lattices in V\_lambda over Q, and similarly at mu. Use the source symmetric bilinear normalized Fischer forms

\[
 [P,Q]_\lambda=\frac{\operatorname{Fischer}(P,Q)}{H_\lambda},
 \qquad H_\lambda=\|p_\lambda\|_{\rm Fischer}^2,
 \tag{R20.1}
\]

where p\_lambda is the primitive extremal, of normalized norm one. These are the exact \nolinkurl{pairBK} conventions of K037 and DO038. The following accepted results are used, not proved anew:

\begin{itemize}
\tightlist
\item
  \nolinkurl{Formalism/LocalRingCardinality/FactDelDualNabla.lean}, theorem \nolinkurl{factDelDualNabla}, line75 in the recovered copy; its statement is \nolinkurl{DataCoeff.lean}126--129. Delta is the full dual lattice of Nabla under (R20.1). Biduality therefore makes Nabla the full dual of Delta.
\item
  \nolinkurl{Formalism/GradedWeylDivisibility/StatementA.lean}34--38 and68--80: Delta is generated from p\_lambda by integral row divided-power words; Nabla is the column-equation lattice, and these row operations preserve it. Its row-weight decomposition is an integral direct sum.
\item
  The multiplicity-one characteristic-zero branching and the literal highest embedding are those used in R019. In particular there is a rational map i:V\_mu to V\_lambda with i(p\_mu)=xiBr. The polynomial xiBr is integral, belongs to Nabla\_lambda, and its last-row coefficient at bBr is exactly p\_mu. R019 proves the required normalization, including its positive rational normalized squared norm kappa.
\item
  \nolinkurl{Columns/FiniteSumNonzeroLargeDepth.lean}118--141 defines B=brStd=cBr\emph{kappa}projVV\emph{coefLast. The accepted branch equivariance is \nolinkurl{InterlaceDifferenceShiftInvariant.lean}3146--3158; its normalized adjoint identity is \nolinkurl{brTranspose}, 304--314, or \nolinkurl{inD_adj},3870--3883. R019 gives B(xiBr)=cBr}kappa\emph{p\_mu and \textbar\textbar B\textbar\textbar\^{}2=\textbar cBr\textbar\^{}2}kappa. None of these identifies an integral image.
\end{itemize}

All integer lattices in this note are the displayed lattices. We do not replace Delta by Nabla when they have different index.

\subsection{2. Exact lattice inclusions for the whole branch}\label{exact-lattice-inclusions-for-the-whole-branch}

Let c=cBr, which is a nonzero integer. Then

\[
 \boxed{\quad
 c\Delta_\mu\ \subset\ B(\Delta_\lambda)
       \ \subset\ c\nabla_\mu
       \ \subset\ \frac{c}{H_\mu}\Delta_\mu .\quad}
 \tag{R20.2}
\]

The inclusions are inclusions of fractional lattices in V\_mu; the sign of c has no effect on them. In particular the whole branch image need not equal c times Delta\_mu.

\textbf{Proof.} Let i\^{}dagger denote the adjoint of i for (R20.1). All coefficients are rational, so this is also the Hermitian adjoint for the corresponding compact metrics. Branching multiplicity one and the normalization on p\_mu give

\[
 i^\dagger i=\kappa I,\qquad B=c\,i^\dagger .
 \tag{R20.3}
\]

Indeed the first equality follows from invariance and the squared norm of i(p\_mu). Both sides of the second equality are the same equivariant projection up to a scalar; applying them to xiBr determines that scalar, since both values are c\emph{kappa}p\_mu. Equivalently (R20.3) is the literal normalized transpose identity cited above with the highest embedding divided by c.

Every integral divided-power word in the rank-n row operators commutes with i and preserves Nabla\_lambda. Applying the word to p\_mu and to xiBr proves

\[
                         i(\Delta_\mu)\subset\nabla_\lambda.
 \tag{R20.4}
\]

Consequently, for v in Delta\_lambda and w in Delta\_mu, {[}i\^{}dagger v,w{]}\_mu={[}v,iw{]}\_lambda is integral. Perfect duality makes i\^{}dagger v an element of Nabla\_mu. This proves the middle inclusion in (R20.2).

For the first inclusion, xiBr is a primitive element of Nabla\_lambda. In fact coefLast(xiBr)=p\_mu, and p\_mu has an integer monomial coefficient one. Thus xiBr also has an integer coefficient one; that coefficient functional on Nabla\_lambda takes value one at xiBr. Perfect duality supplies v in Delta\_lambda satisfying

\[
                              [v,\xi_{\rm Br}]_\lambda=1.
 \tag{R20.5}
\]

Replace v by its integral row-weight component having the weight of xiBr. This component remains in Delta: its divided-power generators are row-weight homogeneous, so the generator span is the direct sum of its row-weight spans. Distinct row weights are Fischer orthogonal, so (R20.5) is preserved. The first n entries of that weight are mu. The image i\^{}dagger v is therefore in the highest weight line Q*p\_mu of V\_mu. Pairing it with p\_mu and using (R20.5) shows that it equals p\_mu. Thus p\_mu belongs to i\^{}dagger(Delta\_lambda).

The latter image is stable under all integral rank-n row divided powers. One way to check this in the polynomial convention is to transpose an intertwining row divided-power identity for i: the adjoint of the (r,s) divided power is the (s,r) divided power for the normalized Fischer form. Both source Weyl lattices are stable under these operators. The divided-power generation of Delta\_mu now gives Delta\_mu contained in i\^{}dagger(Delta\_lambda), proving the first inclusion after multiplication by c.

Finally any two elements of Nabla\_mu are integer polynomials. Their ordinary Fischer pairing is an integer. Equation (R20.1) gives {[}H\_mu*u,w{]}\_mu in Z for all u,w in Nabla\_mu. Perfect duality then gives

\[
                          H_\mu\nabla_\mu\subset\Delta_\mu,
 \tag{R20.6}
\]

which proves the last inclusion. This proof uses no saturation of Delta in the ambient polynomial module. The integral inclusions remain valid after tensoring with the integers of a number field and localizing at any finite place. QED.

\subsection{3. Exact finite-place formulation, tensor labels and frames}\label{exact-finite-place-formulation-tensor-labels-and-frames}

Let O\_w be a DVR, with valuation v\_w, and put I\_w(b)=\{x in Frac(O\_w):v\_w(x) at least b\} for b in its valuation group. Write content\_w(B) for the largest such b for which B(Delta\_lambda,O\_w) is contained in I\_w(b)*Delta\_mu,O\_w. This convention also permits the later normalization v\_w(11)=1 without writing a fractional power of a uniformizer. The image is nonzero, and (R20.2) implies

\[
 v_w(c)-v_w(H_\mu)
                \ \le\ \operatorname{content}_w(B)\ \le\ v_w(c).
 \tag{R20.7}
\]

The left inequality is the image inclusion in the largest displayed lattice. The right inequality follows because c*p\_mu is an actual image and p\_mu is primitive in Delta\_mu. Notice that the inequalities have opposite uses; a lower divisibility certificate is not reversed into the upper inequality.

For several CM labels the full coefficient spaces and maps are tensor products. Put

\[
 C=\prod_\tau c_\tau,\quad H=\prod_\tau H_{\mu_\tau},\quad
 u_T=\prod_\tau\tau(\overline\pi_T)^{-nBr_\tau}.
 \tag{R20.8}
\]

Tensoring (R20.2), and then applying the literal scalar u\_T, gives

\[
 u_TC\Delta_{\boldsymbol\mu}
 \subset \operatorname{brAt}(\Delta_{\boldsymbol\lambda})
 \subset u_TC\nabla_{\boldsymbol\mu}
 \subset \frac{u_TC}{H}\Delta_{\boldsymbol\mu}.
 \tag{R20.9}
\]

Thus (R20.7) has v\_w(u\_T C) and v\_w(H) in this setting. In a finite frame whose source and target changes are integral GL matrices with integral inverses, the same bounds hold, since the Weyl lattices are stable under those changes. A nonintegral frame has its exact source and target lattice factors, as in K037(CR37)--(CR39); they are not silently suppressed. This statement is about standard coefficient lattices and their indicated transports, before stabilizer invariants or restriction to a prior image.

\subsection{4. Evaluation on the actual growing J14 and J56 weights}\label{evaluation-on-the-actual-growing-j14-and-j56-weights}

Use precisely R019's lambda and mu lists, and increase only lambda\_1 at the two designated CM labels by A=11\^{}r and A=121*11\^{}r. The target mu, and hence its natural partition and hook product, stay fixed. R019 proves that both cBr factors at the growing labels have reached their constant saturated values for every r at least13. All other labels have their unchanged baseline values. The T primes of R017 are distinct from11, so every factor in u\_T is an11-unit.

For eta\_i=mu\_i-mu\_n and n=N-1, the exact hook valuation used here is

\[
 v_{11}(H_\eta)=
 \sum_{i=1}^{n}v_{11}((\eta_i+n-i)!)
       -\sum_{1\le i<j\le n}
                        v_{11}(\eta_i-\eta_j+j-i).
 \tag{R20.10}
\]

This is the standard hook-product factorial identity, used with the same normalization as R011. Legendre's floor sum evaluates its first sum; the second sum has finitely many explicit integer factors. The script retains the actual literal RowPar check of the mu list from R019 before evaluating it. The results are:

\begin{longtable}[]{@{}
  >{\raggedright\arraybackslash}p{(\columnwidth - 8\tabcolsep) * \real{0.1579}}
  >{\raggedleft\arraybackslash}p{(\columnwidth - 8\tabcolsep) * \real{0.2105}}
  >{\raggedleft\arraybackslash}p{(\columnwidth - 8\tabcolsep) * \real{0.2105}}
  >{\raggedleft\arraybackslash}p{(\columnwidth - 8\tabcolsep) * \real{0.2105}}
  >{\raggedleft\arraybackslash}p{(\columnwidth - 8\tabcolsep) * \real{0.2105}}@{}}
\toprule\noalign{}
\begin{minipage}[b]{\linewidth}\raggedright
Host
\end{minipage} & \begin{minipage}[b]{\linewidth}\raggedleft
n
\end{minipage} & \begin{minipage}[b]{\linewidth}\raggedleft
Number of CM labels
\end{minipage} & \begin{minipage}[b]{\linewidth}\raggedleft
v11(H\_mu), one label
\end{minipage} & \begin{minipage}[b]{\linewidth}\raggedleft
v11(C), full tensor
\end{minipage} \\
\midrule\noalign{}
\endhead
\bottomrule\noalign{}
\endlastfoot
J14 & 3561 & 42 & 40694631292531 & 61490184399432942 \\
J56 & 3435 & 168 & 144279902901612 & 802279065217743524 \\
\end{longtable}

Therefore the content of the \textbf{whole} literal branch on the standard Weyl lattice satisfies

\[
\begin{array}{ll}
\text{J14:}&
 59781009885146640\ \le\operatorname{content}_{11}(\operatorname{brAt})
                                  \le61490184399432942,\\[2mm]
\text{J56:}&
 778040041530272708\ \le\operatorname{content}_{11}(\operatorname{brAt})
                                  \le802279065217743524.
\end{array}
 \tag{R20.11}
\]

Valuations here are normalized by v(11)=1. The bounds hold at each coefficient-field place above11 with the standard integral tensor model. They are independent of r. Their intersection with the already certified actual fixed-T families is r=13 through55 for J14 and r=15 through90 for the K033 J56 growing witnesses; the branch assertion itself holds for all r at least13. No enlargement of the R018 size intervals follows.

This proves that the normalization/content uncertainty of the whole terminal branching map is bounded independently of the growing first weight. It does not prove that the branch restricted to the actual preceding theta image has bounded content: that image can enter a deeper congruence sublattice or even the branch kernel.

\subsection{5. A complete small example with unequal Smith factors}\label{a-complete-small-example-with-unequal-smith-factors}

Take lambda=(2,1,0), mu=(2,0), and write x\_i and y\_i for the first two entries of row i in a three-row matrix. Then

\[
 p_\lambda=x_1(x_1y_2-x_2y_1),\quad
 \xi=x_1(x_1y_3-x_3y_1),\quad c=4,\quad\kappa=1.
 \tag{R20.12}
\]

The target is Sym\^{}2 in x\_1,x\_2, so projVV is the identity after extracting the coefficient of y\_3 with x\_3-degree zero. Consequently B is four times that extraction. Every source Weyl polynomial is integral, hence

\[
 B(\Delta_{(2,1,0)})\subset
       4\mathbb Zx_1^2+4\mathbb Zx_1x_2+4\mathbb Zx_2^2.
 \tag{R20.13}
\]

Let D\_rs=sum\_j X\_rj partial/partial X\_sj, with divided powers as in the source. The following three source vectors belong to Delta\_lambda:

\[
 \xi=D_{32}p_\lambda,\qquad
 \xi'=D_{21}^{(2)}\xi=x_2(x_2y_3-x_3y_2),\qquad
 v=D_{31}p_\lambda
       =2x_1x_3y_2-x_2x_3y_1-x_1x_2y_3.
 \tag{R20.14}
\]

Their images are 4x\_1\^{}2,4x\_2\^{}2,-4x\_1x\_2, respectively. Thus (R20.13) is equality: the image is 4*Nabla\_(2,0). The integral target Weyl basis is (x\_1\^{}2,2x\_1x\_2,x\_2\^{}2), directly from the divided powers of x\_1\^{}2. The image matrix is diagonal (4,2,4), up to ordering and a sign. Its Smith factors are exactly

\[
                                      2,\quad4,\quad4.
 \tag{R20.15}
\]

The highest-line scalar is4, while the full content is2. This is an explicit counterexample to identifying the whole integral branch with its primitive highest-line scalar, or with a homothetic copy of the target Weyl lattice. It is not a counterexample to R019, whose scope was expressly that highest line and the complex norm.

\subsection{6. Check scope and implications}\label{check-scope-and-implications}

The companion diagnostic constructs the exact divided-power Weyl lattices, applies the literal coefficient extraction and normalized Fischer projection, and computes their full small Smith matrices. It covers70 small branching pairs and finds nine nonhomothetic images. The three-variable proof in section5 is complete without relying on those computations. The actual-host certificate computes the hook and cBr valuation bounds in section4; it does not enumerate the enormous actual matrices.

For DP\_MULT Q3e this isolates another actual terminal factor: the whole branch has a fixed finite content interval in the growing family. It leaves the prior-image restriction, class contraction, terminal detector lattice and full g\_j,e\_j,H comparison. For THETA3 G1, the same B is the branch in DO038's exact mass and real period equation; its nonzero image on a full Weyl module does not imply a nonzero restricted theta period. The polynomial-lattice statement lives over Z and its localizations, and its rational Schur functors have the lambda/K-theory interpretation already recorded at SGA6 VI3.3. No new motive, weight, or RH comparison is asserted.

Author and independent whole-proof readings are logged separately. This file does not claim a Lean build or a new proof of the accepted Weyl duality and characteristic-zero branching theorems.

\clearpage

\section{AM-K038 --- the actual sharp coupled G/Q row and its integral class matrix}\label{am-k038-the-actual-sharp-coupled-gq-row-and-its-integral-class-matrix}

Author: GPT-6 Astra Pro, ALP\_MACHINE /root/kummer\_extension. 22 November 2026 UTC; inherited reasoning effort is not separately exposed.

This finding constructs the two-edge coupled arithmetic row from the actual mixed-CM geometric kernels of K035--037 and the unchanged raw finite functions. It proves their finite level covariance directly. The resulting row belongs to the prescribed evGSD module and is the actual argument of PtilFF. An ordinary ColV row is not substituted for this coupled row.

The statement concerns the sharp m=2 family, with the actual ranks, characters, G/T primes, maximal level at29, and selected partitions already specified in K031/K032, R017, and LIFT26. Every admitted fixed tuple is allowed; its full finite Q set and its entire prescribed Y quotient are retained. This is a constructed arithmetic replacement family on the mixed host. It does not assert an instance of the incompatible old Fields.Kq record, a uniform coefficient field for an unbounded collection of new tuples, or a numerical DP44 gain.

\subsection{1. Data and exact source interfaces}\label{data-and-exact-source-interfaces}

Use the polarized host K\_h and A\_h of K037(CR1), with h=1 or4. Put \[
 a=k,\quad d=t+1,\quad r=a+1,\quad
 b=a+d=r+t,\quad c=a+2d,\quad n_T=c-1=r+2t.
 \tag{CA1}
\] The two concrete rank quadruples (a,d,b,c) are (42,1760,1802,3562) and (168,1634,1802,3436). All three column ranks a,b,c are even. The terminal row has rank n\_T. Keep the actual sharp parameter rS=0. On the coupled chain m2=0 the source forces cOf(0)=1, cOf(1)=0 and h0=h1=1; there are no preceding uncoupled matrix edges.

At every G or Q place the local field is the rational q-adic field, with q a uniformizer. For G this is part of the chosen degree-one split prime data. For Q it follows from literal GoodPl: DegOne requires prime residue norm and q outside the square of the prime. The tuple is tame and unramified outside GoodPl. These are the definitions in \nolinkurl{Tuples.lean64–67,112–114,128,} not an inference from residue degree alone.

Write O for this local valuation ring and kappa for its residue field. Let psi have conductor O and put psi0(z)=psi(q\^{}\{-1\}z). Its finite argument ranges are those of the prescribed LocCst. At G let chi0,chi1 be the two actual residue characters in Ch.chi; the honest splitting characters eta\_j restrict to chi\_j\^{}\{-1\}. At Q let chi be the residue character of the fixed tuple; both eta\_j restrict to chi. These exact conventions are LIFT26(GS7a), LIFT23, and FS1004--1013. No character is conjugated implicitly.

Here and below FS denotes the recovered primary file \nolinkurl{Formalism/LocalToGlobal/Columns/FiniteSumNonzeroLargeDepth.lean.} The other paths in this table have the same recovered source root.

\begin{longtable}[]{@{}
  >{\raggedright\arraybackslash}p{(\columnwidth - 2\tabcolsep) * \real{0.5000}}
  >{\raggedright\arraybackslash}p{(\columnwidth - 2\tabcolsep) * \real{0.5000}}@{}}
\toprule\noalign{}
\begin{minipage}[b]{\linewidth}\raggedright
Source locator
\end{minipage} & \begin{minipage}[b]{\linewidth}\raggedright
Exact interface used
\end{minipage} \\
\midrule\noalign{}
\endhead
\bottomrule\noalign{}
\endlastfoot
FS194--245 & brC, RowPt, tilQ, PtilFF, evRho and evGSD \\
FS270--323 & Cpl, CplKer, its two covariance equations, class-weighted comp and colCpl \\
FS350--445 & DecoF full Q frames, their action, frmRho and frmGSD \\
FS610--639 & ThDist, determinant/rational covariance, and kerOf as an unaveraged finite latSum \\
FS640--665 & inFr, phiT, and liftFr \\
FS684--715 & entB, Q parabolic lattice, lamQB and PhiQ \\
FS738--810 & sharp G flags, lamGBS and PhiGS \\
FS909--941,982--986 & Jat, lamRow, muAt, omPhi and PsiAt \\
FS1019--1026,1062--1083 & zero/last h, initial levels, prescribed middle groups and characters \\
FS1103--1167 & actual frames, XiCpl1, PhiTop, ftY and XiCpl2 \\
FS1180--1210,1220--1275 & Y transport, midGSD, ybCpl, ChKer and its exact Shape \\
FS1408--1418,1451--1454 & initial startC, coupled colC with evGSD codomain, and PerInv \\
\nolinkurl{Columns/RowForms.lean221–231,253–269} & actual rowTyG, rowTyQ, gluing and inverse detector convention \\
\nolinkurl{Columns/RowFormsSym.lean16–20;} Columns/LevGSym.lean11--13 & rPsymF and symIdx, the latter the identity for rS=0 \\
\nolinkurl{Columns/GenericUnipotentCharacter.lean526–557,603–620} & HPsym, chiSym, HQ and the bottom determinant character \\
\nolinkurl{KernelLattices/ThDistC.lean349–455} & exact Cp1Equiv/Cp2Equiv, cp1K/cp2K, and the chKer constructor \\
Columns/RecLevel.lean50--52 & arithmetic finite character and inverse adelic right character \\
\end{longtable}

The companion certificate records SHA256 hashes of the recovered source bytes actually checked here; their mirror/author provenance is in the existing source-recovery records. This proof invokes the already accepted coefficient-module results at K037's stated interfaces. There is no Lean execution.

Fix one number field M containing the actual finite character values and the all-order geometric jet field of K036 for these two edges and this fixed tuple. All equations below are algebraic over M after the common period and source-frame scale of K035 are removed. No effective degree, discriminant or height bound for M is asserted.

\subsection{2. Coordinate and honest frame rules}\label{coordinate-and-honest-frame-rules}

At a local place a decoration frame F changes to F kappa\^{}\{-1\} under a pure finite action. For a simultaneous initial, middle and terminal change kappa0,kappa1,kappaT, the first coordinates and the terminal coordinates therefore change as follows: \[
\begin{split}
 (X,y)&\longmapsto
       (\kappa_1X\kappa_0^{-1},\kappa_1y),\\
 (Z,\eta,y)&\longmapsto
       (\kappa_TZ\kappa_1^{-1},\eta\kappa_1^{-1},\kappa_1y).
\end{split}
\tag{CA2}
\] Here X has size b by a, Z has size n\_T by b, and eta is the final row, including the literal piT\^{}\{-1\} in FS1166. The terminal finite action on the full rank-c target is diag(kappaT,1), so that row is fixed. In particular \[
                 (\eta\kappa_1^{-1})(\kappa_1y)=\eta y.
                 \tag{CA3}
\] This is the exact phase cancellation, with its original minus sign inside ftY. At G/Q piT is a unit; it has not been removed or replaced.

The honest split Weil action is the one checked in LIFT26(GS7a): \[
 [\omega_j(g,h)f](X)
 =\eta_j(\det g)^{-1}\eta_j(\det h)
       |\det g|^{N_{j+1}/2}|\det h|^{-N_j/2}
                    f(h^{-1}Xg).
 \tag{CA4}
\] For integral frame changes as in(CA2), the extra factor is \[
              \eta_j(\det\kappa_{\rm source})
                        \eta_j(\det\kappa_{\rm target})^{-1}.
                        \tag{CA5}
\] All absolute values in(CA5) are one. At Q this is chi(det source)/chi(det target). At G it is chi\_j(det source)\^{}\{-1\}chi\_j(det target). These are also exactly the inverse-frame determinant factors of ThDist.Ok.det with FS1246--1253's rhoS/rhoT. They are necessary: raw polynomial covariance alone leaves the wrong determinant character.

At Q the middle frames in DecoF are full local integral frames, so (CA2) is exact. At G a chosen liftFr is determined only modulo the principal congruence subgroup. Two choices differ by I+qA. The G raw matrix factors below depend only on the residue matrix. The G Y coordinate changes by an integral vector, and eta is integral on ftY's support; hence its additive phase is unchanged. At the terminal Q place the arbitrary liftFr affects only Z, and PhiQ depends only on its residue and integral support. Thus the computations do not depend on a preferred liftFr.

\subsection{3. Individual-Y covariance at Q}\label{individual-y-covariance-at-q}

Let P\_\{r,s\} be the upper block parabolic of sizes(r,s). The actual initial group is GL\_a(O), the actual middle group is the inverse image of P\_\{r,t\}(kappa), and the terminal group is the inverse image of P\_\{r,2t\}(kappa). Write \[
 \bar\kappa_1=
 \begin{pmatrix}A_1&B_1\\0&D_1\end{pmatrix},\qquad
 \bar\kappa_T=
 \begin{pmatrix}A_T&B_T\\0&D_T\end{pmatrix}.
 \tag{CA6}
\] The arithmetic characters are \[
 \lambda_{0,Q}=1,\quad
 \lambda_{1,Q}(\kappa_1)=\chi(\det D_1)^{-1},\quad
 \lambda_{T,Q}(\kappa_T)=\chi(\det D_T)^{-1}.
 \tag{CA7}
\] For the middle group, this follows directly from lamQB and cOf(1)=0. For the terminal it is the inverse lamQP in RowForms230--231.

Write X=(A\_X;C\_X) and y=(z;w), with top block size r. The exact first raw factor, LIFT23(Q5), is \[
 F_Q(X,y)=
 {\bf1}_{X\in M_{b,a}(O),\,C_X\in qM_{t,a}(O),\,
                  z\in O^r,\,w\in qO^t}
                         \chi(\det[\bar A_X,\bar z]).
 \tag{CA8}
\] The character is extended by zero at singular determinants. The r columns are precisely the a columns of X and the last y column. The support is preserved by(CA2) and the stated groups. Modulo q its top square matrix changes to \[
       A_1[\bar A_X,\bar z]\operatorname{diag}
                                      (\bar\kappa_0^{-1},1).
\] Consequently, including zero values, \[
 F_Q(\kappa_1X\kappa_0^{-1},\kappa_1y)
   =\chi(\det A_1)\chi(\det\bar\kappa_0)^{-1}F_Q(X,y).
 \tag{CA9}
\] Multiplication by(CA5) for the first edge gives exactly lambda\_\{1,Q\}(kappa1)lambda\_\{0,Q\}(kappa0)\^{}\{-1\}.

Let H\_Q(Z) be the literal terminal PhiQ. Its support consists of integral Z with lower-left 2t by r block divisible by q and invertible leading r by r block; its value is chi of that leading determinant. The terminal factor is \[
 H_Q(Z)\,q^{-\mathsf M b}
              {\bf1}_{\eta\in q^{-\mathsf M}O^b}
                       \psi(-\eta y),\qquad \mathsf M\ge1.
 \tag{CA10}
\] Its Y parameter lies in q\^{}\{-M\}O\^{}b/q\^{}MO\^{}b. The raw support is preserved under(CA2), and its leading block changes to A\_T Z\_\{11\} A\_1\^{}\{-1\}. The eta support and phase are invariant by(CA3). Thus the raw factor changes by chi(det A\_T)chi(det A\_1)\^{}\{-1\}. Combining with(CA5) for the second edge gives \[
       \chi(\det D_1)\chi(\det D_T)^{-1}
           =\lambda_{1,Q}(\kappa_1)^{-1}
                              \lambda_{T,Q}(\kappa_T).
           \tag{CA11}
\] No Fourier summation or division by the size of the Y set was used. In particular(CA9)--(CA11) hold for every individual Y term.

The statement is independent of the chosen Y representative. In(CA8), adding q\^{}M O\^{}b to y changes neither integral support nor residue determinant; in(CA10), its pairing with q\^{}\{-M\}O\^{}b is integral. This also proves that arbitrary representatives in FS.repY are allowed.

\subsection{4. Individual-Y covariance at G}\label{individual-y-covariance-at-g}

For an upper unipotent a by a matrix u set \[
             \psi_U(u)=\psi_0\left(\sum_{i=1}^{a-1}u_{i,i+1}\right).
             \tag{CA12}
\] It is a group character because the indicated superdiagonal entries add under multiplication. The exact middle residue group is \[
 H^\sharp_{r,t}=
 \left\{\begin{pmatrix}\operatorname{diag}(u,1)&A\\0&D\end{pmatrix}
       :u\in U_a(\kappa),\ D\in GL_t(\kappa)\right\}.
 \tag{CA13}
\] This is the direct flag calculation LIFT22(CG4), not the generic ordinary ColV1 level. The initial residue group is U\_a. The terminal group has upper block form of sizes(r,t,t), with first diagonal block diag(u\_T,1) and other diagonal blocks D\_\{T,0\},D\_\{T,1\}. Its exact characters, from lamGBS and rowTyG(rPsymF), are \[
\begin{split}
 \lambda_{0,G}(u_0)&=\psi_U(u_0),\\
 \lambda_{1,G}(\kappa_1)&=\psi_U(u_1)\chi_0(\det D_1),\\
 \lambda_{T,G}(\kappa_T)&=
          \psi_U(u_T)\chi_0(\det D_{T,0})\chi_1(\det D_{T,1}).
\end{split}
\tag{CA14}
\] For rS=0, symIdx is the identity, the distinguished r-th column in the first r block is fixed, and chiSym omits its joining root. This gives(CA12), rather than a generic character on U\_r.

The first raw factor is zero unless \[
                  y=q^{-1}e_a\pmod{O^b}.
                  \tag{CA15}
\] On that class it is exactly \[
 F_G(X,y)=
 \begin{cases}
  \psi_U(u_X),&X\in M_{b,a}(O),\
                    \bar X=\binom{I_a}{0}u_X,\quad u_X\in U_a(\kappa),\\
  0,&\text{otherwise}.
 \end{cases}
 \tag{CA16}
\] Here e\_a uses zero-based indexing, so is the last vector of the first r block. Both statements are the explicit raw PhiGS/PsiAt calculation LIFT22(CG7--9). They are retained with no projection.

The group(CA13) fixes e\_a, so(CA15) is stable. Under(CA2) the matrix u\_X changes to u\_1 u\_X u\_0\^{}\{-1\}. The raw factor therefore acquires psi\_U(u\_1)psi\_U(u\_0)\^{}\{-1\}. The honest first factor(CA5) is chi0(det kappa0)\^{}\{-1\}chi0(det kappa1)=chi0(det D1). The result is precisely lambda\_\{1,G\}lambda\_\{0,G\}\^{}\{-1\}. This proves the first covariance also off its support.

For the terminal raw factor, write the first b rows of Z as B and the last t rows as Z\_2. Its support is \[
         Z\in M_{n_T,b}(O),\quad \bar Z_2=0,\quad
         \bar B\in H^\sharp_{r,t}.
 \tag{CA17}
\] If the blocks of B are u\_B,D\_B as in(CA13), its value is \[
 H_G(Z)=\psi_U(u_B)\chi_0(\det D_B)\chi_1(\det D_B)^{-1}.
 \tag{CA18}
\] The full terminal local factor multiplies this by 1\_\{eta in O\^{}b\}psi(-eta y). These are the literal PhiGS and ftY: PhiGS has a=b at j=2, its last inverse unitV is chi1(det B)\^{}\{-1\}, and det B=det D\_B on the support.

Right multiplication by kappa1\^{}\{-1\} preserves(CA17) and gives \[
 H_G(Z\kappa_1^{-1})
  =\psi_U(u_1)^{-1}\chi_0(\det D_1)^{-1}
                                      \chi_1(\det D_1)H_G(Z).
 \tag{CA19}
\] For the terminal left action let C\_T be its leading b by b block. It belongs to Hsharp\_\{r,t\}. The zero lower block in(CA17) shows that the new leading block is C\_T B, while the last t rows remain zero. Thus \[
 H_G(\kappa_TZ)
 =\psi_U(u_T)\chi_0(\det D_{T,0})
                           \chi_1(\det D_{T,0})^{-1}H_G(Z).
 \tag{CA20}
\] These are identities of functions, including all zero cases, because the inverse group actions preserve the same supports. The honest second frame factor is \[
             \chi_1(\det D_1)^{-1}
                  \chi_1(\det D_{T,0}\det D_{T,1}).
                  \tag{CA21}
\] Equations(CA19)--(CA21), together with(CA3), leave exactly lambda\_\{1,G\}\^{}\{-1\}lambda\_\{T,G\}. This is the full compact-group covariance, not merely its restriction to a diagonal torus or the identity cell of a quotient model.

Adding an integral vector to y leaves(CA15)--(CA16) unchanged and pairs integrally with eta in the terminal factor. Thus it is an identity on the actual G quotient q\^{}\{-1\}O\^{}b/O\^{}b.

\subsection{5. Rational transport, T, and the actual polynomial kernels}\label{rational-transport-t-and-the-actual-polynomial-kernels}

At T the raw factor is exactly phiT: it says that the linear map sends the first local lattice into the second and pulls its designated covector back to the first covector. Under rational isometries gamma0,gamma1 and x -\textgreater{} gamma1 x gamma0\^{}\{-1\}, this incidence condition is unchanged. The same statement holds on the second edge with gammaT extended to diag(gammaT,1). The literal tilQ/tilH extension and its piT coordinate give exactly this covector transport. The finite identity does not require erasing a character at T. Its prescribed initial type is the actual one of R017/DO035/LIFT24. Away from G,Q,T retain the fixed maximal vacuum and the actual unramified splitting; its compact covariance is already that of K035. In particular no principal29 refinement is introduced.

The Y set is the quotient of the two localized middle lattices specified in FS1125--1129. The CRT description is \[
 {\cal Y}(p)\simeq
 \prod_{G}q^{-1}O^b/O^b\ \times\
 \prod_Qq^{-\mathsf M}O^b/q^{\mathsf M}O^b .
 \tag{CA22}
\] It is finite. A rational gamma1 sends it bijectively to the quotient for gamma1 N; a pure finite decoration action leaves the global coset Y unchanged and changes only its coordinates as in(CA2). This is exactly the action in ybCpl. The source proof FS1180--1210 is the elementary identity gamma1(z-y)=gamma1z-gamma1y.

Define the geometric functionals J\_\{1;p0,p,Y\}, J\_\{2;p,Y,d\} by the finite Mumford translation and induction sums of K035--037, using the literal XiCpl1 and XiCpl2 coefficients just displayed. For an inadmissible Y set the two functionals equal zero. They are defined from algebraic sections, character phases, and finite sums before their analytic theta comparison. No new Schwartz vector is selected.

Let Phi\_J be K037's universal subD operator on the selected coefficient spaces. Define \[
 k_1(p_0,p,Y)=\Phi_{J_{1;p_0,p,Y}},\qquad
 k_2(p,Y,d)=\Phi_{J_{2;p,Y,d}}.
 \tag{CA23}
\] The first map goes from V\_0 to V\_1; the second goes from V\_1 to V\_2. Every V\_j is the source rhoVr polynomial model with its actual selected weight. Its integral lattice is the literal delAtr Weyl lattice, not a replacement by the dual Weyl lattice.

K036 gives exact rational transport of the geometric kernel, including its determinant half-rank scalar. K037(CR18) gives the matching selected polynomial action, with source rhoVr and target rhoVr after the subD contraction. Equivalently, substitute X -\textgreater{} gamma1 X gamma0\^{}\{-1\} in the geometric germ, use K036's actual determinant scalar, and apply the ordinary-derivative chain rule inside K037(CR12). The character exponents in(CR2)--(CR3) cancel that scalar and yield the two displayed polynomial actions. This is the exact rational covariance proved in K036--037; no quadratic elliptic period is inserted.

Combining that rational identity with the finite identities (CA9)--(CA11),(CA14)--(CA21) proves \[
 k_1(g_0p_0,g_1p,g_1Y)
   =\rho_1(g_1)\,k_1(p_0,p,Y)\,\rho_0(g_0^{-1}),
\tag{CA24}
\] and \[
 \boxed{\ k_2(g_1p,g_1Y,g_Td)
       =\rho_{\rm ev}(g_T)k_2(p,Y,d)\rho_1(g_1^{-1}).\ }
 \tag{CA25}
\] Equations(CA24) and(CA25) are exactly CplKer.equiv1/equiv2. Here rho1 is frmRho, including its arithmetic middle character, and \[
 \rho_{\rm ev}(\gamma_T,\kappa_T)
       =\lambda_T(\kappa_T)\,
                      \rho_2^\vee(\operatorname{diag}(\gamma_T,1)).
 \tag{CA26}
\] Indeed evRho takes the inverse of rowTy.inv's character, which is rowTy's original character. Equations(CA7),(CA14) are exactly that character. The zero extension outside the allowed Y set is equivariant because(CA22) is equivariant.

Only the middle group acts on Y in these two equations.

\subsection{6. The actual coupled class map and its terminal domain}\label{the-actual-coupled-class-map-and-its-terminal-domain}

Take the finite decorated class sets C0,C1 for the initial and sharp middle groups, and their fixed representatives p\_\{0,e\},p\_\{1,f\}. Let Gamma0e,Gamma1f be their actual finite stabilizers. Finiteness comes from the definite selfdual lattice-class argument of K036 and finite residue decorations. The odd terminal rank uses the same finite-class argument, whose proof does not require even rank. Even rank was needed only for the subsequent all-order period field transport on the two column edges.

The source module consists of its prescribed equivariant rows P\_e, with its integral lattice defined by intersection with delAtr as in K037(CR21)--(CR22). Define, for every terminal decoration d, \[
 \boxed{\ 
 P_{\rm cpl}(d)=
  \sum_{e\in{\cal C}_0}\frac1{|\Gamma_{0,e}|}
  \sum_{f\in{\cal C}_1}\frac1{|\Gamma_{1,f}|}
  \sum_{Y\in{\cal Y}(p_{1,f})}
       k_2(p_{1,f},Y,d)\,
       k_1(p_{0,e},p_{1,f},Y)P_e .
 \ }
 \tag{CA27}
\] There is no factor 1/\textbar Y\textbar. The scalar q\^{}\{-Mb\} of ftY is already inside k2. This is the actual CplKer.comp followed by the initial class operator in FS298--323. Since m2=0, no preceding kk edge occurs. The singleton initial chart is the source's startC inclusion.

This initial module is the actual MfV=ColV0, not only its finite Steinberg subspace. ColForms.lean93--115 gives its initial G residue group U\_a and arithmetic character psi\_U; RecLevel50--52 gives the inverse adelic right character. At T its finite level fixes the designated nonzero covector, with trivial character on that strict stabilizer. It is maximal at the other places, including 29 and the initial Q places. The ramified source subspace is \[
 {\cal M}_{\rm ram}
       =\bigcap_{P\mid T}\ker\left(\sum_{z\in\kappa_P^\times}
                                      \operatorname{diamond}_P(z)\right)
                                      \subseteq{\cal M}_{f}.
 \tag{CA27a}
\] \nolinkurl{ScaledRamifiedLattice.lean84–91,155–175} is exactly this definition: the norm is an unweighted sum, without an auxiliary character. Record.lean1811--1848 proves its literal transport to the adelic scalar-dilation action. Over C the finite cyclic quotient acts semisimply, so(CA27a) retains precisely the nontrivial dilation characters at each T place. Formula(CA27) is defined on all MfV and restricts to this actual Mram and its inherited integral lattice. There is no extra initial projector in this restriction.

\textbf{Theorem CA28.} Formula(CA27) is an M-linear map from the actual initial row module to the literal evGSD.EqFun module with target polynomial space V\_2 and its prescribed terminal finite character. It is independent of the chosen class representatives, Y representatives, and liftFr choices, with the natural row transport.

\textbf{Proof.} Equations(CA24) and(CA25) cancel the two rho1 factors in each composition. Changing a middle representative therefore only bijects its Y set, and conjugates its stabilizer, whose cardinal remains unchanged. Changing an initial representative has the same effect after applying the equivariance equation for P\_e. This proves representative independence without choosing trivial stabilizers. The earlier periodicity proofs and liftFr argument prove the other two independence assertions.

For a terminal action g\_T, (CA25) left-multiplies every summand of (CA27) by rho\_ev(g\_T). Hence \[
                   P_{\rm cpl}(g_Td)=\rho_{\rm ev}(g_T)P_{\rm cpl}(d).
                   \tag{CA29}
\] This is precisely the defining evGSD equation. The target polynomial space is already correct by K037's integral subD theorem, so no post hoc isotypic projection is used. All sums are finite and all entries lie in M, proving linearity and algebraicity. QED.

This theorem supplies the missing raw covariance input in K037§7 and the admitted coupled-row input of DO038. In particular the exact arithmetic terminal functional is \[
 \operatorname{PtilFF}(P_{\rm cpl},G)
   =\sum_{z\in{\cal C}_T}\frac1{|\Gamma_{T,z}|}
       \operatorname{pairBK}
                 \bigl(\operatorname{brC}P_{\rm cpl}(d_z),G(d_z)\bigr).
 \tag{CA30}
\] Its brC is unchanged, including brConst/kapBr and the specified conjugate piT power. This is FS215--218, not a pairing defined by division of an analytic period. The final column uses rhoVr, the detector rhoV; no transpose or conjugation scalar has been inserted. DO038's actual terminal class transport and top realization therefore apply to this row after their stated hypotheses are retained.

Throughout(CA27)--(CA30), retain R017's maximal29 level. Distinguish the literal MfV level from R017's coarser support realization. At T the former uses the strict covector stabilizer H\_T; R017 uses the larger projective-line stabilizer P\_T with its chosen nontrivial character. At G, R017 keeps the full maximal compact and realizes the finite Steinberg type by its matrix coefficient; literal MfV uses the inverse image of U\_a with right character psi\_U\^{}\{-1\}. R017(R17.9) gives stabilizer1 after the strict generic T decoration even with G still maximal, and stabilizer mu28 for the coarser projective-T/full-G support. No principal29 refinement is made in any of these realizations.

Here is the exact refinement rule. For compact levels H contained in P of index I, the strict classes above a projective class with stabilizer Gamma\_P are the orbits of Gamma\_P on P/H. If their stabilizers are Gamma\_s, orbit-stabilizer gives \[
 \sum_s\frac1{|\Gamma_s|}=\frac I{|\Gamma_P|},\qquad
 v_H=\frac{v_P}{I},\qquad
 \sum_s\frac{v_H}{|\Gamma_s|}
                      =\frac{v_P}{|\Gamma_P|}.
 \tag{CA30a}
\] For the T-only refinement of R017, with G still maximal, the map from mu28 into P\_T/H\_T is injective. Its action is free. Put I\_T=product\_\{P\textbar T\}(q\_P-1). There are I\_T/28 intermediate classes, each with trivial stabilizer. The additional literal G refinement has index \[
 I_G=\prod_{P\in G}
       [GL_a(\kappa_P):U_a(\kappa_P)].
 \tag{CA30c}
\] Because each intermediate stabilizer is trivial, every such class refines into I\_G classes. Thus the full initial strict MfV support has I\_T I\_G/28 classes, each with stabilizer1, and volume v\_P/(I\_T I\_G) per class, where P here is R017's projective-T/full-G support level. Their total mass is v\_P/28. For norms the T character has modulus one, while the actual finite G matrix coefficients have average square1/d\_G by finite Schur orthogonality. Together with the real factor 1/d\_rho this recovers exactly R017(R17.14); it does not assign constant magnitude to the G coefficients.

It does not make the raw theta summands identical. The raw incidence kernel is H\_T-invariant and need not be nu\^{}\{-1\}-eigen under P\_T. Denote its actual right input character by nu (the inverse of R017's nu\_P if that notation is used). On this component coarsening the T integration is the exact identity below, with P=P\_T, H=H\_T and I={[}P\_T:H\_T{]}: \[
 \int_P F(gp)\theta_\phi(gp,h)\,dp
   =v_P F(g)\theta_{\Pi_{\nu^{-1}}\phi}(g,h),\qquad
 \Pi_{\nu^{-1}}\phi=
      \frac1I\sum_{z\in P/H}\nu(z)\omega(z)\phi .
 \tag{CA30b}
\] The character is trivial on H, and the chosen H-fixed finite vector makes this formula well-defined. Change variables in the sum to check that the average is nu\^{}\{-1\}-eigen; its factor 1/I is retained. Formula(CA30b) is an identity obtained from the original integration, not an extra Schwartz replacement. Thus(CA27) always uses the actual strict stabilizers and kernels. Using R017's projective scalar requires its actual transported nu-weighted kernel sum as well. Scalar masses alone do not justify replacing that sum by one representative kernel. At G the analogous coset decomposition retains the finite Steinberg matrix-coefficient weights: the input transforms by psi\_U\^{}\{-1\} and the raw source kernel by psi\_U, so their product descends to K\_G/J\_\{0,G\}. Summing that product over the cosets, with measure vol(J\_\{0,G\}), is exactly its integration over K\_G. It does not replace the raw kernel by an unweighted average. This is the finite G normalization already retained in R017/DO036/LIFT20.

\subsection{7. Exact integral image containment and the finite matrix gate}\label{exact-integral-image-containment-and-the-finite-matrix-gate}

Fix a finite place w of M and write R=O\_\{M,w\}. Use the actual transferred delAtr lattices L0e,L1f,L2d. All are free over R after localization. For each matched triple let j1(e,f,Y;w) and j2(f,Y,d;w) be the fractional ideals generated by the finitely many ordinary derivatives in the exact local frames of K037(CR39) for the selected multidegree on that edge. Include their geometric phase/average coefficients, their cD factors, and the determinant factors of those frames. These are fixed ideals, not the content of an ambient arbitrary integral function.

The integral Weyl-to-Weyl substitution theorem K037(CR8) gives \[
 \begin{split}
 k_1(L_{0,e})&\subseteq j_1(e,f,Y;w)L_{1,f},\\
 k_2(L_{1,f})&\subseteq j_2(f,Y,d;w)L_{2,d}.
 \end{split}
 \tag{CA31}
\] Consequently \[
 \boxed{\ 
 P_{\rm cpl}({\cal L}_0)_d\subseteq
   \left(\sum_{e,f,Y}
       \frac{j_2(f,Y,d;w)\,j_1(e,f,Y;w)}
            {|\Gamma_{0,e}|\,|\Gamma_{1,f}|}\right)L_{2,d}.
 \ }
 \tag{CA32}
\] This statement follows by composition and the triangle inequality for fractional ideals. It does not require an inverse Gram matrix. In valuation form, its lower bound is the minimum of v\_w(j1)+v\_w(j2)-v\_w(\textbar Gamma0\textbar)-v\_w(\textbar Gamma1\textbar) over the indicated terms. Since(CA29) already holds, intersecting the right side with the equivariant target space supplies the actual target row lattice.

When all G and Q residue characteristics differ from11, their raw coefficients are roots of unity, zero, or q\^{}\{-Mb\} times roots of unity. They are integral, and their nonzero scalar factors are 11-units. Thus this raw finite assembly introduces no new 11-denominator beyond the displayed geometric jet/frame ideals and actual stabilizer factors. This conditional statement excludes any tuple having a Q place over11; the proof of covariance itself allows such a tuple. In R017's generic initial strict classes the stabilizer factor is1; in its projective realization28 is an11-unit, but the weighted average(CA30b) has the displayed factor1/I. No claim is made that every intermediate or terminal stabilizer has order prime to11.

For complete precision, choose R-bases of the initial invariant row lattice and target invariant row lattice. Their existence follows from the saturated-intersection description and the DVR property. Formula(CA27), with its literal subD matrices, gives a finite matrix A\_cpl over Frac(R). Formula(CA30) adds the fixed branch and terminal pairing matrix. The least entry valuation of A\_cpl is the exact content of its image in the chosen target lattice; its minors give its other determinantal ideals. On an already produced input-image lattice, use that lattice's basis instead. The preliminary containment(CA32) need not be sharp: Y summation, class summation, and terminal branching can increase valuations or annihilate an image. No saturation assertion follows from nonzero analytic theta lift alone.

The gate is therefore a specific finite algebraic image matrix on fixed arithmetic lattices, with all its entries defined by the actual geometric kernels. It is no longer a missing unspecified coupled covariance or terminal-domain identification. Enumeration and a useful uniform bound for this matrix are not supplied here.

\subsection{8. Analytic comparison, complete type, and fixed highest vectors}\label{analytic-comparison-complete-type-and-fixed-highest-vectors}

Let v0,v1 be the actual strict-covector finite level volumes used for the classes in(CA27), with compact archimedean Haar of mass one. Projective nu-isotypic realizations instead use the precise coarsening(CA30a)--(CA30b). Let D\_\{j,tau\} be the two prescribed selected degrees, and let s\_tau=lambda\_tau/cD\_tau be K035's literal source frame scale. Set \[
 {\cal P}_{01}
   =v_0v_1 C_\theta^{14h(ab+bc)}
            \prod_\tau s_\tau^{D_{0,\tau}+D_{1,\tau}}.
 \tag{CA33}
\] Keep the same actual finite tensor as LIFT26: sum\_Y phi1Y tensor phi0Y. Its pointwise identity(GS7a--c) proves that the moving-lattice Y sum is the action of this one fixed tensor; no h-dependent projection is inserted.

Apply K037's complete coevaluation construction at both real types to this tensor. The coevaluation tensors are the canonical sum v\_i tensor p\_i with pairBK-dual bases. Class integration and pointwise coevaluation give, on the actual terminal restriction, \[
          {\cal T}_{\rm cpl}(P;v)
                      ={\cal P}_{01}\,f_{P_{\rm cpl},v}.
                      \tag{CA34}
\] To check the coefficient directly, first decompose both compact adelic quotients into their actual classes. Their masses are v0/\textbar Gamma0\textbar{} and v1/\textbar Gamma1\textbar. The complete source coevaluation contracts to P\_e; the complete middle coevaluation contracts the resulting V1 vector with the second subD kernel. The remaining sum is precisely(CA27). Formula K035(PF28) supplies the two common period/degree factors in(CA33), independent of Y and of the class. This proves(CA34) without rescaling the arithmetic row.

Complete coevaluation does not create a factor 1/dim V\_j. For a single source coefficient line, compact Schur integration does give that dimension factor, as in K037. In particular the fixed highest Chern creation vectors of LIFT25/LIFT26 are not silently identified with the source's unit-normalized pureTensorO: \[
 \operatorname{pureTensorO}(P_j)
          =\Lambda_j\,r_{P_j},\qquad
 \Lambda_j=\prod_\tau\lambda_\tau^{D_{j,\tau}}\ne0.
 \tag{CA35}
\] Replacing the two Chern vectors by the literal source vectors multiplies the two-edge operator by Lambda0 Lambda1 and its squared norm by \textbar Lambda0 Lambda1\textbar\^{}2. The further single-line Schur and R017 input normalizations remain the explicit factors of LIFT25/LIFT26; they are not absorbed into(CA27) or(CA33). These complex period factors are not assigned finite valuations.

For the actual normalized highest vectors, one can make the relation to the fixed detector exact. Let h\_j be the norm-one highest vector in the rhoV coefficient model and let p\_j be its corresponding dual functional with p\_j(h\_j)=1. The latter is a lowest vector for the dual upper Borel, as in LIFT25(AR6), not the unrelated highest vector of that dual representation. The real polynomial of LIFT25 is iota\_j(p\_j,h\_\{j+1\}), with its explicit Chern/unit-creation conversion(CA35). Compact Schur averaging of h\_j tensor p\_j equals the canonical coevaluation divided by d\_j=dim V\_j: its average is an invariant endomorphism, hence scalar, and its trace is1. Thus, for the same fixed coupled finite tensor and every initial arithmetic row P, \[
 {\cal T}_{\rm cpl}(P;h_2)
    =d_0d_1\,{\cal B}_{\rm source,highest}(f_{P,h_0})
    =d_0d_1\Lambda_0\Lambda_1\,
                   {\cal B}_{\rm Chern,highest}(f_{P,h_0}).
 \tag{CA36}
\] This argument is independent of the finite datum on each summand and so commutes with its prescribed full Y sum. It retains all factors in(CA33) independently. It does not assert that a fixed highest tensor itself is the full coevaluation tensor.

Formula(CA34) identifies the canonical complete-type analytic map with the newly constructed arithmetic row. It does not prove that the nonzero unrestricted full G2 output in LIFT26 has nonzero restriction to this terminal decorated subgroup, or survives brC. That is the actual image/terminal detection question in(CA30).

\subsection{9. The empty-Q reference member and family scope}\label{the-empty-q-reference-member-and-family-scope}

The reference tuple is the identity in tupGrp: Tuples.lean118--120. Its Qprime is empty by definition, and the actual source realization proves this at \nolinkurl{TrueColumnChoicesHeckeCharReferenceEqOne.lean273–279.} Thus Q=empty is an admitted reference branch; it retains the stage G and T characters. Equations(CA22)--(CA34) then simply omit the Q factors. The sharp middle group(CA13) is still the actual coupled middle group and need not equal ordinary ColV1.

This reference branch does not remove the nonreference tuples from the source comparison. DataParams₂.lean108's untwisted condition sets Aj and F to zero at every place whose muAt=mu0. LinksChar.lean386--402 requires every nonzero-index place to lie among the twist places and its tuple there to differ from mu0. Therefore a reference-only example supplies one actual operator and its costs; it does not prove the DP44 family inequalities or the gains carried by the nonreference rows.

The proof here works for every fixed admitted tame tuple with its full Q set. It also works for every fixed word for which the stated two-edge ranks, characters, selected weights and finite T data hold. It does not assert that the explicitly checked finite word/size ranges of R018 are unbounded.

\subsection{10. Checks and residual scope}\label{checks-and-residual-scope}

The companion finite certificate tests the two block covariance identities, including nonzero and zero-support cases, over small prime fields with exact formal character phases. It separately tests the complete ftY pairing modulo q\^{}\{2M\}, the sharp finite group character law, the double class/stabilizer contraction, and the source anchors. These are finite exact form checks; the proofs above establish the arbitrary-rank identities.

New result: the unchanged sharp raw G/Q family supplies actual Cp1Equiv/Cp2Equiv, hence a geometric arithmetic coupled map on the actual Weyl lattices with exact evGSD and PtilFF target. The coefficient matrix, stabilizer weights, common periods and creation scales are specified.

Still required for a numerical improvement: the restricted successive-image content of that matrix, the terminal branch image of the actual theta row, effective field/global height costs, and the complete source family comparison. None is deduced by replacing the image with an ambient integral module or by normalizing away a period, class mass or hook factor.

\clearpage

\section{ALP-LIFT27: injectivity of the unchanged sharp coupled operator on the entire initial ramified type}\label{alp-lift27-injectivity-of-the-unchanged-sharp-coupled-operator-on-the-entire-initial-ramified-type}

Author: GPT-6 Astra Pro, /root/lift\_accounting, ALP\_MACHINE. Date: 22 November 2026 UTC. Inherited reasoning effort is not exposed.

This strengthens LIFT26 from its explicitly chosen R017 input to every vector in the actual initial ramified automorphic row space. The new local point is that finite Steinberg type was unnecessary at G. A direct unitary Godement--Jacquet bound makes the whole-Schwartz quotient map holomorphic on every permitted initial constituent and its actual first theta image. At T, the literal unweighted norm operator is the trivial character projector, so all remaining dilation characters are handled by LIFT24. The source's complete coevaluation operator is also compared exactly with the previously checked highest-kernel operator.

The prescribed finite G/T/Q functions and full coupled Y sum are unchanged. No local Whittaker-genericity, local cuspidality, temperedness, initial finite Steinberg projection, extra finite translation, or new real rotation is assumed. The conclusion is before the literal terminal brC, tilQ, and PtilFF restriction. Injectivity of the unrestricted terminal row does not imply injectivity of that restriction.

\subsection{1. Fixed spaces, initial conditions, and theorem}\label{fixed-spaces-initial-conditions-and-theorem}

Use either actual mixed CM host K031 or K032, the common edge characters and selected real representations of LIFT26, and ranks \[
 (a,d,b,c)=(42,1760,1802,3562)
       \quad\hbox{or}\quad(168,1634,1802,3436),
 \qquad b=a+d,\quad c=b+d.
 \tag{MR1}
\] Write \(G_0=U(I_a),G_1=U(I_b),G_2=U(I_c)\). These groups are positive at every real place. Fix exactly the raw sharp coupled finite tensor \(\Psi_f\) of LIFT26(GS3),(GS6)--(GS7c), the source-style bilinear theta convention, and the unit-normalized source creation vectors of LIFT25(AR27)--(AR30). Their Chern-coordinate versions differ by the fixed nonzero scalar \(\Lambda_0\Lambda_1\).

At a selected split degree-one G place put \[
 K_a=\mathrm{GL}_a(O),\quad K_a^1=\ker(K_a\to\mathrm{GL}_a(\kappa)),
 \quad J_G=\operatorname{red}^{-1}U_a(\kappa).
 \tag{MR2}
\] The initial right condition is \(\psi_U^{-1}\) on \(J_G\), for the literal source generic upper-unipotent character. There is no restriction to the finite Steinberg representation.

At a selected T place let \(u\ne0\) be the actual residue covector, and \[
 H_a=\{g\in K_a:u\bar g=u\},\qquad
 J_a=\{g\in K_a:u\bar g=s(g)u,\ s(g)\in\kappa^\times\}.
 \tag{MR3}
\] The input is \(H_a\)-fixed and killed by the unweighted sum of the \(\kappa^\times\) dilation operators. At Q and all other finite places, including29, it has the actual maximal compact condition of LIFT26. At infinity its right vector is the fixed unit highest vector \(h_0\) of the selected compact representation \(W_0\). Let \(\mathcal H_{\mathrm{ram},h_0}\) be this finite-dimensional automorphic subspace. All rational stabilizer relations, including central roots of unity, are retained. An incompatible tuple of local characters gives a zero summand; it is not artificially admitted.

Let \(\mathcal B\) be the same full finite Y-sum of the two actual highest-kernel theta lifts as in LIFT26. Then \[
 \boxed{\ker(\mathcal B|_{\mathcal H_{\mathrm{ram},h_0}})=0.}
 \tag{MR4}
\] The canonical complete-coevaluation operator from actual coefficient rows to full G2 automorphic functions is injective as well, with the exact factor in section7. Section2 identifies these initial conditions with the source's actual \(M_{\mathrm{ram}}\). The final arithmetic comparison still retains its independently determined periods and class volumes.

\subsection{2. Identification with the literal initial ramified space}\label{identification-with-the-literal-initial-ramified-space}

All Lean locators are relative to work/research/alpoge\_source, commit 7cbf7a0c32fca1badf05931879a37c3f8a2b03b4.

\nolinkurl{Formalism/LocalToGlobal/Columns/ColForms.lean93–115} defines colLev0, MfV, and equivariant forms. At the exact first index \(n=k\), \nolinkurl{GenericUnipotentCharacter.lean437–451} gives \(HX(k,k,t,0)=U_k\) and \(\operatorname{lamX}=\psi_U\). RecLevel.lean50--52 takes the inverse level character for the adelic reading. Thus(MR2) is the entire prescribed G condition.

\nolinkurl{Columns/ScaledRamifiedLattice.lean84–91} defines NmG as the unweighted sum of diamondOperatorOnForms′ over residue units. Lines122--127 define diaOp and NmOp; lines155--169 give \[
 M_{\mathrm{ram}}=\bigcap_{P\in v_T}
       \ker\!\left(\sum_{s\in\kappa_P^\times}\operatorname{dia}_P(s)\right).
 \tag{MR5}
\] The decoration fixes the nonzero covector, not merely its line, giving \(H_a\). Scaling it gives \(J_a/H_a=\kappa^\times\). \nolinkurl{Columns/Record.lean1811–1848} proves transport of precisely that sum under the faithful adelic reading and the exact Mram equivalence. There is no determinant character in that sum. The stage splitting characters are unramified at T and hence trivial on compact determinants. Changing the integral frame conjugates these conditions.

For a unitary representation \(\pi\), the quotient acts unitarily on \(\pi^{H_a}\), so \[
 \pi^{H_a}=\bigoplus_{\nu\in\widehat{\kappa^\times}}\pi^{J_a,\nu},
 \qquad
 \ker\sum_s\pi(\delta(s))
             =\bigoplus_{\nu\ne1}\pi^{J_a,\nu}.
 \tag{MR6}
\] Here \(\delta(s)\) is any lift of s; its action on \(H_a\)-fixed vectors is independent of the lift. Indeed the character projectors are \((q-1)^{-1}\sum_s\nu(s)^{-1}\pi(\delta(s))\); the unweighted sum is \((q-1)\) times the trivial projector. Conditions at different T places commute. Thus the simultaneous kernel is the direct sum over tuples of nontrivial characters. This decomposition is used in the proof; it is not an extra projector in the theta construction.

K037(CR26),(CR32)--(CR33) gives the actual coefficient realization \(P\mapsto f_{P,h_0}\), including arithmetic stabilizers. It is injective: at each finite class a zero scalar function means its coefficient vanishes on the orbit of \(h_0\); that orbit spans the irreducible \(W_0\). With the actual class-row metric, \[
 \|f_{P,h_0}\|^2
   =\frac{v_0}{d_0}\sum_x|\Gamma_x|^{-1}\|P_x\|^2,
 \quad d_0=\dim W_0,\quad v_0=\operatorname{vol}(U_0)>0.
 \tag{MR7}
\] The level volume, stabilizers, source dimension, and arithmetic coefficient lattice are not replaced by unitary-coordinate data. Every actual \(M_{\mathrm{ram}}\) row therefore gives an element of \(\mathcal H_{\mathrm{ram},h_0}\).

Here \(U_0\) is the literal strict covector/unipotent level, not R017's coarser projective-T/full-G support. The exact refinement is important even though it does not affect injectivity. If \(H\subset P\) has finite index I and a P-class has rational stabilizer \(\Gamma_P\), then its H-classes are the \(\Gamma_P\)-orbits on \(P/H\). For their stabilizers \(\Gamma_s\), \[
 \sum_s|\Gamma_s|^{-1}=I/|\Gamma_P|,\qquad
 v_H=v_P/I,\qquad
 \sum_s v_H/|\Gamma_s|=v_P/|\Gamma_P|.
 \tag{MR7a}
\] This is orbit-stabilizer and Haar refinement, with no freeness assumption. In R017's particular generic support, the T refinement has \(I_T/28\) free classes and the further G refinement gives \(I_TI_G/28\) classes of stabilizer1, exactly K038(CA30a),(CA30c). The finite G coefficient squares retain their Schur average \(1/d_G\); they are not all1.

Coarsening raw theta integration also retains a vector average. If the input has actual right character \(\nu\) on P, the H-fixed kernel satisfies \[
 \int_P F(gp)\theta_\phi(gp,h)\,dp
   =v_P F(g)\theta_{\Pi_{\nu^{-1}}\phi}(g,h),\qquad
 \Pi_{\nu^{-1}}\phi=I^{-1}\sum_{z\in P/H}\nu(z)\omega(z)\phi .
 \tag{MR7b}
\] This follows directly by decomposing P into H-cosets; changing variables proves the displayed average is \(\nu^{-1}\)-eigen. It is the average supplied by the original integral, not a replacement of the prescribed Schwartz vector. Thus equal total class masses do not make raw kernel summands identical. All arguments below use the genuine integrals at their actual strict levels. If expressed in R017's projective realization, they use (MR7b) with its factor \(1/I\), as in corrected K038(CA30b). The analogous G refinement retains its finite matrix-coefficient weights. No principal29 level is inserted.

\subsection{3. Uniform local holomorphy for an arbitrary unitary input}\label{uniform-local-holomorphy-for-an-arbitrary-unitary-input}

Let F be nonarchimedean of characteristic zero and residue cardinality q. Normalize additive Haar by \(\operatorname{vol}M_n(O)=1\) and multiplicative Haar by \(\operatorname{vol}K_n=1\). Put \(\kappa_n=\prod_{i=1}^n(1-q^{-i})\), so \(dg=\kappa_n^{-1}|\det g|^{-n}dX\). For an irreducible smooth unitary \(\pi\) of \(\mathrm{GL}_n(F)\), a unitary character \(\eta\), a Schwartz function f on \(M_n(F)\), and a smooth matrix coefficient m, use \[
 Z(z,f,m;\eta)=\int_{\mathrm{GL}_n(F)}
       f(g)m(g)\eta(\det g)^{-1}|\det g|^{z+(n-1)/2}\,dg.
 \tag{MR8}
\] Every such integral is holomorphic in \[
                          \Re z>(n-1)/2.
 \tag{MR9}
\]

\textbf{Proof.} Smooth coefficients of a unitary representation are bounded. A smooth dual functional is represented by a Hilbert vector: take a compact subgroup fixing it, use Riesz on the finite-dimensional invariant space, and compose with the orthogonal compact average. Cauchy--Schwarz then bounds the full coefficient.

The checked determinant integral LIFT16(S8), with these measures, is \[
 \int_{M_n(O)}|\det X|^u\,dX
       =\prod_{i=1}^n\frac{1-q^{-i}}{1-q^{-u-i}},
                       \qquad\Re u>-1.
 \tag{MR10}
\] This accepted matrix-Iwasawa identity and its convergence range are used, not claimed as a new general integration theorem. Substituting \(u=\Re z-(n+1)/2\) gives precisely(MR9). Any Schwartz support is contained in \(\varpi^{-N}M_n(O)\); its determinant integral differs by a finite scaling factor. On every compact subset of(MR9), the same expression gives a locally uniform integrable majorant. Thus the zeta integral converges absolutely and holomorphically. ∎

The published all-irreducibles Godement--Jacquet ideal theorem, as restated by Jiang and Luo(2023), Theorem2.4(1)--(2), pp.313--314, says these integrals generate \(L(z,\pi\eta^{-1})\mathbf C[q^z,q^{-z}]\). The generator is therefore a finite Laurent-polynomial linear combination of such integrals. Consequently \[
 \boxed{L(z,\pi\eta^{-1})\text{ is holomorphic and nonzero for }
                                      \Re z>(n-1)/2.}
 \tag{MR11}
\] Holomorphy follows from the finite combination, and nonvanishing from the standard local expression \(1/P(q^{-z})\), \(P(0)=1\). No temperedness, local cuspidality, or local genericity is assumed.

Now take any unitary irreducible \(\pi_0\) of \(\mathrm{GL}_a(F)\) and unitary characters \(\eta_0,\eta_1\). Set \[
 z_*=(d+1)/2,\quad
 \pi_1=\pi_0\times(\eta_0\circ\det_d),\quad
 \pi_2=\pi_1\times(\eta_1\circ\det_d),\qquad d>a-2.
 \tag{MR12}
\] Both normalized inductions are irreducible and unitary by Tadic's TheoremD: combine the lists of unitary blocks of its factors and use TheoremD(i). The characteristic-zero appendix removes the historical unpublished dependency mentioned on p.339. No disjoint cuspidal support assumption is needed.

The standard parameter of irreducible normalized induction is the direct sum of its factors. The same multiplicativity used in LIFT24(MT27)--(MT30), now without a spherical assumption on \(\pi_0\), therefore gives \[
 L(z_*,\pi_1\eta_1^{-1})
   =L(z_*,\pi_0\eta_1^{-1})
                     \prod_{r=1}^{d}L(r,\eta_0\eta_1^{-1}).
 \tag{MR13}
\] The first factor is holomorphic and nonzero by(MR11). Each rank-one unitary factor at \(r\ge1\) is1 if ramified and otherwise \((1-\eta_0\eta_1^{-1}(\varpi)q^{-r})^{-1}\), so is finite and nonzero. In the actual G application the ratio is ramified, making that product equal1. The first-edge factor \(L(z_*,\pi_0\eta_0^{-1})\) is likewise holomorphic and nonzero. The actual rank margins are \[
                  z_*-(a-1)/2=(d-a+2)/2=860\ \hbox{or}\ 734.
 \tag{MR14}
\] These statements are uniform over all unitary initial constituents, including arbitrary depth-zero cuspidal support at G.

\subsection{4. The actual G quotient and raw norms on the entire generic type}\label{the-actual-g-quotient-and-raw-norms-on-the-entire-generic-type}

The whole-Schwartz construction LIFT24(MT10)--(MT16), at the holomorphic values just proved, gives target-equivariant balanced maps \[
 q_0:(\mathcal S(M_{b,a})\otimes\pi_0)_{G_a}\longrightarrow\pi_1,
 \quad
 q_1:(\mathcal S(M_{c,b})\otimes\pi_1)_{G_b}\longrightarrow\pi_2.
 \tag{MR15}
\] Its exact common-character action is \[
 \omega_\eta(g,h)f(X)=
  \eta(\det g)^{-1}\eta(\det h)
      |\det g|^{m/2}|\det h|^{-n/2}f(h^{-1}Xg).
 \tag{MR16}
\] The exponent at the continued point is \(z_*+(n-1)/2=m/2\), so these are the literal bilinear source coinvariants. MT12--MT15 continues finite-dimensional compact-fixed coordinates, not merely bidual functionals, and proves the full parabolic covariance. A nonzero map(MR15) is the actual small theta quotient by Mínguez(2008), Theorem1 and Tadic irreducibility. No map on an open-rank stratum is simply presumed to extend.

Let \[
 V_0=\pi_0^{J_G,\psi_U^{-1}},\qquad
                 a_q=\frac{|U_a(\kappa)|}{|\mathrm{GL}_a(\kappa)|}.
 \tag{MR17}
\] The unchanged first sharp raw factor is LIFT20(RG6): it is supported on integral matrices with residue \([I_a;0]u\), with value \(\psi_U(u)\). The unique allowed local coupled class is \(y=q^{-1}e_a\bmod O^b\), as in LIFT22(CG7)--(CG8). For every \(v\in V_0\), \[
                             q_0(f_U,v)=a_qs(v),
 \tag{MR18}
\] where s(v) is the identity parabolic-cell section of CG12, with s(v)(1)=v. To check this without Steinberg: the support forces the source variable into \(J_G\), with determinant residue1; the two generic characters cancel, leaving Haar volume times v. The target support is the identity cell and parabolic covariance gives all remaining values. Nothing here uses the cuspidal support of \(\pi_0\).

The section s(v) has the exact sharp middle character \(\lambda^\sharp\) of LIFT22(CG4)--(CG6): on the displayed \(\operatorname{diag}(U_a,1)\times\mathrm{GL}_{d-1}\) blocks it is \(\psi_U^{-1}\eta_0(\det h)\), and is trivial on the specified radical. This is the actual coupled middle level, not the different full generic ColV1 condition.

The terminal raw \(\Phi_y^\sharp\) is exactly LIFT21(TG7)--(TG8). Its top b-by-b block is invertible on its support. Thus its restriction to a fixed full-rank orbit has compact support, and the continued map agrees there with the ordinary compact integral. The calculation of LIFT22 gives \[
 q_1(\Phi_y^\sharp,w)(1)=v_{J^\sharp}w,\quad
 w\in\pi_1^{J^\sharp,\lambda^\sharp},\qquad
 v_{J^\sharp}=\frac{|H^\sharp|}{|\mathrm{GL}_b(\kappa)|}.
 \tag{MR19}
\] It uses only the specified compact character. Consequently \[
 \boxed{[q_1(\Phi_y^\sharp,q_0(f_U,v))(1)](1)
                                    =a_qv_{J^\sharp}v.}
 \tag{MR20}
\] This proves actual small-quotient occurrence and injectivity on all of \(V_0\), without a finite Steinberg restriction.

The exact support-cell norm calculation CG17--CG19 also only uses unitarity. Hence \[
 \|q_1(\Phi_y^\sharp,q_0(f_U,v))\|^2
  =\frac{q^b(a_qv_{J^\sharp})^2}
        {\mathcal J_{a,d}(q)\mathcal J_{b,d}(q)}\|v\|^2,
 \quad
 \mathcal J_{n,d}=|\mathrm{GL}_{n+d}(\kappa)/P_{n,d}(\kappa)|.
 \tag{MR21}
\] The raw squared norms remain \(q^{-ab}|U_a|\) and \(q^{-bc+b}|H^\sharp|\). If both raw vectors are unit-normalized, the squared multiplier is \[
 \frac{q^{ab}|U_a|}{|\mathrm{GL}_a|^2\mathcal J_{a,d}}\,
 \frac{q^{bc}|H^\sharp|}{|\mathrm{GL}_b|^2\mathcal J_{b,d}}.
 \tag{MR22}
\] These are exact local induced-model norms. The globally pinned row normalization is not thereby1, nor is an integral image lattice for arbitrary \(\pi_0\) determined.

\subsection{5. All nontrivial T characters and the other places}\label{all-nontrivial-t-characters-and-the-other-places}

For every nontrivial \(\nu\) in(MR6), checked LIFT24(MT6)--(MT8) applies to every unitary irreducible constituent with that type: \[
 \pi_0=\sigma_0\times\xi,\qquad
 \sigma_0\text{ unitary spherical of rank }a-1,\qquad
                       \xi|_{O^\times}=\nu.
 \tag{MR23}
\] The space \(\pi_0^{J_a,\nu}\) is one-dimensional. The proof uses Roche's tame type theorem, the finite parabolic invariant calculation, and Tadic's unitary blocks; the actual T primes satisfy the residual characteristic condition. It does not assume local cuspidality. An irreducible \(\pi_0\) cannot have two different nontrivial \(\nu\): its inertial support has exactly one nontrivial rank-one character, whose identity is unique by cuspidal-support uniqueness.

The unchanged first raw incidence kernel carries this line to the same \(\nu\)-line in \(\pi_1=(\sigma_0\times\eta_0\det_d)\times\xi\), by MT28. DO036's positive full-rank test applies uniformly to all \(H_a\)-fixed vectors, so it is nonzero on each of these lines. The second raw incidence kernel has exactly the graph value \[
 \lambda_T(\pi_0,\nu)w_1,\qquad
 \lambda_T(\pi_0,\nu)=
  \frac{L((d+1)/2,\sigma_1\eta_1^{-1})}{q^b-1}\ne0,\quad
                    \sigma_1=\sigma_0\times(\eta_0\det_d).
 \tag{MR24}
\] All Euler and graph/frame factors are MT29--MT33. Alternatively (MR11)--(MR13), at initial rank a−1, proves uniform holomorphy; the inserted factors are at1,\ldots,d and are unitary. The prescribed raw function is not replaced by its \(\nu\)-projector. Its actual target covariance preserves the character, so distinct characters in(MR6) cannot cancel. We retain each exact \(\lambda_T\), input norm, and compact coefficient normalization; complex nonvanishing is not a p-adic valuation statement.

DO037's Q calculation already proves injectivity of the full coupled Q quotient on every \(K_a\)-fixed vector, with its exact nonzero \(\Lambda_Q\). All other fixed finite basic vectors are nonzero on the corresponding spherical line by LIFT16. That result distinguishes relative splitting from absolute host ramification, including29. At infinity LIFT25 proves the actual consecutive-highest alignment, with source dual, dimensions, hooks, Gaussian norms and \(\Lambda_j\) conversion. These are carried checked findings used with their existing scope.

\subsection{6. Global injectivity and all automorphic multiplicities}\label{global-injectivity-and-all-automorphic-multiplicities}

The finite-level, fixed-real-type initial space has finitely many irreducible constituents and finite-dimensional multiplicity spaces. All are globally cuspidal and unitary since the positive groups are anisotropic; their finite local representations need not be cuspidal. For each constituent, DO036's first norm proof applies to every embedding and every nonzero vector of its actual initial type: DO36.15--17 already treats the entire inverse-generic G space, its T test treats all \(H_a\)-fixed vectors, and Q and basic tests are unchanged. This first edge has b\textgreater2a and uses GQT7.1 with convergent unfolding, not GQT11.4 outside its dimension window.

The first unramified theta parameter is still the original parameter plus the length-d \(\chi_0\) string. Thus the second global standard factor remains the one in LIFT26: \[
 L_E^S((d+1)/2,\operatorname{BC}_{\rm unr}(\pi_0)\chi_1^{-1})
            \prod_{r=1}^d L_E^S(r,\chi_0\chi_1^{-1}).
 \tag{MR25}
\] Only unramified parameters are used; no automorphic global base-change representation is assumed. The initial factor has the same absolute convergence margin(MR14). The boundary Hecke factor at1 is nonzero by the checked angular-type argument of LIFT14/LIFT26; absolute convergence at1 is not claimed. All fixed stage characters remain unchanged.

Local occurrence is now given at G by(MR15)--(MR20), at T by(MR23)-- (MR24), and elsewhere by the carried tests. The real converse is the explicit compact Schur integral. GQT11.4 applies on the second edge because b\textless c\textless2b and target Witt index is0: the anisotropic convergent case inside that theorem's dimension window. Therefore the second lift is nonzero on every actual embedding of every first constituent.

Apply exactly the multiplicity factorization of LIFT26(GS19)--(GS27): \[
       \Theta_\pi=A_\pi\otimes\bigotimes'_v q_{\pi,v},
                           \qquad\ker A_\pi=0.
 \tag{MR26}
\] Local Howe uniqueness gives the quotient factors; every-embedding nonvanishing makes the actual multiplicity map injective. No automorphic multiplicity-one hypothesis is used. All hypotheses have now been checked on the larger initial space.

Every local fixed coupled map \(q_{1,v}(1\otimes q_{0,v})(\Psi_v\otimes\cdot)\) is injective on the whole allowed initial type, by sections4--5. At basic places it is nonzero on a line, and at infinity it is nonzero on the fixed highest seed. Their restricted tensor product is injective, also on any subspace imposed by rational stabilizer compatibility. LIFT26(GS7a)-- (GS7c) identifies precisely this SAME fixed tensor with the literal moving-frame Y family, including \((\eta_0\eta_1^{-1})(\det h)|\det h|^d\). No middle-dependent projection has replaced it. On each initial constituent the full operator is consequently \[
                       (A_1A_0)\otimes C_{\pi_0},
                              \qquad\ker C_{\pi_0}=0.
 \tag{MR27}
\] It is injective. Distinct initial irreducible representations cannot coalesce after two edges, by the injectivity part of local Howe correspondence. Their target irreducible subspaces are orthogonal, so no cross-constituent cancellation occurs. This proves(MR4) for the whole ramified space, not merely for pure tensors or one class.

\subsection{7. Exact highest-kernel/complete-row normalization}\label{exact-highest-kernelcomplete-row-normalization}

This step is independent of the GJ extension. Use K037's compact metrics and selected Cauchy map \(\iota_j:W_j^\vee\otimes W_{j+1}\to\mathcal F_j\), normalized by the leading coefficient1 of its highest polynomial. For a coefficient row P and target seed v, K037(CR27) defines the actual complete coevaluation lift \[
 \mathcal T_j(P;v)=
   \sum_{i=1}^{d_j}
    \Theta_{j,\operatorname{pureTensorO}
                      (\phi_{j,f},\iota_j(p_i,v))}^{\rm bil}
                                         (f_{P,v_i}),
 \tag{MR28}
\] where \(v_i\) is an orthonormal basis and \(p_i\) its bilinear dual.

For any unit h and its unit dual \(h^\vee\), compact Schur averaging with Haar volume1 gives \[
 \int_{G_{j,\infty}}\rho_j(k)h\otimes\rho_j^\vee(k)h^\vee\,dk
                    =\frac1{d_j}\sum_i v_i\otimes p_i.
 \tag{MR29}
\] Indeed the integral is invariant and has trace1 under the endomorphism identification. The invariant line is the identity, whose trace is \(d_j\). This fixes coefficient and sign.

Apply(MR29) on every finite class of the genuine theta integral. The finite vector, level character, and class mass remain unchanged. Multilinearity gives \[
 \boxed{\mathcal T_j(P;v)=d_j\,
   \Theta_{j,\operatorname{pureTensorO}
                     (\phi_{j,f},\iota_j(h^\vee,v))}^{\rm bil}
                                             (f_{P,h}).}
 \tag{MR30}
\] For consecutive highest vectors h and v, the polynomial \(\iota_j(h^\vee,v)\) is exactly LIFT25's primitive principal-minor polynomial: K037's leading-minor substitution fixes its scalar. No source volume, creation norm, or finite mass is absorbed here.

Apply(MR30) twice to the same coupled finite tensor. The first complete-coevaluation output is a genuine \(W_1\)-coefficient family, so the second use has the correct middle representation and dual. The finite Y-sum commutes with both compact averages. Hence \[
 \boxed{\mathcal T_{10}^{\rm coev}(P;h_2)
                         =d_0d_1\,\mathcal B(f_{P,h_0}).}
 \tag{MR31}
\] Here \(\mathcal B\) uses source-normalized creation. For the Chern operator the right side is instead \(d_0d_1\Lambda_0\Lambda_1\mathcal B^{\rm Chern}(f_{P,h_0})\). These factors are all nonzero. Equations(MR7),(MR4),(MR31) prove injectivity of the complete raw coevaluation operator, valued in full G2 automorphic functions, on every initial \(M_{\mathrm{ram}}\) row.

For the actual restricted arithmetic class map constructed in K038(CA27)--(CA34), the precise comparison scalar is \[
 c_{10}=v_0v_1 C_\theta^{14h(ab+bc)}
     \prod_\tau(\lambda_\tau/cD_\tau)^{D_{0,\tau}+D_{1,\tau}}\ne0.
 \tag{MR32}
\] This is corrected K038(CA33), where \(v_0,v_1\) are the actual strict-covector level volumes at the two stages. They are not the coarser projective-T support volumes. Changing to a projective character realization requires both(MR7a) and the actual weighted kernel average(MR7b). Unequal chart volumes must retain their stipulated ratios, as in K037§4. K038(CA34) identifies the complete coevaluation family with its constructed arithmetic row \textbf{after the literal terminal tilQ restriction}; CA36 agrees with(MR31) after that restriction. Thus that actual row map is fully specified, but its injectivity does NOT follow from(MR4). The latter concerns the full G2-valued operator. Passing to a terminal subgroup or applying brC can still create a kernel. If an unrestricted arithmetic class realization is used instead, its own faithful all-class comparison would transfer (MR4) to characteristic zero. Such a comparison must not be silently substituted for K038's restricted one. No integral elementary-divisor or terminal branch-injectivity assertion is made here.

\subsection{8. Retained norms and the endpoint still to be detected}\label{retained-norms-and-the-endpoint-still-to-be-detected}

For every vector in this larger finite-dimensional ramified space the actual norm identity of LIFT26 remains \[
 \|\mathcal Bf\|^2
   =\sum_{\pi_0}c_{0,\pi_0}c_{1,\pi_1}
              \|(1\otimes C_{\pi_0})f_{\pi_0}\|^2,
                  \qquad c_{0,\pi_0},c_{1,\pi_1}>0.
 \tag{MR33}
\] The c's are actual Rallis and class/Haar factors, not arbitrarily chosen isometries. G retains(MR21)/(MR22), T retains(MR24), Q retains DO037's scalar, and source infinity retains AR27--30. The complete-row value(MR31) has squared norm \((d_0d_1)^2\) times(MR33); comparison with the row input also retains(MR7).

A positive least singular value exists on this fixed finite-level space by injectivity and finite dimensionality. No effective uniform global lower bound is claimed. This does not determine a DP44 gain, a saturated integral prior-image lattice, or finite valuations of complex periods.

For the bottom problem, the precise improvement is that a nonzero initial Mram vector obtained from an actual bottom detector need not be R017's vector: the same unrestricted raw coupled operator survives EVERY such input. Its prescribed terminal restriction must still be detected. The exact source interfaces clarify the remaining issue:

\begin{itemize}
\tightlist
\item
  ValLastX.lean244--251 (MainX) transports top-tail vanishing to bottom vanishing against the SAME downward image dn(ioT F). At m2 StepX is vacuous.
\item
  JointSmooth.lean264--292 (Bot and LastVan) separates all-bottom period detection from testing only that downward image.
\item
  PairForm.lean143--145 (PuNZ) asks for a nonzero Poincare image of some ramified row; it is not an assertion about R017's chosen row.
\item
  \nolinkurl{ExistsCentralBottomDatum.lean48–49} (one\_central\_pin) explicitly assumes \(1\le r_S\). At m2, \(r_S=m/2-1=0\), so this constructor cannot supply the common central pin. BottomRelationOfPin.lean101 instead applies to an already pinned datum. This is a constructor limitation, not an impossibility theorem for m2 theta lifts.
\item
  \nolinkurl{BottomTriangleNonzeroPeriod.lean36–47} is a normalized kernel identity for independently supplied Pair data. \nolinkurl{BottomNonvanishingFromInputs.lean28–31} separately assumes LNe, BotSmall, NormBig and TriEx. Faithfulness and unrestricted G2 nonvanishing do not supply them. Its uniform scalar coordinate condition retains the old host normalization; mixed application needs K035's actual varying comparison.
\end{itemize}

No implication from unrestricted G2 nonvanishing to the actual brC/tilQ/PtilFF period is used here.

\subsection{9. Sources and proof/check scope}\label{sources-and-proofcheck-scope}

New steps proved: uniform unitary holomorphy(MR9)--(MR14); the unchanged G quotient/norm on the entire generic type(MR15)--(MR22); the full-Mram global application(MR25)--(MR27); and the exact highest/coeval comparison(MR29)--(MR31). The finite cyclic decomposition is written out to match the source's unweighted norm.

Primary references:

\begin{itemize}
\tightlist
\item
  Jiang, D., \& Luo, Z. (2023). Certain Fourier operators and their associated Poisson summation formulae on GL1. Pacific Journal of Mathematics,326(2),301--372. §2B,(2-10),Theorem2.4(1)--(2), pp.313--314. \url{https://msp.org/pjm/2023/326-2/pjm-v326-n2-p05-p.pdf} Directly checked copy: work/agents/lift\_refs/jiang\_luo\_2023.pdf. This restates Godement--Jacquet, ChapterI, Theorem3.3; the original monograph is not claimed to have been read in full. Haar constants here are fixed independently.
\item
  Tadic, M. (1986). Classification of unitary representations in irreducible representations of general linear group (non-Archimedean case). Annales scientifiques de l'École Normale Supérieure,19(3),335--382. TheoremD,pp.338--339 and the characteristic-zero appendix. \url{https://www.numdam.org/article/ASENS_1986_4_19_3_335_0.pdf}
\item
  Mínguez, A. (2008). Correspondance de Howe explicite: paires duales de type II. Annales scientifiques de l'École Normale Supérieure, 41(5),717--741. Theorem1,p.718 and normalization(1.4),pp.721--722. \url{https://www.numdam.org/item/ASENS_2008_4_41_5_717_0.pdf}
\item
  Roche, A. (1998). Types and Hecke algebras for principal series representations of split reductive p-adic groups. Annales scientifiques de l'École Normale Supérieure,31(3),361--413. Theorem7.7/Remark7.8,pp.391--392;Theorem9.4,p.404;residue characteristic restriction p.379. Used through checked LIFT24. \url{https://www.numdam.org/article/ASENS_1998_4_31_3_361_0.pdf}
\item
  Gan, W. T., Qiu, Y., \& Takeda, S. The regularized Siegel--Weil formula and the Rallis inner product formula. Theorem7.1 and §§11.3--11.6, with the exact dimension ranges and conjugation dictionary checked in DO028,DO036,LIFT26.
\end{itemize}

Checked run inputs, not reproved: LIFT16(S8 and actual spherical test),LIFT20(RG6),LIFT21(TG7--TG9),LIFT22(CG4--CG21), LIFT24(MT6--MT33),LIFT25(AR5--AR30),LIFT26(GS7a--GS31), DO036/037,K035(PF22--PF29),K037(CR19,CR26--CR36). K038(CA27a,CA30a--CA30c,CA33--CA36) supplies the completed initial-level and restricted arithmetic comparison used in section7; that completed restriction is explicitly distinguished from the full G2 map.

No Lean execution or new finite computation certificate is claimed. Author whole-proof and independent reviews are recorded only after the corresponding complete text has been read.

\clearpage

\section{AM-R-021 --- the current equally spaced family fails the full budget weight windows}\label{am-r-021-the-current-equally-spaced-family-fails-the-full-budget-weight-windows}

Author: GPT-6 Astra Pro, ALP\_MACHINE /root. 22 November 2026 UTC. Reasoning effort is not externally exposed.

The concrete K031/K032 weights satisfy the analytic Hyp3 inequalities, and R018 verifies the four numerical \nolinkurl{Sizes} clauses on explicit finite growing intervals. Neither result claimed the much stronger global \nolinkurl{SizeData₂} record. This finding checks two of that record's additional conditions and proves that the current equally spaced family cannot satisfy them, even after arbitrary common determinant translations at individual CM labels. A different weight profile is required for the actual DP44 budget. This is an obstruction to this specified family, not to all mixed theta chains or all shorter lift counts.

\subsection{1. Exact conditions and a shift-invariant necessary inequality}\label{exact-conditions-and-a-shift-invariant-necessary-inequality}

The source \nolinkurl{Formalism/LocalRingCardinality/DataParams₂.lean}157--164 requires the following, writing S\_c for its positive real parameter Sc. For every CM label tau and every i=0,\ldots,k−2,

\[
                  0\le v_{\tau,i}\le S_c.
                  \tag{R21.1}
\]

For every label outside \nolinkurl{host.topPlaces}, it also requires

\[
                  S_c\le u_{\tau,0}-v_{\tau,0}
                         \le C_c S_c.
                  \tag{R21.2}
\]

Only the first inequality in (R21.2) will be used. There is no dependence on a hypothetical sign of the height term, the chosen local gain, the coefficient field, or an endpoint detector in this argument.

Allow an arbitrary real number s\_tau at each label, and replace

\[
 (u_{\tau,i},v_{\tau,i},b_{\tau,j},e_{\tau,j})
 \longmapsto
 (u_{\tau,i}-s_\tau,\ v_{\tau,i}-s_\tau,
                         b_{\tau,j}+s_\tau,\ e_{\tau,j}+s_\tau).
 \tag{R21.3}
\]

This is the common weight translation induced by a determinant twist, when such a twist exists with the required characters. We permit all real s\_tau here, a larger set than the admissible twists; therefore impossibility in this larger set also excludes every actual integral or character-compatible determinant repair.

\textbf{Necessary inequality.} At every non-top label, any translated profile satisfying (R21.1)--(R21.2) must satisfy

\[
 \boxed{\quad
 \max_i v_{\tau,i}-\min_i v_{\tau,i}
                  \le u_{\tau,0}-v_{\tau,0}.\quad}
                  \tag{R21.4}
\]

\textbf{Proof.} Subtracting the smallest translated v entry from the largest in (R21.1) gives max(v)−min(v) at most S\_c.~The common translation cancels from this difference and from u\_0−v\_0. Apply (R21.2). This proves (R21.4) with no assumption that the translated profile satisfies any of the other budget conditions. QED.

\subsection{2. Actual weights and enough unchanged non-top labels}\label{actual-weights-and-enough-unchanged-non-top-labels}

K031(HC7) and K032's corresponding baseline use

\[
 u_{\tau,i}=(2(k-i)+1)C+\tfrac12,\qquad
 v_{\tau,i}=2(k-i)C.
 \tag{R21.5}
\]

The growing K033 family adds A=11\^{}r to u\_0 at one label and 121A to u\_0 at one other label. It leaves every v entry unchanged and leaves u\_0 unchanged at all other labels. Thus at each unchanged label

\[
 \max v-\min v=2(k-2)C,\qquad u_0-v_0=C+\tfrac12.
 \tag{R21.6}
\]

There is necessarily an unchanged non-top label under any source frame indexing. In fact the literal CM label set is the disjoint union of the two images tau1(A) and tau2(A), each of size f/2. The top set is TT mapped by tau1 for a subset TT of A. Hence there are at least f/2 non-top labels. Removing the at most two growing labels leaves at least f/2−2 unchanged non-top labels. For the actual hosts these numbers are19 and82, respectively.

The exact accepted locators for that counting statement are \nolinkurl{DataHost.lean}100 (TT),129 (topPlaces),190--191 (FactPlaces), and the proven bijection \nolinkurl{tauPl_bij} in \nolinkurl{FactDelDualNabla.lean}26--45. The proof of the last theorem uses the disjoint tau1/tau2 parametrization and the frame's iota bijection. We use these accepted frame statements, not a new audit of the archive. Relabeling the two selected growing places cannot evade the cardinality argument.

Apply (R21.4) at any one of these unchanged non-top labels. It would give

\[
                         2(k-2)C\le C+\tfrac12,
                         \tag{R21.7}
\]

which is false for both concrete hosts. The full exact numbers are:

\begin{longtable}[]{@{}
  >{\raggedright\arraybackslash}p{(\columnwidth - 8\tabcolsep) * \real{0.1579}}
  >{\raggedleft\arraybackslash}p{(\columnwidth - 8\tabcolsep) * \real{0.2105}}
  >{\raggedleft\arraybackslash}p{(\columnwidth - 8\tabcolsep) * \real{0.2105}}
  >{\raggedleft\arraybackslash}p{(\columnwidth - 8\tabcolsep) * \real{0.2105}}
  >{\raggedleft\arraybackslash}p{(\columnwidth - 8\tabcolsep) * \real{0.2105}}@{}}
\toprule\noalign{}
\begin{minipage}[b]{\linewidth}\raggedright
Host
\end{minipage} & \begin{minipage}[b]{\linewidth}\raggedleft
k=f
\end{minipage} & \begin{minipage}[b]{\linewidth}\raggedleft
C
\end{minipage} & \begin{minipage}[b]{\linewidth}\raggedleft
Required lower bound for S\_c
\end{minipage} & \begin{minipage}[b]{\linewidth}\raggedleft
Required upper bound for S\_c
\end{minipage} \\
\midrule\noalign{}
\endhead
\bottomrule\noalign{}
\endlastfoot
J14 & 42 & 45156049293 & 3612483943440 & 90312098587/2 \\
J56 & 168 & 40530339687 & 13456072776084 & 81060679375/2 \\
\end{longtable}

The lower bound exceeds the upper by7134655788293/2 in the first case and26831084872793/2 in the second. This proves:

\textbf{Theorem R21.8.} No member of the specified K031/K032--K033 growing family, including every member of the R018 certified fixed-T ranges, satisfies the full literal \nolinkurl{SizeData₂} weight conditions. This remains true after every relabeling of CM places and every common per-label translation (R21.3) compatible with the indicated source frame. In particular determinant twists alone cannot turn this family into a full DP44 budget instance.

\textbf{Proof.} The counting argument supplies an unchanged non-top label. Equations (R21.6) and the necessary (R21.4) give the contradiction (R21.7), quantified exactly in the table. QED.

\subsection{3. Scope, repair direction, and framework translation}\label{scope-repair-direction-and-framework-translation}

This obstruction is independent of the earlier incompatible old quadratic \nolinkurl{Fields.Kq}/elliptic \nolinkurl{IsTen} wrapper. Even a replacement mixed host that retained the source's two-half CM index and these weight-window inequalities would encounter (R21.7). Conversely, a replacement theorem changing those inequalities would require its own quantitative proof; it could not cite \nolinkurl{SizeData₂} unchanged.

The counterexample to full admissibility does not invalidate the proved Hyp3, numerical Sizes, local selected content, actual raw theta injectivity, geometric class maps, or terminal normalization on this family. Their scopes did not include (R21.1)--(R21.2). In particular R020's fixed whole-branch content interval remains a correct calculation on the stated weights. It is not a budget certificate for a family that fails the additional weight windows.

For a new non-top profile, (R21.4) gives a concrete repair criterion: the v cloud must fit below the first u\_0−v\_0 gap. Changing only the common determinant shift never changes that ratio. One can instead alter the spacings among the v entries and the relative position of u\_0. Top labels also require their distinct \nolinkurl{top_u}, \nolinkurl{top_b}, \nolinkurl{aT_win} and \nolinkurl{win_top} conditions; this note does not assert that the repair criterion alone supplies a full profile. Every such change must be carried through the selected kernels, character compatibility, arithmetic costs and terminal detector.

In the source's representation language the translation (R21.3) is a determinant-character operation. In algebraic geometry it tensors the coefficient representation by a rank-one character, organized by SGA3 I4.7.3/4.7.3.1 and SGA6 VI3.3 as in R008. The two differences in (R21.4) are unchanged by that operation; this is a numerical weight constraint, not an equality of integral images or a new Picard/biextension theorem. The existing formal Cartier--Dieudonne, Picard, motive and ell-adic interpretations of the kernels remain those in the dictionary. No RH implication or universal negative DP44 difference follows.

Primary source citation: Alpöge, L. (2026). \emph{Remark}, step8, with its stipulated archive at commit 7cbf7a0c32fca1badf05931879a37c3f8a2b03b4, exact declarations above. Prior run input: DP\_MULT DP44, which rebuilds every required datum and does not equate the smaller numerical Sizes record with \nolinkurl{SizeData₂}. Whole-proof readings and the arithmetic certificate are recorded separately; no Lean build is claimed.

\clearpage

\section{AM-DO-039: literal split-line tests and the two local Q restriction orbits}\label{am-do-039-literal-split-line-tests-and-the-two-local-q-restriction-orbits}

Author: GPT-6 Astra Pro, ALP\_MACHINE \nolinkurl{/root/definite_oscillator}. 22 November 2026 UTC. Inherited reasoning effort is not separately exposed.

This finding isolates two exact facts about terminal detection. First, varying all integral orthogonal line insertions into the standard terminal lattice does not separate its entire coefficient representation for the actual selected weights. Second, at a Q place the unchanged coupled theta vector is zero on one induced-model H orbit and nonzero on another. The latter has an explicit H-equivariant realization with the correct terminal compact character and a nonzero normalized local test.

Neither statement decides the global terminal period. The first concerns a fixed ambient coefficient space, before its rational-stabilizer and previous-image restrictions. The second concerns functions in the actual local induced representation, not a chosen global automorphic embedding. Their exact scope prevents either an unjustified universal-spanning claim or an unjustified global vanishing claim.

\subsection{1. What the literal inserted line is}\label{what-the-literal-inserted-line-is}

All source locators below refer to commit \nolinkurl{7cbf7a0c32fca1badf05931879a37c3f8a2b03b4} in the recovered source tree. Write E for either K031/K032 host and use the standard Hermitian form \(h(x,y)=\sum_i\overline{x_i}y_i\) on \(E^N\).

\nolinkurl{Columns/StdLatUnimodular.lean}165--176 defines

\[
 \operatorname{tilN}(L)=L\oplus\mathcal O_E e_N,
 \qquad h(e_N,e_N)=1.                         \tag{DO39.1}
\]

Its T covector is defined using \(z_0=\pi_T e_N\), and its last G/Q frame coordinate is \(x_N/\pi_T\). The lattice in (DO39.1) contains the norm-one vector \(e_N\), independently of the size of T. \nolinkurl{FiniteSumNonzeroLargeDepth.lean}204--205 defines the actual \nolinkurl{tilQ} using this \nolinkurl{DecoSet.tilH} construction. Consequently an isometric realization of one of these lattices inside the standard terminal lattice supplies an integral norm-one orthogonal summand there.

This must be distinguished from the original nonstandard-coordinate \nolinkurl{TilTied} record: \nolinkurl{LastCoordinateUniformizerVectorSpansOrthogonalLi.lean} 109--115 supplies a vector of norm T in its fixed Gram matrix and requires an exact orthogonal-sum lattice equation. A collection of arbitrary primitive norm-T vectors in the standard lattice is not automatically a collection of literal (DO39.1) insertions. To give that interpretation in the standard chart, the vector must be \(\pi_T\) times an integral norm-one summand, with the stipulated complement and decorations.

Root R017(R17.7) proves

\[
 U(I_N)(F)\cap\operatorname{GL}_N(\mathcal O_E)
                   =\mu(E)^N\rtimes S_N.       \tag{DO39.2}
\]

We use that settled theorem. Its entrywise argument also proves directly that every integral norm-one vector is a root of unity times one coordinate vector. Indeed every conjugate absolute value of every entry is at most one. A nonzero integral entry has nonzero integral algebraic norm, so all its conjugate absolute values equal one; its powers lie in a bounded subset of the Minkowski lattice, making it a root of unity. The sum of squared absolute values then permits exactly one nonzero entry. Thus there are precisely N such lines in the standard lattice.

\subsection{2. The full coefficient space is not separated by these N lines}\label{the-full-coefficient-space-is-not-separated-by-these-n-lines}

Let dominant integral weights \(\lambda\) of length N and \(\mu\) of length n=N−1 interlace and satisfy the source SameShift condition \(\mu_n=\lambda_N\). Common determinant shifts do not affect the argument. In (DO39.3--9) write the shifted partitions again as lambda and mu. Let \(V_\lambda,V_\mu\) be the characteristic-zero polynomial models after removing that shift, and let

\[
 B:V_\lambda\longrightarrow V_\mu              \tag{DO39.3}
\]

be the nonzero literal branching map, with the scalar computed by R019 and its integral lattice interpretation in R020. Multiplicity-one branching, the source \nolinkurl{BrTranspose/IoHw/CoefXi} interface, and its normalization are accepted carried results. No normalization of B to one is made here; only its kernel is being considered.

For each i choose the coordinate permutation \(g_i\) interchanging \(e_i\) and \(e_N\), and put \(B_i=B\rho(g_i)\). Changing the permutation by an integral isometry preserving the last line does not change its kernel: that isometry has a block \(A\in U(I_n)\) and a root of unity z on the last line. The branch has fixed last-coordinate weight \(|\lambda|-|\mu|\), so

\[
 B\rho(\operatorname{diag}(A,z))
     =z^{|\lambda|-|\mu|}\rho_\mu(A)B.          \tag{DO39.4}
\]

Equation (DO39.4) uses the positive polynomial row action. The literal source \nolinkurl{rhoVr} precomposes that action with inverse transpose (\nolinkurl{Divisibility/CoefMod.lean}320--321); it accordingly replaces A by A\^{}\{-t\} and z by z\^{}\{-1\}. The same kernel identity follows, and coordinate transpositions are unchanged. All factors remain invertible. This shows that changing the root-of-unity generator of a coordinate line or the integral complement basis does not add another independent kernel condition. The \(\overline\pi_T^{-nBr}\) factor of literal \nolinkurl{brC} is nonzero and does not alter this assertion.

The combined map for all integral norm-one line splittings at this fixed standard coefficient value is therefore

\[
 \mathcal B_{\mathrm{split}}:V_\lambda\longrightarrow
       \bigoplus_{i=1}^{N}V_\mu,\qquad
       v\longmapsto(B_i v)_i,\quad
 \operatorname{rank}\mathcal B_{\mathrm{split}}\le N\dim V_\mu.
                                                        \tag{DO39.5}
\]

There is an elementary lower bound for its domain dimension. For every integer s with \(\mu_1\le s\le\lambda_1\), the weight \(\mu(s)=(s,\mu_2,\ldots,\mu_n)\) still interlaces \(\lambda\). These are distinct summands in the multiplicity-free restriction. The Weyl dimension formula gives

\[
 \frac{\dim V_{\mu(s)}}{\dim V_\mu}
       =\prod_{j=2}^{n}
          \frac{s-\mu_j+j-1}{\mu_1-\mu_j+j-1}\ \ge1.
                                                        \tag{DO39.6}
\]

Every denominator is positive, including when the weight has repeated entries. Summing these particular branching summands yields

\[
 \dim V_\lambda\ge(\lambda_1-\mu_1+1)\dim V_\mu.
                                                        \tag{DO39.7}
\]

Thus if \(\lambda_1-\mu_1+1>N\), the common kernel in (DO39.5) is nonzero, with dimension at least

\[
       (\lambda_1-\mu_1+1-N)\dim V_\mu.        \tag{DO39.8}
\]

For the CM tensor coefficient space use bold weights and tensor B at all f CM labels. A global coordinate line has the same coordinate index at every embedding, so there are N tests, not \(N^f\). Root-of-unity entries merely change the nonzero factors in (DO39.4). Therefore

\[
 \dim\ker\mathcal B_{\mathrm{split}}
 \ \ge\left[\prod_\tau
    (\lambda_{\tau,1}-\mu_{\tau,1}+1)-N\right]
                              \dim V_{\boldsymbol\mu}.
                                                        \tag{DO39.9}
\]

This bound concerns the full tensor coefficient module.

\subsubsection{Actual selected weights}\label{actual-selected-weights}

R019(R19.18) gives the actual weights, including their common shift. At every unchanged label,

\[
 \lambda_1-\mu_1=C,\qquad
 \begin{array}{c|r|r|r}
 \text{host}&f&N&C\\ \hline
 J^{14}&42&3562&45\,156\,049\,293\\
 J^{56}&168&3436&40\,530\,339\,687
 \end{array}                                           \tag{DO39.10}
\]

At a label with selected increase A, the difference is \(C+A\). It already exceeds N at A=0. For K033's two growing labels, the exact lower bound in (DO39.9) is

\[
 \left[(C+1)^{f-2}(C+11^r+1)(C+121\cdot11^r+1)-N\right]
                         \dim V_{\boldsymbol\mu}>0.     \tag{DO39.11}
\]

These are the already constructed selected representations. R018 gives the limited numerical \nolinkurl{Sizes} ranges, whereas the independently checked R021 proves that this equally spaced family does not satisfy the full \nolinkurl{SizeData₂} windows. The present algebraic statement neither extends those ranges nor supplies a full budget instance.

\subsubsection{Rational and integral qualifications}\label{rational-and-integral-qualifications}

Ignoring the common nonzero T scalars, the branching and permutation matrices are rational in the fixed polynomial model. Thus their common kernel is defined over the coefficient field. Its intersection with any full source Weyl lattice contains nonzero vectors, by clearing finitely many denominators. At a fixed coefficient DVR this intersection is saturated: if a nonzero scalar times a lattice vector lies in the kernel, the vector already does. Hence a nonzero intersection contains a vector primitive in the ambient local Weyl lattice. No 11-denominator is forced merely by the existence of this ambient kernel.

This does \textbf{not} show that such a vector satisfies all rational-class stabilizer relations, terminal finite characters, or lies in the actual preceding theta image. Varying decorations can also change the value of a coefficient family; (DO39.5) does not replace all such values by one vector. Other terminal lattice classes can have different orthogonal line configurations. Accordingly the result rules out the proposed universal separation of the entire standard coefficient space by its literal split lines. It is neither a counterexample to MainX nor a vanishing theorem for the constructed arithmetic row.

\subsection{3. The actual Q vector and its closed H orbit}\label{the-actual-q-vector-and-its-closed-h-orbit}

Now work at a single actual Q place. By DO037.1 the local field is \(F=\mathbf Q_q\); its nontrivial tame character is \(\chi\). Write

\[
 a=k,\quad d=t+1,\quad b=a+d,\quad N=b+d,
 \quad H=\operatorname{diag}(\operatorname{GL}_{N-1}(F),1).
                                                        \tag{DO39.12}
\]

Use the exact raw functions \(F_y\) and \(H_y\) of LIFT23(Q5)--(Q7), and DO037's normalized whole-Schwartz quotient maps \(q_0,q_1\). Their targets are \(\pi_1=\pi_0\times\eta_0\det_d\) and \(\pi_2=\pi_1\times\eta_1\det_d\), for the actual spherical unitary \(\pi_0\), with \(\eta_0|_{O^\times}
=\eta_1|_{O^\times}=\chi\). Let

\[
 C(v)=q_1\left(\sum_y H_y\otimes q_0(F_y\otimes v)\right)
                      \in\pi_2,\qquad v\in\pi_0^{K_a}.
                                                        \tag{DO39.13}
\]

This is precisely the unchanged full coupled family, not one chosen Y class or a replacement finite vector.

For every first matrix X,

\[
                         \sum_yF_y(X)=0.                \tag{DO39.14}
\]

Indeed on its support the top r=a+1 rows A have full residue rank and \(F_y(X)\) is \(\chi(\ell_A\bar z)\), with uniform higher-lift multiplicity and the stated integral lower coordinates. The nonzero linear form \(\ell_A\) has equally sized fibres, so summing reduces to a nonzero integer times \(\sum_{s\in\mathbf F_q^\times}\chi(s)=0\). Off that support every summand is zero. This is also the zero-frequency case of the checked LIFT23 Fourier formula.

Regard \(\pi_2\) in the normalized parabolic model \(\operatorname{Ind}_{P_{b,d}}^{\operatorname{GL}_N}
(\pi_1\boxtimes\eta_1\det_d)\). Then

\[
                              C(v)(h)=0\quad(h\in H).
                                                        \tag{DO39.15}
\]

To prove this, initially put the second quotient parameter in its convergent half-plane. The matrix \(h^{-1}E_b g\) has last row zero. The actual terminal factor \(H_y\) is therefore independent of y: its Fourier factor \(\psi(-\eta y)\) equals one, and every other factor is independent of y. Its sum with the first output is the operator on \(\sum_y F_y=0\), by linearity. This proves the vector integral is zero in that half-plane. The checked finite-dimensional meromorphic continuation and holomorphy at the actual quotient point then prove (DO39.15). There is no interchange of a divergent integral at the actual point in this argument.

\subsection{4. A nonzero H-equivariant open-orbit restriction}\label{a-nonzero-h-equivariant-open-orbit-restriction}

In normalized coordinates set \(c=q^{-1}e_{a+1}^*\); in the literal unscaled terminal chart set \(c=\pi_Tq^{-1}e_{a+1}^*\). Here \(\pi_T\) is a unit. Relative to block sizes \(b,d-1,1\), put

\[
 g_* =\begin{pmatrix}I_b&0&0\\0&I_{d-1}&0\\-c&0&1\end{pmatrix}.
                                                        \tag{DO39.16}
\]

Its determinant is one. The actual DO037 evaluation, retaining this unit convention, is

\[
 \operatorname{ev}_1(C(v)(g_*))=\Lambda_Q v\ne0
                    \quad\text{if }v\ne0,              \tag{DO39.17}
\]

where \(\operatorname{ev}_1:\pi_1\to\pi_0\) is the first induced model's identity evaluation, and the carried exact scalar is

\[
 \Lambda_Q=q^{-d}G(\chi,\psi_0)\chi(-1)^{-1}
 q^{-t(a+1)}\frac{\kappa_{a+1}\kappa_t}{\kappa_b}
       \prod_{i=1}^{t}(1-(\eta_0\eta_1^{-1})(q)q^{-i})^{-1}.
                                                        \tag{DO39.18}
\]

Here \(\kappa_j=\prod_{i=1}^j(1-q^{-i})\). The Gauss phase and every Haar/Euler factor are retained. Formula (DO39.17) is a carried proved value, not a new normalization.

Define the restriction along this orbit by

\[
 (R_{g_*}F)(h)=F(g_*h),\qquad h\in H.                     \tag{DO39.19}
\]

It is an H-intertwiner for the actual right-translation action. Its left covariance is explicit. Write

\[
 J=H\cap g_*^{-1}P_{b,d}g_*.
\]

A direct matrix calculation gives

\[
 J=\left\{\begin{pmatrix}A&B\\0&D\end{pmatrix}
 \in\operatorname{GL}_{N-1}(F):\ cA=c\right\},
 \quad A\in\operatorname{GL}_b,
 \quad D\in\operatorname{GL}_{d-1}.                     \tag{DO39.20}
\]

To see necessity as well as sufficiency, the b-plane \(g_*^{-1}E_bF^b\) consists of \((x,0,cx)\). An element of H preserves it precisely when its lower-left block is zero and its upper-left block satisfies \(cA=c\). For such j,

\[
 g_*jg_*^{-1}=
 \begin{pmatrix}A&B&0\\0&D&0\\0&-cB&1\end{pmatrix}.
                                                        \tag{DO39.21}
\]

Its lower d-by-d block has determinant \(\det D\). Normalized parabolic covariance therefore gives

\[
 (R_{g_*}F)(jh)=\tau(j)(R_{g_*}F)(h),\qquad
 \tau(j)=|\det A|^{d/2}|\det D|^{-b/2}
                         \eta_1(\det D)\pi_1(A).
                                                        \tag{DO39.22}
\]

Thus (DO39.19) is a specified map into the smooth function space on H with this left covariance. No compact-support condition or unmentioned normalization by the modular function of J is imposed. It is nonzero on the actual vector C(v) by (DO39.17).

This is the open orbit: a b-plane is in it exactly when it does not contain the last coordinate line and is not contained in its complementary hyperplane. Projecting to that hyperplane writes it as the graph of a nonzero functional on a b-dimensional subspace. GL\_(N−1) is transitive on these pairs. These two conditions are open and nonempty. Equivalently, \(\dim H-\dim J=bd\), the dimension of the Grassmannian. The closed orbit used in (DO39.15) consists of b-planes contained in the hyperplane.

\subsection{5. The actual terminal compact character and local mass}\label{the-actual-terminal-compact-character-and-local-mass}

Let \(U_H\subset H\) be the inverse image of \(P_{r,2t}(\mathbf F_q)\), where r=a+1. This is the literal terminal Q compact group in the standard chart. Put

\[
 \chi_H(u)=\chi(\det\overline D_u),\qquad
 v_{U_H}=\frac{|P_{r,2t}(\mathbf F_q)|}
                    {|\operatorname{GL}_{N-1}(\mathbf F_q)|}.
                                                        \tag{DO39.23}
\]

The maximal compact in H has Haar mass one. Then

\[
                 \pi_2(u)C(v)=\chi_H(u)C(v)
                              \quad(u\in U_H).          \tag{DO39.24}
\]

This can be checked directly, without identifying a row by convention. The target Weil action is \(\eta_1(\det u)\phi(u^{-1}X)\) on the compact group. The raw terminal top determinant contributes \(\chi(\det\overline A_u)^{-1}\), the last row and its Fourier factor remain fixed, and the honest Weil factor contributes \(\chi(\det\overline A_u)
\chi(\det\overline D_u)\). Their product is (DO39.24). It is the inverse of the arithmetic \nolinkurl{rowTyQ} in K038, exactly as the source inverse level-character convention requires.

Combining (DO39.17), (DO39.19) and (DO39.24) proves the actual local test

\[
 \boxed{\quad
 \int_{U_H}\chi_H(u)^{-1}
       \operatorname{ev}_1\bigl((R_{g_*}C(v))(u)\bigr)\,du
                         =v_{U_H}\Lambda_Q v.\quad}      \tag{DO39.25}
\]

It has no unknown local scalar. The normalized average divides by \(v_{U_H}\), which need not be an 11-unit. At a coefficient place over11 with q different from11, DO037.33 gives the exact 11-valuation of \(\Lambda_Q\); the terminal parabolic index adds \(v_{11}(v_{U_H})\). These are already the source local volume factors, not a declaration of a primitive integral image.

\subsection{6. Consequence and remaining global gate}\label{consequence-and-remaining-global-gate}

Equations (DO39.15) and (DO39.25) coexist: the same nonzero actual local vector has zero restriction on the closed induced-model orbit and a nonzero open-orbit H-model with the prescribed compact character. An abstract nonzero theta quotient therefore does not identify the terminal functional, but a closed-orbit evaluation zero does not prove that every H-equivariant test is zero either.

A global automorphic realization of \(\pi_2\) is not its local parabolic function model evaluated at one place. Consequently (DO39.15) is not a vanishing assertion for the global function on \(H(F)\backslash H(\mathbf A)\). Conversely (DO39.25) does not show that the actual global brC/tilQ/PtilFF period uses this open-orbit functional or that its global scalar is nonzero. That requires the prescribed global period or a suitable proved seesaw identity, with all actual class insertions and the selected archimedean mu component.

The full actual coupled arithmetic map is now constructed by K038, and DO038 gives its exact terminal mass/branch transport. LIFT27 gives injectivity of the unrestricted full-G2 map on all initial Mram vectors. The remaining issue is survival under this particular global terminal restriction and branch, not the existence of an unspecified arithmetic row map or a free choice of finite Schwartz data.

\subsection{7. Checks and carried references}\label{checks-and-carried-references}

The companion exact certificate checks a nonzero integral common branch kernel in a small Weyl module, the dimension inequalities on the actual weight formulas, the open-orbit matrix conjugation and finite orbit cardinality, and the zero Fourier frequency in small residue fields. The proofs above establish arbitrary-rank statements; the certificate is not a replacement for them or a Lean check.

Carried checked findings used without rederiving their hard inputs: R017(R17.7), R019(R19.18--21), R020(R20.2--8), DO037.1--33, DO038.9--22, K038(CA15--30,CA33--36), LIFT23(Q5--12), LIFT27(MR4,MR32). The classical branching and dimension formulas are the same polynomial representation theory used in the checked R019; see Molev, A. I. (2002), \emph{Gelfand--Tsetlin bases for classical Lie algebras}, arXiv:math/0211289v2, §2 and the interlacing basis. The new results here are the exact scope of (DO39.5--11), the local zero statement (DO39.15), the explicit stabilizer and covariance (DO39.20--22), and the unchanged-character test (DO39.25).

\clearpage

\section{AM-K039 --- an actual coupled-row divisibility bound on the fixed raw family}\label{am-k039-an-actual-coupled-row-divisibility-bound-on-the-fixed-raw-family}

Author: GPT-6 Astra Pro, ALP\_MACHINE \nolinkurl{/root/kummer_extension}, 22 November 2026 UTC; inherited reasoning effort is not separately exposed.

This finding bounds the \textbf{actual arithmetic coupled map} constructed in AM-K038, on its unchanged integral Weyl row lattices and with its actual finite G/T/Q data. It proves the integral-section premise needed to apply the already checked full-module bound AM-K027(EX2--EX4). It then subtracts explicit uniform stabilizer denominators and applies the actual terminal branch AM-R-020 and the source's integral dual-Weyl detector lattice.

The conclusion is a lower bound on image divisibility. It supplies no sharp upper bound on a prior image, no nonvanishing after terminal restriction, and no DP44 gain/error or GRH conclusion. In particular the uniformly spaced HC profiles used below fail the separately required full SizeData2 cloud inequalities, as established in the parent's current profile audit. The numerical Sizes intervals and analytic Hyp3/kernel construction remain useful checked statements at their recorded scope; they are not called full Scale data here.

\subsection{1. Actual objects, source interfaces, and the prime 11}\label{actual-objects-source-interfaces-and-the-prime-11}

There are two principally polarized hosts:

\begin{longtable}[]{@{}
  >{\raggedleft\arraybackslash}p{(\columnwidth - 14\tabcolsep) * \real{0.1333}}
  >{\raggedright\arraybackslash}p{(\columnwidth - 14\tabcolsep) * \real{0.1000}}
  >{\raggedright\arraybackslash}p{(\columnwidth - 14\tabcolsep) * \real{0.1000}}
  >{\raggedleft\arraybackslash}p{(\columnwidth - 14\tabcolsep) * \real{0.1333}}
  >{\raggedleft\arraybackslash}p{(\columnwidth - 14\tabcolsep) * \real{0.1333}}
  >{\raggedleft\arraybackslash}p{(\columnwidth - 14\tabcolsep) * \real{0.1333}}
  >{\raggedleft\arraybackslash}p{(\columnwidth - 14\tabcolsep) * \real{0.1333}}
  >{\raggedleft\arraybackslash}p{(\columnwidth - 14\tabcolsep) * \real{0.1333}}@{}}
\toprule\noalign{}
\begin{minipage}[b]{\linewidth}\raggedleft
h
\end{minipage} & \begin{minipage}[b]{\linewidth}\raggedright
CM field K\_h
\end{minipage} & \begin{minipage}[b]{\linewidth}\raggedright
A\_h
\end{minipage} & \begin{minipage}[b]{\linewidth}\raggedleft
f={[}K\_h:Q{]}/2
\end{minipage} & \begin{minipage}[b]{\linewidth}\raggedleft
a
\end{minipage} & \begin{minipage}[b]{\linewidth}\raggedleft
d
\end{minipage} & \begin{minipage}[b]{\linewidth}\raggedleft
b=a+d
\end{minipage} & \begin{minipage}[b]{\linewidth}\raggedleft
c=a+2d
\end{minipage} \\
\midrule\noalign{}
\endhead
\bottomrule\noalign{}
\endlastfoot
1 & Q(zeta28) B, B the degree-seven subfield of Q(zeta29) & J\^{}14 & 42 & 42 & 1760 & 1802 & 3562 \\
4 & K\_1(sqrt5,sqrt13) & J\^{}56 & 168 & 168 & 1634 & 1802 & 3436 \\
\end{longtable}

Here J is the Jacobian of y\^{}2=x\^{}7−1. The integral product-polarization identifications and actual CM eigenbases are K030/K032; K033 checks that the change from primitive ordinary J coordinates and its inverse are integral at11. Fix a coefficient-field place w above11 and normalize v(11)=1. We may make a finite extension of its complete DVR R to define the finitely many torsion points, phases, and frames. Such an extension is not asserted unramified. All bounds below retain this normalization.

Use the actual two edges a→b→c, the complete selected coefficient spaces of K037, and the fixed raw coupled finite datum of K038. Initial and middle classes use the \textbf{literal strict} levels of K038(CA27--CA33). At T this is the strict covector stabilizer, and at G the actual unipotent-character level. The level at29 stays maximal. Projective-T or full-G realizations, if used to describe an input, retain the actual weighted sums and matching volumes of K038(CA30a--c).

The final row is the actual evGSD row, with rhoVr coefficients. The source terminal detector lattice is exactly

\[
 \mathcal M_T=\mathrm{MtildeL}
       =(\mathrm{mtildeGSD}) .\mathrm{intNab},
 \tag{IG1}
\]

defined in \nolinkurl{Formalism/LocalToGlobal/Columns/RowForms.lean}, lines281--283. Its intDel variant is a separately named object at286. Thus the detector uses the transported dual-Weyl lattice in the opposite rhoV action. The pairing with the branch column is the literal perfect Delta/nabla pairing pairBK, as in K037(CR17--CR19), DO038 section1 and \nolinkurl{Formalism/LocalRingCardinality/FactDelDualNabla.lean}, theorem at75. No orthonormal lattice replaces either integral lattice.

\subsubsection{Lemma IG2. Every admitted Q place has residue characteristic other than11}\label{lemma-ig2.-every-admitted-q-place-has-residue-characteristic-other-than11}

The real field F\_h contains F\_0=Q(zeta28)\^{}+. The powers of11 modulo28 are

\[
                 11,9,15,25,23,1.
 \tag{IG2}
\]

None of the first five is 1 or −1. Hence the Frobenius of11 in Gal(F\_0/Q)=(Z/28)\^{}times/\{±1\} has order6. Every F\_0 place above11 has residue degree6. By the residue-field tower, every F\_h place above11 has residue degree divisible by6. Its absolute norm therefore is not a rational prime. Literal \nolinkurl{DegOne} requires that norm to be prime: \nolinkurl{Formalism/LocalToGlobal/Divisibility/Tuples.lean}, lines64--67. Consequently no such place is GoodPl, and the tuple definition at lines112--114 kills its inertAt subgroup. The exact Qprime definition at128 is the set where that subgroup is not killed. Hence this place cannot belong to the admitted Q ramification set. QED.

For a second exact check, the real cyclotomic polynomial is X\^{}6−7X\^{}4+14X\^{}2−7. Its reduction modulo11 is irreducible. The proof uses the explicit Frobenius order, not an inference from a sample of Q tuples. Both hosts are unramified at11; their ramified rational primes belong to \{2,7,29\} or \{2,5,7,13,29\}. The already fixed G and T primes also differ from11. Thus the conditional Q-over11 caveat in K038 section7 is absent for these two exact hosts and every admitted tuple.

\subsection{2. Representatives and unchanged integral frames}\label{representatives-and-unchanged-integral-frames}

\subsubsection{Lemma IG3. Standard local representatives at11}\label{lemma-ig3.-standard-local-representatives-at11}

For each of the actual initial, middle, and terminal unitary class sets, choose representatives whose finite components are the identity at every place above11. This changes neither the class problem nor its level, scalar masses, or integral coefficient lattices.

\textbf{Proof.} Let U(H) be one of the actual unitary groups over F\_h, and let g\_v be the finitely many components of a representative at v\textbar11. It suffices to find gamma in U(H)(F\_h) with gamma g\_v in K\_v, the unchanged maximal compact subgroup. Then multiply on the right by the element of the allowed compact level whose v-component is (gamma g\_v)\^{}−1 and whose other components are the identity.

Here is a direct weak-approximation argument for the required gamma. On the F\_h-vector space of H-skew-adjoint endomorphisms use

\[
             C(X)=(1-X)(1+X)^{-1},\qquad X^*=-X.
 \tag{IG3}
\]

Matrix multiplication gives C(X)\^{}*C(X)=1. Its inverse on the open set det(1+g)≠0 is X=(1−g)(1+g)\^{}−1. Every nonempty local open subset of U(H)(F\_\{h,v\}) meets this set: multiply any element by a central norm-one scalar z sufficiently close to1. The local norm-one torus has infinitely many elements in each identity neighbourhood, for example from its rational Cayley parameter. The polynomial det(1+zg) has only finitely many zeros, so one can avoid them while staying in the open subset.

Apply this to K\_v g\_v\^{}−1. Choose a Cayley preimage in each such open set. Weak approximation in F\_h, applied to a fixed F\_h-basis of the skew-adjoint vector space, supplies one rational X in all the specified inverse-image neighbourhoods. Its C(X) is the required gamma. The argument does not use a determinant-one constraint, and therefore applies also to the odd terminal rank c−1.

Rational left multiplication and allowed compact right multiplication are precisely the existing class equivalences. Their covariance and lattice transport were proved in K037(CR37--CR39) and K038(CA24--CA29). The \textbf{new representative's} local lattice at11 is standard, so its own transported delAtr or intNab lattice is the standard local Weyl or dual-Weyl lattice. We do not assert that the rational change from an earlier display of this class is integral. Its exact transport is used instead. The final tilQ decoration only adds the specified T element, an11-unit. Thus it also gives integral local frames. QED.

There are finitely many classes for these fixed ranks and levels by K036. The preceding choices can be made once, independently of the growing archimedean weight. No bound for their global heights, no new principal29 level, and no additional arithmetic-field complexity bound is asserted.

\subsection{3. Actual geometric class germs are integral local sections}\label{actual-geometric-class-germs-are-integral-local-sections}

This section proves the hypothesis that cannot be replaced by generic algebraicity or an integral Hodge splitting. In the statement, an integral local section means that its coefficient in an integral algebraic formal frame is integral. The reduced Chern frame itself need not be integral.

\subsubsection{Lemma IG4. Finite character scalars away from11 are11-units}\label{lemma-ig4.-finite-character-scalars-away-from11-are11-units}

Let eta be one of the algebraic angular Hecke characters in the actual stage splitting, including a finite-order tuple twist. If a finite idele x has component1 at all primes above11, then eta\_f(x) is an 11-unit at every chosen algebraic coefficient-field place w.

\textbf{Proof.} Principalize the ideal of x: for some positive integer m, it has m-th power (a). Since that ideal has zero valuation at all primes above11, a is a unit at all those primes. We can write x\^{}m=a u with u a finite unit idele. The restriction of the actual Hecke character to finite units has open conductor and finite image, so eta\_f(u) is a root of unity. The global relation gives

\[
               \eta_f(x)^m=\eta_\infty(a)^{-1}\eta_f(u).
 \tag{IG4}
\]

For the actual angular infinity types, eta\_infty(a) is an algebraic product of embedded a's and roots of their a bar a norms. Every embedded a is an11-unit, since a is a unit at \textbf{all} primes over11. All factors on the right of (IG4) have valuation zero. Hence so does eta\_f(x). Finite-order twists already have root-of-unity values. Determinant roots at a place where g\_v=1 may be chosen to be1; their choice-independence is the actual even-rank restriction retained in DO022/DO026. This proves the assertion for those splitting factors as well. QED.

This lemma does not claim that the finite value of an infinite-order algebraic Hecke character is a root of unity. Only its11-valuation is needed, and (IG4) proves exactly that assertion.

\subsubsection{Lemma IG5. Integral prime-to11 Heisenberg expansion}\label{lemma-ig5.-integral-prime-to11-heisenberg-expansion}

Let phi be a finite Schwartz function in the fixed rational Schrodinger realization. Suppose its component at11 is the standard vacuum and its values are algebraic and integral at w. Then it is a finite linear combination

\[
               \phi=\sum_t b_t W_t\phi_0,
               \qquad b_t\in R,
 \tag{IG5}
\]

where phi\_0 is the actual trace-one vacuum, each rational Heisenberg parameter t represents prime-to11 torsion of the base torus, and all central phases are retained. A finite extension of R is allowed.

\textbf{Proof.} In the standard position lattice Lambda, choose a positive integer D prime to11 such that phi has support in D\^{}−1 Lambda\_hat and is invariant under D Lambda\_hat. This is possible because its 11-component is the standard lattice indicator. On each coset B+Lambda\_hat, B in D\^{}−1 Lambda/Lambda, expand the function in the characters of Lambda/D Lambda. With A running over its dual group, the coefficient is

\[
 b_{A,B}=\frac1{[\Lambda:D\Lambda]}
  \sum_{x\in\Lambda/D\Lambda}
        \phi(B+x)\,\psi_f(-\langle A,x\rangle),
 \tag{IG6}
\]

up to the explicitly prescribed central phase of W\_(A,B). All character values and that phase are roots of unity. The denominator is a power of D and is an11-unit. Thus every coefficient is in R. Finite character orthogonality proves the expansion exactly on each support coset and outside the support.

The actual phi\_0 is a unit phase times a half-lattice Heisenberg translate of the standard indicator: DO016(1), with coordinates alpha=delta/2 and beta=epsilon/2, reads −i 1\_(alpha+Zhat\^{}3)(x) psi\_f(−x\^{}t beta). Its product realization on the actual hosts has the same property. Undoing this fixed translate only multiplies (IG6) by roots of unity and enlarges the denominator by a power of2. Those factors are11-units. Thus (IG5) holds relative to the \textbf{specified} trace-one vacuum, with rational parameters whose denominators divide a power of2D. They represent prime-to11 torsion. No averaging over a group of order q−1 occurs. QED.

\subsubsection{Proposition IG7. The full actual class-kernel module premise}\label{proposition-ig7.-the-full-actual-class-kernel-module-premise}

After the representative choices of Lemma IG3, every normalized geometric germ appearing in either edge of K038(CA27), at every actual class and every allowed Y, is the reduced germ of an integral formal local section of the ordinary-product principal theta line. Its splitting scalar, cD coordinates, and actual local frames add no negative11-valuation.

\textbf{Proof.} At11 the actual raw G/T/Q data are absent by Lemma IG2, and both class representatives are the identity by Lemma IG3. The finite input is therefore precisely the local vacuum. At the other places the raw values in K038 sections3--5 are zero, roots of unity, or powers of the residue cardinality times roots of unity, including the exact ftY factor q\^{}−Mb. All are integral at w. The finite Weil action preserves this coefficient integrality at places of residue characteristic different from11: its Schrodinger generators are linear changes of variable with absolute-determinant half powers, root-of-unity phases, and normalized Fourier transforms. On each compact support/period quotient the Fourier denominator and the measure are powers or square roots of that residue characteristic. These are11-units. Determinant splitting characters have valuation zero by Lemma IG4. This argument applies also at the primes2 and7; it does not assert integrality at their own residue valuations.

Equivalently the prescribed two moved-vacuum averages of DO017(8)/DO026(3) have indices prime to11, because all moved lattices agree with the standard lattice locally at11. Their exact positive square-root normalizers are therefore11-units. This is the same finite Weil realization checked in DO022 and K035(PF28), with no new vacuum normalization.

Lemma IG5 now writes each actual moved input as a sum with integral coefficients of prime-to11 Heisenberg translates of phi\_0. Under the already proved reading K035(PF25--PF28), remove the common complex period C\_theta\^{}(14h n\_j n\_(j+1)). The resulting function is exactly the same linear combination of the corresponding \textbf{actual reduced Mumford translates} of the seed product of sigma\_q's. K035 retains their central phases; it is an all-function equality before taking derivatives, so it identifies the germ already used to define the arithmetic class kernel.

LIFT08 proves that sigma\_q, normalized by leading coefficient u1u3−u2\^{}2, is primitive in the actual integral theta line at11. K034(TG3--TG10) proves that every prime-to11 torsion translate of an integral section, with precisely this reduced translation and its finite Heisenberg fibre scalar, belongs to the same complete coefficient module as an integral local section at0. It compares the normalized torsion connection, not merely the translated line. For even torsion order the same proof uses an exponent2d prime to11; the formal root and the rigidification argument are unchanged. Taking the integral linear combination proves the asserted premise. In the explicit notation of K034(TG9), the resulting coefficient belongs to

\[
            U_0\exp(-\beta_{\rm prod})\,R[[t_1,\ldots,t_g]],
\]

where the t\_i are integral algebraic formal parameters. The exponential is retained when passing to logarithmic derivative coordinates. It is not replaced by an integral Chern Taylor unit.

Finally K033(CS1--CS5) gives integral invertible actual CM coordinates, K035(PF15--PF16) gives cD\_tau an11-unit, and Lemma IG3 gives standard local class frames. The tilQ T-scaling is also an11-unit. Thus the framed functional in K037(CR39--CR40), including determinant twists, has exactly the asserted integral-module bound. QED.

The common complex period, the source unit-Fock creation factors lambda\_tau, and the literal real frame sc\_tau=lambda\_tau/cD\_tau have not been assigned finite valuations. They stay outside the arithmetic functional exactly as in K038(CA33--CA36). This proof does not claim ordinary Taylor integrality of an ample-line reduced frame. Its integral object is the algebraic local-section module, transferred using the checked torsion difference connection.

\subsection{4. A bound for every entry in the complete selected module}\label{a-bound-for-every-entry-in-the-complete-selected-module}

Write D\_(j,tau) for the selected degree at CM label tau on edge j, j=0,1, and D\_j=sum\_tau D\_(j,tau). Put n\_0=a,n\_1=b,n\_2=c and

\[
 d_{11}(n)=\begin{cases}0&n=0,\\1+\lfloor\log_{11}n\rfloor&n>0,\end{cases}
 \qquad
 E_j=n_jn_{j+1}\sum_\tau d_{11}(D_{j,\tau}),
 \qquad
 L_j=\max\{0,\lfloor D_j/1220\rfloor-E_j\}.
 \tag{IG8}
\]

\subsubsection{Theorem IG9. Actual local selected matrix-entry bound}\label{theorem-ig9.-actual-local-selected-matrix-entry-bound}

For every matched actual class pair and Y in K038,

\[
       k_j(\mathcal L_{j,c})
                       \subset 11^{L_j}\mathcal L_{j+1,c'}.
 \tag{IG9}
\]

\textbf{Proof.} By Proposition IG7 its actual germ F satisfies the checked full-module estimate K027(EX2). That estimate remains valid over the present R: its proof uses integral inverse logarithm maps and arbitrary integral formal prefactors, and every coefficientwise inequality is preserved by an integral extension. It makes no claim that this extension is unramified.

For an integral monomial differential operator of the selected multidegree, K027(EX3), obtained there from the exact Legendre formula, gives

\[
 v(\partial^\alpha F(0))
 \ge D_j/1220-\frac1{10}\sum_{\tau,r,s}s_{11}(\alpha_{\tau,r,s}).
 \tag{IG10}
\]

Every exponent at label tau is at most D\_(j,tau), and there are n\_j n\_(j+1) entries at that label. Thus its digit sum is at most 10 d\_11(D\_(j,tau)); (IG10) is at least D\_j/1220−E\_j. The full Hurwitz integrality carried in K023/K027 gives the additional lower bound0. Taking the conservative integer floor yields (IG8). No integer-valued valuation assumption on a ramified R is needed.

All polynomials produced by the actual universal subD substitution have this \textbf{label multidegree} and integral coefficients, on the whole source Weyl lattice. Their source row weights need not be the highest weight. The integral Weyl-to-Weyl theorem K037(CR8--CR14), not an abstract rational Schur projector, therefore applies to this complete family of monomial functionals. It gives (IG9), with the actual transported frames fixed in Proposition IG7. QED.

This argument uses K027(EX2--EX4), whose full proof was checked, and does not invoke the narrower fixed-source-row estimate EX6 on all weights. It introduces neither a hook-product inverse nor a Gram matrix inverse. All hook and compact metric factors remain where the analytic normalization puts them.

\subsection{5. Uniform actual class-stabilizer denominators}\label{uniform-actual-class-stabilizer-denominators}

\subsubsection{Lemma IG11. An elementary uniform11-bound}\label{lemma-ig11.-an-elementary-uniform11-bound}

For every finite stabilizer Gamma of an actual rank-n Hermitian O\_(K\_h)-lattice, including the additional level or row decorations,

\[
 v_{11}(|\Gamma|)\le S_f(n):=2m+v_{11}(m!),
                    \qquad m=\lfloor2fn/5\rfloor.
 \tag{IG11}
\]

\textbf{Proof.} The underlying integral lattice has Z-rank D=2fn, and Gamma embeds faithfully in GL\_D(Z). It is finite because the Hermitian form is positive definite at every real place: preserving the summed positive real trace form bounds all matrix entries of an integral automorphism.

The kernel of reduction GL\_D(Z)→GL\_D(F\_3) is torsion-free. Here is the needed direct proof. A nonidentity matrix in that kernel has the form A=1+3\^{}s B, s≥1, with B not identically zero modulo3. If A has prime order ell different from3, then (A\^{}ell−1)/3\^{}s is ell B modulo3, a contradiction. If ell=3, then (A\^{}3−1)/3\^{}(s+1) is B modulo3, again a contradiction; the remaining two binomial terms are divisible by3 since s≥1. Any nontrivial finite-order element has a power of prime order in the same kernel. Thus reduction is injective on Gamma, and its order divides \textbar GL\_D(F\_3)\textbar.

The exact powers 3,9,5,4,1 modulo11 give ord\_11(3)=5, and 3\^{}5−1=242=2·11\^{}2. Therefore 11 divides 3\^{}i−1 exactly when5 divides i. The elementary lifting formula for odd prime11 gives v\_11(3\^{}(5j)−1)=2+v\_11(j). It follows from the standard product \textbar GL\_D(F\_3)\textbar=3\^{}(D(D−1)/2) product\_(i=1)\^{}D(3\^{}i−1) that

\[
 v_{11}(|GL_D(F_3)|)
       =\sum_{j=1}^{\lfloor D/5\rfloor}(2+v_{11}(j))
       =2m+v_{11}(m!).
 \tag{IG12}
\]

For completeness, the used lifting formula follows by writing A=1+11\^{}e u with e≥1 and 11 not dividing u: raising A to a power prime to11 preserves e, while the binomial expansion of A\^{}11 gives valuation e+1. Factor j into its11-power and prime-to11 parts. This proves (IG12) and the lemma without enumerating classes. QED.

The rank D is the actual \textbf{2fn}, not n or fn. Taking level stabilizers or invariant-row stabilizers only makes the finite group smaller. This estimate does not assert that all actual stabilizers are prime to11, and it makes no principal congruence refinement.

\subsubsection{Theorem IG13. Actual coupled row and actual terminal scalar}\label{theorem-ig13.-actual-coupled-row-and-actual-terminal-scalar}

Let P\_cpl be K038(CA27) on the literal integral MfV lattice, or on its actual Mram sublattice. Then

\[
 \boxed{\quad
 P_{\rm cpl}(\mathcal L_0)
       \subset 11^{G_{\rm cpl}}\mathcal L_{\rm ev},
 \qquad G_{\rm cpl}=L_0+L_1-S_f(a)-S_f(b).
 \quad}
 \tag{IG13}
\]

For the growing family below and r≥13, let B\_h be the actual whole-Weyl branch lower bound of R020(R20.11):

\[
 B_1=59\,781\,009\,885\,146\,640,\qquad
 B_4=778\,040\,041\,530\,272\,708.
 \tag{IG14}
\]

For every actual terminal detector G in MtildeL of (IG1),

\[
 \boxed{\quad
 v\bigl(\operatorname{PtilFF}(P_{\rm cpl}(P),G)\bigr)
       \ge G_{\rm term}:=
          G_{\rm cpl}+B_h-S_f(c-1)
 \quad (P\in\mathcal L_0).
 \quad}
 \tag{IG15}
\]

Zero values have valuation+infinity.

\textbf{Proof.} Apply (IG9) to each of the two actual kernels in every summand of K038(CA27). Their composition maps the initial lattice into 11\^{}(L\_0+L\_1) times the target lattice. Lemma IG11 bounds the two actual denominators \textbar Gamma\_0\textbar{} and \textbar Gamma\_1\textbar, term by term. The Y sum is an actual unnormalized sum: it contributes no inverse cardinality. The nonarchimedean triangle inequality gives the same bound for all class sums. K038(CA29) puts the result in the actual evGSD equivariant subspace. Intersecting the full coefficient lattice with that subspace is saturated over the DVR, so the ambient containment gives exactly (IG13). Restricting to Mram is legitimate because it is a sublattice; no projection onto it is introduced.

R020 proves that the actual brC, including brConst/kapBr, maps the whole final Weyl lattice into 11\^{}B\_h times the terminal Weyl lattice for every r≥13. Its retained conjugate piT power is an 11-unit. Lemma IG3 supplies compatible integral local frames, so this inclusion applies to the preceding image as an inclusion. It does not require that the preceding image fill the whole Weyl lattice. PairBK takes that terminal Weyl lattice and the actual intNab detector lattice (IG1) into R. Finally K038(CA30) divides each terminal term by the actual rank-(c−1) stabilizer. Apply Lemma IG11 one more time and sum. This proves (IG15). QED.

The upper endpoint of the \textbf{whole-Weyl} branch interval in R020 has not been applied to a restricted image. Such an application would be invalid: a previous image can enter a deeper congruence lattice or the branch kernel. No nonzero vector or finite-index surjectivity is obtained from (IG13)--(IG15).

\subsection{6. Exact values on the already constructed growing weights}\label{exact-values-on-the-already-constructed-growing-weights}

Use the numerical HC/NCGaps baseline of K031/K032. At all f CM labels, let

\[
 C=\begin{cases}45\,156\,049\,293&h=1,\\40\,530\,339\,687&h=4,\end{cases}
 \qquad
 \nu_{0,i}=(2(a-i)+1)C+1-a+i-d/2
                    \quad(0\le i<a),
 \tag{IG16}
\]

and Delta=4C−d.~The exact selected recurrence is nu\_1=(nu\_0,0\^{}d)+Delta·1\_b. Define

\[
 \delta_0=\sum_i\nu_{0,i},\qquad
 \delta_1=\delta_0+b\Delta.
 \tag{IG17}
\]

The actual values, independently reconstructed from the checked K031 parameter function, are:

\begin{longtable}[]{@{}
  >{\raggedright\arraybackslash}p{(\columnwidth - 8\tabcolsep) * \real{0.1579}}
  >{\raggedleft\arraybackslash}p{(\columnwidth - 8\tabcolsep) * \real{0.2105}}
  >{\raggedleft\arraybackslash}p{(\columnwidth - 8\tabcolsep) * \real{0.2105}}
  >{\raggedleft\arraybackslash}p{(\columnwidth - 8\tabcolsep) * \real{0.2105}}
  >{\raggedleft\arraybackslash}p{(\columnwidth - 8\tabcolsep) * \real{0.2105}}@{}}
\toprule\noalign{}
\begin{minipage}[b]{\linewidth}\raggedright
Host
\end{minipage} & \begin{minipage}[b]{\linewidth}\raggedleft
delta\_0
\end{minipage} & \begin{minipage}[b]{\linewidth}\raggedleft
f delta\_0
\end{minipage} & \begin{minipage}[b]{\linewidth}\raggedleft
delta\_1
\end{minipage} & \begin{minipage}[b]{\linewidth}\raggedleft
f delta\_1
\end{minipage} \\
\midrule\noalign{}
\endhead
\bottomrule\noalign{}
\endlastfoot
J14 & 83,448,379,055,643 & 3,504,831,920,337,006 & 408,933,179,188,067 & 17,175,193,525,898,814 \\
J56 & 1,157,546,501,309,436 & 194,467,812,219,985,248 & 1,449,689,186,828,864 & 243,547,783,387,249,152 \\
\end{longtable}

The K033 growth increases the first weight at two designated labels by A=11\^{}r and B=121·11\^{}r, respectively. On \textbf{each} edge the selected partition increases in its first row by that same respective amount. Therefore

\[
 D_j(r)=f\delta_j+122\,11^r,\qquad
 \sum_\tau d_{11}(D_{j,\tau}(r))
 =(f-2)d_{11}(\delta_j)
      +d_{11}(\delta_j+11^r)+d_{11}(\delta_j+121\,11^r).
 \tag{IG18}
\]

Substitution in (IG8) gives a completely explicit integer lower bound for every selected coefficient of every actual class kernel. The stabilizer bounds do not depend on r:

\begin{longtable}[]{@{}lrrr@{}}
\toprule\noalign{}
Host & S\_f(a) & S\_f(b) & S\_f(c−1) \\
\midrule\noalign{}
\endhead
\bottomrule\noalign{}
\endlastfoot
J14 & 1,479 & 63,572 & 125,628 \\
J56 & 23,705 & 254,294 & 484,744 \\
\end{longtable}

For example:

\begin{longtable}[]{@{}
  >{\raggedright\arraybackslash}p{(\columnwidth - 10\tabcolsep) * \real{0.1304}}
  >{\raggedleft\arraybackslash}p{(\columnwidth - 10\tabcolsep) * \real{0.1739}}
  >{\raggedleft\arraybackslash}p{(\columnwidth - 10\tabcolsep) * \real{0.1739}}
  >{\raggedleft\arraybackslash}p{(\columnwidth - 10\tabcolsep) * \real{0.1739}}
  >{\raggedleft\arraybackslash}p{(\columnwidth - 10\tabcolsep) * \real{0.1739}}
  >{\raggedleft\arraybackslash}p{(\columnwidth - 10\tabcolsep) * \real{0.1739}}@{}}
\toprule\noalign{}
\begin{minipage}[b]{\linewidth}\raggedright
Host
\end{minipage} & \begin{minipage}[b]{\linewidth}\raggedleft
r
\end{minipage} & \begin{minipage}[b]{\linewidth}\raggedleft
L\_0(r)
\end{minipage} & \begin{minipage}[b]{\linewidth}\raggedleft
L\_1(r)
\end{minipage} & \begin{minipage}[b]{\linewidth}\raggedleft
G\_cpl(r)
\end{minipage} & \begin{minipage}[b]{\linewidth}\raggedleft
G\_term(r)
\end{minipage} \\
\midrule\noalign{}
\endhead
\bottomrule\noalign{}
\endlastfoot
J14 & 13 & 6,325,039,610,289 & 17,526,248,479,794 & 23,851,288,025,032 & 59,804,861,173,046,044 \\
J14 & 15 & 420,597,585,034,725 & 431,798,774,950,794 & 852,396,359,920,468 & 60,633,406,244,941,480 \\
J56 & 13 & 162,851,354,098,892 & 203,065,992,654,567 & 365,917,346,475,460 & 778,405,958,876,263,424 \\
J56 & 15 & 577,123,898,917,856 & 617,338,519,806,723 & 1,194,462,418,446,580 & 779,234,503,948,234,544 \\
\end{longtable}

These numbers are bounds on valuations normalized by v(11)=1, not logarithms of real norms. The exact certificate also records every r in the checked fixed-T numerical Sizes ranges,13--55 for J14 and13--90 for J56. The sharp selected \textbf{upper-witness} argument K033 begins at13 or15 respectively; (IG13) does not need that witness, so its J56 statement at13 is valid as a lower bound.

For the abstract growing kernel family, the digit error is O(r) at fixed ranks, whereas D\_j(r)=122·11\^{}r+O(1). Hence

\[
 \lim_{r\to\infty}
    \frac{G_{\rm cpl}(r)}{D_0(r)+D_1(r)}=\frac1{1220},\qquad
 \lim_{r\to\infty}
    \frac{G_{\rm term}(r)}{D_0(r)+D_1(r)}=\frac1{1220}.
 \tag{IG19}
\]

These are limits of the \textbf{proved lower-bound formulas}, not equalities for actual image content. Actual content can be larger, and the map or terminal scalar can vanish. Moreover fixed T only allows the finite numerical intervals already stated; no unbounded fixed-SizeData2 family follows. The new cloud-profile obstruction means that even those numerical intervals are not full DP44 Scale instances for this baseline.

\subsection{7. Exact remaining cost and image questions}\label{exact-remaining-cost-and-image-questions}

The common analytic comparison remains exactly K038(CA33--CA36):

\[
 v_0v_1 C_\theta^{14h(ab+bc)}
       \prod_\tau s_\tau^{D_{0,\tau}+D_{1,\tau}},
 \tag{IG20}
\]

with literal strict Haar volumes, actual source unit-Fock scales, and complete-coevaluation versus fixed-highest factors d\_0d\_1. Terminal mass is the further actual volume from DO038, not part of the arithmetic PtilFF valuation. These scalar, metric, and period factors have not been replaced by free normalizations.

To deduce a main gain/cost statement one still needs the corrected full parameter profile, a nonzero actual terminal value (or a quantitative statement on the restricted preceding image), and the source's complete norm/height/error accounting in the same normalization. In particular K034's sharp auxiliary first-extraction upper witness uses its explicitly chosen growing29 half-orbit. It does not furnish an upper bound for the fixed raw class map here, whose maximal29 level and complete fixed G/T/Q tensor are retained.

What is gained here is the integral-section comparison for every actual framed class kernel and a uniform bound for the actual composite and terminal maps on their original integral lattices. The previous purely formal or ambient coefficient lower bound is therefore no longer the unresolved premise for this fixed family.

\subsection{Sources and check scope}\label{sources-and-check-scope}

The new arguments are IG2--IG7, the full-label matrix application IG8--IG10, the elementary uniform stabilizer bound, and their actual coupled/terminal consequence IG13--IG19. The accepted primitive common-weight theorem is used through the already fully checked K027(EX2--EX4), not rederived and not promoted from a declaration. The standard Cayley formula, finite character orthogonality, congruence-kernel torsion argument and lifting valuation formula have their full proofs above.

The geometric translation comparison uses the proved K034(TG3--TG10), whose primary source conventions were directly checked against:

Bannai, K., \& Kobayashi, S. (2010). Algebraic theta functions and the p-adic interpolation of Eisenstein--Kronecker numbers. \emph{Duke Mathematical Journal, 153}(2), 229--295. \url{https://doi.org/10.1215/00127094-2010-024}. Author preprint \url{https://arxiv.org/pdf/math/0610163}, version4, section1.3, formula(14), Proposition2.3, Lemma2.7 and Proposition2.8. No ordinary-reduction integrality theorem from that paper is used at11 here.

All source-specific levels, scalar conventions and integral lattices are located above and in the carried K035--K038 reports. The exact certificate verifies the finite residue-degree and Fourier identities, rational Cayley form identities, degree formulas, GL-order/LTE arithmetic, and the numerical bounds. It does not enumerate unitary classes, evaluate theta coefficients or a terminal image matrix, prove nonvanishing, or run Lean. Whole-proof check events and hashes are recorded separately after completion.

\clearpage

\section{AM-R-022: a narrow-cloud replacement of the two-edge weight profile}\label{am-r-022-a-narrow-cloud-replacement-of-the-two-edge-weight-profile}

Author: /root, GPT-6 Astra Pro, ALP\_MACHINE, 22 November2026UTC. Reasoning effort is not externally exposed.

This finding replaces the equally spaced profile excluded by R021. It constructs explicit rational Harish--Chandra data, including every binary interior word, and verifies the \textbf{weight and Scale clauses} listed below. It is not a full \nolinkurl{Host}, \nolinkurl{D7}, \nolinkurl{SizeData₂}, \nolinkurl{SizesAt}, or \nolinkurl{Dat.Ok} instance. In particular the primitive-type/common-field obstruction, global gain and height records, and the terminal nonzero comparison are not removed by a numerically feasible weight profile. The numerical example also needs a new T prime; R017's old fixed primes are too small for its spread.

\subsection{1. Source clauses and parameters}\label{source-clauses-and-parameters}

All source locators use commit \nolinkurl{7cbf7a0c32fca1badf05931879a37c3f8a2b03b4}. \nolinkurl{LocalRingCardinality/Params.lean}6--80 gives the numerical ranks and \nolinkurl{HCParam}. \nolinkurl{DataParams₂.lean}9--49 defines Scale; 79--90 gives \nolinkurl{aT_win}, \nolinkurl{bT_ge}, \nolinkurl{top_u}, \nolinkurl{top_b}; 106 gives \nolinkurl{b_le}; 148--164 gives \nolinkurl{Sc_eq}, the three role windows, \nolinkurl{cloud_gap}, \nolinkurl{cloud_v}, \nolinkurl{cloud_u}, and \nolinkurl{be_close}. The binary interior formula is \nolinkurl{LocalToGlobal/Divisibility/Numbers.lean}, definition \nolinkurl{uW}. The source top/thick sets are disjoint and thin is their complement within the CM labels: \nolinkurl{DataHost.lean}129--136 and193--196.

Use the two already constructed CM hosts and ranks of K031/K032:

\begin{longtable}[]{@{}lrrrrrr@{}}
\toprule\noalign{}
Host & f=k & n=f/2 & t & d=t+1 & S=k+1+2t & N₂=k+2d \\
\midrule\noalign{}
\endhead
\bottomrule\noalign{}
\endlastfoot
J¹⁴ & 42 & 21 & 1759 & 1760 & 3561 & 3562 \\
J⁵⁶ & 168 & 84 & 1633 & 1634 & 3435 & 3436 \\
\end{longtable}

Put

\[
 \ell=28\left\lceil\frac{S^3+k+t+11}{28}\right\rceil,
 \qquad \delta=\ell,\quad
 S_c=4k\delta,\quad C_c=2,\quad h_0=2\ell,
 \quad\eta_{\rm sz}=4\ell=2h_0.                 \tag{R22.1}
\]

In both cases \(28\mid\ell\), \(S^3<\ell<2^{36}\), \(S_c<2^{46}\), and \((t+3)/2<\ell\). These are elementary integer inequalities, recorded exactly in the companion certificate. The source symbol \nolinkurl{Del} is to be assigned \(\delta\); this verifies the equality \nolinkurl{Sc_eq} if and when the remaining source record is built. It does not construct that record.

Choose any disjoint top and thick subsets of the f CM labels and let thin be the remaining labels. This includes every partition arising from the source sets just cited. There can be any number of top labels; all top labels below have the same prescribed extremes.

\subsection{2. An exact Scale whose main weight is an integer}\label{an-exact-scale-whose-main-weight-is-an-integer}

Define

\[
 A=2^{960},\quad B=2^{955}=A/32,\quad
 x=Ae^{-343},\quad
 \epsilon=\frac{\log343}{960\log2-343},\quad J=1.
                                                        \tag{R22.2}
\]

Then this is a Scale and its derived values are exactly

\[
 L=343,\quad b=e^{343},\quad a=A,\quad X'=e^7,
 \quad b'=Ae^{-7},\quad b'/x=e^{336},\quad
 w''=e^{336}/n.                                        \tag{R22.3}
\]

\textbf{Proof.} The elementary bounds \(0.69<\log2<0.7\), \(\log343<6\), proved for example by the convergent exponential series with its positive remainder bound, imply \(\log x>960(0.69)-343=319.4>52\cdot6>52\log343\). Thus \(x>1\) and \(0<\epsilon<1/52\). The definition gives \(x^\epsilon=343\); since \(343=7^3\), all the displayed Scale identities follow from their definitions. Also \(J\ge1\). Taking \(a_T=A\) satisfies both ends of the arbitrarily short \nolinkurl{aT_win} because its lower endpoint equals A exactly. Taking \(b_T=B\) satisfies \nolinkurl{bT_ge}, since \(e^7>2^7>32\) and hence \(Ae^{-7}<A/32\). No assertion about rounding a generic irrational a into a short interval is needed. ∎

For each label let the exact real role centre be

\[
 w_\tau=\begin{cases}
 B+e^{336}&\tau\text{ top},\\
 e^{343}&\tau\text{ thick},\\
 e^{336}/n&\tau\text{ thin}.
 \end{cases}                                           \tag{R22.4}
\]

Put

\[
 E_\tau=28\lceil w_\tau/28\rceil,\quad
 (e_{\tau,0},e_{\tau,1})=(E_\tau,E_\tau+2\ell),\quad
 (b_{\tau,0},b_{\tau,1})=(E_\tau-\ell,E_\tau+\ell).
                                                        \tag{R22.5}
\]

These are actual integers, specified uniquely by the displayed formula. The certificate determines their complete decimal values using rational upper and lower exponential-series bounds; it does not use a floating point rounding decision. Both e entries are multiples28 and their difference is the same \(2\ell\) at every CM label. Existence and finite conductors of compatible nonuniform stage characters are a separate character construction, not inferred from the integer formula.

\subsection{3. The full binary interior cloud}\label{the-full-binary-interior-cloud}

For \(0\le i\le k-2\) set

\[
 v_{\tau,i}=4\delta(k-1-i).                             \tag{R22.6}
\]

For every binary word \(\omega\in\{0,1\}^{k-2}\), define the interior entries, with \(1\le i\le k-2\), by

\[
 u^\omega_{\tau,i}=v_{\tau,i-1}-\tfrac12-
                             \delta\omega_{i-1}.        \tag{R22.7}
\]

At the endpoints use

\[
 (u^\omega_{\tau,0},u^\omega_{\tau,k-1})=
 \begin{cases}
 (A+(k-1)/2,-B-(k-1)/2)&\tau\text{ top},\\
 (v_{\tau,0}+S_c+1/2,2\delta+1/2)&\tau\text{ not top}.
 \end{cases}                                           \tag{R22.8}
\]

These formulas are compatible with the literal \nolinkurl{uW} rule, simultaneously at all labels. They include all words, not just one choice of each interior coefficient. If a later application uses a subset of the words, all the assertions below continue to hold on that subset.

\subsubsection{Proposition R22.9: HC, centredness, and noncompact gaps}\label{proposition-r22.9-hc-centredness-and-noncompact-gaps}

For every word and every label, (R22.5--8) is an \nolinkurl{HCParam} with the same numerical ranks as K031/K032. It satisfies the centredness and \nolinkurl{NCGaps} inequalities used in K031(HC5--HC8), including the actual natural-subtraction sign rule of \nolinkurl{Dpar}.

\textbf{Proof.} The interior interlacing gaps are exactly \(v_{i-1}-u_i=1/2+\delta\omega_{i-1}>0\) and \(u_i-v_i=4\delta-1/2-\delta\omega_{i-1}>0\). At the non-top endpoints the first gap is \(S_c+1/2>0\) and the last gap is \(2\delta-1/2>0\). At the top endpoints the first gap is positive since \(A>2^{46}>v_0\), and the last is positive since \(v_{k-2}>0>u_{k-1}\). Parity follows because k is even, all v,b,e are integers, and all u are strict half-integers. The be and eb gaps are all \(\ell\), larger than \((t+3)/2\).

Every role centre is at least \(e^{336}/84>2^{329}\), while all rank and cloud constants used here are less than \(2^{47}\). At a top label, \(u_{k-1}+b_0\ge e^{336}-\ell-(k-1)/2>0\). At other labels \(u_{k-1}>0\) and \(b_0>0\). This is the remaining HC inequality.

For centredness the distances from \(-e_j\) to u and v are larger than \((t+3)/2\); the only cancellation requiring an explicit check is the top last u entry, where \(u_{k-1}+e_0\ge e^{336}-(k-1)/2\). The distances to the two b entries are \(\ell\) or \(3\ell\). The block for the other e entry has distances at least \(2\ell-(t-1)/2>(t+3)/2\). These are precisely the complete centredness comparisons recorded in K031, and they hold for every interior word because those entries remain positive.

At stage0 the sole negative-sign entry of \nolinkurl{Dpar} is \(-b_1\). The positive-sign entries are the v entries and \(-b_0\). Their distances from \(-b_1\) are at least \(\min(2\ell,4\delta+b_1)>S^3\). At stages1 and2 the literal sign rule makes all entries positive; there are no opposite-sign comparisons. The odd-stage common shift \(-1/2\) changes no distance. This verifies every \nolinkurl{NCGaps} clause. No analytic nonvanishing is deduced just from this numerical check. ∎

\subsubsection{Proposition R22.10: all listed weight-window clauses}\label{proposition-r22.10-all-listed-weight-window-clauses}

The Scale (R22.2), top integers A,B, and (R22.1,5--8) satisfy \nolinkurl{aT_win}, \nolinkurl{bT_ge}, \nolinkurl{top_u}, \nolinkurl{top_b}, \nolinkurl{b_le}, \nolinkurl{Sc_eq}, \nolinkurl{win_thick}, \nolinkurl{win_top}, \nolinkurl{win_thin}, \nolinkurl{cloud_gap}, \nolinkurl{cloud_v}, \nolinkurl{cloud_u}, and \nolinkurl{be_close}, with positive Sc and etaSz.

\textbf{Proof.} The first three clauses were checked above. For \nolinkurl{top_b}, \(b_0-B\ge e^{336}-\ell>(k-1)/2+3\). Every e value lies between its role centre and that centre plus \(2\ell+28\le4\ell=\eta_{\rm sz}\); this gives each role window, including the allowed additional1 in the thick upper end. Every b/e difference has absolute value at most \(3\ell\), giving \nolinkurl{be_close}. The v entries lie between \(4\delta\) and \(4\delta(k-1)<S_c\). Every interior u lies strictly between0 and that same upper bound, for both choices of its word bit. At every non-top label \(u_0-v_0=S_c+1/2\in[S_c,2S_c]\). This proves every cloud clause and removes exactly R021's contradiction.

For the uniform b upper bound, \(e<4\) gives \(e^{343}<2^{686}\) and \(e^{336}<2^{672}\). At top labels \(b_1<B+2^{672}+28+2^{36}<2^{956}\); at other labels \(b_1<2^{687}\). But \(\log2<0.7\) gives \(e^{686}>2^{980}\). Thus every \(b_j<e^{686}=b^2\). The remaining equalities and positivity assertions follow directly from (R22.1). ∎

\subsection{4. Selected partitions and the numerical size of a new prime}\label{selected-partitions-and-the-numerical-size-of-a-new-prime}

Use the already checked source recurrence (R018.2), now with the actual label-dependent e entries. Write

\[
 w_{0,i}=u_i-(k-1)/2+i,\quad
 \nu_{0,i}=w_{0,i}+e_0-(k+d)/2,\quad
 \Delta=2\ell-d>0.                                    \tag{R22.11}
\]

Every \(\nu_0\) is an integral decreasing partition with positive last entry. Indeed at top labels its last entry is \(E_\tau-B-(k+d)/2>0\); at other labels it is \(E_\tau+2\delta-d/2>0\). Interlacing gives the required dominance and the exact adjacent differences are integers. The next selected partition is exactly

\[
                 \nu_1=(\nu_0,0^d)+\Delta\mathbf1_{k+d}.
                                                        \tag{R22.12}
\]

This is the same two-edge recurrence as before with its new e values; it retains every determinant shift.

The literal spread is first minus last plus rank. At top labels its three values are exactly

\[
 A+B+k,\qquad A+E_\tau+(k+d)/2,\qquad
 A+E_\tau+2\ell+(k+d)/2.                               \tag{R22.13}
\]

At a non-top label put \(W_0=v_0+S_c+1-k/2\). Then they are

\[
 u_0-u_{k-1}+1,\qquad W_0+E_\tau+(k+d)/2,\qquad
 W_0+E_\tau+2\ell+(k+d)/2.                            \tag{R22.14}
\]

The displayed selected last entries verify that each new constant block is the minimum at its stage; the first entry remains the shifted old first entry. Therefore the recurrence proves (R22.13--14) directly. They do not depend on the binary interior word. Each is strictly less than \(2A\), by the powers-of-two bounds used above, and \(2fN_2+1<4A\). The old G-size inequalities are unchanged. Consequently a norm/congruence T-prime with \(T>4A\) suffices for all four clauses of the numerical \nolinkurl{Sizes} record R018.8. This statement is conditional on choosing that new prime; it does not recycle the old fixed T values or prove the stronger \nolinkurl{AdmT}/height window.

\subsection{5. Remaining source and cost interfaces}\label{remaining-source-and-cost-interfaces}

This replacement proves the explicitly listed HC, centredness, gap, cloud and role-window conditions. The globally unramified character extension for its nonuniform e entries, a suitably large T prime, and the actual analytic/arithmetic transport must be supplied separately. Even after those steps, no full \nolinkurl{SizeData₂} is asserted here. Among its other fields are \nolinkurl{o_tie}, \nolinkurl{QS_tie}, \nolinkurl{split_tie}, the gain/error/support data, the logarithmic height equality, and the prescribed common host. The source's additional \nolinkurl{Host.o_pow}, \nolinkurl{qd_big}, and common quadratic CM-type requirements are not consequences of this weight calculation.

For example at the explicit Scale above, \(x^{\epsilon/100}=343^{1/100}<2\). Thus this particular Scale cannot satisfy \nolinkurl{o_tie} for a frame with o at least2. It is deliberately a concrete \textbf{weight-window repair}, not an assertion that all the source's global size records follow. Increasing the scale and rebuilding the host is a separate question. The terminal bottom detector and all actual DP44 gains, errors, heights and cross-family endpoint costs remain to be established in their common normalization.

The geometric interpretation is a replacement of coefficient weights in the same integral Schur functors; it changes neither the CM abelian scheme nor its Picard/biextension line. Exact programme locators remain SGA3I4.7.3/4.7.3.1, SGA6VI3.3, FGA232 Theorem3.1 pp148--149, FGA236 Corollary1.7 p229/2.3 p232, SGA7I VII2.7/2.9 pp169--174 andVIII7.4(b) pp310--311, EGAIV4 16.11.2--3(ii) pp54--55. No integral image, motivic correspondence, ell-adic weight theorem, or Kummer RH result is implied by this numerical repair.

\clearpage

\section{ALP-LIFT-28: a fixed-level holomorphic bottom detector with the exact selected jet degree}\label{alp-lift-28-a-fixed-level-holomorphic-bottom-detector-with-the-exact-selected-jet-degree}

Author: /root/lift\_accounting, GPT-6 Astra Pro, ALP\_MACHINE. 22 November 2026 UTC. Reasoning effort is not separately exposed.

This is a qualitative bottom-period construction. For either actual mixed CM host K031/K032, the literal initial compact level and its finite character extend to an explicit level on \(U(a,1)\). Its compact \(U(a)\) packet has distinct points and precisely the original rational stabilizers. The selected initial compact representation occurs in a single explicitly computed holomorphic jet degree. After increasing the last bottom parameter through a fixed arithmetic progression, cusp forms at this same finite level realize all these selected jets. Consequently every nonzero initial coefficient row has a nonzero ordinary compact bottom period against such a cusp form.

The theorem does not identify this constructed level with the literal rJ/finite frame in a completed source RealPack. It does not identify the new cusp forms with the prescribed downward theta image. There is no effective weight threshold, cusp-form norm bound, or DP44 saving here. The source's central-pin constructor requires positive rS, whereas the current m=2 problem has rS=0; that constructor is not supplied by this argument. These distinctions remain even though the finite packet and the selected real type below are actual, rather than parameterized existence interfaces.

\subsection{1. Fixed groups, initial rows, and conventions}\label{fixed-groups-initial-rows-and-conventions}

Let \(E/F\) be either K031's degree-84 CM host or K032's degree-336 CM host. Its chosen CM type has size f=42 or f=168, respectively. Let \(a=k=42\) or 168 and put

\[
 H=U(I_a),\qquad G^+=U(I_a\oplus(-1)),\qquad
 H\hookrightarrow G^+,\quad h\longmapsto\operatorname{diag}(h,1).
                                                        \tag{HJ1}
\]

At all real places H is compact and \(G^+\) has signature \((a,1)\). The Hermitian domain D is the product of the a-dimensional balls of negative lines, with base point \(o=[e_{a+1}]\).

Retain the actual initial strict compact level \(J_0\subset H(\mathbf
A_f)\) and actual right character \(\chi_0\) of LIFT27(MR4--7) and K038(CA27a,CA30a--c). At G its residue group is the full upper unipotent \(U_a(\kappa)\), with inverse generic character. At T it is the strict covector stabilizer, with trivial character. At Q, 29, and the other finite places it is the stipulated maximal basic compact group, with trivial character. This is the strict level, not the projective T level used to describe an individual nontrivial dilation eigenspace.

The positive compact highest weight at a CM label is denoted by \(\lambda_\tau\), and \(W_\lambda=\bigotimes_\tau W_{\lambda_\tau}\). Equip it with a fixed invariant positive Hermitian norm. Let \(\mathcal A_\lambda(J_0,\chi_0)\) be its actual coefficient-row space, equivalently the multiplicity space for \(W_\lambda\) in functions on \([H]=H(F)\backslash H(\mathbf A)\) with right character \(\chi_0\). All rational-stabilizer relations are imposed. Initial Mram is its specified subspace, including the unweighted T norm-kernel conditions.

For a unit vector w in \(W_\lambda\), write \(f_{P,w}\) for the scalar automorphic realization of a row P. Use real compact Haar mass one and the same finite Haar as K038/LIFT27. If \(x\) runs over the actual finite classes and \(\Gamma_x=H(F)\cap xJ_0x^{-1}\), then

\[
 \|f_{P,w}\|_{[H]}^2
   ={v_0\over d_\lambda}\sum_x{ \|P_x\|^2\over|\Gamma_x|},
 \quad v_0=\operatorname{vol}(J_0),\quad
 d_\lambda=\dim W_\lambda.                         \tag{HJ2}
\]

This is the carried compact Schur/mass identity, not a normalization of the future bottom cusp form. In the bilinear period the bottom finite character must restrict to \(\chi_0^{-1}\).

\subsection{2. The exact bottom Harish--Chandra list and jet order}\label{the-exact-bottom-harishchandra-list-and-jet-order}

Use zero-based indices. At a fixed label suppress tau and assume the source interlacing and parity conditions. In particular,

\[
 u_i>v_i>u_{i+1}\quad(0\le i<a-1),\qquad
 u_{a-1}>-b_0,\qquad b_0<e_0<b_1.
                                                        \tag{HJ3}
\]

The relevant half-integer/integer differences are at least \(1/2\). The exact definitions are \nolinkurl{Formalism/LocalRingCardinality/DataTheta.lean,} TPar.Interlace 529--533, Dpar 538--543, and Xpar 545--548, in the recovered commit 7cbf7a0c32fca1badf05931879a37c3f8a2b03b4. For m=2, j=0, the signed bottom parameter is

\[
       (p_0,\ldots,p_{a-1};q)
       =(v_0,\ldots,v_{a-2},-b_0;-b_1).             \tag{HJ4}
\]

Every positive-slot entry exceeds the negative-slot entry. There is no eta shift at j=0. Let \(\sigma\) be the discrete series with this parameter. For \(K=U(a)\times U(1)\), its minimal K type is

\[
 \tau_{\min}=W_\alpha\boxtimes z^\beta,\quad
 \alpha_i=p_i-{a-2\over2}+i,\qquad
 \beta=-b_1-{a\over2}.                             \tag{HJ5}
\]

Indeed the compact half-sum is \(\rho_c=((a-1)/2-i;0)\), while the noncompact half-sum for this chamber is \(\rho_n=((1/2)^{[a]};-a/2)\). Thus the Blattner parameter is the HC parameter plus \(\rho_n-\rho_c\), giving (HJ5). This is the usual minimal-type formula; see Harris, Kobayashi, and Speh (2025), Definition 2.7 and (2.24), printed p.16.

The initial compact type has highest weight

\[
                \lambda_i=u_i-{a-1\over2}+i.       \tag{HJ6}
\]

This is also the literal initial lamTrue: Columns/BaseLemmas.lean, lamN\_eq and lamTrue\_zero, 475--494, together with \nolinkurl{LocalRingCardinality/DataCoeff.lean41,} hcToWeightQ.

At o the holomorphic tangent representation is \(T_oD=\mathbf C^a\boxtimes z^{-1}\). The Harish--Chandra module of \(\sigma\), in the polarization with this tangent representation, has K-module structure

\[
       \tau_{\min}\otimes
       \operatorname{Sym}^{\bullet}
          (\mathbf C^a\boxtimes z^{-1}).            \tag{HJ7}
\]

This is the standard free generalized-Verma description for a holomorphic or antiholomorphic discrete series, with the name depending on the complex-structure convention. We use the explicit module in (HJ7). Its dual is realized by holomorphic sections on the negative-line ball with fibre \(\tau_{\min}^{\vee}\). Howe and Kraft (1997), opening discussion and equation (3), give this general description before their additional symplectic-group theorem. We do not invoke that latter theorem for unitary groups.

The required Pieri inequalities are exact:

\[
 \lambda_i-\alpha_i=u_i-v_i-\tfrac12\ge0
       \quad(i<a-1),\qquad
 \alpha_i-\lambda_{i+1}=v_i-u_{i+1}-\tfrac12\ge0,
\] \[
 \lambda_{a-1}-\alpha_{a-1}
                  =u_{a-1}+b_0-\tfrac12\ge0.      \tag{HJ8}
\]

Consequently \(\lambda/\alpha\) is a horizontal strip after any common determinant shift making these weights partitions. Pieri's formula therefore gives one copy of \(W_\lambda\), in precisely degree

\[
 \boxed{\quad N=|\lambda|-|\alpha|
       =\sum_{i=0}^{a-1}u_i-\sum_{i=0}^{a-2}v_i
                          +b_0-{a\over2}\in\mathbf Z_{\ge0}.\quad}
                                                        \tag{HJ9}
\]

Its extra \(U(1)\) weight is \(\beta-N\). The degree is unique because the scalar centre of \(U(a)\) counts polynomial degree. The multiplicity is one by the horizontal-strip Pieri rule (Macdonald, 1995, Chapter I, equation (5.16), printed p.73). Apply this at every label; the selected multidegree is \(\mathbf N=(N_\tau)_\tau\).

In particular, \(N_\tau\) does not involve \(b_{1,\tau}\). It is also unchanged by a common HC determinant shift \((u,v,-b_0)\mapsto(u+s,v+s,-b_0+s)\). No equal-spacing assumption on v occurs anywhere in (HJ3--9). These formulas therefore apply to a future narrow-cloud weight profile satisfying the same interlacing.

\subsection{3. An explicit bottom finite level extending the initial level}\label{an-explicit-bottom-finite-level-extending-the-initial-level}

At a relative split G place use the compatible identifications \(H(F_v)=\operatorname{GL}_a(F_v)\) and \(G^+(F_v)=\operatorname{GL}_{a+1}(F_v)\). Set

\[
 J_v^+=\operatorname{red}^{-1}U_{a+1}(\kappa_v).
                                                        \tag{HJ10}
\]

Extend the initial inverse generic character by the same additive character on all successive upper-diagonal entries, including the new last one. Denote the resulting character by \(\chi_v^+\). Then \(\chi_v^+|_{J_{0,v}}=\chi_{0,v}\).

At T, if the actual initial covector is u, take the strict stabilizer of \((u,0)\) modulo the same prime, inside \(\operatorname{GL}_{a+1}(O_v)\), with trivial character. At every other place take the stabilizer of the standard selfdual lattice in \(I_a\oplus(-1)\), again with trivial character. Denote the product by \((J^+,\chi^+)\). The form is integral unimodular at all these places; the relative CM extension in K031/K032 is unramified. At a split place this just means the displayed integral general linear compact. We have the literal equality

\[
             J^+\cap H(\mathbf A_f)=J_0,\qquad
                         \chi^+|_{J_0}=\chi_0.     \tag{HJ11}
\]

The bottom forms constructed below have right character \((\chi^+)^{-1}\), so the bilinear H period is invariant at the initial level. This inverse is essential.

At least one selected G prime has residue characteristic prime to \(|\mu(E)|\). The actual primes of K031/K032 satisfy this, and root reduction is injective there. For every finite component g, put \(\Gamma_g=G^+(F)\cap gJ^+g^{-1}\). Then

\[
                         \det\Gamma_g=\{1\}.       \tag{HJ12}
\]

To prove it, \(\det\gamma\) is a norm-one E-unit: every finite valuation is zero because every component of J+ has integral-unit determinant. At all complex embeddings it has absolute value one. Its powers are bounded in the Minkowski lattice of \(O_E\), so it is a root of unity. At the selected G place its residue is one, because the determinant of an upper unipotent residue is one; determinant is unchanged by the conjugating component g. Injective reduction of roots of unity makes the root equal to one. This proof works on every adelic component.

The construction is an explicit genuine level, not a claim that an uninstantiated source rJ record has these values. Section 8 states the remaining comparison precisely.

\subsection{4. The compact packet is embedded, with the correct stabilizers}\label{the-compact-packet-is-embedded-with-the-correct-stabilizers}

Map the finite class set

\[
        Z=H(F)\backslash H(\mathbf A_f)/J_0
      \longrightarrow
        G^+(F)\backslash(D\times G^+(\mathbf A_f))/J^+,
              \qquad [x]\longmapsto[(o,x)].         \tag{HJ13}
\]

It is finite because H is definite. It is injective. Suppose two displayed points agree. The relating rational element \(\gamma\) fixes o at every real embedding; hence it preserves the E-line \(Ee_{a+1}\), and has the form \(\operatorname{diag}(h,t)\), with \(h\in H(F)\) and \(t\in E^1(F)\). The finite level condition makes t an integral unit at every finite place. It is a root of unity by the same bounded-powers argument. At G the last diagonal residue of an element of (HJ10) is one, so t=1. Thus h gives exactly the original H double-coset equivalence, proving injectivity.

The same argument with the two representatives equal identifies the orbifold stabilizer at \([(o,x)]\) with

\[
                   \Gamma_x=H(F)\cap xJ_0x^{-1},   \tag{HJ14}
\]

including its action on the fibre and the finite character. There is no extra central U(1) stabilizer and no quotient by an additional root-of-unity group. In particular all compatible coefficient rows are legitimate equivariant jet data at these packet points. This establishes the actual finite input fit for the detector, rather than assuming that independently chosen packet values descend.

\subsection{5. The weight direction is an ample canonical twist}\label{the-weight-direction-is-an-ample-canonical-twist}

Hold u, v and b0 fixed at every label. Set, for an integer \(\ell\ge0\),

\[
 b_{1,\tau}(\ell)=b_{1,\tau}(0)+(a+1)\ell,\qquad
 \beta_\tau(\ell)=\beta_\tau(0)-(a+1)\ell.          \tag{HJ15}
\]

The compact type \(\lambda\), minimal positive weight \(\alpha\), and all jet degrees \(N_\tau\) stay fixed. The separation in (HJ4) only improves. If one also wants the elementary interlacing \(b_1<e_1\), increase e1 by the same amount. This statement alone is not a full SizeData2 assertion.

At a real place write a maximal compact element as \((A,z)\). Its determinant on the tangent space is \(\det A\,z^{-a}\); the canonical line has character

\[
       \kappa(A,z)=\det A^{-1}z^a,\qquad
       z^{a+1}=\det(\operatorname{diag}(A,z))\,\kappa(A,z).
                                                        \tag{HJ16}
\]

The bundle with fibre \((W_\alpha\boxtimes z^{\beta(\ell)})^\vee\) consequently changes by \(\kappa^\ell\) and the flat determinant character \(\det_{G^+}^{\,\ell}\). By (HJ12) the latter character is trivial on every arithmetic component group. Hence, as a holomorphic bundle on that component, the change is its canonical automorphic line to the \(\ell\)-th power. This is a statement about the quotient bundle, not an equality of maximal-compact representation weights: the flat determinant factor in (HJ16) must be retained when reading the real type. It is precisely why \(\alpha\) stays fixed.

On a neat finite cover choose a smooth projective toroidal compactification \(\overline X\) with reduced boundary \(D_\infty\) and morphism \(\pi:\overline X\to X^*\) to Baily--Borel. The canonical automorphic line extends to \(K_{\overline X}(D_\infty)\); it is the pullback of an ample line bundle on \(X^*\), after a fixed positive power if necessary for descent to a coarse quotient. These are Mumford's Main Theorem 3.1 and Proposition 3.4(a),(b), printed pp.256--257 and 264--265. Take such a fixed power q and henceforth allow \(\ell\) in one progression modulo q. Nothing depends on altering J+ or J0. The neat cover is used to prove descent statements; the constructed forms will have the original level.

\subsection{6. Cusp sections separate the selected finite jets}\label{cusp-sections-separate-the-selected-finite-jets}

Here is the finite-jet argument, including the cuspidal condition. On the neat cover let \(\overline{\mathcal E}_0\) be the canonical extension of the initial automorphic vector bundle, including the finite character after passage to its kernel. Its subcanonical extension is

\[
       \mathcal E_0^{\rm sub}
                  =\overline{\mathcal E}_0(-D_\infty).
                                                        \tag{HJ17}
\]

Its global sections are vector-valued holomorphic cusp forms. Canonical extension and the boundary condition are the established automorphic-bundle results of Mumford (1977), Proposition 3.3 and the proof of Corollary 3.5, and Lan and Stroh (2014), introduction and Section 1. In the present F-rank-one case this condition also has a direct Fourier interpretation: vanishing on the toroidal boundary removes the constant term at each rational cusp. We need no higher direct-image vanishing theorem.

Let L be the ample line on \(X^*\) furnished in Section 5 and put \(\mathcal F=\pi_*\mathcal E_0^{\rm sub}\). It is coherent by properness. Over the interior, \(\pi\) is an isomorphism, so \(\mathcal F\) is the desired automorphic vector bundle there. Let \(Z_{\mathbf N}\) be the finite infinitesimal scheme at all packet preimages imposing all Taylor coefficients through degree \(N_\tau\) in each ball variable. One may instead use all total degrees at most \(\sum_\tau N_\tau\), which contains these data. It is disjoint from the boundary.

Let \(\mathcal K\) be the kernel of the coherent restriction \(\mathcal F\to\mathcal F|_{Z_{\mathbf N}}\). For all sufficiently large n, Serre vanishing gives

\[
 H^1(X^*,\mathcal K\otimes L^n)=0.
                                                        \tag{HJ18}
\]

The long exact sequence gives the surjection

\[
 H^0(X^*,\mathcal F\otimes L^n)
       \longrightarrow
 H^0(Z_{\mathbf N},\mathcal F\otimes L^n|_{Z_{\mathbf N}}).
                                                        \tag{HJ19}
\]

The projection formula identifies the source with the cusp sections of the appropriate weight on the toroidal compactification. Serre vanishing is used only in this elementary coherent-sheaf form; see Grothendieck (1961), EGA III.1, Theorem (2.2.1)(ii), printed p.100. The same formulation is The Stacks Project, Lemma 30.17.1, equivalence (1)--(2).

Take the neat cover normal and deep enough to kill the finite character. Its finite deck group acts on (HJ19), on the line linearization, and on all the packet preimages. Averaging in characteristic zero makes passage to its specified character and invariant subspace exact. Thus (HJ19) descends to the original level and gives every jet satisfying precisely the stabilizer conditions (HJ14). This does not refine the actual bottom level.

At a packet point the selected homogeneous quotient of the jet space is

\[
 \bigotimes_\tau
 \left[(W_{\alpha_\tau}\boxtimes z^{\beta_\tau})
       \otimes\operatorname{Sym}^{N_\tau}
                 (\mathbf C^a\boxtimes z^{-1})\right]^\vee
    \supset W_\lambda^\vee\boxtimes
                       z^{-\boldsymbol\beta+\mathbf N}.
                                                        \tag{HJ20}
\]

The direct summand has multiplicity one by (HJ8--9). Its restriction to the actual stabilizers is exactly the dual initial coefficient module. Projecting (HJ19) to this summand remains surjective. For all sufficiently large n, therefore, all its compatible packet values arise from genuine cusp forms of the required discrete series \(\sigma_\ell^\vee\).

For completeness, the last type assertion uses the holomorphic discrete-series realization, not a dimensional guess. The minimal K type in the section construction is \(\tau_{\min,\ell}^\vee\), and the antiholomorphic lowering algebra annihilates its section vectors. Their generated Harish--Chandra modules are the corresponding generalized-Verma quotients. In the regular discrete-series chamber of (HJ4) this is the irreducible holomorphic discrete-series module, whose free polynomial K structure is (HJ7) dual. The cusp forms are in \(L^2\) and decompose into copies of that representation. Large ell remains in this chamber. We impose no Hecke eigenvector or irreducible finite-adelic constituent restriction.

\subsection{7. From jets to the actual compact restriction and period}\label{from-jets-to-the-actual-compact-restriction-and-period}

One must distinguish a jet of the minimal-type section from evaluation of its selected higher K-type vector. This distinction does not spoil surjectivity, but carries a real normalization factor.

Fix at each label a nonzero Pieri inclusion \(j_\tau:W_{\lambda_\tau}\to
W_{\alpha_\tau}\otimes\operatorname{Sym}^{N_\tau}T_oD\). The Poincaré--Birkhoff--Witt map into \(\sigma_\ell\) is injective in this discrete-series range by (HJ7). Its restriction to \(j_\tau\) is therefore nonzero. After choosing invariant Hilbert norms it is a nonzero scalar times the isometric selected K-type inclusion. Denote the scalar by \(b_{\ell,\tau}\ne0\); its absolute value is the norm of that fixed PBW image of a unit vector. No value one is assigned.

The transpose operation on holomorphic cusp sections is a differential operator of order \(N_\tau\). Its leading symbol is the transpose of \(j_\tau\), times the inverse of the selected normalization scalar (and the sign dictated by the chosen Lie-action convention). It is nonzero. This can be checked in Harish--Chandra coordinates: each noncompact Lie generator acts by a first-order operator whose principal part is the corresponding constant directional derivative at o; a product of \(N_\tau\) generators has the symmetric product as principal symbol. Lower-order terms come only from differentiating the frame multiplier. The prescribed Pieri projection gives precisely the indicated selected symbol. This argument also proves that no different derivative order can substitute for (HJ9).

Prescribe all lower jets to be zero, and prescribe the selected multidegree jet with the inverse of this nonzero symbol factor. Equation (HJ19) then realizes any desired selected higher-K-type coefficient value at every packet point. More generally its triangular lower-jet correction can be inverted because the top symbol is surjective onto this summand. Thus the actual restriction map

\[
 \operatorname{Hom}_{G^+_\infty}
       \bigl(\sigma_\ell^\vee,
         \mathcal A_{\rm cusp}(G^+,J^+,(\chi^+)^{-1})\bigr)
  \longrightarrow
        \mathcal A_{\lambda^\vee}(J_0,\chi_0^{-1})
                                                        \tag{HJ21}
\]

is surjective for all ell sufficiently large in the progression of Section 5. The map uses the fixed isometric inclusion of the selected \(W_\lambda^\vee\) in \(\sigma_\ell^\vee\), then the literal restriction along (HJ1). It does not insert a rational or real translation to force a nonzero value.

In particular let \(P\ne0\) be any initial row, including any nonzero Mram row, and w any unit seed. Choose the target row in (HJ21) to be the Hermitian conjugate of P. The resulting actual cusp form F in the selected \(W_\lambda^\vee\) has

\[
                F|_{[H]}=\overline{f_{P,w}}.       \tag{HJ22}
\]

Its bilinear bottom period is exactly

\[
 \boxed{\quad
   \int_{[H]}f_{P,w}(h)F(\operatorname{diag}(h,1))\,dh
       ={v_0\over d_\lambda}
             \sum_x{ \|P_x\|^2\over|\Gamma_x|}
       >0.\quad}                                  \tag{HJ23}
\]

This equality is an interpolation prescription on a freely chosen bottom cusp form; it is not a change of normalization of either prescribed theta kernel. The factors \(b_{\ell,\tau}\), the norm of a right inverse of (HJ19), and the Petersson norm of F are not estimated. Normalizing F to norm one divides (HJ23) by its actual positive norm; there is no uniform lower bound for that normalized period asserted here.

For one sufficiently large ell, (HJ21) is surjective on the entire finite initial coefficient space, so the same ell works for all initial rows, not merely one chosen vector. Taking Hilbert adjoints, the corresponding period/Poincaré operator is injective on that finite coefficient space. This last statement refers to the explicitly constructed period operator (HJ21), not automatically to the source operator D.s.Pu.

The same conclusion is uniform over any fixed finite collection of admissible word profiles: use the union of their finite jet conditions and the maximum of their finitely many Serre thresholds. When v and b0 are fixed across the words, the minimal-type bundle is already common and only the selected jet degrees/components vary. This is a qualitative uniform statement over the source's finite word set; it does not estimate that maximum.

\subsection{8. Exact remaining gate to the source endpoint}\label{exact-remaining-gate-to-the-source-endpoint}

The current result supplies an actual bottom representation, finite level, finite character, compact packet, selected type, and nonzero bottom period. It removes the need to guess that some initial vector is compatible with this particular geometric detector: every compatible row is detected.

The following identifications have not been proved here.

\begin{enumerate}
\def\labelenumi{\arabic{enumi}.}
\tightlist
\item
  The displayed \((J^+,(\chi^+)^{-1})\), lattice frame and coordinate embedding have not been identified with rJ/rChi and the prescribed finite translate in a completed actual RealPack. PairForm.lean42--45 defines rJ via dat0, betaY and levR; these pins contain more data than (HJ11).
\item
  A cusp form furnished by (HJ21) has not been shown to lie in the actual prescribed downward theta image from the terminal input. R016's free-data seesaw and LIFT27's full-G2 injectivity do not supply this.
\item
  Even an abstract identification of the full bottom cusp space does not construct BotDat.Pin, ValLastX, or the same-family ChainOnto/LastX transport needed by MainX. The existing \nolinkurl{ExistsCentralBottomDatum.lean,} \nolinkurl{MatchingTranslatesBottom.lean} and \nolinkurl{PlacedBesselCompatibility.lean} constructors used in this route explicitly assume \(1\le rS\); m=2 has rS=0.
\item
  The threshold for ell is ineffective in this proof. It has not been fitted into the b/e and SizeData2 windows. Root R021 separately rules out the previously equally spaced family as a full SizeData2 instance; no consequence here repairs that obstruction. No derivative gain, saturation, global row content, or complete H cost is inferred from (HJ23).
\end{enumerate}

Thus a source PuNZ conclusion would require the actual detector identification in item 1 together with the fixed real/Poincaré normalizations; a terminal nonzero conclusion additionally requires items 2--3. These are concrete comparisons, not an assumption that all bottom cusp forms are automatically in the stipulated image.

\subsection{9. Sources, checks, and status}\label{sources-checks-and-status}

Primary mathematical inputs used with the indicated scope:

\begin{itemize}
\tightlist
\item
  Harris, M., Kobayashi, T., \& Speh, B. (2025, September 21). Translation functors, branching problems, and applications to the restriction of coherent cohomology of Shimura varieties. Definition 2.7, equation (2.24), printed p.16; the opposite holomorphic convention and contragredient rule are explicit in Sections 10.3.1 and 11.2. \url{https://www.ms.u-tokyo.ac.jp/~toshi/texpdf/tk2025l-HKS_250921.pdf}
\item
  Howe, R., \& Kraft, H. (1997). Principal covariants, multiplicity-free actions, and the K-types of holomorphic discrete series. In Geometry and representation theory of real and p-adic groups (pp.147--161). Birkhäuser. Opening general discussion, equation (3), preprint pp.1--2. \url{https://doi.org/10.1007/978-1-4612-4162-1_8} The original paper's opening text was directly inspected in the transcription hosted by Academia.
\item
  Macdonald, I. G. (1995). Symmetric functions and Hall polynomials (2nd ed.). Oxford University Press. Chapter I, (5.16), p.73. \url{https://math.berkeley.edu/~corteel/MATH249/macdonald.pdf}
\item
  Mumford, D. (1977). Hirzebruch's proportionality theorem in the non-compact case. Inventiones Mathematicae, 42, 239--272. Main Theorem 3.1; Propositions 3.3--3.4; proof of Corollary 3.5. \url{https://www.dam.brown.edu/people/mumford/alg_geom/papers/1977b--HirzProp-GT.pdf} The relevant original text was inspected through its paperzz transcription because the author and repository PDF hosts returned errors in this session; no PDF inspection is claimed.
\item
  Lan, K.-W., \& Stroh, B. (2014). Relative cohomology of cuspidal forms on PEL-type Shimura varieties. Algebra \& Number Theory, 8(8), 1787--1799. Introduction and Section 1 for the subcanonical cusp-form interpretation; its stronger direct-image theorem is not needed. \url{https://msp.org/ant/2014/8-8/ant-v8-n8-p01-s.pdf}
\item
  The Stacks Project Authors. Ample invertible sheaves and cohomology, Lemma 30.17.1 (Tag 0B5U; contained in Section 01XO). \url{https://stacks.math.columbia.edu/tag/0B5U}
\item
  Grothendieck, A. (1961). Éléments de géométrie algébrique: III. Étude cohomologique des faisceaux cohérents, première partie. Publications Mathématiques de l'IHÉS, 11, 5--167. Theorem (2.2.1)(ii), p.100. \url{https://www.numdam.org/item/PMIHES_1961__11__5_0.pdf} The original PDF was recovered and the theorem directly inspected in its extracted text, work/agents/lift\_refs/ega\_iii\_1.txt.
\end{itemize}

Carried accepted run inputs: K031/K032 actual hosts and finite primes; R015's source signed HC convention; K038's strict initial row level and character; LIFT27's realization of full Mram and exact compact mass; R021's separate failure of the previous weight cloud.

Author whole-proof review completed 2026-11-22T02:29:00+00:00: PASS in the stated qualitative fixed-level bottom-detector scope. Every section and HJ1--23 was read; the HC signs, exact Pieri degree, root-of-unity packet/stabilizer argument, determinant-flat correction, coherent jet exact sequence, finite-level descent, and higher-K-type principal symbol were checked. The actual source initial lamTrue formula was also directly checked. No numerical Serre threshold, finite norm bound, source rJ identification, or Lean verification is claimed.

\clearpage

\section{AM-K040 --- nonuniform CM stage characters with fixed finite data}\label{am-k040-nonuniform-cm-stage-characters-with-fixed-finite-data}

Author: GPT-6 Astra Pro, ALP\_MACHINE \nolinkurl{/root/kummer_extension}, 22 November 2026 UTC; inherited reasoning effort is not separately exposed.

This finding removes the character-existence restriction to uniform archimedean weights in K031/K032. On either actual host, every vector e in (28 Z)\^{}Psi has an everywhere finite-unramified unitary Hecke character alpha\_e, trivial on the full base-field idele group, with angular infinity exponents −2e\_tau. The construction is coherent in e and has the exact uniform-direction identity

\[
                  \alpha_{c\mathbf1}=\eta^{-2c}
                     \qquad(c\in28\mathbb Z).
 \tag{NC1}
\]

Multiplying by the already constructed zero-infinity G character therefore supplies the actual stage characters for arbitrary such nonuniform profiles, without adding finite ramification or changing the prescribed tame G characters. One finite Galois coefficient field works for all profiles. The proof explicitly handles compatibility with principal ideals extended from the real field; checking units alone would not establish the required base-field restriction.

This is an arithmetic character theorem. It does not establish the new profile's cloud inequalities, full SizeData2, selected content, terminal nonvanishing, or the remaining DP44 cost window.

\subsection{1. Hosts and literal infinity convention}\label{hosts-and-literal-infinity-convention}

Use the fields and actual CM types of K029--K032:

\[
 E_0=\mathbb Q(\zeta_{28}),\quad
 K_1=E_0B,\quad [B:\mathbb Q]=7,\quad B\subset\mathbb Q(\zeta_{29}),
 \qquad K_4=K_1(\sqrt5,\sqrt{13}).
 \tag{NC2}
\]

Let F\_h=K\_h\^{}+ and let Psi\_h be the actual induced CM type of the principal variety J\^{}(14h), with f=42h labels. Both CM extensions are unramified at every finite place. Retain the actual odd character eta\_h from K029 and its norm pullback in K032:

\[
 (\eta_h)_\infty(z)=\prod_{\tau\in\Psi_h}z_\tau/|z_\tau|,
 \qquad \eta_h|_{\mathbb A_{F_h}^\times}=\eta_{K_h/F_h}.
 \tag{NC3}
\]

The right side is quadratic, and eta\_h is everywhere finite-unramified. Its normalized etaVal values belong to the already constructed fixed field. This accepted result is not rederived here.

The exact source definitions are \nolinkurl{AbelianHostFieldConditions.lean}, lines314--331: IsMu0 requires unitary, TrivOn F, HasInfQ(−e), unramified away from the two G primes, tame everywhere, and the prescribed inverse residue character at each selected G prime. Chars.Tied additionally requires all unramified finite values in one coefficient field and retains the same etaT/etaVal. In \nolinkurl{Formalism/LocalToGlobal/UnitaryNorms/IsUnitaryNormOf.lean}, lines90--94, HasInfQ(a) means the \textbf{angular integer exponent 2a}. Thus the needed infinity component for e is

\[
 b_e(z)=\prod_{\tau\in\Psi_h}
              (z_\tau/|z_\tau|)^{-2e_\tau}.
 \tag{NC4}
\]

All e\_tau below may be positive, negative, or zero. They are multiples of28. No uniformity over labels is imposed.

\subsubsection{Root-of-unity fact}\label{root-of-unity-fact}

\[
                         \mu(K_h)=\mu_{28}\quad(h=1,4).
 \tag{NC5}
\]

For K\_1 this is K029(ER4--ER5). For K\_4, the added ramification indices at5 and13 are2, whereas the indices at2,7,29 remain2,6,7. A primitive ell\^{}a-th root would force phi(ell\^{}a) to divide the rational ramification index at ell. Thus neither fifth nor thirteenth roots can occur; an eighth, forty-ninth, or twenty-ninth root is excluded by the same test. All other odd primes are unramified. The order therefore divides28, and E\_0 already contains mu28. This proves (NC5).

\subsection{2. Exact principal/base-ideal compatibility}\label{exact-principalbase-ideal-compatibility}

Write I\_K for the fractional ideal group, P\_K for its principal subgroup, and j:I\_F→I\_K for extension of ideals. Define

\[
 H=P_Kj(I_F),\qquad
 C_{K/F}=I_K/H=\operatorname{Cl}(K)/j\operatorname{Cl}(F).
 \tag{NC6}
\]

This quotient is finite; no value of either class number is assumed. For a in K\^{}times put

\[
 T_e(a)=\prod_{\tau\in\Psi}
                    \tau(a/\bar a)^{e_\tau}\in K^\times.
 \tag{NC7}
\]

Here K is Galois over Q, so all its embeddings have image K after the fixed complex embedding. Each complex value of T\_e(a) has absolute value1. Equation (NC4) gives T\_e(a)=b\_e(a)\^{}−1.

\subsubsection{Lemma NC8}\label{lemma-nc8}

If (a) belongs to j(I\_F), then T\_e(a)=1. Consequently

\[
           c_e((a)j(\mathfrak b))=T_e(a)
                    \qquad((a)j(\mathfrak b)\in H)
 \tag{NC8}
\]

is a well-defined homomorphism H→K\^{}times, trivial on j(I\_F), with values of complex absolute value1.

\textbf{Proof.} If (a)=j(b), that ideal is fixed by complex conjugation. Thus (a/bar a)=1, and a/bar a is a global unit. Every complex embedding of this unit has absolute value1. The elementary CM unit argument, proved in K029 section1, makes it a root of unity: the monic polynomials of its positive powers have bounded integer coefficients, so two powers coincide. By (NC5) it belongs to mu28. Each exponent e\_tau is divisible by28; every factor of T\_e(a) is therefore1, proving the first assertion.

If (a)j(b)=(a')j(b'), then (a/a')=j(b'/b). Apply the first assertion to a/a' to obtain T\_e(a)=T\_e(a'). This proves well-definedness of (NC8). Multiplicativity follows from (NC7), and a=1 proves its triviality on j(I\_F). QED.

In particular this handles any capitulating real ideal class. It does not first extend a character on P\_K and then presume that the resulting restriction to F is trivial. The compatibility is checked on the full intersection P\_K intersect j(I\_F) before extension.

\subsection{3. Explicit extension across the finite relative class quotient}\label{explicit-extension-across-the-finite-relative-class-quotient}

Choose an invariant-factor decomposition

\[
 C_{K/F}=\bigoplus_{s=1}^q\mathbb Z/d_s\mathbb Z,
 \qquad d_s>1,
 \qquad H_{\rm rel}=\prod_s d_s,
 \quad d_{\rm rel}=\operatorname{lcm}_s d_s.
 \tag{NC9}
\]

For the trivial group use H\_rel=d\_rel=1. Choose fractional ideals A\_s representing the generators, and fixed relations

\[
                  A_s^{d_s}=(a_s)j(\mathfrak b_s).
 \tag{NC10}
\]

Choose z\_(s,e) with z\_(s,e)\^{}(d\_s)=T\_e(a\_s). Its chosen complex value has absolute value1, as do all d\_s-th roots. Every ideal can be written uniquely in quotient-exponent ranges as

\[
 I=h\prod_s A_s^{r_s},\qquad h\in H,\quad0\le r_s<d_s.
\]

Define

\[
                 \widetilde c_e(I)
                    =c_e(h)\prod_s z_{s,e}^{r_s}.
 \tag{NC11}
\]

This is a homomorphism: when multiplying two ideals, each carry of d\_s moves A\_s\^{}(d\_s) into H and contributes exactly the same factor T\_e(a\_s)=z\_(s,e)\^{}(d\_s). It extends c\_e and is trivial on j(I\_F). Every value belongs to the field generated by K and the finitely many chosen roots.

For a finite idele x\_f set I(x\_f)=product\_P P\^{}(v\_P(x\_P)). Then

\[
             \alpha_e(x)=b_e(x_\infty)\widetilde c_e(I(x_f))
 \tag{NC12}
\]

is continuous, unitary, everywhere finite-unramified, and trivial on K\^{}times. The last assertion follows from c\_e((a))=b\_e(a)\^{}−1. Continuity follows because the finite ideal map is locally constant and b\_e is continuous. For an F-idele, its finite ideal extends from F and so has value1. At infinity all its components are real, and the even angular exponents in (NC4) make b\_e identically1 there. Thus

\[
                   \alpha_e|_{\mathbb A_F^\times}=1.
 \tag{NC13}
\]

Equations (NC8)--(NC13) prove existence with the exact base restriction. They use only the finite relative ideal quotient, not class number one of the enlarged host or a character extension with unverified capitulation compatibility.

\subsection{4. A coherent family with the prescribed odd character}\label{a-coherent-family-with-the-prescribed-odd-character}

Independent choices in (NC11) would not automatically respect addition of weight vectors. We now make coherent choices and enforce the exact uniform relation (NC1).

Enumerate Psi=\{tau\_1,\ldots,tau\_f\}. For i\textless f, construct once a character xi\_i using (NC12) for the basis vector e=28 e\_i. Put

\[
             \xi_f=\eta_h^{-56}\prod_{i<f}\xi_i^{-1},
 \qquad
             \alpha_e=\prod_{i=1}^f\xi_i^{e_{\tau_i}/28}.
 \tag{NC14}
\]

Since eta\_h has the quadratic base restriction (NC3), its even power eta\_h\^{}−56 is trivial on the full F-idele group. Therefore xi\_f is also trivial there. All xi\_i are finite-unramified and unitary. The infinity exponent of eta\_h\^{}−56 is−56 at every label; dividing by the first f−1 basis characters leaves exponent−56 exactly at tau\_f.~Thus xi\_f has the required last basis type, and alpha\_e has (NC4) at every label. All finite conditions remain as proved above.

The construction is a homomorphism from the free weight lattice:

\[
       \alpha_{e+e'}=\alpha_e\alpha_{e'},\qquad
       \prod_i\xi_i=\eta_h^{-56},\qquad
       \alpha_{c\mathbf1}=\eta_h^{-2c}\quad(c\in28\mathbb Z).
 \tag{NC15}
\]

In particular, for two nonuniform profiles with

\[
       e_{1,\tau}-e_{0,\tau}=2\ell
                     \quad\text{for every tau},\qquad2\ell\in28\mathbb Z,
\]

one has the \textbf{identity of full Hecke characters}

\[
                         \alpha_{e_1}/\alpha_{e_0}=\eta_h^{-4\ell}.
 \tag{NC16}
\]

It is not merely an equality of their infinity components up to an unspecified finite class character.

\subsection{5. The actual fixed-G stage characters}\label{the-actual-fixed-g-stage-characters}

Let nu\_j be the already constructed zero-infinity finite-order character of E\_0 with the chosen G data, and let

\[
       \widetilde\nu_j=\nu_j\circ N_{K_h/E_0},\qquad
       \boxed{\ \mu_j=\widetilde\nu_j\alpha_{e_j}.\ }
 \tag{NC17}
\]

Use exactly the existing K031 or K032 G primes and local characters. Their completely split degree-one completions make each local norm the identity. Thus at every selected G place P,

\[
                 \mu_{j,P}|_{\mathcal O_{K_h,P}^\times}
                                      =\chi_{j,P}^{-1}.
 \tag{NC18}
\]

The conjugate place has the inverse character, as required by triviality on A\_F. Nothing in alpha\_e changes a compact finite restriction, because alpha\_e is unramified at every finite place. Consequently the exact tame conductor remains

\[
       \mathfrak f(\mu_j)=g_0g_1\mathcal O_{K_h},\qquad
       N\mathfrak f(\mu_j)=(g_0g_1)^{2f}
                       \quad(j=0,1),
 \tag{NC19}
\]

for the original nontrivial distinct G characters. In particular no ramification is added at29,11,T,2,7,5, or13. Unitarity, TrivOn F, HasInfQ(−e\_j), the exact atg clause and every finite unramified/tame clause in literal IsMu0 all follow. The same eta\_h and etaVal are retained in Chars.Tied. The numerical sum e\_even is automatic since each e\_(j,tau) is a multiple of28.

For the constant stage difference above, (NC16) gives

\[
          \mu_1\mu_0^{-1}
          =(\widetilde\nu_1\widetilde\nu_0^{-1})\eta_h^{-4\ell}.
 \tag{NC20}
\]

Thus the exact common relative character used in the two-stage analytic comparison is preserved. A common finite tuple twist, if subsequently applied as in the original construction, cancels from this ratio and retains its own separately prescribed Q data.

\subsection{6. One coefficient field for every profile, with an explicit bound}\label{one-coefficient-field-for-every-profile-with-an-explicit-bound}

Let M\_base be a finite Galois field containing K\_h, the existing normalized etaVal field, and the finite G roots. It is fixed before the profiles e\_j are chosen. Write L\_h=lcm(28,g\_0−1,g\_1−1) for the latter root exponent.

For each relative class relation (NC10) and every selected tau, choose a root, using the root already chosen for the corresponding basis character when that character is among the first f−1:

\[
          z_{s,\tau}^{d_s}=\tau(a_s/\bar a_s)^{28}.
 \tag{NC21}
\]

All finite values of the first f−1 basis characters lie in K\_h(z\_(s,tau):s,tau), by (NC11). The last basis character in (NC14) also takes values in the compositum with M\_base: on an unramified prime P, eta\_h(P)\^{}−56 equals

\[
       \bigl(\eta_h(P)(NP)^{-1/2}\bigr)^{-56}(NP)^{-28},
 \tag{NC22}
\]

which is in the old etaVal field. Multiplication extends this to every ideal. Hence

\[
 M=M_{\rm base}\bigl(\mu_{d_{\rm rel}},
                             z_{s,\tau}:s,\tau\in\Psi\bigr)
 \tag{NC23}
\]

contains every finite value of every alpha\_e and mu\_j, for \textbf{all} e in (28 Z)\^{}Psi. It therefore supplies the literal muVal field, not merely a different algebraic extension for each prime or weight.

This M is Galois over Q. Indeed K\_h and M\_base are Galois, and complex conjugation is central on K\_h. A Q-automorphism sends a radicand in (NC21) to the corresponding radicand for another selected embedding or its inverse. All its d\_s-th roots are already present because mu\_(d\_rel) is present. Every conjugate of each generator therefore lies in M. Character coherence is not changed by adjoining the unused last-label roots.

The field degree satisfies

\[
       [M:M_{\rm base}]
       \le\varphi(d_{\rm rel})\prod_s d_s^f
       =\varphi(d_{\rm rel})H_{\rm rel}^{f}
       \le h(K_h)^{f+1}.
 \tag{NC24}
\]

This gives a profile-independent bound without computing the large class group. A wholly numerical bound follows from the known discriminants and the standard Minkowski theorem. The exact absolute discriminants are

\[
 |\Delta_{K_1}|=2^{84}7^{70}29^{72},\qquad
 |\Delta_{K_4}|=2^{336}7^{280}29^{288}5^{168}13^{168}.
 \tag{NC25}
\]

The first follows from \textbar Delta\_(E0)\textbar=2\^{}12·7\^{}10, \textbar Delta\_B\textbar=29\^{}6 and the coprime-discriminant compositum. The second uses \textbar Delta\_(Q(sqrt5,sqrt13))\textbar=5\^{}2·13\^{}2 and the same formula. These are the integral composita already proved in K030/K032.

Minkowski's ideal-class representative bound is (4/pi)\^{}f (2f)!/(2f)\^{}(2f) sqrt(\textbar Delta\_K\textbar); see Milne (2020), Theorem4.3, printed p.70. In particular every class has an integral ideal of norm at most

\[
 B_1=2^{84}7^{35}29^{36},\qquad
 B_4=2^{336}7^{140}29^{144}65^{84},
 \tag{NC26}
\]

because these are exactly 2\^{}f sqrt(\textbar Delta\_(K\_h)\textbar). The number of integral ideals of norm m in a degree-n field is at most d\_n(m), the number of ordered n-tuples of positive integers with product m. To prove the bound, at a rational prime p the ideal-counting Euler factor is product\_(P\textbar p)(1−T\^{}(f\_P))\^{}−1. Coefficientwise it is bounded by (1−T)\^{}−n, since there are at most n primes and all f\_P≥1. Multiplication over p proves the assertion. Also d\_n(m)≤m\^{}(n−1), by choosing its first n−1 entries. Therefore

\[
                  h(K_h)\le H_h:=B_h^{2f},\qquad
                  [M:M_{\rm base}]\le H_h^{f+1}.
 \tag{NC27}
\]

These crude bounds are finite, explicit, and independent of the new weight profile. They are not proposed as economical arithmetic costs.

One can also bound the previously existing M\_base rather than treating its degree as a fixed input. In K029(ER25--ER26), the etaVal construction adjoins, for each cyclic generator of Cl(K\_1), one root of degree its class order, and the correcting class character has order dividing h(K\_1). For each root relation, adjoin one root of every conjugate radicand under the84 automorphisms of K\_1, and also mu\_(h(K\_1)); this contains every conjugate of each chosen root. The resulting normal field has degree at most h(K\_1)\^{}85 over K\_1. The normalized etaVal field for K\_4 is contained in the same field by its norm-pullback formula. Hence the deliberately loose completely numerical bound

\[
             [M:\mathbb Q]
                  \le 2f\,L_h\,H_1^{85}\,H_h^{f+1}
 \tag{NC28}
\]

is available. In this formula H\_1=B\_1\^{}84 refers to the original degree84 field in (NC26), even when h=4. No class number, regulator, or discriminant of the coefficient field has been computed here.

The bounded \textbf{field degree} does not bound all finite values by a constant height. Profile-dependent values are powers of finitely many fixed basis characters. At each fixed finite place P their logarithmic heights are bounded above by the sum of \textbar e\_tau\textbar/28 times the heights of those basis values, plus the fixed finite G factor. An effective global height/error statement would still need those arithmetic constants. No such cost is suppressed in (NC24)--(NC28).

\subsection{7. What this repairs in the new profile}\label{what-this-repairs-in-the-new-profile}

Arbitrary top/thick/thin assignments of e\_tau in28Z are covered, including negative assignments, independently at all42 or168 CM labels. If the second stage is obtained by the same increment2ell at every label, with2ell in28Z, the exact stage-ratio formula(NC20) is available. The G primes, their generic tame characters and residue fields, the odd etaT, the maximal29 level, and unramified 11 components are unchanged.

Thus nonuniform stage-character existence and a common algebraic value field impose no additional obstruction on that proposed profile. The field/type constructor's old common-quadratic failure is not altered. The profile still needs its own numerical proof and the full norm/height/terminal-image analysis; this theorem does not certify those independent requirements.

\subsection{References and check scope}\label{references-and-check-scope}

Milne, J. S. (2020). \emph{Algebraic number theory} (Version3.08). Theorem4.3, printed p.70; Theorem4.4, printed pp.70--71; Lemma6.5 and Remark6.6(c), printed pp.100--101. \url{https://www.jmilne.org/math/CourseNotes/ANT.pdf}. Theorem4.3 was directly read from the primary author PDF in this run. The two latter compositum statements are retained from the checked K030/K032 application. The finite relative-ideal extension, base restriction, coherent eta relation and all field-degree deductions are proved above.

Source conventions: \nolinkurl{AbelianHostFieldConditions.lean},314--331; \nolinkurl{IsUnitaryNormOf.lean},90--94. Carried inputs: K029(ER2,ER5,ER25--26), K031(HC16--21), K032(QH27--31), and LIFT18's zero-infinity finite character. No new class-number-one assertion or analytic L-function theorem is used. The companion exact certificate checks root-order exclusions, finite quotient carry identities, coherent type exponents, conductor and discriminant arithmetic, the exact fixed G root exponents, and the stated field bounds as formulas. It does not compute the large class groups, their ideals, Hecke values or coefficient heights, and it does not run Lean.

\clearpage

\section{AM-DO-040: two certified replacement T-primes for the large weight profile}\label{am-do-040-two-certified-replacement-t-primes-for-the-large-weight-profile}

Author: GPT-6 Astra Pro, ALP\_MACHINE \nolinkurl{/root/definite_oscillator}. 22 November 2026 UTC. Inherited reasoning effort is not separately exposed.

This finding replaces the two auxiliary T-primes of R017 by explicit larger primes, retaining the same CM hosts and G-primes. They satisfy T\textgreater16 A for A=2\^{}960, the exact norm and full source congruence, and an explicit projective residue-vector condition. The latter uses powers of the integer 2, with a certified order bound; no primitive generator of the whole residue-field multiplicative group or complete factorization of T−1 is assumed.

A standalone replay proves primality using only integer arithmetic, trial division of the specified factor primes, modular powers, and gcds. The probable-prime discovery was a separate calculation and is not an input to that proof. The result does not claim that these T choices alone fill a complete SizeData₂ record or repair the original quadratic host-field wrapper.

\subsection{1. Exact parameters}\label{exact-parameters}

Retain the checked hosts, conductors and G-primes of K031/K032:

\begin{longtable}[]{@{}lrrrrrr@{}}
\toprule\noalign{}
Host & f=a & c & g0 & g1 & B & m \\
\midrule\noalign{}
\endhead
\bottomrule\noalign{}
\endlastfoot
J\^{}14 & 42 & 812 & 4615409 & 5276377 & 392 & 263 \\
J\^{}56 & 168 & 52780 & 22285932761 & 64072175441 & 334 & 54 \\
\end{longtable}

Let delta=sqrt(−7) in their common external quadratic subfield and set

\[
 D=c(g_0g_1)^2\,2^B m,\qquad
 \pi_T=1+2D+D\delta,\qquad
 T=11D^2+4D+1,\qquad F=4D.                    \tag{DO40.1}
\]

The quantities in this formula are integers except for the explicitly specified quadratic generator in piT. Thus it completely specifies the prime and generator without a search oracle. The exact decimal primes are recorded below and in the replay JSON.

\textbf{J14} (1001 bits, 302 decimal digits):

\begin{verbatim}
17952167255185635051321154351674895613987764293769673456729204696357336089011562929338665498585310931065658058642583455560352833168574896704942121915304013574614983189345694485846080474422936718464345734525547995223578746517703781220872222478207325026191727208735314247959847795300609580930248088748033
\end{verbatim}

\textbf{J56} (996 bits, 300 decimal digits):

\begin{verbatim}
454943247013778122512187780897314598574665539127335259605458539964582491987337758099655723658178473308877448900224099993773653967422266630524027386542260738612598969923877926849959335045299704408598932706573786754076850608752614606900351396549840892191161193206683853263893019307714081695789188382721
\end{verbatim}

The replay verifies D\textgreater=2\^{}490 and 4 divides D in each case. In fact

\[
 \left\lfloor\frac{T}{4\cdot2^{960}}\right\rfloor
 =\begin{cases}
 460\,533\,319\,762&J^{14},\\
 11\,670\,820\,624&J^{56}.
 \end{cases}                                           \tag{DO40.2}
\]

Hence both T\textgreater16 A, more than the requested T\textgreater4 A. This inequality is an exact integer comparison. No floating-point exponential estimate enters the certificate.

\subsection{2. Deterministic primality proof}\label{deterministic-primality-proof}

The exact known factorizations are

\[
\begin{aligned}
 F_{14}&=2^{396}\,7\,29\,263\,4615409^2\,5276377^2,\\
 F_{56}&=2^{339}\,3^3\,5\,7\,13\,29\,
                      22285932761^2\,64072175441^2.
\end{aligned}                                          \tag{DO40.3}
\]

The replay checks each equality with F=4D. It proves every listed factor prime by testing all odd divisors through its integer square root (and handles 2 directly). The largest trial-division bound is only the integer square root of 64072175441. Primality of the large T is proved by the following complete factored-divisor argument.

First,

\[
 T-1=4D\left(1+\frac{11D}{4}\right),\qquad
 F^2-T=(5D+1)(D-1)>0.                                  \tag{DO40.4}
\]

Thus F divides T−1 and F\textgreater sqrt(T). For every distinct prime q dividing F, use the specified witness z in the following table:

\begin{longtable}[]{@{}lll@{}}
\toprule\noalign{}
Host & q & z \\
\midrule\noalign{}
\endhead
\bottomrule\noalign{}
\endlastfoot
J14 & 2 & 5 \\
J14 & 7,29,263,4615409,5276377 & 2 \\
J56 & 2 & 11 \\
J56 & 3 & 3 \\
J56 & 5,7,13,29,22285932761,64072175441 & 2 \\
\end{longtable}

Every one of the exact congruences and gcd identities

\[
 z^{T-1}\equiv1\pmod T,\qquad
       \gcd\left(z^{(T-1)/q}-1,T\right)=1               \tag{DO40.5}
\]

is replayed. The complete quotient-power residues, not only PASS flags, are saved in the JSON.

Here is why these checks prove primality. If p is any prime divisor of T, the multiplicative order of the q-specific z modulo p divides T−1. The second condition in (DO40.5) says it does not divide (T−1)/q. Its q-adic order therefore contains the full q-primary part of T−1, in particular the q-primary part of F. Since that multiplicative order divides p−1, this holds for each q, and F divides p−1. Consequently every prime divisor p of T satisfies p\textgreater=F+1\textgreater sqrt(T). A composite positive integer has a prime divisor at most its square root. Hence T is prime. This also excludes prime powers of exponent at least two.

This is the elementary Pocklington factored-divisor criterion used in R017, with its proof included. No unverified property of the complementary factor 1+11D/4 is needed; it is not asserted prime or fully factored.

\subsection{3. Exact source norm, congruence and residue conditions}\label{exact-source-norm-congruence-and-residue-conditions}

Direct multiplication in Q(delta) gives

\[
 \pi_T\overline\pi_T=(1+2D)^2+7D^2=T,\qquad
 \pi_T\equiv1\pmod{c(g_0g_1)^2}.                       \tag{DO40.6}
\]

Both generators are algebraic integers. The replay also checks

\[
 \gcd(T,2(1+2D))=1,\quad
 \gcd(T,11cg_0g_1)=1,\quad T\equiv1\pmod c.            \tag{DO40.7}
\]

A prime ideal dividing both piT and its conjugate would divide their product T and sum 2(1+2D), contradicting the first equality. Their ideals are therefore coprime, also after extending to either host.

Each checked host is contained in the c-th cyclotomic field. The prime T is coprime to c and congruent to one modulo c, so it splits completely in that cyclotomic field and therefore in the host. At all host primes over T its valuation is one. Equation (DO40.6) and coprimality show that precisely one prime of each complex-conjugate pair divides piT. There are exactly f such primes, each with residue field F\_T and piT-valuation one. These are the norm, congruence, coprimality, residue-norm and half-cardinality conditions used in R017. The G-primes and their previously checked congruences are unchanged.

The literal old \nolinkurl{piT_Kq} clause is not claimed: as K031 recorded, its original source-defined common field is Q on this induced mixed host, whereas piT belongs to the actual external Q(sqrt(−7)). The replacement geometric family retains the latter explicit quadratic subfield. Constructing the remaining generalized source data remains separate from the elementary prime certificate.

\subsection{4. A completely specified residue covector}\label{a-completely-specified-residue-covector}

Both hosts have roots of unity mu28, by the carried R017 host check. Reduction at T is injective on this group. Since T is congruent to one modulo28, its image is the unique subgroup of order28 in F\_T\^{}times. The replay checks, for every integer 1\textless=j\textless=a, that

\[
                         2^{28j}\not\equiv1\pmod T.     \tag{DO40.8}
\]

It follows that the image of 2 in the cyclic quotient F\_T\^{}times/mu28 has order h\textgreater a. In particular the residue covector

\[
                         u=(1,2,2^2,\ldots,2^{a-1})    \tag{DO40.9}
\]

has nonzero coordinates in distinct mu28 orbits. This condition has been proved without factoring all of T−1 or claiming that 2 is a primitive residue-field generator.

The projective stabilizer in the integral monomial group is exactly central mu28. For completeness, suppose a monomial matrix with entries in mu28 maps u to c u. Passing to F\_T\^{}times/mu28 says that translation by the class of c preserves the set I=\{0,1,\ldots,a−1\} in the cyclic group generated by the class of 2. That class of c lies in this cyclic subgroup because one member of I is zero. Around its cyclic ordering, I has a−1 successive gaps of length one and one gap of length h−a+1\textgreater1. The unique larger gap is preserved by a translation stabilizing I, so that translation is zero. Each coordinate is therefore fixed, and every monomial entry equals the same c in mu28. Conversely these scalar matrices clearly preserve the projective line. If u itself is fixed rather than its line, c=1. Thus

\[
 \Gamma\cap P_u=\mu_{28}I_a,\qquad
 \Gamma\cap H_u=1,\quad
 \Gamma=U(I_a)(F)\cap\operatorname{GL}_a(\mathcal O_K).
                                                        \tag{DO40.10}
\]

Here the exact monomial description of Gamma is R017(R17.7). A single T-place already forces these conclusions; imposing all f T decorations does not enlarge either group. This is the same class support needed for R017's matrix-coefficient input, now with the explicit covector (DO40.9). No principal29 refinement is introduced.

The restriction map from residue characters of F\_T\^{}times to characters of mu28 is surjective. Every fibre has (T−1)/28 members, which is larger than one. Even the trivial root-character therefore admits a nontrivial residue-character extension. Choose such nontrivial T characters at f−1 places; at the last choose a nontrivial extension of the root-character needed to cancel the actual archimedean and G central character. This gives precisely R017's central compatibility for any fixed selected archimedean coefficient representation. If these particular character-valued inputs are used as algebraic rows, their finite character values must be included in the coefficient field; adjoining mu\_(T−1) suffices. No cost-free field enlargement is claimed.

\subsection{5. Exact inherited class mass and its T-part at11}\label{exact-inherited-class-mass-and-its-t-part-at11}

Write

\[
 d_T=\frac{T^a-1}{T-1},\qquad
 d_G=\prod_{P\in G}|U_a(\mathbf F_{g_P})|,
 \qquad d_\rho=\dim W_0.
                                                        \tag{DO40.11}
\]

For the actual R017 projective-T/full-G matrix-coefficient support, with its chosen unit real and finite Steinberg vectors, the squared norm is still exactly

\[
                        \frac1{28\,d_T^f d_\rho d_G}.
                                                        \tag{DO40.12}
\]

This is an application of R017's checked mass and Schur-normalization formula to the newly verified stabilizer and level; it does not declare the raw kernel constant on this support. The literal strict refinement remains K038(CA30a--c): with

\[
 I_T=(T-1)^f,\qquad I_G=\prod_{P\in G}
                    [\operatorname{GL}_a(\mathbf F_{g_P}):U_a],
                                                        \tag{DO40.13}
\]

there are I\_T I\_G/28 strict classes of stabilizer one, each with volume v\_P/(I\_T I\_G). Both actual weighted kernel sums are retained. The dimensions and G factors have not been dropped from the norm.

The new T contributions have particularly simple exact11-valuations:

\begin{longtable}[]{@{}lrrrrr@{}}
\toprule\noalign{}
Host & T mod11 & order of T mod11 & v11(T−1) & v11(d\_T) & v11(d\_T\^{}f) \\
\midrule\noalign{}
\endhead
\bottomrule\noalign{}
\endlastfoot
J14 & 9 & 5 & 0 & 0 & 0 \\
J56 & 10 & 2 & 0 & 2 & 336 \\
\end{longtable}

For J14, 5 does not divide a=42, so T\^{}a−1 is nonzero modulo11. For J56, the replay gives v11(T+1)=2; since a=168, the odd-prime lifting formula gives v11(T\^{}a−1)=2+v11(84)=2. Both denominators T−1 are units at11. Thus the T-character coarsening factor 1/I\_T is an11-unit on both hosts, whereas the J56 projective-line mass retains its valuation −336. The unchanged G refinement and variable real dimension factors must still be accounted separately.

\subsection{6. Sources, checks and scope}\label{sources-checks-and-scope}

The full deterministic replay is \nolinkurl{definite_oscillator_large_t_certificate.py}; its JSON contains the exact primes, generators, factorization, every witness residue and gcd, all210 projective-order tests, and the exact11 mass valuations. The separate discovery script is explicitly labelled as such. The proof above explains why the replay establishes primality and the geometric level conditions; numerical checks are not substituted for those implications.

Carried findings: K031/K032's fields, conductors and G-primes; R017(R17.7--14) for the monomial group and matrix-coefficient mass; K038(CA30a--c) for exact strict refinement and weighted averages. The literal source conditions are \nolinkurl{Columns/AbelianHostFieldConditions.lean}206--238 (\nolinkurl{Primes.Ok}) and 486--494 (\nolinkurl{Sizes}), with the common-field clause separated as above.

The prime criterion is the factored-divisor argument of Pocklington, H. C. (1914--1916), ``The determination of the prime or composite nature of large numbers by Fermat's theorem,'' \emph{Proceedings of the Cambridge Philosophical Society, 18}, 29--30, already recorded in R017. Its complete required implication is proved in section2, so no additional inaccessible literature assertion is used here.

These data supply the requested large auxiliary prime and actual class support for the new weight-profile construction. They do not establish all its weight windows, a terminal nonzero period, or the complete arithmetic gain/cost inequality. No Lean execution is claimed.

\clearpage

\section{ALP-LIFT-29: the rank-zero sharp bottom level and an exact T translation detector}\label{alp-lift-29-the-rank-zero-sharp-bottom-level-and-an-exact-t-translation-detector}

Author: /root/lift\_accounting, GPT-6 Astra Pro, ALP\_MACHINE. 22 November 2026 UTC. Reasoning effort is not separately exposed.

This refines LIFT28's convenient bottom level. In the actual m=2, rS=0 ranks, the source bottom sharp residue group is exactly diag(U\_a,1), with its generic character on U\_a. The standard-frame dat0 level is maximal at T. An explicit lower unipotent at each T place makes its intersection with the compact subgroup H exactly the original strict covector level. The translated compact packet has exactly the original classes and stabilizers. LIFT28's holomorphic interpolation argument therefore applies at this source-shaped bottom level, without a new T level restriction.

This gives a concrete candidate for the finite bottom level and translate. The fixed tS, betaY, real model, and downward theta family of a completed source RealPack have not thereby been identified. No assertion about their terminal endpoint is made.

\subsection{1. Exact specialization of the source sharp subgroup}\label{exact-specialization-of-the-source-sharp-subgroup}

Use the recovered source commit 7cbf7a0c32fca1badf05931879a37c3f8a2b03b4. Let m=2, a=k, r=k+m-1=a+1, and rS=0. All finite residue statements below concern the actual selected relative split G places.

\nolinkurl{Divisibility/DefGroup.lean186–189} defines unipCol(r,A,C) as the upper unipotent subgroup with all entries (i,C) equal to zero for A\textless=i\textless C. Lines258--261 set

\[
 \operatorname{unipSym}(r,\rho,k)
       =\operatorname{unipCol}(r,\rho,\rho+k).
                                                        \tag{BT1}
\]

Thus, at r=a+1, rho=0, k=a, all entries above the last diagonal entry in the last column are zero. Upper triangularity already forces the last row below the diagonal to vanish. Consequently

\[
 \operatorname{unipSym}(a+1,0,a)
                =\{\operatorname{diag}(u,1):u\in U_a\}.
                                                        \tag{BT2}
\]

The symEdge definition at lines218--220 removes the adjacent edge ending at column a; its extra bridge edge would require i+1=0, which is impossible. All and only the adjacent edges within the first a coordinates remain. Its additive character is therefore exactly the usual generic character on U\_a. PsharpSym is defined from these local groups and characters at lines608--609.

The permutation in LinksChar.lean109--110 satisfies IsSymOrd (defined at lines75--76). At rS=0 every one of its cases gives pi(i)=i. There is no residual reordering in this specialization.

Further, LinksDefs.lean430--431 and \nolinkurl{SubgroupCharExtensionOneApply.lean32–43} show that dat0 has n=r=a+1, j=s=0, Q empty and Frm.B. At G the block function blkOf(r,t) is constant, so parLev adds no condition beyond integrality. The character product over inserted blocks is empty, as is the Q product; see DefGroup.lean635--637 and704--707. With the standard lattice frame beta=1, the exact finite data are

\[
 J_{B,v}=
 \begin{cases}
   \operatorname{red}^{-1}\operatorname{diag}(U_a(\kappa_v),1),
                                                 &v\in G,\\
   K_{a+1,v},                                   &v\notin G,
 \end{cases}
\quad
 \chi_{B,v}=
 \begin{cases}\psi_{U_a},&v\in G,\\1,&v\notin G.\end{cases}
                                                        \tag{BT3}
\]

Here \(K_{a+1,v}\) is the unitary stabilizer of the standard selfdual lattice; at a relative split place it is the ordinary general linear maximal compact. The character at G is the inverse of the actual initial right character. This statement specializes the source level formula; it does not assert that every other RealPack tie has been constructed.

\subsection{2. The exact T conjugation}\label{the-exact-t-conjugation}

Work at one actual T place. It is relative split and unramified of residue characteristic different from11 in the current constructions. Let \(O\), \(\varpi\), and \(\kappa\) be its local integer ring, uniformizer, and residue field. Let \(u\in O^{1\times a}\) be a lift of the actual nonzero covector defining the initial strict T level. Set

\[
 \tau_u=
 \begin{pmatrix}I_a&0\\ \varpi^{-1}u&1\end{pmatrix}
           \in\operatorname{GL}_{a+1}(F_v)
                     \simeq G^+(F_v).             \tag{BT4}
\]

The relative split identification is the one compatible with H embedded as diag(h,1). It sends every displayed matrix to an actual unitary point by its paired component \(J^{-1}g^{-t}J\), where \(J=\operatorname{diag}(I_a,-1)\); no rational global unitary matrix is being postulated.

Direct multiplication gives, for h in GL\_a(F\_v),

\[
 \tau_u^{-1}\operatorname{diag}(h,1)\tau_u
   =\begin{pmatrix}
      h&0\\ \varpi^{-1}u(1-h)&1
     \end{pmatrix}.                              \tag{BT5}
\]

This matrix and its inverse are integral if and only if \(h\in\operatorname{GL}_a(O)\) and \(\overline u\,\overline h=\overline u\). Hence

\[
 H(F_v)\cap\tau_u K_{a+1,v}\tau_u^{-1}
   =J_{0,v}
   =\{h\in\operatorname{GL}_a(O):
                         \overline u\,\overline h=\overline u\}.
                                                        \tag{BT6}
\]

This is the strict covector stabilizer, not the projective stabilizer and not an additional principal congruence group. There is no dilation average in (BT6). Both relevant characters are trivial on this group.

For later use, the identical calculation for an extra scalar t is

\[
 \tau_u^{-1}\operatorname{diag}(A,t)\tau_u
       =\begin{pmatrix}
            A&0\\
            \varpi^{-1}(t u-uA)&t
         \end{pmatrix}.                          \tag{BT7}
\]

Membership in the maximal compact implies both \(A\in\operatorname{GL}_a(O)\) and \(t\in O^\times\); this follows from the upper-left and lower-right blocks of the matrix and its inverse. Once t=1 it is exactly the original condition (BT6).

Take the finite adele \(\tau\) with components (BT4) at all T places, identity elsewhere, and identity at infinity. It is a genuine finite unitary adele of determinant one. From (BT2--6),

\[
 H(\mathbf A_f)\cap\tau J_B\tau^{-1}=J_0,\qquad
 \chi_B(\tau^{-1}h\tau)=\chi_0(h)^{-1}
                                         \quad(h\in J_0).
                                                        \tag{BT8}
\]

In particular, if F has bottom right level (J\_B,chi\_B), then h maps to F(diag(h,1)tau) has precisely the required inverse initial finite character for the bilinear period.

\subsection{3. The translated packet and all rational stabilizers}\label{the-translated-packet-and-all-rational-stabilizers}

Map

\[
 Z=H(F)\backslash H(\mathbf A_f)/J_0
   \longrightarrow
 G^+(F)\backslash(D\times G^+(\mathbf A_f))/J_B,
 \qquad [x]\longmapsto[(o,x\tau)].                 \tag{BT9}
\]

It is well defined by (BT8). Suppose two points agree. The relating rational element fixes the negative line at every real place and is diag(h,t), with \(h\in H(F)\), \(t\in E^1(F)\). At a place away T the finite level condition makes t an integral unit; at T this follows from (BT7). Thus t is an integral norm-one E-unit, hence a root of unity. At a selected G place, (BT2) forces its last diagonal residue to be one. The residue characteristic is prime to the root order, so t=1. Now (BT8) says exactly that the original two H classes agree. Thus (BT9) is injective.

Applying the same argument to a single packet point identifies its entire orbifold stabilizer, including its finite character, with

\[
                       \Gamma_x=H(F)\cap xJ_0x^{-1}.
                                                        \tag{BT10}
\]

This retains the actual strict-level stabilizers. It introduces no extra factor from the complementary U(1).

For an arbitrary arithmetic component \(\Gamma_g=G^+(F)\cap gJ_Bg^{-1}\), its determinant is a global norm-one unit. The G determinant in (BT2) is one modulo the selected prime, so the same root argument gives

\[
                              \det\Gamma_g=1.     \tag{BT11}
\]

Therefore the determinant-flat correction in LIFT28(HJ16) is still trivial on every component. The proof of jet separation does not change when the packet points move from x to x tau.

\subsection{4. The actual translated bottom detector}\label{the-actual-translated-bottom-detector}

Keep the actual selected initial HC u, bottom positive list \((v,-b_0)\), and the selected degrees

\[
 N_\tau=\sum_i u_{\tau,i}-\sum_i v_{\tau,i}
                              +b_{0,\tau}-a/2.    \tag{BT12}
\]

Vary b1 as in LIFT28(HJ15). For all sufficiently large ell in its fixed progression, the cusp restriction map at the actual finite level (BT3), using the fixed selected K-type embedding, is onto:

\[
 \operatorname{Hom}_{G^+_\infty}
    \bigl(\sigma_\ell^\vee,
       \mathcal A_{\rm cusp}(G^+,J_B,\chi_B)\bigr)
 \longrightarrow
       \mathcal A_{\lambda^\vee}(J_0,\chi_0^{-1}),
       \qquad F\longmapsto[h\mapsto F(\operatorname{diag}(h,1)\tau)].
                                                        \tag{BT13}
\]

Proof: repeat LIFT28 Sections5--7 with the finite level (BT3) and packet (BT9). Equation(BT10) supplies exactly the same compatible jet target, and(BT11) supplies the same ample canonical twist. The PBW-to-isometric factors are unchanged: tau is finite and the differential operators are real. No new finite Schwartz input is used in this argument.

For every nonzero initial row P, including every nonzero Mram row, and every unit seed w, there is consequently an actual bottom cusp form F such that

\[
 F(\operatorname{diag}(h,1)\tau)=\overline{f_{P,w}(h)}
                                             \quad(h\in[H]),
\] \[
 \int_{[H]} f_{P,w}(h)F(\operatorname{diag}(h,1)\tau)\,dh
       ={v_0\over d_\lambda}
           \sum_x{\|P_x\|^2\over|\Gamma_x|}>0.     \tag{BT14}
\]

All Haar and stabilizer factors are those of the original strict initial level. Right translation by tau is unitary on the ambient automorphic L2 space, so it does not itself change the Petersson norm of F. This is not an assertion that tau preserves an integral coefficient lattice or has zero arithmetic height cost. The norm of the interpolating cusp form remains uncontrolled, as does the sufficient ell.

\subsection{5. What has and has not been matched}\label{what-has-and-has-not-been-matched}

The new conclusions are (i) the exact rS=0 source residue specialization (BT2--3), and (ii) the explicit maximal-T translation and packet comparison (BT4--11). They refine the bottom detector to the standard-frame dat0 level formula, while preserving the literal initial strict level.

The source RealPins.tSpin condition is stronger. It identifies the entire conjugated bottom level and character with the fixed rowJ/rowChi in a particular real/arithmetic model; see \nolinkurl{BottomHypothesisMatchingCentral.lean22–44} and PairForm.lean66--73. It also uses the prescribed tS, not a translate freely redefined after the kernel construction. The equality of this tau with that fixed tS, up to the allowed embedded/rational/central/level corrections, has not yet been proved. Nor has betaY0 been tied to the standard frame while preserving the actual two vertical raw kernels and their terminal lattice family.

Thus (BT13--14) does not assert ValLastX, BotDat.Pin, or that the interpolated cusp form belongs to the prescribed downward theta image. Those remain the endpoint comparison. The conclusion is nevertheless stronger than choosing an unrelated bottom finite level: its exact residue subgroup, maximal T condition, character, and initial-level intersection are now fixed explicitly.

\subsection{6. Verification}\label{verification-1}

The source specialization was read directly in DefGroup.lean 138--275, 601--609, 620--707, LinksDefs.lean430--431, and \nolinkurl{SubgroupCharExtensionOneApply.lean32–43;} the reordering was checked in LinksChar.lean75--76 and109--110. These files were recovered from the pinned source; retrieval manifests are source\_recovery\_20261122T022749Z.json and source\_recovery\_20261122T022927Z.json. Recovery is not a Lean check.

LIFT28's whole checked interpolation proof is a carried input. The displayed block identities, compact intersection in both directions, root-unit packet argument, character inverse and Haar scope are the new proof here. Author whole-proof review completed 2026-11-22T02:34:29+00:00: PASS after reading every section and BT1--14, with direct source specialization and both directions of the block intersection checked. No effective image-cost or terminal-nonzero conclusion is claimed.

\clearpage

\section{AM-R-023: literal Host rank and Scale obstructions for the repaired CM fields}\label{am-r-023-literal-host-rank-and-scale-obstructions-for-the-repaired-cm-fields}

Author: GPT-6 Astra Pro, ALP\_MACHINE \nolinkurl{/root}, 22 November 2026 UTC. Reasoning effort is not separately exposed.

This finding fixes the scope of the remaining global parameter adapter. The actual mixed theta families on J14 and J56 exist in the senses proved in K029--K040 and LIFT26--29. Neither underlying CM field can, however, be the field of a literal source \nolinkurl{HostRaw}, for any choice of its abstract Frame. Changing only the weights, Scale, or T prime cannot repair this. Even if one discards the already proved CM-type obstruction and enlarges the field, the unchanged source Frame/Host inequalities impose the exact rank and Scale lower bounds below. These are necessary bounds, not a construction of the enlarged host or a sign computation for DP44.

\subsection{1. Source data and exact scope}\label{source-data-and-exact-scope}

Write n for Frame.ord, A for its additive group of order n, e for its cyclotomic parameter, a0 for its distinguished generator, and

\[
 n_0=\operatorname{ord}(a_0),\qquad o=n/n_0.
                                                        \tag{R23.1}
\]

The literal Frame contains a surjective homomorphism

\[
 h:A\longrightarrow H_e,
 \qquad H_e=(\mathbb Z/2^{e+2}\mathbb Z)^\times/\{\pm1\},
                                                        \tag{R23.2}
\]

and requires that h be injective on the cyclic subgroup generated by a0. The target is cyclic of order 2\^{}e. This last standard cyclotomic group fact is the source interface in \nolinkurl{FrobeniusOrbits/Statement.lean99;} it is used here, not presented as a new cyclotomic theorem.

A \nolinkurl{HostRaw} further requires e\textgreater=2 and an isomorphism

\[
 \operatorname{Gal}(E/\mathbb Q)\simeq C_2^2\times A.
                                                        \tag{R23.3}
\]

The full \nolinkurl{Host} imposes o\^{}64\textless=2\^{}e. The source rank and Scale conditions include

\[
 f=2n\le k,\quad k+m+4\le t,\qquad
 o\le x^{\epsilon/100},\quad x>1,\quad0<\epsilon<1/52,
                                                        \tag{R23.4}
\] \[
 L=x^\epsilon,\qquad \mathfrak a=x\exp L.
                                                        \tag{R23.5}
\]

The different symbol mathfrak a is used for the source Scale.a, to avoid confusing it with a unitary-group rank. The conclusion below uses only these displayed clauses. It does not use, weaken, or audit the source's accepted theorem.

\subsection{2. No literal HostRaw on J14 or J56}\label{no-literal-hostraw-on-j14-or-j56}

The actual fields already constructed are

\[
 K_{14}=\mathbb Q(\zeta_{28})B,\qquad
 K_{56}=K_{14}(\sqrt5,\sqrt{13}),
                                                        \tag{R23.6}
\]

where B is the cyclic degree-seven subfield of Q(zeta29). Their checked Galois groups are

\[
 \operatorname{Gal}(K_{14}/\mathbb Q)=C_2^2\times C_{21},
 \qquad
 \operatorname{Gal}(K_{56}/\mathbb Q)=C_2^4\times C_{21}.
                                                        \tag{R23.7}
\]

Both have exponent 42. Therefore neither field supports a literal HostRaw, under any relabeling or choice of Frame.

\textbf{Proof.} Equations (R23.2--3) would give a surjection from its Galois group to H\_e, hence to C4 because e\textgreater=2. Every quotient of a group of exponent 42 has exponent dividing 42. Since 4 does not divide 42, C4 cannot be such a quotient. This is a contradiction. The argument does not depend on the previously chosen e=0 Frames; it excludes every possible Frame on either fixed field. Their earlier numerical \nolinkurl{Ranks} records did not claim to instantiate HostRaw and remain valid. ∎

The same proof excludes every compositum of either field with finitely many quadratic extensions of Q: its Galois group embeds in the product of the original group with a group of exponent two and still has exponent dividing 42. This assertion concerns quadratic extensions of Q, not arbitrary iterated quadratic towers, which can have cyclic quartic quotients.

\subsection{3. Necessary rank of any enlargement retaining the literal inequalities}\label{necessary-rank-of-any-enlargement-retaining-the-literal-inequalities}

Let q0 be the odd part of n.~Every Frame satisfies

\[
 2^e\mid n,\qquad n_0\mid2^e,
 \qquad q_0\mid o.
                                                        \tag{R23.8}
\]

If it satisfies the full Host inequality, set

\[
 e_{\min}(q_0)=\min\{j\ge0:q_0^{64}\le2^j\}.
\] Then

\[
 e\ge e_{\min}(q_0),\qquad
 n\ge q_0\,2^{e_{\min}(q_0)},\qquad
 k\ge2q_0\,2^{e_{\min}(q_0)}.
                                                        \tag{R23.9}
\]

\textbf{Proof.} Surjectivity in (R23.2) gives 2\^{}e\textbar n by the first isomorphism theorem. Injectivity on the cyclic subgroup gives n0\textbar2\^{}e by Lagrange's theorem. In particular n0 is a power of two, so the integer o=n/n0 is divisible by the odd part q0 of n.~Thus 2\^{}e\textgreater=o\^{}64\textgreater=q0\^{}64. The integer 2\^{}e q0 divides n because its two factors are coprime and both divide n. These observations prove (R23.9), including k\textgreater=2n. ∎

In particular, if E contains K14 or K56 and has the literal degree {[}E:Q{]}=4n, the tower formula gives 21\textbar n.~One may apply the same argument directly with the odd divisor 21, even if the enlargement introduces additional odd factors. Exact integer arithmetic gives

\[
 2^{281}<21^{64}<2^{282},\qquad
 e\ge282,\quad n\ge21\,2^{282},\quad
 f,k\ge42\,2^{282},\quad t\ge42\,2^{282}+m+4.
                                                        \tag{R23.10}
\]

The numerical lower bound for the Frame plus o\_pow is sharp: take A=C21 x C\_(2\^{}282), let h be projection to the second factor, and choose a0=(0,1). Then n0=2\^{}282, o=21 and the displayed Frame conditions and o\_pow all hold. This verifies sharpness only for these group inequalities; it does not construct a number field, CM type, polarization, full Host, or admissible theta family.

For comparison, suppose a principally polarized host is isogenous to an ordinary power J\^{}N and has a full CM field of degree 2dim(J\^{}N)=6N. If its degree also equals 4n, then N is even, n=3N/2 is divisible by3, and

\[
 2^{101}<3^{64}<2^{102},\qquad
 n\ge3\,2^{102},\quad k\ge6\,2^{102},\quad N\ge2^{103}.
                                                        \tag{R23.11}
\]

This is a conditional numerical bound obtained before imposing the CM type. AM-R-003 already excludes the actual quadratic-induced Host type on every nonzero J power. The two statements must not be confused: (R23.11) is not a construction that avoids R003.

\subsection{4. The Scale cost cannot be repaired at fixed a=2\^{}960}\label{the-scale-cost-cannot-be-repaired-at-fixed-a2960}

For any o\textgreater=1 satisfying (R23.4--5),

\[
 L\ge o^{100},\qquad
 \log\mathfrak a>L+52\log L
                    \ge o^{100}+5200\log o.
                                                        \tag{R23.12}
\]

\textbf{Proof.} The o condition is o\textless=L\^{}(1/100), and taking its hundredth power gives the first inequality. Since x\textgreater1 and epsilon\textgreater0, L\textgreater1 and log L=epsilon log x. The strict upper bound epsilon\textless1/52 gives log x\textgreater52log L. Use log mathfrak a=log x+L. The function y+52log y is increasing for y\textgreater0, proving the final inequality. ∎

Consequently every enlarged host still containing K14 or K56 requires

\[
 L\ge21^{100},\qquad
 \mathfrak a>\exp(21^{100})\,21^{5200}.
                                                        \tag{R23.13}
\]

For any o\textgreater=2, there is no Scale with mathfrak a=2\^{}960 satisfying the literal o\_tie clause, regardless of how x, epsilon or J are changed. Indeed (R23.12) would give log mathfrak a\textgreater2\^{}100, whereas log(2\^{}960)=960log2\textless960\textless2\^{}100. Thus R022's specified weight profile and DO040's corresponding primes solve their stated weight/numerical Sizes problems; reparametrizing that same fixed Scale.a cannot complete the literal global record. A successful replacement must change the scale and coefficient profile, or prove an appropriately modified global budget statement for the mixed host.

\subsection{5. Consequences and geometric language}\label{consequences-and-geometric-language}

For DP\_MULT Q3e, this converts the fixed-host realization issue into an explicit negative statement: neither K14 nor K56 can be inserted in the literal HostRaw interface. The unchanged global Host inequalities also make the rank and Scale costs of a field-enlargement attempt explicit. They do not evaluate the actual g\_j, e\_j, H or endpoint terms in R001, and do not prove that every adapted mixed construction has negative DP44 difference. The mixed Picard/section and actual two-edge kernel constructions remain available independently of this literal interface.

Over Spec(Q) the quotient obstruction is a statement about finite etale Galois realizations and their constant finite quotient groups. The CM-type distinction belongs to the simple weight-one isogeny constituent h1(J), exactly as in R003; it is not a statement that every higher tensor or exterior-power correspondence vanishes. The previously fixed schemes, Picard functors and biextension lines retain the locators FGA232 Theorem3.1, pp148--149; FGA236 Corollary1.7, p229, and2.3, p232; SGA7I VII2.7/2.9, pp169--174, andVIII7.4(b), pp310--311. Representation lambda-rings/K-theory (SGA6 VI3.3) do not alter these scalar inequalities. The cote4/cote54 Cartier--Dieudonne realization is unchanged; no new ell-adic purity, motivic comparison, or Kummer RH assertion follows.

\subsection{6. Exact locators and verification}\label{exact-locators-and-verification}

\begin{itemize}
\tightlist
\item
  Alpöge, L. (n.d.). \emph{Remark}, step8 and the associated Lean archive, pinned commit7cbf7a0c32fca1badf05931879a37c3f8a2b03b4: \nolinkurl{Formalism/SharpRestriction/RamificationFrobenius/TotientGaloisClasses/ZsmulGeneratorEqZeroOfProjection.lean}, Frame16--24, n0/o68--70; \nolinkurl{Formalism/FrobeniusOrbits/Statement.lean}, HGrp27--28 and cyclotomic group interface97--101; \nolinkurl{.../TotientGaloisClasses/ConjEMemHostFIff.lean}, HostRaw12--43, cmType66, HostDer degree91--92; \nolinkurl{Formalism/LocalRingCardinality/DataHost.lean}, o\_pow50; \nolinkurl{Params.lean}, Ranks6--22 and f/r28--32; \nolinkurl{DataParams₂.lean}, Scale9--38 and o\_tie131.
\item
  ALP\_MACHINE AM-K029/030 and AM-K032, especially K032(QH22) and its preceding Galois decomposition: actual fields and ordinary J powers. AM-R-003: the independent isogeny/CM-type obstruction, used as a settled result without rederiving its proof.
\item
  AM-R-022 and AM-DO-040: the exact scope of the fixed-a replacement weights and prime certificates; none claimed a full Host.
\end{itemize}

The accompanying certificate uses only integer powers, divisibility, group exponents, and the displayed rank arithmetic. It does not certify a Host, a number-field extension, or a transcendental equality. The Scale proof above is symbolic and independent of floating-point logarithm approximations. Whole-proof checks are recorded separately in CHECK\_EVENTS.tsv.

\clearpage

\section{AM-K041 --- actual coupled divisibility and full-theta survival for the narrow cloud}\label{am-k041-actual-coupled-divisibility-and-full-theta-survival-for-the-narrow-cloud}

Author: GPT-6 Astra Pro, ALP\_MACHINE /root/kummer\_extension, 22 November 2026 UTC. Inherited reasoning effort is not separately exposed.

This finding specializes the \textbf{actual} class-kernel estimate K039 to the repaired profile R022, K040's actual nonuniform stage characters, and DO040's actual new T-prime. It covers every binary interior word and every top/thick/thin assignment. It also checks the changed hypotheses needed to transport LIFT25--27's full-theta injectivity. No old uniform numerical Schwartz norm or coefficient dimension is carried over.

The results are an arithmetic coupled-image lower bound and, separately, analytic injectivity before terminal restriction. They do not assert nonvanishing of the restricted row, PtilFF, or MainX. They do not fill Fields.Kq, Host, D7, SizeData₂, a downward-image condition, or a gain/error and height record. R022's stated o\_tie limitation remains. The old R020 numerical branch bounds are not used for the new partitions.

\subsection{1. Actual family and integral modules}\label{actual-family-and-integral-modules}

\begin{longtable}[]{@{}lrrrrrr@{}}
\toprule\noalign{}
Host & h & f=a=k & d & b=a+d & c=b+d & ell \\
\midrule\noalign{}
\endhead
\bottomrule\noalign{}
\endlastfoot
J\^{}14, K031 & 1 & 42 & 1760 & 1802 & 3562 & 45,156,049,316 \\
J\^{}56, K032 & 4 & 168 & 1634 & 1802 & 3436 & 40,530,339,696 \\
\end{longtable}

Here J is the Jacobian of y²=x⁷−1, and the CM type has f labels. Keep the product-polarized integral models, actual standard CM eigenbases, source creation operators, and period of K030/K032/K033/K035. Choose any partition of the labels into top, thick, thin, with counts \(N_T,N_H,N_L\) summing to f.~Set

\[
 A=2^{960},\quad B=2^{955},\quad n=f/2,\qquad
 E_T=28\left\lceil\frac{B+e^{336}}{28}\right\rceil,\quad
 E_H=28\left\lceil\frac{e^{343}}{28}\right\rceil,\quad
 E_L=28\left\lceil\frac{e^{336}/n}{28}\right\rceil.
 \tag{NW1}
\]

These are R022's exact integers. The replay independently encloses the exponentials by rational Taylor bounds and checks each ceiling. At a label of role beta use \(e_{0,\tau}=E_\beta,\ e_{1,\tau}=E_\beta+2\ell\). K040 constructs actual unitary characters with these infinity types, the prescribed real-adele restriction, the unchanged tame G-unit characters, and unramified components at29,11,T. Its exact relation is

\[
 \mu_1\mu_0^{-1}
    =(\widetilde\nu_1\widetilde\nu_0^{-1})\eta_T^{-4\ell}.
 \tag{NW2}
\]

Use DO040's new T and piT, the unchanged G-primes, and any fixed admitted common tuple character xi and its Q set. Both members of edge j have splitting character \(\chi_j=\mu_j\xi\). Keep K038's exact raw G/T/Q tensor, unnormalized Y sum, literal strict initial and middle levels, actual class stabilizers, and terminal tilQ decoration. The level at29 stays maximal. The changed numerical T and residue covector are exactly DO040's; no auxiliary Heisenberg translation or principal29 refinement is introduced.

For a word \(\omega\in\{0,1\}^{k-2}\), write \(\mathcal L_0^\omega,\mathcal L_{\rm ev}^\omega\) for K037/K038's actual integral Weyl coefficient lattices at these weights, with their stabilizer relations and transported class frames. The target is the coupled evGSD row with rhoVr action, not an arbitrary ordinary ColV row. The ramified source is its actual Mram sublattice. Allow a finite extension of the coefficient DVR R at a place over11; always normalize v(11)=1. This extension includes the actual root-of-unity fields of the new T input characters, with no claim that their field cost is trivial.

\subsection{2. Exact all-word selected degrees}\label{exact-all-word-selected-degrees}

Put \(W=\sum_i\omega_i\). The literal source weights and selected partitions of R022 are, in zero-based indices,

\[
 w_{0,i}=u_i-\frac{k-1}{2}+i,\qquad
 \nu_{0,i}=w_{0,i}+E_\beta-\frac{k+d}{2},\qquad
 \Delta=2\ell-d>2,\qquad
 \nu_1=(\nu_0,0^d)+\Delta\mathbf1_b.
 \tag{NW3}
\]

Let \(\delta_{j,\beta}(W)=|\nu_{j,\tau}|\) for role beta. Then

\[
\begin{split}
 \delta_{0,T}(W)
  &=A-B+kE_T+\ell\{2k(k-1)-4-W\}
       -\frac{k-2}{2}-\frac{k(k+d)}2,\\
 \delta_{0,\beta}(W)
  &=kE_\beta+\ell\{2k^2+6k-6-W\}
       +\frac{4-k}{2}-\frac{k(k+d)}2
                      \quad(\beta=H,L),\\
 \delta_{1,\beta}(W)&=\delta_{0,\beta}(W)+b\Delta.
\end{split}
 \tag{NW4}
\]

\textbf{Proof.} The rho shifts in w0 sum to zero. The R022 interior sum is \[
 \sum_{i=1}^{k-2}u_i
   =2\ell k(k-1)-4\ell-\frac{k-2}{2}-\ell W.       \tag{NW5}
\] The endpoints sum to A−B at a top label, and \(\ell(8k-2)+1\) otherwise. Add these sums and then \(kE_\beta-k(k+d)/2\). This proves the first two equations; the last follows from (NW3). All terms are integers since k,d are even. The arrangement of the bits can change the partition and its dimension, but its degree depends only on W. QED.

For \(d_{11}(m)=1+\lfloor\log_{11}m\rfloor\), the exact digit lengths on \textbf{both} edges, for every \(0\le W\le k-2\), are:

\begin{longtable}[]{@{}lrrr@{}}
\toprule\noalign{}
Host & top & thick & thin \\
\midrule\noalign{}
\endhead
\bottomrule\noalign{}
\endlastfoot
J14 & 278 & 145 & 141 \\
J56 & 279 & 146 & 141 \\
\end{longtable}

Indeed (NW4) decreases by ell at each increment of W. The exact endpoint checks are \[
 11^{z_\beta-1}\le\delta_{j,\beta}(k-2)
        \le\delta_{j,\beta}(0)<11^{z_\beta}.        \tag{NW6}
\] No floating point logarithm enters. Put \[
 Z=N_Tz_T+N_Hz_H+N_Lz_L,\qquad
 D_j(W)=\sum_\beta N_\beta\delta_{j,\beta}(W)
                         =D_j(0)-f\ell W.         \tag{NW7}
\]

\subsection{3. The actual local-section premise transports}\label{the-actual-local-section-premise-transports}

\subsubsection{Proposition NW8}\label{proposition-nw8}

Every actual normalized framed class germ in the two changed-weight K038 kernels satisfies K039(IG7)'s integral algebraic local-section premise in the same normalized Chern/CM gauge. Thus K039(IG9) applies with (NW7).

\textbf{Proof.} Only the infinity types, selected polynomials, and numerical T/residue covector have changed. The CM varieties, theta line, primitive sigma section, and source/logarithm comparison are unchanged. K039(IG2) excludes11 from every admitted Q set by the literal DegOne condition and residue degree6 in Q(zeta28)\^{}+. DO040's T differs from11 and splits completely; G is unchanged; K040's stages are unramified at11.

Apply K039(IG3)'s Cayley weak approximation to the new finite-level class sets. Representatives have component one at11 and their \textbf{own transported} standard integral Weyl frames. A rational change from an earlier display is not presumed integral. These choices depend on the fixed ranks and levels, not the archimedean word. The piT in tilQ is an11-unit.

K039(IG4)'s principal-ideal proof applies to the new algebraic angular characters: their finite splitting values on ideles trivial over11 are11-units. That proof never assumed uniform infinity type or finite-order global stages. The raw G/T/Q values and finite Weil transforms remain integral at11, since their Fourier/measure denominators involve only residue characteristics different from11. K039(IG5) expresses them as integral linear combinations of prime-to11 Heisenberg translates of the specified vacuum. There is no division by T−1, by a residue character group, or by the cardinality of Y.

K035's all-polynomial identity then gives the actual reduced torsion translates of the sigma product, after removing the fixed period. LIFT08's primitive normalization and K034's normalized torsion difference connection place them in the complete module \(U_0\exp(-\beta_{\rm prod})R[[t]]\). The exponential is retained: reduced Chern Taylor integrality is not claimed. K033's CM coordinate change and inverse, and K035's algebraic cD factors, are still integral and invertible at11.

Finally K037's integral Weyl-to-Weyl subD theorem applies to every actual dominant partition (NW3). Every resulting monomial has the selected degree at each label. Therefore the full-module bound K027(EX2--EX4), for arbitrary integral algebraic prefactors in the retained gauge, applies to the \textbf{whole} selected module. No rational isotypic or Gram projector is introduced. These are exactly the hypotheses of K039(IG9). QED.

\subsection{4. Explicit actual coupled-row divisibility}\label{explicit-actual-coupled-row-divisibility}

Define \[
 S_f(r)=2m+v_{11}(m!),\qquad m=\lfloor2fr/5\rfloor,
 \quad E_0=abZ,\quad E_1=bcZ,
\] \[
 L_j(W)=\left\lfloor D_j(W)/1220\right\rfloor-E_j,\qquad
 G(W;N_T,N_H,N_L)=L_0(W)+L_1(W)-S_f(a)-S_f(b).       \tag{NW9}
\] All L\_j here are positive, as proved below, so K039's maximum with zero is unnecessary. The stabilizer bounds are:

\begin{longtable}[]{@{}lrr@{}}
\toprule\noalign{}
Host & S\_f(a) & S\_f(b) \\
\midrule\noalign{}
\endhead
\bottomrule\noalign{}
\endlastfoot
J14 & 1,479 & 63,572 \\
J56 & 23,705 & 254,294 \\
\end{longtable}

\subsubsection{Theorem NW10}\label{theorem-nw10}

For every word and admitted fixed tuple, the actual K038 arithmetic map at the new profile satisfies \[
 \boxed{\ 
 P_{\rm cpl}^\omega(\mathcal L_0^\omega)
   \subset11^{G(W;N_T,N_H,N_L)}\mathcal L_{\rm ev}^\omega.
 \ }                                                     \tag{NW10}
\] The same holds on Mram. Uniformly in all words at fixed roles, use \(G(k-2;N_T,N_H,N_L)\).

\textbf{Proof.} NW8 and K039(IG8--IG10) give the two entrywise bounds, with digit errors abZ and bcZ. K039(IG11)'s reduction-mod3 bound uses the actual underlying integral rank 2fr and applies to every stabilizer of the new strict classes without a level refinement. Divide by the two actual stabilizers in the class sum. The Y sum is unnormalized. The nonarchimedean triangle inequality gives (NW10), exactly K039(IG13). K038 covariance applies because its finite-unit characters are unchanged and its integral subD identities are valid at the new weights. Intersection with the actual evGSD equivariant subspace is saturated; Mram is a sublattice, with no added projector. Finally Z is independent of W and D\_j decreases with W, proving the uniform-word version. QED.

A bound independent even of the role assignment is \[
\begin{split}
 G_{\min,h}={}&
 \left\lfloor f\delta_{0,L}(k-2)/1220\right\rfloor
 +\left\lfloor f\delta_{1,L}(k-2)/1220\right\rfloor\\
 &-141f(ab+bc)-S_f(a)-S_f(b).
\end{split}                                              \tag{NW11}
\] It is the minimum of the \textbf{lower-bound formula}, not a claim about the actual image's minimum content. Exact integers are in the replay; useful rigorous sizes are \[
                  G_{\min,1}\ge2^{481},\qquad
                  G_{\min,4}\ge2^{483}.                 \tag{NW12}
\]

\textbf{Proof of minimization and positivity.} For beta=H,T the difference \(R_\beta=\delta_{0,\beta}(W)-\delta_{0,L}(W)\) is independent of W and identical on edge1. The exact certificate verifies \[
 \lfloor R_\beta/1220\rfloor
         >\max(ab,bc)(z_\beta-141)>0.                    \tag{NW13}
\] Replacing a thin label by beta increases either floor in (NW9) by at least floor(R\_beta/1220), and its digit error by only ab(z\_beta−141) or bc(z\_beta−141). Thus each L\_j increases. Both L\_j and G are smallest for all thin and W=k−2. Direct substitution proves (NW12) and positivity of each L\_j there. This minimizes over a possibly larger role set than any completed source Host, which is valid for a uniform lower bound; no all-thin Host is asserted. QED.

There is also an exact finite degree comparison. For \(D=D_0(W)+D_1(W)\), \[
 \boxed{\qquad
     G(W;N_T,N_H,N_L)/D>1/1220-2^{-455}.
 \qquad}                                                 \tag{NW14}
\] The two floors lose less than2. With zmax=278 or279, the whole error is at most f(ab+bc)zmax+2+S\_f(a)+S\_f(b), respectively 75,828,772,861 or304,406,169,777. These are less than2\^{}37 and2\^{}39. The minimum D is at least2\^{}492 or2\^{}494. The error/D is therefore strictly less than2\^{}-455. This is a ratio of the proved arithmetic lower bound to raw derivative degree, not exact image content or a normalized DP44 gain/cost ratio.

\subsection{5. Actual selected highest nonvanishing transports}\label{actual-selected-highest-nonvanishing-transports}

\subsubsection{Proposition NW15: both compact weight ties and the middle line}\label{proposition-nw15-both-compact-weight-ties-and-the-middle-line}

The new compact oscillator on each edge has the actual summand \(W_{j+1}\otimes W_j^\vee\), with the source dual convention of LIFT25. The two primitive highest kernels pair the same middle line with value one.

\textbf{Proof.} At each label the literal recurrence is \[
 w_1=(w_0-\tfrac d2\mathbf1_a,
                   (-E_\beta+\tfrac a2)^{[d]}),\qquad
 w_2=(w_1-\tfrac d2\mathbf1_b,
                   (-E_\beta-2\ell+\tfrac b2)^{[d]}).
 \tag{NW16}
\] This is the accepted source lamTrue\_old, lamTrue\_new and nuW\_eq\_true, \nolinkurl{Columns/BaseLemmas.lean496–552,} at the new e entries. Thus, for j=0,1, \[
 \nu_j=w_j+e_j-\frac{N_{j+1}}2\mathbf1,\qquad
 w_{j+1}=(\nu_j,0^d)-e_j+\frac{N_j}2\mathbf1.             \tag{NW17}
\] These identities match \textbf{both} group actions coordinate by coordinate, for either value of every interior bit. The last nu0 entry is \(E_T-B-b/2>0\) at top, and \(E_\beta+2\ell-d/2>0\) otherwise. The last nu1 entry is Delta\textgreater2. Writing eta\_j=w\_j−min(w\_j), one obtains \(\eta_1=(\nu_0,0^d)\), \(\eta_2=(\nu_1,0^d)\), and \(\nu_1=\eta_1+\Delta\mathbf1_b\).

The principal-minor polynomials consequently satisfy \(p_{\nu_1}=\det^\Delta p_{\eta_1}\), with leading diagonal coefficient one. The determinant-insertion isometry in the fixed row metrics, LIFT25(AR8--AR9), gives \(h_1^\vee(h_1)=1\). Its proof is labelwise; neither the partitions nor e0 need be uniform across labels. Compact Schur averaging therefore gives the same nonzero consecutive highest contraction, with its unchanged dimension factors. QED.

\subsubsection{Theorem NW18: full analytic coupled injectivity}\label{theorem-nw18-full-analytic-coupled-injectivity}

For each word, let \(\mathcal H_{\rm ram,h_0}^\omega\) be LIFT27's actual initial ramified automorphic space, at the new T, K040 stages, and present compact weight. The raw coupled operator with these two primitive highest polynomials and source-normalized creation is injective on that space. The complete-coevaluation operator from its coefficient rows to full G2 functions is also injective.

\textbf{Proof by checking all changed hypotheses of LIFT27.} The local GJ holomorphy and rank margins use a,d and unitarity: \((d-a+2)/2=860\) or734 is unchanged. At G, K040 preserves the distinct nontrivial tame unit characters. The exact generic and sharp-middle quotient calculations use these characters and ranks, so LIFT27(MR15--MR22) applies without a finite Steinberg assumption.

At the new T the stages are unramified, the source has the same strict covector condition, and its norm-kernel is the sum of nontrivial residue dilation types. DO040 gives degree-one splitting and a residue characteristic exceeding the type bounds. The first tail test uses only b\textgreater=4a and q\textgreater=2; the second whole-Schwartz quotient uses unitary unramified stages and the nontrivial dilation type. Hence LIFT27(MR23--MR24) applies to every permitted T component, without recycling an old T mass.

At Q the two stage unit characters agree. Their ratio is unramified and unitary, so DO037's Gauss/graph quotient remains nonzero for the new values. At other finite places the relative field, selfdual lattices, and basic vectors are unchanged. Their tested spherical quotients allow arbitrary unitary unramified stage values. In particular29 remains maximal. NW15 proves the exact real selected occurrence and consecutive-highest alignment. The compact local parameters and angular twists give the same finite nonzero real gamma factors.

Globally put \(\upsilon=\chi_0\chi_1^{-1}=\mu_0\mu_1^{-1}\); the tuple cancels. HasInfQ(−e\_j) means angular exponent −2e\_j. Thus upsilon has angular exponent4ell at \textbf{every} label, and its square has exponent8ell. Both are nontrivial. LIFT26's exact boundary lemma at1, carried in LIFT27(MR25), therefore applies. Higher Hecke factors are at integers at least2, and the initial Euler product retains absolute-convergence margin860 or734. Neither e0 uniformity nor the old numerical equation e1−e0=4C is used once these hypotheses are checked.

The first global edge still has b\textgreater2a and the second b\textless c\textless2b with anisotropic target. Thus the accepted norm identities and every-embedding nonvanishing have exactly the same valid scope. LIFT27's multiplicity factorization, fixed moving-frame tensor identity, and local injectivity prove the assertion. The complete coevaluation comparison is its exact factor d0 d1. All arguments use actual characters and their stated local properties, not an old uniform numerical specialization. QED.

The initial space is nonzero for each word. DO040's explicit covector gives projective class stabilizer mu28 and strict stabilizer one on the identity support. Every character of mu28 extends to a \textbf{nontrivial} residue character at T, since each restriction fibre has (T−1)/28\textgreater1 elements. Choose the T characters with the last cancelling the actual compact/G central character, as in R017/DO040. The class-supported highest coefficient with the finite G Steinberg generic vector is a nonzero Mram input. This provides one input; NW18 itself covers the full Mram space. Its norm uses the new T and actual word-dependent compact dimension.

\subsection{6. Normalizations and the unresolved terminal step}\label{normalizations-and-the-unresolved-terminal-step}

Write d\_j=dim W\_j, D\_(j,tau)=\textbar nu\_(j,tau)\textbar, and H\_nu for the hook product. With LIFT25(AR22)'s actual positive gamma\_tau, the Chern highest-vector squared Schwartz norm is \[
 A_j^2=\det(2Y)^{-7hN_jN_{j+1}}
       \prod_\tau H_{\nu_{j,\tau}}
                         (\pi\gamma_\tau)^{D_{j,\tau}}. \tag{NW19}
\] Source creation multiplies by \(\Lambda_j=\prod_\tau\lambda_\tau^{D_{j,\tau}}\); its squared norm is \(\det(2Y)^{-7hN_jN_{j+1}}\prod_\tau H_{\nu_{j,\tau}}\). The literal compact coefficient integrals have 1/(d0 d1). The auxiliary isometric quotient has the corresponding square-root dimension factors. Hooks and dimensions use the actual partitions; LIFT25(AR26)'s old uniform numerical tables are inapplicable.

The arithmetic-to-analytic comparison, only after literal tilQ, is \[
 v_0v_1 C_\theta^{14h(ab+bc)}
       \prod_\tau(\lambda_\tau/cD_\tau)^{
                         D_{0,\tau}+D_{1,\tau}}.        \tag{NW20}
\] Here v0,v1 are the actual strict level volumes at the new T. Projective-T/full-G descriptions retain K038's weighted averages and matching volume refinements. The cD factors are11-units. The complex periods, source-creation factors, and real metrics are not assigned finite valuations in NW10. The finite input character field remains an arithmetic cost.

NW18 concerns full G2 functions. Restriction by tilQ, terminal brC, and PtilFF can still create a kernel. NW10 is valid even if the restricted map is zero. An upper content witness on the preceding image, a terminal nonzero scalar, and full e/H and gain/error accounting remain separate. This is actual new-profile image divisibility, not a source gainD value or GRH conclusion.

\subsection{Sources and check scope}\label{sources-and-check-scope-1}

Carried checked inputs are K027(EX2--EX4), K033(CS1--CS5), K034(TG3--TG10), K035(PF22--PF29), K037(CR8--CR14,CR35--CR40), corrected K038(CA27--CA36), K039(IG2--IG13), K040(NC17--NC28), R022(R22.1--R22.14), DO040, LIFT25(AR2--AR30), and LIFT27(MR8--MR32). R017 supplies the class-supported construction with DO040's new prime/covector and central-character extension. No accepted global theta theorem is rederived.

Direct pinned source anchors: \nolinkurl{Formalism/LocalToGlobal/Columns/BaseLemmas.lean496–552} (lamTrue\_old, lamTrue\_new, nuW\_eq\_true); \nolinkurl{Formalism/LocalRingCardinality/DataCoeff.lean21–23,39–41,51–55,83–84} (determinant shift, HC conversion, Weyl extremal, row norm); \nolinkurl{Formalism/LocalToGlobal/UnitaryNorms/IsUnitaryNormOf.lean90–94} (HasInf, HasInfQ), all in recovered commit 7cbf7a0c32fca1badf05931879a37c3f8a2b03b4.

The new proof is the all-word/role degree specialization, the actual geometric and character premise check, and the coupled-image bound. The replay independently encloses the defining exponentials, checks each Hamming weight by an explicit word, checks all local bit choices needed for arbitrary words, verifies both oscillator weight ties and digit endpoints, evaluates the bounds, and enumerates every role-count triple as a redundant arithmetic check. It does not enumerate 2\^{}(k−2) words, calculate a theta coefficient or image Smith form, prove a Rallis theorem by arithmetic tests, or build Lean.

\clearpage

\section{AM-R-024: actual terminal divisibility for the repaired profile}\label{am-r-024-actual-terminal-divisibility-for-the-repaired-profile}

Author: GPT-6 Astra Pro, ALP\_MACHINE \nolinkurl{/root}, 22 November 2026 UTC. Reasoning effort is not separately exposed.

This finding transports the terminal arithmetic lower bound to the R022/K040/DO040 family whose actual coupled kernel is K041. It uses the already proved whole-Weyl inclusions R020, rather than identifying the full branch with its highest line. There are two conclusions. For the literal Weyl detector sublattice MtildeLdel, every role assignment has a large terminal divisibility bound without a hook denominator. For the full literal MtildeL=intNab detector lattice, every assignment with at least one top label also has a positive explicit bound, now obtained from a long interval in the actual integer spectral factor. Neither conclusion asserts that a terminal value is nonzero or that the normalization is cheap in the DP44 budget.

\subsection{1. Actual terminal weights, with the word retained}\label{actual-terminal-weights-with-the-word-retained}

Keep the notation of K041: k=f is42 or168, d is1760 or1634, c=k+2d, n=c−1, ell is45,156,049,316 or40,530,339,696, A=2\^{}960, B=2\^{}955, and the role integer is E=E\_T,E\_H,E\_L. Write W for the Hamming weight of the binary word. At one CM label,

\[
 e_0=E,\quad e_1=E+2\ell,\qquad b_0=E-\ell,\quad b_1=E+\ell,
 \qquad v_i=4\ell(k-1-i)\quad(0\le i\le k-2).
                                                        \tag{R24.1}
\]

The literal column weight lambda is K041's w2, not its selected polynomial partition nu1. The literal row weight mu is RowPar.sgTrue. In zero-based initial indices their complete lists are

\[
\begin{split}
 \lambda={}&
 (u_i-(k-1)/2+i-d)_{i=0}^{k-1},\quad
 (-E+(k-d)/2)^{[d]},\quad
 (-E-2\ell+(k+d)/2)^{[d]},\\
 \mu={}&
 (v_i-(c-2)/2+i)_{i=0}^{k-2},\quad
 -b_0+k/2-d,\quad
 (-E+(k-d)/2)^{[d-1]},\quad
 -b_1+k/2,\quad
 (-E-2\ell+(k+d)/2)^{[d-1]}.
\end{split}                                             \tag{R24.2}
\]

Commas between displayed groups mean concatenation. These formulas follow by substituting (R24.1) in the source lamTrue\_old/new and RowPar.sgHC/sgTrue\_eq formulas already used in R019 and K041. All entries are integers. They interlace and have the common last entry s=−E−2ell+(k+d)/2. Set eta=mu−s1\_n and D\_mu(E)=sum eta\_i. Direct summation gives the word-independent integer

\[
 D_\mu(E)=(k-1)
       \left(E+2(k+1)\ell-\frac{k+3d}{2}\right)
                  +(d+1)(2\ell-d).
                                                        \tag{R24.3}
\]

For clarity, put a\_i=mu\_i−lambda\_(i+1), b\_i\^{}Br=lambda\_i−mu\_i, using zero-based indices 0\textless=i\textless n.~The nonzero a entries are

\[
 a_i=\ell\omega_i\ (0\le i\le k-3),\qquad
 a_{k-2}=
 \begin{cases}
 B+4\ell+(k-2)/2,&\text{top},\\
 2\ell-1,&\text{thick or thin},
 \end{cases}
\] \[
 a_{k-1}=a_{k+d-1}=\ell-d/2,
                                                        \tag{R24.4}
\] and all others vanish. The initial b\^{}Br entries are u\_i−v\_i−1/2 for i\textless k−1, the next is u\_(k−1)+b0−1/2, and b\^{}Br\_(k+d−1)=ell−d/2; all remaining entries are zero. R022 makes all these quantities nonnegative for every arrangement of the bits. This also checks interlacing without replacing arbitrary words by a Hamming representative.

Consequently R019's d0=sum a\_i and nBr=sum lambda−sum mu are

\[
\begin{array}{c|cc}
 &d_0&n_{\rm Br}\\ \hline
\text{top}&B+(W+6)\ell+(k-2)/2-d
          &A-B-(W+6)\ell+1-k/2\\
\text{thick/thin}&(W+4)\ell-d-1
          &(8k-8-W)\ell+2-k/2.
\end{array}                                             \tag{R24.5}
\]

These are exact changed-weight formulas. In particular the old uniform numerical R019/R020 tables do not apply.

\subsection{2. The two actual detector lattices}\label{the-two-actual-detector-lattices}

At each coefficient-field place above11 use v(11)=1, and choose the actual class representatives and transported integral frames of K039(IG3)/K041(NW8). The source explicitly distinguishes

\[
 \mathcal M_\nabla=\mathrm{MtildeL}
             =\mathrm{mtildeGSD.intNab},\qquad
 \mathcal M_\Delta=\mathrm{MtildeLdel}
             =\mathrm{mtildeGSD.intDel}.
                                                        \tag{R24.6}
\]

They have the same EqFun character, residue-action and stabilizer conditions. RowForms.lean281--286 gives these exact definitions. The coefficient inclusion Delta\_mu subset Nabla\_mu implies M\_Delta subset M\_Nabla; thus M\_Delta is an allowed sublattice of the literal detector lattice, not an invented unconstrained row space.

Let C=product\_tau cBr\_tau, computed with the actual weights (R24.2), and H=product\_tau H\_(eta\_tau). R019 proves every cBr\_tau is a nonzero integer. R020(R20.9) gives, with uT the literal piT power,

\[
 B_{\rm term}(\Delta_{\boldsymbol\lambda})
       \subset u_TC\nabla_{\boldsymbol\mu}
       \subset (u_TC/H)\Delta_{\boldsymbol\mu}.
                                                        \tag{R24.7}
\]

The new piT is an11-unit, so v(uT)=0. The normalized Fischer pairing is integral on Nabla\_mu x Delta\_mu. If G\_cpl is K041(NW9), the exact class formula DO038(DO38.9)/K038(CA30) therefore gives

\[
 v\bigl(\operatorname{PtilFF}(P_{\rm cpl}(P),Z)\bigr)
       \ge G_{\rm cpl}+v(C)-S_f(n)
                         \qquad(Z\in\mathcal M_\Delta),
                                                        \tag{R24.8}
\] \[
 v\bigl(\operatorname{PtilFF}(P_{\rm cpl}(P),Z)\bigr)
       \ge G_{\rm cpl}+v(C)-v(H)-S_f(n)
                         \qquad(Z\in\mathcal M_\nabla).
                                                        \tag{R24.9}
\]

Here P is in the actual initial integral lattice or Mram, and S\_f(n)=2floor(2fn/5)+v11(floor(2fn/5)!) is K039's proved stabilizer bound. It is125,628 on J14 and484,744 on J56.

\textbf{Proof.} K041 places every coupled coefficient in 11\^{}G\_cpl times the actual final Weyl lattice, including all initial and middle stabilizer denominators. Apply the first inclusion of (R24.7) and pair with the actual Weyl detector coefficient. Perfect Weyl/dual-Weyl integrality gives (R24.8), after dividing by each terminal class stabilizer and applying K039's bound. For the full dual-Weyl detector instead use the second inclusion of (R24.7). The same pairing and class sum prove (R24.9). K041's local frames and piT-unit transport preserve all these inclusions. No projection onto an invariant row has been applied with an unrecorded denominator; the rows already satisfy their exact EqFun constraints. ∎

Since v(C)\textgreater=0, the exact K041 lower bound minus S\_f(n) is still at least2\^{}481 on J14 and2\^{}483 on J56. Thus (R24.8) gives those uniform terminal divisibility exponents for every role assignment and every binary word. This result does not require a top label.

\subsection{3. A large spectral interval controls the full detector lattice}\label{a-large-spectral-interval-controls-the-full-detector-lattice}

Assume there is at least one top label. Put

\[
 a_* = B+4\ell+(k-2)/2,\qquad
 C_*=n\,v_{11}(a_*!),\qquad
 D_T=D_\mu(E_T),\quad D_H=D_\mu(E_H),\quad D_L=D_\mu(E_L).
                                                        \tag{R24.10}
\]

At every top label, the \textbf{actual integer} spectral factor satisfies

\[
                         v_{11}(c_{\rm Br})\ge C_*.
                                                        \tag{R24.11}
\]

\textbf{Proof.} Use the spectral product R019(R19.12) at index i=k−2. By (R24.4), its lower endpoint is L\_i=mu\_i−a\_* and its upper endpoint U\_i is at least mu\_i. Its product therefore contains, for each of the n column indices j, every factor

\[
     \prod_{v=\mu_i-a_*}^{\mu_i-1}
                         (v-(\mu_j+i-j)).
                                                        \tag{R24.12}
\]

The omitted zero in the full product is v=mu\_i,j=i; it is absent from (R24.12). For j=i the absolute product is a\_\emph{!. For j\textless i all terms are negative, and for j\textgreater i all terms are positive: these are the strict endpoint inequalities proved in R019 for interlacing weights. In either case the absolute values are a\_} consecutive positive integers. Their product is a\_\emph{! times an integer binomial coefficient. Each of the n products is thus divisible by a\_}!, while all remaining spectral factors are nonzero integers. This proves (R24.11), without enumerating the long interval. ∎

For every partition eta of D, the carried hook-length identity D!/H\_eta=dim Specht\_eta is a positive integer (R011, with Fulton and Harris's exact locator below). Therefore

\[
                    v_{11}(H_\eta)\le v_{11}(D!)\le D/10.
                                                        \tag{R24.13}
\]

For the exact two profiles, D\_T\textgreater=D\_H\textgreater=D\_L and 10C\_*\textgreater D\_T−D\_H. Define the explicit integer

\[
 B_* = C_*-\left\lfloor\frac{D_T+(f-1)D_H}{10}\right\rfloor.
                                                        \tag{R24.14}
\]

Then, for every assignment with N\_T\textgreater=1 top labels and every word,

\[
 \boxed{\quad v(C)-v(H)\ge B_*\ge2^{963}.\quad}
                                                        \tag{R24.15}
\]

\textbf{Proof.} All cBr factors are integral, and each top factor contributes at least C\_\emph{. The hook loss is at most floor((N\_T D\_T+N\_H D\_H+N\_L D\_L)/10), by (R24.13). Replacing each non-top label by a thick label only increases this upper bound. Increasing N\_T by one increases the difference between the proposed lower and upper bounds by at least C\_}−ceil((D\_T−D\_H)/10), which is nonnegative because 10C\_\emph{\textgreater D\_T−D\_H. Thus its minimum for N\_T\textgreater=1 is at N\_T=1, with all remaining labels thick. This gives (R24.14). The exact integers E\_T,E\_H from R022/K041, formula (R24.3), and Legendre's floor sum give B\_}\textgreater=2\^{}963 on both hosts. These last inequalities are exact integer comparisons in the saved certificate; there is no asymptotic approximation to a factorial or exponential. The one-top/all-thick comparison is a numerical worst-case bound, not an assertion that this role assignment supplies a full Host. ∎

Combining (R24.9), (R24.15), and the positive K041 lower bound minus the terminal stabilizer loss gives

\[
 \boxed{\quad
 v\bigl(\operatorname{PtilFF}(P_{\rm cpl}(P),Z)\bigr)
       \ge G_{\rm cpl}+B_*-S_f(n)\ge2^{963}
       \quad(Z\in\mathcal M_\nabla,\ N_T\ge1).
 \quad}                                                  \tag{R24.16}
\]

The coefficient inclusion B\_term(Delta\_lambda) subset 11\^{}B\_* Delta\_mu is itself valid on the whole standard tensor Weyl module in these local frames. It is an inclusion before restricting to the preceding image, and remains an inclusion after restriction. No upper content estimate on that restricted image follows.

\subsection{4. Nonzero values and costs remain distinct}\label{nonzero-values-and-costs-remain-distinct}

All displayed valuation bounds include zero, whose valuation is infinite. K041's injectivity into full G2 functions does not prove that the literal lowered class family is nonzero. If that lowered family is nonzero and the full compatible terminal row space of DO038 is admitted, its nondegenerate class pairing has a detecting row in M\_Delta after clearing denominators: M\_Delta is a full lattice in that compatible field-vector space. This is the usual field-spanning argument for a lattice; it is not a unit-valued or bounded-norm detector construction.

The large number in (R24.16) includes the mandatory integer cBr. R019(R19.15--16) retains the corresponding real norm

\[
 \log\|B_{\rm term}\|
 =-\tfrac12\log T\sum_\tau n_{{\rm Br},\tau}
    +\sum_\tau\log|c_{{\rm Br},\tau}|
    +\tfrac12\sum_\tau\log\kappa_\tau.
                                                        \tag{R24.17}
\]

Here nBr is the new formula (R24.5), and cBr/kappa are the exact factorial products R019(R19.4,R19.13) at (R24.2). Thus the integer gain is not inserted without its archimedean normalization, nor counted twice as a new local slope. The global mixed Host/budget replacement isolated by R023, terminal image nonvanishing, a detector of adequate norm, and the actual complete DP44 comparison remain open.

There is an exact way to separate this scalar, rather than only a warning about it. For the fixed word and role assignment, C and uT are independent of the terminal class. Put

\[
 s_B=u_TC,\qquad
 \widehat B=s_B^{-1}B_{\rm term}
                      =\bigotimes_\tau i_\tau^\dagger.
                                                        \tag{R24.18}
\]

The last equality is the already proved R020(R20.3). Its full integral inclusions are

\[
 \Delta_{\boldsymbol\mu}
       \subset\widehat B(\Delta_{\boldsymbol\lambda})
       \subset\nabla_{\boldsymbol\mu},\qquad
 -v(H)\le\operatorname{content}_v(\widehat B)\le0.
                                                        \tag{R24.19}
\]

Thus the normalized whole branch has no positive compulsory divisibility on its full Weyl source; its image contains the primitive target Weyl lattice. This does not bound its content on the preceding theta image above, because that image need not contain the vectors realizing the left inclusion.

Choose a number field M containing the actual algebraic scalars. For every complex embedding sigma of M, all conjugates of piT have absolute value sqrt(T), because piT belongs to Q(sqrt(−7)) and has rational norm T. Hence, with D\_Br=sum\_tau nBr\_tau,

\[
 \log|\sigma(s_B)|=\log|C|-\frac{D_{\rm Br}}2\log T,
 \qquad \log\|\widehat B\|=\frac12\sum_\tau\log\kappa_\tau.
                                                        \tag{R24.20}
\]

With ord\_w normalized by a local uniformizer, the product formula and R002's scalar-invariant normalization balance give the exact scalar cancellation

\[
 \frac1{[M:\mathbb Q]}
   \sum_w\operatorname{ord}_w(s_B)\log Nw
 =\frac1{[M:\mathbb Q]}
   \sum_\sigma\log|\sigma(s_B)|
 =\log|C|-\frac{D_{\rm Br}}2\log T.
                                                        \tag{R24.21}
\]

This is an application of R002 and R020, not a new product-formula theorem. It computes the signed contribution of this actual scalar as zero after both finite and real terms are included. It does not declare the entire terminal restriction cost zero, nor assign the remaining coupled-image, detector, field or error terms of DP44. In particular the large cBr contribution in (R24.16) cannot by itself be a net lift saving.

\subsection{5. Framework, sources and verification}\label{framework-sources-and-verification}

This is an inclusion of actual coefficient-image lattices over the DVR above11, paired with one of the source's two specified integral row lattices. It refines the integral representation/flag picture of SGA3I4.7.3/4.7.3.1 and SGA6VI3.3, rather than replacing lattice indices by equality of rational representation classes. The actual theta section retains the Picard and biextension locators FGA232 Theorem3.1, pp148--149; FGA236 Corollary1.7, p229, and2.3, p232; SGA7I VII2.7/2.9, pp169--174, andVIII7.4(b), pp310--311. Its formal jet realization retains EGAIV4 16.11.2--3(ii), pp54--55, and the Cartier--Dieudonne context at cote4/cote54 scans5--7 already fixed in the dictionary. There is no new motivic, ell-adic-weight, cycle nonvanishing, or Kummer RH conclusion.

Exact inputs and locators:

\begin{itemize}
\tightlist
\item
  AM-K041(NW1--20): actual all-word coupled kernels, frames and lower bounds, with R022, K040 and DO040's specified data.
\item
  AM-R-020(R20.2,R20.7,R20.9): whole Weyl/dual-Weyl inclusions and normalized pairing; AM-R-019(R19.1,R19.4,R19.12--16): exact spectral product, real norm and same-shift normalization.
\item
  AM-K039(IG3,IG11--13), AM-K038(CA27--30) and AM-DO-038(DO38.9): actual local frames, class stabilizers and literal PtilFF formula.
\item
  Alpöge, L. (n.d.). \emph{Remark}, step8, associated archive commit 7cbf7a0c32fca1badf05931879a37c3f8a2b03b4: \nolinkurl{Formalism/LocalToGlobal/Columns/RowForms.lean},281--286; \nolinkurl{AbelianHostFieldConditions.lean},408--421; \nolinkurl{BaseLemmas.lean},496--552,575--582; \nolinkurl{FiniteSumNonzeroLargeDepth.lean},41--55,118--151; \nolinkurl{Formalism/LocalRingCardinality/FactDelDualNabla.lean},75.
\item
  Fulton, W., \& Harris, J. (1991). \emph{Representation theory: A first course}. Springer. Theorem6.3(2)--(4), pp77--78, and Exercise6.11(b), p80, as already checked and carried in AM-R-011. The hook identity is an input; it is not rederived here.
\end{itemize}

The companion certificate checks the literal row formula independently, the complete terminal interlacing lists, exact d0/nBr and hook degrees, 624 Hamming representatives with the bitwise identities (R24.4), the long top spectral interval, every displayed finite integer bound, and the terminal stabilizer loss after the carried coupled-map bound. It does not compute a theta coefficient or decide its nonvanishing. Whole-proof readings are recorded separately in CHECK\_EVENTS.tsv.

\clearpage

\section{ALP-LIFT-30: the primitive norm-minus-T line behind the bottom translation}\label{alp-lift-30-the-primitive-norm-minus-t-line-behind-the-bottom-translation}

Author: /root/lift\_accounting, GPT-6 Astra Pro, ALP\_MACHINE. 22 November 2026 UTC.

This identifies LIFT29's finite translation with an explicit rational primitive negative line. It also identifies the necessary real-vector transport. The finite translation is not the source's tS merely because its intersection with the compact subgroup has the desired T level: the source tS conjugates the entire bottom level to the row level. The norm-minus-T Bessel embedding is a separate part of that source.

\subsection{1. Actual arithmetic inputs}\label{actual-arithmetic-inputs}

Let E/F be either checked mixed host K031 or K032, and let a be its chosen initial rank. Work on the actual rational Hermitian space

\[
 V=E^{a+1},\qquad J=\operatorname{diag}(I_a,-1),\qquad
 H=U(I_a)\subset G=U(V).
 \tag{NL1}
\]

Let T be one of the actual primes of DO040, or any odd rational prime with the same norm and congruence properties. Thus

\[
 \pi\bar\pi=T,\qquad (\pi,\bar\pi)=1,\qquad
 \pi\equiv1\pmod{(g_0g_1)^2},
 \tag{NL2}
\]

with pi integral and T completely split in E. These properties for the specified large primes are carried from DO040.1 and DO040.6--7 in definite\_oscillator\_large\_t\_primes.md; their primality is not reproved. At a selected T place, pi is a uniformizer and its conjugate is a unit. The other component of the relative split unitary group is always the paired inverse-transpose component; a displayed GL matrix below refers to the selected component.

Set

\[
 A=(T+1)/2,\qquad B=(T-1)/2.
 \tag{NL3}
\]

Both are rational integers, A-B=1, and A squared minus B squared is T. In the plane spanned by the first positive coordinate and the negative coordinate, define

\[
 g=
 \begin{pmatrix}
  A/\pi&B/\bar\pi\\
  B/\pi&A/\bar\pi
 \end{pmatrix}.
 \tag{NL4}
\]

On every other coordinate g is the identity.

\textbf{Proposition.} The element g lies in SU(V)(F). The integral vector

\[
 z_T=B e_1+A e_{a+1}
       =g(\bar\pi e_{a+1})
 \tag{NL5}
\]

is primitive, has norm minus T, and its stabilizer is

\[
 H_T=gHg^{-1}.
 \tag{NL6}
\]

Its Hermitian orthogonal complement is therefore rationally isometric to I\_a, with the explicit isometry obtained by restricting g.

\textbf{Proof.} The first column of the displayed two-plane matrix has norm (A squared minus B squared)/(pi times conjugate pi)=1. The second has norm minus one. Their Hermitian cross product is AB/(conjugate pi squared) minus BA/(conjugate pi squared), hence zero. The determinant is (A squared minus B squared)/(pi times conjugate pi)=1. This proves the special-unitary assertion and (NL5). The norm of z\_T is B squared minus A squared, or minus T. Because A-B=1, the coordinates generate the unit ideal of O\_E; equivalently the map O\_E to O\_E\^{}(a+1) with value z\_T has an integral left inverse. Its image is a saturated rank-one summand. Stabilizing the nonzero scalar multiple bar(pi)e\_last is the same as stabilizing e\_last. Conjugation by g proves (NL6). The restriction of g to the orthogonal complement of e\_last gives the last assertion. QED.

This supplies the actual primitive negative vector and positive rational complement. It does not by itself assert every extra field of the source's FlagS, FixedHyp, SBes, or transport records.

\subsection{2. Exact finite factorization}\label{exact-finite-factorization}

At a selected T place put c=bar(pi) and

\[
 \tau_c=
 \begin{pmatrix}I_a&0\\\pi^{-1}c e_1^*&1\end{pmatrix}.
 \tag{NL7}
\]

On the distinguished two-plane, direct multiplication gives

\[
 \begin{split}
 g^{-1}&=\tau_c k,\\
 k&=
 \begin{pmatrix}
 A/\bar\pi&-B/\bar\pi\\
 -\bar\pi&\bar\pi
 \end{pmatrix}\in\operatorname{SL}_2(\mathcal O_{F_v}).
 \end{split}
 \tag{NL8}
\]

Indeed the lower-left entry of tau\_c inverse times g inverse is minus (A+B)/pi=minus bar(pi), and its lower-right entry is bar(pi). The determinant of k is A-B=1. All entries are integral because bar(pi) is a unit. This proves membership, including integrality of the inverse, not merely a determinant valuation.

Away from T both pi and bar(pi) are units, so g and its inverse are integral. At each selected G place, (NL2) gives T congruent to one, A congruent to one, and B congruent to zero modulo g\_0g\_1. The prime 2 is distinct from both G primes. Consequently g is identity modulo the relevant G prime.

Let J\_B be LIFT29's standard-frame dat0-shaped bottom compact: at G its reduction is diag(U\_a,1), and elsewhere it is maximal. Let chi\_B be its generic G character and be trivial elsewhere. Let tau\_f be the product of the elements (NL7) at T and identity elsewhere. Then (NL8) and the preceding integrality assertions prove

\[
       g_f^{-1}=\tau_f k_f,\qquad
       k_f\in J_B,\qquad \chi_B(k_f)=1.
 \tag{NL9}
\]

There is no finite scalar in this identity for automorphic functions of right type chi\_B inverse.

\subsection{3. An arbitrary actual primitive T covector}\label{an-arbitrary-actual-primitive-t-covector}

The chosen covector of DO040 is generally not e\_1. Let u\_v be any primitive integral row at the selected T place, including its actual row (1,2,\ldots,2\^{}(a-1)). Choose

\[
 a_v\in\operatorname{GL}_a(\mathcal O_{F_v}),
       \qquad c e_1^*a_v=u_v.
 \tag{NL10}
\]

Such a matrix exists: a primitive row extends to an integral basis, and c is a unit. Complete its relative unitary pair in H(F\_v). Let a\_f be these matrices at T and identity elsewhere, and put hat(a\_f)=diag(a\_f,1). There is the exact identity

\[
 \tau_{u_v}
   =\widehat a_v^{-1}\tau_c\widehat a_v.
 \tag{NL11}
\]

The rightmost hat(a\_v) is in J\_B and has trivial chi\_B. Therefore, for any right-chi\_B-inverse automorphic function F and any initial automorphic row function P,

\[
 \begin{split}
 \int_{[H]}P(h)F(h\tau_u)\,dh
 &=\int_{[H]}P(h a_f)F(h\tau_f)\,dh.
 \end{split}
 \tag{NL12}
\]

Here a\_f is the H element, without its ambient block embedding. The equality is right invariance of Haar measure after the substitution h\_old=h\_new a\_f, followed by (NL11). The strict groups are related by J\_u=a\_f inverse J\_c a\_f.~Hence P\_c(h)=P(h a\_f) is exactly the corresponding initial row at the coordinate-covector level. This is an isometric change of finite class presentation. Rational stabilizers, class masses and the unweighted T diamond-norm kernels are transported by conjugation, so the whole Mram subspace is carried to the corresponding Mram subspace. No projective-level coarsening or character average is involved.

The chosen a\_f is an actual finite coordinate transport, not a claim that the unchanged source transport tr already has this value.

\subsection{4. The real component must also be transported}\label{the-real-component-must-also-be-transported}

A rational element is not a finite adelic element. Let sigma be the actual unitary bottom discrete-series representation and let

\[
 {\cal W}:\sigma^\infty\longrightarrow{\cal A}_{\rm cusp}(G)
\]

be an equivariant automorphic intertwiner, with right action R(x)F(y)=F(yx). For any smooth vector v put

\[
                  v_T=\sigma(g_\infty)v.
 \tag{NL13}
\]

Then, with h in H(A\_F),

\[
 \begin{split}
 {\cal W}(v_T)(g h g^{-1})
 &= {\cal W}(v)(g h g^{-1}(g_\infty,1))\\
 &= {\cal W}(v)(h(1,g_f^{-1}))\\
 &= {\cal W}(v)(h\tau_f).
 \end{split}
 \tag{NL14}
\]

The first equality is equivariance, the second is left G(F) invariance and the separation of finite and infinite components, and the last is (NL9). This proves the exact norm-minus-T Bessel interpretation of the translated detector.

If v belongs to the selected compact H type at the standard negative line, v\_T belongs to that same type under H\_T at the conjugate negative line. Unitarity gives norm(v\_T)=norm(v). Equivalently, the selected compact-type embedding is conjugated together with the negative line. It is incorrect to omit (NL13) while keeping the original fixed compact vector.

At every real place the two nontrivial singular values of g are sqrt(T) and 1/sqrt(T). Indeed its columns differ by unit phases from the real matrix with entries A/sqrt(T),B/sqrt(T), whose eigenvalues are (A+B)/sqrt(T)=sqrt(T) and (A-B)/sqrt(T)=1/sqrt(T). Thus the equality of Hilbert norms is a unitary-representation fact; it is not a bound of one for algebraic matrix entries, coordinate heights, finite polynomial coefficients, or selected jet frames. Those distinct costs have not been discarded.

At 11 the entire matrix g and its inverse are integral because T is different from 11. In particular the ambient algebraic lattice change is invertible at 11. This assertion does not identify an integral theta trivialization or an integral discrete-series model.

\subsection{5. Exact comparison with the source pins}\label{exact-comparison-with-the-source-pins}

The following source statements were read in this run at recovered commit 7cbf7a0c32fca1badf05931879a37c3f8a2b03b4.

\begin{itemize}
\tightlist
\item
  \nolinkurl{HaarTransversals/HidaSetting.lean117–130} defines K2 by congruence modulo g=g0g1 inside the full lattice compact. Hence its localization away G is maximal. Lines210--226 specify the norm-minus-T vector, its lattice and flag conditions. Lines267--276 put xu in that lattice compact and require OneOff away G.
\item
  \nolinkurl{Definitions/LatticeFrames/ThetaIdentityOfMapsAll.lean68–73} defines rowJ, rowChi and Matches. Matches conjugates the entire compact and its character, not only their intersection with H. Lines150--153 define the additional finite shift using xu.
\item
  \nolinkurl{IndefiniteReal/RealPins.lean47–52} requires a finite tS and the full conjugacy together with RowLevelMatchAt. Its betaStd field at line56 constrains betaV and betaX; it does not set betaY to one.
\item
  \nolinkurl{FlagStabilizerTensors/RowLevelMatch.lean22–47} divides the further pin into LevelMatchA, LevelMatchB and LevelMatchC. A is the compact character match. B also requires the actual full residue-group action on the automorphic multiplicity module, and C equates the selected idempotent coefficients. These latter clauses are not consequences of a compact intersection calculation.
\end{itemize}

For our standard-frame J\_B, tS therefore cannot be identified with the nonintegral T element tau merely from LIFT29: the row compact is maximal at T, whereas tau conjugates that maximal compact. Equations (NL5)--(NL14) place tau in the Bessel embedding and its finite/real coordinate transport instead.

They provide an explicit primitive negative line, its rational positive complement, and the exact transported detector. They do not identify the completed source's chosen betaY, tS, tr, SBes, RowLevelMatchB/C, or its idempotent-selected bottom space. The holomorphic detector still has to be placed in the actual prescribed downward theta image. No terminal period, DP44 gain, or full D7 construction is claimed here.

\clearpage
\part{Proofs awaiting completed whole-proof review}

The following contributions are retained as work in progress at this build. Their status is not upgraded by inclusion.

\clearpage

\section{AM-K042 --- the existing strict levels are neat, and the bottom twist has an exact finite window}\label{am-k042-the-existing-strict-levels-are-neat-and-the-bottom-twist-has-an-exact-finite-window}

Author: GPT-6 Astra Pro, ALP\_MACHINE /root/kummer\_extension, 22 November2026UTC. Inherited reasoning effort is not separately exposed.

This finding proves that the \textbf{existing} strict initial level and the dat0 bottom level have neat arithmetic groups. Consequently every initial definite class stabilizer is trivial, improving the actual K041 bound without a level refinement. It also computes the exact selected bottom jet degree and the entire canonical-twist interval that retains R022's windows and K040's character construction.

The comparison with the jet degree is not an interpolation-impossibility theorem. An effective coherent-jet or selected-period estimate remains required. No terminal nonzero value, full Host/SizeData₂, or DP44 cost sign is asserted.

\subsection{1. A direct local neatness criterion}\label{a-direct-local-neatness-criterion}

Call a characteristic-zero matrix neat if the subgroup of the algebraic multiplicative group generated by all its eigenvalues contains no root of unity other than one.

\subsubsection{Lemma NB1}\label{lemma-nb1}

Let q be a rational prime and \(M\in{\rm GL}_D(\mathbb Q_q)\). Suppose its characteristic polynomial lies in \(\mathbb Z_q[X]\) and reduces to \((X-1)^D\). If q\textgreater D+1, then M is neat.

\textbf{Proof.} If lambda is an eigenvalue, y=lambda−1 satisfies \[
 y^D+c_{D-1}y^{D-1}+\cdots+c_0=0,\qquad c_i\in q\mathbb Z_q.
 \tag{NB1}
\] Normalize v(q)=1. If v(y)\textless1/D, the leading term has strictly smaller valuation than every other term. For v(y)\textless=0 this follows from \(Dv(y)<1+iv(y)\); for \(0<v(y)<1/D\) it follows from \(Dv(y)<1\le1+iv(y)\). Both contradict (NB1). Therefore \[
                       v(\lambda-1)\ge1/D.              \tag{NB2}
\] This includes lambda=1. Each integer power of lambda, including negative powers, has the same lower bound on its difference from one. Multiplying finitely many such units preserves the lower bound. Every element of the eigenvalue-generated group is consequently in \(1+\{z:v(z)\ge1/D\}\).

For completeness, a nontrivial root of unity zeta with v(zeta−1)\textgreater0 has q-power order. Indeed a root of order m prime to q cannot be nontrivially near one: dividing (1+y)\^{}m−1 by y gives a unit when v(y)\textgreater0. The prime-to-q component of an arbitrary root must therefore be one. If zeta has exact order q\^{}s, then \[
 v(\zeta-1)=\frac1{q^{s-1}(q-1)}\le\frac1{q-1}.           \tag{NB3}
\] To prove the equality, the cyclotomic identity \(\prod_{(j,q)=1}(1-\zeta^j)=q\) has q\^{}(s−1)(q−1) factors. Their valuations agree because \((1-\zeta^j)/(1-\zeta)\) reduces to the q-unit j. Taking valuations proves (NB3). Roots with a nontrivial prime-to-q component instead have valuation zero. Since 1/(q−1)\textless1/D, no nontrivial root can satisfy (NB2). QED.

This proof controls \textbf{products} of eigenvalues, not merely the absence of finite-order matrices. It is stronger than the usual torsion-free conclusion needed for finite stabilizers.

\subsection{2. Application to the actual initial and bottom levels}\label{application-to-the-actual-initial-and-bottom-levels}

Let K be K031 or K032, with {[}K:Q{]}=2f and f=42 or168. For rank n, use the faithful rational representation on the underlying vector space of K\^{}n, of dimension \[
                              D=2fn.                    \tag{NB4}
\] The actual G data contain one chosen K prime from every conjugate pair above each rational g0,g1. Those rational primes split completely. Thus they contain every place of the real field above each g\_i. K031(HC12--HC13) gives42 chosen primes per rational prime; K032(QH27--28) gives168. This all-place condition matters when controlling products of eigenvalues from different embeddings.

At the literal initial level, the reduction at every selected G place is \(U_a(\mathbb F_g)\). The exact source is \nolinkurl{Columns/ColForms.lean93–115} together with \nolinkurl{Columns/GenericUnipotentCharacter.lean437–451:} HX(k,k,t,0)=HU(k), the full upper unipotent subgroup. At the dat0 bottom level of LIFT29(BT2--BT3), reduction is \[
                  {\rm diag}(U_a(\mathbb F_g),1).         \tag{NB5}
\] This is the exact m=2,rS=0 specialization of \nolinkurl{Divisibility/DefGroup.lean186–189,258–261,608–609.} It is not a projective T level or a full maximal G compact.

At a split place the second K component of a unitary matrix is the conjugate inverse transpose of the first. Unipotence is preserved by that operation. Therefore, in the rational representation (NB4), every block of every level element is unipotent modulo g. Rational conjugation by a finite class representative preserves the characteristic polynomial. Hence each arithmetic group \(G(\mathbb Q)\cap xJx^{-1}\) satisfies NB1 in dimension D.

The sufficient numerical bounds hold with a large margin:

\begin{longtable}[]{@{}lrrr@{}}
\toprule\noalign{}
Host & g0 & initial D=2fa & bottom D=2f(a+1) \\
\midrule\noalign{}
\endhead
\bottomrule\noalign{}
\endlastfoot
J14 & 4,615,409 & 3,528 & 3,612 \\
J56 & 22,285,932,761 & 56,448 & 56,784 \\
\end{longtable}

In particular g0\textgreater D+1 in both columns.

\subsubsection{Theorem NB6}\label{theorem-nb6}

Every arithmetic component group for the existing strict initial level and the existing dat0 bottom level is neat. Every initial definite class stabilizer is therefore trivial: \[
                         \Gamma_{0,x}=\{1\}.             \tag{NB6}
\] The statement remains true for every word, role assignment, new T-prime of DO040, and admitted tuple used in K041.

\textbf{Proof.} Apply NB1 to the characteristic polynomial just described. This works for every element and every class representative, in the faithful rational representation, so it is neatness of the actual restriction-of-scalars group. A finite subgroup of a neat group is trivial: a finite-order matrix has root-of-unity eigenvalues, hence all eigenvalues one; it is diagonalizable in characteristic zero, so it is the identity. Initial stabilizers are finite because the unitary group is definite. This proves (NB6). Neither T nor the archimedean weight changes the G-level premise. QED.

No principal29 level or new level at G was introduced. The statement does not say that R017's \textbf{coarser} projective-T/full-G support has trivial stabilizer: that support still has mu28, and its norm formula retains the factor28 and its G Schur denominator. Its refinement to the literal strict level still uses K038's weighted finite averages.

Nor does NB6 apply to the actual sharp middle level, whose reduction contains a GL\_(d−1) block. The terminal level has its own blocks. Their previous stabilizer estimates are not removed.

\subsubsection{Corollary NB7: exact improvement of K041}\label{corollary-nb7-exact-improvement-of-k041}

With the notation of K041(NW7--NW10), the actual arithmetic coupled map satisfies \[
 P_{\rm cpl}^\omega(\mathcal L_0^\omega)
       \subset11^{\,G(W;N_T,N_H,N_L)+S_f(a)}
                                  \mathcal L_{\rm ev}^\omega.
 \tag{NB7}
\] Thus one removes the loss1479 for J14 and23705 for J56.

\textbf{Proof.} In the literal K038 class sum the initial denominator is \(|\Gamma_{0,x}|=1\) in every summand. The middle denominator, the unnormalized Y sum, and every framed jet bound stay unchanged. Repeat K041(NW10)'s proof with no initial stabilizer subtraction. This also applies to its Mram sublattice. QED.

For LIFT29's translated compact packet the point stabilizers equal the initial ones, so they are all trivial as well. The bottom arithmetic quotient is already neat; a further cover is not needed to remove orbifold stabilizers. This does not supply an effective ample power, a jet interpolation bound, or a norm of a right inverse.

\subsection{3. Exact selected bottom jet order on the repaired profile}\label{exact-selected-bottom-jet-order-on-the-repaired-profile}

Retain R022's a=k, ell, A,B, E\_beta and binary word of Hamming weight W. LIFT28(HJ9) gives the exact selected degree \[
 N_\tau=\sum_i u_{\tau,i}-\sum_i v_{\tau,i}
                                      +b_{\tau,0}-a/2.
 \tag{NB8}
\] Since b0=E\_beta−ell and \(\sum_i v_i=2\ell k(k-1)\), substitution of the already proved sums in K041(NW4--NW5) gives \[
 \begin{split}
 N_T(W)&=A-B+E_T-\ell(W+5)-k+1,\\
 N_\beta(W)&=E_\beta+\ell(8k-7-W)+2-k
                                      \quad(\beta=H,L).
 \end{split}                                             \tag{NB9}
\] This is an exact formula for all words, not the degree of an unrelated ambient polynomial. The selected Pieri component has multiplicity one and this degree is unique, by the carried LIFT28 argument. A common HC determinant shift leaves NB8 unchanged.

The minimum occurs at W=k−2. Exact certified lower bounds are:

\begin{longtable}[]{@{}lrrr@{}}
\toprule\noalign{}
Host & top N\_min & thick N\_min & thin N\_min \\
\midrule\noalign{}
\endhead
\bottomrule\noalign{}
\endlastfoot
J14 & \textgreater2\^{}960 & \textgreater2\^{}494 & \textgreater2\^{}480 \\
J56 & \textgreater2\^{}960 & \textgreater2\^{}494 & \textgreater2\^{}478 \\
\end{longtable}

Full integer values and direct checks against NB8 are in the companion certificate. Large N by itself is not an obstruction: the minimal automorphic vector bundle also depends on the profile.

\subsection{4. The canonical-twist interval preserving these windows}\label{the-canonical-twist-interval-preserving-these-windows}

Use a new variable s to distinguish the canonical twist from the fixed gap ell. Hold u,v,b0,e0 fixed and put, at every label, \[
 q_s=(a+1)s,\qquad
 b_1(s)=E_\beta+\ell+q_s,\qquad
 e_1(s)=E_\beta+2\ell+q_s,\qquad s\in\mathbb Z_{\ge0}.
 \tag{NB10}
\] This is precisely LIFT28's ample canonical direction, with the retained flat determinant factor. It leaves the selected degree NB8 and the initial compact type unchanged. K040's specified character construction still applies exactly when e1(s) is a multiple28. Since gcd(a+1,28)=1 for a=42,168, this is 28\textbar s.

\subsubsection{Theorem NB11}\label{theorem-nb11}

Within the direction (NB10), retaining R022's fixed Scale, role centres, etaSz=4ell, and other listed window parameters, all the named R022 HC/centred/gap/window clauses hold exactly for \[
 s\in28\mathbb Z_{\ge0},\qquad
        s\le s_{\max}=28\left\lfloor
                           \frac{\ell}{28(a+1)}\right\rfloor.
 \tag{NB11}
\] Here ``exactly'' also includes the required multiple28 character condition; no assertion of a full SizeData₂ record is made.

\begin{longtable}[]{@{}lrrrr@{}}
\toprule\noalign{}
Host & s\_max & q\_smax & ell−q\_smax & number of candidate s \\
\midrule\noalign{}
\endhead
\bottomrule\noalign{}
\endlastfoot
J14 & 1,050,140,672 & 45,156,048,896 & 420 & 37,505,025 \\
J56 & 239,824,480 & 40,530,337,120 & 2,576 & 8,565,161 \\
\end{longtable}

\textbf{Proof.} The literal be\_close clause quantifies \textbf{all} pairs b\_i,e\_j: DataParams₂.lean164. The pair b0,e1(s) has difference \[
                          e_1(s)-b_0=3\ell+q_s.
 \tag{NB12}
\] Thus be\_close forces q\_s\textless=ell. Conversely, for 0\textless=q\_s\textless=ell, the four differences are ell,3ell+q\_s,ell+q\_s,ell, all at most4ell. This proves the exact restrictive inequality.

Write E\_beta=w\_beta+r\_beta with 0\textless=r\_beta\textless28. Then e1(s)\textless=w\_beta+3ell+28\textless w\_beta+4ell, so every role window still holds; the thick window allows an additional1 anyway. The earlier e0 is unchanged. The cloud clauses and top\_u/top\_b are unchanged, and the small increment of b1 preserves R022's proved b² bound.

All three HC interlacing gaps b0\textless e0\textless b1(s)\textless e1(s) remain positive, and the last positive/negative HC comparison still uses b0. The original centred comparisons either stay at ell, increase, or stay at least their previous lower bound; in particular the old e0 block is separated from e1(s) by 2ell+q\_s. At stage0 the opposite-sign gap between −b0 and −b1 increases. At stages1,2 the literal Dpar rule still gives no negative-sign entries. Thus the checked NCGaps condition is preserved.

The stage0 and1 compact spreads stay unchanged, and the stage2 spread increases by q\_s\textless=ell\textless2\^{}36. R022's elementary bounds still give every spread\textless2A, so DO040's same T\textgreater16A supplies the numerical Sizes inequalities. Positivity and the fixed eta/cloud equalities are unchanged. Finally gcd(a+1,28)=1 gives the displayed progression. Its next member violates NB12; the exact excesses are784 and2156. QED.

The two selected theta partitions become \[
 \nu_0(s)=\nu_0,\qquad
 \nu_1(s)=(\nu_0,0^d)+(\Delta+q_s)\mathbf1_b.             \tag{NB13}
\] Thus D0 is fixed, and D1 increases by fbq\_s. The digit lengths in K041 stay unchanged throughout the interval, as checked at its two endpoints. Every actual algebraic/gauge argument of K041 is unchanged, so NB7 holds with this new D1. The actual compact highest ties remain NB13, and the global stage ratio now has angular type4ell+2q\_s, which is nonzero. Hence K041's full-theta injectivity proof transports to the entire interval, with the new hook, dimension, and creation exponents retained.

If LIFT28's construction uses an additional descent power q\_can, one must intersect NB11 with its stipulated progression. Neatness alone is not used to set q\_can=1. The finite interval above is an upper allowance, not a proof of interpolation for any member.

\subsection{5. The remaining effective interpolation problem}\label{the-remaining-effective-interpolation-problem}

At the actual translated packet Z, NB6 removes every point stabilizer. Let \(\mathcal F_s\) be LIFT28/LIFT29's subcanonical automorphic sheaf, pushed to the minimal compactification, with its actual finite character. Let Z\_N impose the full jets through degree N\_tau at each ball factor at every packet point. The restriction kernel \(\mathcal K_s\) fits into \[
 0\longrightarrow\mathcal K_s\longrightarrow\mathcal F_s
  \longrightarrow\mathcal F_s|_{Z_{\mathbf N}}\longrightarrow0.
 \tag{NB14}
\] A sufficient explicit target is H¹(K\_s)=0 for one s in NB11's allowed progression. A sharper target is surjectivity only after projection to the selected Pieri quotient. Neither vanishing nor surjectivity has been proved in this interval.

For clarity, with packet cardinality z and minimal fibre rank r0, the unconstrained full-jet target has exact dimension \[
 z\,r_0\prod_\tau {N_\tau+a\choose a}.                    \tag{NB15}
\] The selected homogeneous target instead has dimension z dim W\_lambda. These give necessary dimension comparisons with the actual cusp section space, not estimates of that space. The neatness proof justifies the absence of an orbifold-invariant correction in these dimensions.

The generic Serre argument establishes a threshold for unbounded s; NB11 shows precisely the bounded interval in which that threshold must be controlled here. Neither N\_tau≫s\_max nor the size of the full-jet target alone proves failure: the sheaf's rank and weight also vary with the original profile. A proof of an upper section- dimension bound or a quantitative jet-separation theorem is still required.

I checked Lan's primary automorphic-cohomology vanishing theorem: Lan (2016), Theorem4.1(2), printed p.18, in the corrected author PDF \url{https://www.kwlan.org/articles/van-coh-aut.pdf}. It treats subcanonical automorphic bundles under weight conditions. The point-jet kernel in NB14 is an additional ideal modification and is not itself one of those homogeneous bundles. Thus that theorem alone does not give H¹(K\_s)=0. No such implication is used.

\subsection{6. Geometric interpretation and scope}\label{geometric-interpretation-and-scope}

The neatness result concerns the arithmetic quotient of the fixed unitary group; it changes neither the CM abelian scheme nor its Poincare biextension. NB7 is an improved coherent-module inclusion over Spec(R) for the same actual Schur coefficient map. NB8--NB15 concern principal-parts sheaves and coherent interpolation on the actual bottom Shimura variety, with its fixed finite character.

The inherited framework locators are EGA IV.4,16.11.2--3(ii) for jets/principal parts and EGA III.1,2.2.1(ii), printed p.100, for the qualitative coherent vanishing used by LIFT28; SGA6 VI3.3 for the lambda-ring description of Schur functors. Picard and biextension objects retain FGA232 Theorem3.1, pp148--149; FGA236 Corollary1.7,p229 and2.3,p232; SGA7I VII2.7/2.9, pp169--174 andVIII7.4(b),pp310--311. The carried formal-group inputs remain DP\_MULT's cote4/cote54,scans5--7 dictionary. No new ell-adic weight theorem, motivic purity statement, or Kummer RH assertion follows from these level and jet calculations.

\subsection{Check scope}\label{check-scope}

The companion script checks the exact source level anchors, field degrees and G-prime inequalities, small characteristic-polynomial and cyclotomic identities used in NB1, all displayed jet formulas, the complete interval endpoints, and preservation of all changed R022 comparisons. Endpoint and local-bit checks support the affine/monotone all-word and all-twist proofs; they do not enumerate all words or every integer in the twist interval.

The mathematical proof of neatness and its actual level application are checked separately from those finite diagnostics. No arithmetic class enumeration, cusp dimension calculation, successful jet interpolation, or Lean build is claimed.

\clearpage
\part{Source audit and geometric dictionary}

\clearpage

\section{Source audit: why CMData plus generic thetaG does not instantiate the mixed family}\label{source-audit-why-cmdata-plus-generic-thetag-does-not-instantiate-the-mixed-family}

Run ALP\_MACHINE, \nolinkurl{/root/lift_accounting}, 21 November 2026 UTC. This is an exact source audit and refinement of the current gate, not a new theorem replacing accepted DP52 or DP57. The generic-interface distinction is already stated immediately after the proof of DP52 in the intake. No Lean execution is claimed.

\subsection{The field repair is real but has a deliberately narrower target}\label{the-field-repair-is-real-but-has-a-deliberately-narrower-target}

Accepted DP57 proves that \(E'=\mathbb Q(\zeta_{28})\), \(F'=(E')^+\), and \(A=J^2\) furnish the indicated integral CM action, principal product polarization, unramified CM field condition, and an actual generic adelic endpoint. In particular it really meets \nolinkurl{CMData}, whose four fields are \nolinkurl{cmUnram}, \nolinkurl{totReal}, \nolinkurl{deg2}, and \nolinkurl{fixed} in \nolinkurl{Formalism/LocalRingCardinality/DataSpaces.lean:14–19}. The geometric type is induced from no imaginary quadratic subfield. DP57 explicitly does not construct the old \nolinkurl{Datum.dom}/\nolinkurl{CMObj.HasType} input; its separate old \nolinkurl{8 ∣ f} wrapper also fails since \(f=6\). Those are carried findings, not new obstructions derived here.

\nolinkurl{CMData} is sufficient for the generic \nolinkurl{Spaces} layer, subject to its other geometric inputs. It is not the type of the input of the source \nolinkurl{thetaG_exists} theorem. That theorem begins with \nolinkurl{variable (D : D7)}.

\subsection{The actual input to thetaG is substantially larger}\label{the-actual-input-to-thetag-is-substantially-larger}

\nolinkurl{Formalism/LocalRingCardinality/RhoDetPrimeEqMuTildeMulChiPow.lean} defines \nolinkurl{D7} from line 76. Besides \nolinkurl{Ranks} and \nolinkurl{Spaces}, it includes:

\begin{itemize}
\tightlist
\item
  The full old \nolinkurl{Host}, \nolinkurl{Rows.Data}, all member column chains, their periods, fixed coefficient fields and integral row/column lattices.
\item
  Exact ties between \nolinkurl{perCol} and the member \nolinkurl{PtilCol} maps, both at the base tuple and at twisted tuples; actual lattice embeddings and compatibility with archimedean norms.
\item
  The local theory, local comparison systems, theta-jet data, finite places, stage characters, their arithmetic and Galois compatibilities, and the remaining global data.
\end{itemize}

This is not an inference from the name \nolinkurl{D7}: these are visible fields in its declaration. For example \nolinkurl{perCol_base_eq_PtilCol} appears at lines 169--171, \nolinkurl{perCol_tie} at 181--183, and the theta-jet block starts at line 311 in the recovered copy.

The old \nolinkurl{Host} itself extends \nolinkurl{Formalism.HC.HostAx}, rather than \nolinkurl{CMData} (\nolinkurl{DataHost.lean:39}). Its predecessor \nolinkurl{HostRaw} includes the special Galois/cyclotomic field package and \(\sqrt{-d}\). Its CM type is literally the set of embeddings taking that element to \(i\sqrt d\): \nolinkurl{Formalism/SharpRestriction/RamificationFrobenius/TotientGaloisClasses/ConjEMemHostFIff.lean:66}. This is an induced imaginary-quadratic type. Supplying DP57's generic field and geometric CM action therefore does not, by itself, supply this old arithmetic \nolinkurl{D7} package.

\subsection{What the generic constructor actually proves}\label{what-the-generic-constructor-actually-proves}

\nolinkurl{Formalism/LocalRingCardinality/SpanPowerMembership/ThetaFamilyAdmissibleThree.lean} defines one auxiliary Hilbert space \(HG=HR\oplus HC\oplus\mathbb C\) at line 16. Its row embeddings land in the first factor, its column embeddings in the second, and \nolinkurl{v0} in the last (\nolinkurl{rowJ}, \nolinkurl{colJ}, \nolinkurl{v0}, lines 18--33). The \nolinkurl{thetaG} structure is at lines 71--153: its automorphic spaces all equal \nolinkurl{HG}, the cusp/type spaces are the full space, group actions are identity, Schwartz spaces are all endomorphisms, and periods are inner products. The specified lowering operator is \nolinkurl{rowJ ∘ io ∘ rowQ}.

For an already supplied \nolinkurl{D7}, the final theorem at lines 277--279 gives exactly

\begin{verbatim}
∃ X : ThetaX D,
  TieRowIsom X ∧ TieIo X ∧ AdmS D X ∧
  TieRowTyp X ∧ TypStable X.
\end{verbatim}

It proves all six components of \nolinkurl{AdmS} at lines 215--274. The Rallis representation used for this construction is the line spanned by the extra vector \nolinkurl{v0} (line 216). In particular the construction's existence does not prove that the actual abelian theta germs occur in its prescribed row/column images. This is an auxiliary realization of the interface, already recognized as such by accepted DP52.

There is a further explicit omission: the theorem does \textbf{not} supply \nolinkurl{TieBm}. In \nolinkurl{DataTheta.lean:747–752}, that predicate identifies the actual terminal bilinear period \nolinkurl{D.what₂.perCol} with the specific \nolinkurl{Btop} expression on the conjugated column and embedded terminal row, including its \nolinkurl{mY} factor. \nolinkurl{TieRowIsom} and \nolinkurl{TieIo} are distinct predicates at lines 739--745. Their proofs do not establish this terminal identification. No claim that \nolinkurl{TieBm} is identically zero or impossible for the generic construction is made: its noncanonical conjugation is not silently assumed to preserve the three factors.

\subsection{Exact conclusion for the proposed shortcut}\label{exact-conclusion-for-the-proposed-shortcut}

The proposed implication ``DP57 supplies CMData for J², and thetaG supplies the six identities, so only one finite coefficient matrix remains'' does not follow from these sources. The theorem providing thetaG requires the still-unconstructed arithmetic \nolinkurl{D7} first and its conclusion omits \nolinkurl{TieBm}; more importantly its theta maps have not been identified with the classical mixed-CM correspondence or the abelian theta kernels. The numerical source \nolinkurl{SizeData₂} also remains tied to the selected arithmetic derivatives and lattice errors, rather than to the bare CM orbit certificate.

There are useful concrete pieces: accepted DP52/DP57 provide the generic endpoint, AM-R-004 supplies the actual analytic CM torus and theta germ, ALP-LIFT-05 gives its full stable Hodge normalization, ALP-LIFT-06 computes the directional bare-certificate/Cauchy equality, and ALP-LIFT-07 fixes its exact characteristic and theta-null finite normalization. They reduce the unknown analytic normalization and identify a genuine source-line mismatch. They do not yet fill the mixed arithmetic constructor, the actual six identities on its tied theta family, or the final quantitative selected-value comparison.

The smallest honest next local target is to construct the prescribed trace-one section and allowed odd-level translates in explicit CM coordinates, compare their reduced gauges to the integral connection lattices at the selected primes, and identify the resulting finite coefficient matrix with the actual theta-lift row/column map. Even that local matrix identification is not a substitute for all global constructor and norm hypotheses. This source audit makes no universal nonexistence assertion about mixed-type theta correspondences.

\clearpage

\section{Geometric dictionary and exact Grothendieck locators}\label{geometric-dictionary-and-exact-grothendieck-locators}

Author: \nolinkurl{/root}, GPT-6 Astra Pro; ALP\_MACHINE, 21 November 2026 UTC. This is a dictionary of the proved constructions and their remaining comparison problems. It does not attribute the new numerical formulas or theta identities to Grothendieck. Statements carried from the supplied runs keep their recorded intake status. The source-reading events below distinguish fresh primary inspection from carried references.

\subsection{Current additions through the selected even chain and raw finite datum}\label{current-additions-through-the-selected-even-chain-and-raw-finite-datum}

The following table extends the coverage below. Its locators refer to the exact primary passages already identified and qualified in this dictionary; it does not attribute the new calculations to those texts.

\begin{longtable}[]{@{}
  >{\raggedright\arraybackslash}p{(\columnwidth - 4\tabcolsep) * \real{0.3333}}
  >{\raggedright\arraybackslash}p{(\columnwidth - 4\tabcolsep) * \real{0.3333}}
  >{\raggedright\arraybackslash}p{(\columnwidth - 4\tabcolsep) * \real{0.3333}}@{}}
\toprule\noalign{}
\begin{minipage}[b]{\linewidth}\raggedright
Findings
\end{minipage} & \begin{minipage}[b]{\linewidth}\raggedright
Geometric object and exact comparison
\end{minipage} & \begin{minipage}[b]{\linewidth}\raggedright
Grothendieck framework and locator
\end{minipage} \\
\midrule\noalign{}
\endhead
\bottomrule\noalign{}
\endlastfoot
ALP-LIFT-23 & Exact Fourier contraction of the entire coupled Q theta-group family, supported on the cofactor line, with an explicit Gauss line and integral content & SGA7 I VII2.7/2.9, pp.169--174, retains the finite theta group and phase; SGA3 I4.7.3 organizes its character eigenspace. The finite character line has its actual arithmetic scalar, whose11 content and Weil height are computed. The finite Fourier contraction is a map of coefficient modules, not the subsequent global automorphic pushforward or its terminal mass. \\
ALP-LIFT-22 & The exact composition of two finite theta-section correspondences through their prescribed intermediate character module & SGA7 I VII2.7/2.9 retains the theta-group line and phase; SGA3 I4.7.3 and SGA6 VI3.3 organize the actual induced flag coefficients. The composition is a saturated integral injection after the explicit Haar scalar, a stronger assertion than an equality in K0. The global arithmetic row lattice and period still need comparison. \\
AM-K034 & A rigidified prime-to11 torsion Picard-zero line with its unique integral finite-holonomy connection, and equality of complete translated principal-parts modules & FGA236 Corollaries1.7/2.3, pp.229/232, and SGA7 I VII2.7/2.9 give the Picard/biextension setting; cote54 scans003/005/007 distinguishes universal-extension data from a flat structure on the ample line. EGA IV4 16.11.2/16.11.3(ii) retains ordinary differential factorials. The actual auxiliary ideal equality is proved by the integral torsion frame and independent jets, not by the virtual line class alone. \\
AM-R-018 & Exact dominant-weight extremes and finite numerical bounds on the actual selected coefficient bundles & SGA3 I4.7.3 organizes weights, SGA6 VI3.3 the flag/Schur bundle. The source spread includes its rank; this is explicit numerical representation data, not a consequence of the motivic weight. The same crystalline selected-ideal computation of K033 is compatible with the displayed finite source Sizes intervals, without a global height-window conclusion. \\
AM-R-017 & Explicit split prime divisors of a norm element, a projective residue flag with exact central rational stabilizer, and descent of an actual automorphic coefficient section & SGA3 I4.7.3/4.7.3.1 gives character gradings and SGA6 VI3.3 the flag coefficient module. The integral polarized lattice is in the FGA232 Theorem3.1, pp.148--149, setting. Its exact class mass retains the stabilizer28; the projective cover degree is an11-unit but has large real size. A finite-etale realization of the characters records their monodromy, not theta nonvanishing or a height-window theorem. \\
ALP-LIFT-21 & The actual finite theta-group section mapped through the open matrix-rank correspondence to a parabolically induced coefficient module & SGA7 I VII2.7/2.9, pp.169--174, retains the theta line and phase; SGA3 I4.7.3 and SGA6 VI3.3 organize the character and flag-module action. The normalized map is an actual saturated integral injection with the computed content and metric. Minguez's analytic boundary theorem is a separate input; the geometric coefficient map does not calibrate the global terminal mass. \\
AM-DO-035 & Exact local character cancellation, a nontrivial dilation eigensheaf on the projective residue orbit, and its integral finite Fourier theta section & SGA3 I4.7.3 retains the character grading; SGA7 I VII2.7/2.9 retains the Heisenberg line and phase. FGA236 Corollary2.3, p.232, and EGA IV4 Theorem16.11.2/Remark16.11.3(ii), pp.54--55, organize the Picard section and its ordinary jets. The raw and filtered norms and class indices are actual data. Linear lattice motion and affine theta translation are distinct operations on the section. \\
AM-K033 & Frobenius-coupled CM coordinates on the actual product principal polarization, a primitive Schur differential, and the exact valuation of its permitted section image ideal & Cote4 scan003 sections8.1.1--8.1.2/scans005--006 sections8.4.1--8.4.3 and cote54 scans005--007 give the crystalline/duality framework. EGA IV4 16.11.2/16.11.3(ii), pp.54--55, retains ordinary derivative factorials; SGA6 VI3.3 supplies the selected coefficient module; SGA7 I VII2.7/2.9 retains the primitive biextension. The sharp slope1/1220 is proved for this actual image ideal with fixed-baseline constants; the Newton polygon or virtual motive alone would not establish it. \\
AM-K032 & Positive weighted trace modules over two real quadratic integer rings, their actual product Picard polarization and a full integral CM action on J56 & FGA232 Theorem3.1, pp.148--149, and FGA236 Corollary2.3, p.232, give the polarized/Picard setting; SGA7 I VII2.7/2.9 retains the actual symmetric theta line. The weighted trace identity constructs the ppav itself. Its8-divisible real degree does not change the primitive Hodge-type common-field obstruction; the5/13 coordinate denominators and arch metric remain explicit arithmetic data. \\
ALP-LIFT-20 & The literal initial coupled arithmetic section as a finite theta-group element, its primitive affine coefficient vector and an exact finite Steinberg contraction & SGA7 I VII2.7/2.9, pp.169--174, retains the theta-group phases and line; SGA6 VI3.3 supplies the flag coefficient module and SGA3 I4.7.3 its character grading. The exact projector scalar is d\_St²/GL-order and can have11-adic denominators even though the raw section is primitive. This is a nonzero first contraction for the actual G factor, not a statement that later arithmetic pushforwards are saturated. \\
AM-K031 & Actual norm-pulled rank-one Hecke characters, their exact tame conductor divisor and integral Hodge grading, and a preserved polarized alpha-isogeny family & SGA3 I4.7.3 supplies character gradings; SGA7 I VII2.7/2.9 and FGA236 Corollary2.3, p.232, organize the symmetric theta line and allowed isogeny pullback. The numerical ranks and finite residue fields are explicit arithmetic data. The common CM field obstruction remains a property of the actual Hodge representation, not a defect removed by a virtual class. \\
AM-DO-034 & A new primitive Schur differential on the actual J14 Picard line, integral principal parts at11, and an explicitly nonzero automorphic image with its metric & EGA IV4 Theorem16.11.2/Remark16.11.3(ii), pp.54--55, distinguish ordinary derivative factorials; FGA236 Corollary1.7, p.229, gives the Picard setting; SGA6 VI3.3 supplies the flag/Schur module; SGA7 I VII2.7/2.9 retains the transported theta line. The actual42-block selected section survives the first map. Its hook metric, pi\_B coordinate scalar and ramified29 index remain additional data, and do not determine later integral images. \\
AM-R-016 & Orthogonal factorization of the polarized Hom space, conjugate and dual theta lines, and actual compact period contractions & SGA7 I VII2.7/2.9, pp.169--174, organizes the biextension and theta-group tensor products; SGA3 I4.7.3 retains the splitting-character grading. The final two period transfers are exact integrals with the same Haar measures, and at m2 they complete the six analytic predicates. No integral arithmetic row identity or bounded free-data witness follows from this tensor identity alone. \\
ALP-LIFT-19 & Distinct tame character lines, a finite induced coefficient module and an explicit primitive column-type vector with nontrivial metric & SGA3 I4.7.3 gives the grading and SGA6 VI3.3 the flag/representation-module setting. The Mackey summand is an actual submodule over the displayed coefficient ring; its Hermitian denominator is the calculated finite flag index. Its virtual K-class and its type occurrence do not identify the global arithmetic image lattice or the raw theta operator. \\
AM-K030 & An explicit trace-self-dual fractional ideal in the real degree-seven field, integral symmetric endomorphisms of the product ppav J\^{}14, and the actual polarized CM lattice & FGA232 Theorem3.1, pp.148--149, and FGA236 Corollary2.3, p.232, supply the Picard/abelian-scheme setting; SGA7 I VII2.7/2.9, pp.169--174, records the transported theta line and semicharacter. The inverse different, xi\_K=xi\_0*pi\_B\^{}(-6), trace-one diagonal, and integral11 eigenbasis are calculated explicitly. The crystalline direct sum survives this prime-to11 integral transport, while a selected highest-weight upper witness and the source's stronger host record require separate proofs. \\
AM-DO-033 & A primitive integral finite-Steinberg coefficient vector, its actual finite theta-group realization, and an explicit automorphic matrix-coefficient section with exact metric & SGA3 I4.7.3 organizes the finite character and grading; SGA6 VI3.3 supplies the representation-bundle setting; SGA7 I VII2.7/2.9 retains the exact theta-group line and finite phase. The class mass, Steinberg dimension and non-unit Hermitian norm are actual numerical data. A nonzero coefficient section is shown to survive this first theta operator, but no equality of virtual classes identifies the later prescribed image lattice. \\
AM-K029 & Descent of the induced angular CM character through every integral unit, extension from principal ideals, and exact correction of its quadratic restriction through an exponent-two capitulation kernel & SGA3 I Proposition4.7.3/Corollary4.7.3.1 supplies the character grading; FGA232 Theorem3.1, pp.148--149, and FGA236 Corollary2.3, p.232, supply the CM/Picard setting. The resulting normalized algebraic Hecke character has rank-one realizations; the unit, class-group and coefficient-field arguments are explicit arithmetic input. No effective height or number-field RH follows from the corresponding motivic or l-adic class. \\
ALP-LIFT-18 & Actual tame character lines on the norm-one torus with an exact integral-unit descent condition and fixed conductor divisor; a nonempty selected-weight locus & SGA3 I4.7.3 identifies the torus character lattice and gradings; SGA6 VI3.3 records the induced coefficient-bundle class. Its finite character is ramified on precisely the displayed G divisor. Cote4 scan012 section8.6 supplies a monodromy question, not a theorem converting this explicit tame character into an analytic nonvanishing bound. \\
AM-DO-032 & The source-defined terminal Fourier family, its finite character quotient/radical, a saturated integral coefficient submodule and a non-unimodular restricted pairing & SGA7 I VII2.7/2.9, pp.169--174, organizes the finite theta-group action and line; SGA6 VI3.3 supplies the representation-module language. The orthogonal projector denominator is additional integral metric data. Saturation in the coefficient lattice does not make the restricted pairing perfect, and its finite Gram is not the global arithmetic terminal mass. \\
AM-R-015 & The exact signed real weights, opposing character shifts, compact U(t) Levi and first-occurrence block of the actual vertical correspondence & SGA3 I4.7.3 and SGA6 VI3.3 organize weights and flag/representation classes; SGA7 I VII2.7/2.9 retains the polarized kernel line and its orientation. Li's regular real theta theorem, Atobe's first-occurrence theorem and Howe duality are separate analytic inputs. No equality in K-theory or of crystalline slopes substitutes for those representation-level correspondences. \\
AM-K028 & A prescribed infinity character fails to descend through the integral units; the actual CM Hodge representation of an ordinary J-power imposes the divisibility7 condition on multiplicity & SGA3 I Proposition4.7.3 and Corollary4.7.3.1 organize character gradings, while FGA232 Theorem3.1, pp.148--149, and FGA236 Corollary2.3, p.232, give the CM/Picard setting. The global idele-unit descent obstruction is proved explicitly here; it is not removed by passing to a virtual motive or by an integral change of theta frame. \\
ALP-LIFT-17 & An absent finite representation type in a characteristic-zero coefficient bundle, proved using flag stabilizers; exact tame conductor divisor & SGA3 I4.7.3 describes the character decomposition and SGA6 VI Theorem3.3 the flag/representation setting. The projector is zero on the actual finite flag modules, hence on all compact-generated spherical local modules. Its conductor records ramification of the corresponding rank-one character sheaf; the dictionary does not identify this divisor with the complete arithmetic height cost. \\
ALP-LIFT-16 & Actual unramified local tests, a full theta-group translate supported at29, exact principal-level degrees and arithmetic class masses & SGA7 I VII2.7 and2.9, pp.169--174, describe the theta-group/biextension lines; SGA3 I4.7.3 describes the character actions. The class index is an actual finite cover degree. Neither its local zeta test nor its unknown depth bound is a consequence of a virtual K-class. \\
AM-DO-030 & Simultaneous ordinary jets at finitely many distinct torsion sections, translated Picard-zero lines, and a finite theta-group operator with its exact norm & FGA236 Corollary1.7, p.229, and Corollary2.3, p.232, give the Picard setting; EGA IV4 Theorem16.11.2 and Remark16.11.3(ii), pp.54--55, distinguish these ordinary derivatives from divided differential operators; SGA7 I VII2.7/2.9 organizes the transported line and its finite phase. A new29 torsion choice is an actual change of section. \\
AM-DO-031 & Complete terminal finite Fourier datum, explicit lattice windows, a primitive coefficient vector away11 and a rational point with exact raw value & SGA7 I VII2.7/2.9 and VIII Proposition7.4(b), pp.310--311, in its stated extension scope. The finite group-algebra element, its character phase, and its integral content are retained; geometric nonzero sections need not survive a specified theta projection. \\
AM-K026--027 & Explicit full-Hyp3 Schur weights, exact hook metrics, and the common-weight crystalline lower bound on the actual formal theta line & SGA6 V and VI Theorem3.3 give the Schur and lambda/K-theory language; SGA3 I4.7.3 its weight grading. Cote4 scan003, sections8.1.1--8.1.2, scans005--006, sections8.4.1--8.4.3, and cote54 scans005--007 give the crystalline/duality setting. The selected slope1/1220 and retained factorial cost are calculations in the actual line and integral differential module, not just its Newton polygon. \\
AM-R-014 & Genuine rational Witt-lower Hermitian spaces, parabolic constant terms and stability of full smooth real isotypic realizations & Reductive groups and their character gradings belong to SGA3 I4.7.3; coefficient-bundle classes to SGA6 VI3.3; theta lines to SGA7 I VII2.7/2.9. Wu's Proposition3.1 is the separate analytic tower theorem. None of these framework identifications supplies the remaining source identities or arithmetic host pins. \\
AM-K019 & Quasi-periodic theta logarithmic derivatives and the polynomial auxiliary-variable connection; a full ordinary differential lattice has unbounded creation powers & EGA IV4, Theorem16.11.2 and Remark16.11.3(ii), pp.54--55, distinguish divided differential operators from the ordinary Weyl action. Cote54 scans003,005,007 supply the universal extension and duality setting, not a nilpotence theorem for this entire connection. \\
AM-K020; ALP-LIFT-13 & Exact evaluated theta derivatives and the primitive direction at the height-three primes3 and5 & Cote4 scan003, section8.1.1--8.1.2; scans005--006, section8.4.1--8.4.3; cote54 scans005--007. The computed scalar slopes concern the evaluated section, not just the isocrystal Newton slopes. \\
AM-R-009--010 & A principal theta line on the actual product A=J squared, its E-prime action, norm-one torus characters, and permitted odd-isogeny translates & FGA232 Theorem3.1, pp.148--149; FGA236 Corollaries1.7 and2.3, pp.229/232; SGA7 I VII2.7 and2.9.3, pp.169--174. The explicit period and its Gaussian metric are additional analytic data. \\
ALP-LIFT-11--12 & Tensoring the automorphic coefficient bundle by an actual unitary determinant line, then a nonzero second theta integral & SGA3 I Proposition4.7.3 and Corollary4.7.3.1 describe the character grading; SGA6 VI Theorem3.3 describes the corresponding representation/flag class. Rallis nonvanishing and its measures are separate automorphic input, not consequences of K-theory. \\
AM-DO-025--026 & Selected integral highest weights, all-polynomial theta-group kernel, real covariance, adjoint, and weighted class matrix & SGA3 I Proposition4.7.3; SGA6 V lambda-operations and VI Theorem3.3; SGA7 I VII2.7/2.9 for the biextension line. The arithmetic terminal matrix is an actual map with a metric and cannot be replaced by its virtual class. \\
AM-DO-027 & An allowed scalar normalization of the literal arithmetic datum; equal-volume chart phases are integral roots of unity & Rank-one character bundles in SGA3 I4.7.3 and SGA6 VI3.3. This is an actual integral isomorphism of each existing input line. It does not construct a missing raw arithmetic section or terminal functional. \\
AM-R-011 & Integral Schur modules, determinant insertion, hook-normalized Fischer metrics and a rational central idempotent with an explicit denominator & SGA6 V and VI3.3, with the carried-source status stated below. Over Z, the denominator primes matter; characteristic-zero decomposition alone does not provide an integral splitting. \\
AM-K022; AM-R-012 & Ordinary principal parts at the zero section, integral jet surjectivity at explicit level, determinantal ideals of the allowed section map, and a concrete first automorphic class function & EGA IV4 section16, especially Theorem16.11.2 and Remark16.11.3(ii); FGA236 Corollary1.7; SGA7 I VII2.7/2.9. The finite matrix uses ordinary jets; evaluating a derivative introduces the displayed factorial. The class-function mass is additional adelic data. \\
AM-K021 & A finite Artin character component, CM-type height classes factoring through a central-conjugation quotient, and a ray-theta Mellin period carrying L-prime/L & The Artin projector is a finite-etale, weight-zero motivic summand and can also occur inside a CM Hom motive. SGA4--5 supply its sheaf/trace language; the finite-field weight results cited below do not give the required arithmetic height comparison or number-field RH. \\
AM-DO-028; ALP-LIFT-14 & The polarized Hom space as a tensor product with its dual skew-Hermitian factor; determinant character transport; actual selected nonzero theta iteration & The character and dual-representation operations lie in SGA3 I4.7.3 and SGA6 VI3.3. The tensor/Hom symplectic map retains an actual biextension normalization in SGA7 I VII2.7/2.9. The Rallis formula, its analytic L-values and Haar factors are additional automorphic theorems, not conclusions of these geometric frameworks. \\
ALP-LIFT-15 & Full finite theta-group eigenspaces, their orthonormal Weyl basis, and one nonzero translated section at each analytic edge & SGA7 I VII2.7, pp.169--170, and VII2.9.1/Proposition2.9.3, pp.172--174, organize the biextension and theta-group lines; FGA236 Corollary2.3, p.232, supplies the Picard-functor setting. The varying finite level, norm and orbit index must be retained in any arithmetic comparison. \\
AM-R-013 & Algebraic tensor-factorized theta images, branching of the compact coefficient bundle to an actual codimension-one Hermitian subspace, and its positive restriction pairing & SGA3 I4.7.3 and SGA6 VI3.3 give the representation-bundle language. FGA232 Theorem3.1, pp.148--149, and SGA7 I VIII Theorem7.1, pp.299--300, describe the polarized/Picard and extension settings with their stated hypotheses. They do not identify the analytic restriction norm with the source's arithmetic terminal functional. \\
AM-DO-029 & A finite theta-group Fourier expansion of the literal nonterminal raw arithmetic Xi, its support/period lattices, exact phases and norm & SGA7 I VII2.7/2.9, with VIII Proposition7.4(b), pp.310--311, in its stated fully faithful extension scope. The explicit integral group-algebra element and its denominators are stronger data than the class of its line. Its selected jet evaluation can still vanish. \\
AM-K023; AM-K025 & The full integral formal section module, restricted differential operators, the finite mod11 evaluation quotient and adjacent-gap factorial content & Cote54 scans003,005,007 and cote4 scan003 section8.1.1--8.1.2, scans005--006 section8.4.1--8.4.3 supply the universal-extension/crystalline framework. EGA IV4 Theorem16.11.2 and Remark16.11.3(ii), pp.54--55, distinguish the ordinary derivative factorials from divided operators. The finite quotient and its unit Smith rank are computed here; no higher divided-lattice stability is inferred. \\
AM-K024 & A primitive integral polynomial already in its selected source/target Schur summand, with exact ordinary-CM gauge and permitted jet realization & SGA6 V and VI3.3 describe its Schur class; SGA3 I4.7.3 its highest-weight grading. EGA IV4 section16 and FGA236 Corollary1.7, p.229, organize the actual finite principal-parts map. Its 1/1220 upper local slope is a computed evaluation property, not a consequence of a virtual representation identity. \\
\end{longtable}

For every entry involving a metric, the period, archimedean norm, finite lattice and actual map remain explicit data. A statement in a representation lambda-ring, a Dieudonne crystal, or an l-adic Grothendieck group does not determine those missing normalizations.

\subsection{1. What the archive pages actually supply}\label{what-the-archive-pages-actually-supply}

\textbf{Cote 54.} The inspected manuscript concerns the universal vector extension of an abelian variety, its de Rham realization, Cartier duality, and the restricted Lie structure in characteristic \(p\). In the official PDF, scan 003 (manuscript page 2) displays the Hodge exact sequence and the vector extension; scan 005 (manuscript page 4) discusses the Cartier dual construction; scan 007 (manuscript page 6) treats Frobenius, Verschiebung, and the restricted operation on the vector kernel. The zero restricted operation there belongs to that vector kernel. It is not a statement that every invariant derivation of an abelian variety has zero \(p\)-th power. AM-K013 and AM-K017 compute precisely the different restricted powers that occur on the actual Jacobian.

\textbf{Cote 4.} Scan 003 (printed page 2), §8.1.1, relates crystalline and de Rham cohomology through the stated derived comparison; §8.1.2 records the associated Tor issue. Scan 005 (printed page 4), §8.4.1--2, singles out the Frobenius-bijective stable part. The manuscript's terminology ``semi-simple'' in this passage refers to that Frobenius-stable part; it must not be substituted for ordinary linear diagonalizability. Scan 006 (printed page 5), Lemma 8.4.3, compares the stable parts of de Rham and coherent cohomology and uses the Frobenius action on the filtration. The subsequent §8.5 contains conditional ingredients. Scan 012 (printed page 11), §8.6, asks a monodromy question for the indicated finite-field local systems and explains a possible trace/congruence consequence. It is a proposed route in the manuscript, not a theorem about all number-field Hecke \(L\)-functions.

These are the exact archive loci used here. Their content organizes the realization and monodromy questions; no newly computed theta slope or automorphic nonvanishing is inferred from an archival question.

\subsection{2. Formal groups, crystals, and integral differential operators}\label{formal-groups-crystals-and-integral-differential-operators}

Over a finite extension where the Jacobian has good reduction, its \(p\)-divisible group and Cartier dual determine dual Dieudonné objects. The accepted DP\_MULT DP30--34 and DP49 distinguish the one-dimensional height-three factor from its two-dimensional dual. The three labelled tangent directions are not three independent height-three formal lines. Frobenius permutes the CM labels and Verschiebung controls the restricted derivations. This is why the restricted nilpotence exponents in AM-K013 are direction dependent.

The exact finite-level framework is SGA 3, Exposé VII\_A, §§6.2--6.4, 7.2--7.4, and 8.1, as carried in DP\_MULT. For the de Rham/finite-group comparison, the carried primary locators are Oda (1969), Theorem 3.19, p.~89, and Corollary 5.11, p.~131. DP\_MULT's source audit additionally fixes Tate (1967), §2.2, Proposition 1, p.~162, equation (3), p.~163, equation (4), p.~164, and §2.3, Proposition 3, p.~166; Zink's dual display operators are Definition 19, equations (22)--(24), pp.~143--144, and the formal-biextension injection is Proposition 121 and equation (243), pp.~241 and 246. Those established results are cited, not proved again in this run.

The new theta statements concern the completed local rings and actual sections of a polarized abelian scheme. In AM-K012 the explicit ordered divisor chart is étale on a dense open and supplies a ring stable under the invariant derivations. AM-K017 extends a logarithmic Hessian across the theta divisor by checking every horizontal and vertical divisor and using normality. It then uses properness to turn a global regular function on the special fibre into a constant. This is a scheme-theoretic argument on the actual integral model. It does not replace a missing integral reduced theta frame by a complex analytic one.

The resulting Hurwitz series use the ordinary invariant operators (D\_i\^{}n), with factorials in the multiplication law. Their reductions and restricted-power identities are explicit. A Chern connection on an ample theta line is not thereby made into a global flat algebraic connection or a Dieudonné crystal. To make that additional assertion one would have to construct the relevant stratification, integrable connection, and Frobenius structure. The present proof does not do so.

For the oscillator, the same distinction has an exact EGA locator: EGA IV\_4, Theorem 16.11.2, p.~54, and Remarks 16.11.3(ii), p.~55. Over an integral base the full sheaf of differential operators contains divided derivatives and is generally larger than the algebra generated by coordinates and ordinary derivations. AM-DO-003 constructs an integral Weyl-algebra module, becoming the familiar delta \(D\)-module after rationalization. It does not claim an integral action of the full EGA differential-operator sheaf.

\subsection{3. Picard schemes, the theta line, and finite models}\label{picard-schemes-the-theta-line-and-finite-models}

The Jacobian is a Picard variety. The precise carried FGA references are Bourbaki exposé 232, Theorem 3.1, pp.~148--149; exposé 236, Corollary 1.7, p.~229, Corollary 2.3, p.~232, and Proposition 2.10(ii), p.~236. They give the representability and base-change framework used in DP\_MULT DP22 and DP49. The actual divisor and section calculations in the present run are additional calculations on that Picard scheme.

ALP-LIFT-07 distinguishes the non-effective characteristic (D\_\{135\}) from the invariant trace-one theta characteristic. They differ by a nonzero element of (J{[}2{]}). This is a distinction in the relative Picard scheme, so an analytic normalization or a scalar change cannot remove it. Pullback by an odd integer preserves the nontrivial two-torsion difference. AM-DO-015 instead uses the specified Mumford translation; AM-R-006 computes its actual scalar.

ALP-LIFT-08 computes the leading determinant section on an integral Picard chart, including residue characteristic two. Its first nonzero jet is the determinant of the cup-product matrix, (u\_1u\_3-u\_2\^{}2). The Hodge-basis transformations in ALP-LIFT-05 then determine the vertical contents at 2 and 7. The objects are line bundles, their primitive sections, and their integral jet images over specified discrete valuation rings. An equality after tensoring with the fraction field would not determine these contents.

AM-R-004--006 and ALP-LIFT-02--06 give the corresponding Betti lattice, de Rham basis, complex metric, determinant metric, and exact scalar comparison. Thus the height identities live on the metrized Hodge determinant line, with the finite integral lattices retained. The archimedean Cauchy estimate and the finite derivative estimate are different realizations of the same normalized section; neither may be changed independently in the cost calculation.

\subsection{4. Biextensions and the phase problem}\label{biextensions-and-the-phase-problem}

The polarized Poincaré line on \(J\times J^\vee\) is a \(\mathbf G_m\)-biextension, rigidified along both zero sections. Pullbacks, duality, and transposition organize the theta-group and trace-one constructions. The exact inspected SGA 7 I loci are:

\begin{itemize}
\tightlist
\item
  Exposé VII, §2.7, printed pp.~169--170: transposition of a biextension. These are PDF pages 174--175 in the original scan.
\item
  Exposé VII, equation (2.9.1) and Proposition 2.9.3, printed pp.~172--174: the forgetful functor to birigidified torsors and its fully faithful form under the displayed hypotheses. Example 2.9.5, p.~175, gives the abelian/Picard setting.
\item
  Exposé VIII, Theorem 7.1(a--b), printed pp.~299--300: extension over a trait. In part (b), the component groups enter a canonical obstruction pairing into \(\mathbf Q/\mathbf Z\); under the stated bounded-torsion hypothesis its vanishing is equivalent to extension. The connected-fibre case is explicitly separated.
\item
  Exposé VIII, Proposition 7.4(b), printed pp.~310--311: full faithfulness over a reduced base for the indicated smooth group schemes with connected generic fibres; over a normal base, with one group having connected fibres, the stated generic-fibre criterion describes the essential image.
\end{itemize}

The local finite Gauss factors of AM-DO-017 are additional data in a specified analytic/Heisenberg realization. The uniqueness of a birigidified biextension structure does not say that every chosen intertwiner of its realizations has scalar one. AM-DO-020 proves that scalar assertion at relative split places by evaluating positive normalized lattice indicators. AM-DO-021 exhibits the actual residual compact sign at the ramified place 7. This is a normalization issue inside an existing realization, not a counterexample to the SGA 7 existence or uniqueness theorems.

The common period in AM-DO-018 is retained for every moved lattice. Its coefficient denominators are intersection indices, and its norm is computed with specified probability Haar measure. In geometric language the Heisenberg eigenspaces are the one-dimensional section spaces of principal theta lines on isogenous complex tori, compared over common finite covers. Their phase and arithmetic integral structure are separate from the isomorphism class of the line bundle.

\subsection{5. Group schemes, weights, and the oscillator}\label{group-schemes-weights-and-the-oscillator}

The exact fresh SGA 3 locator for a diagonalizable-group action is Exposé I, Proposition 4.7.3; Corollary 4.7.3.1 gives the corresponding graded-algebra/equivariant-affine-scheme statement. Accordingly the charges in T-052 are actual characters of a torus, and its degree pieces are equivariant locally free modules. The nonnegative equal-charge construction of the accepted T-053 is a specified graded coordinate subalgebra. It is not merely a dimension count.

AM-DO-002 and AM-DO-007 compare the actual character weights of the canonical ball representation with the required polynomial degrees. AM-DO-006 proves an obstruction for the entire creation-operator weight set. AM-DO-003 changes polarization by a Fourier/Weyl algebra map and records why the resulting delta module is not the positive Fock vacuum. AM-DO-005 supplies a selected bicanonical comparison with an exact local Haar pairing, so the obstruction must not be generalized to every representation or determinant twist.

AM-R-007 belongs to the reductive-group side: the exact determinant sequence separates (SU\(V\)) from the norm-one torus in (U\(V\)), and strong approximation gives representatives for a Hermitian lattice genus with level structure. SGA 3 supplies the group-scheme framework; the approximation theorem actually used is Rapinchuk (2013), Theorem 2.3, p.~281. The source locator is not relabelled as a theorem of SGA 3. ALP-LIFT-09 then constructs an actual automorphic character and theta integral on the norm-one torus. That analytic integral requires its Haar measure and splitting, beyond the algebraic group scheme and its representation class.

\subsection{6. Lambda rings, K-theory, and what a class does not determine}\label{lambda-rings-k-theory-and-what-a-class-does-not-determine}

For a finite locally free representation \(V\), its exterior powers define the lambda operations on its exact Grothendieck group. The carried exact locators are SGA 6, Exposé V, §2 (lambda operations), §3 (gamma operations and filtration), Proposition 4.2, Definition 4.4, Proposition 4.8, and Corollary 4.9 (the universal representation-ring description); Exposé VI, Theorem 3.3, is the associated Schur-functor framework used by GCT\_PLETH. Cote 133, scan 006, records the archival pre-lambda/classifying-topos framework, as carried in that run's source study. The SGA 6 download attempt in this run failed; these numbered references remain carried references, not fresh primary inspections.

Each finite degree of the oscillator polynomial algebra has a class formed by symmetric and Schur operations. A tensor-to-Hom map such as AM-DO-009 induces the corresponding equality of representation classes after dualizing the source factor. The map itself retains the contragredient action, the alternating form and the determinant roots. The representation class alone forgets the positive metric, the integral section lattice, finite Weil phases and the actual map on automorphic forms.

Likewise AM-R-002's content ideal is data of an actual integral linear map applied to an actual vector. Its least denominator and dual metric can be expressed through determinant lines and metrized lattices; they are not recovered from a virtual rank. AM-R-001's cost subtraction and ALP-LIFT-06's directional cancellation compare those retained numerical realizations. No lambda-ring identity is used to erase a height, a covolume or an endpoint denominator.

\subsection{7. CM motives, Artin motives, and the RH boundary}\label{cm-motives-artin-motives-and-the-rh-boundary}

AM-R-003 uses the primitive CM type to separate the simple rational weight-one Hodge structure of (J) from an elliptic induced type. The corresponding CM realizations are those of (h\^{}1(J)), with its endomorphism action and polarization. Products and field extensions do not turn that simple isotypy class into the source's elliptic Serre-tensor class. AM-DO-010 instead identifies specified CM eigensummands with actual Galois twists up to isogeny. It does not identify their integral models without the separate lattice data.

For AM-K001--004 the order-(r) Kummer character is a rank-one Artin representation over the cyclotomic base. Over the open arithmetic base excluding its ramified primes it is a lisse rank-one sheaf with finite monodromy. Its induced representation is the stated augmentation representation over \(\mathbf Q\). The exact representation identity therefore carries the ramified Euler factors as recorded in the proof. For (r=3), the explicit isomorphism of two CM elliptic curves over the Kummer extension identifies the character with a line in their Hom realization. This gives an actual descent cocycle, rather than an analogy based on dimension.

The proper/smooth realization framework is SGA 4, Exposé XII, Theorems 5.1--5.2; Exposé XV, Theorem 2.1; Exposé XVI, Theorem 2.1, and SGA 5, Exposé VI, §§1.2.1, 1.4.2--1.4.3, pp.~265--266, as carried in CX\_GRH/GRH\_NIGHT. For finite-field weights the exact carried locators are Deligne's \emph{Weil I}, Theorem 1.6, p.~276, and \emph{Weil II}, Corollary 3.3.9, p.~207. They control Frobenius eigenvalues in their stated finite-field cohomological settings. In particular an Artin line has weight zero. These statements do not put the zeros of its number-field Hecke \(L\)-function on the critical line.

AM-K005 isolates the remaining analytic implication through an explicit every-epsilon arithmetic window. AM-K007--011 supply ray characters, splitting twists and CM hosts while retaining ramified conductors, prime norms and host costs. Cote 4 §8.6's monodromy question is relevant background for such realization questions; it does not supply the missing arithmetic window. The accepted Dirichlet theorem of Alpöge remains an input with its original scope. No statement of Masullo's draft is used as an unconditional theorem.

\subsection{8. Coverage of findings and the programme's rows}\label{coverage-of-findings-and-the-programmes-rows}

\begin{longtable}[]{@{}
  >{\raggedright\arraybackslash}p{(\columnwidth - 4\tabcolsep) * \real{0.3333}}
  >{\raggedright\arraybackslash}p{(\columnwidth - 4\tabcolsep) * \real{0.3333}}
  >{\raggedright\arraybackslash}p{(\columnwidth - 4\tabcolsep) * \real{0.3333}}@{}}
\toprule\noalign{}
\begin{minipage}[b]{\linewidth}\raggedright
Findings in this run
\end{minipage} & \begin{minipage}[b]{\linewidth}\raggedright
Concrete geometric content
\end{minipage} & \begin{minipage}[b]{\linewidth}\raggedright
Original row and remaining boundary
\end{minipage} \\
\midrule\noalign{}
\endhead
\bottomrule\noalign{}
\endlastfoot
AM-R-001--002; ALP-LIFT-01 & Metrized integral row/endpoint maps, their content ideals, and exact changed-family accounting & DP\_MULT DP44, DP52, Q3e; the actual terminal vector and all costs remain required. \\
AM-R-003; ALP-LIFT-02--06 & Primitive CM isotypy; polarized Hodge lattices and determinant metrics; finite normalization and directional balance & DP\_MULT DP49, DP57, DP64, Q3e; a bare-certificate equality is not a selected-gain upper bound. \\
AM-R-004--006; ALP-LIFT-07--08; AM-DO-011--016 & Actual Picard/theta lines, characteristic translation, integral leading determinant, periods, scalar contents and archimedean bounds & DP\_MULT DP22, DP49, DP54--65; source-line identification and normalization are explicit. \\
AM-K012--018 & Integral Abel charts, algebraic logarithmic Hessians, restricted powers, explicit Hurwitz jets, growth, mixed-jet counterexamples, and exact full-CM canonical-gauge matching & DP\_MULT Q3e and the local step-2 input; the two fixed mixed-line slopes now apply to the actual section, while conjugate-action products retain a separate gauge condition. \\
AM-DO-017--022; AM-R-007 & Principal theta eigenspaces on commensurable lattices, exact normalized averages, split-place phase, ramified compact character, and all-finite phase repair over the unramified CM host & THETA3 G1 and DP\_MULT Q3e; full arithmetic rows, later images and terminal mass remain distinct equations. \\
ALP-LIFT-09--10 & Genuine first theta integrals from an anisotropic torus, with actual splitting, permitted half-torsion orbit, nonvanishing, exact finite-index input norms and character projections & DP\_MULT Q3e; later definite lifts and the prescribed chain are not supplied by this first image. The first-lift host E and the repaired all-finite host E-prime still require a common-tower comparison. \\
AM-DO-001, 004, 009--010 & Signed Hermitian spaces, positive multiplicity space, exact tensor-to-Hom map and CM isogeny realization & THETA3 A7 and G1; the common positive arithmetic pin and global MainX value need their own comparison. \\
AM-DO-002--003, 005--008 & Equivariant oscillator modules, integral Weyl/Fourier map, degree and charge obstructions, selected bicanonical pairing, finite scalar equations & GCT\_PLETH T-052/T-053, Q-ALPOGE; ALPOGE\_GCT Q-ALPOGE-METHOD; scoped obstructions coexist with a selected local comparison. \\
AM-K000--006 & Artin--Schreier finite-field construction, Kummer Artin motive, cubic conductor, CM Hom line, analytic pole/window reduction & CX\_GRH CX-ALG-024; the general Kummer RH and factoring implication remain conditional on the missing window. \\
AM-K007--011 & Explicit finite etale/ray characters, determinant-root twists and arithmetic CM-host lattices & CX\_GRH and GRH\_NIGHT Kummer gate; relative prime/conductor and quantitative tower adapters remain. \\
\end{longtable}

New findings added after this version must be added to this table before the final package. The table is a coverage record, not a new proof of the findings it indexes.

\subsection{References and source status}\label{references-and-source-status}

Grothendieck, A. (n.d.). \emph{Manuscript cote 4} {[}Digitized manuscript{]}. Archives Grothendieck, Université de Montpellier. \url{https://grothendieck.umontpellier.fr/4.pdf} . Scans and paragraph locators are listed in §1; primary PDF and the relevant page images were inspected.

Grothendieck, A. (n.d.). \emph{Manuscript cote 54} {[}Digitized manuscript{]}. Archives Grothendieck, Université de Montpellier. \url{https://grothendieck.umontpellier.fr/54.pdf} . Exact used scans 003, 005, 007; primary inspection in this run.

Grothendieck, A. (1962). \emph{Techniques de construction et théorèmes d'existence en géométrie algébrique V--VI: Les schémas de Picard}. Séminaire Bourbaki, exposés 232 and 236. Exact theorem/page locators in §3, carried through the supplied DP\_MULT DP22 and its intake.

Grothendieck, A., with Dieudonné, J. (1967). Éléments de géométrie algébrique IV, quatrième partie. \emph{Publications Mathématiques de l'IHÉS, 32}, 5--361. Exact used locus IV\_4, 16.11.2--16.11.3(ii), pp.~54--55; freshly read by the oscillator agent.

Demazure, M., \& Grothendieck, A. (Eds.). (1970). \emph{Schémas en groupes (SGA 3), Tome I} (Lecture Notes in Mathematics, Vol. 151). Springer. Exposé I, 4.7.3 and 4.7.3.1 were freshly read in the recomposed French text; VII\_A locators in §2 are carried.

Berthelot, P., Grothendieck, A., \& Illusie, L. (Eds.). (1971). \emph{Théorie des intersections et théorème de Riemann--Roch (SGA 6)} (Lecture Notes in Mathematics, Vol. 225). Springer. Exact carried lambda/K-theory locators in §6; no successful fresh PDF retrieval claimed.

Grothendieck, A., Raynaud, M., \& Rim, D. S. (1972). \emph{Groupes de monodromie en géométrie algébrique (SGA 7 I)} (Lecture Notes in Mathematics, Vol. 288). Springer. Exact freshly inspected VII--VIII locators in §4. The original-scan PDF has five front-matter pages: PDF page number equals printed page plus five at these loci.

Oda, T. (1969). The first de Rham cohomology group and Dieudonné modules. \emph{Annales Scientifiques de l'École Normale Supérieure, 2}(1), 63--135. Theorem 3.19 and Corollary 5.11 are used through the accepted DP\_MULT realization results.

Deligne, P. (1974). La conjecture de Weil. I. \emph{Publications Mathématiques de l'IHÉS, 43}, 273--307. Theorem 1.6, p.~276.

Deligne, P. (1980). La conjecture de Weil. II. \emph{Publications Mathématiques de l'IHÉS, 52}, 137--252. Corollary 3.3.9, p.~207. Both weight references are carried framework inputs, with their finite-field scope retained.

\subsection{10. Additional checked pairing and twist refinements}\label{additional-checked-pairing-and-twist-refinements}

AM-DO-024 supplies a positive comparison between the actual finite Weil translate and the direct Heisenberg average on the E-prime host, together with the exact probability-Haar pairing and adjoints. These are identities of metrized section lines in the biextension/Heisenberg realization of section4, stronger data than the isomorphism class of a Picard line. They do not determine an automorphic quotient mass. AM-R-008 supplies actual integral diagonal maps for tensoring compatible coefficient modules by a unitary determinant character, within the SGA3 character and SGA6 representation K-theory frameworks of sections5--6. Its all-prime content preservation comes from root-of-unity scalars, not from equality in K0. The first-lift target-level membership and its refinement cost remain additional obligations. AM-DO-023 is held pending reconciliation of the general character convention; its unchanged-character local vanishing has a separate narrower scope.

\subsection{11. Unchanged first raw kernel and its mass}\label{unchanged-first-raw-kernel-and-its-mass}

AM-DO-036 proves positivity for the actual unchanged first raw theta tensor on the principally polarized J14/J56 hosts. Its finite lattice realization is the section/Heisenberg construction organized by the rigidified Poincare biextension of SGA7 I, VII2.7/2.9 (pp169--174) and VIII7.4(b) (pp310--311), and the Picard representability interfaces FGA232 Theorem3.1 (pp148--149) and FGA236 Corollary1.7 (p229),2.3 (p232). Those scheme results provide the object and functoriality, not this analytic positivity. The latter is a convergent doubled theta norm calculation with anisotropic target, using GQT7.1 (p33); its target quotient mass remains an explicit scalar. It therefore removes the first fixed-kernel vanishing gate of DP\_MULT Q3e and narrows THETA3 G1, without supplying an integral later image lattice or terminal mass identity. The common CM weight-one motive is retained; its Dieudonne and Cartier-dual realization remains the cote4/cote54 correspondence recorded in sections2--3. No finite-field weight statement or equality in representation K-theory supplies the missing global mass or number-field RH.

\subsection{12. Mixed periods, embedding ideals and local quotient scalars}\label{mixed-periods-embedding-ideals-and-local-quotient-scalars}

AM-K035 expresses the polarized Picard variety through its Betti/de Rham comparison with an actual per-embedding period vector. Its coordinate ideals are integral module intersections over Spec(O\_M), and its reduced theta jets are formal sections of the rigidified line in the Picard/biextension setting of FGA232 Theorem3.1 (pp148--149), FGA236 Corollary1.7 (p229),2.3 (p232), SGA7 I VII2.7/2.9 (pp169--174) and VIII7.4(b) (pp310--311). The all-order algebraic formal coordinates use the EGA IV4 16.11.2--16.11.3(ii) interface (pp54--55) already recorded above. The explicit analytic period and unit-Fock constants are new comparisons, not consequences of representability alone. At11 its cD is a unit, so it preserves the previously proved Dieudonne/Cartier-dual selected-jet content; cote4 and cote54 scans5--7 remain the exact historical framework locators of sections2--3.

ALP-LIFT-24 and AM-DO-037 calculate actual local quotient maps and their Haar/L/Gauss scalars. Their finite inputs are section representations of the same Heisenberg extension; the literal DegOne condition is the statement e=f=1 for the relevant finite localization. An identity in the representation lambda-ring or K-group (SGA6 VI3.3) records types but does not determine these periods, integral content or the global terminal mass. Likewise finite-field weights of the CM motive do not supply number-field RH for the Kummer character. These scope distinctions carry forward DP\_MULT Q3e and THETA3 G1 without changing CX\_GRH CX-ALG-024.

\subsection{13. Fixed-edge descent and compact contractions}\label{fixed-edge-descent-and-compact-contractions}

K036 is a descent statement for complete formal sections and their torsion translations on the fixed Picard/biextension family. It uses the exact algebraic rational-isometry scalar and finite-class reduction, not representability alone; its FGA/SGA7/EGA locators are those in sections4 and12. The explicit nc measures a finite isogeny/Heisenberg induction coefficient, while Cartier/Dieudonne slopes (cote4; cote54 scans5--7) concern the separate formal section content. LIFT25 computes an actual contraction in the compact representation category and its Fock metric, refining the SGA3 I4.7.3/4.7.3.1 and SGA6 VI3.3 class information. The dimension denominator and invariant-tensor norm are distinct. No virtual representation or motivic equality erases these scalar costs.

\subsection{14. Integral branching with its metric}\label{integral-branching-with-its-metric}

R019 makes the SGA3/SGA6 representation/flag restriction of section13 quantitative: the actual integral minor vector, spectral integer and metrized branch have exact factorial expressions. SGA3 I4.7.3/4.7.3.1 and SGA6 VI3.3 organize the representations, while Molev Proposition2.4 supplies the classical contravariant norm input and the new path calculation fixes its source normalization. The integer scalar is not an integral splitting theorem for every chosen lattice. The Picard/biextension and cote4/cote54 formal realizations remain those already identified; no new motivic or number-field RH consequence follows from this metric refinement.

\subsection{15. ALP-LIFT-26: fixed coupled analytic image}\label{alp-lift-26-fixed-coupled-analytic-image}

The two selected correspondences have a nonzero composition on their actual complex automorphic realization, with prescribed local sections and every multiplicity copy retained. In the Picard/biextension interpretation of sections12--14, the tensor identity uses the same transported line-bundle frames; its character and volume factors are part of the realization, not an identification of integral lattices. The exact Grothendieck locators remain FGA232 Theorem3.1 pp148--149 and236 Corollary1.7 p229/2.3 p232 for Picard geometry, SGA7I VII2.7/2.9 pp169--174 and VIII7.4(b) pp310--311 for biextensions, and EGAIV4 16.11.2--3(ii) pp54--55 for jets. None of these citations identifies the image of the arithmetic class map or its terminal branching projection. In lambda/K-theory the selected Schur factors are the same factors already recorded in section14; in motivic or ell-adic language no new algebraic correspondence comparison or RH statement follows from this analytic nonvanishing theorem.

\subsection{16. AM-K037: integral polynomial functor and class-map realization}\label{am-k037-integral-polynomial-functor-and-class-map-realization}

The map CR7 is a morphism of the specified integral Weyl polynomial functors over Z; the monic-determinant descent proves it extends from GL to the full matrix scheme on the relevant image. Its evaluation against the independently fixed algebraic theta jet is the actual arithmetic class map. This places the new fact directly in the polynomial-functor/lambda-ring language of the earlier SGA6 VI3.3 comparison, while its jet realization uses EGAIV4 16.11.2--3(ii), pp54--55. The fixed line bundle and reduced theta section retain the Picard and biextension locators of sections12--15 (FGA232/236 and SGA7I VII--VIII). Coevaluation is the usual duality morphism for the selected coefficient representation; its norm and finite stabilizers are part of the realization and are not removed by the K-theory description. No new identification of an integral motivic image, ell-adic weights or Kummer RH is claimed.

\subsection{17. AM-DO-038: terminal contraction and mass}\label{am-do-038-terminal-contraction-and-mass}

The fixed branch is a morphism between the selected coefficient realizations; the finite class pairing contracts it with the dual row and weights each class by its actual finite stabilizer. The top automorphic realization realizes this same pairing and introduces the exact adelic level volume. This is a comparison at the complex/arithmetic coefficient level, not a new identity of schemes or motives. Its polynomial/jet part has the same SGA6 VI3.3 and EGAIV4 16.11.2--3(ii) pp54--55 locators; its fixed line/frame transports retain FGA232 Theorem3.1 pp148--149, FGA236 Corollary1.7 p229 and2.3 p232, and SGA7I VII2.7/2.9 pp169--174 andVIII7.4(b) pp310--311. Stabilizer cancellation and compact Schur integration do not imply nonzero algebraic cycle classes or ell-adic purity. The latter realizations require the independently stated correspondence comparison, still open in THETA3 G1.

\subsection{18. AM-R-020: integral Weyl-image sandwich}\label{am-r-020-integral-weyl-image-sandwich}

The new inclusions are statements about finite integral representation modules over Spec(Z) and their localizations, and they survive extension to Spec(O\_M). They refine the rational restriction/Schur-functor information organized by SGA3I4.7.3/4.7.3.1 and SGA6VI3.3: equality of rational representation classes does not determine an integral image or its elementary divisors. The exact finite-place bounds use the specified source divided-power lattice and normalized duality. Their application to the Picard/biextension theta section retains the existing FGA232 Theorem3.1 pp148--149, FGA236 Corollary1.7 p229/2.3 p232, SGA7I VII2.7/2.9 pp169--174 andVIII7.4(b) pp310--311 locators; R020 adds no claim about that earlier image. The Dieudonne/Cartier-dual framework remains cote4 and cote54 scans5--7 as recorded in sections2--3. No new motivic correspondence, ell-adic weight or Kummer RH implication is asserted.

\subsection{19. AM-K038 and ALP-LIFT-27: coupled section descent and faithful analytic image}\label{am-k038-and-alp-lift-27-coupled-section-descent-and-faithful-analytic-image}

K038 is an explicit descent construction of the two composed theta-section differential operators on the finite decorated lattice quotient, with coefficient modules over O\_M and actual stabilizer weights. Its polynomial scheme maps use the existing SGA3I4.7.3/4.7.3.1, SGA6VI3.3 and EGAIV4 16.11.2--3(ii) pp54--55 interfaces. The section/translation geometry retains FGA232 Theorem3.1 pp148--149, FGA236 Corollary1.7 p229/2.3 p232, SGA7I VII2.7/2.9 pp169--174 andVIII7.4(b) pp310--311. The strict/projective refinement is an actual quotient mass identity, not a change in the underlying Picard or motivic object. LIFT27 proves faithfulness of the specified full analytic image on its initial ramified space; it does not identify a faithful algebraic correspondence after terminal restriction. The formal selected content remains within the cote4/cote54 Cartier--Dieudonne framework already recorded. Neither analytic injectivity nor equality of lambda-ring classes proves a nonzero terminal motive, an ell-adic weight assertion or Kummer RH.

\subsection{20. AM-R-021: character twists and numerical weight windows}\label{am-r-021-character-twists-and-numerical-weight-windows}

Tensoring by a determinant character shifts both u and v equally; the v-width and first u−v gap are invariant. The full arithmetic budget requires an inequality between these invariants, violated by the current equally spaced family. This is a numerical restriction on coefficient representations within the SGA3I4.7.3/4.7.3.1 and SGA6VI3.3 framework, not a new geometric obstruction to all theta correspondences. It neither changes the existing Picard/biextension/formal crystal constructions nor supplies an ell-adic or motivic RH conclusion. The exact source and prior-run locators are R021§1--3.

\subsection{21. AM-DO-039: finite integral tests and relative representation models}\label{am-do-039-finite-integral-tests-and-relative-representation-models}

The common split-line kernel is a kernel of morphisms of the specified finite polynomial representation modules over Z and their localizations. Its characteristic-zero dimension is not determined by a single highest line. This refines the coefficient-restriction picture of SGA3I4.7.3/4.7.3.1 and SGA6VI3.3; it does not assert a kernel in the actual theta-correspondence image. The local Q open/closed orbits are relative representation models of one fixed analytic vector and are not different motivic realizations. The Picard, biextension and jet comparisons retain FGA232 Theorem3.1 pp148--149, FGA236 Corollary1.7 p229/2.3 p232, SGA7I VII2.7/2.9 pp169--174 andVIII7.4(b) pp310--311, and EGAIV4 16.11.2--3(ii) pp54--55. No new cycle nonvanishing, ell-adic weight assertion or L-function theorem follows.

\subsection{22. AM-K039: actual section integrality and class descent}\label{am-k039-actual-section-integrality-and-class-descent}

The new divisibility theorem is about the integral local-section module of the fixed principal theta line and its actual class descent over a DVR above11. Prime-to11 Heisenberg translation is identified through the normalized torsion difference connection, so the reduced exponential is retained. Its Picard/biextension/jet setting has the exact FGA232 Theorem3.1 pp148--149, FGA236 Corollary1.7 p229/2.3 p232, SGA7I VII2.7/2.9 pp169--174 andVIII7.4(b) pp310--311, EGAIV4 16.11.2--3(ii) pp54--55 locators already recorded. The selected differential functor acts on integral Weyl modules over Z (SGA3I4.7.3/4.7.3.1; SGA6VI3.3), while the numerical11 bound uses the established Cartier--Dieudonne comparison (cote4; cote54 scans5--7). Finite class stabilizers are explicit descent denominators; a K-theory class or a motivic description does not remove them. These inclusions give no nonzero terminal cycle, ell-adic weight theorem or Kummer RH.

\subsection{23. AM-R-022: a feasible coefficient-weight cloud}\label{am-r-022-a-feasible-coefficient-weight-cloud}

The new profile supplies explicit highest weights and binary interlacing grids in the same integral representation framework (SGA3I4.7.3/4.7.3.1; SGA6VI3.3), with exact real Scale parameters and rational coefficient weights. It changes the coefficient representations used by the jets, not the underlying J-power, Picard scheme or biextension line; the exact FGA/SGA7/EGA jet locators are those in R022§5 and sections12--22 above. Numerical window feasibility does not identify the integral theta image or a motivic correspondence and implies no ell-adic weight or RH assertion.

\subsection{24. ALP-LIFT-28: coherent jets and a bottom detector}\label{alp-lift-28-coherent-jets-and-a-bottom-detector}

The new bottom detector is an application of coherent finite-jet interpolation on a projective minimal compactification. Its exact Grothendieck input is EGAIII.1 Theorem2.2.1(ii), p100: sufficiently ample twists kill higher cohomology of the coherent kernel of finite-jet restriction. The finite jet space retains EGAIV4 16.11.2--3(ii), pp54--55. Mumford1977 Theorem3.1 and Proposition3.4(a),(b) supply the canonical extension and ample logarithmic-canonical line. These results construct an actual analytic/coherent detector at the explicit level, with stabilizer-equivariant coefficients; they do not identify it with a prescribed theta correspondence or bound its norm. The Picard/biextension locators remain FGA232/236 and SGA7I VII--VIII as above. No ell-adic cycle or motivic L-function conclusion follows.

\subsection{25. AM-K040: coherent CM character types and arithmetic coefficient fields}\label{am-k040-coherent-cm-character-types-and-arithmetic-coefficient-fields}

The principal/base-ideal compatibility builds actual algebraic angular characters for every weight in a fixed lattice, with one finite value field and preserved finite conductors. In the CM-motive language this is the infinity-type/finite-realization datum of rank-one character realizations; it is not a proof of RH or an automatic comparison with the mixed theta motive. The underlying CM Picard scheme and biextension line retain FGA232 Theorem3.1 pp148--149, FGA236 Corollary1.7 p229/2.3 p232 and SGA7I VII2.7/2.9 pp169--174, VIII7.4(b) pp310--311. The new theorem is an ideal/character construction, not a new theorem from those geometric locators. Its finite radical-field bound controls degree only; no unrecorded discriminant or height bound is supplied.

\subsection{26. AM-DO-040: certified prime and decorated lattice data}\label{am-do-040-certified-prime-and-decorated-lattice-data}

The explicit algebraic integers piT give finite etale split residue fibres and the actual decorated integral lattice support for the replacement coefficient weights. Their deterministic primality, norm and stabilizer assertions are arithmetic inputs to the existing scheme/jet comparison, not new theorems about Picard schemes. The exact geometric locators remain FGA232/236, SGA7I VII--VIII and EGAIV4 16.11.2--3(ii) stated above. Strict/projective level masses and finite character-field enlargements remain numerical costs; the lambda-ring or motivic description does not cancel them. No new nonvanishing or L-function conclusion follows.

\subsection{27. ALP-LIFT-29: translated finite packet and the same coherent jets}\label{alp-lift-29-translated-finite-packet-and-the-same-coherent-jets}

The finite level and translation define an actual collection of special points with exactly the original isotropy groups; the displayed local matrix calculation fixes the intersection of their integral models at T. The coherent detector therefore uses the same EGAIII.1 Theorem2.2.1(ii), p100, and EGAIV4 16.11.2--3(ii), pp54--55, finite-jet interfaces as LIFT28. The translation is an isometry of the complex automorphic realization; it is not asserted to preserve an integral section lattice or have zero arithmetic height. The Picard/biextension locators remain FGA232 Theorem3.1 pp148--149, FGA236 Corollary1.7 p229/2.3 p232 and SGA7I VII2.7/2.9 pp169--174, VIII7.4(b) pp310--311. No identification of the prescribed correspondence, motive, ell-adic cycle or L-function follows from the packet comparison alone.

\subsection{28. AM-R-023: finite quotient and scale constraints}\label{am-r-023-finite-quotient-and-scale-constraints}

The fixed-field obstruction is a finite Galois quotient statement over Spec(Q), independent of the valid mixed Picard/section construction. The conditional rank bounds concern the size of a full CM realization of h1(J\^{}N); they do not remove R003's separate simple-isogeny/type obstruction. The exact FGA/SGA7 Picard and biextension locators, SGA6VI3.3 representation framework and cote4/cote54 crystal realizations are retained in R023§5. No change in a virtual representation, motive or ell-adic realization relaxes the displayed integer and Scale inequalities. No L-function theorem is obtained.

\subsection{29. AM-K041: actual new-profile integral image}\label{am-k041-actual-new-profile-integral-image}

The all-word bound is an inclusion of the actual finite-class coherent Weyl coefficient modules over the DVR above11, with both stabilizer denominators. The primitive theta section is transported through the normalized torsion/Poincare connection on its formal Picard completion, retaining the Chern exponential. Schur operations retain SGA6VI3.3; the exact Picard/biextension and jet locators are FGA232 Theorem3.1 pp148--149, FGA236 Corollary1.7 p229/2.3 p232, SGA7I VII2.7/2.9 pp169--174 andVIII7.4(b) pp310--311, EGAIV4 16.11.2--3(ii) pp54--55. The separate unrestricted analytic injectivity does not imply a nonzero restricted motivic or ell-adic correspondence. No Kummer RH result follows.

\subsection{30. AM-R-024 and ALP-LIFT-30: terminal integral rows and negative-line transport}\label{am-r-024-and-alp-lift-30-terminal-integral-rows-and-negative-line-transport}

R024 gives inclusions of the actual image in Weyl/dual-Weyl coefficient lattices over the DVR above11, retaining normalized pairing, finite stabilizers and every EqFun constraint. The exact SGA3I4.7.3/4.7.3.1 and SGA6VI3.3 frameworks do not remove these indices. LIFT30 gives a primitive integral negative line and rational isometry of its positive complement, followed by a separately required real-vector transport; a unitary Hilbert norm does not imply bounded algebraic coordinate height. The fixed Picard/biextension/formal jet locators remain FGA232 Theorem3.1 pp148--149, FGA236 Corollary1.7 p229/2.3 p232, SGA7I VII2.7/2.9 pp169--174 andVIII7.4(b) pp310--311, EGAIV4 16.11.2--3(ii) pp54--55. No algebraic-cycle nonvanishing, ell-adic weight or Kummer RH statement is inferred.

R024.18--21 applies R002's metrized image-line balance to the actual class-independent branching scalar. Its principal finite divisor and all-embedding real logarithmic norm agree, so the corresponding scalar has degree zero. This is the product formula on the number-field base, not a cancellation of the remaining integral correspondence, restricted image, or detector norm.

\end{document}
